% Auto-generated from data/segments.json + layout.json (+ transforms) — do not edit by hand.
% volume: 40Abhi12
% mode: reading
% segments: 5196
% transform_rules: 0
% toc_marks: 16

% Volume layout defaults (layout.json ``layout``).
\csromanlayoutapply{1}{1.5}{21.6pt}{5pt}{6.3pt}{65pt}{2.5em}%

% Reading physical-page layout (layout.json ``page_layout_reading_mode``).
\csromanreadingpagelayoutvolume{1}{1.5}{21.6pt}{5pt}{6.3pt}{65pt}{2.5em}%
\csromanreadingpagelayoutenable

\pitaka{\textbf{อภิธมฺมปิฏก}}
\csromanheader{2}{\textbf{ธมฺมานุโลม}}
\csromanmatikaanchor{mtk.0}
\csromantocmarknopage{chapter}{ปฏฺฐานปาฬิ}{mtk.0}
\bookmark[dest={mtk.0},level=0]{ปฏฺฐานปาฬิ}
\gambhira{\textbf{ติกติกปฏฺฐานปาฬิ}}
\namakkaram{นโม ตสฺส ภควโต อรหโต สมฺมาสมฺพุทฺธสฺส.}
\csromanheader{2}{1. กุสลตฺติก 1. เวทนาตฺติก}
\csromanheader{2}{1. สุขายเวทนายสมฺปยุตฺตปท 1-7. ปฏิจฺจาทิวาร}
\csromanheader{2}{\textbf{ปจฺจยจตุกฺก-เหตุ}}
\proseitem{1}{\csromanfolio{1}กุสลํ สุขาย เวทนาย สมฺปยุตฺตํ ธมฺมํ ปฏิจฺจ กุสโล สุขาย เวทนาย สมฺปยุตฺโต ธมฺโม อุปฺปชฺชติ เหตุปจฺจยา. (1)}

\prose{อกุสลํ สุขาย เวทนาย สมฺปยุตฺตํ ธมฺมํ ปฏิจฺจ อกุสโล สุขาย เวทนาย สมฺปยุตฺโต ธมฺโม อุปฺปชฺชติ เหตุปจฺจยา. (1)}

\prose{อพฺยากตํ สุขาย เวทนาย สมฺปยุตฺตํ ธมฺมํ ปฏิจฺจ อพฺยากโต สุขาย เวทนาย สมฺปยุตฺโต ธมฺโม อุปฺปชฺชติ เหตุปจฺจยา. (1)}

\proseitem{2}{เหตุยา ตีณิ, อารมฺมเณ ตีณิ ฯเปฯ อวิคเต ตีณิ. (สํขิตฺตํ. ฯเปฯ สหชาตวารมฺปิ ฯเปฯ สมฺปยุตฺตวารมฺปิ ปฏิจฺจวารสทิสํ.)}

\proseitem{3}{กุสโล สุขาย เวทนาย สมฺปยุตฺโต ธมฺโม กุสลสฺส สุขาย เวทนาย สมฺปยุตฺตสฺส ธมฺมสฺส เหตุปจฺจเยน ปจฺจโย. (1)}

\prose{อกุสโล สุขาย เวทนาย สมฺปยุตฺโต ธมฺโม อกุสลสฺส สุขาย เวทนาย สมฺปยุตฺตสฺส ธมฺมสฺส เหตุปจฺจเยน ปจฺจโย. (1)}

\prose{\csromanfolio{2}อพฺยากโต สุขาย เวทนาย สมฺปยุตฺโต ธมฺโม อพฺยากตสฺส สุขาย เวทนาย สมฺปยุตฺตสฺส ธมฺมสฺส เหตุปจฺจเยน ปจฺจโย. (1)}

\prose{กุสโล สุขาย เวทนาย สมฺปยุตฺโต ธมฺโม กุสลสฺส สุขาย เวทนาย สมฺปยุตฺตสฺส ธมฺมสฺส อารมฺมณปจฺจเยน ปจฺจโย.}

\proseitem{4}{เหตุยา ตีณิ, อารมฺมเณ นว. (สํขิตฺตํ.) (ยถา กุสลตฺติเก ปญฺหาวารํ, เอวํ วิตฺถาเรตพฺพํ.)}

\csromanheader{2}{2. ทุกฺขายเวทนายสมฺปยุตฺตปท 1-7. ปฏิจฺจาทิวาร}
\csromanheader{2}{\textbf{ปจฺจยจตุกฺก-เหตุ-อารมฺมณ}}
\proseitem{5}{อกุสลํ ทุกฺขาย เวทนาย สมฺปยุตฺตํ ธมฺมํ ปฏิจฺจ อกุสโล ทุกฺขาย เวทนาย สมฺปยุตฺโต ธมฺโม อุปฺปชฺชติ เหตุปจฺจยา. (1)}

\prose{อกุสลํ ทุกฺขาย เวทนาย สมฺปยุตฺตํ ธมฺมํ ปฏิจฺจ อกุสโล ทุกฺขาย เวทนาย สมฺปยุตฺโต ธมฺโม อุปฺปชฺชติ อารมฺมณปจฺจยา. (1)}

\prose{อพฺยากตํ ทุกฺขาย เวทนาย สมฺปยุตฺตํ ธมฺมํ ปฏิจฺจ อพฺยากโต ทุกฺขาย เวทนาย สมฺปยุตฺโต ธมฺโม อุปฺปชฺชติ อารมฺมณปจฺจยา. (1) (สํขิตฺตํ.)}

\proseitem{6}{เหตุยา เอกํ, อารมฺมเณ เทฺว ฯเปฯ อวิคเต เทฺว. (สํขิตฺตํ.) (สหชาตวารมฺปิ ฯเปฯ สมฺปยุตฺตวารมฺปิ ปฏิจฺจวารสทิสํ.)}

\proseitem{7}{อกุสโล ทุกฺขาย เวทนาย สมฺปยุตฺโต ธมฺโม อกุสลสฺส ทุกฺขาย เวทนาย สมฺปยุตฺตสฺส ธมฺมสฺส เหตุปจฺจเยน ปจฺจโย. (1)}

\prose{อกุสโล ทุกฺขาย เวทนาย สมฺปยุตฺโต ธมฺโม อกุสลสฺส ทุกฺขาย เวทนาย สมฺปยุตฺตสฺส ธมฺมสฺส อารมฺมณปจฺจเยน ปจฺจโย. (1)}

\prose{อพฺยากโต ทุกฺขาย เวทนาย สมฺปยุตฺโต ธมฺโม อกุสลสฺส ทุกฺขาย เวทนาย สมฺปยุตฺตสฺส ธมฺมสฺส อารมฺมณปจฺจเยน ปจฺจโย. (1) (สํขิตฺตํ.)}

\csromanheader{2}{1. กุสลตฺติก 1. เวทนาตฺติก}
\proseitem{8}{\csromanfolio{3}เหตุยา เอกํ, อารมฺมเณ เทฺว. (สํขิตฺตํ.\csromanspacer{} ยถา กุสลตฺติเก ปญฺหาวารํ, เอวํ วิตฺถาเรตพฺพํ.)}

\csromanheader{2}{3. อทุกฺขมสุขเวทนายสมฺปยุตฺตปท 1-7. ปฏิจฺจาทิวาร}
\csromanheader{2}{\textbf{ปจฺจยจตุกฺก-เหตุ}}
\proseitem{9}{กุสลํ อทุกฺขมสุขาย เวทนาย สมฺปยุตฺตํ ธมฺมํ ปฏิจฺจ กุสโล อทุกฺขมสุขาย เวทนาย สมฺปยุตฺโต ธมฺโม อุปฺปชฺชติ}

\prose{อกุสลํ อทุกฺขมสุขาย เวทนาย สมฺปยุตฺตํ ธมฺมํ ปฏิจฺจ อกุสโล อทุกฺขมสุขาย เวทนาย สมฺปยุตฺโต ธมฺโม อุปฺปชฺชติ}

\prose{อพฺยากตํ อทุกฺขมสุขาย เวทนาย สมฺปยุตฺตํ ธมฺมํ ปฏิจฺจ อพฺยากโต อทุกฺขมสุขาย เวทนาย สมฺปยุตฺโต ธมฺโม อุปฺปชฺชติ เหตุปจฺจยา. (1)}

\proseitem{10}{เหตุยา ตีณิ, อารมฺมเณ ตีณิ ฯเปฯ อวิคเต ตีณิ. (สํขิตฺตํ.\csromanspacer{} สหชาตวารมฺปิ ฯเปฯ สมฺปยุตฺตวารมฺปิ ปฏิจฺจวารสทิสํ.)}

\proseitem{11}{กุสโล อทุกฺขมสุขาย เวทนาย สมฺปยุตฺโต ธมฺโม กุสลสฺส อทุกฺขมสุขาย เวทนาย สมฺปยุตฺตสฺส ธมฺมสฺส}

\prose{อกุสโล อทุกฺขมสุขาย เวทนาย สมฺปยุตฺโต ธมฺโม อกุสลสฺส อทุกฺขมสุขาย เวทนาย สมฺปยุตฺตสฺส ธมฺมสฺส เหตุปจฺจเยน ปจฺจโย. (1)}

\prose{อพฺยากโต อทุกฺขมสุขาย เวทนาย สมฺปยุตฺโต ธมฺโม อพฺยากตสฺส อทุกฺขมสุขาย เวทนาย สมฺปยุตฺตสฺส ธมฺมสฺส เหตุปจฺจเยน ปจฺจโย. (1) (สํขิตฺตํ.)}

\proseitem{12}{เหตุยา ตีณิ, อารมฺมเณ นว, อธิปติยา สตฺต, อนนฺตเร สตฺต ฯเปฯ อุปนิสฺสเย นว, อวิคเต ตีณิ. (สํขิตฺตํ.\csromanspacer{} ยถา กุสลตฺติเก ปญฺหาวารํ, เอวํ วิตฺถาเรตพฺพํ.)}

\csromanheader{2}{1. กุสลตฺติก 2. วิปากตฺติก}
\proseitem{13}{\csromanfolio{4}อพฺยากตํ วิปากํ ธมฺมํ ปฏิจฺจ อพฺยากโต วิปาโก ธมฺโม อุปฺปชฺชติ เหตุปจฺจยา. (สํขิตฺตํ.)}

\prose{เหตุยา เอกํ, อารมฺมเณ เอกํ ฯเปฯ อวิคเต เอกํ. (สํขิตฺตํ.) (สหชาตวาเรปิ ฯเปฯ ปญฺหาวาเรปิ สพฺพตฺถ เอกํ.)}

\proseitem{14}{กุสลํ วิปากธมฺมธมฺมํ ปฏิจฺจ กุสโล วิปากธมฺมธมฺโม อุปฺปชฺชติ เหตุปจฺจยา. (1)}

\prose{อกุสลํ วิปากธมฺมธมฺมํ ปฏิจฺจ อกุสโล วิปากธมฺมธมฺโม อุปฺปชฺชติ เหตุปจฺจยา. (1) (สํขิตฺตํ.)}

\prose{เหตุยา เทฺว, อารมฺมเณ เทฺว ฯเปฯ อวิคเต เทฺว. (สํขิตฺตํ.)}

\prose{(สหชาตวารมฺปิ ฯเปฯ สมฺปยุตฺตวารมฺปิ ปฏิจฺจวารสทิสํ.)}

\proseitem{15}{กุสโล วิปากธมฺมธมฺโม กุสลสฺส วิปากธมฺมธมฺมสฺส เหตุปจฺจเยน ปจฺจโย. (1)}

\csromanheader{2}{อกุสโล วิปากธมฺมธมฺโม อกุสลสฺส วิปากธมฺมธมฺมสฺส}
\prose{กุสโล วิปากธมฺมธมฺโม กุสลสฺส วิปากธมฺมธมฺมสฺส อารมฺมณปจฺจเยน ปจฺจโย.\csromanspacer{}\csromanspacer{}กุสโล วิปากธมฺมธมฺโม อกุสลสฺส วิปากธมฺมธมฺมสฺส อารมฺมณปจฺจเยน ปจฺจโย. (2)}

\prose{อกุสโล วิปากธมฺมธมฺโม อกุสลสฺส วิปากธมฺมธมฺมสฺส อารมฺมณปจฺจเยน ปจฺจโย.\csromanspacer{}\csromanspacer{}อกุสโล วิปากธมฺมธมฺโม กุสลสฺส วิปากธมฺมธมฺมสฺส อารมฺมณปจฺจเยน ปจฺจโย. (2) (สํขิตฺตํ.)}

\proseitem{16}{เหตุยา เทฺว, อารมฺมเณ จตฺตาริ, อธิปติยา ตีณิ, อนนฺตเร เทฺว ฯเปฯ สหชาเต เทฺว, อุปนิสฺสเย จตฺตาริ ฯเปฯ อวิคเต เทฺว. (สํขิตฺตํ.\csromanspacer{} ยถา กุสลตฺติเก ปญฺหาวารํ, เอวํ วิตฺถาเรตพฺพํ.)}

\csromanheader{2}{1. กุสลตฺติก 3. อุปาทินฺนตฺติก}
\proseitem{17}{\csromanfolio{5}อพฺยากตํ เนววิปากนวิปากธมฺมธมฺมํ ปฏิจฺจ อพฺยากโต เนว วิปากนวิปากธมฺมธมฺโม อุปฺปชฺชติ เหตุปจฺจยา. (สํขิตฺตํ.)}

\prose{เหตุยา เอกํ, อารมฺมเณ เอกํ ฯเปฯ อวิคเต เอกํ. (สํขิตฺตํ.)}

\vspace{\tipitakaparskip}
\csromancenter{(สหชาตวาเรปิ ฯเปฯ ปญฺหาวาเรปิ สพฺพตฺถ เอกํ.)}
\csromanheader{2}{1. กุสลตฺติก 3. อุปาทินฺนตฺติก}
\proseitem{18}{อพฺยากตํ อุปาทินฺนุปาทานิยํ ธมฺมํ ปฏิจฺจ อพฺยากโต อุปาทินฺนุปาทานิโย ธมฺโม อุปฺปชฺชติ เหตุปจฺจยา. (สํขิตฺตํ.)}

\csromanheader{2}{เหตุยา เอกํ ฯเปฯ อวิคเต เอกํ. (สํขิตฺตํ.)}
\vspace{\tipitakaparskip}
\csromancenter{(สหชาตวาเรปิ ฯเปฯ ปญฺหาวาเรปิ สพฺพตฺถ เอกํ.)}
\proseitem{19}{กุสลํ อนุปาทินฺนุปาทานิยํ ธมฺมํ ปฏิจฺจ กุสโล อนุปาทินฺนุปาทานิโย ธมฺโม อุปฺปชฺชติ เหตุปจฺจยา.\csromanspacer{} ตีณิ.}

\prose{อกุสลํ อนุปาทินฺนุปาทานิยํ ธมฺมํ ปฏิจฺจ อกุสโล อนุปาทินฺนุปาทานิโย ธมฺโม อุปฺปชฺชติ เหตุปจฺจยา.\csromanspacer{} ตีณิ.}

\prose{อพฺยากตํ อนุปาทินฺนุปาทานิยํ ธมฺมํ ปฏิจฺจ อพฺยากโต อนุปาทินฺนุปาทานิโย ธมฺโม อุปฺปชฺชติ เหตุปจฺจยา. (1)}

\prose{กุสลํ อนุปาทินฺนุปาทานิยญฺจ อพฺยากตํ อนุปาทินฺนุปาทานิยญฺจ ธมฺมํ ปฏิจฺจ อพฺยากโต อนุปาทินฺนุปาทานิโย ธมฺโม อุปฺปชฺชติ}

\prose{อกุสลํ อนุปาทินฺนุปาทานิยญฺจ อพฺยากตํ อนุปาทินฺนุปาทานิยญฺจ ธมฺมํ ปฏิจฺจ อพฺยากโต อนุปาทินฺนุปาทานิโย ธมฺโม อุปฺปชฺชติ เหตุปจฺจยา. (1) (สํขิตฺตํ.)}

\proseitem{20}{เหตุยา นว, อวิคเต นว. (สํขิตฺตํ.\csromanspacer{} สหชาตวารมฺปิ ฯเปฯ สมฺปยุตฺตวารมฺปิ ปฏิจฺจวารสทิสํ.)}

\proseitem{21}{\csromanfolio{6}กุสโล อนุปาทินฺนุปาทานิโย ธมฺโม กุสลสฺส อนุปาทินฺนุปาทานิยสฺส ธมฺมสฺส เหตุปจฺจเยน ปจฺจโย.\csromanspacer{} ตีณิ.}

\prose{อกุสโล อนุปาทินฺนุปาทานิโย ธมฺโม อกุสลสฺส อนุปาทินฺนุปาทานิยสฺส ธมฺมสฺส เหตุปจฺจเยน ปจฺจโย.\csromanspacer{} ตีณิ.}

\prose{อพฺยากโต อนุปาทินฺนุปาทานิโย ธมฺโม อพฺยากตสฺส อนุปาทินฺนุปาทานิยสฺส ธมฺมสฺส เหตุปจฺจเยน ปจฺจโย. (1)}

\proseitem{22}{กุสโล อนุปาทินฺนุปาทานิโย ธมฺโม กุสลสฺส อนุปาทินฺนุปาทานิยสฺส ธมฺมสฺส อารมฺมณปจฺจเยน ปจฺจโย.\csromanspacer{} ตีณิ.}

\prose{อกุสโล อนุปาทินฺนุปาทานิโย ธมฺโม อกุสลสฺส อนุปาทินฺนุปาทานิยสฺส ธมฺมสฺส อารมฺมณปจฺจเยน ปจฺจโย.\csromanspacer{} ตีณิ.}

\prose{อพฺยากโต อนุปาทินฺนุปาทานิโย ธมฺโม อพฺยากตสฺส อนุปาทินฺนุปาทานิยสฺส ธมฺมสฺส อารมฺมณปจฺจเยน ปจฺจโย.\csromanspacer{} ตีณิ.}

\proseitem{23}{เหตุยา สตฺต, อารมฺมเณ นว, อวิคเต เอกาทส. (สํขิตฺตํ.)}

\prose{(ยถา กุสลตฺติเก ปญฺหาวารํ, เอวํ วิตฺถาเรตพฺพํ.)}

\proseitem{24}{กุสลํ อนุปาทินฺน-อนุปาทานิยํ ธมฺมํ ปฏิจฺจ กุสโล อนุปาทินฺนุปาทานิโย ธมฺโม อุปฺปชฺชติ เหตุปจฺจยา. (1)}

\prose{อพฺยากตํ อนุปาทินฺน-อนุปาทานิยํ ธมฺมํ ปฏิจฺจ อพฺยากโต อนุปาทินฺน-อนุปาทานิโย ธมฺโม อุปฺปชฺชติ เหตุปจฺจยา. (1) (สํขิตฺตํ.)}

\prose{เหตุยา เทฺว, อวิคเต เทฺว. (สํขิตฺตํ.\csromanspacer{} สหชาตวารมฺปิ ฯเปฯ ปญฺหาวารมฺปิ สพฺพตฺถ วิตฺถาโร.)}

\csromanheader{2}{1. กุสลตฺติก 4. สํกิลิฏฺฐตฺติก}
\proseitem{25}{อกุสลํ สํกิลิฏฺฐสํกิเลสิกํ ธมฺมํ ปฏิจฺจ อกุสโล สํกิลิฏฺฐสํกิเลสิโก ธมฺโม อุปฺปชฺชติ เหตุปจฺจยา. (สํขิตฺตํ.)}

\csromanheader{2}{1. กุสลตฺติก 5. วิตกฺกตฺติก}
\csromanheader{2}{เหตุยา เอกํ ฯเปฯ อวิคเต เอกํ. (สํขิตฺตํ.)}
\prose{\csromanfolio{7}(สหชาตวาเรปิ ฯเปฯ ปญฺหาวาเรปิ สพฺพตฺถ เอกํ.)}

\proseitem{26}{กุสลํ อสํกิลิฏฺฐสํกิเลสิกํ ธมฺมํ ปฏิจฺจ กุสโล อสํกิลิฏฺฐสํกิเลสิโก ธมฺโม อุปฺปชฺชติ เหตุปจฺจยา.\csromanspacer{} ตีณิ.}

\prose{อพฺยากตํ อสํกิลิฏฺฐสํกิเลสิกํ ธมฺมํ ปฏิจฺจ อพฺยากโต อสํกิลิฏฺฐสํกิเลสิโก ธมฺโม อุปฺปชฺชติ เหตุปจฺจยา. (1)}

\prose{กุสลํ อสํกิลิฏฺฐสํกิเลสิกญฺจ อพฺยากตํ อสํกิลิฏฺฐสํกิเลสิกญฺจ ธมฺมํ ปฏิจฺจ อพฺยากโต อสํกิลิฏฺฐสํกิเลสิโก ธมฺโม อุปฺปชฺชติ เหตุปจฺจยา. (1)}

\proseitem{27}{เหตุยา ปญฺจ, อวิคเต ปญฺจ. (สํขิตฺตํ.\csromanspacer{} สหชาตวาเรปิ ฯเปฯ ปญฺหาวาเรปิ สพฺพตฺถ วิตฺถาโร.)}

\proseitem{28}{กุสลํ อสํกิลิฏฺฐ-อสํกิเลสิกํ ธมฺมํ ปฏิจฺจ กุสโล อสํกิลิฏฺฐ-อสํกิเลสิโก ธมฺโม อุปฺปชฺชติ เหตุปจฺจยา. (1)}

\prose{อพฺยากตํ อสํกิลิฏฺฐ-อสํกิเลสิกํ ธมฺมํ ปฏิจฺจ อพฺยากโต อสํกิลิฏฺฐ-อสํกิเลสิโก ธมฺโม อุปฺปชฺชติ เหตุปจฺจยา. (1) (สํขิตฺตํ.)}

\prose{เหตุยา เทฺว, อวิคเต เทฺว. (สํขิตฺตํ.\csromanspacer{} สหชาตวาเรปิ ฯเปฯ ปญฺหาวาเรปิ สพฺพตฺถ วิตฺถาเรตพฺพํ.)}

\csromanheader{2}{1. กุสลตฺติก 5. วิตกฺกตฺติก}
\csromanheader{2}{1-7. ปฏิจฺจาทิวาร ปจฺจยจตุกฺก-เหตุ}
\proseitem{29}{กุสลํ สวิตกฺกสวิจารํ ธมฺมํ ปฏิจฺจ กุสโล สวิตกฺกสวิจาโร ธมฺโม อุปฺปชฺชติ เหตุปจฺจยา. (1)}

\prose{อกุสลํ สวิตกฺกสวิจารํ ธมฺมํ ปฏิจฺจ อกุสโล สวิตกฺกสวิจาโร ธมฺโม อุปฺปชฺชติ เหตุปจฺจยา. (1)}

\prose{\csromanfolio{8}อพฺยากตํ สวิตกฺกสวิจารํ ธมฺมํ ปฏิจฺจ อพฺยากโต สวิตกฺกสวิจาโร ธมฺโม อุปฺปชฺชติ เหตุปจฺจยา. (1) (สํขิตฺตํ.)}

\prose{เหตุยา ตีณิ, อารมฺมเณ ตีณิ, อวิคเต ตีณิ. (สํขิตฺตํ.\csromanspacer{} สหชาตวารมฺปิ ฯเปฯ สมฺปยุตฺตวารมฺปิ ปฏิจฺจวารสทิสํ.)}

\proseitem{30}{กุสโล สวิตกฺกสวิจาโร ธมฺโม กุสลสฺส สวิตกฺกสวิจารสฺส ธมฺมสฺส เหตุปจฺจเยน ปจฺจโย. (1)}

\prose{อกุสโล สวิตกฺกสวิจาโร ธมฺโม อกุสลสฺส สวิตกฺกสวิจารสฺส ธมฺมสฺส เหตุปจฺจเยน ปจฺจโย. (1)}

\prose{อพฺยากโต สวิตกฺกสวิจาโร ธมฺโม อพฺยากตสฺส สวิตกฺกสวิจารสฺส ธมฺมสฺส เหตุปจฺจเยน ปจฺจโย. (1) (สํขิตฺตํ.)}

\proseitem{31}{เหตุยา ตีณิ, อารมฺมเณ นว, อวิคเต ตีณิ. (สํขิตฺตํ.) (ยถา กุสลตฺติเก ปญฺหาวารํ, เอวํ วิตฺถาเรตพฺพํ.)}

\proseitem{32}{กุสลํ อวิตกฺกวิจารมตฺตํ ธมฺมํ ปฏิจฺจ กุสโล อวิตกฺกวิจารมตฺโต ธมฺโม อุปฺปชฺชติ เหตุปจฺจยา. (1)}

\prose{อพฺยากตํ อวิตกฺกวิจารมตฺตํ ธมฺมํ ปฏิจฺจ อพฺยากโต อวิตกฺกวิจารมตฺโต ธมฺโม อุปฺปชฺชติ เหตุปจฺจยา. (1) (สํขิตฺตํ.)}

\prose{เหตุยา เทฺว, อวิคเต เทฺว. (สํขิตฺตํ.)}

\prose{(สหชาตวาเรปิ ฯเปฯ สมฺปยุตฺตวาเรปิ ปญฺหาวาเรปิ สพฺพตฺถ วิตฺถาโร.)}

\proseitem{33}{กุสลํ อวิตกฺก-อวิจารํ ธมฺมํ ปฏิจฺจ กุสโล อวิตกฺก-อวิจาโร ธมฺโม อุปฺปชฺชติ เหตุปจฺจยา.\csromanspacer{} ตีณิ.}

\prose{อพฺยากตํ อวิตกฺก-อวิจารํ ธมฺมํ ปฏิจฺจ อพฺยากโต อวิตกฺก-อวิจาโร ธมฺโม อุปฺปชฺชติ เหตุปจฺจยา. (1)}

\prose{กุสลํ อวิตกฺก-อวิจารญฺจ อพฺยากตํ อวิตกฺก-อวิจารญฺจ ธมฺมํ ปฏิจฺจ อพฺยากโต อวิตกฺก-อวิจาโร ธมฺโม อุปฺปชฺชติ เหตุปจฺจยา. (1) (สํขิตฺตํ.)}

\csromanheader{2}{1. กุสลตฺติก 6. ปีติตฺติก}
\prose{\csromanfolio{9}เหตุยา ปญฺจ, อารมฺมเณ เทฺว, อวิคเต ปญฺจ. (สํขิตฺตํ.\csromanspacer{} สหชาตวารมฺปิ ฯเปฯ ปญฺหาวารมฺปิ วิตฺถาเรตพฺพํ.)}

\csromanheader{2}{1. กุสลตฺติก 6. ปีติตฺติก}
\csromanheader{2}{1-7. ปฏิจฺจาทิวาร ปจฺจยจตุกฺก-เหตุ}
\csromanheader{2}{34. กุสลํ ปีติสหคตํ ธมฺมํ ปฏิจฺจ กุสโล ปีติสหคโต}
\csromanheader{2}{อกุสลํ ปีติสหคตํ ธมฺมํ ปฏิจฺจ อกุสโล ปีติสหคโต}
\prose{อพฺยากตํ ปีติสหคตํ ธมฺมํ ปฏิจฺจ อพฺยากโต ปีติสหคโต ธมฺโม อุปฺปชฺชติ เหตุปจฺจยา. (1) (สํขิตฺตํ.)}

\prose{เหตุยา ตีณิ, อารมฺมเณ ตีณิ, อวิคเต ตีณิ. (สํขิตฺตํ.\csromanspacer{} สหชาตวารมฺปิ ฯเปฯ สมฺปยุตฺตวารมฺปิ ปฏิจฺจวารสทิสํ.)}

\proseitem{35}{กุสโล ปีติสหคโต ธมฺโม กุสลสฺส ปีติสหคตสฺส ธมฺมสฺส}

\prose{อกุสโล ปีติสหคโต ธมฺโม อกุสลสฺส ปีติสหคตสฺส ธมฺมสฺส}

\prose{อพฺยากโต ปีติสหคโต ธมฺโม อพฺยากตสฺส ปีติสหคตสฺส ธมฺมสฺส เหตุปจฺจเยน ปจฺจโย. (1) (สํขิตฺตํ.)}

\prose{เหตุยา ตีณิ, อารมฺมเณ นว, อธิปติยา สตฺต, อนนฺตเร ปญฺจ, อวิคเต ตีณิ. (สํขิตฺตํ.) (ยถา กุสลตฺติเก ปญฺหาวารํ, เอวํ วิตฺถาเรตพฺพํ.)}

\csromanheader{2}{36. กุสลํ สุขสหคตํ ธมฺมํ ปฏิจฺจ กุสโล สุขสหคโต}
\prose{อกุสลํ สุขสหคตํ ธมฺมํ ปฏิจฺจ อกุสโล สุขสหคโต}

\prose{\csromanfolio{10}อพฺยากตํ สุขสหคตํ ธมฺมํ ปฏิจฺจ อพฺยากโต สุขสหคโต ธมฺโม อุปฺปชฺชติ เหตุปจฺจยา. (1) (สํขิตฺตํ.)}

\prose{เหตุยา ตีณิ, อารมฺมเณ ตีณิ, อวิคเต ตีณิ. (สํขิตฺตํ.\csromanspacer{} สหชาตวารมฺปิ ฯเปฯ ปญฺหาวารมฺปิ วิตฺถาเรตพฺพํ.)}

\proseitem{37}{กุสลํ อุเปกฺขาสหคตํ ธมฺมํ ปฏิจฺจ กุสโล อุเปกฺขาสหคโต ธมฺโม อุปฺปชฺชติ เหตุปจฺจยา. (1)}

\prose{อกุสลํ อุเปกฺขาสหคตํ ธมฺมํ ปฏิจฺจ อกุสโล อุเปกฺขาสหคโต ธมฺโม อุปฺปชฺชติ เหตุปจฺจยา. (1)}

\prose{อพฺยากตํ อุเปกฺขาสหคตํ ธมฺมํ ปฏิจฺจ อพฺยากโต อุเปกฺขาสหคโต ธมฺโม อุปฺปชฺชติ เหตุปจฺจยา. (1) (สํขิตฺตํ.)}

\prose{เหตุยา ตีณิ, อวิคเต ตีณิ. (สํขิตฺตํ.)}

\vspace{\tipitakaparskip}
\csromancenter{(สหชาตวารมฺปิ ฯเปฯ ปญฺหาวารมฺปิ วิตฺถาเรตพฺพํ.)}
\csromanheader{2}{1. กุสลตฺติก 7. ทสฺสนตฺติก}
\proseitem{38}{อกุสลํ ทสฺสเนน ปหาตพฺพํ ธมฺมํ ปฏิจฺจ อกุสโล ทสฺสเนน ปหาตพฺโพ ธมฺโม อุปฺปชฺชติ เหตุปจฺจยา. (สํขิตฺตํ.)}

\prose{เหตุยา เอกํ ฯเปฯ อวิคเต เอกํ. (สํขิตฺตํ.) (สหชาตวาเรปิ ฯเปฯ ปญฺหาวาเรปิ สพฺพตฺถ เอกํ.)}

\proseitem{39}{อกุสลํ ภาวนาย ปหาตพฺพํ ธมฺมํ ปฏิจฺจ อกุสโล ภาวนาย ปหาตพฺโพ ธมฺโม อุปฺปชฺชติ เหตุปจฺจยา. (1) (สํขิตฺตํ.)}

\csromanheader{2}{เหตุยา เอกํ ฯเปฯ อวิคเต เอกํ. (สํขิตฺตํ.)}
\proseitem{40}{กุสลํ เนวทสฺสเนน นภาวนาย ปหาตพฺพํ ธมฺมํ ปฏิจฺจ กุสโล เนวทสฺสเนน นภาวนาย ปหาตพฺโพ ธมฺโม อุปฺปชฺชติ เหตุปจฺจยา.\csromanspacer{} ตีณิ.}

\csromanheader{2}{1. กุสลตฺติก 8. ทสฺสนเหตุกตฺติก}
\prose{\csromanfolio{11}อพฺยากตํ เนวทสฺสเนน นภาวนาย ปหาตพฺพํ ธมฺมํ ปฏิจฺจ อพฺยากโต เนวทสฺสเนน นภาวนาย ปหาตพฺโพ ธมฺโม อุปฺปชฺชติ}

\prose{กุสลํ เนวทสฺสเนน นภาวนาย ปหาตพฺพญฺจ อพฺยากตํ เนวทสฺสเนน นภาวนาย ปหาตพฺพญฺจ ธมฺมํ ปฏิจฺจ อพฺยากโต เนวทสฺสเนน นภาวนาย ปหาตพฺโพ ธมฺโม อุปฺปชฺชติ เหตุปจฺจยา. (1)}

\prose{เหตุยา ปญฺจ, อวิคเต ปญฺจ. (สํขิตฺตํ.\csromanspacer{} สหชาตวาเรปิ ฯเปฯ ปญฺหาวาเรปิ สพฺพตฺถ วิตฺถาโร.)}

\csromanheader{2}{1. กุสลตฺติก 8. ทสฺสนเหตุกตฺติก}
\proseitem{41}{อกุสลํ ทสฺสเนน ปหาตพฺพเหตุกํ ธมฺมํ ปฏิจฺจ อกุสโล ทสฺสเนน ปหาตพฺพเหตุโก ธมฺโม อุปฺปชฺชติ เหตุปจฺจยา.}

\csromanheader{2}{เหตุยา เอกํ ฯเปฯ อวิคเต เอกํ. (สํขิตฺตํ.)}
\proseitem{42}{อกุสลํ ภาวนาย ปหาตพฺพเหตุกํ ธมฺมํ ปฏิจฺจ อกุสโล ภาวนาย ปหาตพฺพเหตุโก ธมฺโม อุปฺปชฺชติ เหตุปจฺจยา.}

\csromanheader{2}{เหตุยา เอกํ ฯเปฯ อวิคเต เอกํ. (สํขิตฺตํ.)}
\proseitem{43}{กุสลํ เนวทสฺสเนน นภาวนาย ปหาตพฺพเหตุกํ ธมฺมํ ปฏิจฺจ กุสโล เนวทสฺสเนน นภาวนาย ปหาตพฺพเหตุโก ธมฺโม อุปฺปชฺชติ เหตุปจฺจยา. (สํขิตฺตํ.)}

\prose{เหตุยา สตฺต, อารมฺมเณ เทฺว, อวิคเต สตฺต. (สํขิตฺตํ.\csromanspacer{} สหชาตวาเรปิ ฯเปฯ ปญฺหาวาเรปิ สพฺพตฺถ วิตฺถาโร.)}

\csromanheader{2}{1. กุสลตฺติก 9. อาจยคามิตฺติก}
\csromanheader{2}{44. กุสลํ อาจยคามิํ ธมฺมํ ปฏิจฺจ กุสโล อาจยคามี}
\prose{\csromanfolio{12}อกุสลํ อาจยคามิํ ธมฺม ปฏิจฺจ อกุสโล อาจยคามี ธมฺโม อุปฺปชฺชติ เหตุปจฺจยา. (1) (สํขิตฺตํ.)}

\prose{เหตุยา เทฺว ฯเปฯ อวิคเต เทฺว. (สํขิตฺตํ.) (สหชาตวาเรปิ ฯเปฯ ปญฺหาวาเรปิ สพฺพตฺถ วิตฺถาโร.)}

\proseitem{45}{กุสลํ อปจยคามิํ ธมฺมํ ปฏิจฺจ กุสโล อปจยคามี ธมฺโม อุปฺปชฺชติ เหตุปจฺจยา. (สํขิตฺตํ.\csromanspacer{} ปฏิจฺจวาเรปิ ฯเปฯ ปญฺหาวาเรปิ สพฺพตฺถ เอกํ.)}

\proseitem{46}{อพฺยากตํ เนวาจยคามินาปจยคามิํ ธมฺมํ ปฏิจฺจ อพฺยากโต เนวาจยคามินาปจยคามี ธมฺโม อุปฺปชฺชติ เหตุปจฺจยา.}

\prose{เหตุยา เอกํ ฯเปฯ อวิคเต เอกํ. (สํขิตฺตํ.) (สหชาตวารมฺปิ ฯเปฯ ปญฺหาวารมฺปิ วิตฺถาเรตพฺพํ.)}

\csromanheader{2}{1. กุสลตฺติก 10. เสกฺขตฺติก}
\proseitem{47}{กุสลํ เสกฺขํ ธมฺมํ ปฏิจฺจ กุสโล เสกฺโข ธมฺโม อุปฺปชฺชติ เหตุปจฺจยา. (1)}

\prose{อพฺยากตํ เสกฺขํ ธมฺมํ ปฏิจฺจ อพฺยากโต เสกฺโข ธมฺโม อุปฺปชฺชติ เหตุปจฺจยา. (1) (สํขิตฺตํ.)}

\prose{เหตุยา เทฺว, อวิคเต เทฺว. (สํขิตฺตํ.\csromanspacer{} สหชาตวาเรปิ ฯเปฯ ปญฺหาวาเรปิ สพฺพตฺถ วิตฺถาโร.)}

\csromanheader{2}{1. กุสลตฺติก 11. ปริตฺตตฺติก}
\proseitem{48}{\csromanfolio{13}อพฺยากตํ อเสกฺขํ ธมฺมํ ปฏิจฺจ อพฺยากโต อเสกฺโข ธมฺโม อุปฺปชฺชติ เหตุปจฺจยา. (สํขิตฺตํ.)}

\prose{เหตุยา เอกํ ฯเปฯ อวิคเต เอกํ. (สํขิตฺตํ.) (สหชาตวาเรปิ ฯเปฯ ปญฺหาวาเรปิ สพฺพตฺถ เอกํ.)}

\proseitem{49}{กุสลํ เนวเสกฺขนาเสกฺขํ ธมฺมํ ปฏิจฺจ กุสโล เนวเสกฺขนาเสกฺโข ธมฺโม อุปฺปชฺชติ เหตุปจฺจยา.\csromanspacer{} ตีณิ.}

\prose{อกุสลํ เนวเสกฺขนาเสกฺขํ ธมฺมํ ปฏิจฺจ อกุสโล เนวเสกฺขนาเสกฺโข ธมฺโม อุปฺปชฺชติ เหตุปจฺจยา.\csromanspacer{} ตีณิ.}

\prose{อพฺยากตํ เนวเสกฺขนาเสกฺขํ ธมฺมํ ปฏิจฺจ อพฺยากโต เนวเสกฺขนาเสกฺโข ธมฺโม อุปฺปชฺชติ เหตุปจฺจยา.\csromanspacer{} ตีณิ.}

\prose{เหตุยา นว, อวิคเต นว. (สํขิตฺตํ.\csromanspacer{} สหชาตวารมฺปิ ฯเปฯ สมฺปยุตฺตวารมฺปิ ปฏิจฺจวารสทิสํ.)}

\proseitem{50}{กุสโล เนวเสกฺขนาเสกฺโข ธมฺโม กุสลสฺส เนวเสกฺขนาเสกฺขสฺส ธมฺมสฺส เหตุปจฺจเยน ปจฺจโย. (สํขิตฺตํ.)}

\prose{เหตุยา สตฺต, อารมฺมเณ นว, อวิคเต เตรส. (สํขิตฺตํ.\csromanspacer{} ยถา กุสลตฺติเก ปญฺหาวารํ, เอวํ วิตฺถาเรตพฺพํ.)}

\csromanheader{2}{1. กุสลตฺติก 11. ปริตฺตตฺติก}
\proseitem{51}{กุสลํ ปริตฺตํ ธมฺมํ ปฏิจฺจ กุสโล ปริตฺโต ธมฺโม อุปฺปชฺชติ เหตุปจฺจยา.\csromanspacer{}\csromanspacer{}กุสลํ ปริตฺตํ ธมฺมํ ปฏิจฺจ อพฺยากโต ปริตฺโต ธมฺโม อุปฺปชฺชติ เหตุปจฺจยา.\csromanspacer{}\csromanspacer{}กุสลํ ปริตฺตํ ธมฺมํ ปฏิจฺจ กุสโล ปริตฺโต จ อพฺยากโต ปริตฺโต จ ธมฺมา อุปฺปชฺชนฺติ เหตุปจฺจยา. (3)}

\prose{อกุสลํ ปริตฺตํ ธมฺมํ ปฏิจฺจ อกุสโล ปริตฺโต ธมฺโม อุปฺปชฺชติ เหตุปจฺจยา.\csromanspacer{} ตีณิ.}

\prose{อพฺยากตํ ปริตฺตํ ธมฺมํ ปฏิจฺจ อพฺยากโต ปริตฺโต ธมฺโม อุปฺปชฺชติ เหตุปจฺจยา. (1)}

\prose{\csromanfolio{14}กุสลํ ปริตฺตญฺจ อพฺยากตํ ปริตฺตญฺจ ธมฺมํ ปฏิจฺจ อพฺยากโต ปริตฺโต ธมฺโม อุปฺปชฺชติ เหตุปจฺจยา. (1)}

\prose{อกุสลํ ปริตฺตญฺจ อพฺยากตํ ปริตฺตญฺจ ธมฺมํ ปฏิจฺจ อพฺยากโต ปริตฺโต ธมฺโม อุปฺปชฺชติ เหตุปจฺจยา. (1) (สํขิตฺตํ.)}

\prose{เหตุยา นว, อารมฺมเณ ตีณิ, อวิคเต นว. (สํขิตฺตํ.\csromanspacer{} สหชาตวารมฺปิ ฯเปฯ สมฺปยุตฺตวารมฺปิ ปฏิจฺจวารสทิสํ.)}

\proseitem{52}{กุสโล ปริตฺโต ธมฺโม กุสลสฺส ปริตฺตสฺส ธมฺมสฺส เหตุปจฺจเยน ปจฺจโย.\csromanspacer{} ตีณิ.}

\prose{อกุสโล ปริตฺโต ธมฺโม อกุสลสฺส ปริตฺตสฺส ธมฺมสฺส เหตุปจฺจเยน ปจฺจโย.\csromanspacer{} ตีณิ.}

\csromanheader{2}{อพฺยากโต ปริตฺโต ธมฺโม อพฺยากตสฺส ปริตฺตสฺส ธมฺมสฺส}
\prose{กุสโล ปริตฺโต ธมฺโม กุสลสฺส ปริตฺตสฺส ธมฺมสฺส อารมฺมณปจฺจเยน ปจฺจโย. (สํขิตฺตํ.)}

\prose{เหตุยา สตฺต, อารมฺมเณ นว, อวิคเต เตรส. (สํขิตฺตํ.\csromanspacer{} ยถา กุสลตฺติเก ปญฺหาวารํ, เอวํ วิตฺถาเรตพฺพํ.)}

\csromanheader{2}{\textbf{มหคฺคตาทิปท-เหตุ}}
\csromanheader{2}{53. กุสลํ มหคฺคตํ ธมฺมํ ปฏิจฺจ กุสโล มหคฺคโต}
\prose{อพฺยากตํ มหคฺคตํ ธมฺมํ ปฏิจฺจ อพฺยากโต มหคฺคโต ธมฺโม อุปฺปชฺชติ เหตุปจฺจยา. (1) (สํขิตฺตํ.)}

\prose{เหตุยา เทฺว, อารมฺมเณ เทฺว, อวิคเต เทฺว. (สํขิตฺตํ.\csromanspacer{} สหชาตวาเรปิ ฯเปฯ ปญฺหาวาเรปิ สพฺพตฺถ วิตฺถาโร.)}

\proseitem{54}{กุสลํ อปฺปมาณํ ธมฺมํ ปฏิจฺจ กุสโล อปฺปมาโณ ธมฺโม อุปฺปชฺชติ เหตุปจฺจยา. (1)}

\prose{อพฺยากตํ อปฺปมาณํ ธมฺมํ ปฏิจฺจ อพฺยากโต อปฺปมาโณ ธมฺโม อุปฺปชฺชติ เหตุปจฺจยา. (1) (สํขิตฺตํ.)}

\csromanheader{2}{1. กุสลตฺติก 12. ปริตฺตารมฺมณตฺติก}
\prose{\csromanfolio{15}เหตุยา เทฺว, อวิคเต เทฺว. (สํขิตฺตํ.\csromanspacer{} สหชาตวาเรปิ ฯเปฯ ปญฺหาวาเรปิ สพฺพตฺถ วิตฺถาโร.)}

\csromanheader{2}{1. กุสลตฺติก 12. ปริตฺตารมฺมณตฺติก}
\proseitem{55}{กุสลํ ปริตฺตารมฺมณํ ธมฺมํ ปฏิจฺจ กุสโล ปริตฺตารมฺมโณ ธมฺโม อุปฺปชฺชติ เหตุปจฺจยา. (1)}

\prose{อกุสลํ ปริตฺตารมฺมณํ ธมฺมํ ปฏิจฺจ อกุสโล ปริตฺตารมฺมโณ}

\prose{อพฺยากตํ ปริตฺตารมฺมณํ ธมฺมํ ปฏิจฺจ อพฺยากโต ปริตฺตารมฺมโณ ธมฺโม อุปฺปชฺชติ เหตุปจฺจยา. (1) (สํขิตฺตํ.)}

\prose{เหตุยา ตีณิ, อวิคเต ตีณิ. (สํขิตฺตํ.) (สหชาตวารมฺปิ ฯเปฯ ปญฺหาวารมฺปิ วิตฺถาเรตพฺพํ.)}

\proseitem{56}{กุสลํ มหคฺคตารมฺมณํ ธมฺมํ ปฏิจฺจ กุสโล มหคฺคตารมฺมโณ ธมฺโม อุปฺปชฺชติ เหตุปจฺจยา. (1)}

\prose{อกุสลํ มหคฺคตารมฺมณํ ธมฺมํ ปฏิจฺจ อกุสโล มหคฺคตารมฺมโณ ธมฺโม อุปฺปชฺชติ เหตุปจฺจยา. (1)}

\prose{อพฺยากตํ มหคฺคตารมฺมณํ ธมฺมํ ปฏิจฺจ อพฺยากโต มหคฺคตารมฺมโณ ธมฺโม อุปฺปชฺชติ เหตุปจฺจยา. (1) (สํขิตฺตํ.)}

\prose{เหตุยา ตีณิ, อวิคเต ตีณิ. (สํขิตฺตํ.) (สหชาตวาเรปิ ฯเปฯ ปญฺหาวาเรปิ สพฺพตฺถ วิตฺถาโร.)}

\proseitem{57}{กุสลํ อปฺปมาณารมฺมณํ ธมฺมํ ปฏิจฺจ กุสโล อปฺปมาณารมฺมโณ ธมฺโม อุปฺปชฺชติ เหตุปจฺจยา. (1)}

\prose{อพฺยากตํ อปฺปมาณารมฺมณํ ธมฺมํ ปฏิจฺจ อพฺยากโต อปฺปมาณารมฺมโณ ธมฺโม อุปฺปชฺชติ เหตุปจฺจยา. (1) (สํขิตฺตํ.)}

\prose{เหตุยา เทฺว, อวิคเต เทฺว. (สํขิตฺตํ.) (สหชาตวาเรปิ ฯเปฯ ปญฺหาวาเรปิ สพฺพตฺถ วิตฺถาโร.)}

\csromanheader{2}{1. กุสลตฺติก 13. หีนตฺติก}
\proseitem{58}{\csromanfolio{16}อกุสลํ หีนํ ธมฺมํ ปฏิจฺจ อกุสโล หีโน ธมฺโม อุปฺปชฺชติ เหตุปจฺจยา. (สํขิตฺตํ.)}

\prose{เหตุยา เอกํ ฯเปฯ อวิคเต เอกํ. (สํขิตฺตํ.\csromanspacer{} สพฺพตฺถ วิตฺถาโร.)}

\proseitem{59}{กุสลํ มชฺฌิมํ ธมฺมํ ปฏิจฺจ กุสโล มชฺฌิโม ธมฺโม อุปฺปชฺชติ เหตุปจฺจยา.\csromanspacer{}\csromanspacer{}กุสลํ มชฺฌิมํ ธมฺมํ ปฏิจฺจ อพฺยากโต มชฺฌิโม ธมฺโม อุปฺปชฺชติ เหตุปจฺจยา.\csromanspacer{}\csromanspacer{}กุสลํ มชฺฌิมํ ธมฺมํ ปฏิจฺจ กุสโล มชฺฌิโม จ อพฺยากโต มชฺฌิโม จ ธมฺมา อุปฺปชฺชนฺติ เหตุปจฺจยา. (3)}

\prose{อพฺยากตํ มชฺฌิมํ ธมฺมํ ปฏิจฺจ อพฺยากโต มชฺฌิโม ธมฺโม อุปฺปชฺชติ เหตุปจฺจยา. (1)}

\prose{กุสลํ มชฺฌิมญฺจ อพฺยากตํ มชฺฌิมญฺจ ธมฺมํ ปฏิจฺจ อพฺยากโต มชฺฌิโม ธมฺโม อุปฺปชฺชติ เหตุปจฺจยา. (1)}

\prose{กุสลํ มชฺฌิมํ ธมฺมํ ปฏิจฺจ กุสโล มชฺฌิโม ธมฺโม อุปฺปชฺชติ อารมฺมณปจฺจยา. (1)}

\prose{อพฺยากตํ มชฺฌิมํ ธมฺโม ปฏิจฺจ อพฺยากโต มชฺฌิโม ธมฺโม อุปฺปชฺชติ อารมฺมณปจฺจยา. (1) (สํขิตฺตํ.)}

\prose{เหตุยา ปญฺจ, อารมฺมเณ เทฺว, อวิคเต ปญฺจ. (สํขิตฺตํ.\csromanspacer{} สหชาตวารมฺปิ ฯเปฯ สมฺปยุตฺตวารมฺปิ ปฏิจฺจวารสทิสํ.)}

\proseitem{60}{กุสโล มชฺฌิโม ธมฺโม กุสลสฺส มชฺฌิมสฺส ธมฺมสฺส เหตุปจฺจเยน ปจฺจโย.\csromanspacer{} ตีณิ.}

\csromanheader{2}{อพฺยากโต มชฺฌิโม ธมฺโม อพฺยากตสฺส มชฺฌิมสฺส ธมฺมสฺส}
\prose{กุสโล มชฺฌิโม ธมฺโม กุสลสฺส มชฺฌิมสฺส ธมฺมสฺส อารมฺมณปจฺจเยน ปจฺจโย.\csromanspacer{}\csromanspacer{}กุสโล มชฺฌิโม ธมฺโม อพฺยากตสฺส มชฺฌิมสฺส ธมฺมสฺส อารมฺมณปจฺจเยน ปจฺจโย. (2)}

\vspace{\tipitakaparskip}
\csromancenter{กุสลตฺติก 14.\csromanspacer{} มิจฺฉตฺตนิยตตฺติก อพฺยากโต มชฺฌิโม ธมฺโม อพฺยากตสฺส มชฺฌิมสฺส ธมฺมสฺส อารมฺมณปจฺจเยน ปจฺจโย.\csromanspacer{}\csromanspacer{}อพฺยากโต มชฺฌิโม ธมฺโม กุสลสฺส มชฺฌิมสฺส ธมฺมสฺส อารมฺมณปจฺจเยน ปจฺจโย. (2) (สํขิตฺตํ.)}
\proseitem{61}{\csromanfolio{17}เหตุยา จตฺตาริ, อารมฺมเณ จตฺตาริ, อวิคเต สตฺต. (สํขิตฺตํ.\csromanspacer{} ยถา กุสลตฺติเก ปญฺหาวารํ, เอวํ วิตฺถาเรตพฺพํ.)}

\csromanheader{2}{\textbf{ปณีตปท-เหตุ}}
\proseitem{62}{กุสลํ ปณีตํ ธมฺมํ ปฏิจฺจ กุสโล ปณีโต ธมฺโม อุปฺปชฺชติ เหตุปจฺจยา. (1)}

\prose{อพฺยากตํ ปณีตํ ธมฺมํ ปฏิจฺจ อพฺยากโต ปณีโต ธมฺโม อุปฺปชฺชติ เหตุปจฺจยา. (1) (สํขิตฺตํ.)}

\prose{เหตุยา เทฺว, อวิคเต เทฺว. (สํขิตฺตํ.\csromanspacer{} สหชาตวาเรปิ ฯเปฯ ปญฺหาวาเรปิ สพฺพตฺถ วิตฺถาโร.)}

\csromanheader{2}{1. กุสลตฺติก 14. มิจฺฉตฺตนิยตตฺติก}
\proseitem{63}{อกุสลํ มิจฺฉตฺตนิยตํ ธมฺมํ ปฏิจฺจ อกุสโล มิจฺฉตฺตนิยโต ธมฺโม อุปฺปชฺชติ เหตุปจฺจยา. (สํขิตฺตํ.)}

\prose{เหตุยา เอกํ ฯเปฯ อวิคเต เอกํ. (สํขิตฺตํ.\csromanspacer{} สพฺพตฺถ วิตฺถาโร.)}

\proseitem{64}{กุสลํ สมฺมตฺตนิยตํ ธมฺมํ ปฏิจฺจ กุสโล สมฺมตฺตนิยโต ธมฺโม อุปฺปชฺชติ เหตุปจฺจยา. (สํขิตฺตํ.)}

\prose{เหตุยา เอกํ ฯเปฯ อวิคเต เอกํ. (สํขิตฺตํ.\csromanspacer{} สหชาตวาเรปิ ฯเปฯ ปญฺหาวาเรปิ สพฺพตฺถ เอกํ.)}

\proseitem{65}{กุสลํ อนิยตํ ธมฺมํ ปฏิจฺจ กุสโล อนิยโต ธมฺโม อุปฺปชฺชติ เหตุปจฺจยา.\csromanspacer{}\csromanspacer{}กุสลํ อนิยตํ ธมฺมํ ปฏิจฺจ อพฺยากโต อนิยโต ธมฺโม อุปฺปชฺชติ เหตุปจฺจยา.\csromanspacer{}\csromanspacer{}กุสลํ อนิยตํ ธมฺมํ \csromanfolio{18}ปฏิจฺจ กุสโล อนิยโต จ อพฺยากโต อนิยโต จ ธมฺมา อุปฺปชฺชนฺติ เหตุปจฺจยา. (3)}

\prose{อกุสลํ อนิยตํ ธมฺมํ ปฏิจฺจ อกุสโล อนิยโต ธมฺโม อุปฺปชฺชติ เหตุปจฺจยา.\csromanspacer{} ตีณิ.}

\prose{อพฺยากตํ อนิยตํ ธมฺมํ ปฏิจฺจ อพฺยากโต อนิยโต ธมฺโม อุปฺปชฺชติ เหตุปจฺจยา. (1)}

\prose{กุสลํ อนิยตญฺจ อพฺยากตํ อนิยตญฺจ ธมฺมํ ปฏิจฺจ อพฺยากโต อนิยโต ธมฺโม อุปฺปชฺชติ เหตุปจฺจยา. (1)}

\prose{อกุสลํ อนิยตญฺจ อพฺยากตํ อนิยตญฺจ ธมฺมํ ปฏิจฺจ อพฺยากโต อนิยโต ธมฺโม อุปฺปชฺชติ เหตุปจฺจยา. (1) (สํขิตฺตํ.)}

\prose{เหตุยา นว, อารมฺมเณ ตีณิ, อวิคเต นว. (สํขิตฺตํ.\csromanspacer{} สหชาตวารมฺปิ ฯเปฯ สมฺปยุตฺตวารมฺปิ ปฏิจฺจวารสทิสํ.)}

\proseitem{66}{กุสโล อนิยโต ธมฺโม กุสลสฺส อนิยตสฺส ธมฺมสฺส เหตุปจฺจเยน ปจฺจโย.\csromanspacer{} ตีณิ.}

\prose{อกุสโล อนิยโต ธมฺโม อกุสลสฺส อนิยตสฺส ธมฺมสฺส เหตุปจฺจเยน ปจฺจโย.\csromanspacer{} ตีณิ.}

\prose{อพฺยากโต อนิยโต ธมฺโม อพฺยากตสฺส อนิยตสฺส ธมฺมสฺส เหตุปจฺจเยน ปจฺจโย. (1) (สํขิตฺตํ.)}

\prose{เหตุยา สตฺต, อารมฺมเณ นว, อวิคเต เตรส. (สํขิตฺตํ.\csromanspacer{} ยถา กุสลตฺติเก ปญฺหาวารํ, เอวํ วิตฺถาเรตพฺพํ.)}

\csromanheader{2}{1. กุสลตฺติก 15. มคฺคารมฺมณตฺติก}
\proseitem{67}{กุสลํ มคฺคารมฺมณํ ธมฺมํ ปฏิจฺจ กุสโล มคฺคารมฺมโณ ธมฺโม อุปฺปชฺชติ เหตุปจฺจยา. (1)}

\csromanheader{2}{1. กุสลตฺติก 17. อตีตตฺติก}
\prose{\csromanfolio{19}อพฺยากตํ มคฺคารมฺมณํ ธมฺมํ ปฏิจฺจ อพฺยากโต มคฺคารมฺมโณ ธมฺโม อุปฺปชฺชติ เหตุปจฺจยา. (1) (สํขิตฺตํ.)}

\prose{เหตุยา เทฺว ฯเปฯ อวิคเต เทฺว. (สํขิตฺตํ.\csromanspacer{} สพฺพตฺถ วิตฺถาโร.)}

\proseitem{68}{กุสลํ มคฺคเหตุกํ ธมฺมํ ปฏิจฺจ กุสโล มคฺคเหตุโก ธมฺโม อุปฺปชฺชติ เหตุปจฺจยา. (สํขิตฺตํ.)}

\csromanheader{2}{เหตุยา เอกํ ฯเปฯ อวิคเต เอกํ. (สพฺพตฺถ เอกํ. สํขิตฺตํ.)}
\csromanheader{2}{69. กุสลํ มคฺคาธิปติํ ธมฺมํ ปฏิจฺจ กุสโล มคฺคาธิปติ}
\prose{อพฺยากตํ มคฺคาธิปติํ ธมฺมํ ปฏิจฺจ อพฺยากโต มคฺคาธิปติ ธมฺโม อุปฺปชฺชติ เหตุปจฺจยา. (1) (สํขิตฺตํ.)}

\prose{เหตุยา เทฺว ฯเปฯ อวิคเต เทฺว. (สํขิตฺตํ.\csromanspacer{} สพฺพตฺถ วิตฺถาโร.)}

\csromanheader{2}{1. กุสลตฺติก 16. อุปฺปนฺนตฺติก}
\csromanheader{2}{7. ปญฺหาวาร ปจฺจยจตุกฺก}
\proseitem{70}{กุสโล อุปฺปนฺโน ธมฺโม กุสลสฺส อุปฺปนฺนสฺส ธมฺมสฺส เหตุปจฺจเยน ปจฺจโย. (สํขิตฺตํ.)}

\csromanheader{2}{เหตุยา สตฺต. (สํขิตฺตํ.)}
\csromanheader{2}{1. กุสลตฺติก 17. อตีตตฺติก}
\csromanheader{2}{7. ปญฺหาวาร ปจฺจยจตุกฺก}
\proseitem{71}{กุสโล ปจฺจุปฺปนฺโน ธมฺโม กุสลสฺส ปจฺจุปฺปนฺนสฺส ธมฺมสฺส เหตุปจฺจเยน ปจฺจโย. (สํขิตฺตํ.)}

\csromanheader{2}{เหตุยา สตฺต. (สํขิตฺตํ.)}
\csromanheader{2}{1. กุสลตฺติก 18. อตีตารมฺมณตฺติก}
\csromanheader{2}{72. กุสลํ อตีตารมฺมณํ ธมฺมํ ปฏิจฺจ กุสโล อตีตารมฺมโณ}
\csromanheader{2}{อกุสลํ อตีตารมฺมณํ ธมฺมํ ปฏิจฺจ อกุสโล อตีตารมฺมโณ}
\prose{\csromanfolio{20}อพฺยากตํ อตีตารมฺมณํ ธมฺมํ ปฏิจฺจ อพฺยากโต อตีตารมฺมโณ ธมฺโม อุปฺปชฺชติ เหตุปจฺจยา. (1) (สํขิตฺตํ.)}

\prose{เหตุยา ตีณิ, อวิคเต ตีณิ. (สํขิตฺตํ.) (สหชาตวารมฺปิ ฯเปฯ สมฺปยุตฺตวารมฺปิ ปฏิจฺจวารสทิสํ.)}

\proseitem{73}{กุสโล อตีตารมฺมโณ ธมฺโม กุสลสฺส อตีตารมฺมณสฺส ธมฺมสฺส เหตุปจฺจเยน ปจฺจโย. (1)}

\prose{อกุสโล อตีตารมฺมโณ ธมฺโม อกุสลสฺส อตีตารมฺมณสฺส ธมฺมสฺส เหตุปจฺจเยน ปจฺจโย. (1)}

\prose{อพฺยากโต อตีตารมฺมโณ ธมฺโม อพฺยากตสฺส อตีตารมฺมณสฺส ธมฺมสฺส เหตุปจฺจเยน ปจฺจโย. (1) (สํขิตฺตํ.)}

\prose{เหตุยา ตีณิ, อารมฺมเณ นว, อวิคเต ตีณิ. (สํขิตฺตํ.\csromanspacer{} ยถา กุสลตฺติเก ปญฺหาวารํ, เอวํ วิตฺถาเรตพฺพํ.)}

\csromanheader{2}{\textbf{อนาคตารมฺมณปท-เหตุ}}
\proseitem{74}{กุสลํ อนาคตารมฺมณํ ธมฺมํ ปฏิจฺจ กุสโล อนาคตารมฺมโณ ธมฺโม อุปฺปชฺชติ เหตุปจฺจยา. (สํขิตฺตํ.)}

\prose{เหตุยา ตีณิ, อารมฺมเณ ตีณิ, อวิคเต ตีณิ. (สํขิตฺตํ.\csromanspacer{} สหชาตวารมฺปิ ฯเปฯ สมฺปยุตฺตวารมฺปิ ปฏิจฺจวารสทิสํ.)}

\csromanheader{2}{1. กุสลตฺติก 19. อชฺฌตฺตตฺติก}
\proseitem{75}{\csromanfolio{21}กุสโล อนาคตารมฺมโณ ธมฺโม กุสลสฺส อนาคตารมฺมณสฺส ธมฺมสฺส เหตุปจฺจเยน ปจฺจโย. (สํขิตฺตํ.)}

\prose{เหตุยา ตีณิ, อารมฺมเณ นว, อวิคเต ตีณิ. (สํขิตฺตํ.\csromanspacer{} ยถา กุสลตฺติเก ปญฺหาวารํ, เอวํ วิตฺถาเรตพฺพํ.)}

\csromanheader{2}{\textbf{ปจฺจุปฺปนฺนารมฺมณปท-เหตุ}}
\proseitem{76}{กุสลํ ปจฺจุปฺปนฺนารมฺมณํ ธมฺมํ ปฏิจฺจ กุสโล ปจฺจุปฺปนฺนารมฺมโณ ธมฺโม อุปฺปชฺชติ เหตุปจฺจยา. (สํขิตฺตํ.)}

\prose{เหตุยา ตีณิ, อารมฺมเณ ตีณิ, อวิคเต ตีณิ. (สํขิตฺตํ.\csromanspacer{} สหชาตวารมฺปิ ฯเปฯ สมฺปยุตฺตวารมฺปิ ปฏิจฺจวารสทิสํ.)}

\proseitem{77}{กุสโล ปจฺจุปฺปนฺนารมฺมโณ ธมฺโม กุสลสฺส ปจฺจุปฺปนฺนารมฺมณสฺส ธมฺมสฺส เหตุปจฺจเยน ปจฺจโย. (1)}

\prose{อกุสโล ปจฺจุปฺปนฺนารมฺมโณ ธมฺโม อกุสลสฺส ปจฺจุปฺปนฺนารมฺมณสฺส ธมฺมสฺส เหตุปจฺจเยน ปจฺจโย. (1)}

\prose{อพฺยากโต ปจฺจุปฺปนฺนารมฺมโณ ธมฺโม อพฺยากตสฺส ปจฺจุปฺปนฺนารมฺมณสฺส ธมฺมสฺส เหตุปจฺจเยน ปจฺจโย. (1)}

\prose{เหตุยา ตีณิ, อารมฺมเณ ฉ, อวิคเต ตีณิ. (สํขิตฺตํ.\csromanspacer{} ยถา กุสลตฺติเก ปญฺหาวารํ, เอวํ วิตฺถาเรตพฺพํ.)}

\csromanheader{2}{1. กุสลตฺติก 19. อชฺฌตฺตตฺติก}
\proseitem{78}{กุสลํ อชฺฌตฺตํ ธมฺมํ ปฏิจฺจ กุสโล อชฺฌตฺโต ธมฺโม อุปฺปชฺชติ เหตุปจฺจยา.\csromanspacer{} ตีณิ.}

\prose{อกุสลํ อชฺฌตฺตํ ธมฺมํ ปฏิจฺจ อกุสโล อชฺฌตฺโต ธมฺโม อุปฺปชฺชติ เหตุปจฺจยา.\csromanspacer{} ตีณิ.}

\prose{อพฺยากตํ อชฺฌตฺตํ ธมฺมํ ปฏิจฺจ อพฺยากโต อชฺฌตฺโต ธมฺโม อุปฺปชฺชติ เหตุปจฺจยา.\csromanspacer{} ตีณิ.}

\prose{\csromanfolio{22}เหตุยา นว, อวิคเต นว. (สํขิตฺตํ.) (สหชาตวารมฺปิ ฯเปฯ สมฺปยุตฺตวารมฺปิ ปฏิจฺจวารสทิสํ, เอวํ วิตฺถาเรตพฺพํ.)}

\proseitem{79}{กุสโล อชฺฌตฺโต ธมฺโม กุสลสฺส อชฺฌตฺตสฺส ธมฺมสฺส เหตุปจฺจเยน ปจฺจโย.\csromanspacer{} ตีณิ.}

\prose{อกุสโล อชฺฌตฺโต ธมฺโม อกุสลสฺส อชฺฌตฺตสฺส ธมฺมสฺส เหตุปจฺจเยน ปจฺจโย.\csromanspacer{} ตีณิ.}

\prose{อพฺยากโต อชฺฌตฺโต ธมฺโม อพฺยากตสฺส อชฺฌตฺตสฺส ธมฺมสฺส เหตุปจฺจเยน ปจฺจโย. (1) (สํขิตฺตํ.)}

\prose{เหตุยา สตฺต, อารมฺมเณ นว, อวิคเต เตรส. (สํขิตฺตํ.\csromanspacer{} ยถา กุสลตฺติเก ปญฺหาวารํ, เอวํ วิตฺถาเรตพฺพํ.)}

\csromanheader{2}{\textbf{พหิทฺธาปท-เหตุ}}
\proseitem{80}{กุสลํ พหิทฺธา ธมฺมํ ปฏิจฺจ กุสโล พหิทฺธา ธมฺโม อุปฺปชฺชติ เหตุปจฺจยา.\csromanspacer{} ตีณิ.}

\prose{อกุสลํ พหิทฺธา ธมฺมํ ปฏิจฺจ อกุสโล พหิทฺธา ธมฺโม อุปฺปชฺชติ เหตุปจฺจยา.\csromanspacer{} ตีณิ.}

\prose{อพฺยากตํ พหิทฺธา ธมฺมํ ปฏิจฺจ อพฺยากโต พหิทฺธา ธมฺโม อุปฺปชฺชติ เหตุปจฺจยา. (1)}

\prose{กุสลํ พหิทฺธา จ อพฺยากตํ พหิทฺธา จ ธมฺมํ ปฏิจฺจ อพฺยากโต พหิทฺธา ธมฺโม อุปฺปชฺชติ เหตุปจฺจยา. (1)}

\prose{อกุสลํ พหิทฺธา จ อพฺยากตํ พหิทฺธา จ ธมฺมํ ปฏิจฺจ อพฺยากโต พหิทฺธา ธมฺโม อุปฺปชฺชติ เหตุปจฺจยา. (1) (สํขิตฺตํ.)}

\prose{เหตุยา นว ฯเปฯ วิปาเก เอกํ ฯเปฯ อวิคเต นว. (สํขิตฺตํ.) (สหชาตวารมฺปิ ฯเปฯ สมฺปยุตฺตวารมฺปิ ปฏิจฺจวารสทิสํ.)}

\proseitem{81}{กุสโล พหิทฺธา ธมฺโม กุสลสฺส พหิทฺธา ธมฺมสฺส เหตุปจฺจเยน ปจฺจโย.\csromanspacer{} ตีณิ.}

\csromanheader{2}{1. กุสลตฺติก 21. นิทสฺสนตฺติก}
\prose{\csromanfolio{23}อกุสโล พหิทฺธา ธมฺโม อกุสลสฺส พหิทฺธา ธมฺมสฺส เหตุปจฺจเยน ปจฺจโย.\csromanspacer{} ตีณิ.}

\prose{อพฺยากโต พหิทฺธา ธมฺโม อพฺยากตสฺส พหิทฺธา ธมฺมสฺส เหตุปจฺจเยน ปจฺจโย. (1) (สํขิตฺตํ.)}

\prose{เหตุยา สตฺต, อารมฺมเณ นว, อธิปติยา ทส ฯเปฯ อวิคเต เตรส.}

\prose{(สํขิตฺตํ.\csromanspacer{} ยถา กุสลตฺติเก ปญฺหาวารํ, เอวํ วิตฺถาเรตพฺพํ.)}

\csromanheader{2}{1. กุสลตฺติก 20. อชฺฌตฺตารมฺมณตฺติก}
\proseitem{82}{กุสลํ อชฺฌตฺตารมฺมณํ ธมฺมํ ปฏิจฺจ กุสโล อชฺฌตฺตารมฺมโณ ธมฺโม อุปฺปชฺชติ เหตุปจฺจยา. (สํขิตฺตํ.)}

\prose{เหตุยา ตีณิ ฯเปฯ วิปาเก เอกํ ฯเปฯ อวิคเต ตีณิ. (สํขิตฺตํ.)}

\proseitem{83}{กุสลํ พหิทฺธารมฺมณํ ธมฺมํ ปฏิจฺจ กุสโล พหิทฺธารมฺมโณ ธมฺโม อุปฺปชฺชติ เหตุปจฺจยา. (สํขิตฺตํ.)}

\prose{เหตุยา ตีณิ ฯเปฯ วิปาเก เอกํ ฯเปฯ อวิคเต ตีณิ. (สํขิตฺตํ.)}

\csromanheader{2}{1. กุสลตฺติก 21. สนิทสฺสนตฺติก}
\proseitem{84}{อพฺยากตํ อนิทสฺสนสปฺปฏิฆํ ธมฺมํ ปฏิจฺจ อพฺยากโต อนิทสฺสนสปฺปฏิโฆ ธมฺโม อุปฺปชฺชติ เหตุปจฺจยา. (สํขิตฺตํ.)}

\prose{เหตุยา เอกํ ฯเปฯ อวิคเต เอกํ. (สํขิตฺตํ.) (สหชาตวารมฺปิ ฯเปฯ ปญฺหาวารมฺปิ วิตฺถาเรตพฺพํ.)}

\proseitem{85}{กุสลํ อนิทสฺสน-อปฺปฏิฆํ ธมฺมํ ปฏิจฺจ กุสโล อนิทสฺสน-อปฺปฏิโฆ ธมฺโม อุปฺปชฺชติ เหตุปจฺจยา.\csromanspacer{}\csromanspacer{}กุสลํ อนิทสฺสน-อปฺปฏิฆํ ธมฺมํ \csromanfolio{24}ปฏิจฺจ อพฺยากโต อนิทสฺสน-อปฺปฏิโฆ ธมฺโม อุปฺปชฺชติ เหตุปจฺจยา.\csromanspacer{}\csromanspacer{}กุสลํ อนิทสฺสน-อปฺปฏิฆํ ธมฺมํ ปฏิจฺจ กุสโล อนิทสฺสน-อปฺปฏิโฆ จ อพฺยากโต อนิทสฺสน-อปฺปฏิโฆ จ ธมฺมา อุปฺปชฺชนฺติ เหตุปจฺจยา. (3)}

\prose{อกุสลํ อนิทสฺสน-อปฺปฏิฆํ ธมฺมํ ปฏิจฺจ อกุสโล อนิทสฺสน-อปฺปฏิโฆ ธมฺโม อุปฺปชฺชติ เหตุปจฺจยา.\csromanspacer{} ตีณิ.}

\prose{อพฺยากตํ อนิทสฺสน-อปฺปฏิฆํ ธมฺมํ ปฏิจฺจ อพฺยากโต อนิทสฺสน-อปฺปฏิโฆ ธมฺโม อุปฺปชฺชติ เหตุปจฺจยา. (1) กุสลํ อนิทสฺสน-อปฺปฏิฆญฺจ อพฺยากตํ อนิทสฺสน-อปฺปฏิฆญฺจ ธมฺมํ ปฏิจฺจ อพฺยากโต อนิทสฺสน-อปฺปฏิโฆ ธมฺโม อุปฺปชฺชติ}

\prose{อกุสลํ อนิทสฺสน-อปฺปฏิฆญฺจ อพฺยากตํ อนิทสฺสน-อปฺปฏิฆญฺจ ธมฺมํ ปฏิจฺจ อพฺยากโต อนิทสฺสน-อปฺปฏิโฆ ธมฺโม อุปฺปชฺชติ}

\prose{เหตุยา นว, อารมฺมเณ ตีณิ ฯเปฯ วิปาเก เอกํ ฯเปฯ อวิคเต นว. (สํขิตฺตํ.\csromanspacer{} สหชาตวารมฺปิ ฯเปฯ สมฺปยุตฺตวารมฺปิ ปญฺหาวารมฺปิ วิตฺถาเรตพฺพํ.) กุสลตฺติกสนิทสฺสนตฺติกํ นิฏฺฐิตํ.}

\csromanheader{2}{2. เวทนาตฺติก 1. กุสลตฺติก}
\proseitem{86}{สุขาย เวทนาย สมฺปยุตฺตํ กุสลํ ธมฺมํ ปฏิจฺจ สุขาย เวทนาย สมฺปยุตฺโต กุสโล ธมฺโม อุปฺปชฺชติ เหตุปจฺจยา. (1)}

\prose{อทุกฺขมสุขาย เวทนาย สมฺปยุตฺตํ กุสลํ ธมฺมํ ปฏิจฺจ อทุกฺขมสุขาย เวทนาย สมฺปยุตฺโต กุสโล ธมฺโม อุปฺปชฺชติ}

\csromanheader{2}{เหตุยา เทฺว ฯเปฯ อวิคเต เทฺว. (สํขิตฺตํ.)}
\prose{(สหชาตวารมฺปิ ฯเปฯ สมฺปยุตฺตวารมฺปิ ปฏิจฺจวารสทิสํ.)}

\proseitem{87}{สุขาย เวทนาย สมฺปยุตฺโต กุสโล ธมฺโม สุขาย เวทนาย สมฺปยุตฺตสฺส กุสลสฺส ธมฺมสฺส เหตุปจฺจเยน ปจฺจโย. (1)}

\csromanheader{2}{2. เวทนาตฺติก 1. กุสลตฺติก}
\prose{\csromanfolio{25}อทุกฺขมสุขาย เวทนาย สมฺปยุตฺโต กุสโล ธมฺโม อทุกฺขมสุขาย เวทนาย สมฺปยุตฺตสฺส กุสลสฺส ธมฺมสฺส เหตุปจฺจเยน ปจฺจโย. (1) (สํขิตฺตํ.)}

\prose{เหตุยา เทฺว, อวิคเต เทฺว. (สํขิตฺตํ.\csromanspacer{} สพฺพตฺถ วิตฺถาโร.)}

\proseitem{88}{สุขาย เวทนาย สมฺปยุตฺตํ อกุสลํ ธมฺมํ ปฏิจฺจ สุขาย เวทนาย สมฺปยุตฺโต อกุสโล ธมฺโม อุปฺปชฺชติ เหตุปจฺจยา. (1)}

\prose{ทุกฺขาย เวทนาย สมฺปยุตฺตํ อกุสลํ ธมฺมํ ปฏิจฺจ ทุกฺขาย เวทนาย สมฺปยุตฺโต อกุสโล ธมฺโม อุปฺปชฺชติ เหตุปจฺจยา. (1)}

\prose{อทุกฺขมสุขาย เวทนาย สมฺปยุตฺตํ อกุสลํ ธมฺมํ ปฏิจฺจ อทุกฺขมสุขาย เวทนาย สมฺปยุตฺโต อกุสโล ธมฺโม อุปฺปชฺชติ เหตุปจฺจยา. (1) (สํขิตฺตํ.)}

\prose{เหตุยา ตีณิ ฯเปฯ อวิคเต ตีณิ. (สํขิตฺตํ.\csromanspacer{} สพฺพตฺถ วิตฺถาโร.)}

\proseitem{89}{สุขาย เวทนาย สมฺปยุตฺตํ อพฺยากตํ ธมฺมํ ปฏิจฺจ สุขาย เวทนาย สมฺปยุตฺโต อพฺยากโต ธมฺโม อุปฺปชฺชติ เหตุปจฺจยา. (1)}

\prose{อทุกฺขมสุขาย เวทนาย สมฺปยุตฺตํ อพฺยากตํ ธมฺมํ ปฏิจฺจ อทุกฺขมสุขาย เวทนาย สมฺปยุตฺโต อพฺยากโต ธมฺโม อุปฺปชฺชติ เหตุปจฺจยา. (1) (สํขิตฺตํ.)}

\prose{เหตุยา เทฺว, อวิคเต ตีณิ. (สํขิตฺตํ.)}

\csromanheader{2}{3. วิปากตฺติก 1. กุสลตฺติก}
\csromanheader{2}{1. ปฏิจฺจวาร-เหตุ}
\proseitem{90}{วิปากธมฺมธมฺมํ กุสลํ ธมฺมํ ปฏิจฺจ วิปากธมฺมธมฺโม กุสโล ธมฺโม อุปฺปชฺชติ เหตุปจฺจยา.}

\csromanheader{2}{เหตุยา เอกํ ฯเปฯ อวิคเต เอกํ. (สํขิตฺตํ.)}
\proseitem{91}{\csromanfolio{26}วิปากธมฺมธมฺมํ อกุสลํ ธมฺมํ ปฏิจฺจ วิปากธมฺมธมฺโม อกุสโล ธมฺโม อุปฺปชฺชติ เหตุปจฺจยา.}

\csromanheader{2}{เหตุยา เอกํ. (สพฺพตฺถ เอกํ.)}
\proseitem{92}{วิปากํ อพฺยากตํ ธมฺมํ ปฏิจฺจ วิปาโก อพฺยากโต ธมฺโม อุปฺปชฺชติ เหตุปจฺจยา.\csromanspacer{} ตีณิ.}

\prose{เนววิปากนวิปากธมฺมธมฺมํ อพฺยากตํ ธมฺมํ ปฏิจฺจ เนววิปากนวิปากธมฺมธมฺโม อพฺยากโต ธมฺโม อุปฺปชฺชติ เหตุปจฺจยา.\csromanspacer{} ตีณิ.}

\prose{วิปากํ อพฺยากตญฺจ เนววิปากนวิปากธมฺมธมฺมํ อพฺยากตญฺจ ธมฺมํ ปฏิจฺจ วิปาโก อพฺยากโต ธมฺโม อุปฺปชฺชติ เหตุปจฺจยา.\csromanspacer{} ตีณิ. (สํขิตฺตํ.)}

\prose{เหตุยา นว, อวิคเต นว. (สํขิตฺตํ.\csromanspacer{} สพฺพตฺถ วิตฺถาโร.)}

\csromanheader{2}{4. อุปาทินฺนตฺติก 1. กุสลตฺติก}
\csromanheader{2}{1. ปฏิจฺจวาร-เหตุ}
\proseitem{93}{อนุปาทินฺนุปาทานิยํ กุสลํ ธมฺมํ ปฏิจฺจ อนุปาทินฺนุปาทานิโย กุสโล ธมฺโม อุปฺปชฺชติ เหตุปจฺจยา. (1)}

\prose{อนุปาทินฺน-อนุปาทานิยํ กุสลํ ธมฺมํ ปฏิจฺจ อนุปาทินฺน-อนุปาทานิโย กุสโล ธมฺโม อุปฺปชฺชติ เหตุปจฺจยา. (1)}

\csromanheader{2}{เหตุยา เทฺว. (สพฺพตฺถ วิตฺถาโร.)}
\prose{อนุปาทินฺนุปาทานิยํ อกุสลํ ธมฺมํ ปฏิจฺจ อนุปาทินฺนุปาทานิโย อกุสโล ธมฺโม อุปฺปชฺชติ เหตุปจฺจยา. (สพฺพตฺถ เอกํ.\csromanspacer{} สพฺพตฺถ วิตฺถาโร.)}

\proseitem{94}{อุปาทินฺนุปาทานิยํ อพฺยากตํ ธมฺมํ ปฏิจฺจ อุปาทินฺนุปาทานิโย อพฺยากโต ธมฺโม อุปฺปชฺชติ เหตุปจฺจยา.\csromanspacer{} ตีณิ. (สพฺพตฺถ วิตฺถาโร.)}

\prose{อนุปาทินฺนุปาทานิยํ อพฺยากตํ ธมฺมํ ปฏิจฺจ อนุปาทินฺนุปาทานิโย อพฺยากโต ธมฺโม อุปฺปชฺชติ เหตุปจฺจยา.\csromanspacer{} ตีณิ.}

\csromanheader{2}{6. วิตกฺกตฺติก 1. กุสลตฺติก}
\prose{\csromanfolio{27}อนุปาทินฺน-อนุปาทานิยํ อพฺยากตํ ธมฺมํ ปฏิจฺจ อนุปาทินฺน-อนุปาทานิโย อพฺยากโต ธมฺโม อุปฺปชฺชติ เหตุปจฺจยา. (1)}

\prose{อนุปาทินฺนุปาทานิยํ อพฺยากตญฺจ อนุปาทินฺน-อนุปาทานิยํ อพฺยากตญฺจ ธมฺมํ ปฏิจฺจ อนุปาทินฺนุปาทานิโย อพฺยากโต ธมฺโม อุปฺปชฺชติ เหตุปจฺจยา. (1)}

\prose{อุปาทินฺนุปาทานิยํ อพฺยากตญฺจ อนุปาทินฺนุปาทานิยํ อพฺยากตญฺจ ธมฺมํ ปฏิจฺจ อนุปาทินฺนุปาทานิโย อพฺยากโต ธมฺโม อุปฺปชฺชติ เหตุปจฺจยา. (1) (สํขิตฺตํ.)}

\csromanheader{2}{เหตุยา นว. (สพฺพตฺถ วิตฺถาโร.)}
\csromanheader{2}{5. สํกิลิฏฺฐตฺติก 1. กุสลตฺติก}
\csromanheader{2}{1. ปฏิจฺจวาร-เหตุ}
\proseitem{95}{อสํกิลิฏฺฐสํกิเลสิกํ กุสลํ ธมฺมํ ปฏิจฺจ อสํกิลิฏฺฐสํกิเลสิโก กุสโล ธมฺโม อุปฺปชฺชติ เหตุปจฺจยา.}

\prose{อสํกิลิฏฺฐ-อสํกิเลสิกํ กุสลํ ธมฺมํ ปฏิจฺจ อสํกิลิฏฺฐ-อสํกิเลสิโก กุสโล ธมฺโม อุปฺปชฺชติ เหตุปจฺจยา. (สพฺพตฺถ เทฺว.)}

\prose{สํกิลิฏฺฐสํกิเลสิกํ อกุสลํ ธมฺมํ ปฏิจฺจ สํกิลิฏฺฐสํกิเลสิโก อกุสโล ธมฺโม อุปฺปชฺชติ เหตุปจฺจยา. (สพฺพตฺถ เอกํ.)}

\prose{อสํกิลิฏฺฐสํกิเลสิกํ อพฺยากตํ ธมฺมํ ปฏิจฺจ อสํกิลิฏฺฐสํกิเลสิโก อพฺยากโต ธมฺโม อุปฺปชฺชติ เหตุปจฺจยา.}

\prose{เหตุยา ปญฺจ.}

\csromanheader{2}{6. วิตกฺกตฺติก 1. กุสลตฺติก}
\csromanheader{2}{1. ปฏิจฺจวาร-เหตุ}
\proseitem{96}{สวิตกฺกสวิจารํ กุสลํ ธมฺมํ ปฏิจฺจ สวิตกฺกสวิจาโร กุสโล ธมฺโม อุปฺปชฺชติ เหตุปจฺจยา.\csromanspacer{} ตีณิ.}

\prose{\csromanfolio{28}อวิตกฺกวิจารมตฺตํ กุสลํ ธมฺมํ ปฏิจฺจ อวิตกฺกวิจารมตฺโต กุสโล ธมฺโม อุปฺปชฺชติ เหตุปจฺจยา.\csromanspacer{} จตฺตาริ.}

\prose{อวิตกฺก-อวิจารํ กุสลํ ธมฺมํ ปฏิจฺจ อวิตกฺก-อวิจาโร กุสโล}

\prose{อวิตกฺก-อวิจารํ กุสลํ ธมฺมํ ปฏิจฺจ อวิตกฺกวิจารมตฺโต กุสโล ธมฺโม อุปฺปชฺชติ เหตุปจฺจยา. (1)}

\prose{อวิตกฺกวิจารมตฺตํ กุสลญฺจ อวิตกฺก-อวิจารํ กุสลญฺจ ธมฺมํ ฯเปฯ.\csromanspacer{} สวิตกฺกสวิจารํ กุสลญฺจ อวิตกฺกวิจารมตฺตํ กุสลญฺจ ธมฺมํ ปฏิจฺจ สวิตกฺกสวิจาโร กุสโล ธมฺโม อุปฺปชฺชติ}

\prose{เหตุยา เอกาทส.}

\proseitem{97}{สวิตกฺกสวิจารํ อกุสลํ ธมฺมํ ปฏิจฺจ สวิตกฺกสวิจาโร อกุสโล ธมฺโม อุปฺปชฺชติ เหตุปจฺจยา.\csromanspacer{} ตีณิ.}

\prose{อวิตกฺกวิจารมตฺตํ อกุสลํ ธมฺมํ ปฏิจฺจ สวิตกฺกสวิจาโร อกุสโล ธมฺโม อุปฺปชฺชติ เหตุปจฺจยา. (1)}

\prose{สวิตกฺกสวิจารํ อกุสลญฺจ อวิตกฺกวิจารมตฺตํ อกุสลญฺจ ธมฺมํ ปฏิจฺจ สวิตกฺกสวิจาโร อกุสโล ธมฺโม อุปฺปชฺชติ}

\prose{เหตุยา ปญฺจ.}

\proseitem{98}{สวิตกฺกสวิจารํ อพฺยากตํ ธมฺมํ ปฏิจฺจ สวิตกฺกสวิจาโร อพฺยากโต ธมฺโม อุปฺปชฺชติ เหตุปจฺจยา.}

\prose{เหตุยา สตฺตติํส.}

\csromanheader{2}{7. ปีติตฺติก 1. กุสลตฺติก}
\csromanheader{2}{1. ปฏิจฺจวาร-เหตุ}
\proseitem{99}{ปีติสหคตํ กุสลํ ธมฺมํ ปฏิจฺจ ปีติสหคโต กุสโล ธมฺโม อุปฺปชฺชติ เหตุปจฺจยา.\csromanspacer{} ตีณิ.}

\csromanheader{2}{7. ปีติตฺติก 1. กุสลตฺติก}
\prose{\csromanfolio{29}สุขสหคตํ กุสลํ ธมฺมํ ปฏิจฺจ สุขสหคโต กุสโล ธมฺโม อุปฺปชฺชติ เหตุปจฺจยา.\csromanspacer{} ตีณิ.}

\prose{อุเปกฺขาสหคตํ กุสลํ ธมฺมํ ปฏิจฺจ อุเปกฺขาสหคโต กุสโล}

\prose{ปีติสหคตํ กุสลญฺจ สุขสหคตํ กุสลญฺจ ธมฺมํ ปฏิจฺจ ปีติสหคโต กุสโล ธมฺโม อุปฺปชฺชติ เหตุปจฺจยา.\csromanspacer{} ตีณิ. (สพฺพตฺถ ทส.\csromanspacer{} สพฺพตฺถ วิตฺถาโร.)}

\proseitem{100}{ปีติสหคตํ อกุสลํ ธมฺมํ ปฏิจฺจ ปีติสหคโต อกุสโล ธมฺโม อุปฺปชฺชติ เหตุปจฺจยา.\csromanspacer{} ตีณิ.}

\prose{สุขสหคตํ อกุสลํ ธมฺมํ ปฏิจฺจ สุขสหคโต อกุสโล ธมฺโม อุปฺปชฺชติ เหตุปจฺจยา.\csromanspacer{} ตีณิ.}

\prose{อุเปกฺขาสหคตํ อกุสลํ ธมฺมํ ปฏิจฺจ อุเปกฺขาสหคโต อกุสโล ธมฺโม อุปฺปชฺชติ เหตุปจฺจยา. (1)}

\prose{ปีติสหคตํ อกุสลญฺจ สุขสหคตํ อกุสลญฺจ ธมฺมํ ปฏิจฺจ ปีติสหคโต อกุสโล ธมฺโม อุปฺปชฺชติ เหตุปจฺจยา.\csromanspacer{} ตีณิ. (สพฺพตฺถ ทส.\csromanspacer{} สพฺพตฺถ วิตฺถาโร.)}

\proseitem{101}{ปีติสหคตํ อพฺยากตํ ธมฺมํ ปฏิจฺจ ปีติสหคโต อพฺยากโต ธมฺโม อุปฺปชฺชติ เหตุปจฺจยา.\csromanspacer{} ตีณิ.}

\prose{สุขสหคตํ อพฺยากตํ ธมฺมํ ปฏิจฺจ สุขสหคโต อพฺยากโต ธมฺโม อุปฺปชฺชติ เหตุปจฺจยา.\csromanspacer{} ตีณิ.}

\prose{อุเปกฺขาสหคตํ อพฺยากตํ ธมฺมํ ปฏิจฺจ อุเปกฺขาสหคโต อพฺยากโต ธมฺโม อุปฺปชฺชติ เหตุปจฺจยา. (1)}

\prose{ปีติสหคตํ อพฺยากตญฺจ สุขสหคตํ อพฺยากตญฺจ ธมฺมํ ปฏิจฺจ ปีติสหคโต อพฺยากโต ธมฺโม อุปฺปชฺชติ เหตุปจฺจยา.\csromanspacer{} ตีณิ. (สพฺพตฺถ ทส.\csromanspacer{} สพฺพตฺถ วิตฺถาโร.)}

\csromanheader{2}{8. ทสฺสนตฺติก 1. กุสลตฺติก}
\csromanheader{2}{1. ปฏิจฺจวาร-เหตุ}
\proseitem{102}{\csromanfolio{30}เนวทสฺสเนน นภาวนาย ปหาตพฺพํ กุสลํ ธมฺมํ ปฏิจฺจ เนวทสฺสเนน นภาวนาย ปหาตพฺโพ กุสโล ธมฺโม อุปฺปชฺชติ เหตุปจฺจยา. (สพฺพตฺถ เอกํ.)}

\proseitem{103}{ทสฺสเนน ปหาตพฺพํ อกุสลํ ธมฺมํ ปฏิจฺจ ทสฺสเนน ปหาตพฺโพ อกุสโล ธมฺโม อุปฺปชฺชติ เหตุปจฺจยา. (1)}

\prose{ภาวนาย ปหาตพฺพํ อกุสลํ ธมฺมํ ปฏิจฺจ ภาวนาย ปหาตพฺโพ อกุสโล ธมฺโม อุปฺปชฺชติ เหตุปจฺจยา. (1) (สพฺพตฺถ เทฺว.\csromanspacer{} สพฺพตฺถ วิตฺถาโร.)}

\proseitem{104}{เนวทสฺสเนน นภาวนาย ปหาตพฺพํ อพฺยากตํ ธมฺมํ ปฏิจฺจ เนวทสฺสเนน นภาวนาย ปหาตพฺโพ อพฺยากโต ธมฺโม อุปฺปชฺชติ เหตุปจฺจยา. (สพฺพตฺถ เอกํ.)}

\csromanheader{2}{9. ทสฺสนเหตุตฺติก 1. กุสลตฺติก}
\csromanheader{2}{1. ปฏิจฺจวาร-เหตุ}
\proseitem{105}{เนวทสฺสเนน นภาวนาย ปหาตพฺพเหตุกํ กุสลํ ธมฺมํ ปฏิจฺจ เนวทสฺสเนน นภาวนาย ปหาตพฺพเหตุโก กุสโล ธมฺโม อุปฺปชฺชติ เหตุปจฺจยา. (สพฺพตฺถ เอกํ.\csromanspacer{} สพฺพตฺถ วิตฺถาโร.)}

\proseitem{106}{ทสฺสเนน ปหาตพฺพเหตุกํ อกุสลํ ธมฺมํ ปฏิจฺจ ทสฺสเนน ปหาตพฺพเหตุโก อกุสโล ธมฺโม อุปฺปชฺชติ เหตุปจฺจยา.}

\prose{ภาวนาย ปหาตพฺพเหตุกํ อกุสลํ ธมฺมํ ปฏิจฺจ ภาวนาย ปหาตพฺพเหตุโก อกุสโล ธมฺโม อุปฺปชฺชติ เหตุปจฺจยา. (สํขิตฺตํ.)}

\prose{เหตุยา ฉ, อารมฺมเณ ทส, อธิปติยา เทฺว ฯเปฯ อวิคเต ทส. (สพฺพตฺถ วิตฺถาโร.)}

\csromanheader{2}{11. เสกฺขตฺติก 1. กุสลตฺติก}
\proseitem{107}{\csromanfolio{31}เนวทสฺสเนน นภาวนาย ปหาตพฺพเหตุกํ อพฺยากตํ ธมฺมํ ปฏิจฺจ เนวทสฺสเนน นภาวนาย ปหาตพฺพเหตุโก อพฺยากโต ธมฺโม อุปฺปชฺชติ เหตุปจฺจยา. (สพฺพตฺถ เอกํ.)}

\csromanheader{2}{10. อาจยคามิตฺติก 1. กุสลตฺติก}
\csromanheader{2}{1. ปฏิจฺจวาร-เหตุ}
\csromanheader{2}{108. อาจยคามิํ กุสลํ ธมฺมํ ปฏิจฺจ อาจยคามี กุสโล}
\prose{อปจยคามิํ กุสลํ ธมฺมํ ปฏิจฺจ อปจยคามี กุสโล ธมฺโม อุปฺปชฺชติ เหตุปจฺจยา. (1) (สพฺพตฺถ เทฺว.\csromanspacer{} สพฺพตฺถ วิตฺถาโร.)}

\prose{อาจยคามิํ อกุสลํ ธมฺมํ ปฏิจฺจ อาจยคามี อกุสโล ธมฺโม อุปฺปชฺชติ เหตุปจฺจยา. (สพฺพตฺถ เอกํ.)}

\prose{เนวาจยคามินาปจยคามิํ อพฺยากตํ ธมฺมํ ปฏิจฺจ เนวาจยคามินาปจยคามี อพฺยากโต ธมฺโม อุปฺปชฺชติ เหตุปจฺจยา. (สพฺพตฺถ เอกํ.)}

\csromanheader{2}{11. เสกฺขตฺติก 1. กุสลตฺติก}
\csromanheader{2}{1. ปฏิจฺจวาร-เหตุ}
\proseitem{109}{เสกฺขํ กุสลํ ธมฺมํ ปฏิจฺจ เสกฺโข กุสโล ธมฺโม อุปฺปชฺชติ เหตุปจฺจยา. (1)}

\prose{เนวเสกฺขนาเสกฺขํ กุสลํ ธมฺมํ ปฏิจฺจ เนวเสกฺขนาเสกฺโข กุสโล ธมฺโม อุปฺปชฺชติ เหตุปจฺจยา. (1) (สพฺพตฺถ เทฺว.\csromanspacer{} สพฺพตฺถ วิตฺถาโร.)}

\prose{เนวเสกฺขนาเสกฺขํ อกุสลํ ธมฺมํ ปฏิจฺจ เนวเสกฺขนาเสกฺโข อกุสโล ธมฺโม อุปฺปชฺชติ เหตุปจฺจยา. (สพฺพตฺถ เอกํ.)}

\prose{เสกฺขํ อพฺยากตํ ธมฺมํ ปฏิจฺจ เสกฺโข อพฺยากโต ธมฺโม อุปฺปชฺชติ เหตุปจฺจยา. (สํขิตฺตํ.)}

\csromanheader{2}{เหตุยา นว. (สพฺพตฺถ วิตฺถาเรตพฺพํ.)}
\csromanheader{2}{12. ปริตฺตติก 1. กุสลตฺติก}
\csromanheader{2}{1. ปฏิจฺจวาร-เหตุ}
\proseitem{110}{\csromanfolio{32}ปริตฺตํ กุสลํ ธมฺมํ ปฏิจฺจ ปริตฺโต กุสโล ธมฺโม อุปฺปชฺชติ เหตุปจฺจยา. (1)}

\prose{มหคฺคตํ กุสลํ ธมฺมํ ปฏิจฺจ มหคฺคโต กุสโล ธมฺโม อุปฺปชฺชติ เหตุปจฺจยา. (1)}

\prose{อปฺปมาณํ กุสลํ ธมฺมํ ปฏิจฺจ อปฺปมาโณ กุสโล ธมฺโม อุปฺปชฺชติ เหตุปจฺจยา. (1) (สพฺพตฺถ ตีณิ.\csromanspacer{} สพฺพตฺถ วิตฺถาโร.)}

\proseitem{111}{ปริตฺตํ อกุสลํ ธมฺมํ ปฏิจฺจ ปริตฺโต อกุสโล ธมฺโม อุปฺปชฺชติ เหตุปจฺจยา. (สพฺพตฺถ เอกํ.)}

\prose{ปริตฺตํ อพฺยากตํ ธมฺมํ ปฏิจฺจ ปริตฺโต อพฺยากโต ธมฺโม อุปฺปชฺชติ เหตุปจฺจยา.\csromanspacer{} ตีณิ.}

\prose{มหคฺคตํ อพฺยากตํ ธมฺมํ ปฏิจฺจ มหคฺคโต อพฺยากโต ธมฺโม อุปฺปชฺชติ เหตุปจฺจยา.\csromanspacer{} ตีณิ.}

\prose{อปฺปมาณํ อพฺยากตํ ธมฺมํ ปฏิจฺจ อปฺปมาโณ อพฺยากโต ธมฺโม อุปฺปชฺชติ เหตุปจฺจยา. (สํขิตฺตํ.)}

\prose{เหตุยา เตรส.}

\csromanheader{2}{13. ปริตฺตารมฺมณตฺติก 1. กุสลตฺติก}
\csromanheader{2}{1. ปฏิจฺจวาร-เหตุ}
\proseitem{112}{ปริตฺตารมฺมณํ กุสลํ ธมฺมํ ปฏิจฺจ ปริตฺตารมฺมโณ กุสโล ธมฺโม อุปฺปชฺชติ เหตุปจฺจยา. (1)}

\prose{มหคฺคตารมฺมณํ กุสลํ ธมฺมํ ปฏิจฺจ มหคฺคตารมฺมโณ กุสโล ธมฺโม อุปฺปชฺชติ เหตุปจฺจยา. (1)}

\csromanheader{2}{14. หีนตฺติก 1. กุสลตฺติก}
\prose{\csromanfolio{33}อปฺปมาณารมฺมณํ กุสลํ ธมฺมํ ปฏิจฺจ อปฺปมาณารมฺมโณ กุสโล ธมฺโม อุปฺปชฺชติ เหตุปจฺจยา. (1) (สพฺพตฺถ ตีณิ.\csromanspacer{} สพฺพตฺถ วิตฺถาโร.)}

\proseitem{113}{ปริตฺตารมฺมณํ อกุสลํ ธมฺมํ ปฏิจฺจ ปริตฺตารมฺมโณ อกุสโล ธมฺโม อุปฺปชฺชติ เหตุปจฺจยา. (1)}

\prose{มหคฺคตารมฺมณํ อกุสลํ ธมฺมํ ปฏิจฺจ มหคฺคตารมฺมโณ อกุสโล ธมฺโม อุปฺปชฺชติ เหตุปจฺจยา. (1) (สพฺพตฺถ เทฺว.\csromanspacer{} สพฺพตฺถ วิตฺถาโร.)}

\proseitem{114}{ปริตฺตารมฺมณํ อพฺยากตํ ธมฺมํ ปฏิจฺจ ปริตฺตารมฺมโณ อพฺยากโต ธมฺโม อุปฺปชฺชติ เหตุปจฺจยา. (1)}

\prose{มหคฺคตารมฺมณํ อพฺยากตํ ธมฺมํ ปฏิจฺจ มหคฺคตารมฺมโณ อพฺยากโต ธมฺโม อุปฺปชฺชติ เหตุปจฺจยา. (1)}

\prose{อปฺปมาณารมฺมณํ อพฺยากตํ ธมฺมํ ปฏิจฺจ อปฺปมาณารมฺมโณ อพฺยากโต ธมฺโม อุปฺปชฺชติ เหตุปจฺจยา. (1) (สพฺพตฺถ ตีณิ.\csromanspacer{} สพฺพตฺถ วิตฺถาโร.)}

\csromanheader{2}{14. หีนตฺติก 1. กุสลตฺติก}
\csromanheader{2}{1. ปฏิจฺจวาร-เหตุ}
\proseitem{115}{มชฺฌิมํ กุสลํ ธมฺมํ ปฏิจฺจ มชฺฌิโม กุสโล ธมฺโม อุปฺปชฺชติ เหตุปจฺจยา. (1)}

\prose{ปณีตํ กุสลํ ธมฺมํ ปฏิจฺจ ปณีโต กุสโล ธมฺโม อุปฺปชฺชติ เหตุปจฺจยา. (1) (สพฺพตฺถ เทฺว.\csromanspacer{} สพฺพตฺถ วิตฺถาโร.)}

\proseitem{116}{หีนํ อกุสลํ ธมฺมํ ปฏิจฺจ หีโน อกุสโล ธมฺโม อุปฺปชฺชติ เหตุปจฺจยา. (สพฺพตฺถ เอกํ.)}

\csromanheader{2}{117. มชฺฌิมํ อพฺยากตํ ธมฺมํ ปฏิจฺจ มชฺฌิโม อพฺยากโต}
\prose{ปณีตํ อพฺยากตํ ธมฺมํ ปฏิจฺจ ปณีโต อพฺยากโต ธมฺโม อุปฺปชฺชติ เหตุปจฺจยา.\csromanspacer{} ตีณิ.}

\prose{\csromanfolio{34}มชฺฌิมํ อพฺยากตญฺจ ปณีตํ อพฺยากตญฺจ ธมฺมํ ปฏิจฺจ มชฺฌิโม อพฺยากโต ธมฺโม อุปฺปชฺชติ เหตุปจฺจยา. (1) (สํขิตฺตํ.)}

\csromanheader{2}{เหตุยา ปญฺจ. (สพฺพตฺถ วิตฺถาโร.)}
\csromanheader{2}{15. มิจฺฉตฺตตฺติก 1. กุสลตฺติก}
\csromanheader{2}{1. ปฏิจฺจวาร-เหตุ}
\proseitem{118}{สมฺมตฺตนิยตํ กุสลํ ธมฺมํ ปฏิจฺจ สมฺมตฺตนิยโต กุสโล ธมฺโม อุปฺปชฺชติ เหตุปจฺจยา. (1)}

\prose{อนิยตํ กุสลํ ธมฺมํ ปฏิจฺจ อนิยโต กุสโล ธมฺโม อุปฺปชฺชติ เหตุปจฺจยา. (1) (สพฺพตฺถ เทฺว.\csromanspacer{} สพฺพตฺถ วิตฺถาโร.)}

\proseitem{119}{มิจฺฉตฺตนิยตํ อกุสลํ ธมฺมํ ปฏิจฺจ มิจฺฉตฺตนิยโต อกุสโล ธมฺโม อุปฺปชฺชติ เหตุปจฺจยา. (1)}

\prose{อนิยตํ อกุสลํ ธมฺมํ ปฏิจฺจ อนิยโต อกุสโล ธมฺโม อุปฺปชฺชติ เหตุปจฺจยา. (1) (สพฺพตฺถ เทฺว.\csromanspacer{} สพฺพตฺถ วิตฺถาโร.)}

\prose{อนิยตํ อพฺยากตํ ธมฺมํ ปฏิจฺจ อนิยโต อพฺยากโต ธมฺโม อุปฺปชฺชติ เหตุปจฺจยา. (สพฺพตฺถ เอกํ.)}

\csromanheader{2}{16. มคฺคารมฺมณตฺติก 1. กุสลตฺติก}
\csromanheader{2}{1. ปฏิจฺจวาร-เหตุ}
\proseitem{120}{มคฺคารมฺมณํ กุสลํ ธมฺมํ ปฏิจฺจ มคฺคารมฺมโณ กุสโล ธมฺโม อุปฺปชฺชติ เหตุปจฺจยา.\csromanspacer{} ตีณิ.}

\prose{มคฺคเหตุกํ กุสลํ ธมฺมํ ปฏิจฺจ มคฺคเหตุโก กุสโล ธมฺโม อุปฺปชฺชติ เหตุปจฺจยา.\csromanspacer{} ตีณิ.}

\prose{มคฺคาธิปติํ กุสลํ ธมฺมํ ปฏิจฺจ มคฺคาธิปติ กุสโล ธมฺโม อุปฺปชฺชติ เหตุปจฺจยา.\csromanspacer{} ปญฺจ.}

\csromanheader{2}{18. อตีตตฺติก 1. กุสลตฺติก}
\prose{\csromanfolio{35}มคฺคารมฺมณํ กุสลญฺจ มคฺคาธิปติํ กุสลญฺจ ธมฺมํ ปฏิจฺจ มคฺคารมฺมโณ กุสโล ธมฺโม อุปฺปชฺชติ เหตุปจฺจยา.\csromanspacer{} ตีณิ.}

\prose{มคฺคเหตุกํ กุสลญฺจ มคฺคาธิปติํ กุสลญฺจ ธมฺมํ ปฏิจฺจ มคฺคเหตุโก กุสโล ธมฺโม อุปฺปชฺชติ เหตุปจฺจยา.\csromanspacer{} ตีณิ. (สํขิตฺตํ.)}

\csromanheader{2}{เหตุยา สตฺตรส ฯเปฯ อวิคเต สตฺตรส. (สํขิตฺตํ.)}
\proseitem{121}{มคฺคารมฺมณํ อพฺยากตํ ธมฺมํ ปฏิจฺจ มคฺคารมฺมโณ อพฺยากโต ธมฺโม อุปฺปชฺชติ เหตุปจฺจยา.\csromanspacer{} ตีณิ.}

\prose{มคฺคาธิปติํ อพฺยากตํ ธมฺมํ ปฏิจฺจ มคฺคาธิปติ อพฺยากโต ธมฺโม อุปฺปชฺชติ เหตุปจฺจยา.\csromanspacer{} ตีณิ.}

\prose{(สพฺพตฺถ วิตฺถาโร.)}

\csromanheader{2}{17. อุปฺปนฺนตฺติก 1. กุสลตฺติก}
\csromanheader{2}{7. ปญฺหาวาร-เหตุ}
\proseitem{122}{อุปฺปนฺโน กุสโล ธมฺโม อุปฺปนฺนสฺส กุสลสฺส ธมฺมสฺส เหตุปจฺจเยน ปจฺจโย. (สพฺพตฺถ เอกํ.\csromanspacer{} สพฺพตฺถ วิตฺถาโร.)}

\prose{อุปฺปนฺโน อกุสโล ธมฺโม อุปฺปนฺนสฺส อกุสลสฺส ธมฺมสฺส เหตุปจฺจเยน ปจฺจโย. (สพฺพตฺถ เอกํ.\csromanspacer{} สพฺพตฺถ วิตฺถาโร.)}

\prose{อุปฺปนฺโน อพฺยากโต ธมฺโม อุปฺปนฺนสฺส อพฺยากตสฺส ธมฺมสฺส เหตุปจฺจเยน ปจฺจโย. (สํขิตฺตํ.)}

\prose{เหตุยา เอกํ, อารมฺมเณ ตีณิ ฯเปฯ อุปนิสฺสเย ตีณิ ฯเปฯ อวิคเต เอกํ. (สพฺพตฺถ วิตฺถาโร.)}

\csromanheader{2}{18. อตีตตฺติก 1. กุสลตฺติก}
\csromanheader{2}{7. ปญฺหาวาร-เหตุ}
\proseitem{123}{ปจฺจุปฺปนฺโน กุสโล ธมฺโม ปจฺจุปฺปนฺนสฺส กุสลสฺส ธมฺมสฺส เหตุปจฺจเยน ปจฺจโย. (สพฺพตฺถ เอกํ.)}

\prose{\csromanfolio{36}ปจฺจุปฺปนฺโน อกุสโล ธมฺโม ปจฺจุปฺปนฺนสฺส อกุสลสฺส ธมฺมสฺส เหตุปจฺจเยน ปจฺจโย. (สพฺพตฺถ เอกํ.\csromanspacer{} สพฺพตฺถ วิตฺถาโร.)}

\prose{ปจฺจุปฺปนฺโน อพฺยากโต ธมฺโม ปจฺจุปฺปนฺนสฺส อพฺยากตสฺส ธมฺมสฺส เหตุปจฺจเยน ปจฺจโย.}

\csromanheader{2}{เหตุยา เอกํ. (สพฺพตฺถ วิตฺถาโร.)}
\csromanheader{2}{19. อตีตารมฺมณตฺติก 1. กุสลตฺติก}
\csromanheader{2}{1. ปฏิจฺจวาร-เหตุ}
\csromanheader{2}{124. อตีตารมฺมณํ กุสลํ ธมฺมํ ปฏิจฺจ อตีตารมฺมโณ กุสโล}
\prose{อนาคตารมฺมณํ กุสลํ ธมฺมํ ปฏิจฺจ อนาคตารมฺมโณ กุสโล}

\prose{ปจฺจุปฺปนฺนารมฺมณํ กุสลํ ธมฺมํ ปฏิจฺจ ปจฺจุปฺปนฺนารมฺมโณ กุสโล ธมฺโม อุปฺปชฺชติ เหตุปจฺจยา. (1) (สพฺพตฺถ ตีณิ.\csromanspacer{} สพฺพตฺถ วิตฺถาโร.)}

\proseitem{125}{อตีตารมฺมณํ อกุสลํ ธมฺมํ ปฏิจฺจ อตีตารมฺมโณ อกุสโล ธมฺโม อุปฺปชฺชติ เหตุปจฺจยา. (1)}

\prose{อนาคตารมฺมณํ อกุสลํ ธมฺมํ ปฏิจฺจ อนาคตารมฺมโณ อกุสโล ธมฺโม อุปฺปชฺชติ เหตุปจฺจยา. (1)}

\prose{ปจฺจุปฺปนฺนารมฺมณํ อกุสลํ ธมฺมํ ปฏิจฺจ ปจฺจุปฺปนฺนารมฺมโณ อกุสโล ธมฺโม อุปฺปชฺชติ เหตุปจฺจยา. (1) (สพฺพตฺถ ตีณิ.\csromanspacer{} สพฺพตฺถ วิตฺถาโร.)}

\proseitem{126}{อตีตารมฺมณํ อพฺยากตํ ธมฺมํ ปฏิจฺจ อตีตารมฺมโณ อพฺยากโต ธมฺโม อุปฺปชฺชติ เหตุปจฺจยา. (1)}

\prose{อนาคตารมฺมณํ อพฺยากตํ ธมฺมํ ปฏิจฺจ อนาคตารมฺมโณ อพฺยากโต ธมฺโม อุปฺปชฺชติ เหตุปจฺจยา. (1)}

\prose{ปจฺจุปฺปนฺนารมฺมณํ อพฺยากตํ ธมฺมํ ปฏิจฺจ ปจฺจุปฺปนฺนารมฺมโณ อพฺยากโต ธมฺโม อุปฺปชฺชติ เหตุปจฺจยา. (1) (สพฺพตฺถ ตีณิ.\csromanspacer{} สพฺพตฺถ วิตฺถาโร.)}

\vspace{\tipitakaparskip}
\csromancenter{สนิทสฺสนตฺติก 1.\csromanspacer{} กุสลตฺติก \textbf{20.}\csromanspacer{}\textbf{ อชฺฌตฺตตฺติก 1.}\csromanspacer{}\textbf{ กุสลตฺติก}}
\csromanheader{2}{1. ปฏิจฺจวาร-เหตุ}
\proseitem{127}{\csromanfolio{37}อชฺฌตฺตํ กุสลํ ธมฺมํ ปฏิจฺจ อชฺฌตฺโต กุสโล ธมฺโม อุปฺปชฺชติ เหตุปจฺจยา. (1)}

\prose{พหิทฺธา กุสลํ ธมฺมํ ปฏิจฺจ พหิทฺธา กุสโล ธมฺโม อุปฺปชฺชติ เหตุปจฺจยา. (1) (สพฺพตฺถ เทฺว.\csromanspacer{} สพฺพตฺถ วิตฺถาโร.)}

\proseitem{128}{อชฺฌตฺตํ อกุสลํ ธมฺมํ ปฏิจฺจ อชฺฌตฺโต อกุสโล ธมฺโม อุปฺปชฺชติ เหตุปจฺจยา. (1)}

\prose{พหิทฺธา อกุสลํ ธมฺมํ ปฏิจฺจ พหิทฺธา อกุสโล ธมฺโม อุปฺปชฺชติ เหตุปจฺจยา. (1) (สพฺพตฺถ เทฺว.\csromanspacer{} สพฺพตฺถ วิตฺถาโร.)}

\proseitem{129}{อชฺฌตฺตํ อพฺยากตํ ธมฺมํ ปฏิจฺจ อชฺฌตฺโต อพฺยากโต ธมฺโม อุปฺปชฺชติ เหตุปจฺจยา. (1)}

\prose{พหิทฺธา อพฺยากตํ ธมฺมํ ปฏิจฺจ พหิทฺธา อพฺยากโต ธมฺโม อุปฺปชฺชติ เหตุปจฺจยา. (1) (สพฺพตฺถ เทฺว.\csromanspacer{} สพฺพตฺถ วิตฺถาโร.)}

\csromanheader{2}{21. อชฺฌตฺตารมฺมณตฺติก 1. กุสลตฺติก}
\csromanheader{2}{1. ปฏิจฺจวาร-เหตุ}
\proseitem{130}{อชฺฌตฺตารมฺมณํ กุสลํ ธมฺมํ ปฏิจฺจ อชฺฌตฺตารมฺมโณ กุสโล ธมฺโม อุปฺปชฺชติ เหตุปจฺจยา. (สพฺพตฺถ เทฺว.\csromanspacer{} สพฺพตฺถ วิตฺถาโร.)}

\prose{อชฺฌตฺตารมฺมณํ อกุสลํ ธมฺมํ ปฏิจฺจ อชฺฌตฺตารมฺมโณ อกุสโล ธมฺโม อุปฺปชฺชติ เหตุปจฺจยา. (สพฺพตฺถ เทฺว.\csromanspacer{} สพฺพตฺถ วิตฺถาโร.)}

\prose{อชฺฌตฺตารมฺมณํ อพฺยากตํ ธมฺมํ ปฏิจฺจ อชฺฌตฺตารมฺมโณ อพฺยากโต ธมฺโม อุปฺปชฺชติ เหตุปจฺจยา. (สพฺพตฺถ เทฺว.\csromanspacer{} สพฺพตฺถ วิตฺถาโร.)}

\csromanheader{2}{22. สนิทสฺสนตฺติก 1. กุสลตฺติก}
\csromanheader{2}{1. ปฏิจฺจวาร-เหตุ}
\proseitem{131}{อนิทสฺสน-อปฺปฏิฆํ กุสลํ ธมฺมํ ปฏิจฺจ อนิทสฺสน-อปฺปฏิโฆ กุสโล ธมฺโม อุปฺปชฺชติ เหตุปจฺจยา. (สพฺพตฺถ เอกํ.)}

\prose{\csromanfolio{38}อนิทสฺสน-อปฺปฏิฆํ อกุสลํ ธมฺมํ ปฏิจฺจ อนิทสฺสน-อปฺปฏิโฆ อกุสโล ธมฺโม อุปฺปชฺชติ เหตุปจฺจยา. (สพฺพตฺถ เอกํ.)}

\proseitem{132}{อนิทสฺสนสปฺปฏิฆํ อพฺยากตํ ธมฺมํ ปฏิจฺจ อนิทสฺสนสปฺปฏิโฆ อพฺยากโต ธมฺโม อุปฺปชฺชติ เหตุปจฺจยา. (เอกํ.) .\csromanspacer{} อนิทสฺสนสปฺปฏิฆํ อพฺยากตํ ธมฺมํ ปฏิจฺจ สนิทสฺสนสปฺปฏิโฆ อพฺยากโต ธมฺโม อุปฺปชฺชติ เหตุปจฺจยา. (เทฺว.) .\csromanspacer{} อนิทสฺสนสปฺปฏิฆํ อพฺยากตํ ธมฺมํ ปฏิจฺจ อนิทสฺสน-อปฺปฏิโฆ อพฺยากโต ธมฺโม อุปฺปชฺชติ เหตุปจฺจยา. (ตีณิ.) .\csromanspacer{} อนิทสฺสนสปฺปฏิฆํ อพฺยากตํ ธมฺมํ ปฏิจฺจ สนิทสฺสนสปฺปฏิโฆ อพฺยากโต จ อนิทสฺสน-อปฺปฏิโฆ อพฺยากโต จ ธมฺมา อุปฺปชฺชนฺติ เหตุปจฺจยา. (จตฺตาริ.) .\csromanspacer{} อนิทสฺสนสปฺปฏิฆํ อพฺยากตํ ธมฺมํ ปฏิจฺจ อนิทสฺสนสปฺปฏิโฆ อพฺยากโต จ อนิทสฺสน-อปฺปฏิโฆ อพฺยากโต จ ธมฺมา อุปฺปชฺชนฺติ เหตุปจฺจยา. (ปญฺจ.) .\csromanspacer{} อนิทสฺสนสปฺปฏิฆํ อพฺยากตํ ธมฺมํ ปฏิจฺจ สนิทสฺสนสปฺปฏิโฆ อพฺยากโต จ อนิทสฺสนสปฺปฏิโฆ อพฺยากโต จ ธมฺมา อุปฺปชฺชนฺติ เหตุปจฺจยา. (ฉ.) .\csromanspacer{} อนิทสฺสนสปฺปฏิฆํ อพฺยากตํ ธมฺมํ ปฏิจฺจ สนิทสฺสนสปฺปฏิโฆ อพฺยากโต จ อนิทสฺสนสปฺปฏิโฆ อพฺยากโต จ อนิทสฺสน-อปฺปฏิโฆ อพฺยากโต จ ธมฺมา อุปฺปชฺชนฺติ เหตุปจฺจยา. (สตฺต.)}

\proseitem{133}{อนิทสฺสน-อปฺปฏิฆํ อพฺยากตํ ธมฺมํ ปฏิจฺจ อนิทสฺสน-อปฺปฏิโฆ อพฺยากโต ธมฺโม อุปฺปชฺชติ เหตุปจฺจยา. (สตฺต.) .\csromanspacer{} อนิทสฺสน-อปฺปฏิฆํ อพฺยากตญฺจ อนิทสฺสนสปฺปฏิฆํ อพฺยากตญฺจ ธมฺมํ ปฏิจฺจ สนิทสฺสนสปฺปฏิโฆ อพฺยากโต ธมฺโม อุปฺปชฺชติ เหตุปจฺจยา. (สตฺต.) (สํขิตฺตํ.)}

\proseitem{134}{เหตุยา เอกวีส, อวิคเต เอกวีส. (สํขิตฺตํ.) (สหชาตวารมฺปิ ฯเปฯ ปญฺหาวารมฺปิ วิตฺถาเรตพฺพํ.)}

\nitthitam{\textbf{ธมฺมานุโลเม ติกติกปฏฺฐานํ นิฏฺฐิตํ.}}
\pitaka{\textbf{อภิธมฺมปิฏก}}
\csromanheader{2}{\textbf{ธมฺมานุโลม}}
\csromanmatikaanchor{mtk.1}
\csromanmatikaanchor{mtk.3}
\csromantocmarknopage{section}{ปญฺจมภาค}{mtk.1}
\bookmark[dest={mtk.1},level=1]{ปญฺจมภาค}
\gambhira{\textbf{ทุกทุกปฏฺฐานปาฬิ}}
\namakkaram{นโม ตสฺส ภควโต อรหโต สมฺมาสมฺพุทฺธสฺส.}
\csromanheader{2}{1. เหตุทุก 1. สเหตุกทุก}
\csromanheader{2}{\textbf{สเหตุกปท 1-7. ปฏิจฺจาทิวาร}}
\csromanheader{2}{\textbf{ปจฺจยจตุกฺก-เหตุ}}
\proseitem{1}{\csromanfolio{39}เหตุํ สเหตุกํ ธมฺมํ ปฏิจฺจ เหตุ สเหตุโก ธมฺโม อุปฺปชฺชติ เหตุปจฺจยา.\csromanspacer{}\csromanspacer{}เหตุํ สเหตุกํ ธมฺมํ ปฏิจฺจ นเหตุ สเหตุโก ธมฺโม อุปฺปชฺชติ เหตุปจฺจยา.\csromanspacer{}\csromanspacer{}เหตุํ สเหตุกํ ธมฺมํ ปฏิจฺจ เหตุ สเหตุโก จ นเหตุ สเหตุโก จ ธมฺมา อุปฺปชฺชนฺติ เหตุปจฺจยา. (3)}

\proseitem{2}{นเหตุํ สเหตุกํ ธมฺมํ ปฏิจฺจ นเหตุ สเหตุโก ธมฺโม อุปฺปชฺชติ เหตุปจฺจยา.\csromanspacer{}\csromanspacer{}นเหตุํ สเหตุกํ ธมฺมํ ปฏิจฺจ เหตุ สเหตุโก ธมฺโม อุปฺปชฺชติ เหตุปจฺจยา.\csromanspacer{}\csromanspacer{}นเหตุํ สเหตุกํ ธมฺมํ ปฏิจฺจ เหตุ สเหตุโก จ นเหตุ สเหตุโก จ ธมฺมา อุปฺปชฺชนฺติ เหตุปจฺจยา. (3)}

\proseitem{3}{เหตุํ สเหตุกญฺจ นเหตุํ สเหตุกญฺจ ธมฺมํ ปฏิจฺจ เหตุ สเหตุโก ธมฺโม อุปฺปชฺชติ เหตุปจฺจยา.\csromanspacer{}\csromanspacer{}เหตุํ สเหตุกญฺจ นเหตุํ สเหตุกญฺจ ธมฺมํ ปฏิจฺจ นเหตุ สเหตุโก ธมฺโม อุปฺปชฺชติ \csromanfolio{40}เหตุปจฺจยา.\csromanspacer{}\csromanspacer{}เหตุํ สเหตุกญฺจ นเหตุํ สเหตุกญฺจ ธมฺมํ ปฏิจฺจ เหตุ สเหตุโก จ นเหตุ สเหตุโก จ ธมฺมา อุปฺปชฺชนฺติ เหตุปจฺจยา. (3)}

\proseitem{4}{เหตุยา นว, อารมฺมเณ นว ฯเปฯ อวิคเต นว. (สํขิตฺตํ.)}

\csromanheader{2}{\textbf{ปจฺจนีย-น-อธิปติ}}
\proseitem{5}{เหตุํ สเหตุกํ ธมฺมํ ปฏิจฺจ เหตุ สเหตุโก ธมฺโม อุปฺปชฺชติ น-อธิปติปจฺจยา. (สํขิตฺตํ.)}

\prose{น-อธิปติยา นว, นปุเรชาเต นว, นปจฺฉาชาเต นว, นกมฺเม ตีณิ, นวิปาเก นว, นวิปฺปยุตฺเต นว. (สํขิตฺตํ.)}

\csromanheader{2}{เหตุปจฺจยา น-อธิปติยา นว. (สํขิตฺตํ.)}
\csromanheader{2}{น-อธิปติปจฺจยา เหตุยา นว. (สํขิตฺตํ.)}
\prose{(สหชาตวารมฺปิ ปจฺจยวารมฺปิ นิสฺสยวารมฺปิ สํสฏฺฐวารมฺปิ สมฺปยุตฺตวารมฺปิ ปฏิจฺจวารสทิสํ.)}

\csromanheader{2}{\textbf{เหตุ อารมฺมณ}}
\proseitem{6}{เหตุ สเหตุโก ธมฺโม เหตุสฺส สเหตุกสฺส ธมฺมสฺส เหตุปจฺจเยน ปจฺจโย.\csromanspacer{}\csromanspacer{}เหตุ สเหตุโก ธมฺโม นเหตุสฺส สเหตุกสฺส ธมฺมสฺส เหตุปจฺจเยน ปจฺจโย.\csromanspacer{}\csromanspacer{}เหตุ สเหตุโก ธมฺโม เหตุสฺส สเหตุกสฺส จ นเหตุสฺส สเหตุกสฺส จ ธมฺมสฺส เหตุปจฺจเยน ปจฺจโย. (3)}

\proseitem{7}{เหตุ สเหตุโก ธมฺโม เหตุสฺส สเหตุกสฺส ธมฺมสฺส อารมฺมณปจฺจเยน ปจฺจโย.\csromanspacer{}\csromanspacer{}เหตุ สเหตุโก ธมฺโม นเหตุสฺส สเหตุกสฺส ธมฺมสฺส อารมฺมณปจฺจเยน ปจฺจโย.\csromanspacer{}\csromanspacer{}เหตุ สเหตุโก ธมฺโม เหตุสฺส สเหตุกสฺส จ นเหตุสฺส สเหตุกสฺส จ ธมฺมสฺส อารมฺมณปจฺจเยน ปจฺจโย. (3)}

\prose{นเหตุ สเหตุโก ธมฺโม นเหตุสฺส สเหตุกสฺส ธมฺมสฺส อารมฺมณปจฺจเยน ปจฺจโย.\csromanspacer{}\csromanspacer{}นเหตุ สเหตุโก ธมฺโม เหตุสฺส สเหตุกสฺส ธมฺมสฺส อารมฺมณปจฺจเยน ปจฺจโย.\csromanspacer{}\csromanspacer{}นเหตุ สเหตุโก ธมฺโม เหตุสฺส สเหตุกสฺส จ นเหตุสฺส สเหตุกสฺส จ ธมฺมสฺส อารมฺมณปจฺจเยน ปจฺจโย. (3)}

\csromanheader{2}{1. เหตุทุก 1. สเหตุกทุก}
\prose{\csromanfolio{41}เหตุ สเหตุโก จ นเหตุ สเหตุโก จ ธมฺมา เหตุสฺส สเหตุกสฺส ธมฺมสฺส อารมฺมณปจฺจเยน ปจฺจโย.\csromanspacer{}\csromanspacer{}เหตุ สเหตุโก จ นเหตุ สเหตุโก จ ธมฺมา นเหตุสฺส สเหตุกสฺส ธมฺมสฺส อารมฺมณปจฺจเยน ปจฺจโย.\csromanspacer{}\csromanspacer{}เหตุ สเหตุโก จ นเหตุ สเหตุโก จ ธมฺมา เหตุสฺส สเหตุกสฺส จ นเหตุสฺส สเหตุกสฺส จ ธมฺมสฺส อารมฺมณปจฺจเยน ปจฺจโย. (3) (สํขิตฺตํ.)}

\proseitem{8}{เหตุยา ตีณิ, อารมฺมเณ นว ฯเปฯ อุปนิสฺสเย นว, อวิคเต นว. (ยถา กุสลตฺติเก ปญฺหาวารํ, เอวํ วิตฺถาเรตพฺพํ.)}

\csromanheader{2}{\textbf{อเหตุกปท ปจฺจยจตุกฺก}}
\proseitem{9}{เหตุํ อเหตุกํ ธมฺมํ ปฏิจฺจ นเหตุ อเหตุโก ธมฺโม อุปฺปชฺชติ เหตุปจฺจยา. (1)}

\prose{นเหตุํ อเหตุกํ ธมฺมํ ปฏิจฺจ นเหตุ อเหตุโก ธมฺโม อุปฺปชฺชติ เหตุปจฺจยา. (1)}

\prose{เหตุํ อเหตุกญฺจ นเหตุํ อเหตุกญฺจ ธมฺมํ ปฏิจฺจ นเหตุ อเหตุโก ธมฺโม อุปฺปชฺชติ เหตุปจฺจยา. (1) (สํขิตฺตํ.)}

\prose{เหตุยา ตีณิ, อารมฺมเณ เอกํ, อวิคเต ตีณิ. (สํขิตฺตํ.\csromanspacer{} สหชาตวารมฺปิ ฯเปฯ สมฺปยุตฺตวารมฺปิ ปฏิจฺจวารสทิสํ.)}

\csromanheader{2}{10. เหตุ อเหตุโก ธมฺโม นเหตุสฺส อเหตุกสฺส ธมฺมสฺส}
\prose{เหตุ อเหตุโก ธมฺโม เหตุสฺส อเหตุกสฺส ธมฺมสฺส อารมฺมณปจฺจเยน ปจฺจโย.\csromanspacer{}\csromanspacer{}เหตุ อเหตุโก ธมฺโม นเหตุสฺส อเหตุกสฺส ธมฺมสฺส อารมฺมณปจฺจเยน ปจฺจโย. (2)}

\prose{นเหตุ อเหตุโก ธมฺโม นเหตุสฺส อเหตุกสฺส ธมฺมสฺส อารมฺมณปจฺจเยน ปจฺจโย.\csromanspacer{}\csromanspacer{}นเหตุ อเหตุโก ธมฺโม เหตุสฺส อเหตุกสฺส ธมฺมสฺส อารมฺมณปจฺจเยน ปจฺจโย. (2) (สํขิตฺตํ.)}

\proseitem{11}{\csromanfolio{42}เหตุยา เอกํ, อารมฺมเณ จตฺตาริ, อวิคเต จตฺตาริ. (สํขิตฺตํ.) (ยถา กุสลตฺติเก ปญฺหาวารํ, เอวํ วิตฺถาเรตพฺพํ.)}

\csromanheader{2}{1. เหตุทุก 2. เหตุสมฺปยุตฺตทุก}
\csromanheader{2}{1-7. ปฏิจฺจาทิวาร ปจฺจยจตุกฺก-เหตุ}
\proseitem{12}{เหตุํ เหตุสมฺปยุตฺตํ ธมฺมํ ปฏิจฺจ เหตุ เหตุสมฺปยุตฺโต ธมฺโม อุปฺปชฺชติ เหตุปจฺจยา.\csromanspacer{}\csromanspacer{}เหตุํ เหตุสมฺปยุตฺตํ ธมฺมํ ปฏิจฺจ นเหตุ เหตุสมฺปยุตฺโต ธมฺโม อุปฺปชฺชติ เหตุปจฺจยา.\csromanspacer{}\csromanspacer{}เหตุํ เหตุสมฺปยุตฺตํ ธมฺมํ ปฏิจฺจ เหตุ เหตุสมฺปยุตฺโต จ นเหตุ เหตุสมฺปยุตฺโต จ ธมฺมา อุปฺปชฺชนฺติ เหตุปจฺจยา. (3)}

\prose{นเหตุํ เหตุสมฺปยุตฺตํ ธมฺมํ ปฏิจฺจ นเหตุ เหตุสมฺปยุตฺโต ธมฺโม อุปฺปชฺชติ เหตุปจฺจยา.\csromanspacer{} ตีณิ.}

\prose{เหตุํ เหตุสมฺปยุตฺตญฺจ นเหตุํ เหตุสมฺปยุตฺตญฺจ ธมฺมํ ปฏิจฺจ เหตุ เหตุสมฺปยุตฺโต ธมฺโม อุปฺปชฺชติ เหตุปจฺจยา.\csromanspacer{} ตีณิ.}

\proseitem{13}{เหตุยา นว, อารมฺมเณ นว ฯเปฯ อวิคเต นว. (สํขิตฺตํ.)}

\proseitem{14}{เหตุ เหตุสมฺปยุตฺโต ธมฺโม เหตุสฺส เหตุสมฺปยุตฺตสฺส ธมฺมสฺส เหตุปจฺจเยน ปจฺจโย.\csromanspacer{} ตีณิ. (สํขิตฺตํ.)}

\csromanheader{2}{เหตุยา ตีณิ ฯเปฯ อวิคเต นว. (สํขิตฺตํ.)}
\prose{(ยถา กุสลตฺติเก ปญฺหาวารํ, เอวํ วิตฺถาเรตพฺพํ.)}

\csromanheader{2}{15. เหตุํ เหตุวิปฺปยุตฺตํ ธมฺมํ ปฏิจฺจ นเหตุ เหตุวิปฺปยุตฺโต}
\prose{นเหตุํ เหตุวิปฺปยุตฺตํ ธมฺมํ ปฏิจฺจ นเหตุ เหตุวิปฺปยุตฺโต}

\csromanheader{2}{1. เหตุทุก 5. นเหตุสเหตุกทุก}
\prose{\csromanfolio{43}เหตุํ เหตุวิปฺปยุตฺตญฺจ นเหตุํ เหตุวิปฺปยุตฺตญฺจ ธมฺมํ ปฏิจฺจ นเหตุ เหตุวิปฺปยุตฺโต ธมฺโม อุปฺปชฺชติ เหตุปจฺจยา. (1)}

\prose{เหตุยา ตีณิ, อารมฺมเณ เอกํ ฯเปฯ อวิคเต ตีณิ.}

\prose{(สหชาตวารมฺปิ ฯเปฯ สมฺปยุตฺตวารมฺปิ ปฏิจฺจวารสทิสํ วิตฺถาเรตพฺพํ.)}

\proseitem{16}{เหตุ เหตุวิปฺปยุตฺโต ธมฺโม นเหตุสฺส เหตุวิปฺปยุตฺตสฺส ธมฺมสฺส เหตุปจฺจเยน ปจฺจโย. (สํขิตฺตํ.)}

\prose{เหตุยา เอกํ, อารมฺมเณ จตฺตาริ ฯเปฯ อวิคเต จตฺตาริ. (สํขิตฺตํ.) (ยถา กุสลตฺติเก ปญฺหาวารํ, เอวํ วิตฺถาเรตพฺพํ.)}

\csromanheader{2}{1. เหตุทุก 3. เหตุสเหตุกทุก}
\csromanheader{2}{1. ปฏิจฺจวาร ปจฺจยจตุกฺก-เหตุ}
\proseitem{17}{เหตุํ เหตุญฺเจว สเหตุกญฺจ ธมฺมํ ปฏิจฺจ เหตุ เหตุ เจว สเหตุโก ธมฺโม อุปฺปชฺชติ เหตุปจฺจยา. (สพฺพตฺถ เอกํ.)}

\prose{นเหตุํ สเหตุกญฺเจว น จ เหตุํ ธมฺมํ ปฏิจฺจ นเหตุ สเหตุโก เจว น จ เหตุ ธมฺโม อุปฺปชฺชติ เหตุปจฺจยา. (สพฺพตฺถ เอกํ.)}

\csromanheader{2}{1. เหตุทุก 4. เหตุเหตุสมฺปยุตฺตทุก}
\csromanheader{2}{1. ปฏิจฺจวาร ปจฺจยจตุกฺก-เหตุ}
\proseitem{18}{เหตุํ เหตุญฺเจว เหตุสมฺปยุตฺตญฺจ ธมฺมํ ปฏิจฺจ เหตุ เหตุ เจว เหตุสมฺปยุตฺโต จ ธมฺโม อุปฺปชฺชติ เหตุปจฺจยา. (สพฺพตฺถ เอกํ.)}

\prose{นเหตุํ เหตุสมฺปยุตฺตญฺเจว น จ เหตุํ ธมฺมํ ปฏิจฺจ นเหตุ เหตุสมฺปยุตฺโต เจว น จ เหตุ ธมฺโม อุปฺปชฺชติ เหตุปจฺจยา. (สพฺพตฺถ เอกํ.)}

\csromanheader{2}{1. เหตุทุก 5. นเหตุสเหตุกทุก}
\proseitem{19}{นเหตุํ สเหตุกํ ธมฺมํ ปฏิจฺจ นเหตุ สเหตุโก ธมฺโม อุปฺปชฺชติ เหตุปจฺจยา. (สพฺพตฺถ เอกํ.)}

\prose{\csromanfolio{44}นเหตุํ อเหตุกํ ธมฺมํ ปฏิจฺจ นเหตุ อเหตุโก ธมฺโม อุปฺปชฺชติ เหตุปจฺจยา. (สพฺพตฺถ เอกํ.)}

\nitthitam{เหตุโคจฺฉกํ นิฏฺฐิตํ.}
\csromanheader{2}{1. เหตุทุก 6. สปฺปจฺจยทุก}
\proseitem{20}{เหตุํ สปฺปจฺจยํ ธมฺมํ ปฏิจฺจ เหตุ สปฺปจฺจโย ธมฺโม อุปฺปชฺชติ เหตุปจฺจยา.\csromanspacer{} ตีณิ.}

\prose{นเหตุํ สปฺปจฺจยํ ธมฺมํ ปฏิจฺจ นเหตุ สปฺปจฺจโย ธมฺโม อุปฺปชฺชติ เหตุปจฺจยา.\csromanspacer{} ตีณิ.}

\prose{เหตุํ สปฺปจฺจยญฺจ นเหตุํ สปฺปจฺจยญฺจ ธมฺมํ ปฏิจฺจ เหตุ สปฺปจฺจโย ธมฺโม อุปฺปชฺชติ เหตุปจฺจยา.\csromanspacer{} ตีณิ. (สํขิตฺตํ.)}

\prose{เหตุยา นว, อารมฺมเณ นว ฯเปฯ อวิคเต นว. (สํขิตฺตํ.) (สหชาตวารมฺปิ ฯเปฯ สมฺปยุตฺตวารมฺปิ ปฏิจฺจวารสทิสํ.)}

\proseitem{21}{เหตุ สปฺปจฺจโย ธมฺโม เหตุสฺส สปฺปจฺจยสฺส ธมฺมสฺส เหตุปจฺจเยน ปจฺจโย. (สํขิตฺตํ.)}

\csromanheader{2}{เหตุยา ตีณิ. (สํขิตฺตํ. ปญฺหาวารมฺปิ เอวํ วิตฺถาเรตพฺพํ.)}
\csromanheader{2}{1. เหตุทุก 7. สงฺขตทุก}
\proseitem{22}{เหตุํ สงฺขตํ ธมฺมํ ปฏิจฺจ เหตุ สงฺขโต ธมฺโม อุปฺปชฺชติ เหตุปจฺจยา.}

\csromanheader{2}{เหตุยา นว ฯเปฯ อวิคเต นว. (สปฺปจฺจยทุกสทิสํ.)}
\csromanheader{2}{1. เหตุทุก 8. สนิทสฺสนทุก}
\proseitem{23}{เหตุํ อนิทสฺสนํ ธมฺมํ ปฏิจฺจ เหตุ อนิทสฺสโน ธมฺโม อุปฺปชฺชติ เหตุปจฺจยา.\csromanspacer{} ตีณิ.}

\csromanheader{2}{1. เหตุทุก 10. รูปีทุก}
\prose{\csromanfolio{45}นเหตุํ อนิทสฺสนํ ธมฺมํ ปฏิจฺจ นเหตุ อนิทสฺสโน ธมฺโม อุปฺปชฺชติ เหตุปจฺจยา.\csromanspacer{} ตีณิ.}

\prose{เหตุํ อนิทสฺสนญฺจ นเหตุํ อนิทสฺสนญฺจ ธมฺมํ ปฏิจฺจ เหตุ อนิทสฺสโน ธมฺโม อุปฺปชฺชติ เหตุปจฺจยา.\csromanspacer{} ตีณิ. (สํขิตฺตํ.)}

\prose{เหตุยา นว ฯเปฯ อวิคเต นว. (สํขิตฺตํ.) (สหชาตวารมฺปิ ฯเปฯ ปญฺหาวารมฺปิ วิตฺถาเรตพฺพํ.)}

\csromanheader{2}{1. เหตุทุก 9. สปฺปฏิฆทุก}
\proseitem{24}{นเหตุํ สปฺปฏิฆํ ธมฺมํ ปฏิจฺจ นเหตุ สปฺปฏิโฆ ธมฺโม อุปฺปชฺชติ เหตุปจฺจยา. (สพฺพตฺถ เอกํ.)}

\proseitem{25}{เหตุํ อปฺปฏิฆํ ธมฺมํ ปฏิจฺจ เหตุ อปฺปฏิโฆ ธมฺโม อุปฺปชฺชติ เหตุปจฺจยา.\csromanspacer{}\csromanspacer{}เหตุํ อปฺปฏิฆํ ธมฺมํ ปฏิจฺจ นเหตุ อปฺปฏิโฆ ธมฺโม อุปฺปชฺชติ เหตุปจฺจยา.\csromanspacer{}\csromanspacer{}เหตุํ อปฺปฏิฆํ ธมฺมํ ปฏิจฺจ เหตุ อปฺปฏิโฆ จ นเหตุ อปฺปฏิโฆ จ ธมฺมา อุปฺปชฺชนฺติ เหตุปจฺจยา. (3)}

\prose{นเหตุํ อปฺปฏิฆํ ธมฺมํ ปฏิจฺจ นเหตุ อปฺปฏิโฆ ธมฺโม อุปฺปชฺชติ เหตุปจฺจยา.\csromanspacer{} ตีณิ.}

\prose{เหตุํ อปฺปฏิฆญฺจ นเหตุํ อปฺปฏิฆญฺจ ธมฺมํ ปฏิจฺจ เหตุ อปฺปฏิโฆ ธมฺโม อุปฺปชฺชติ เหตุปจฺจยา.\csromanspacer{} ตีณิ. (สํขิตฺตํ.)}

\prose{เหตุยา นว, อารมฺมเณ นว ฯเปฯ อวิคเต นว. (สํขิตฺตํ.) (สหชาตวารมฺปิ ฯเปฯ ปญฺหาวารมฺปิ วิตฺถาเรตพฺพํ.)}

\csromanheader{2}{1. เหตุทุก 10. รูปีทุก}
\proseitem{26}{นเหตุํ รูปิํ ธมฺมํ ปฏิจฺจ นเหตุ รูปี ธมฺโม อุปฺปชฺชติ เหตุปจฺจยา. (สพฺพตฺถ เอกํ.)}

\proseitem{27}{เหตุํ อรูปิํ ธมฺมํ ปฏิจฺจ เหตุ อรูปี ธมฺโม อุปฺปชฺชติ เหตุปจฺจยา.\csromanspacer{} ตีณิ.}

\prose{นเหตุํ อรูปิํ ธมฺมํ ปฏิจฺจ นเหตุ อรูปี ธมฺโม อุปฺปชฺชติ เหตุปจฺจยา.\csromanspacer{} ตีณิ.}

\prose{\csromanfolio{46}เหตุํ อรูปิญฺจ นเหตุํ อรูปิญฺจ ธมฺมํ ปฏิจฺจ เหตุ อรูปี ธมฺโม อุปฺปชฺชติ เหตุปจฺจยา.\csromanspacer{} ตีณิ. (สํขิตฺตํ.)}

\prose{เหตุยา นว, อารมฺมเณ นว ฯเปฯ อวิคเต นว. (สํขิตฺตํ.) (สหชาตวารมฺปิ ฯเปฯ ปญฺหาวารมฺปิ วิตฺถาเรตพฺพํ.)}

\csromanheader{2}{1. เหตุทุก 11. โลกิยทุก}
\proseitem{28}{เหตุํ โลกิยํ ธมฺมํ ปฏิจฺจ เหตุ โลกิโย ธมฺโม อุปฺปชฺชติ เหตุปจฺจยา.\csromanspacer{} ตีณิ.}

\prose{นเหตุํ โลกิยํ ธมฺมํ ปฏิจฺจ นเหตุ โลกิโย ธมฺโม อุปฺปชฺชติ เหตุปจฺจยา.\csromanspacer{} ตีณิ.}

\prose{เหตุํ โลกิยญฺจ นเหตุํ โลกิยญฺจ ธมฺมํ ปฏิจฺจ เหตุ โลกิโย ธมฺโม อุปฺปชฺชติ เหตุปจฺจยา.\csromanspacer{} ตีณิ. (สํขิตฺตํ.)}

\prose{เหตุยา นว, อารมฺมเณ นว ฯเปฯ อวิคเต นว. (สํขิตฺตํ.) (สหชาตวารมฺปิ ฯเปฯ ปญฺหาวารมฺปิ เอวํ วิตฺถาเรตพฺพํ.)}

\proseitem{29}{เหตุํ โลกุตฺตรํ ธมฺมํ ปฏิจฺจ เหตุ โลกุตฺตโร ธมฺโม อุปฺปชฺชติ เหตุปจฺจยา.\csromanspacer{} ตีณิ.}

\prose{นเหตุํ โลกุตฺตรํ ธมฺมํ ปฏิจฺจ นเหตุ โลกุตฺตโร ธมฺโม อุปฺปชฺชติ เหตุปจฺจยา.\csromanspacer{} ตีณิ.}

\prose{เหตุํ โลกุตฺตรญฺจ นเหตุํ โลกุตฺตรญฺจ ธมฺมํ ปฏิจฺจ เหตุ โลกุตฺตโร ธมฺโม อุปฺปชฺชติ เหตุปจฺจยา.\csromanspacer{} ตีณิ. (สํขิตฺตํ.)}

\prose{เหตุยา นว, อารมฺมเณ นว ฯเปฯ อวิคเต นว. (สํขิตฺตํ.)}

\prose{(สหชาตวารมฺปิ ฯเปฯ สมฺปยุตฺตวารมฺปิ ปญฺหาวารมฺปิ ปฏิจฺจวารสทิสํ วิตฺถาเรตพฺพํ.)}

\csromanheader{2}{1. เหตุทุก 12. เกนจิวิญฺเญยฺยทุก}
\proseitem{30}{เหตุํ เกนจิ วิญฺเญยฺยํ ธมฺมํ ปฏิจฺจ เหตุ เกนจิ วิญฺเญยฺโย ธมฺโม อุปฺปชฺชติ เหตุปจฺจยา. (สํขิตฺตํ.)}

\csromanheader{2}{เหตุยา นว ฯเปฯ อวิคเต นว. (สํขิตฺตํ.)}
\vspace{\tipitakaparskip}
\csromancenter{(สหชาตวารมฺปิ ฯเปฯ ปญฺหาวารมฺปิ วิตฺถาเรตพฺพํ.)}
\csromanheader{2}{1. เหตุทุก 14. สาสวทุก}
\proseitem{31}{\csromanfolio{47}เหตุํ เกนจิ นวิญฺเญยฺยํ ธมฺมํ ปฏิจฺจ เหตุ เกนจิ นวิญฺเญยฺโย ธมฺโม อุปฺปชฺชติ เหตุปจฺจยา. (สํขิตฺตํ.)}

\csromanheader{2}{เหตุยา นว ฯเปฯ อวิคเต นว. (สพฺพตฺถ วิตฺถาโร.)}
\nitthitam{จูฬนฺตรทุกํ นิฏฺฐิตํ.}
\csromanheader{2}{1. เหตุทุก 13. อาสวทุก}
\proseitem{32}{เหตุํ อาสวํ ธมฺมํ ปฏิจฺจ เหตุ อาสโว ธมฺโม อุปฺปชฺชติ เหตุปจฺจยา.\csromanspacer{}\csromanspacer{}เหตุํ อาสวํ ธมฺมํ ปฏิจฺจ นเหตุ อาสโว ธมฺโม อุปฺปชฺชติ เหตุปจฺจยา.\csromanspacer{}\csromanspacer{}เหตุํ อาสวํ ธมฺมํ ปฏิจฺจ เหตุ อาสโว จ นเหตุ อาสโว จ ธมฺมา อุปฺปชฺชนฺติ เหตุปจฺจยา. (3)}

\csromanheader{2}{นเหตุํ อาสวํ ธมฺมํ ปฏิจฺจ นเหตุ อาสโว ธมฺโม อุปฺปชฺชติ}
\prose{เหตุํ อาสวญฺจ นเหตุํ อาสวญฺจ ธมฺมํ ปฏิจฺจ เหตุ อาสโว ธมฺโม อุปฺปชฺชติ เหตุปจฺจยา. (1) (สํขิตฺตํ.)}

\prose{เหตุยา ปญฺจ, อารมฺมเณ ปญฺจ ฯเปฯ อวิคเต ปญฺจ. (สพฺพตฺถ วิตฺถาโร.)}

\proseitem{33}{เหตุํ โน-อาสวํ ธมฺมํ ปฏิจฺจ เหตุ โน-อาสโว ธมฺโม อุปฺปชฺชติ เหตุปจฺจยา.\csromanspacer{} ตีณิ.}

\prose{นเหตุํ โน-อาสวํ ธมฺมํ ปฏิจฺจ นเหตุ โน-อาสโว ธมฺโม อุปฺปชฺชติ เหตุปจฺจยา.\csromanspacer{} ตีณิ.}

\prose{เหตุํ โน-อาสวญฺจ นเหตุํ โน-อาสวญฺจ ธมฺมํ ปฏิจฺจ เหตุ โน-อาสโว ธมฺโม อุปฺปชฺชติ เหตุปจฺจยา.\csromanspacer{} ตีณิ. (สํขิตฺตํ.)}

\csromanheader{2}{เหตุยา นว ฯเปฯ อวิคเต นว. (สพฺพตฺถ วิตฺถาโร.)}
\csromanheader{2}{1. เหตุทุก 14. สาสวทุก}
\proseitem{34}{เหตุํ สาสวํ ธมฺมํ ปฏิจฺจ เหตุ สาสโว ธมฺโม อุปฺปชฺชติ เหตุปจฺจยา.\csromanspacer{} ตีณิ.}

\prose{นเหตุํ สาสวํ ธมฺมํ ปฏิจฺจ นเหตุ สาสโว ธมฺโม อุปฺปชฺชติ เหตุปจฺจยา.\csromanspacer{} ตีณิ.}

\prose{\csromanfolio{48}เหตุํ สาสวญฺจ นเหตุํ สาสวญฺจ ธมฺมํ ปฏิจฺจ เหตุ สาสโว ธมฺโม อุปฺปชฺชติ เหตุปจฺจยา.\csromanspacer{} ตีณิ. (สํขิตฺตํ.)}

\csromanheader{2}{เหตุยา นว ฯเปฯ อวิคเต นว. (สพฺพตฺถ วิตฺถาโร.)}
\proseitem{35}{เหตุํ อนาสวํ ธมฺมํ ปฏิจฺจ เหตุ อนาสโว ธมฺโม อุปฺปชฺชติ เหตุปจฺจยา.\csromanspacer{} ตีณิ.}

\prose{นเหตุํ อนาสวํ ธมฺมํ ปฏิจฺจ นเหตุ อนาสโว ธมฺโม อุปฺปชฺชติ เหตุปจฺจยา.\csromanspacer{} ตีณิ.}

\prose{เหตุํ อนาสวญฺจ นเหตุํ อนาสวญฺจ ธมฺมํ ปฏิจฺจ เหตุ อนาสโว ธมฺโม อุปฺปชฺชติ เหตุปจฺจยา.\csromanspacer{} ตีณิ. (สํขิตฺตํ.)}

\csromanheader{2}{เหตุยา นว ฯเปฯ อวิคเต นว. (สพฺพตฺถ วิตฺถาโร.)}
\csromanheader{2}{1. เหตุทุก 15. อาสวสมฺปยุตฺตทุก}
\proseitem{36}{เหตุํ อาสวสมฺปยุตฺตํ ธมฺมํ ปฏิจฺจ เหตุ อาสวสมฺปยุตฺโต ธมฺโม อุปฺปชฺชติ เหตุปจฺจยา. (สํขิตฺตํ.)}

\csromanheader{2}{เหตุยา นว ฯเปฯ อวิคเต นว. (สพฺพตฺถ วิตฺถาโร.)}
\proseitem{37}{เหตุํ อาสววิปฺปยุตฺตํ ธมฺมํ ปฏิจฺจ เหตุ อาสววิปฺปยุตฺโต ธมฺโม อุปฺปชฺชติ เหตุปจฺจยา. (สํขิตฺตํ.)}

\csromanheader{2}{เหตุยา นว ฯเปฯ อวิคเต นว. (สพฺพตฺถ วิตฺถาโร.)}
\csromanheader{2}{1. เหตุทุก 16. อาสวสาสวทุก}
\proseitem{38}{เหตุํ อาสวญฺเจว สาสวญฺจ ธมฺมํ ปฏิจฺจ เหตุ อาสโว เจว สาสโว จ ธมฺโม อุปฺปชฺชติ เหตุปจฺจยา. (สํขิตฺตํ.)}

\csromanheader{2}{เหตุยา ปญฺจ ฯเปฯ อวิคเต ปญฺจ. (สพฺพตฺถ วิตฺถาโร.)}
\proseitem{39}{เหตุํ สาสวญฺเจว โน จ อาสวํ ธมฺมํ ปฏิจฺจ เหตุ สาสโว เจว โน จ อาสโว ธมฺโม อุปฺปชฺชติ เหตุปจฺจยา. (สํขิตฺตํ.)}

\csromanheader{2}{เหตุยา นว ฯเปฯ อวิคเต นว. (สพฺพตฺถ วิตฺถาโร.)}
\vspace{\tipitakaparskip}
\csromancenter{เหตุทุก 15-53.\csromanspacer{} สญฺโญชนาทิทุก \textbf{1.}\csromanspacer{}\textbf{ เหตุทุก 17.}\csromanspacer{}\textbf{ อาสว-อาสวสมฺปยุตฺตทุก}}
\proseitem{40}{\csromanfolio{49}เหตุํ อาสวญฺเจว อาสวสมฺปยุตฺตญฺจ ธมฺมํ ปฏิจฺจ เหตุ อาสโว เจว อาสวสมฺปยุตฺโต จ ธมฺโม อุปฺปชฺชติ เหตุปจฺจยา.}

\csromanheader{2}{เหตุยา ปญฺจ ฯเปฯ อวิคเต ปญฺจ. (สพฺพตฺถ วิตฺถาโร.)}
\proseitem{41}{เหตุํ อาสวสมฺปยุตฺตญฺเจว โน จ อาสวํ ธมฺมํ ปฏิจฺจ เหตุ อาสวสมฺปยุตฺโต เจว โน จ อาสโว ธมฺโม อุปฺปชฺชติ เหตุปจฺจยา.}

\csromanheader{2}{เหตุยา ปญฺจ ฯเปฯ อวิคเต ปญฺจ. (สพฺพตฺถ วิตฺถาโร.)}
\csromanheader{2}{1. เหตุทุก 18. อาสววิปฺปยุตฺตสาสวทุก}
\proseitem{42}{เหตุํ อาสววิปฺปยุตฺตํ สาสวํ ธมฺมํ ปฏิจฺจ เหตุ อาสววิปฺปยุตฺโต สาสโว ธมฺโม อุปฺปชฺชติ เหตุปจฺจยา. (สํขิตฺตํ.)}

\csromanheader{2}{เหตุยา นว ฯเปฯ อวิคเต นว. (สพฺพตฺถ วิตฺถาโร.)}
\prose{เหตุํ อาสววิปฺปยุตฺตํ อนาสวํ ธมฺมํ ปฏิจฺจ เหตุ อาสววิปฺปยุตฺโต อนาสโว ธมฺโม อุปฺปชฺชติ เหตุปจฺจยา. (สํขิตฺตํ.)}

\csromanheader{2}{เหตุยา นว ฯเปฯ อวิคเต นว. (สพฺพตฺถ วิตฺถาโร.)}
\nitthitam{เหตุทุก-อาสวโคจฺฉกํ นิฏฺฐิตํ.}
\csromanheader{2}{1. เหตุทุก 19-53 สญฺโญชนาทิทุก}
\proseitem{43}{เหตุํ สญฺโญชนํ ธมฺมํ ปฏิจฺจ ฯเปฯ.\csromanspacer{} เหตุํ คนฺถํ ธมฺมํ ปฏิจฺจ ฯเปฯ.\csromanspacer{} เหตุํ โอฆํ ธมฺมํ ปฏิจฺจ ฯเปฯ.\csromanspacer{} เหตุํ โยคํ ธมฺมํ ปฏิจฺจ ฯเปฯ.\csromanspacer{} เหตุํ นีวรณํ ธมฺมํ ปฏิจฺจ ฯเปฯ.\csromanspacer{} เหตุํ โนปรามาสํ ธมฺมํ ปฏิจฺจ เหตุ โน ปรามาโส ธมฺโม อุปฺปชฺชติ เหตุปจฺจยา. (สํขิตฺตํ.)}

\prose{เหตุยา นว ฯเปฯ อวิคเต นว. (สพฺพตฺถ โคจฺฉกํ วิตฺถาเรตพฺพํ.)}

\nitthitam{เหตุทุกปรามาสโคจฺฉกํ นิฏฺฐิตํ.}
\csromanheader{2}{1. เหตุทุก 54. สารมฺมณทุก}
\proseitem{44}{\csromanfolio{50}เหตุํ สารมฺมณํ ธมฺมํ ปฏิจฺจ เหตุ สารมฺมโณ ธมฺโม อุปฺปชฺชติ เหตุปจฺจยา.\csromanspacer{} ตีณิ.}

\prose{นเหตุํ สารมฺมณํ ธมฺมํ ปฏิจฺจ นเหตุ สารมฺมโณ ธมฺโม อุปฺปชฺชติ เหตุปจฺจยา.\csromanspacer{} ตีณิ.}

\prose{เหตุํ สารมฺมณญฺจ นเหตุํ สารมฺมณญฺจ ธมฺมํ ปฏิจฺจ เหตุ สารมฺมโณ ธมฺโม อุปฺปชฺชติ เหตุปจฺจยา.\csromanspacer{} ตีณิ. (สํขิตฺตํ.)}

\csromanheader{2}{เหตุยา นว ฯเปฯ อวิคเต นว. (สพฺพตฺถ วิตฺถาโร.)}
\prose{นเหตุํ อนารมฺมณํ ธมฺมํ ปฏิจฺจ นเหตุ อนารมฺมโณ ธมฺโม อุปฺปชฺชติ เหตุปจฺจยา. (สพฺพตฺถ เอกํ.)}

\csromanheader{2}{1. เหตุทุก 55. จิตฺตทุก}
\proseitem{45}{เหตุํ โนจิตฺตํ ธมฺมํ ปฏิจฺจ เหตุ โนจิตฺโต ธมฺโม อุปฺปชฺชติ เหตุปจฺจยา. (สํขิตฺตํ.)}

\csromanheader{2}{เหตุยา นว ฯเปฯ อวิคเต นว. (สพฺพตฺถ วิตฺถาโร.)}
\csromanheader{2}{1. เหตุทุก 56. เจตสิกทุก}
\proseitem{46}{เหตุํ เจตสิกํ ธมฺมํ ปฏิจฺจ เหตุ เจตสิโก ธมฺโม อุปฺปชฺชติ เหตุปจฺจยา. (สํขิตฺตํ.)}

\csromanheader{2}{เหตุยา นว ฯเปฯ อวิคเต นว. (สพฺพตฺถ วิตฺถาโร.)}
\prose{นเหตุํ อเจตสิกํ ธมฺมํ ปฏิจฺจ นเหตุ อเจตสิโก ธมฺโม อุปฺปชฺชติ เหตุปจฺจยา. (สพฺพตฺถ เอกํ.)}

\csromanheader{2}{1. เหตุทุก 57. จิตฺตสมฺปยุตฺตทุก}
\proseitem{47}{เหตุํ จิตฺตสมฺปยุตฺตํ ธมฺมํ ปฏิจฺจ เหตุ จิตฺตสมฺปยุตฺโต ธมฺโม อุปฺปชฺชติ เหตุปจฺจยา. (สํขิตฺตํ.)}

\csromanheader{2}{เหตุยา นว ฯเปฯ อวิคเต นว. (สพฺพตฺถ วิตฺถาโร.)}
\prose{นเหตุํ จิตฺตวิปฺปยุตฺตํ ธมฺมํ ปฏิจฺจ นเหตุ จิตฺตวิปฺปยุตฺโต ธมฺโม อุปฺปชฺชติ เหตุปจฺจยา. (สพฺพตฺถ เอกํ.)}

\vspace{\tipitakaparskip}
\csromancenter{เหตุทุก 62.\csromanspacer{} จิตฺตสํสฏฺฐสมุฏฺฐานทุก \textbf{1.}\csromanspacer{}\textbf{ เหตุทุก 58.}\csromanspacer{}\textbf{ จิตฺตสํสฏฺฐทุก}}
\proseitem{48}{\csromanfolio{51}เหตุํ จิตฺตสํสฏฺฐํ ธมฺมํ ปฏิจฺจ เหตุ จิตฺตสํสฏฺโฐ ธมฺโม อุปฺปชฺชติ เหตุปจฺจยา. (สํขิตฺตํ.)}

\csromanheader{2}{เหตุยา นว ฯเปฯ อวิคเต นว. (สพฺพตฺถ วิตฺถาโร.)}
\prose{นเหตุํ จิตฺตวิสํสฏฺฐํ ธมฺมํ ปฏิจฺจ นเหตุ จิตฺตวิสํสฏฺโฐ ธมฺโม อุปฺปชฺชติ เหตุปจฺจยา. (สพฺพตฺถ เอกํ.)}

\csromanheader{2}{1. เหตุทุก 59. จิตฺตสมุฏฺฐานทุก}
\proseitem{49}{เหตุํ จิตฺตสมุฏฺฐานํ ธมฺมํ ปฏิจฺจ เหตุ จิตฺตสมุฏฺฐาโน ธมฺโม อุปฺปชฺชติ เหตุปจฺจยา. (สํขิตฺตํ.)}

\csromanheader{2}{เหตุยา นว ฯเปฯ อวิคเต นว. (สพฺพตฺถ วิตฺถาโร.)}
\prose{นเหตุํ โนจิตฺตสมุฏฺฐานํ ธมฺมํ ปฏิจฺจ นเหตุ โนจิตฺตสมุฏฺฐาโน ธมฺโม อุปฺปชฺชติ เหตุปจฺจยา. (สพฺพตฺถ เอกํ.)}

\csromanheader{2}{1. เหตุทุก 60. จิตฺตสหภูทุก}
\proseitem{50}{เหตุํ จิตฺตสหภุํ ธมฺมํ ปฏิจฺจ เหตุ จิตฺตสหภู ธมฺโม อุปฺปชฺชติ เหตุปจฺจยา. (สพฺพตฺถ นว.)}

\prose{นเหตุํ โนจิตฺตสหภุํ ธมฺมํ ปฏิจฺจ นเหตุ โนจิตฺตสหภู ธมฺโม อุปฺปชฺชติ เหตุปจฺจยา. (สพฺพตฺถ เอกํ.\csromanspacer{} สพฺพตฺถ วิตฺถาโร.)}

\csromanheader{2}{1. เหตุทุก 61. จิตฺตานุปริวตฺติทุก}
\proseitem{51}{เหตุํ จิตฺตานุปริวตฺติํ ธมฺมํ ปฏิจฺจ เหตุ จิตฺตานุปริวตฺตี ธมฺโม อุปฺปชฺชติ เหตุปจฺจยา. (สพฺพตฺถ นว.)}

\prose{นเหตุํ โนจิตฺตานุปริวตฺติํ ธมฺมํ ปฏิจฺจ นเหตุ โนจิตฺตานุปริวตฺตี ธมฺโม อุปฺปชฺชติ เหตุปจฺจยา. (สพฺพตฺถ เอกํ.\csromanspacer{} สพฺพตฺถ วิตฺถาโร.)}

\csromanheader{2}{1. เหตุทุก 62. จิตฺตสํสฏฺฐสมุฏฺฐานทุก}
\proseitem{52}{เหตุํ จิตฺตสํสฏฺฐสมุฏฺฐานํ ธมฺมํ ปฏิจฺจ เหตุ จิตฺตสํสฏฺฐสมุฏฺฐาโน ธมฺโม อุปฺปชฺชติ เหตุปจฺจยา. (สพฺพตฺถ นว.)}

\prose{นเหตุํ โนจิตฺตสํสฏฺฐสมุฏฺฐานํ ธมฺมํ ปฏิจฺจ นเหตุ โนจิตฺตสํสฏฺฐสมุฏฺฐาโน ธมฺโม อุปฺปชฺชติ เหตุปจฺจยา. (สพฺพตฺถ เอกํ.\csromanspacer{} สพฺพตฺถ วิตฺถาโร.)}

\csromanheader{2}{1. เหตุทุก 63. จิตฺตสํสฏฺฐสมุฏฺฐานสหภูทุก}
\proseitem{53}{\csromanfolio{52}เหตุํ จิตฺตสํสฏฺฐสมุฏฺฐานสหภุํ ธมฺมํ ปฏิจฺจ เหตุ จิตฺตสํสฏฺฐสมุฏฺฐานสหภู ธมฺโม อุปฺปชฺชติ เหตุปจฺจยา. (สพฺพตฺถ นว.)}

\prose{นเหตุํ โนจิตฺตสํสฏฺฐสมุฏฺฐานสหภุํ ธมฺมํ ปฏิจฺจ นเหตุ โนจิตฺตสํสฏฺฐสมุฏฺฐานสหภู ธมฺโม อุปฺปชฺชติ เหตุปจฺจยา. (สพฺพตฺถ เอกํ.\csromanspacer{} สพฺพตฺถ วิตฺถาโร.)}

\csromanheader{2}{1. เหตุทุก 64. จิตฺตสํสฏฺฐสมุฏฺฐานานุปริวตฺติทุก}
\proseitem{54}{เหตุํ จิตฺตสํสฏฺฐสมุฏฺฐานานุปริวตฺติํ ธมฺมํ ปฏิจฺจ เหตุ จิตฺตสํสฏฺฐสมุฏฺฐานานุปริวตฺตี ธมฺโม อุปฺปชฺชติ เหตุปจฺจยา. (สพฺพตฺถ นว.)}

\prose{นเหตุํ โนจิตฺตสํสฏฺฐสมุฏฺฐานานุปริวตฺติํ ธมฺมํ ปฏิจฺจ นเหตุ โนจิตฺตสํสฏฺฐสมุฏฺฐานานุปริวตฺตี ธมฺโม อุปฺปชฺชติ เหตุปจฺจยา. (สพฺพตฺถ เอกํ.\csromanspacer{} สพฺพตฺถ วิตฺถาโร.)}

\csromanheader{2}{1. เหตุทุก 65. อชฺฌตฺติกทุก}
\proseitem{55}{นเหตุํ อชฺฌตฺติกํ ธมฺมํ ปฏิจฺจ นเหตุ อชฺฌตฺติโก ธมฺโม อุปฺปชฺชติ เหตุปจฺจยา. (สพฺพตฺถ เอกํ.\csromanspacer{} สพฺพตฺถ วิตฺถาโร.)}

\prose{เหตุํ พาหิรํ ธมฺมํ ปฏิจฺจ เหตุ พาหิโร ธมฺโม อุปฺปชฺชติ เหตุปจฺจยา. (สพฺพตฺถ นว.\csromanspacer{} สพฺพตฺถ วิตฺถาโร.)}

\csromanheader{2}{1. เหตุทุก 66. อุปาทาทุก}
\proseitem{56}{เหตุํ โน-อุปาทา ธมฺมํ ปฏิจฺจ เหตุ โน-อุปาทา ธมฺโม อุปฺปชฺชติ เหตุปจฺจยา. (สํขิตฺตํ.)}

\prose{เหตุยา นว ฯเปฯ อวิคเต นว. (สํขิตฺตํ.\csromanspacer{} สพฺพตฺถ วิตฺถาโร.)}

\csromanheader{2}{1. เหตุทุก 67. อุปาทินฺนทุก}
\proseitem{57}{เหตุํ อุปาทินฺนํ ธมฺมํ ปฏิจฺจ เหตุ อุปาทินฺโน ธมฺโม อุปฺปชฺชติ เหตุปจฺจยา. (สํขิตฺตํ.)}

\csromanheader{2}{เหตุยา นว ฯเปฯ อวิคเต นว. (สพฺพตฺถ วิตฺถาโร.)}
\csromanheader{2}{1. เหตุทุก 82. ปิฏฺฐิทุก}
\prose{\csromanfolio{53}เหตุํ อนุปาทินฺนํ ธมฺมํ ปฏิจฺจ เหตุ อนุปาทินฺโน ธมฺโม อุปฺปชฺชติ เหตุปจฺจยา. (สํขิตฺตํ.)}

\csromanheader{2}{เหตุยา นว ฯเปฯ อวิคเต นว. (สพฺพตฺถ วิตฺถาโร.)}
\nitthitam{เหตุทุกมหนฺตรทุกํ นิฏฺฐิตํ.}
\csromanheader{2}{1. เหตุทุก 68. อุปาทานโคจฺฉก}
\proseitem{58}{เหตุํ อุปาทานํ ธมฺมํ ปฏิจฺจ นเหตุ อุปาทาโน ธมฺโม อุปฺปชฺชติ เหตุปจฺจยา.}

\csromanheader{2}{เหตุยา เทฺว ฯเปฯ อวิคเต เทฺว. (สํขิตฺตํ.)}
\csromanheader{2}{1. เหตุทุก 74. กิเลสโคจฺฉก}
\proseitem{59}{เหตุํ กิเลสํ ธมฺมํ ปฏิจฺจ เหตุ กิเลโส ธมฺโม อุปฺปชฺชติ เหตุปจฺจยา.}

\csromanheader{2}{เหตุยา นว ฯเปฯ อวิคเต นว. (สํขิตฺตํ.)}
\csromanheader{2}{1. เหตุทุก 82. ปิฏฺฐิทุก}
\proseitem{60}{เหตุํ ทสฺสเนน ปหาตพฺพํ ธมฺมํ ปฏิจฺจ เหตุ ทสฺสเนน ปหาตพฺโพ ธมฺโม อุปฺปชฺชติ เหตุปจฺจยา. (สพฺพตฺถ นว.\csromanspacer{} วิปากํ นตฺถิ.)}

\prose{เหตุํ นทสฺสเนน ปหาตพฺพํ ธมฺมํ ปฏิจฺจ เหตุ นทสฺสเนน ปหาตพฺโพ ธมฺโม อุปฺปชฺชติ เหตุ ปจฺจยา. (สพฺพตฺถ นว.\csromanspacer{} วิปากํ นตฺถิ.)}

\proseitem{61}{เหตุํ ภาวนาย ปหาตพฺพํ ธมฺมํ ปฏิจฺจ เหตุ ภาวนาย ปหาตพฺโพ ธมฺโม อุปฺปชฺชติ เหตุปจฺจยา. (สพฺพตฺถ นว.\csromanspacer{} วิปากํ นตฺถิ.)}

\prose{เหตุํ นภาวนาย ปหาตพฺพํ ธมฺมํ ปฏิจฺจ เหตุ นภาวนาย ปหาตพฺโพ ธมฺโม อุปฺปชฺชติ เหตุปจฺจยา. (สพฺพตฺถ นว.\csromanspacer{} วิปากํ นตฺถิ.)}

\proseitem{62}{เหตุํ ทสฺสเนน ปหาตพฺพเหตุกํ ธมฺมํ ปฏิจฺจ เหตุ ทสฺสเนน ปหาตพฺพเหตุโก ธมฺโม อุปฺปชฺชติ เหตุปจฺจยา. (สพฺพตฺถ นว.)}

\prose{เหตุํ นทสฺสเนน ปหาตพฺพเหตุกํ ธมฺมํ ปฏิจฺจ เหตุ นทสฺสเนน ปหาตพฺพเหตุโก ธมฺโม อุปฺปชฺชติ เหตุปจฺจยา. (สพฺพตฺถ นว.)}

\proseitem{63}{\csromanfolio{54}เหตุํ ภาวนาย ปหาตพฺพเหตุกํ ธมฺมํ ปฏิจฺจ เหตุ ภาวนาย ปหาตพฺพเหตุโก ธมฺโม อุปฺปชฺชติ เหตุปจฺจยา. (สพฺพตฺถ นว.)}

\prose{เหตุํ นภาวนาย ปหาตพฺพเหตุกํ ธมฺมํ ปฏิจฺจ เหตุ นภาวนาย ปหาตพฺพเหตุโก ธมฺโม อุปฺปชฺชติ เหตุปจฺจยา. (สพฺพตฺถ นว.)}

\proseitem{64}{เหตุํ สวิตกฺกํ ธมฺมํ ปฏิจฺจ เหตุ สวิตกฺโก ธมฺโม อุปฺปชฺชติ เหตุปจฺจยา. (สพฺพตฺถ นว.)}

\prose{เหตุํ อวิตกฺกํ ธมฺมํ ปฏิจฺจ เหตุ อวิตกฺโก ธมฺโม อุปฺปชฺชติ เหตุปจฺจยา. (สพฺพตฺถ นว.)}

\proseitem{65}{เหตุํ สวิจารํ ธมฺมํ ปฏิจฺจ เหตุ สวิจาโร ธมฺโม อุปฺปชฺชติ เหตุปจฺจยา. (สพฺพตฺถ นว.)}

\prose{เหตุํ อวิจารํ ธมฺมํ ปฏิจฺจ เหตุ อวิจาโร ธมฺโม อุปฺปชฺชติ เหตุปจฺจยา. (สพฺพตฺถ นว.)}

\proseitem{66}{เหตุํ สปฺปีติกํ ธมฺมํ ปฏิจฺจ ฯเปฯ.\csromanspacer{} เหตุํ อปฺปีติกํ ธมฺมํ ปฏิจฺจ ฯเปฯ.}

\prose{เหตุํ ปีติสหคตํ ธมฺมํ ปฏิจฺจ ฯเปฯ.\csromanspacer{} เหตุํ นปีติสหคตํ ธมฺมํ ปฏิจฺจ ฯเปฯ.}

\prose{เหตุํ สุขสหคตํ ธมฺมํ ปฏิจฺจ ฯเปฯ.\csromanspacer{} เหตุํ นสุขสหคตํ ธมฺมํ ปฏิจฺจ ฯเปฯ.}

\prose{เหตุํ อุเปกฺขาสหคตํ ธมฺมํ ปฏิจฺจ ฯเปฯ.\csromanspacer{} เหตุํ น-อุเปกฺขาสหคตํ ธมฺมํ ปฏิจฺจ ฯเปฯ.}

\prose{เหตุํ กามาวจรํ ธมฺมํ ปฏิจฺจ ฯเปฯ.\csromanspacer{} เหตุํ นกามาวจรํ ธมฺมํ ปฏิจฺจ ฯเปฯ.}

\prose{เหตุํ รูปาวจรํ ธมฺมํ ปฏิจฺจ ฯเปฯ.\csromanspacer{} เหตุํ นรูปาวจรํ ธมฺมํ ปฏิจฺจ ฯเปฯ.}

\prose{เหตุํ อรูปาวจรํ ธมฺมํ ปฏิจฺจ ฯเปฯ.\csromanspacer{} เหตุํ น-อรูปาวจรํ ธมฺมํ ปฏิจฺจ ฯเปฯ.}

\prose{เหตุํ ปริยาปนฺนํ ธมฺมํ ปฏิจฺจ ฯเปฯ.\csromanspacer{} เหตุํ อปริยาปนฺนํ ธมฺมํ ปฏิจฺจ ฯเปฯ.}

\prose{เหตุํ นิยฺยานิกํ ธมฺมํ ปฏิจฺจ ฯเปฯ.\csromanspacer{} เหตุํ อนิยฺยานิกํ ธมฺมํ ปฏิจฺจ ฯเปฯ.}

\prose{เหตุํ นิยตํ ธมฺมํ ปฏิจฺจ ฯเปฯ.\csromanspacer{} เหตุํ อนิยตํ ธมฺมํ ปฏิจฺจ ฯเปฯ.}

\prose{เหตุํ ส-อุตฺตรํ ธมฺมํ ปฏิจฺจ ฯเปฯ.\csromanspacer{} เหตุํ อนุตฺตรํ ธมฺมํ ปฏิจฺจ ฯเปฯ.}

\csromanheader{2}{5. เหตุเหตุสมฺปยุตฺตทุก 1. เหตุทุก}
\prose{\csromanfolio{55}เหตุํ สรณํ ธมฺมํ ปฏิจฺจ ฯเปฯ.\csromanspacer{} เหตุํ อรณํ ธมฺมํ ปฏิจฺจ เหตุ อรโณ ธมฺโม อุปฺปชฺชติ เหตุปจฺจยา. (สํขิตฺตํ.)}

\prose{เหตุยา นว ฯเปฯ อวิคเต นว. (สํขิตฺตํ.\csromanspacer{} สพฺพตฺถ วิตฺถาโร.)}

\nitthitam{เหตุทุกปิฏฺฐิทุกํ นิฏฺฐิตํ.}
\csromanheader{2}{2. สเหตุกทุก 1. เหตุทุก}
\proseitem{67}{สเหตุกํ เหตุํ ธมฺมํ ปฏิจฺจ สเหตุโก เหตุ ธมฺโม อุปฺปชฺชติ เหตุปจฺจยา. (สพฺพตฺถ เอกํ.)}

\prose{สเหตุกํ นเหตุํ ธมฺมํ ปฏิจฺจ สเหตุโก นเหตุ ธมฺโม อุปฺปชฺชติ เหตุปจฺจยา.}

\prose{เหตุยา นว.}

\csromanheader{2}{3. เหตุสมฺปยุตฺตทุก 1. เหตุทุก}
\proseitem{68}{เหตุสมฺปยุตฺตํ เหตุํ ธมฺมํ ปฏิจฺจ เหตุสมฺปยุตฺโต เหตุ ธมฺโม อุปฺปชฺชติ เหตุปจฺจยา. (สพฺพตฺถ เอกํ.)}

\prose{เหตุสมฺปยุตฺตํ นเหตุํ ธมฺมํ ปฏิจฺจ เหตุสมฺปยุตฺโต นเหตุ ธมฺโม อุปฺปชฺชติ เหตุปจฺจยา. (สํขิตฺตํ.)}

\prose{เหตุยา นว.}

\csromanheader{2}{4. เหตุสเหตุกทุก 1. เหตุทุก}
\proseitem{69}{เหตุญฺเจว สเหตุกญฺจ เหตุํ ธมฺมํ ปฏิจฺจ เหตุ เจว สเหตุโก จ เหตุ ธมฺโม อุปฺปชฺชติ เหตุปจฺจยา. (สพฺพตฺถ เอกํ.)}

\prose{สเหตุกญฺเจว น จ เหตุํ นเหตุํ ธมฺมํ ปฏิจฺจ สเหตุโก เจว น จ เหตุ นเหตุ ธมฺโม อุปฺปชฺชติ เหตุปจฺจยา. (สพฺพตฺถ เอกํ.)}

\csromanheader{2}{5. เหตุเหตุสมฺปยุตฺตทุก 1. เหตุทุก}
\proseitem{70}{เหตุญฺเจว เหตุสมฺปยุตฺตญฺจ เหตุํ ธมฺมํ ปฏิจฺจ เหตุ เจว เหตุสมฺปยุตฺโต จ เหตุ ธมฺโม อุปฺปชฺชติ เหตุปจฺจยา. (สพฺพตฺถ เอกํ.)}

\prose{\csromanfolio{56}เหตุสมฺปยุตฺตญฺเจว น จ เหตุํ นเหตุํ ธมฺมํ ปฏิจฺจ เหตุสมฺปยุตฺโต เจว น จ เหตุ นเหตุ ธมฺโม อุปฺปชฺชติ เหตุปจฺจยา. (สพฺพตฺถ เอกํ.)}

\csromanheader{2}{6. นเหตุสเหตุกทุก 1. เหตุทุก}
\proseitem{71}{นเหตุสเหตุกํ นเหตุํ ธมฺมํ ปฏิจฺจ นเหตุสเหตุโก นเหตุ ธมฺโม อุปฺปชฺชติ เหตุปจฺจยา.}

\prose{นเหตุํ อเหตุกํ นเหตุํ ธมฺมํ ปฏิจฺจ นเหตุ อเหตุโก นเหตุ ธมฺโม อุปฺปชฺชติ เหตุปจฺจยา.}

\csromanheader{2}{7. จูฬนฺตรทุก 1. เหตุทุก}
\proseitem{72}{สปฺปจฺจยํ เหตุํ ธมฺมํ ปฏิจฺจ สปฺปจฺจโย เหตุ ธมฺโม อุปฺปชฺชติ เหตุปจฺจยา. (สพฺพตฺถ เอกํ.)}

\prose{สปฺปจฺจยํ นเหตุํ ธมฺมํ ปฏิจฺจ สปฺปจฺจโย นเหตุ ธมฺโม อุปฺปชฺชติ เหตุปจฺจยา. (สพฺพตฺถ เอกํ.)}

\proseitem{73}{สงฺขตํ เหตุํ ธมฺมํ ปฏิจฺจ สงฺขโต เหตุ ธมฺโม อุปฺปชฺชติ เหตุปจฺจยา. (สพฺพตฺถ เอกํ.)}

\prose{สงฺขตํ นเหตุํ ธมฺมํ ปฏิจฺจ สงฺขโต นเหตุ ธมฺโม อุปฺปชฺชติ เหตุปจฺจยา. (สพฺพตฺถ เอกํ.)}

\proseitem{74}{อนิทสฺสนํ เหตุํ ธมฺมํ ปฏิจฺจ อนิทสฺสโน เหตุ ธมฺโม อุปฺปชฺชติ เหตุปจฺจยา. (สพฺพตฺถ เอกํ.)}

\prose{อนิทสฺสนํ นเหตุํ ธมฺมํ ปฏิจฺจ อนิทสฺสโน นเหตุ ธมฺโม อุปฺปชฺชติ เหตุปจฺจยา. (สํขิตฺตํ.)}

\prose{เหตุยา ตีณิ, อารมฺมเณ เอกํ ฯเปฯ อวิคเต ตีณิ. (สํขิตฺตํ.)}

\proseitem{75}{อปฺปฏิฆํ เหตุํ ธมฺมํ ปฏิจฺจ สปฺปฏิโฆ เหตุ ธมฺโม อุปฺปชฺชติ เหตุปจฺจยา. (สพฺพตฺถ เอกํ.)}

\csromanheader{2}{7. จูฬนฺตรทุก 1. เหตุทุก}
\prose{\csromanfolio{57}สปฺปฏิฆํ นเหตุํ ธมฺมํ ปฏิจฺจ สปฺปฏิโฆ นเหตุ ธมฺโม อุปฺปชฺชติ เหตุปจฺจยา.\csromanspacer{} ตีณิ.}

\prose{อปฺปฏิฆํ นเหตุํ ธมฺมํ ปฏิจฺจ อปฺปฏิโฆ นเหตุ ธมฺโม อุปฺปชฺชติ เหตุปจฺจยา.\csromanspacer{} ตีณิ.}

\prose{สปฺปฏิฆํ นเหตุญฺจ อปฺปฏิฆํ นเหตุญฺจ ธมฺมํ ปฏิจฺจ สปฺปฏิโฆ นเหตุ ธมฺโม อุปฺปชฺชติ เหตุปจฺจยา.\csromanspacer{} ตีณิ. (สํขิตฺตํ.)}

\prose{เหตุยา นว, อารมฺมเณ เอกํ ฯเปฯ อญฺญมญฺเญ ฉ ฯเปฯ อวิคเต นว.}

\proseitem{76}{อรูปิํ เหตุํ ธมฺมํ ปฏิจฺจ อรูปี เหตุ ธมฺโม อุปฺปชฺชติ เหตุปจฺจยา. (สพฺพตฺถ เอกํ.)}

\prose{รูปิํ นเหตุํ ธมฺมํ ปฏิจฺจ รูปี นเหตุ ธมฺโม อุปฺปชฺชติ เหตุปจฺจยา. (สํขิตฺตํ.)}

\prose{เหตุยา นว.}

\csromanheader{2}{77. โลกิยํ เหตุํ ธมฺมํ ปฏิจฺจ โลกิโย เหตุ ธมฺโม อุปฺปชฺชติ}
\prose{โลกุตฺตรํ เหตุํ ธมฺมํ ปฏิจฺจ โลกุตฺตโร เหตุ ธมฺโม อุปฺปชฺชติ เหตุปจฺจยา. (1) (สพฺพตฺถ เทฺว.)}

\prose{โลกิยํ นเหตุํ ธมฺมํ ปฏิจฺจ โลกิโย นเหตุ ธมฺโม อุปฺปชฺชติ}

\prose{โลกุตฺตรํ นเหตุํ ธมฺมํ ปฏิจฺจ โลกุตฺตโร นเหตุ ธมฺโม อุปฺปชฺชติ เหตุปจฺจยา.\csromanspacer{} ตีณิ.}

\prose{โลกิยํ นเหตุญฺจ โลกุตฺตรํ นเหตุญฺจ ธมฺมํ ปฏิจฺจ โลกิโย นเหตุ ธมฺโม อุปฺปชฺชติ เหตุปจฺจยา. (1) (สํขิตฺตํ.)}

\prose{เหตุยา ปญฺจ.}

\proseitem{78}{เกนจิ วิญฺเญยฺยํ เหตุํ ธมฺมํ ปฏิจฺจ ฯเปฯ.\csromanspacer{} นเกนจิ วิญฺเญยฺยํ เหตุํ ธมฺมํ ปฏิจฺจ. (สพฺพตฺถ นว.)}

\prose{เกนจิ วิญฺเญยฺยํ นเหตุํ ธมฺมํ ปฏิจฺจ ฯเปฯ.\csromanspacer{} นเกนจิ วิญฺเญยฺยํ นเหตุํ ธมฺมํ ปฏิจฺจ. (สพฺพตฺถ นว.)}

\csromanheader{2}{14. อาสวโคจฺฉก 1. เหตุทุก}
\proseitem{79}{\csromanfolio{58}อาสวํ เหตุํ ธมฺมํ ปฏิจฺจ ฯเปฯ.\csromanspacer{} โน-อาสวํ เหตุํ ธมฺมํ ปฏิจฺจ. (สํขิตฺตํ.)}

\prose{อาสวํ นเหตุํ ธมฺมํ ปฏิจฺจ ฯเปฯ.\csromanspacer{} โน-อาสวํ นเหตุํ ธมฺมํ ปฏิจฺจ. (สํขิตฺตํ.)}

\proseitem{80}{สาสวํ เหตุํ ธมฺมํ ปฏิจฺจ ฯเปฯ.\csromanspacer{} อนาสวํ เหตุํ ธมฺมํ ปฏิจฺจ. (สพฺพตฺถ เทฺว.)}

\prose{สาสวํ นเหตุํ ธมฺมํ ปฏิจฺจ ฯเปฯ.\csromanspacer{} อนาสวํ นเหตุํ ธมฺมํ ปฏิจฺจ ฯเปฯ.}

\proseitem{81}{อาสวสมฺปยุตฺตํ เหตุํ ธมฺมํ ปฏิจฺจ ฯเปฯ.\csromanspacer{} อาสววิปฺปยุตฺตํ เหตุํ ธมฺมํ ปฏิจฺจ. (สํขิตฺตํ.)}

\prose{อาสวสมฺปยุตฺตํ นเหตุํ ธมฺมํ ปฏิจฺจ ฯเปฯ.\csromanspacer{} อาสววิปฺปยุตฺตํ นเหตุํ ธมฺมํ ปฏิจฺจ. (สํขิตฺตํ.)}

\proseitem{82}{อาสวญฺเจว สาสวญฺจ เหตุํ ธมฺมํ ปฏิจฺจ ฯเปฯ.\csromanspacer{} สาสวญฺเจว โน จ อาสวํ เหตุํ ธมฺมํ ปฏิจฺจ.}

\prose{อาสวญฺเจว สาสวญฺจ นเหตุํ ธมฺมํ ปฏิจฺจ ฯเปฯ.\csromanspacer{} สาสวญฺเจว โน จ อาสวํ นเหตุํ ธมฺมํ ปฏิจฺจ. (สํขิตฺตํ.)}

\proseitem{83}{อาสวญฺเจว อาสวสมฺปยุตฺตญฺจ เหตุํ ธมฺมํ ปฏิจฺจ ฯเปฯ.\csromanspacer{} อาสวสมฺปยุตฺตญฺเจว โน จ อาสวํ เหตุํ ธมฺมํ ปฏิจฺจ.}

\prose{อาสวญฺเจว อาสวสมฺปยุตฺตญฺจ นเหตุํ ธมฺมํ ปฏิจฺจ ฯเปฯ.\csromanspacer{} อาสวสมฺปยุตฺตญฺเจว โน จ อาสวํ นเหตุํ ธมฺมํ ปฏิจฺจ.}

\proseitem{84}{อาสววิปฺปยุตฺตํ สาสวํ เหตุํ ธมฺมํ ปฏิจฺจ ฯเปฯ.\csromanspacer{} อาสววิปฺปยุตฺตํ อนาสวํ เหตุํ ธมฺมํ ปฏิจฺจ.}

\prose{อาสววิปฺปยุตฺตํ สาสวํ นเหตุํ ธมฺมํ ปฏิจฺจ ฯเปฯ.\csromanspacer{} อาสววิปฺปยุตฺตํ อนาสวํ นเหตุํ ธมฺมํ ปฏิจฺจ.}

\csromanheader{2}{55. มหนฺตรทุก 1. เหตุทุก \textbf{20. สญฺโญชนาทิทุก 1. เหตุทุก}}
\proseitem{85}{\csromanfolio{59}สญฺโญชนํ เหตุํ ธมฺมํ ปฏิจฺจ ฯเปฯ.\csromanspacer{} คนฺถํ เหตุํ ธมฺมํ ปฏิจฺจ ฯเปฯ.}

\prose{โอฆํ เหตุํ ธมฺมํ ปฏิจฺจ ฯเปฯ.\csromanspacer{} โยคํ เหตุํ ธมฺมํ ปฏิจฺจ ฯเปฯ.}

\csromanheader{2}{นีวรณํ เหตุํ ธมฺมํ ปฏิจฺจ. (สํขิตฺตํ.)}
\prose{โนปรามาสํ เหตุํ ธมฺมํ ปฏิจฺจ. (สพฺพตฺถ เอกํ.\csromanspacer{} สํขิตฺตํ.)}

\csromanheader{2}{55. มหนฺตรทุก 1. เหตุทุก}
\proseitem{86}{สารมฺมณํ เหตุํ ธมฺมํ ปฏิจฺจ. (สพฺพตฺถ เอกํ.) สารมฺมณํ นเหตุํ ธมฺมํ ปฏิจฺจ. (สํขิตฺตํ.)}

\proseitem{87}{โนจิตฺตํ เหตุํ ธมฺมํ ปฏิจฺจ.}

\prose{เจตสิกํ เหตุํ ธมฺมํ ปฏิจฺจ.}

\prose{จิตฺตสมฺปยุตฺตํ เหตุํ ธมฺมํ ปฏิจฺจ.}

\prose{จิตฺตสํสฏฺฐํ เหตุํ ธมฺมํ ปฏิจฺจ.}

\prose{จิตฺตสมุฏฺฐานํ เหตุํ ธมฺมํ ปฏิจฺจ.}

\prose{จิตฺตสหภุํ เหตุํ ธมฺมํ ปฏิจฺจ.}

\prose{จิตฺตานุปริวตฺติํ เหตุํ ธมฺมํ ปฏิจฺจ.}

\prose{จิตฺตสํสฏฺฐสมุฏฺฐานํ เหตุํ ธมฺมํ ปฏิจฺจ.}

\prose{จิตฺตสํสฏฺฐสมุฏฺฐานสหภุํ เหตุํ ธมฺมํ ปฏิจฺจ.}

\prose{จิตฺตสํสฏฺฐสมุฏฺฐานานุปริวตฺติํ เหตุํ ธมฺมํ ปฏิจฺจ.}

\prose{พาหิรํ เหตุํ ธมฺมํ ปฏิจฺจ.}

\prose{โน-อุปาทา เหตุํ ธมฺมํ ปฏิจฺจ.}

\prose{อุปาทินฺนํ เหตุํ ธมฺมํ ปฏิจฺจ ฯเปฯ.\csromanspacer{} อนุปาทินฺนํ เหตุํ ธมฺมํ ปฏิจฺจ.}

\csromanheader{2}{69. อุปาทานโคจฺฉก 1. เหตุทุก}
\vspace{\tipitakaparskip}
\csromancenter{อุปาทานํ เหตุํ ธมฺมํ ปฏิจฺจ ฯเปฯ.}
\csromanheader{2}{75. กิเลสโคจฺฉก 1. เหตุทุก}
\vspace{\tipitakaparskip}
\csromancenter{กิเลสํ เหตุํ ธมฺมํ ปฏิจฺจ ฯเปฯ.}
\csromanheader{2}{83. ปิฏฺฐิทุก 1. เหตุทุก}
\proseitem{90}{\csromanfolio{60}ทสฺสเนน ปหาตพฺพํ เหตุํ ธมฺมํ ปฏิจฺจ ฯเปฯ.\csromanspacer{} นทสฺสเนน ปหาตพฺพํ เหตุํ ธมฺมํ ปฏิจฺจ.}

\proseitem{91}{ภาวนาย ปหาตพฺพํ เหตุํ ธมฺมํ ปฏิจฺจ ฯเปฯ.\csromanspacer{} นภาวนาย ปหาตพฺพํ เหตุํ ธมฺมํ ปฏิจฺจ.}

\proseitem{92}{ทสฺสเนน ปหาตพฺพเหตุกํ เหตุํ ธมฺมํ ปฏิจฺจ ฯเปฯ.\csromanspacer{} นทสฺสเนน ปหาตพฺพเหตุกํ เหตุํ ธมฺมํ ปฏิจฺจ.}

\proseitem{93}{ภาวนาย ปหาตพฺพเหตุกํ เหตุํ ธมฺมํ ปฏิจฺจ ฯเปฯ.\csromanspacer{} นภาวนาย ปหาตพฺพเหตุกํ เหตุํ ธมฺมํ ปฏิจฺจ.}

\proseitem{94}{สวิตกฺกํ เหตุํ ธมฺมํ ปฏิจฺจ ฯเปฯ.\csromanspacer{} อวิตกฺกํ เหตุํ ธมฺมํ ปฏิจฺจ ฯเปฯ.}

\prose{สวิจารํ เหตุํ ธมฺมํ ปฏิจฺจ ฯเปฯ.\csromanspacer{} อวิจารํ เหตุํ ธมฺมํ ปฏิจฺจ ฯเปฯ.}

\prose{สปฺปีติกํ เหตุํ ธมฺมํ ปฏิจฺจ ฯเปฯ.\csromanspacer{} อปฺปีติกํ เหตุํ ธมฺมํ ปฏิจฺจ. (สํขิตฺตํ.)}

\prose{ปีติสหคตํ เหตุํ ธมฺมํ ปฏิจฺจ ฯเปฯ.\csromanspacer{} นปีติสหคตํ เหตุํ ธมฺมํ ปฏิจฺจ ฯเปฯ.}

\prose{สุขสหคตํ เหตุํ ธมฺมํ ปฏิจฺจ ฯเปฯ.\csromanspacer{} นสุขสหคตํ เหตุํ ธมฺมํ ปฏิจฺจ ฯเปฯ.}

\prose{อุเปกฺขาสหคตํ เหตุํ ธมฺมํ ปฏิจฺจ ฯเปฯ.\csromanspacer{} น-อุเปกฺขาสหคตํ เหตุํ ธมฺมํ ปฏิจฺจ ฯเปฯ.}

\prose{กามาวจรํ เหตุํ ธมฺมํ ปฏิจฺจ ฯเปฯ.\csromanspacer{} นกามาวจรํ เหตุํ ธมฺมํ ปฏิจฺจ ฯเปฯ.}

\csromanheader{2}{83. ปิฏฺฐิทุก 1. เหตุทุก}
\prose{\csromanfolio{61}รูปาวจรํ เหตุํ ธมฺมํ ปฏิจฺจ ฯเปฯ.\csromanspacer{} นรูปาวจรํ เหตุํ ธมฺมํ ปฏิจฺจ ฯเปฯ.}

\prose{อรูปาวจรํ เหตุํ ธมฺมํ ปฏิจฺจ ฯเปฯ.\csromanspacer{} น-อรูปาวจรํ เหตุํ ธมฺมํ ปฏิจฺจ ฯเปฯ.}

\prose{ปริยาปนฺนํ เหตุํ ธมฺมํ ปฏิจฺจ ฯเปฯ.\csromanspacer{} อปริยาปนฺนํ เหตุํ ธมฺมํ ปฏิจฺจ ฯเปฯ.}

\prose{นิยฺยานิกํ เหตุํ ธมฺมํ ปฏิจฺจ ฯเปฯ.\csromanspacer{} อนิยฺยานิกํ เหตุํ ธมฺมํ ปฏิจฺจ ฯเปฯ.}

\prose{นิยตํ เหตุํ ธมฺมํ ปฏิจฺจ ฯเปฯ.\csromanspacer{} อนิยตํ เหตุํ ธมฺมํ ปฏิจฺจ ฯเปฯ.}

\prose{ส-อุตฺตรํ เหตุํ ธมฺมํ ปฏิจฺจ ฯเปฯ.\csromanspacer{} อนุตฺตรํ เหตุํ ธมฺมํ ปฏิจฺจ ฯเปฯ. (สพฺพตฺถ เทฺว.)}

\prose{สรณํ เหตุํ ธมฺมํ ปฏิจฺจ ฯเปฯ.\csromanspacer{} อรณํ เหตุํ ธมฺมํ ปฏิจฺจ อรโณ เหตุ ธมฺโม อุปฺปชฺชติ เหตุปจฺจยา. (สํขิตฺตํ.)}

\prose{อรณํ นเหตุํ ธมฺมํ ปฏิจฺจ อรโณ นเหตุ ธมฺโม อุปฺปชฺชติ เหตุปจฺจยา.}

\prose{(สพฺพตฺถ วิตฺถาโร.)}

\nitthitam{ธมฺมานุโลเม ทุกทุกปฏฺฐานํ นิฏฺฐิตํ.}
\nitthitam{\textbf{อนุโลมปฏฺฐานํ นิฏฺฐิตํ.}}
\pitaka{\textbf{อภิธมฺมปิฏก}}
\csromanheader{2}{\textbf{ธมฺมปจฺจนีย}}
\csromanmatikaanchor{mtk.2}
\csromanmatikaanchor{mtk.6}
\csromantocmarknopage{chapter}{ติกติกปฏฺฐานปาฬิ}{mtk.2}
\bookmark[dest={mtk.2},level=0]{ติกติกปฏฺฐานปาฬิ}
\csromantocmarknopage{chapter}{ทุกทุกปฏฺฐานปาฬิ}{mtk.3}
\bookmark[dest={mtk.3},level=0]{ทุกทุกปฏฺฐานปาฬิ}
\gambhira{\textbf{ติกปฏฺฐานปาฬิ}}
\namakkaram{นโม ตสฺส ภควโต อรหโต สมฺมาสมฺพุทฺธสฺส.}
\csromanheader{2}{1. กุสลตฺติก 1. ปฏิจฺจวาร}
\csromanheader{2}{\textbf{ปจฺจยจตุกฺก-เหตุ-อารมฺมณ}}
\proseitem{1}{\csromanfolio{63}นกุสลํ ธมฺมํ ปฏิจฺจ นกุสโล ธมฺโม อุปฺปชฺชติ เหตุปจฺจยา.\csromanspacer{} อกุสลํ อพฺยากตํ เอกํ ขนฺธํ ปฏิจฺจ อกุสลา อพฺยากตา ตโย ขนฺธา จิตฺตสมุฏฺฐานญฺจ รูปํ ฯเปฯ เทฺว ขนฺเธ ปฏิจฺจ เทฺว ขนฺธา จิตฺตสมุฏฺฐานญฺจ รูปํ.\csromanspacer{}\csromanspacer{}นกุสลํ ธมฺมํ ปฏิจฺจ น-อกุสโล ธมฺโม อุปฺปชฺชติ เหตุปจฺจยา.\csromanspacer{}\csromanspacer{}นกุสลํ ธมฺมํ ปฏิจฺจ น-อพฺยากโต ธมฺโม อุปฺปชฺชติ เหตุปจฺจยา.\csromanspacer{}\csromanspacer{}นกุสลํ ธมฺมํ ปฏิจฺจ นกุสโล จ น-อพฺยากโต จ ธมฺมา อุปฺปชฺชนฺติ เหตุปจฺจยา.\csromanspacer{}\csromanspacer{}นกุสลํ ธมฺมํ ปฏิจฺจ นกุสโล จ น-อกุสโล จ ธมฺมา อุปฺปชฺชนฺติ เหตุปจฺจยา.\csromanspacer{} ปญฺจ.}

\proseitem{2}{น-อกุสลํ ธมฺมํ ปฏิจฺจ น-อกุสโล ธมฺโม อุปฺปชฺชติ เหตุปจฺจยา.\csromanspacer{}\csromanspacer{}น-อกุสลํ ธมฺมํ ปฏิจฺจ นกุสโล ธมฺโม อุปฺปชฺชติ เหตุปจฺจยา.\csromanspacer{}\csromanspacer{}น-อกุสลํ ธมฺมํ ปฏิจฺจ น-อพฺยากโต ธมฺโม อุปฺปชฺชติ เหตุปจฺจยา.\csromanspacer{}\csromanspacer{}น-อกุสลํ ธมฺมํ ปฏิจฺจ น-อกุสโล จ น-อพฺยากโต จ ธมฺมา อุปฺปชฺชนฺติ เหตุปจฺจยา.\csromanspacer{}\csromanspacer{}น-อกุสลํ ธมฺมํ ปฏิจฺจ นกุสโล จ น-อกุสโล จ ธมฺมา อุปฺปชฺชนฺติ เหตุปจฺจยา.\csromanspacer{} ปญฺจ.}

\csromanmatikaanchor{mtk.4}
\csromantocmarknopage{chapter}{ติกปฏฺฐานปาฬิ}{mtk.4}
\bookmark[dest={mtk.4},level=0]{ติกปฏฺฐานปาฬิ}
\gambhira{ธมฺมปจฺจนีย ติกปฏฺฐานปาฬิ}
\proseitem{3}{\csromanfolio{64}น-อพฺยากตํ ธมฺมํ ปฏิจฺจ น-อพฺยากโต ธมฺโม อุปฺปชฺชติ เหตุปจฺจยา.\csromanspacer{}\csromanspacer{}น-อพฺยากตํ ธมฺมํ ปฏิจฺจ นกุสโล ธมฺโม อุปฺปชฺชติ เหตุปจฺจยา.\csromanspacer{}\csromanspacer{}น-อพฺยากตํ ธมฺมํ ปฏิจฺจ น-อกุสโล ธมฺโม อุปฺปชฺชติ เหตุปจฺจยา.\csromanspacer{}\csromanspacer{}น-อพฺยากตํ ธมฺมํ ปฏิจฺจ นกุสโล จ น-อพฺยากโต จ ธมฺมา อุปฺปชฺชนฺติ เหตุปจฺจยา.\csromanspacer{}\csromanspacer{}น-อพฺยากตํ ธมฺมํ ปฏิจฺจ น-อกุสโล จ น-อพฺยากโต จ ธมฺมา อุปฺปชฺชนฺติ เหตุปจฺจยา.\csromanspacer{}\csromanspacer{}น-อพฺยากตํ ธมฺมํ ปฏิจฺจ นกุสโล จ น-อกุสโล จ ธมฺมา อุปฺปชฺชนฺติ เหตุปจฺจยา.\csromanspacer{} ฉ.}

\proseitem{4}{นกุสลญฺจ น-อพฺยากตญฺจ ธมฺมํ ปฏิจฺจ นกุสโล ธมฺโม อุปฺปชฺชติ เหตุปจฺจยา.\csromanspacer{}\csromanspacer{}นกุสลญฺจ น-อพฺยากตญฺจ ธมฺมํ ปฏิจฺจ น-อกุสโล ธมฺโม อุปฺปชฺชติ เหตุปจฺจยา.\csromanspacer{}\csromanspacer{}นกุสลญฺจ น-อพฺยากตญฺจ ธมฺมํ ปฏิจฺจ น-อพฺยากโต ธมฺโม อุปฺปชฺชติ เหตุปจฺจยา.\csromanspacer{}\csromanspacer{}นกุสลญฺจ น-อพฺยากตญฺจ ธมฺมํ ปฏิจฺจ นกุสโล จ น-อพฺยากโต จ ธมฺมา อุปฺปชฺชนฺติ เหตุปจฺจยา.\csromanspacer{}\csromanspacer{}นกุสลญฺจ น-อพฺยากตญฺจ ธมฺมํ ปฏิจฺจ นกุสโล จ น-อกุสโล จ ธมฺมา อุปฺปชฺชนฺติ เหตุปจฺจยา.\csromanspacer{} ปญฺจ.}

\proseitem{5}{น-อกุสลญฺจ น-อพฺยากตญฺจ ธมฺมํ ปฏิจฺจ นกุสโล ธมฺโม อุปฺปชฺชติ เหตุปจฺจยา.\csromanspacer{}\csromanspacer{}น-อกุสลญฺจ น-อพฺยากตญฺจ ธมฺมํ ปฏิจฺจ น-อกุสโล ธมฺโม อุปฺปชฺชติ เหตุปจฺจยา.\csromanspacer{}\csromanspacer{}น-อกุสลญฺจ น-อพฺยากตญฺจ ธมฺมํ ปฏิจฺจ น-อพฺยากโต ธมฺโม อุปฺปชฺชติ เหตุปจฺจยา.\csromanspacer{}\csromanspacer{}น-อกุสลญฺจ น-อพฺยากตญฺจ ธมฺมํ ปฏิจฺจ น-อกุสโล จ น-อพฺยากโต จ ธมฺมา อุปฺปชฺชนฺติ เหตุปจฺจยา.\csromanspacer{}\csromanspacer{}น-อกุสลญฺจ น-อพฺยากตญฺจ ธมฺมํ ปฏิจฺจ นกุสโล จ น-อกุสโล จ ธมฺมา อุปฺปชฺชนฺติ เหตุปจฺจยา.\csromanspacer{} ปญฺจ.}

\proseitem{6}{นกุสลญฺจ น-อกุสลญฺจ ธมฺมํ ปฏิจฺจ นกุสโล ธมฺโม อุปฺปชฺชติ เหตุปจฺจยา.\csromanspacer{}\csromanspacer{}นกุสลญฺจ น-อกุสลญฺจ ธมฺมํ ปฏิจฺจ น-อกุสโล ธมฺโม อุปฺปชฺชติ เหตุปจฺจยา.\csromanspacer{}\csromanspacer{}นกุสลญฺจ น-อกุสลญฺจ ธมฺมํ ปฏิจฺจ นกุสโล จ น-อกุสโล จ ธมฺมา อุปฺปชฺชนฺติ เหตุปจฺจยา.\csromanspacer{} ตีณิ.}

\prose{นกุสลํ ธมฺมํ ปฏิจฺจ นกุสโล ธมฺโม อุปฺปชฺชติ อารมฺมณปจฺจยา. (สํขิตฺตํ.)}

\proseitem{7}{เหตุยา เอกูนติํส, อารมฺมเณ จตุวีส ฯเปฯ อวิคเต เอกูนติํส.}

\csromanheader{2}{1. กุสลตฺติก}
\prose{\csromanfolio{65}(สหชาตวารมฺปิ ปจฺจยวารมฺปิ นิสฺสยวารมฺปิ สํสฏฺฐวารมฺปิ สมฺปยุตฺตวารมฺปิ ปฏิจฺจวารสทิสํ.)}

\csromanheader{2}{1. กุสลตฺติก 7. ปญฺหาวาร}
\csromanheader{2}{\textbf{ปจฺจยจตุกฺก-เหตุ-อารมฺมณาทิ}}
\proseitem{8}{นกุสโล ธมฺโม นกุสลสฺส ธมฺมสฺส เหตุปจฺจเยน ปจฺจโย ฯเปฯ.\csromanspacer{} นกุสโล ธมฺโม นกุสลสฺส ธมฺมสฺส อารมฺมณปจฺจเยน ปจฺจโย.\csromanspacer{}\csromanspacer{}นกุสโล ธมฺโม น-อกุสลสฺส ธมฺมสฺส อารมฺมณปจฺจเยน ปจฺจโย.\csromanspacer{}\csromanspacer{}นกุสโล ธมฺโม น-อพฺยากตสฺส ธมฺมสฺส อารมฺมณปจฺจเยน ปจฺจโย.\csromanspacer{}\csromanspacer{}นกุสโล ธมฺโม นกุสลสฺส จ น-อพฺยากตสฺส จ ธมฺมสฺส อารมฺมณปจฺจเยน ปจฺจโย.\csromanspacer{}\csromanspacer{}นกุสโล ธมฺโม น-อกุสลสฺส จ น-อพฺยากตสฺส จ ธมฺมสฺส อารมฺมณปจฺจเยน ปจฺจโย.\csromanspacer{}\csromanspacer{}นกุสโล ธมฺโม นกุสลสฺส จ น-อกุสลสฺส จ ธมฺมสฺส อารมฺมณปจฺจเยน ปจฺจโย.\csromanspacer{} ฉ.}

\prose{น-อกุสโล ธมฺโม น-อกุสลสฺส ธมฺมสฺส อารมฺมณปจฺจเยน ปจฺจโย.\csromanspacer{} ฉ.}

\prose{น-อพฺยากโต ธมฺโม น-อพฺยากตสฺส ธมฺมสฺส อารมฺมณปจฺจเยน ปจฺจโย.\csromanspacer{} ฉ.}

\prose{นกุสโล จ น-อพฺยากโต จ ธมฺมา นกุสลสฺส ธมฺมสฺส อารมฺมณปจฺจเยน ปจฺจโย.\csromanspacer{} ฉ.}

\prose{น-อกุสโล จ น-อพฺยากโต จ ธมฺมา นกุสลสฺส ธมฺมสฺส อารมฺมณปจฺจเยน ปจฺจโย.\csromanspacer{} ฉ.}

\prose{นกุสโล จ น-อกุสโล จ ธมฺมา นกุสลสฺส ธมฺมสฺส อารมฺมณปจฺจเยน ปจฺจโย.\csromanspacer{} ฉ.}

\proseitem{9}{นกุสโล ธมฺโม นกุสลสฺส ธมฺมสฺส อธิปติปจฺจเยน ปจฺจโย.\csromanspacer{} อนนฺตรปจฺจเยน ปจฺจโย.\csromanspacer{} สมนนฺตรปจฺจเยน ปจฺจโย.\csromanspacer{} สหชาตปจฺจเยน ปจฺจโย.\csromanspacer{} อญฺญมญฺญปจฺจเยน ปจฺจโย.\csromanspacer{} นิสฺสยปจฺจเยน ปจฺจโย.\csromanspacer{} อุปนิสฺสยปจฺจเยน ปจฺจโย.\csromanspacer{} ปุเรชาตปจฺจเยน ปจฺจโย.\csromanspacer{} นกุสโล ธมฺโม น-อกุสลสฺส ธมฺมสฺส ปุเรชาตปจฺจเยน ปจฺจโย ฯเปฯ.\csromanspacer{}\csromanspacer{}นกุสโล ธมฺโม นกุสลสฺส จ น-อกุสลสฺส จ ธมฺมสฺส ปุเรชาตปจฺจเยน ปจฺจโย.\csromanspacer{} ฉ.}

\gambhira{ธมฺมปจฺจนีย ติกปฏฺฐานปาฬิ}
\prose{\csromanfolio{66}น-อกุสโล ธมฺโม น-อกุสลสฺส ธมฺมสฺส ปุเรชาตปจฺจเยน ปจฺจโย.\csromanspacer{}\csromanspacer{}น-อกุสโล ธมฺโม นกุสลสฺส ธมฺมสฺส ปุเรชาตปจฺจเยน ปจฺจโย ฯเปฯ.\csromanspacer{}\csromanspacer{}น-อกุสโล ธมฺโม นกุสลสฺส จ น-อกุสลสฺส จ ธมฺมสฺส ปุเรชาตปจฺจเยน ปจฺจโย.\csromanspacer{} ฉ.}

\prose{นกุสโล จ น-อกุสโล จ ธมฺมา นกุสลสฺส ธมฺมสฺส ปุเรชาตปจฺจเยน ปจฺจโย.\csromanspacer{}\csromanspacer{}นกุสโล จ น-อกุสโล จ ธมฺมา น-อกุสลสฺส ธมฺมสฺส ปุเรชาตปจฺจเยน ปจฺจโย ฯเปฯ.\csromanspacer{}\csromanspacer{}นกุสโล จ น-อกุสโล จ ธมฺมา นกุสลสฺส จ น-อกุสลสฺส จ ธมฺมสฺส ปุเรชาตปจฺจเยน ปจฺจโย.\csromanspacer{} ฉ. (สํขิตฺตํ.)}

\proseitem{10}{เหตุยา เอกูนติํส, อารมฺมเณ ฉตฺติํส, อธิปติยา ปญฺจติํส, อนนฺตเร จตุตฺติํส, สมนนฺตเร จตุตฺติํส, สหชาเต เอกูนติํส, อญฺญมญฺเญ จตุวีส, นิสฺสเย จตุตฺติํส, อุปนิสฺสเย ฉตฺติํส, ปุเรชาเต อฏฺฐารส, ปจฺฉาชาเต อฏฺฐารส, อาเสวเน จตุวีส, กมฺเม เอกูนติํส, วิปาเก นว, อาหาเร เอกูนติํส ฯเปฯ อวิคเต จตุตฺติํส.}

\prose{(ยถา กุสลตฺติเก ปญฺหาวารสฺส อนุโลมมฺปิ ปจฺจนียมฺปิ อนุโลมปจฺจนียมฺปิ ปจฺจนียานุโลมมฺปิ คณิตํ, เอวํ คเณตพฺพํ.)}

\csromanheader{2}{2. เวทนาตฺติก 1. ปฏิจฺจวาร}
\proseitem{11}{นสุขาย เวทนาย สมฺปยุตฺตํ ธมฺมํ ปฏิจฺจ นสุขาย เวทนาย สมฺปยุตฺโต ธมฺโม อุปฺปชฺชติ เหตุปจฺจยา.\csromanspacer{}\csromanspacer{}นสุขาย เวทนาย สมฺปยุตฺตํ ธมฺมํ ปฏิจฺจ นทุกฺขาย เวทนาย สมฺปยุตฺโต ธมฺโม อุปฺปชฺชติ เหตุปจฺจยา.\csromanspacer{}\csromanspacer{}นสุขาย เวทนาย สมฺปยุตฺตํ ธมฺมํ ปฏิจฺจ น-อทุกฺขมสุขาย เวทนาย สมฺปยุตฺโต ธมฺโม อุปฺปชฺชติ เหตุปจฺจยา.\csromanspacer{}\csromanspacer{}นสุขาย เวทนาย สมฺปยุตฺตํ ธมฺมํ ปฏิจฺจ นสุขาย เวทนาย สมฺปยุตฺโต จ น-อทุกฺขมสุขาย เวทนาย สมฺปยุตฺโต จ ธมฺมา อุปฺปชฺชนฺติ เหตุปจฺจยา.\csromanspacer{}\csromanspacer{}นสุขาย เวทนาย สมฺปยุตฺตํ ธมฺมํ ปฏิจฺจ นทุกฺขาย เวทนาย สมฺปยุตฺโต จ น-อทุกฺขมสุขาย เวทนาย สมฺปยุตฺโต จ ธมฺมา อุปฺปชฺชนฺติ เหตุปจฺจยา.\csromanspacer{}\csromanspacer{}นสุขาย เวทนาย สมฺปยุตฺตํ ธมฺมํ ปฏิจฺจ นสุขาย เวทนาย สมฺปยุตฺโต จ นทุกฺขาย เวทนาย สมฺปยุตฺโต จ ธมฺมา อุปฺปชฺชนฺติ เหตุปจฺจยา.\csromanspacer{}\csromanspacer{}}

\vspace{\tipitakaparskip}
\csromancenter{สํกิลิฏฺฐตฺติก นสุขาย เวทนาย สมฺปยุตฺตํ ธมฺมํ ปฏิจฺจ นสุขาย เวทนาย สมฺปยุตฺโต จ นทุกฺขาย เวทนาย สมฺปยุตฺโต จ น-อทุกฺขมสุขาย เวทนาย สมฺปยุตฺโต จ ธมฺมา อุปฺปชฺชนฺติ เหตุปจฺจยา. (7)}
\prose{\csromanfolio{67}นทุกฺขาย เวทนาย สมฺปยุตฺตํ ธมฺมํ ปฏิจฺจ นทุกฺขาย เวทนาย สมฺปยุตฺโต ธมฺโม อุปฺปชฺชติ เหตุปจฺจยา.\csromanspacer{} สตฺต.}

\prose{น-อทุกฺขมสุขาย เวทนาย สมฺปยุตฺตํ ธมฺมํ ปฏิจฺจ น-อทุกฺขมสุขาย เวทนาย สมฺปยุตฺโต ธมฺโม อุปฺปชฺชติ เหตุปจฺจยา.\csromanspacer{} สตฺต. (สํขิตฺตํ.)}

\csromanheader{2}{3. วิปากตฺติก 1. ปฏิจฺจวาร}
\proseitem{12}{นวิปากํ ธมฺมํ ปฏิจฺจ นวิปาโก ธมฺโม อุปฺปชฺชติ เหตุปจฺจยา ฯเปฯ.\csromanspacer{} นวิปากธมฺมธมฺมํ ปฏิจฺจ นวิปากธมฺมธมฺโม อุปฺปชฺชติ เหตุปจฺจยา ฯเปฯ.}

\prose{นเนววิปากนวิปากธมฺมธมฺมํ ปฏิจฺจ นเนววิปากนวิปากธมฺมธมฺโม อุปฺปชฺชติ เหตุปจฺจยา.}

\csromanheader{2}{4. อุปาทินฺนตฺติก 1. ปฏิจฺจวาร}
\proseitem{13}{น-อุปาทินฺนุปาทานิยํ ธมฺมํ ปฏิจฺจ น-อุปาทินฺนุปาทานิโย ธมฺโม อุปฺปชฺชติ เหตุปจฺจยา ฯเปฯ.}

\prose{น-อนุปาทินฺนุปาทานิยํ ธมฺมํ ปฏิจฺจ น-อนุปาทินฺนุปาทานิโย ธมฺโม อุปฺปชฺชติ เหตุปจฺจยา ฯเปฯ.}

\prose{น-อนุปาทินฺน-อนุปาทานิยํ ธมฺมํ ปฏิจฺจ น-อนุปาทินฺน-อนุปาทานิโย ธมฺโม อุปฺปชฺชติ เหตุปจฺจยา. (สํขิตฺตํ.)}

\csromanheader{2}{5. สํกิลิฏฺฐตฺติก 1. ปฏิจฺจวาร}
\proseitem{14}{นสํกิลิฏฺฐสํกิเลสิกํ ธมฺมํ ปฏิจฺจ นสํกิลิฏฺฐสํกิเลสิโก ธมฺโม อุปฺปชฺชติ เหตุปจฺจยา ฯเปฯ.}

\prose{น-อสํกิลิฏฺฐสํกิเลสิกํ ธมฺมํ ปฏิจฺจ น-อสํกิลิฏฺฐสํกิเลสิโก ธมฺโม อุปฺปชฺชติ เหตุปจฺจยา ฯเปฯ.}

\prose{น-อสํกิลิฏฺฐ-อสํกิเลสิกํ ธมฺมํ ปฏิจฺจ น-อสํกิลิฏฺฐ-อสํกิเลสิโก ธมฺโม อุปฺปชฺชติ เหตุปจฺจยา. (สํขิตฺตํ.)}

\csromanheader{2}{ธมฺมปจฺจนีย ติกปฏฺฐานปาฬิ \textbf{6. วิตกฺกตฺติก 1. ปฏิจฺจวาร}}
\proseitem{15}{\csromanfolio{68}นสวิตกฺกสวิจารํ ธมฺมํ ปฏิจฺจ นสวิตกฺกสวิจาโร ธมฺโม อุปฺปชฺชติ เหตุปจฺจยา ฯเปฯ.}

\prose{น-อวิตกฺกวิจารมตฺตํ ธมฺมํ ปฏิจฺจ น-อวิตกฺกวิจารมตฺโต ธมฺโม อุปฺปชฺชติ เหตุปจฺจยา ฯเปฯ.}

\prose{น-อวิตกฺก-อวิจารํ ธมฺมํ ปฏิจฺจ น-อวิตกฺก-อวิจาโร ธมฺโม อุปฺปชฺชติ เหตุปจฺจยา. (สํขิตฺตํ.) \textbf{7.}\csromanspacer{}\textbf{ ปีติตฺติก 1.}\csromanspacer{}\textbf{ ปฏิจฺจวาร}}

\proseitem{16}{นปีติสหคตํ ธมฺมํ ปฏิจฺจ นปีติสหคโต ธมฺโม อุปฺปชฺชติ เหตุปจฺจยา ฯเปฯ.}

\prose{นสุขสหคตํ ธมฺมํ ปฏิจฺจ ฯเปฯ.\csromanspacer{}\csromanspacer{}น-อุเปกฺขาสหคตํ ธมฺมํ ปฏิจฺจ. (สํขิตฺตํ.) \textbf{8.}\csromanspacer{}\textbf{ ทสฺสนตฺติก 1.}\csromanspacer{}\textbf{ ปฏิจฺจวาร}}

\proseitem{17}{นทสฺสเนน ปหาตพฺพํ ธมฺมํ ปฏิจฺจ ฯเปฯ.\csromanspacer{}\csromanspacer{}นภาวนาย ปหาตพฺพํ ธมฺมํ ปฏิจฺจ ฯเปฯ.\csromanspacer{}\csromanspacer{}นเนวทสฺสเนน นภาวนาย ปหาตพฺพํ ธมฺมํ ปฏิจฺจ. (สํขิตฺตํ.) \textbf{9.}\csromanspacer{}\textbf{ ทสฺสนเหตุตฺติก 1.}\csromanspacer{}\textbf{ ปฏิจฺจวาร}}

\proseitem{18}{นทสฺสเนน ปหาตพฺพเหตุกํ ธมฺมํ ปฏิจฺจ ฯเปฯ.\csromanspacer{}\csromanspacer{}นภาวนาย ปหาตพฺพเหตุกํ ธมฺมํ ปฏิจฺจ ฯเปฯ.\csromanspacer{}\csromanspacer{}นเนวทสฺสเนน นภาวนาย ปหาตพฺพเหตุกํ ธมฺมํ ปฏิจฺจ. \textbf{10.}\csromanspacer{}\textbf{ อาจยคามิตฺติก 1.}\csromanspacer{}\textbf{ ปฏิจฺจวาร}}

\proseitem{19}{น-อาจยคามิํ ธมฺมํ ปฏิจฺจ ฯเปฯ.\csromanspacer{} น-อปจยคามิํ ธมฺมํ ปฏิจฺจ ฯเปฯ.\csromanspacer{}\csromanspacer{}นเนวาจยคามินาปจยคามิํ ธมฺมํ ปฏิจฺจ. (สํขิตฺตํ.) \textbf{11.}\csromanspacer{}\textbf{ เสกฺขตฺติก 1.}\csromanspacer{}\textbf{ ปฏิจฺจวาร}}

\proseitem{20}{นเสกฺขํ ธมฺมํ ปฏิจฺจ ฯเปฯ.\csromanspacer{}\csromanspacer{}น-อเสกฺขํ ธมฺมํ ปฏิจฺจ ฯเปฯ.\csromanspacer{}\csromanspacer{}นเนวเสกฺขนาเสกฺขํ ธมฺมํ ปฏิจฺจ. (สํขิตฺตํ.)}

\csromanheader{2}{17. อุปฺปนฺนตฺติก \textbf{12. ปริตฺตตฺติก 1. ปฏิจฺจวาร}}
\proseitem{21}{\csromanfolio{69}นปริตฺตํ ธมฺมํ ปฏิจฺจ ฯเปฯ.\csromanspacer{}\csromanspacer{}นมหคฺคตํ ธมฺมํ ปฏิจฺจ ฯเปฯ.\csromanspacer{}\csromanspacer{}น-อปฺปมาณํ ธมฺมํ ปฏิจฺจ. (สํขิตฺตํ.) \textbf{13.}\csromanspacer{}\textbf{ ปริตฺตารมฺมณตฺติก 1.}\csromanspacer{}\textbf{ ปฏิจฺจวาร}}

\proseitem{22}{นปริตฺตารมฺมณํ ธมฺมํ ปฏิจฺจ ฯเปฯ.\csromanspacer{}\csromanspacer{}นมหคฺคตารมฺมณํ ธมฺมํ ปฏิจฺจ ฯเปฯ.\csromanspacer{}\csromanspacer{}น-อปฺปมาณารมฺมณํ ธมฺมํ ปฏิจฺจ. (สํขิตฺตํ.) \textbf{14.}\csromanspacer{}\textbf{ หีนตฺติก 1.}\csromanspacer{}\textbf{ ปฏิจฺจวาร}}

\proseitem{23}{นหีนํ ธมฺมํ ปฏิจฺจ ฯเปฯ.\csromanspacer{}\csromanspacer{}นมชฺฌิมํ ธมฺมํ ปฏิจฺจ ฯเปฯ.\csromanspacer{}\csromanspacer{}นปณีตํ ธมฺมํ ปฏิจฺจ. (สํขิตฺตํ.) \textbf{15.}\csromanspacer{}\textbf{ มิจฺฉตฺตตฺติก 1.}\csromanspacer{}\textbf{ ปฏิจฺจวาร}}

\proseitem{24}{นมิจฺฉตฺตนิยตํ ธมฺมํ ปฏิจฺจ ฯเปฯ.\csromanspacer{}\csromanspacer{}นสมฺมตฺตนิยตํ ธมฺมํ ปฏิจฺจ ฯเปฯ.\csromanspacer{}\csromanspacer{}น-อนิยตํ ธมฺมํ ปฏิจฺจ. (สํขิตฺตํ.) \textbf{16.}\csromanspacer{}\textbf{ มคฺคารมฺมณตฺติก 1.}\csromanspacer{}\textbf{ ปฏิจฺจวาร}}

\proseitem{25}{นมคฺคารมฺมณํ ธมฺมํ ปฏิจฺจ ฯเปฯ.\csromanspacer{}\csromanspacer{}นมคฺคเหตุกํ ธมฺมํ ปฏิจฺจ ฯเปฯ.\csromanspacer{}\csromanspacer{}นมคฺคาธิปติํ ธมฺมํ ปฏิจฺจ. \textbf{17.}\csromanspacer{}\textbf{ อุปฺปนฺนตฺติก 1.}\csromanspacer{}\textbf{ ปฏิจฺจวาร}}

\proseitem{26}{น-อนุปฺปนฺนํ ธมฺมํ ปฏิจฺจ น-อนุปฺปนฺโน ธมฺโม อุปฺปชฺชติ เหตุปจฺจยา.\csromanspacer{}\csromanspacer{}น-อนุปฺปนฺนํ ธมฺมํ ปฏิจฺจ น-อุปฺปาที ธมฺโม อุปฺปชฺชติ เหตุปจฺจยา.\csromanspacer{}\csromanspacer{}น-อนุปฺปนฺนํ ธมฺมํ ปฏิจฺจ น-อนุปฺปนฺโน จ น-อุปฺปาที จ ธมฺมา อุปฺปชฺชนฺติ เหตุปจฺจยา. (3)}

\prose{น-อุปฺปาทิํ ธมฺมํ ปฏิจฺจ น-อุปฺปาที ธมฺโม อุปฺปชฺชติ เหตุปจฺจยา.\csromanspacer{}\csromanspacer{}น-อุปฺปาทิํ ธมฺมํ ปฏิจฺจ น-อนุปฺปนฺโน ธมฺโม อุปฺปชฺชติ เหตุปจฺจยา.\csromanspacer{}\csromanspacer{}น-อุปฺปาทิํ ธมฺมํ ปฏิจฺจ น-อนุปฺปนฺโน จ น-อุปฺปาที จ ธมฺมา อุปฺปชฺชนฺติ เหตุปจฺจยา. (3)}

\prose{น-อนุปฺปนฺนญฺจ น-อุปฺปาทิญฺจ ธมฺมํ ปฏิจฺจ น-อนุปฺปนฺโน ธมฺโม อุปฺปชฺชติ เหตุปจฺจยา.\csromanspacer{}\csromanspacer{}น-อนุปฺปนฺนญฺจ น-อุปฺปาทิญฺจ ธมฺมํ ปฏิจฺจ น-อุปฺปาที ธมฺโม อุปฺปชฺชติ เหตุปจฺจยา.\csromanspacer{}\csromanspacer{}น-อนุปฺปนฺนญฺจ น-อุปฺปาทิญฺจ ธมฺมํ ปฏิจฺจ น-อนุปฺปนฺโน จ น-อุปฺปาที จ ธมฺมา อุปฺปชฺชนฺติ เหตุปจฺจยา. (3) (สํขิตฺตํ.)}

\csromanheader{2}{ธมฺมปจฺจนีย ติกปฏฺฐานปาฬิ \textbf{18. อตีตตฺติก 1-7. ปฏิจฺจาทิวาร}}
\proseitem{27}{\csromanfolio{70}น-อตีตํ ธมฺมํ ปฏิจฺจ น-อตีโต ธมฺโม อุปฺปชฺชติ เหตุปจฺจยา. (สํขิตฺตํ.)}

\prose{น-อตีโต ธมฺโม น-อตีตสฺส ธมฺมสฺส เหตุปจฺจเยน ปจฺจโย.\csromanspacer{}\csromanspacer{}น-อตีโต ธมฺโม น-อนาคตสฺส ธมฺมสฺส เหตุปจฺจเยน ปจฺจโย.\csromanspacer{}\csromanspacer{}น-อตีโต ธมฺโม น-อตีตสฺส จ น-อนาคตสฺส จ ธมฺมสฺส เหตุปจฺจเยน ปจฺจโย. (3)}

\prose{น-อนาคโต ธมฺโม น-อนาคตสฺส ธมฺมสฺส เหตุปจฺจเยน ปจฺจโย.\csromanspacer{}\csromanspacer{}น-อนาคโต ธมฺโม น-อตีตสฺส ธมฺมสฺส เหตุปจฺจเยน ปจฺจโย.\csromanspacer{}\csromanspacer{}น-อนาคโต ธมฺโม น-อตีตสฺส จ น-อนาคตสฺส จ ธมฺมสฺส เหตุปจฺจเยน ปจฺจโย. (3)}

\prose{น-อตีโต จ น-อนาคโต จ ธมฺมา น-อตีตสฺส ธมฺมสฺส เหตุปจฺจเยน ปจฺจโย.\csromanspacer{}\csromanspacer{}น-อตีโต จ น-อนาคโต จ ธมฺมา น-อนาคตสฺส ธมฺมสฺส เหตุปจฺจเยน ปจฺจโย.\csromanspacer{}\csromanspacer{}น-อตีโต จ น-อนาคโต จ ธมฺมา น-อตีตสฺส จ น-อนาคตสฺส จ ธมฺมสฺส เหตุปจฺจเยน ปจฺจโย. (3)}

\csromanheader{2}{19. อตีตารมฺมณตฺติก 1. ปฏิจฺจวาร}
\proseitem{28}{น-อตีตารมฺมณํ ธมฺมํ ปฏิจฺจ ฯเปฯ.\csromanspacer{} น-อนาคตารมฺมณํ ธมฺมํ ปฏิจฺจ ฯเปฯ.\csromanspacer{} นปจฺจุปฺปนฺนารมฺมณํ ธมฺมํ ปฏิจฺจ.}

\csromanheader{2}{20. อชฺฌตฺตตฺติก 1. ปฏิจฺจวาร}
\proseitem{29}{น-อชฺฌตฺตํ ธมฺมํ ปฏิจฺจ น-อชฺฌตฺโต ธมฺโม อุปฺปชฺชติ เหตุปจฺจยา.\csromanspacer{} นพหิทฺธา ธมฺมํ ปฏิจฺจ นพหิทฺธา ธมฺโม อุปฺปชฺชติ เหตุปจฺจยา. (สํขิตฺตํ.)}

\csromanheader{2}{21. อชฺฌตฺตารมฺมณตฺติก 1. ปฏิจฺจวาร}
\proseitem{30}{น-อชฺฌตฺตารมฺมณํ ธมฺมํ ปฏิจฺจ น-อชฺฌตฺตารมฺมโณ ธมฺโม อุปฺปชฺชติ เหตุปจฺจยา.\csromanspacer{}\csromanspacer{}น-อชฺฌตฺตารมฺมณํ ธมฺมํ ปฏิจฺจ นพหิทฺธารมฺมโณ ธมฺโม อุปฺปชฺชติ เหตุปจฺจยา. (2)}

\csromanheader{2}{22. สนิทสฺสนตฺติก}
\prose{\csromanfolio{71}นพหิทฺธารมฺมณํ ธมฺมํ ปฏิจฺจ นพหิทฺธารมฺมโณ ธมฺโม อุปฺปชฺชติ เหตุปจฺจยา.\csromanspacer{}\csromanspacer{}นพหิทฺธารมฺมณํ ธมฺมํ ปฏิจฺจ น-อชฺฌตฺตารมฺมโณ ธมฺโม อุปฺปชฺชติ เหตุปจฺจยา. (2) (สํขิตฺตํ.)}

\csromanheader{2}{22. สนิทสฺสนตฺติก 1. ปฏิจฺจวาร}
\proseitem{31}{นสนิทสฺสนสปฺปฏิฆํ ธมฺมํ ปฏิจฺจ นสนิทสฺสนสปฺปฏิโฆ ธมฺโม อุปฺปชฺชติ เหตุปจฺจยา.\csromanspacer{}\csromanspacer{}นสนิทสฺสนสปฺปฏิฆํ ธมฺมํ ปฏิจฺจ น-อนิทสฺสนสปฺปฏิโฆ ธมฺโม อุปฺปชฺชติ เหตุปจฺจยา.\csromanspacer{}\csromanspacer{}นสนิทสฺสนสปฺปฏิฆํ ธมฺมํ ปฏิจฺจ น-อนิทสฺสน-อปฺปฏิโฆ ธมฺโม อุปฺปชฺชติ เหตุปจฺจยา.\csromanspacer{}\csromanspacer{}นสนิทสฺสนสปฺปฏิฆํ ธมฺมํ ปฏิจฺจ นสนิทสฺสนสปฺปฏิโฆ จ น-อนิทสฺสน-อปฺปฏิโฆ จ ธมฺมา อุปฺปชฺชนฺติ เหตุปจฺจยา.\csromanspacer{}\csromanspacer{}นสนิทสฺสนสปฺปฏิฆํ ธมฺมํ ปฏิจฺจ น-อนิทสฺสนสปฺปฏิโฆ จ น-อนิทสฺสน-อปฺปฏิโฆ จ ธมฺมา อุปฺปชฺชนฺติ เหตุปจฺจยา.\csromanspacer{}\csromanspacer{}นสนิทสฺสนสปฺปฏิฆํ ธมฺมํ ปฏิจฺจ นสนิทสฺสนสปฺปฏิโฆ จ น-อนิทสฺสนสปฺปฏิโฆ จ ธมฺมา อุปฺปชฺชนฺติ เหตุปจฺจยา. (6)}

\proseitem{32}{น-อนิทสฺสนสปฺปฏิฆํ ธมฺมํ ปฏิจฺจ น-อนิทสฺสนสปฺปฏิโฆ ธมฺโม อุปฺปชฺชติ เหตุปจฺจยา.\csromanspacer{}\csromanspacer{}น-อนิทสฺสนสปฺปฏิฆํ ธมฺมํ ปฏิจฺจ นสนิทสฺสนสปฺปฏิโฆ ธมฺโม อุปฺปชฺชติ เหตุปจฺจยา.\csromanspacer{}\csromanspacer{}น-อนิทสฺสนสปฺปฏิฆํ ธมฺมํ ปฏิจฺจ น-อนิทสฺสน-อปฺปฏิโฆ ธมฺโม อุปฺปชฺชติ เหตุปจฺจยา.\csromanspacer{}\csromanspacer{}น-อนิทสฺสนสปฺปฏิฆํ ธมฺมํ ปฏิจฺจ นสนิทสฺสนสปฺปฏิโฆ จ น-อนิทสฺสน-อปฺปฏิโฆ จ ธมฺมา อุปฺปชฺชนฺติ เหตุปจฺจยา.\csromanspacer{}\csromanspacer{}น-อนิทสฺสนสปฺปฏิฆํ ธมฺมํ ปฏิจฺจ น-อนิทสฺสนสปฺปฏิโฆ จ น-อนิทสฺสน-อปฺปฏิโฆ จ ธมฺมา อุปฺปชฺชนฺติ เหตุปจฺจยา.\csromanspacer{}\csromanspacer{}น-อนิทสฺสนสปฺปฏิฆํ ธมฺมํ ปฏิจฺจ นสนิทสฺสนสปฺปฏิโฆ จ น-อนิทสฺสนสปฺปฏิโฆ จ ธมฺมา อุปฺปชฺชนฺติ เหตุปจฺจยา. (6)}

\proseitem{33}{น-อนิทสฺสน-อปฺปฏิฆํ ธมฺมํ ปฏิจฺจ น-อนิทสฺสน-อปฺปฏิโฆ ธมฺโม อุปฺปชฺชติ เหตุปจฺจยา.\csromanspacer{}\csromanspacer{}น-อนิทสฺสน-อปฺปฏิฆํ ธมฺมํ ปฏิจฺจ นสนิทสฺสนสปฺปฏิโฆ ธมฺโม อุปฺปชฺชติ เหตุปจฺจยา.\csromanspacer{}\csromanspacer{}น-อนิทสฺสน-อปฺปฏิฆํ ธมฺมํ ปฏิจฺจ น-อนิทสฺสนสปฺปฏิโฆ ธมฺโม อุปฺปชฺชติ เหตุปจฺจยา.\csromanspacer{} ฉ.}

\prose{นสนิทสฺสนสปฺปฏิฆญฺจ น-อนิทสฺสน-อปฺปฏิฆญฺจ ธมฺมํ ปฏิจฺจ นสนิทสฺสนสปฺปฏิโฆ ธมฺโม อุปฺปชฺชติ เหตุปจฺจยา.\csromanspacer{} ฉ.}

\gambhira{ธมฺมปจฺจนีย ติกปฏฺฐานปาฬิ}
\prose{\csromanfolio{72}นสนิทสฺสนสปฺปฏิฆญฺจ น-อนิทสฺสนสปฺปฏิฆญฺจ ธมฺมํ ปฏิจฺจ นสนิทสฺสนสปฺปฏิโฆ ธมฺโม อุปฺปชฺชติ เหตุปจฺจยา.\csromanspacer{} ฉ.}

\proseitem{34}{เหตุยา ติํส ฯเปฯ อวิคเต ติํส.}

\prose{(ยถา กุสลตฺติเก สหชาตวารมฺปิ ปจฺจยวารมฺปิ นิสฺสยวารมฺปิ สํสฏฺฐวารมฺปิ สมฺปยุตฺตวารมฺปิ ปญฺหาวารมฺปิ, คณิตํ, เอวํ คเณตพฺพํ.)}

\nitthitam{\textbf{ธมฺมปจฺจนีเย ติกปฏฺฐานํ นิฏฺฐิตํ.}}
\pitaka{\textbf{อภิธมฺมปิฏก}}
\csromanheader{2}{\textbf{ธมฺมปจฺจนีย}}
\csromanmatikaanchor{mtk.5}
\csromantocmarknopage{chapter}{ทุกปฏฺฐานปาฬิ}{mtk.5}
\bookmark[dest={mtk.5},level=0]{ทุกปฏฺฐานปาฬิ}
\csromantocmarknopage{chapter}{ทุกติกปฏฺฐานปาฬิ}{mtk.6}
\bookmark[dest={mtk.6},level=0]{ทุกติกปฏฺฐานปาฬิ}
\gambhira{\textbf{ทุกปฏฺฐานปาฬิ}}
\namakkaram{นโม ตสฺส ภควโต อรหโต สมฺมาสมฺพุทฺธสฺส.}
\csromanheader{2}{1. เหตุทุก 1-7. ปฏิจฺจาทิวาร}
\csromanheader{2}{\textbf{ปจฺจยจตุกฺก-เหตุ}}
\proseitem{1}{\csromanfolio{73}นเหตุํ ธมฺมํ ปฏิจฺจ นเหตุ ธมฺโม อุปฺปชฺชติ เหตุปจฺจยา.\csromanspacer{} นเหตุํ เอกํ ขนฺธํ ปฏิจฺจ ตโย ขนฺธา จิตฺตสมุฏฺฐานญฺจ รูปํ ฯเปฯ เทฺว ขนฺเธ ปฏิจฺจ เทฺว ขนฺธา จิตฺตสมุฏฺฐานญฺจ รูปํ ปฏิสนฺธิกฺขเณ ฯเปฯ. (ยาว มหาภูตา.) .\csromanspacer{} นเหตุํ ธมฺมํ ปฏิจฺจ น นเหตุ ธมฺโม อุปฺปชฺชติ เหตุปจฺจยา.\csromanspacer{} นเหตู ขนฺเธ ปฏิจฺจ เหตู.\csromanspacer{} ปฏิสนฺธิกฺขเณ ฯเปฯ.\csromanspacer{}\csromanspacer{}นเหตุํ ธมฺมํ ปฏิจฺจ นเหตุ จ น นเหตุ จ ธมฺมา อุปฺปชฺชนฺติ เหตุปจฺจยา. (3)}

\prose{น นเหตุํ ธมฺมํ ปฏิจฺจ น นเหตุ ธมฺโม อุปฺปชฺชติ เหตุปจฺจยา.\csromanspacer{} น นเหตุํ ธมฺมํ ปฏิจฺจ นเหตุ ธมฺโม อุปฺปชฺชติ เหตุปจฺจยา.\csromanspacer{}\csromanspacer{}น นเหตุํ ธมฺมํ ปฏิจฺจ นเหตุ จ น นเหตุ จ ธมฺมา อุปฺปชฺชนฺติ เหตุปจฺจยา. (3)}

\prose{นเหตุญฺจ น นเหตุญฺจ ธมฺมํ ปฏิจฺจ นเหตุ ธมฺโม อุปฺปชฺชติ เหตุปจฺจยา.\csromanspacer{}\csromanspacer{}นเหตุญฺจ น นเหตุญฺจ ธมฺมํ ปฏิจฺจ น นเหตุ ธมฺโม อุปฺปชฺชติ เหตุปจฺจยา.\csromanspacer{}\csromanspacer{}นเหตุญฺจ น นเหตุญฺจ ธมฺมํ ปฏิจฺจ นเหตุ จ น นเหตุ จ ธมฺมา อุปฺปชฺชนฺติ เหตุปจฺจยา. (3)}

\proseitem{2}{เหตุยา นว, อารมฺมเณ นว ฯเปฯ อวิคเต นว.}

\prose{(สหชาตวารมฺปิ ฯเปฯ สมฺปยุตฺตวารมฺปิ ปฏิจฺจวารสทิสํ, เอวํ วิตฺถาเรตพฺพํ.)}

\csromanheader{2}{ธมฺมปจฺจนีย ทุกปฏฺฐานปาฬิ \textbf{เหตุ-อารมฺมณ}}
\proseitem{3}{\csromanfolio{74}น นเหตุ ธมฺโม น นเหตุสฺส ธมฺมสฺส เหตุปจฺจเยน ปจฺจโย.\csromanspacer{}\csromanspacer{}น นเหตุ ธมฺโม นเหตุสฺส ธมฺมสฺส เหตุปจฺจเยน ปจฺจโย.\csromanspacer{}\csromanspacer{}น นเหตุ ธมฺโม นเหตุสฺส จ น นเหตุสฺส จ ธมฺมสฺส เหตุปจฺจเยน ปจฺจโย. (3)}

\prose{นเหตุ ธมฺโม นเหตุสฺส ธมฺมสฺส อารมฺมณปจฺจเยน ปจฺจโย.\csromanspacer{} นเหตุ ธมฺโม น นเหตุสฺส ธมฺมสฺส อารมฺมณปจฺจเยน ปจฺจโย.\csromanspacer{}\csromanspacer{}นเหตุ ธมฺโม นเหตุสฺส จ น นเหตุสฺส จ ธมฺมสฺส อารมฺมณปจฺจเยน ปจฺจโย. (3)}

\prose{น นเหตุ ธมฺโม น นเหตุสฺส ธมฺมสฺส อารมฺมณปจฺจเยน ปจฺจโย.\csromanspacer{}\csromanspacer{}น นเหตุ ธมฺโม นเหตุสฺส ธมฺมสฺส อารมฺมณปจฺจเยน ปจฺจโย.\csromanspacer{}\csromanspacer{}น นเหตุ ธมฺโม นเหตุสฺส จ น นเหตุสฺส จ ธมฺมสฺส อารมฺมณปจฺจเยน ปจฺจโย.\csromanspacer{}\csromanspacer{}(3)}

\prose{นเหตุ จ น นเหตุ จ ธมฺมา นเหตุสฺส ธมฺมสฺส อารมฺมณปจฺจเยน ปจฺจโย.\csromanspacer{}\csromanspacer{}นเหตุ จ น นเหตุ จ ธมฺมา น นเหตุสฺส ธมฺมสฺส อารมฺมณปจฺจเยน ปจฺจโย.\csromanspacer{}\csromanspacer{}นเหตุ จ น นเหตุ จ ธมฺมา นเหตุสฺส จ น นเหตุสฺส จ ธมฺมสฺส อารมฺมณปจฺจเยน ปจฺจโย. (3) (สํขิตฺตํ.)}

\proseitem{4}{เหตุยา ตีณิ, อารมฺมเณ นว ฯเปฯ อวิคเต นว. (ยถา กุสลตฺติเก ปญฺหาวารํ, เอวํ วิตฺถาเรตพฺพํ.)}

\csromanheader{2}{2. สเหตุกทุก 1-7. ปฏิจฺจาทิวาร}
\csromanheader{2}{\textbf{ปจฺจยจตุกฺก-เหตุ}}
\proseitem{5}{นสเหตุกํ ธมฺมํ ปฏิจฺจ น สเหตุโก ธมฺโม อุปฺปชฺชติ เหตุปจฺจยา.\csromanspacer{} วิจิกิจฺฉาสหคตํ อุทฺธจฺจสหคตํ โมหํ ปฏิจฺจ จิตฺตสมุฏฺฐานํ รูปํ.\csromanspacer{} เอกํ มหาภูตํ ปฏิจฺจ ตโย มหาภูตา ฯเปฯ เทฺว มหาภูเต ปฏิจฺจ เทฺว มหาภูตา.\csromanspacer{} มหาภูเต ปฏิจฺจ จิตฺตสมุฏฺฐานํ รูปํ, กฏตฺตารูปํ, อุปาทารูปํ.\csromanspacer{}\csromanspacer{}นสเหตุกํ ธมฺมํ ปฏิจฺจ น-อเหตุโก ธมฺโม อุปฺปชฺชติ เหตุปจฺจยา.\csromanspacer{} วิจิกิจฺฉาสหคตํ อุทฺธจฺจสหคตํ โมหํ ปฏิจฺจ สมฺปยุตฺตกา ขนฺธา.\csromanspacer{} ปฏิสนฺธิกฺขเณ วตฺถุํ ปฏิจฺจ สเหตุกา ขนฺธา.\csromanspacer{}\csromanspacer{}นสเหตุกํ ธมฺมํ ปฏิจฺจ นสเหตุโก จ น-อเหตุโก จ ธมฺมา อุปฺปชฺชนฺติ เหตุปจฺจยา. (3)}

\csromanheader{2}{3. เหตุสมฺปยุตฺตทุก}
\prose{\csromanfolio{75}น-อเหตุกํ ธมฺมํ ปฏิจฺจ น-อเหตุโก ธมฺโม อุปฺปชฺชติ เหตุปจฺจยา.\csromanspacer{} ตีณิ.}

\prose{นสเหตุกญฺจ น-อเหตุกญฺจ ธมฺมํ ปฏิจฺจ นสเหตุโก ธมฺโม อุปฺปชฺชติ เหตุปจฺจยา.\csromanspacer{}\csromanspacer{}นสเหตุกญฺจ น-อเหตุกญฺจ ธมฺมํ ปฏิจฺจ น-อเหตุโก ธมฺโม อุปฺปชฺชติ เหตุปจฺจยา.\csromanspacer{}\csromanspacer{}นสเหตุกญฺจ น-อเหตุกญฺจ ธมฺมํ ปฏิจฺจ นสเหตุโก จ น-อเหตุโก จ ธมฺมา อุปฺปชฺชนฺติ เหตุปจฺจยา. (3) (สํขิตฺตํ.)}

\proseitem{6}{เหตุยา นว, อารมฺมเณ ฉ ฯเปฯ อวิคเต นว. (สหชาตวารมฺปิ ฯเปฯ สมฺปยุตฺตวารมฺปิ ปฏิจฺจวารสทิสํ, เอวํ วิตฺถาเรตพฺพํ.)}

\csromanheader{2}{\textbf{เหตุ-อารมฺมณ}}
\proseitem{7}{นสเหตุโก ธมฺโม นสเหตุกสฺส ธมฺมสฺส เหตุปจฺจเยน ปจฺจโย.\csromanspacer{}\csromanspacer{}นสเหตุโก ธมฺโม น-อเหตุกสฺส ธมฺมสฺส เหตุปจฺจเยน ปจฺจโย.\csromanspacer{}\csromanspacer{}นสเหตุโก ธมฺโม นสเหตุกสฺส จ น-อเหตุกสฺส จ ธมฺมสฺส เหตุปจฺจเยน ปจฺจโย. (3)}

\prose{น-อเหตุโก ธมฺโม น-อเหตุกสฺส ธมฺมสฺส เหตุปจฺจเยน ปจฺจโย.\csromanspacer{} ตีณิ.}

\prose{นสเหตุโก ธมฺโม นสเหตุกสฺส ธมฺมสฺส อารมฺมณปจฺจเยน ปจฺจโย. (สํขิตฺตํ.)}

\proseitem{8}{เหตุยา ฉ, อารมฺมเณ นว ฯเปฯ อวิคเต นว.}

\prose{(ยถา กุสลตฺติเก ปญฺหาวารํ, เอวํ วิตฺถาเรตพฺพํ.)}

\csromanheader{2}{3. เหตุสมฺปยุตฺตทุก 1. ปฏิจฺจวาร}
\proseitem{9}{นเหตุสมฺปยุตฺตํ ธมฺมํ ปฏิจฺจ นเหตุสมฺปยุตฺโต ธมฺโม อุปฺปชฺชติ เหตุปจฺจยา.\csromanspacer{}\csromanspacer{}นเหตุสมฺปยุตฺตํ ธมฺมํ ปฏิจฺจ นเหตุวิปฺปยุตฺโต ธมฺโม อุปฺปชฺชติ เหตุปจฺจยา.\csromanspacer{}\csromanspacer{}นเหตุสมฺปยุตฺตํ ธมฺมํ ปฏิจฺจ นเหตุสมฺปยุตฺโต จ นเหตุวิปฺปยุตฺโต จ ธมฺมา อุปฺปชฺชนฺติ เหตุปจฺจยา. (3)}

\prose{นเหตุวิปฺปยุตฺตํ ธมฺมํ ปฏิจฺจ นเหตุวิปฺปยุตฺโต ธมฺโม อุปฺปชฺชติ เหตุปจฺจยา.\csromanspacer{}\csromanspacer{}นเหตุวิปฺปยุตฺตํ ธมฺมํ ปฏิจฺจ นเหตุสมฺปยุตฺโต ธมฺโม อุปฺปชฺชติ}

\vspace{\tipitakaparskip}
\csromancenter{ธมฺมปจฺจนีย ทุกปฏฺฐานปาฬิ เหตุปจฺจยา.\csromanspacer{}\csromanspacer{}นเหตุวิปฺปยุตฺตํ ธมฺมํ ปฏิจฺจ นเหตุสมฺปยุตฺโต จ นเหตุวิปฺปยุตฺโต จ ธมฺมา อุปฺปชฺชนฺติ เหตุปจฺจยา. (3)}
\prose{\csromanfolio{76}นเหตุสมฺปยุตฺตญฺจ นเหตุวิปฺปยุตฺตญฺจ ธมฺมํ ปฏิจฺจ นเหตุสมฺปยุตฺโต ธมฺโม อุปฺปชฺชติ เหตุปจฺจยา.\csromanspacer{} ตีณิ. (สํขิตฺตํ.)}

\proseitem{10}{เหตุยา นว, อารมฺมเณ ฉ ฯเปฯ อวิคเต นว. (สหชาตวาเรปิ ฯเปฯ ปญฺหาวาเรปิ สพฺพตฺถ วิตฺถาเรตพฺพํ.)}

\csromanheader{2}{4. เหตุสเหตุกทุก 1. ปฏิจฺจวาร}
\proseitem{11}{นเหตุญฺเจว น-อเหตุกญฺจ ธมฺมํ ปฏิจฺจ นเหตุ เจว น-อเหตุโก จ ธมฺโม อุปฺปชฺชติ เหตุปจฺจยา.\csromanspacer{}\csromanspacer{}นเหตุญฺเจว น-อเหตุกญฺจ ธมฺมํ ปฏิจฺจ น-อเหตุโก เจว น นเหตุ จ ธมฺโม อุปฺปชฺชติ เหตุปจฺจยา.\csromanspacer{}\csromanspacer{}นเหตุญฺเจว น-อเหตุกญฺจ ธมฺมํ ปฏิจฺจ นเหตุ เจว น-อเหตุโก จ น-อเหตุโก เจว น นเหตุ จ ธมฺมา อุปฺปชฺชนฺติ เหตุปจฺจยา. (3)}

\prose{น-อเหตุกญฺเจว น นเหตุญฺจ ธมฺมํ ปฏิจฺจ น-อเหตุโก เจว น นเหตุ จ ธมฺโม อุปฺปชฺชติ เหตุปจฺจยา.\csromanspacer{}\csromanspacer{}น-อเหตุกญฺเจว น นเหตุญฺจ ธมฺมํ ปฏิจฺจ นเหตุ เจว น-อเหตุโก จ ธมฺโม อุปฺปชฺชติ เหตุปจฺจยา.\csromanspacer{}\csromanspacer{}น-อเหตุกญฺเจว น นเหตุญฺจ ธมฺมํ ปฏิจฺจ นเหตุ เจว น-อเหตุโก จ น-อเหตุโก เจว น นเหตุ จ ธมฺมา อุปฺปชฺชนฺติ เหตุปจฺจยา. (3)}

\prose{นเหตุญฺเจว น-อเหตุกญฺจ น-อเหตุกญฺเจว น นเหตุญฺจ ธมฺมํ ปฏิจฺจ นเหตุ เจว น-อเหตุโก จ ธมฺโม อุปฺปชฺชติ เหตุปจฺจยา.\csromanspacer{} ตีณิ.}

\prose{เหตุยา นว, อารมฺมเณ นว ฯเปฯ อวิคเต นว. (สหชาตวาเรปิ ฯเปฯ ปญฺหาวาเรปิ สพฺพตฺถ วิตฺถาโร.)}

\csromanheader{2}{5. เหตุเหตุสมฺปยุตฺตทุก 1. ปฏิจฺจวาร}
\proseitem{12}{นเหตุญฺเจว นเหตุวิปฺปยุตฺตญฺจ ธมฺมํ ปฏิจฺจ นเหตุ เจว นเหตุวิปฺปยุตฺโต จ ธมฺโม อุปฺปชฺชติ เหตุปจฺจยา.\csromanspacer{}\csromanspacer{}นเหตุญฺเจว นเหตุวิปฺปยุตฺตญฺจ ธมฺมํ ปฏิจฺจ นเหตุวิปฺปยุตฺโต เจว น นเหตุ จ ธมฺโม อุปฺปชฺชติ}

\vspace{\tipitakaparskip}
\csromancenter{นเหตุสเหตุกทุก เหตุปจฺจยา.\csromanspacer{}\csromanspacer{}นเหตุญฺเจว นเหตุวิปฺปยุตฺตญฺจ ธมฺมํ ปฏิจฺจ นเหตุ เจว นเหตุวิปฺปยุตฺโต จ นเหตุวิปฺปยุตฺโต เจว น นเหตุ จ ธมฺมา อุปฺปชฺชนฺติ เหตุปจฺจยา. (3)}
\prose{\csromanfolio{77}นเหตุวิปฺปยุตฺตญฺเจว น นเหตุญฺจ ธมฺมํ ปฏิจฺจ นเหตุวิปฺปยุตฺโต เจว น นเหตุ จ ธมฺโม อุปฺปชฺชติ เหตุปจฺจยา.\csromanspacer{}\csromanspacer{}นเหตุวิปฺปยุตฺตญฺเจว น นเหตุญฺจ ธมฺมํ ปฏิจฺจ นเหตุ เจว นเหตุวิปฺปยุตฺโต จ ธมฺโม อุปฺปชฺชติ เหตุปจฺจยา.\csromanspacer{}\csromanspacer{}นเหตุวิปฺปยุตฺตญฺเจว น นเหตุญฺจ ธมฺมํ ปฏิจฺจ นเหตุ เจว นเหตุวิปฺปยุตฺโต จ นเหตุวิปฺปยุตฺโต เจว น นเหตุ จ ธมฺมา อุปฺปชฺชนฺติ เหตุปจฺจยา. (3) นเหตุญฺเจว นเหตุวิปฺปยุตฺตญฺจ นเหตุวิปฺปยุตฺตญฺเจว น นเหตุญฺจ ธมฺมํ ปฏิจฺจ นเหตุ เจว นเหตุวิปฺปยุตฺโต จ ธมฺโม อุปฺปชฺชติ เหตุปจฺจยา.\csromanspacer{} ตีณิ. (สํขิตฺตํ.)}

\csromanheader{2}{เหตุยา นว ฯเปฯ อวิคเต นว. (สพฺพตฺถ วิตฺถาโร.)}
\csromanheader{2}{6. นเหตุสเหตุกทุก 1-7. ปฏิจฺจาทิวาร}
\proseitem{13}{นเหตุํ นสเหตุกํ ธมฺมํ ปฏิจฺจ นเหตุ นสเหตุโก ธมฺโม อุปฺปชฺชติ เหตุปจฺจยา.\csromanspacer{} เอกํ มหาภูตํ ปฏิจฺจ ตโย มหาภูตา ฯเปฯ.\csromanspacer{} นเหตุํ นสเหตุกํ ธมฺมํ ปฏิจฺจ นเหตุ น-อเหตุโก ธมฺโม อุปฺปชฺชติ เหตุปจฺจยา.\csromanspacer{} ปฏิสนฺธิกฺขเณ วตฺถุํ ปฏิจฺจ นเหตุสเหตุกา ขนฺธา.\csromanspacer{}\csromanspacer{}นเหตุํ นสเหตุกํ ธมฺมํ ปฏิจฺจ นเหตุ นสเหตุโก จ นเหตุ น-อเหตุโก จ ธมฺมา อุปฺปชฺชนฺติ เหตุปจฺจยา. (3)}

\prose{นเหตุํ น-อเหตุกํ ธมฺมํ ปฏิจฺจ นเหตุ น-อเหตุโก ธมฺโม อุปฺปชฺชติ เหตุปจฺจยา.\csromanspacer{}\csromanspacer{}นเหตุํ น-อเหตุกํ ธมฺมํ ปฏิจฺจ นเหตุ นสเหตุโก ธมฺโม อุปฺปชฺชติ เหตุปจฺจยา.\csromanspacer{}\csromanspacer{}นเหตุํ น-อเหตุกํ ธมฺมํ ปฏิจฺจ นเหตุ นสเหตุโก จ นเหตุ น-อเหตุโก จ ธมฺมา อุปฺปชฺชนฺติ เหตุปจฺจยา. (3)}

\prose{นเหตุํ นสเหตุกญฺจ นเหตุํ น-อเหตุกญฺจ ธมฺมํ ปฏิจฺจ นเหตุ นสเหตุโก ธมฺโม อุปฺปชฺชติ เหตุปจฺจยา.\csromanspacer{}\csromanspacer{}นเหตุํ นสเหตุกญฺจ นเหตุํ น-อเหตุกญฺจ ธมฺมํ ปฏิจฺจ นเหตุ น-อเหตุโก ธมฺโม อุปฺปชฺชติ เหตุปจฺจยา.\csromanspacer{}\csromanspacer{}นเหตุํ นสเหตุกญฺจ นเหตุํ น-อเหตุกญฺจ ธมฺมํ}

\vspace{\tipitakaparskip}
\csromancenter{ธมฺมปจฺจนีย ทุกปฏฺฐานปาฬิ ปฏิจฺจ นเหตุ นสเหตุโก จ นเหตุ น-อเหตุโก จ ธมฺมา อุปฺปชฺชนฺติ เหตุปจฺจยา. (3) (สํขิตฺตํ.)}
\proseitem{14}{\csromanfolio{78}เหตุยา นว, อารมฺมเณ จตฺตาริ ฯเปฯ อวิคเต นว.}

\prose{(สหชาตวารมฺปิ ฯเปฯ สมฺปยุตฺตวารมฺปิ ปฏิจฺจวารสทิสํ.)}

\csromanheader{2}{\textbf{อารมฺมณ}}
\proseitem{15}{นเหตุ นสเหตุโก ธมฺโม นเหตุสฺส นสเหตุกสฺส ธมฺมสฺส อารมฺมณปจฺจเยน ปจฺจโย.\csromanspacer{}\csromanspacer{}นเหตุ นสเหตุโก ธมฺโม นเหตุสฺส น-อเหตุกสฺส ธมฺมสฺส อารมฺมณปจฺจเยน ปจฺจโย. (2)}

\prose{นเหตุ น-อเหตุโก ธมฺโม นเหตุสฺส น-อเหตุกสฺส ธมฺมสฺส อารมฺมณปจฺจเยน ปจฺจโย.\csromanspacer{}\csromanspacer{}นเหตุ น-อเหตุโก ธมฺโม นเหตุสฺส นสเหตุกสฺส ธมฺมสฺส อารมฺมณปจฺจเยน ปจฺจโย. (2) (สํขิตฺตํ.)}

\nitthitam{16. อารมฺมเณ จตฺตาริ ฯเปฯ อวิคเต สตฺต. เหตุโคจฺฉกํ นิฏฺฐิตํ.}
\csromanheader{2}{7. สปฺปจฺจยทุก 1-7. ปฏิจฺจาทิวาร}
\proseitem{17}{น-อปฺปจฺจยํ ธมฺมํ ปฏิจฺจ น-อปฺปจฺจโย ธมฺโม อุปฺปชฺชติ เหตุปจฺจยา.}

\csromanheader{2}{เหตุยา เอกํ. (สพฺพตฺถ วิตฺถาโร.)}
\proseitem{18}{น-อปฺปจฺจโย ธมฺโม น-อปฺปจฺจยสฺส ธมฺมสฺส เหตุปจฺจเยน ปจฺจโย.}

\prose{นสปฺปจฺจโย ธมฺโม น-อปฺปจฺจยสฺส ธมฺมสฺส อารมฺมณปจฺจเยน ปจฺจโย. (สํขิตฺตํ.)}

\prose{เหตุยา เอกํ, อารมฺมเณ เทฺว. (สพฺพตฺถ วิตฺถาโร.)}

\csromanheader{2}{8. สงฺขตทุก 1. ปฏิจฺจวาร}
\proseitem{19}{น-อสงฺขตํ ธมฺมํ ปฏิจฺจ น-อสงฺขโต ธมฺโม อุปฺปชฺชติ เหตุปจฺจยา.\csromanspacer{} เหตุยา เอกํ. (สปฺปจฺจยทุกสทิสํ.)}

\csromanheader{2}{10. สปฺปฏิฆทุก \textbf{9. สนิทสฺสนทุก 1. ปฏิจฺจวาร}}
\proseitem{20}{\csromanfolio{79}นสนิทสฺสนํ ธมฺมํ ปฏิจฺจ นสนิทสฺสโน ธมฺโม อุปฺปชฺชติ เหตุปจฺจยา.\csromanspacer{}\csromanspacer{}นสนิทสฺสนํ ธมฺมํ ปฏิจฺจ น-อนิทสฺสโน ธมฺโม อุปฺปชฺชติ เหตุปจฺจยา.\csromanspacer{}\csromanspacer{}นสนิทสฺสนํ ธมฺมํ ปฏิจฺจ นสนิทสฺสโน จ น-อนิทสฺสโน จ ธมฺมา อุปฺปชฺชนฺติ เหตุปจฺจยา. (3)}

\prose{เหตุยา ตีณิ.}

\csromanheader{2}{10. สปฺปฏิฆทุก 1-7. ปฏิจฺจาทิวาร}
\proseitem{21}{นสปฺปฏิฆํ ธมฺมํ ปฏิจฺจ นสปฺปฏิโฆ ธมฺโม อุปฺปชฺชติ เหตุปจฺจยา.\csromanspacer{}\csromanspacer{}นสปฺปฏิฆํ ธมฺมํ ปฏิจฺจ น-อปฺปฏิโฆ ธมฺโม อุปฺปชฺชติ เหตุปจฺจยา.\csromanspacer{}\csromanspacer{}นสปฺปฏิฆํ ธมฺมํ ปฏิจฺจ นสปฺปฏิโฆ จ น-อปฺปฏิโฆ จ ธมฺมา อุปฺปชฺชนฺติ เหตุปจฺจยา. (3)}

\prose{น-อปฺปฏิฆํ ธมฺมํ ปฏิจฺจ น-อปฺปฏิโฆ ธมฺโม อุปฺปชฺชติ เหตุปจฺจยา.\csromanspacer{} ตีณิ.}

\prose{นสปฺปฏิฆญฺจ น-อปฺปฏิฆญฺจ ธมฺมํ ปฏิจฺจ นสปฺปฏิโฆ ธมฺโม อุปฺปชฺชติ เหตุปจฺจยา.\csromanspacer{} ตีณิ. (สํขิตฺตํ.)}

\prose{เหตุยา นว, อารมฺมเณ เอกํ. (สหชาตวารมฺปิ ฯเปฯ สมฺปยุตฺตวารมฺปิ ปฏิจฺจวารสทิสํ.)}

\proseitem{22}{นสปฺปฏิโฆ ธมฺโม นสปฺปฏิฆสฺส ธมฺมสฺส เหตุปจฺจเยน ปจฺจโย.\csromanspacer{}\csromanspacer{}นสปฺปฏิโฆ ธมฺโม น-อปฺปฏิฆสฺส ธมฺมสฺส เหตุปจฺจเยน ปจฺจโย.\csromanspacer{}\csromanspacer{}นสปฺปฏิโฆ ธมฺโม นสปฺปฏิฆสฺส จ น-อปฺปฏิฆสฺส จ ธมฺมสฺส เหตุปจฺจเยน ปจฺจโย. (3)}

\prose{นสปฺปฏิโฆ ธมฺโม นสปฺปฏิฆสฺส ธมฺมสฺส อารมฺมณปจฺจเยน ปจฺจโย.\csromanspacer{}\csromanspacer{}น-อปฺปฏิโฆ ธมฺโม นสปฺปฏิฆสฺส ธมฺมสฺส อารมฺมณปจฺจเยน ปจฺจโย. (2) (สํขิตฺตํ.)}

\proseitem{23}{เหตุยา ตีณิ, อารมฺมเณ เทฺว, อธิปติยา จตฺตาริ ฯเปฯ อวิคเต นว.}

\csromanheader{2}{ธมฺมปจฺจนีย ทุกปฏฺฐานปาฬิ \textbf{11. รูปีทุก 1. ปฏิจฺจวาร}}
\proseitem{24}{\csromanfolio{80}นรูปิํ ธมฺมํ ปฏิจฺจ นรูปี ธมฺโม อุปฺปชฺชติ เหตุปจฺจยา.\csromanspacer{}\csromanspacer{}นรูปิํ ธมฺมํ ปฏิจฺจ น-อรูปี ธมฺโม อุปฺปชฺชติ เหตุปจฺจยา.\csromanspacer{}\csromanspacer{}นรูปิํ ธมฺมํ ปฏิจฺจ นรูปี จ น-อรูปี จ ธมฺมา อุปฺปชฺชนฺติ เหตุปจฺจยา. (3)}

\prose{น-อรูปิํ ธมฺมํ ปฏิจฺจ น-อรูปี ธมฺโม อุปฺปชฺชติ เหตุปจฺจยา.\csromanspacer{}\csromanspacer{}น-อรูปิํ ธมฺมํ ปฏิจฺจ นรูปี ธมฺโม อุปฺปชฺชติ เหตุปจฺจยา.\csromanspacer{}\csromanspacer{}น-อรูปิํ ธมฺมํ ปฏิจฺจ นรูปี จ น-อรูปี จ ธมฺมา อุปฺปชฺชนฺติ เหตุปจฺจยา. (3)}

\prose{นรูปิญฺจ น-อรูปิญฺจ ธมฺมํ ปฏิจฺจ นรูปี ธมฺโม อุปฺปชฺชติ เหตุปจฺจยา.\csromanspacer{}\csromanspacer{}นรูปิญฺจ น-อรูปิญฺจ ธมฺมํ ปฏิจฺจ น-อรูปี ธมฺโม อุปฺปชฺชติ เหตุปจฺจยา.\csromanspacer{}\csromanspacer{}นรูปิญฺจ น-อรูปิญฺจ ธมฺมํ ปฏิจฺจ นรูปี จ น-อรูปี จ ธมฺมา อุปฺปชฺชนฺติ เหตุปจฺจยา. (3) (สํขิตฺตํ.)}

\proseitem{25}{เหตุยา นว, อารมฺมเณ ตีณิ ฯเปฯ อวิคเต นว. (สหชาตวาเรปิ ฯเปฯ ปญฺหาวาเรปิ สพฺพตฺถ วิตฺถาโร.)}

\csromanheader{2}{12. โลกิยทุก 1. ปฏิจฺจวาร}
\proseitem{26}{นโลกิยํ ธมฺมํ ปฏิจฺจ นโลกิโย ธมฺโม อุปฺปชฺชติ เหตุปจฺจยา.\csromanspacer{}\csromanspacer{}นโลกิยํ ธมฺมํ ปฏิจฺจ นโลกุตฺตโร ธมฺโม อุปฺปชฺชติ เหตุปจฺจยา.\csromanspacer{}\csromanspacer{}นโลกิยํ ธมฺมํ ปฏิจฺจ นโลกิโย จ นโลกุตฺตโร จ ธมฺมา อุปฺปชฺชนฺติ เหตุปจฺจยา. (3)}

\prose{นโลกุตฺตรํ ธมฺมํ ปฏิจฺจ นโลกุตฺตโร ธมฺโม อุปฺปชฺชติ เหตุปจฺจยา.}

\prose{นโลกิยญฺจ นโลกุตฺตรญฺจ ธมฺมํ ปฏิจฺจ นโลกุตฺตโร ธมฺโม อุปฺปชฺชติ เหตุปจฺจยา. (2) (สํขิตฺตํ.)}

\prose{เหตุยา ปญฺจ, อารมฺมเณ เทฺว ฯเปฯ อวิคเต ปญฺจ. (สพฺพตฺถ วิตฺถาโร.)}

\csromanheader{2}{13. เกนจิวิญฺเญยฺยทุก 1. ปฏิจฺจวาร}
\proseitem{27}{นเกนจิ วิญฺเญยฺยํ ธมฺมํ ปฏิจฺจ นเกนจิ วิญฺเญยฺโย ธมฺโม อุปฺปชฺชติ เหตุปจฺจยา.\csromanspacer{}\csromanspacer{}นเกนจิ วิญฺเญยฺยํ ธมฺมํ ปฏิจฺจ นนเกนจิ วิญฺเญยฺโย ธมฺโม อุปฺปชฺชติ เหตุปจฺจยา.\csromanspacer{}\csromanspacer{}นเกนจิ วิญฺเญยฺยํ ธมฺมํ ปฏิจฺจ นเกนจิ วิญฺเญยฺโย จ นนเกนจิ วิญฺเญยฺโย จ ธมฺมา อุปฺปชฺชนฺติ เหตุปจฺจยา. (3)}

\csromanheader{2}{15. สาสวทุก}
\prose{\csromanfolio{81}นนเกนจิ วิญฺเญยฺยํ ธมฺมํ ปฏิจฺจ นนเกนจิ วิญฺเญยฺโย ธมฺโม อุปฺปชฺชติ เหตุปจฺจยา.\csromanspacer{} ตีณิ.}

\prose{นเกนจิ วิญฺเญยฺยญฺจ นนเกนจิ วิญฺเญยฺยญฺจ ธมฺมํ ปฏิจฺจ นเกนจิ วิญฺเญยฺโย ธมฺโม อุปฺปชฺชติ เหตุปจฺจยา.\csromanspacer{} ตีณิ. (สํขิตฺตํ.)}

\csromanheader{2}{เหตุยา นว ฯเปฯ อวิคเต นว. (สพฺพตฺถ วิตฺถาโร.)}
\nitthitam{จูฬนฺตรทุกํ นิฏฺฐิตํ.}
\csromanheader{2}{14. อาสวทุก 1. ปฏิจฺจวาร}
\proseitem{28}{โน-อาสวํ ธมฺมํ ปฏิจฺจ โน-อาสโว ธมฺโม อุปฺปชฺชติ เหตุปจฺจยา.\csromanspacer{}\csromanspacer{}โน-อาสวํ ธมฺมํ ปฏิจฺจ นโน-อาสโว ธมฺโม อุปฺปชฺชติ เหตุปจฺจยา.\csromanspacer{}\csromanspacer{}โน-อาสวํ ธมฺมํ ปฏิจฺจ โน-อาสโว จ นโน-อาสโว จ ธมฺมา อุปฺปชฺชนฺติ เหตุปจฺจยา. (3)}

\prose{นโน-อาสวํ ธมฺมํ ปฏิจฺจ นโน-อาสโว ธมฺโม อุปฺปชฺชติ เหตุปจฺจยา.\csromanspacer{} ตีณิ.}

\prose{โน-อาสวญฺจ นโน-อาสวญฺจ ธมฺมํ ปฏิจฺจ โน-อาสโว ธมฺโม อุปฺปชฺชติ เหตุปจฺจยา.\csromanspacer{}\csromanspacer{}โน-อาสวญฺจ นโน-อาสวญฺจ ธมฺมํ ปฏิจฺจ นโน-อาสโว ธมฺโม อุปฺปชฺชติ เหตุปจฺจยา.\csromanspacer{} โน-อาสวญฺจ นโน-อาสวญฺจ ธมฺมํ ปฏิจฺจ โน-อาสโว จ นโน-อาสโว จ ธมฺมา อุปฺปชฺชนฺติ เหตุปจฺจยา. (3) (สํขิตฺตํ.)}

\prose{เหตุยา นว, อารมฺมเณ นว ฯเปฯ อวิคเต นว. (สพฺพตฺถ วิตฺถาโร.)}

\csromanheader{2}{15. สาสวทุก 1. ปฏิจฺจวาร}
\proseitem{29}{นสาสวํ ธมฺมํ ปฏิจฺจ นสาสโว ธมฺโม อุปฺปชฺชติ เหตุปจฺจยา.\csromanspacer{}\csromanspacer{}นสาสวํ ธมฺมํ ปฏิจฺจ น-อนาสโว ธมฺโม อุปฺปชฺชติ เหตุปจฺจยา.\csromanspacer{}\csromanspacer{}นสาสวํ ธมฺมํ ปฏิจฺจ นสาสโว จ น-อนาสโว จ ธมฺมา อุปฺปชฺชนฺติ เหตุปจฺจยา. (3)}

\csromanheader{2}{น-อนาสวํ ธมฺมํ ปฏิจฺจ น-อนาสโว ธมฺโม อุปฺปชฺชติ}
\gambhira{ธมฺมปจฺจนีย ทุกปฏฺฐานปาฬิ}
\prose{\csromanfolio{82}นสาสวญฺจ น-อนาสวญฺจ ธมฺมํ ปฏิจฺจ น-อนาสโว ธมฺโม อุปฺปชฺชติ เหตุปจฺจยา. (1) (สํขิตฺตํ.)}

\prose{เหตุยา ปญฺจ, อารมฺมเณ เทฺว ฯเปฯ อวิคเต ปญฺจ. (สพฺพตฺถ วิตฺถาโร.)}

\csromanheader{2}{16. อาสวสมฺปยุตฺตทุก 1. ปฏิจฺจวาร}
\proseitem{30}{น-อาสวสมฺปยุตฺตํ ธมฺมํ ปฏิจฺจ น-อาสวสมฺปยุตฺโต ธมฺโม อุปฺปชฺชติ เหตุปจฺจยา.\csromanspacer{}\csromanspacer{}น-อาสวสมฺปยุตฺตํ ธมฺมํ ปฏิจฺจ น-อาสววิปฺปยุตฺโต ธมฺโม อุปฺปชฺชติ เหตุปจฺจยา.\csromanspacer{}\csromanspacer{}น-อาสวสมฺปยุตฺตํ ธมฺมํ ปฏิจฺจ น-อาสวสมฺปยุตฺโต จ น-อาสววิปฺปยุตฺโต จ ธมฺมา อุปฺปชฺชนฺติ เหตุปจฺจยา. (3)}

\prose{น-อาสววิปฺปยุตฺตํ ธมฺมํ ปฏิจฺจ น-อาสววิปฺปยุตฺโต ธมฺโม อุปฺปชฺชติ เหตุปจฺจยา.\csromanspacer{}\csromanspacer{}น-อาสววิปฺปยุตฺตํ ธมฺมํ ปฏิจฺจ น-อาสวสมฺปยุตฺโต ธมฺโม อุปฺปชฺชติ เหตุปจฺจยา.\csromanspacer{}\csromanspacer{}น-อาสววิปฺปยุตฺตํ ธมฺมํ ปฏิจฺจ น-อาสวสมฺปยุตฺโต จ น-อาสววิปฺปยุตฺโต จ ธมฺมา อุปฺปชฺชนฺติ เหตุปจฺจยา. (3)}

\prose{น-อาสวสมฺปยุตฺตญฺจ น-อาสววิปฺปยุตฺตญฺจ ธมฺมํ ปฏิจฺจ น-อาสวสมฺปยุตฺโต ธมฺโม อุปฺปชฺชติ เหตุปจฺจยา.\csromanspacer{}\csromanspacer{}น-อาสวสมฺปยุตฺตญฺจ น-อาสววิปฺปยุตฺตญฺจ ธมฺมํ ปฏิจฺจ น-อาสววิปฺปยุตฺโต ธมฺโม อุปฺปชฺชติ เหตุปจฺจยา.\csromanspacer{}\csromanspacer{}น-อาสวสมฺปยุตฺตญฺจ น-อาสววิปฺปยุตฺตญฺจ ธมฺมํ ปฏิจฺจ น-อาสวสมฺปยุตฺโต จ น-อาสววิปฺปยุตฺโต จ ธมฺมา อุปฺปชฺชนฺติ เหตุปจฺจยา. (3) (สํขิตฺตํ.)}

\prose{เหตุยา นว, อารมฺมเณ ฉ ฯเปฯ อวิคเต นว. (สพฺพตฺถ วิตฺถาโร.)}

\csromanheader{2}{17. อาสวสาสวทุก 1. ปฏิจฺจวาร}
\proseitem{31}{น-อาสวญฺเจว น-อนาสวญฺจ ธมฺมํ ปฏิจฺจ น-อาสโว เจว น-อนาสโว จ ธมฺโม อุปฺปชฺชติ เหตุปจฺจยา.\csromanspacer{}\csromanspacer{}น-อาสวญฺเจว น-อนาสวญฺจ ธมฺมํ ปฏิจฺจ น-อนาสโว เจว นโน จ อาสโว ธมฺโม อุปฺปชฺชติ เหตุปจฺจยา.\csromanspacer{}\csromanspacer{}น-อาสวญฺเจว น-อนาสวญฺจ ธมฺมํ ปฏิจฺจ โน-อาสโว เจว น-อนาสโว จ น-อนาสโว เจว นโน จ อาสโว จ ธมฺมา อุปฺปชฺชนฺติ เหตุปจฺจยา. (3)}

\prose{น-อนาสวญฺเจว นโน จ อาสวํ ธมฺมํ ปฏิจฺจ น-อนาสโว เจว นโน จ อาสโว ธมฺโม อุปฺปชฺชติ เหตุปจฺจยา.\csromanspacer{}\csromanspacer{}น-อนาสวญฺเจว นโน จ อาสวํ ธมฺมํ ปฏิจฺจ น-อาสโว เจว น-อนาสโว จ ธมฺโม อุปฺปชฺชติ เหตุปจฺจยา.\csromanspacer{}\csromanspacer{}น-อนาสวญฺเจว นโน จ อาสวํ ธมฺมํ ปฏิจฺจ น-อาสโว เจว}

\vspace{\tipitakaparskip}
\csromancenter{อาสว-อาสวสมฺปยุตฺตทุก น-อนาสโว จ น-อนาสโว เจว นโน จ อาสโว จ ธมฺมา อุปฺปชฺชนฺติ เหตุปจฺจยา. (3)}
\prose{\csromanfolio{83}น-อาสวญฺเจว น น-อนาสวญฺจ น-อนาสวญฺเจว นโน-อาสวญฺจ ธมฺมํ ปฏิจฺจ น-อาสโว เจว น-อนาสโว จ ธมฺโม อุปฺปชฺชติ เหตุปจฺจยา.\csromanspacer{}\csromanspacer{}น-อาสวญฺเจว น-อนาสวญฺจ น-อนาสวญฺเจว นโน จ อาสวํ ธมฺมํ ปฏิจฺจ น-อนาสโว เจว นโน จ อาสโว ธมฺโม อุปฺปชฺชติ เหตุปจฺจยา.\csromanspacer{}\csromanspacer{}น-อาสวญฺเจว น-อนาสวญฺจ น-อนาสวญฺเจว นโน-อาสวญฺจ ธมฺมํ ปฏิจฺจ น-อาสโว เจว น-อนาสโว จ น-อนาสโว เจว นโน จ อาสโว จ ธมฺมา อุปฺปชฺชนฺติ เหตุปจฺจยา. (3) (สํขิตฺตํ.)}

\prose{เหตุยา นว, อารมฺมเณ นว ฯเปฯ อวิคเต นว. (สพฺพตฺถ วิตฺถาโร.)}

\csromanheader{2}{18. อาสว-อาสวสมฺปยุตฺตทุก 1. ปฏิจฺจวาร}
\proseitem{32}{น-อาสวญฺเจว น-อาสววิปฺปยุตฺตญฺจ ธมฺมํ ปฏิจฺจ น-อาสโว เจว น-อาสววิปฺปยุตฺโต จ ธมฺโม อุปฺปชฺชติ เหตุปจฺจยา.\csromanspacer{}\csromanspacer{}น-อาสวญฺเจว น-อาสววิปฺปยุตฺตญฺจ ธมฺมํ ปฏิจฺจ น-อาสววิปฺปยุตฺโต เจว นโน จ อาสโว ธมฺโม อุปฺปชฺชติ เหตุปจฺจยา.\csromanspacer{}\csromanspacer{}น-อาสวญฺเจว น-อาสววิปฺปยุตฺตญฺจ ธมฺมํ ปฏิจฺจ น-อาสโว เจว น-อาสววิปฺปยุตฺโต จ น-อาสววิปฺปยุตฺโต เจว นโน-อาสโว จ ธมฺมา อุปฺปชฺชนฺติ เหตุปจฺจยา. (3)}

\prose{น-อาสววิปฺปยุตฺตญฺเจว นโน-อาสวญฺจ ธมฺมํ ปฏิจฺจ น-อาสววิปฺปยุตฺโต เจว นโน-อาสโว จ ธมฺโม อุปฺปชฺชติ เหตุปจฺจยา.\csromanspacer{}\csromanspacer{}น-อาสววิปฺปยุตฺตญฺเจว นโน-อาสวญฺจ ธมฺมํ ปฏิจฺจ น-อาสโว เจว น-อาสววิปฺปยุตฺโต จ ธมฺโม อุปฺปชฺชติ เหตุปจฺจยา.\csromanspacer{}\csromanspacer{}น-อาสววิปฺปยุตฺตญฺเจว นโน จ อาสวํ ธมฺมํ ปฏิจฺจ น-อาสโว เจว น-อาสววิปฺปยุตฺโต จ น-อาสววิปฺปยุตฺโต เจว นโน จ อาสโว จ ธมฺมา อุปฺปชฺชนฺติ เหตุปจฺจยา. (3)}

\prose{น-อาสวญฺเจว น-อาสววิปฺปยุตฺตญฺจ น-อาสววิปฺปยุตฺตญฺเจว นโน-อาสวญฺจ ธมฺมํ ปฏิจฺจ น-อาสโว เจว น-อาสววิปฺปยุตฺโต จ ธมฺโม อุปฺปชฺชติ เหตุปจฺจยา.\csromanspacer{}\csromanspacer{}น-อาสวญฺเจว น-อาสววิปฺปยุตฺตญฺจ น-อาสววิปฺปยุตฺตญฺเจว นโน จ อาสวญฺจ ธมฺมํ ปฏิจฺจ น-อาสววิปฺปยุตฺโต เจว นโน จ อาสโว ธมฺโม อุปฺปชฺชติ เหตุปจฺจยา.\csromanspacer{}\csromanspacer{}น-อาสวญฺเจว น-อาสววิปฺปยุตฺตญฺจ น-อาสววิปฺปยุตฺตญฺเจว นโน จ อาสวญฺจ ธมฺมํ ปฏิจฺจ}

\vspace{\tipitakaparskip}
\csromancenter{ธมฺมปจฺจนีย ทุกปฏฺฐานปาฬิ น-อาสโว เจว น-อาสววิปฺปยุตฺโต จ น-อาสววิปฺปยุตฺโต เจว นโน จ อาสโว จ ธมฺมา อุปฺปชฺชนฺติ เหตุปจฺจยา. (3) (สํขิตฺตํ.)}
\prose{\csromanfolio{84}เหตุยา นว, อารมฺมเณ นว ฯเปฯ อวิคเต นว. (สพฺพตฺถ วิตฺถาโร.)}

\csromanheader{2}{19. อาสววิปฺปยุตฺตสาสวทุก 1. ปฏิจฺจวาร}
\proseitem{33}{อาสววิปฺปยุตฺตํ นสาสวํ ธมฺมํ ปฏิจฺจ อาสววิปฺปยุตฺโต นสาสโว ธมฺโม อุปฺปชฺชติ เหตุปจฺจยา.\csromanspacer{}\csromanspacer{}อาสววิปฺปยุตฺตํ นสาสวํ ธมฺมํ ปฏิจฺจ อาสววิปฺปยุตฺโต น-อนาสโว ธมฺโม อุปฺปชฺชติ เหตุปจฺจยา.\csromanspacer{}\csromanspacer{}อาสววิปฺปยุตฺตํ นสาสวํ ธมฺมํ ปฏิจฺจ อาสววิปฺปยุตฺโต นสาสโว จ อาสววิปฺปยุตฺโต น-อนาสโว จ ธมฺมา อุปฺปชฺชนฺติ เหตุปจฺจยา. (3)}

\prose{อาสววิปฺปยุตฺตํ น-อนาสวํ ธมฺมํ ปฏิจฺจ อาสววิปฺปยุตฺโต น-อนาสโว ธมฺโม อุปฺปชฺชติ เหตุปจฺจยา. (1)}

\prose{อาสววิปฺปยุตฺตํ นสาสวญฺจ อาสววิปฺปยุตฺตํ น-อนาสวญฺจ ธมฺมํ ปฏิจฺจ อาสววิปฺปยุตฺโต น-อนาสโว ธมฺโม อุปฺปชฺชติ เหตุปจฺจยา. (สํขิตฺตํ.)}

\prose{เหตุยา ปญฺจ, อารมฺมเณ เทฺว ฯเปฯ อวิคเต ปญฺจ. (สพฺพตฺถ วิตฺถาโร.)}

\nitthitam{อาสวโคจฺฉกํ นิฏฺฐิตํ.}
\csromanheader{2}{20. สญฺโญชนทุก}
\proseitem{34}{โนสญฺโญชนํ ธมฺมํ ปฏิจฺจ โนสญฺโญชโน ธมฺโม อุปฺปชฺชติ เหตุปจฺจยา.\csromanspacer{}\csromanspacer{}โนสญฺโญชนํ ธมฺมํ ปฏิจฺจ นโนสญฺโญชโน ธมฺโม อุปฺปชฺชติ เหตุปจฺจยา.\csromanspacer{}\csromanspacer{}โนสญฺโญชนํ ธมฺมํ ปฏิจฺจ โนสญฺโญชโน จ นโนสญฺโญชโน จ ธมฺมา อุปฺปชฺชนฺติ เหตุปจฺจยา. (3)}

\prose{นโนสญฺโญชนํ ธมฺมํ ปฏิจฺจ นโนสญฺโญชโน ธมฺโม อุปฺปชฺชติ เหตุปจฺจยา.\csromanspacer{}\csromanspacer{}นโนสญฺโญชนํ ธมฺมํ ปฏิจฺจ โนสญฺโญชโน ธมฺโม อุปฺปชฺชติ เหตุปจฺจยา.\csromanspacer{}\csromanspacer{}นโนสญฺโญชนํ ธมฺมํ ปฏิจฺจ โนสญฺโญชโน จ นโนสญฺโญชโน จ ธมฺมา อุปฺปชฺชนฺติ เหตุปจฺจยา. (3)}

\prose{โนสญฺโญชนญฺจ นโนสญฺโญชนญฺจ ธมฺมํ ปฏิจฺจ โนสญฺโญชโน ธมฺโม อุปฺปชฺชติ เหตุปจฺจยา.\csromanspacer{}\csromanspacer{}โนสญฺโญชนญฺจ นโนสญฺโญชนญฺจ}

\vspace{\tipitakaparskip}
\csromancenter{สญฺโญชนสมฺปยุตฺตทุก ธมฺมํ ปฏิจฺจ นโนสญฺโญชโน ธมฺโม อุปฺปชฺชติ เหตุปจฺจยา.\csromanspacer{}\csromanspacer{}โนสญฺโญชนญฺจ นโนสญฺโญชนญฺจ ธมฺมํ ปฏิจฺจ โนสญฺโญชโน จ นโนสญฺโญชโน จ ธมฺมา อุปฺปชฺชนฺติ เหตุปจฺจยา. (3) (สํขิตฺตํ.)}
\prose{\csromanfolio{85}เหตุยา นว, อารมฺมเณ นว ฯเปฯ อวิคเต นว. (สพฺพตฺถ วิตฺถาโร.) \textbf{21.}\csromanspacer{}\textbf{ สญฺโญชนิยทุก}}

\proseitem{35}{นสญฺโญชนิยํ ธมฺมํ ปฏิจฺจ นสญฺโญชนิโย ธมฺโม อุปฺปชฺชติ เหตุปจฺจยา.\csromanspacer{}\csromanspacer{}นสญฺโญชนิยํ ธมฺมํ ปฏิจฺจ น-อสญฺโญชนิโย ธมฺโม อุปฺปชฺชติ เหตุปจฺจยา.\csromanspacer{}\csromanspacer{}นสญฺโญชนิยํ ธมฺมํ ปฏิจฺจ นสญฺโญชนิโย จ น-อสญฺโญชนิโย จ ธมฺมา อุปฺปชฺชนฺติ เหตุปจฺจยา. (3)}

\prose{น-อสญฺโญชนิยํ ธมฺมํ ปฏิจฺจ น-อสญฺโญชนิโย ธมฺโม อุปฺปชฺชติ เหตุปจฺจยา. (1)}

\prose{นสญฺโญชนิยญฺจ น-อสญฺโญชนิยญฺจ ธมฺมํ ปฏิจฺจ น-อสญฺโญชนิโย ธมฺโม อุปฺปชฺชติ เหตุปจฺจยา. (1) (สํขิตฺตํ.)}

\prose{เหตุยา ปญฺจ, อารมฺมเณ เทฺว ฯเปฯ อวิคเต ปญฺจ. (สพฺพตฺถ วิตฺถาโร.) \textbf{22.}\csromanspacer{}\textbf{ สญฺโญชนสมฺปยุตฺตทุก}}

\proseitem{36}{นสญฺโญชนสมฺปยุตฺตํ ธมฺมํ ปฏิจฺจ นสญฺโญชนสมฺปยุตฺโต ธมฺโม อุปฺปชฺชติ เหตุปจฺจยา.\csromanspacer{}\csromanspacer{}นสญฺโญชนสมฺปยุตฺตํ ธมฺมํ ปฏิจฺจ นสญฺโญชนวิปฺปยุตฺโต ธมฺโม อุปฺปชฺชติ เหตุปจฺจยา.\csromanspacer{}\csromanspacer{}นสญฺโญชนสมฺปยุตฺตํ ธมฺมํ ปฏิจฺจ นสญฺโญชนสมฺปยุตฺโต จ นสญฺโญชนวิปฺปยุตฺโต จ ธมฺมา อุปฺปชฺชนฺติ เหตุปจฺจยา. (3)}

\prose{นสญฺโญชนวิปฺปยุตฺตํ ธมฺมํ ปฏิจฺจ นสญฺโญชนวิปฺปยุตฺโต ธมฺโม อุปฺปชฺชติ เหตุปจฺจยา.\csromanspacer{}\csromanspacer{}นสญฺโญชนวิปฺปยุตฺตํ ธมฺมํ ปฏิจฺจ นสญฺโญชนสมฺปยุตฺโต ธมฺโม อุปฺปชฺชติ เหตุปจฺจยา.\csromanspacer{}\csromanspacer{}นสญฺโญชนวิปฺปยุตฺตํ ธมฺมํ ปฏิจฺจ นสญฺโญชนสมฺปยุตฺโต จ นสญฺโญชนวิปฺปยุตฺโต จ ธมฺมา อุปฺปชฺชนฺติ เหตุปจฺจยา. (3)}

\prose{นสญฺโญชนสมฺปยุตฺตญฺจ นสญฺโญชนวิปฺปยุตฺตญฺจ ธมฺมํ ปฏิจฺจ นสญฺโญชนสมฺปยุตฺโต ธมฺโม อุปฺปชฺชติ เหตุปจฺจยา.\csromanspacer{}\csromanspacer{}นสญฺโญชนสมฺปยุตฺตญฺจ นสญฺโญชนวิปฺปยุตฺตญฺจ ธมฺมํ ปฏิจฺจ นสญฺโญชนวิปฺปยุตฺโต ธมฺโม อุปฺปชฺชติ}

\vspace{\tipitakaparskip}
\csromancenter{ธมฺมปจฺจนีย ทุกปฏฺฐานปาฬิ เหตุปจฺจยา.\csromanspacer{}\csromanspacer{}นสญฺโญชนสมฺปยุตฺตญฺจ นสญฺโญชนวิปฺปยุตฺตญฺจ ธมฺมํ ปฏิจฺจ นสญฺโญชนสมฺปยุตฺโต จ นสญฺโญชนวิปฺปยุตฺโต จ ธมฺมา อุปฺปชฺชนฺติ เหตุปจฺจยา. (3) (สํขิตฺตํ.)}
\prose{\csromanfolio{86}เหตุยา นว, อารมฺมเณ ฉ ฯเปฯ อวิคเต นว. (สพฺพตฺถ วิตฺถาโร.)}

\csromanheader{2}{23. สญฺโญชนสญฺโญชนิยทุก}
\proseitem{37}{นสญฺโญชนญฺเจว น-อสญฺโญชนิยญฺจ ธมฺมํ ปฏิจฺจ นสญฺโญชโน เจว น-อสญฺโญชนิโย จ ธมฺโม อุปฺปชฺชติ เหตุปจฺจยา.\csromanspacer{}\csromanspacer{}นสญฺโญชนญฺเจว น-อสญฺโญชนิยญฺจ ธมฺมํ ปฏิจฺจ น-อสญฺโญชนิโย เจว นโน-อสญฺโญชโน จ ธมฺโม อุปฺปชฺชติ เหตุปจฺจยา.\csromanspacer{}\csromanspacer{}นสญฺโญชนญฺเจว น-อสญฺโญชนิยญฺจ ธมฺมํ ปฏิจฺจ นสญฺโญชโน เจว น-อสญฺโญชนิโย จ น-อสญฺโญชนิโย เจว นโน จ สญฺโญชโน จ ธมฺมา อุปฺปชฺชนฺติ เหตุปจฺจยา. (3)}

\prose{น-อสญฺโญชนิยญฺเจว นโน จ สญฺโญชนํ ธมฺมํ ปฏิจฺจ น-อสญฺโญชนิโย เจว นโน จ สญฺโญชโน ธมฺโม อุปฺปชฺชติ เหตุปจฺจยา.\csromanspacer{} ตีณิ.}

\prose{นสญฺโญชนญฺเจว น-อสญฺโญชนิยญฺจ น-อสญฺโญชนิยญฺเจว นโนสญฺโญชนญฺจ ธมฺมํ ปฏิจฺจ นสญฺโญชโน เจว น-อสญฺโญชนิโย จ ธมฺโม อุปฺปชฺชติ เหตุปจฺจยา.\csromanspacer{} ตีณิ. (สํขิตฺตํ.)}

\prose{เหตุยา นว, อารมฺมเณ นว ฯเปฯ อวิคเต นว. (สพฺพตฺถ วิตฺถาโร.)}

\csromanheader{2}{24. สญฺโญชนสญฺโญชนสมฺปยุตฺตทุก}
\proseitem{38}{นสญฺโญชนญฺเจว นสญฺโญชนวิปฺปยุตฺตญฺจ ธมฺมํ ปฏิจฺจ นสญฺโญชโน เจว นสญฺโญชนวิปฺปยุตฺโต จ ธมฺโม อุปฺปชฺชติ เหตุปจฺจยา.\csromanspacer{}\csromanspacer{}นสญฺโญชนญฺเจว นสญฺโญชนวิปฺปยุตฺตญฺจ ธมฺมํ ปฏิจฺจ นสญฺโญชนวิปฺปยุตฺโต เจว นโนสญฺโญชโน จ ธมฺโม อุปฺปชฺชติ เหตุปจฺจยา.\csromanspacer{}\csromanspacer{}นสญฺโญชนญฺเจว นสญฺโญชนวิปฺปยุตฺตญฺจ ธมฺมํ ปฏิจฺจ นสญฺโญชโน เจว นสญฺโญชนวิปฺปยุตฺโต จ นสญฺโญชนวิปฺปยุตฺโต เจว นโน จ สญฺโญชโน จ ธมฺมา อุปฺปชฺชนฺติ เหตุปจฺจยา. (3)}

\csromanheader{2}{25. สญฺโญชนวิปฺปยุตฺตสญฺโญชนิยทุก}
\prose{\csromanfolio{87}นสญฺโญชนวิปฺปยุตฺตญฺเจว นโน จ สญฺโญชนํ ธมฺมํ ปฏิจฺจ นสญฺโญชนวิปฺปยุตฺโต เจว นโน จ สญฺโญชโน ธมฺโม อุปฺปชฺชติ เหตุปจฺจยา.\csromanspacer{}\csromanspacer{}นสญฺโญชนวิปฺปยุตฺตญฺเจว นโน จ สญฺโญชนํ ธมฺมํ ปฏิจฺจ นสญฺโญชโน เจว นสญฺโญชนวิปฺปยุตฺโต จ ธมฺโม อุปฺปชฺชติ เหตุปจฺจยา.\csromanspacer{}\csromanspacer{}นสญฺโญชนวิปฺปยุตฺตญฺเจว นโน จ สญฺโญชนํ ธมฺมํ ปฏิจฺจ นสญฺโญชโน เจว นสญฺโญชนวิปฺปยุตฺโต จ นสญฺโญชนวิปฺปยุตฺโต เจว นโน จ สญฺโญชโน จ ธมฺมา อุปฺปชฺชนฺติ เหตุปจฺจยา. (3)}

\prose{นสญฺโญชนญฺเจว นสญฺโญชนวิปฺปยุตฺตญฺจ นสญฺโญชนวิปฺปยุตฺตญฺเจว นโน จ สญฺโญชนญฺจ ธมฺมํ ปฏิจฺจ นสญฺโญชโน เจว นสญฺโญชนวิปฺปยุตฺโต จ ธมฺโม อุปฺปชฺชติ เหตุปจฺจยา.\csromanspacer{} ตีณิ. (สํขิตฺตํ.)}

\csromanheader{2}{เหตุยา นว ฯเปฯ อวิคเต นว. (สพฺพตฺถ วิตฺถาโร.)}
\csromanheader{2}{25. สญฺโญชนวิปฺปยุตฺตสญฺโญชนิยทุก}
\proseitem{39}{สญฺโญชนวิปฺปยุตฺตํ นสญฺโญชนิยํ ธมฺมํ ปฏิจฺจ สญฺโญชนวิปฺปยุตฺโต นสญฺโญชนิโย ธมฺโม อุปฺปชฺชติ เหตุปจฺจยา.\csromanspacer{}\csromanspacer{}สญฺโญชนวิปฺปยุตฺตํ นสญฺโญชนิยํ ธมฺมํ ปฏิจฺจ สญฺโญชนวิปฺปยุตฺโต น-อสญฺโญชนิโย ธมฺโม อุปฺปชฺชติ เหตุปจฺจยา.\csromanspacer{}\csromanspacer{}สญฺโญชนวิปฺปยุตฺตํ นสญฺโญชนิยํ ธมฺมํ ปฏิจฺจ สญฺโญชนวิปฺปยุตฺโต นสญฺโญชนิโย จ สญฺโญชนวิปฺปยุตฺโต น-อสญฺโญชนิโย จ ธมฺมา อุปฺปชฺชนฺติ เหตุปจฺจยา. (3)}

\prose{สญฺโญชนวิปฺปยุตฺตํ น-อสญฺโญชนิยํ ธมฺมํ ปฏิจฺจ สญฺโญชนวิปฺปยุตฺโต น-อสญฺโญชนิโย ธมฺโม อุปฺปชฺชติ เหตุปจฺจยา. (1)}

\prose{สญฺโญชนวิปฺปยุตฺตํ นสญฺโญชนิยญฺจ สญฺโญชนวิปฺปยุตฺตํ น-อสญฺโญชนิยญฺจ ธมฺมํ ปฏิจฺจ สญฺโญชนวิปฺปยุตฺโต น-อสญฺโญชนิโย ธมฺโม อุปฺปชฺชติ เหตุปจฺจยา. (1) (สํขิตฺตํ.)}

\csromanheader{2}{เหตุยา ปญฺจ ฯเปฯ อวิคเต ปญฺจ. (สพฺพตฺถ วิตฺถาโร.)}
\nitthitam{สญฺโญชนโคจฺฉกํ นิฏฺฐิตํ.}
\csromanheader{2}{ธมฺมปจฺจนีย ทุกปฏฺฐานปาฬิ \textbf{26. คนฺถทุก}}
\proseitem{40}{\csromanfolio{88}โนคนฺถํ ธมฺมํ ปฏิจฺจ โนคนฺโถ ธมฺโม อุปฺปชฺชติ เหตุปจฺจยา.\csromanspacer{}\csromanspacer{}โนคนฺถํ ธมฺมํ ปฏิจฺจ นโนคนฺโถ ธมฺโม อุปฺปชฺชติ เหตุปจฺจยา.\csromanspacer{}\csromanspacer{}โนคนฺถํ ธมฺมํ ปฏิจฺจ โนคนฺโถ จ นโนคนฺโถ จ ธมฺมา อุปฺปชฺชนฺติ เหตุปจฺจยา. (3)}

\prose{นโนคนฺถํ ธมฺมํ ปฏิจฺจ นโนคนฺโถ ธมฺโม อุปฺปชฺชติ เหตุปจฺจยา.\csromanspacer{}\csromanspacer{}นโนคนฺถํ ธมฺมํ ปฏิจฺจ โนคนฺโถ ธมฺโม อุปฺปชฺชติ เหตุปจฺจยา.\csromanspacer{}\csromanspacer{}นโนคนฺถํ ธมฺมํ ปฏิจฺจ โนคนฺโถ จ นโนคนฺโถ จ ธมฺมา อุปฺปชฺชนฺติ เหตุปจฺจยา. (3)}

\prose{โนคนฺถญฺจ นโนคนฺถญฺจ ธมฺมํ ปฏิจฺจ โนคนฺโถ ธมฺโม อุปฺปชฺชติ เหตุปจฺจยา.\csromanspacer{}\csromanspacer{}โนคนฺถญฺจ นโนคนฺถญฺจ ธมฺมํ ปฏิจฺจ นโนคนฺโถ ธมฺโม อุปฺปชฺชติ เหตุปจฺจยา.\csromanspacer{}\csromanspacer{}โนคนฺถญฺจ นโนคนฺถญฺจ ธมฺมํ ปฏิจฺจ โนคนฺโถ จ นโนคนฺโถ จ ธมฺมา อุปฺปชฺชนฺติ เหตุปจฺจยา. (3) (สํขิตฺตํ.)}

\csromanheader{2}{เหตุยา นว ฯเปฯ อวิคเต นว. (สพฺพตฺถ นว.)}
\csromanheader{2}{27. คนฺถนิยทุก}
\proseitem{41}{นคนฺถนิยํ ธมฺมํ ปฏิจฺจ นคนฺถนิโย ธมฺโม อุปฺปชฺชติ เหตุปจฺจยา.\csromanspacer{}\csromanspacer{}นคนฺถนิยํ ธมฺมํ ปฏิจฺจ น-อคนฺถนิโย ธมฺโม อุปฺปชฺชติ เหตุปจฺจยา.\csromanspacer{}\csromanspacer{}นคนฺถนิยํ ธมฺมํ ปฏิจฺจ นคนฺถนิโย จ น-อคนฺถนิโย จ ธมฺมา อุปฺปชฺชนฺติ เหตุปจฺจยา. (3)}

\prose{น-อคนฺถนิยํ ธมฺมํ ปฏิจฺจ น-อคนฺถนิโย ธมฺโม อุปฺปชฺชติ}

\prose{นคนฺถนิยญฺจ น-อคนฺถนิยญฺจ ธมฺมํ ปฏิจฺจ น-อคนฺถนิโย ธมฺโม อุปฺปชฺชติ เหตุปจฺจยา. (1) (สํขิตฺตํ.)}

\csromanheader{2}{เหตุยา ปญฺจ ฯเปฯ อวิคเต ปญฺจ. (สพฺพตฺถ ปญฺจ.)}
\csromanheader{2}{28. คนฺถสมฺปยุตฺตทุก}
\proseitem{42}{นคนฺถสมฺปยุตฺตํ ธมฺมํ ปฏิจฺจ นคนฺถสมฺปยุตฺโต ธมฺโม อุปฺปชฺชติ เหตุปจฺจยา.\csromanspacer{}\csromanspacer{}นคนฺถสมฺปยุตฺตํ ธมฺมํ ปฏิจฺจ นคนฺถวิปฺปยุตฺโต ธมฺโม อุปฺปชฺชติ เหตุปจฺจยา.\csromanspacer{}\csromanspacer{}นคนฺถสมฺปยุตฺตํ ธมฺมํ ปฏิจฺจ นคนฺถสมฺปยุตฺโต จ นคนฺถวิปฺปยุตฺโต จ ธมฺมา อุปฺปชฺชนฺติ เหตุปจฺจยา. (3)}

\csromanheader{2}{29. คนฺถคนฺถนิยทุก}
\prose{\csromanfolio{89}นคนฺถวิปฺปยุตฺตํ ธมฺมํ ปฏิจฺจ นคนฺถวิปฺปยุตฺโต ธมฺโม อุปฺปชฺชติ เหตุปจฺจยา.\csromanspacer{}\csromanspacer{}นคนฺถวิปฺปยุตฺตํ ธมฺมํ ปฏิจฺจ นคนฺถสมฺปยุตฺโต ธมฺโม อุปฺปชฺชติ เหตุปจฺจยา.\csromanspacer{}\csromanspacer{}นคนฺถวิปฺปยุตฺตํ ธมฺมํ ปฏิจฺจ นคนฺถสมฺปยุตฺโต จ นคนฺถวิปฺปยุตฺโต จ ธมฺมา อุปฺปชฺชนฺติ เหตุปจฺจยา. (3)}

\prose{นคนฺถสมฺปยุตฺตญฺจ นคนฺถวิปฺปยุตฺตญฺจ ธมฺมํ ปฏิจฺจ นคนฺถสมฺปยุตฺโต ธมฺโม อุปฺปชฺชติ เหตุปจฺจยา.\csromanspacer{}\csromanspacer{}นคนฺถสมฺปยุตฺตญฺจ นคนฺถวิปฺปยุตฺตญฺจ ธมฺมํ ปฏิจฺจ นคนฺถวิปฺปยุตฺโต ธมฺโม อุปฺปชฺชติ เหตุปจฺจยา.\csromanspacer{}\csromanspacer{}นคนฺถสมฺปยุตฺตญฺจ นคนฺถวิปฺปยุตฺตญฺจ ธมฺมํ ปฏิจฺจ นคนฺถสมฺปยุตฺโต จ นคนฺถวิปฺปยุตฺโต จ ธมฺมา อุปฺปชฺชนฺติ เหตุปจฺจยา. (3) (สํขิตฺตํ.)}

\prose{เหตุยา นว, อารมฺมเณ ฉ, อธิปติยา นว ฯเปฯ อวิคเต นว. (สพฺพตฺถ วิตฺถาโร.)}

\csromanheader{2}{29. คนฺถคนฺถนิยทุก}
\proseitem{43}{นคนฺถญฺเจว น-อคนฺถนิยญฺจ ธมฺมํ ปฏิจฺจ นคนฺโถ เจว น-อคนฺถนิโย จ ธมฺโม อุปฺปชฺชติ เหตุปจฺจยา.\csromanspacer{}\csromanspacer{}นคนฺถญฺเจว น-อคนฺถนิยญฺจ ธมฺมํ ปฏิจฺจ น-อคนฺถนิโย เจว นโน จ คนฺโถ ธมฺโม อุปฺปชฺชติ เหตุปจฺจยา.\csromanspacer{}\csromanspacer{}นคนฺถญฺเจว น-อคนฺถนิยญฺจ ธมฺมํ ปฏิจฺจ นคนฺโถ เจว น-อคนฺถนิโย จ น-อคนฺถนิโย เจว นโน จ คนฺโถ จ ธมฺมา อุปฺปชฺชนฺติ เหตุปจฺจยา. (3)}

\prose{น-อคนฺถนิยญฺเจว นโน จ คนฺถํ ธมฺมํ ปฏิจฺจ น-อคนฺถนิโย เจว นโนคนฺโถ จ ธมฺโม อุปฺปชฺชติ เหตุปจฺจยา.\csromanspacer{}\csromanspacer{}น-อคนฺถนิยญฺเจว นโน จ คนฺถํ ธมฺมํ ปฏิจฺจ นคนฺโถ เจว น-อคนฺถนิโย จ ธมฺโม อุปฺปชฺชติ เหตุปจฺจยา.\csromanspacer{}\csromanspacer{}น-อคนฺถนิยญฺเจว นโน จ คนฺถํ ธมฺมํ ปฏิจฺจ นคนฺโถ เจว น-อคนฺถนิโย จ น-อคนฺถนิโย เจว นโน จ คนฺโถ จ ธมฺมา อุปฺปชฺชนฺติ เหตุปจฺจยา. (3)}

\prose{นคนฺถญฺเจว น-อคนฺถนิยญฺจ น-อคนฺถนิยญฺเจว นโน จ คนฺถญฺจ ธมฺมํ ปฏิจฺจ นคนฺโถ เจว น-อคนฺถนิโย จ ธมฺโม อุปฺปชฺชติ เหตุปจฺจยา.\csromanspacer{} ตีณิ. (สํขิตฺตํ.)}

\csromanheader{2}{เหตุยา นว ฯเปฯ อวิคเต นว. (สพฺพตฺถ วิตฺถาโร.)}
\csromanheader{2}{ธมฺมปจฺจนีย ทุกปฏฺฐานปาฬิ \textbf{30. คนฺถคนฺถสมฺปยุตฺตทุก}}
\proseitem{44}{\csromanfolio{90}นคนฺถญฺเจว นคนฺถวิปฺปยุตฺตญฺจ ธมฺมํ ปฏิจฺจ นคนฺโถ เจว นคนฺถวิปฺปยุตฺโต จ ธมฺโม อุปฺปชฺชติ เหตุปจฺจยา.\csromanspacer{}\csromanspacer{}นคนฺถญฺเจว นคนฺถวิปฺปยุตฺตญฺจ ธมฺมํ ปฏิจฺจ นคนฺถวิปฺปยุตฺโต เจว นโน จ คนฺโถ ธมฺโม อุปฺปชฺชติ เหตุปจฺจยา.\csromanspacer{}\csromanspacer{}นคนฺถญฺเจว นคนฺถวิปฺปยุตฺตญฺจ ธมฺมํ ปฏิจฺจ นคนฺโถ เจว นคนฺถวิปฺปยุตฺโต จ นคนฺถวิปฺปยุตฺโต เจว นโน จ คนฺโถ จ ธมฺมา อุปฺปชฺชนฺติ เหตุปจฺจยา. (3)}

\prose{นคนฺถวิปฺปยุตฺตญฺเจว นโน จ คนฺถํ ธมฺมํ ปฏิจฺจ นคนฺถวิปฺปยุตฺโต เจว นโน จ คนฺโถ ธมฺโม อุปฺปชฺชติ เหตุปจฺจยา.\csromanspacer{}\csromanspacer{}นคนฺถวิปฺปยุตฺตญฺเจว นโน จ คนฺถํ ธมฺมํ ปฏิจฺจ นคนฺโถ เจว นคนฺถวิปฺปยุตฺโต จ ธมฺโม อุปฺปชฺชติ เหตุปจฺจยา.\csromanspacer{}\csromanspacer{}นคนฺถวิปฺปยุตฺตญฺเจว นโน จ คนฺถํ ธมฺมํ ปฏิจฺจ นคนฺโถ เจว นคนฺถวิปฺปยุตฺโต จ นคนฺถวิปฺปยุตฺโต เจว นโน จ คนฺโถ จ ธมฺมา อุปฺปชฺชนฺติ เหตุปจฺจยา. (3)}

\prose{นคนฺถญฺเจว นคนฺถวิปฺปยุตฺตญฺจ นคนฺถวิปฺปยุตฺตญฺเจว นโน จ คนฺถญฺจ ธมฺมํ ปฏิจฺจ นคนฺโถ เจว นคนฺถวิปฺปยุตฺโต จ ธมฺโม อุปฺปชฺชติ เหตุปจฺจยา.\csromanspacer{}\csromanspacer{}นคนฺถญฺเจว นคนฺถวิปฺปยุตฺตญฺจ นคนฺถวิปฺปยุตฺตญฺเจว นโน จ คนฺถญฺจ ธมฺมํ ปฏิจฺจ นคนฺถวิปฺปยุตฺโต เจว นโน จ คนฺโถ ธมฺโม อุปฺปชฺชติ เหตุปจฺจยา.\csromanspacer{}\csromanspacer{}นคนฺถญฺเจว นคนฺถวิปฺปยุตฺตญฺจ นคนฺถวิปฺปยุตฺตญฺเจว นโน จ คนฺถญฺจ ธมฺมํ ปฏิจฺจ นคนฺโถ เจว นคนฺถวิปฺปยุตฺโต จ นคนฺถวิปฺปยุตฺโต เจว นโน จ คนฺโถ จ ธมฺมา อุปฺปชฺชนฺติ เหตุปจฺจยา. (3) (สํขิตฺตํ.)}

\csromanheader{2}{เหตุยา นว ฯเปฯ อวิคเต นว. (สพฺพตฺถ วิตฺถาโร.)}
\csromanheader{2}{31. คนฺถวิปฺปยุตฺตคนฺถนิยทุก}
\proseitem{45}{คนฺถวิปฺปยุตฺตํ นคนฺถนิยํ ธมฺมํ ปฏิจฺจ คนฺถวิปฺปยุตฺโต นคนฺถนิโย ธมฺโม อุปฺปชฺชติ เหตุปจฺจยา.\csromanspacer{}\csromanspacer{}คนฺถวิปฺปยุตฺตํ นคนฺถนิยํ ธมฺมํ ปฏิจฺจ คนฺถวิปฺปยุตฺโต น-อคนฺถนิโย ธมฺโม อุปฺปชฺชติ เหตุปจฺจยา.\csromanspacer{}\csromanspacer{}คนฺถวิปฺปยุตฺตํ นคนฺถนิยํ ธมฺมํ ปฏิจฺจ คนฺถวิปฺปยุตฺโต นคนฺถนิโย จ คนฺถวิปฺปยุตฺโต น-อคนฺถนิโย จ ธมฺมา อุปฺปชฺชนฺติ เหตุปจฺจยา. (3)}

\prose{คนฺถวิปฺปยุตฺตํ น-อคนฺถนิยํ ธมฺมํ ปฏิจฺจ คนฺถวิปฺปยุตฺโต น-อคนฺถนิโย ธมฺโม อุปฺปชฺชติ เหตุปจฺจยา. (1)}

\csromanheader{2}{50-54. ปรามาสทุกาทิ}
\prose{\csromanfolio{91}คนฺถวิปฺปยุตฺตํ นคนฺถนิยญฺจ คนฺถวิปฺปยุตฺตํ น-อคนฺถนิยญฺจ ธมฺมํ ปฏิจฺจ คนฺถวิปฺปยุตฺโต น-อคนฺถนิโย ธมฺโม อุปฺปชฺชติ เหตุปจฺจยา. (1) (สํขิตฺตํ.)}

\prose{เหตุยา ปญฺจ, อารมฺมเณ เทฺว ฯเปฯ อวิคเต ปญฺจ. (สพฺพตฺถ วิตฺถาโร.)}

\nitthitam{คนฺถโคจฺฉกํ นิฏฺฐิตํ.}
\csromanheader{2}{32-49. โอฆาทิทุก}
\proseitem{46}{โน-โอฆํ ธมฺมํ ปฏิจฺจ ฯเปฯ.\csromanspacer{} โนโยคํ ธมฺมํ ปฏิจฺจ ฯเปฯ. (โอฆโคจฺฉกโยคโคจฺฉเกสุ อามสนํ อาสวโคจฺฉเก อามสนสทิสํ.)}

\proseitem{47}{โนนีวรณํ ธมฺมํ ปฏิจฺจ ฯเปฯ. (นีวรณโคจฺฉเก อามสนํ สญฺโญชนโคจฺฉเก อามสนสทิสํ.)}

\nitthitam{นีวรณโคจฺฉกํ นิฏฺฐิตํ.}
\csromanheader{2}{50-54. ปรามาสทุกาทิ}
\proseitem{48}{โนปรามาสํ ธมฺมํ ปฏิจฺจ โนปรามาโส ธมฺโม อุปฺปชฺชติ เหตุปจฺจยา. (สํขิตฺตํ.)}

\proseitem{49}{นปรามฏฺฐํ ธมฺมํ ปฏิจฺจ นปรามฏฺโฐ ธมฺโม อุปฺปชฺชติ เหตุปจฺจยา. (สํขิตฺตํ.)}

\proseitem{50}{นปรามาสสมฺปยุตฺตํ ธมฺมํ ปฏิจฺจ นปรามาสสมฺปยุตฺโต ธมฺโม อุปฺปชฺชติ เหตุปจฺจยา.}

\proseitem{51}{นปรามาสญฺเจว น-อปรามฏฺฐญฺจ ธมฺมํ นปรามาโส เจว น-อปรามฏฺโฐ จ ธมฺโม อุปฺปชฺชติ เหตุปจฺจยา.\csromanspacer{}\csromanspacer{}นปรามาสญฺเจว น-อปรามฏฺฐญฺจ ธมฺมํ ปฏิจฺจ น-อปรามฏฺโฐ เจว นโน จ ปรามาโส ธมฺโม อุปฺปชฺชติ เหตุปจฺจยา.\csromanspacer{}\csromanspacer{}นปรามาสญฺเจว น-อปรามฏฺฐญฺจ ธมฺมํ ปฏิจฺจ}

\vspace{\tipitakaparskip}
\csromancenter{ธมฺมปจฺจนีย ทุกปฏฺฐานปาฬิ นปรามาโส เจว น-อปรามฏฺโฐ จ น-อปรามฏฺโฐ เจว นโนปรามาโส จ ธมฺมา อุปฺปชฺชนฺติ เหตุปจฺจยา. (3)}
\prose{\csromanfolio{92}น-อปรามฏฺฐญฺเจว นโน จ ปรามาสํ ธมฺมํ ปฏิจฺจ นปรามาโส เจว น-อปรามฏฺโฐ จ ธมฺโม อุปฺปชฺชติ เหตุปจฺจยา. (1)}

\prose{นปรามาสญฺเจว น-อปรามฏฺฐญฺจ น-อปรามฏฺฐญฺเจว นโน จ ปรามาสญฺจ ธมฺมํ ปฏิจฺจ นปรามาโส เจว น-อปรามฏฺโฐ จ ธมฺโม อุปฺปชฺชติ เหตุปจฺจยา. (1) (สํขิตฺตํ.)}

\proseitem{52}{ปรามาสวิปฺปยุตฺตํ นปรามฏฺฐํ ธมฺมํ ปฏิจฺจ ปรามาสวิปฺปยุตฺโต นปรามฏฺโฐ ธมฺโม อุปฺปชฺชติ เหตุปจฺจยา.\csromanspacer{}\csromanspacer{}ปรามาสวิปฺปยุตฺตํ นปรามฏฺฐํ ธมฺมํ ปฏิจฺจ ปรามาสวิปฺปยุตฺโต น-อปรามฏฺโฐ ธมฺโม อุปฺปชฺชติ เหตุปจฺจยา.\csromanspacer{}\csromanspacer{}ปรามาสวิปฺปยุตฺตํ นปรามฏฺฐํ ธมฺมํ ปฏิจฺจ ปรามาสวิปฺปยุตฺโต นปรามฏฺโฐ จ ปรามาสวิปฺปยุตฺโต น-อปรามฏฺโฐ จ ธมฺมา อุปฺปชฺชนฺติ เหตุปจฺจยา. (3)}

\prose{ปรามาสวิปฺปยุตฺตํ น-อปรามฏฺฐํ ธมฺมํ ปฏิจฺจ ปรามาสวิปฺปยุตฺโต น-อปรามฏฺโฐ ธมฺโม อุปฺปชฺชติ เหตุปจฺจยา. (1)}

\prose{ปรามาสวิปฺปยุตฺตํ นปรามฏฺฐญฺจ ปรามาสวิปฺปยุตฺตํ น-อปรามฏฺฐญฺจ ธมฺมํ ปฏิจฺจ ปรามาสวิปฺปยุตฺโต น-อปรามฏฺโฐ}

\prose{เหตุยา ปญฺจ, อารมฺมเณ เทฺว ฯเปฯ อวิคเต ปญฺจ.}

\nitthitam{ปรามาสโคจฺฉกํ นิฏฺฐิตํ.}
\csromanheader{2}{55. สารมฺมณทุก}
\proseitem{53}{นสารมฺมณํ ธมฺมํ ปฏิจฺจ นสารมฺมโณ ธมฺโม อุปฺปชฺชติ เหตุปจฺจยา.\csromanspacer{}\csromanspacer{}นสารมฺมณํ ธมฺมํ ปฏิจฺจ น-อนารมฺมโณ ธมฺโม อุปฺปชฺชติ เหตุปจฺจยา.\csromanspacer{}\csromanspacer{}นสารมฺมณํ ธมฺมํ ปฏิจฺจ นสารมฺมโณ จ น-อนารมฺมโณ จ ธมฺมา อุปฺปชฺชนฺติ เหตุปจฺจยา. (3)}

\prose{น-อนารมฺมณํ ธมฺมํ ปฏิจฺจ น-อนารมฺมโณ ธมฺโม อุปฺปชฺชติ เหตุปจฺจยา.\csromanspacer{}\csromanspacer{}น-อนารมฺมณํ ธมฺมํ ปฏิจฺจ นสารมฺมโณ ธมฺโม อุปฺปชฺชติ}

\vspace{\tipitakaparskip}
\csromancenter{เจตสิกทุก เหตุปจฺจยา.\csromanspacer{}\csromanspacer{}น-อนารมฺมณํ ธมฺมํ ปฏิจฺจ นสารมฺมโณ จ น-อนารมฺมโณ จ ธมฺมา อุปฺปชฺชนฺติ เหตุปจฺจยา. (3)}
\prose{\csromanfolio{93}นสารมฺมณญฺจ น-อนารมฺมณญฺจ ธมฺมํ ปฏิจฺจ นสารมฺมโณ ธมฺโม อุปฺปชฺชติ เหตุปจฺจยา.\csromanspacer{}\csromanspacer{}นสารมฺมณญฺจ น-อนารมฺมณญฺจ ธมฺมํ ปฏิจฺจ น-อนารมฺมโณ ธมฺโม อุปฺปชฺชติ เหตุปจฺจยา.\csromanspacer{}\csromanspacer{}นสารมฺมณญฺจ น-อนารมฺมณญฺจ ธมฺมํ ปฏิจฺจ นสารมฺมโณ จ น-อนารมฺมโณ จ ธมฺมา อุปฺปชฺชนฺติ เหตุปจฺจยา. (3) (สํขิตฺตํ.)}

\csromanheader{2}{เหตุยา นว ฯเปฯ อวิคเต นว. (สพฺพตฺถ วิตฺถาโร.)}
\csromanheader{2}{56. จิตฺตทุก}
\proseitem{54}{นจิตฺตํ ธมฺมํ ปฏิจฺจ นจิตฺโต ธมฺโม อุปฺปชฺชติ เหตุปจฺจยา.\csromanspacer{}\csromanspacer{}นจิตฺตํ ธมฺมํ ปฏิจฺจ นโนจิตฺโต ธมฺโม อุปฺปชฺชติ เหตุปจฺจยา.\csromanspacer{}\csromanspacer{}นจิตฺตํ ธมฺมํ ปฏิจฺจ นจิตฺโต จ นโนจิตฺโต จ ธมฺมา อุปฺปชฺชนฺติ เหตุปจฺจยา. (3)}

\prose{นโนจิตฺตํ ธมฺมํ ปฏิจฺจ นจิตฺโต ธมฺโม อุปฺปชฺชติ เหตุปจฺจยา. (1)}

\prose{นจิตฺตญฺจ นโนจิตฺตญฺจ ธมฺมํ ปฏิจฺจ นจิตฺโต ธมฺโม อุปฺปชฺชติ เหตุปจฺจยา. (1) (สํขิตฺตํ.)}

\csromanheader{2}{เหตุยา ปญฺจ. (สพฺพตฺถ ปญฺจ.)}
\csromanheader{2}{57. เจตสิกทุก}
\proseitem{55}{นเจตสิกํ ธมฺมํ ปฏิจฺจ นเจตสิโก ธมฺโม อุปฺปชฺชติ เหตุปจฺจยา.\csromanspacer{}\csromanspacer{}นเจตสิกํ ธมฺมํ ปฏิจฺจ น-อเจตสิโก ธมฺโม อุปฺปชฺชติ เหตุปจฺจยา.\csromanspacer{}\csromanspacer{}นเจตสิกํ ธมฺมํ ปฏิจฺจ นเจตสิโก จ น-อเจตสิโก จ ธมฺมา อุปฺปชฺชนฺติ เหตุปจฺจยา. (3)}

\prose{น-อเจตสิกํ ธมฺมํ ปฏิจฺจ น-อเจตสิโก ธมฺโม อุปฺปชฺชติ เหตุปจฺจยา.\csromanspacer{}\csromanspacer{}น-อเจตสิกํ ธมฺมํ ปฏิจฺจ นเจตสิโก ธมฺโม อุปฺปชฺชติ เหตุปจฺจยา.\csromanspacer{}\csromanspacer{}น-อเจตสิกํ ธมฺมํ ปฏิจฺจ นเจตสิโก จ น-อเจตสิโก จ ธมฺมา อุปฺปชฺชนฺติ เหตุปจฺจยา. (3)}

\prose{นเจตสิกญฺจ น-อเจตสิกญฺจ ธมฺมํ ปฏิจฺจ นเจตสิโก ธมฺโม อุปฺปชฺชติ เหตุปจฺจยา.\csromanspacer{}\csromanspacer{}นเจตสิกญฺจ น-อเจตสิกญฺจ ธมฺมํ ปฏิจฺจ น-อเจตสิโก}

\vspace{\tipitakaparskip}
\csromancenter{ธมฺมปจฺจนีย ทุกปฏฺฐานปาฬิ ธมฺโม อุปฺปชฺชติ เหตุปจฺจยา.\csromanspacer{}\csromanspacer{}นเจตสิกญฺจ น-อเจตสิกญฺจ ธมฺมํ ปฏิจฺจ นเจตสิโก จ น-อเจตสิโก จ ธมฺมา อุปฺปชฺชนฺติ เหตุปจฺจยา. (3)}
\csromanheader{2}{เหตุยา นว ฯเปฯ อวิคเต นว. (สพฺพตฺถ วิตฺถาโร.)}
\csromanheader{2}{58. จิตฺตสมฺปยุตฺตทุก}
\proseitem{56}{\csromanfolio{94}นจิตฺตสมฺปยุตฺตํ ธมฺมํ ปฏิจฺจ นจิตฺตสมฺปยุตฺโต ธมฺโม อุปฺปชฺชติ เหตุปจฺจยา.\csromanspacer{} จิตฺตํ ปฏิจฺจ จิตฺตสมุฏฺฐานํ รูปํ.\csromanspacer{}\csromanspacer{}นจิตฺตสมฺปยุตฺตํ ธมฺมํ ปฏิจฺจ นจิตฺตวิปฺปยุตฺโต ธมฺโม อุปฺปชฺชติ เหตุปจฺจยา.\csromanspacer{} จิตฺตํ ปฏิจฺจ จิตฺตสมฺปยุตฺตกา ขนฺธา.}

\prose{(ปจฺจนียทุกมตฺตานิ น วตฺตพฺพํ ธมฺมํ ปูเรตุํ นว นว ปญฺหานิ กโรตุ.)}

\csromanheader{2}{59. จิตฺตสํสฏฺฐทุก}
\proseitem{57}{นจิตฺตสํสฏฺฐํ ธมฺมํ ปฏิจฺจ นจิตฺตสํสฏฺโฐ ธมฺโม อุปฺปชฺชติ เหตุปจฺจยา.\csromanspacer{} เหตุยา นว.}

\csromanheader{2}{60. จิตฺตสมุฏฺฐานทุก}
\proseitem{58}{นจิตฺตสมุฏฺฐานํ ธมฺมํ ปฏิจฺจ นจิตฺตสมุฏฺฐาโน ธมฺโม อุปฺปชฺชติ เหตุปจฺจยา. เหตุยา นว.}

\csromanheader{2}{61. จิตฺตสหภูทุก}
\proseitem{59}{นจิตฺตสหภุํ ธมฺมํ ปฏิจฺจ นจิตฺตสหภู ธมฺโม อุปฺปชฺชติ เหตุปจฺจยา. เหตุยา นว.}

\csromanheader{2}{62. จิตฺตานุปริวตฺติทุก}
\proseitem{60}{นจิตฺตานุปริวตฺติํ ธมฺมํ ปฏิจฺจ นจิตฺตานุปริวตฺตี ธมฺโม อุปฺปชฺชติ เหตุปจฺจยา.\csromanspacer{} เหตุยา นว.}

\csromanheader{2}{63. จิตฺตสํสฏฺฐสมุฏฺฐานทุก}
\proseitem{61}{นจิตฺตสํสฏฺฐสมุฏฺฐานํ ธมฺมํ ปฏิจฺจ นจิตฺตสํสฏฺฐสมุฏฺฐาโน ธมฺโม อุปฺปชฺชติ เหตุปจฺจยา.\csromanspacer{}\csromanspacer{}เหตุยา นว.}

\csromanheader{2}{68. อุปาทินฺนทุก \textbf{64. จิตฺตสํสฏฺฐสมุฏฺฐานสหภูทุก}}
\proseitem{62}{\csromanfolio{95}นจิตฺตสํสฏฺฐสมุฏฺฐานสหภุํ ธมฺมํ ปฏิจฺจ นจิตฺตสํสฏฺฐสมุฏฺฐานสหภู ธมฺโม อุปฺปชฺชติ เหตุปจฺจยา.\csromanspacer{}\csromanspacer{}เหตุยา นว.}

\csromanheader{2}{65. จิตฺตสํสฏฺฐสมุฏฺฐานานุปริวตฺติทุก}
\proseitem{63}{นจิตฺตสํสฏฺฐสมุฏฺฐานานุปริวตฺติํ ธมฺมํ ปฏิจฺจ นจิตฺตสํสฏฺฐสมุฏฺฐานานุปริวตฺตี ธมฺโม อุปฺปชฺชติ เหตุปจฺจยา.}

\csromanheader{2}{เหตุยา นว ฯเปฯ อวิคเต นว. (สพฺพตฺถ วิตฺถาโร.)}
\csromanheader{2}{66. อชฺฌตฺติกทุก}
\proseitem{64}{น-อชฺฌตฺติกํ ธมฺมํ ปฏิจฺจ น-อชฺฌตฺติโก ธมฺโม อุปฺปชฺชติ เหตุปจฺจยา.\csromanspacer{} ตีณิ.}

\prose{นพาหิรํ ธมฺมํ ปฏิจฺจ นพาหิโร ธมฺโม อุปฺปชฺชติ เหตุปจฺจยา.\csromanspacer{} ปฏิสนฺธิกฺขเณ จิตฺตํ ปฏิจฺจ อชฺฌตฺติกํ กฏตฺตารูปํ. (สํขิตฺตํ.)}

\csromanheader{2}{เหตุยา นว ฯเปฯ อวิคเต นว. (สพฺพตฺถ วิตฺถาโร.)}
\csromanheader{2}{67. อุปาทาทุก}
\proseitem{65}{น-อุปาทา ธมฺมํ ปฏิจฺจ น-อุปาทา ธมฺโม อุปฺปชฺชติ เหตุปจฺจยา.\csromanspacer{}\csromanspacer{}น-อุปาทา ธมฺมํ ปฏิจฺจ นโน-อุปาทา ธมฺโม อุปฺปชฺชติ เหตุปจฺจยา.\csromanspacer{}\csromanspacer{}น-อุปาทา ธมฺมํ ปฏิจฺจ น-อุปาทา จ นโน-อุปาทา จ ธมฺมา อุปฺปชฺชนฺติ เหตุปจฺจยา. (3)}

\csromanheader{2}{นโน-อุปาทา ธมฺมํ ปฏิจฺจ น-อุปาทา ธมฺโม อุปฺปชฺชติ}
\prose{น-อุปาทา จ นโน-อุปาทา จ ธมฺมํ ปฏิจฺจ น-อุปาทา ธมฺโม อุปฺปชฺชติ เหตุปจฺจยา. (1) (สํขิตฺตํ.)}

\prose{เหตุยา ปญฺจ.}

\csromanheader{2}{68. อุปาทินฺนทุก}
\csromanheader{2}{66. น-อุปาทินฺนํ ธมฺมํ ปฏิจฺจ น-อุปาทินฺโน ธมฺโม อุปฺปชฺชติ}
\prose{น-อนุปาทินฺนํ ธมฺมํ ปฏิจฺจ น-อนุปาทินฺโน ธมฺโม อุปฺปชฺชติ เหตุปจฺจยา.\csromanspacer{}\csromanspacer{}น-อนุปาทินฺนํ ธมฺมํ ปฏิจฺจ น-อุปาทินฺโน ธมฺโม อุปฺปชฺชติ เหตุปจฺจยา.\csromanspacer{}\csromanspacer{}}

\vspace{\tipitakaparskip}
\csromancenter{ธมฺมปจฺจนีย ทุกปฏฺฐานปาฬิ น-อนุปาทินฺนํ ธมฺมํ ปฏิจฺจ น-อุปาทินฺโน จ น-อนุปาทินฺโน จ ธมฺมา อุปฺปชฺชนฺติ เหตุปจฺจยา. (3)}
\prose{\csromanfolio{96}น-อุปาทินฺนญฺจ น-อนุปาทินฺนญฺจ ธมฺมํ ปฏิจฺจ น-อุปาทินฺโน ธมฺโม อุปฺปชฺชติ เหตุปจฺจยา. (1) (สํขิตฺตํ.)}

\nitthitam{เหตุยา ปญฺจ. (สพฺพตฺถ วิตฺถาโร.) มหนฺตรทุกํ นิฏฺฐิตํ.}
\csromanheader{2}{69. อุปาทานโคจฺฉก}
\proseitem{67}{น-อุปาทานํ ธมฺมํ ปฏิจฺจ น-อุปาทาโน ธมฺโม อุปฺปชฺชติ เหตุปจฺจยา. (สํขิตฺตํ.) \textbf{75.}\csromanspacer{}\textbf{ กิเลสทุก}}

\proseitem{68}{โนกิเลสํ ธมฺมํ ปฏิจฺจ โนกิเลโส ธมฺโม อุปฺปชฺชติ เหตุปจฺจยา.\csromanspacer{}\csromanspacer{}โนกิเลสํ ธมฺมํ ปฏิจฺจ นโนกิเลโส ธมฺโม อุปฺปชฺชติ เหตุปจฺจยา.\csromanspacer{}\csromanspacer{}โนกิเลสํ ธมฺมํ ปฏิจฺจ โนกิเลโส จ นโนกิเลโส จ ธมฺมา อุปฺปชฺชนฺติ เหตุปจฺจยา. (3)}

\prose{นโนกิเลสํ ธมฺมํ ปฏิจฺจ นโนกิเลโส ธมฺโม อุปฺปชฺชติ เหตุปจฺจยา.\csromanspacer{} ตีณิ.}

\prose{โนกิเลสญฺจ นโนกิเลสญฺจ ธมฺมํ ปฏิจฺจ โนกิเลโส ธมฺโม อุปฺปชฺชติ เหตุปจฺจยา.\csromanspacer{} ตีณิ. (สํขิตฺตํ.)}

\csromanheader{2}{เหตุยา นว. (สพฺพตฺถ นว.) \textbf{76. สํกิเลสิกทุก}}
\proseitem{69}{นสํกิเลสิกํ ธมฺมํ ปฏิจฺจ นสํกิเลสิโก ธมฺโม อุปฺปชฺชติ เหตุปจฺจยา.\csromanspacer{}\csromanspacer{}นสํกิเลสิกํ ธมฺมํ ปฏิจฺจ น-อสํกิเลสิโก ธมฺโม อุปฺปชฺชติ เหตุปจฺจยา.\csromanspacer{}\csromanspacer{}นสํกิเลสิกํ ธมฺมํ ปฏิจฺจ นสํกิเลสิโก จ น-อสํกิเลสิโก จ ธมฺมา อุปฺปชฺชนฺติ เหตุปจฺจยา. (3)}

\prose{น-อสํกิเลสิกํ ธมฺมํ ปฏิจฺจ น-อสํกิเลสิโก ธมฺโม อุปฺปชฺชติ เหตุปจฺจยา. (1)}

\csromanheader{2}{79. กิเลสสํกิเลสิกทุก}
\prose{\csromanfolio{97}นสํกิเลสิกญฺจ น-อสํกิเลสิกญฺจ ธมฺมํ ปฏิจฺจ น-อสํกิเลสิโก ธมฺโม อุปฺปชฺชติ เหตุปจฺจยา. (1) (สํขิตฺตํ.)}

\csromanheader{2}{เหตุยา ปญฺจ ฯเปฯ อวิคเต ปญฺจ. (สพฺพตฺถ วิตฺถาโร.)}
\csromanheader{2}{77. สํกิลิฏฺฐทุก}
\proseitem{70}{นสํกิลิฏฺฐํ ธมฺมํ ปฏิจฺจ นสํกิลิฏฺโฐ ธมฺโม อุปฺปชฺชติ}

\prose{น-อสํกิลิฏฺฐํ ธมฺมํ ปฏิจฺจ น-อสํกิลิฏฺโฐ ธมฺโม อุปฺปชฺชติ เหตุปจฺจยา.\csromanspacer{}\csromanspacer{}น-อสํกิลิฏฺฐํ ธมฺมํ ปฏิจฺจ นสํกิลิฏฺโฐ ธมฺโม อุปฺปชฺชติ เหตุปจฺจยา.\csromanspacer{}\csromanspacer{}น-อสํกิลิฏฺฐํ ธมฺมํ ปฏิจฺจ นสํกิลิฏฺโฐ จ น-อสํกิลิฏฺโฐ จ ธมฺมา อุปฺปชฺชนฺติ เหตุปจฺจยา. (3)}

\prose{นสํกิลิฏฺฐญฺจ น-อสํกิลิฏฺฐญฺจ ธมฺมํ ปฏิจฺจ นสํกิลิฏฺโฐ ธมฺโม อุปฺปชฺชติ เหตุปจฺจยา. (1) (สํขิตฺตํ.)}

\csromanheader{2}{เหตุยา ปญฺจ ฯเปฯ อวิคเต ปญฺจ. (สพฺพตฺถ วิตฺถาโร.)}
\csromanheader{2}{78. กิเลสสมฺปยุตฺตทุก}
\csromanheader{2}{71. นกิเลสสมฺปยุตฺตํ ธมฺมํ ปฏิจฺจ นกิเลสสมฺปยุตฺโต}
\prose{นกิเลสวิปฺปยุตฺตํ ธมฺมํ ปฏิจฺจ นกิเลสวิปฺปยุตฺโต ธมฺโม อุปฺปชฺชติ เหตุปจฺจยา.\csromanspacer{} ตีณิ.}

\prose{นกิเลสสมฺปยุตฺตญฺจ นกิเลสวิปฺปยุตฺตญฺจ ธมฺมํ ปฏิจฺจ นกิเลสสมฺปยุตฺโต ธมฺโม อุปฺปชฺชติ เหตุปจฺจยา. (1) (สํขิตฺตํ.)}

\prose{เหตุยา ปญฺจ, อวิคเต ปญฺจ. (สพฺพตฺถ วิตฺถาโร.)}

\csromanheader{2}{79. กิเลสสํกิเลสิกทุก}
\proseitem{72}{นกิเลสญฺเจว น-อสํกิเลสิกญฺจ ธมฺมํ ปฏิจฺจ นกิเลโส เจว น-อสํกิเลสิโก จ ธมฺโม อุปฺปชฺชติ เหตุปจฺจยา.\csromanspacer{} ตีณิ.}

\prose{น-อสํกิเลสิกญฺเจว นโน จ กิเลสํ ธมฺมํ ปฏิจฺจ น-อสํกิเลสิโก เจว นโน กิเลโส ธมฺโม อุปฺปชฺชติ เหตุปจฺจยา.\csromanspacer{} ตีณิ.}

\gambhira{ธมฺมปจฺจนีย ทุกปฏฺฐานปาฬิ}
\prose{\csromanfolio{98}นกิเลสญฺเจว น-อสํกิเลสิกญฺจ น-อสํกิเลสิกญฺเจว นโน จ กิเลสญฺจ ธมฺมํ ปฏิจฺจ นกิเลโส เจว น-อสํกิเลสิโก จ ธมฺโม อุปฺปชฺชติ เหตุปจฺจยา.\csromanspacer{} ตีณิ. (สํขิตฺตํ.)}

\prose{เหตุยา นว.}

\csromanheader{2}{80. กิเลสสํกิลิฏฺฐทุก}
\proseitem{73}{นกิเลสญฺเจว น-อสํกิลิฏฺฐญฺจ ธมฺมํ ปฏิจฺจ นกิเลโส เจว น-อสํกิลิฏฺโฐ จ ธมฺโม อุปฺปชฺชติ เหตุปจฺจยา.\csromanspacer{} ตีณิ.}

\prose{น-อสํกิลิฏฺฐญฺเจว นโน จ กิเลสํ ธมฺมํ ปฏิจฺจ น-อสํกิลิฏฺโฐ เจว นโน จ กิเลโส ธมฺโม อุปฺปชฺชติ เหตุปจฺจยา.\csromanspacer{} ตีณิ.}

\prose{นกิเลสญฺเจว น-อสํกิลิฏฺฐญฺจ น-อสํกิลิฏฺฐญฺเจว นโน จ กิเลสญฺจ ธมฺมํ ปฏิจฺจ นกิเลโส เจว น-อสํกิลิฏฺโฐ จ ธมฺโม อุปฺปชฺชติ เหตุปจฺจยา.\csromanspacer{} ตีณิ. (สํขิตฺตํ.)}

\prose{เหตุยา นว.}

\csromanheader{2}{81. กิเลสกิเลสสมฺปยุตฺตทุก}
\proseitem{74}{นกิเลสญฺเจว นกิเลสวิปฺปยุตฺตญฺจ ธมฺมํ ปฏิจฺจ นกิเลโส เจว นกิเลสวิปฺปยุตฺโต จ ธมฺโม อุปฺปชฺชติ เหตุปจฺจยา.\csromanspacer{} ตีณิ.}

\prose{นกิเลสวิปฺปยุตฺตญฺเจว นโน จ กิเลสํ ธมฺมํ ปฏิจฺจ นกิเลสวิปฺปยุตฺโต เจว นโน จ กิเลโส ธมฺโม อุปฺปชฺชติ เหตุปจฺจยา.\csromanspacer{} ตีณิ.}

\prose{นกิเลสญฺเจว นกิเลสวิปฺปยุตฺตญฺจ นกิเลสวิปฺปยุตฺตญฺเจว นโน จ กิเลสญฺจ ธมฺมํ ปฏิจฺจ นกิเลโส เจว นกิเลสวิปฺปยุตฺโต จ ธมฺโม อุปฺปชฺชติ เหตุปจฺจยา.\csromanspacer{} ตีณิ. (สํขิตฺตํ.)}

\prose{เหตุยา นว.}

\csromanheader{2}{82. กิเลสวิปฺปยุตฺตสํกิเลสิกทุก}
\proseitem{75}{กิเลสวิปฺปยุตฺตํ นสํกิเลสิกํ ธมฺมํ ปฏิจฺจ กิเลสวิปฺปยุตฺโต นสํกิเลสิโก ธมฺโม อุปฺปชฺชติ เหตุปจฺจยา.\csromanspacer{}\csromanspacer{}กิเลสวิปฺปยุตฺตํ นสํกิเลสิกํ ธมฺมํ ปฏิจฺจ กิเลสวิปฺปยุตฺโต น-อสํกิเลสิโก ธมฺโม อุปฺปชฺชติ เหตุปจฺจยา.\csromanspacer{}\csromanspacer{}กิเลสวิปฺปยุตฺตํ นสํกิเลสิกํ ธมฺมํ ปฏิจฺจ}

\vspace{\tipitakaparskip}
\csromancenter{ทสฺสเนนปหาตพฺพเหตุกทุก กิเลสวิปฺปยุตฺโต นสํกิเลสิโก จ กิเลสวิปฺปยุตฺโต น-อสํกิเลสิโก จ ธมฺมา อุปฺปชฺชนฺติ เหตุปจฺจยา. (3)}
\prose{\csromanfolio{99}กิเลสวิปฺปยุตฺตํ น-อสํกิเลสิกํ ธมฺมํ ปฏิจฺจ กิเลสวิปฺปยุตฺโต น-อสํกิเลสิโก ธมฺโม อุปฺปชฺชติ เหตุปจฺจยา. (1)}

\prose{กิเลสวิปฺปยุตฺตํ นสํกิเลสิกญฺจ กิเลสวิปฺปยุตฺตํ น-อสํกิเลสิกญฺจ ธมฺมํ ปฏิจฺจ กิเลสวิปฺปยุตฺโต น-อสํกิเลสิโก ธมฺโม อุปฺปชฺชติ เหตุปจฺจยา. (1) (สํขิตฺตํ.)}

\csromanheader{2}{เหตุยา ปญฺจ. (สพฺพตฺถ วิตฺถาโร.)}
\csromanheader{2}{83. ทสฺสเนนปหาตพฺพทุก}
\csromanheader{2}{76. นทสฺสเนน ปหาตพฺพํ ธมฺมํ ปฏิจฺจ นทสฺสเนน ปหาตพฺโพ}
\prose{นนทสฺสเนน ปหาตพฺพํ ธมฺมํ ปฏิจฺจ นนทสฺสเนน ปหาตพฺโพ ธมฺโม อุปฺปชฺชติ เหตุปจฺจยา.\csromanspacer{} ตีณิ.}

\prose{นทสฺสเนน ปหาตพฺพญฺจ นนทสฺสเนน ปหาตพฺพญฺจ ธมฺมํ ปฏิจฺจ นทสฺสเนน ปหาตพฺโพ ธมฺโม อุปฺปชฺชติ เหตุปจฺจยา. (1)}

\prose{เหตุยา ปญฺจ, อารมฺมเณ เทฺว.}

\csromanheader{2}{84. ภาวนายปหาตพฺพทุก}
\proseitem{77}{นภาวนาย ปหาตพฺพํ ธมฺมํ ปฏิจฺจ นภาวนาย ปหาตพฺโพ ธมฺโม อุปฺปชฺชติ เหตุปจฺจยา. (1)}

\prose{นนภาวนาย ปหาตพฺพํ ธมฺมํ ปฏิจฺจ นนภาวนาย ปหาตพฺโพ ธมฺโม อุปฺปชฺชติ เหตุปจฺจยา.\csromanspacer{} ตีณิ.}

\prose{นภาวนาย ปหาตพฺพญฺจ นนภาวนาย ปหาตพฺพญฺจ ธมฺมํ ปฏิจฺจ นภาวนาย ปหาตพฺโพ ธมฺโม อุปฺปชฺชติ เหตุปจฺจยา. (1)}

\prose{เหตุยา ปญฺจ.}

\csromanheader{2}{85. ทสฺสเนนปหาตพฺพเหตุกทุก}
\proseitem{78}{นทสฺสเนน ปหาตพฺพเหตุกํ ธมฺมํ ปฏิจฺจ นทสฺสเนน ปหาตพฺพเหตุโก ธมฺโม อุปฺปชฺชติ เหตุปจฺจยา.\csromanspacer{} ตีณิ.}

\gambhira{ธมฺมปจฺจนีย ทุกปฏฺฐานปาฬิ}
\prose{\csromanfolio{100}นนทสฺสเนน ปหาตพฺพเหตุกํ ธมฺมํ ปฏิจฺจ นนทสฺสเนน ปหาตพฺพเหตุโก ธมฺโม อุปฺปชฺชติ เหตุปจฺจยา.\csromanspacer{} ตีณิ.}

\prose{นทสฺสเนน ปหาตพฺพเหตุกญฺจ นนทสฺสเนน ปหาตพฺพเหตุกญฺจ ธมฺมํ ปฏิจฺจ นทสฺสเนน ปหาตพฺพเหตุโก ธมฺโม อุปฺปชฺชติ เหตุปจฺจยา.\csromanspacer{} ตีณิ. (สํขิตฺตํ.)}

\prose{เหตุยา นว.}

\csromanheader{2}{86. ภาวนายปหาตพฺพเหตุกทุก}
\proseitem{79}{นภาวนาย ปหาตพฺพเหตุกํ ธมฺมํ ปฏิจฺจ นภาวนาย ปหาตพฺพเหตุโก ธมฺโม อุปฺปชฺชติ เหตุปจฺจยา.\csromanspacer{} ตีณิ.}

\prose{นนภาวนาย ปหาตพฺพเหตุกํ ธมฺมํ ปฏิจฺจ นนภาวนาย ปหาตพฺพเหตุโก ธมฺโม อุปฺปชฺชติ เหตุปจฺจยา.\csromanspacer{} ตีณิ.}

\prose{นภาวนาย ปหาตพฺพเหตุกญฺจ นนภาวนาย ปหาตพฺพเหตุกญฺจ ธมฺมํ ปฏิจฺจ นภาวนาย ปหาตพฺพเหตุโก ธมฺโม อุปฺปชฺชติ เหตุปจฺจยา.\csromanspacer{} ตีณิ. (สํขิตฺตํ.)}

\prose{เหตุยา นว.}

\csromanheader{2}{87. สวิตกฺกทุก}
\proseitem{80}{นสวิตกฺกํ ธมฺมํ ปฏิจฺจ นสวิตกฺโก ธมฺโม อุปฺปชฺชติ เหตุปจฺจยา.\csromanspacer{} ตีณิ.}

\prose{น-อวิตกฺกํ ธมฺมํ ปฏิจฺจ น-อวิตกฺโก ธมฺโม อุปฺปชฺชติ เหตุปจฺจยา.\csromanspacer{} ตีณิ.}

\prose{นสวิตกฺกญฺจ น-อวิตกฺกญฺจ ธมฺมํ ปฏิจฺจ นสวิตกฺโก ธมฺโม อุปฺปชฺชติ เหตุปจฺจยา.\csromanspacer{} ตีณิ. (สํขิตฺตํ.)}

\prose{เหตุยา นว.}

\csromanheader{2}{88. สวิจารทุก}
\proseitem{81}{นสวิจารํ ธมฺมํ ปฏิจฺจ นสวิจาโร ธมฺโม อุปฺปชฺชติ เหตุปจฺจยา.\csromanspacer{} ตีณิ.}

\prose{น-อวิจารํ ธมฺมํ ปฏิจฺจ น-อวิจาโร ธมฺโม อุปฺปชฺชติ เหตุปจฺจยา.\csromanspacer{} ตีณิ.}

\csromanheader{2}{91. สุขสหคตทุก}
\prose{\csromanfolio{101}นสวิจารญฺจ น-อวิจารญฺจ ธมฺมํ ปฏิจฺจ นสวิจาโร ธมฺโม อุปฺปชฺชติ เหตุปจฺจยา.\csromanspacer{} ตีณิ. (สํขิตฺตํ.)}

\prose{เหตุยา นว.}

\csromanheader{2}{89. สปฺปีติกทุก}
\proseitem{82}{นสปฺปีติกํ ธมฺมํ ปฏิจฺจ นสปฺปีติโก ธมฺโม อุปฺปชฺชติ เหตุปจฺจยา.\csromanspacer{} ตีณิ.}

\prose{น-อปฺปีติกํ ธมฺมํ ปฏิจฺจ น-อปฺปีติโก ธมฺโม อุปฺปชฺชติ เหตุปจฺจยา.\csromanspacer{} ตีณิ.}

\prose{นสปฺปีติกญฺจ น-อปฺปีติกญฺจ ธมฺมํ ปฏิจฺจ นสปฺปีติโก ธมฺโม อุปฺปชฺชติ เหตุปจฺจยา.\csromanspacer{} ตีณิ. (สํขิตฺตํ.)}

\prose{เหตุยา นว.}

\csromanheader{2}{90. ปีติสหคตทุก}
\proseitem{83}{นปีติสหคตํ ธมฺมํ ปฏิจฺจ นปีติสหคโต ธมฺโม อุปฺปชฺชติ เหตุปจฺจยา.\csromanspacer{} ตีณิ.}

\prose{นนปีติสหคตํ ธมฺมํ ปฏิจฺจ นนปีติสหคโต ธมฺโม อุปฺปชฺชติ เหตุปจฺจยา.\csromanspacer{} ตีณิ.}

\prose{นปีติสหคตญฺจ นนปีติสหคตญฺจ ธมฺมํ ปฏิจฺจ นปีติสหคโต ธมฺโม อุปฺปชฺชติ เหตุปจฺจยา.\csromanspacer{} ตีณิ. (สํขิตฺตํ.)}

\prose{เหตุยา นว.}

\csromanheader{2}{91. สุขสหคตทุก}
\proseitem{84}{นสุขสหคตํ ธมฺมํ ปฏิจฺจ นสุขสหคโต ธมฺโม อุปฺปชฺชติ เหตุปจฺจยา.\csromanspacer{} ตีณิ.}

\prose{นนสุขสหคตํ ธมฺมํ ปฏิจฺจ นนสุขสหคโต ธมฺโม อุปฺปชฺชติ เหตุปจฺจยา.\csromanspacer{} ตีณิ.}

\prose{นสุขสหคตญฺจ นนสุขสหคตญฺจ ธมฺมํ ปฏิจฺจ นสุขสหคโต ธมฺโม อุปฺปชฺชติ เหตุปจฺจยา.\csromanspacer{} ตีณิ. (สํขิตฺตํ.)}

\prose{เหตุยา นว.}

\csromanheader{2}{ธมฺมปจฺจนีย ทุกปฏฺฐานปาฬิ \textbf{92. อุเปกฺขาสหคตทุก}}
\proseitem{85}{\csromanfolio{102}น-อุเปกฺขาสหคตํ ธมฺมํ ปฏิจฺจ น-อุเปกฺขาสหคโต ธมฺโม อุปฺปชฺชติ เหตุปจฺจยา.\csromanspacer{} ตีณิ.}

\prose{นน-อุเปกฺขาสหคตํ ธมฺมํ ปฏิจฺจ นน-อุเปกฺขาสหคโต ธมฺโม อุปฺปชฺชติ เหตุปจฺจยา.\csromanspacer{} ตีณิ.}

\prose{น-อุเปกฺขาสหคตญฺจ นน-อุเปกฺขาสหคตญฺจ ธมฺมํ ปฏิจฺจ น-อุเปกฺขาสหคโต ธมฺโม อุปฺปชฺชติ เหตุปจฺจยา.\csromanspacer{} ตีณิ. (สํขิตฺตํ.)}

\prose{เหตุยา นว.}

\csromanheader{2}{93. กามาวจรทุก}
\proseitem{86}{นกามาวจรํ ธมฺมํ ปฏิจฺจ นกามาวจโร ธมฺโม อุปฺปชฺชติ เหตุปจฺจยา.\csromanspacer{} ตีณิ.}

\prose{นนกามาวจรํ ธมฺมํ ปฏิจฺจ นนกามาวจโร ธมฺโม อุปฺปชฺชติ เหตุปจฺจยา.\csromanspacer{} ตีณิ.}

\prose{นกามาวจรญฺจ นนกามาวจรญฺจ ธมฺมํ ปฏิจฺจ นกามาวจโร ธมฺโม อุปฺปชฺชติ เหตุปจฺจยา.\csromanspacer{} ตีณิ. (สํขิตฺตํ.)}

\prose{เหตุยา นว.}

\csromanheader{2}{94. รูปาวจรทุก}
\proseitem{87}{นรูปาวจรํ ธมฺมํ ปฏิจฺจ นรูปาวจโร ธมฺโม อุปฺปชฺชติ เหตุปจฺจยา.\csromanspacer{} ตีณิ.}

\prose{นนรูปาวจรํ ธมฺมํ ปฏิจฺจ นนรูปาวจโร ธมฺโม อุปฺปชฺชติ เหตุปจฺจยา.\csromanspacer{} ตีณิ.}

\prose{นรูปาวจรญฺจ นนรูปาวจรญฺจ ธมฺมํ ปฏิจฺจ นรูปาวจโร ธมฺโม อุปฺปชฺชติ เหตุปจฺจยา.\csromanspacer{} ตีณิ. (สํขิตฺตํ.)}

\prose{เหตุยา นว.}

\csromanheader{2}{95. อรูปาวจรทุก}
\proseitem{88}{น-อรูปาวจรํ ธมฺมํ ปฏิจฺจ น-อรูปาวจโร ธมฺโม อุปฺปชฺชติ เหตุปจฺจยา.\csromanspacer{} เอกํ.}

\csromanheader{2}{98. นิยตทุก}
\prose{\csromanfolio{103}นน-อรูปาวจรํ ธมฺมํ ปฏิจฺจ นน-อรูปาวจโร ธมฺโม อุปฺปชฺชติ เหตุปจฺจยา.\csromanspacer{} ตีณิ.}

\prose{น-อรูปาวจรญฺจ นน-อรูปาวจรญฺจ ธมฺมํ ปฏิจฺจ น-อรูปาวจโร ธมฺโม อุปฺปชฺชติ เหตุปจฺจยา.\csromanspacer{} เอกํ. (สํขิตฺตํ.)}

\prose{เหตุยา ปญฺจ.}

\csromanheader{2}{96. ปริยาปนฺนทุก}
\proseitem{89}{นปริยาปนฺนํ ธมฺมํ ปฏิจฺจ นปริยาปนฺโน ธมฺโม อุปฺปชฺชติ เหตุปจฺจยา.\csromanspacer{} ตีณิ.}

\prose{น-อปริยาปนฺนํ ธมฺมํ ปฏิจฺจ น-อปริยาปนฺโน ธมฺโม อุปฺปชฺชติ เหตุปจฺจยา.\csromanspacer{} เอกํ.}

\prose{นปริยาปนฺนญฺจ น-อปริยาปนฺนญฺจ ธมฺมํ ปฏิจฺจ น-อปริยาปนฺโน ธมฺโม อุปฺปชฺชติ เหตุปจฺจยา.\csromanspacer{} เอกํ. (สํขิตฺตํ.)}

\prose{เหตุยา ปญฺจ.}

\csromanheader{2}{97. นิยฺยานิกทุก}
\proseitem{90}{นนิยฺยานิกํ ธมฺมํ ปฏิจฺจ นนิยฺยานิโก ธมฺโม อุปฺปชฺชติ เหตุปจฺจยา.\csromanspacer{} เอกํ.}

\prose{น-อนิยฺยานิกํ ธมฺมํ ปฏิจฺจ น-อนิยฺยานิโก ธมฺโม อุปฺปชฺชติ เหตุปจฺจยา.\csromanspacer{} ตีณิ.}

\prose{นนิยฺยานิกญฺจ น-อนิยฺยานิกญฺจ ธมฺมํ ปฏิจฺจ นนิยฺยานิโก ธมฺโม อุปฺปชฺชติ เหตุปจฺจยา.\csromanspacer{} เอกํ. (สํขิตฺตํ.)}

\prose{เหตุยา ปญฺจ.}

\csromanheader{2}{98. นิยตทุก}
\proseitem{91}{นนิยตํ ธมฺมํ ปฏิจฺจ นนิยโต ธมฺโม อุปฺปชฺชติ เหตุปจฺจยา.\csromanspacer{} เอกํ.}

\prose{น-อนิยตํ ธมฺมํ ปฏิจฺจ น-อนิยโต ธมฺโม อุปฺปชฺชติ เหตุปจฺจยา.\csromanspacer{} ตีณิ.}

\gambhira{ธมฺมปจฺจนีย ทุกปฏฺฐานปาฬิ}
\prose{\csromanfolio{104}นนิยตญฺจ น-อนิยตญฺจ ธมฺมํ ปฏิจฺจ นนิยโต ธมฺโม อุปฺปชฺชติ เหตุปจฺจยา.\csromanspacer{} เอกํ. (สํขิตฺตํ.)}

\prose{เหตุยา ปญฺจ.}

\csromanheader{2}{99. ส-อุตฺตรทุก}
\proseitem{92}{นส-อุตฺตรํ ธมฺมํ ปฏิจฺจ นส-อุตฺตโร ธมฺโม อุปฺปชฺชติ เหตุปจฺจยา.\csromanspacer{} ตีณิ.}

\prose{น-อนุตฺตรํ ธมฺมํ ปฏิจฺจ น-อนุตฺตโร ธมฺโม อุปฺปชฺชติ เหตุปจฺจยา.\csromanspacer{} เอกํ.}

\prose{นส-อุตฺตรญฺจ น-อนุตฺตรญฺจ ธมฺมํ ปฏิจฺจ น-อนุตฺตโร ธมฺโม อุปฺปชฺชติ เหตุปจฺจยา.\csromanspacer{} เอกํ. (สํขิตฺตํ.)}

\prose{เหตุยา ปญฺจ.}

\csromanheader{2}{100. สรณทุก}
\proseitem{93}{นสรณํ ธมฺมํ ปฏิจฺจ นสรโณ ธมฺโม อุปฺปชฺชติ เหตุปจฺจยา.\csromanspacer{} เอกํ.}

\prose{น-อรณํ ธมฺมํ ปฏิจฺจ น-อรโณ ธมฺโม อุปฺปชฺชติ เหตุปจฺจยา.\csromanspacer{} ตีณิ.}

\prose{นสรณญฺจ น-อรณญฺจ ธมฺมํ ปฏิจฺจ นสรโณ ธมฺโม อุปฺปชฺชติ เหตุปจฺจยา.\csromanspacer{} เอกํ. (สํขิตฺตํ.)}

\prose{เหตุยา ปญฺจ.}

\nitthitam{ปิฏฺฐิทุกํ นิฏฺฐิตํ.}
\nitthitam{\textbf{ธมฺมปจฺจนีเย ทุกปฏฺฐานํ นิฏฺฐิตํ.}}
\pitaka{\textbf{อภิธมฺมปิฏก}}
\csromanheader{2}{\textbf{ธมฺมปจฺจนีย}}
\gambhira{\textbf{ทุกติกปฏฺฐานปาฬิ}}
\namakkaram{นโม ตสฺส ภควโต อรหโต สมฺมาสมฺพุทฺธสฺส.}
\csromanheader{2}{1. เหตุทุก 1. กุสลตฺติก}
\csromanheader{2}{1-7. ปฏิจฺจาทิวาร ปจฺจยจตุกฺก-เหตุ}
\proseitem{1}{\csromanfolio{105}นเหตุํ นกุสลํ ธมฺมํ ปฏิจฺจ นเหตุ นกุสโล ธมฺโม อุปฺปชฺชติ เหตุปจฺจยา.\csromanspacer{} นเหตุํ อกุสลํ อพฺยากตํ เอกํ ขนฺธํ ปฏิจฺจ ตโย ขนฺธา จิตฺตสมุฏฺฐานญฺจ รูปํ.\csromanspacer{}\csromanspacer{}นเหตุํ นกุสลํ ธมฺมํ ปฏิจฺจ นนเหตุํ นกุสโล ธมฺโม อุปฺปชฺชติ เหตุปจฺจยา.\csromanspacer{} นเหตู อกุสเล อพฺยากเต ขนฺเธ ปฏิจฺจ เหตู.\csromanspacer{}\csromanspacer{}นเหตุํ นกุสลํ ธมฺมํ ปฏิจฺจ นเหตุ นกุสโล จ นนเหตุ นกุสโล จ ธมฺมา อุปฺปชฺชนฺติ เหตุปจฺจยา. (3)}

\prose{นนเหตุํ นกุสลํ ธมฺมํ ปฏิจฺจ นนเหตุ นกุสโล ธมฺโม อุปฺปชฺชติ เหตุปจฺจยา.\csromanspacer{}\csromanspacer{}นนเหตุํ นกุสลํ ธมฺมํ ปฏิจฺจ นเหตุ นกุสโล ธมฺโม อุปฺปชฺชติ เหตุปจฺจยา.\csromanspacer{}\csromanspacer{}นนเหตุํ นกุสลํ ธมฺมํ ปฏิจฺจ นเหตุ นกุสโล จ นนเหตุ นกุสโล จ ธมฺมา อุปฺปชฺชนฺติ เหตุปจฺจยา. (3)}

\prose{นเหตุํ นกุสลญฺจ นนเหตุํ นกุสลญฺจ ธมฺมํ ปฏิจฺจ นเหตุ นกุสโล ธมฺโม อุปฺปชฺชติ เหตุปจฺจยา.\csromanspacer{}\csromanspacer{}นเหตุํ นกุสลญฺจ นนเหตุํ นกุสลญฺจ ธมฺมํ ปฏิจฺจ นนเหตุ นกุสโล ธมฺโม อุปฺปชฺชติ เหตุปจฺจยา.\csromanspacer{}\csromanspacer{}นเหตุํ นกุสลญฺจ นนเหตุํ นกุสลญฺจ ธมฺมํ ปฏิจฺจ นเหตุ นกุสโล จ นนเหตุ นกุสโล จ ธมฺมา อุปฺปชฺชนฺติ เหตุปจฺจยา. (3) (สํขิตฺตํ.)}

\gambhira{ธมฺมปจฺจนีย ทุกติกปฏฺฐานปาฬิ}
\proseitem{2}{\csromanfolio{106}เหตุยา นว ฯเปฯ อวิคเต นว.}

\prose{นเหตุยา เทฺว.}

\prose{(สหชาตวาเรปิ ฯเปฯ สมฺปยุตฺตวาเรปิ สพฺพตฺถ วิตฺถาโร.)}

\csromanheader{2}{\textbf{เหตุ อารมฺมณ}}
\proseitem{3}{นนเหตุ นกุสโล ธมฺโม นนเหตุสฺส นกุสลสฺส ธมฺมสฺส เหตุปจฺจเยน ปจฺจโย.\csromanspacer{}\csromanspacer{}นนเหตุ นกุสโล ธมฺโม นเหตุสฺส นกุสลสฺส ธมฺมสฺส เหตุปจฺจเยน ปจฺจโย.\csromanspacer{}\csromanspacer{}นนเหตุ นกุสโล ธมฺโม นเหตุสฺส นกุสลสฺส จ นนเหตุสฺส นกุสลสฺส จ ธมฺมสฺส เหตุปจฺจเยน ปจฺจโย. (3)}

\prose{นเหตุ นกุสโล ธมฺโม นเหตุสฺส นกุสลสฺส ธมฺมสฺส อารมฺมณปจฺจเยน ปจฺจโย. (สํขิตฺตํ.)}

\proseitem{4}{เหตุยา ตีณิ, อารมฺมเณ นว ฯเปฯ อวิคเต นว. (ปญฺหาวารํ เอวํ วิตฺถาเรตพฺพํ.)}

\proseitem{5}{นเหตุํ น-อกุสลํ ธมฺมํ ปฏิจฺจ นเหตุ น-อกุสโล ธมฺโม อุปฺปชฺชติ เหตุปจฺจยา.\csromanspacer{} นเหตุํ กุสลํ อพฺยากตํ เอกํ ขนฺธํ ปฏิจฺจ ตโย ขนฺธา จิตฺตสมุฏฺฐานญฺจ รูปํ ฯเปฯ เทฺว ขนฺเธ ฯเปฯ.\csromanspacer{}\csromanspacer{}นเหตุํ น-อกุสลํ ธมฺมํ ปฏิจฺจ นนเหตุ น-อกุสโล ธมฺโม อุปฺปชฺชติ เหตุปจฺจยา.\csromanspacer{} นเหตู กุสเล อพฺยากเต ขนฺเธ ปฏิจฺจ เหตู.\csromanspacer{}\csromanspacer{}นเหตุํ น-อกุสลํ ธมฺมํ ปฏิจฺจ นเหตุ น-อกุสโล จ นนเหตุ น-อกุสโล จ ธมฺมา อุปฺปชฺชนฺติ เหตุปจฺจยา. (3)}

\prose{นนเหตุํ น-อกุสลํ ธมฺมํ ปฏิจฺจ นนเหตุ น-อกุสโล ธมฺโม อุปฺปชฺชติ เหตุปจฺจยา.\csromanspacer{}\csromanspacer{}นนเหตุํ น-อกุสลํ ธมฺมํ ปฏิจฺจ นเหตุ น-อกุสโล ธมฺโม อุปฺปชฺชติ เหตุปจฺจยา.\csromanspacer{}\csromanspacer{}นนเหตุํ น-อกุสลํ ธมฺมํ ปฏิจฺจ นเหตุ น-อกุสโล จ นนเหตุ น-อกุสโล จ ธมฺมา อุปฺปชฺชนฺติ เหตุปจฺจยา. (3)}

\prose{นเหตุํ น-อกุสลญฺจ นนเหตุํ น-อกุสลญฺจ ธมฺมํ ปฏิจฺจ นเหตุ น-อกุสโล ธมฺโม อุปฺปชฺชติ เหตุปจฺจยา.\csromanspacer{}\csromanspacer{}นเหตุํ น-อกุสลญฺจ นนเหตุํ น-อกุสลญฺจ ธมฺมํ ปฏิจฺจ นนเหตุ น-อกุสโล ธมฺโม อุปฺปชฺชติ}

\vspace{\tipitakaparskip}
\csromancenter{เหตุทุก 2.\csromanspacer{} เวทนาตฺติก เหตุปจฺจยา.\csromanspacer{}\csromanspacer{}นเหตุํ น-อกุสลญฺจ นนเหตุํ น-อกุสลญฺจ ธมฺมํ ปฏิจฺจ นเหตุ น-อกุสโล จ นนเหตุ น-อกุสโล จ ธมฺมา อุปฺปชฺชนฺติ เหตุปจฺจยา. (3) (สํขิตฺตํ.)}
\csromanheader{2}{เหตุยา นว ฯเปฯ อวิคเต นว. (สพฺพตฺถ นว.)}
\proseitem{6}{\csromanfolio{107}นเหตุํ น-อพฺยากตํ ธมฺมํ ปฏิจฺจ นเหตุ น-อพฺยากโต ธมฺโม อุปฺปชฺชติ เหตุปจฺจยา.\csromanspacer{} กุสลากุสลํ นเหตุํ เอกํ ขนฺธํ ปฏิจฺจ ตโย ขนฺธา.\csromanspacer{} ตโย ขนฺเธ ปฏิจฺจ เอโก ขนฺธา ฯเปฯ.\csromanspacer{}\csromanspacer{}นเหตุํ น-อพฺยากตํ ธมฺมํ ปฏิจฺจ นนเหตุ น-อพฺยากโต ธมฺโม อุปฺปชฺชติ เหตุปจฺจยา.\csromanspacer{}\csromanspacer{}นเหตุํ น-อพฺยากตํ ธมฺมํ ปฏิจฺจ นเหตุ น-อพฺยากโต จ นนเหตุ น-อพฺยากโต จ ธมฺมา อุปฺปชฺชนฺติ เหตุปจฺจยา. (3)}

\prose{นนเหตุํ น-อพฺยากตํ ธมฺมํ ปฏิจฺจ นนเหตุ น-อพฺยากโต ธมฺโม อุปฺปชฺชติ เหตุปจฺจยา.\csromanspacer{} ตีณิ.}

\prose{นเหตุํ น-อพฺยากตญฺจ นนเหตุํ น-อพฺยากตญฺจ ธมฺมํ ปฏิจฺจ นเหตุ น-อพฺยากโต ธมฺโม อุปฺปชฺชติ เหตุปจฺจยา.\csromanspacer{}\csromanspacer{}นเหตุํ น-อพฺยากตญฺจ นนเหตุํ น-อพฺยากตญฺจ ธมฺมํ ปฏิจฺจ นนเหตุ น-อพฺยากโต ธมฺโม อุปฺปชฺชติ เหตุปจฺจยา.\csromanspacer{}\csromanspacer{}นเหตุํ น-อพฺยากตญฺจ นนเหตุํ น-อพฺยากตญฺจ ธมฺมํ ปฏิจฺจ นเหตุ น-อพฺยากโต จ นนเหตุ น-อพฺยากโต จ ธมฺมา อุปฺปชฺชนฺติ เหตุปจฺจยา. (3) (สํขิตฺตํ.)}

\csromanheader{2}{เหตุยา นว ฯเปฯ อวิคเต นว. (สพฺพตฺถ นว.)}
\csromanheader{2}{1. เหตุทุก 2. เวทนาตฺติก}
\proseitem{7}{นเหตุํ นสุขาย เวทนาย สมฺปยุตฺตํ ธมฺมํ ปฏิจฺจ นเหตุ นสุขาย เวทนาย สมฺปยุตฺโต ธมฺโม อุปฺปชฺชติ เหตุปจฺจยา.\csromanspacer{}\csromanspacer{}นเหตุํ นสุขาย เวทนาย สมฺปยุตฺตํ ธมฺมํ ปฏิจฺจ นนเหตุ นสุขาย เวทนาย สมฺปยุตฺโต ธมฺโม อุปฺปชฺชติ เหตุปจฺจยา.}

\csromanheader{2}{เหตุยา นว. (สพฺพตฺถ นว.)}
\proseitem{8}{นเหตุํ นทุกฺขาย เวทนาย สมฺปยุตฺตํ ธมฺมํ ปฏิจฺจ นเหตุ นทุกฺขาย เวทนาย สมฺปยุตฺโต ธมฺโม อุปฺปชฺชติ เหตุปจฺจยา.}

\csromanheader{2}{เหตุยา นว. (สพฺพตฺถ นว.)}
\gambhira{ธมฺมปจฺจนีย ทุกติกปฏฺฐานปาฬิ}
\proseitem{9}{\csromanfolio{108}นเหตุํ น-อทุกฺขมสุขาย เวทนาย สมฺปยุตฺตํ ธมฺมํ ปฏิจฺจ นเหตุ น-อทุกฺขมสุขาย เวทนาย สมฺปยุตฺโต ธมฺโม อุปฺปชฺชติ เหตุปจฺจยา. (สํขิตฺตํ.)}

\csromanheader{2}{เหตุยา นว. (สพฺพตฺถ นว.)}
\csromanheader{2}{1. เหตุทุก 3-9. วิปากตฺติกาทิ}
\proseitem{10}{นเหตุํ นวิปากํ ธมฺมํ ปฏิจฺจ ฯเปฯ.}

\prose{นเหตุํ นวิปากธมฺมธมฺมํ ปฏิจฺจ ฯเปฯ.}

\prose{นเหตุํ นเนววิปากนวิปากธมฺมธมฺมํ ปฏิจฺจ ฯเปฯ.}

\proseitem{11}{นเหตุํ น-อุปาทินฺนุปาทานิยํ ธมฺมํ ปฏิจฺจ ฯเปฯ.}

\prose{นเหตุํ น-อนุปาทินฺนุปาทานิยํ ธมฺมํ ปฏิจฺจ ฯเปฯ.}

\prose{นเหตุํ น-อนุปาทินฺน-อนุปาทานิยํ ธมฺมํ ปฏิจฺจ ฯเปฯ.}

\proseitem{12}{นเหตุํ นสํกิลิฏฺฐสํกิเลสิกํ ธมฺมํ ปฏิจฺจ ฯเปฯ.}

\prose{นเหตุํ น-อสํกิลิฏฺฐสํกิเลสิกํ ธมฺมํ ปฏิจฺจ ฯเปฯ.}

\prose{นเหตุํ น-อสํกิลิฏฺฐ-อสํกิเลสิกํ ธมฺมํ ปฏิจฺจ ฯเปฯ.}

\proseitem{13}{นเหตุํ นสวิตกฺกสวิจารํ ธมฺมํ ปฏิจฺจ ฯเปฯ.}

\prose{นเหตุํ น-อวิตกฺกวิจารมตฺตํ ธมฺมํ ปฏิจฺจ ฯเปฯ.}

\prose{นเหตุํ น-อวิตกฺก-อวิจารํ ธมฺมํ ปฏิจฺจ ฯเปฯ.}

\proseitem{14}{นเหตุํ นปีติสหคตํ ธมฺมํ ปฏิจฺจ ฯเปฯ.}

\prose{นเหตุํ นสุขสหคตํ ธมฺมํ ปฏิจฺจ ฯเปฯ.}

\prose{นเหตุํ น-อุเปกฺขาสหคตํ ธมฺมํ ปฏิจฺจ ฯเปฯ.}

\proseitem{15}{นเหตุํ นทสฺสเนน ปหาตพฺพํ ธมฺมํ ปฏิจฺจ ฯเปฯ.}

\prose{นเหตุํ นภาวนาย ปหาตพฺพํ ธมฺมํ ปฏิจฺจ ฯเปฯ.}

\prose{นเหตุํ นเนวทสฺสเนน นภาวนาย ปหาตพฺพํ ธมฺมํ ปฏิจฺจ ฯเปฯ.}

\csromanheader{2}{1. เหตุทุก 10-21. อาจยคามิตฺติกาทิ}
\proseitem{16}{\csromanfolio{109}นเหตุํ นทสฺสเนน ปหาตพฺพเหตุกํ ธมฺมํ ปฏิจฺจ ฯเปฯ.}

\prose{นเหตุํ นภาวนาย ปหาตพฺพเหตุกํ ธมฺมํ ปฏิจฺจ ฯเปฯ.}

\prose{นเหตุํ นเนวทสฺสเนน นภาวนาย ปหาตพฺพเหตุกํ ธมฺมํ ปฏิจฺจ ฯเปฯ.}

\csromanheader{2}{1. เหตุทุก 10-21. อาจยคามิตฺติกาทิ}
\proseitem{17}{นเหตุํ น-อาจยคามิํ ธมฺมํ ปฏิจฺจ ฯเปฯ.}

\prose{นเหตุํ น-อปจยคามิํ ธมฺมํ ปฏิจฺจ ฯเปฯ.}

\prose{นเหตุํ นเนว-อาจยคามิํ นาปจยคามิํ ธมฺมํ ปฏิจฺจ ฯเปฯ.}

\proseitem{18}{นเหตุํ นเสกฺขํ ธมฺมํ ปฏิจฺจ ฯเปฯ.}

\prose{นเหตุํ น-อเสกฺขํ ธมฺมํ ปฏิจฺจ ฯเปฯ.}

\prose{นเหตุํ นเนวเสกฺขนาเสกฺขํ ธมฺมํ ปฏิจฺจ ฯเปฯ.}

\proseitem{19}{นเหตุํ นปริตฺตํ ธมฺมํ ปฏิจฺจ ฯเปฯ.}

\prose{นเหตุํ นมหคฺคตํ ธมฺมํ ปฏิจฺจ ฯเปฯ.}

\prose{นเหตุํ น-อปฺปมาณํ ธมฺมํ ปฏิจฺจ ฯเปฯ.}

\proseitem{20}{นเหตุํ นปริตฺตารมฺมณํ ธมฺมํ ปฏิจฺจ ฯเปฯ.}

\prose{นเหตุํ นมหคฺคตารมฺมณํ ธมฺมํ ปฏิจฺจ ฯเปฯ.}

\prose{นเหตุํ น-อปฺปมาณารมฺมณํ ธมฺมํ ปฏิจฺจ ฯเปฯ.}

\proseitem{21}{นเหตุํ นหีนํ ธมฺมํ ปฏิจฺจ ฯเปฯ.}

\prose{นเหตุํ นมชฺฌิมํ ธมฺมํ ปฏิจฺจ ฯเปฯ.}

\prose{นเหตุํ นปณีตํ ธมฺมํ ปฏิจฺจ ฯเปฯ.}

\proseitem{22}{นเหตุํ นมิจฺฉตฺตนิยตํ ธมฺมํ ปฏิจฺจ ฯเปฯ.}

\prose{นเหตุํ นสมฺมตฺตนิยตํ ธมฺมํ ปฏิจฺจ ฯเปฯ.}

\prose{นเหตุํ น-อนิยตํ ธมฺมํ ปฏิจฺจ ฯเปฯ.}

\gambhira{ธมฺมปจฺจนีย ทุกติกปฏฺฐานปาฬิ}
\proseitem{23}{\csromanfolio{110}นเหตุํ นมคฺคารมฺมณํ ธมฺมํ ปฏิจฺจ ฯเปฯ.}

\prose{นเหตุํ นมคฺคเหตุกํ ธมฺมํ ปฏิจฺจ ฯเปฯ.}

\prose{นเหตุํ นมคฺคาธิปติํ ธมฺมํ ปฏิจฺจ ฯเปฯ.}

\proseitem{24}{นเหตุํ น-อนุปฺปนฺนํ ธมฺมํ ปฏิจฺจ ฯเปฯ.}

\prose{นเหตุํ น-อุปฺปาทิํ ธมฺมํ ปฏิจฺจ ฯเปฯ.}

\proseitem{25}{นเหตุํ น-อตีตํ ธมฺมํ ปฏิจฺจ ฯเปฯ.}

\prose{นเหตุํ น-อนาคตํ ธมฺมํ ปฏิจฺจ ฯเปฯ.}

\proseitem{26}{นเหตุํ น-อตีตารมฺมณํ ธมฺมํ ปฏิจฺจ ฯเปฯ.}

\prose{นเหตุํ น-อนาคตารมฺมณํ ธมฺมํ ปฏิจฺจ ฯเปฯ.}

\prose{นเหตุํ นปจฺจุปฺปนฺนารมฺมณํ ธมฺมํ ปฏิจฺจ ฯเปฯ.}

\proseitem{27}{นเหตุํ น-อชฺฌตฺตํ ธมฺมํ ปฏิจฺจ ฯเปฯ.}

\prose{นเหตุํ นพหิทฺธา ธมฺมํ ปฏิจฺจ ฯเปฯ.}

\proseitem{28}{นเหตุํ น-อชฺฌตฺตารมฺมณํ ธมฺมํ ปฏิจฺจ ฯเปฯ.}

\prose{นเหตุํ นพหิทฺธารมฺมณํ ธมฺมํ ปฏิจฺจ ฯเปฯ.}

\csromanheader{2}{1. เหตุทุก 22. สนิทสฺสนตฺติก}
\csromanheader{2}{1-7. ปฏิจฺจาทิวาร ปจฺจยจตุกฺก-เหตุ}
\proseitem{29}{นเหตุํ นสนิทสฺสนสปฺปฏิฆํ ธมฺมํ ปฏิจฺจ นเหตุ นสนิทสฺสนสปฺปฏิโฆ ธมฺโม อุปฺปชฺชติ เหตุปจฺจยา.\csromanspacer{}\csromanspacer{}นเหตุํ นสนิทสฺสนสปฺปฏิฆํ ธมฺมํ ปฏิจฺจ นนเหตุ นสนิทสฺสนสปฺปฏิโฆ ธมฺโม อุปฺปชฺชติ เหตุปจฺจยา.\csromanspacer{}\csromanspacer{}นเหตุํ นสนิทสฺสนสปฺปฏิฆํ ธมฺมํ ปฏิจฺจ นเหตุ นสนิทสฺสนสปฺปฏิโฆ จ นนเหตุ นสนิทสฺสนสปฺปฏิโฆ จ ธมฺมา อุปฺปชฺชนฺติ เหตุปจฺจยา. (3)}

\prose{นนเหตุํ นสนิทสฺสนสปฺปฏิฆํ ธมฺมํ ปฏิจฺจ นนเหตุ นสนิทสฺสนสปฺปฏิโฆ ธมฺโม อุปฺปชฺชติ เหตุปจฺจยา.\csromanspacer{}\csromanspacer{}นนเหตุํ นสนิทสฺสนสปฺปฏิฆํ ธมฺมํ}

\vspace{\tipitakaparskip}
\csromancenter{เหตุทุก 22.\csromanspacer{} สนิทสฺสนตฺติก ปฏิจฺจ นเหตุ นสนิทสฺสนสปฺปฏิโฆ ธมฺโม อุปฺปชฺชติ เหตุปจฺจยา.\csromanspacer{}\csromanspacer{}นนเหตุํ นสนิทสฺสนสปฺปฏิฆํ ธมฺมํ ปฏิจฺจ นเหตุ นสนิทสฺสนสปฺปฏิโฆ จ นนเหตุ นสนิทสฺสนสปฺปฏิโฆ จ ธมฺมา อุปฺปชฺชนฺติ เหตุปจฺจยา. (3)}
\prose{\csromanfolio{111}นเหตุํ นสนิทสฺสนสปฺปฏิฆญฺจ นนเหตุํ นสนิทสฺสนสปฺปฏิฆญฺจ ธมฺมํ ปฏิจฺจ นเหตุ นสนิทสฺสนสปฺปฏิโฆ ธมฺโม อุปฺปชฺชติ เหตุปจฺจยา.\csromanspacer{}\csromanspacer{}นเหตุํ นสนิทสฺสนสปฺปฏิฆญฺจ นนเหตุํ นสนิทสฺสนสปฺปฏิฆญฺจ ธมฺมํ ปฏิจฺจ นนเหตุ นสนิทสฺสนสปฺปฏิโฆ ธมฺโม อุปฺปชฺชติ เหตุปจฺจยา.\csromanspacer{}\csromanspacer{}นเหตุํ นสนิทสฺสนสปฺปฏิฆญฺจ นนเหตุํ นสนิทสฺสนสปฺปฏิฆญฺจ ธมฺมํ ปฏิจฺจ นเหตุ นสนิทสฺสนสปฺปฏิโฆ จ นนเหตุ นสนิทสฺสนสปฺปฏิโฆ จ ธมฺมา อุปฺปชฺชนฺติ เหตุปจฺจยา. (3) (สํขิตฺตํ.)}

\csromanheader{2}{เหตุยา นว ฯเปฯ อวิคเต นว. (สพฺพตฺถ นว.)}
\proseitem{30}{นเหตุํ น-อนิทสฺสนสปฺปฏิฆํ ธมฺมํ ปฏิจฺจ นเหตุ น-อนิทสฺสนสปฺปฏิโฆ ธมฺโม อุปฺปชฺชติ เหตุปจฺจยา.\csromanspacer{}\csromanspacer{}นเหตุํ น-อนิทสฺสนสปฺปฏิฆํ ธมฺมํ ปฏิจฺจ นนเหตุ น-อนิทสฺสนสปฺปฏิโฆ ธมฺโม อุปฺปชฺชติ เหตุปจฺจยา.\csromanspacer{}\csromanspacer{}นเหตุํ น-อนิทสฺสนสปฺปฏิฆํ ธมฺมํ ปฏิจฺจ นเหตุ น-อนิทสฺสนสปฺปฏิโฆ จ นนเหตุ น-อนิทสฺสนสปฺปฏิโฆ จ ธมฺมา อุปฺปชฺชนฺติ เหตุปจฺจยา. (3)}

\prose{นนเหตุํ น-อนิทสฺสนสปฺปฏิฆํ ธมฺมํ ปฏิจฺจ นนเหตุ น-อนิทสฺสนสปฺปฏิโฆ ธมฺโม อุปฺปชฺชติ เหตุปจฺจยา.\csromanspacer{}\csromanspacer{}นนเหตุํ น-อนิทสฺสนสปฺปฏิฆํ ธมฺมํ ปฏิจฺจ นเหตุ น-อนิทสฺสนสปฺปฏิโฆ ธมฺโม อุปฺปชฺชติ เหตุปจฺจยา.\csromanspacer{}\csromanspacer{}นนเหตุํ น-อนิทสฺสนสปฺปฏิฆํ ธมฺมํ ปฏิจฺจ นเหตุ น-อนิทสฺสนสปฺปฏิโฆ จ นนเหตุ น-อนิทสฺสนสปฺปฏิโฆ จ ธมฺมา อุปฺปชฺชนฺติ เหตุปจฺจยา. (3)}

\prose{นเหตุํ น-อนิทสฺสนสปฺปฏิฆญฺจ นนเหตุํ น-อนิทสฺสนสปฺปฏิฆญฺจ ธมฺมํ ปฏิจฺจ นเหตุ น-อนิทสฺสนสปฺปฏิโฆ ธมฺโม อุปฺปชฺชติ เหตุปจฺจยา.\csromanspacer{}\csromanspacer{}นเหตุํ น-อนิทสฺสนสปฺปฏิฆญฺจ นนเหตุํ น-อนิทสฺสนสปฺปฏิฆญฺจ ธมฺมํ ปฏิจฺจ นนเหตุ น-อนิทสฺสนสปฺปฏิโฆ ธมฺโม อุปฺปชฺชติ เหตุปจฺจยา.\csromanspacer{}\csromanspacer{}นเหตุํ น-อนิทสฺสนสปฺปฏิฆญฺจ นนเหตุํ น-อนิทสฺสนสปฺปฏิฆญฺจ ธมฺมํ ปฏิจฺจ นเหตุ น-อนิทสฺสนสปฺปฏิโฆ จ นนเหตุ น-อนิทสฺสนสปฺปฏิโฆ จ ธมฺมา อุปฺปชฺชนฺติ เหตุปจฺจยา. (3) (สํขิตฺตํ.)}

\prose{เหตุยา นว ฯเปฯ อวิคเต นว. (สหชาตวาเรปิ.\csromanspacer{} สมฺปยุตฺตวาเรปิ สพฺพตฺถ วิตฺถาโร.)}

\csromanheader{2}{ธมฺมปจฺจนีย ทุกติกปฏฺฐานปาฬิ \textbf{เหตุ-อารมฺมณ}}
\proseitem{31}{\csromanfolio{112}นนเหตุ น-อนิทสฺสนสปฺปฏิโฆ ธมฺโม นนเหตุสฺส น-อนิทสฺสนสปฺปฏิฆสฺส ธมฺมสฺส เหตุปจฺจเยน ปจฺจโย.\csromanspacer{}\csromanspacer{}นนเหตุ น-อนิทสฺสนสปฺปฏิโฆ ธมฺโม นเหตุสฺส น-อนิทสฺสนสปฺปฏิฆสฺส ธมฺมสฺส เหตุปจฺจเยน ปจฺจโย.\csromanspacer{}\csromanspacer{}นนเหตุ น-อนิทสฺสนสปฺปฏิโฆ ธมฺโม นเหตุสฺส น-อนิทสฺสนสปฺปฏิฆสฺส จ นนเหตุสฺส น-อนิทสฺสนสปฺปฏิฆสฺส จ ธมฺมสฺส เหตุปจฺจเยน ปจฺจโย. (3)}

\prose{นเหตุ น-อนิทสฺสนสปฺปฏิโฆ ธมฺโม นเหตุสฺส น-อนิทสฺสนสปฺปฏิฆสฺส ธมฺมสฺส อารมฺมณปจฺจเยน ปจฺจโย.}

\prose{เหตุยา ตีณิ, อารมฺมเณ นว ฯเปฯ อุปนิสฺสเย นว, ปุเรชาเต ปจฺฉาชาเต ตีณิ, อาเสวเน นว, กมฺเม ตีณิ, วิปาเก นว, อาหาเร ตีณิ, อินฺทฺริเย นว, ฌาเน ตีณิ, มคฺเค สมฺปยุตฺเต นว, วิปฺปยุตฺเต ปญฺจ ฯเปฯ อวิคเต นว. (ปญฺหาวารํ เอวํ วิตฺถาเรตพฺพํ.)}

\proseitem{32}{นเหตุํ น-อนิทสฺสน-อปฺปฏิฆํ ธมฺมํ ปฏิจฺจ นเหตุ น-อนิทสฺสน-อปฺปฏิโฆ ธมฺโม อุปฺปชฺชติ เหตุปจฺจยา. (1) (สํขิตฺตํ.)}

\csromanheader{2}{เหตุยา เอกํ. (สพฺพตฺถ วิตฺถาโร.)}
\csromanheader{2}{2. สเหตุกทุก 1. กุสลตฺติก}
\proseitem{33}{นสเหตุกํ นกุสลํ ธมฺมํ ปฏิจฺจ นสเหตุโก นกุสโล ธมฺโม อุปฺปชฺชติ เหตุปจฺจยา.\csromanspacer{} วิจิกิจฺฉาสหคตํ อุทฺธจฺจสหคตํ โมหํ ปฏิจฺจ จิตฺตสมุฏฺฐานํ รูปํ.\csromanspacer{} เอกํ มหาภูตํ ปฏิจฺจ ตโย มหาภูตา ฯเปฯ เทฺว มหาภูเต ฯเปฯ.\csromanspacer{}\csromanspacer{}นสเหตุกํ นกุสลํ ธมฺมํ ปฏิจฺจ น-อเหตุโก นกุสโล ธมฺโม อุปฺปชฺชติ เหตุปจฺจยา.\csromanspacer{} วิจิกิจฺฉาสหคตํ อุทฺธจฺจสหคตํ โมหํ ปฏิจฺจ สมฺปยุตฺตกา ขนฺธา.\csromanspacer{} ปฏิสนฺธิกฺขเณ วตฺถุํ ปฏิจฺจ สเหตุกา ขนฺธา.\csromanspacer{}\csromanspacer{}นสเหตุกํ นกุสลํ ธมฺมํ ปฏิจฺจ นสเหตุโก นกุสโล จ น-อเหตุโก นกุสโล จ ธมฺมา อุปฺปชฺชนฺติ เหตุปจฺจยา. (3)}

\prose{น-อเหตุกํ นกุสลํ ธมฺมํ น-อเหตุโก นกุสโล ธมฺโม อุปฺปชฺชติ เหตุปจฺจยา.\csromanspacer{} ตีณิ.}

\csromanheader{2}{4. เหตุสเหตุกทุก 1. กุสลตฺติก}
\prose{\csromanfolio{113}นสเหตุกํ นกุสลญฺจ น-อเหตุกํ นกุสลญฺจ ธมฺมํ ปฏิจฺจ นสเหตุโก นกุสโล ธมฺโม อุปฺปชฺชติ เหตุปจฺจยา.\csromanspacer{} ตีณิ. (สํขิตฺตํ.)}

\prose{เหตุยา นว.}

\proseitem{34}{นสเหตุกํ น-อกุสลํ ธมฺมํ ปฏิจฺจ นสเหตุโก น-อกุสโล ธมฺโม อุปฺปชฺชติ เหตุปจฺจยา.\csromanspacer{} ตีณิ.}

\prose{น-อเหตุกํ น-อกุสลํ ธมฺมํ ปฏิจฺจ น-อเหตุโก น-อกุสโล ธมฺโม อุปฺปชฺชติ เหตุปจฺจยา.\csromanspacer{} ตีณิ.}

\prose{นสเหตุกํ น-อกุสลญฺจ น-อเหตุกํ น-อกุสลญฺจ ธมฺมํ ปฏิจฺจ นสเหตุโก น-อกุสโล ธมฺโม อุปฺปชฺชติ เหตุปจฺจยา.\csromanspacer{} ตีณิ.}

\prose{เหตุยา นว.}

\proseitem{35}{นสเหตุกํ น-อพฺยากตํ ธมฺมํ ปฏิจฺจ น-อเหตุโก น-อพฺยากโต ธมฺโม อุปฺปชฺชติ เหตุปจฺจยา. (1)}

\prose{น-อเหตุกํ น-อพฺยากตํ ธมฺมํ ปฏิจฺจ น-อเหตุโก น-อพฺยากโต}

\prose{นสเหตุกํ น-อพฺยากตญฺจ น-อเหตุกํ น-อพฺยากตญฺจ ธมฺมํ ปฏิจฺจ น-อเหตุโก น-อพฺยากโต ธมฺโม อุปฺปชฺชติ เหตุปจฺจยา. (1)}

\csromanheader{2}{เหตุยา ตีณิ. (เหตุสมฺปยุตฺตทุกํ สเหตุกทุกสทิสํ.)}
\csromanheader{2}{4. เหตุสเหตุกทุก 1. กุสลตฺติก}
\proseitem{36}{นเหตุญฺเจว น-อเหตุกญฺจ นกุสลํ ธมฺมํ ปฏิจฺจ นเหตุ เจว น-อเหตุโก จ นกุสโล ธมฺโม อุปฺปชฺชติ เหตุปจฺจยา.\csromanspacer{}\csromanspacer{}นเหตุญฺเจว น-อเหตุกญฺจ นกุสลํ ธมฺมํ ปฏิจฺจ น-อเหตุโก เจว นน จ เหตุ นกุสโล ธมฺโม อุปฺปชฺชติ เหตุปจฺจยา.\csromanspacer{}\csromanspacer{}นเหตุญฺเจว น-อเหตุกญฺจ นกุสลํ ธมฺมํ ปฏิจฺจ นเหตุ เจว น-อเหตุโก จ นกุสโล จ น-อเหตุโก เจว นน จ เหตุ นกุสโล จ ธมฺมา อุปฺปชฺชนฺติ เหตุปจฺจยา. (3)}

\gambhira{ธมฺมปจฺจนีย ทุกติกปฏฺฐานปาฬิ}
\prose{\csromanfolio{114}น-อเหตุกญฺเจว นน จ เหตุํ นกุสลํ ธมฺมํ ปฏิจฺจ น-อเหตุโก เจว นน จ เหตุ นกุสโล ธมฺโม อุปฺปชฺชติ เหตุปจฺจยา.\csromanspacer{} ตีณิ.}

\prose{นเหตุญฺเจว น-อเหตุกญฺจ นกุสลญฺจ น-อเหตุกญฺเจว นน จ เหตุํ นกุสลญฺจ ธมฺมํ ปฏิจฺจ นเหตุ เจว น-อเหตุโก จ นกุสโล ธมฺโม อุปฺปชฺชติ เหตุปจฺจยา.\csromanspacer{} ตีณิ. (สํขิตฺตํ.)}

\prose{เหตุยา นว.}

\proseitem{37}{นเหตุญฺเจว น-อเหตุกญฺจ น-อกุสลํ ธมฺมํ ปฏิจฺจ นเหตุ เจว น-อเหตุโก จ น-อกุสโล ธมฺโม อุปฺปชฺชติ เหตุปจฺจยา.}

\prose{เหตุยา นว.}

\proseitem{38}{นเหตุญฺเจว น-อเหตุกญฺจ น-อพฺยากตํ ธมฺมํ ปฏิจฺจ นเหตุ เจว น-อเหตุโก จ น-อพฺยากโต ธมฺโม อุปฺปชฺชติ เหตุปจฺจยา.}

\prose{เหตุยา นว. (เหตุเหตุสมฺปยุตฺตทุกํ เหตุสเหตุกทุกสทิสํ.)}

\csromanheader{2}{6. นเหตุสเหตุกทุก 1. กุสลตฺติก}
\proseitem{39}{นเหตุํ นสเหตุกํ นกุสลํ ธมฺมํ ปฏิจฺจ นเหตุ นสเหตุโก นกุสโล ธมฺโม อุปฺปชฺชติ เหตุปจฺจยา. (สํขิตฺตํ.)}

\prose{เหตุยา นว.}

\proseitem{40}{นเหตุํ นสเหตุกํ น-อกุสลํ ธมฺมํ ปฏิจฺจ นเหตุ นสเหตุโก น-อกุสโล ธมฺโม อุปฺปชฺชติ เหตุปจฺจยา. (สํขิตฺตํ.)}

\prose{เหตุยา นว.}

\proseitem{41}{นเหตุํ น-อเหตุกํ น-อพฺยากตํ ธมฺมํ ปฏิจฺจ นเหตุ น-อเหตุโก น-อพฺยากโต ธมฺโม อุปฺปชฺชติ เหตุปจฺจยา. (เอกปญฺหเมว.\csromanspacer{} เอเตน อุปาเยน กุสลากุสลทุกํ กุสลตฺติกเมว เอตฺถ วฏฺฏติ.)}

\csromanheader{2}{14. อาสวโคจฺฉก 1. กุสลตฺติก \textbf{7. จูฬนฺตรทุก 1. กุสลตฺติก}}
\proseitem{42}{\csromanfolio{115}น-อปจฺจยํ นกุสลํ ธมฺมํ ปฏิจฺจ ฯเปฯ. (สพฺพตฺถ เอกํ.) น-อสงฺขตํ นกุสลํ ธมฺมํ ปฏิจฺจ ฯเปฯ.}

\proseitem{43}{นสนิทสฺสนํ นกุสลํ ธมฺมํ ปฏิจฺจ.}

\proseitem{44}{นสปฺปฏิฆํ นกุสลํ ธมฺมํ ปฏิจฺจ ฯเปฯ.\csromanspacer{}\csromanspacer{}น-อปฺปฏิฆํ นกุสลํ ธมฺมํ ปฏิจฺจ.}

\proseitem{45}{นรูปิํ นกุสลํ ธมฺมํ ปฏิจฺจ ฯเปฯ.\csromanspacer{}\csromanspacer{}น-อรูปิํ นกุสลํ ธมฺมํ ปฏิจฺจ.}

\proseitem{46}{นโลกิยํ นกุสลํ ธมฺมํ ปฏิจฺจ ฯเปฯ.\csromanspacer{}\csromanspacer{}นโลกุตฺตรํ นกุสลํ ธมฺมํ ปฏิจฺจ.}

\proseitem{47}{นเกนจิ วิญฺเญยฺยํ นกุสลํ ธมฺมํ ปฏิจฺจ ฯเปฯ.\csromanspacer{}\csromanspacer{}นนเกนจิ วิญฺเญยฺยํ นกุสลํ ธมฺมํ ปฏิจฺจ.}

\csromanheader{2}{14. อาสวโคจฺฉก 1. กุสลตฺติก}
\proseitem{48}{น-อาสวํ นกุสลํ ธมฺมํ ปฏิจฺจ ฯเปฯ.\csromanspacer{}\csromanspacer{}นโน-อาสวํ นกุสลํ ธมฺมํ ปฏิจฺจ.}

\proseitem{49}{นสาสวํ นกุสลํ ธมฺมํ ปฏิจฺจ ฯเปฯ.\csromanspacer{}\csromanspacer{}น-อนาสวํ กุสลํ ธมฺมํ ปฏิจฺจ.}

\proseitem{50}{น-อาสวสมฺปยุตฺตํ นกุสลํ ธมฺมํ ปฏิจฺจ ฯเปฯ.\csromanspacer{}\csromanspacer{}น-อาสววิปฺปยุตฺตํ นกุสลํ ธมฺมํ ปฏิจฺจ.}

\proseitem{51}{น-อาสวญฺเจว น-อนาสวญฺจ นกุสลํ ธมฺมํ ปฏิจฺจ ฯเปฯ.\csromanspacer{}\csromanspacer{}น-อนาสวญฺเจว นโน จ อาสวํ นกุสลํ ธมฺมํ ปฏิจฺจ.}

\proseitem{52}{น-อาสวญฺเจว น-อาสววิปฺปยุตฺตญฺจ นกุสลํ ธมฺมํ ปฏิจฺจ ฯเปฯ.\csromanspacer{}\csromanspacer{}น-อาสววิปฺปยุตฺตญฺเจว นโน จ อาสวํ นกุสลํ ธมฺมํ ปฏิจฺจ.}

\proseitem{53}{อาสววิปฺปยุตฺตํ นสาสวํ นกุสลํ ธมฺมํ ปฏิจฺจ ฯเปฯ.\csromanspacer{}\csromanspacer{}อาสววิปฺปยุตฺตํ น-อนาสวํ นกุสลํ ธมฺมํ ปฏิจฺจ.}

\csromanheader{2}{ธมฺมปจฺจนีย ทุกติกปฏฺฐานปาฬิ \textbf{20-54. ฉโคจฺฉก 1. กุสลตฺติก}}
\proseitem{54}{\csromanfolio{116}โนสญฺโญชนํ นกุสลํ ธมฺมํ ปฏิจฺจ ฯเปฯ.\csromanspacer{}\csromanspacer{}นโนสญฺโญชนํ นกุสลํ ธมฺมํ ปฏิจฺจ.}

\proseitem{55}{โนคนฺถํ นกุสลํ ธมฺมํ ปฏิจฺจ ฯเปฯ.\csromanspacer{}\csromanspacer{}โน-โอฆํ ฯเปฯ.\csromanspacer{}\csromanspacer{}โนโยคํ ฯเปฯ.\csromanspacer{}\csromanspacer{}โนนีวรณํ.}

\prose{โนปรามาสํ นกุสลํ ธมฺมํ ปฏิจฺจ.}

\csromanheader{2}{55. มหนฺตรทุก 1. กุสลตฺติก}
\proseitem{56}{นสารมฺมณํ นกุสลํ ธมฺมํ ปฏิจฺจ ฯเปฯ.\csromanspacer{}\csromanspacer{}น-อนารมฺมณํ นกุสลํ ธมฺมํ ปฏิจฺจ.}

\proseitem{57}{นจิตฺตํ นกุสลํ ธมฺมํ ปฏิจฺจ ฯเปฯ.\csromanspacer{}\csromanspacer{}นโนจิตฺตํ นกุสลํ ธมฺมํ ปฏิจฺจ.}

\csromanheader{2}{58. นเจตสิกํ นกุสลํ ธมฺมํ ปฏิจฺจ. (สํขิตฺตํ.)}
\csromanheader{2}{59. นจิตฺตสมฺปยุตฺตํ นกุสลํ ธมฺมํ ปฏิจฺจ. (สํขิตฺตํ.)}
\proseitem{60}{นจิตฺตสํสฏฺฐํ นกุสลํ ธมฺมํ ปฏิจฺจ.}

\proseitem{61}{โนจิตฺตสมุฏฺฐานํ นกุสลํ ธมฺมํ ปฏิจฺจ.}

\proseitem{62}{โนจิตฺตสหภุํ นกุสลํ ธมฺมํ ปฏิจฺจ.}

\proseitem{63}{โนจิตฺตานุปริวตฺติํ นกุสลํ ธมฺมํ ปฏิจฺจ ฯเปฯ.\csromanspacer{}\csromanspacer{}โนจิตฺตสํสฏฺฐสมุฏฺฐานํ นกุสลํ ธมฺมํ ปฏิจฺจ ฯเปฯ.\csromanspacer{}\csromanspacer{}โนจิตฺตสํสฏฺฐสมุฏฺฐานสหภุํ นกุสลํ ธมฺมํ ปฏิจฺจ ฯเปฯ.\csromanspacer{}\csromanspacer{}โนจิตฺตสํสฏฺฐสมุฏฺฐานานุปริวตฺติํ นกุสลํ ธมฺมํ ปฏิจฺจ.}

\proseitem{64}{น-อชฺฌตฺติกํ นกุสลํ ธมฺมํ ปฏิจฺจ ฯเปฯ.\csromanspacer{}\csromanspacer{}นพาหิรํ นกุสลํ ธมฺมํ ปฏิจฺจ.}

\proseitem{65}{น-อุปาทา นกุสลํ ธมฺมํ ปฏิจฺจ.}

\proseitem{66}{น-อุปาทินฺนํ นกุสลํ ธมฺมํ ปฏิจฺจ.}

\csromanheader{2}{83. ปิฏฺฐิทุก 1. กุสลตฺติก \textbf{69. อุปาทานโคจฺฉก 1. กุสลตฺติก}}
\vspace{\tipitakaparskip}
\csromancenter{โน-อุปาทานํ นกุสลํ ธมฺมํ ปฏิจฺจ.}
\csromanheader{2}{75. กิเลสโคจฺฉก 1. กุสลตฺติก}
\proseitem{68}{\csromanfolio{117}โนกิเลสํ นกุสลํ ธมฺมํ ปฏิจฺจ ฯเปฯ.\csromanspacer{}\csromanspacer{}นโนกิเลสํ นกุสลํ ธมฺมํ ปฏิจฺจ.}

\proseitem{69}{นสํกิเลสิกํ นกุสลํ ธมฺมํ ปฏิจฺจ ฯเปฯ.\csromanspacer{}\csromanspacer{}น-อสํกิเลสิกํ นกุสลํ ธมฺมํ ปฏิจฺจ.}

\proseitem{70}{นสํกิลิฏฺฐํ นกุสลํ ธมฺมํ ปฏิจฺจ ฯเปฯ.\csromanspacer{}\csromanspacer{}น-อสํกิลิฏฺฐํ นกุสลํ ธมฺมํ ปฏิจฺจ.}

\proseitem{71}{นกิเลสสมฺปยุตฺตํ นกุสลํ ธมฺมํ ปฏิจฺจ ฯเปฯ.\csromanspacer{}\csromanspacer{}นกิเลสวิปฺปยุตฺตํ นกุสลํ ธมฺมํ ปฏิจฺจ.}

\proseitem{72}{นกิเลสญฺเจว น-อสํกิเลสิกญฺจ นกุสลํ ธมฺมํ ปฏิจฺจ ฯเปฯ.\csromanspacer{}\csromanspacer{}น-อสํกิเลสิกญฺเจว นโน จ กิเลสํ นกุสลํ ธมฺมํ ปฏิจฺจ.}

\proseitem{73}{นกิเลสญฺเจว น-อสํกิลิฏฺฐญฺจ นกุสลํ ธมฺมํ ปฏิจฺจ ฯเปฯ.\csromanspacer{} น-อสํกิลิฏฺฐญฺเจว นโน จ กิเลสํ นกุสลํ ธมฺมํ ปฏิจฺจ ฯเปฯ.\csromanspacer{} นกิเลสญฺเจว นกิเลสวิปฺปยุตฺตญฺจ นกุสลํ ธมฺมํ ปฏิจฺจ ฯเปฯ.\csromanspacer{} นกิเลสวิปฺปยุตฺตญฺเจว นโน จ กิเลสํ นกุสลํ ธมฺมํ ปฏิจฺจ ฯเปฯ. (สพฺพตฺถ นว.\csromanspacer{} วิปากํ นตฺถิ.)}

\prose{กิเลสวิปฺปยุตฺตํ นสํกิเลสิกํ นกุสลํ ธมฺมํ ปฏิจฺจ ฯเปฯ.\csromanspacer{} กิเลสวิปฺปยุตฺตํ น-อสํกิเลสิกํ นกุสลํ ธมฺมํ ปฏิจฺจ.}

\csromanheader{2}{83. ปิฏฺฐิทุก 1. กุสลตฺติก}
\proseitem{74}{นทสฺสเนน ปหาตพฺพํ นกุสลํ ธมฺมํ ปฏิจฺจ ฯเปฯ.\csromanspacer{}\csromanspacer{}นนทสฺสเนน ปหาตพฺพํ นกุสลํ ธมฺมํ ปฏิจฺจ.}

\proseitem{75}{นภาวนาย ปหาตพฺพํ นกุสลํ ธมฺมํ ปฏิจฺจ ฯเปฯ.\csromanspacer{}\csromanspacer{}นนภาวนาย ปหาตพฺพํ นกุสลํ ธมฺมํ ปฏิจฺจ.}

\gambhira{ธมฺมปจฺจนีย ทุกติกปฏฺฐานปาฬิ}
\proseitem{76}{\csromanfolio{118}นทสฺสเนน ปหาตพฺพเหตุกํ นกุสลํ ธมฺมํ ปฏิจฺจ ฯเปฯ.\csromanspacer{} นนทสฺสเนน ปหาตพฺพเหตุกํ นกุสลํ ธมฺมํ ปฏิจฺจ ฯเปฯ.\csromanspacer{} นภาวนาย ปหาตพฺพเหตุกํ นกุสลํ ธมฺมํ ปฏิจฺจ ฯเปฯ.\csromanspacer{} นนภาวนาย ปหาตพฺพเหตุกํ นกุสลํ ธมฺมํ ปฏิจฺจ.}

\proseitem{77}{นสวิตกฺกํ นกุสลํ ธมฺมํ ปฏิจฺจ ฯเปฯ.\csromanspacer{} น-อวิตกฺกํ นกุสลํ ธมฺมํ ปฏิจฺจ ฯเปฯ.\csromanspacer{} นสวิจารํ นกุสลํ ธมฺมํ ปฏิจฺจ ฯเปฯ.\csromanspacer{} น-อวิจารํ นกุสลํ ธมฺมํ ปฏิจฺจ ฯเปฯ.\csromanspacer{} นสปฺปีติกํ นกุสลํ ธมฺมํ ปฏิจฺจ ฯเปฯ.\csromanspacer{} น-อปฺปีติกํ นกุสลํ ธมฺมํ ปฏิจฺจ ฯเปฯ.\csromanspacer{} นปีติสหคตํ นกุสลํ ธมฺมํ ปฏิจฺจ ฯเปฯ.\csromanspacer{} นนปีติสหคตํ นกุสลํ ธมฺมํ ปฏิจฺจ ฯเปฯ.\csromanspacer{} นสุขสหคตํ นกุสลํ ธมฺมํ ปฏิจฺจ ฯเปฯ.\csromanspacer{} นนสุขสหคตํ นกุสลํ ธมฺมํ ปฏิจฺจ ฯเปฯ.\csromanspacer{} น-อุเปกฺขาสหคตํ นกุสลํ ธมฺมํ ปฏิจฺจ ฯเปฯ.\csromanspacer{} นน-อุเปกฺขาสหคตํ นกุสลํ ธมฺมํ ปฏิจฺจ.}

\proseitem{78}{นกามาวจรํ นกุสลํ ธมฺมํ ปฏิจฺจ ฯเปฯ.\csromanspacer{} นนกามาวจรํ นกุสลํ ธมฺมํ ปฏิจฺจ ฯเปฯ.\csromanspacer{} นรูปาวจรํ นกุสลํ ธมฺมํ ปฏิจฺจ ฯเปฯ.\csromanspacer{} นนรูปาวจรํ นกุสลํ ธมฺมํ ปฏิจฺจ.}

\proseitem{79}{น-อรูปาวจรํ นกุสลํ ธมฺมํ ปฏิจฺจ ฯเปฯ.\csromanspacer{} นน-อรูปาวจรํ นกุสลํ ธมฺมํ ปฏิจฺจ.}

\proseitem{80}{นปริยาปนฺนํ นกุสลํ ธมฺมํ ปฏิจฺจ ฯเปฯ.\csromanspacer{} น-อปริยาปนฺนํ นกุสลํ ธมฺมํ ปฏิจฺจ.}

\proseitem{81}{นนิยฺยานิกํ นกุสลํ ธมฺมํ ปฏิจฺจ ฯเปฯ.\csromanspacer{} น-อนิยฺยานิกํ นกุสลํ ธมฺมํ ปฏิจฺจ. (สํขิตฺตํ.\csromanspacer{} สพฺพตฺถ เอกํ.)}

\prose{นนิยตํ นกุสลํ ธมฺมํ ปฏิจฺจ ฯเปฯ.\csromanspacer{} น-อนิยตํ นกุสลํ ธมฺมํ ปฏิจฺจ. (สํขิตฺตํ.)}

\proseitem{82}{นส-อุตฺตรํ นกุสลํ ธมฺมํ ปฏิจฺจ ฯเปฯ.\csromanspacer{} น-อนุตฺตรํ นกุสลํ ธมฺมํ ปฏิจฺจ. (สํขิตฺตํ.)}

\csromanheader{2}{83. นสรณํ นกุสลํ ธมฺมํ ปฏิจฺจ นสรโณ นกุสโล}
\csromanheader{2}{100. สรณทุก 2. เวทนาตฺติกาทิ}
\prose{\csromanfolio{119}น-อรณํ นกุสลํ ธมฺมํ ปฏิจฺจ น-อรโณ นกุสโล ธมฺโม อุปฺปชฺชติ เหตุปจฺจยา.\csromanspacer{}\csromanspacer{}น-อรณํ นกุสลํ ธมฺมํ ปฏิจฺจ นสรโณ นกุสโล ธมฺโม อุปฺปชฺชติ เหตุปจฺจยา.\csromanspacer{}\csromanspacer{}น-อรณํ นกุสลํ ธมฺมํ ปฏิจฺจ นสรโณ นกุสโล จ น-อรโณ นกุสโล จ ธมฺมา อุปฺปชฺชนฺติ เหตุปจฺจยา. (3)}

\prose{นสรณํ นกุสลญฺจ น-อรณํ นกุสลญฺจ ธมฺมํ ปฏิจฺจ นสรโณ นกุสโล ธมฺโม อุปฺปชฺชติ เหตุปจฺจยา. (1) (สํขิตฺตํ.) เหตุยา ปญฺจ ฯเปฯ อวิคเต ปญฺจ.}

\proseitem{84}{นสรณํ น-อกุสลํ ธมฺมํ ปฏิจฺจ นสรโณ น-อกุสโล ธมฺโม อุปฺปชฺชติ เหตุปจฺจยา. (สํขิตฺตํ.) เหตุยา เอกํ. (สพฺพตฺถ เอกํ.)}

\csromanheader{2}{85. นสรณํ น-อพฺยากตํ ธมฺมํ ปฏิจฺจ นสรโณ น-อพฺยากโต}
\prose{น-อรณํ น-อพฺยากตํ ธมฺมํ ปฏิจฺจ น-อรโณ น-อพฺยากโต ธมฺโม อุปฺปชฺชติ เหตุปจฺจยา. (1) (สํขิตฺตํ.)}

\prose{เหตุยา เทฺว, อารมฺมเณ เทฺว ฯเปฯ อวิคเต เทฺว. (สหชาตวาเรปิ ฯเปฯ ปญฺหาวาเรปิ สพฺพตฺถ วิตฺถาโร.)}

\csromanheader{2}{100. สรณทุก 2. เวทนาตฺติกาทิ}
\csromanheader{2}{\textbf{ปจฺจยจตุกฺก-เหตุ}}
\proseitem{86}{นสรณํ นสุขาย เวทนาย สมฺปยุตฺตํ ธมฺมํ ปฏิจฺจ ฯเปฯ.\csromanspacer{} นสรณํ นทุกฺขาย เวทนาย สมฺปยุตฺตํ ธมฺมํ ปฏิจฺจ ฯเปฯ.\csromanspacer{} นสรณํ น-อทุกฺขมสุขาย เวทนาย สมฺปยุตฺตํ ธมฺมํ ปฏิจฺจ.}

\proseitem{87}{นสรณํ นวิปากํ ธมฺมํ ปฏิจฺจ.}

\prose{นสรณํ นวิปากธมฺมธมฺมํ ปฏิจฺจ.}

\prose{นสรณํ นเนววิปากนวิปากธมฺมธมฺมํ ปฏิจฺจ.}

\gambhira{ธมฺมปจฺจนีย ทุกติกปฏฺฐานปาฬิ}
\proseitem{88}{\csromanfolio{120}นสรณํ น-อุปาทินฺนุปาทานิยํ ธมฺมํ ปฏิจฺจ.}

\prose{นสรณํ น-อนุปาทินฺนุปาทานิยํ ธมฺมํ ปฏิจฺจ.}

\prose{นสรณํ น-อนุปาทินฺน-อนุปาทานิยํ ธมฺมํ ปฏิจฺจ.}

\csromanheader{2}{100. สรณทุก 22. สนิทสฺสนตฺติก}
\proseitem{89}{นสรณํ นสนิทสฺสนสปฺปฏิฆํ ธมฺมํ ปฏิจฺจ นสรโณ นสนิทสฺสนสปฺปฏิโฆ ธมฺโม อุปฺปชฺชติ เหตุปจฺจยา. (1)}

\prose{น-อรณํ นสนิทสฺสนสปฺปฏิฆํ ธมฺมํ ปฏิจฺจ น-อรโณ นสนิทสฺสนสปฺปฏิโฆ ธมฺโม อุปฺปชฺชติ เหตุปจฺจยา.\csromanspacer{}\csromanspacer{}น-อรณํ นสนิทสฺสนสปฺปฏิฆํ ธมฺมํ ปฏิจฺจ นสรโณ นสนิทสฺสนสปฺปฏิโฆ ธมฺโม อุปฺปชฺชติ เหตุปจฺจยา.\csromanspacer{}\csromanspacer{}น-อรณํ นสนิทสฺสนสปฺปฏิฆํ ธมฺมํ ปฏิจฺจ นสรโณ นสนิทสฺสนสปฺปฏิโฆ จ น-อรโณ นสนิทสฺสนสปฺปฏิโฆ จ ธมฺมา อุปฺปชฺชนฺติ เหตุปจฺจยา. (3)}

\prose{นสรณํ นสนิทสฺสนสปฺปฏิฆญฺจ น-อรณํ นสนิทสฺสนสปฺปฏิฆญฺจ ธมฺมํ ปฏิจฺจ นสรโณ นสนิทสฺสนสปฺปฏิโฆ ธมฺโม อุปฺปชฺชติ เหตุปจฺจยา. (1) (สํขิตฺตํ.)}

\prose{เหตุยา ปญฺจ.}

\proseitem{90}{นสรณํ น-อนิทสฺสนสปฺปฏิฆํ ธมฺมํ ปฏิจฺจ นสรโณ น-อนิทสฺสนสปฺปฏิโฆ ธมฺโม อุปฺปชฺชติ เหตุปจฺจยา. (1)}

\prose{น-อรณํ น-อนิทสฺสนสปฺปฏิฆํ ธมฺมํ ปฏิจฺจ นสรโณ น-อนิทสฺสนสปฺปฏิโฆ ธมฺโม อุปฺปชฺชติ เหตุปจฺจยา.\csromanspacer{}\csromanspacer{}น-อรณํ น-อนิทสฺสนสปฺปฏิฆํ ธมฺมํ ปฏิจฺจ นสรโณ น-อนิทสฺสนสปฺปฏิโฆ จ น-อรโณ น-อนิทสฺสนสปฺปฏิโฆ จ ธมฺมา อุปฺปชฺชนฺติ เหตุปจฺจยา. (3)}

\csromanheader{2}{100. สรณทุก 22. สนิทสฺสนตฺติก}
\prose{\csromanfolio{121}นสรณํ น-อนิทสฺสนสปฺปฏิฆญฺจ น-อรณํ น-อนิทสฺสนสปฺปฏิฆญฺจ ธมฺมํ ปฏิจฺจ นสรโณ น-อนิทสฺสนสปฺปฏิโฆ ธมฺโม อุปฺปชฺชติ เหตุปจฺจยา. (1) (สํขิตฺตํ.)}

\prose{เหตุยา ปญฺจ, อารมฺมเณ เทฺว ฯเปฯ อวิคเต ปญฺจ. (สหชาตวาเรปิ ฯเปฯ ปญฺหาวาเรปิ สพฺพตฺถ วิตฺถาเรตพฺพํ.)}

\proseitem{91}{นสรณํ น-อนิทสฺสน-อปฺปฏิฆํ ธมฺมํ ปฏิจฺจ นสรโณ น-อนิทสฺสน-อปฺปฏิโฆ ธมฺโม อุปฺปชฺชติ เหตุปจฺจยา. (สํขิตฺตํ.)}

\csromanheader{2}{เหตุยา เอกํ. (สพฺพตฺถ วิตฺถาโร.)}
\nitthitam{\textbf{ธมฺมปจฺจนีเย ทุกติกปฏฺฐานํ นิฏฺฐิตํ.}}
\pitaka{\textbf{อภิธมฺมปิฏก}}
\csromanheader{2}{\textbf{ธมฺมปจฺจนีย}}
\csromanmatikaanchor{mtk.7}
\csromantocmarknopage{chapter}{ติกทุกปฏฺฐานปาฬิ}{mtk.7}
\bookmark[dest={mtk.7},level=0]{ติกทุกปฏฺฐานปาฬิ}
\gambhira{\textbf{ติกทุกปฏฺฐานปาฬิ}}
\namakkaram{นโม ตสฺส ภควโต อรหโต สมฺมาสมฺพุทฺธสฺส.}
\csromanheader{2}{1. กุสลตฺติก 1. เหตุทุก}
\csromanheader{2}{1. ปฏิจฺจวาร-เหตุ}
\proseitem{1}{\csromanfolio{123}นกุสลํ นเหตุํ ธมฺมํ ปฏิจฺจ นกุสโล นเหตุ ธมฺโม อุปฺปชฺชติ เหตุปจฺจยา.\csromanspacer{}\csromanspacer{}นกุสลํ นเหตุํ ธมฺมํ ปฏิจฺจ น-อกุสโล นเหตุ ธมฺโม อุปฺปชฺชติ เหตุปจฺจยา.\csromanspacer{}\csromanspacer{}นกุสลํ นเหตุํ ธมฺมํ ปฏิจฺจ น-อพฺยากโต นเหตุ ธมฺโม อุปฺปชฺชติ เหตุปจฺจยา.\csromanspacer{}\csromanspacer{}นกุสลํ นเหตุํ ธมฺมํ ปฏิจฺจ นกุสโล นเหตุ จ น-อพฺยากโต นเหตุ จ ธมฺมา อุปฺปชฺชนฺติ เหตุปจฺจยา.\csromanspacer{}\csromanspacer{}นกุสลํ นเหตุํ ธมฺมํ ปฏิจฺจ นกุสโล นเหตุ จ น-อกุสโล นเหตุ จ ธมฺมา อุปฺปชฺชนฺติ เหตุปจฺจยา.\csromanspacer{} ปญฺจ.}

\proseitem{2}{น-อกุสลํ นเหตุํ ธมฺมํ ปฏิจฺจ น-อกุสโล นเหตุ ธมฺโม อุปฺปชฺชติ เหตุปจฺจยา.\csromanspacer{}\csromanspacer{}น-อกุสลํ นเหตุํ ธมฺมํ ปฏิจฺจ นกุสโล นเหตุ ธมฺโม อุปฺปชฺชติ เหตุปจฺจยา.\csromanspacer{}\csromanspacer{}น-อกุสลํ นเหตุํ ธมฺมํ ปฏิจฺจ น-อพฺยากโต นเหตุ ธมฺโม อุปฺปชฺชติ เหตุปจฺจยา.\csromanspacer{}\csromanspacer{}น-อกุสลํ นเหตุํ ธมฺมํ ปฏิจฺจ น-อกุสโล นเหตุ จ น-อพฺยากโต นเหตุ จ ธมฺมา อุปฺปชฺชนฺติ เหตุปจฺจยา.\csromanspacer{}\csromanspacer{}น-อกุสลํ นเหตุํ ธมฺมํ ปฏิจฺจ นกุสโล นเหตุ จ น-อกุสโล นเหตุ จ ธมฺมา อุปฺปชฺชนฺติ เหตุปจฺจยา.\csromanspacer{} ปญฺจ.}

\gambhira{ธมฺมปจฺจนีย ติกทุกปฏฺฐานปาฬิ}
\proseitem{3}{\csromanfolio{124}น-อพฺยากตํ นเหตุํ ธมฺมํ ปฏิจฺจ น-อพฺยากโต นเหตุ ธมฺโม อุปฺปชฺชติ เหตุปจฺจยา.\csromanspacer{}\csromanspacer{}น-อพฺยากตํ นเหตุํ ธมฺมํ ปฏิจฺจ นกุสโล นเหตุ ธมฺโม อุปฺปชฺชติ เหตุปจฺจยา.\csromanspacer{}\csromanspacer{}น-อพฺยากตํ นเหตุํ ธมฺมํ ปฏิจฺจ น-อกุสโล นเหตุ ธมฺโม อุปฺปชฺชติ เหตุปจฺจยา.\csromanspacer{}\csromanspacer{}น-อพฺยากตํ นเหตุํ ธมฺมํ ปฏิจฺจ นกุสโล นเหตุ จ น-อพฺยากโต นเหตุ จ ธมฺมา อุปฺปชฺชนฺติ เหตุปจฺจยา.\csromanspacer{}\csromanspacer{}น-อพฺยากตํ นเหตุํ ธมฺมํ ปฏิจฺจ น-อกุสโล นเหตุ จ น-อพฺยากโต นเหตุ จ ธมฺมา อุปฺปชฺชนฺติ เหตุปจฺจยา.\csromanspacer{}\csromanspacer{}น-อพฺยากตํ นเหตุํ ธมฺมํ ปฏิจฺจ นกุสโล นเหตุ จ น-อกุสโล นเหตุ จ ธมฺมา อุปฺปชฺชนฺติ เหตุปจฺจยา.\csromanspacer{} ฉ.}

\proseitem{4}{นกุสลํ นเหตุญฺจ น-อพฺยากตํ นเหตุญฺจ ธมฺมํ ปฏิจฺจ นกุสโล นเหตุ ธมฺโม อุปฺปชฺชติ เหตุปจฺจยา.\csromanspacer{}\csromanspacer{}นกุสลํ นเหตุญฺจ น-อพฺยากตํ นเหตุญฺจ ธมฺมํ ปฏิจฺจ น-อกุสโล นเหตุ ธมฺโม อุปฺปชฺชติ เหตุปจฺจยา.\csromanspacer{}\csromanspacer{}นกุสลํ นเหตุญฺจ น-อพฺยากตํ นเหตุญฺจ ธมฺมํ ปฏิจฺจ น-อพฺยากโต นเหตุ ธมฺโม อุปฺปชฺชติ เหตุปจฺจยา.\csromanspacer{}\csromanspacer{}นกุสลํ นเหตุญฺจ น-อพฺยากตํ นเหตุญฺจ ธมฺมํ ปฏิจฺจ นกุสโล นเหตุ จ น-อพฺยากโต นเหตุ จ ธมฺมา อุปฺปชฺชนฺติ เหตุปจฺจยา.\csromanspacer{}\csromanspacer{}นกุสลํ นเหตุญฺจ น-อพฺยากตํ นเหตุญฺจ ธมฺมํ ปฏิจฺจ นกุสโล นเหตุ จ น-อกุสโล นเหตุ จ ธมฺมา อุปฺปชฺชนฺติ เหตุปจฺจยา.\csromanspacer{} ปญฺจ.}

\proseitem{5}{น-อกุสลํ นเหตุญฺจ น-อพฺยากตํ นเหตุญฺจ ธมฺมํ ปฏิจฺจ นกุสโล นเหตุ ธมฺโม อุปฺปชฺชติ เหตุปจฺจยา.\csromanspacer{}\csromanspacer{}น-อกุสลํ นเหตุญฺจ น-อพฺยากตํ นเหตุญฺจ ธมฺมํ ปฏิจฺจ น-อกุสโล นเหตุ ธมฺโม อุปฺปชฺชติ เหตุปจฺจยา.\csromanspacer{}\csromanspacer{}น-อกุสลํ นเหตุญฺจ น-อพฺยากตํ นเหตุญฺจ ธมฺมํ ปฏิจฺจ น-อพฺยากโต นเหตุ ธมฺโม อุปฺปชฺชติ เหตุปจฺจยา.\csromanspacer{}\csromanspacer{}น-อกุสลํ นเหตุญฺจ น-อพฺยากตํ นเหตุญฺจ ธมฺมํ ปฏิจฺจ น-อกุสโล นเหตุ จ น-อพฺยากโต นเหตุ จ ธมฺมา อุปฺปชฺชนฺติ เหตุปจฺจยา.\csromanspacer{}\csromanspacer{}น-อกุสลํ นเหตุญฺจ น-อพฺยากตํ นเหตุญฺจ ธมฺมํ ปฏิจฺจ นกุสโล นเหตุ จ น-อกุสโล นเหตุ จ ธมฺมา อุปฺปชฺชนฺติ เหตุปจฺจยา.\csromanspacer{} ปญฺจ.}

\proseitem{6}{นกุสลํ นเหตุญฺจ น-อกุสลํ นเหตุญฺจ ธมฺมํ ปฏิจฺจ นกุสโล นเหตุ ธมฺโม อุปฺปชฺชติ เหตุปจฺจยา.\csromanspacer{}\csromanspacer{}นกุสลํ นเหตุญฺจ น-อกุสลํ นเหตุญฺจ ธมฺมํ ปฏิจฺจ น-อกุสโล นเหตุ ธมฺโม อุปฺปชฺชติ เหตุปจฺจยา.\csromanspacer{}\csromanspacer{}นกุสลํ นเหตุญฺจ น-อกุสลํ นเหตุญฺจ ธมฺมํ ปฏิจฺจ}

\vspace{\tipitakaparskip}
\csromancenter{กุสลตฺติก 2.\csromanspacer{} สเหตุกทุก นกุสโล นเหตุ จ น-อกุสโล นเหตุ จ ธมฺมา อุปฺปชฺชนฺติ เหตุปจฺจยา.\csromanspacer{} ตีณิ. (สํขิตฺตํ.)}
\prose{\csromanfolio{125}เหตุยา เอกูนติํส, อารมฺมเณ จตุวีส ฯเปฯ อวิคเต เอกูนติํส. (สพฺพตฺถ วิตฺถาโร.)}

\proseitem{7}{นกุสลํ นนเหตุํ ธมฺมํ ปฏิจฺจ นกุสโล นนเหตุ ธมฺโม อุปฺปชฺชติ เหตุปจฺจยา.\csromanspacer{}\csromanspacer{}นกุสลํ นนเหตุํ ธมฺมํ ปฏิจฺจ น-อกุสโล นนเหตุ ธมฺโม อุปฺปชฺชติ เหตุปจฺจยา.\csromanspacer{}\csromanspacer{}นกุสลํ นนเหตุํ ธมฺมํ ปฏิจฺจ น-อพฺยากโต นนเหตุ ธมฺโม อุปฺปชฺชติ เหตุปจฺจยา.\csromanspacer{}\csromanspacer{}นกุสลํ นนเหตุํ ธมฺมํ ปฏิจฺจ นกุสโล นนเหตุ จ น-อพฺยากโต นนเหตุ จ ธมฺมา อุปฺปชฺชนฺติ เหตุปจฺจยา.\csromanspacer{}\csromanspacer{}นกุสลํ นนเหตุํ ธมฺมํ ปฏิจฺจ นกุสโล นนเหตุ จ น-อกุสโล นนเหตุ จ ธมฺมา อุปฺปชฺชนฺติ เหตุปจฺจยา. (5)}

\prose{น-อกุสลํ นนเหตุํ ธมฺมํ ปฏิจฺจ น-อกุสโล นนเหตุ ธมฺโม อุปฺปชฺชติ เหตุปจฺจยา.\csromanspacer{} ปญฺจ.}

\prose{น-อพฺยากตํ นนเหตุํ ธมฺมํ ปฏิจฺจ น-อพฺยากโต นนเหตุ ธมฺโม อุปฺปชฺชติ เหตุปจฺจยา.\csromanspacer{} ปญฺจ.}

\prose{นกุสลํ นนเหตุญฺจ น-อพฺยากตํ นนเหตุญฺจ ธมฺมํ ปฏิจฺจ นกุสโล นนเหตุ ธมฺโม อุปฺปชฺชติ เหตุปจฺจยา.\csromanspacer{} ตีณิ.}

\prose{น-อกุสลํ นนเหตุญฺจ น-อพฺยากตํ นนเหตุญฺจ ธมฺมํ ปฏิจฺจ น-อกุสโล นนเหตุ ธมฺโม อุปฺปชฺชติ เหตุปจฺจยา.\csromanspacer{} ตีณิ.}

\prose{นกุสลํ นนเหตุญฺจ น-อกุสลํ นนเหตุญฺจ ธมฺมํ ปฏิจฺจ นกุสโล นนเหตุ ธมฺโม อุปฺปชฺชติ เหตุปจฺจยา.\csromanspacer{} ตีณิ. (สํขิตฺตํ.)}

\prose{เหตุยา จตุวีส, อารมฺมเณ จตุวีส ฯเปฯ อวิคเต จตุวีส. (สพฺพตฺถ วิตฺถาโร.)}

\csromanheader{2}{1. กุสลตฺติก 2. สเหตุกทุก}
\proseitem{8}{นกุสลํ นสเหตุกํ ธมฺมํ ปฏิจฺจ นกุสโล นสเหตุโก ธมฺโม อุปฺปชฺชติ เหตุปจฺจยา.\csromanspacer{}\csromanspacer{}นกุสลํ นสเหตุกํ ธมฺมํ ปฏิจฺจ น-อกุสโล นสเหตุโก ธมฺโม อุปฺปชฺชติ เหตุปจฺจยา.\csromanspacer{}\csromanspacer{}นกุสลํ นสเหตุกํ}

\vspace{\tipitakaparskip}
\csromancenter{ธมฺมปจฺจนีย ติกทุกปฏฺฐานปาฬิ ธมฺมํ ปฏิจฺจ นกุสโล นสเหตุโก จ น-อกุสโล นสเหตุโก จ ธมฺมา อุปฺปชฺชนฺติ เหตุปจฺจยา. (3)}
\prose{\csromanfolio{126}น-อกุสลํ นสเหตุกํ ธมฺมํ ปฏิจฺจ น-อกุสโล นสเหตุโก ธมฺโม อุปฺปชฺชติ เหตุปจฺจยา.\csromanspacer{}\csromanspacer{}น-อกุสลํ นสเหตุกํ ธมฺมํ ปฏิจฺจ นกุสโล นสเหตุโก ธมฺโม อุปฺปชฺชติ เหตุปจฺจยา.\csromanspacer{}\csromanspacer{}น-อกุสลํ นสเหตุกํ ธมฺมํ ปฏิจฺจ นกุสโล นสเหตุโก จ น-อกุสโล นสเหตุโก จ ธมฺมา อุปฺปชฺชนฺติ เหตุปจฺจยา. (3)}

\prose{น-อพฺยากตํ นสเหตุกํ ธมฺมํ ปฏิจฺจ นกุสโล นสเหตุโก ธมฺโม อุปฺปชฺชติ เหตุปจฺจยา.\csromanspacer{} ตีณิ.}

\prose{นกุสลํ นสเหตุกญฺจ น-อพฺยากตํ นสเหตุกญฺจ ธมฺมํ ปฏิจฺจ นกุสโล นสเหตุโก ธมฺโม อุปฺปชฺชติ เหตุปจฺจยา.\csromanspacer{} ตีณิ.}

\prose{นกุสลํ นสเหตุกญฺจ น-อกุสลํ นสเหตุกญฺจ ธมฺมํ ปฏิจฺจ นกุสโล นสเหตุโก ธมฺโม อุปฺปชฺชติ เหตุปจฺจยา.\csromanspacer{} ตีณิ.}

\prose{เหตุยา ปนฺนรส, อารมฺมเณ นว ฯเปฯ อธิปติยา นว ฯเปฯ วิปาเก นว ฯเปฯ อวิคเต ปนฺนรส. (สพฺพตฺถ วิตฺถาโร.)}

\proseitem{9}{นกุสลํ น-อเหตุกํ ธมฺมํ ปฏิจฺจ นกุสโล น-อเหตุโก ธมฺโม อุปฺปชฺชติ เหตุปจฺจยา.\csromanspacer{} ปญฺจ.}

\prose{น-อกุสลํ น-อเหตุกํ ธมฺมํ ปฏิจฺจ น-อกุสโล น-อเหตุโก ธมฺโม อุปฺปชฺชติ เหตุปจฺจยา.\csromanspacer{} ปญฺจ.}

\prose{น-อพฺยากตํ น-อเหตุกํ ธมฺมํ ปฏิจฺจ น-อพฺยากโต น-อเหตุโก ธมฺโม อุปฺปชฺชติ เหตุปจฺจยา.\csromanspacer{} ปญฺจ.}

\prose{นกุสลํ น-อเหตุกญฺจ น-อพฺยากตํ น-อเหตุกญฺจ ธมฺมํ ปฏิจฺจ นกุสโล น-อเหตุโก ธมฺโม อุปฺปชฺชติ เหตุปจฺจยา.\csromanspacer{} ตีณิ.}

\prose{น-อกุสลํ น-อเหตุกญฺจ น-อพฺยากตํ น-อเหตุกญฺจ ธมฺมํ ปฏิจฺจ น-อกุสโล น-อเหตุโก ธมฺโม อุปฺปชฺชติ เหตุปจฺจยา.\csromanspacer{} ตีณิ.}

\prose{นกุสลํ น-อเหตุกญฺจ น-อกุสลํ น-อเหตุกญฺจ ธมฺมํ ปฏิจฺจ นกุสโล น-อเหตุโก ธมฺโม อุปฺปชฺชติ เหตุปจฺจยา.\csromanspacer{} ตีณิ.}

\csromanheader{2}{เหตุยา จตุวีส ฯเปฯ อวิคเต จตุวีส. (สพฺพตฺถ วิตฺถาโร.)}
\vspace{\tipitakaparskip}
\csromancenter{กุสลตฺติก 14.\csromanspacer{} อาสวโคจฺฉก \textbf{1.}\csromanspacer{}\textbf{ กุสลตฺติก 3.}\csromanspacer{}\textbf{ เหตุสมฺปยุตฺตทุกาทิ}}
\proseitem{10}{\csromanfolio{127}นกุสลํ นเหตุสมฺปยุตฺตํ ธมฺมํ ปฏิจฺจ.}

\proseitem{11}{นกุสลํ นเหตุวิปฺปยุตฺตํ ธมฺมํ ปฏิจฺจ.}

\proseitem{12}{นกุสลํ นเหตุญฺเจว น-อเหตุกญฺจ ธมฺมํ ปฏิจฺจ ฯเปฯ.\csromanspacer{} นกุสลํ น-อเหตุกญฺเจว นนเหตุญฺจ ธมฺมํ ปฏิจฺจ ฯเปฯ.\csromanspacer{} นกุสลํ นเหตุญฺเจว นเหตุวิปฺปยุตฺตญฺจ ธมฺมํ ปฏิจฺจ ฯเปฯ.\csromanspacer{} นกุสลํ นเหตุวิปฺปยุตฺตญฺเจว นนเหตุญฺจ ธมฺมํ ปฏิจฺจ.}

\proseitem{13}{นกุสลํ นเหตุํ นสเหตุกํ ธมฺมํ ปฏิจฺจ.}

\prose{นกุสลํ นเหตุํ น-อเหตุกํ ธมฺมํ ปฏิจฺจ.}

\csromanheader{2}{1. กุสลตฺติก 7. จูฬนฺตรทุก}
\proseitem{14}{นกุสลํ น-อปฺปจฺจยํ ธมฺมํ ปฏิจฺจ ฯเปฯ.\csromanspacer{} นกุสลํ น-อสงฺขตํ ธมฺมํ ปฏิจฺจ.}

\proseitem{15}{นกุสลํ นสนิทสฺสนํ ธมฺมํ ปฏิจฺจ.}

\proseitem{16}{นกุสลํ นสปฺปฏิฆํ ธมฺมํ ปฏิจฺจ.}

\proseitem{17}{นกุสลํ น-อปฺปฏิฆํ ธมฺมํ ปฏิจฺจ.}

\proseitem{18}{นกุสลํ นรูปิํ ธมฺมํ ปฏิจฺจ.}

\proseitem{19}{นกุสลํ น-อรูปิํ ธมฺมํ ปฏิจฺจ.}

\proseitem{20}{นกุสลํ นโลกิยํ ธมฺมํ ปฏิจฺจ.}

\proseitem{21}{นกุสลํ นโลกุตฺตรํ ธมฺมํ ปฏิจฺจ.}

\proseitem{22}{นกุสลํ นเกนจิ วิญฺเญยฺยํ ธมฺมํ ปฏิจฺจ ฯเปฯ.\csromanspacer{} นกุสลํ นนเกนจิ วิญฺเญยฺยํ ธมฺมํ ปฏิจฺจ.}

\csromanheader{2}{1. กุสลตฺติก 14. อาสวโคจฺฉก}
\proseitem{23}{นกุสลํ โน-อาสวํ ธมฺมํ ปฏิจฺจ.}

\prose{นกุสลํ นโน-อาสวํ ธมฺมํ ปฏิจฺจ.}

\gambhira{ธมฺมปจฺจนีย ติกทุกปฏฺฐานปาฬิ}
\proseitem{24}{\csromanfolio{128}นกุสลํ นสาสวํ ธมฺมํ ปฏิจฺจ.}

\prose{นกุสลํ น-อนาสวํ ธมฺมํ ปฏิจฺจ.}

\proseitem{25}{นกุสลํ น-อาสวสมฺปยุตฺตํ ธมฺมํ ปฏิจฺจ.}

\prose{นกุสลํ น-อาสววิปฺปยุตฺตํ ธมฺมํ ปฏิจฺจ.}

\proseitem{26}{นกุสลํ โน-อาสวญฺเจว น-อนาสวญฺจ ธมฺมํ ปฏิจฺจ.}

\prose{นกุสลํ น-อนาสวญฺเจว นโน จ อาสวํ ธมฺมํ ปฏิจฺจ.}

\proseitem{27}{นกุสลํ น-อนาสวญฺเจว น-อาสววิปฺปยุตฺตญฺจ ธมฺมํ ปฏิจฺจ.}

\prose{นกุสลํ น-อาสววิปฺปยุตฺตญฺเจว นโน จ อาสวํ ธมฺมํ ปฏิจฺจ.}

\proseitem{28}{นกุสลํ อาสววิปฺปยุตฺตํ นสาสวํ ธมฺมํ ปฏิจฺจ.}

\prose{นกุสลํ อาสววิปฺปยุตฺตํ น-อนาสวํ ธมฺมํ ปฏิจฺจ.}

\csromanheader{2}{1. กุสลตฺติก 20. สญฺโญชนโคจฺฉกาทิ}
\proseitem{29}{นกุสลํ โนสญฺโญชนํ ธมฺมํ ปฏิจฺจ.}

\proseitem{30}{นกุสลํ โนคนฺถํ ธมฺมํ ปฏิจฺจ ฯเปฯ.\csromanspacer{} นกุสลํ โน-โอฆํ ธมฺมํ ปฏิจฺจ ฯเปฯ.\csromanspacer{} นกุสลํ โนโยคํ ธมฺมํ ปฏิจฺจ ฯเปฯ.\csromanspacer{} นกุสลํ โนนีวรณํ ธมฺมํ ปฏิจฺจ ฯเปฯ.\csromanspacer{} นกุสลํ โนปรามาสํ ธมฺมํ ปฏิจฺจ.}

\csromanheader{2}{1. กุสลตฺติก 55. มหนฺตรทุก}
\proseitem{31}{นกุสลํ นสารมฺมณํ ธมฺมํ ปฏิจฺจ.}

\csromanheader{2}{นกุสลํ น-อนารมฺมณํ ธมฺมํ ปฏิจฺจ. (สํขิตฺตํ.)}
\proseitem{32}{นกุสลํ นจิตฺตํ ธมฺมํ ปฏิจฺจ. (สํขิตฺตํ.\csromanspacer{} นโนจิตฺตปทํ น ลพฺภติ.)}

\proseitem{33}{นกุสลํ นเจตสิกํ ธมฺมํ ปฏิจฺจ.}

\proseitem{34}{นกุสลํ นจิตฺตสมฺปยุตฺตํ ธมฺมํ ปฏิจฺจ ฯเปฯ.\csromanspacer{} นกุสลํ นจิตฺตสํสฏฺฐํ ธมฺมํ ปฏิจฺจ ฯเปฯ.\csromanspacer{} นกุสลํ นจิตฺตสมุฏฺฐานํ ธมฺมํ ปฏิจฺจ.}

\csromanheader{2}{1. กุสลตฺติก 83. ปิฏฺฐิทุก}
\proseitem{35}{\csromanfolio{129}นกุสลํ นจิตฺตสหภุํ ธมฺมํ ปฏิจฺจ.}

\proseitem{36}{นกุสลํ นจิตฺตานุปริวตฺติํ ธมฺมํ ปฏิจฺจ ฯเปฯ.\csromanspacer{} นกุสลํ นจิตฺตสํสฏฺฐสมุฏฺฐานํ ธมฺมํ ปฏิจฺจ ฯเปฯ.\csromanspacer{} นกุสลํ นจิตฺตสํสฏฺฐสมุฏฺฐานสหภุํ ธมฺมํ ปฏิจฺจ ฯเปฯ.\csromanspacer{} นกุสลํ นจิตฺตสํสฏฺฐสมุฏฺฐานานุปริวตฺติํ ธมฺมํ ปฏิจฺจ.}

\proseitem{37}{นกุสลํ น-อชฺฌตฺติกํ ธมฺมํ ปฏิจฺจ.}

\prose{นกุสลํ นพาหิรํ ธมฺมํ ปฏิจฺจ.}

\proseitem{38}{นกุสลํ น-อุปาทา ธมฺมํ ปฏิจฺจ.}

\proseitem{39}{นกุสลํ น-อุปาทินฺนํ ธมฺมํ ปฏิจฺจ.}

\prose{นกุสลํ น-อนุปาทินฺนํ ธมฺมํ ปฏิจฺจ.}

\csromanheader{2}{1. กุสลตฺติก 69. อุปาทานทุกาทิ}
\proseitem{40}{นกุสลํ โน-อุปาทานํ ธมฺมํ ปฏิจฺจ.}

\prose{นกุสลํ นโน-อุปาทานํ ธมฺมํ ปฏิจฺจ.}

\proseitem{41}{นกุสลํ โนกิเลสํ ธมฺมํ ปฏิจฺจ.}

\prose{นกุสลํ นโนกิเลสํ ธมฺมํ ปฏิจฺจ.}

\csromanheader{2}{1. กุสลตฺติก 83. ปิฏฺฐิทุก}
\proseitem{42}{นกุสลํ นทสฺสเนน ปหาตพฺพํ ธมฺมํ ปฏิจฺจ ฯเปฯ.\csromanspacer{} นกุสลํ นนทสฺสเนน ปหาตพฺพํ ธมฺมํ ปฏิจฺจ ฯเปฯ.\csromanspacer{} นกุสลํ นภาวนาย ปหาตพฺพํ ธมฺมํ ปฏิจฺจ ฯเปฯ.\csromanspacer{} นกุสลํ นภาวนาย ปหาตพฺพํ ธมฺมํ ปฏิจฺจ ฯเปฯ.\csromanspacer{} นกุสลํ นทสฺสเนน ปหาตพฺพเหตุกํ ธมฺมํ ปฏิจฺจ ฯเปฯ.\csromanspacer{} นกุสลํ นนทสฺสเนน ปหาตพฺพเหตุกํ ธมฺมํ ปฏิจฺจ ฯเปฯ.\csromanspacer{} นกุสลํ นภาวนาย ปหาตพฺพเหตุกํ ธมฺมํ ปฏิจฺจ.}

\prose{นกุสลํ นนภาวนาย ปหาตพฺพเหตุกํ ธมฺมํ ปฏิจฺจ.}

\proseitem{43}{นกุสลํ นสวิตกฺกํ ธมฺมํ ปฏิจฺจ ฯเปฯ.\csromanspacer{} นกุสลํ น-อวิตกฺกํ ธมฺมํ ปฏิจฺจ ฯเปฯ.\csromanspacer{} นกุสลํ นสวิจารํ ธมฺมํ ปฏิจฺจ.}

\prose{นกุสลํ น-อวิจารํ ธมฺมํ ปฏิจฺจ.}

\gambhira{ธมฺมปจฺจนีย ติกทุกปฏฺฐานปาฬิ}
\proseitem{44}{\csromanfolio{130}นกุสลํ นสปฺปีติกํ ธมฺมํ ปฏิจฺจ ฯเปฯ.\csromanspacer{} นกุสลํ น-อปฺปีติกํ ธมฺมํ ปฏิจฺจ ฯเปฯ.\csromanspacer{} นกุสลํ นปีติสหคตํ ธมฺมํ ปฏิจฺจ ฯเปฯ.\csromanspacer{} นกุสลํ นนปีติสหคตํ ธมฺมํ ปฏิจฺจ ฯเปฯ.\csromanspacer{} นกุสลํ นสุขสหคตํ ธมฺมํ ปฏิจฺจ ฯเปฯ.\csromanspacer{} นกุสลํ นนสุขสหคตํ ธมฺมํ ปฏิจฺจ ฯเปฯ.\csromanspacer{} นกุสลํ น-อุเปกฺขาสหคตํ ธมฺมํ ปฏิจฺจ.}

\prose{นกุสลํ นน-อุเปกฺขาสหคตํ ธมฺมํ ปฏิจฺจ.}

\proseitem{45}{นกุสลํ นกามาวจรํ ธมฺมํ ปฏิจฺจ.}

\prose{นกุสลํ นนกามาวจรํ ธมฺมํ ปฏิจฺจ.}

\proseitem{46}{นกุสลํ นรูปาวจรํ ธมฺมํ ปฏิจฺจ ฯเปฯ.\csromanspacer{} นกุสลํ นนรูปาวจรํ ธมฺมํ ปฏิจฺจ ฯเปฯ.\csromanspacer{} นกุสลํ น-อรูปาวจรํ ธมฺมํ ปฏิจฺจ.}

\prose{นกุสลํ นน-อรูปาวจรํ ธมฺมํ ปฏิจฺจ.}

\proseitem{47}{นกุสลํ นปริยาปนฺนํ ธมฺมํ ปฏิจฺจ.}

\prose{นกุสลํ น-อปริยาปนฺนํ ธมฺมํ ปฏิจฺจ.}

\proseitem{48}{นกุสลํ นนิยฺยานิกํ ธมฺมํ ปฏิจฺจ.}

\prose{นกุสลํ น-อนิยฺยานิกํ ธมฺมํ ปฏิจฺจ.}

\proseitem{49}{นกุสลํ นนิยตํ ธมฺมํ ปฏิจฺจ.}

\prose{นกุสลํ น-อนิยตํ ธมฺมํ ปฏิจฺจ.}

\proseitem{50}{นกุสลํ นส-อุตฺตรํ ธมฺมํ ปฏิจฺจ.}

\prose{นกุสลํ น-อนุตฺตรํ ธมฺมํ ปฏิจฺจ.}

\proseitem{51}{นกุสลํ นสรณํ ธมฺมํ ปฏิจฺจ นกุสโล นสรโณ ธมฺโม อุปฺปชฺชติ เหตุปจฺจยา.\csromanspacer{}\csromanspacer{}นกุสลํ นสรณํ ธมฺมํ ปฏิจฺจ น-อกุสโล นสรโณ ธมฺโม อุปฺปชฺชติ เหตุปจฺจยา.\csromanspacer{}\csromanspacer{}นกุสลํ นสรณํ ธมฺมํ ปฏิจฺจ นกุสโล นสรโณ จ น-อกุสโล นสรโณ จ ธมฺมา อุปฺปชฺชนฺติ เหตุปจฺจยา. (3)}

\prose{น-อกุสลํ นสรณํ ธมฺมํ ปฏิจฺจ น-อกุสโล นสรโณ ธมฺโม อุปฺปชฺชติ เหตุปจฺจยา.\csromanspacer{} ปญฺจ.}

\csromanheader{2}{2. เวทนาตฺติก 1. เหตุทุก}
\prose{\csromanfolio{131}น-อพฺยากตํ นสรณํ ธมฺมํ ปฏิจฺจ น-อพฺยากโต นสรโณ ธมฺโม อุปฺปชฺชติ เหตุปจฺจยา.\csromanspacer{} ปญฺจ.}

\prose{น-อกุสลํ นสรณญฺจ น-อพฺยากตํ นสรณญฺจ ธมฺมํ ปฏิจฺจ นกุสโล นสรโณ ธมฺโม อุปฺปชฺชติ เหตุปจฺจยา.\csromanspacer{} ปญฺจ.}

\prose{นกุสลํ นสรณญฺจ น-อกุสลํ นสรณญฺจ ธมฺมํ ปฏิจฺจ นกุสโล นสรโณ ธมฺโม อุปฺปชฺชติ เหตุปจฺจยา.\csromanspacer{} ตีณิ. (สํขิตฺตํ.)}

\prose{เหตุยา เอกวีส, อารมฺมเณ สตฺตรส ฯเปฯ อวิคเต เอกวีส. (สพฺพตฺถ วิตฺถาโร.)}

\prose{นกุสลํ น-อรณํ ธมฺมํ ปฏิจฺจ นกุสโล น-อรโณ ธมฺโม อุปฺปชฺชติ เหตุปจฺจยา. (สํขิตฺตํ.) เหตุยา นว.}

\nitthitam{กุสลตฺติกปิฏฺฐิทุกํ นิฏฺฐิตํ.}
\csromanheader{2}{2. เวทนาตฺติก 1. เหตุทุก}
\proseitem{52}{นสุขาย เวทนาย สมฺปยุตฺตํ นเหตุํ ธมฺมํ ปฏิจฺจ นสุขาย เวทนาย สมฺปยุตฺโต นเหตุ ธมฺโม อุปฺปชฺชติ เหตุปจฺจยา.\csromanspacer{}\csromanspacer{}นสุขาย เวทนาย สมฺปยุตฺตํ นเหตุํ ธมฺมํ ปฏิจฺจ นทุกฺขาย เวทนาย สมฺปยุตฺโต นเหตุ ธมฺโม อุปฺปชฺชติ เหตุปจฺจยา.\csromanspacer{}\csromanspacer{}นสุขาย เวทนาย สมฺปยุตฺตํ นเหตุํ ธมฺมํ ปฏิจฺจ น-อทุกฺขมสุขาย เวทนาย สมฺปยุตฺโต นเหตุ ธมฺโม อุปฺปชฺชติ เหตุปจฺจยา ฯเปฯ.\csromanspacer{} สตฺต.}

\prose{นทุกฺขาย เวทนาย สมฺปยุตฺตํ นเหตุํ ธมฺมํ ปฏิจฺจ นทุกฺขาย เวทนาย สมฺปยุตฺโต นเหตุ ธมฺโม อุปฺปชฺชติ เหตุปจฺจยา.\csromanspacer{} สตฺต. (สํขิตฺตํ.)}

\prose{น-อทุกฺขมสุขาย เวทนาย สมฺปยุตฺตํ นเหตุํ ธมฺมํ ปฏิจฺจ.\csromanspacer{} สตฺต. (สํขิตฺตํ.)}

\prose{นสุขาย เวทนาย สมฺปยุตฺตํ นนเหตุํ ธมฺมํ ปฏิจฺจ นสุขาย เวทนาย สมฺปยุตฺโต นนเหตุ ธมฺโม อุปฺปชฺชติ เหตุปจฺจยา.\csromanspacer{} ปญฺจ.}

\csromanheader{2}{ธมฺมปจฺจนีย ติกทุกปฏฺฐานปาฬิ \textbf{3. วิปากตฺติกาทิ 1. เหตุทุก}}
\proseitem{53}{\csromanfolio{132}นวิปากํ นเหตุํ ธมฺมํ ปฏิจฺจ ฯเปฯ.\csromanspacer{} นวิปากธมฺมธมฺมํ นเหตุํ ธมฺมํ ปฏิจฺจ ฯเปฯ.\csromanspacer{} นเนววิปากนวิปากธมฺมธมฺมํ นเหตุํ ธมฺมํ ปฏิจฺจ.}

\proseitem{54}{น-อุปาทินฺนุปาทานิยํ นเหตุํ ธมฺมํ ปฏิจฺจ ฯเปฯ.\csromanspacer{} น-อนุปาทินฺนุปาทานิยํ นเหตุํ ธมฺมํ ปฏิจฺจ ฯเปฯ.\csromanspacer{} น-อนุปาทินฺน-อนุปาทานิยํ นเหตุํ ธมฺมํ ปฏิจฺจ.}

\proseitem{55}{นสํกิลิฏฺฐสํกิเลสิกํ นเหตุํ ธมฺมํ ปฏิจฺจ ฯเปฯ.\csromanspacer{} น-อสํกิลิฏฺฐสํกิเลสิกํ นเหตุํ ธมฺมํ ปฏิจฺจ ฯเปฯ.\csromanspacer{} น-อสํกิลิฏฺฐ-อสํกิเลสิกํ นเหตุํ ธมฺมํ ปฏิจฺจ.}

\proseitem{56}{นสวิตกฺกสวิจารํ นเหตุํ ธมฺมํ ปฏิจฺจ ฯเปฯ.\csromanspacer{} น-อวิตกฺกวิจารมตฺตํ นเหตุํ ธมฺมํ ปฏิจฺจ ฯเปฯ.\csromanspacer{} น-อวิตกฺก-อวิจารํ นเหตุํ ธมฺมํ ปฏิจฺจ.}

\proseitem{57}{นปีติสหคตํ นเหตุํ ธมฺมํ ปฏิจฺจ ฯเปฯ.\csromanspacer{} นสุขสหคตํ นเหตุํ ธมฺมํ ปฏิจฺจ ฯเปฯ.\csromanspacer{} น-อุเปกฺขาสหคตํ นเหตุํ ธมฺมํ ปฏิจฺจ.}

\proseitem{58}{นทสฺสเนน ปหาตพฺพํ นเหตุํ ธมฺมํ ปฏิจฺจ ฯเปฯ.\csromanspacer{} นภาวนาย ปหาตพฺพํ นเหตุํ ธมฺมํ ปฏิจฺจ ฯเปฯ.\csromanspacer{} นเนวทสฺสเนน นภาวนาย ปหาตพฺพํ นเหตุํ ธมฺมํ ปฏิจฺจ ฯเปฯ.\csromanspacer{} นทสฺสเนน ปหาตพฺพเหตุกํ นเหตุํ ธมฺมํ ปฏิจฺจ ฯเปฯ.\csromanspacer{} นภาวนาย ปหาตพฺพเหตุกํ นเหตุํ ธมฺมํ ปฏิจฺจ ฯเปฯ.\csromanspacer{} นเนวทสฺสเนน นภาวนาย ปหาตพฺพเหตุกํ นเหตุํ ธมฺมํ ปฏิจฺจ.}

\proseitem{59}{น-อาจยคามิํ นเหตุํ ธมฺมํ ปฏิจฺจ ฯเปฯ.\csromanspacer{} น-อปจยคามิํ นเหตุํ ธมฺมํ ปฏิจฺจ ฯเปฯ.\csromanspacer{} นเนวาจยคามินาปจยคามิํ นเหตุํ ธมฺมํ ปฏิจฺจ.}

\prose{นเสกฺขํ นเหตุํ ธมฺมํ ปฏิจฺจ ฯเปฯ.\csromanspacer{} น-อเสกฺขํ นเหตุํ ธมฺมํ ปฏิจฺจ ฯเปฯ.\csromanspacer{} นเนวเสกฺขนาเสกฺขํ นเหตุํ ธมฺมํ ปฏิจฺจ.}

\proseitem{60}{นปริตฺตํ นเหตุํ ธมฺมํ ปฏิจฺจ ฯเปฯ.\csromanspacer{} นมหคฺคตํ นเหตุํ ธมฺมํ ปฏิจฺจ ฯเปฯ.\csromanspacer{} น-อปฺปมาณํ นเหตุํ ธมฺมํ ปฏิจฺจ.}

\proseitem{61}{นปริตฺตารมฺมณํ นเหตุํ ธมฺมํ ปฏิจฺจ ฯเปฯ.\csromanspacer{} นมหคฺคตารมฺมณํ นเหตุํ ธมฺมํ ปฏิจฺจ ฯเปฯ.\csromanspacer{} น-อปฺปมาณารมฺมณํ นเหตุํ ธมฺมํ ปฏิจฺจ.}

\csromanheader{2}{22. สนิทสฺสนตฺติก 1. เหตุทุกาทิ}
\proseitem{62}{\csromanfolio{133}นหีนํ นเหตุํ ธมฺมํ ปฏิจฺจ ฯเปฯ.\csromanspacer{} นมชฺฌิมํ นเหตุํ ธมฺมํ ปฏิจฺจ ฯเปฯ.\csromanspacer{} นปณีตํ นเหตุํ ธมฺมํ ปฏิจฺจ.}

\proseitem{63}{นมิจฺฉตฺตนิยตํ นเหตุํ ธมฺมํ ปฏิจฺจ ฯเปฯ.\csromanspacer{} นสมฺมตฺตนิยตํ นเหตุํ ธมฺมํ ปฏิจฺจ ฯเปฯ.\csromanspacer{} น-อนิยตํ นเหตุํ ธมฺมํ ปฏิจฺจ.}

\proseitem{64}{นมคฺคารมฺมณํ นเหตุํ ธมฺมํ ปฏิจฺจ ฯเปฯ.\csromanspacer{} นมคฺคเหตุกํ นเหตุํ ธมฺมํ ปฏิจฺจ ฯเปฯ.\csromanspacer{} นมคฺคาธิปติํ นเหตุํ ธมฺมํ ปฏิจฺจ.}

\proseitem{65}{น-อนุปฺปนฺนํ นเหตุํ ธมฺมํ ปฏิจฺจ ฯเปฯ.\csromanspacer{} น-อุปฺปาทิํ นเหตุํ ธมฺมํ ปฏิจฺจ.}

\proseitem{66}{น-อตีตํ นเหตุํ ธมฺมํ ปฏิจฺจ ฯเปฯ.\csromanspacer{} น-อนาคตํ นเหตุํ ธมฺมํ ปฏิจฺจ.}

\proseitem{67}{น-อตีตารมฺมณํ นเหตุํ ธมฺมํ ปฏิจฺจ ฯเปฯ.\csromanspacer{} น-อนาคตารมฺมณํ นเหตุํ ธมฺมํ ปฏิจฺจ ฯเปฯ.\csromanspacer{} นปจฺจุปฺปนฺนารมฺมณํ นเหตุํ ธมฺมํ ปฏิจฺจ.}

\proseitem{68}{น-อชฺฌตฺตํ นเหตุํ ธมฺมํ ปฏิจฺจ ฯเปฯ.\csromanspacer{} นพหิทฺธา นเหตุํ ธมฺมํ ปฏิจฺจ.}

\prose{น-อชฺฌตฺตารมฺมณํ นเหตุํ ธมฺมํ ปฏิจฺจ ฯเปฯ.\csromanspacer{} นพหิทฺธารมฺมณํ นเหตุํ ธมฺมํ ปฏิจฺจ.}

\csromanheader{2}{22. สนิทสฺสนตฺติก 1. เหตุทุกาทิ}
\proseitem{69}{นสนิทสฺสนสปฺปฏิฆํ นเหตุํ ธมฺมํ ปฏิจฺจ นสนิทสฺสนสปฺปฏิโฆ นเหตุ ธมฺโม อุปฺปชฺชติ เหตุปจฺจยา.\csromanspacer{}\csromanspacer{}นสนิทสฺสนสปฺปฏิฆํ นเหตุํ ธมฺมํ ปฏิจฺจ น-อนิทสฺสนสปฺปฏิโฆ นเหตุ ธมฺโม อุปฺปชฺชติ เหตุปจฺจยา.\csromanspacer{}\csromanspacer{}นสนิทสฺสนสปฺปฏิฆํ นเหตุํ ธมฺมํ ปฏิจฺจ น-อนิทสฺสน-อปฺปฏิโฆ นเหตุ ธมฺโม อุปฺปชฺชติ เหตุปจฺจยา.\csromanspacer{}\csromanspacer{}นสนิทสฺสนสปฺปฏิฆํ นเหตุํ ธมฺมํ ปฏิจฺจ นสนิทสฺสนสปฺปฏิโฆ นเหตุ จ น-อนิทสฺสน-อปฺปฏิโฆ นเหตุ จ ธมฺมา อุปฺปชฺชนฺติ เหตุปจฺจยา.\csromanspacer{}\csromanspacer{}นสนิทสฺสนสปฺปฏิฆํ นเหตุํ ธมฺมํ ปฏิจฺจ น-อนิทสฺสนสปฺปฏิโฆ นเหตุ จ น-อนิทสฺสน-อปฺปฏิโฆ นเหตุ จ ธมฺมา อุปฺปชฺชนฺติ เหตุปจฺจยา.\csromanspacer{}\csromanspacer{}นสนิทสฺสนสปฺปฏิฆํ นเหตุํ ธมฺมํ ปฏิจฺจ นสนิทสฺสนสปฺปฏิโฆ นเหตุ จ น-อนิทสฺสนสปฺปฏิโฆ นเหตุ จ ธมฺมา อุปฺปชฺชนฺติ เหตุปจฺจยา. (6)}

\prose{น-อนิทสฺสนสปฺปฏิฆํ นเหตุํ ธมฺมํ ปฏิจฺจ น-อนิทสฺสนสปฺปฏิโฆ นเหตุ ธมฺโม อุปฺปชฺชติ เหตุปจฺจยา.\csromanspacer{}\csromanspacer{}น-อนิทสฺสนสปฺปฏิฆํ นเหตุํ ธมฺมํ ปฏิจฺจ}

\vspace{\tipitakaparskip}
\csromancenter{ธมฺมปจฺจนีย ติกทุกปฏฺฐานปาฬิ นสนิทสฺสนสปฺปฏิโฆ นเหตุ ธมฺโม อุปฺปชฺชติ เหตุปจฺจยา.\csromanspacer{}\csromanspacer{}น-อนิทสฺสนสปฺปฏิฆํ นเหตุํ ธมฺมํ ปฏิจฺจ น-อนิทสฺสน-อปฺปฏิโฆ นเหตุ ธมฺโม อุปฺปชฺชติ เหตุปจฺจยา.\csromanspacer{}\csromanspacer{}น-อนิทสฺสนสปฺปฏิฆํ นเหตุํ ธมฺมํ ปฏิจฺจ นสนิทสฺสนสปฺปฏิโฆ นเหตุ จ น-อนิทสฺสน-อปฺปฏิโฆ นเหตุ จ ธมฺมา อุปฺปชฺชนฺติ เหตุปจฺจยา.\csromanspacer{}\csromanspacer{}น-อนิทสฺสนสปฺปฏิฆํ นเหตุํ ธมฺมํ ปฏิจฺจ น-อนิทสฺสนสปฺปฏิโฆ นเหตุ จ น-อนิทสฺสน-อปฺปฏิโฆ นเหตุ จ ธมฺมา อุปฺปชฺชนฺติ เหตุปจฺจยา.\csromanspacer{}\csromanspacer{}น-อนิทสฺสน-อปฺปฏิฆํ นเหตุํ ธมฺมํ ปฏิจฺจ นสนิทสฺสนสปฺปฏิโฆ นเหตุ จ น-อนิทสฺสนสปฺปฏิโฆ นเหตุ จ ธมฺมา อุปฺปชฺชนฺติ เหตุปจฺจยา. (6)}
\prosecont{\csromanfolio{134}น-อนิทสฺสน-อปฺปฏิฆํ นเหตุํ ธมฺมํ ปฏิจฺจ น-อนิทสฺสน-อปฺปฏิโฆ นเหตุ ธมฺโม อุปฺปชฺชติ เหตุปจฺจยา.\csromanspacer{} ฉ.}

\prose{นสนิทสฺสนสปฺปฏิฆํ นเหตุญฺจ น-อนิทสฺสน-อปฺปฏิฆํ นเหตุญฺจ ธมฺมํ ปฏิจฺจ นสนิทสฺสนสปฺปฏิโฆ นเหตุ ธมฺโม อุปฺปชฺชติ เหตุปจฺจยา.\csromanspacer{} ฉ.}

\prose{นสนิทสฺสนสปฺปฏิฆํ นเหตุญฺจ น-อนิทสฺสนสปฺปฏิฆํ นเหตุญฺจ ธมฺมํ ปฏิจฺจ นสนิทสฺสนสปฺปฏิโฆ นเหตุ ธมฺโม อุปฺปชฺชติ เหตุปจฺจยา.\csromanspacer{} ฉ. (สํขิตฺตํ.)}

\prose{เหตุยา ติํส, อารมฺมเณ นว. (สพฺพตฺถ วิตฺถาโร.)}

\proseitem{70}{นสนิทสฺสนสปฺปฏิฆํ นนเหตุํ ธมฺมํ ปฏิจฺจ นสนิทสฺสนสปฺปฏิโฆ นนเหตุ ธมฺโม อุปฺปชฺชติ เหตุปจฺจยา. (นว.)}

\prose{นสนิทสฺสนสปฺปฏิฆํ นสเหตุกํ ธมฺมํ ปฏิจฺจ.}

\prose{นสนิทสฺสนสปฺปฏิฆํ น-อเหตุกํ ธมฺมํ ปฏิจฺจ.}

\proseitem{71}{นสนิทสฺสนสปฺปฏิฆํ นเหตุสมฺปยุตฺตํ ธมฺมํ ปฏิจฺจ.}

\prose{นสนิทสฺสนสปฺปฏิฆํ นเหตุวิปฺปยุตฺตํ ธมฺมํ ปฏิจฺจ.}

\proseitem{72}{นสนิทสฺสนสปฺปฏิฆํ นเหตุญฺเจว น-อเหตุกญฺจ ธมฺมํ ปฏิจฺจ ฯเปฯ.\csromanspacer{} นสนิทสฺสนสปฺปฏิฆํ น-อเหตุกญฺเจว นนเหตุญฺจ ธมฺมํ ปฏิจฺจ ฯเปฯ.\csromanspacer{} นสนิทสฺสนสปฺปฏิฆํ นเหตุญฺเจว นเหตุวิปฺปยุตฺตญฺจ ธมฺมํ ปฏิจฺจ.}

\prose{นสนิทสฺสนสปฺปฏิฆํ นเหตุวิปฺปยุตฺตญฺเจว นนเหตุญฺจ ธมฺมํ ปฏิจฺจ.}

\csromanheader{2}{22. สนิทสฺสนตฺติก 55. มหนฺตรทุก}
\prose{\csromanfolio{135}นสนิทสฺสนสปฺปฏิฆํ นเหตุํ นสเหตุกํ ธมฺมํ ปฏิจฺจ.}

\prose{นสนิทสฺสนสปฺปฏิฆํ นเหตุํ น-อเหตุกํ ธมฺมํ ปฏิจฺจ.}

\csromanheader{2}{22. สนิทสฺสนตฺติก 7. จูฬนฺตรทุก}
\proseitem{73}{นสนิทสฺสนสปฺปฏิฆํ น-อปฺปจฺจยํ ธมฺมํ ปฏิจฺจ ฯเปฯ.\csromanspacer{} นสนิทสฺสนสปฺปฏิฆํ น-อสงฺขตํ ธมฺมํ ปฏิจฺจ.}

\proseitem{74}{นสนิทสฺสนสปฺปฏิฆํ นสนิทสฺสนํ ธมฺมํ ปฏิจฺจ.}

\proseitem{75}{นสนิทสฺสนสปฺปฏิฆํ นสปฺปฏิฆํ ธมฺมํ ปฏิจฺจ.}

\prose{นสนิทสฺสนสปฺปฏิฆํ น-อปฺปฏิฆํ ธมฺมํ ปฏิจฺจ.}

\proseitem{76}{นสนิทสฺสนสปฺปฏิฆํ นรูปิํ ธมฺมํ ปฏิจฺจ.\csromanspacer{} นสนิทสฺสนสปฺปฏิฆํ น-อรูปิํ ธมฺมํ ปฏิจฺจ.}

\proseitem{77}{นสนิทสนสปฺปฏิฆํ นโลกิยํ ธมฺมํ ปฏิจฺจ.}

\prose{นสนิทสฺสนสปฺปฏิฆํ นโลกุตฺตรํ ธมฺมํ ปฏิจฺจ.}

\proseitem{78}{นสนิทสฺสนสปฺปฏิฆํ นเกนจิ วิญฺเญยฺยํ ธมฺมํ ปฏิจฺจ ฯเปฯ.\csromanspacer{} นสนิทสฺสนสปฺปฏิฆํ นนเกนจิ วิญฺเญยฺยํ ธมฺมํ ปฏิจฺจ.}

\csromanheader{2}{22. สนิทสฺสนตฺติก 14. อาสวโคจฺฉกาทิ}
\proseitem{79}{นสนิทสฺสนสปฺปฏิฆํ โน-อาสวํ ธมฺมํ ปฏิจฺจ.}

\proseitem{80}{นสนิทสฺสนสปฺปฏิฆํ โนสญฺโญชนํ ธมฺมํ ปฏิจฺจ ฯเปฯ.\csromanspacer{} นสนิทสฺสนสปฺปฏิฆํ โนคนฺถํ ธมฺมํ ปฏิจฺจ ฯเปฯ.\csromanspacer{} นสนิทสฺสนสปฺปฏิฆํ โน-โอฆํ ธมฺมํ ปฏิจฺจ ฯเปฯ.\csromanspacer{} นสนิทสฺสนสปฺปฏิฆํ โนโยคํ ธมฺมํ ปฏิจฺจ ฯเปฯ.\csromanspacer{} นสนิทสฺสนสปฺปฏิฆํ โนนีวรณํ ธมฺมํ ปฏิจฺจ ฯเปฯ.\csromanspacer{} นสนิทสฺสนสปฺปฏิฆํ โนปรามาสํ ธมฺมํ ปฏิจฺจ.}

\csromanheader{2}{22. สนิทสฺสนตฺติก 55. มหนฺตรทุก}
\proseitem{81}{นสนิทสฺสนสปฺปฏิฆํ นสารมฺมณํ ธมฺมํ ปฏิจฺจ.}

\gambhira{ธมฺมปจฺจนีย ติกทุกปฏฺฐานปาฬิ}
\proseitem{82}{\csromanfolio{136}นสนิทสฺสนสปฺปฏิฆํ นจิตฺตํ ธมฺมํ ปฏิจฺจ.}

\prose{(นโนจิตฺตปทํ น ลพฺภติ.)}

\proseitem{83}{นสนิทสฺสนสปฺปฏิฆํ นเจตสิกํ ธมฺมํ ปฏิจฺจ.}

\proseitem{84}{นสนิทสฺสนสปฺปฏิฆํ นจิตฺตสมฺปยุตฺตํ ธมฺมํ ปฏิจฺจ ฯเปฯ.\csromanspacer{} นสนิทสฺสนสปฺปฏิฆํ นจิตฺตสํสฏฺฐํ ธมฺมํ ปฏิจฺจ.}

\proseitem{85}{นสนิทสฺสนสปฺปฏิฆํ นจิตฺตสมุฏฺฐานํ ธมฺมํ ปฏิจฺจ ฯเปฯ.\csromanspacer{} นสนิทสฺสนสปฺปฏิฆํ นจิตฺตสหภุํ ธมฺมํ ปฏิจฺจ ฯเปฯ.\csromanspacer{} นสนิทสฺสนสปฺปฏิฆํ นจิตฺตานุปริวตฺติํ ธมฺมํ ปฏิจฺจ ฯเปฯ.\csromanspacer{} นสนิทสฺสนสปฺปฏิฆํ นจิตฺตสํสฏฺฐสมุฏฺฐานํ ธมฺมํ ปฏิจฺจ ฯเปฯ.\csromanspacer{} นสนิทสฺสนสปฺปฏิฆํ นจิตฺตสํสฏฺฐสมุฏฺฐานสหภุํ ธมฺมํ ปฏิจฺจ ฯเปฯ.\csromanspacer{} นสนิทสฺสนสปฺปฏิฆํ นจิตฺตสํสฏฺฐสมุฏฺฐานานุปริวตฺติํ ธมฺมํ ปฏิจฺจ.}

\proseitem{86}{นสนิทสฺสนสปฺปฏิฆํ น-อชฺฌตฺติกํ ธมฺมํ ปฏิจฺจ.}

\prose{นสนิทสฺสนสปฺปฏิฆํ นพาหิรํ ธมฺมํ ปฏิจฺจ.}

\proseitem{87}{นสนิทสฺสนสปฺปฏิฆํ น-อุปาทา ธมฺมํ ปฏิจฺจ.}

\proseitem{88}{นสนิทสฺสนสปฺปฏิฆํ น-อุปาทินฺนํ ธมฺมํ ปฏิจฺจ ฯเปฯ.\csromanspacer{} นสนิทสฺสนสปฺปฏิฆํ น-อนุปาทินฺนํ ธมฺมํ ปฏิจฺจ.}

\csromanheader{2}{22. สนิทสฺสนตฺติก 69. อุปาทานาทิทุก}
\proseitem{89}{นสนิทสฺสนสปฺปฏิฆํ โน-อุปาทานํ ธมฺมํ ปฏิจฺจ ฯเปฯ.\csromanspacer{} นสนิทสฺสนสปฺปฏิฆํ โนกิเลสํ ธมฺมํ ปฏิจฺจ.}

\csromanheader{2}{22. สนิทสฺสนตฺติก 83. ปิฏฺฐิทุก}
\proseitem{90}{นสนิทสฺสนสปฺปฏิฆํ นทสฺสเนน ปหาตพฺพํ ธมฺมํ ปฏิจฺจ.}

\prose{เหตุยา ติํส, อารมฺมเณ นว ฯเปฯ อญฺญมญฺเญ ปญฺจวีส ฯเปฯ อวิคเต ติํส. (ปิฏฺฐิทุกํ วิตฺถาเรตพฺพํ.)}

\proseitem{91}{นสนิทสฺสนสปฺปฏิฆํ นสรณํ ธมฺมํ ปฏิจฺจ นสนิทสฺสนสปฺปฏิโฆ นสรโณ ธมฺโม อุปฺปชฺชติ เหตุปจฺจยา.\csromanspacer{}\csromanspacer{}นสนิทสฺสนสปฺปฏิฆํ นสรณํ}

\vspace{\tipitakaparskip}
\csromancenter{สนิทสฺสนตฺติก 83.\csromanspacer{} ปิฏฺฐิทุก ธมฺมํ ปฏิจฺจ น-อนิทสฺสนสปฺปฏิโฆ นสรโณ ธมฺโม อุปฺปชฺชติ เหตุปจฺจยา.\csromanspacer{}\csromanspacer{}นสนิทสฺสนสปฺปฏิฆํ นสรณํ ธมฺมํ ปฏิจฺจ น-อนิทสฺสน-อปฺปฏิโฆ นสรโณ ธมฺโม อุปฺปชฺชติ เหตุปจฺจยา ฯเปฯ.\csromanspacer{} ฉ.}
\prose{\csromanfolio{137}น-อนิทสฺสนสปฺปฏิฆํ นสรณํ ธมฺมํ ปฏิจฺจ น-อนิทสฺสนสปฺปฏิโฆ นสรโณ ธมฺโม อุปฺปชฺชติ เหตุปจฺจยา.\csromanspacer{}\csromanspacer{}น-อนิทสฺสนสปฺปฏิฆํ นสรณํ ธมฺมํ ปฏิจฺจ นสนิทสฺสนสปฺปฏิโฆ นสรโณ ธมฺโม อุปฺปชฺชติ เหตุปจฺจยา ฯเปฯ.\csromanspacer{} ฉ.}

\prose{น-อนิทสฺสนสปฺปฏิฆํ นสรณํ ธมฺมํ ปฏิจฺจ น-อนิทสฺสน-อปฺปฏิโฆ นสรโณ ธมฺโม อุปฺปชฺชติ เหตุปจฺจยา.\csromanspacer{} ฉ. นสนิทสฺสนสปฺปฏิฆํ นสรณญฺจ น-อนิทสฺสน-อปฺปฏิฆํ นสรณญฺจ ธมฺมํ ปฏิจฺจ นสนิทสฺสนสปฺปฏิโฆ นสรโณ ธมฺโม อุปฺปชฺชติ เหตุปจฺจยา.\csromanspacer{} ฉ. นสนิทสฺสนสปฺปฏิฆํ นสรณญฺจ น-อนิทสฺสนสปฺปฏิฆํ นสรณญฺจ ธมฺมํ ปฏิจฺจ นสนิทสฺสนสปฺปฏิโฆ นสรโณ ธมฺโม อุปฺปชฺชติ เหตุปจฺจยา.\csromanspacer{} ฉ.}

\prose{เหตุยา ติํส, อารมฺมเณ นว ฯเปฯ อวิคเต ติํส. (สพฺพตฺถ วิตฺถาโร.)}

\proseitem{92}{นสนิทสฺสนสปฺปฏิฆํ น-อรณํ ธมฺมํ ปฏิจฺจ นสนิทสฺสนสปฺปฏิโฆ น-อรโณ ธมฺโม อุปฺปชฺชติ เหตุปจฺจยา.\csromanspacer{}\csromanspacer{}นสนิทสฺสนสปฺปฏิฆํ น-อรณํ ธมฺมํ ปฏิจฺจ น-อนิทสฺสนสปฺปฏิโฆ น-อรโณ ธมฺโม อุปฺปชฺชติ เหตุปจฺจยา.\csromanspacer{}\csromanspacer{}นสนิทสฺสนสปฺปฏิฆํ น-อรณํ ธมฺมํ ปฏิจฺจ นสนิทสฺสนสปฺปฏิโฆ น-อรโณ จ น-อนิทสฺสนสปฺปฏิโฆ น-อรโณ จ ธมฺมา อุปฺปชฺชนฺติ เหตุปจฺจยา. (3)}

\proseitem{93}{น-อนิทสฺสนสปฺปฏิฆํ น-อรณํ ธมฺมํ ปฏิจฺจ น-อนิทสฺสนสปฺปฏิโฆ น-อรโณ ธมฺโม อุปฺปชฺชติ เหตุปจฺจยา.\csromanspacer{}\csromanspacer{}น-อนิทสฺสนสปฺปฏิฆํ น-อรณํ ธมฺมํ ปฏิจฺจ นสนิทสฺสนสปฺปฏิโฆ น-อรโณ ธมฺโม อุปฺปชฺชติ เหตุปจฺจยา.\csromanspacer{}\csromanspacer{}น-อนิทสฺสนสปฺปฏิฆํ น-อรณํ ธมฺมํ ปฏิจฺจ นสนิทสฺสนสปฺปฏิโฆ น-อรโณ จ น-อนิทสฺสนสปฺปฏิโฆ น-อรโณ จ ธมฺมา อุปฺปชฺชนฺติ เหตุปจฺจยา. (3)}

\gambhira{ธมฺมปจฺจนีย ติกทุกปฏฺฐานปาฬิ}
\proseitem{94}{\csromanfolio{138}นสนิทสฺสนสปฺปฏิฆํ น-อรณญฺจ น-อนิทสฺสนสปฺปฏิฆํ น-อรณญฺจ ธมฺมํ ปฏิจฺจ นสนิทสฺสนสปฺปฏิโฆ น-อรโณ ธมฺโม อุปฺปชฺชติ เหตุปจฺจยา.\csromanspacer{}\csromanspacer{}นสนิทสฺสนสปฺปฏิฆํ น-อรณญฺจ น-อนิทสฺสนสปฺปฏิฆํ น-อรณญฺจ ธมฺมํ ปฏิจฺจ น-อนิทสฺสนสปฺปฏิโฆ น-อรโณ ธมฺโม อุปฺปชฺชติ เหตุปจฺจยา.\csromanspacer{}\csromanspacer{}นสนิทสฺสนสปฺปฏิฆํ น-อรณญฺจ น-อนิทสฺสนสปฺปฏิฆํ น-อรณญฺจ ธมฺมํ ปฏิจฺจ นสนิทสฺสนสปฺปฏิโฆ น-อรโณ จ น-อนิทสฺสนสปฺปฏิโฆ น-อรโณ จ ธมฺมา อุปฺปชฺชนฺติ เหตุปจฺจยา. (3)}

\prose{เหตุยา นว, อารมฺมเณ นว ฯเปฯ อวิคเต นว. (สหชาตวาเรปิ ปจฺจยวาเรปิ นิสฺสยวาเรปิ สํสฏฺฐวาเรปิ สมฺปยุตฺตวาเรปิ สพฺพตฺถ นว.)}

\nitthitam{\textbf{ธมฺมปจฺจนีเย ติกทุกปฏฺฐานํ นิฏฺฐิตํ.}}
\pitaka{\textbf{อภิธมฺมปิฏก}}
\csromanheader{2}{\textbf{ธมฺมปจฺจนีย}}
\gambhira{\textbf{ติกติกปฏฺฐานปาฬิ}}
\namakkaram{นโม ตสฺส ภควโต อรหโต สมฺมาสมฺพุทฺธสฺส.}
\csromanheader{2}{1. กุสลตฺติก 1. เวทนาตฺติก}
\csromanheader{2}{1. ปฏิจฺจวาร-เหตุ}
\proseitem{1}{\csromanfolio{139}นกุสลํ นสุขาย เวทนาย สมฺปยุตฺตํ ธมฺมํ ปฏิจฺจ นกุสโล นสุขาย เวทนาย สมฺปยุตฺโต ธมฺโม อุปฺปชฺชติ เหตุปจฺจยา.\csromanspacer{}\csromanspacer{}นกุสลํ นสุขาย เวทนาย สมฺปยุตฺตํ ธมฺมํ ปฏิจฺจ น-อกุสโล นสุขาย เวทนาย สมฺปยุตฺโต ธมฺโม อุปฺปชฺชติ เหตุปจฺจยา.\csromanspacer{}\csromanspacer{}นกุสลํ นสุขาย เวทนาย สมฺปยุตฺตํ ธมฺมํ ปฏิจฺจ น-อพฺยากโต นสุขาย เวทนาย สมฺปยุตฺโต ธมฺโม อุปฺปชฺชติ เหตุปจฺจยา.\csromanspacer{}\csromanspacer{}นกุสลํ นสุขาย เวทนาย สมฺปยุตฺตํ ธมฺมํ ปฏิจฺจ นกุสโล นสุขาย เวทนาย สมฺปยุตฺโต จ น-อพฺยากโต นสุขาย เวทนาย สมฺปยุตฺโต จ ธมฺมา อุปฺปชฺชนฺติ เหตุปจฺจยา.\csromanspacer{}\csromanspacer{}นกุสลํ นสุขาย เวทนาย สมฺปยุตฺตํ ธมฺมํ ปฏิจฺจ นกุสโล นสุขาย เวทนาย สมฺปยุตฺโต จ น-อกุสโล นสุขาย เวทนาย สมฺปยุตฺโต จ ธมฺมา อุปฺปชฺชนฺติ เหตุปจฺจยา. (5)}

\proseitem{2}{น-อกุสลํ นสุขาย เวทนาย สมฺปยุตฺตํ ธมฺมํ ปฏิจฺจ น-อกุสโล นสุขาย เวทนาย สมฺปยุตฺโต ธมฺโม อุปฺปชฺชติ เหตุปจฺจยา.\csromanspacer{}\csromanspacer{}น-อกุสลํ นสุขาย เวทนาย สมฺปยุตฺตํ ธมฺมํ ปฏิจฺจ นกุสโล นสุขาย เวทนาย สมฺปยุตฺโต ธมฺโม อุปฺปชฺชติ เหตุปจฺจยา.\csromanspacer{}\csromanspacer{}น-อกุสลํ นสุขาย เวทนาย สมฺปยุตฺตํ ธมฺมํ ปฏิจฺจ น-อพฺยากโต นสุขาย เวทนาย}

\vspace{\tipitakaparskip}
\csromancenter{ธมฺมปจฺจนีย ติกติกปฏฺฐานปาฬิ สมฺปยุตฺโต ธมฺโม อุปฺปชฺชติ เหตุปจฺจยา.\csromanspacer{}\csromanspacer{}น-อกุสลํ นสุขาย เวทนาย สมฺปยุตฺตํ ธมฺมํ ปฏิจฺจ น-อกุสโล นสุขาย เวทนาย สมฺปยุตฺโต จ น-อพฺยากโต นสุขาย เวทนาย สมฺปยุตฺโต จ ธมฺมา อุปฺปชฺชนฺติ เหตุปจฺจยา.\csromanspacer{}\csromanspacer{}น-อกุสลํ นสุขาย เวทนาย สมฺปยุตฺตํ ธมฺมํ ปฏิจฺจ นกุสโล นสุขาย เวทนาย สมฺปยุตฺโต จ น-อกุสโล นสุขาย เวทนาย สมฺปยุตฺโต จ ธมฺมา อุปฺปชฺชนฺติ เหตุปจฺจยา. (5)}
\proseitem{3}{\csromanfolio{140}น-อพฺยากตํ นสุขาย เวทนาย สมฺปยุตฺตํ ธมฺมํ ปฏิจฺจ น-อพฺยากโต นสุขาย เวทนาย สมฺปยุตฺโต ธมฺโม อุปฺปชฺชติ เหตุปจฺจยา.\csromanspacer{}\csromanspacer{}น-อพฺยากตํ นสุขาย เวทนาย สมฺปยุตฺตํ ธมฺมํ ปฏิจฺจ นกุสโล นสุขาย เวทนาย สมฺปยุตฺโต ธมฺโม อุปฺปชฺชติ เหตุปจฺจยา.\csromanspacer{}\csromanspacer{}น-อพฺยากตํ นสุขาย เวทนาย สมฺปยุตฺตํ ธมฺมํ ปฏิจฺจ น-อกุสโล นสุขาย เวทนาย สมฺปยุตฺโต ธมฺโม อุปฺปชฺชติ เหตุปจฺจยา.\csromanspacer{}\csromanspacer{}น-อพฺยากตํ นสุขาย เวทนาย สมฺปยุตฺตํ ธมฺมํ ปฏิจฺจ นกุสโล นสุขาย เวทนาย สมฺปยุตฺโต จ น-อพฺยากโต นสุขาย เวทนาย สมฺปยุตฺโต จ ธมฺมา อุปฺปชฺชนฺติ เหตุปจฺจยา.\csromanspacer{}\csromanspacer{}น-อพฺยากตํ นสุขาย เวทนาย สมฺปยุตฺตํ ธมฺมํ ปฏิจฺจ น-อกุสโล นสุขาย เวทนาย สมฺปยุตฺโต จ น-อพฺยากโต นสุขาย เวทนาย สมฺปยุตฺโต จ ธมฺมา อุปฺปชฺชนฺติ เหตุปจฺจยา.\csromanspacer{}\csromanspacer{}น-อพฺยากตํ นสุขาย เวทนาย สมฺปยุตฺตํ ธมฺมํ ปฏิจฺจ นกุสโล นสุขาย เวทนาย สมฺปยุตฺโต จ น-อกุสโล นสุขาย เวทนาย สมฺปยุตฺโต จ ธมฺมา อุปฺปชฺชนฺติ เหตุปจฺจยา. (6)}

\proseitem{4}{นกุสลํ นสุขาย เวทนาย สมฺปยุตฺตญฺจ น-อพฺยากตํ นสุขาย เวทนาย สมฺปยุตฺตญฺจ ธมฺมํ ปฏิจฺจ นกุสโล นสุขาย เวทนาย สมฺปยุตฺโต ธมฺโม อุปฺปชฺชติ เหตุปจฺจยา.\csromanspacer{}\csromanspacer{}นกุสลํ นสุขาย เวทนาย สมฺปยุตฺตญฺจ น-อพฺยากตํ นสุขาย เวทนาย สมฺปยุตฺตญฺจ ธมฺมํ ปฏิจฺจ น-อกุสโล นสุขาย เวทนาย สมฺปยุตฺโต ธมฺโม อุปฺปชฺชติ เหตุปจฺจยา.\csromanspacer{}\csromanspacer{}นกุสลํ นสุขาย เวทนาย สมฺปยุตฺตญฺจ น-อพฺยากตํ นสุขาย เวทนาย สมฺปยุตฺตญฺจ ธมฺมํ ปฏิจฺจ น-อพฺยากโต นสุขาย เวทนาย สมฺปยุตฺโต ธมฺโม อุปฺปชฺชติ เหตุปจฺจยา.\csromanspacer{}\csromanspacer{}นกุสลํ นสุขาย เวทนาย สมฺปยุตฺตญฺจ น-อพฺยากตํ นสุขาย เวทนาย สมฺปยุตฺตญฺจ ธมฺมํ ปฏิจฺจ นกุสโล นสุขาย เวทนาย สมฺปยุตฺโต จ น-อพฺยากโต นสุขาย เวทนาย สมฺปยุตฺโต จ ธมฺมา อุปฺปชฺชนฺติ}

\vspace{\tipitakaparskip}
\csromancenter{กุสลตฺติก 1.\csromanspacer{} เวทนาตฺติก เหตุปจฺจยา.\csromanspacer{}\csromanspacer{}นกุสลํ นสุขาย เวทนาย สมฺปยุตฺตญฺจ น-อพฺยากตํ นสุขาย เวทนาย สมฺปยุตฺตญฺจ ธมฺมํ ปฏิจฺจ นกุสโล นสุขาย เวทนาย สมฺปยุตฺโต จ น-อกุสโล นสุขาย เวทนาย สมฺปยุตฺโต จ ธมฺมา อุปฺปชฺชนฺติ เหตุปจฺจยา. (5)}
\proseitem{5}{\csromanfolio{141}น-อกุสลํ นสุขาย เวทนาย สมฺปยุตฺตญฺจ น-อพฺยากตํ นสุขาย เวทนาย สมฺปยุตฺตญฺจ ธมฺมํ ปฏิจฺจ นกุสโล นสุขาย เวทนาย สมฺปยุตฺโต ธมฺโม อุปฺปชฺชติ เหตุปจฺจยา.\csromanspacer{}\csromanspacer{}น-อกุสลํ นสุขาย เวทนาย สมฺปยุตฺตญฺจ น-อพฺยากตํ นสุขาย เวทนาย สมฺปยุตฺตญฺจ ธมฺมํ ปฏิจฺจ น-อกุสโล นสุขาย เวทนาย สมฺปยุตฺโต ธมฺโม อุปฺปชฺชติ เหตุปจฺจยา.\csromanspacer{}\csromanspacer{}น-อกุสลํ นสุขาย เวทนาย สมฺปยุตฺตญฺจ น-อพฺยากตํ นสุขาย เวทนาย สมฺปยุตฺตญฺจ ธมฺมํ ปฏิจฺจ น-อพฺยากโต นสุขาย เวทนาย สมฺปยุตฺโต ธมฺโม อุปฺปชฺชติ เหตุปจฺจยา.\csromanspacer{}\csromanspacer{}น-อกุสลํ นสุขาย เวทนาย สมฺปยุตฺตญฺจ น-อพฺยากตํ นสุขาย เวทนาย สมฺปยุตฺตญฺจ ธมฺมํ ปฏิจฺจ น-อกุสโล นสุขาย เวทนาย สมฺปยุตฺโต จ น-อพฺยากโต นสุขาย เวทนาย สมฺปยุตฺโต จ ธมฺมา อุปฺปชฺชนฺติ เหตุปจฺจยา.\csromanspacer{}\csromanspacer{}น-อกุสลํ นสุขาย เวทนาย สมฺปยุตฺตญฺจ น-อพฺยากตํ นสุขาย เวทนาย สมฺปยุตฺตญฺจ ธมฺมํ ปฏิจฺจ นกุสโล นสุขาย เวทนาย สมฺปยุตฺโต จ น-อกุสโล นสุขาย เวทนาย สมฺปยุตฺโต จ ธมฺมา อุปฺปชฺชนฺติ เหตุปจฺจยา. (5) นกุสลํ นสุขาย เวทนาย สมฺปยุตฺตญฺจ น-อกุสลํ นสุขาย เวทนาย สมฺปยุตฺตญฺจ ธมฺมํ ปฏิจฺจ นกุสโล นสุขาย เวทนาย สมฺปยุตฺโต ธมฺโม อุปฺปชฺชติ เหตุปจฺจยา.\csromanspacer{}\csromanspacer{}นกุสลํ นสุขาย เวทนาย สมฺปยุตฺตญฺจ น-อกุสลํ นสุขาย เวทนาย สมฺปยุตฺตญฺจ ธมฺมํ ปฏิจฺจ น-อกุสโล นสุขาย เวทนาย สมฺปยุตฺโต ธมฺโม อุปฺปชฺชติ เหตุปจฺจยา.\csromanspacer{}\csromanspacer{}นกุสลํ นสุขาย เวทนาย สมฺปยุตฺตญฺจ น-อกุสลํ นสุขาย เวทนาย สมฺปยุตฺตญฺจ ธมฺมํ ปฏิจฺจ นกุสโล นสุขาย เวทนาย สมฺปยุตฺโต จ น-อกุสโล นสุขาย เวทนาย สมฺปยุตฺโต จ ธมฺมา อุปฺปชฺชนฺติ เหตุปจฺจยา. (3)}

\proseitem{6}{เหตุยา เอกูนติํส, อารมฺมเณ จตุวีส ฯเปฯ อวิคเต เอกูนติํส. (สหชาตวาเรปิ ฯเปฯ ปญฺหาวาเรปิ วิตฺถาโร.)}

\csromanheader{2}{ธมฺมปจฺจนีย ติกติกปฏฺฐานปาฬิ \textbf{ทุกฺขายเวทนายสมฺปยุตฺตปท-เหตุ}}
\proseitem{7}{\csromanfolio{142}นกุสลํ นทุกฺขาย เวทนาย สมฺปยุตฺตํ ธมฺมํ ปฏิจฺจ นกุสโล นทุกฺขาย เวทนาย สมฺปยุตฺโต ธมฺโม อุปฺปชฺชติ เหตุปจฺจยา.\csromanspacer{}\csromanspacer{}นกุสลํ นทุกฺขาย เวทนาย สมฺปยุตฺตํ ธมฺมํ ปฏิจฺจ น-อกุสโล นทุกฺขาย เวทนาย สมฺปยุตฺโต ธมฺโม อุปฺปชฺชติ เหตุปจฺจยา.\csromanspacer{}\csromanspacer{}นกุสลํ นทุกฺขาย เวทนาย สมฺปยุตฺตํ ธมฺมํ ปฏิจฺจ น-อพฺยากโต นทุกฺขาย เวทนาย สมฺปยุตฺโต ธมฺโม อุปฺปชฺชติ เหตุปจฺจยา.\csromanspacer{}\csromanspacer{}นกุสลํ นทุกฺขาย เวทนาย สมฺปยุตฺตํ ธมฺมํ ปฏิจฺจ นกุสโล นทุกฺขาย เวทนาย สมฺปยุตฺโต จ น-อพฺยากโต นทุกฺขาย เวทนาย สมฺปยุตฺโต จ ธมฺมา อุปฺปชฺชนฺติ เหตุปจฺจยา.\csromanspacer{}\csromanspacer{}นกุสลํ นทุกฺขาย เวทนาย สมฺปยุตฺตํ ธมฺมํ ปฏิจฺจ นกุสโล นทุกฺขาย เวทนาย สมฺปยุตฺโต จ น-อกุสโล นทุกฺขาย เวทนาย สมฺปยุตฺโต จ ธมฺมา อุปฺปชฺชนฺติ เหตุปจฺจยา. (5)}

\prose{น-อกุสลํ นทุกฺขาย เวทนาย สมฺปยุตฺตํ ธมฺมํ ปฏิจฺจ น-อกุสโล นทุกฺขาย เวทนาย สมฺปยุตฺโต ธมฺโม อุปฺปชฺชติ เหตุปจฺจยา.\csromanspacer{} ปญฺจ.}

\prose{น-อพฺยากตํ นทุกฺขาย เวทนาย สมฺปยุตฺตํ ธมฺมํ ปฏิจฺจ น-อพฺยากโต นทุกฺขาย เวทนาย สมฺปยุตฺโต ธมฺโม อุปฺปชฺชติ เหตุปจฺจยา.\csromanspacer{} ฉ.}

\prose{นกุสลํ นทุกฺขาย เวทนาย สมฺปยุตฺตญฺจ น-อพฺยากตํ นทุกฺขาย เวทนาย สมฺปยุตฺตญฺจ ธมฺมํ ปฏิจฺจ นกุสโล นทุกฺขาย เวทนาย สมฺปยุตฺโต ธมฺโม อุปฺปชฺชติ เหตุปจฺจยา.\csromanspacer{} ปญฺจ.}

\prose{น-อกุสลํ นทุกฺขาย เวทนาย สมฺปยุตฺตญฺจ น-อพฺยากตํ นทุกฺขาย เวทนาย สมฺปยุตฺตญฺจ ธมฺมํ ปฏิจฺจ นกุสโล นทุกฺขาย เวทนาย สมฺปยุตฺโต ธมฺโม อุปฺปชฺชติ เหตุปจฺจยา.\csromanspacer{} ปญฺจ.}

\prose{นกุสลํ นทุกฺขาย เวทนาย สมฺปยุตฺตญฺจ น-อกุสลํ นทุกฺขาย เวทนาย สมฺปยุตฺตญฺจ ธมฺมํ ปฏิจฺจ นกุสโล นทุกฺขาย เวทนาย สมฺปยุตฺโต ธมฺโม อุปฺปชฺชติ เหตุปจฺจยา.\csromanspacer{} ตีณิ.}

\prose{เหตุยา เอกูนติํส ฯเปฯ อวิคเต เอกูนติํส. (สพฺพตฺถ วิตฺถาเรตพฺพํ.)}

\csromanheader{2}{1. กุสลตฺติก 2. วิปากตฺติก \textbf{ตติยปท-เหตุ}}
\proseitem{8}{\csromanfolio{143}นกุสลํ น-อทุกฺขมสุขาย เวทนาย สมฺปยุตฺตํ ธมฺมํ ปฏิจฺจ นกุสโล น-อทุกฺขมสุขาย เวทนาย สมฺปยุตฺโต ธมฺโม อุปฺปชฺชติ เหตุปจฺจยา.\csromanspacer{} ปญฺจ.}

\prose{น-อกุสลํ น-อทุกฺขมสุขาย เวทนาย สมฺปยุตฺตํ ธมฺมํ ปฏิจฺจ น-อกุสโล น-อทุกฺขมสุขาย เวทนาย สมฺปยุตฺโต ธมฺโม อุปฺปชฺชติ เหตุปจฺจยา.\csromanspacer{} ปญฺจ.}

\prose{น-อพฺยากตํ น-อทุกฺขมสุขาย เวทนาย สมฺปยุตฺตํ ธมฺมํ ปฏิจฺจ น-อพฺยากโต น-อทุกฺขมสุขาย เวทนาย สมฺปยุตฺโต ธมฺโม อุปฺปชฺชติ เหตุปจฺจยา.\csromanspacer{} ฉ.}

\prose{นกุสลํ น-อทุกฺขมสุขาย เวทนาย สมฺปยุตฺตญฺจ น-อพฺยากตํ น-อทุกฺขมสุขาย เวทนาย สมฺปยุตฺตญฺจ ธมฺมํ ปฏิจฺจ นกุสโล น-อทุกฺขมสุขาย เวทนาย สมฺปยุตฺโต ธมฺโม อุปฺปชฺชติ เหตุปจฺจยา.\csromanspacer{} ปญฺจ.}

\prose{น-อกุสลํ น-อทุกฺขมสุขาย เวทนาย สมฺปยุตฺตญฺจ น-อพฺยากตํ น-อทุกฺขมสุขาย เวทนาย สมฺปยุตฺตญฺจ ธมฺมํ ปฏิจฺจ นกุสโล น-อทุกฺขมสุขาย เวทนาย สมฺปยุตฺโต ธมฺโม อุปฺปชฺชติ เหตุปจฺจยา.\csromanspacer{} ปญฺจ.}

\prose{นกุสลํ น-อทุกฺขมสุขาย เวทนาย สมฺปยุตฺตญฺจ น-อกุสลํ น-อทุกฺขมสุขาย เวทนาย สมฺปยุตฺตญฺจ ธมฺมํ ปฏิจฺจ นกุสโล น-อทุกฺขมสุขาย เวทนาย สมฺปยุตฺโต ธมฺโม อุปฺปชฺชติ เหตุปจฺจยา.\csromanspacer{} ตีณิ. (สํขิตฺตํ.)}

\prose{เหตุยา เอกูนติํส ฯเปฯ อวิคเต เอกูนติํส. (สพฺพตฺถ วิตฺถาโร.)}

\csromanheader{2}{1. กุสลตฺติก 2. วิปากตฺติก}
\proseitem{9}{นกุสลํ นวิปากํ ธมฺมํ ปฏิจฺจ นกุสโล นวิปาโก ธมฺโม อุปฺปชฺชติ เหตุปจฺจยา.\csromanspacer{}\csromanspacer{}นกุสลํ นวิปากํ ธมฺมํ ปฏิจฺจ น-อกุสโล นวิปาโก ธมฺโม อุปฺปชฺชติ เหตุปจฺจยา.\csromanspacer{}\csromanspacer{}นกุสลํ นวิปากํ ธมฺมํ ปฏิจฺจ น-อพฺยากโต นวิปาโก ธมฺโม อุปฺปชฺชติ เหตุปจฺจยา.\csromanspacer{}\csromanspacer{}นกุสลํ นวิปากํ ธมฺมํ ปฏิจฺจ นกุสโล นวิปาโก จ น-อพฺยากโต นวิปาโก จ ธมฺมา อุปฺปชฺชนฺติ}

\vspace{\tipitakaparskip}
\csromancenter{ธมฺมปจฺจนีย ติกติกปฏฺฐานปาฬิ เหตุปจฺจยา.\csromanspacer{}\csromanspacer{}นกุสลํ นวิปากํ ธมฺมํ ปฏิจฺจ นกุสโล นวิปาโก จ น-อกุสโล นวิปาโก จ ธมฺมา อุปฺปชฺชนฺติ เหตุปจฺจยา. (5)}
\prose{\csromanfolio{144}น-อกุสลํ นวิปากํ ธมฺมํ ปฏิจฺจ น-อกุสโล นวิปาโก ธมฺโม อุปฺปชฺชติ เหตุปจฺจยา.\csromanspacer{} ปญฺจ.}

\prose{น-อพฺยากตํ นวิปากํ ธมฺมํ ปฏิจฺจ น-อพฺยากโต นวิปาโก ธมฺโม อุปฺปชฺชติ เหตุปจฺจยา.\csromanspacer{} ฉ.}

\prose{นกุสลํ นวิปากญฺจ น-อพฺยากตํ นวิปากญฺจ ธมฺมํ ปฏิจฺจ นกุสโล นวิปาโก ธมฺโม อุปฺปชฺชติ เหตุปจฺจยา.\csromanspacer{} ปญฺจ.}

\prose{น-อกุสลํ นวิปากญฺจ น-อพฺยากตํ นวิปากญฺจ ธมฺมํ ปฏิจฺจ นกุสโล นวิปาโก ธมฺโม อุปฺปชฺชติ เหตุปจฺจยา.\csromanspacer{} ปญฺจ.}

\prose{นกุสลํ นวิปากญฺจ น-อกุสลํ นวิปากญฺจ ธมฺมํ ปฏิจฺจ นกุสโล นวิปาโก ธมฺโม อุปฺปชฺชติ เหตุปจฺจยา.\csromanspacer{} ตีณิ. (สํขิตฺตํ.)}

\prose{เหตุยา เอกูนติํส ฯเปฯ อวิคเต เอกูนติํส. (สพฺพตฺถ วิตฺถาโร.)}

\proseitem{10}{นกุสลํ นวิปากธมฺมธมฺมํ ปฏิจฺจ นกุสโล นวิปากธมฺมธมฺโม อุปฺปชฺชติ เหตุปจฺจยา.\csromanspacer{}\csromanspacer{}นกุสลํ นวิปากธมฺมธมฺมํ ปฏิจฺจ น-อกุสโล นวิปากธมฺมธมฺโม อุปฺปชฺชติ เหตุปจฺจยา.\csromanspacer{}\csromanspacer{}นกุสลํ นวิปากธมฺมธมฺมํ ปฏิจฺจ นกุสโล นวิปากธมฺมธมฺโม จ น-อกุสโล นวิปากธมฺมธมฺโม จ ธมฺมา อุปฺปชฺชนฺติ เหตุปจฺจยา. (3)}

\prose{น-อกุสลํ นวิปากธมฺมธมฺมํ ปฏิจฺจ น-อกุสโล นวิปากธมฺมธมฺโม อุปฺปชฺชติ เหตุปจฺจยา.\csromanspacer{}\csromanspacer{}น-อกุสลํ นวิปากธมฺมธมฺมํ ปฏิจฺจ นกุสโล นวิปากธมฺมธมฺโม อุปฺปชฺชติ เหตุปจฺจยา.\csromanspacer{}\csromanspacer{}น-อกุสลํ นวิปากธมฺมธมฺมํ ปฏิจฺจ นกุสโล นวิปากธมฺมธมฺโม จ น-อกุสโล นวิปากธมฺมธมฺโม จ ธมฺมา อุปฺปชฺชนฺติ เหตุปจฺจยา. (3)}

\prose{นกุสลํ นวิปากธมฺมธมฺมญฺจ น-อกุสลํ นวิปากธมฺมธมฺมญฺจ ธมฺมํ ปฏิจฺจ นกุสโล นวิปากธมฺมธมฺโม อุปฺปชฺชติ เหตุปจฺจยา.\csromanspacer{}\csromanspacer{}นกุสลํ นวิปากธมฺมธมฺมญฺจ น-อกุสลํ นวิปากธมฺมธมฺมญฺจ ธมฺมํ ปฏิจฺจ น-อกุสโล นวิปากธมฺมธมฺโม อุปฺปชฺชติ เหตุปจฺจยา.\csromanspacer{}\csromanspacer{}นกุสลํ นวิปากธมฺมธมฺมญฺจ}

\vspace{\tipitakaparskip}
\csromancenter{กุสลตฺติก 6.\csromanspacer{} ปีติตฺติก น-อกุสลํ นวิปากธมฺมธมฺมญฺจ ธมฺมํ ปฏิจฺจ นกุสโล นวิปากธมฺมธมฺโม จ น-อกุสโล นวิปากธมฺมธมฺโม จ ธมฺมา อุปฺปชฺชนฺติ เหตุปจฺจยา. (3) (สํขิตฺตํ.)}
\csromanheader{2}{เหตุยา นว ฯเปฯ อวิคเต นว. (สพฺพตฺถ นว.)}
\proseitem{11}{\csromanfolio{145}นกุสลํ นเนววิปากนวิปากธมฺมธมฺมํ ปฏิจฺจ นกุสโล นเนววิปากนวิปากธมฺมธมฺโม อุปฺปชฺชติ เหตุปจฺจยา.}

\csromanheader{2}{1. กุสลตฺติก 3. อุปาทินฺนตฺติก}
\proseitem{12}{นกุสลํ น-อุปาทินฺนุปาทานิยํ ธมฺมํ ปฏิจฺจ.}

\proseitem{13}{นกุสลํ น-อนุปาทินฺนุปาทานิยํ ธมฺมํ ปฏิจฺจ.}

\proseitem{14}{นกุสลํ น-อนุปาทินฺน-อนุปาทานิยํ ธมฺมํ ปฏิจฺจ.}

\csromanheader{2}{1. กุสลตฺติก 4. สํกิลิฏฺฐตฺติก}
\proseitem{15}{นกุสลํ นสํกิลิฏฺฐสํกิเลสิกํ ธมฺมํ ปฏิจฺจ.}

\proseitem{16}{นกุสลํ น-อสํกิลิฏฺฐสํกิเลสิกํ ธมฺมํ ปฏิจฺจ.}

\proseitem{17}{นกุสลํ น-อสํกิลิฏฺฐ-อสํกิเลสิกํ ธมฺมํ ปฏิจฺจ.}

\csromanheader{2}{1. กุสลตฺติก 5. วิตกฺกตฺติก}
\proseitem{18}{นกุสลํ นสวิตกฺกสวิจารํ ธมฺมํ ปฏิจฺจ.}

\proseitem{19}{นกุสลํ น-อวิตกฺกวิจารมตฺตํ ธมฺมํ ปฏิจฺจ.}

\proseitem{20}{นกุสลํ น-อวิตกฺก-อวิจารํ ธมฺมํ ปฏิจฺจ.}

\csromanheader{2}{1. กุสลตฺติก 6. ปีติตฺติก}
\proseitem{21}{นกุสลํ นปีติสหคตํ ธมฺมํ ปฏิจฺจ ฯเปฯ.\csromanspacer{} นกุสลํ นสุขสหคตํ ธมฺมํ ปฏิจฺจ ฯเปฯ.\csromanspacer{} นกุสลํ น-อุเปกฺขาสหคตํ ธมฺมํ ปฏิจฺจ.}

\csromanheader{2}{ธมฺมปจฺจนีย ติกติกปฏฺฐานปาฬิ \textbf{1. กุสลตฺติก 7-12. ทสฺสนตฺติกาทิ}}
\proseitem{22}{\csromanfolio{146}นกุสลํ นทสฺสเนน ปหาตพฺพํ ธมฺมํ ปฏิจฺจ ฯเปฯ.\csromanspacer{} นกุสลํ นภาวนาย ปหาตพฺพํ ธมฺมํ ปฏิจฺจ.}

\prose{นกุสลํ นเนวทสฺสเนน นภาวนาย ปหาตพฺพํ ธมฺมํ ปฏิจฺจ.}

\proseitem{23}{นกุสลํ นทสฺสเนน ปหาตพฺพเหตุกํ ธมฺมํ ปฏิจฺจ ฯเปฯ.\csromanspacer{} นกุสลํ นภาวนาย ปหาตพฺพเหตุกํ ธมฺมํ ปฏิจฺจ.}

\prose{นกุสลํ นเนวทสฺสเนน นภาวนาย ปหาตพฺพเหตุกํ ธมฺมํ ปฏิจฺจ.}

\proseitem{24}{นกุสลํ น-อาจยคามิํ ธมฺมํ ปฏิจฺจ.}

\prose{นกุสลํ น-อปจยคามิํ ธมฺมํ ปฏิจฺจ.}

\prose{นกุสลํ นเนวาจยคามินาปจยคามิํ ธมฺมํ ปฏิจฺจ.}

\proseitem{25}{นกุสลํ นเสกฺขํ ธมฺมํ ปฏิจฺจ ฯเปฯ.\csromanspacer{} นกุสลํ น-อเสกฺขํ ธมฺมํ ปฏิจฺจ.\csromanspacer{} นกุสลํ นเนวเสกฺขนาเสกฺขํ ธมฺมํ ปฏิจฺจ.}

\proseitem{26}{นกุสลํ นปริตฺตํ ธมฺมํ ปฏิจฺจ.}

\prose{นกุสลํ นมหคฺคตํ ธมฺมํ ปฏิจฺจ ฯเปฯ.\csromanspacer{} นกุสลํ น-อปฺปมาณํ ธมฺมํ ปฏิจฺจ.}

\proseitem{27}{นกุสลํ นปริตฺตารมฺมณํ ธมฺมํ ปฏิจฺจ ฯเปฯ.\csromanspacer{} นกุสลํ นมหคฺคตารมฺมณํ ธมฺมํ ปฏิจฺจ ฯเปฯ.\csromanspacer{} นกุสลํ น-อปฺปมาณารมฺมณํ ธมฺมํ ปฏิจฺจ.}

\csromanheader{2}{1. กุสลตฺติก 13-20. หีนตฺติกาทิ}
\proseitem{28}{นกุสลํ นหีนํ ธมฺมํ ปฏิจฺจ.}

\prose{นกุสลํ นมชฺฌิมํ ธมฺมํ ปฏิจฺจ.}

\prose{นกุสลํ นปณีตํ ธมฺมํ ปฏิจฺจ.}

\proseitem{29}{นกุสลํ นมิจฺฉตฺตนิยตํ ธมฺมํ ปฏิจฺจ ฯเปฯ.\csromanspacer{} นกุสลํ นสมฺมตฺตนิยตํ ธมฺมํ ปฏิจฺจ.}

\prose{นกุสลํ น-อนิยตํ ธมฺมํ ปฏิจฺจ.}

\csromanheader{2}{1. กุสลตฺติก 21. สนิทสฺสนตฺติก}
\proseitem{30}{\csromanfolio{147}นกุสลํ นมคฺคารมฺมณํ ธมฺมํ ปฏิจฺจ ฯเปฯ.\csromanspacer{} นกุสลํ นมคฺคเหตุกํ ธมฺมํ ปฏิจฺจ ฯเปฯ.\csromanspacer{} นกุสลํ นมคฺคาธิปติํ ธมฺมํ ปฏิจฺจ.}

\proseitem{31}{นกุสลํ น-อนุปฺปนฺนํ ธมฺมํ ปฏิจฺจ ฯเปฯ.\csromanspacer{} นกุสลํ น-อุปฺปาทิํ ธมฺมํ ปฏิจฺจ.}

\proseitem{32}{นกุสลํ น-อตีตํ ธมฺมํ ปฏิจฺจ ฯเปฯ.\csromanspacer{} นกุสลํ น-อนาคตํ ธมฺมํ ปฏิจฺจ.}

\proseitem{33}{นกุสลํ น-อตีตารมฺมณํ ธมฺมํ ปฏิจฺจ ฯเปฯ.\csromanspacer{} นกุสลํ น-อนาคตารมฺมณํ ธมฺมํ ปฏิจฺจ ฯเปฯ.\csromanspacer{} นกุสลํ นปจฺจุนฺนารมฺมณํ ธมฺมํ ปฏิจฺจ.}

\proseitem{34}{นกุสลํ น-อชฺฌตฺตํ ธมฺมํ ปฏิจฺจ ฯเปฯ.\csromanspacer{} นกุสลํ นพหิทฺธา ธมฺมํ ปฏิจฺจ.}

\proseitem{35}{นกุสลํ น-อชฺฌตฺตารมฺมณํ ธมฺมํ ปฏิจฺจ ฯเปฯ.\csromanspacer{} นกุสลํ นพหิทฺธารมฺมณํ ธมฺมํ ปฏิจฺจ.}

\csromanheader{2}{1. กุสลตฺติก 21. สนิทสฺสนตฺติก}
\proseitem{36}{นกุสลํ นสนิทสฺสนสปฺปฏิฆํ ธมฺมํ ปฏิจฺจ นกุสโล นสนิทสฺสนสปฺปฏิโฆ ธมฺโม อุปฺปชฺชติ เหตุปจฺจยา.\csromanspacer{}\csromanspacer{}นกุสลํ นสนิทสฺสนสปฺปฏิฆํ ธมฺมํ ปฏิจฺจ น-อกุสโล นสนิทสฺสนสปฺปฏิโฆ ธมฺโม อุปฺปชฺชติ เหตุปจฺจยา.\csromanspacer{}\csromanspacer{}นกุสลํ นสนิทสฺสนสปฺปฏิฆํ ธมฺมํ ปฏิจฺจ น-อพฺยากโต นสนิทสฺสนสปฺปฏิโฆ ธมฺโม อุปฺปชฺชติ เหตุปจฺจยา.\csromanspacer{}\csromanspacer{}นกุสลํ นสนิทสฺสนสปฺปฏิฆํ ธมฺมํ ปฏิจฺจ นกุสโล นสนิทสฺสนสปฺปฏิโฆ จ น-อพฺยากโต นสนิทสฺสนสปฺปฏิโฆ จ ธมฺมา อุปฺปชฺชนฺติ เหตุปจฺจยา.\csromanspacer{}\csromanspacer{}นกุสลํ นสนิทสฺสนสปฺปฏิฆํ ธมฺมํ ปฏิจฺจ นกุสโล นสนิทสฺสนสปฺปฏิโฆ จ น-อกุสโล นสนิทสฺสนสปฺปฏิโฆ จ ธมฺมา อุปฺปชฺชนฺติ เหตุปจฺจยา. (5)}

\prose{น-อกุสลํ นสนิทสฺสนสปฺปฏิฆํ ธมฺมํ ปฏิจฺจ น-อกุสโล นสนิทสฺสนสปฺปฏิโฆ ธมฺโม อุปฺปชฺชติ เหตุปจฺจยา.\csromanspacer{}\csromanspacer{}น-อกุสลํ นสนิทสฺสนสปฺปฏิฆํ ธมฺมํ ปฏิจฺจ นกุสโล นสนิทสฺสนสปฺปฏิโฆ ธมฺโม อุปฺปชฺชติ เหตุปจฺจยา.\csromanspacer{}\csromanspacer{}น-อกุสลํ นสนิทสฺสนสปฺปฏิฆํ ธมฺมํ ปฏิจฺจ น-อพฺยากโต นสนิทสฺสนสปฺปฏิโฆ ธมฺโม อุปฺปชฺชติ เหตุปจฺจยา.\csromanspacer{}\csromanspacer{}น-อกุสลํ นสนิทสฺสนสปฺปฏิฆํ ธมฺมํ ปฏิจฺจ น-อกุสโล นสนิทสฺสนสปฺปฏิโฆ จ น-อพฺยากโต นสนิทสฺสนสปฺปฏิโฆ จ ธมฺมา อุปฺปชฺชนฺติ เหตุปจฺจยา.\csromanspacer{}\csromanspacer{}}

\vspace{\tipitakaparskip}
\csromancenter{ธมฺมปจฺจนีย ติกติกปฏฺฐานปาฬิ น-อกุสลํ นสนิทสฺสนสปฺปฏิฆํ ธมฺมํ ปฏิจฺจ นกุสโล นสนิทสฺสนสปฺปฏิโฆ จ น-อกุสโล นสนิทสฺสนสปฺปฏิโฆ จ ธมฺมา อุปฺปชฺชนฺติ เหตุปจฺจยา. (5)}
\prose{\csromanfolio{148}น-อพฺยากตํ นสนิทสฺสนสปฺปฏิฆํ ธมฺมํ ปฏิจฺจ น-อพฺยากโต นสนิทสฺสนสปฺปฏิโฆ ธมฺโม อุปฺปชฺชติ เหตุปจฺจยา.\csromanspacer{}\csromanspacer{}น-อพฺยากตํ นสนิทสฺสนสปฺปฏิฆํ ธมฺมํ ปฏิจฺจ นกุสโล นสนิทสฺสนสปฺปฏิโฆ ธมฺโม อุปฺปชฺชติ เหตุปจฺจยา.\csromanspacer{}\csromanspacer{}น-อพฺยากตํ นสนิทสฺสนสปฺปฏิฆํ ธมฺมํ ปฏิจฺจ น-อกุสโล นสนิทสฺสนสปฺปฏิโฆ ธมฺโม อุปฺปชฺชติ เหตุปจฺจยา.\csromanspacer{}\csromanspacer{}น-อพฺยากตํ นสนิทสฺสนสปฺปฏิฆํ ธมฺมํ ปฏิจฺจ นกุสโล นสนิทสฺสนสปฺปฏิโฆ จ น-อพฺยากโต นสนิทสฺสนสปฺปฏิโฆ จ ธมฺมา อุปฺปชฺชนฺติ เหตุปจฺจยา.\csromanspacer{}\csromanspacer{}น-อพฺยากตํ นสนิทสฺสนสปฺปฏิฆํ ธมฺมํ ปฏิจฺจ น-อกุสโล นสนิทสฺสนสปฺปฏิโฆ จ น-อพฺยากโต นสนิทสฺสนสปฺปฏิโฆ จ ธมฺมา อุปฺปชฺชนฺติ เหตุปจฺจยา.\csromanspacer{}\csromanspacer{}น-อพฺยากตํ นสนิทสฺสนสปฺปฏิฆํ ธมฺมํ ปฏิจฺจ นกุสโล นสนิทสฺสนสปฺปฏิโฆ จ น-อกุสโล นสนิทสฺสนสปฺปฏิโฆ จ ธมฺมา อุปฺปชฺชนฺติ เหตุปจฺจยา. (6)}

\prose{นกุสลํ นสนิทสฺสนสปฺปฏิฆญฺจ น-อพฺยากตํ นสนิทสฺสนสปฺปฏิฆญฺจ ธมฺมํ ปฏิจฺจ นกุสโล นสนิทสฺสนสปฺปฏิโฆ ธมฺโม อุปฺปชฺชติ เหตุปจฺจยา.\csromanspacer{} ปญฺจ. น-อกุสลํ นสนิทสฺสนสปฺปฏิฆญฺจ น-อพฺยากตํ นสนิทสฺสนสปฺปฏิฆญฺจ ธมฺมํ ปฏิจฺจ นกุสโล นสนิทสฺสนสปฺปฏิโฆ ธมฺโม อุปฺปชฺชติ เหตุปจฺจยา.\csromanspacer{} ปญฺจ.}

\prose{นกุสลํ นสนิทสฺสนสปฺปฏิฆญฺจ น-อกุสลํ นสนิทสฺสนสปฺปฏิฆญฺจ ธมฺมํ ปฏิจฺจ นกุสโล นสนิทสฺสนสปฺปฏิโฆ ธมฺโม อุปฺปชฺชติ เหตุปจฺจยา.\csromanspacer{}\csromanspacer{}นกุสลํ นสนิทสฺสนสปฺปฏิฆญฺจ น-อกุสลํ นสนิทสฺสนสปฺปฏิฆญฺจ ธมฺมํ ปฏิจฺจ น-อกุสโล นสนิทสฺสนสปฺปฏิโฆ ธมฺโม อุปฺปชฺชติ เหตุปจฺจยา.\csromanspacer{}\csromanspacer{}นกุสลํ นสนิทสฺสนสปฺปฏิฆญฺจ น-อกุสลํ นสนิทสฺสนสปฺปฏิฆญฺจ ธมฺมํ ปฏิจฺจ นกุสโล นสนิทสฺสนสปฺปฏิโฆ จ น-อกุสโล นสนิทสฺสนสปฺปฏิโฆ จ ธมฺมา อุปฺปชฺชนฺติ เหตุปจฺจยา. (3) (สํขิตฺตํ.)}

\prose{เหตุยา เอกูนติํส ฯเปฯ อวิคเต เอกูนติํส. (สพฺพตฺถ วิตฺถาโร.)}

\csromanheader{2}{2. เวทนาตฺติก 1. กุสลตฺติก}
\proseitem{37}{\csromanfolio{149}นกุสลํ น-อนิทสฺสนสปฺปฏิฆํ ธมฺมํ ปฏิจฺจ นกุสโล น-อนิทสฺสนสปฺปฏิโฆ ธมฺโม อุปฺปชฺชติ เหตุปจฺจยา. (สํขิตฺตํ.)}

\prose{เหตุยา เอกูนติํส ฯเปฯ อวิคเต เอกูนติํส. (สพฺพตฺถ วิตฺถาโร.)}

\proseitem{38}{นกุสลํ น-อนิทสฺสน-อปฺปฏิฆํ ธมฺมํ ปฏิจฺจ นกุสโล น-อนิทสฺสน-อปฺปฏิโฆ ธมฺโม อุปฺปชฺชติ เหตุปจฺจยา. (สํขิตฺตํ.)}

\csromanheader{2}{เหตุยา นว. (สพฺพตฺถ วิตฺถาโร.)}
\csromanheader{2}{2. เวทนาตฺติก 1. กุสลตฺติก}
\proseitem{39}{นสุขาย เวทนาย สมฺปยุตฺตํ นกุสลํ ธมฺมํ ปฏิจฺจ นสุขาย เวทนาย สมฺปยุตฺโต นกุสโล ธมฺโม อุปฺปชฺชติ เหตุปจฺจยา.\csromanspacer{}\csromanspacer{}นสุขาย เวทนาย สมฺปยุตฺตํ นกุสลํ ธมฺมํ ปฏิจฺจ นทุกฺขาย เวทนาย สมฺปยุตฺโต นกุสโล ธมฺโม อุปฺปชฺชติ เหตุปจฺจยา.\csromanspacer{}\csromanspacer{}นสุขาย เวทนาย สมฺปยุตฺตํ นกุสลํ ธมฺมํ ปฏิจฺจ น-อทุกฺขมสุขาย เวทนาย สมฺปยุตฺโต นกุสโล ธมฺโม อุปฺปชฺชติ เหตุปจฺจยา.\csromanspacer{}\csromanspacer{}นสุขาย เวทนาย สมฺปยุตฺตํ นกุสลํ ธมฺมํ ปฏิจฺจ นสุขาย เวทนาย สมฺปยุตฺโต นกุสโล จ น-อทุกฺขมสุขาย เวทนาย สมฺปยุตฺโต นกุสโล จ ธมฺมา อุปฺปชฺชนฺติ เหตุปจฺจยา ฯเปฯ.\csromanspacer{} สตฺต.}

\prose{นทุกฺขาย เวทนาย สมฺปยุตฺตํ นกุสลํ ธมฺมํ ปฏิจฺจ นทุกฺขาย เวทนาย สมฺปยุตฺโต นกุสโล ธมฺโม อุปฺปชฺชติ เหตุปจฺจยา.\csromanspacer{}\csromanspacer{}นทุกฺขาย เวทนาย สมฺปยุตฺตํ นกุสลํ ธมฺมํ ปฏิจฺจ นสุขาย เวทนาย สมฺปยุตฺโต นกุสโล ธมฺโม อุปฺปชฺชติ เหตุปจฺจยา.\csromanspacer{}\csromanspacer{}นทุกฺขาย เวทนาย สมฺปยุตฺตํ นกุสลํ ธมฺมํ ปฏิจฺจ น-อทุกฺขมสุขาย เวทนาย สมฺปยุตฺโต นกุสโล ธมฺโม อุปฺปชฺชติ เหตุปจฺจยา.\csromanspacer{}\csromanspacer{}นทุกฺขาย เวทนาย สมฺปยุตฺตํ นกุสลํ ธมฺมํ ปฏิจฺจ นสุขาย เวทนาย สมฺปยุตฺโต นกุสโล จ น-อทุกฺขมสุขาย เวทนาย สมฺปยุตฺโต นกุสโล จ ธมฺมา อุปฺปชฺชนฺติ เหตุปจฺจยา ฯเปฯ.\csromanspacer{} สตฺต.}

\prose{น-อทุกฺขมสุขาย เวทนาย สมฺปยุตฺตํ นกุสลํ ธมฺมํ ปฏิจฺจ น-อทุกฺขมสุขาย เวทนาย สมฺปยุตฺโต นกุสโล ธมฺโม อุปฺปชฺชติ เหตุปจฺจยา.\csromanspacer{}\csromanspacer{}น-อทุกฺขมสุขาย เวทนาย สมฺปยุตฺตํ นกุสลํ ธมฺมํ ปฏิจฺจ}

\vspace{\tipitakaparskip}
\csromancenter{ธมฺมปจฺจนีย ติกติกปฏฺฐานปาฬิ นสุขาย เวทนาย สมฺปยุตฺโต นกุสโล ธมฺโม อุปฺปชฺชติ เหตุปจฺจยา.\csromanspacer{}\csromanspacer{}น-อทุกฺขมสุขาย เวทนาย สมฺปยุตฺตํ นกุสลํ ธมฺมํ ปฏิจฺจ นทุกฺขาย เวทนาย สมฺปยุตฺโต นกุสโล ธมฺโม อุปฺปชฺชติ เหตุปจฺจยา.\csromanspacer{}\csromanspacer{}น-อทุกฺขมสุขาย เวทนาย สมฺปยุตฺตํ นกุสลํ ธมฺมํ ปฏิจฺจ นสุขาย เวทนาย สมฺปยุตฺโต นกุสโล จ น-อทุกฺขมสุขาย เวทนาย สมฺปยุตฺโต นกุสโล จ ธมฺมา อุปฺปชฺชนฺติ เหตุปจฺจยา ฯเปฯ.\csromanspacer{} สตฺต.}
\proseitem{40}{\csromanfolio{150}นสุขาย เวทนาย สมฺปยุตฺตํ นกุสลญฺจ น-อทุกฺขมสุขาย เวทนาย สมฺปยุตฺตํ นกุสลญฺจ ธมฺมํ ปฏิจฺจ นสุขาย เวทนาย สมฺปยุตฺโต นกุสโล ธมฺโม อุปฺปชฺชติ เหตุปจฺจยา.\csromanspacer{}\csromanspacer{}นสุขาย เวทนาย สมฺปยุตฺตํ นกุสลญฺจ น- อทุกฺขมสุขาย เวทนาย สมฺปยุตฺตํ นกุสลญฺจ ธมฺมํ ปฏิจฺจ นทุกฺขาย เวทนาย สมฺปยุตฺโต นกุสโล ธมฺโม อุปฺปชฺชติ เหตุปจฺจยา.\csromanspacer{}\csromanspacer{}นสุขาย เวทนาย สมฺปยุตฺตํ นกุสลญฺจ น-อทุกฺขมสุขาย เวทนาย สมฺปยุตฺตํ นกุสลญฺจ ธมฺมํ ปฏิจฺจ น-อทุกฺขมสุขาย เวทนาย สมฺปยุตฺโต นกุสโล ธมฺโม อุปฺปชฺชติ เหตุปจฺจยา.\csromanspacer{}\csromanspacer{}นสุขาย เวทนาย สมฺปยุตฺตํ นกุสลญฺจ น-อทุกฺขมสุขาย เวทนาย สมฺปยุตฺตํ นกุสลญฺจ ธมฺมํ ปฏิจฺจ นสุขาย เวทนาย สมฺปยุตฺโต นกุสโล จ น-อทุกฺขมสุขาย เวทนาย สมฺปยุตฺโต นกุสโล จ ธมฺมา อุปฺปชฺชนฺติ เหตุปจฺจยา ฯเปฯ.\csromanspacer{} สตฺต.}

\prose{นทุกฺขาย เวทนาย สมฺปยุตฺตํ นกุสลญฺจ น-อทุกฺขมสุขาย เวทนาย สมฺปยุตฺตํ นกุสลญฺจ ธมฺมํ ปฏิจฺจ นสุขาย เวทนาย สมฺปยุตฺโต นกุสโล ธมฺโม อุปฺปชฺชติ เหตุปจฺจยา.\csromanspacer{}\csromanspacer{}นทุกฺขาย เวทนาย สมฺปยุตฺตํ นกุสลญฺจ น- อทุกฺขมสุขาย เวทนาย สมฺปยุตฺตํ นกุสลญฺจ ธมฺมํ ปฏิจฺจ นทุกฺขาย เวทนาย สมฺปยุตฺโต นกุสโล ธมฺโม อุปฺปชฺชติ เหตุปจฺจยา.\csromanspacer{}\csromanspacer{}นทุกฺขาย เวทนาย สมฺปยุตฺตํ นกุสลญฺจ น-อทุกฺขมสุขาย เวทนาย สมฺปยุตฺตํ นกุสลญฺจ ธมฺมํ ปฏิจฺจ น-อทุกฺขมสุขาย เวทนาย สมฺปยุตฺโต นกุสโล ธมฺโม อุปฺปชฺชติ เหตุปจฺจยา ฯเปฯ.\csromanspacer{} สตฺต. (สํขิตฺตํ.)}

\prose{เหตุยา เอกูนปญฺญาส, อารมฺมเณ เอกูนปญฺญาส ฯเปฯ อวิคเต เอกูนปญฺญาส.}

\csromanheader{2}{6-21. วิตกฺกตฺติกาทิ 1. กุสลตฺติก}
\proseitem{41}{\csromanfolio{151}นสุขาย เวทนาย สมฺปยุตฺตํ น-อกุสลํ ธมฺมํ ปฏิจฺจ นสุขาย เวทนาย สมฺปยุตฺโต น-อกุสโล ธมฺโม อุปฺปชฺชติ เหตุปจฺจยา. (สํขิตฺตํ.) .\csromanspacer{} เหตุยา เอกูนปญฺญาส.}

\proseitem{42}{นสุขาย เวทนาย สมฺปยุตฺตํ น-อพฺยากตํ ธมฺมํ ปฏิจฺจ นสุขาย เวทนาย สมฺปยุตฺโต น-อพฺยากโต ธมฺโม อุปฺปชฺชติ เหตุปจฺจยา. (สํขิตฺตํ.) .\csromanspacer{} เหตุยา อฏฺฐจตฺตาลีส. (สพฺพตฺถ วิตฺถาโร.)}

\csromanheader{2}{2. เวทนาตฺติก 2. วิปากตฺติก}
\proseitem{43}{นสุขาย เวทนาย สมฺปยุตฺตํ นวิปากํ ธมฺมํ ปฏิจฺจ นสุขาย เวทนาย สมฺปยุตฺโต นวิปาโก ธมฺโม อุปฺปชฺชติ เหตุปจฺจยา.}

\csromanheader{2}{เหตุยา เอกูนปญฺญาส ฯเปฯ วิคเต เอกูนปญฺญาส. (สํขิตฺตํ.)}
\csromanheader{2}{3. วิปากตฺติก 1. กุสลตฺติก}
\proseitem{44}{นวิปากํ นกุสลํ ธมฺมํ ปฏิจฺจ ฯเปฯ.\csromanspacer{} นวิปากธมฺมธมฺมํ น-อกุสลํ ธมฺมํ ปฏิจฺจ.}

\prose{นเนววิปากนวิปากธมฺมธมฺมํ น-อพฺยากตํ ธมฺมํ ปฏิจฺจ.}

\csromanheader{2}{4. อุปาทินฺนตฺติก 1. กุสลตฺติก}
\proseitem{45}{น-อุปาทินฺนุปาทานิยํ นกุสลํ\footnote{อิโต ปฏฺฐาย ยาว อชฺฌตฺตารมฺมณตฺติกา สพฺพโปตฺถเกสุ อกุสลาพฺยากตปทานิ น ทิสฺสนฺติ, เวทนาตฺติก สนิทสฺสนตฺติเกสุ วิย นนุ เตหิปิ ภวิตพฺพํ.} ธมฺมํ ปฏิจฺจ ฯเปฯ.\csromanspacer{} น-อนุปาทินฺนุปาทานิยํ นกุสลํ ธมฺมํ ปฏิจฺจ ฯเปฯ.\csromanspacer{} น-อนุปาทินฺน-อนุปาทานิยํ นกุสลํ ธมฺมํ ปฏิจฺจ.}

\csromanheader{2}{5. สํกิลิฏฺฐตฺติก 1. กุสลตฺติก}
\proseitem{46}{นสํกิลิฏฺฐสํกิเลสิกํ นกุสลํ ธมฺมํ ปฏิจฺจ ฯเปฯ.\csromanspacer{} น-อสํกิลิฏฺฐสํกิเลสิกํ นกุสลํ ธมฺมํ ปฏิจฺจ ฯเปฯ.\csromanspacer{} น-อสํกิลิฏฺฐ-อสํกิเลสิกํ นกุสลํ ธมฺมํ ปฏิจฺจ.}

\csromanheader{2}{6-21. วิตกฺกตฺติกาทิ 1. กุสลตฺติก}
\proseitem{47}{นสวิตกฺกสวิจารํ นกุสลํ ธมฺมํ ปฏิจฺจ ฯเปฯ.\csromanspacer{} น-อวิตกฺกวิจารมตฺตํ นกุสลํ ธมฺมํ ปฏิจฺจ ฯเปฯ.\csromanspacer{} น-อวิตกฺก-อวิจารํ นกุสลํ ธมฺมํ ปฏิจฺจ.}

\gambhira{ธมฺมปจฺจนีย ติกติกปฏฺฐานปาฬิ}
\proseitem{48}{\csromanfolio{152}นปีติสหคตํ นกุสลํ ธมฺมํ ปฏิจฺจ ฯเปฯ.\csromanspacer{} นสุขสหคตํ นกุสลํ ธมฺมํ ปฏิจฺจ ฯเปฯ.\csromanspacer{} น-อุเปกฺขาสหคตํ นกุสลํ ธมฺมํ ปฏิจฺจ.}

\proseitem{49}{นทสฺสเนน ปหาตพฺพํ นกุสลํ ธมฺมํ ปฏิจฺจ ฯเปฯ.\csromanspacer{} นภาวนาย ปหาตพฺพํ นกุสลํ ธมฺมํ ปฏิจฺจ ฯเปฯ.\csromanspacer{} นเนวทสฺสเนน นภาวนาย ปหาตพฺพํ นกุสลํ ธมฺมํ ปฏิจฺจ.}

\proseitem{50}{นทสฺสเนน ปหาตพฺพเหตุกํ นกุสลํ ธมฺมํ ปฏิจฺจ ฯเปฯ.\csromanspacer{} นภาวนาย ปหาตพฺพเหตุกํ นกุสลํ ธมฺมํ ปฏิจฺจ ฯเปฯ.\csromanspacer{} นเนวทสฺสเนน นภาวนาย ปหาตพฺพเหตุกํ นกุสลํ ธมฺมํ ปฏิจฺจ.}

\proseitem{51}{น-อาจยคามิํ นกุสลํ ธมฺมํ ปฏิจฺจ ฯเปฯ.\csromanspacer{} น-อปจยคามิํ นกุสลํ ธมฺมํ ปฏิจฺจ ฯเปฯ.\csromanspacer{} นเนวาจยคามินาปจยคามิํ นกุสลํ ธมฺมํ ปฏิจฺจ.}

\proseitem{52}{นเสกฺขํ นกุสลํ ธมฺมํ ปฏิจฺจ ฯเปฯ.\csromanspacer{} น-อเสกฺขํ นกุสลํ ธมฺมํ ปฏิจฺจ ฯเปฯ.\csromanspacer{} นเนวาเสกฺขนาเสกฺขํ นกุสลํ ธมฺมํ ปฏิจฺจ.}

\proseitem{53}{นปริตฺตํ นกุสลํ ธมฺมํ ปฏิจฺจ ฯเปฯ.\csromanspacer{} นมหคฺคตํ นกุสลํ ธมฺมํ ปฏิจฺจ ฯเปฯ.\csromanspacer{} น-อปฺปมาณํ นกุสลํ ธมฺมํ ปฏิจฺจ.}

\proseitem{54}{นปริตฺตารมฺมณํ นกุสลํ ธมฺมํ ปฏิจฺจ ฯเปฯ.\csromanspacer{} นมหคฺคตารมฺมณํ นกุสลํ ธมฺมํ ปฏิจฺจ ฯเปฯ.\csromanspacer{} น-อปฺปมาณารมฺมณํ นกุสลํ ธมฺมํ ปฏิจฺจ.}

\proseitem{55}{นหีนํ นกุสลํ ธมฺมํ ปฏิจฺจ ฯเปฯ.\csromanspacer{} นมชฺฌิมํ นกุสลํ ธมฺมํ ปฏิจฺจ ฯเปฯ.\csromanspacer{} นปณีตํ นกุสลํ ธมฺมํ ปฏิจฺจ.}

\proseitem{56}{นมิจฺฉตฺตนิยตํ นกุสลํ ธมฺมํ ปฏิจฺจ ฯเปฯ.\csromanspacer{} นสมฺมตฺตนิยตํ นกุสลํ ธมฺมํ ปฏิจฺจ ฯเปฯ.\csromanspacer{} น-อนิยตํ นกุสลํ ธมฺมํ ปฏิจฺจ.}

\proseitem{57}{นมคฺคารมฺมณํ นกุสลํ ธมฺมํ ปฏิจฺจ ฯเปฯ.\csromanspacer{} นมคฺคเหตุกํ นกุสลํ ธมฺมํ ปฏิจฺจ ฯเปฯ.\csromanspacer{} นมคฺคาธิปติํ นกุสลํ ธมฺมํ ปฏิจฺจ.}

\proseitem{58}{น-อนุปฺปนฺนํ นกุสลํ ธมฺมํ ปฏิจฺจ ฯเปฯ.\csromanspacer{} น-อุปฺปาทิํ นกุสลํ ธมฺมํ ปฏิจฺจ.}

\proseitem{59}{น-อตีตํ นกุสลํ ธมฺมํ ปฏิจฺจ ฯเปฯ.\csromanspacer{} น-อนาคตํ นกุสลํ ธมฺมํ ปฏิจฺจ.}

\proseitem{60}{น-อตีตารมฺมณํ นกุสลํ ธมฺมํ ปฏิจฺจ ฯเปฯ.\csromanspacer{} น-อนาคตารมฺมณํ นกุสลํ ธมฺมํ ปฏิจฺจ ฯเปฯ.\csromanspacer{} นปจฺจุปฺปนฺนารมฺมณํ นกุสลํ ธมฺมํ ปฏิจฺจ.}

\csromanheader{2}{22. สนิทสฺสนตฺติก 1. กุสลตฺติก}
\proseitem{61}{\csromanfolio{153}น-อชฺฌตฺตํ น-อกุสลํ ธมฺมํ ปฏิจฺจ ฯเปฯ.\csromanspacer{} นพหิทฺธา นกุสลํ ธมฺมํ ปฏิจฺจ.}

\proseitem{62}{น-อชฺฌตฺตารมฺมณํ นกุสลํ ธมฺมํ ปฏิจฺจ ฯเปฯ.\csromanspacer{} นพหิทฺธารมฺมณํ นกุสลํ ธมฺมํ ปฏิจฺจ.}

\csromanheader{2}{22. สนิทสฺสนตฺติก 1. กุสลตฺติก}
\proseitem{63}{นสนิทสฺสนสปฺปฏิฆํ นกุสลํ ธมฺมํ ปฏิจฺจ นสนิทสฺสนสปฺปฏิโฆ นกุสโล ธมฺโม อุปฺปชฺชติ เหตุปจฺจยา.\csromanspacer{}\csromanspacer{}นสนิทสฺสนสปฺปฏิฆํ นกุสลํ ธมฺมํ ปฏิจฺจ น-อนิทสฺสนสปฺปฏิโฆ นกุสโล ธมฺโม อุปฺปชฺชติ เหตุปจฺจยา.\csromanspacer{}\csromanspacer{}นสนิทสฺสนสปฺปฏิฆํ นกุสลํ ธมฺมํ ปฏิจฺจ น-อนิทสฺสน-อปฺปฏิโฆ นกุสโล ธมฺโม อุปฺปชฺชติ เหตุปจฺจยา ฯเปฯ.\csromanspacer{} ฉ.}

\prose{น-อนิทสฺสนสปฺปฏิฆํ นกุสลํ ธมฺมํ ปฏิจฺจ น-อนิทสฺสนสปฺปฏิโฆ นกุสโล ธมฺโม อุปฺปชฺชติ เหตุปจฺจยา.\csromanspacer{}\csromanspacer{}น-อนิทสฺสนสปฺปฏิฆํ นกุสลํ ธมฺมํ ปฏิจฺจ นสนิทสฺสนสปฺปฏิโฆ นกุสโล ธมฺโม อุปฺปชฺชติ เหตุปจฺจยา.\csromanspacer{}\csromanspacer{}น- อนิทสฺสนสปฺปฏิฆํ นกุสลํ ธมฺมํ ปฏิจฺจ น-อนิทสฺสน-อปฺปฏิโฆ นกุสโล ธมฺโม อุปฺปชฺชติ เหตุปจฺจยา ฯเปฯ.\csromanspacer{} ฉ.}

\prose{น-อนิทสฺสน-อปฺปฏิฆํ นกุสลํ ธมฺมํ ปฏิจฺจ น-อนิทสฺสน-อปฺปฏิโฆ นกุสโล ธมฺโม อุปฺปชฺชติ เหตุปจฺจยา.\csromanspacer{} ฉ.}

\prose{นสนิทสฺสนสปฺปฏิฆํ นกุสลญฺจ น-อนิทสฺสน-อปฺปฏิฆํ นกุสลญฺจ ธมฺมํ ปฏิจฺจ นสนิทสฺสนสปฺปฏิโฆ นกุสโล ธมฺโม อุปฺปชฺชติ เหตุปจฺจยา.\csromanspacer{} ฉ.}

\prose{น-อนิทสฺสนสปฺปฏิฆํ นกุสลญฺจ น-อนิทสฺสน-อปฺปฏิฆํ นกุสลญฺจ ธมฺมํ ปฏิจฺจ นสนิทสฺสนสปฺปฏิโฆ นกุสโล ธมฺโม อุปฺปชฺชติ เหตุปจฺจยา.\csromanspacer{} ฉ. (สํขิตฺตํ.) .\csromanspacer{} เหตุยา ติํส, อารมฺมเณ นว ฯเปฯ อวิคเต ติํส. (สพฺพตฺถ วิตฺถาโร.)}

\proseitem{64}{นสนิทสฺสนสปฺปฏิฆํ น-อกุสลํ ธมฺมํ ปฏิจฺจ นสนิทสฺสนสปฺปฏิโฆ น-อกุสโล ธมฺโม อุปฺปชฺชติ เหตุปจฺจยา.\csromanspacer{}\csromanspacer{}เหตุยา ติํส, อารมฺมเณ นว ฯเปฯ อวิคเต ติํส. (สพฺพตฺถ วิตฺถาโร.)}

\proseitem{65}{นสนิทสฺสนสปฺปฏิฆํ น-อพฺยากตํ ธมฺมํ ปฏิจฺจ นสนิทสฺสนสปฺปฏิโฆ น-อพฺยากโต ธมฺโม อุปฺปชฺชติ เหตุปจฺจยา.}

\csromanheader{2}{เหตุยา นว. (สพฺพตฺถ นว.)}
\csromanheader{2}{ธมฺมปจฺจนีย ติกติกปฏฺฐานปาฬิ \textbf{22. สนิทสฺสนตฺติก 2-21. เวทนาตฺติกาทิ}}
\proseitem{66}{\csromanfolio{154}นสนิทสฺสนสปฺปฏิฆํ นสุขาย เวทนาย สมฺปยุตฺตํ ธมฺมํ ปฏิจฺจ.\csromanspacer{} เหตุยา ติํส.}

\prose{นสนิทสฺสนสปฺปฏิฆํ นทุกฺขาย เวทนาย สมฺปยุตฺตํ ธมฺมํ ปฏิจฺจ.\csromanspacer{} เหตุยา ติํส.}

\prose{นสนิทสฺสนสปฺปฏิฆํ น-อทุกฺขมสุขาย เวทนาย สมฺปยุตฺตํ ธมฺมํ ปฏิจฺจ.\csromanspacer{} เหตุยา ติํส.}

\proseitem{67}{นสนิทสฺสนสปฺปฏิฆํ นวิปากํ ธมฺมํ ปฏิจฺจ ฯเปฯ.}

\prose{นสนิทสฺสนสปฺปฏิฆํ นวิปากธมฺมธมฺมํ ปฏิจฺจ.\csromanspacer{} เหตุยา ติํส.}

\prose{นสนิทสฺสนสปฺปฏิฆํ นเนววิปากนวิปากธมฺมธมฺมํ ปฏิจฺจ.\csromanspacer{} เหตุยา นว.}

\proseitem{68}{นสนิทสฺสนสปฺปฏิฆํ น-อุปาทินฺนุปาทานิยํ ธมฺมํ ปฏิจฺจ ฯเปฯ.}

\prose{นสนิทสฺสนสปฺปฏิฆํ น-อนุปาทินฺนุปาทานิยํ ธมฺมํ ปฏิจฺจ ฯเปฯ.}

\prose{นสนิทสฺสนสปฺปฏิฆํ น-อนุปาทินฺน-อนุปาทานิยํ ธมฺมํ ปฏิจฺจ. เหตุยา ติํส.}

\proseitem{69}{นสนิทสฺสนสปฺปฏิฆํ นสํกิลิฏฺฐสํกิเลสิกํ ธมฺมํ ปฏิจฺจ. เหตุยา ติํส.}

\prose{นสนิทสฺสนสปฺปฏิฆํ น-อสํกิลิฏฺฐสํกิเลสิกํ ธมฺมํ ปฏิจฺจ ฯเปฯ.}

\prose{นสนิทสฺสนสปฺปฏิฆํ น-อสํกิลิฏฺฐ-อสํกิเลสิกํ ธมฺมํ ปฏิจฺจ. เหตุยา ติํส.}

\proseitem{70}{นสนิทสฺสนสปฺปฏิฆํ นสวิตกฺกสวิจารํ ธมฺมํ ปฏิจฺจ ฯเปฯ.}

\prose{นสนิทสฺสนสปฺปฏิฆํ น-อวิตกฺกวิจารมตฺตํ ธมฺมํ ปฏิจฺจ.}

\prose{นสนิทสฺสนสปฺปฏิฆํ น-อวิตกฺก-อวิจารํ ธมฺมํ ปฏิจฺจ ฯเปฯ.}

\proseitem{71}{นสนิทสฺสนสปฺปฏิฆํ นปีติสหคตํ ธมฺมํ ปฏิจฺจ ฯเปฯ.\csromanspacer{} นสนิทสฺสนสปฺปฏิฆํ นสุขสหคตํ ธมฺมํ ปฏิจฺจ ฯเปฯ.\csromanspacer{} นสนิทสฺสนสปฺปฏิฆํ น-อุเปกฺขาสหคตํ ธมฺมํ ปฏิจฺจ.}

\csromanheader{2}{22. สนิทสฺสนตฺติก 2-21. เวทนาตฺติกาทิ}
\proseitem{72}{\csromanfolio{155}นสนิทสฺสนสปฺปฏิฆํ นทสฺสเนน ปหาตพฺพํ ธมฺมํ ปฏิจฺจ ฯเปฯ.\csromanspacer{} นสนิทสฺสนสปฺปฏิฆํ นภาวนาย ปหาตพฺพํ ธมฺมํ ปฏิจฺจ.}

\prose{นสนิทสฺสนสปฺปฏิฆํ นเนวทสฺสเนน นภาวนาย ปหาตพฺพํ ธมฺมํ ปฏิจฺจ.}

\proseitem{73}{นสนิทสฺสนสปฺปฏิฆํ นทสฺสเนน ปหาตพฺพเหตุกํ ธมฺมํ ปฏิจฺจ ฯเปฯ.\csromanspacer{} นสนิทสฺสนสปฺปฏิฆํ นภาวนาย ปหาตพฺพเหตุกํ ธมฺมํ ปฏิจฺจ.}

\prose{นสนิทสฺสนสปฺปฏิฆํ นเนวทสฺสเนน นภาวนาย ปหาตพฺพเหตุกํ ธมฺมํ ปฏิจฺจ.}

\proseitem{74}{นสนิทสฺสนสปฺปฏิฆํ น-อาจยคามิํ ธมฺมํ ปฏิจฺจ ฯเปฯ.\csromanspacer{} นสนิทสฺสนสปฺปฏิฆํ น-อปจยคามิํ ธมฺมํ ปฏิจฺจ.}

\prose{นสนิทสฺสนสปฺปฏิฆํ นเนวาจยคามินาปจยคามิํ ธมฺมํ ปฏิจฺจ.}

\proseitem{75}{นสนิทสฺสนสปฺปฏิฆํ นเสกฺขํ ธมฺมํ ปฏิจฺจ ฯเปฯ.\csromanspacer{} นสนิทสฺสนสปฺปฏิฆํ น-อเสกฺขํ ธมฺมํ ปฏิจฺจ.}

\prose{นสนิทสฺสนสปฺปฏิฆํ นเนวาเสกฺขนาเสกฺขํ ธมฺมํ ปฏิจฺจ ฯเปฯ.}

\proseitem{76}{นสนิทสฺสนสปฺปฏิฆํ นปริตฺตํ ธมฺมํ ปฏิจฺจ.}

\prose{นสนิทสฺสนสปฺปฏิฆํ นมหคฺคตํ ธมฺมํ ปฏิจฺจ ฯเปฯ.\csromanspacer{} นสนิทสฺสนสปฺปฏิฆํ น-อปฺปมาณํ ธมฺมํ ปฏิจฺจ.}

\proseitem{77}{นสนิทสฺสนสปฺปฏิฆํ นปริตฺตารมฺมณํ ธมฺมํ ปฏิจฺจ ฯเปฯ.\csromanspacer{} นสนิทสฺสนสปฺปฏิฆํ นมหคฺคตารมฺมณํ ธมฺมํ ปฏิจฺจ ฯเปฯ.\csromanspacer{} นสนิทสฺสนสปฺปฏิฆํ น-อปฺปมาณารมฺมณํ ธมฺมํ ปฏิจฺจ.}

\proseitem{78}{นสนิทสฺสนสปฺปฏิฆํ นหีนํ ธมฺมํ ปฏิจฺจ.}

\prose{นสนิทสฺสนสปฺปฏิฆํ นมชฺฌิมํ ธมฺมํ ปฏิจฺจ.}

\prose{นสนิทสฺสนสปฺปฏิฆํ นปณีตํ ธมฺมํ ปฏิจฺจ.}

\proseitem{79}{นสนิทสฺสนสปฺปฏิฆํ นมิจฺฉตฺตนิยตํ ธมฺมํ ปฏิจฺจ ฯเปฯ.\csromanspacer{} นสนิทสฺสนสปฺปฏิฆํ นสมฺมตฺตนิยตํ ธมฺมํ ปฏิจฺจ.}

\prose{นสนิทสฺสนสปฺปฏิฆํ น-อนิยตํ ธมฺมํ ปฏิจฺจ ฯเปฯ.}

\gambhira{ธมฺมปจฺจนีย ติกติกปฏฺฐานปาฬิ}
\proseitem{80}{\csromanfolio{156}นสนิทสฺสนสปฺปฏิฆํ นมคฺคารมฺมณํ ธมฺมํ ปฏิจฺจ ฯเปฯ.\csromanspacer{} นสนิทสฺสนสปฺปฏิฆํ นมคฺคเหตุกํ ธมฺมํ ปฏิจฺจ ฯเปฯ.\csromanspacer{} นสนิทสฺสนสปฺปฏิฆํ นมคฺคาธิปติํ ธมฺมํ ปฏิจฺจ.}

\proseitem{81}{นสนิทสฺสนสปฺปฏิฆํ น-อนุปฺปนฺนํ ธมฺมํ ปฏิจฺจ ฯเปฯ.\csromanspacer{} นสนิทสฺสนสปฺปฏิฆํ น-อุปฺปาทิํ ธมฺมํ ปฏิจฺจ.}

\proseitem{82}{นสนิทสฺสนสปฺปฏิฆํ น-อตีตํ ธมฺมํ ปฏิจฺจ ฯเปฯ.\csromanspacer{} นสนิทสฺสนสปฺปฏิฆํ น-อนาคตํ ธมฺมํ ปฏิจฺจ.}

\proseitem{83}{นสนิทสฺสนสปฺปฏิฆํ น-อตีตารมฺมณํ ธมฺมํ ปฏิจฺจ ฯเปฯ.\csromanspacer{} นสนิทสฺสนสปฺปฏิฆํ น-อนาคตารมฺมณํ ธมฺมํ ปฏิจฺจ ฯเปฯ.\csromanspacer{} นสนิทสฺสนสปฺปฏิฆํ นปจฺจุปฺปนฺนารมฺมณํ ธมฺมํ ปฏิจฺจ.}

\proseitem{84}{นสนิทสฺสนสปฺปฏิฆํ น-อชฺฌตฺตํ ธมฺมํ ปฏิจฺจ ฯเปฯ.\csromanspacer{} นสนิทสฺสนสปฺปฏิฆํ นพหิทฺธา ธมฺมํ ปฏิจฺจ.}

\proseitem{85}{นสนิทสฺสนสปฺปฏิฆํ น-อชฺฌตฺตารมฺมณํ ธมฺมํ ปฏิจฺจ นสนิทสฺสนสปฺปฏิโฆ น-อชฺฌตฺตารมฺมโณ ธมฺโม อุปฺปชฺชติ เหตุปจฺจยา.\csromanspacer{}\csromanspacer{}นสนิทสฺสนสปฺปฏิฆํ น-อชฺฌตฺตารมฺมณํ ธมฺมํ ปฏิจฺจ น-อนิทสฺสนสปฺปฏิโฆ น-อชฺฌตฺตารมฺมโณ ธมฺโม อุปฺปชฺชติ เหตุปจฺจยา.}

\csromanheader{2}{เหตุยา ติํส ฯเปฯ อวิคเต ติํส. (สพฺพตฺถ วิตฺถาโร.)}
\prosecont{นสนิทสฺสนสปฺปฏิฆํ นพหิทฺธารมฺมณํ ธมฺมํ ปฏิจฺจ นสนิทสฺสนสปฺปฏิโฆ นพหิทฺธารมฺมโณ ธมฺโม อุปฺปชฺชติ เหตุปจฺจยา.\csromanspacer{}\csromanspacer{}นสนิทสฺสนสปฺปฏิฆํ นพหิทฺธารมฺมณํ ธมฺมํ ปฏิจฺจ น-อนิทสฺสนสปฺปฏิโฆ นพหิทฺธารมฺมโณ ธมฺโม อุปฺปชฺชติ เหตุปจฺจยา.\csromanspacer{}\csromanspacer{}นสนิทสฺสนสปฺปฏิฆํ นพหิทฺธารมฺมณํ ธมฺมํ ปฏิจฺจ น-อนิทสฺสน-อปฺปฏิโฆ นพหิทฺธารมฺมโณ ธมฺโม อุปฺปชฺชติ เหตุปจฺจยา ฯเปฯ.\csromanspacer{} ฉ.}

\proseitem{86}{น-อนิทสฺสนสปฺปฏิฆํ นพหิทฺธารมฺมณํ ธมฺมํ ปฏิจฺจ น-อนิทสฺสนสปฺปฏิโฆ นพหิทฺธารมฺมโณ ธมฺโม อุปฺปชฺชติ เหตุปจฺจยา.\csromanspacer{} ฉ.}

\prose{น-อนิทสฺสน-อปฺปฏิฆํ นพหิทฺธารมฺมณํ ธมฺมํ ปฏิจฺจ น-อนิทสฺสน-อปฺปฏิโฆ นพหิทฺธารมฺมโณ ธมฺโม อุปฺปชฺชติ เหตุปจฺจยา.\csromanspacer{} ฉ.}

\csromanheader{2}{22. สนิทสฺสนตฺติก 2-21. เวทนาตฺติกาทิ}
\prose{\csromanfolio{157}นสนิทสฺสนสปฺปฏิฆํ นพหิทฺธารมฺมณญฺจ น-อนิทสฺสน-อปฺปฏิฆํ นพหิทฺธารมฺมณญฺจ ธมฺมํ ปฏิจฺจ นสนิทสฺสนสปฺปฏิโฆ นพหิทฺธารมฺมโณ ธมฺโม อุปฺปชฺชติ เหตุปจฺจยา.\csromanspacer{} ฉ.}

\prose{น-อนิทสฺสนสปฺปฏิฆํ นพหิทฺธารมฺมณญฺจ น-อนิทสฺสน-อปฺปฏิฆํ นพหิทฺธารมฺมณญฺจ ธมฺมํ ปฏิจฺจ นสนิทสฺสนสปฺปฏิโฆ นพหิทฺธารมฺมโณ ธมฺโม อุปฺปชฺชติ เหตุปจฺจยา.\csromanspacer{} ฉ. (สํขิตฺตํ.)}

\prose{เหตุยา ติํส, อารมฺมเณ นว ฯเปฯ อวิคเต ติํส.}

\prose{(สหชาตวาเรปิ ปจฺจยวาเรปิ นิสฺสยวาเรปิ สํสฏฺฐวาเรปิ สมฺปยุตฺตวาเรปิ ปญฺหาวาเรปิ วิตฺถาโร.)}

\nitthitam{\textbf{ธมฺมปจฺจนีเย ติกติกปฏฺฐานํ นิฏฺฐิตํ.}}
\pitaka{\textbf{อภิธมฺมปิฏก}}
\csromanheader{2}{\textbf{ธมฺมปจฺจนีย}}
\gambhira{\textbf{ทุกทุกปฏฺฐานปาฬิ}}
\namakkaram{นโม ตสฺส ภควโต อรหโต สมฺมาสมฺพุทฺธสฺส.}
\csromanheader{2}{1. เหตุทุก 1-5. สเหตุกทุกาทิ}
\proseitem{1}{\csromanfolio{159}นเหตุํ นสเหตุกํ ธมฺมํ ปฏิจฺจ นเหตุ นสเหตุโก ธมฺโม อุปฺปชฺชติ เหตุปจฺจยา. (สํขิตฺตํ.)}

\prose{เหตุยา ตีณิ, อารมฺมเณ เอกํ ฯเปฯ อวิคเต ตีณิ.}

\proseitem{2}{นเหตุํ นสเหตุกํ ธมฺมํ ปฏิจฺจ นเหตุ นสเหตุโก ธมฺโม อุปฺปชฺชติ นเหตุปจฺจยา. (1)}

\prose{นเหตุํ นสเหตุกํ ธมฺมํ ปฏิจฺจ นเหตุ นสเหตุโก ธมฺโม อุปฺปชฺชติ น-อารมฺมณปจฺจยา. (สํขิตฺตํ.)}

\prose{นเหตุยา เอกํ, น-อารมฺมเณ ตีณิ, น-อธิปติยา ตีณิ.}

\csromanheader{2}{\textbf{เหตุ-อารมฺมณ}}
\csromanheader{2}{3. นนเหตุ นสเหตุโก ธมฺโม นเหตุสฺส นสเหตุกสฺส ธมฺมสฺส}
\prose{นเหตุ นสเหตุโก ธมฺโม นเหตุสฺส นสเหตุกสฺส ธมฺมสฺส อารมฺมณปจฺจเยน ปจฺจโย.\csromanspacer{}\csromanspacer{}นเหตุ นสเหตุโก ธมฺโม นนเหตุสฺส นสเหตุกสฺส ธมฺมสฺส อารมฺมณปจฺจเยน ปจฺจโย.}

\prose{เหตุยา เอกํ, อารมฺมเณ จตฺตาริ ฯเปฯ อวิคเต จตฺตาริ.}

\gambhira{ธมฺมปจฺจนีย ทุกทุกปฏฺฐานปาฬิ}
\proseitem{4}{\csromanfolio{160}นเหตุํ น-อเหตุกํ ธมฺมํ ปฏิจฺจ นเหตุ น-อเหตุโก ธมฺโม อุปฺปชฺชติ เหตุปจฺจยา.\csromanspacer{}\csromanspacer{}นเหตุํ น-อเหตุกํ ธมฺมํ ปฏิจฺจ นนเหตุ น-อเหตุโก ธมฺโม อุปฺปชฺชติ เหตุปจฺจยา.\csromanspacer{}\csromanspacer{}นเหตุํ น-อเหตุกํ ธมฺมํ ปฏิจฺจ นเหตุ น-อเหตุโก จ นนเหตุ น-อเหตุโก จ ธมฺมา อุปฺปชฺชนฺติ เหตุปจฺจยา. (3)}

\prose{นนเหตุํ น-อเหตุกํ ธมฺมํ ปฏิจฺจ นนเหตุ น-อเหตุโก ธมฺโม อุปฺปชฺชติ เหตุปจฺจยา.\csromanspacer{}\csromanspacer{}นนเหตุํ น-อเหตุกํ ธมฺมํ ปฏิจฺจ นเหตุ น-อเหตุโก ธมฺโม อุปฺปชฺชติ เหตุปจฺจยา.\csromanspacer{}\csromanspacer{}นนเหตุํ น-อเหตุกํ ธมฺมํ ปฏิจฺจ นเหตุ น-อเหตุโก จ นนเหตุ น-อเหตุโก จ ธมฺมา อุปฺปชฺชนฺติ เหตุปจฺจยา. (3)}

\prose{นเหตุํ น-อเหตุกญฺจ นนเหตุํ น-อเหตุกญฺจ ธมฺมํ ปฏิจฺจ นเหตุ น-อเหตุโก ธมฺโม อุปฺปชฺชติ เหตุปจฺจยา.\csromanspacer{}\csromanspacer{}นเหตุํ น-อเหตุกญฺจ นนเหตุํ น-อเหตุกญฺจ ธมฺมํ ปฏิจฺจ นนเหตุ น-อเหตุโก ธมฺโม อุปฺปชฺชติ เหตุปจฺจยา.\csromanspacer{}\csromanspacer{}นเหตุํ น-อเหตุกญฺจ นนเหตุํ น-อเหตุกญฺจ ธมฺมํ ปฏิจฺจ นเหตุ น-อเหตุโก จ นนเหตุ น-อเหตุโก จ ธมฺมา อุปฺปชฺชนฺติ เหตุปจฺจยา. (3) (สํขิตฺตํ.)}

\prose{เหตุยา นว, อารมฺมเณ นว ฯเปฯ อวิคเต นว. (สพฺพตฺถ นว.)}

\proseitem{5}{นเหตุํ นเหตุสมฺปยุตฺตํ ธมฺมํ ปฏิจฺจ นเหตุ นเหตุสมฺปยุตฺโต ธมฺโม อุปฺปชฺชติ เหตุปจฺจยา.\csromanspacer{} ตีณิ.}

\proseitem{6}{นเหตุํ นเหตุวิปฺปยุตฺตํ ธมฺมํ ปฏิจฺจ นเหตุ นเหตุวิปฺปยุตฺโต ธมฺโม อุปฺปชฺชติ เหตุปจฺจยา.\csromanspacer{} ตีณิ.}

\prose{นนเหตุํ นเหตุวิปฺปยุตฺตํ ธมฺมํ ปฏิจฺจ นนเหตุ นเหตุวิปฺปยุตฺโต ธมฺโม อุปฺปชฺชติ เหตุปจฺจยา.\csromanspacer{} ตีณิ.}

\prose{นเหตุํ นเหตุวิปฺปยุตฺตญฺจ นนเหตุํ นเหตุวิปฺปยุตฺตญฺจ ธมฺมํ ปฏิจฺจ นเหตุ นเหตุวิปฺปยุตฺโต ธมฺโม อุปฺปชฺชติ เหตุปจฺจยา.\csromanspacer{} ตีณิ.}

\csromanheader{2}{เหตุยา นว. (สพฺพตฺถ วิตฺถาโร.)}
\csromanheader{2}{1. เหตุทุก 6. จูฬนฺตรทุก}
\proseitem{7}{\csromanfolio{161}นเหตุํ นเหตุญฺเจว น-อเหตุกญฺจ ธมฺมํ ปฏิจฺจ นเหตุ นเหตุ เจว น-อเหตุโก จ ธมฺโม อุปฺปชฺชติ เหตุปจฺจยา. (ยาว ปญฺหาวาเรปิ เอกํ.)}

\prose{นนเหตุํ น-อเหตุกญฺเจว นน จ เหตุํ ธมฺมํ ปฏิจฺจ นนเหตุ น-อเหตุโก เจว นน จ เหตุ ธมฺโม อุปฺปชฺชติ เหตุปจฺจยา.}

\prose{นเหตุํ นเหตุญฺเจว นเหตุวิปฺปยุตฺตญฺจ ธมฺมํ ปฏิจฺจ นเหตุ นเหตุ เจว นเหตุวิปฺปยุตฺโต จ ธมฺโม อุปฺปชฺชติ เหตุปจฺจยา.}

\prose{นนเหตุํ นเหตุวิปฺปยุตฺตญฺเจว นน จ เหตุํ ธมฺมํ ปฏิจฺจ นนเหตุ นเหตุวิปฺปยุตฺโต เจว นน จ เหตุ ธมฺโม อุปฺปชฺชติ เหตุปจฺจยา.}

\prose{นเหตุํ นสเหตุกํ ธมฺมํ ปฏิจฺจ นเหตุ นสเหตุโก ธมฺโม อุปฺปชฺชติ เหตุปจฺจยา.}

\prose{นเหตุํ น-อเหตุกํ ธมฺมํ ปฏิจฺจ นเหตุ น-อเหตุโก ธมฺโม อุปฺปชฺชติ เหตุปจฺจยา. (สพฺพตฺถ เอกํ.)}

\csromanheader{2}{1. เหตุทุก 6. จูฬนฺตรทุก}
\proseitem{8}{นเหตุํ น-อปฺปจฺจยํ ธมฺมํ ปฏิจฺจ นเหตุ น-อปฺปจฺจโย ธมฺโม อุปฺปชฺชติ เหตุปจฺจยา.\csromanspacer{}\csromanspacer{}นเหตุํ น-อปฺปจฺจยํ ธมฺมํ ปฏิจฺจ นนเหตุ น-อปฺปจฺจโย ธมฺโม อุปฺปชฺชติ เหตุปจฺจยา.\csromanspacer{}\csromanspacer{}นเหตุํ น-อปฺปจฺจยํ ธมฺมํ ปฏิจฺจ นเหตุ น-อปฺปจฺจโย จ นนเหตุ น-อปฺปจฺจโย จ ธมฺมา อุปฺปชฺชนฺติ เหตุปจฺจยา. (3)}

\prose{นนเหตุํ น-อปฺปจฺจยํ ธมฺมํ ปฏิจฺจ นนเหตุ น-อปฺปจฺจโย ธมฺโม อุปฺปชฺชติ เหตุปจฺจยา.\csromanspacer{}\csromanspacer{}นนเหตุํ น-อปฺปจฺจยํ ธมฺมํ ปฏิจฺจ นเหตุ น-อปฺปจฺจโย ธมฺโม อุปฺปชฺชติ เหตุปจฺจยา.\csromanspacer{}\csromanspacer{}นนเหตุํ น-อปฺปจฺจยํ ธมฺมํ ปฏิจฺจ นเหตุ น-อปฺปจฺจโย จ นนเหตุ น-อปฺปจฺจโย จ ธมฺมา อุปฺปชฺชนฺติ เหตุปจฺจยา. (3)}

\prose{นเหตุํ น-อปฺปจฺจยญฺจ นนเหตุํ น-อปฺปจฺจยญฺจ ธมฺมํ ปฏิจฺจ นเหตุ น-อปฺปจฺจโย ธมฺโม อุปฺปชฺชติ เหตุปจฺจยา.\csromanspacer{}\csromanspacer{}นเหตุํ น-อปฺปจฺจยญฺจ นนเหตุํ น-อปฺปจฺจยญฺจ ธมฺมํ ปฏิจฺจ นนเหตุ น-อปฺปจฺจโย ธมฺโม อุปฺปชฺชติ เหตุปจฺจยา.\csromanspacer{}\csromanspacer{}นเหตุํ น-อปฺปจฺจยญฺจ นนเหตุํ น-อปฺปจฺจยญฺจ ธมฺมํ ปฏิจฺจ นเหตุ น-อปฺปจฺจโย จ นนเหตุ น-อปฺปจฺจโย จ ธมฺมา อุปฺปชฺชนฺติ เหตุปจฺจยา. (3)}

\gambhira{ธมฺมปจฺจนีย ทุกทุกปฏฺฐานปาฬิ}
\prose{\csromanfolio{162}เหตุยา นว.}

\prose{นเหตุํ น-อสงฺขตํ ธมฺมํ ปฏิจฺจ.}

\prose{นเหตุํ นสนิทสฺสนํ ธมฺมํ ปฏิจฺจ.}

\prose{นเหตุํ นสปฺปฏิฆํ ธมฺมํ ปฏิจฺจ.}

\prose{นเหตุํ น-อปฺปฏิฆํ ธมฺมํ ปฏิจฺจ.}

\prose{นเหตุํ นรูปิํ ธมฺมํ ปฏิจฺจ.}

\prose{นเหตุํ น-อรูปิํ ธมฺมํ ปฏิจฺจ.}

\prose{นเหตุํ นโลกิยํ ธมฺมํ ปฏิจฺจ ฯเปฯ.\csromanspacer{} นเหตุํ นโลกุตฺตรํ ธมฺมํ ปฏิจฺจ.}

\prose{นเหตุํ นเกนจิ วิญฺเญยฺยํ ธมฺมํ ปฏิจฺจ ฯเปฯ.\csromanspacer{} นเหตุํ นนเกนจิ วิญฺเญยฺยํ ธมฺมํ ปฏิจฺจ.}

\csromanheader{2}{1. เหตุทุก 13-18. อาสวโคจฺฉก}
\proseitem{9}{นเหตุํ โน-อาสวํ ธมฺมํ ปฏิจฺจ.}

\prose{นเหตุํ นโน-อาสวํ ธมฺมํ ปฏิจฺจ.}

\prose{นเหตุํ นสาสวํ ธมฺมํ ปฏิจฺจ ฯเปฯ.\csromanspacer{} นเหตุํ น-อนาสวํ ธมฺมํ ปฏิจฺจ.}

\prose{นเหตุํ น-อาสวสมฺปยุตฺตํ ธมฺมํ ปฏิจฺจ.}

\prose{นเหตุํ น-อาสววิปฺปยุตฺตํ ธมฺมํ ปฏิจฺจ.}

\prose{นเหตุํ โน-อาสวญฺเจว น-อนาสวญฺจ ธมฺมํ ปฏิจฺจ.}

\prose{นเหตุํ น-อนาสวญฺเจว นโน จ อาสวํ ธมฺมํ ปฏิจฺจ.}

\prose{นเหตุํ น-อาสวญฺเจว น-อาสววิปฺปยุตฺตญฺจ ธมฺมํ ปฏิจฺจ.}

\prose{นเหตุํ น-อาสววิปฺปยุตฺตญฺเจว นโน จ อาสวํ ธมฺมํ ปฏิจฺจ.}

\prose{นเหตุํ อาสววิปฺปยุตฺตํ นสาสวํ ธมฺมํ ปฏิจฺจ ฯเปฯ.\csromanspacer{} นเหตุํ อาสววิปฺปยุตฺตํ น-อนาสวํ ธมฺมํ ปฏิจฺจ.}

\csromanheader{2}{1. เหตุทุก 19-53. สญฺโญชนาทิทุก}
\proseitem{10}{นเหตุํ โนสญฺโญชนํ ธมฺมํ ปฏิจฺจ ฯเปฯ.\csromanspacer{} นเหตุํ โนคนฺถํ ธมฺมํ ปฏิจฺจ.\csromanspacer{} นเหตุํ โน-โอฆํ ธมฺมํ ปฏิจฺจ ฯเปฯ.\csromanspacer{} นเหตุํ โนโยคํ ธมฺมํ ปฏิจฺจ.}

\csromanheader{2}{1. เหตุทุก 82. ปิฏฺฐิทุก}
\prose{\csromanfolio{163}นเหตุํ โนนีวรณํ ธมฺมํ ปฏิจฺจ.}

\prose{นเหตุํ โนปรามาสํ ธมฺมํ ปฏิจฺจ.}

\csromanheader{2}{1. เหตุทุก 54. มหนฺตรทุก}
\csromanheader{2}{11. นเหตุํ นสารมฺมณํ ธมฺมํ ปฏิจฺจ. (สํขิตฺตํ.)}
\prose{นเหตุํ นจิตฺตํ ธมฺมํ ปฏิจฺจ.}

\prose{นเหตุํ นเจตสิกํ ธมฺมํ ปฏิจฺจ.}

\prose{นเหตุํ นจิตฺตสมฺปยุตฺตํ ธมฺมํ ปฏิจฺจ ฯเปฯ.\csromanspacer{} นเหตุํ นจิตฺตสํสฏฺฐํ ธมฺมํ ปฏิจฺจ ฯเปฯ.}

\prose{นเหตุํ นจิตฺตสมุฏฺฐานํ ธมฺมํ ปฏิจฺจ.}

\prose{นเหตุํ นจิตฺตสหภุํ ธมฺมํ ปฏิจฺจ ฯเปฯ.\csromanspacer{} นเหตุํ นจิตฺตานุปริวตฺติํ ธมฺมํ ปฏิจฺจ.}

\prose{นเหตุํ นจิตฺตสํสฏฺฐสมุฏฺฐานํ ธมฺมํ ปฏิจฺจ ฯเปฯ.\csromanspacer{} นเหตุํ นจิตฺตสํสฏฺฐสมุฏฺฐานสหภุํ ธมฺมํ ปฏิจฺจ ฯเปฯ.\csromanspacer{} นเหตุํ นจิตฺตสํสฏฺฐสมุฏฺฐานานุปริวตฺติํ ธมฺมํ ปฏิจฺจ.}

\prose{นเหตุํ น-อชฺฌตฺติกํ ธมฺมํ ปฏิจฺจ.}

\prose{นเหตุํ นพาหิรํ ธมฺมํ ปฏิจฺจ.}

\prose{นเหตุํ น-อุปาทาธมฺมํ ปฏิจฺจ.}

\prose{นเหตุํ น-อุปาทินฺนํ ธมฺมํ ปฏิจฺจ.}

\prose{นเหตุํ น-อนุปาทินฺนํ ธมฺมํ ปฏิจฺจ.}

\prose{นเหตุํ โน-อุปาทานํ ธมฺมํ ปฏิจฺจ ฯเปฯ.\csromanspacer{} นเหตุํ โนกิเลสํ ธมฺมํ ปฏิจฺจ ฯเปฯ.}

\csromanheader{2}{1. เหตุทุก 82. ปิฏฺฐิทุก}
\proseitem{12}{นเหตุํ นทสฺสเนน ปหาตพฺพํ ธมฺมํ ปฏิจฺจ.}

\prose{นเหตุํ นนทสฺสเนน ปหาตพฺพํ ธมฺมํ ปฏิจฺจ.}

\proseitem{13}{นเหตุํ นภาวนาย ปหาตพฺพํ ธมฺมํ ปฏิจฺจ ฯเปฯ.\csromanspacer{} นเหตุํ นนภาวนาย ปหาตพฺพํ ธมฺมํ ปฏิจฺจ ฯเปฯ.\csromanspacer{} นเหตุํ นทสฺสเนน ปหาตพฺพเหตุกํ ธมฺมํ ปฏิจฺจ ฯเปฯ.\csromanspacer{} นเหตุํ นนทสฺสเนน ปหาตพฺพเหตุกํ ธมฺมํ ปฏิจฺจ ฯเปฯ.\csromanspacer{} นเหตุํ นภาวนาย ปหาตพฺพเหตุกํ ธมฺมํ ปฏิจฺจ. นเหตุํ นนภาวนาย ปหาตพฺพเหตุกํ ธมฺมํ ปฏิจฺจ.}

\gambhira{ธมฺมปจฺจนีย ทุกทุกปฏฺฐานปาฬิ}
\proseitem{14}{\csromanfolio{164}นเหตุํ นสวิตกฺกํ ธมฺมํ ปฏิจฺจ ฯเปฯ.\csromanspacer{} นเหตุํ น-อวิตกฺกํ ธมฺมํ ปฏิจฺจ ฯเปฯ.\csromanspacer{} นเหตุํ นสวิจารํ ธมฺมํ ปฏิจฺจ ฯเปฯ.\csromanspacer{} นเหตุํ น-อวิจารํ ธมฺมํ ปฏิจฺจ ฯเปฯ.\csromanspacer{} นเหตุํ นสปฺปีติกํ ธมฺมํ ปฏิจฺจ ฯเปฯ.\csromanspacer{} นเหตุํ น-อปฺปีติกํ ธมฺมํ ปฏิจฺจ ฯเปฯ.\csromanspacer{} นเหตุํ นปีติสหคตํ ธมฺมํ ปฏิจฺจ ฯเปฯ.\csromanspacer{} นเหตุํ นนปีติสหคตํ ธมฺมํ ปฏิจฺจ ฯเปฯ.\csromanspacer{} นเหตุํ นสุขสหคตํ ธมฺมํ ปฏิจฺจ ฯเปฯ.\csromanspacer{} นเหตุํ นนสุขสหคตํ ธมฺมํ ปฏิจฺจ ฯเปฯ.\csromanspacer{} นเหตุํ น-อุเปกฺขาสหคตํ ธมฺมํ ปฏิจฺจ ฯเปฯ.\csromanspacer{} นเหตุํ นน-อุเปกฺขาสหคตํ ธมฺมํ ปฏิจฺจ ฯเปฯ.}

\proseitem{15}{นเหตุํ นกามาวจรํ ธมฺมํ ปฏิจฺจ ฯเปฯ.\csromanspacer{} นเหตุํ นนกามาวจรํ ธมฺมํ ปฏิจฺจ ฯเปฯ.\csromanspacer{} นเหตุํ นรูปาวจรํ ธมฺมํ ปฏิจฺจ ฯเปฯ.\csromanspacer{} นเหตุํ นนรูปาวจรํ ธมฺมํ ปฏิจฺจ ฯเปฯ.\csromanspacer{} นเหตุํ น-อรูปาวจรํ ธมฺมํ ปฏิจฺจ ฯเปฯ.\csromanspacer{} นเหตุํ นน-อรูปาวจรํ ธมฺมํ ปฏิจฺจ ฯเปฯ.}

\proseitem{16}{นเหตุํ นปริยาปนฺนํ ธมฺมํ ปฏิจฺจ ฯเปฯ.\csromanspacer{} นเหตุํ น-อปริยาปนฺนํ ธมฺมํ ปฏิจฺจ ฯเปฯ.\csromanspacer{} นเหตุํ นนิยฺยานิกํ ธมฺมํ ปฏิจฺจ ฯเปฯ.\csromanspacer{} นเหตุํ น-อนิยฺยานิกํ ธมฺมํ ปฏิจฺจ ฯเปฯ.\csromanspacer{} นเหตุํ นนิยตํ ธมฺมํ ปฏิจฺจ ฯเปฯ.\csromanspacer{} นเหตุํ น-อนิยตํ ธมฺมํ ปฏิจฺจ ฯเปฯ.}

\prose{นเหตุํ นส-อุตฺตรํ ธมฺมํ ปฏิจฺจ ฯเปฯ.\csromanspacer{} นเหตุํ น-อนุตฺตรํ ธมฺมํ ปฏิจฺจ ฯเปฯ.}

\proseitem{17}{นเหตุํ นสรณํ ธมฺมํ ปฏิจฺจ นเหตุ นสรโณ ธมฺโม อุปฺปชฺชติ เหตุปจฺจยา. (นว.)}

\prose{นเหตุํ น-อรณํ ธมฺมํ ปฏิจฺจ นเหตุ น-อรโณ ธมฺโม อุปฺปชฺชติ เหตุปจฺจยา.\csromanspacer{}\csromanspacer{}นเหตุํ น-อรณํ ธมฺมํ ปฏิจฺจ นนเหตุ น-อรโณ ธมฺโม อุปฺปชฺชติ เหตุปจฺจยา.\csromanspacer{}\csromanspacer{}นเหตุํ น-อรณํ ธมฺมํ ปฏิจฺจ นเหตุ น-อรโณ จ นนเหตุ น-อรโณ จ ธมฺมา อุปฺปชฺชนฺติ เหตุปจฺจยา. (3)}

\prose{นนเหตุํ น-อรณํ ธมฺมํ ปฏิจฺจ นนเหตุ น-อรโณ ธมฺโม อุปฺปชฺชติ เหตุปจฺจยา.\csromanspacer{}\csromanspacer{}นนเหตุํ น-อรณํ ธมฺมํ ปฏิจฺจ นเหตุ น-อรโณ ธมฺโม อุปฺปชฺชติ เหตุปจฺจยา.\csromanspacer{}\csromanspacer{}นนเหตุํ น-อรณํ ธมฺมํ ปฏิจฺจ นเหตุ น-อรโณ จ นนเหตุ น-อรโณ จ ธมฺมา อุปฺปชฺชนฺติ เหตุปจฺจยา. (3)}

\prose{นเหตุํ น-อรณญฺจ นนเหตุํ น-อรณญฺจ ธมฺมํ ปฏิจฺจ นเหตุ น-อรโณ ธมฺโม อุปฺปชฺชติ เหตุปจฺจยา.\csromanspacer{}\csromanspacer{}นเหตุํ น-อรณญฺจ นนเหตุํ น-อรณญฺจ ธมฺมํ ปฏิจฺจ นนเหตุ น-อรโณ ธมฺโม อุปฺปชฺชติ เหตุปจฺจยา.\csromanspacer{}\csromanspacer{}นเหตุํ น-อรณญฺจ นนเหตุํ น-อรณญฺจ ธมฺมํ ปฏิจฺจ นเหตุ น-อรโณ จ นนเหตุ}

\vspace{\tipitakaparskip}
\csromancenter{สเหตุกาทิทุก 1.\csromanspacer{} เหตุทุก น-อรโณ จ ธมฺมา อุปฺปชฺชนฺติ เหตุปจฺจยา. (3) (สํขิตฺตํ.) เหตุยา นว. (สพฺพตฺถ นว.)}
\csromanheader{2}{2. สเหตุกาทิทุก 1. เหตุทุก}
\proseitem{18}{\csromanfolio{165}นสเหตุกํ นเหตุํ ธมฺมํ ปฏิจฺจ นสเหตุโก นเหตุ ธมฺโม อุปฺปชฺชติ เหตุปจฺจยา.\csromanspacer{}\csromanspacer{}นสเหตุกํ นเหตุํ ธมฺมํ ปฏิจฺจ น-อเหตุโก นเหตุ ธมฺโม อุปฺปชฺชติ เหตุปจฺจยา.\csromanspacer{}\csromanspacer{}นสเหตุกํ นเหตุํ ธมฺมํ ปฏิจฺจ นสเหตุโก นเหตุ จ น-อเหตุโก นเหตุ จ ธมฺมา อุปฺปชฺชนฺติ เหตุปจฺจยา. (3)}

\prose{น-อเหตุกํ นเหตุํ ธมฺมํ ปฏิจฺจ น-อเหตุโก นเหตุ ธมฺโม อุปฺปชฺชติ เหตุปจฺจยา.\csromanspacer{}\csromanspacer{}น-อเหตุกํ นเหตุํ ธมฺมํ ปฏิจฺจ นสเหตุโก นเหตุ ธมฺโม อุปฺปชฺชติ เหตุปจฺจยา.\csromanspacer{}\csromanspacer{}น-อเหตุกํ นเหตุํ ธมฺมํ ปฏิจฺจ นสเหตุโก นเหตุ จ น-อเหตุโก นเหตุ จ ธมฺมา อุปฺปชฺชนฺติ เหตุปจฺจยา. (3)}

\prose{นสเหตุกํ นเหตุญฺจ น-อเหตุกํ นเหตุญฺจ ธมฺมํ ปฏิจฺจ นสเหตุโก นเหตุ ธมฺโม อุปฺปชฺชติ เหตุปจฺจยา.\csromanspacer{}\csromanspacer{}นสเหตุกํ นเหตุญฺจ น-อเหตุกํ นเหตุญฺจ ธมฺมํ ปฏิจฺจ น-อเหตุโก นเหตุ ธมฺโม อุปฺปชฺชติ เหตุปจฺจยา.\csromanspacer{}\csromanspacer{}นสเหตุกํ นเหตุญฺจ น-อเหตุกํ นเหตุญฺจ ธมฺมํ ปฏิจฺจ นสเหตุโก นเหตุ จ น-อเหตุโก นเหตุ จ ธมฺมา อุปฺปชฺชนฺติ เหตุปจฺจยา. (3) (สํขิตฺตํ.)}

\csromanheader{2}{เหตุยา นว. (สพฺพตฺถ วิตฺถาโร.)}
\proseitem{19}{น-อเหตุกํ นนเหตุํ ธมฺมํ ปฏิจฺจ น-อเหตุโก นนเหตุ ธมฺโม อุปฺปชฺชติ เหตุปจฺจยา. (สํขิตฺตํ.)}

\csromanheader{2}{เหตุยา เอกํ ฯเปฯ อวิคเต เอกํ. (สพฺพตฺถ เอกํ.)}
\proseitem{20}{นเหตุสมฺปยุตฺตํ นเหตุํ ธมฺมํ ปฏิจฺจ.}

\prose{นเหตุวิปฺปยุตฺตํ นนเหตุํ ธมฺมํ ปฏิจฺจ.}

\proseitem{21}{นเหตุญฺเจว น-อเหตุกญฺจ นเหตุํ ธมฺมํ ปฏิจฺจ ฯเปฯ.\csromanspacer{} น-อเหตุกญฺเจว นน จ เหตุํ ธมฺมํ ปฏิจฺจ ฯเปฯ.\csromanspacer{} นเหตุญฺเจว นเหตุวิปฺปยุตฺตญฺจ นเหตุํ ธมฺมํ ปฏิจฺจ ฯเปฯ.\csromanspacer{} นเหตุวิปฺปยุตฺตญฺเจว นน จ เหตุํ ธมฺมํ ปฏิจฺจ.}

\gambhira{ธมฺมปจฺจนีย ทุกทุกปฏฺฐานปาฬิ}
\proseitem{22}{\csromanfolio{166}นเหตุํ นสเหตุกํ นเหตุํ ธมฺมํ ปฏิจฺจ ฯเปฯ.\csromanspacer{} นเหตุํ น-อเหตุกํ นเหตุํ ธมฺมํ ปฏิจฺจ.}

\csromanheader{2}{7. จูฬนฺตรทุก 1. เหตุทุก}
\proseitem{23}{น-อปฺปจฺจยํ นเหตุํ ธมฺมํ ปฏิจฺจ ฯเปฯ.\csromanspacer{} น-อสงฺขตํ นเหตุํ ธมฺมํ ปฏิจฺจ.}

\prose{นสนิทสฺสนํ นเหตุํ ธมฺมํ ปฏิจฺจ.}

\prose{นสปฺปฏิฆํ นเหตุํ ธมฺมํ ปฏิจฺจ ฯเปฯ.\csromanspacer{} น-อปฺปฏิฆํ นเหตุํ ธมฺมํ ปฏิจฺจ.}

\prose{นรูปิํ นเหตุํ ธมฺมํ ปฏิจฺจ ฯเปฯ.\csromanspacer{} น-อรูปิํ นเหตุํ ธมฺมํ ปฏิจฺจ.}

\prose{นโลกิยํ นเหตุํ ธมฺมํ ปฏิจฺจ ฯเปฯ.\csromanspacer{} นโลกุตฺตรํ นเหตุํ ธมฺมํ ปฏิจฺจ.}

\prose{นเกนจิ วิญฺเญยฺยํ นเหตุํ ธมฺมํ ปฏิจฺจ ฯเปฯ.\csromanspacer{} นนเกนจิ วิญฺเญยฺยํ นเหตุํ ธมฺมํ ปฏิจฺจ.}

\csromanheader{2}{14. อาสวโคจฺฉก 1. เหตุทุก}
\proseitem{24}{โน-อาสวํ นเหตุํ ธมฺมํ ปฏิจฺจ ฯเปฯ.\csromanspacer{} นโน-อาสวํ นเหตุํ ธมฺมํ ปฏิจฺจ.}

\prose{นสาสวํ นเหตุํ ธมฺมํ ปฏิจฺจ ฯเปฯ.\csromanspacer{} น-อนาสวํ นเหตุํ ธมฺมํ ปฏิจฺจ.}

\prose{น-อาสวสมฺปยุตฺตํ นเหตุํ ธมฺมํ ปฏิจฺจ ฯเปฯ.\csromanspacer{} น-อาสววิปฺปยุตฺตํ นเหตุํ ธมฺมํ ปฏิจฺจ.}

\prose{น-อาสวญฺเจว น-อนาสวํ นเหตุํ ธมฺมํ ปฏิจฺจ ฯเปฯ.\csromanspacer{} น-อนาสวญฺเจว นโน จ อาสวํ นเหตุํ ธมฺมํ ปฏิจฺจ.}

\prose{น-อาสวญฺเจว น-อาสววิปฺปยุตฺตญฺจ นเหตุํ ธมฺมํ ปฏิจฺจ.}

\prose{น-อาสววิปฺปยุตฺตญฺเจว นโน จ อาสวํ นเหตุํ ธมฺมํ ปฏิจฺจ ฯเปฯ.\csromanspacer{} อาสววิปฺปยุตฺตํ นสาสวํ นเหตุํ ธมฺมํ ปฏิจฺจ ฯเปฯ.\csromanspacer{} อาสววิปฺปยุตฺตํ น-อนาสวํ นเหตุํ ธมฺมํ ปฏิจฺจ.}

\csromanheader{2}{20. สญฺโญชนาทิทุก 1. เหตุทุก}
\proseitem{25}{โนสญฺโญชนํ นเหตุํ ธมฺมํ ปฏิจฺจ.}

\proseitem{26}{โนคนฺถํ นเหตุํ ธมฺมํ ปฏิจฺจ.}

\csromanheader{2}{83. ปิฏฺฐิทุก 1. เหตุทุก}
\proseitem{27}{\csromanfolio{167}โน-โอฆํ นเหตุํ ธมฺมํ ปฏิจฺจ ฯเปฯ.\csromanspacer{} โนโยคํ นเหตุํ ธมฺมํ ปฏิจฺจ.}

\proseitem{28}{โนนีวรณํ นเหตุํ ธมฺมํ ปฏิจฺจ.}

\proseitem{29}{โนปรามาสํ นเหตุํ ธมฺมํ ปฏิจฺจ.}

\csromanheader{2}{55. มหนฺตรทุก 1. เหตุทุก}
\csromanheader{2}{30. นสารมฺมณํ นเหตุํ ธมฺมํ ปฏิจฺจ. (สํขิตฺตํ.)}
\proseitem{31}{นจิตฺตํ นเหตุํ ธมฺมํ ปฏิจฺจ.}

\prose{นเจตสิกํ นเหตุํ ธมฺมํ ปฏิจฺจ.}

\prose{นจิตฺตสมฺปยุตฺตํ นเหตุํ ธมฺมํ ปฏิจฺจ ฯเปฯ.\csromanspacer{} นจิตฺตสํสฏฺฐํ นเหตุํ ธมฺมํ ปฏิจฺจ.}

\prose{นจิตฺตสมุฏฺฐานํ นเหตุํ ธมฺมํ ปฏิจฺจ.}

\proseitem{32}{นจิตฺตสหภุํ นเหตุํ ธมฺมํ ปฏิจฺจ ฯเปฯ.\csromanspacer{} นจิตฺตานุปริวตฺติํ นเหตุํ ธมฺมํ ปฏิจฺจ ฯเปฯ.\csromanspacer{} นจิตฺตสํสฏฺฐสมุฏฺฐานํ นเหตุํ ธมฺมํ ปฏิจฺจ ฯเปฯ.\csromanspacer{} นจิตฺตสํสฏฺฐสมุฏฺฐานสหภุํ นเหตุํ ธมฺมํ ปฏิจฺจ ฯเปฯ.\csromanspacer{} นจิตฺตสํสฏฺฐสมุฏฺฐานานุปริวตฺติํ นเหตุํ ธมฺมํ ปฏิจฺจ.}

\proseitem{33}{น-อชฺฌตฺติกํ นเหตุํ ธมฺมํ ปฏิจฺจ ฯเปฯ.\csromanspacer{} นพาหิรํ นเหตุํ ธมฺมํ ปฏิจฺจ.}

\prose{โน-อุปาทา นเหตุํ ธมฺมํ ปฏิจฺจ.}

\proseitem{34}{น-อุปาทินฺนํ นเหตุํ ธมฺมํ ปฏิจฺจ ฯเปฯ.\csromanspacer{} น-อนุปาทินฺนํ นเหตุํ ธมฺมํ ปฏิจฺจ.}

\csromanheader{2}{69. อุปาทานทุกาทิ 1. เหตุทุก}
\proseitem{35}{โน-อุปาทานํ นเหตุํ ธมฺมํ ปฏิจฺจ.}

\proseitem{36}{โนกิเลสํ นเหตุํ ธมฺมํ ปฏิจฺจ ฯเปฯ.\csromanspacer{} นโนกิเลสํ นเหตุํ ธมฺมํ ปฏิจฺจ.}

\csromanheader{2}{83. ปิฏฺฐิทุก 1. เหตุทุก}
\proseitem{37}{นทสฺสเนน ปหาตพฺพํ นเหตุํ ธมฺมํ ปฏิจฺจ ฯเปฯ.\csromanspacer{} นนทสฺสเนน ปหาตพฺพํ นเหตุํ ธมฺมํ ปฏิจฺจ ฯเปฯ.\csromanspacer{} นภาวนาย ปหาตพฺพํ นเหตุํ ธมฺมํ}

\vspace{\tipitakaparskip}
\csromancenter{ธมฺมปจฺจนีย ทุกทุกปฏฺฐานปาฬิ ปฏิจฺจ ฯเปฯ.\csromanspacer{} นนภาวนาย ปหาตพฺพํ นเหตุํ ธมฺมํ ปฏิจฺจ ฯเปฯ.\csromanspacer{} นทสฺสเนน ปหาตพฺพเหตุกํ นเหตุํ ธมฺมํ ปฏิจฺจ ฯเปฯ.\csromanspacer{} นนทสฺสเนน ปหาตพฺพเหตุกํ นเหตุํ ธมฺมํ ปฏิจฺจ ฯเปฯ.\csromanspacer{} นภาวนาย ปหาตพฺพเหตุกํ นเหตุํ ธมฺมํ ปฏิจฺจ ฯเปฯ.\csromanspacer{} นนภาวนาย ปหาตพฺพเหตุกํ นเหตุํ ธมฺมํ ปฏิจฺจ.}
\proseitem{38}{\csromanfolio{168}นสวิตกฺกํ นเหตุํ ธมฺมํ ปฏิจฺจ ฯเปฯ.\csromanspacer{} น-อวิตกฺกํ นเหตุํ ธมฺมํ ปฏิจฺจ ฯเปฯ.\csromanspacer{} นสวิจารํ นเหตุํ ธมฺมํ ปฏิจฺจ ฯเปฯ.\csromanspacer{} น-อวิจารํ นเหตุํ ธมฺมํ ปฏิจฺจ ฯเปฯ.\csromanspacer{} นสปฺปีติกํ นเหตุํ ธมฺมํ ปฏิจฺจ ฯเปฯ.\csromanspacer{} น-อปฺปีติกํ นเหตุํ ธมฺมํ ปฏิจฺจ ฯเปฯ.\csromanspacer{} นปีติสหคตํ นเหตุํ ธมฺมํ ปฏิจฺจ ฯเปฯ.\csromanspacer{} นนปีติสหคตํ นเหตุํ ธมฺมํ ปฏิจฺจ ฯเปฯ.\csromanspacer{} นสุขสหคตํ นเหตุํ ธมฺมํ ปฏิจฺจ ฯเปฯ.\csromanspacer{} นนสุขสหคตํ นเหตุํ ธมฺมํ ปฏิจฺจ ฯเปฯ.\csromanspacer{} น-อุเปกฺขาสหคตํ นเหตุํ ธมฺมํ ปฏิจฺจ ฯเปฯ.\csromanspacer{} นน-อุเปกฺขาสหคตํ นเหตุํ ธมฺมํ ปฏิจฺจ.}

\proseitem{39}{นกามาวจรํ นเหตุํ ธมฺมํ ปฏิจฺจ ฯเปฯ.\csromanspacer{} นนกามาวจรํ นเหตุํ ธมฺมํ ปฏิจฺจ ฯเปฯ.\csromanspacer{} นรูปาวจรํ นเหตุํ ธมฺมํ ปฏิจฺจ ฯเปฯ.\csromanspacer{} นนรูปาวจรํ นเหตุํ ธมฺมํ ปฏิจฺจ.}

\proseitem{40}{น-อรูปาวจรํ นเหตุํ ธมฺมํ ปฏิจฺจ ฯเปฯ.\csromanspacer{} นน-อรูปาวจรํ นเหตุํ ธมฺมํ ปฏิจฺจ.}

\prose{นปริยาปนฺนํ นเหตุํ ธมฺมํ ปฏิจฺจ ฯเปฯ.\csromanspacer{} น-อปริยาปนฺนํ นเหตุํ ธมฺมํ ปฏิจฺจ.}

\proseitem{41}{นนิยฺยานิกํ นเหตุํ ธมฺมํ ปฏิจฺจ ฯเปฯ.\csromanspacer{} น-อนิยฺยานิกํ นเหตุํ ธมฺมํ ปฏิจฺจ ฯเปฯ.\csromanspacer{} นนิยตํ นเหตุํ ธมฺมํ ปฏิจฺจ ฯเปฯ.\csromanspacer{} น-อนิยตํ นเหตุํ ธมฺมํ ปฏิจฺจ.}

\prose{นส-อุตฺตรํ นเหตุํ ธมฺมํ ปฏิจฺจ ฯเปฯ.\csromanspacer{} น-อนุตฺตรํ นเหตุํ ธมฺมํ ปฏิจฺจ.}

\proseitem{42}{นสรณํ นเหตุํ ธมฺมํ ปฏิจฺจ นสรโณ นเหตุ ธมฺโม อุปฺปชฺชติ เหตุปจฺจยา. (1)}

\prose{น-อรณํ นเหตุํ ธมฺมํ ปฏิจฺจ น-อรโณ นเหตุ ธมฺโม อุปฺปชฺชติ เหตุปจฺจยา.\csromanspacer{}\csromanspacer{}น-อรณํ นเหตุํ ธมฺมํ ปฏิจฺจ นสรโณ นเหตุ ธมฺโม อุปฺปชฺชติ เหตุปจฺจยา.\csromanspacer{}\csromanspacer{}น-อรณํ นเหตุํ ธมฺมํ ปฏิจฺจ นสรโณ นเหตุ จ น-อรโณ นเหตุ จ ธมฺมา อุปฺปชฺชนฺติ เหตุปจฺจยา. (3)}

\prose{นสรณํ นเหตุญฺจ น-อรณํ นเหตุญฺจ ธมฺมํ ปฏิจฺจ นสรโณ นเหตุ ธมฺโม อุปฺปชฺชติ เหตุปจฺจยา. (1)}

\csromanheader{2}{เหตุยา ปญฺจ ฯเปฯ อวิคเต ปญฺจ. (สพฺพตฺถ วิตฺถาโร.)}
\csromanheader{2}{100. สรณทุก 7. จูฬนฺตรทุก}
\csromanheader{2}{43. นสรณํ นนเหตุํ ธมฺมํ ปฏิจฺจ นสรโณ นนเหตุ}
\prose{\csromanfolio{169}น-อรณํ นนเหตุํ ธมฺมํ ปฏิจฺจ น-อรโณ นนเหตุ ธมฺโม อุปฺปชฺชติ เหตุปจฺจยา. (1)}

\csromanheader{2}{เหตุยา เทฺว ฯเปฯ อวิคเต เทฺว. (สพฺพตฺถ วิตฺถาโร.)}
\csromanheader{2}{100. สรณทุก 2. สเหตุกทุกาทิ}
\proseitem{44}{นสรณํ นสเหตุกํ ธมฺมํ ปฏิจฺจ ฯเปฯ.\csromanspacer{} นสรณํ น-อเหตุกํ ธมฺมํ ปฏิจฺจ ฯเปฯ.\csromanspacer{} นสรณํ นเหตุสมฺปยุตฺตํ ธมฺมํ ปฏิจฺจ.}

\prose{นสรณํ นเหตุวิปฺปยุตฺตํ ธมฺมํ ปฏิจฺจ.}

\proseitem{45}{นสรณํ นเหตุญฺเจว น-อเหตุกญฺจ ธมฺมํ ปฏิจฺจ ฯเปฯ.\csromanspacer{} นสรณํ น-อเหตุกญฺเจว นน จ เหตุํ ธมฺมํ ปฏิจฺจ ฯเปฯ.\csromanspacer{} นสรณํ นเหตุญฺเจว นเหตุวิปฺปยุตฺตญฺจ ธมฺมํ ปฏิจฺจ ฯเปฯ.\csromanspacer{} นสรณํ นเหตุวิปฺปยุตฺตญฺเจว นน จ เหตุํ ธมฺมํ ปฏิจฺจ.}

\prose{นสรณํ นเหตุํ นสเหตุกํ ธมฺมํ ปฏิจฺจ.}

\prose{นสรณํ นเหตุํ น-อเหตุกํ ธมฺมํ ปฏิจฺจ.}

\csromanheader{2}{100. สรณทุก 7. จูฬนฺตรทุก}
\proseitem{46}{นสรณํ น-อปฺปจฺจยํ ธมฺมํ ปฏิจฺจ ฯเปฯ.\csromanspacer{} นสรณํ น-อสงฺขตํ ธมฺมํ ปฏิจฺจ.}

\prose{นสรณํ นสนิทสฺสนํ ธมฺมํ ปฏิจฺจ.}

\prose{นสรณํ นสปฺปฏิฆํ ธมฺมํ ปฏิจฺจ.}

\prose{นสรณํ น-อปฺปฏิฆํ ธมฺมํ ปฏิจฺจ.}

\prose{นสรณํ นรูปิํ ธมฺมํ ปฏิจฺจ.}

\prose{นสรณํ น-อรูปิํ ธมฺมํ ปฏิจฺจ.}

\prose{นสรณํ นโลกุตฺตรํ ธมฺมํ ปฏิจฺจ.}

\prose{นสรณํ นเกนจิ วิญฺเญยฺยํ ธมฺมํ ปฏิจฺจ ฯเปฯ.\csromanspacer{} นสรณํ นนเกนจิ วิญฺเญยฺยํ ธมฺมํ ปฏิจฺจ.}

\csromanheader{2}{ธมฺมปจฺจนีย ทุกทุกปฏฺฐานปาฬิ \textbf{100. สรณทุก 14. อาสวโคจฺฉก}}
\proseitem{47}{\csromanfolio{170}นสรณํ โน-อาสวํ ธมฺมํ ปฏิจฺจ.}

\prose{นสรณํ นโน-อาสวํ ธมฺมํ ปฏิจฺจ.}

\prose{นสรณํ นสาสวํ ธมฺมํ ปฏิจฺจ.}

\prose{นสรณํ น-อนาสวํ ธมฺมํ ปฏิจฺจ.}

\prose{นสรณํ น-อาสวสมฺปยุตฺตํ ธมฺมํ ปฏิจฺจ.}

\prose{นสรณํ น-อาสววิปฺปยุตฺตํ ธมฺมํ ปฏิจฺจ.}

\prose{นสรณํ น-อาสวญฺเจว น-อนาสวญฺจ ธมฺมํ ปฏิจฺจ.}

\prose{นสรณํ น-อนาสวญฺเจว นโน จ อาสวํ ธมฺมํ ปฏิจฺจ.}

\prose{นสรณํ น-อาสวญฺเจว น-อาสววิปฺปยุตฺตญฺจ ธมฺมํ ปฏิจฺจ ฯเปฯ.\csromanspacer{} นสรณํ น-อาสววิปฺปยุตฺตญฺเจว นโน จ อาสวํ ธมฺมํ ปฏิจฺจ.}

\prose{นสรณํ น-อาสววิปฺปยุตฺตํ นสาสวํ ธมฺมํ ปฏิจฺจ.}

\prose{นสรณํ น-อาสววิปฺปยุตฺตํ น-อนาสวํ ธมฺมํ ปฏิจฺจ.}

\csromanheader{2}{100. สรณทุก 20. สญฺโญชนทุกาทิ}
\proseitem{48}{นสรณํ โนสญฺโญชนํ ธมฺมํ ปฏิจฺจ.}

\proseitem{49}{นสรณํ โนคนฺถํ ธมฺมํ ปฏิจฺจ ฯเปฯ.\csromanspacer{} นสรณํ โน-โอฆํ ธมฺมํ ปฏิจฺจ ฯเปฯ.\csromanspacer{} นสรณํ โนโยคํ ธมฺมํ ปฏิจฺจ ฯเปฯ.\csromanspacer{} นสรณํ โนนีวรณํ ธมฺมํ ปฏิจฺจ ฯเปฯ.\csromanspacer{} นสรณํ โนปรามาสํ ธมฺมํ ปฏิจฺจ.}

\csromanheader{2}{100. สรณทุก 55. มหนฺตรทุก}
\csromanheader{2}{50. นสรณํ นสารมฺมณํ ธมฺมํ ปฏิจฺจ. (สํขิตฺตํ.)}
\prose{นสรณํ นจิตฺตํ ธมฺมํ ปฏิจฺจ.}

\prose{นสรณํ นเจตสิกํ ธมฺมํ ปฏิจฺจ.}

\prose{นสรณํ นจิตฺตสมฺปยุตฺตํ ธมฺมํ ปฏิจฺจ ฯเปฯ.\csromanspacer{} นสรณํ นจิตฺตสํสฏฺฐํ ธมฺมํ ปฏิจฺจ.}

\csromanheader{2}{100. สรณทุก 83. ทสฺสนทุกาทิ}
\prose{\csromanfolio{171}นสรณํ นจิตฺตสมุฏฺฐานํ ธมฺมํ ปฏิจฺจ.}

\prose{นสรณํ นจิตฺตสหภุํ ธมฺมํ ปฏิจฺจ ฯเปฯ.\csromanspacer{} นสรณํ นจิตฺตานุปริวตฺติํ ธมฺมํ ปฏิจฺจ.}

\proseitem{51}{นสรณํ นจิตฺตสํสฏฺฐสมุฏฺฐานํ ธมฺมํ ปฏิจฺจ ฯเปฯ.\csromanspacer{} นสรณํ นจิตฺตสํสฏฺฐสมุฏฺฐานสหภุํ ธมฺมํ ปฏิจฺจ ฯเปฯ.\csromanspacer{} นสรณํ นจิตฺตสํสฏฺฐสมุฏฺฐานานุปริวตฺติํ ธมฺมํ ปฏิจฺจ.}

\prose{นสรณํ น-อชฺฌตฺติกํ ธมฺมํ ปฏิจฺจ.}

\prose{นสรณํ นพาหิรํ ธมฺมํ ปฏิจฺจ.}

\prose{นสรณํ น-อุปาทา ธมฺมํ ปฏิจฺจ.}

\prose{นสรณํ น-อุปาทินฺนํ ธมฺมํ ปฏิจฺจ.}

\prose{นสรณํ น-อนุปาทินฺนํ ธมฺมํ ปฏิจฺจ.}

\csromanheader{2}{100. สรณทุก 69. อุปาทานาทิทุก}
\proseitem{52}{นสรณํ น-อุปาทานํ ธมฺมํ ปฏิจฺจ.}

\proseitem{53}{นสรณํ โนกิเลสํ ธมฺมํ ปฏิจฺจ.}

\csromanheader{2}{100. สรณทุก 83. ทสฺสนทุกาทิ}
\proseitem{54}{นสรณํ นทสฺสเนน ปหาตพฺพํ ธมฺมํ ปฏิจฺจ ฯเปฯ.\csromanspacer{} นสรณํ นนทสฺสเนน ปหาตพฺพํ ธมฺมํ ปฏิจฺจ ฯเปฯ.\csromanspacer{} นสรณํ นภาวนาย ปหาตพฺพํ ธมฺมํ ปฏิจฺจ ฯเปฯ.\csromanspacer{} นสรณํ นนภาวนาย ปหาตพฺพํ ธมฺมํ ปฏิจฺจ ฯเปฯ.\csromanspacer{} นสรณํ นทสฺสเนน ปหาตพฺพเหตุกํ ธมฺมํ ปฏิจฺจ ฯเปฯ.\csromanspacer{} นสรณํ นนทสฺสเนน ปหาตพฺพเหตุกํ ธมฺมํ ปฏิจฺจ ฯเปฯ.\csromanspacer{} นสรณํ นภาวนาย ปหาตพฺพเหตุกํ ธมฺมํ ปฏิจฺจ.}

\prose{นสรณํ นนภาวนาย ปหาตพฺพเหตุกํ ธมฺมํ ปฏิจฺจ.}

\proseitem{55}{นสรณํ นสวิตกฺกํ ธมฺมํ ปฏิจฺจ ฯเปฯ.\csromanspacer{} นสรณํ น-อวิตกฺกํ ธมฺมํ ปฏิจฺจ ฯเปฯ.\csromanspacer{} นสรณํ นสวิจารํ ธมฺมํ ปฏิจฺจ.}

\prose{นสรณํ น-อวิจารํ ธมฺมํ ปฏิจฺจ.}

\proseitem{56}{นสรณํ นสปฺปีติกํ ธมฺมํ ปฏิจฺจ ฯเปฯ.\csromanspacer{} นสรณํ น-อปฺปีติกํ ธมฺมํ ปฏิจฺจ ฯเปฯ.\csromanspacer{} นสรณํ นปีติสหคตํ ธมฺมํ ปฏิจฺจ ฯเปฯ.\csromanspacer{} นสรณํ นนปีติสหคตํ ธมฺมํ ปฏิจฺจ ฯเปฯ.\csromanspacer{} นสรณํ นสุขสหคตํ ธมฺมํ ปฏิจฺจ ฯเปฯ.\csromanspacer{} นสรณํ นนสุขสหคตํ ธมฺมํ ปฏิจฺจ ฯเปฯ.\csromanspacer{} นสรณํ น- อุเปกฺขาสหคตํ ธมฺมํ ปฏิจฺจ.}

\gambhira{ธมฺมปจฺจนีย ทุกทุกปฏฺฐานปาฬิ}
\proseitem{57}{\csromanfolio{172}นสรณํ นน-อุเปกฺขาสหคตํ ธมฺมํ ปฏิจฺจ.}

\prose{นสรณํ นกามาวจรํ ธมฺมํ ปฏิจฺจ.}

\prose{นสรณํ นนกามาวจรํ ธมฺมํ ปฏิจฺจ.}

\proseitem{58}{นสรณํ นรูปาวจรํ ธมฺมํ ปฏิจฺจ ฯเปฯ.\csromanspacer{} นสรณํ นนรูปาวจรํ ธมฺมํ ปฏิจฺจ ฯเปฯ.\csromanspacer{} นสรณํ น-อรูปาวจรํ ธมฺมํ ปฏิจฺจ.}

\prose{นสรณํ นน-อรูปาวจรํ ธมฺมํ ปฏิจฺจ.}

\prose{นสรณํ นปริยาปนฺนํ ธมฺมํ ปฏิจฺจ.}

\prose{นสรณํ น-อปริยาปนฺนํ ธมฺมํ ปฏิจฺจ.}

\proseitem{59}{นสรณํ นนิยฺยานิกํ ธมฺมํ ปฏิจฺจ.}

\prose{นสรณํ น-อนิยฺยานิกํ ธมฺมํ ปฏิจฺจ.}

\prose{นสรณํ นนิยตํ ธมฺมํ ปฏิจฺจ.}

\prose{นสรณํ น-อนิยตํ ธมฺมํ ปฏิจฺจ.}

\proseitem{60}{นสรณํ นส-อุตฺตรํ ธมฺมํ ปฏิจฺจ นสรโณ นส-อุตฺตโร ธมฺโม อุปฺปชฺชติ เหตุปจฺจยา.}

\csromanheader{2}{นสรณํ น-อนุตฺตรํ ธมฺมํ ปฏิจฺจ นสรโณ น-อนุตฺตโร}
\prose{น-อรณํ น-อนุตฺตรํ ธมฺมํ ปฏิจฺจ น-อรโณ น-อนุตฺตโร ธมฺโม อุปฺปชฺชติ เหตุปจฺจยา.\csromanspacer{}\csromanspacer{}น-อรณํ น-อนุตฺตรํ ธมฺมํ ปฏิจฺจ นสรโณ น-อนุตฺตโร ธมฺโม อุปฺปชฺชติ เหตุปจฺจยา.\csromanspacer{}\csromanspacer{}น-อรณํ น-อนุตฺตรํ ธมฺมํ ปฏิจฺจ นสรโณ น-อนุตฺตโร จ น-อรโณ น-อนุตฺตโร จ ธมฺมา อุปฺปชฺชนฺติ เหตุปจฺจยา. (3)}

\prose{นสรณํ น-อนุตฺตรญฺจ น-อรณํ น-อนุตฺตรญฺจ ธมฺมํ ปฏิจฺจ นสรโณ น-อนุตฺตโร ธมฺโม อุปฺปชฺชติ เหตุปจฺจยา. (1) (สํขิตฺตํ.)}

\prose{เหตุยา ปญฺจ. (สพฺพตฺถ วิตฺตาโร.\csromanspacer{} สหชาตวารมฺปิ ปจฺจยวารมฺปิ นิสฺสยวารมฺปิ สํสฏฺฐวารมฺปิ สมฺปยุตฺตวารมฺปิ ปญฺหาวารมฺปิ วิตฺถาเรตพฺพํ.)}

\nitthitam{ธมฺมปจฺจนีเย ทุกทุกปฏฺฐานํ นิฏฺฐิตํ.}
\nitthitam{\textbf{ปจฺจนียปฏฺฐานํ นิฏฺฐิตํ.}}
\csromanheader{2}{\textbf{อภิธมฺมปิฏิก}}
\csromanheader{2}{\textbf{ธมฺมานุโลมปจฺจนีย}}
\gambhira{\textbf{ติกปฏฺฐานปาฬิ}}
\namakkaram{นโม ตสฺส ภควโต อรหโต สมฺมาสมฺพุทฺธสฺส.}
\csromanheader{2}{1. กุสลตฺติก 1. ปฏิจฺจวาร}
\csromanheader{2}{\textbf{ปจฺจยจตุกฺก-เหตุ}}
\proseitem{1}{\csromanfolio{173}กุสลํ ธมฺมํ ปฏิจฺจ นกุสโล ธมฺโม อุปฺปชฺชติ เหตุปจฺจยา.\csromanspacer{} กุสเล ขนฺเธ ปฏิจฺจ จิตฺตสมุฏฺฐานํ รูปํ.\csromanspacer{}\csromanspacer{}กุสลํ ธมฺมํ ปฏิจฺจ น-อกุสโล ธมฺโม อุปฺปชฺชติ เหตุปจฺจยา.\csromanspacer{} กุสลํ เอกํ ขนฺธํ ปฏิจฺจ ตโย ขนฺธา จิตฺตสมุฏฺฐานญฺจ รูปํ ฯเปฯ.\csromanspacer{}\csromanspacer{}กุสลํ ธมฺมํ ปฏิจฺจ น-อพฺยากโต ธมฺโม อุปฺปชฺชติ เหตุปจฺจยา.\csromanspacer{} กุสลํ เอกํ ขนฺธํ ปฏิจฺจ ตโย ขนฺธา ฯเปฯ.\csromanspacer{}\csromanspacer{}กุสลํ ธมฺมํ ปฏิจฺจ น-อกุสโล จ น-อพฺยากโต จ ธมฺมา อุปฺปชฺชนฺติ เหตุปจฺจยา.\csromanspacer{} กุสลํ เอกํ ขนฺธํ ปฏิจฺจ ตโย ขนฺธา ฯเปฯ.\csromanspacer{}\csromanspacer{}กุสลํ ธมฺมํ ปฏิจฺจ นกุสโล จ น-อกุสโล จ ธมฺมา อุปฺปชฺชนฺติ เหตุปจฺจยา.\csromanspacer{} กุสเล ขนฺเธ ปฏิจฺจ จิตฺตสมุฏฺฐานํ รูปํ. (5)}

\proseitem{2}{อกุสลํ ธมฺมํ ปฏิจฺจ น-อกุสโล ธมฺโม อุปฺปชฺชติ เหตุปจฺจยา.\csromanspacer{} อกุสเล ขนฺเธ ปฏิจฺจ จิตฺตสมุฏฺฐานํ รูปํ.\csromanspacer{}\csromanspacer{}อกุสลํ ธมฺมํ ปฏิจฺจ นกุสโล ธมฺโม อุปฺปชฺชติ เหตุปจฺจยา.\csromanspacer{} อกุสลํ เอกํ ขนฺธํ ปฏิจฺจ ตโย ขนฺธา จิตฺตสมุฏฺฐานญฺจ รูปํ ฯเปฯ.\csromanspacer{}\csromanspacer{}อกุสลํ ธมฺมํ ปฏิจฺจ น-อพฺยากโต ธมฺโม อุปฺปชฺชติ เหตุปจฺจยา.\csromanspacer{}\csromanspacer{}อกุสลํ ธมฺมํ ปฏิจฺจ นกุสโล จ น-อพฺยากโต จ ธมฺมา อุปฺปชฺชนฺติ เหตุปจฺจยา.\csromanspacer{}\csromanspacer{}อกุสลํ ธมฺมํ ปฏิจฺจ นกุสโล จ น-อกุสโล จ ธมฺมา อุปฺปชฺชนฺติ เหตุปจฺจยา. (5)}

\gambhira{ธมฺมานุโลมปจฺจนีย ติกปฏฺฐานปาฬิ}
\proseitem{3}{\csromanfolio{174}อพฺยากตํ ธมฺมํ ปฏิจฺจ นกุสโล ธมฺโม อุปฺปชฺชติ เหตุปจฺจยา.\csromanspacer{} วิปากาพฺยากตํ กิริยาพฺยากตํ เอกํ ขนฺธํ ปฏิจฺจ ตโย ขนฺธา จิตฺตสมุฏฺฐานญฺจ รูปํ ฯเปฯ.\csromanspacer{} ปฏิสนฺธิกฺขเณ ฯเปฯ เอกํ มหาภูตํ ปฏิจฺจ ตโย มหาภูตา ฯเปฯ เทฺว มหาภูเต ปฏิจฺจ เทฺว มหาภูตา, มหาภูเต ปฏิจฺจ จิตฺตสมุฏฺฐานํ รูปํ.\csromanspacer{}\csromanspacer{}อพฺยากตํ ธมฺมํ ปฏิจฺจ น-อกุสโล ธมฺโม อุปฺปชฺชติ เหตุปจฺจยา.\csromanspacer{} วิปากาพฺยากตํ ฯเปฯ.\csromanspacer{} ปฏิสนฺธิกฺขเณ ฯเปฯ มหาภูตํ ฯเปฯ.\csromanspacer{}\csromanspacer{}อพฺยากตํ ธมฺมํ ปฏิจฺจ นกุสโล จ น-อกุสโล จ ธมฺมา อุปฺปชฺชนฺติ เหตุปจฺจยา.\csromanspacer{} วิปากาพฺยากตํ กิริยาพฺยากตํ ฯเปฯ. (3) (สํขิตฺตํ.)}

\proseitem{4}{กุสลญฺจ อพฺยากตญฺจ ธมฺมํ ปฏิจฺจ นกุสโล ธมฺโม อุปฺปชฺชติ เหตุปจฺจยา.\csromanspacer{} กุสเล ขนฺเธ จ มหาภูเต จ ปฏิจฺจ จิตฺตสมุฏฺฐานํ รูปํ.\csromanspacer{}\csromanspacer{}กุสลญฺจ อพฺยากตญฺจ ธมฺมํ ปฏิจฺจ น-อกุสโล ธมฺโม อุปฺปชฺชติ เหตุปจฺจยา.\csromanspacer{}\csromanspacer{}กุสลญฺจ อพฺยากตญฺจ ธมฺมํ ปฏิจฺจ นกุสโล จ น-อกุสโล จ ธมฺมา อุปฺปชฺชนฺติ เหตุปจฺจยา. (3)}

\prose{อกุสลญฺจ อพฺยากตญฺจ ธมฺมํ ปฏิจฺจ นกุสโล ธมฺโม อุปฺปชฺชติ เหตุปจฺจยา.\csromanspacer{}\csromanspacer{}อกุสลญฺจ อพฺยากตญฺจ ธมฺมํ ปฏิจฺจ น-อกุสโล ธมฺโม อุปฺปชฺชติ เหตุปจฺจยา.\csromanspacer{}\csromanspacer{}อกุสลญฺจ อพฺยากตญฺจ ธมฺมํ ปฏิจฺจ นกุสโล จ น-อกุสโล จ ธมฺมา อุปฺปชฺชนฺติ เหตุปจฺจยา. (3) (จิตฺตสมุฏฺฐานรูปเมว เอตฺถ วตฺตติ.\csromanspacer{} เอกูนวีสติ ปญฺหา กาตพฺพา.)}

\csromanheader{2}{\textbf{อารมฺมณ}}
\proseitem{5}{กุสลํ ธมฺมํ ปฏิจฺจ น-อกุสโล ธมฺโม อุปฺปชฺชติ อารมฺมณปจฺจยา.\csromanspacer{}\csromanspacer{}กุสลํ ธมฺมํ ปฏิจฺจ น-อพฺยากโต ธมฺโม อุปฺปชฺชติ อารมฺมณปจฺจยา.\csromanspacer{}\csromanspacer{}กุสลํ ธมฺมํ ปฏิจฺจ น-อกุสโล จ น-อพฺยากโต จ ธมฺมา อุปฺปชฺชนฺติ อารมฺมณปจฺจยา. (3)}

\prose{อกุสลํ ธมฺมํ ปฏิจฺจ นกุสโล ธมฺโม อุปฺปชฺชติ อารมฺมณปจฺจยา.\csromanspacer{}\csromanspacer{}อกุสลํ ธมฺมํ ปฏิจฺจ น-อพฺยากโต ธมฺโม อุปฺปชฺชติ อารมฺมณปจฺจยา.\csromanspacer{}\csromanspacer{}อกุสลํ ธมฺมํ ปฏิจฺจ นกุสโล จ น-อพฺยากโต จ ธมฺมา อุปฺปชฺชนฺติ อารมฺมณปจฺจยา. (3)}

\prose{อพฺยากตํ ธมฺมํ ปฏิจฺจ นกุสโล ธมฺโม อุปฺปชฺชติ อารมฺมณปจฺจยา.\csromanspacer{}\csromanspacer{}อพฺยากตํ ธมฺมํ ปฏิจฺจ น-อกุสโล ธมฺโม อุปฺปชฺชติ อารมฺมณปจฺจยา.\csromanspacer{}\csromanspacer{}}

\vspace{\tipitakaparskip}
\csromancenter{กุสลตฺติก อพฺยากตํ ธมฺมํ ปฏิจฺจ นกุสโล จ น-อกุสโล จ ธมฺมา อุปฺปชฺชนฺติ อารมฺมณปจฺจยา. (3) (สํขิตฺตํ.)}
\prose{\csromanfolio{175}เหตุยา เอกูนวีส, อารมฺมเณ นว, อธิปติยา เอกูนวีส, อนนฺตเร นว, สมนนฺตเร นว, สหชาเต เอกูนวีส ฯเปฯ อวิคเต เอกูนวีส.}

\csromanheader{2}{\textbf{ปจฺจนีย-นเหตุน-อารมฺมณ}}
\proseitem{6}{อกุสลํ ธมฺมํ ปฏิจฺจ นกุสโล ธมฺโม อุปฺปชฺชติ นเหตุปจฺจยา.\csromanspacer{}\csromanspacer{}อกุสลํ ธมฺมํ ปฏิจฺจ น-อพฺยากโต ธมฺโม อุปฺปชฺชติ นเหตุปจฺจยา.\csromanspacer{}\csromanspacer{}อกุสลํ ธมฺมํ ปฏิจฺจ นกุสโล จ น-อพฺยากโต จ ธมฺมา อุปฺปชฺชนฺติ นเหตุปจฺจยา. (3) อพฺยากตํ ธมฺมํ ปฏิจฺจ นกุสโล ธมฺโม อุปฺปชฺชติ นเหตุปจฺจยา.\csromanspacer{}\csromanspacer{}อพฺยากตํ ธมฺมํ ปฏิจฺจ น-อกุสโล ธมฺโม อุปฺปชฺชติ นเหตุปจฺจยา.\csromanspacer{}\csromanspacer{}อพฺยากตํ ธมฺมํ ปฏิจฺจ นกุสโล จ น-อกุสโล จ ธมฺมา อุปฺปชฺชนฺติ นเหตุปจฺจยา. (3)}

\proseitem{7}{กุสลํ ธมฺมํ ปฏิจฺจ นกุสโล ธมฺโม อุปฺปชฺชติ น-อารมฺมณปจฺจยา.\csromanspacer{}\csromanspacer{}กุสลํ ธมฺมํ ปฏิจฺจ น-อกุสโล ธมฺโม อุปฺปชฺชติ น-อารมฺมณปจฺจยา.\csromanspacer{}\csromanspacer{}กุสลํ ธมฺมํ ปฏิจฺจ นกุสโล จ น-อกุสโล จ ธมฺมา อุปฺปชฺชนฺติ น-อารมฺมณปจฺจยา. (3) อกุสลํ ธมฺมํ ปฏิจฺจ น-อกุสโล ธมฺโม อุปฺปชฺชติ น-อารมฺมณปจฺจยา.\csromanspacer{}\csromanspacer{}อกุสลํ ธมฺมํ ปฏิจฺจ นกุสโล ธมฺโม อุปฺปชฺชติ น-อารมฺมณปจฺจยา.\csromanspacer{}\csromanspacer{}อกุสลํ ธมฺมํ ปฏิจฺจ นกุสโล จ น-อกุสโล จ ธมฺมา อุปฺปชฺชนฺติ น-อารมฺมณปจฺจยา. (3) (สํขิตฺตํ.)}

\prose{นเหตุยา ฉ, น-อารมฺมเณ ปนฺนรส, น-อธิปติยา เอกูนวีส ฯเปฯ โนวิคเต ปนฺนรส.}

\prose{(ปจฺจนียํ วิตฺถาเรตพฺพํ.\csromanspacer{} สหชาตวารมฺปิ ปจฺจยวารมฺปิ วิตฺถาเรตพฺพํ.\csromanspacer{} ปจฺจยวาเรปิ เหตุยา ฉพฺพีส, อารมฺมเณ อฏฺฐารส ฯเปฯ อวิคเต ฉพฺพีส.\csromanspacer{} นิสฺสยวารมฺปิ สํสฏฺฐวารมฺปิ สมฺปยุตฺตวารมฺปิ ปฏิจฺจวารสทิสํ.)}

\csromanheader{2}{ธมฺมานุโลมปจฺจนีย ติกปฏฺฐานปาฬิ \textbf{1. กุสลตฺติก 7. ปญฺหาวาร}}
\csromanheader{2}{\textbf{ปจฺจยจตุกฺก-เหตุ-อารมฺมณ}}
\proseitem{8}{\csromanfolio{176}กุสโล ธมฺโม นกุสลสฺส ธมฺมสฺส เหตุปจฺจเยน ปจฺจโย.\csromanspacer{}\csromanspacer{}กุสโล ธมฺโม น-อกุสลสฺส ธมฺมสฺส เหตุปจฺจเยน ปจฺจโย.\csromanspacer{}\csromanspacer{}กุสโล ธมฺโม น-อพฺยากตสฺส ธมฺมสฺส เหตุปจฺจเยน ปจฺจโย.\csromanspacer{}\csromanspacer{}กุสโล ธมฺโม น-อกุสลสฺส จ น-อพฺยากตสฺส จ ธมฺมสฺส เหตุปจฺจเยน ปจฺจโย.\csromanspacer{}\csromanspacer{}กุสโล ธมฺโม นกุสลสฺส จ น-อกุสลสฺส จ ธมฺมสฺส เหตุปจฺจเยน ปจฺจโย. (5)}

\proseitem{9}{อกุสโล ธมฺโม น-อกุสลสฺส ธมฺมสฺส เหตุปจฺจเยน ปจฺจโย.\csromanspacer{}\csromanspacer{}อกุสโล ธมฺโม นกุสลสฺส ธมฺมสฺส เหตุปจฺจเยน ปจฺจโย.\csromanspacer{}\csromanspacer{}อกุสโล ธมฺโม น-อพฺยากตสฺส ธมฺมสฺส เหตุปจฺจเยน ปจฺจโย.\csromanspacer{}\csromanspacer{}อกุสโล ธมฺโม นกุสลสฺส จ น-อพฺยากตสฺส จ ธมฺมสฺส เหตุปจฺจเยน ปจฺจโย.\csromanspacer{}\csromanspacer{}อกุสโล ธมฺโม นกุสลสฺส จ น-อกุสลสฺส จ ธมฺมสฺส เหตุปจฺจเยน ปจฺจโย. (5)}

\prose{อพฺยากโต ธมฺโม นกุสลสฺส ธมฺมสฺส เหตุปจฺจเยน ปจฺจโย.\csromanspacer{}\csromanspacer{}อพฺยากโต ธมฺโม น-อกุสลสฺส ธมฺมสฺส เหตุปจฺจเยน ปจฺจโย.\csromanspacer{}\csromanspacer{}อพฺยากโต ธมฺโม นกุสลสฺส จ น-อกุสลสฺส จ ธมฺมสฺส เหตุปจฺจเยน ปจฺจโย. (3)}

\proseitem{10}{กุสโล ธมฺโม นกุสลสฺส ธมฺมสฺส อารมฺมณปจฺจเยน ปจฺจโย. (ฉ ปญฺหา.)}

\prose{อกุสโล ธมฺโม น-อกุสลสฺส ธมฺมสฺส อารมฺมณปจฺจเยน ปจฺจโย. (ฉ ปญฺหา.)}

\prose{อพฺยากโต ธมฺโม น-อพฺยากตสฺส ธมฺมสฺส อารมฺมณปจฺจเยน ปจฺจโย. (ฉ ปญฺหา.\csromanspacer{} สํขิตฺตํ.)}

\proseitem{11}{เหตุยา เตรส, อารมฺมเณ อฏฺฐารส, อธิปติยา สตฺตรส, อนนฺตเร โสฬส, สมนนฺตเร โสฬส, สหชาเต เอกูนวีส, อญฺญมญฺเญ นว, นิสฺสเย ฉพฺพีส, อุปนิสฺสเย อฏฺฐารส, ปุเรชาเต ฉ, ปจฺฉาชาเต นว, อาเสวเน นว, กมฺเม เตรส, วิปาเก ตีณิ, อาหาเร เตรส ฯเปฯ มคฺเค เตรส, สมฺปยุตฺเต นว, วิปฺปยุตฺเต ทฺวาทส ฯเปฯ อวิคเต ฉพฺพีส. (ปญฺหาวารํ วิตฺถาเรตพฺพํ.)}

\csromanheader{2}{2. เวทนาตฺติก \textbf{2. เวทนาตฺติก 1. ปฏิจฺจวาร}}
\csromanheader{2}{\textbf{ปจฺจยจตุกฺก-เหตุ}}
\proseitem{12}{\csromanfolio{177}สุขาย เวทนาย สมฺปยุตฺตํ ธมฺมํ ปฏิจฺจ นสุขาย เวทนาย สมฺปยุตฺโต ธมฺโม อุปฺปชฺชติ เหตุปจฺจยา.\csromanspacer{} สุขาย เวทนาย สมฺปยุตฺเต ขนฺเธ ปฏิจฺจ สุขเวทนา จิตฺตสมุฏฺฐานญฺจ รูปํ.\csromanspacer{} ปฏิสนฺธิกฺขเณ ฯเปฯ. (มหาภูตา นตฺถิ.) .\csromanspacer{} สุขาย เวทนาย สมฺปยุตฺตํ ธมฺมํ ปฏิจฺจ นทุกฺขาย เวทนาย สมฺปยุตฺโต ธมฺโม อุปฺปชฺชติ เหตุปจฺจยา.\csromanspacer{} สุขาย เวทนาย สมฺปยุตฺตํ เอกํ ขนฺธํ ปฏิจฺจ ตโย ขนฺธา จิตฺตสมุฏฺฐานญฺจ รูปํ ฯเปฯ.\csromanspacer{} สุขาย เวทนาย สมฺปยุตฺตํ ธมฺมํ ปฏิจฺจ น-อทุกฺขมสุขาย เวทนาย สมฺปยุตฺโต ธมฺโม อุปฺปชฺชติ เหตุปจฺจยา.\csromanspacer{}\csromanspacer{}สุขาย เวทนาย สมฺปยุตฺตํ ธมฺมํ ปฏิจฺจ นสุขาย เวทนาย สมฺปยุตฺโต จ น-อทุกฺขมสุขาย เวทนาย สมฺปยุตฺโต จ ธมฺมา อุปฺปชฺชนฺติ เหตุปจฺจยา.\csromanspacer{}\csromanspacer{}สุขาย เวทนาย สมฺปยุตฺตํ ธมฺมํ ปฏิจฺจ นทุกฺขาย เวทนาย สมฺปยุตฺโต จ น-อทุกฺขมสุขาย เวทนาย สมฺปยุตฺโต จ ธมฺมา อุปฺปชฺชนฺติ เหตุปจฺจยา.\csromanspacer{}\csromanspacer{}สุขาย เวทนาย สมฺปยุตฺตํ ธมฺมํ ปฏิจฺจ นสุขาย เวทนาย สมฺปยุตฺโต จ นทุกฺขาย เวทนาย สมฺปยุตฺโต จ ธมฺมา อุปฺปชฺชนฺติ เหตุปจฺจยา.\csromanspacer{}\csromanspacer{}สุขาย เวทนาย สมฺปยุตฺตํ ธมฺมํ ปฏิจฺจ นสุขาย เวทนาย สมฺปยุตฺโต จ นทุกฺขาย เวทนาย สมฺปยุตฺโต จ น-อทุกฺขมสุขาย เวทนาย สมฺปยุตฺโต จ ธมฺมา อุปฺปชฺชนฺติ เหตุปจฺจยา. (7)}

\proseitem{13}{ทุกฺขาย เวทนาย สมฺปยุตฺตํ ธมฺมํ ปฏิจฺจ นทุกฺขาย เวทนาย สมฺปยุตฺโต ธมฺโม อุปฺปชฺชติ เหตุปจฺจยา.\csromanspacer{}\csromanspacer{}ทุกฺขาย เวทนาย สมฺปยุตฺตํ ธมฺมํ ปฏิจฺจ นสุขาย เวทนาย สมฺปยุตฺโต ธมฺโม อุปฺปชฺชติ เหตุปจฺจยา.\csromanspacer{}\csromanspacer{}ทุกฺขาย เวทนาย สมฺปยุตฺตํ ธมฺมํ ปฏิจฺจ น-อทุกฺขมสุขาย เวทนาย สมฺปยุตฺโต ธมฺโม อุปฺปชฺชติ เหตุปจฺจยา.\csromanspacer{}\csromanspacer{}ทุกฺขาย เวทนาย สมฺปยุตฺตํ ธมฺมํ ปฏิจฺจ นสุขาย เวทนาย สมฺปยุตฺโต จ น-อทุกฺขมสุขาย เวทนาย สมฺปยุตฺโต จ ธมฺมา อุปฺปชฺชนฺติ เหตุปจฺจยา.\csromanspacer{}\csromanspacer{}ทุกฺขาย เวทนาย สมฺปยุตฺตํ ธมฺมํ ปฏิจฺจ นทุกฺขาย เวทนาย สมฺปยุตฺโต จ น-อทุกฺขมสุขาย เวทนาย สมฺปยุตฺโต จ ธมฺมา อุปฺปชฺชนฺติ เหตุปจฺจยา.\csromanspacer{}\csromanspacer{}ทุกฺขาย เวทนาย สมฺปยุตฺตํ ธมฺมํ ปฏิจฺจ นสุขาย เวทนาย สมฺปยุตฺโต จ นทุกฺขาย เวทนาย สมฺปยุตฺโต จ ธมฺมา อุปฺปชฺชนฺติ เหตุปจฺจยา.\csromanspacer{}\csromanspacer{}ทุกฺขาย เวทนาย สมฺปยุตฺตํ ธมฺมํ ปฏิจฺจ นสุขาย เวทนาย สมฺปยุตฺโต จ นทุกฺขาย เวทนาย}

\vspace{\tipitakaparskip}
\csromancenter{ธมฺมานุโลมปจฺจนีย ติกปฏฺฐานปาฬิ สมฺปยุตฺโต จ น-อทุกฺขมสุขาย เวทนาย สมฺปยุตฺโต จ ธมฺมา อุปฺปชฺชนฺติ เหตุปจฺจยา. (7)}
\proseitem{14}{\csromanfolio{178}อทุกฺขมสุขาย เวทนาย สมฺปยุตฺตํ ธมฺมํ ปฏิจฺจ น-อทุกฺขมสุขาย เวทนาย สมฺปยุตฺโต ธมฺโม อุปฺปชฺชติ เหตุปจฺจยา.\csromanspacer{}\csromanspacer{}อทุกฺขมสุขาย เวทนาย สมฺปยุตฺตํ ธมฺมํ ปฏิจฺจ นสุขาย เวทนาย สมฺปยุตฺโต ธมฺโม อุปฺปชฺชติ เหตุปจฺจยา.\csromanspacer{}\csromanspacer{}อทุกฺขมสุขาย เวทนาย สมฺปยุตฺตํ ธมฺมํ ปฏิจฺจ นทุกฺขาย เวทนาย สมฺปยุตฺโต ธมฺโม อุปฺปชฺชติ เหตุปจฺจยา.\csromanspacer{}\csromanspacer{}อทุกฺขมสุขาย เวทนาย สมฺปยุตฺตํ ธมฺมํ ปฏิจฺจ นสุขาย เวทนาย สมฺปยุตฺโต จ น-อทุกฺขมสุขาย เวทนาย สมฺปยุตฺโต จ ธมฺมา อุปฺปชฺชนฺติ เหตุปจฺจยา.\csromanspacer{}\csromanspacer{}อทุกฺขมสุขาย เวทนาย สมฺปยุตฺตํ ธมฺมํ ปฏิจฺจ นทุกฺขาย เวทนาย สมฺปยุตฺโต จ น-อทุกฺขมสุขาย เวทนาย สมฺปยุตฺโต จ ธมฺมา อุปฺปชฺชนฺติ เหตุปจฺจยา.\csromanspacer{}\csromanspacer{}อทุกฺขมสุขาย เวทนาย สมฺปยุตฺตํ ธมฺมํ ปฏิจฺจ นสุขาย เวทนาย สมฺปยุตฺโต จ นทุกฺขาย เวทนาย สมฺปยุตฺโต จ ธมฺมา อุปฺปชฺชนฺติ เหตุปจฺจยา.\csromanspacer{}\csromanspacer{}อทุกฺขมสุขาย เวทนาย สมฺปยุตฺตํ ธมฺมํ ปฏิจฺจ นสุขาย เวทนาย สมฺปยุตฺโต จ นทุกฺขาย เวทนาย สมฺปยุตฺโต จ น-อทุกฺขมสุขาย เวทนาย สมฺปยุตฺโต จ ธมฺมา อุปฺปชฺชนฺติ เหตุปจฺจยา. (7) (สํขิตฺตํ.)}

\prose{เหตุยา เอกวีส, อารมฺมเณ เอกวีส ฯเปฯ อวิคเต เอกวีส.}

\csromanheader{2}{\textbf{ปจฺจนีย-นเหตุ}}
\proseitem{15}{สุขาย เวทนาย สมฺปยุตฺตํ ธมฺมํ ปฏิจฺจ นสุขาย เวทนาย สมฺปยุตฺโต ธมฺโม อุปฺปชฺชติ นเหตุปจฺจยา. (สํขิตฺตํ.)}

\prose{นเหตุยา เอกวีส, น-อารมฺมเณ เอกวีส ฯเปฯ นวิปฺปยุตฺโต จุทฺทส ฯเปฯ โนวิคเต เอกวีส.}

\prose{(สหชาตวารมฺปิ ปจฺจยวารมฺปิ นิสฺสยวารมฺปิ สํสฏฺฐวารมฺปิ สมฺปยุตฺตวารมฺปิ ปฏิจฺจวารสทิสํ.)}

\csromanheader{2}{2. เวทนาตฺติก 7. ปญฺหาวาร}
\csromanheader{2}{\textbf{ปจฺจยจตุกฺก-เหตุ}}
\proseitem{16}{สุขาย เวทนาย สมฺปยุตฺโต ธมฺโม นสุขาย เวทนาย สมฺปยุตฺตสฺส ธมฺมสฺส เหตุปจฺจเยน ปจฺจโย.\csromanspacer{}\csromanspacer{}สุขาย เวทนาย}

\vspace{\tipitakaparskip}
\csromancenter{วิปากตฺติก สมฺปยุตฺโต ธมฺโม นทุกฺขาย เวทนาย สมฺปยุตฺตสฺส ธมฺมสฺส เหตุปจฺจเยน ปจฺจโย.\csromanspacer{}\csromanspacer{}สุขาย เวทนาย สมฺปยุตฺโต ธมฺโม น-อทุกฺขมสุขาย เวทนาย สมฺปยุตฺตสฺส ธมฺมสฺส เหตุปจฺจเยน ปจฺจโย.\csromanspacer{}\csromanspacer{}สุขาย เวทนาย สมฺปยุตฺโต ธมฺโม นสุขาย เวทนาย สมฺปยุตฺตสฺส จ น-อทุกฺขมสุขาย เวทนาย สมฺปยุตฺตสฺส จ ธมฺมสฺส เหตุปจฺจเยน ปจฺจโย.\csromanspacer{}\csromanspacer{}สุขาย เวทนาย สมฺปยุตฺโต ธมฺโม นทุกฺขาย เวทนาย สมฺปยุตฺตสฺส จ น-อทุกฺขมสุขาย เวทนาย สมฺปยุตฺตสฺส จ ธมฺมสฺส เหตุปจฺจเยน ปจฺจโย.\csromanspacer{}\csromanspacer{}สุขาย เวทนาย สมฺปยุตฺโต ธมฺโม นสุขาย เวทนาย สมฺปยุตฺตสฺส จ นทุกฺขาย เวทนาย สมฺปยุตฺตสฺส จ ธมฺมสฺส เหตุปจฺจเยน ปจฺจโย.\csromanspacer{}\csromanspacer{}สุขาย เวทนาย สมฺปยุตฺโต ธมฺโม นสุขาย เวทนาย สมฺปยุตฺตสฺส จ นทุกฺขาย เวทนาย สมฺปยุตฺตสฺส จ น-อทุกฺขมสุขาย เวทนาย สมฺปยุตฺตสฺส จ ธมฺมสฺส เหตุปจฺจเยน ปจฺจโย. (7) (สํขิตฺตํ.)}
\prose{\csromanfolio{179}เหตุยา เอกวีส, อารมฺมเณ เอกวีส ฯเปฯ อวิคเต เอกวีส. (ปญฺหาวารํ วิตฺถาเรตพฺพํ.)}

\csromanheader{2}{3. วิปากตฺติก 1. ปฏิจฺจวาร}
\csromanheader{2}{\textbf{ปจฺจยจตุกฺก-เหตุ-อารมฺมณ}}
\proseitem{17}{วิปากํ ธมฺมํ ปฏิจฺจ นวิปาโก ธมฺโม อุปฺปชฺชติ เหตุปจฺจยา.\csromanspacer{}\csromanspacer{}วิปากํ ธมฺมํ ปฏิจฺจ นวิปากธมฺมธมฺโม อุปฺปชฺชติ เหตุปจฺจยา.\csromanspacer{}\csromanspacer{}วิปากํ ธมฺมํ ปฏิจฺจ นเนววิปากนวิปากธมฺมธมฺโม อุปฺปชฺชติ เหตุปจฺจยา.\csromanspacer{}\csromanspacer{}วิปากํ ธมฺมํ ปฏิจฺจ นวิปากธมฺมธมฺโม จ นเนววิปากนวิปากธมฺมธมฺโม จ ธมฺมา อุปฺปชฺชนฺติ เหตุปจฺจยา.\csromanspacer{}\csromanspacer{}วิปากํ ธมฺมํ ปฏิจฺจ นวิปาโก จ นวิปากธมฺมธมฺโม จ ธมฺมา อุปฺปชฺชนฺติ เหตุปจฺจยา. (ปญฺจ ปญฺหา.)}

\proseitem{18}{วิปากธมฺมธมฺมํ ปฏิจฺจ นวิปากธมฺมธมฺโม อุปฺปชฺชติ เหตุปจฺจยา.\csromanspacer{}\csromanspacer{}วิปากธมฺมธมฺมํ ปฏิจฺจ นวิปาโก ธมฺโม อุปฺปชฺชติ เหตุปจฺจยา.\csromanspacer{}\csromanspacer{}วิปากธมฺมธมฺมํ ปฏิจฺจ นเนววิปากนวิปากธมฺมธมฺโม อุปฺปชฺชติ เหตุปจฺจยา.\csromanspacer{}\csromanspacer{}วิปากธมฺมธมฺมํ ปฏิจฺจ นวิปาโก จ นเนววิปากนวิปากธมฺมธมฺโม จ ธมฺมา อุปฺปชฺชนฺติ เหตุปจฺจยา.\csromanspacer{}\csromanspacer{}วิปากธมฺมธมฺมํ ปฏิจฺจ นวิปาโก จ นวิปากธมฺมธมฺโม จ ธมฺมา อุปฺปชฺชนฺติ เหตุปจฺจยา. (ปญฺจ ปญฺหา.)}

\gambhira{ธมฺมานุโลมปจฺจนีย ติกปฏฺฐานปาฬิ}
\proseitem{19}{\csromanfolio{180}เนววิปากนวิปากธมฺมธมฺมํ ปฏิจฺจ นเนววิปากนวิปากธมฺมธมฺโม อุปฺปชฺชติ เหตุปจฺจยา.\csromanspacer{}\csromanspacer{}เนววิปากนวิปากธมฺมธมฺมํ ปฏิจฺจ นวิปาโก ธมฺโม อุปฺปชฺชติ เหตุปจฺจยา.\csromanspacer{}\csromanspacer{}เนววิปากนวิปากธมฺมธมฺมํ ปฏิจฺจ นวิปากธมฺมธมฺโม อุปฺปชฺชติ เหตุปจฺจยา.\csromanspacer{}\csromanspacer{}เนววิปากนวิปากธมฺมธมฺมํ ปฏิจฺจ นวิปากธมฺมธมฺโม จ นเนววิปากนวิปากธมฺมธมฺโม จ ธมฺมา อุปฺปชฺชนฺติ เหตุปจฺจยา.\csromanspacer{}\csromanspacer{}เนววิปากนวิปากธมฺมธมฺมํ ปฏิจฺจ นวิปาโก จ นวิปากธมฺมธมฺโม จ ธมฺมา อุปฺปชฺชนฺติ เหตุปจฺจยา. (ปญฺจ ปญฺหา.)}

\proseitem{20}{วิปากญฺจ เนววิปากนวิปากธมฺมธมฺมญฺจ ธมฺมํ ปฏิจฺจ นวิปาโก ธมฺโม อุปฺปชฺชติ เหตุปจฺจยา.\csromanspacer{}\csromanspacer{}วิปากญฺจ เนววิปากนวิปากธมฺมธมฺมญฺจ ธมฺมํ ปฏิจฺจ นวิปากธมฺมธมฺโม อุปฺปชฺชติ เหตุปจฺจยา.\csromanspacer{}\csromanspacer{}วิปากญฺจ เนววิปากนวิปากธมฺมธมฺมญฺจ ธมฺมํ ปฏิจฺจ นเนววิปากนวิปากธมฺมธมฺโม อุปฺปชฺชติ เหตุปจฺจยา.\csromanspacer{}\csromanspacer{}วิปากญฺจ เนววิปากนวิปากธมฺมธมฺมญฺจ ธมฺมํ ปฏิจฺจ นวิปากธมฺมธมฺโม จ นเนววิปากนวิปากธมฺมธมฺโม จ ธมฺมา อุปฺปชฺชนฺติ เหตุปจฺจยา.\csromanspacer{}\csromanspacer{}วิปากญฺจ เนววิปากนวิปากธมฺมธมฺมญฺจ ธมฺมํ ปฏิจฺจ นวิปาโก จ นวิปากธมฺมธมฺโม จ ธมฺมา อุปฺปชฺชนฺติ เหตุปจฺจยา. (5)}

\prose{วิปากธมฺมธมฺมญฺจ เนววิปากนวิปากธมฺมธมฺมญฺจ ธมฺมํ ปฏิจฺจ นวิปาโก ธมฺโม อุปฺปชฺชติ เหตุปจฺจยา.\csromanspacer{}\csromanspacer{}วิปากธมฺมธมฺมญฺจ เนววิปากนวิปากธมฺมธมฺมญฺจ ธมฺมํ ปฏิจฺจ นวิปากธมฺมธมฺโม อุปฺปชฺชติ เหตุปจฺจยา.\csromanspacer{}\csromanspacer{}วิปากธมฺมธมฺมญฺจ เนววิปากนวิปากธมฺมธมฺมญฺจ ธมฺมํ ปฏิจฺจ นวิปาโก จ นวิปากธมฺมธมฺโม จ ธมฺมา อุปฺปชฺชนฺติ เหตุปจฺจยา. (3)}

\proseitem{21}{วิปากํ ธมฺมํ ปฏิจฺจ นวิปากธมฺมธมฺโม อุปฺปชฺชติ อารมฺมณปจฺจยา.\csromanspacer{} ตีณิ.}

\prose{วิปากธมฺมธมฺมํ ปฏิจฺจ นวิปาโก ธมฺโม อุปฺปชฺชติ อารมฺมณปจฺจยา.\csromanspacer{} ตีณิ.}

\prose{เนววิปากนวิปากธมฺมธมฺมํ ปฏิจฺจ นเนววิปากนวิปากธมฺมธมฺโม อุปฺปชฺชติ อารมฺมณปจฺจยา. (ปญฺจ ปญฺหา.)}

\prose{วิปากญฺจ เนววิปากนวิปากธมฺมธมฺมญฺจ ธมฺมํ ปฏิจฺจ นวิปากธมฺมธมฺโม อุปฺปชฺชติ อารมฺมณปจฺจยา.\csromanspacer{} ตีณิ. (สํขิตฺตํ.)}

\proseitem{22}{เหตุยา เตวีส, อารมฺมเณ จุทฺทส ฯเปฯ อวิคเต เตวีส. (สํขิตฺตํ.)}

\prose{นเหตุยา อฏฺฐารส, น-อารมฺมเณ ปนฺนรส.}

\csromanheader{2}{3. วิปากตฺติก}
\prose{\csromanfolio{181}(สหชาตวารมฺปิ ปจฺจยวารมฺปิ นิสฺสยวารมฺปิ สํสฏฺฐวารมฺปิ สมฺปยุตฺตวารมฺปิ วิตฺถาเรตพฺพํ.)}

\csromanheader{2}{3. วิปากตฺติก 7. ปญฺหาวาร}
\csromanheader{2}{\textbf{ปจฺจยจตุกฺก-เหตุ-อารมฺมณ}}
\proseitem{23}{วิปาโก ธมฺโม นวิปากสฺส ธมฺมสฺส เหตุปจฺจเยน ปจฺจโย.\csromanspacer{}\csromanspacer{}วิปาโก ธมฺโม นวิปากธมฺมธมฺมสฺส เหตุปจฺจเยน ปจฺจโย.\csromanspacer{}\csromanspacer{}วิปาโก ธมฺโม นเนววิปากนวิปากธมฺมธมฺมสฺส เหตุปจฺจเยน ปจฺจโย.\csromanspacer{}\csromanspacer{}วิปาโก ธมฺโม นวิปากธมฺมธมฺมสฺส จ นเนววิปากนวิปากธมฺมธมฺมสฺส จ ธมฺมสฺส เหตุปจฺจเยน ปจฺจโย.\csromanspacer{}\csromanspacer{}วิปาโก ธมฺโม นวิปากสฺส จ นวิปากธมฺมธมฺมสฺส จ ธมฺมสฺส เหตุปจฺจเยน ปจฺจโย. (5)}

\proseitem{24}{วิปาโก ธมฺโม นวิปากสฺส ธมฺมสฺส อารมฺมณปจฺจเยน ปจฺจโย.\csromanspacer{}\csromanspacer{}วิปาโก ธมฺโม นวิปากธมฺมธมฺมสฺส อารมฺมณปจฺจเยน ปจฺจโย.\csromanspacer{}\csromanspacer{}วิปาโก ธมฺโม นเนววิปากนวิปากธมฺมธมฺมสฺส อารมฺมณปจฺจเยน ปจฺจโย.\csromanspacer{}\csromanspacer{}วิปาโก ธมฺโม นวิปากสฺส จ นเนววิปากนวิปากธมฺมธมฺมสฺส จ ธมฺมสฺส อารมฺมณปจฺจเยน ปจฺจโย.\csromanspacer{}\csromanspacer{}วิปาโก ธมฺโม นวิปากธมฺมธมฺมสฺส จ นเนววิปากนวิปากธมฺมธมฺมสฺส จ ธมฺมสฺส อารมฺมณปจฺจเยน ปจฺจโย.\csromanspacer{}\csromanspacer{}วิปาโก ธมฺโม นวิปากสฺส จ นวิปากธมฺมธมฺมสฺส จ ธมฺมสฺส อารมฺมณปจฺจเยน ปจฺจโย. (6)}

\prose{วิปากธมฺมธมฺโม นวิปากธมฺมธมฺมสฺส อารมฺมณปจฺจเยน ปจฺจโย.\csromanspacer{}\csromanspacer{}วิปากธมฺมธมฺโม นวิปากสฺส ธมฺมสฺส อารมฺมณปจฺจเยน ปจฺจโย.\csromanspacer{}\csromanspacer{}วิปากธมฺมธมฺโม นเนววิปากนวิปากธมฺมธมฺมสฺส อารมฺมณปจฺจเยน ปจฺจโย.\csromanspacer{}\csromanspacer{}วิปากธมฺมธมฺโม นวิปากสฺส จ นเนววิปากนวิปากธมฺมธมฺมสฺส จ ธมฺมสฺส อารมฺมณปจฺจเยน ปจฺจโย.\csromanspacer{}\csromanspacer{}วิปากธมฺมธมฺโม นวิปากธมฺมธมฺมสฺส จ นเนววิปากนวิปากธมฺมธมฺมสฺส จ ธมฺมสฺส อารมฺมณปจฺจเยน ปจฺจโย.\csromanspacer{}\csromanspacer{}วิปากธมฺมธมฺโม นวิปากสฺส จ นวิปากธมฺมธมฺมสฺส จ ธมฺมสฺส อารมฺมณปจฺจเยน ปจฺจโย. (6)}

\prose{เนววิปากนวิปากธมฺมธมฺโม นเนววิปากนวิปากธมฺมธมฺมสฺส อารมฺมณปจฺจเยน ปจฺจโย ฯเปฯ.\csromanspacer{} ฉ. (สํขิตฺตํ.)}

\gambhira{ธมฺมานุโลมปจฺจนีย ติกปฏฺฐานปาฬิ}
\proseitem{25}{\csromanfolio{182}เหตุยา เตรส, อารมฺมเณ อฏฺฐารส, อธิปติยา สตฺตรส, อนนฺตเร โสฬส ฯเปฯ ปุเรชาเต ฉ, ปจฺฉาชาเต นว, อาเสวเน ฉ, กมฺเม จุทฺทส, วิปาเก ปญฺจ, อินฺทฺริเย อฏฺฐารส ฯเปฯ วิปฺปยุตฺเต ทฺวาทส ฯเปฯ อวิคเต ฉพฺพีส. (ปญฺหาวารํ วิตฺถาเรตพฺพํ.)}

\csromanheader{2}{4. อุปาทินฺนตฺติก 1. ปฏิจฺจวาร}
\csromanheader{2}{\textbf{ปจฺจยจตุกฺก-เหตุ}}
\proseitem{26}{อุปาทินฺนุปาทานิยํ ธมฺมํ ปฏิจฺจ น-อุปาทินฺนุปาทานิโย ธมฺโม อุปฺปชฺชติ เหตุปจฺจยา.\csromanspacer{}\csromanspacer{}อุปาทินฺนุปาทานิยํ ธมฺมํ ปฏิจฺจ น-อนุปาทินฺนุปาทานิโย ธมฺโม อุปฺปชฺชติ เหตุปจฺจยา.\csromanspacer{}\csromanspacer{}อุปาทินฺนุปาทานิยํ ธมฺมํ ปฏิจฺจ น-อนุปาทินฺน-อนุปาทานิโย ธมฺโม อุปฺปชฺชติ เหตุปจฺจยา.\csromanspacer{}\csromanspacer{}อุปาทินฺนุปาทานิยํ ธมฺมํ ปฏิจฺจ น-อุปาทินฺนุปาทานิโย จ น-อนุปาทินฺน-อนุปาทานิโย จ ธมฺมา อุปฺปชฺชนฺติ เหตุปจฺจยา.\csromanspacer{}\csromanspacer{}อุปาทินฺนุปาทานิยํ ธมฺมํ ปฏิจฺจ น-อนุปาทินฺนุปาทานิโย จ น-อนุปาทินฺน-อนุปาทานิโย จ ธมฺมา อุปฺปชฺชนฺติ เหตุปจฺจยา. (5)}

\prose{อนุปาทินฺนุปาทานิยํ ธมฺมํ ปฏิจฺจ น-อุปาทินฺนุปาทานิโย ธมฺโม อุปฺปชฺชติ เหตุปจฺจยา.\csromanspacer{}\csromanspacer{}อนุปาทินฺนุปาทานิยํ ธมฺมํ ปฏิจฺจ น-อนุปาทินฺน-อนุปาทานิโย ธมฺโม อุปฺปชฺชติ เหตุปจฺจยา.\csromanspacer{}\csromanspacer{}อนุปาทินฺนุปาทานิยํ ธมฺมํ ปฏิจฺจ น-อุปาทินฺนุปาทานิโย จ น-อนุปาทินฺน-อนุปาทานิโย จ ธมฺมา อุปฺปชฺชนฺติ เหตุปจฺจยา. (3)}

\prose{อนุปาทินฺน-อนุปาทานิยํ ธมฺมํ ปฏิจฺจ น-อนุปาทินฺน-อนุปาทานิโย ธมฺโม อุปฺปชฺชติ เหตุปจฺจยา.\csromanspacer{}\csromanspacer{}อนุปาทินฺน-อนุปาทานิยํ ธมฺมํ ปฏิจฺจ น-อุปาทินฺนุปาทานิโย ธมฺโม อุปฺปชฺชติ เหตุปจฺจยา.\csromanspacer{}\csromanspacer{}อนุปาทินฺน-อนุปาทานิยํ ธมฺมํ ปฏิจฺจ น-อนุปาทินฺนุปาทานิโย ธมฺโม อุปฺปชฺชติ เหตุปจฺจยา.\csromanspacer{}\csromanspacer{}อนุปาทินฺน-อนุปาทานิยํ ธมฺมํ ปฏิจฺจ น-อุปาทินฺนุปาทานิโย จ น-อนุปาทินฺน-อนุปาทานิโย จ ธมฺมา อุปฺปชฺชนฺติ เหตุปจฺจยา.\csromanspacer{}\csromanspacer{}อนุปาทินฺน-อนุปาทานิยํ ธมฺมํ ปฏิจฺจ น-อุปาทินฺนุปาทานิโย จ น-อนุปาทินฺนุปาทานิโย จ ธมฺมา อุปฺปชฺชนฺติ เหตุปจฺจยา. (5)}

\prose{อนุปาทินฺนุปาทานิยญฺจ อนุปาทินฺน-อนุปาทานิยญฺจ ธมฺมํ ปฏิจฺจ น-อุปาทินฺนุปาทานิโย ธมฺโม อุปฺปชฺชติ เหตุปจฺจยา.\csromanspacer{}\csromanspacer{}อนุปาทินฺนุปาทานิยญฺจ อนุปาทินฺน-อนุปาทานิยญฺจ ธมฺมํ ปฏิจฺจ น-อนุปาทินฺน-อนุปาทานิโย ธมฺโม อุปฺปชฺชติ เหตุปจฺจยา.\csromanspacer{}\csromanspacer{}อนุปาทินฺนุปาทานิยญฺจ อนุปาทินฺน-อนุปาทานิยญฺจ ธมฺมํ ปฏิจฺจ}

\vspace{\tipitakaparskip}
\csromancenter{อุปาทินฺนตฺติก น-อุปาทินฺนุปาทานิโย จ น-อนุปาทินฺน-อนุปาทานิโย จ ธมฺมา อุปฺปชฺชนฺติ เหตุปจฺจยา. (3)}
\prose{\csromanfolio{183}อุปาทินฺนุปาทานิยญฺจ อนุปาทินฺนุปาทานิยญฺจ ธมฺมํ ปฏิจฺจ น-อุปาทินฺนุปาทานิโย ธมฺโม อุปฺปชฺชติ เหตุปจฺจยา.\csromanspacer{}\csromanspacer{}อุปาทินฺนุปาทานิยญฺจ อนุปาทินฺนุปาทานิยญฺจ ธมฺมํ ปฏิจฺจ น-อนุปาทินฺน-อนุปาทานิโย ธมฺโม อุปฺปชฺชติ เหตุปจฺจยา.\csromanspacer{}\csromanspacer{}อุปาทินฺนุปาทานิยญฺจ อนุปาทินฺนุปาทานิยญฺจ ธมฺมํ ปฏิจฺจ น-อุปาทินฺนุปาทานิโย จ น-อนุปาทินฺน-อนุปาทานิโย จ ธมฺมา อุปฺปชฺชนฺติ เหตุปจฺจยา. (3) (สํขิตฺตํ.)}

\proseitem{27}{เหตุยา เอกูนวีส, อารมฺมเณ นว, อธิปติยา เอกาทส ฯเปฯ สหชาเต เอกูนวีส. (สหชาตวารมฺปิ ปจฺจยวารมฺปิ นิสฺสยวารมฺปิ สํสฏฺฐวารมฺปิ สมฺปยุตฺตวารมฺปิ วิตฺถาเรตพฺพํ.)}

\csromanheader{2}{4. อุปาทินฺนตฺติก 7. ปญฺหาวาร}
\csromanheader{2}{\textbf{ปจฺจยจตุกฺก-เหตุ-อารมฺมณ}}
\proseitem{28}{อุปาทินฺนุปาทานิโย ธมฺโม น-อุปาทินฺนุปาทานิยสฺส ธมฺมสฺส เหตุปจฺจเยน ปจฺจโย.\csromanspacer{}\csromanspacer{}อุปาทินฺนุปาทานิโย ธมฺโม น-อนุปาทินฺนุปาทานิยสฺส ธมฺมสฺส เหตุปจฺจเยน ปจฺจโย ฯเปฯ.}

\proseitem{29}{อุปาทินฺนุปาทานิโย ธมฺโม น-อุปาทินฺนุปาทานิยสฺส ธมฺมสฺส อารมฺมณปจฺจเยน ปจฺจโย.\csromanspacer{}\csromanspacer{}อุปาทินฺนุปาทานิโย ธมฺโม น-อนุปาทินฺนุปาทานิยสฺส ธมฺมสฺส อารมฺมณปจฺจเยน ปจฺจโย.\csromanspacer{}\csromanspacer{}อุปาทินฺนุปาทานิโย ธมฺโม น-อนุปาทินฺน-อนุปาทานิยสฺส ธมฺมสฺส อารมฺมณปจฺจเยน ปจฺจโย.\csromanspacer{}\csromanspacer{}อุปาทินฺนุปาทานิโย ธมฺโม น-อุปาทินฺนุปาทานิยสฺส จ น-อนุปาทินฺน-อนุปาทานิยสฺส จ ธมฺมสฺส อารมฺมณปจฺจเยน ปจฺจโย.\csromanspacer{}\csromanspacer{}อุปาทินฺนุปาทานิโย ธมฺโม น-อนุปาทินฺนุปาทานิยสฺส จ น-อนุปาทินฺน-อนุปาทานิยสฺส จ ธมฺมสฺส อารมฺมณปจฺจเยน ปจฺจโย. (5)}

\prose{อนุปาทินฺนุปาทานิโย ธมฺโม น-อนุปาทินฺนุปาทานิยสฺส ธมฺมสฺส อารมฺมณปจฺจเยน ปจฺจโย.\csromanspacer{}\csromanspacer{}อนุปาทินฺนุปาทานิโย ธมฺโม น-อุปาทินฺนุปาทานิยสฺส ธมฺมสฺส อารมฺมณปจฺจเยน ปจฺจโย.\csromanspacer{}\csromanspacer{}อนุปาทินฺนุปาทานิโย ธมฺโม น-อนุปาทินฺน-อนุปาทานิยสฺส ธมฺมสฺส อารมฺมณปจฺจเยน ปจฺจโย.\csromanspacer{}\csromanspacer{}}

\vspace{\tipitakaparskip}
\csromancenter{ธมฺมานุโลมปจฺจนีย ติกปฏฺฐานปาฬิ อนุปาทินฺนุปาทานิโย ธมฺโม น-อุปาทินฺนุปาทานิยสฺส จ น-อนุปาทินฺน-อนุปาทานิยสฺส จ ธมฺมสฺส อารมฺมณปจฺจเยน ปจฺจโย.\csromanspacer{}\csromanspacer{}อนุปาทินฺนุปาทานิโย ธมฺโม น-อนุปาทินฺนุปาทานิยสฺส จ น-อนุปาทินฺน-อนุปาทานิยสฺส จ ธมฺมสฺส อารมฺมณปจฺจเยน ปจฺจโย. (5)}
\prose{\csromanfolio{184}อนุปาทินฺน-อนุปาทานิโย ธมฺโม น-อนุปาทินฺน-อนุปาทานิยสฺส ธมฺมสฺส อารมฺมณปจฺจเยน ปจฺจโย ฯเปฯ.\csromanspacer{} ปญฺจ. (สํขิตฺตํ.)}

\proseitem{30}{เหตุยา เตรส, อารมฺมเณ ปนฺนรส, อธิปติยา เอกาทส. (ปญฺหาวารมฺปิ วิตฺถาเรตพฺพํ.)}

\csromanheader{2}{5. สํกิลิฏฺฐตฺติก 1. ปฏิจฺจวาร}
\csromanheader{2}{1. ปจฺจยานุโลม-เหตุ}
\proseitem{31}{สํกิลิฏฺฐสํกิเลสิกํ ธมฺมํ ปฏิจฺจ นสํกิลิฏฺฐสํกิเลสิโก ธมฺโม อุปฺปชฺชติ เหตุปจฺจยา.\csromanspacer{}\csromanspacer{}สํกิลิฏฺฐสํกิเลสิกํ ธมฺมํ ปฏิจฺจ น-อสํกิลิฏฺฐสํกิเลสิโก ธมฺโม อุปฺปชฺชติ เหตุปจฺจยา.\csromanspacer{}\csromanspacer{}สํกิลิฏฺฐสํกิเลสิกํ ธมฺมํ ปฏิจฺจ น-อสํกิลิฏฺฐ-อสํกิเลสิโก ธมฺโม อุปฺปชฺชติ เหตุปจฺจยา.\csromanspacer{}\csromanspacer{}สํกิลิฏฺฐสํกิเลสิกํ ธมฺมํ ปฏิจฺจ นสํกิลิฏฺฐสํกิเลสิโก จ น-อสํกิลิฏฺฐ-อสํกิเลสิโก จ ธมฺมา อุปฺปชฺชนฺติ เหตุปจฺจยา.\csromanspacer{}\csromanspacer{}สํกิลิฏฺฐสํกิเลสิกํ ธมฺมํ ปฏิจฺจ น-อสํกิลิฏฺฐสํกิเลสิโก จ น-อสํกิลิฏฺฐ-อสํกิเลสิโก จ ธมฺมา อุปฺปชฺชนฺติ เหตุปจฺจยา. (5)}

\prose{อสํกิลิฏฺฐสํกิเลสิกํ ธมฺมํ ปฏิจฺจ นสํกิลิฏฺฐสํกิเลสิโก ธมฺโม อุปฺปชฺชติ เหตุปจฺจยา.\csromanspacer{}\csromanspacer{}อสํกิลิฏฺฐสํกิเลสิกํ ธมฺมํ ปฏิจฺจ น-อสํกิลิฏฺฐ-อสํกิเลสิโก ธมฺโม อุปฺปชฺชติ เหตุปจฺจยา.\csromanspacer{}\csromanspacer{}อสํกิลิฏฺฐสํกิเลสิกํ ธมฺมํ ปฏิจฺจ นสํกิลิฏฺฐสํกิเลสิโก จ น-อสํกิลิฏฺฐ-อสํกิเลสิโก จ ธมฺมา อุปฺปชฺชนฺติ เหตุปจฺจยา. (3)}

\prose{อสํกิลิฏฺฐ-อสํกิเลสิกํ ธมฺมํ ปฏิจฺจ น-อสํกิลิฏฺฐ-อสํกิเลสิโก ธมฺโม อุปฺปชฺชติ เหตุปจฺจยา.\csromanspacer{}\csromanspacer{}อสํกิลิฏฺฐ-อสํกิเลสิกํ ธมฺมํ ปฏิจฺจ นสํกิลิฏฺฐสํกิเลสิโก ธมฺโม อุปฺปชฺชติ เหตุปจฺจยา.\csromanspacer{}\csromanspacer{}อสํกิลิฏฺฐ-อสํกิเลสิกํ ธมฺมํ ปฏิจฺจ น-อสํกิลิฏฺฐสํกิเลสิโก ธมฺโม อุปฺปชฺชติ}

\vspace{\tipitakaparskip}
\csromancenter{วิตกฺกตฺติก เหตุปจฺจยา.\csromanspacer{}\csromanspacer{}อสํกิลิฏฺฐ-อสํกิเลสิกํ ธมฺมํ ปฏิจฺจ นสํกิลิฏฺฐสํกิเลสิโก จ น-อสํกิลิฏฺฐ-อสํกิเลสิโก จ ธมฺมา อุปฺปชฺชนฺติ เหตุปจฺจยา.\csromanspacer{}\csromanspacer{}อสํกิลิฏฺฐ-อสํกิเลสิกํ ธมฺมํ ปฏิจฺจ นสํกิลิฏฺฐสํกิเลสิโก จ น-อสํกิลิฏฺฐสํกิเลสิโก จ ธมฺมา อุปฺปชฺชนฺติ เหตุปจฺจยา. (5)}
\prose{\csromanfolio{185}อสํกิลิฏฺฐสํกิเลสิกญฺจ อสํกิลิฏฺฐ-อสํกิเลสิกญฺจ ธมฺมํ ปฏิจฺจ นสํกิลิฏฺฐสํกิเลสิโก ธมฺโม อุปฺปชฺชติ เหตุปจฺจยา.\csromanspacer{} ตีณิ.}

\prose{สํกิลิฏฺฐสํกิเลสิกญฺจ อสํกิลิฏฺฐสํกิเลสิกญฺจ ธมฺมํ ปฏิจฺจ นสํกิลิฏฺฐสํกิเลสิโก ธมฺโม อุปฺปชฺชติ เหตุปจฺจยา ฯเปฯ.\csromanspacer{} ตีณิ.}

\prose{เหตุยา เอกูนวีส, อารมฺมเณ นว ฯเปฯ อวิคเต เอกูนวีส. (สหชาตวารมฺปิ ปจฺจยวารมฺปิ นิสฺสยวารมฺปิ สํสฏฺฐวารมฺปิ สมฺปยุตฺตวารมฺปิ วิตฺถาเรตพฺพํ.\csromanspacer{} ปญฺหาวาเร น สทิสํ.)}

\proseitem{32}{เหตุยา เตรส, อารมฺมเณ ปนฺนรส, อธิปติยา ปนฺนรส, อนนฺตเร โสฬส ฯเปฯ ปุเรชาเต ฉ, ปจฺฉาชาเต นว, อาเสวเน อฏฺฐ, กมฺเม เตรส, วิปาเก อฏฺฐ, อาหาเร เตรส ฯเปฯ มคฺเค เตรส, วิปฺปยุตฺเต ทฺวาทส ฯเปฯ อวิคเต ฉพฺพีส.}

\csromanheader{2}{6. วิตกฺกตฺติก 1. ปฏิจฺจวาร}
\csromanheader{2}{1. ปจฺจยานุโลม-เหตุ}
\proseitem{33}{สวิตกฺกสวิจารํ ธมฺมํ ปฏิจฺจ นสวิตกฺกสวิจาโร ธมฺโม อุปฺปชฺชติ เหตุปจฺจยา.\csromanspacer{}\csromanspacer{}สวิตกฺกสวิจารํ ธมฺมํ ปฏิจฺจ น-อวิตกฺกวิจารมตฺโต ธมฺโม อุปฺปชฺชติ เหตุปจฺจยา.\csromanspacer{}\csromanspacer{}สวิตกฺกสวิจารํ ธมฺมํ ปฏิจฺจ น-อวิตกฺก-อวิจาโร ธมฺโม อุปฺปชฺชติ เหตุปจฺจยา.\csromanspacer{}\csromanspacer{}สวิตกฺกสวิจารํ ธมฺมํ ปฏิจฺจ นสวิตกฺกสวิจาโร จ น-อวิตกฺก-อวิจาโร จ ธมฺมา อุปฺปชฺชนฺติ เหตุปจฺจยา.\csromanspacer{}\csromanspacer{}สวิตกฺกสวิจารํ ธมฺมํ ปฏิจฺจ น-อวิตกฺกวิจารมตฺโต จ น-อวิตกฺก-อวิจาโร จ ธมฺมา อุปฺปชฺชนฺติ เหตุปจฺจยา.\csromanspacer{}\csromanspacer{}สวิตกฺกสวิจารํ ธมฺมํ ปฏิจฺจ นสวิตกฺกสวิจาโร จ น-อวิตกฺกวิจารมตฺโต จ ธมฺมา อุปฺปชฺชนฺติ เหตุปจฺจยา.\csromanspacer{}\csromanspacer{}สวิตกฺกสวิจารํ ธมฺมํ ปฏิจฺจ นสวิตกฺกสวิจาโร จ น-อวิตกฺกวิจารมตฺโต จ น-อวิตกฺก-อวิจาโร จ ธมฺมา อุปฺปชฺชนฺติ เหตุปจฺจยา. (7)}

\gambhira{ธมฺมานุโลมปจฺจนีย ติกปฏฺฐานปาฬิ}
\proseitem{34}{\csromanfolio{186}อวิตกฺกวิจารมตฺตํ ธมฺมํ ปฏิจฺจ น-อวิตกฺกวิจารมตฺโต ธมฺโม อุปฺปชฺชติ เหตุปจฺจยา.\csromanspacer{}\csromanspacer{}อวิตกฺกวิจารมตฺตํ ธมฺมํ ปฏิจฺจ นสวิตกฺกสวิจาโร ธมฺโม อุปฺปชฺชติ เหตุปจฺจยา.\csromanspacer{}\csromanspacer{}อวิตกฺกวิจารมตฺตํ ธมฺมํ ปฏิจฺจ น-อวิตกฺก-อวิจาโร ธมฺโม อุปฺปชฺชติ เหตุปจฺจยา.\csromanspacer{}\csromanspacer{}อวิตกฺกวิจารมตฺตํ ธมฺมํ ปฏิจฺจ นสวิตกฺกสวิจาโร จ น-อวิตกฺก-อวิจาโร จ ธมฺมา อุปฺปชฺชนฺติ เหตุปจฺจยา.\csromanspacer{}\csromanspacer{}อวิตกฺกวิจารมตฺตํ ธมฺมํ ปฏิจฺจ น-อวิตกฺกวิจารมตฺโต จ น-อวิตกฺก-อวิจาโร จ ธมฺมา อุปฺปชฺชนฺติ เหตุปจฺจยา.\csromanspacer{}\csromanspacer{}อวิตกฺกวิจารมตฺตํ ธมฺมํ ปฏิจฺจ นสวิตกฺกสวิจาโร จ น-อวิตกฺกวิจารมตฺโต จ ธมฺมา อุปฺปชฺชนฺติ เหตุปจฺจยา.\csromanspacer{}\csromanspacer{}อวิตกฺกวิจารมตฺตํ ธมฺมํ ปฏิจฺจ นสวิตกฺกสวิจาโร จ น-อวิตกฺกวิจารมตฺโต จ น-อวิตกฺก-อวิจาโร จ ธมฺมา อุปฺปชฺชนฺติ เหตุปจฺจยา. (7)}

\prose{อวิตกฺก-อวิจารํ ธมฺมํ ปฏิจฺจ น-อวิตกฺก-อวิจาโร ธมฺโม อุปฺปชฺชติ เหตุปจฺจยา ฯเปฯ.\csromanspacer{} สตฺต.}

\prose{สวิตกฺกสวิจารญฺจ อวิตกฺก-อวิจารญฺจ ธมฺมํ ปฏิจฺจ นสวิตกฺกสวิจาโร ธมฺโม อุปฺปชฺชติ เหตุปจฺจยา ฯเปฯ.\csromanspacer{} สตฺต.}

\prose{อวิตกฺกวิจารมตฺตญฺจ อวิตกฺก-อวิจารญฺจ ธมฺมํ ปฏิจฺจ นสวิตกฺกสวิจาโร ธมฺโม อุปฺปชฺชติ เหตุปจฺจยา ฯเปฯ.\csromanspacer{} สตฺต.}

\prose{สวิตกฺกสวิจารญฺจ อวิตกฺกวิจารมตฺตญฺจ ธมฺมํ ปฏิจฺจ นสวิตกฺกสวิจาโร ธมฺโม อุปฺปชฺชติ เหตุปจฺจยา ฯเปฯ.\csromanspacer{} สตฺต.}

\prose{สวิตกฺกสวิจารญฺจ อวิตกฺกวิจารมตฺตญฺจ อวิตกฺก-อวิจารญฺจ ธมฺมํ ปฏิจฺจ นสวิตกฺกสวิจาโร ธมฺโม อุปฺปชฺชติ เหตุปจฺจยา ฯเปฯ.\csromanspacer{} สตฺต. (สํขิตฺตํ.)}

\prose{เหตุยา เอกูนปญฺญาส, อารมฺมเณ เอกูนปญฺญาส ฯเปฯ อวิคเต เอกูนปญฺญาส. (สหชาตวารมฺปิ ฯเปฯ ปญฺหาวารมฺปิ วิตฺถาเรตพฺพํ.)}

\csromanheader{2}{7. ปีติตฺติก 1. ปฏิจฺจวาร}
\csromanheader{2}{1. ปจฺจยานุโลม-เหตุ}
\proseitem{35}{ปีติสหคตํ ธมฺมํ ปฏิจฺจ นปีติสหคโต ธมฺโม อุปฺปชฺชติ เหตุปจฺจยา.\csromanspacer{}\csromanspacer{}ปีติสหคตํ ธมฺมํ ปฏิจฺจ นสุขสหคโต ธมฺโม อุปฺปชฺชติ เหตุปจฺจยา.\csromanspacer{}\csromanspacer{}ปีติสหคตํ ธมฺมํ ปฏิจฺจ น-อุเปกฺขาสหคโต ธมฺโม}

\vspace{\tipitakaparskip}
\csromancenter{ทสฺสนตฺติก อุปฺปชฺชติ เหตุปจฺจยา.\csromanspacer{}\csromanspacer{}ปีติสหคตํ ธมฺมํ ปฏิจฺจ นปีติสหคโต จ น-อุเปกฺขาสหคโต จ ธมฺมา อุปฺปชฺชนฺติ เหตุปจฺจยา.\csromanspacer{}\csromanspacer{}ปีติสหคตํ ธมฺมํ ปฏิจฺจ นสุขสหคโต จ น-อุเปกฺขาสหคโต จ ธมฺมา อุปฺปชฺชนฺติ เหตุปจฺจยา.\csromanspacer{}\csromanspacer{}ปีติสหคตํ ธมฺมํ ปฏิจฺจ นปีติสหคโต จ นสุขสหคโต จ ธมฺมา อุปฺปชฺชนฺติ เหตุปจฺจยา.\csromanspacer{}\csromanspacer{}ปีติสหคตํ ธมฺมํ ปฏิจฺจ นปีติสหคโต จ นสุขสหคโต จ น-อุเปกฺขาสหคโต จ ธมฺมา อุปฺปชฺชนฺติ เหตุปจฺจยา. (7)}
\prose{\csromanfolio{187}สุขสหคตํ ธมฺมํ ปฏิจฺจ นสุขสหคโต ธมฺโม อุปฺปชฺชติ เหตุปจฺจยา ฯเปฯ.\csromanspacer{} สตฺต.}

\prose{อุเปกฺขาสหคตํ ธมฺมํ ปฏิจฺจ น-อุเปกฺขาสหคโต ธมฺโม อุปฺปชฺชติ เหตุปจฺจยา ฯเปฯ.\csromanspacer{} สตฺต.}

\prose{ปีติสหคตญฺจ สุขสหคตญฺจ ธมฺมํ ปฏิจฺจ นปีติสหคโต ธมฺโม อุปฺปชฺชติ เหตุปจฺจยา ฯเปฯ.\csromanspacer{} สตฺต. (สํขิตฺตํ.)}

\prose{เหตุยา อฏฺฐวีส, อารมฺมเณ จตุวีส ฯเปฯ อวิคเต อฏฺฐวีส.}

\prose{(สหชาตวารมฺปิ ฯเปฯ สมฺปยุตฺตวารมฺปิ ปญฺหาวารมฺปิ สพฺพตฺถ วิตฺถาเรตพฺพํ.)}

\csromanheader{2}{8. ทสฺสนตฺติก 1. ปฏิจฺจวาร}
\csromanheader{2}{1. ปจฺจยานุโลม-เหตุ}
\proseitem{36}{ทสฺสเนน ปหาตพฺพํ ธมฺมํ ปฏิจฺจ นทสฺสเนน ปหาตพฺโพ ธมฺโม อุปฺปชฺชติ เหตุปจฺจยา.\csromanspacer{}\csromanspacer{}ทสฺสเนน ปหาตพฺพํ ธมฺมํ ปฏิจฺจ นภาวนาย ปหาตพฺโพ ธมฺโม อุปฺปชฺชติ เหตุปจฺจยา.\csromanspacer{}\csromanspacer{}ทสฺสเนน ปหาตพฺพํ ธมฺมํ ปฏิจฺจ นเนวทสฺสเนน นภาวนาย ปหาตพฺโพ ธมฺโม อุปฺปชฺชติ เหตุปจฺจยา.\csromanspacer{}\csromanspacer{}ทสฺสเนน ปหาตพฺพํ ธมฺมํ ปฏิจฺจ นภาวนาย ปหาตพฺโพ จ นเนวทสฺสเนน นภาวนาย ปหาตพฺโพ จ ธมฺมา อุปฺปชฺชนฺติ เหตุปจฺจยา.\csromanspacer{}\csromanspacer{}ทสฺสเนน ปหาตพฺพํ ธมฺมํ ปฏิจฺจ นทสฺสเนน ปหาตพฺโพ จ นภาวนาย ปหาตพฺโพ จ ธมฺมา อุปฺปชฺชนฺติ เหตุปจฺจยา. (5)}

\prose{ภาวนาย ปหาตพฺพํ ธมฺมํ ปฏิจฺจ นภาวนาย ปหาตพฺโพ ธมฺโม อุปฺปชฺชติ เหตุปจฺจยา.\csromanspacer{} ปญฺจ.}

\gambhira{ธมฺมานุโลมปจฺจนีย ติกปฏฺฐานปาฬิ}
\prose{\csromanfolio{188}เนวทสฺสเนน นภาวนาย ปหาตพฺพํ ธมฺมํ ปฏิจฺจ นทสฺสเนน ปหาตพฺโพ ธมฺโม อุปฺปชฺชติ เหตุปจฺจยา.\csromanspacer{} ตีณิ.}

\prose{ทสฺสเนน ปหาตพฺพญฺจ เนวทสฺสเนน นภาวนาย ปหาตพฺพญฺจ ธมฺมํ ปฏิจฺจ นทสฺสเนน ปหาตพฺโพ ธมฺโม อุปฺปชฺชติ เหตุปจฺจยา.\csromanspacer{} ตีณิ.}

\prose{ภาวนาย ปหาตพฺพญฺจ เนวทสฺสเนน นภาวนาย ปหาตพฺพญฺจ ธมฺมํ ปฏิจฺจ นทสฺสเนน ปหาตพฺโพ ธมฺโม อุปฺปชฺชติ เหตุปจฺจยา.\csromanspacer{} ตีณิ. (สํขิตฺตํ.)}

\prose{เหตุยา เอกูนวีส ฯเปฯ อวิคเต เอกูนวีส. (สหชาตวารมฺปิ ฯเปฯ สมฺปยุตฺตวารมฺปิ ปญฺหาวารมฺปิ วิตฺถาเรตพฺพํ.)}

\csromanheader{2}{9. ทสฺสนเหตุตฺติก 1. ปฏิจฺจวาร}
\proseitem{37}{ทสฺสเนน ปหาตพฺพเหตุกํ ธมฺมํ ปฏิจฺจ นทสฺสเนน ปหาตพฺพเหตุโก ธมฺโม อุปฺปชฺชติ เหตุปจฺจยา. (สํขิตฺตํ.)}

\csromanheader{2}{เหตุยา ฉพฺพีส ฯเปฯ อวิคเต ฉพฺพีส. (วิตฺถาเรตพฺพํ.)}
\csromanheader{2}{10. อาจยคามิตฺติก 1. ปฏิจฺจวาร}
\csromanheader{2}{1. ปจฺจยานุโลม-เหตุ}
\proseitem{38}{อาจยคามิํ ธมฺมํ ปฏิจฺจ น-อาจยคามี ธมฺโม อุปฺปชฺชติ เหตุปจฺจยา.\csromanspacer{}\csromanspacer{}อาจยคามิํ ธมฺมํ ปฏิจฺจ น-อปจยคามี ธมฺโม อุปฺปชฺชติ เหตุปจฺจยา.\csromanspacer{}\csromanspacer{}อาจยคามิํ ธมฺมํ ปฏิจฺจ นเนวาจยคามินาปจยคามี ธมฺโม อุปฺปชฺชติ เหตุปจฺจยา.\csromanspacer{}\csromanspacer{}อาจยคามิํ ธมฺมํ ปฏิจฺจ น-อปจยคามี จ นเนวาจยคามินาปจยคามี จ ธมฺมา อุปฺปชฺชนฺติ เหตุปจฺจยา.\csromanspacer{}\csromanspacer{}อาจยคามิํ ธมฺมํ ปฏิจฺจ น-อาจยคามี จ น-อปจยคามี จ ธมฺมา อุปฺปชฺชนฺติ เหตุปจฺจยา. (5)}

\prose{อปจยคามิํ ธมฺมํ ปฏิจฺจ น-อปจยคามี ธมฺโม อุปฺปชฺชติ เหตุปจฺจยา.\csromanspacer{}\csromanspacer{}อปจยคามิํ ธมฺมํ ปฏิจฺจ น-อาจยคามี ธมฺโม อุปฺปชฺชติ เหตุปจฺจยา.\csromanspacer{}\csromanspacer{}อปจยคามิํ ธมฺมํ ปฏิจฺจ นเนวาจยคามินาปจยคามี ธมฺโม อุปฺปชฺชติ เหตุปจฺจยา.\csromanspacer{}\csromanspacer{}อปจยคามิํ ธมฺมํ ปฏิจฺจ น-อาจยคามี จ นเนวาจยคามินาปจยคามี จ ธมฺมา อุปฺปชฺชนฺติ เหตุปจฺจยา.\csromanspacer{}\csromanspacer{}อปจยคามิํ ธมฺมํ ปฏิจฺจ น-อาจยคามี จ น-อปจยคามี จ ธมฺมา อุปฺปชฺชนฺติ เหตุปจฺจยา. (5)}

\csromanheader{2}{11. เสกฺขตฺติก}
\prose{\csromanfolio{189}เนวาจยคามินาปจยคามิํ ธมฺมํ ปฏิจฺจ น-อาจยคามี ธมฺโม อุปฺปชฺชติ เหตุปจฺจยา.\csromanspacer{}\csromanspacer{}เนวาจยคามินาปจยคามิํ ธมฺมํ ปฏิจฺจ น-อปจยคามี ธมฺโม อุปฺปชฺชติ เหตุปจฺจยา.\csromanspacer{}\csromanspacer{}เนวาจยคามินาปจยคามิํ ธมฺมํ ปฏิจฺจ น-อาจยคามี จ น-อปจยคามี จ ธมฺมา อุปฺปชฺชนฺติ เหตุปจฺจยา. (3)}

\prose{อาจยคามิญฺจ เนวาจยคามินาปจยคามิญฺจ ธมฺมํ ปฏิจฺจ น-อาจยคามี ธมฺโม อุปฺปชฺชติ เหตุปจฺจยา.\csromanspacer{}\csromanspacer{}อาจยคามิญฺจ เนวาจยคามินาปจยคามิญฺจ ธมฺมํ ปฏิจฺจ น-อปจยคามี ธมฺโม อุปฺปชฺชติ เหตุปจฺจยา.\csromanspacer{}\csromanspacer{}อาจยคามิญฺจ เนวาจยคามินาปจยคามิญฺจ ธมฺมํ ปฏิจฺจ น-อาจยคามี จ น-อปจยคามี จ ธมฺมา อุปฺปชฺชนฺติ เหตุปจฺจยา. (3)}

\prose{อปจยคามิญฺจ เนวาจยคามินาปจยคามิญฺจ ธมฺมํ ปฏิจฺจ น-อาจยคามี ธมฺโม อุปฺปชฺชติ เหตุปจฺจยา.\csromanspacer{}\csromanspacer{}อปจยคามิญฺจ เนวาจยคามินาปจยคามิญฺจ ธมฺมํ ปฏิจฺจ น-อปจยคามี ธมฺโม อุปฺปชฺชติ เหตุปจฺจยา.\csromanspacer{}\csromanspacer{}อปจยคามิญฺจ เนวาจยคามินาปจยคามิญฺจ ธมฺมํ ปฏิจฺจ น-อาจยคามี จ น-อปจยคามี จ ธมฺมา อุปฺปชฺชนฺติ เหตุปจฺจยา. (3) (สํขิตฺตํ.)}

\csromanheader{2}{เหตุยา เอกูนวีส. (สพฺพตฺถ วิตฺถาโร.)}
\csromanheader{2}{11. เสกฺขตฺติก 1. ปฏิจฺจวาร}
\csromanheader{2}{1. ปจฺจยานุโลม-เหตุ}
\proseitem{39}{เสกฺขํ ธมฺมํ ปฏิจฺจ นเสกฺโข ธมฺโม อุปฺปชฺชติ เหตุปจฺจยา.\csromanspacer{}\csromanspacer{}เสกฺขํ ธมฺมํ ปฏิจฺจ น-อเสกฺโข ธมฺโม อุปฺปชฺชติ เหตุปจฺจยา.\csromanspacer{}\csromanspacer{}เสกฺขํ ธมฺมํ ปฏิจฺจ นเนวเสกฺขนาเสกฺโข ธมฺโม อุปฺปชฺชติ เหตุปจฺจยา.\csromanspacer{}\csromanspacer{}เสกฺขํ ธมฺมํ ปฏิจฺจ น-อเสกฺโข จ นเนวเสกฺขนาเสกฺโข จ ธมฺมา อุปฺปชฺชนฺติ เหตุปจฺจยา.\csromanspacer{}\csromanspacer{}เสกฺขํ ธมฺมํ ปฏิจฺจ นเสกฺโข จ น-อเสกฺโข จ ธมฺมา อุปฺปชฺชนฺติ เหตุปจฺจยา. (5)}

\prose{อเสกฺขํ ธมฺมํ ปฏิจฺจ น-อเสกฺโข ธมฺโม อุปฺปชฺชติ เหตุปจฺจยา.\csromanspacer{}\csromanspacer{}อเสกฺขํ ธมฺมํ ปฏิจฺจ นเสกฺโข ธมฺโม อุปฺปชฺชติ เหตุปจฺจยา.\csromanspacer{}\csromanspacer{}อเสกฺขํ ธมฺมํ ปฏิจฺจ นเนวเสกฺขนาเสกฺโข ธมฺโม อุปฺปชฺชติ เหตุปจฺจยา.\csromanspacer{}\csromanspacer{}อเสกฺขํ ธมฺมํ ปฏิจฺจ นเสกฺโข จ นเนวเสกฺขนาเสกฺโข จ ธมฺมา อุปฺปชฺชนฺติ เหตุปจฺจยา.\csromanspacer{}\csromanspacer{}อเสกฺขํ ธมฺมํ ปฏิจฺจ นเสกฺโข จ น-อเสกฺโข จ ธมฺมา อุปฺปชฺชนฺติ เหตุปจฺจยา. (5)}

\gambhira{ธมฺมานุโลมปจฺจนีย ติกปฏฺฐานปาฬิ}
\prose{\csromanfolio{190}เนวเสกฺขนาเสกฺขํ ธมฺมํ ปฏิจฺจ นเสกฺโข ธมฺโม อุปฺปชฺชติ เหตุปจฺจยา.\csromanspacer{}\csromanspacer{}เนวเสกฺขนาเสกฺขํ ธมฺมํ ปฏิจฺจ น-อเสกฺโข ธมฺโม อุปฺปชฺชติ เหตุปจฺจยา.\csromanspacer{}\csromanspacer{}เนวเสกฺขนาเสกฺขํ ธมฺมํ ปฏิจฺจ นเสกฺโข จ น-อเสกฺโข จ ธมฺมา อุปฺปชฺชนฺติ เหตุปจฺจยา. (3)}

\prose{เสกฺขญฺจ เนวเสกฺขนาเสกฺขญฺจ ธมฺมํ ปฏิจฺจ นเสกฺโข ธมฺโม อุปฺปชฺชติ เหตุปจฺจยา.\csromanspacer{}\csromanspacer{}เสกฺขญฺจ เนวเสกฺขนาเสกฺขญฺจ ธมฺมํ ปฏิจฺจ น-อเสกฺโข ธมฺโม อุปฺปชฺชติ เหตุปจฺจยา.\csromanspacer{}\csromanspacer{}เสกฺขญฺจ เนวเสกฺขนาเสกฺขญฺจ ธมฺมํ ปฏิจฺจ นเสกฺโข จ น-อเสกฺโข จ ธมฺมา อุปฺปชฺชนฺติ เหตุปจฺจยา. (3)}

\prose{อเสกฺขญฺจ เนวเสกฺขนาเสกฺขญฺจ ธมฺมํ ปฏิจฺจ นเสกฺโข ธมฺโม อุปฺปชฺชติ เหตุปจฺจยา.\csromanspacer{}\csromanspacer{}อเสกฺขญฺจ เนวเสกฺขนาเสกฺขญฺจ ธมฺมํ ปฏิจฺจ น-อเสกฺโข ธมฺโม อุปฺปชฺชติ เหตุปจฺจยา.\csromanspacer{}\csromanspacer{}อเสกฺขญฺจ เนวเสกฺขนาเสกฺขญฺจ ธมฺมํ ปฏิจฺจ นเสกฺโข จ น-อเสกฺโข จ ธมฺมา อุปฺปชฺชนฺติ เหตุปจฺจยา. (3) (สํขิตฺตํ.)}

\csromanheader{2}{เหตุยา เอกูนวีส. (สพฺพตฺถ วิตฺถาโร.)}
\csromanheader{2}{12. ปริตฺตตฺติก 1. ปฏิจฺจวาร}
\csromanheader{2}{1. ปจฺจยานุโลม-เหตุ}
\proseitem{40}{ปริตฺตํ ธมฺมํ ปฏิจฺจ นปริตฺโต ธมฺโม อุปฺปชฺชติ เหตุปจฺจยา.\csromanspacer{}\csromanspacer{}ปริตฺตํ ธมฺมํ ปฏิจฺจ นมหคฺคโต ธมฺโม อุปฺปชฺชติ เหตุปจฺจยา.\csromanspacer{}\csromanspacer{}ปริตฺตํ ธมฺมํ ปฏิจฺจ น-อปฺปมาโณ ธมฺโม อุปฺปชฺชติ เหตุปจฺจยา.\csromanspacer{}\csromanspacer{}ปริตฺตํ ธมฺมํ ปฏิจฺจ นปริตฺโต จ น-อปฺปมาโณ จ ธมฺมา อุปฺปชฺชนฺติ เหตุปจฺจยา.\csromanspacer{}\csromanspacer{}ปริตฺตํ ธมฺมํ ปฏิจฺจ นมหคฺคโต จ น-อปฺปมาโณ จ ธมฺมา อุปฺปชฺชนฺติ เหตุปจฺจยา. (5)}

\prose{มหคฺคตํ ธมฺมํ ปฏิจฺจ นมหคฺคโต ธมฺโม อุปฺปชฺชติ เหตุปจฺจยา.\csromanspacer{}\csromanspacer{}มหคฺคตํ ธมฺมํ ปฏิจฺจ นปริตฺโต ธมฺโม อุปฺปชฺชติ เหตุปจฺจยา.\csromanspacer{}\csromanspacer{}มหคฺคตํ ธมฺมํ ปฏิจฺจ น-อปฺปมาโณ ธมฺโม อุปฺปชฺชติ เหตุปจฺจยา.\csromanspacer{}\csromanspacer{}มหคฺคตํ ธมฺมํ ปฏิจฺจ นปริตฺโต จ น-อปฺปมาโณ จ ธมฺมา อุปฺปชฺชนฺติ เหตุปจฺจยา.\csromanspacer{}\csromanspacer{}มหคฺคตํ ธมฺมํ ปฏิจฺจ นมหคฺคโต จ น-อปฺปมาโณ จ ธมฺมา อุปฺปชฺชนฺติ เหตุปจฺจยา. (5)}

\prose{อปฺปมาณํ ธมฺมํ ปฏิจฺจ น-อปฺปมาโณ ธมฺโม อุปฺปชฺชติ เหตุปจฺจยา.\csromanspacer{}\csromanspacer{}อปฺปมาณํ ธมฺมํ ปฏิจฺจ นปริตฺโต ธมฺโม อุปฺปชฺชติ เหตุปจฺจยา.\csromanspacer{}\csromanspacer{}อปฺปมาณํ ธมฺมํ}

\vspace{\tipitakaparskip}
\csromancenter{มิจฺฉตฺตตฺติก ปฏิจฺจ นมหคฺคโต ธมฺโม อุปฺปชฺชติ เหตุปจฺจยา.\csromanspacer{}\csromanspacer{}อปฺปมาณํ ธมฺมํ ปฏิจฺจ นมหคฺคโต จ น-อปฺปมาโณ จ ธมฺมา อุปฺปชฺชนฺติ เหตุปจฺจยา.\csromanspacer{}\csromanspacer{}อปฺปมาณํ ธมฺมํ ปฏิจฺจ นปริตฺโต จ นมหคฺคโต จ ธมฺมา อุปฺปชฺชนฺติ เหตุปจฺจยา. (5)}
\prose{\csromanfolio{191}ปริตฺตญฺจ อปฺปมาณญฺจ ธมฺมํ ปฏิจฺจ นมหคฺคโต ธมฺโม อุปฺปชฺชติ เหตุปจฺจยา.\csromanspacer{}\csromanspacer{}ปริตฺตญฺจ อปฺปมาณญฺจ ธมฺมํ ปฏิจฺจ น-อปฺปมาโณ ธมฺโม อุปฺปชฺชติ เหตุปจฺจยา.\csromanspacer{}\csromanspacer{}ปริตฺตญฺจ อปฺปมาณญฺจ ธมฺมํ ปฏิจฺจ นมหคฺคโต จ น-อปฺปมาโณ จ ธมฺมา อุปฺปชฺชนฺติ เหตุปจฺจยา. (3)}

\prose{ปริตฺตญฺจ มหคฺคตญฺจ ธมฺมํ ปฏิจฺจ นปริตฺโต ธมฺโม อุปฺปชฺชติ เหตุปจฺจยา.\csromanspacer{}\csromanspacer{}ปริตฺตญฺจ มหคฺคตญฺจ ธมฺมํ ปฏิจฺจ นมหคฺคโต ธมฺโม อุปฺปชฺชติ เหตุปจฺจยา.\csromanspacer{}\csromanspacer{}ปริตฺตญฺจ มหคฺคตญฺจ ธมฺมํ ปฏิจฺจ น-อปฺปมาโณ ธมฺโม อุปฺปชฺชติ เหตุปจฺจยา.\csromanspacer{}\csromanspacer{}ปริตฺตญฺจ มหคฺคตญฺจ ธมฺมํ ปฏิจฺจ นปริตฺโต จ น-อปฺปมาโณ จ ธมฺมา อุปฺปชฺชนฺติ เหตุปจฺจยา.\csromanspacer{}\csromanspacer{}ปริตฺตญฺจ มหคฺคตญฺจ ธมฺมํ ปฏิจฺจ นมหคฺคโต จ น-อปฺปมาโณ จ ธมฺมา อุปฺปชฺชนฺติ เหตุปจฺจยา. (5)}

\prose{เหตุยา เตวีส, อารมฺมเณ จุทฺทส. (สพฺพตฺถ วิตฺถาเรตพฺพํ.)}

\csromanheader{2}{13. ปริตฺตารมฺมณตฺติก 1. ปฏิจฺจวาร}
\proseitem{41}{ปริตฺตารมฺมณํ ธมฺมํ ปฏิจฺจ นปริตฺตารมฺมโณ ธมฺโม อุปฺปชฺชติ เหตุปจฺจยา.}

\prose{เหตุยา เอกวีส ฯเปฯ อวิคเต เอกวีส.}

\csromanheader{2}{14. หีนตฺติก 1. ปฏิจฺจวาร}
\proseitem{42}{หีนํ ธมฺมํ ปฏิจฺจ นหีโน ธมฺโม อุปฺปชฺชติ เหตุปจฺจยา. (สํกิลิฏฺฐสํกิเลสิกตฺติกสทิสํ.)}

\prose{เหตุยา เอกูนวีส ฯเปฯ อวิคเต เอกูนวีส.}

\csromanheader{2}{15. มิจฺฉตฺตตฺติก 1. ปฏิจฺจวาร}
\csromanheader{2}{1. ปจฺจยานุโลม-เหตุ}
\proseitem{43}{มิจฺฉตฺตนิยตํ ธมฺมํ ปฏิจฺจ นมิจฺฉตฺตนิยโต ธมฺโม อุปฺปชฺชติ เหตุปจฺจยา.\csromanspacer{}\csromanspacer{}มิจฺฉตฺตนิยตํ ธมฺมํ ปฏิจฺจ นสมฺมตฺตนิยโต ธมฺโม อุปฺปชฺชติ}

\vspace{\tipitakaparskip}
\csromancenter{ธมฺมานุโลมปจฺจนีย ติกปฏฺฐานปาฬิ เหตุปจฺจยา.\csromanspacer{}\csromanspacer{}มิจฺฉตฺตนิยตํ ธมฺมํ ปฏิจฺจ น-อนิยโต ธมฺโม อุปฺปชฺชติ เหตุปจฺจยา.\csromanspacer{}\csromanspacer{}มิจฺฉตฺตนิยตํ ธมฺมํ ปฏิจฺจ นสมฺมตฺตนิยโต จ น-อนิยโต จ ธมฺมา อุปฺปชฺชนฺติ เหตุปจฺจยา.\csromanspacer{}\csromanspacer{}มิจฺฉตฺตนิยตํ ธมฺมํ ปฏิจฺจ นมิจฺฉตฺตนิยโต จ นสมฺมตฺตนิยโต จ ธมฺมา อุปฺปชฺชนฺติ เหตุปจฺจยา. (5)}
\prose{\csromanfolio{192}สมฺมตฺตนิยตํ ธมฺมํ ปฏิจฺจ นสมฺมตฺตนิยโต ธมฺโม อุปฺปชฺชติ เหตุปจฺจยา.\csromanspacer{}\csromanspacer{}สมฺมตฺตนิยตํ ธมฺมํ ปฏิจฺจ นมิจฺฉตฺตนิยโต ธมฺโม อุปฺปชฺชติ เหตุปจฺจยา.\csromanspacer{}\csromanspacer{}สมฺมตฺตนิยตํ ธมฺมํ ปฏิจฺจ น-อนิยโต ธมฺโม อุปฺปชฺชติ เหตุปจฺจยา.\csromanspacer{}\csromanspacer{}สมฺมตฺตนิยตํ ธมฺมํ ปฏิจฺจ นมิจฺฉตฺตนิยโต จ น-อนิยโต จ ธมฺมา อุปฺปชฺชนฺติ เหตุปจฺจยา.\csromanspacer{}\csromanspacer{}สมฺมตฺตนิยตํ ธมฺมํ ปฏิจฺจ นมิจฺฉตฺตนิยโต จ นสมฺมตฺตนิยโต จ ธมฺมา อุปฺปชฺชนฺติ เหตุปจฺจยา. (5)}

\prose{อนิยตํ ธมฺมํ ปฏิจฺจ นมิจฺฉตฺตนิยโต ธมฺโม อุปฺปชฺชติ เหตุปจฺจยา.\csromanspacer{} ตีณิ.}

\prose{มิจฺฉตฺตนิยตญฺจ อนิยตญฺจ ธมฺมํ ปฏิจฺจ นมิจฺฉตฺตนิยโต ธมฺโม อุปฺปชฺชติ เหตุปจฺจยา.\csromanspacer{} ตีณิ.}

\prose{สมฺมตฺตนิยตญฺจ อนิยตญฺจ ธมฺมํ ปฏิจฺจ นมิจฺฉตฺตนิยโต ธมฺโม อุปฺปชฺชติ เหตุปจฺจยา.\csromanspacer{} ตีณิ. (สํขิตฺตํ.)}

\csromanheader{2}{เหตุยา เอกูนวีส. (สพฺพตฺถ วิตฺถาเรตพฺพํ.)}
\csromanheader{2}{16. มคฺคารมฺมณตฺติก 1. ปฏิจฺจวาร}
\proseitem{44}{มคฺคารมฺมณํ ธมฺมํ ปฏิจฺจ นมคฺคารมฺมโณ ธมฺโม อุปฺปชฺชติ เหตุปจฺจยา.\csromanspacer{}\csromanspacer{}มคฺคารมฺมณํ ธมฺมํ ปฏิจฺจ นมคฺคเหตุโก ธมฺโม อุปฺปชฺชติ เหตุปจฺจยา.\csromanspacer{}\csromanspacer{}มคฺคารมฺมณํ ธมฺมํ ปฏิจฺจ นมคฺคาธิปติ ธมฺโม อุปฺปชฺชติ เหตุปจฺจยา.\csromanspacer{}\csromanspacer{}มคฺคารมฺมณํ ธมฺมํ ปฏิจฺจ นมคฺคารมฺมโณ จ นมคฺคาธิปติ จ ธมฺมา อุปฺปชฺชนฺติ เหตุปจฺจยา.\csromanspacer{}\csromanspacer{}มคฺคารมฺมณํ ธมฺมํ ปฏิจฺจ นมคฺคเหตุโก จ นมคฺคาธิปติ จ ธมฺมา อุปฺปชฺชนฺติ เหตุปจฺจยา.\csromanspacer{}\csromanspacer{}มคฺคารมฺมณํ ธมฺมํ ปฏิจฺจ นมคฺคารมฺมโณ จ นมคฺคเหตุโก จ ธมฺมา อุปฺปชฺชนฺติ เหตุปจฺจยา.\csromanspacer{}\csromanspacer{}มคฺคารมฺมณํ ธมฺมํ ปฏิจฺจ นมคฺคารมฺมโณ จ นมคฺคเหตุโก จ นมคฺคาธิปติ จ ธมฺมา อุปฺปชฺชนฺติ เหตุปจฺจยา. (7) (สํขิตฺตํ.)}

\prose{เหตุยา ปญฺจติํส ฯเปฯ อวิคเต ปญฺจติํส. (สพฺพตฺถ วิตฺถาเรตพฺพํ.)}

\csromanheader{2}{20. อชฺฌตฺตตฺติก \textbf{17. อุปฺปนฺนตฺติก 7. ปญฺหาวาร}}
\proseitem{45}{\csromanfolio{193}อุปฺปนฺโน ธมฺโม น-อนุปฺปนฺนสฺส ธมฺมสฺส เหตุปจฺจเยน ปจฺจโย. (สํขิตฺตํ.)}

\prose{เหตุยา ตีณิ, อารมฺมเณ นว.}

\csromanheader{2}{18. อตีตตฺติก 7. ปญฺหาวาร}
\proseitem{46}{ปจฺจุปฺปนฺโน ธมฺโม น-อตีตสฺส ธมฺมสฺส เหตุปจฺจเยน ปจฺจโย. (สํขิตฺตํ.)}

\prose{เหตุยา ตีณิ, อารมฺมเณ นว.}

\csromanheader{2}{19. อตีตารมฺมณตฺติก 1. ปฏิจฺจวาร}
\proseitem{47}{อตีตารมฺมณํ ธมฺมํ ปฏิจฺจ น-อตีตารมฺมโณ ธมฺโม อุปฺปชฺชติ เหตุปจฺจยา.\csromanspacer{}\csromanspacer{}อตีตารมฺมณํ ธมฺมํ ปฏิจฺจ น-อนาคตารมฺมโณ ธมฺโม อุปฺปชฺชติ เหตุปจฺจยา.\csromanspacer{}\csromanspacer{}อตีตารมฺมณํ ธมฺมํ ปฏิจฺจ นปจฺจุปฺปนฺนารมฺมโณ ธมฺโม อุปฺปชฺชติ เหตุปจฺจยา.\csromanspacer{}\csromanspacer{}อตีตารมฺมณํ ธมฺมํ ปฏิจฺจ น-อตีตารมฺมโณ จ นปจฺจุปฺปนฺนารมฺมโณ จ ธมฺมา อุปฺปชฺชนฺติ เหตุปจฺจยา.\csromanspacer{}\csromanspacer{}อตีตารมฺมณํ ธมฺมํ ปฏิจฺจ น-อนาคตารมฺมโณ จ นปจฺจุปฺปนฺนารมฺมโณ จ ธมฺมา อุปฺปชฺชนฺติ เหตุปจฺจยา.\csromanspacer{}\csromanspacer{}อตีตารมฺมณํ ธมฺมํ ปฏิจฺจ น-อตีตารมฺมโณ จ น-อนาคตารมฺมโณ จ ธมฺมา อุปฺปชฺชนฺติ เหตุปจฺจยา.\csromanspacer{}\csromanspacer{}อตีตารมฺมณํ ธมฺมํ ปฏิจฺจ น-อตีตารมฺมโณ จ น-อนาคตารมฺมโณ จ นปจฺจุปฺปนฺนารมฺมโณ จ ธมฺมา อุปฺปชฺชนฺติ เหตุปจฺจยา. (7) (สํขิตฺตํ.)}

\prose{เหตุยา เอกวีส ฯเปฯ อวิคเต เอกวีส.}

\csromanheader{2}{20. อชฺฌตฺตตฺติก 1. ปฏิจฺจวาร}
\csromanheader{2}{48. อชฺฌตฺตํ ธมฺมํ ปฏิจฺจ นพหิทฺธา ธมฺโม อุปฺปชฺชติ}
\prose{พหิทฺธา ธมฺมํ ปฏิจฺจ น-อชฺฌตฺโต ธมฺโม อุปฺปชฺชติ เหตุปจฺจยา. (1) เหตุยา เทฺว. (สพฺพตฺถ วิตฺถาโร.)}

\csromanheader{2}{ธมฺมานุโลมปจฺจนีย ติกปฏฺฐานปาฬิ \textbf{21. อชฺฌตฺตารมฺมณตฺติก 1. ปฏิจฺจวาร}}
\proseitem{49}{\csromanfolio{194}อชฺฌตฺตารมฺมณํ ธมฺมํ ปฏิจฺจ น-อชฺฌตฺตารมฺมโณ ธมฺโม อุปฺปชฺชติ เหตุปจฺจยา.}

\prose{เหตุยา ฉ.}

\csromanheader{2}{22. สนิทสฺสนตฺติก 1-7. ปฏิจฺจาทิวาร}
\csromanheader{2}{1. ปจฺจยานุโลม-เหตุ}
\proseitem{50}{อนิทสฺสนสปฺปฏิฆํ ธมฺมํ ปฏิจฺจ น-อนิทสฺสนสปฺปฏิโฆ ธมฺโม อุปฺปชฺชติ เหตุปจฺจยา.\csromanspacer{}\csromanspacer{}อนิทสฺสนสปฺปฏิฆํ ธมฺมํ ปฏิจฺจ นสนิทสฺสนสปฺปฏิโฆ ธมฺโม อุปฺปชฺชติ เหตุปจฺจยา.\csromanspacer{}\csromanspacer{}อนิทสฺสนสปฺปฏิฆํ ธมฺมํ ปฏิจฺจ น-อนิทสฺสน-อปฺปฏิโฆ ธมฺโม อุปฺปชฺชติ เหตุปจฺจยา.\csromanspacer{}\csromanspacer{}อนิทสฺสนสปฺปฏิฆํ ธมฺมํ ปฏิจฺจ นสนิทสฺสนสปฺปฏิโฆ จ น-อนิทสฺสน-อปฺปฏิโฆ จ ธมฺมา อุปฺปชฺชนฺติ เหตุปจฺจยา.\csromanspacer{}\csromanspacer{}อนิทสฺสนสปฺปฏิฆํ ธมฺมํ ปฏิจฺจ น-อนิทสฺสนสปฺปฏิโฆ จ น-อนิทสฺสน-อปฺปฏิโฆ จ ธมฺมา อุปฺปชฺชนฺติ เหตุปจฺจยา.\csromanspacer{}\csromanspacer{}อนิทสฺสนสปฺปฏิฆํ ธมฺมํ ปฏิจฺจ นสนิทสฺสนสปฺปฏิโฆ จ น-อนิทสฺสนสปฺปฏิโฆ จ ธมฺมา อุปฺปชฺชนฺติ เหตุปจฺจยา. (6)}

\prose{อนิทสฺสน-อปฺปฏิฆํ ธมฺมํ ปฏิจฺจ น-อนิทสฺสน-อปฺปฏิโฆ ธมฺโม อุปฺปชฺชติ เหตุปจฺจยา.\csromanspacer{} ฉ.}

\prose{อนิทสฺสนสปฺปฏิฆญฺจ อนิทสฺสน-อปฺปฏิฆญฺจ ธมฺมํ ปฏิจฺจ นสนิทสฺสนสปฺปฏิโฆ ธมฺโม อุปฺปชฺชติ เหตุปจฺจยา.\csromanspacer{} ฉ.}

\prose{เหตุยา อฏฺฐารส, อารมฺมเณ ตีณิ ฯเปฯ อวิคเต อฏฺฐารส. (สพฺพตฺถ วิตฺถาโร.\csromanspacer{} สหชาตวารมฺปิ ฯเปฯ สมฺปยุตฺตวารมฺปิ วิตฺถาเรตพฺพํ.)}

\csromanheader{2}{\textbf{เหตุ-อารมฺมณ}}
\proseitem{51}{อนิทสฺสน-อปฺปฏิโฆ ธมฺโม น-อนิทสฺสน-อปฺปฏิฆสฺส ธมฺมสฺส เหตุปจฺจเยน ปจฺจโย.\csromanspacer{}\csromanspacer{}อนิทสฺสน-อปฺปฏิโฆ ธมฺโม นสนิทสฺสนสปฺปฏิฆสฺส ธมฺมสฺส เหตุปจฺจเยน ปจฺจโย.\csromanspacer{}\csromanspacer{}อนิทสฺสน-อปฺปฏิโฆ ธมฺโม น-อนิทสฺสนสปฺปฏิฆสฺส ธมฺมสฺส เหตุปจฺจเยน ปจฺจโย.\csromanspacer{}\csromanspacer{}}

\vspace{\tipitakaparskip}
\csromancenter{สนิทสฺสนตฺติก อนิทสฺสน-อปฺปฏิโฆ ธมฺโม นสนิทสฺสนสปฺปฏิฆสฺส จ น-อนิทสฺสน-อปฺปฏิฆสฺส จ ธมฺมสฺส เหตุปจฺจเยน ปจฺจโย.\csromanspacer{}\csromanspacer{}อนิทสฺสน-อปฺปฏิโฆ ธมฺโม น-อนิทสฺสนสปฺปฏิฆสฺส จ น-อนิทสฺสน-อปฺปฏิฆสฺส จ ธมฺมสฺส เหตุปจฺจเยน ปจฺจโย.\csromanspacer{}\csromanspacer{}อนิทสฺสน-อปฺปฏิโฆ ธมฺโม นสนิทสฺสนสปฺปฏิฆสฺส จ น-อนิทสฺสนสปฺปฏิฆสฺส จ ธมฺมสฺส เหตุปจฺจเยน ปจฺจโย. (6)}
\prose{\csromanfolio{195}สนิทสฺสนสปฺปฏิโฆ ธมฺโม นสนิทสฺสนสปฺปฏิฆสฺส ธมฺมสฺส อารมฺมณปจฺจเยน ปจฺจโย.\csromanspacer{} ตีณิ. (สํขิตฺตํ.)}

\proseitem{52}{เหตุยา ฉ, อารมฺมเณ นว. (ปญฺหาวารํ วิตฺถาเรตพฺพํ.)}

\nitthitam{\textbf{ธมฺมานุโลมปจฺจนีเย ติกปฏฺฐานํ นิฏฺฐิตํ.}}
\pitaka{\textbf{อภิธมฺมปิฏก}}
\csromanheader{2}{\textbf{ธมฺมานุโลมปจฺจนีย}}
\gambhira{\textbf{ทุกปฏฺฐานปาฬิ}}
\namakkaram{นโม ตสฺส ภควโต อรหโต สมฺมาสมฺพุทฺธสฺส.}
\csromanheader{2}{1. เหตุทุก 1. ปฏิจฺจาทิวาร}
\csromanheader{2}{\textbf{ปจฺจยานุโลม-เหตุ}}
\proseitem{1}{\csromanfolio{197}เหตุํ ธมฺมํ ปฏิจฺจ นเหตุ ธมฺโม อุปฺปชฺชติ เหตุปจฺจยา.\csromanspacer{} เหตุํ ธมฺมํ ปฏิจฺจ สมฺปยุตฺตกา ขนฺธา จิตฺตสมุฏฺฐานญฺจ รูปํ.\csromanspacer{} ปฏิสนฺธิกฺขเณ ฯเปฯ.\csromanspacer{} เหตุํ ธมฺมํ ปฏิจฺจ นนเหตุ ธมฺโม อุปฺปชฺชติ เหตุปจฺจยา.\csromanspacer{}\csromanspacer{}เหตุํ ธมฺมํ ปฏิจฺจ นเหตุ จ นนเหตุ จ ธมฺมา อุปฺปชฺชนฺติ เหตุปจฺจยา. (3)}

\prose{นเหตุํ ธมฺมํ ปฏิจฺจ นนเหตุ ธมฺโม อุปฺปชฺชติ เหตุปจฺจยา.\csromanspacer{}\csromanspacer{}นเหตุํ ธมฺมํ ปฏิจฺจ นเหตุ ธมฺโม อุปฺปชฺชติ เหตุปจฺจยา.\csromanspacer{}\csromanspacer{}นเหตุํ ธมฺมํ ปฏิจฺจ นเหตุ จ นนเหตุ จ ธมฺมา อุปฺปชฺชนฺติ เหตุปจฺจยา. (3)}

\prose{เหตุญฺจ นเหตุญฺจ ธมฺมํ ปฏิจฺจ นเหตุ ธมฺโม อุปฺปชฺชติ เหตุปจฺจยา.\csromanspacer{}\csromanspacer{}เหตุญฺจ นเหตุญฺจ ธมฺมํ ปฏิจฺจ นนเหตุ ธมฺโม อุปฺปชฺชติ เหตุปจฺจยา.\csromanspacer{}\csromanspacer{}เหตุญฺจ นเหตุญฺจ ธมฺมํ ปฏิจฺจ นเหตุ จ นนเหตุ จ ธมฺมา อุปฺปชฺชนฺติ เหตุปจฺจยา. (3) (สํขิตฺตํ.)}

\prose{เหตุยา นว, อารมฺมเณ นว ฯเปฯ อวิคเต นว. (สหชาตวารมฺปิ ฯเปฯ สมฺปยุตฺตวารมฺปิ ปฏิจฺจวารสทิสํ.)}

\csromanheader{2}{ธมฺมานุโลมปจฺจนีย ทุกปฏฺฐานปาฬิ \textbf{เหตุ-อารมฺมณ}}
\proseitem{2}{\csromanfolio{198}เหตุ ธมฺโม นเหตุสฺส ธมฺมสฺส เหตุปจฺจเยน ปจฺจโย.\csromanspacer{} ตีณิ.}

\prose{เหตุ ธมฺโม นเหตุสฺส ธมฺมสฺส อารมฺมณปจฺจเยน ปจฺจโย.\csromanspacer{} ตีณิ.}

\prose{นเหตุ ธมฺโม นนเหตุสฺส ธมฺมสฺส อารมฺมณปจฺจเยน ปจฺจโย.\csromanspacer{} ตีณิ.}

\prose{เหตุ จ นเหตุ จ ธมฺมา นเหตุสฺส ธมฺมสฺส อารมฺมณปจฺจเยน ปจฺจโย.\csromanspacer{} ตีณิ. (สํขิตฺตํ.)}

\proseitem{3}{เหตุยา ตีณิ, อารมฺมเณ นว ฯเปฯ อวิคเต นว. (ปญฺหาวารมฺปิ เอวํ วิตฺถาเรตพฺพํ.)}

\csromanheader{2}{2. สเหตุกทุก 1. ปฏิจฺจวาร}
\proseitem{4}{สเหตุกํ ธมฺมํ ปฏิจฺจ นสเหตุโก ธมฺโม อุปฺปชฺชติ เหตุปจฺจยา.\csromanspacer{}\csromanspacer{}สเหตุกํ ธมฺมํ ปฏิจฺจ น-อเหตุโก ธมฺโม อุปฺปชฺชติ เหตุปจฺจยา.\csromanspacer{}\csromanspacer{}สเหตุกํ ธมฺมํ ปฏิจฺจ นสเหตุโก จ น-อเหตุโก จ ธมฺมา อุปฺปชฺชนฺติ เหตุปจฺจยา. (3)}

\prose{อเหตุกํ ธมฺมํ ปฏิจฺจ น-อเหตุโก ธมฺโม อุปฺปชฺชติ เหตุปจฺจยา.\csromanspacer{}\csromanspacer{}อเหตุกํ ธมฺมํ ปฏิจฺจ นสเหตุโก ธมฺโม อุปฺปชฺชติ เหตุปจฺจยา.\csromanspacer{}\csromanspacer{}อเหตุกํ ธมฺมํ ปฏิจฺจ นสเหตุโก จ น-อเหตุโก จ ธมฺมา อุปฺปชฺชนฺติ เหตุปจฺจยา. (3)}

\prose{สเหตุกญฺจ อเหตุกญฺจ ธมฺมํ ปฏิจฺจ นสเหตุโก ธมฺโม อุปฺปชฺชติ เหตุปจฺจยา.\csromanspacer{}\csromanspacer{}สเหตุกญฺจ อเหตุกญฺจ ธมฺมํ ปฏิจฺจ น-อเหตุโก ธมฺโม อุปฺปชฺชติ เหตุปจฺจยา.\csromanspacer{}\csromanspacer{}สเหตุกญฺจ อเหตุกญฺจ ธมฺมํ ปฏิจฺจ นสเหตุโก จ น-อเหตุโก จ ธมฺมา อุปฺปชฺชนฺติ เหตุปจฺจยา. (3) (สํขิตฺตํ.)}

\prose{เหตุยา นว, อารมฺมเณ ฉ ฯเปฯ อวิคเต นว. (สหชาตวารมฺปิ ฯเปฯ ปญฺหาวารมฺปิ วิตฺถาเรตพฺพํ.)}

\csromanheader{2}{4. เหตุสเหตุกทุก \textbf{3. เหตุสมฺปยุตฺตทุก 1. ปฏิจฺจวาร}}
\proseitem{5}{\csromanfolio{199}เหตุสมฺปยุตฺตํ ธมฺมํ ปฏิจฺจ นเหตุสมฺปยุตฺโต ธมฺโม อุปฺปชฺชติ เหตุปจฺจยา.\csromanspacer{}\csromanspacer{}เหตุสมฺปยุตฺตํ ธมฺมํ ปฏิจฺจ นเหตุวิปฺปยุตฺโต ธมฺโม อุปฺปชฺชติ เหตุปจฺจยา.\csromanspacer{}\csromanspacer{}เหตุสมฺปยุตฺตํ ธมฺมํ ปฏิจฺจ นเหตุสมฺปยุตฺโต จ นเหตุวิปฺปยุตฺโต จ ธมฺมา อุปฺปชฺชนฺติ เหตุปจฺจยา. (3)}

\prose{เหตุวิปฺปยุตฺตํ ธมฺมํ ปฏิจฺจ นเหตุวิปฺปยุตฺโต ธมฺโม อุปฺปชฺชติ เหตุปจฺจยา.\csromanspacer{}\csromanspacer{}เหตุวิปฺปยุตฺตํ ธมฺมํ ปฏิจฺจ นเหตุสมฺปยุตฺโต ธมฺโม อุปฺปชฺชติ เหตุปจฺจยา.\csromanspacer{}\csromanspacer{}เหตุวิปฺปยุตฺตํ ธมฺมํ ปฏิจฺจ นเหตุสมฺปยุตฺโต จ นเหตุวิปฺปยุตฺโต จ ธมฺมา อุปฺปชฺชนฺติ เหตุปจฺจยา. (3)}

\prose{เหตุสมฺปยุตฺตญฺจ เหตุวิปฺปยุตฺตญฺจ ธมฺมํ ปฏิจฺจ นเหตุสมฺปยุตฺโต ธมฺโม อุปฺปชฺชติ เหตุปจฺจยา.\csromanspacer{}\csromanspacer{}เหตุสมฺปยุตฺตญฺจ เหตุวิปฺปยุตฺตญฺจ ธมฺมํ ปฏิจฺจ นเหตุวิปฺปยุตฺโต ธมฺโม อุปฺปชฺชติ เหตุปจฺจยา.\csromanspacer{}\csromanspacer{}เหตุสมฺปยุตฺตญฺจ เหตุวิปฺปยุตฺตญฺจ ธมฺมํ ปฏิจฺจ นเหตุสมฺปยุตฺโต จ นเหตุวิปฺปยุตฺโต จ ธมฺมา อุปฺปชฺชนฺติ เหตุปจฺจยา. (3) (สํขิตฺตํ.)}

\prose{เหตุยา นว. (สหชาตวารมฺปิ ฯเปฯ ปญฺหาวารมฺปิ วิตฺถาเรตพฺพํ.)}

\csromanheader{2}{4. เหตุสเหตุกทุก 1. ปฏิจฺจวาร}
\proseitem{6}{เหตุญฺเจว สเหตุกญฺจ ธมฺมํ ปฏิจฺจ นเหตุ เจว น-อเหตุโก จ ธมฺโม อุปฺปชฺชติ เหตุปจฺจยา.\csromanspacer{}\csromanspacer{}เหตุญฺเจว สเหตุกญฺจ ธมฺมํ ปฏิจฺจ น-อเหตุโก เจว นนเหตุ จ ธมฺโม อุปฺปชฺชติ เหตุปจฺจยา.\csromanspacer{}\csromanspacer{}เหตุญฺเจว สเหตุกญฺจ ธมฺมํ ปฏิจฺจ นเหตุ เจว น-อเหตุโก จ น-อเหตุโก เจว นนเหตุ จ ธมฺมา อุปฺปชฺชนฺติ เหตุปจฺจยา. (3)}

\prose{สเหตุกญฺเจว น จ เหตุํ ธมฺมํ ปฏิจฺจ น-อเหตุโก เจว นนเหตุ จ ธมฺโม อุปฺปชฺชติ เหตุปจฺจยา.\csromanspacer{}\csromanspacer{}สเหตุกญฺเจว น จ เหตุํ ธมฺมํ ปฏิจฺจ นเหตุ เจว น-อเหตุโก จ ธมฺโม อุปฺปชฺชติ เหตุปจฺจยา.\csromanspacer{}\csromanspacer{}สเหตุกญฺเจว น จ เหตุํ ธมฺมํ ปฏิจฺจ นเหตุ เจว น-อเหตุโก จ น-อเหตุโก เจว นนเหตุ จ ธมฺมา อุปฺปชฺชนฺติ เหตุปจฺจยา. (3)}

\prose{เหตุญฺเจว สเหตุกญฺจ สเหตุกญฺเจว น จ เหตุํ ธมฺมํ ปฏิจฺจ นเหตุ เจว น-อเหตุโก จ ธมฺโม อุปฺปชฺชติ เหตุปจฺจยา.\csromanspacer{}\csromanspacer{}เหตุญฺเจว สเหตุกญฺจ สเหตุกญฺเจว น จ เหตุญฺจ ธมฺมํ ปฏิจฺจ น-อเหตุโก}

\vspace{\tipitakaparskip}
\csromancenter{ธมฺมานุโลมปจฺจนีย ทุกปฏฺฐานปาฬิ เจว นนเหตุ จ ธมฺโม อุปฺปชฺชติ เหตุปจฺจยา.\csromanspacer{}\csromanspacer{}เหตุญฺเจว สเหตุกญฺจ สเหตุกญฺเจว น จ เหตุํ ธมฺมํ ปฏิจฺจ นเหตุ เจว น-อเหตุโก จ น-อเหตุโก เจว นนเหตุ จ ธมฺมา อุปฺปชฺชนฺติ เหตุปจฺจยา. (3)}
\prose{\csromanfolio{200}เหตุยา นว. (สหชาตวารมฺปิ ฯเปฯ ปญฺหาวารมฺปิ วิตฺถาเรตพฺพํ.)}

\csromanheader{2}{5. เหตุเหตุสมฺปยุตฺตทุก 1. ปฏิจฺจวาร}
\proseitem{7}{เหตุญฺเจว เหตุสมฺปยุตฺตญฺจ ธมฺมํ ปฏิจฺจ นเหตุ เจว นเหตุวิปฺปยุตฺโต จ ธมฺโม อุปฺปชฺชติ เหตุปจฺจยา.\csromanspacer{}\csromanspacer{}เหตุญฺเจว เหตุสมฺปยุตฺตญฺจ ธมฺมํ ปฏิจฺจ นเหตุวิปฺปยุตฺโต เจว นนเหตุ จ ธมฺโม อุปฺปชฺชติ เหตุปจฺจยา.\csromanspacer{}\csromanspacer{}เหตุญฺเจว เหตุสมฺปยุตฺตญฺจ ธมฺมํ ปฏิจฺจ นเหตุ เจว นเหตุวิปฺปยุตฺโต จ นเหตุวิปฺปยุตฺโต เจว นนเหตุ จ ธมฺมา อุปฺปชฺชนฺติ เหตุปจฺจยา. (3)}

\prose{เหตุสมฺปยุตฺตญฺเจว น จ เหตุํ ธมฺมํ ปฏิจฺจ นเหตุวิปฺปยุตฺโต เจว นนเหตุ จ ธมฺโม อุปฺปชฺชติ เหตุปจฺจยา.\csromanspacer{}\csromanspacer{}เหตุสมฺปยุตฺตญฺเจว น จ เหตุํ ธมฺมํ ปฏิจฺจ นเหตุ เจว นเหตุวิปฺปยุตฺโต จ ธมฺโม อุปฺปชฺชติ เหตุปจฺจยา.\csromanspacer{}\csromanspacer{}เหตุสมฺปยุตฺตญฺเจว น จ เหตุํ ธมฺมํ ปฏิจฺจ นเหตุ เจว นเหตุวิปฺปยุตฺโต จ นเหตุวิปฺปยุตฺโต เจว นนเหตุ จ ธมฺมา อุปฺปชฺชนฺติ เหตุปจฺจยา. (3)}

\prose{เหตุญฺเจว เหตุสมฺปยุตฺตญฺจ เหตุสมฺปยุตฺตญฺเจว น จ เหตุญฺจ ธมฺมํ ปฏิจฺจ นเหตุ เจว นเหตุวิปฺปยุตฺโต จ ธมฺโม อุปฺปชฺชติ เหตุปจฺจยา.\csromanspacer{}\csromanspacer{}เหตุญฺเจว เหตุสมฺปยุตฺตญฺจ เหตุสมฺปยุตฺตญฺเจว น จ เหตุญฺจ ธมฺมํ ปฏิจฺจ นเหตุวิปฺปยุตฺโต เจว นนเหตุ ธมฺโม จ อุปฺปชฺชติ เหตุปจฺจยา.\csromanspacer{}\csromanspacer{}เหตุญฺเจว เหตุสมฺปยุตฺตญฺจ เหตุสมฺปยุตฺตญฺเจว น จ เหตุญฺจ ธมฺมํ ปฏิจฺจ นเหตุ เจว นเหตุวิปฺปยุตฺโต จ นเหตุวิปฺปยุตฺโต เจว นนเหตุ จ ธมฺมา อุปฺปชฺชนฺติ เหตุปจฺจยา. (3) (สํขิตฺตํ.)}

\prose{เหตุยา นว. (สหชาตวารมฺปิ ฯเปฯ ปญฺหาวารมฺปิ วิตฺถาเรตพฺพํ.)}

\csromanheader{2}{10. สปฺปฏิฆทุก \textbf{6. นเหตุสเหตุกทุก 1. ปฏิจฺจวาร}}
\proseitem{8}{\csromanfolio{201}นเหตุํ สเหตุกํ ธมฺมํ ปฏิจฺจ นเหตุ นสเหตุโก ธมฺโม อุปฺปชฺชติ เหตุปจฺจยา. (สํขิตฺตํ.)}

\nitthitamruled{เหตุยา นว. (สหชาตวารมฺปิ ฯเปฯ ปญฺหาวารมฺปิ วิตฺถาเรตพฺพํ.) เหตุโคจฺฉกํ นิฏฺฐิตํ.}
\csromanheader{2}{7. สปฺปจฺจยทุก 1. ปฏิจฺจาทิวาร}
\proseitem{9}{สปฺปจฺจยํ ธมฺมํ ปฏิจฺจ น-อปฺปจฺจโย ธมฺโม อุปฺปชฺชติ เหตุปจฺจยา. (สํขิตฺตํ.) .\csromanspacer{} เหตุยา เอกํ.}

\proseitem{10}{สปฺปจฺจโย ธมฺโม น-อปฺปจฺจยสฺส ธมฺมสฺส อารมฺมณปจฺจเยน ปจฺจโย.}

\prose{อปฺปจฺจโย ธมฺโม น-อปฺปจฺจยสฺส ธมฺมสฺส อารมฺมณปจฺจเยน ปจฺจโย. (สงฺขตํ สปฺปจฺจยสทิสํ.)}

\csromanheader{2}{9. สนิทสฺสนทุก 1. ปฏิจฺจวาร}
\proseitem{11}{อนิทสฺสนํ ธมฺมํ ปฏิจฺจ น-อนิทสฺสโน ธมฺโม อุปฺปชฺชติ เหตุปจฺจยา.\csromanspacer{}\csromanspacer{}อนิทสฺสนํ ธมฺมํ ปฏิจฺจ นสนิทสฺสโน ธมฺโม อุปฺปชฺชติ เหตุปจฺจยา.\csromanspacer{}\csromanspacer{}อนิทสฺสนํ ธมฺมํ ปฏิจฺจ นสนิทสฺสโน จ น-อนิทสฺสโน จ ธมฺมา อุปฺปชฺชนฺติ เหตุปจฺจยา. (3) (สํขิตฺตํ.)}

\csromanheader{2}{เหตุยา ตีณิ. (สพฺพตฺถ วิตฺถาโร.)}
\csromanheader{2}{10. สปฺปฏิฆทุก 1. ปฏิจฺจวาร}
\proseitem{12}{สปฺปฏิฆํ ธมฺมํ ปฏิจฺจ นสปฺปฏิโฆ ธมฺโม อุปฺปชฺชติ เหตุปจฺจยา.\csromanspacer{}\csromanspacer{}สปฺปฏิฆํ ธมฺมํ ปฏิจฺจ น-อปฺปฏิโฆ ธมฺโม อุปฺปชฺชติ เหตุปจฺจยา.\csromanspacer{}\csromanspacer{}สปฺปฏิฆํ ธมฺมํ ปฏิจฺจ นสปฺปฏิโฆ จ น-อปฺปฏิโฆ จ ธมฺมา อุปฺปชฺชนฺติ เหตุปจฺจยา. (3)}

\prose{อปฺปฏิฆํ ธมฺมํ ปฏิจฺจ น-อปฺปฏิโฆ ธมฺโม อุปฺปชฺชติ เหตุปจฺจยา.\csromanspacer{}\csromanspacer{}อปฺปฏิฆํ ธมฺมํ ปฏิจฺจ นสปฺปฏิโฆ ธมฺโม อุปฺปชฺชติ เหตุปจฺจยา.\csromanspacer{}\csromanspacer{}อปฺปฏิฆํ ธมฺมํ ปฏิจฺจ นสปฺปฏิโฆ จ น-อปฺปฏิโฆ จ ธมฺมา อุปฺปชฺชนฺติ เหตุปจฺจยา. (3)}

\gambhira{ธมฺมานุโลมปจฺจนีย ทุกปฏฺฐานปาฬิ}
\prose{\csromanfolio{202}สปฺปฏิฆญฺจ อปฺปฏิฆญฺจ ธมฺมํ ปฏิจฺจ นสปฺปฏิโฆ ธมฺโม อุปฺปชฺชติ เหตุปจฺจยา.\csromanspacer{}\csromanspacer{}สปฺปฏิฆญฺจ อปฺปฏิฆญฺจ ธมฺมํ ปฏิจฺจ น-อปฺปฏิโฆ ธมฺโม อุปฺปชฺชติ เหตุปจฺจยา.\csromanspacer{}\csromanspacer{}สปฺปฏิฆญฺจ อปฺปฏิฆญฺจ ธมฺมํ ปฏิจฺจ นสปฺปฏิโฆ จ น-อปฺปฏิโฆ จ ธมฺมา อุปฺปชฺชนฺติ เหตุปจฺจยา. (3) (สํขิตฺตํ.)}

\csromanheader{2}{เหตุยา นว. (สพฺพตฺถ วิตฺถาโร.)}
\csromanheader{2}{11. รูปีทุก 1. ปฏิจฺจวาร}
\proseitem{13}{รูปิํ ธมฺมํ ปฏิจฺจ นรูปี ธมฺโม อุปฺปชฺชติ เหตุปจฺจยา.}

\csromanheader{2}{เหตุยา นว. (สพฺพตฺถ วิตฺถาโร.)}
\csromanheader{2}{12. โลกิยทุก 1. ปฏิจฺจวาร}
\csromanheader{2}{14. โลกิยํ ธมฺมํ ปฏิจฺจ นโลกุตฺตโร ธมฺโม อุปฺปชฺชติ}
\prose{โลกุตฺตรํ ธมฺมํ ปฏิจฺจ นโลกิโย ธมฺโม อุปฺปชฺชติ เหตุปจฺจยา.\csromanspacer{}\csromanspacer{}โลกุตฺตรํ ธมฺมํ ปฏิจฺจ นโลกุตฺตโร ธมฺโม อุปฺปชฺชติ เหตุปจฺจยา.\csromanspacer{}\csromanspacer{}(โลกุตฺตรํ ธมฺมํ ปฏิจฺจ นโลกิโย จ นโลกุตฺตโร จ ธมฺมา อุปฺปชฺชนฺติ เหตุปจฺจยา.)* (3)}

\prose{โลกิยญฺจ โลกุตฺตรญฺจ ธมฺมํ ปฏิจฺจ นโลกุตฺตโร ธมฺโม อุปฺปชฺชติ เหตุปจฺจยา. (1) (สํขิตฺตํ.)}

\csromanheader{2}{เหตุยา ปญฺจ. (สพฺพตฺถ ปญฺจ.)}
\csromanheader{2}{13. เกนจิวิญฺเญยฺยทุก 1. ปฏิจฺจวาร}
\proseitem{15}{เกนจิ วิญฺเญยฺยํ ธมฺมํ ปฏิจฺจ นเกนจิ วิญฺเญยฺโย ธมฺโม อุปฺปชฺชติ เหตุปจฺจยา.\csromanspacer{}\csromanspacer{}เกนจิ วิญฺเญยฺยํ ธมฺมํ ปฏิจฺจ นนเกนจิ วิญฺเญยฺโย ธมฺโม อุปฺปชฺชติ เหตุปจฺจยา.\csromanspacer{}\csromanspacer{}เกนจิ วิญฺเญยฺยํ ธมฺมํ ปฏิจฺจ นเกนจิ วิญฺเญยฺโย จ นนเกนจิ วิญฺเญยฺโย จ ธมฺมา อุปฺปชฺชนฺติ เหตุปจฺจยา. (3)}

\prose{นเกนจิ วิญฺเญยฺยํ ธมฺมํ ปฏิจฺจ นนเกนจิ วิญฺเญยฺโย ธมฺโม อุปฺปชฺชติ เหตุปจฺจยา.\csromanspacer{} ตีณิ.}

\csromanheader{2}{15. สาสวทุก}
\prose{\csromanfolio{203}เกนจิ วิญฺเญยฺยญฺจ นเกนจิ วิญฺเญยฺยญฺจ ธมฺมํ ปฏิจฺจ นเกนจิ วิญฺเญยฺโย ธมฺโม อุปฺปชฺชติ เหตุปจฺจยา.\csromanspacer{} ตีณิ. (สํขิตฺตํ.) เหตุยา นว. (สพฺพตฺถ นว.)}

\nitthitam{จูฬนฺตรทุกํ นิฏฺฐิตํ.}
\csromanheader{2}{14. อาสวทุก 1. ปฏิจฺจวาร}
\proseitem{16}{อาสวํ ธมฺมํ ปฏิจฺจ โน-อาสโว ธมฺโม อุปฺปชฺชติ เหตุปจฺจยา.\csromanspacer{}\csromanspacer{}อาสวํ ธมฺมํ ปฏิจฺจ นโน-อาสโว ธมฺโม อุปฺปชฺชติ เหตุปจฺจยา.\csromanspacer{}\csromanspacer{}อาสวํ ธมฺมํ ปฏิจฺจ โน-อาสโว จ นโน-อาสโว จ ธมฺมา อุปฺปชฺชนฺติ เหตุปจฺจยา. (3)}

\prose{โน-อาสวํ ธมฺมํ ปฏิจฺจ นโน-อาสโว ธมฺโม อุปฺปชฺชติ เหตุปจฺจยา.\csromanspacer{}\csromanspacer{}โน-อาสวํ ธมฺมํ ปฏิจฺจ โน-อาสโว ธมฺโม อุปฺปชฺชติ เหตุปจฺจยา.\csromanspacer{}\csromanspacer{}โน-อาสวํ ธมฺมํ ปฏิจฺจ โน-อาสโว จ นโน-อาสโว จ ธมฺมา อุปฺปชฺชนฺติ เหตุปจฺจยา. (3)}

\prose{อาสวญฺจ โน-อาสวญฺจ ธมฺมํ ปฏิจฺจ โน-อาสโว ธมฺโม อุปฺปชฺชติ เหตุปจฺจยา.\csromanspacer{}\csromanspacer{}อาสวญฺจ โน-อาสวญฺจ ธมฺมํ ปฏิจฺจ นโน-อาสโว ธมฺโม อุปฺปชฺชติ เหตุปจฺจยา.\csromanspacer{}\csromanspacer{}อาสวญฺจ โน-อาสวญฺจ ธมฺมํ ปฏิจฺจ โน-อาสโว จ นโน-อาสโว จ ธมฺมา อุปฺปชฺชนฺติ เหตุปจฺจยา. (3) (สํขิตฺตํ.) เหตุยา นว. (สพฺพตฺถ นว.)}

\csromanheader{2}{15. สาสวทุก 1. ปฏิจฺจวาร}
\csromanheader{2}{17. สาสวํ ธมฺมํ ปฏิจฺจ น-อนาสโว ธมฺโม อุปฺปชฺชติ}
\prose{อนาสวํ ธมฺมํ ปฏิจฺจ น-อนาสโว ธมฺโม อุปฺปชฺชติ เหตุปจฺจยา.\csromanspacer{}\csromanspacer{}อนาสวํ ธมฺมํ ปฏิจฺจ นสาสโว ธมฺโม อุปฺปชฺชติ เหตุปจฺจยา.\csromanspacer{}\csromanspacer{}อนาสวํ ธมฺมํ ปฏิจฺจ นสาสโว จ น-อนาสโว จ ธมฺมา อุปฺปชฺชนฺติ เหตุปจฺจยา. (3)}

\prose{สาสวญฺจ อนาสวญฺจ ธมฺมํ ปฏิจฺจ น-อนาสโว ธมฺโม อุปฺปชฺชติ เหตุปจฺจยา. (1) (สํขิตฺตํ.)}

\csromanheader{2}{เหตุยา ปญฺจ. (สพฺพตฺถ วิตฺตาโร.)}
\csromanheader{2}{ธมฺมานุโลมปจฺจนีย ทุกปฏฺฐานปาฬิ \textbf{16. อาสวสมฺปยุตฺตทุก 1. ปฏิจฺจวาร}}
\proseitem{18}{\csromanfolio{204}อาสวสมฺปยุตฺตํ ธมฺมํ ปฏิจฺจ น-อาสวสมฺปยุตฺโต ธมฺโม อุปฺปชฺชติ เหตุปจฺจยา.\csromanspacer{}\csromanspacer{}อาสวสมฺปยุตฺตํ ธมฺมํ ปฏิจฺจ น-อาสววิปฺปยุตฺโต ธมฺโม อุปฺปชฺชติ เหตุปจฺจยา.\csromanspacer{}\csromanspacer{}อาสวสมฺปยุตฺตํ ธมฺมํ ปฏิจฺจ น-อาสวสมฺปยุตฺโต จ น-อาสววิปฺปยุตฺโต จ ธมฺมา อุปฺปชฺชนฺติ เหตุปจฺจยา. (3)}

\prose{อาสววิปฺปยุตฺตํ ธมฺมํ ปฏิจฺจ น-อาสววิปฺปยุตฺโต ธมฺโม อุปฺปชฺชติ เหตุปจฺจยา.\csromanspacer{} อาสววิปฺปยุตฺตํ ธมฺมํ ปฏิจฺจ น-อาสวสมฺปยุตฺโต ธมฺโม อุปฺปชฺชติ เหตุปจฺจยา.\csromanspacer{}\csromanspacer{}อาสววิปฺปยุตฺตํ ธมฺมํ ปฏิจฺจ น-อาสวสมฺปยุตฺโต จ น-อาสววิปฺปยุตฺโต จ ธมฺมา อุปฺปชฺชนฺติ เหตุปจฺจยา. (3)}

\prose{อาสวสมฺปยุตฺตญฺจ อาสววิปฺปยุตฺตญฺจ ธมฺมํ ปฏิจฺจ น-อาสวสมฺปยุตฺโต ธมฺโม อุปฺปชฺชติ เหตุปจฺจยา.\csromanspacer{}\csromanspacer{}อาสวสมฺปยุตฺตญฺจ อาสววิปฺปยุตฺตญฺจ ธมฺมํ ปฏิจฺจ น-อาสววิปฺปยุตฺโต ธมฺโม อุปฺปชฺชติ เหตุปจฺจยา.\csromanspacer{}\csromanspacer{}อาสวสมฺปยุตฺตญฺจ อาสววิปฺปยุตฺตญฺจ ธมฺมํ ปฏิจฺจ น-อาสวสมฺปยุตฺโต จ น-อาสววิปฺปยุตฺโต จ ธมฺมา อุปฺปชฺชนฺติ เหตุปจฺจยา. (3) (สํขิตฺตํ.)}

\csromanheader{2}{เหตุยา นว. (สพฺพตฺถ นว.)}
\csromanheader{2}{17. อาสวสาสวทุก 1. ปฏิจฺจวาร}
\proseitem{19}{อาสวญฺเจว สาสวญฺจ ธมฺมํ ปฏิจฺจ น-อาสโว เจว น-อนาสโว จ ธมฺโม อุปฺปชฺชติ เหตุปจฺจยา.\csromanspacer{}\csromanspacer{}อาสวญฺเจว สาสวญฺจ ธมฺมํ ปฏิจฺจ น-อนาสโว เจว นโน จ อาสโว ธมฺโม อุปฺปชฺชติ เหตุปจฺจยา.\csromanspacer{}\csromanspacer{}อาสวญฺเจว สาสวญฺจ ธมฺมํ ปฏิจฺจ น-อาสโว เจว น-อนาสโว จ น-อนาสโว เจว นโน จ อาสโว จ ธมฺมา อุปฺปชฺชนฺติ เหตุปจฺจยา. (3)}

\prose{สาสวญฺเจว โน จ อาสวํ ธมฺมํ ปฏิจฺจ น-อนาสโว เจว นโน จ อาสโว ธมฺโม อุปฺปชฺชติ เหตุปจฺจยา.\csromanspacer{} ตีณิ.}

\prose{อาสวญฺเจว สาสวญฺจ สาสวญฺเจว โน จ อาสวญฺจ ธมฺมํ ปฏิจฺจ น-อาสโว เจว น-อนาสโว จ ธมฺโม อุปฺปชฺชติ เหตุปจฺจยา.\csromanspacer{} ตีณิ.\csromanspacer{} เหตุยา นว. (สพฺพตฺถ นว.)}

\csromanheader{2}{50. ปรามาสทุก \textbf{18. อาสว-อาสวสมฺปยุตฺตทุก 1. ปฏิจฺจวาร}}
\proseitem{20}{\csromanfolio{205}อาสวญฺเจว อาสวสมฺปยุตฺตญฺจ ธมฺมํ ปฏิจฺจ น-อาสโว เจว น-อาสววิปฺปยุตฺโต จ ธมฺโม อุปฺปชฺชติ เหตุปจฺจยา.\csromanspacer{} ตีณิ.}

\prose{อาสวสมฺปยุตฺตญฺเจว โน จ อาสวํ ธมฺมํ ปฏิจฺจ น-อาสววิปฺปยุตฺโต เจว นโน จ อาสโว ธมฺโม อุปฺปชฺชติ เหตุปจฺจยา.\csromanspacer{} ตีณิ.}

\prose{อาสวญฺเจว อาสวสมฺปยุตฺตญฺจ อาสววิปฺปยุตฺตญฺเจว โน จ อาสวญฺจ ธมฺมํ ปฏิจฺจ โน-อาสโว เจว น-อาสววิปฺปยุตฺโต จ ธมฺโม อุปฺปชฺชติ เหตุปจฺจยา.\csromanspacer{} ตีณิ. (สํขิตฺตํ.) เหตุยา นว. (สพฺพตฺถ นว.)}

\csromanheader{2}{19. อาสววิปฺปยุตฺตสาสวทุก 1. ปฏิจฺจวาร}
\proseitem{21}{อาสววิปฺปยุตฺตํ สาสวํ ธมฺมํ ปฏิจฺจ อาสววิปฺปยุตฺโต น-อนาสโว ธมฺโม อุปฺปชฺชติ เหตุปจฺจยา. (1)}

\prose{อาสววิปฺปยุตฺตํ อนาสวํ ธมฺมํ ปฏิจฺจ อาสววิปฺปยุตฺโต น-อนาสโว ธมฺโม อุปฺปชฺชติ เหตุปจฺจยา.\csromanspacer{} ตีณิ.}

\prose{อาสววิปฺปยุตฺตํ สาสวญฺจ อาสววิปฺปยุตฺตํ อนาสวญฺจ ธมฺมํ ปฏิจฺจ อาสววิปฺปยุตฺโต นโน-อนาสโว ธมฺโม อุปฺปชฺชติ เหตุปจฺจยา. (1)}

\csromanheader{2}{เหตุยา ปญฺจ ฯเปฯ อวิคเต ปญฺจ. (สพฺพตฺถ วิตฺถาโร.)}
\nitthitam{อาสวโคจฺฉกํ นิฏฺฐิตํ.}
\csromanheader{2}{20. สญฺโญชนทุกาทิ 1. ปฏิจฺจวาร}
\proseitem{22}{สญฺโญชนํ ธมฺมํ ปฏิจฺจ โนสญฺโญชโน ธมฺโม อุปฺปชฺชติ เหตุปจฺจยา ฯเปฯ.\csromanspacer{} คนฺถํ ธมฺมํ ปฏิจฺจ โนคนฺโถ ธมฺโม อุปฺปชฺชติ เหตุปจฺจยา ฯเปฯ.\csromanspacer{} โอฆํ ธมฺมํ ปฏิจฺจ โน-โอโฆ ธมฺโม อุปฺปชฺชติ เหตุปจฺจยา ฯเปฯ.\csromanspacer{} โยคํ ธมฺมํ ปฏิจฺจ โนโยโค ธมฺโม อุปฺปชฺชติ เหตุปจฺจยา ฯเปฯ.\csromanspacer{} นีวรณํ ธมฺมํ ปฏิจฺจ โนนีวรโณ ธมฺโม อุปฺปชฺชติ เหตุปจฺจยา.}

\csromanheader{2}{50. ปรามาสทุก 1. ปฏิจฺจวาร}
\proseitem{23}{ปรามาสํ ธมฺมํ ปฏิจฺจ โนปรามาโส ธมฺโม อุปฺปชฺชติ เหตุปจฺจยา. (อาสวโคจฺฉกสทิสํ.)}

\csromanheader{2}{ธมฺมานุโลมปจฺจนีย ทุกปฏฺฐานปาฬิ \textbf{55. สารมฺมณทุก 1. ปฏิจฺจวาร}}
\proseitem{24}{\csromanfolio{206}สารมฺมณํ ธมฺมํ ปฏิจฺจ นสารมฺมโณ ธมฺโม อุปฺปชฺชติ เหตุปจฺจยา.\csromanspacer{} ตีณิ.}

\prose{อนารมฺมณํ ธมฺมํ ปฏิจฺจ น-อนารมฺมโณ ธมฺโม อุปฺปชฺชติ เหตุปจฺจยา.\csromanspacer{} ตีณิ.}

\prose{สารมฺมณญฺจ อนารมฺมณญฺจ ธมฺมํ ปฏิจฺจ นสารมฺมโณ ธมฺโม อุปฺปชฺชติ เหตุปจฺจยา.\csromanspacer{} ตีณิ. (สํขิตฺตํ.)}

\csromanheader{2}{เหตุยา นว ฯเปฯ อวิคเต นว. (สพฺพตฺถ วิตฺถาโร.)}
\csromanheader{2}{56. จิตฺตทุก 1. ปฏิจฺจวาร}
\proseitem{25}{จิตฺตํ ธมฺมํ ปฏิจฺจ โนจิตฺโต ธมฺโม อุปฺปชฺชติ เหตุปจฺจยา.\csromanspacer{} เอกํ.}

\prose{โนจิตฺตํ ธมฺมํ ปฏิจฺจ นโนจิตฺโต ธมฺโม อุปฺปชฺชติ เหตุปจฺจยา.\csromanspacer{} ตีณิ.}

\prose{จิตฺตญฺจ โนจิตฺตญฺจ ธมฺมํ ปฏิจฺจ โนจิตฺโต ธมฺโม อุปฺปชฺชติ เหตุปจฺจยา. (1) (สํขิตฺตํ.\csromanspacer{} ปญฺจ.)}

\csromanheader{2}{57. เจตสิกทุกาทิ 1. ปฏิจฺจวาร}
\proseitem{26}{เจตสิกํ ธมฺมํ ปฏิจฺจ นเจตสิโก ธมฺโม อุปฺปชฺชติ เหตุปจฺจยา.}

\proseitem{27}{จิตฺตสมฺปยุตฺตํ ธมฺมํ ปฏิจฺจ นจิตฺตสมฺปยุตฺโต ธมฺโม อุปฺปชฺชติ เหตุปจฺจยา ฯเปฯ.}

\prose{จิตฺตสํสฏฺฐํ ธมฺมํ ปฏิจฺจ นจิตฺตสํสฏฺโฐ ธมฺโม อุปฺปชฺชติ เหตุปจฺจยา.}

\proseitem{28}{จิตฺตสมุฏฺฐานํ ธมฺมํ ปฏิจฺจ นจิตฺตสมุฏฺฐาโน ธมฺโม อุปฺปชฺชติ เหตุปจฺจยา.}

\proseitem{29}{จิตฺตสหภุํ ธมฺมํ ปฏิจฺจ นจิตฺตสหภู ธมฺโม อุปฺปชฺชติ เหตุปจฺจยา ฯเปฯ.\csromanspacer{} จิตฺตานุปริวตฺติํ ธมฺมํ ปฏิจฺจ นจิตฺตานุปริวตฺตี ธมฺโม อุปฺปชฺชติ เหตุปจฺจยา ฯเปฯ.\csromanspacer{} จิตฺตสํสฏฺฐสมุฏฺฐานํ ธมฺมํ ปฏิจฺจ นจิตฺตสํสฏฺฐสมุฏฺฐาโน ธมฺโม อุปฺปชฺชติ เหตุปจฺจยา ฯเปฯ.\csromanspacer{} จิตฺตสํสฏฺฐสมุฏฺฐานสหภุํ ธมฺมํ ปฏิจฺจ นจิตฺตสํสฏฺฐสมุฏฺฐานสหภู ธมฺโม อุปฺปชฺชติ เหตุปจฺจยา ฯเปฯ.\csromanspacer{} จิตฺตสํสฏฺฐสมุฏฺฐานานุปริวตฺติํ ธมฺมํ ปฏิจฺจ นจิตฺตสํสฏฺฐสมุฏฺฐานานุปริวตฺตี ธมฺโม อุปฺปชฺชติ เหตุปจฺจยา.}

\csromanheader{2}{83. ทสฺสเนนปหาตพฺพทุก}
\proseitem{30}{\csromanfolio{207}อชฺฌตฺติกํ ธมฺมํ ปฏิจฺจ น-อชฺฌตฺติโก ธมฺโม อุปฺปชฺชติ เหตุปจฺจยา.\csromanspacer{} ตีณิ.}

\prose{พาหิรํ ธมฺมํ ปฏิจฺจ นพาหิโร ธมฺโม อุปฺปชฺชติ เหตุปจฺจยา.\csromanspacer{} ตีณิ.}

\prose{อชฺฌตฺติกญฺจ พาหิรญฺจ ธมฺมํ ปฏิจฺจ น-อชฺฌตฺติโก ธมฺโม อุปฺปชฺชติ เหตุปจฺจยา.\csromanspacer{} ตีณิ.}

\proseitem{31}{อุปาทา ธมฺมํ ปฏิจฺจ โน-อุปาทา ธมฺโม อุปฺปชฺชติ เหตุปจฺจยา. (1)}

\prose{โน-อุปาทา ธมฺมํ ปฏิจฺจ นโน-อุปาทา ธมฺโม อุปฺปชฺชติ เหตุปจฺจยา.\csromanspacer{} ตีณิ.}

\prose{อุปาทา จ โน-อุปาทา จ ธมฺมํ ปฏิจฺจ โน-อุปาทา ธมฺโม อุปฺปชฺชติ}

\proseitem{32}{อุปาทินฺนํ ธมฺมํ ปฏิจฺจ น-อุปาทินฺโน ธมฺโม อุปฺปชฺชติ เหตุปจฺจยา.\csromanspacer{} ตีณิ.}

\csromanheader{2}{อนุปาทินฺนํ ธมฺมํ ปฏิจฺจ น-อนุปาทินฺโน ธมฺโม อุปฺปชฺชติ}
\prose{อุปาทินฺนญฺจ อนุปาทินฺนญฺจ ธมฺมํ ปฏิจฺจ น-อุปาทินฺโน ธมฺโม อุปฺปชฺชติ เหตุปจฺจยา. (1) (สํขิตฺตํ.)}

\csromanheader{2}{เหตุยา ปญฺจ ฯเปฯ อวิคเต ปญฺจ. (สพฺพตฺถ วิตฺถาโร.)}
\nitthitam{มหนฺตรทุกํ นิฏฺฐิตํ.}
\csromanheader{2}{69. อุปาทานโคจฺฉก}
\vspace{\tipitakaparskip}
\csromancenter{อุปาทานํ ธมฺมํ ปฏิจฺจ โน-อุปาทาโน ธมฺโม อุปฺปชฺชติ เหตุปจฺจยา.\csromanspacer{} เหตุยา นว.}
\csromanheader{2}{75. กิเลสโคจฺฉก}
\vspace{\tipitakaparskip}
\csromancenter{กิเลสํ ธมฺมํ ปฏิจฺจ โนกิเลโส ธมฺโม อุปฺปชฺชติ เหตุปจฺจยา.\csromanspacer{} เหตุยา นว.}
\csromanheader{2}{83. ทสฺสเนนปหาตพฺพทุก}
\proseitem{35}{ทสฺสเนน ปหาตพฺพํ ธมฺมํ ปฏิจฺจ นทสฺสเนน ปหาตพฺโพ ธมฺโม อุปฺปชฺชติ เหตุปจฺจยา.\csromanspacer{}\csromanspacer{}ทสฺสเนน ปหาตพฺพํ ธมฺมํ ปฏิจฺจ นนทสฺสเนน}

\vspace{\tipitakaparskip}
\csromancenter{ธมฺมานุโลมปจฺจนีย ทุกปฏฺฐานปาฬิ ปหาตพฺโพ ธมฺโม อุปฺปชฺชติ เหตุปจฺจยา.\csromanspacer{}\csromanspacer{}ทสฺสเนน ปหาตพฺพํ ธมฺมํ ปฏิจฺจ นทสฺสเนน ปหาตพฺโพ จ นนทสฺสเนน ปหาตพฺโพ จ ธมฺมา อุปฺปชฺชนฺติ เหตุปจฺจยา. (3)}
\csromanheader{2}{นทสฺสเนน ปหาตพฺพํ ธมฺมํ ปฏิจฺจ นทสฺสเนน ปหาตพฺโพ}
\prose{\csromanfolio{208}ทสฺสเนน ปหาตพฺพญฺจ นทสฺสเนน ปหาตพฺพญฺจ ธมฺมํ ปฏิจฺจ นทสฺสเนน ปหาตพฺโพ ธมฺโม อุปฺปชฺชติ เหตุปจฺจยา. (1) (สํขิตฺตํ.)}

\prose{เหตุยา ปญฺจ ฯเปฯ อวิคเต ปญฺจ.}

\csromanheader{2}{84. ภาวนายปหาตพฺพทุก}
\proseitem{36}{ภาวนาย ปหาตพฺพํ ธมฺมํ ปฏิจฺจ นภาวนาย ปหาตพฺโพ ธมฺโม อุปฺปชฺชติ เหตุปจฺจยา.\csromanspacer{}\csromanspacer{}ภาวนาย ปหาตพฺพํ ธมฺมํ ปฏิจฺจ นนภาวนาย ปหาตพฺโพ ธมฺโม อุปฺปชฺชติ เหตุปจฺจยา.\csromanspacer{}\csromanspacer{}ภาวนาย ปหาตพฺพํ ธมฺมํ ปฏิจฺจ นภาวนาย ปหาตพฺโพ จ นนภาวนาย ปหาตพฺโพ จ ธมฺมา อุปฺปชฺชนฺติ เหตุปจฺจยา. (3)}

\csromanheader{2}{นภาวนาย ปหาตพฺพํ ธมฺมํ ปฏิจฺจ นภาวนาย ปหาตพฺโพ}
\prose{ภาวนาย ปหาตพฺพญฺจ นภาวนาย ปหาตพฺพญฺจ ธมฺมํ ปฏิจฺจ นภาวนาย ปหาตพฺโพ ธมฺโม อุปฺปชฺชติ เหตุปจฺจยา. (1)}

\prose{เหตุยา ปญฺจ.}

\csromanheader{2}{85. ทสฺสเนนปหาตพฺพเหตุกทุก}
\proseitem{37}{ทสฺสเนน ปหาตพฺพเหตุกํ ธมฺมํ ปฏิจฺจ นทสฺสเนน ปหาตพฺพเหตุโก ธมฺโม อุปฺปชฺชติ เหตุปจฺจยา.\csromanspacer{}\csromanspacer{}ทสฺสเนน ปหาตพฺพเหตุกํ ธมฺมํ ปฏิจฺจ นนทสฺสเนน ปหาตพฺพเหตุโก ธมฺโม อุปฺปชฺชติ เหตุปจฺจยา. (สํขิตฺตํ.)}

\prose{เหตุยา นว.}

\csromanheader{2}{89-92. สปฺปีติกาทิทุก \textbf{86. ภาวนายปหาตพฺพเหตุกทุก}}
\proseitem{38}{\csromanfolio{209}ภาวนาย ปหาตพฺพเหตุกํ ธมฺมํ ปฏิจฺจ นภาวนาย ปหาตพฺพเหตุโก ธมฺโม อุปฺปชฺชติ เหตุปจฺจยา.\csromanspacer{}\csromanspacer{}ภาวนาย ปหาตพฺพเหตุกํ ธมฺมํ ปฏิจฺจ นนภาวนาย ปหาตพฺพเหตุโก ธมฺโม อุปฺปชฺชติ เหตุปจฺจยา.\csromanspacer{}\csromanspacer{}เหตุยา นว.}

\csromanheader{2}{87-88. สวิตกฺกาทิทุก}
\proseitem{39}{สวิตกฺกํ ธมฺมํ ปฏิจฺจ นสวิตกฺโก ธมฺโม อุปฺปชฺชติ เหตุปจฺจยา.\csromanspacer{} ตีณิ.}

\prose{อวิตกฺกํ ธมฺมํ ปฏิจฺจ น-อวิตกฺโก ธมฺโม อุปฺปชฺชติ เหตุปจฺจยา. (สํขิตฺตํ.) เหตุยา นว.}

\proseitem{40}{สวิจารํ ธมฺมํ ปฏิจฺจ นสวิจาโร ธมฺโม อุปฺปชฺชติ เหตุปจฺจยา.\csromanspacer{} ตีณิ.}

\prose{อวิจารํ ธมฺมํ ปฏิจฺจ น-อวิจาโร ธมฺโม อุปฺปชฺชติ เหตุปจฺจยา. (สํขิตฺตํ.) เหตุยา นว.}

\csromanheader{2}{89-92. สปฺปีติกาทิทุก}
\proseitem{41}{สปฺปีติกํ ธมฺมํ ปฏิจฺจ นสปฺปีติโก ธมฺโม อุปฺปชฺชติ เหตุปจฺจยา.\csromanspacer{} ตีณิ.}

\prose{อปฺปีติกํ ธมฺมํ ปฏิจฺจ น-อปฺปีติโก ธมฺโม อุปฺปชฺชติ เหตุปจฺจยา.\csromanspacer{}\csromanspacer{}เหตุยา นว.}

\proseitem{42}{ปีติสหคตํ ธมฺมํ ปฏิจฺจ นปีติสหคโต ธมฺโม อุปฺปชฺชติ เหตุปจฺจยา.\csromanspacer{} ตีณิ.}

\prose{นปีติสหคตํ ธมฺมํ ปฏิจฺจ นนปีติสหคโต ธมฺโม อุปฺปชฺชติ เหตุปจฺจยา.\csromanspacer{}\csromanspacer{}เหตุยา นว.}

\gambhira{ธมฺมานุโลมปจฺจนีย ทุกปฏฺฐานปาฬิ}
\proseitem{43}{\csromanfolio{210}สุขสหคตํ ธมฺมํ ปฏิจฺจ นสุขสหคโต ธมฺโม อุปฺปชฺชติ เหตุปจฺจยา.\csromanspacer{} ตีณิ.}

\prose{นสุขสหคตํ ธมฺมํ ปฏิจฺจ นนสุขสหคโต ธมฺโม อุปฺปชฺชติ เหตุปจฺจยา.\csromanspacer{}\csromanspacer{}เหตุยา นว.}

\proseitem{44}{อุเปกฺขาสหคตํ ธมฺมํ ปฏิจฺจ น-อุเปกฺขาสหคโต ธมฺโม อุปฺปชฺชติ เหตุปจฺจยา.\csromanspacer{} ตีณิ.}

\prose{น-อุเปกฺขาสหคตํ ธมฺมํ ปฏิจฺจ นน-อุเปกฺขาสหคโต ธมฺโม อุปฺปชฺชติ เหตุปจฺจยา.\csromanspacer{}\csromanspacer{}เหตุยา นว.}

\csromanheader{2}{93-95. กามาวจราทิทุก}
\proseitem{45}{กามาวจรํ ธมฺมํ ปฏิจฺจ นกามาวจโร ธมฺโม อุปฺปชฺชติ เหตุปจฺจยา.\csromanspacer{} ตีณิ.}

\prose{นกามาวจรํ ธมฺมํ ปฏิจฺจ นนกามาวจโร ธมฺโม อุปฺปชฺชติ เหตุปจฺจยา.\csromanspacer{} เหตุยา นว.}

\proseitem{46}{รูปาวจรํ ธมฺมํ ปฏิจฺจ นรูปาวจโร ธมฺโม อุปฺปชฺชติ เหตุปจฺจยา.\csromanspacer{} ตีณิ.}

\prose{นรูปาวจรํ ธมฺมํ ปฏิจฺจ นนรูปาวจโร ธมฺโม อุปฺปชฺชติ เหตุปจฺจยา.\csromanspacer{} เหตุยา นว.}

\proseitem{47}{อรูปาวจรํ ธมฺมํ ปฏิจฺจ น-อรูปาวจโร ธมฺโม อุปฺปชฺชติ เหตุปจฺจยา.\csromanspacer{} ตีณิ.}

\prose{น-อรูปาวจรํ ธมฺมํ ปฏิจฺจ นน-อรูปาวจโร ธมฺโม อุปฺปชฺชติ เหตุปจฺจยา. (1)}

\prose{อรูปาวจรญฺจ น-อรูปาวจรญฺจ ธมฺมํ ปฏิจฺจ น-อรูปาวจโร ธมฺโม อุปฺปชฺชติ เหตุปจฺจยา. (1).\csromanspacer{} เหตุยา ปญฺจ.}

\csromanheader{2}{99. ส-อุตฺตรทุก \textbf{96. ปริยาปนฺนทุก}}
\proseitem{48}{\csromanfolio{211}ปริยาปนฺนํ ธมฺมํ ปฏิจฺจ น-อปริยาปนฺโน ธมฺโม อุปฺปชฺชติ เหตุปจฺจยา. (1)}

\prose{อปริยาปนฺนํ ธมฺมํ ปฏิจฺจ น-อปริยาปนฺโน ธมฺโม อุปฺปชฺชติ เหตุปจฺจยา.\csromanspacer{}\csromanspacer{}อปริยาปนฺนํ ธมฺมํ ปฏิจฺจ นปริยาปนฺโน ธมฺโม อุปฺปชฺชติ เหตุปจฺจยา.\csromanspacer{}\csromanspacer{}อปริยาปนฺนํ ธมฺมํ ปฏิจฺจ นปริยาปนฺโน จ น-อปริยาปนฺโน จ ธมฺมา อุปฺปชฺชนฺติ เหตุปจฺจยา. (3)}

\prose{ปริยาปนฺนญฺจ อปริยาปนฺนญฺจ ธมฺมํ ปฏิจฺจ น-อปริยาปนฺโน ธมฺโม อุปฺปชฺชติ เหตุปจฺจยา. (1) (สํขิตฺตํ.) .\csromanspacer{} เหตุยา ปญฺจ.}

\csromanheader{2}{97. นิยฺยานิกทุก}
\proseitem{49}{นิยฺยานิกํ ธมฺมํ ปฏิจฺจ นนิยฺยานิโก ธมฺโม อุปฺปชฺชติ เหตุปจฺจยา.\csromanspacer{} ตีณิ.}

\csromanheader{2}{อนิยฺยานิกํ ธมฺมํ ปฏิจฺจ นนิยฺยานิโก ธมฺโม อุปฺปชฺชติ}
\prose{นิยฺยานิกญฺจ อนิยฺยานิกญฺจ ธมฺมํ ปฏิจฺจ นนิยฺยานิโก ธมฺโม อุปฺปชฺชติ เหตุปจฺจยา. (1).\csromanspacer{} เหตุยา ปญฺจ.}

\csromanheader{2}{98. นิยตทุก}
\proseitem{50}{นิยตํ ธมฺมํ ปฏิจฺจ นนิยโต ธมฺโม อุปฺปชฺชติ เหตุปจฺจยา.\csromanspacer{} ตีณิ.}

\prose{อนิยตํ ธมฺมํ ปฏิจฺจ นนิยโต ธมฺโม อุปฺปชฺชติ เหตุปจฺจยา. (1)}

\prose{นิยตญฺจ อนิยตญฺจ ธมฺมํ ปฏิจฺจ นนิยโต ธมฺโม อุปฺปชฺชติ เหตุปจฺจยา. (1).\csromanspacer{} เหตุยา ปญฺจ.}

\csromanheader{2}{99. ส-อุตฺตรทุก}
\csromanheader{2}{51. ส-อุตฺตรํ ธมฺมํ ปฏิจฺจ น-อนุตฺตโร ธมฺโม อุปฺปชฺชติ}
\prose{อนุตฺตรํ ธมฺมํ ปฏิจฺจ น-อนุตฺตโร ธมฺโม อุปฺปชฺชติ เหตุปจฺจยา.\csromanspacer{}\csromanspacer{}อนุตฺตรํ ธมฺมํ ปฏิจฺจ นส-อุตฺตโร ธมฺโม อุปฺปชฺชติ เหตุปจฺจยา.\csromanspacer{}\csromanspacer{}อนุตฺตรํ ธมฺมํ ปฏิจฺจ นส-อุตฺตโร จ น-อนุตฺตโร จ ธมฺมา อุปฺปชฺชนฺติ เหตุปจฺจยา. (3)}

\gambhira{ธมฺมานุโลมปจฺจนีย ทุกปฏฺฐานปาฬิ}
\prose{\csromanfolio{212}ส-อุตฺตรญฺจ อนุตฺตรญฺจ ธมฺมํ ปฏิจฺจ น-อนุตฺตโร ธมฺโม อุปฺปชฺชติ เหตุปจฺจยา. (1).\csromanspacer{} เหตุยา ปญฺจ.}

\csromanheader{2}{100. สรณทุก 1. ปฏิจฺจวาร}
\proseitem{52}{สรณํ ธมฺมํ ปฏิจฺจ นสรโณ ธมฺโม อุปฺปชฺชติ เหตุปจฺจยา.\csromanspacer{}\csromanspacer{}สรณํ ธมฺมํ ปฏิจฺจ น-อรโณ ธมฺโม อุปฺปชฺชติ เหตุปจฺจยา.\csromanspacer{}\csromanspacer{}สรณํ ธมฺมํ ปฏิจฺจ นสรโณ จ น-อรโณ จ ธมฺมา อุปฺปชฺชนฺติ เหตุปจฺจยา. (3)}

\prose{อรณํ ธมฺมํ ปฏิจฺจ น-อรโณ ธมฺโม อุปฺปชฺชติ เหตุปจฺจยา. (1)}

\prose{สรณญฺจ อรณญฺจ ธมฺมํ ปฏิจฺจ นสรโณ ธมฺโม อุปฺปชฺชติ เหตุปจฺจยา. (1).\csromanspacer{} เหตุยา ปญฺจ, อารมฺมเณ เทฺว ฯเปฯ อวิคเต ปญฺจ.}

\csromanheader{2}{\textbf{ปจฺจนีย นเหตุ}}
\prosecont{สรณํ ธมฺมํ ปฏิจฺจ น-อรโณ ธมฺโม อุปฺปชฺชติ นเหตุปจฺจยา. (1)}

\proseitem{53}{อรณํ ธมฺมํ ปฏิจฺจ นน-อรโณ ธมฺโม อุปฺปชฺชติ นเหตุปจฺจยา. (สํขิตฺตํ.)}

\prose{นเหตุยา เทฺว, น-อารมฺมเณ ตีณิ ฯเปฯ โนวิคเต ตีณิ.}

\prose{(สหชาตวารมฺปิ ปจฺจยวารมฺปิ นิสฺสยวารมฺปิ สํสฏฺฐวารมฺปิ สมฺปยุตฺตวารมฺปิ ปฏิจฺจวารสทิสํ วิตฺถาเรตพฺพํ.)}

\csromanheader{2}{100. สรณทุก 7. ปญฺหาวาร}
\proseitem{54}{สรโณ ธมฺโม นสรณสฺส ธมฺมสฺส เหตุปจฺจเยน ปจฺจโย.\csromanspacer{}\csromanspacer{}สรโณ ธมฺโม น-อรณสฺส ธมฺมสฺส เหตุปจฺจเยน ปจฺจโย.\csromanspacer{}\csromanspacer{}สรโณ ธมฺโม นสรณสฺส จ น-อรณสฺส จ ธมฺมสฺส เหตุปจฺจเยน ปจฺจโย. (3)}

\prose{อรโณ ธมฺโม นสรณสฺส ธมฺมสฺส เหตุปจฺจเยน ปจฺจโย. (1)}

\proseitem{55}{สรโณ ธมฺโม นสรณสฺส ธมฺมสฺส อารมฺมณปจฺจเยน ปจฺจโย.\csromanspacer{}\csromanspacer{}สรโณ ธมฺโม น-อรณสฺส ธมฺมสฺส อารมฺมณปจฺจเยน ปจฺจโย. (สํขิตฺตํ.)}

\prose{เหตุยา จตฺตาริ, อารมฺมเณ จตฺตาริ, อธิปติยา ปญฺจ, อนนฺตเร จตฺตาริ ฯเปฯ อวิคเต สตฺต.}

\csromanheader{2}{\textbf{ปจฺจนียุทฺธาร} \textbf{ปจฺจนียุทฺธาร}}
\proseitem{56}{\csromanfolio{213}สรโณ ธมฺโม นสรณสฺส ธมฺมสฺส อารมฺมณปจฺจเยน ปจฺจโย.\csromanspacer{} สหชาตปจฺจเยน ปจฺจโย.\csromanspacer{} อุปนิสฺสยปจฺจเยน ปจฺจโย.\csromanspacer{} ปจฺฉาชาตปจฺจเยน ปจฺจโย.\csromanspacer{} กมฺมปจฺจเยน ปจฺจโย.}

\prose{สรโณ ธมฺโม น-อรณสฺส ธมฺมสฺส อารมฺมณปจฺจเยน ปจฺจโย.\csromanspacer{} สหชาตปจฺจเยน ปจฺจโย.\csromanspacer{} อุปนิสฺสยปจฺจเยน ปจฺจโย. (สํขิตฺตํ.)}

\proseitem{57}{นเหตุยา สตฺต, น-อารมฺมเณ สตฺต, น-อธิปติยา สตฺต, น-อนนฺตเร สตฺต ฯเปฯ โน-อวิคเต จตฺตาริ.}

\csromanheader{2}{เหตุปจฺจยา น-อารมฺมเณ จตฺตาริ. (สํขิตฺตํ.)}
\csromanheader{2}{นเหตุปจฺจยา อารมฺมเณ จตฺตาริ. (สํขิตฺตํ.)}
\prose{(ยถา กุสลตฺติเก ปญฺหาวารํ, เอวํ วิตฺถาเรตพฺพํ.)}

\nitthitam{\textbf{ธมฺมานุโลมปจฺจนีเย ทุกปฏฺฐานํ นิฏฺฐิตํ.}}
\pitaka{\textbf{อภิธมฺมปิฏก}}
\csromanheader{2}{\textbf{ธมฺมานุโลมปจฺจนีย}}
\gambhira{\textbf{ทุกติกปฏฺฐานปาฬิ}}
\namakkaram{นโม ตสฺส ภควโต อรหโต สมฺมาสมฺพุทฺธสฺส.}
\csromanheader{2}{1. เหตุทุก 1. กุสลตฺติก}
\csromanheader{2}{1. กุสลปท 1-7. ปฏิจฺจาทิวาร}
\proseitem{1}{\csromanfolio{215}เหตุํ กุสลํ ธมฺมํ ปฏิจฺจ นเหตุ นกุสโล ธมฺโม อุปฺปชฺชติ เหตุปจฺจยา.\csromanspacer{}\csromanspacer{}นเหตุํ กุสลํ ธมฺมํ ปฏิจฺจ นนเหตุ นกุสโล ธมฺโม อุปฺปชฺชติ เหตุปจฺจยา.\csromanspacer{}\csromanspacer{}เหตุํ กุสลญฺจ นเหตุํ กุสลญฺจ ธมฺมํ ปฏิจฺจ นเหตุ นกุสโล ธมฺโม อุปฺปชฺชติ เหตุปจฺจยา.\csromanspacer{} ตีณิ. (จิตฺตสมุฏฺฐานเมว, อารมฺมณํ นตฺถิ.)}

\prose{เหตุยา ตีณิ, อธิปติยา ตีณิ ฯเปฯ อวิคเต ตีณิ. (สหชาตวารมฺปิ ฯเปฯ นิสฺสยวารมฺปิ ปฏิจฺจวารสทิสํ.)}

\csromanheader{2}{2. เหตุ กุสโล ธมฺโม นเหตุสฺส นกุสลสฺส ธมฺมสฺส}
\prose{เหตุ กุสโล ธมฺโม นเหตุสฺส นกุสลสฺส ธมฺมสฺส อารมฺมณปจฺจเยน ปจฺจโย.\csromanspacer{} ตีณิ.}

\prose{นเหตุ กุสโล ธมฺโม นนเหตุสฺส นกุสลสฺส ธมฺมสฺส อารมฺมณปจฺจเยน ปจฺจโย.\csromanspacer{} ตีณิ.}

\prose{เหตุ กุสโล จ นเหตุ กุสโล จ ธมฺมา นเหตุสฺส นกุสลสฺส ธมฺมสฺส อารมฺมณปจฺจเยน ปจฺจโย.\csromanspacer{} ตีณิ. (สํขิตฺตํ.)}

\csromanheader{2}{ธมฺมานุโลมปจฺจนีย ทุกติกปฏฺฐานปาฬิ}
\proseitem{3}{\csromanfolio{216}เหตุยา เอกํ, อารมฺมเณ นว, อธิปติยา นว ฯเปฯ อวิคเต ตีณิ. (ปญฺหาวารมฺปิ วิตฺถาเรตพฺพํ.)}

\csromanheader{2}{2. อกุสลปท 1-7. ปฏิจฺจาทิวาร}
\proseitem{4}{เหตุํ อกุสลํ ธมฺมํ ปฏิจฺจ นเหตุ น-อกุสโล ธมฺโม อุปฺปชฺชติ เหตุปจฺจยา. (1)}

\prose{นเหตุํ อกุสลํ ธมฺมํ ปฏิจฺจ นเหตุ น-อกุสโล ธมฺโม อุปฺปชฺชติ เหตุปจฺจยา. (1)}

\prose{เหตุํ อกุสลญฺจ นเหตุํ อกุสลญฺจ ธมฺมํ ปฏิจฺจ นเหตุ น-อกุสโล ธมฺโม อุปฺปชฺชติ เหตุปจฺจยา. (1) (สํขิตฺตํ.) .\csromanspacer{} เหตุยา ตีณิ, อธิปติยา ตีณิ ฯเปฯ อวิคเต ตีณิ. (สหชาตวารมฺปิ ฯเปฯ นิสฺสยวารมฺปิ ปฏิจฺจวารสทิสํ.)}

\csromanheader{2}{5. เหตุ อกุสโล ธมฺโม นเหตุสฺส น-อกุสลสฺส ธมฺมสฺส}
\prose{เหตุ อกุสโล ธมฺโม นเหตุสฺส น-อกุสลสฺส ธมฺมสฺส อารมฺมณปจฺจเยน ปจฺจโย.\csromanspacer{} ตีณิ.}

\prose{นเหตุ อกุสโล ธมฺโม นนเหตุสฺส น-อกุสลสฺส ธมฺมสฺส อารมฺมณปจฺจเยน ปจฺจโย.\csromanspacer{} ตีณิ.}

\prose{เหตุ อกุสโล จ นเหตุ อกุสโล จ ธมฺมา นเหตุสฺส น-อกุสลสฺส ธมฺมสฺส อารมฺมณปจฺจเยน ปจฺจโย.\csromanspacer{} ตีณิ. (สํขิตฺตํ.)}

\proseitem{6}{เหตุยา เอกํ, อารมฺมเณ นว, อธิปติยา เอกํ ฯเปฯ อวิคเต ตีณิ.}

\csromanheader{2}{3. อพฺยากตปท 3. ปจฺจยวาร}
\proseitem{7}{นเหตุํ อพฺยากตํ ธมฺมํ ปจฺจยา นนเหตุ น-อพฺยากโต ธมฺโม อุปฺปชฺชติ เหตุปจฺจยา.\csromanspacer{}\csromanspacer{}เหตุยา ตีณิ. (นิสฺสยวารมฺปิ ฯเปฯ ปญฺหาวารมฺปิ วิตฺถาเรตพฺพํ.)}

\csromanheader{2}{1. เหตุทุก 4. อุปาทินฺนตฺติก \textbf{1. เหตุทุก 2. เวทนาตฺติก}}
\proseitem{8}{\csromanfolio{217}เหตุํ สุขาย เวทนาย สมฺปยุตฺตํ ธมฺมํ ปฏิจฺจ นเหตุ นสุขาย เวทนาย สมฺปยุตฺโต ธมฺโม อุปฺปชฺชติ เหตุปจฺจยา. (1)}

\prose{นเหตุํ สุขาย เวทนาย สมฺปยุตฺตํ ธมฺมํ ปฏิจฺจ นเหตุ นสุขาย เวทนาย สมฺปยุตฺโต ธมฺโม อุปฺปชฺชติ เหตุปจฺจยา. (1)}

\prose{เหตุํ สุขาย เวทนาย สมฺปยุตฺตญฺจ นเหตุํ สุขาย เวทนาย สมฺปยุตฺตญฺจ ธมฺมํ ปฏิจฺจ นเหตุ นสุขาย เวทนาย สมฺปยุตฺโต ธมฺโม อุปฺปชฺชติ เหตุปจฺจยา. (1) (สํขิตฺตํ.) .\csromanspacer{} เหตุยา ตีณิ, อารมฺมเณ ตีณิ ฯเปฯ อวิคเต ตีณิ. (สหชาตวารมฺปิ ฯเปฯ ปญฺหาวารมฺปิ วิตฺถาเรตพฺพํ.)}

\proseitem{9}{เหตุํ ทุกฺขาย เวทนาย สมฺปยุตฺตํ ธมฺมํ ปฏิจฺจ นเหตุ นทุกฺขาย เวทนาย สมฺปยุตฺโต ธมฺโม อุปฺปชฺชติ เหตุปจฺจยา. (สํขิตฺตํ.) .\csromanspacer{} เหตุยา ตีณิ. (สพฺพตฺถ วิตฺถาโร.)}

\proseitem{10}{เหตุํ อทุกฺขมสุขาย เวทนาย สมฺปยุตฺตํ ธมฺมํ ปฏิจฺจ นเหตุ น-อทุกฺขมสุขาย เวทนาย สมฺปยุตฺโต ธมฺโม อุปฺปชฺชติ เหตุปจฺจยา. (สํขิตฺตํ.) .\csromanspacer{} เหตุยา ตีณิ. (สพฺพตฺถ วิตฺถาโร.)}

\csromanheader{2}{1. เหตุทุก 3. วิปากตฺติก}
\proseitem{11}{เหตุํ วิปากํ ธมฺมํ ปฏิจฺจ นเหตุ นวิปาโก ธมฺโม อุปฺปชฺชติ เหตุปจฺจยา.\csromanspacer{}\csromanspacer{}เหตุยา ตีณิ.}

\prose{เหตุํ วิปากธมฺมธมฺมํ ปฏิจฺจ นเหตุ นวิปากธมฺมธมฺโม อุปฺปชฺชติ เหตุปจฺจยา.\csromanspacer{}\csromanspacer{}เหตุยา ตีณิ.}

\prose{นเหตุํ เนววิปากนวิปากธมฺมธมฺมํ ปฏิจฺจ นนเหตุ นเนววิปากนวิปากธมฺมธมฺโม อุปฺปชฺชติ เหตุปจฺจยา.\csromanspacer{}\csromanspacer{}เหตุยา ตีณิ.}

\csromanheader{2}{1. เหตุทุก 4. อุปาทินฺนตฺติก}
\proseitem{12}{เหตุํ อุปาทินฺนุปาทานิยํ ธมฺมํ ปฏิจฺจ นเหตุ น-อุปาทินฺนุปาทานิโย ธมฺโม อุปฺปชฺชติ เหตุปจฺจยา.\csromanspacer{}\csromanspacer{}เหตุยา ตีณิ.}

\csromanheader{2}{ธมฺมานุโลมปจฺจนีย ทุกติกปฏฺฐานปาฬิ}
\prose{\csromanfolio{218}นเหตุํ อนุปาทินฺนุปาทานิยํ ธมฺมํ ปฏิจฺจ นนเหตุ อนุปาทินฺนุปาทานิโย ธมฺโม อุปฺปชฺชติ เหตุปจฺจยา.}

\prose{เหตุํ อนุปาทินฺน-อนุปาทานิยํ ธมฺมํ ปฏิจฺจ นเหตุ น-อนุปาทินฺน-อนุปาทานิโย ธมฺโม อุปฺปชฺชติ เหตุปจฺจยา.\csromanspacer{}\csromanspacer{}เหตุยา ตีณิ.}

\csromanheader{2}{1. เหตุทุก 5. สํกิลิฏฺฐตฺติก}
\proseitem{13}{เหตุํ สํกิลิฏฺฐสํกิเลสิกํ ธมฺมํ ปฏิจฺจ นเหตุ นสํกิลิฏฺฐสํกิเลสิโก ธมฺโม อุปฺปชฺชติ เหตุปจฺจยา.\csromanspacer{}\csromanspacer{}เหตุยา ตีณิ.}

\prose{นเหตุํ อสํกิลิฏฺฐสํกิเลสิกํ ธมฺมํ ปจฺจยา นนเหตุ น-อสํกิลิฏฺฐสํกิเลสิโก ธมฺโม อุปฺปชฺชติ เหตุปจฺจยา.\csromanspacer{}\csromanspacer{}เหตุยา ตีณิ.}

\prose{เหตุํ อสํกิลิฏฺฐ-อสํกิเลสิกํ ธมฺมํ ปฏิจฺจ นเหตุ น-อสํกิลิฏฺฐ-อสํกิเลสิโก ธมฺโม อุปฺปชฺชติ เหตุปจฺจยา.\csromanspacer{}\csromanspacer{}เหตุยา ตีณิ.}

\csromanheader{2}{1. เหตุทุก 6. วิตกฺกตฺติก}
\proseitem{14}{เหตุํ สวิตกฺกสวิจารํ ธมฺมํ ปฏิจฺจ นเหตุ นสวิตกฺกสวิจาโร ธมฺโม อุปฺปชฺชติ เหตุปจฺจยา.\csromanspacer{}\csromanspacer{}เหตุยา ตีณิ.}

\prose{เหตุํ อวิตกฺกวิจารมตฺตํ ธมฺมํ ปฏิจฺจ นเหตุ น-อวิตกฺกวิจารมตฺโต ธมฺโม อุปฺปชฺชติ เหตุปจฺจยา.\csromanspacer{}\csromanspacer{}เหตุยา ตีณิ.}

\prose{นเหตุํ อวิตกฺก-อวิจารํ ธมฺมํ ปฏิจฺจ นนเหตุ น-อวิตกฺก-อวิจาโร ธมฺโม อุปฺปชฺชติ เหตุปจฺจยา.\csromanspacer{}\csromanspacer{}เหตุยา ตีณิ.}

\csromanheader{2}{1. เหตุทุก 7. ปีติตฺติก}
\proseitem{15}{เหตุํ ปีติสหคตํ ธมฺมํ ปฏิจฺจ นเหตุ นปีติสหคโต ธมฺโม อุปฺปชฺชติ เหตุปจฺจยา.\csromanspacer{}\csromanspacer{}เหตุยา ตีณิ.}

\prose{เหตุํ สุขสหคตํ ธมฺมํ ปฏิจฺจ นเหตุ นสุขสหคโต ธมฺโม อุปฺปชฺชติ เหตุปจฺจยา.\csromanspacer{}\csromanspacer{}เหตุยา ตีณิ.}

\prose{เหตุํ อุเปกฺขาสหคตํ ธมฺมํ ปฏิจฺจ นเหตุ น-อุเปกฺขาสหคโต ธมฺโม อุปฺปชฺชติ เหตุปจฺจยา.\csromanspacer{}\csromanspacer{}เหตุยา ตีณิ.}

\csromanheader{2}{1. เหตุทุก 11. เสกฺขตฺติก \textbf{1. เหตุทุก 8. ทสฺสนตฺติก}}
\proseitem{16}{\csromanfolio{219}เหตุํ ทสฺสเนน ปหาตพฺพํ ธมฺมํ ปฏิจฺจ นเหตุ นทสฺสเนน ปหาตพฺโพ ธมฺโม อุปฺปชฺชติ เหตุปจฺจยา.\csromanspacer{}\csromanspacer{}เหตุยา ตีณิ.}

\prose{เหตุํ ภาวนาย ปหาตพฺพํ ธมฺมํ ปฏิจฺจ นเหตุ นภาวนาย ปหาตพฺโพ ธมฺโม อุปฺปชฺชติ เหตุปจฺจยา.\csromanspacer{}\csromanspacer{}เหตุยา ตีณิ.}

\prose{นเหตุํ เนวทสฺสเนน นภาวนาย ปหาตพฺพํ ธมฺมํ ปจฺจยา นนเหตุ นเนวทสฺสเนน นภาวนาย ปหาตพฺโพ ธมฺโม อุปฺปชฺชติ เหตุปจฺจยา.\csromanspacer{}\csromanspacer{}เหตุยา ตีณิ.}

\csromanheader{2}{1. เหตุทุก 9. ทสฺสนเหตุตฺติก}
\proseitem{17}{เหตุํ ทสฺสเนน ปหาตพฺพเหตุกํ ธมฺมํ ปฏิจฺจ นเหตุ นทสฺสเนน ปหาตพฺพเหตุโก ธมฺโม อุปฺปชฺชติ เหตุปจฺจยา.\csromanspacer{}\csromanspacer{}เหตุยา ตีณิ.}

\prose{เหตุํ ภาวนาย ปหาตพฺพเหตุกํ ธมฺมํ ปฏิจฺจ นเหตุ นภาวนาย ปหาตพฺพเหตุโก ธมฺโม อุปฺปชฺชติ เหตุปจฺจยา.\csromanspacer{}\csromanspacer{}เหตุยา ตีณิ.}

\prose{เหตุํ เนวทสฺสเนน นภาวนาย ปหาตพฺพเหตุกํ ธมฺมํ ปฏิจฺจ นเหตุ นเนวทสฺสเนน นภาวนาย ปหาตพฺพเหตุโก ธมฺโม อุปฺปชฺชติ เหตุปจฺจยา.\csromanspacer{}\csromanspacer{}เหตุยา เอกํ.}

\csromanheader{2}{1. เหตุทุก 10. อาจยคามิตฺติก}
\proseitem{18}{เหตุํ อาจยคามิํ ธมฺมํ ปฏิจฺจ นเหตุ น-อาจยคามี ธมฺโม อุปฺปชฺชติ เหตุปจฺจยา.\csromanspacer{}\csromanspacer{}เหตุยา ตีณิ.}

\prose{เหตุํ อปจยคามิํ ธมฺมํ ปฏิจฺจ นเหตุ น-อปจยคามี ธมฺโม อุปฺปชฺชติ เหตุปจฺจยา.\csromanspacer{}\csromanspacer{}เหตุยา ตีณิ.}

\prose{นเหตุํ เนวาจยคามินาปจยคามิํ ธมฺมํ ปจฺจยา นนเหตุ นเนวาจยคามินาปจยคามี ธมฺโม อุปฺปชฺชติ เหตุปจฺจยา.\csromanspacer{}\csromanspacer{}เหตุยา ตีณิ.}

\csromanheader{2}{1. เหตุทุก 11. เสกฺขตฺติก}
\proseitem{19}{เหตุํ เสกฺขํ ธมฺมํ ปฏิจฺจ นเหตุ นเสกฺโข ธมฺโม อุปฺปชฺชติ เหตุปจฺจยา.\csromanspacer{}\csromanspacer{}เหตุยา ตีณิ.}

\csromanheader{2}{ธมฺมานุโลมปจฺจนีย ทุกติกปฏฺฐานปาฬิ}
\prose{\csromanfolio{220}เหตุํ อเสกฺขํ ธมฺมํ ปฏิจฺจ นเหตุ น-อเสกฺโข ธมฺโม อุปฺปชฺชติ เหตุปจฺจยา.\csromanspacer{}\csromanspacer{}เหตุยา ตีณิ.}

\prose{นเหตุํ เนวเสกฺขนาเสกฺขํ ธมฺมํ ปจฺจยา นนเหตุ นเนวเสกฺขนาเสกฺโข ธมฺโม อุปฺปชฺชติ เหตุปจฺจยา.\csromanspacer{}\csromanspacer{}เหตุยา ตีณิ.}

\csromanheader{2}{1. เหตุทุก 12. ปริตฺตตฺติก}
\proseitem{20}{นเหตุํ ปริตฺตํ ธมฺมํ ปฏิจฺจ นนเหตุ นปริตฺโต ธมฺโม อุปฺปชฺชติ เหตุปจฺจยา. (ตีณิ.)}

\prose{เหตุํ มหคฺคตํ ธมฺมํ ปฏิจฺจ นเหตุ นมหคฺคโต ธมฺโม อุปฺปชฺชติ เหตุปจฺจยา.\csromanspacer{}\csromanspacer{}เหตุยา ตีณิ.}

\prose{เหตุํ อปฺปมาณํ ธมฺมํ ปฏิจฺจ นเหตุ น-อปฺปมาโณ ธมฺโม อุปฺปชฺชติ เหตุปจฺจยา.\csromanspacer{}\csromanspacer{}เหตุยา ตีณิ.}

\csromanheader{2}{1. เหตุทุก 13. ปริตฺตารมฺมณตฺติก}
\proseitem{21}{เหตุํ ปริตฺตารมฺมณํ ธมฺมํ ปฏิจฺจ นเหตุ นปริตฺตารมฺมโณ ธมฺโม อุปฺปชฺชติ เหตุปจฺจยา. (ตีณิ.)}

\prose{เหตุํ มหคฺคตารมฺมณํ ธมฺมํ ปฏิจฺจ นเหตุ นมหคฺคตารมฺมโณ ธมฺโม อุปฺปชฺชติ เหตุปจฺจยา.\csromanspacer{} ตีณิ.}

\prose{เหตุํ อปฺปมาณารมฺมณํ ธมฺมํ ปฏิจฺจ นเหตุ น-อปฺปมาณารมฺมโณ ธมฺโม อุปฺปชฺชติ เหตุปจฺจยา.\csromanspacer{}\csromanspacer{}เหตุยา ตีณิ.}

\csromanheader{2}{1. เหตุทุก 14. หีนตฺติก}
\proseitem{22}{เหตุํ หีนํ ธมฺมํ ปฏิจฺจ นเหตุ นหีโน ธมฺโม อุปฺปชฺชติ เหตุปจฺจยา.\csromanspacer{}\csromanspacer{}เหตุยา ตีณิ.}

\prose{นเหตุํ มชฺฌิมํ ธมฺมํ ปจฺจยา นนเหตุ นมชฺฌิโม ธมฺโม อุปฺปชฺชติ เหตุปจฺจยา.\csromanspacer{}\csromanspacer{}เหตุยา ตีณิ.}

\prose{เหตุํ ปณีตํ ธมฺมํ ปฏิจฺจ นเหตุ นปณีโต ธมฺโม อุปฺปชฺชติ เหตุปจฺจยา.\csromanspacer{}\csromanspacer{}เหตุยา ตีณิ.}

\csromanheader{2}{1. เหตุทุก 18. อตีตตฺติก \textbf{1. เหตุทุก 15. มิจฺฉตฺตนิยตตฺติก}}
\proseitem{23}{\csromanfolio{221}เหตุํ มิจฺฉตฺตนิยตํ ธมฺมํ ปฏิจฺจ นเหตุ นมิจฺฉตฺตนิยโต ธมฺโม อุปฺปชฺชติ เหตุปจฺจยา. (ตีณิ.)}

\prose{เหตุํ สมฺมตฺตนิยตํ ธมฺมํ ปฏิจฺจ นเหตุ นสมฺมตฺตนิยโต ธมฺโม อุปฺปชฺชติ เหตุปจฺจยา.\csromanspacer{} ตีณิ.}

\prose{นเหตุํ อนิยตํ ธมฺมํ ปจฺจยา นนเหตุ น-อนิยโต ธมฺโม อุปฺปชฺชติ เหตุปจฺจยา.\csromanspacer{}\csromanspacer{}เหตุยา ตีณิ.}

\csromanheader{2}{1. เหตุทุก 16. มคฺคารมฺมณตฺติก}
\proseitem{24}{เหตุํ มคฺคารมฺมณํ ธมฺมํ ปฏิจฺจ นเหตุ นมคฺคารมฺมโณ ธมฺโม อุปฺปชฺชติ เหตุปจฺจยา.\csromanspacer{} ตีณิ.}

\prose{เหตุํ มคฺคเหตุกํ ธมฺมํ ปฏิจฺจ นเหตุ นมคฺคเหตุโก ธมฺโม อุปฺปชฺชติ เหตุปจฺจยา.\csromanspacer{} ตีณิ.}

\prose{เหตุํ มคฺคาธิปติํ ธมฺมํ ปฏิจฺจ นเหตุ นมคฺคาธิปติ ธมฺโม อุปฺปชฺชติ เหตุปจฺจยา.\csromanspacer{}\csromanspacer{}เหตุยา ตีณิ.}

\csromanheader{2}{1. เหตุทุก 17. อุปฺปนฺนตฺติก}
\proseitem{25}{เหตุ อนุปฺปนฺโน ธมฺโม นเหตุสฺส น-อนุปฺปนฺนสฺส ธมฺมสฺส อารมฺมณปจฺจเยน ปจฺจโย.\csromanspacer{}\csromanspacer{}อารมฺมเณ นว, อธิปติยา อุปนิสฺสเย นว.}

\prose{เหตุ อุปฺปาที ธมฺโม นเหตุสฺส น-อุปฺปาทิสฺส ธมฺมสฺส อารมฺมณปจฺจเยน ปจฺจโย.\csromanspacer{}\csromanspacer{}อารมฺมเณ นว.}

\csromanheader{2}{1. เหตุทุก 18. อตีตตฺติก}
\proseitem{26}{เหตุ อตีโต ธมฺโม นเหตุสฺส น-อตีตสฺส ธมฺมสฺส อารมฺมณปจฺจเยน ปจฺจโย.\csromanspacer{}\csromanspacer{}อารมฺมเณ นว.}

\prose{เหตุ อนาคโต ธมฺโม นเหตุสฺส น-อนาคตสฺส ธมฺมสฺส อารมฺมณปจฺจเยน ปจฺจโย.\csromanspacer{}\csromanspacer{}อารมฺมเณ นว.}

\csromanheader{2}{ธมฺมานุโลมปจฺจนีย ทุกติกปฏฺฐานปาฬิ \textbf{1. เหตุทุก 19. อตีตารมฺมณตฺติก}}
\proseitem{27}{\csromanfolio{222}เหตุํ อตีตารมฺมณํ ธมฺมํ ปฏิจฺจ นเหตุ น-อตีตารมฺมโณ ธมฺโม อุปฺปชฺชติ เหตุปจฺจยา.\csromanspacer{}\csromanspacer{}เหตุยา ตีณิ.}

\prose{เหตุํ อนาคตารมฺมณํ ธมฺมํ ปฏิจฺจ นเหตุ น-อนาคตารมฺมโณ ธมฺโม อุปฺปชฺชติ เหตุปจฺจยา.\csromanspacer{}\csromanspacer{}เหตุยา ตีณิ.}

\prose{เหตุํ ปจฺจุปฺปนฺนารมฺมณํ ธมฺมํ ปฏิจฺจ นเหตุ นปจฺจุปฺปนฺนารมฺมโณ ธมฺโม อุปฺปชฺชติ เหตุปจฺจยา.\csromanspacer{}\csromanspacer{}เหตุยา ตีณิ.}

\csromanheader{2}{1. เหตุทุก 20. อชฺฌตฺตารมฺมณตฺติก}
\proseitem{28}{เหตุํ อชฺฌตฺตารมฺมณํ ธมฺมํ ปฏิจฺจ นเหตุ น-อชฺฌตฺตารมฺมโณ ธมฺโม อุปฺปชฺชติ เหตุปจฺจยา. (ตีณิ.)}

\prose{เหตุํ พหิทฺธารมฺมณํ ธมฺมํ ปฏิจฺจ นเหตุ นพหิทฺธารมฺมโณ ธมฺโม อุปฺปชฺชติ เหตุปจฺจยา.\csromanspacer{}\csromanspacer{}เหตุยา ตีณิ.}

\csromanheader{2}{1. เหตุทุก 21. สนิทสฺสนตฺติก}
\proseitem{29}{นเหตุ สนิทสฺสนสปฺปฏิโฆ ธมฺโม นนเหตุสฺส นสนิทสฺสนสปฺปฏิฆสฺส ธมฺมสฺส อารมฺมณปจฺจเยน ปจฺจโย ฯเปฯ.\csromanspacer{} นเหตุ สนิทสฺสนสปฺปฏิโฆ ธมฺโม นเหตุสฺส นสนิทสฺสนสปฺปฏิฆสฺส จ นนเหตุสฺส นสนิทสฺสนสปฺปฏิฆสฺส จ ธมฺมสฺส อารมฺมณปจฺจเยน ปจฺจโย.\csromanspacer{}\csromanspacer{}อารมฺมเณ ตีณิ, อธิปติยา อุปนิสฺสเย ปุเรชาเต อตฺถิยา อวิคเต ตีณิ.}

\prose{นเหตุํ อนิทสฺสนสปฺปฏิฆํ ธมฺมํ ปฏิจฺจ นเหตุ น-อนิทสฺสนสปฺปฏิโฆ ธมฺโม อุปฺปชฺชติ เหตุปจฺจยา.\csromanspacer{}\csromanspacer{}เหตุยา เอกํ ฯเปฯ อวิคเต เอกํ.}

\prose{เหตุํ อนิทสฺสน-อปฺปฏิฆํ ธมฺมํ ปฏิจฺจ นเหตุ น-อนิทสฺสน-อปฺปฏิโฆ ธมฺโม อุปฺปชฺชติ เหตุปจฺจยา.\csromanspacer{}\csromanspacer{}นเหตุํ อนิทสฺสน-อปฺปฏิฆํ ธมฺมํ ปฏิจฺจ นเหตุ น-อนิทสฺสน-อปฺปฏิโฆ ธมฺโม อุปฺปชฺชติ เหตุปจฺจยา.\csromanspacer{}\csromanspacer{}เหตุํ อนิทสฺสน-อปฺปฏิฆญฺจ นเหตุํ อนิทสฺสน-อปฺปฏิฆญฺจ ธมฺมํ ปฏิจฺจ นเหตุ น-อนิทสฺสน-อปฺปฏิโฆ ธมฺโม อุปฺปชฺชติ เหตุปจฺจยา. (3)}

\prose{เหตุยา ตีณิ.}

\vspace{\tipitakaparskip}
\csromancenter{เหตุสเหตุกทุก 1.\csromanspacer{} กุสลตฺติก \textbf{2.}\csromanspacer{}\textbf{ สเหตุกทุก 1.}\csromanspacer{}\textbf{ กุสลตฺติก}}
\proseitem{30}{\csromanfolio{223}สเหตุกํ กุสลํ ธมฺมํ ปฏิจฺจ นสเหตุโก นกุสโล ธมฺโม อุปฺปชฺชติ เหตุปจฺจยา.\csromanspacer{}\csromanspacer{}เหตุยา เอกํ.}

\prose{สเหตุกํ อกุสลํ ธมฺมํ ปฏิจฺจ นสเหตุโก น-อกุสโล ธมฺโม อุปฺปชฺชติ เหตุปจฺจยา.\csromanspacer{}\csromanspacer{}เหตุยา ตีณิ.}

\prose{อเหตุกํ อพฺยากตํ ธมฺมํ ปจฺจยา น-อเหตุโก น-อพฺยากโต ธมฺโม อุปฺปชฺชติ เหตุปจฺจยา.\csromanspacer{}\csromanspacer{}เหตุยา เอกํ.}

\csromanheader{2}{3. เหตุสมฺปยุตฺตทุก 1. กุสลตฺติก}
\proseitem{31}{เหตุสมฺปยุตฺตํ กุสลํ ธมฺมํ ปฏิจฺจ นเหตุสมฺปยุตฺโต นกุสโล ธมฺโม อุปฺปชฺชติ เหตุปจฺจยา.\csromanspacer{}\csromanspacer{}เหตุยา เอกํ.}

\prose{เหตุสมฺปยุตฺตํ อกุสลํ ธมฺมํ ปฏิจฺจ นเหตุสมฺปยุตฺโต น-อกุสโล ธมฺโม อุปฺปชฺชติ เหตุปจฺจยา.\csromanspacer{}\csromanspacer{}เหตุยา ตีณิ.}

\prose{เหตุวิปฺปยุตฺตํ อพฺยากตํ ธมฺมํ ปจฺจยา นเหตุวิปฺปยุตฺโต น-อพฺยากโต ธมฺโม อุปฺปชฺชติ เหตุปจฺจยา.\csromanspacer{}\csromanspacer{}เหตุยา เอกํ.}

\csromanheader{2}{4. เหตุสเหตุกทุก 1. กุสลตฺติก}
\proseitem{32}{เหตุ เจว สเหตุโก จ กุสโล ธมฺโม นเหตุสฺส เจว น-อเหตุกสฺส จ นกุสลสฺส ธมฺมสฺส อารมฺมณปจฺจเยน ปจฺจโย.\csromanspacer{}\csromanspacer{}เหตุ เจว สเหตุโก จ กุสโล ธมฺโม น-อเหตุกสฺส เจว นน จ เหตุสฺส นกุสลสฺส ธมฺมสฺส อารมฺมณปจฺจเยน ปจฺจโย.\csromanspacer{}\csromanspacer{}เหตุ เจว สเหตุโก จ กุสโล ธมฺโม นเหตุสฺส เจว น-อเหตุกสฺส จ นกุสลสฺส จ น-อเหตุกสฺส เจว นน จ เหตุสฺส นกุสลสฺส จ ธมฺมสฺส อารมฺมณปจฺจเยน ปจฺจโย. (3)}

\prose{สเหตุโก เจว น จ เหตุ กุสโล ธมฺโม น-อเหตุกสฺส เจว นน จ เหตุสฺส นกุสลสฺส ธมฺมสฺส อารมฺมณปจฺจเยน ปจฺจโย.\csromanspacer{}\csromanspacer{}สเหตุโก เจว น จ เหตุ กุสโล ธมฺโม นเหตุสฺส เจว น-อเหตุกสฺส จ นกุสลสฺส ธมฺมสฺส อารมฺมณปจฺจเยน ปจฺจโย.\csromanspacer{}\csromanspacer{}}

\vspace{\tipitakaparskip}
\csromancenter{ธมฺมานุโลมปจฺจนีย ทุกติกปฏฺฐานปาฬิ สเหตุโก เจว น จ เหตุ กุสโล ธมฺโม นเหตุสฺส เจว น-อเหตุกสฺส จ นกุสลสฺส จ น-อเหตุกสฺส เจว นน จ เหตุสฺส นกุสลสฺส จ ธมฺมสฺส อารมฺมณปจฺจเยน ปจฺจโย. (3)}
\prose{\csromanfolio{224}เหตุ เจว สเหตุโก กุสโล จ สเหตุโก เจว น จ เหตุ กุสโล จ ธมฺมา นเหตุสฺส เจว น-อเหตุกสฺส นกุสลสฺส ธมฺมสฺส อารมฺมณปจฺจเยน ปจฺจโย.\csromanspacer{}\csromanspacer{}เหตุ เจว สเหตุโก กุสโล จ สเหตุโก เจว น จ เหตุ กุสโล จ ธมฺมา น-อเหตุกสฺส เจว นน จ เหตุสฺส นกุสลสฺส ธมฺมสฺส อารมฺมณปจฺจเยน ปจฺจโย.\csromanspacer{}\csromanspacer{}เหตุ เจว สเหตุโก กุสโล จ สเหตุโก เจว น จ เหตุ กุสโล จ ธมฺมา นเหตุสฺส เจว น-อเหตุกสฺส นกุสลสฺส จ น-อเหตุกสฺส เจว นน เหตุสฺส นกุสลสฺส จ ธมฺมสฺส อารมฺมณปจฺจเยน ปจฺจโย. (3)}

\prose{อารมฺมเณ นว.}

\proseitem{33}{เหตุ เจว สเหตุโก จ อกุสโล ธมฺโม นเหตุสฺส เจว น-อเหตุกสฺส น-อกุสลสฺส จ ธมฺมสฺส อารมฺมณปจฺจเยน ปจฺจโย. (เอเตน อุปาเยน นว ปญฺหา กาตพฺพา.)}

\proseitem{34}{เหตุ เจว สเหตุโก จ อพฺยากโต ธมฺโม นเหตุสฺส เจว น-อเหตุกสฺส จ น-อพฺยากตสฺส ธมฺมสฺส อารมฺมณปจฺจเยน ปจฺจโย. (นว ปญฺหา กาตพฺพา.) (สํขิตฺตํ.) (เหตุเหตุสมฺปยุตฺตทุกํ เหตุสเหตุกทุกสทิสํ.\csromanspacer{} สํขิตฺตํ.\csromanspacer{} นว ปญฺหา.)}

\csromanheader{2}{6. นเหตุสเหตุกทุก 1. กุสลตฺติก}
\proseitem{35}{นเหตุํ สเหตุกํ กุสลํ ธมฺมํ ปฏิจฺจ นเหตุ นสเหตุโก นกุสโล ธมฺโม อุปฺปชฺชติ เหตุปจฺจยา.\csromanspacer{}\csromanspacer{}เหตุยา เอกํ.}

\prose{นเหตุํ สเหตุกํ อกุสลํ ธมฺมํ ปฏิจฺจ นเหตุ นสเหตุโก น-อกุสโล ธมฺโม อุปฺปชฺชติ เหตุปจฺจยา.\csromanspacer{}\csromanspacer{}เหตุยา เอกํ.}

\prose{นเหตุํ อเหตุกํ อพฺยากตํ ธมฺมํ ปจฺจยา นเหตุ น-อเหตุโก น-อพฺยากโต ธมฺโม อุปฺปชฺชติ เหตุปจฺจยา.\csromanspacer{}\csromanspacer{}เหตุยา เอกํ.}

\nitthitam{เหตุโคจฺฉกํ นิฏฺฐิตํ.}
\csromanheader{2}{12. โลกิยทุก 1. กุสลตฺติก \textbf{7. สปฺปจฺจยทุก 1. กุสลตฺติก}}
\proseitem{36}{\csromanfolio{225}สปฺปจฺจยํ กุสลํ ธมฺมํ ปฏิจฺจ น-อปฺปจฺจโย นกุสโล ธมฺโม อุปฺปชฺชติ เหตุปจฺจยา.\csromanspacer{}\csromanspacer{}เหตุยา เอกํ.}

\prose{สปฺปจฺจยํ อกุสลํ ธมฺมํ ปฏิจฺจ น-อปฺปจฺจโย น-อกุสโล ธมฺโม อุปฺปชฺชติ เหตุปจฺจยา.\csromanspacer{}\csromanspacer{}เหตุยา เอกํ.}

\prose{สปฺปจฺจยํ อพฺยากตํ ธมฺมํ ปจฺจยา น-อปฺปจฺจโย น-อพฺยากโต ธมฺโม อุปฺปชฺชติ เหตุปจฺจยา.\csromanspacer{}\csromanspacer{}เหตุยา เอกํ ฯเปฯ. (สํขิตฺตํ.\csromanspacer{} สปฺปจฺจยสทิสํ.)}

\csromanheader{2}{9. สนิทสฺสนทุก 1. กุสลตฺติก}
\proseitem{37}{อนิทสฺสนํ กุสลํ ธมฺมํ ปฏิจฺจ น-อนิทสฺสโน นกุสโล ธมฺโม อุปฺปชฺชติ เหตุปจฺจยา.\csromanspacer{}\csromanspacer{}เหตุยา ตีณิ. (อกุสลํ กุสลสทิสํ.)}

\prose{อนิทสฺสนํ อพฺยากตํ ธมฺมํ ปจฺจยา นสนิทสฺสโน น-อพฺยากโต ธมฺโม อุปฺปชฺชติ เหตุปจฺจยา.\csromanspacer{}\csromanspacer{}เหตุยา เอกํ.}

\prose{อปฺปฏิฆํ กุสลํ ธมฺมํ ปฏิจฺจ น-อปฺปฏิโฆ นกุสโล ธมฺโม อุปฺปชฺชติ เหตุปจฺจยา.\csromanspacer{}\csromanspacer{}เหตุยา ตีณิ.}

\prose{อปฺปฏิฆํ อกุสลํ ธมฺมํ ปฏิจฺจ น-อปฺปฏิโฆ น-อกุสโล ธมฺโม อุปฺปชฺชติ เหตุปจฺจยา.\csromanspacer{}\csromanspacer{}เหตุยา ตีณิ.}

\prose{อพฺยากเต เอกํ.}

\csromanheader{2}{11. รูปีทุก 1. กุสลตฺติก}
\proseitem{38}{อรูปิํ กุสลํ ธมฺมํ ปฏิจฺจ น-อรูปี นกุสโล ธมฺโม อุปฺปชฺชติ เหตุปจฺจยา.\csromanspacer{}\csromanspacer{}เหตุยา เอกํ.}

\prose{อรูปิํ อกุสลํ ธมฺมํ ปฏิจฺจ น-อรูปี น-อกุสโล ธมฺโม อุปฺปชฺชติ เหตุปจฺจยา.\csromanspacer{}\csromanspacer{}เหตุยา เอกํ.}

\prose{รูปิํ อพฺยากตํ ธมฺมํ ปจฺจยา นรูปี น-อพฺยากโต ธมฺโม อุปฺปชฺชติ เหตุปจฺจยา.\csromanspacer{}\csromanspacer{}เหตุยา เอกํ.}

\csromanheader{2}{12. โลกิยทุก 1. กุสลตฺติก}
\proseitem{39}{โลกิยํ กุสลํ ธมฺมํ ปฏิจฺจ นโลกุตฺตโร นกุสโล ธมฺโม อุปฺปชฺชติ เหตุปจฺจยา.\csromanspacer{}\csromanspacer{}เหตุยา เทฺว.}

\csromanheader{2}{ธมฺมานุโลมปจฺจนีย ทุกติกปฏฺฐานปาฬิ}
\prose{\csromanfolio{226}โลกิยํ อกุสลํ ธมฺมํ ปฏิจฺจ นโลกุตฺตโร น-อกุสโล ธมฺโม อุปฺปชฺชติ เหตุปจฺจยา.\csromanspacer{}\csromanspacer{}เหตุยา เอกํ.}

\prose{โลกิยํ อพฺยากตํ ธมฺมํ ปจฺจยา นโลกิโย น-อพฺยากโต ธมฺโม อุปฺปชฺชติ เหตุปจฺจยา.\csromanspacer{}\csromanspacer{}เหตุยา เทฺว.}

\csromanheader{2}{13. เกนจิวิญฺเญยฺยทุก 1. กุสลตฺติก}
\proseitem{40}{เกนจิ วิญฺเญยฺยํ กุสลํ ธมฺมํ ปฏิจฺจ นเกนจิ วิญฺเญยฺโย นกุสโล ธมฺโม อุปฺปชฺชติ เหตุปจฺจยา.\csromanspacer{}\csromanspacer{}เหตุยา นว.}

\prose{เกนจิ วิญฺเญยฺยํ อกุสลํ ธมฺมํ ปฏิจฺจ นเกนจิ วิญฺเญยฺโย น-อกุสโล ธมฺโม อุปฺปชฺชติ เหตุปจฺจยา.\csromanspacer{}\csromanspacer{}เหตุยา นว.}

\prose{เกนจิ วิญฺเญยฺยํ อพฺยากตํ ธมฺมํ ปจฺจยา นเกนจิ วิญฺเญยฺโย น-อพฺยากโต ธมฺโม อุปฺปชฺชติ เหตุปจฺจยา.\csromanspacer{}\csromanspacer{}เหตุยา นว.}

\nitthitam{จูฬนฺตรทุกํ นิฏฺฐิตํ.}
\csromanheader{2}{14. อาสวทุก 1. กุสลตฺติก}
\proseitem{41}{โน-อาสวํ กุสลํ ธมฺมํ ปฏิจฺจ โน-อาสโว นกุสโล ธมฺโม อุปฺปชฺชติ เหตุปจฺจยา.\csromanspacer{}\csromanspacer{}เหตุยา เอกํ.}

\prose{อาสวํ อกุสลํ ธมฺมํ ปฏิจฺจ โน-อาสโว น-อกุสโล ธมฺโม อุปฺปชฺชติ เหตุปจฺจยา.\csromanspacer{}\csromanspacer{}เหตุยา ตีณิ.}

\prose{โน-อาสวํ อพฺยากตํ ธมฺมํ ปจฺจยา นโน-อาสโว น-อพฺยากโต ธมฺโม อุปฺปชฺชติ เหตุปจฺจยา.\csromanspacer{}\csromanspacer{}เหตุยา ตีณิ.}

\csromanheader{2}{15. สาสวทุก 1. กุสลตฺติก}
\proseitem{42}{สาสวํ กุสลํ ธมฺมํ ปฏิจฺจ น-อนาสโว นกุสโล ธมฺโม อุปฺปชฺชติ เหตุปจฺจยา.\csromanspacer{}\csromanspacer{}อนาสวํ กุสลํ ธมฺมํ ปฏิจฺจ น-อนาสโว นกุสโล ธมฺโม อุปฺปชฺชติ เหตุปจฺจยา.\csromanspacer{}\csromanspacer{}เหตุยา เทฺว.}

\prose{สาสวํ อกุสลํ ธมฺมํ ปฏิจฺจ น-อนาสโว น-อกุสโล ธมฺโม อุปฺปชฺชติ เหตุปจฺจยา.\csromanspacer{}\csromanspacer{}เหตุยา เอกํ.}

\csromanheader{2}{19. อาสววิปฺปยุตฺตสาสวทุก 1. กุสลตฺติก}
\prose{\csromanfolio{227}สาสวํ อพฺยากตํ ธมฺมํ ปจฺจยา นสาสโว น-อพฺยากโต ธมฺโม อุปฺปชฺชติ เหตุปจฺจยา.\csromanspacer{}\csromanspacer{}เหตุยา เทฺว.}

\csromanheader{2}{16. อาสวสมฺปยุตฺตทุก 1. กุสลตฺติก}
\proseitem{43}{อาสววิปฺปยุตฺตํ กุสลํ ธมฺมํ ปฏิจฺจ น-อาสวสมฺปยุตฺโต นกุสโล ธมฺโม อุปฺปชฺชติ เหตุปจฺจยา.\csromanspacer{}\csromanspacer{}เหตุยา เอกํ.}

\prose{อาสวสมฺปยุตฺตํ อกุสลํ ธมฺมํ ปฏิจฺจ น-อาสวสมฺปยุตฺโต น-อกุสโล ธมฺโม อุปฺปชฺชติ เหตุปจฺจยา.\csromanspacer{}\csromanspacer{}เหตุยา ตีณิ.}

\prose{อาสววิปฺปยุตฺตํ อพฺยากตํ ธมฺมํ ปจฺจยา น-อาสวสมฺปยุตฺโต น-อพฺยากโต ธมฺโม อุปฺปชฺชติ เหตุปจฺจยา.\csromanspacer{}\csromanspacer{}เหตุยา ตีณิ.}

\csromanheader{2}{17. อาสวสาสวทุก 1. กุสลตฺติก}
\proseitem{44}{สาสวญฺเจว โน จ อาสวํ กุสลํ ธมฺมํ ปฏิจฺจ น-อาสโว เจว น-อนาสโว จ นกุสโล ธมฺโม อุปฺปชฺชติ เหตุปจฺจยา.\csromanspacer{}\csromanspacer{}เหตุยา เอกํ.}

\prose{อาสวญฺเจว สาสวญฺจ อกุสลํ ธมฺมํ ปฏิจฺจ โน-อาสโว เจว น-อนาสโว จ น-อกุสโล ธมฺโม อุปฺปชฺชติ เหตุปจฺจยา.\csromanspacer{}\csromanspacer{}เหตุยา ตีณิ.}

\prose{สาสวญฺเจว โน จ อาสวํ อพฺยากตํ ธมฺมํ ปจฺจยา น-อนาสโว เจว นโน จ อาสโว น-อพฺยากโต ธมฺโม อุปฺปชฺชติ เหตุปจฺจยา.\csromanspacer{}\csromanspacer{}เหตุยา ตีณิ. (อาสว-อาสวสมฺปยุตฺตทุกํ อาสวสาสวทุกสทิสํ.)}

\csromanheader{2}{19. อาสววิปฺปยุตฺตสาสวทุก 1. กุสลตฺติก}
\proseitem{45}{อาสววิปฺปยุตฺตํ สาสวํ กุสลํ ธมฺมํ ปฏิจฺจ อาสววิปฺปยุตฺโต น-อนาสโว นกุสโล ธมฺโม อุปฺปชฺชติ เหตุปจฺจยา.\csromanspacer{} เอกํ.}

\prose{อาสววิปฺปยุตฺตํ อนาสวํ กุสลํ ธมฺมํ ปฏิจฺจ อาสววิปฺปยุตฺโต น-อนาสโว นกุสโล ธมฺโม อุปฺปชฺชติ เหตุปจฺจยา.\csromanspacer{} เอกํ.}

\csromanheader{2}{ธมฺมานุโลมปจฺจนีย ทุกติกปฏฺฐานปาฬิ}
\prose{\csromanfolio{228}อาสววิปฺปยุตฺตํ สาสวํ อกุสลํ ธมฺมํ ปฏิจฺจ อาสววิปฺปยุตฺโต น-อนาสโว น-อกุสโล ธมฺโม อุปฺปชฺชติ เหตุปจฺจยา.\csromanspacer{}\csromanspacer{}เหตุยา เอกํ.}

\prose{อาสววิปฺปยุตฺตํ สาสวํ อพฺยากตํ ธมฺมํ ปจฺจยา อาสววิปฺปยุตฺโต นสาสโว น-อพฺยากโต ธมฺโม อุปฺปชฺชติ เหตุปจฺจยา.\csromanspacer{}\csromanspacer{}อาสววิปฺปยุตฺตํ สาสวํ อพฺยากตํ ธมฺมํ ปจฺจยา อาสววิปฺปยุตฺโต น-อนาสโว น-อพฺยากโต ธมฺโม อุปฺปชฺชติ เหตุปจฺจยา.\csromanspacer{}\csromanspacer{}เหตุยา เทฺว.}

\nitthitam{อาสวโคจฺฉกํ นิฏฺฐิตํ.}
\csromanheader{2}{20. สญฺโญชนาทิทุก 1. กุสลตฺติก}
\proseitem{46}{โนสญฺโญชนํ กุสลํ ธมฺมํ ปฏิจฺจ โนสญฺโญชโน นกุสโล ธมฺโม อุปฺปชฺชติ เหตุปจฺจยา ฯเปฯ.\csromanspacer{} โนคนฺถํ กุสลํ ธมฺมํ ปฏิจฺจ โนคนฺโถ นกุสโล ธมฺโม อุปฺปชฺชติ เหตุปจฺจยา ฯเปฯ.\csromanspacer{} โน-โอฆํ กุสลํ ธมฺมํ ปฏิจฺจ โน-โอโฆ นกุสโล ธมฺโม อุปฺปชฺชติ เหตุปจฺจยา ฯเปฯ.\csromanspacer{} โนโยคํ กุสลํ ธมฺมํ ปฏิจฺจ โนโยโค นกุสโล ธมฺโม อุปฺปชฺชติ เหตุปจฺจยา ฯเปฯ.\csromanspacer{} โนนีวรณํ กุสลํ ธมฺมํ ปฏิจฺจ โนนีวรโณ นกุสโล ธมฺโม อุปฺปชฺชติ เหตุปจฺจยา ฯเปฯ.\csromanspacer{} โนปรามาสํ กุสลํ ธมฺมํ ปฏิจฺจ โนปรามาโส นกุสโล ธมฺโม อุปฺปชฺชติ เหตุปจฺจยา.}

\csromanheader{2}{55. สารมฺมณทุก 1. กุสลตฺติก}
\proseitem{47}{สารมฺมณํ กุสลํ ธมฺมํ ปฏิจฺจ นสารมฺมโณ นกุสโล ธมฺโม อุปฺปชฺชติ เหตุปจฺจยา.\csromanspacer{}\csromanspacer{}เหตุยา เอกํ.}

\prose{สารมฺมณํ อกุสลํ ธมฺมํ ปฏิจฺจ นสารมฺมโณ น-อกุสโล ธมฺโม อุปฺปชฺชติ เหตุปจฺจยา.\csromanspacer{}\csromanspacer{}เหตุยา เอกํ.}

\prose{อนารมฺมณํ อพฺยากตํ ธมฺมํ ปจฺจยา น-อนารมฺมโณ น-อพฺยากโต ธมฺโม อุปฺปชฺชติ เหตุปจฺจยา.\csromanspacer{}\csromanspacer{}เหตุยา เอกํ.}

\csromanheader{2}{56. จิตฺตทุก 1. กุสลตฺติก}
\proseitem{48}{จิตฺตํ กุสลํ ธมฺมํ ปฏิจฺจ นจิตฺโต นกุสโล ธมฺโม อุปฺปชฺชติ เหตุปจฺจยา.\csromanspacer{}\csromanspacer{}เหตุยา ตีณิ.}

\csromanheader{2}{60. จิตฺตสมุฏฺฐานทุก 1. กุสลตฺติก}
\prose{\csromanfolio{229}จิตฺตํ อกุสลํ ธมฺมํ ปฏิจฺจ นจิตฺโต น-อกุสโล ธมฺโม อุปฺปชฺชติ เหตุปจฺจยา.\csromanspacer{}\csromanspacer{}เหตุยา ตีณิ.}

\prose{โนจิตฺตํ อพฺยากตํ ธมฺมํ ปจฺจยา นโนจิตฺโต น-อพฺยากโต ธมฺโม อุปฺปชฺชติ เหตุปจฺจยา.\csromanspacer{}\csromanspacer{}เหตุยา ตีณิ.}

\csromanheader{2}{57. เจตสิกทุก 1. กุสลตฺติก}
\proseitem{49}{เจตสิกํ กุสลํ ธมฺมํ ปฏิจฺจ นเจตสิโก นกุสโล ธมฺโม อุปฺปชฺชติ เหตุปจฺจยา.\csromanspacer{}\csromanspacer{}เหตุยา ตีณิ.}

\prose{เจตสิกํ อกุสลํ ธมฺมํ ปฏิจฺจ นเจตสิโก น-อกุสโล ธมฺโม อุปฺปชฺชติ เหตุปจฺจยา.\csromanspacer{}\csromanspacer{}เหตุยา ตีณิ.}

\prose{อเจตสิกํ อพฺยากตํ ธมฺมํ ปจฺจยา น-อเจตสิโก น-อพฺยากโต ธมฺโม อุปฺปชฺชติ เหตุปจฺจยา.\csromanspacer{}\csromanspacer{}เหตุยา ตีณิ.}

\csromanheader{2}{58. จิตฺตสมฺปยุตฺตทุก 1. กุสลตฺติก}
\proseitem{50}{จิตฺตสมฺปยุตฺตํ กุสลํ ธมฺมํ ปฏิจฺจ นจิตฺตสมฺปยุตฺโต นกุสโล ธมฺโม อุปฺปชฺชติ เหตุปจฺจยา.\csromanspacer{}\csromanspacer{}เหตุยา เอกํ.}

\prose{จิตฺตสมฺปยุตฺตํ อกุสลํ ธมฺมํ ปฏิจฺจ นจิตฺตสมฺปยุตฺโต น-อกุสโล ธมฺโม อุปฺปชฺชติ เหตุปจฺจยา.\csromanspacer{}\csromanspacer{}เหตุยา เอกํ. (อพฺยากตมูลํ ตีณิเยว ปจฺจยวเสน.)}

\csromanheader{2}{59. จิตฺตสํสฏฺฐทุก 1. กุสลตฺติก}
\proseitem{51}{จิตฺตสํสฏฺฐํ กุสลํ ธมฺมํ ปฏิจฺจ นจิตฺตสํสฏฺโฐ นกุสโล ธมฺโม อุปฺปชฺชติ เหตุปจฺจยา.\csromanspacer{}\csromanspacer{}เหตุยา เอกํ.}

\prose{จิตฺตสํสฏฺฐํ อกุสลํ ธมฺมํ ปฏิจฺจ นจิตฺตสํสฏฺโฐ น-อกุสโล ธมฺโม อุปฺปชฺชติ เหตุปจฺจยา.\csromanspacer{}\csromanspacer{}เหตุยา เอกํ. (อพฺยากตมูลํ ตีณิเยว ปจฺจยวเสน.)}

\csromanheader{2}{60. จิตฺตสมุฏฺฐานทุก 1. กุสลตฺติก}
\proseitem{52}{จิตฺตสมุฏฺฐานํ กุสลํ ธมฺมํ ปฏิจฺจ นโนจิตฺตสมุฏฺฐาโน นกุสโล ธมฺโม อุปฺปชฺชติ เหตุปจฺจยา.\csromanspacer{}\csromanspacer{}เหตุยา ตีณิ.}

\csromanheader{2}{ธมฺมานุโลมปจฺจนีย ทุกติกปฏฺฐานปาฬิ}
\prose{\csromanfolio{230}จิตฺตสมุฏฺฐานํ อกุสลํ ธมฺมํ ปฏิจฺจ นโนจิตฺตสมุฏฺฐาโน น-อกุสโล ธมฺโม อุปฺปชฺชติ เหตุปจฺจยา.\csromanspacer{}\csromanspacer{}เหตุยา ตีณิ. (อพฺยากตมูลํ ตีณิเยว ปจฺจยวเสน.)}

\csromanheader{2}{61. จิตฺตสหภูทุก 1. กุสลตฺติก}
\proseitem{53}{จิตฺตสหภุํ กุสลํ ธมฺมํ ปฏิจฺจ โนจิตฺตสหภู นกุสโล ธมฺโม อุปฺปชฺชติ เหตุปจฺจยา.\csromanspacer{}\csromanspacer{}เหตุยา นว.}

\prose{จิตฺตสหภุํ อกุสลํ ธมฺมํ ปฏิจฺจ โนจิตฺตสหภู น-อกุสโล ธมฺโม อุปฺปชฺชติ เหตุปจฺจยา.\csromanspacer{}\csromanspacer{}เหตุยา นว. (อพฺยากตมูลํ ตีณิเยว ปจฺจยวเสน.)}

\csromanheader{2}{62. จิตฺตานุปริวตฺติทุก 1. กุสลตฺติก}
\proseitem{54}{จิตฺตานุปริวตฺติํ กุสลํ ธมฺมํ ปฏิจฺจ โนจิตฺตานุปริวตฺตี นกุสโล ธมฺโม อุปฺปชฺชติ เหตุปจฺจยา.\csromanspacer{}\csromanspacer{}เหตุยา นว.}

\prose{จิตฺตานุปริวตฺติํ อกุสลํ ธมฺมํ ปฏิจฺจ โนจิตฺตานุปริวตฺตี น-อกุสโล ธมฺโม อุปฺปชฺชติ เหตุปจฺจยา.\csromanspacer{}\csromanspacer{}เหตุยา นว. (อพฺยากตมูลํ ตีณิเยว ปจฺจยวเสน.)}

\csromanheader{2}{63. จิตฺตสํสฏฺฐสมุฏฺฐานทุก 1. กุสลตฺติก}
\proseitem{55}{จิตฺตสํสฏฺฐสมุฏฺฐานํ กุสลํ ธมฺมํ ปฏิจฺจ โนจิตฺตสํสฏฺฐสมุฏฺฐาโน นกุสโล ธมฺโม อุปฺปชฺชติ เหตุปจฺจยา.\csromanspacer{}\csromanspacer{}เหตุยา ตีณิ.}

\prose{จิตฺตสํสฏฺฐสมุฏฺฐานํ อกุสลํ ธมฺมํ ปฏิจฺจ โนจิตฺตสํสฏฺฐสมุฏฺฐาโน น-อกุสโล ธมฺโม อุปฺปชฺชติ เหตุปจฺจยา.\csromanspacer{}\csromanspacer{}เหตุยา ตีณิ. (อพฺยากตมูลํ ตีณิเยว ปจฺจยวเสน.)}

\csromanheader{2}{64. จิตฺตสํสฏฺฐสมุฏฺฐานสหภูทุก 1. กุสลตฺติก}
\proseitem{56}{จิตฺตสํสฏฺฐสมุฏฺฐานสหภุํ กุสลํ ธมฺมํ ปฏิจฺจ โนจิตฺตสํสฏฺฐสมุฏฺฐานสหภู นกุสโล ธมฺโม อุปฺปชฺชติ เหตุปจฺจยา.\csromanspacer{}\csromanspacer{}เหตุยา ตีณิ.}

\prose{จิตฺตสํสฏฺฐสมุฏฺฐานสหภุํ อกุสลํ ธมฺมํ ปฏิจฺจ โนจิตฺตสํสฏฺฐสมุฏฺฐานสหภู น-อกุสโล ธมฺโม อุปฺปชฺชติ เหตุปจฺจยา.\csromanspacer{}\csromanspacer{}เหตุยา ตีณิ. (อพฺยากตมูลํ ตีณิเยว ปจฺจยวเสน.)}

\vspace{\tipitakaparskip}
\csromancenter{อุปาทินฺนทุก 1.\csromanspacer{} กุสลตฺติก \textbf{65.}\csromanspacer{}\textbf{ จิตฺตสํสฏฺฐสมุฏฺฐานานุปริวตฺติทุก 1.}\csromanspacer{}\textbf{ กุสลตฺติก}}
\proseitem{57}{\csromanfolio{231}จิตฺตสํสฏฺฐสมุฏฺฐานานุปริวตฺติํ กุสลํ ธมฺมํ ปฏิจฺจ โนจิตฺตสํสฏฺฐสมุฏฺฐานานุปริวตฺตี นกุสโล ธมฺโม อุปฺปชฺชติ เหตุปจฺจยา.\csromanspacer{}\csromanspacer{}เหตุยา ตีณิ.}

\prose{จิตฺตสํสฏฺฐสมุฏฺฐานานุปริวตฺติํ อกุสลํ ธมฺมํ ปฏิจฺจ โนจิตฺตสํสฏฺฐสมุฏฺฐานานุปริวตฺตี น-อกุสโล ธมฺโม อุปฺปชฺชติ เหตุปจฺจยา.\csromanspacer{}\csromanspacer{}เหตุยา ตีณิ. (อพฺยากตมูลํ ตีณิเยว.)}

\csromanheader{2}{66. อชฺฌตฺติกทุก 1. กุสลตฺติก}
\proseitem{58}{อชฺฌตฺติกํ กุสลํ ธมฺมํ ปฏิจฺจ น-อชฺฌตฺติโก นกุสโล ธมฺโม อุปฺปชฺชติ เหตุปจฺจยา.\csromanspacer{}\csromanspacer{}พาหิรํ กุสลํ ธมฺมํ ปฏิจฺจ น-อชฺฌตฺติโก นกุสโล ธมฺโม อุปฺปชฺชติ เหตุปจฺจยา.\csromanspacer{}\csromanspacer{}อชฺฌตฺติกํ กุสลญฺจ พาหิรํ กุสลญฺจ ธมฺมํ ปฏิจฺจ น-อชฺฌตฺติโก นกุสโล ธมฺโม อุปฺปชฺชติ เหตุปจฺจยา.\csromanspacer{}\csromanspacer{}เหตุยา ตีณิ.}

\prose{อชฺฌตฺติกํ อกุสลํ ธมฺมํ ปฏิจฺจ น-อชฺฌตฺติโก น-อกุสโล ธมฺโม อุปฺปชฺชติ เหตุปจฺจยา.\csromanspacer{}\csromanspacer{}พาหิรํ อกุสลํ ธมฺมํ ปฏิจฺจ น-อชฺฌตฺติโก น-อกุสโล ธมฺโม อุปฺปชฺชติ เหตุปจฺจยา.\csromanspacer{}\csromanspacer{}อชฺฌตฺติกํ อกุสลญฺจ พาหิรํ อกุสลญฺจ ธมฺมํ ปฏิจฺจ น-อชฺฌตฺติโก น-อกุสโล ธมฺโม อุปฺปชฺชติ เหตุปจฺจยา.\csromanspacer{}\csromanspacer{}เหตุยา ตีณิ. (อพฺยากตมูลํ ตีณิเยว.)}

\csromanheader{2}{67. อุปาทาทุก 1. กุสลตฺติก}
\proseitem{59}{โน-อุปาทา กุสลํ ธมฺมํ ปฏิจฺจ นโน-อุปาทา นกุสโล ธมฺโม อุปฺปชฺชติ เหตุปจฺจยา.\csromanspacer{}\csromanspacer{}เหตุยา ตีณิ.}

\prose{โน-อุปาทา อกุสลํ ธมฺมํ ปฏิจฺจ นโน-อุปาทา น-อกุสโล ธมฺโม อุปฺปชฺชติ เหตุปจฺจยา.\csromanspacer{}\csromanspacer{}เหตุยา ตีณิ. (อพฺยากตมูลํ เอกํเยว.)}

\csromanheader{2}{68. อุปาทินฺนทุก 1. กุสลตฺติก}
\proseitem{60}{อนุปาทินฺนํ กุสลํ ธมฺมํ ปฏิจฺจ น-อุปาทินฺโน นกุสโล ธมฺโม อุปฺปชฺชติ เหตุปจฺจยา.\csromanspacer{}\csromanspacer{}เหตุยา เอกํ.}

\csromanheader{2}{ธมฺมานุโลมปจฺจนีย ทุกติกปฏฺฐานปาฬิ}
\prose{\csromanfolio{232}อนุปาทินฺนํ อกุสลํ ธมฺมํ ปฏิจฺจ น-อุปาทินฺโน น-อกุสโล ธมฺโม อุปฺปชฺชติ เหตุปจฺจยา.\csromanspacer{}\csromanspacer{}เหตุยา เอกํ. (อพฺยากตมูลํ เอกํเยว.)}

\nitthitam{มหนฺตรทุกํ นิฏฺฐิตํ.}
\csromanheader{2}{69. ทฺวิโคจฺฉก 1. กุสลตฺติก}
\proseitem{61}{โน-อุปาทานํ กุสลํ ธมฺมํ ปฏิจฺจ โน-อุปาทาโน นกุสโล ธมฺโม อุปฺปชฺชติ เหตุปจฺจยา.\csromanspacer{}\csromanspacer{}เหตุยา เอกํ.}

\prose{อุปาทานํ อกุสลํ ธมฺมํ ปฏิจฺจ โน-อุปาทาโน น-อกุสโล ธมฺโม อุปฺปชฺชติ เหตุปจฺจยา.\csromanspacer{}\csromanspacer{}เหตุยา ตีณิ.}

\prose{โนกิเลสํ กุสลํ ธมฺมํ ปฏิจฺจ โนกิเลโส นกุสโล ธมฺโม อุปฺปชฺชติ เหตุปจฺจยา.\csromanspacer{}\csromanspacer{}เหตุยา เอกํ.}

\csromanheader{2}{83. ทสฺสเนนปหาตพฺพทุก 1. กุสลตฺติก}
\proseitem{62}{นทสฺสเนน ปหาตพฺพํ กุสลํ ธมฺมํ ปฏิจฺจ นทสฺสเนน ปหาตพฺโพ นกุสโล ธมฺโม อุปฺปชฺชติ เหตุปจฺจยา.\csromanspacer{}\csromanspacer{}เหตุยา เอกํ.}

\prose{ทสฺสเนน ปหาตพฺพํ อกุสลํ ธมฺมํ ปฏิจฺจ นทสฺสเนน ปหาตพฺโพ น-อกุสโล ธมฺโม อุปฺปชฺชติ เหตุปจฺจยา.\csromanspacer{}\csromanspacer{}นทสฺสเนน ปหาตพฺพํ อกุสลํ ธมฺมํ ปฏิจฺจ นทสฺสเนน ปหาตพฺโพ น-อกุสโล ธมฺโม อุปฺปชฺชติ เหตุปจฺจยา.\csromanspacer{}\csromanspacer{}เหตุยา เทฺว.}

\prose{นทสฺสเนน ปหาตพฺพํ อพฺยากตํ ธมฺมํ ปจฺจยา นนทสฺสเนน ปหาตพฺโพ น-อพฺยากโต ธมฺโม อุปฺปชฺชติ เหตุปจฺจยา.\csromanspacer{}\csromanspacer{}เหตุยา เทฺว.}

\csromanheader{2}{84. ภาวนายปหาตพฺพทุก 1. กุสลตฺติก}
\proseitem{63}{นภาวนาย ปหาตพฺพํ กุสลํ ธมฺมํ ปฏิจฺจ นภาวนาย ปหาตพฺโพ นกุสโล ธมฺโม อุปฺปชฺชติ เหตุปจฺจยา.\csromanspacer{}\csromanspacer{}เหตุยา เอกํ.}

\prose{ภาวนาย ปหาตพฺพํ อกุสลํ ธมฺมํ ปฏิจฺจ นภาวนาย ปหาตพฺโพ น-อกุสโล ธมฺโม อุปฺปชฺชติ เหตุปจฺจยา.\csromanspacer{}\csromanspacer{}เหตุยา เทฺว.}

\prose{นภาวนาย ปหาตพฺพํ อพฺยากตํ ธมฺมํ ปจฺจยา นนภาวนาย ปหาตพฺโพ น-อพฺยากโต ธมฺโม อุปฺปชฺชติ เหตุปจฺจยา.\csromanspacer{}\csromanspacer{}เหตุยา เทฺว.}

\vspace{\tipitakaparskip}
\csromancenter{สวิจารทุก 1.\csromanspacer{} กุสลตฺติก \textbf{85.}\csromanspacer{}\textbf{ ทสฺสเนนปหาตพฺพเหตุกทุก 1.}\csromanspacer{}\textbf{ กุสลตฺติก}}
\proseitem{64}{\csromanfolio{233}นทสฺสเนน ปหาตพฺพเหตุกํ กุสลํ ธมฺมํ ปฏิจฺจ นทสฺสเนน ปหาตพฺพเหตุโก นกุสโล ธมฺโม อุปฺปชฺชติ เหตุปจฺจยา.\csromanspacer{}\csromanspacer{}เหตุยา เอกํ.}

\prose{ทสฺสเนน ปหาตพฺพเหตุกํ อกุสลํ ธมฺมํ ปฏิจฺจ นทสฺสเนน ปหาตพฺพเหตุโก น-อกุสโล ธมฺโม อุปฺปชฺชติ เหตุปจฺจยา.\csromanspacer{}\csromanspacer{}เหตุยา ตีณิ.}

\prose{นทสฺสเนน ปหาตพฺพเหตุกํ อพฺยากตํ ธมฺมํ ปจฺจยา นนทสฺสเนน ปหาตพฺพเหตุโก น-อพฺยากโต ธมฺโม อุปฺปชฺชติ เหตุปจฺจยา.\csromanspacer{}\csromanspacer{}เหตุยา เทฺว. (อพฺยากตวาเร สพฺพตฺถ ปจฺจยวเสน คเณตพฺพํ.)}

\csromanheader{2}{86. ภาวนายปหาตพฺพเหตุกทุก 1. กุสลตฺติก}
\proseitem{65}{นภาวนาย ปหาตพฺพเหตุกํ กุสลํ ธมฺมํ ปฏิจฺจ นภาวนาย ปหาตพฺพเหตุโก นกุสโล ธมฺโม อุปฺปชฺชติ เหตุปจฺจยา.\csromanspacer{}\csromanspacer{}เหตุยา เอกํ.}

\prose{ภาวนาย ปหาตพฺพเหตุกํ อกุสลํ ธมฺมํ ปฏิจฺจ นภาวนาย ปหาตพฺพเหตุโก น-อกุสโล ธมฺโม อุปฺปชฺชติ เหตุปจฺจยา.\csromanspacer{}\csromanspacer{}เหตุยา ตีณิ. (อพฺยากตมูลํ เทฺว.)}

\csromanheader{2}{87. สวิตกฺกทุก 1. กุสลตฺติก}
\proseitem{66}{สวิตกฺกํ กุสลํ ธมฺมํ ปฏิจฺจ นสวิตกฺโก นกุสโล ธมฺโม อุปฺปชฺชติ เหตุปจฺจยา.\csromanspacer{}\csromanspacer{}เหตุยา ตีณิ.}

\prose{สวิตกฺกํ อกุสลํ ธมฺมํ ปฏิจฺจ นสวิตกฺโก น-อกุสโล ธมฺโม อุปฺปชฺชติ เหตุปจฺจยา.\csromanspacer{}\csromanspacer{}เหตุยา ตีณิ. (อพฺยากตมูลํ ตีณิเยว.)}

\csromanheader{2}{88. สวิจารทุก 1. กุสลตฺติก}
\proseitem{67}{สวิจารํ กุสลํ ธมฺมํ ปฏิจฺจ นสวิจาโร นกุสโล ธมฺโม อุปฺปชฺชติ เหตุปจฺจยา.\csromanspacer{}\csromanspacer{}เหตุยา ตีณิ.}

\prose{สวิจารํ อกุสลํ ธมฺมํ ปฏิจฺจ นสวิจาโร น-อกุสโล ธมฺโม อุปฺปชฺชติ เหตุปจฺจยา.\csromanspacer{}\csromanspacer{}เหตุยา ตีณิ. (อพฺยากตมูลํ ตีณิเยว.)}

\csromanheader{2}{ธมฺมานุโลมปจฺจนีย ทุกติกปฏฺฐานปาฬิ \textbf{89. สปฺปีติกทุก 1. กุสลตฺติก}}
\proseitem{68}{\csromanfolio{234}สปฺปีติกํ กุสลํ ธมฺมํ ปฏิจฺจ นสปฺปีติโก นกุสโล ธมฺโม อุปฺปชฺชติ เหตุปจฺจยา.\csromanspacer{}\csromanspacer{}อปฺปีติกํ กุสลํ ธมฺมํ ปฏิจฺจ นสปฺปีติโก นกุสโล ธมฺโม อุปฺปชฺชติ เหตุปจฺจยา.\csromanspacer{}\csromanspacer{}สปฺปีติกํ กุสลญฺจ อปฺปีติกํ กุสลญฺจ ธมฺมํ ปฏิจฺจ นสปฺปีติโก นกุสโล ธมฺโม อุปฺปชฺชติ เหตุปจฺจยา.\csromanspacer{}\csromanspacer{}เหตุยา ตีณิ.}

\prose{สปฺปีติกํ อกุสลํ ธมฺมํ ปฏิจฺจ นสปฺปีติโก น-อกุสโล ธมฺโม อุปฺปชฺชติ เหตุปจฺจยา.\csromanspacer{}\csromanspacer{}อปฺปีติกํ อกุสลํ ธมฺมํ ปฏิจฺจ นสปฺปีติโก น-อกุสโล ธมฺโม อุปฺปชฺชติ เหตุปจฺจยา.\csromanspacer{}\csromanspacer{}สปฺปีติกํ อกุสลญฺจ อปฺปีติกํ อกุสลญฺจ ธมฺมํ ปฏิจฺจ นสปฺปีติโก อกุสโล ธมฺโม อุปฺปชฺชติ เหตุปจฺจยา.\csromanspacer{}\csromanspacer{}เหตุยา ตีณิ. (อพฺยากตมูลํ ตีณิเยว.)}

\csromanheader{2}{90. ปีติสหคตทุก 1. กุสลตฺติก}
\proseitem{69}{ปีติสหคตํ กุสลํ ธมฺมํ ปฏิจฺจ นปีติสหคโต นกุสโล ธมฺโม อุปฺปชฺชติ เหตุปจฺจยา.\csromanspacer{}\csromanspacer{}เหตุยา ตีณิ.}

\prose{ปีติสหคตํ อกุสลํ ธมฺมํ ปฏิจฺจ นปีติสหคโต น-อกุสโล ธมฺโม อุปฺปชฺชติ เหตุปจฺจยา.\csromanspacer{}\csromanspacer{}เหตุยา ตีณิ. (อพฺยากตมูลํ ตีณิเยว.)}

\csromanheader{2}{91. สุขสหคตทุก 1. กุสลตฺติก}
\proseitem{70}{สุขสหคตํ กุสลํ ธมฺมํ ปฏิจฺจ นสุขสหคโต นกุสโล ธมฺโม อุปฺปชฺชติ เหตุปจฺจยา.\csromanspacer{}\csromanspacer{}นสุขสหคตํ กุสลํ ธมฺมํ ปฏิจฺจ นสุขสหคโต นกุสโล ธมฺโม อุปฺปชฺชติ เหตุปจฺจยา.\csromanspacer{}\csromanspacer{}สุขสหคตํ กุสลญฺจ นสุขสหคตํ กุสลญฺจ ธมฺมํ ปฏิจฺจ นสุขสหคโต นกุสโล ธมฺโม อุปฺปชฺชติ เหตุปจฺจยา.\csromanspacer{}\csromanspacer{}เหตุยา ตีณิ.}

\prose{สุขสหคตํ อกุสลํ ธมฺมํ ปฏิจฺจ นสุขสหคโต น-อกุสโล ธมฺโม อุปฺปชฺชติ เหตุปจฺจยา.\csromanspacer{}\csromanspacer{}นสุขสหคตํ อกุสลํ ธมฺมํ ปฏิจฺจ นสุขสหคโต น-อกุสโล ธมฺโม อุปฺปชฺชติ เหตุปจฺจยา.\csromanspacer{}\csromanspacer{}สุขสหคตํ อกุสลญฺจ นสุขสหคตํ อกุสลญฺจ ธมฺมํ ปฏิจฺจ นสุขสหคโต น-อกุสโล ธมฺโม อุปฺปชฺชติ เหตุปจฺจยา.\csromanspacer{}\csromanspacer{}เหตุยา ตีณิ. (อพฺยากตมูลํ ตีณิเยว.)}

\vspace{\tipitakaparskip}
\csromancenter{อรูปาวจรทุก 1.\csromanspacer{} กุสลตฺติก \textbf{92.}\csromanspacer{}\textbf{ อุเปกฺขาสหคตทุก 1.}\csromanspacer{}\textbf{ กุสลตฺติก}}
\proseitem{71}{\csromanfolio{235}อุเปกฺขาสหคตํ กุสลํ ธมฺมํ ปฏิจฺจ น-อุเปกฺขาสหคโต นกุสโล ธมฺโม อุปฺปชฺชติ เหตุปจฺจยา.\csromanspacer{}\csromanspacer{}น-อุเปกฺขาสหคตํ กุสลํ ธมฺมํ ปฏิจฺจ น-อุเปกฺขาสหคโต นกุสโล ธมฺโม อุปฺปชฺชติ เหตุปจฺจยา.\csromanspacer{}\csromanspacer{}อุเปกฺขาสหคตํ กุสลญฺจ น-อุเปกฺขาสหคตํ กุสลญฺจ ธมฺมํ ปฏิจฺจ น-อุเปกฺขาสหคโต นกุสโล ธมฺโม อุปฺปชฺชติ เหตุปจฺจยา.\csromanspacer{}\csromanspacer{}เหตุยา ตีณิ.}

\prose{อุเปกฺขาสหคตํ อกุสลํ ธมฺมํ ปฏิจฺจ น-อุเปกฺขาสหคโต น-อกุสโล ธมฺโม อุปฺปชฺชติ เหตุปจฺจยา.\csromanspacer{}\csromanspacer{}น-อุเปกฺขาสหคตํ อกุสลํ ธมฺมํ ปฏิจฺจ น-อุเปกฺขาสหคโต น-อกุสโล ธมฺโม อุปฺปชฺชติ เหตุปจฺจยา.\csromanspacer{}\csromanspacer{}อุเปกฺขาสหคตํ อกุสลญฺจ น-อุเปกฺขาสหคตํ อกุสลญฺจ ธมฺมํ ปฏิจฺจ น-อุเปกฺขาสหคโต น-อกุสโล ธมฺโม อุปฺปชฺชติ เหตุปจฺจยา.\csromanspacer{}\csromanspacer{}เหตุยา ตีณิ. (อพฺยากตมูลํ ตีณิเยว.)}

\csromanheader{2}{93. กามาวจรทุก 1. กุสลตฺติก}
\proseitem{72}{กามาวจรํ กุสลํ ธมฺมํ ปฏิจฺจ นนกามาวจโร นกุสโล ธมฺโม อุปฺปชฺชติ เหตุปจฺจยา.\csromanspacer{}\csromanspacer{}นกามาวจรํ กุสลํ ธมฺมํ ปฏิจฺจ นนกามาวจโร นกุสโล ธมฺโม อุปฺปชฺชติ เหตุปจฺจยา.\csromanspacer{}\csromanspacer{}เหตุยา เทฺว.}

\prose{กามาวจรํ อกุสลํ ธมฺมํ ปฏิจฺจ นนกามาวจโร น-อกุสโล ธมฺโม อุปฺปชฺชติ เหตุปจฺจยา.\csromanspacer{}\csromanspacer{}เหตุยา เอกํ. (อพฺยากตมูลํ เทฺว.)}

\csromanheader{2}{94. รูปาวจรทุก 1. กุสลตฺติก}
\proseitem{73}{รูปาวจรํ กุสลํ ธมฺมํ ปฏิจฺจ นรูปาวจโร นกุสโล ธมฺโม อุปฺปชฺชติ เหตุปจฺจยา.\csromanspacer{}\csromanspacer{}นรูปาวจรํ กุสลํ ธมฺมํ ปฏิจฺจ นรูปาวจโร นกุสโล ธมฺโม อุปฺปชฺชติ เหตุปจฺจยา.\csromanspacer{}\csromanspacer{}เหตุยา เทฺว.}

\prose{นรูปาวจรํ อกุสลํ ธมฺมํ ปฏิจฺจ นรูปาวจโร น-อกุสโล ธมฺโม อุปฺปชฺชติ เหตุปจฺจยา.\csromanspacer{}\csromanspacer{}เหตุยา เอกํ. (อพฺยากตมูลํ เทฺว.)}

\csromanheader{2}{95. อรูปาวจรทุก 1. กุสลตฺติก}
\proseitem{74}{อรูปาวจรํ กุสลํ ธมฺมํ ปฏิจฺจ น-อรูปาวจโร นกุสโล ธมฺโม อุปฺปชฺชติ เหตุปจฺจยา.\csromanspacer{}\csromanspacer{}น-อรูปาวจรํ กุสลํ ธมฺมํ ปฏิจฺจ น-อรูปาวจโร นกุสโล ธมฺโม อุปฺปชฺชติ เหตุปจฺจยา.\csromanspacer{}\csromanspacer{}เหตุยา เทฺว.}

\csromanheader{2}{ธมฺมานุโลมปจฺจนีย ทุกติกปฏฺฐานปาฬิ}
\prose{\csromanfolio{236}น-อรูปาวจรํ อกุสลํ ธมฺมํ ปฏิจฺจ น-อรูปาวจโร น-อกุสโล ธมฺโม อุปฺปชฺชติ เหตุปจฺจยา.\csromanspacer{} เหตุยา เอกํ. (อพฺยากตมูลํ เทฺว.)}

\csromanheader{2}{96. ปริยาปนฺนทุก 1. กุสลตฺติก}
\proseitem{75}{ปริยาปนฺนํ กุสลํ ธมฺมํ ปฏิจฺจ น-อปริยาปนฺโน นกุสโล ธมฺโม อุปฺปชฺชติ เหตุปจฺจยา.\csromanspacer{}\csromanspacer{}อปริยาปนฺนํ กุสลํ ธมฺมํ ปฏิจฺจ น-อปริยาปนฺโน นกุสโล ธมฺโม อุปฺปชฺชติ เหตุปจฺจยา.\csromanspacer{}\csromanspacer{}เหตุยา เทฺว.}

\prose{ปริยาปนฺนํ อกุสลํ ธมฺมํ ปฏิจฺจ น-อปริยาปนฺโน น-อกุสโล ธมฺโม อุปฺปชฺชติ เหตุปจฺจยา.\csromanspacer{}\csromanspacer{}เหตุยา เอกํ. (อพฺยากตมูลํ เทฺว.)}

\csromanheader{2}{97. นิยฺยานิกทุก 1. กุสลตฺติก}
\proseitem{76}{นิยฺยานิกํ กุสลํ ธมฺมํ ปฏิจฺจ นนิยฺยานิโก นกุสโล ธมฺโม อุปฺปชฺชติ เหตุปจฺจยา.\csromanspacer{}\csromanspacer{}อนิยฺยานิกํ กุสลํ ธมฺมํ ปฏิจฺจ นนิยฺยานิโก นกุสโล ธมฺโม อุปฺปชฺชติ เหตุปจฺจยา.\csromanspacer{}\csromanspacer{}เหตุยา เทฺว.}

\prose{อนิยฺยานิกํ อกุสลํ ธมฺมํ ปฏิจฺจ นนิยฺยานิโก น-อกุสโล ธมฺโม อุปฺปชฺชติ เหตุปจฺจยา.\csromanspacer{}\csromanspacer{}เหตุยา เอกํ. (อพฺยากตมูลํ เทฺว.)}

\csromanheader{2}{98. นิยตทุก 1. กุสลตฺติก}
\proseitem{77}{นิยตํ กุสลํ ธมฺมํ ปฏิจฺจ นนิยโต นกุสโล ธมฺโม อุปฺปชฺชติ เหตุปจฺจยา.\csromanspacer{}\csromanspacer{}อนิยตํ กุสลํ ธมฺมํ ปฏิจฺจ นนิยโต นกุสโล ธมฺโม อุปฺปชฺชติ เหตุปจฺจยา.\csromanspacer{}\csromanspacer{}เหตุยา เทฺว.}

\prose{นิยตํ อกุสลํ ธมฺมํ ปฏิจฺจ นนิยโต น-อกุสโล ธมฺโม อุปฺปชฺชติ เหตุปจฺจยา.\csromanspacer{}\csromanspacer{}อนิยตํ อกุสลํ ธมฺมํ ปฏิจฺจ นนิยโต น-อกุสโล ธมฺโม อุปฺปชฺชติ เหตุปจฺจยา.\csromanspacer{}\csromanspacer{}เหตุยา เทฺว. (อพฺยากตมูเล เทฺว.)}

\csromanheader{2}{99. ส-อุตฺตรทุก 1. กุสลตฺติก}
\proseitem{78}{ส-อุตฺตรํ กุสลํ ธมฺมํ ปฏิจฺจ น-อนุตฺตโร นกุสโล ธมฺโม อุปฺปชฺชติ เหตุปจฺจยา.\csromanspacer{}\csromanspacer{}อนุตฺตรํ กุสลํ ธมฺมํ ปฏิจฺจ น-อนุตฺตโร นกุสโล ธมฺโม อุปฺปชฺชติ เหตุปจฺจยา.\csromanspacer{}\csromanspacer{}เหตุยา เทฺว.}

\prose{ส-อุตฺตรํ อกุสลํ ธมฺมํ ปฏิจฺจ น-อนุตฺตโร น-อกุสโล ธมฺโม อุปฺปชฺชติ เหตุปจฺจยา.\csromanspacer{}\csromanspacer{}เหตุยา เอกํ. (อพฺยากตมูเล เทฺว.)}

\csromanheader{2}{100. สรณทุก 2. เวทนาตฺติก \textbf{100. สรณทุก 1. กุสลตฺติก}}
\proseitem{79}{\csromanfolio{237}อรณํ กุสลํ ธมฺมํ ปฏิจฺจ นสรโณ นกุสโล ธมฺโม อุปฺปชฺชติ เหตุปจฺจยา.\csromanspacer{}\csromanspacer{}เหตุยา เอกํ.}

\prose{สรณํ อกุสลํ ธมฺมํ ปฏิจฺจ นสรโณ น-อกุสโล ธมฺโม อุปฺปชฺชติ เหตุปจฺจยา.\csromanspacer{}\csromanspacer{}เหตุยา เอกํ.}

\prose{อรณํ อพฺยากตํ ธมฺมํ ปจฺจยา น-อรโณ น-อพฺยากโต ธมฺโม อุปฺปชฺชติ เหตุปจฺจยา.\csromanspacer{}\csromanspacer{}อรณํ อพฺยากตํ ธมฺมํ ปจฺจยา นสรโณ น-อพฺยากโต ธมฺโม อุปฺปชฺชติ เหตุปจฺจยา.}

\prose{เหตุยา เทฺว. (สหชาตวารมฺปิ ฯเปฯ สมฺปยุตฺตวารมฺปิ ปฏิจฺจวารสทิสํ.)}

\proseitem{80}{อรโณ อพฺยากโต ธมฺโม น-อรณสฺส น-อพฺยากตสฺส ธมฺมสฺส อารมฺมณปจฺจเยน ปจฺจโย.\csromanspacer{}\csromanspacer{}อรโณ อพฺยากโต ธมฺโม นสรณสฺส น-อพฺยากตสฺส ธมฺมสฺส อารมฺมณปจฺจเยน ปจฺจโย. (2)}

\prose{อารมฺมเณ เทฺว ฯเปฯ อวิคเต เทฺว.}

\csromanheader{2}{100. สรณทุก 2. เวทนาตฺติก}
\proseitem{81}{สรณํ สุขาย เวทนาย สมฺปยุตฺตํ ธมฺมํ ปฏิจฺจ นสรโณ นสุขาย เวทนาย สมฺปยุตฺโต ธมฺโม อุปฺปชฺชติ เหตุปจฺจยา.\csromanspacer{}\csromanspacer{}สรณํ สุขาย เวทนาย สมฺปยุตฺตํ ธมฺมํ ปฏิจฺจ น-อรโณ นสุขาย เวทนาย สมฺปยุตฺโต ธมฺโม อุปฺปชฺชติ เหตุปจฺจยา.\csromanspacer{}\csromanspacer{}สรณํ สุขาย เวทนาย สมฺปยุตฺตํ ธมฺมํ ปฏิจฺจ นสรโณ นสุขาย เวทนาย สมฺปยุตฺโต จ น-อรโณ นสุขาย เวทนาย สมฺปยุตฺโต จ ธมฺมา อุปฺปชฺชนฺติ เหตุปจฺจยา. (3)}

\prose{อรณํ สุขาย เวทนาย สมฺปยุตฺตํ ธมฺมํ ปฏิจฺจ นสรโณ นสุขาย เวทนาย สมฺปยุตฺโต ธมฺโม อุปฺปชฺชติ เหตุปจฺจยา. (1)}

\prose{เหตุยา จตฺตาริ.}

\proseitem{82}{สรณํ ทุกฺขาย เวทนาย สมฺปยุตฺตํ ธมฺมํ ปฏิจฺจ นสรโณ นทุกฺขาย เวทนาย สมฺปยุตฺโต ธมฺโม อุปฺปชฺชติ เหตุปจฺจยา.}

\prose{เหตุยา ตีณิ.}

\csromanheader{2}{ธมฺมานุโลมปจฺจนีย ทุกติกปฏฺฐานปาฬิ}
\prose{\csromanfolio{238}สรณํ อทุกฺขมสุขาย เวทนาย สมฺปยุตฺตํ ธมฺมํ ปฏิจฺจ นสรโณ น-อทุกฺขมสุขาย เวทนาย สมฺปยุตฺโต ธมฺโม อุปฺปชฺชติ เหตุปจฺจยา.}

\prose{เหตุยา จตฺตาริ.}

\csromanheader{2}{100. สรณทุก 3. วิปากตฺติก}
\proseitem{83}{อรณํ วิปากํ ธมฺมํ ปฏิจฺจ นสรโณ นวิปาโก ธมฺโม อุปฺปชฺชติ เหตุปจฺจยา.\csromanspacer{}\csromanspacer{}เหตุยา เอกํ.}

\prose{สรณํ วิปากธมฺมธมฺมํ ปฏิจฺจ นสรโณ นวิปากธมฺมธมฺโม อุปฺปชฺชติ เหตุปจฺจยา.\csromanspacer{}\csromanspacer{}อรณํ วิปากธมฺมธมฺมํ ปฏิจฺจ นสรโณ นวิปากธมฺมธมฺโม อุปฺปชฺชติ เหตุปจฺจยา.\csromanspacer{}\csromanspacer{}เหตุยา เทฺว.}

\prose{อรณํ เนววิปากนวิปากธมฺมธมฺมํ ปฏิจฺจ นสรโณ นเนววิปากนวิปากธมฺมธมฺโม อุปฺปชฺชติ เหตุปจฺจยา.\csromanspacer{}\csromanspacer{}เหตุยา เอกํ.}

\csromanheader{2}{100. สรณทุก 4. อุปาทินฺนตฺติก}
\proseitem{84}{อรณํ อุปาทินฺนุปาทานิยํ ธมฺมํ ปฏิจฺจ นสรโณ น-อุปาทินฺนุปาทานิโย ธมฺโม อุปฺปชฺชติ เหตุปจฺจยา.\csromanspacer{}\csromanspacer{}เหตุยา เอกํ.}

\prose{อรณํ อนุปาทินฺน-อนุปาทานิยํ ธมฺมํ ปฏิจฺจ นสรโณ น-อนุปาทินฺน-อนุปาทานิโย ธมฺโม อุปฺปชฺชติ เหตุปจฺจยา.\csromanspacer{}\csromanspacer{}เหตุยา เอกํ.}

\csromanheader{2}{100. สรณทุก 5. สํกิลิฏฺฐตฺติก}
\proseitem{85}{สรณํ สํกิลิฏฺฐสํกิเลสิกํ ธมฺมํ ปฏิจฺจ นสรโณ นสํกิลิฏฺฐสํกิเลสิโก ธมฺโม อุปฺปชฺชติ เหตุปจฺจยา.\csromanspacer{}\csromanspacer{}เหตุยา เอกํ.}

\prose{อรณํ อสํกิลิฏฺฐสํกิเลสิกํ ธมฺมํ ปจฺจยา น-อรโณ น-อสํกิลิฏฺฐสํกิเลสิโก ธมฺโม อุปฺปชฺชติ เหตุปจฺจยา.\csromanspacer{}\csromanspacer{}อรณํ อสํกิลิฏฺฐสํกิเลสิกํ ธมฺมํ ปจฺจยา นสรโณ น-อสํกิลิฏฺฐสํกิเลสิโก ธมฺโม อุปฺปชฺชติ เหตุปจฺจยา.\csromanspacer{}\csromanspacer{}เหตุยา เทฺว.}

\prose{อรณํ อสํกิลิฏฺฐ-อสํกิเลสิกํ ธมฺมํ ปฏิจฺจ นสรโณ น-อสํกิลิฏฺฐ-อสํกิเลสิโก ธมฺโม อุปฺปชฺชติ เหตุปจฺจยา.\csromanspacer{}\csromanspacer{}เหตุยา เอกํ.}

\csromanheader{2}{100. สรณทุก 8. ทสฺสนตฺติก \textbf{100. สรณทุก 6. วิตกฺกตฺติก}}
\proseitem{86}{\csromanfolio{239}สรณํ สวิตกฺกสวิจารํ ธมฺมํ ปฏิจฺจ นสรโณ นสวิตกฺกสวิจาโร ธมฺโม อุปฺปชฺชติ เหตุปจฺจยา.\csromanspacer{} ตีณิ.}

\prose{อรณํ สวิตกฺกสวิจารํ ธมฺมํ ปฏิจฺจ น-อรโณ นสวิตกฺกสวิจาโร ธมฺโม อุปฺปชฺชติ เหตุปจฺจยา.\csromanspacer{} เอกํ.}

\proseitem{87}{สรณํ อวิตกฺกวิจารมตฺตํ ธมฺมํ ปฏิจฺจ นสรโณ น-อวิตกฺกวิจารมตฺโต ธมฺโม อุปฺปชฺชติ เหตุปจฺจยา.\csromanspacer{}\csromanspacer{}เหตุยา จตฺตาริ.}

\prose{อรณํ อวิตกฺก-อวิจารํ ธมฺมํ ปฏิจฺจ นสรโณ น-อวิตกฺก-อวิจาโร ธมฺโม อุปฺปชฺชติ เหตุปจฺจยา.\csromanspacer{}\csromanspacer{}เหตุยา เอกํ.}

\csromanheader{2}{100. สรณทุก 7. ปีติตฺติก}
\proseitem{88}{สรณํ ปีติสหคตํ ธมฺมํ ปฏิจฺจ นสรโณ นปีติสหคโต ธมฺโม อุปฺปชฺชติ เหตุปจฺจยา.\csromanspacer{} ตีณิ.}

\prose{อรณํ ปีติสหคตํ ธมฺมํ ปฏิจฺจ นสรโณ นปีติสหคโต ธมฺโม อุปฺปชฺชติ เหตุปจฺจยา.\csromanspacer{} เอกํ.}

\prose{สรณํ สุขสหคตํ ธมฺมํ ปฏิจฺจ นสรโณ นสุขสหคโต ธมฺโม อุปฺปชฺชติ เหตุปจฺจยา.\csromanspacer{} ตีณิ.}

\prose{อรณํ สุขสหคตํ ธมฺมํ ปฏิจฺจ นสรโณ นสุขสหคโต ธมฺโม อุปฺปชฺชติ เหตุปจฺจยา.\csromanspacer{} เอกํ.}

\prose{สรณํ อุเปกฺขาสหคตํ ธมฺมํ ปฏิจฺจ นสรโณ น-อุเปกฺขาสหคโต ธมฺโม อุปฺปชฺชติ เหตุปจฺจยา.\csromanspacer{}\csromanspacer{}เหตุยา จตฺตาริ.}

\csromanheader{2}{100. สรณทุก 8. ทสฺสนตฺติก}
\proseitem{89}{สรณํ ทสฺสเนน ปหาตพฺพํ ธมฺมํ ปฏิจฺจ นสรโณ นทสฺสเนน ปหาตพฺโพ ธมฺโม อุปฺปชฺชติ เหตุปจฺจยา.\csromanspacer{}\csromanspacer{}เหตุยา เอกํ.}

\prose{สรณํ ภาวนาย ปหาตพฺพํ ธมฺมํ ปฏิจฺจ นสรโณ นภาวนาย ปหาตพฺโพ ธมฺโม อุปฺปชฺชติ เหตุปจฺจยา.\csromanspacer{}\csromanspacer{}เหตุยา เอกํ.}

\csromanheader{2}{ธมฺมานุโลมปจฺจนีย ทุกติกปฏฺฐานปาฬิ}
\prose{\csromanfolio{240}อรณํ เนวทสฺสเนน นภาวนาย ปหาตพฺพํ ธมฺมํ ปจฺจยา น-อรโณ นเนวทสฺสเนน นภาวนาย ปหาตพฺโพ ธมฺโม อุปฺปชฺชติ เหตุปจฺจยา.\csromanspacer{}\csromanspacer{}เหตุยา เอกํ.}

\csromanheader{2}{100. สรณทุก 9. ทสฺสนเหตุตฺติก}
\proseitem{90}{สรณํ ทสฺสเนน ปหาตพฺพเหตุกํ ธมฺมํ ปจฺจยา นสรโณ นทสฺสเนน ปหาตพฺพเหตุโก ธมฺโม อุปฺปชฺชติ เหตุปจฺจยา.\csromanspacer{}\csromanspacer{}เหตุยา เอกํ.}

\prose{สรณํ ภาวนาย ปหาตพฺพเหตุกํ ธมฺมํ ปฏิจฺจ นสรโณ นภาวนาย ปหาตพฺพเหตุโก ธมฺโม อุปฺปชฺชติ เหตุปจฺจยา.\csromanspacer{}\csromanspacer{}เหตุยา เอกํ.}

\prose{อรณํ เนวทสฺสเนน นภาวนาย ปหาตพฺพเหตุกํ ธมฺมํ ปฏิจฺจ น-อรโณ นเนวทสฺสเนน นภาวนาย ปหาตพฺพเหตุโก ธมฺโม อุปฺปชฺชติ เหตุปจฺจยา.\csromanspacer{}\csromanspacer{}เหตุยา เอกํ.}

\csromanheader{2}{100. สรณทุก 10. อาจยคามิตฺติก}
\proseitem{91}{สรณํ อาจยคามิํ ธมฺมํ ปฏิจฺจ นสรโณ น-อาจยคามี ธมฺโม อุปฺปชฺชติ เหตุปจฺจยา.\csromanspacer{}\csromanspacer{}เหตุยา เทฺว.}

\prose{อรณํ อปจยคามิํ ธมฺมํ ปฏิจฺจ นสรโณ น-อปจยคามี ธมฺโม อุปฺปชฺชติ เหตุปจฺจยา.\csromanspacer{}\csromanspacer{}เหตุยา เอกํ.}

\prose{อรณํ เนวาจยคามินาปจยคามิํ ธมฺมํ ปจฺจยา น-อรโณ นเนวาจยคามินาปจยคามี ธมฺโม อุปฺปชฺชติ เหตุปจฺจยา.}

\csromanheader{2}{100. สรณทุก 11. เสกฺขตฺติก}
\proseitem{92}{อรณํ เสกฺขํ ธมฺมํ ปฏิจฺจ นสรโณ นเสกฺโข ธมฺโม อุปฺปชฺชติ เหตุปจฺจยา.\csromanspacer{}\csromanspacer{}เหตุยา เอกํ.}

\prose{อรณํ อเสกฺขํ ธมฺมํ ปฏิจฺจ นสรโณ น-อเสกฺโข ธมฺโม อุปฺปชฺชติ เหตุปจฺจยา.\csromanspacer{}\csromanspacer{}เหตุยา เอกํ.}

\csromanheader{2}{100. สรณทุก 14. หีนตฺติก}
\prose{\csromanfolio{241}อรณํ เนวเสกฺขนาเสกฺขํ ธมฺมํ ปจฺจยา นสรโณ นเนวเสกฺขนาเสกฺโข ธมฺโม อุปฺปชฺชติ เหตุปจฺจยา.\csromanspacer{}\csromanspacer{}เหตุยา เอกํ.}

\csromanheader{2}{100. สรณทุก 12. ปริตฺตตฺติก}
\proseitem{93}{อรณํ ปริตฺตํ ธมฺมํ ปฏิจฺจ นสรโณ นปริตฺโต ธมฺโม อุปฺปชฺชติ เหตุปจฺจยา.\csromanspacer{}\csromanspacer{}เหตุยา เอกํ.}

\prose{อรณํ มหคฺคตํ ธมฺมํ ปฏิจฺจ นสรโณ นมหคฺคโต ธมฺโม อุปฺปชฺชติ เหตุปจฺจยา.\csromanspacer{}\csromanspacer{}เหตุยา เอกํ.}

\prose{อรณํ อปฺปมาณํ ธมฺมํ ปฏิจฺจ นสรโณ น-อปฺปมาโณ ธมฺโม อุปฺปชฺชติ เหตุปจฺจยา.\csromanspacer{}\csromanspacer{}เหตุยา เอกํ.}

\csromanheader{2}{100. สรณทุก 13. ปริตฺตารมฺมณตฺติก}
\proseitem{94}{สรณํ ปริตฺตารมฺมณํ ธมฺมํ ปฏิจฺจ นสรโณ นปริตฺตารมฺมโณ ธมฺโม อุปฺปชฺชติ เหตุปจฺจยา.\csromanspacer{}\csromanspacer{}อรณํ ปริตฺตารมฺมณํ ธมฺมํ ปฏิจฺจ นสรโณ นปริตฺตารมฺมโณ ธมฺโม อุปฺปชฺชติ เหตุปจฺจยา.\csromanspacer{}\csromanspacer{}เหตุยา เทฺว.}

\prose{สรณํ มหคฺคตารมฺมณํ ธมฺมํ ปฏิจฺจ นสรโณ นมหคฺคตารมฺมโณ ธมฺโม อุปฺปชฺชติ เหตุปจฺจยา.\csromanspacer{}\csromanspacer{}อรณํ มหคฺคตารมฺมณํ ธมฺมํ ปฏิจฺจ นสรโณ นมหคฺคตารมฺมโณ ธมฺโม อุปฺปชฺชติ เหตุปจฺจยา.\csromanspacer{}\csromanspacer{}เหตุยา เทฺว.}

\prose{อรณํ อปฺปมาณารมฺมณํ ธมฺมํ ปฏิจฺจ นสรโณ น-อปฺปมาณารมฺมโณ ธมฺโม อุปฺปชฺชติ เหตุปจฺจยา.\csromanspacer{}\csromanspacer{}เหตุยา เอกํ.}

\csromanheader{2}{100. สรณทุก 14. หีนตฺติก}
\proseitem{95}{สรณํ หีนํ ธมฺมํ ปฏิจฺจ นสรโณ นหีโน ธมฺโม อุปฺปชฺชติ เหตุปจฺจยา.\csromanspacer{}\csromanspacer{}เหตุยา เอกํ.}

\prose{อรณํ มชฺฌิมํ ธมฺมํ ปจฺจยา น-อรโณ นมชฺฌิโม ธมฺโม อุปฺปชฺชติ เหตุปจฺจยา.\csromanspacer{}\csromanspacer{}อรณํ มชฺฌิมํ ธมฺมํ ปจฺจยา นสรโณ นมชฺฌิโม ธมฺโม อุปฺปชฺชติ เหตุปจฺจยา.\csromanspacer{}\csromanspacer{}เหตุยา เทฺว.}

\prose{อรณํ ปณีตํ ธมฺมํ ปฏิจฺจ นสรโณ นปณีโต ธมฺโม อุปฺปชฺชติ เหตุปจฺจยา.\csromanspacer{}\csromanspacer{}เหตุยา เอกํ.}

\vspace{\tipitakaparskip}
\csromancenter{ธมฺมานุโลมปจฺจนีย ทุกติกปฏฺฐานปาฬิ \textbf{100.}\csromanspacer{}\textbf{ สรณทุก 15.}\csromanspacer{}\textbf{ มิจฺฉตฺตนิยตตฺติก}}
\proseitem{96}{\csromanfolio{242}สรณํ มิจฺฉตฺตนิยตํ ธมฺมํ ปฏิจฺจ นสรโณ นมิจฺฉตฺตนิยโต ธมฺโม อุปฺปชฺชติ เหตุปจฺจยา.\csromanspacer{}\csromanspacer{}เหตุยา เอกํ.}

\prose{อรณํ สมฺมตฺตนิยตํ ธมฺมํ ปฏิจฺจ นสรโณ นสมฺมตฺตนิยโต ธมฺโม อุปฺปชฺชติ เหตุปจฺจยา.\csromanspacer{}\csromanspacer{}เหตุยา เอกํ.}

\prose{อรณํ อนิยตํ ธมฺมํ ปจฺจยา น-อรโณ น-อนิยโต ธมฺโม อุปฺปชฺชติ เหตุปจฺจยา.\csromanspacer{}\csromanspacer{}อรณํ อนิยตํ ธมฺมํ ปจฺจยา นสรโณ น-อนิยโต ธมฺโม อุปฺปชฺชติ เหตุปจฺจยา.\csromanspacer{}\csromanspacer{}เหตุยา เทฺว.}

\csromanheader{2}{100. สรณทุก 16. มคฺคารมฺมณตฺติก}
\proseitem{97}{อรณํ มคฺคารมฺมณํ ธมฺมํ ปฏิจฺจ นสรโณ นมคฺคารมฺมโณ ธมฺโม อุปฺปชฺชติ เหตุปจฺจยา.\csromanspacer{}\csromanspacer{}เหตุยา เอกํ.}

\prose{อรณํ มคฺคเหตุกํ ธมฺมํ ปฏิจฺจ นสรโณ นมคฺคเหตุโก ธมฺโม อุปฺปชฺชติ เหตุปจฺจยา.\csromanspacer{}\csromanspacer{}เหตุยา เอกํ.}

\prose{อรณํ มคฺคาธิปติํ ธมฺมํ ปฏิจฺจ นสรโณ นมคฺคาธิปติ ธมฺโม อุปฺปชฺชติ เหตุปจฺจยา.\csromanspacer{}\csromanspacer{}เหตุยา เอกํ.}

\csromanheader{2}{100. สรณทุก 17. อุปฺปนฺนตฺติก}
\proseitem{98}{สรโณ อนุปฺปนฺโน ธมฺโม นสรณสฺส น-อนุปฺปนฺนสฺส ธมฺมสฺส อารมฺมณปจฺจเยน ปจฺจโย.\csromanspacer{}\csromanspacer{}อารมฺมเณ จตฺตาริ.}

\prose{อรโณ อุปฺปาที ธมฺโม น-อรณสฺส น-อุปฺปาทิสฺส ธมฺมสฺส อารมฺมณปจฺจเยน ปจฺจโย.\csromanspacer{}\csromanspacer{}อรโณ อุปฺปาที ธมฺโม นสรณสฺส น-อุปฺปาทิสฺส ธมฺมสฺส อารมฺมณปจฺจเยน ปจฺจโย.\csromanspacer{}\csromanspacer{}อารมฺมเณ เทฺว.}

\csromanheader{2}{100. สรณทุก 18. อตีตตฺติก}
\proseitem{99}{สรโณ อตีโต ธมฺโม นสรณสฺส น-อตีตสฺส ธมฺมสฺส อารมฺมณปจฺจเยน ปจฺจโย.\csromanspacer{}\csromanspacer{}อารมฺมเณ จตฺตาริ. (อนาคโต อตีตสทิโส.)}

\vspace{\tipitakaparskip}
\csromancenter{สรณทุก 22.\csromanspacer{} สนิทสฺสนตฺติก \textbf{100.}\csromanspacer{}\textbf{ สรณทุก 19.}\csromanspacer{}\textbf{ อตีตารมฺมณตฺติก}}
\prose{\csromanfolio{243}สรณํ อตีตารมฺมณํ ธมฺมํ ปฏิจฺจ นสรโณ น-อตีตารมฺมโณ ธมฺโม อุปฺปชฺชติ เหตุปจฺจยา.\csromanspacer{}\csromanspacer{}อรณํ อตีตารมฺมณํ ธมฺมํ ปฏิจฺจ นสรโณ น-อตีตารมฺมโณ ธมฺโม อุปฺปชฺชติ เหตุปจฺจยา.\csromanspacer{}\csromanspacer{}เหตุยา เทฺว.}

\prose{สรณํ อนาคตารมฺมณํ ธมฺมํ ปฏิจฺจ นสรโณ น-อนาคตารมฺมโณ ธมฺโม อุปฺปชฺชติ เหตุปจฺจยา.\csromanspacer{}\csromanspacer{}อรณํ อนาคตารมฺมณํ ธมฺมํ ปฏิจฺจ นสรโณ น-อนาคตารมฺมโณ ธมฺโม อุปฺปชฺชติ เหตุปจฺจยา.\csromanspacer{}\csromanspacer{}เหตุยา เทฺว.}

\prose{สรณํ ปจฺจุปฺปนฺนารมฺมณํ ธมฺมํ ปฏิจฺจ นสรโณ นปจฺจุปฺปนฺนารมฺมโณ ธมฺโม อุปฺปชฺชติ เหตุปจฺจยา.\csromanspacer{}\csromanspacer{}อรณํ ปจฺจุปฺปนฺนารมฺมณํ ธมฺมํ ปฏิจฺจ นสรโณ นปจฺจุปฺปนฺนารมฺมโณ ธมฺโม อุปฺปชฺชติ เหตุปจฺจยา.\csromanspacer{}\csromanspacer{}เหตุยา เทฺว.}

\csromanheader{2}{100. สรณทุก 21. อชฺฌตฺตารมฺมณตฺติก}
\proseitem{101}{สรณํ อชฺฌตฺตารมฺมณํ ธมฺมํ ปฏิจฺจ นสรโณ น-อชฺฌตฺตารมฺมโณ ธมฺโม อุปฺปชฺชติ เหตุปจฺจยา.\csromanspacer{}\csromanspacer{}อรณํ อชฺฌตฺตารมฺมณํ ธมฺมํ ปฏิจฺจ นสรโณ น-อชฺฌตฺตารมฺมโณ ธมฺโม อุปฺปชฺชติ เหตุปจฺจยา.\csromanspacer{}\csromanspacer{}เหตุยา เทฺว.}

\prose{สรณํ พหิทฺธารมฺมณํ ธมฺมํ ปฏิจฺจ นสรโณ นพหิทฺธารมฺมโณ ธมฺโม อุปฺปชฺชติ เหตุปจฺจยา.\csromanspacer{}\csromanspacer{}อรณํ พหิทฺธารมฺมณํ ธมฺมํ ปฏิจฺจ นสรโณ นพหิทฺธารมฺมโณ ธมฺโม อุปฺปชฺชติ เหตุปจฺจยา.\csromanspacer{}\csromanspacer{}เหตุยา เทฺว.}

\csromanheader{2}{100. สรณทุก 22. สนิทสฺสนตฺติก}
\proseitem{102}{อรณํ อนิทสฺสนสปฺปฏิฆํ ธมฺมํ ปฏิจฺจ นสรโณ น-อนิทสฺสนสปฺปฏิโฆ ธมฺโม อุปฺปชฺชติ เหตุปจฺจยา.\csromanspacer{}\csromanspacer{}เอกํ.}

\prose{สรณํ อนิทสฺสน-อปฺปฏิฆํ ธมฺมํ ปฏิจฺจ นสรโณ น-อนิทสฺสน-อปฺปฏิโฆ ธมฺโม อุปฺปชฺชติ เหตุปจฺจยา.\csromanspacer{}\csromanspacer{}อรณํ อนิทสฺสน-อปฺปฏิฆํ ธมฺมํ ปฏิจฺจ นสรโณ น-อนิทสฺสน-อปฺปฏิโฆ ธมฺโม อุปฺปชฺชติ เหตุปจฺจยา.\csromanspacer{}\csromanspacer{}สรณํ อนิทสฺสน-อปฺปฏิฆญฺจ อรณํ อนิทสฺสน-อปฺปฏิฆญฺจ ธมฺมํ ปฏิจฺจ นสรโณ น-อนิทสฺสน-อปฺปฏิโฆ ธมฺโม อุปฺปชฺชติ เหตุปจฺจยา. (3)}

\prose{เหตุยา ตีณิ ฯเปฯ อวิคเต ตีณิ.}

\csromanheader{2}{ธมฺมานุโลมปจฺจนีย ทุกติกปฏฺฐานปาฬิ \textbf{นเหตุ น-อารมฺมณ}}
\proseitem{103}{\csromanfolio{244}อรณํ อนิทสฺสน-อปฺปฏิฆํ ธมฺมํ ปฏิจฺจ นสรโณ น-อนิทสฺสน-อปฺปฏิโฆ ธมฺโม อุปฺปชฺชติ นเหตุปจฺจยา. (1)}

\prose{สรณํ อนิทสฺสน-อปฺปฏิฆํ ธมฺมํ ปฏิจฺจ นสรโณ น-อนิทสฺสน-อปฺปฏิโฆ ธมฺโม อุปฺปชฺชติ น-อารมฺมณปจฺจยา.}

\prose{นเหตุยา เอกํ, น-อารมฺมเณ ตีณิ ฯเปฯ โนวิคเต ตีณิ.}

\prose{(สหชาตวารมฺปิ ฯเปฯ สมฺปยุตฺตวารมฺปิ ปฏิจฺจวารสทิสํ.)}

\proseitem{104}{สรโณ อนิทสฺสน-อปฺปฏิโฆ ธมฺโม นสรณสฺส น-อนิทสฺสน-อปฺปฏิฆสฺส ธมฺมสฺส เหตุปจฺจเยน ปจฺจโย. (1)}

\prose{อรโณ อนิทสฺสน-อปฺปฏิโฆ ธมฺโม นสรณสฺส น-อนิทสฺสน-อปฺปฏิฆสฺส ธมฺมสฺส เหตุปจฺจเยน ปจฺจโย. (1)}

\prose{เหตุยา เทฺว, อธิปติยา เทฺว ฯเปฯ อวิคเต ตีณิ.}

\csromanheader{2}{\textbf{ปจฺจนียุทฺธาร}}
\proseitem{105}{สรโณ อนิทสฺสน-อปฺปฏิโฆ ธมฺโม นสรณสฺส น-อนิทสฺสน-อปฺปฏิฆสฺส ธมฺมสฺส สหชาตปจฺจเยน ปจฺจโย.\csromanspacer{} ปจฺฉาชาตปจฺจเยน ปจฺจโย.\csromanspacer{} กมฺมปจฺจเยน ปจฺจโย. (สํขิตฺตํ.)}

\prose{นเหตุยา ตีณิ, น-อารมฺมเณ ตีณิ ฯเปฯ โน-อวิคเต ตีณิ.}

\prose{(ยถา กุสลตฺติเก ปญฺหาวารํ, เอวํ วิตฺถาเรตพฺพํ.)}

\nitthitam{\textbf{ธมฺมานุโลมปจฺจนีเย ทุกติกปฏฺฐานํ นิฏฺฐิตํ.}}
\pitaka{\textbf{อภิธมฺมปิฏก}}
\csromanheader{2}{\textbf{ธมฺมานุโลมปจฺจนีย}}
\gambhira{\textbf{ติกทุกปฏฺฐานปาฬิ}}
\namakkaram{นโม ตสฺส ภควโต อรหโต สมฺมาสมฺพุทฺธสฺส.}
\csromanheader{2}{1. กุสลตฺติก 1. เหตุทุก}
\csromanheader{2}{1-7. ปฏิจฺจาทิวาร ปจฺจยจตุกฺก-เหตุ}
\proseitem{1}{\csromanfolio{245}กุสลํ เหตุํ ธมฺมํ ปฏิจฺจ นกุสโล นเหตุ ธมฺโม อุปฺปชฺชติ เหตุปจฺจยา.\csromanspacer{}\csromanspacer{}กุสลํ เหตุํ ธมฺมํ ปฏิจฺจ น-อกุสโล นเหตุ ธมฺโม อุปฺปชฺชติ เหตุปจฺจยา.\csromanspacer{}\csromanspacer{}กุสลํ เหตุํ ธมฺมํ ปฏิจฺจ น-อพฺยากโต นเหตุ ธมฺโม อุปฺปชฺชติ เหตุปจฺจยา.\csromanspacer{}\csromanspacer{}กุสลํ เหตุํ ธมฺมํ ปฏิจฺจ น-อกุสโล นเหตุ จ น-อพฺยากโต นเหตุ จ ธมฺมา อุปฺปชฺชนฺติ เหตุปจฺจยา.\csromanspacer{}\csromanspacer{}กุสลํ เหตุํ ธมฺมํ ปฏิจฺจ นกุสโล นเหตุ จ น-อกุสโล นเหตุ จ ธมฺมา อุปฺปชฺชนฺติ เหตุปจฺจยา. (5)}

\prose{อกุสลํ เหตุํ ธมฺมํ ปฏิจฺจ น-อกุสโล นเหตุ ธมฺโม อุปฺปชฺชติ เหตุปจฺจยา.\csromanspacer{}\csromanspacer{}อกุสลํ เหตุํ ธมฺมํ ปฏิจฺจ นกุสโล นเหตุ ธมฺโม อุปฺปชฺชติ เหตุปจฺจยา.\csromanspacer{}\csromanspacer{}อกุสลํ เหตุํ ธมฺมํ ปฏิจฺจ น-อพฺยากโต นเหตุ ธมฺโม อุปฺปชฺชติ เหตุปจฺจยา.\csromanspacer{}\csromanspacer{}อกุสลํ เหตุํ ธมฺมํ ปฏิจฺจ นกุสโล นเหตุ จ น-อพฺยากโต นเหตุ จ ธมฺมา อุปฺปชฺชนฺติ เหตุปจฺจยา.\csromanspacer{}\csromanspacer{}อกุสลํ เหตุํ ธมฺมํ ปฏิจฺจ นกุสโล นเหตุ จ น-อกุสโล นเหตุ จ ธมฺมา อุปฺปชฺชนฺติ เหตุปจฺจยา. (5)}

\csromanheader{2}{ธมฺมานุโลมปจฺจนีย ติกทุกปฏฺฐานปาฬิ}
\prose{\csromanfolio{246}อพฺยากตํ เหตุํ ธมฺมํ ปฏิจฺจ นกุสโล นเหตุ ธมฺโม อุปฺปชฺชติ เหตุปจฺจยา.\csromanspacer{}\csromanspacer{}อพฺยากตํ เหตุํ ธมฺมํ ปฏิจฺจ น-อกุสโล นเหตุ ธมฺโม อุปฺปชฺชติ เหตุปจฺจยา.\csromanspacer{}\csromanspacer{}อพฺยากตํ เหตุํ ธมฺมํ ปฏิจฺจ นกุสโล นเหตุ จ น-อกุสโล นเหตุ จ ธมฺมา อุปฺปชฺชนฺติ เหตุปจฺจยา. (3) (สํขิตฺตํ.)}

\prose{เหตุยา เตรส, อารมฺมเณ นว ฯเปฯ อวิคเต เตรส.}

\csromanheader{2}{\textbf{ปจฺจนีย-น-อารมฺมณ}}
\proseitem{2}{กุสลํ เหตุํ ธมฺมํ ปฏิจฺจ นกุสโล นเหตุ ธมฺโม อุปฺปชฺชติ น-อารมฺมณปจฺจยา. (สํขิตฺตํ.)}

\prose{น-อารมฺมเณ นว, น-อธิปติยา เตรส ฯเปฯ นวิปฺปยุตฺเต นว.}

\prose{(สหชาตวารมฺปิ ฯเปฯ สมฺปยุตฺตวารมฺปิ ปฏิจฺจวารสทิสํ วิตฺถาเรตพฺพํ.)}

\csromanheader{2}{\textbf{เหตุ อารมฺมณ}}
\proseitem{3}{กุสโล เหตุ ธมฺโม นกุสลสฺส นเหตุสฺส ธมฺมสฺส เหตุปจฺจเยน ปจฺจโย.\csromanspacer{}\csromanspacer{}กุสโล เหตุ ธมฺโม น-อกุสลสฺส นเหตุสฺส ธมฺมสฺส เหตุปจฺจเยน ปจฺจโย.\csromanspacer{}\csromanspacer{}กุสโล เหตุ ธมฺโม น-อพฺยากตสฺส นเหตุสฺส ธมฺมสฺส เหตุปจฺจเยน ปจฺจโย.\csromanspacer{} กุสโล เหตุ ธมฺโม น-อกุสลสฺส นเหตุสฺส จ น-อพฺยากตสฺส นเหตุสฺส จ ธมฺมสฺส เหตุปจฺจเยน ปจฺจโย.\csromanspacer{}\csromanspacer{}กุสโล เหตุ ธมฺโม นกุสลสฺส นเหตุสฺส จ น-อกุสลสฺส นเหตุสฺส จ ธมฺมสฺส เหตุปจฺจเยน ปจฺจโย. (5)}

\prose{อกุสโล เหตุ ธมฺโม น-อกุสลสฺส นเหตุสฺส ธมฺมสฺส เหตุปจฺจเยน ปจฺจโย.\csromanspacer{} ปญฺจ.}

\prose{อพฺยากโต เหตุ ธมฺโม นกุสลสฺส นเหตุสฺส ธมฺมสฺส เหตุปจฺจเยน ปจฺจโย.\csromanspacer{} ตีณิ.}

\prose{กุสโล เหตุ ธมฺโม นกุสลสฺส นเหตุสฺส ธมฺมสฺส อารมฺมณปจฺจเยน ปจฺจโย. (สํขิตฺตํ.)}

\proseitem{4}{เหตุยา เตรส, อารมฺมเณ อฏฺฐารส ฯเปฯ อวิคเต เตรส. (ปญฺหาวารํ วิตฺถาเรตพฺพํ.)}

\csromanheader{2}{1. กุสลตฺติก 2. สเหตุกทุก}
\proseitem{5}{\csromanfolio{247}กุสลํ นเหตุํ ธมฺมํ ปฏิจฺจ น-อกุสโล นนเหตุ ธมฺโม อุปฺปชฺชติ เหตุปจฺจยา.\csromanspacer{}\csromanspacer{}กุสลํ นเหตุํ ธมฺมํ ปฏิจฺจ น-อพฺยากโต นนเหตุ ธมฺโม อุปฺปชฺชติ เหตุปจฺจยา. (สํขิตฺตํ.)}

\prose{เหตุยา นว, อารมฺมเณ นว ฯเปฯ อวิคเต นว. (สพฺพตฺถ นว.)}

\csromanheader{2}{1. กุสลตฺติก 2. สเหตุกทุก}
\csromanheader{2}{1-7. ปฏิจฺจาทิวาร ปจฺจยจตุกฺก-เหตุ}
\proseitem{6}{กุสลํ สเหตุกํ ธมฺมํ ปฏิจฺจ นกุสโล นสเหตุโก ธมฺโม อุปฺปชฺชติ เหตุปจฺจยา.\csromanspacer{}\csromanspacer{}กุสลํ สเหตุกํ ธมฺมํ ปฏิจฺจ น-อกุสโล นสเหตุโก ธมฺโม อุปฺปชฺชติ เหตุปจฺจยา.\csromanspacer{}\csromanspacer{}กุสลํ สเหตุกํ ธมฺมํ ปฏิจฺจ นกุสโล นสเหตุโก จ น-อกุสโล นสเหตุโก จ ธมฺมา อุปฺปชฺชนฺติ เหตุปจฺจยา.\csromanspacer{} ตีณิ.}

\prose{อกุสลํ สเหตุกํ ธมฺมํ ปฏิจฺจ น-อกุสโล นสเหตุโก ธมฺโม อุปฺปชฺชติ เหตุปจฺจยา.\csromanspacer{}\csromanspacer{}อกุสลํ สเหตุกํ ธมฺมํ ปฏิจฺจ นกุสโล นสเหตุโก ธมฺโม อุปฺปชฺชติ เหตุปจฺจยา.\csromanspacer{}\csromanspacer{}อกุสลํ สเหตุกํ ธมฺมํ ปฏิจฺจ นกุสโล นสเหตุโก จ น-อกุสโล นสเหตุโก จ ธมฺมา อุปฺปชฺชนฺติ เหตุปจฺจยา.\csromanspacer{} ตีณิ.}

\prose{อพฺยากตํ สเหตุกํ ธมฺมํ ปฏิจฺจ นกุสโล นสเหตุโก ธมฺโม อุปฺปชฺชติ เหตุปจฺจยา.\csromanspacer{}\csromanspacer{}อพฺยากตํ สเหตุกํ ธมฺมํ ปฏิจฺจ น-อกุสโล นสเหตุโก ธมฺโม อุปฺปชฺชติ เหตุปจฺจยา.\csromanspacer{}\csromanspacer{}อพฺยากตํ สเหตุกํ ธมฺมํ ปฏิจฺจ นกุสโล นสเหตุโก จ น-อกุสโล นสเหตุโก จ ธมฺมา อุปฺปชฺชนฺติ เหตุปจฺจยา.\csromanspacer{} ตีณิ.}

\prose{เหตุยา นว, อารมฺมเณ ตีณิ ฯเปฯ อวิคเต เอกาทส. (สหชาตวารมฺปิ ฯเปฯ สมฺปยุตฺตวารมฺปิ ปฏิจฺจวารสทิสํ.)}

\proseitem{7}{กุสโล สเหตุโก ธมฺโม นกุสลสฺส นสเหตุกสฺส ธมฺมสฺส เหตุปจฺจเยน ปจฺจโย.\csromanspacer{}\csromanspacer{}กุสโล สเหตุโก ธมฺโม น-อกุสลสฺส นสเหตุกสฺส ธมฺมสฺส เหตุปจฺจเยน ปจฺจโย.\csromanspacer{}\csromanspacer{}กุสโล สเหตุโก ธมฺโม นกุสลสฺส นสเหตุกสฺส จ น-อกุสลสฺส นสเหตุกสฺส จ ธมฺมสฺส เหตุปจฺจเยน ปจฺจโย.\csromanspacer{} ตีณิ.}

\csromanheader{2}{ธมฺมานุโลมปจฺจนีย ติกทุกปฏฺฐานปาฬิ}
\prose{\csromanfolio{248}อกุสโล สเหตุโก ธมฺโม น-อกุสลสฺส นสเหตุกสฺส ธมฺมสฺส เหตุปจฺจเยน ปจฺจโย.\csromanspacer{}\csromanspacer{}อกุสโล สเหตุโก ธมฺโม นกุสลสฺส นสเหตุกสฺส ธมฺมสฺส เหตุปจฺจเยน ปจฺจโย.\csromanspacer{}\csromanspacer{}อกุสโล สเหตุโก ธมฺโม นกุสลสฺส นสเหตุกสฺส จ น-อกุสลสฺส นสเหตุกสฺส จ ธมฺมสฺส เหตุปจฺจเยน ปจฺจโย.\csromanspacer{} ตีณิ.}

\prose{อพฺยากโต สเหตุโก ธมฺโม นกุสลสฺส นสเหตุกสฺส ธมฺมสฺส เหตุปจฺจเยน ปจฺจโย.\csromanspacer{}\csromanspacer{}อพฺยากโต สเหตุโก ธมฺโม น-อกุสลสฺส นสเหตุกสฺส ธมฺมสฺส เหตุปจฺจเยน ปจฺจโย.\csromanspacer{}\csromanspacer{}อพฺยากโต สเหตุโก ธมฺโม นกุสลสฺส นสเหตุกสฺส จ น-อกุสลสฺส นสเหตุกสฺส จ ธมฺมสฺส เหตุปจฺจเยน ปจฺจโย.\csromanspacer{} ตีณิ. (สํขิตฺตํ.)}

\proseitem{8}{เหตุยา นว, อารมฺมเณ ปนฺนรส ฯเปฯ อวิคเต เอกาทส.}

\proseitem{9}{อกุสลํ อเหตุกํ ธมฺมํ ปฏิจฺจ นกุสโล น-อเหตุโก ธมฺโม อุปฺปชฺชติ เหตุปจฺจยา.\csromanspacer{}\csromanspacer{}อกุสลํ อเหตุกํ ธมฺมํ ปฏิจฺจ น-อพฺยากโต น-อเหตุโก ธมฺโม อุปฺปชฺชติ เหตุปจฺจยา.\csromanspacer{}\csromanspacer{}อกุสลํ อเหตุกํ ธมฺมํ ปฏิจฺจ นกุสโล น-อเหตุโก จ น-อพฺยากโต น-อเหตุโก จ ธมฺมา อุปฺปชฺชนฺติ เหตุปจฺจยา.\csromanspacer{} ตีณิ.}

\prose{อพฺยากตํ อเหตุกํ ธมฺมํ ปฏิจฺจ นกุสโล น-อเหตุโก ธมฺโม อุปฺปชฺชติ เหตุปจฺจยา.\csromanspacer{}\csromanspacer{}อพฺยากตํ อเหตุกํ ธมฺมํ ปฏิจฺจ น-อกุสโล น-อเหตุโก ธมฺโม อุปฺปชฺชติ เหตุปจฺจยา.\csromanspacer{}\csromanspacer{}อพฺยากตํ อเหตุกํ ธมฺมํ ปฏิจฺจ นกุสโล น-อเหตุโก จ น-อกุสโล น-อเหตุโก จ ธมฺมา อุปฺปชฺชนฺติ เหตุปจฺจยา.\csromanspacer{} ตีณิ.}

\prose{เหตุยา ฉ, อารมฺมเณ ฉ ฯเปฯ อวิคเต ฉ.}

\csromanheader{2}{1. กุสลตฺติก 3. เหตุสมฺปยุตฺตทุก}
\proseitem{10}{กุสลํ เหตุสมฺปยุตฺตํ ธมฺมํ ปฏิจฺจ นกุสโล นเหตุสมฺปยุตฺโต ธมฺโม อุปฺปชฺชติ เหตุปจฺจยา. (สเหตุกทุกสทิสํ.)}

\csromanheader{2}{1. กุสลตฺติก 4. เหตุสเหตุกทุก}
\proseitem{11}{กุสลํ เหตุญฺเจว สเหตุกญฺจ ธมฺมํ ปฏิจฺจ น-อกุสโล นเหตุ เจว น-อเหตุโก จ ธมฺโม อุปฺปชฺชติ เหตุปจฺจยา.\csromanspacer{}\csromanspacer{}กุสลํ}

\vspace{\tipitakaparskip}
\csromancenter{กุสลตฺติก 4.\csromanspacer{} เหตุสเหตุกทุก เหตุญฺเจว สเหตุกญฺจ ธมฺมํ ปฏิจฺจ น-อพฺยากโต นเหตุ เจว น-อเหตุโก จ ธมฺโม อุปฺปชฺชติ เหตุปจฺจยา.\csromanspacer{}\csromanspacer{}กุสลํ เหตุญฺเจว สเหตุกญฺจ ธมฺมํ ปฏิจฺจ น-อกุสโล นเหตุ เจว น-อเหตุโก จ น-อพฺยากโต นเหตุ เจว น-อเหตุโก จ ธมฺมา อุปฺปชฺชนฺติ เหตุปจฺจยา. (3)}
\prose{\csromanfolio{249}อกุสลํ เหตุญฺเจว สเหตุญฺจ ธมฺมํ ปฏิจฺจ นกุสโล นเหตุ เจว น-อเหตุโก จ ธมฺโม อุปฺปชฺชติ เหตุปจฺจยา.\csromanspacer{} ตีณิ.}

\prose{อพฺยากตํ เหตุญฺเจว สเหตุกญฺจ ธมฺมํ ปฏิจฺจ นกุสโล นเหตุ เจว น-อเหตุโก จ ธมฺโม อุปฺปชฺชติ เหตุปจฺจยา.\csromanspacer{}\csromanspacer{}อพฺยากตํ เหตุญฺเจว สเหตุกญฺจ ธมฺมํ ปฏิจฺจ น-อกุสโล นเหตุ เจว น-อเหตุโก จ ธมฺโม อุปฺปชฺชติ เหตุปจฺจยา.\csromanspacer{}\csromanspacer{}อพฺยากตํ เหตุญฺเจว สเหตุกญฺจ ธมฺมํ ปฏิจฺจ นกุสโล นเหตุ เจว น-อเหตุโก จ น-อกุสโล นเหตุ เจว น-อเหตุโก จ ธมฺมา อุปฺปชฺชนฺติ เหตุปจฺจยา. (3) (สํขิตฺตํ.)}

\prose{เหตุยา นว.}

\proseitem{12}{กุสลํ สเหตุกญฺเจว น จ เหตุํ ธมฺมํ ปฏิจฺจ น-อกุสโล น-อเหตุโก เจว นนเหตุ จ ธมฺโม อุปฺปชฺชติ เหตุปจฺจยา.\csromanspacer{}\csromanspacer{}กุสลํ สเหตุกญฺเจว น จ เหตุํ ธมฺมํ ปฏิจฺจ น-อพฺยากโต น-อเหตุโก เจว นนเหตุ จ ธมฺโม อุปฺปชฺชติ เหตุปจฺจยา.\csromanspacer{}\csromanspacer{}กุสลํ สเหตุกญฺเจว น จ เหตุํ ธมฺมํ ปฏิจฺจ น-อกุสโล น-อเหตุโก เจว นนเหตุ จ น-อพฺยากโต น-อเหตุโก เจว นนเหตุ จ ธมฺมา อุปฺปชฺชนฺติ เหตุปจฺจยา.\csromanspacer{} ตีณิ.}

\prose{อกุสลํ สเหตุกญฺเจว น จ เหตุํ ธมฺมํ ปฏิจฺจ นกุสโล น-อเหตุโก เจว นนเหตุ จ ธมฺโม อุปฺปชฺชติ เหตุปจฺจยา.\csromanspacer{}\csromanspacer{}อกุสลํ สเหตุกญฺเจว น จ เหตุํ ธมฺมํ ปฏิจฺจ น-อพฺยากโต น-อเหตุโก เจว นนเหตุ จ ธมฺโม อุปฺปชฺชติ เหตุปจฺจยา.\csromanspacer{}\csromanspacer{}อกุสลํ สเหตุกญฺเจว น จ เหตุํ ธมฺมํ ปฏิจฺจ น-อกุสโล น-อเหตุโก เจว นนเหตุ จ น-อพฺยากโต น-อเหตุโก เจว นนเหตุ จ ธมฺมา อุปฺปชฺชนฺติ เหตุปจฺจยา.\csromanspacer{} ตีณิ.}

\prose{อพฺยากตํ สเหตุกญฺเจว น จ เหตุํ ธมฺมํ ปฏิจฺจ นกุสโล น-อเหตุโก เจว นนเหตุ จ ธมฺโม อุปฺปชฺชติ เหตุปจฺจยา.\csromanspacer{}\csromanspacer{}}

\vspace{\tipitakaparskip}
\csromancenter{ธมฺมานุโลมปจฺจนีย ติกทุกปฏฺฐานปาฬิ อพฺยากตํ สเหตุกญฺเจว น จ เหตุํ ธมฺมํ ปฏิจฺจ น-อกุสโล น-อเหตุโก เจว นนเหตุ จ ธมฺโม อุปฺปชฺชติ เหตุปจฺจยา.\csromanspacer{}\csromanspacer{}อพฺยากตํ สเหตุกญฺเจว น จ เหตุํ ธมฺมํ ปฏิจฺจ นกุสโล น-อเหตุโก เจว นนเหตุ จ น-อกุสโล น-อเหตุโก เจว นนเหตุ จ ธมฺมา อุปฺปชฺชนฺติ เหตุปจฺจยา.\csromanspacer{} ตีณิ.}
\csromanheader{2}{1. กุสลตฺติก 5. เหตุเหตุสมฺปยุตฺตทุก}
\proseitem{13}{\csromanfolio{250}กุสลํ เหตุญฺเจว เหตุสมฺปยุตฺตญฺจ ธมฺมํ ปฏิจฺจ น-อกุสโล นเหตุ เจว นเหตุวิปฺปยุตฺโต จ ธมฺโม อุปฺปชฺชติ เหตุปจฺจยา.\csromanspacer{}\csromanspacer{}กุสลํ เหตุญฺเจว เหตุสมฺปยุตฺตญฺจ ธมฺมํ ปฏิจฺจ น-อพฺยากโต นเหตุ เจว นเหตุวิปฺปยุตฺโต จ ธมฺโม อุปฺปชฺชติ เหตุปจฺจยา.\csromanspacer{}\csromanspacer{}กุสลํ เหตุญฺเจว เหตุสมฺปยุตฺตญฺจ ธมฺมํ ปฏิจฺจ น-อกุสโล นเหตุ เจว นเหตุวิปฺปยุตฺโต จ น-อพฺยากโต นเหตุ เจว นเหตุวิปฺปยุตฺโต จ ธมฺมา อุปฺปชฺชนฺติ เหตุปจฺจยา.\csromanspacer{} ตีณิ.}

\prose{(อกุสลํ ตีณิ กาตพฺพํ.\csromanspacer{} อพฺยากตํ ตีณิ กาตพฺพํ.)}

\prose{เหตุยา นว.}

\proseitem{14}{กุสลํ เหตุสมฺปยุตฺตญฺเจว น จ เหตุํ ธมฺมํ ปฏิจฺจ น-อกุสโล นเหตุวิปฺปยุตฺโต เจว นนเหตุ จ ธมฺโม อุปฺปชฺชติ เหตุปจฺจยา.\csromanspacer{}\csromanspacer{}กุสลํ เหตุสมฺปยุตฺตญฺเจว น จ เหตุํ ธมฺมํ ปฏิจฺจ น-อพฺยากโต นเหตุวิปฺปยุตฺโต เจว นนเหตุ จ ธมฺโม อุปฺปชฺชติ เหตุปจฺจยา.\csromanspacer{}\csromanspacer{}กุสลํ เหตุสมฺปยุตฺตญฺเจว น จ เหตุํ ธมฺมํ ปฏิจฺจ น-อกุสโล นเหตุวิปฺปยุตฺโต เจว นนเหตุ จ น-อพฺยากโต นเหตุวิปฺปยุตฺโต เจว นนเหตุ จ ธมฺมา อุปฺปชฺชนฺติ เหตุปจฺจยา.\csromanspacer{} ตีณิ.}

\prose{อกุสลํ เหตุสมฺปยุตฺตญฺเจว น จ เหตุํ ธมฺมํ ปฏิจฺจ.\csromanspacer{} ตีณิ.\csromanspacer{} อพฺยากตํ ตีณิ.\csromanspacer{}\csromanspacer{}เหตุยา นว.}

\csromanheader{2}{1. กุสลตฺติก 6. เหตุสเหตุกทุก}
\proseitem{15}{กุสลํ นเหตุํ สเหตุกํ ธมฺมํ ปฏิจฺจ นกุสโล นเหตุ นสเหตุโก ธมฺโม อุปฺปชฺชติ เหตุปจฺจยา.\csromanspacer{}\csromanspacer{}กุสลํ นเหตุํ สเหตุกํ ธมฺมํ ปฏิจฺจ น-อกุสโล นเหตุ นสเหตุโก ธมฺโม}

\vspace{\tipitakaparskip}
\csromancenter{กุสลตฺติก 7.\csromanspacer{} สปฺปจฺจยทุก อุปฺปชฺชติ เหตุปจฺจยา.\csromanspacer{}\csromanspacer{}กุสลํ นเหตุํ สเหตุกํ ธมฺมํ ปฏิจฺจ นกุสโล นเหตุ นสเหตุโก จ น-อกุสโล นเหตุ นสเหตุโก จ ธมฺมา อุปฺปชฺชนฺติ เหตุปจฺจยา.\csromanspacer{} ตีณิ.}
\prose{\csromanfolio{251}อกุสลํ นเหตุํ สเหตุกํ.\csromanspacer{} ตีณิ.}

\prose{อพฺยากตํ นเหตุํ สเหตุกํ ธมฺมํ ปฏิจฺจ นกุสโล ฯเปฯ.\csromanspacer{} น-อกุสโล ฯเปฯ.\csromanspacer{} นกุสโล นเหตุ นสเหตุโก จ น-อกุสโล นเหตุ นสเหตุโก จ ธมฺมา อุปฺปชฺชนฺติ เหตุปจฺจยา.\csromanspacer{} ตีณิ.}

\proseitem{16}{อพฺยากตํ นเหตุํ อเหตุกํ ธมฺมํ ปฏิจฺจ นกุสโล นเหตุ น-อเหตุโก ธมฺโม อุปฺปชฺชติ เหตุปจฺจยา.\csromanspacer{}\csromanspacer{}อพฺยากตํ นเหตุํ อเหตุกํ ธมฺมํ ปฏิจฺจ น-อกุสโล นเหตุ น-อเหตุโก ธมฺโม อุปฺปชฺชติ เหตุปจฺจยา.\csromanspacer{}\csromanspacer{}อพฺยากตํ นเหตุํ อเหตุกํ ธมฺมํ ปฏิจฺจ นกุสโล นเหตุ น-อเหตุโก จ น-อกุสโล นเหตุ น-อเหตุโก จ ธมฺมา อุปฺปชฺชนฺติ เหตุปจฺจยา. (3)}

\nitthitam{เหตุโคจฺฉกํ นิฏฺฐิตํ.}
\csromanheader{2}{1. กุสลตฺติก 7. สปฺปจฺจยทุก}
\proseitem{17}{อพฺยากโต อปฺปจฺจโย ธมฺโม น-อพฺยากตสฺส น-อปฺปจฺจยสฺส ธมฺมสฺส อารมฺมณปจฺจเยน ปจฺจโย.\csromanspacer{}\csromanspacer{}อพฺยากโต อปฺปจฺจโย ธมฺโม นกุสลสฺส น-อปฺปจฺจยสฺส ธมฺมสฺส อารมฺมณปจฺจเยน ปจฺจโย.\csromanspacer{}\csromanspacer{}อพฺยากโต อปฺปจฺจโย ธมฺโม น-อกุสลสฺส น-อปฺปจฺจยสฺส ธมฺมสฺส อารมฺมณปจฺจเยน ปจฺจโย.\csromanspacer{}\csromanspacer{}อพฺยากโต อปฺปจฺจโย ธมฺโม น-อกุสลสฺส น อปฺปจฺจยสฺส จ น-อพฺยากตสฺส น-อปฺปจฺจยสฺส จ ธมฺมสฺส อารมฺมณปจฺจเยน ปจฺจโย.\csromanspacer{}\csromanspacer{}อพฺยากโต อปฺปจฺจโย ธมฺโม นกุสลสฺส น-อปฺปจฺจยสฺส จ น-อกุสลสฺส น-อปฺปจฺจยสฺส จ ธมฺมสฺส อารมฺมณปจฺจเยน ปจฺจโย. (5)}

\csromanheader{2}{อารมฺมเณ ปญฺจ. (อสงฺขตํ อปฺปจฺจยสทิสํ.)}
\csromanheader{2}{ธมฺมานุโลมปจฺจนีย ติกทุกปฏฺฐานปาฬิ \textbf{1. กุสลตฺติก 9. สนิทสฺสนทุก}}
\proseitem{18}{\csromanfolio{252}อพฺยากโต สนิทสฺสโน ธมฺโม น-อพฺยากตสฺส นสนิทสฺสนสฺส ธมฺมสฺส อารมฺมณปจฺจเยน ปจฺจโย.\csromanspacer{}\csromanspacer{}อพฺยากโต สนิทสฺสโน ธมฺโม นกุสลสฺส นสนิทสฺสนสฺส ฯเปฯ. (ฉ ปญฺหา กาตพฺพา.)}

\prose{กุสลํ อนิทสฺสนํ ธมฺมํ ปฏิจฺจ นกุสโล น-อนิทสฺสโน ธมฺโม อุปฺปชฺชติ เหตุปจฺจยา.\csromanspacer{}\csromanspacer{}อกุสเลน ตีณิเยว.\csromanspacer{} อพฺยากเตน ตีณิเยว.\csromanspacer{}\csromanspacer{}เหตุยา นว.}

\csromanheader{2}{1. กุสลตฺติก 10. สปฺปฏิฆทุก}
\proseitem{19}{อพฺยากตํ สปฺปฏิฆํ ธมฺมํ ปฏิจฺจ นกุสโล นสปฺปฏิโฆ ธมฺโม อุปฺปชฺชติ เหตุปจฺจยา.\csromanspacer{}\csromanspacer{}อพฺยากตํ สปฺปฏิฆํ ธมฺมํ ปฏิจฺจ น-อกุสโล นสปฺปฏิโฆ ธมฺโม อุปฺปชฺชติ เหตุปจฺจยา.\csromanspacer{}\csromanspacer{}อพฺยากตํ สปฺปฏิฆํ ธมฺมํ ปฏิจฺจ นกุสโล นสปฺปฏิโฆ จ น-อกุสโล นสปฺปฏิโฆ จ ธมฺมา อุปฺปชฺชนฺติ เหตุปจฺจยา.\csromanspacer{}\csromanspacer{}เหตุยา ตีณิ.}

\prose{กุสลํ อปฺปฏิเฆน ตีณิ.\csromanspacer{} อกุสลํ อปฺปฏิเฆน ตีณิ.\csromanspacer{} อพฺยากตํ อปฺปฏิเฆน ตีณิ.\csromanspacer{} กุสลํ อปฺปฏิฆญฺจ อพฺยากตํ อปฺปฏิฆญฺจ ตีณิ.\csromanspacer{} อกุสลํ อปฺปฏิฆญฺจ อพฺยากตํ อปฺปฏิฆญฺจ ตีณิ.\csromanspacer{}\csromanspacer{}เหตุยา ปนฺนรส.}

\csromanheader{2}{1. กุสลตฺติก 11. รูปีทุก}
\proseitem{20}{อพฺยากตํ รูปิํ ธมฺมํ ปฏิจฺจ นกุสโล นรูปี ธมฺโม อุปฺปชฺชติ เหตุปจฺจยา.\csromanspacer{}\csromanspacer{}อพฺยากตํ รูปิํ ธมฺมํ ปฏิจฺจ น-อกุสโล นรูปี ธมฺโม อุปฺปชฺชติ เหตุปจฺจยา.\csromanspacer{}\csromanspacer{}อพฺยากตํ รูปิํ ธมฺมํ ปฏิจฺจ นกุสโล นรูปี จ น-อกุสโล น รูปี จ ธมฺมา อุปฺปชฺชนฺติ เหตุปจฺจยา.\csromanspacer{}\csromanspacer{}เหตุยา ตีณิ.}

\prose{กุสลํ อรูปิํ ธมฺมํ ปฏิจฺจ ตีณิ.\csromanspacer{} อกุสลํ อรูปิํ ธมฺมํ ปฏิจฺจ ตีณิ.\csromanspacer{} อพฺยากตํ อรูปิํ ธมฺมํ ปฏิจฺจ ตีณิ.\csromanspacer{}\csromanspacer{}เหตุยา นว.}

\csromanheader{2}{1. กุสลตฺติก 12. โลกิยทุก}
\proseitem{21}{อพฺยากตํ โลกิยํ ธมฺมํ ปจฺจยา น-อพฺยากโต นโลกิโย ธมฺโม อุปฺปชฺชติ เหตุปจฺจยา.\csromanspacer{}\csromanspacer{}อพฺยากตํ โลกิยํ ธมฺมํ ปจฺจยา นกุสโล นโลกิโย ธมฺโม อุปฺปชฺชติ เหตุปจฺจยา.\csromanspacer{}\csromanspacer{}อพฺยากตํ}

\vspace{\tipitakaparskip}
\csromancenter{กุสลตฺติก 13.\csromanspacer{} เกนจิวิญฺเญยฺยทุก โลกิยํ ธมฺมํ ปจฺจยา น-อกุสโล นโลกิโย ธมฺโม อุปฺปชฺชติ เหตุปจฺจยา.\csromanspacer{}\csromanspacer{}อพฺยากตํ โลกิยํ ธมฺมํ ปจฺจยา น-อกุสโล นโลกิโย จ น-อพฺยากโต นโลกิโย จ ธมฺมา อุปฺปชฺชนฺติ เหตุปจฺจยา.\csromanspacer{}\csromanspacer{}อพฺยากตํ โลกิยํ ธมฺมํ ปจฺจยา นกุสโล นโลกิโย จ น-อกุสโล นโลกิโย จ ธมฺมา อุปฺปชฺชนฺติ เหตุปจฺจยา.\csromanspacer{}\csromanspacer{}เหตุยา ปญฺจ.}
\prose{\csromanfolio{253}กุสลํ โลกุตฺตรํ ธมฺมํ ปฏิจฺจ นกุสโล นโลกุตฺตโร ธมฺโม อุปฺปชฺชติ เหตุปจฺจยา.\csromanspacer{}\csromanspacer{}กุสลํ โลกุตฺตรํ ธมฺมํ ปฏิจฺจ น-อกุสโล นโลกุตฺตโร ธมฺโม อุปฺปชฺชติ เหตุปจฺจยา.\csromanspacer{}\csromanspacer{}กุสลํ โลกุตฺตรํ ธมฺมํ ปฏิจฺจ นกุสโล นโลกุตฺตโร จ น-อกุสโล นโลกุตฺตโร จ ธมฺมา อุปฺปชฺชนฺติ เหตุปจฺจยา.\csromanspacer{} ตีณิ.}

\prose{อพฺยากตํ โลกุตฺตรํ ธมฺมํ ปฏิจฺจ นกุสโล นโลกุตฺตโร ธมฺโม อุปฺปชฺชติ เหตุปจฺจยา.\csromanspacer{}\csromanspacer{}อพฺยากตํ โลกุตฺตรํ ธมฺมํ ปฏิจฺจ น-อกุสโล นโลกุตฺตโร ธมฺโม อุปฺปชฺชติ เหตุปจฺจยา.\csromanspacer{}\csromanspacer{}อพฺยากตํ โลกุตฺตรํ ธมฺมํ ปฏิจฺจ นกุสโล นโลกุตฺตโร จ น-อกุสโล นโลกุตฺตโร จ ธมฺมา อุปฺปชฺชนฺติ เหตุปจฺจยา.\csromanspacer{} ตีณิ.}

\csromanheader{2}{1. กุสลตฺติก 13. เกนจิวิญฺเญยฺยทุก}
\proseitem{22}{กุสลํ เกนจิ วิญฺเญยฺยํ ธมฺมํ ปฏิจฺจ นกุสโล นเกนจิ วิญฺเญยฺโย ธมฺโม อุปฺปชฺชติ เหตุปจฺจยา.\csromanspacer{}\csromanspacer{}กุสลํ เกนจิ วิญฺเญยฺยํ ธมฺมํ ปฏิจฺจ น-อกุสโล นเกนจิ วิญฺเญยฺโย ธมฺโม อุปฺปชฺชติ เหตุปจฺจยา.\csromanspacer{}\csromanspacer{}กุสลํ เกนจิ วิญฺเญยฺยํ ธมฺมํ ปฏิจฺจ น-อพฺยากโต นเกนจิ วิญฺเญยฺโย ธมฺโม อุปฺปชฺชติ เหตุปจฺจยา. (สํขิตฺตํ.) .\csromanspacer{} เหตุยา เอกูนวีสติ.}

\proseitem{23}{กุสลํ เกนจิ นวิญฺเญยฺยํ ธมฺมํ ปฏิจฺจ นกุสโล นเกนจิ นวิญฺเญยฺโย ธมฺโม อุปฺปชฺชติ เหตุปจฺจยา.\csromanspacer{}\csromanspacer{}กุสลํ เกนจิ นวิญฺเญยฺยํ ธมฺมํ ปฏิจฺจ น-อกุสโล นเกนจิ นวิญฺเญยฺโย ธมฺโม อุปฺปชฺชติ เหตุปจฺจยา.\csromanspacer{}\csromanspacer{}กุสลํ เกนจิ นวิญฺเญยฺยํ ธมฺมํ ปฏิจฺจ น-อพฺยากโต นเกนจิ นวิญฺเญยฺโย ธมฺโม อุปฺปชฺชติ เหตุปจฺจยา.}

\prose{(เอเตน อุปาเยน เกนจิ นวิญฺเญยฺเย เอกูนวีสติ ปญฺหา กาตพฺพา.)}

\csromanheader{2}{ธมฺมานุโลมปจฺจนีย ติกทุกปฏฺฐานปาฬิ \textbf{1. กุสลตฺติก 14. อาสวโคจฺฉก}}
\proseitem{24}{\csromanfolio{254}อกุสลํ อาสวํ ธมฺมํ ปฏิจฺจ น-อกุสโล น-อาสโว ธมฺโม ฯเปฯ.\csromanspacer{}\csromanspacer{}นกุสโล น-อาสโว ธมฺโม ฯเปฯ.\csromanspacer{}\csromanspacer{}น-อพฺยากโต น-อาสโว ธมฺโม ฯเปฯ.\csromanspacer{}\csromanspacer{}นกุสโล น-อาสโว จ น-อพฺยากโต น-อาสโว จ ธมฺมา ฯเปฯ.\csromanspacer{}\csromanspacer{}นกุสโล น-อาสโว จ น-อกุสโล น-อาสโว จ ธมฺมา อุปฺปชฺชนฺติ เหตุปจฺจยา.\csromanspacer{}\csromanspacer{}เหตุยา ปญฺจ, อารมฺมเณ ตีณิ ฯเปฯ อวิคเต ปญฺจ.}

\prose{อกุสลํ โน-อาสวํ ธมฺมํ ปฏิจฺจ นกุสโล นโน-อาสโว ธมฺโม อุปฺปชฺชติ เหตุปจฺจยา.\csromanspacer{}\csromanspacer{}อกุสลํ โน-อาสวํ ธมฺมํ ปฏิจฺจ น-อพฺยากโต นโน-อาสโว ธมฺโม ฯเปฯ.\csromanspacer{}\csromanspacer{}นกุสโล นโน-อาสโว จ น-อพฺยากโต นโน-อาสโว จ ธมฺมา อุปฺปชฺชนฺติ เหตุปจฺจยา.\csromanspacer{}\csromanspacer{}เหตุยา ตีณิ ฯเปฯ อวิคเต ตีณิ.}

\csromanheader{2}{25. อพฺยากตํ สาสวํ ธมฺมํ ปฏิจฺจ ฯเปฯ. (โลกิยทุกสทิสํ.)}
\proseitem{26}{อกุสลํ อาสวสมฺปยุตฺตํ ธมฺมํ ปฏิจฺจ น-อกุสโล น-อาสวสมฺปยุตฺโต ธมฺโม ฯเปฯ.\csromanspacer{}\csromanspacer{}นกุสโล น-อาสวสมฺปยุตฺโต ธมฺโม ฯเปฯ.\csromanspacer{}\csromanspacer{}น-อพฺยากโต น-อาสวสมฺปยุตฺโต ธมฺโม ฯเปฯ.\csromanspacer{}\csromanspacer{}นกุสโล น-อาสวสมฺปยุตฺโต จ น-อพฺยากโต น-อาสวสมฺปยุตฺโต จ ธมฺมา ฯเปฯ.\csromanspacer{}\csromanspacer{}นกุสโล น-อาสวสมฺปยุตฺโต จ น-อกุสโล น-อาสวสมฺปยุตฺโต จ ธมฺมา อุปฺปชฺชนฺติ เหตุปจฺจยา.\csromanspacer{}\csromanspacer{}เหตุยา ปญฺจ.}

\prose{อกุสลํ อาสววิปฺปยุตฺตํ ธมฺมํ ปฏิจฺจ นกุสโล น-อาสววิปฺปยุตฺโต ธมฺโม ฯเปฯ.\csromanspacer{}\csromanspacer{}น-อพฺยากโต น-อาสววิปฺปยุตฺโต ธมฺโม ฯเปฯ.\csromanspacer{}\csromanspacer{}นกุสโล น-อาสววิปฺปยุตฺโต จ น-อพฺยากโต น-อาสววิปฺปยุตฺโต จ ธมฺมา อุปฺปชฺชนฺติ เหตุปจฺจยา.\csromanspacer{}\csromanspacer{}เหตุยา ตีณิ.}

\proseitem{27}{อกุสลํ อาสวญฺเจว สาสวญฺจ ธมฺมํ ปฏิจฺจ น-อกุสโล น-อาสโว เจว น-อนาสโว จ ธมฺโม ฯเปฯ.\csromanspacer{}\csromanspacer{}นกุสโล น-อาสโว เจว น-อนาสโว จ ธมฺโม ฯเปฯ.\csromanspacer{}\csromanspacer{}น-อพฺยากโต น-อาสโว เจว น-อนาสโว จ ธมฺโม ฯเปฯ.\csromanspacer{}\csromanspacer{}นกุสโล น-อาสโว เจว น-อนาสโว จ น-อพฺยากโต น-อาสโว เจว น-อนาสโว จ ธมฺมา ฯเปฯ.\csromanspacer{}\csromanspacer{}นกุสโล}

\vspace{\tipitakaparskip}
\csromancenter{กุสลตฺติก 55.\csromanspacer{} สารมฺมณทุก น-อาสโว เจว น-อนาสโว จ น-อกุสโล น-อาสโว เจว น-อนาสโว จ ธมฺมา อุปฺปชฺชนฺติ เหตุปจฺจยา.\csromanspacer{}\csromanspacer{}เหตุยา ปญฺจ.}
\prose{\csromanfolio{255}อกุสลํ สาสวญฺเจว โน จ อาสวํ ธมฺมํ ปฏิจฺจ นกุสโล น-อนาสโว เจว นโน จ อาสโว ธมฺโม ฯเปฯ.\csromanspacer{}\csromanspacer{}น-อพฺยากโต น-อนาสโว เจว นโน จ อาสโว ธมฺโม ฯเปฯ.\csromanspacer{}\csromanspacer{}นกุสโล น-อนาสโว เจว นโน จ อาสโว น-อพฺยากโต น-อนาสโว เจว นโน-อาสโว จ ธมฺมา อุปฺปชฺชนฺติ เหตุปจฺจยา.\csromanspacer{}\csromanspacer{}เหตุยา ตีณิ.}

\proseitem{28}{อกุสลํ อาสวญฺเจว อาสวสมฺปยุตฺตญฺจ ธมฺมํ ปฏิจฺจ นกุสโล น-อาสโว เจว น-อาสววิปฺปยุตฺโต จ ธมฺโม ฯเปฯ.\csromanspacer{}\csromanspacer{}น-อพฺยากโต น-อาสโว เจว น-อาสววิปฺปยุตฺโต จ ธมฺโม ฯเปฯ.\csromanspacer{}\csromanspacer{}นกุสโล น-อาสโว เจว น-อาสววิปฺปยุตฺโต จ น-อพฺยากโต น-อาสโว เจว น-อาสววิปฺปยุตฺโต จ ธมฺมา อุปฺปชฺชนฺติ เหตุปจฺจยา.\csromanspacer{}\csromanspacer{}เหตุยา ตีณิ.}

\prose{อกุสลํ อาสวสมฺปยุตฺตญฺเจว โน จ อาสวํ ธมฺมํ ปฏิจฺจ นกุสโล น-อาสววิปฺปยุตฺโต เจว นโน จ อาสโว ธมฺโม ฯเปฯ.\csromanspacer{} น-อพฺยากโต น-อาสววิปฺปยุตฺโต เจว นโน จ อาสโว ธมฺโม ฯเปฯ.\csromanspacer{} นกุสโล น-อาสววิปฺปยุตฺโต เจว นโน อาสโว จ น-อพฺยากโต น-อาสววิปฺปยุตฺโต เจว นโน จ อาสโว จ ธมฺมา อุปฺปชฺชนฺติ เหตุปจฺจยา.\csromanspacer{}\csromanspacer{}เหตุยา ตีณิ.}

\proseitem{29}{อพฺยากตํ อาสววิปฺปยุตฺตํ สาสวํ ธมฺมํ ปฏิจฺจ ฯเปฯ. (โลกิยทุกสทิสํ.)}

\csromanheader{2}{1. กุสลตฺติก 20. ฉโคจฺฉกทุก}
\proseitem{30}{อกุสลํ สญฺโญชนํ ธมฺมํ ปฏิจฺจ ฯเปฯ.\csromanspacer{} คนฺถํ ฯเปฯ.\csromanspacer{} โอฆํ ฯเปฯ.\csromanspacer{} โยคํ ฯเปฯ.\csromanspacer{} นีวรณํ ฯเปฯ.\csromanspacer{} ปรามาสํ. (สํขิตฺตํ.)}

\csromanheader{2}{1. กุสลตฺติก 55. สารมฺมณทุก}
\proseitem{31}{กุสลํ สารมฺมณํ ธมฺมํ ปฏิจฺจ นกุสโล นสารมฺมโณ ธมฺโม อุปฺปชฺชติ เหตุปจฺจยา.\csromanspacer{}\csromanspacer{}กุสลํ สารมฺมณํ ธมฺมํ ปฏิจฺจ น-อกุสโล นสารมฺมโณ ธมฺโม อุปฺปชฺชติ เหตุปจฺจยา.\csromanspacer{}\csromanspacer{}กุสลํ สารมฺมณํ ธมฺมํ}

\vspace{\tipitakaparskip}
\csromancenter{ธมฺมานุโลมปจฺจนีย ติกทุกปฏฺฐานปาฬิ ปฏิจฺจ นกุสโล นสารมฺมโณ จ น-อกุสโล นสารมฺมโณ จ ธมฺมา อุปฺปชฺชนฺติ เหตุปจฺจยา.\csromanspacer{} ตีณิ.}
\prose{\csromanfolio{256}อกุสลํ สารมฺมณํ ธมฺมํ ปฏิจฺจ น-อกุสโล นสารมฺมโณ ธมฺโม อุปฺปชฺชติ เหตุปจฺจยา.\csromanspacer{}\csromanspacer{}อกุสลํ สารมฺมณํ ธมฺมํ ปฏิจฺจ นกุสโล นสารมฺมโณ ธมฺโม อุปฺปชฺชติ เหตุปจฺจยา.\csromanspacer{}\csromanspacer{}อกุสลํ สารมฺมณํ ธมฺมํ ปฏิจฺจ นกุสโล นสารมฺมโณ จ น-อกุสโล นสารมฺมโณ จ ธมฺมา อุปฺปชฺชนฺติ เหตุปจฺจยา.\csromanspacer{} ตีณิ.}

\prose{อพฺยากตํ สารมฺมณํ ธมฺมํ ปฏิจฺจ นกุสโล นสารมฺมโณ ธมฺโม อุปฺปชฺชติ เหตุปจฺจยา.\csromanspacer{} ตีณิ.}

\prose{เหตุยา นว.}

\proseitem{32}{อพฺยากตํ อนารมฺมณํ ธมฺมํ ปฏิจฺจ นกุสโล น-อนารมฺมโณ ธมฺโม อุปฺปชฺชติ เหตุปจฺจยา.\csromanspacer{}\csromanspacer{}อพฺยากตํ อนารมฺมณํ ธมฺมํ ปฏิจฺจ น-อกุสโล น-อนารมฺมโณ ธมฺโม อุปฺปชฺชติ เหตุปจฺจยา.\csromanspacer{}\csromanspacer{}อพฺยากตํ อนารมฺมณํ ธมฺมํ ปฏิจฺจ นกุสโล น-อนารมฺมโณ จ น-อกุสโล น-อนารมฺมโณ จ ธมฺมา อุปฺปชฺชนฺติ เหตุปจฺจยา. (สํขิตฺตํ.) .\csromanspacer{} เหตุยา ตีณิ.}

\csromanheader{2}{1. กุสลตฺติก 56. จิตฺตทุก}
\proseitem{33}{กุสลํ จิตฺตํ ธมฺมํ ปฏิจฺจ นกุสโล นจิตฺโต ธมฺโม อุปฺปชฺชติ เหตุปจฺจยา.\csromanspacer{}\csromanspacer{}กุสลํ จิตฺตํ ธมฺมํ ปฏิจฺจ น-อกุสโล นจิตฺโต ธมฺโม อุปฺปชฺชติ เหตุปจฺจยา.\csromanspacer{}\csromanspacer{}กุสลํ จิตฺตํ ธมฺมํ ปฏิจฺจ น-อพฺยากโต นจิตฺโต ธมฺโม อุปฺปชฺชติ เหตุปจฺจยา.\csromanspacer{}\csromanspacer{}กุสลํ จิตฺตํ ธมฺมํ ปฏิจฺจ น-อกุสโล นจิตฺโต จ น-อพฺยากโต นจิตฺโต จ ธมฺมา อุปฺปชฺชนฺติ เหตุปจฺจยา.\csromanspacer{}\csromanspacer{}กุสลํ จิตฺตํ ธมฺมํ ปฏิจฺจ นกุสโล นจิตฺโต จ น-อกุสโล นจิตฺโต จ ธมฺมา อุปฺปชฺชนฺติ เหตุปจฺจยา.\csromanspacer{} ปญฺจ.}

\prose{อกุสลํ จิตฺตํ ธมฺมํ ปฏิจฺจ น-อกุสโล นจิตฺโต ธมฺโม อุปฺปชฺชติ เหตุปจฺจยา.\csromanspacer{}\csromanspacer{}อกุสลํ จิตฺตํ ธมฺมํ ปฏิจฺจ นกุสโล นจิตฺโต ธมฺโม อุปฺปชฺชติ เหตุปจฺจยา.\csromanspacer{}\csromanspacer{}อกุสลํ จิตฺตํ ธมฺมํ ปฏิจฺจ น-อพฺยากโต นจิตฺโต ธมฺโม อุปฺปชฺชติ เหตุปจฺจยา.\csromanspacer{}\csromanspacer{}อกุสลํ จิตฺตํ ธมฺมํ ปฏิจฺจ นกุสโล นจิตฺโต จ น-อพฺยากโต นจิตฺโต จ ธมฺมา อุปฺปชฺชนฺติ เหตุปจฺจยา.\csromanspacer{}\csromanspacer{}}

\vspace{\tipitakaparskip}
\csromancenter{กุสลตฺติก 57.\csromanspacer{} เจตสิกทุก อกุสลํ จิตฺตํ ธมฺมํ ปฏิจฺจ นกุสโล นจิตฺโต จ น-อกุสโล นจิตฺโต จ ธมฺมา อุปฺปชฺชนฺติ เหตุปจฺจยา.\csromanspacer{} ปญฺจ.}
\prose{\csromanfolio{257}อพฺยากตํ จิตฺตํ ธมฺมํ ปฏิจฺจ นกุสโล นจิตฺโต ธมฺโม อุปฺปชฺชติ เหตุปจฺจยา.\csromanspacer{}\csromanspacer{}อพฺยากตํ จิตฺตํ ธมฺมํ ปฏิจฺจ น-อกุสโล นจิตฺโต ธมฺโม อุปฺปชฺชติ เหตุปจฺจยา.\csromanspacer{}\csromanspacer{}อพฺยากตํ จิตฺตํ ธมฺมํ ปฏิจฺจ นกุสโล นจิตฺโต จ น-อกุสโล นจิตฺโต จ ธมฺมา อุปฺปชฺชนฺติ เหตุปจฺจยา.\csromanspacer{} ตีณิ.\csromanspacer{} เหตุยา เตรส.}

\proseitem{34}{กุสลํ โนจิตฺตํ ธมฺมํ ปฏิจฺจ น-อกุสโล นโนจิตฺโต ธมฺโม อุปฺปชฺชติ เหตุปจฺจยา.\csromanspacer{}\csromanspacer{}กุสลํ โนจิตฺตํ ธมฺมํ ปฏิจฺจ น-อพฺยากโต นโนจิตฺโต ธมฺโม อุปฺปชฺชติ เหตุปจฺจยา.\csromanspacer{}\csromanspacer{}กุสลํ โนจิตฺตํ ธมฺมํ ปฏิจฺจ น-อกุสโล นโนจิตฺโต จ น-อพฺยากโต นโนจิตฺโต จ ธมฺมา อุปฺปชฺชนฺติ เหตุปจฺจยา.\csromanspacer{} ตีณิ.}

\prose{อกุสลํ โนจิตฺตํ ธมฺมํ ปฏิจฺจ นกุสโล นโนจิตฺโต ธมฺโม อุปฺปชฺชติ เหตุปจฺจยา.\csromanspacer{}\csromanspacer{}อกุสลํ โนจิตฺตํ ธมฺมํ ปฏิจฺจ น-อพฺยากโต นโนจิตฺโต ธมฺโม อุปฺปชฺชติ เหตุปจฺจยา.\csromanspacer{}\csromanspacer{}อกุสลํ โนจิตฺตํ ธมฺมํ ปฏิจฺจ นกุสโล นโนจิตฺโต จ น-อพฺยากโต นโนจิตฺโต จ ธมฺมา อุปฺปชฺชนฺติ เหตุปจฺจยา.\csromanspacer{} ตีณิ.}

\prose{อพฺยากตานิ ตีณิ.\csromanspacer{}\csromanspacer{}เหตุยา นว.}

\csromanheader{2}{1. กุสลตฺติก 57. เจตสิกทุก}
\proseitem{35}{กุสลํ เจตสิกํ ธมฺมํ ปฏิจฺจ นกุสโล นเจตสิโก ธมฺโม อุปฺปชฺชติ เหตุปจฺจยา.\csromanspacer{}\csromanspacer{}กุสลํ เจตสิกํ ธมฺมํ ปฏิจฺจ น-อกุสโล นเจตสิโก ธมฺโม อุปฺปชฺชติ เหตุปจฺจยา.\csromanspacer{}\csromanspacer{}กุสลํ เจตสิกํ ธมฺมํ ปฏิจฺจ น-อพฺยากโต นเจตสิโก ธมฺโม อุปฺปชฺชติ เหตุปจฺจยา.\csromanspacer{}\csromanspacer{}กุสลํ เจตสิกํ ธมฺมํ ปฏิจฺจ น-อกุสโล นเจตสิโก จ น-อพฺยากโต นเจตสิโก จ ธมฺมา อุปฺปชฺชนฺติ เหตุปจฺจยา.\csromanspacer{}\csromanspacer{}กุสลํ เจตสิกํ ธมฺมํ ปฏิจฺจ นกุสโล นเจตสิโก จ น-อกุสโล นเจตสิโก จ ธมฺมา อุปฺปชฺชนฺติ เหตุปจฺจยา.\csromanspacer{} ปญฺจ.}

\prose{อกุสลํ เจตสิกํ ธมฺมํ ปฏิจฺจ น-อกุสโล นเจตสิโก ธมฺโม อุปฺปชฺชติ เหตุปจฺจยา.\csromanspacer{}\csromanspacer{}อกุสลํ ธมฺมํ ปฏิจฺจ นกุสโล}

\vspace{\tipitakaparskip}
\csromancenter{ธมฺมานุโลมปจฺจนีย ติกทุกปฏฺฐานปาฬิ นเจตสิโก ธมฺโม อุปฺปชฺชติ เหตุปจฺจยา.\csromanspacer{}\csromanspacer{}อกุสลํ เจตสิกํ ธมฺมํ ปฏิจฺจ น-อพฺยากโต นเจตสิโก ธมฺโม อุปฺปชฺชติ เหตุปจฺจยา.\csromanspacer{}\csromanspacer{}อกุสลํ เจตสิกํ ธมฺมํ ปฏิจฺจ นกุสโล นเจตสิโก จ น-อพฺยากโต นเจตสิโก จ ธมฺมา อุปฺปชฺชนฺติ เหตุปจฺจยา.\csromanspacer{}\csromanspacer{}อกุสลํ เจตสิกํ ธมฺมํ ปฏิจฺจ นกุสโล นเจตสิโก จ น-อกุสโล นเจตสิโก จ ธมฺมา อุปฺปชฺชนฺติ เหตุปจฺจยา.\csromanspacer{} ปญฺจ.}
\prose{\csromanfolio{258}อพฺยากตํ เจตสิกํ ธมฺมํ ปฏิจฺจ นกุสโล นเจตสิโก ธมฺโม อุปฺปชฺชติ เหตุปจฺจยา.\csromanspacer{}\csromanspacer{}อพฺยากตํ เจตสิกํ ธมฺมํ ปฏิจฺจ น-อกุสโล นเจตสิโก ธมฺโม อุปฺปชฺชติ เหตุปจฺจยา.\csromanspacer{}\csromanspacer{}อพฺยากตํ เจตสิกํ ธมฺมํ ปฏิจฺจ นกุสโล นเจตสิโก จ น-อกุสโล นเจตสิโก จ ธมฺมา อุปฺปชฺชนฺติ เหตุปจฺจยา.\csromanspacer{} ตีณิ.}

\prose{เหตุยา เตรส.\csromanspacer{}\csromanspacer{}อเจตสิกานิ นว.}

\csromanheader{2}{1. กุสลตฺติก 58. จิตฺตสมฺปยุตฺตทุก}
\proseitem{36}{กุสลํ จิตฺตสมฺปยุตฺตํ ธมฺมํ ปฏิจฺจ นกุสโล นจิตฺตสมฺปยุตฺโต ธมฺโม อุปฺปชฺชติ เหตุปจฺจยา.\csromanspacer{}\csromanspacer{}กุสลํ จิตฺตสมฺปยุตฺตํ ธมฺมํ ปฏิจฺจ น-อกุสโล นจิตฺตสมฺปยุตฺโต ธมฺโม อุปฺปชฺชติ เหตุปจฺจยา.\csromanspacer{}\csromanspacer{}กุสลํ จิตฺตสมฺปยุตฺตํ ธมฺมํ ปฏิจฺจ น-อพฺยากโต นจิตฺตสมฺปยุตฺโต ธมฺโม อุปฺปชฺชติ เหตุปจฺจยา.\csromanspacer{}\csromanspacer{}กุสลํ จิตฺตสมฺปยุตฺตํ ธมฺมํ ปฏิจฺจ น-อกุสโล นจิตฺตสมฺปยุตฺโต จ น-อพฺยากโต นจิตฺตสมฺปยุตฺโต จ ธมฺมา อุปฺปชฺชนฺติ เหตุปจฺจยา.\csromanspacer{}\csromanspacer{}กุสลํ จิตฺตสมฺปยุตฺตํ ธมฺมํ ปฏิจฺจ นกุสโล นจิตฺตสมฺปยุตฺโต จ น-อกุสโล นจิตฺตสมฺปยุตฺโต จ ธมฺมา อุปฺปชฺชนฺติ เหตุปจฺจยา.\csromanspacer{} ปญฺจ.}

\prose{อกุสลํ จิตฺตสมฺปยุตฺตํ ธมฺมํ ปฏิจฺจ น-อกุสโล นจิตฺตสมฺปยุตฺโต ธมฺโม อุปฺปชฺชติ เหตุปจฺจยา.\csromanspacer{} ปญฺจ.}

\prose{อพฺยากตํ จิตฺตสมฺปยุตฺตํ ธมฺมํ ปฏิจฺจ นกุสโล นจิตฺตสมฺปยุตฺโต ธมฺโม อุปฺปชฺชติ เหตุปจฺจยา.\csromanspacer{}\csromanspacer{}อพฺยากตํ จิตฺตสมฺปยุตฺตํ ธมฺมํ ปฏิจฺจ น-อกุสโล นจิตฺตสมฺปยุตฺโต ธมฺโม อุปฺปชฺชติ เหตุปจฺจยา.\csromanspacer{}\csromanspacer{}อพฺยากตํ จิตฺตสมฺปยุตฺตํ ธมฺมํ ปฏิจฺจ นกุสโล นจิตฺตสมฺปยุตฺโต จ น-อกุสโล นจิตฺตสมฺปยุตฺโต จ ธมฺมา อุปฺปชฺชนฺติ เหตุปจฺจยา.\csromanspacer{} ตีณิ. (สํขิตฺตํ.) .\csromanspacer{} เหตุยา เตรส.\csromanspacer{}\csromanspacer{}จิตฺตวิปฺปยุตฺเต ตีณิ.}

\vspace{\tipitakaparskip}
\csromancenter{กุสลตฺติก 60.\csromanspacer{} จิตฺตสมุฏฺฐานทุก \textbf{1.}\csromanspacer{}\textbf{ กุสลตฺติก 59.}\csromanspacer{}\textbf{ จิตฺตสํสฏฺฐทุก}}
\proseitem{37}{\csromanfolio{259}กุสลํ จิตฺตสํสฏฺฐํ ธมฺมํ ปฏิจฺจ นกุสโล นจิตฺตสํสฏฺโฐ ธมฺโม อุปฺปชฺชติ เหตุปจฺจยา.\csromanspacer{}\csromanspacer{}กุสลํ จิตฺตสํสฏฺฐํ ธมฺมํ ปฏิจฺจ น-อกุสโล นจิตฺตสํสฏฺโฐ ธมฺโม อุปฺปชฺชติ เหตุปจฺจยา.\csromanspacer{}\csromanspacer{}กุสลํ จิตฺตสํสฏฺฐํ ธมฺมํ ปฏิจฺจ น-อพฺยากโต นจิตฺตสํสฏฺโฐ ธมฺโม อุปฺปชฺชติ เหตุปจฺจยา.\csromanspacer{}\csromanspacer{}กุสลํ จิตฺตสํสฏฺฐํ ธมฺมํ ปฏิจฺจ น-อกุสโล นจิตฺตสํสฏฺโฐ จ น-อพฺยากโต นจิตฺตสํสฏฺโฐ จ ธมฺมา อุปฺปชฺชนฺติ เหตุปจฺจยา.\csromanspacer{}\csromanspacer{}กุสลํ จิตฺตสํสฏฺฐํ ธมฺมํ ปฏิจฺจ นกุสโล นจิตฺตสํสฏฺโฐ จ น-อกุสโล นจิตฺตสํสฏฺโฐ จ ธมฺมา อุปฺปชฺชนฺติ เหตุปจฺจยา.\csromanspacer{} ปญฺจ.}

\prose{อกุสลํ จิตฺตสํสฏฺฐํ ธมฺมํ ปฏิจฺจ น-อกุสโล นจิตฺตสํสฏฺโฐ ธมฺโม อุปฺปชฺชติ เหตุปจฺจยา.\csromanspacer{}\csromanspacer{}อกุสลํ จิตฺตสํสฏฺฐํ ธมฺมํ ปฏิจฺจ นกุสโล นจิตฺตสํสฏฺโฐ ธมฺโม อุปฺปชฺชติ เหตุปจฺจยา.\csromanspacer{}\csromanspacer{}อกุสลํ จิตฺตสํสฏฺฐํ ธมฺมํ ปฏิจฺจ น-อพฺยากโต นจิตฺตสํสฏฺโฐ ธมฺโม อุปฺปชฺชติ เหตุปจฺจยา.\csromanspacer{}\csromanspacer{}อกุสลํ จิตฺตสํสฏฺฐํ ธมฺมํ ปฏิจฺจ นกุสโล นจิตฺตสํสฏฺโฐ จ น-อพฺยากโต นจิตฺตสํสฏฺโฐ จ ธมฺมา อุปฺปชฺชนฺติ เหตุปจฺจยา.\csromanspacer{}\csromanspacer{}อกุสลํ จิตฺตสํสฏฺฐํ ธมฺมํ ปฏิจฺจ นกุสโล นจิตฺตสํสฏฺโฐ จ น-อกุสโล นจิตฺตสํสฏฺโฐ จ ธมฺมา อุปฺปชฺชนฺติ เหตุปจฺจยา.\csromanspacer{} ปญฺจ.}

\prose{อพฺยากตํ จิตฺตสํสฏฺฐํ ธมฺมํ ปฏิจฺจ นกุสโล นจิตฺตสํสฏฺโฐ ธมฺโม อุปฺปชฺชติ เหตุปจฺจยา.\csromanspacer{}\csromanspacer{}อพฺยากตํ จิตฺตสํสฏฺฐํ ธมฺมํ ปฏิจฺจ น-อกุสโล นจิตฺตสํสฏฺโฐ ธมฺโม อุปฺปชฺชติ เหตุปจฺจยา.\csromanspacer{}\csromanspacer{}อพฺยากตํ จิตฺตสํสฏฺฐํ ธมฺมํ ปฏิจฺจ นกุสโล นจิตฺตสํสฏฺโฐ จ น-อกุสโล นจิตฺตสํสฏฺโฐ จ ธมฺมา อุปฺปชฺชนฺติ เหตุปจฺจยา.\csromanspacer{} ตีณิ. (สํขิตฺตํ.) .\csromanspacer{} เหตุยา เตรส.}

\csromanheader{2}{1. กุสลตฺติก 60. จิตฺตสมุฏฺฐานทุก}
\proseitem{38}{กุสลํ จิตฺตสมุฏฺฐานํ ธมฺมํ ปฏิจฺจ น-อกุสโล นจิตฺตสมุฏฺฐาโน ธมฺโม อุปฺปชฺชติ เหตุปจฺจยา.\csromanspacer{}\csromanspacer{}กุสลํ จิตฺตสมุฏฺฐานํ ธมฺมํ ปฏิจฺจ น-อพฺยากโต นจิตฺตสมุฏฺฐาโน ธมฺโม อุปฺปชฺชติ เหตุปจฺจยา.\csromanspacer{}\csromanspacer{}กุสลํ จิตฺตสมุฏฺฐานํ ธมฺมํ ปฏิจฺจ น-อกุสโล นจิตฺตสมุฏฺฐาโน จ น-อพฺยากโต นจิตฺตสมุฏฺฐาโน จ ธมฺมา อุปฺปชฺชนฺติ เหตุปจฺจยา.\csromanspacer{} ตีณิ.}

\csromanheader{2}{ธมฺมานุโลมปจฺจนีย ติกทุกปฏฺฐานปาฬิ}
\prose{\csromanfolio{260}อกุสลํ จิตฺตสมุฏฺฐานํ ธมฺมํ ปฏิจฺจ นกุสโล นจิตฺตสมุฏฺฐาโน ธมฺโม อุปฺปชฺชติ เหตุปจฺจยา.\csromanspacer{}\csromanspacer{}อกุสลํ จิตฺตสมุฏฺฐานํ ธมฺมํ ปฏิจฺจ น-อพฺยากโต นจิตฺตสมุฏฺฐาโน ธมฺโม อุปฺปชฺชติ เหตุปจฺจยา.\csromanspacer{}\csromanspacer{}อกุสลํ จิตฺตสมุฏฺฐานํ ธมฺมํ ปฏิจฺจ นกุสโล นจิตฺตสมุฏฺฐาโน จ น-อพฺยากโต นจิตฺตสมุฏฺฐาโน จ ธมฺมา อุปฺปชฺชนฺติ เหตุปจฺจยา.\csromanspacer{} ตีณิ.}

\prose{อพฺยากตํ จิตฺตสมุฏฺฐานํ ธมฺมํ ปฏิจฺจ นกุสโล นจิตฺตสมุฏฺฐาโน ธมฺโม อุปฺปชฺชติ เหตุปจฺจยา.\csromanspacer{} ตีณิ.\csromanspacer{}\csromanspacer{}เหตุยา นว ฯเปฯ อวิคเต นว.}

\proseitem{39}{กุสลํ โนจิตฺตสมุฏฺฐานํ ธมฺมํ ปฏิจฺจ นกุสโล นโนจิตฺตสมุฏฺฐาโน ธมฺโม อุปฺปชฺชติ เหตุปจฺจยา.\csromanspacer{} ปญฺจ.}

\prose{อกุสลํ โนจิตฺตสมุฏฺฐานํ ธมฺมํ ปฏิจฺจ น-อกุสโล นโนจิตฺตสมุฏฺฐาโน ธมฺโม อุปฺปชฺชติ เหตุปจฺจยา.\csromanspacer{} ปญฺจ.}

\prose{อพฺยากตํ โนจิตฺตสมุฏฺฐานํ ธมฺมํ ปฏิจฺจ นกุสโล นโนจิตฺตสมุฏฺฐาโน ธมฺโม อุปฺปชฺชติ เหตุปจฺจยา.\csromanspacer{} ตีณิ.\csromanspacer{}\csromanspacer{}เหตุยา เตรส ฯเปฯ อวิคเต เตรส.}

\csromanheader{2}{1. กุสลตฺติก 61. จิตฺตสหภูทุก}
\proseitem{40}{กุสลํ จิตฺตสหภุํ ธมฺมํ ปฏิจฺจ นกุสโล นจิตฺตสหภู ธมฺโม อุปฺปชฺชติ เหตุปจฺจยา.\csromanspacer{} ปญฺจ.}

\prose{อกุสลํ จิตฺตสหภุํ ธมฺมํ ปฏิจฺจ น-อกุสโล นจิตฺตสหภู ธมฺโม อุปฺปชฺชติ เหตุปจฺจยา.\csromanspacer{} ปญฺจ.}

\prose{อพฺยากตํ จิตฺตสหภุํ ธมฺมํ ปฏิจฺจ นกุสโล นจิตฺตสหภู ธมฺโม อุปฺปชฺชติ เหตุปจฺจยา.\csromanspacer{} ตีณิ.\csromanspacer{}\csromanspacer{}เหตุยา เตรส ฯเปฯ อวิคเต เตรส.}

\proseitem{41}{กุสลํ โนจิตฺตสหภุํ ธมฺมํ ปฏิจฺจ นกุสโล นโนจิตฺตสหภู ธมฺโม อุปฺปชฺชติ เหตุปจฺจยา.\csromanspacer{} ปญฺจ.}

\prose{อกุสลํ โนจิตฺตสหภุํ ธมฺมํ ปฏิจฺจ น-อกุสโล นโนจิตฺตสหภู ธมฺโม อุปฺปชฺชติ เหตุปจฺจยา.\csromanspacer{} ปญฺจ.}

\csromanheader{2}{1. กุสลตฺติก 69. อุปาทานทุกาทิ}
\prose{\csromanfolio{261}อพฺยากตํ โนจิตฺตสหภุํ ธมฺมํ ปฏิจฺจ นกุสโล นโนจิตฺตสหภู ธมฺโม อุปฺปชฺชติ เหตุปจฺจยา.\csromanspacer{} ตีณิ. (สํขิตฺตํ.) .\csromanspacer{} เหตุยา เตรส ฯเปฯ อวิคเต เตรส.}

\csromanheader{2}{1. กุสลตฺติก 62. จิตฺตานุปริวตฺติทุก}
\proseitem{42}{กุสลํ จิตฺตานุปริวตฺติํ ธมฺมํ ปฏิจฺจ. (สํขิตฺตํ.) .\csromanspacer{} เหตุยา เตรส.}

\prose{กุสลํ โนจิตฺตานุปริวตฺติํ ธมฺมํ ปฏิจฺจ. (สํขิตฺตํ.) .\csromanspacer{} เหตุยา เตรส. (เอเต สํขิตฺตา.\csromanspacer{} ทุกตฺตยํ จิตฺตทุกสทิสํ.)}

\csromanheader{2}{1. กุสลตฺติก 66. อชฺฌตฺติกทุกาทิ}
\proseitem{43}{กุสลํ อชฺฌตฺติกํ ธมฺมํ ปฏิจฺจ. (สํขิตฺตํ.\csromanspacer{} จิตฺตทุกสทิสํ.)}

\prose{กุสลํ พาหิรํ ธมฺมํ ปฏิจฺจ น-อกุสโล นพาหิโร ธมฺโม อุปฺปชฺชติ เหตุปจฺจยา.\csromanspacer{} ตีณิ.}

\proseitem{44}{อพฺยากตํ อุปาทา ธมฺมํ ปฏิจฺจ นกุสโล น-อุปาทา ธมฺโม อุปฺปชฺชติ เหตุปจฺจยา.\csromanspacer{}\csromanspacer{}เหตุยา ตีณิ.}

\prose{กุสลํ โน-อุปาทา ธมฺมํ ปฏิจฺจ นกุสโล นโน-อุปาทา ธมฺโม อุปฺปชฺชติ เหตุปจฺจยา.\csromanspacer{} ตีณิ.\csromanspacer{}\csromanspacer{}นว.}

\proseitem{45}{อพฺยากตํ อุปาทินฺนํ ธมฺมํ ปฏิจฺจ นกุสโล น-อุปาทินฺโน ธมฺโม อุปฺปชฺชติ เหตุปจฺจยา.\csromanspacer{}\csromanspacer{}อพฺยากตํ อุปาทินฺนํ ธมฺมํ ปฏิจฺจ น-อกุสโล น-อุปาทินฺโน ธมฺโม อุปฺปชฺชติ เหตุปจฺจยา.\csromanspacer{}\csromanspacer{}อพฺยากตํ อุปาทินฺนํ ธมฺมํ ปฏิจฺจ นกุสโล น-อุปาทินฺโน จ น-อกุสโล น-อุปาทินฺโน จ ธมฺมา อุปฺปชฺชนฺติ เหตุปจฺจยา.\csromanspacer{}\csromanspacer{}เหตุยา ตีณิ.}

\csromanheader{2}{1. กุสลตฺติก 69. อุปาทานทุกาทิ}
\proseitem{46}{อกุสลํ อุปาทานํ ธมฺมํ ปฏิจฺจ น-อกุสโล โน-อุปาทาโน ธมฺโม อุปฺปชฺชติ เหตุปจฺจยา.}

\prose{อกุสลํ กิเลสํ ธมฺมํ ปฏิจฺจ น-อกุสโล โนกิเลโส ธมฺโม อุปฺปชฺชติ เหตุปจฺจยา.}

\csromanheader{2}{ธมฺมานุโลมปจฺจนีย ติกทุกปฏฺฐานปาฬิ \textbf{1. กุสลตฺติก 83. ปิฏฺฐิทุก}}
\proseitem{47}{\csromanfolio{262}อกุสลํ ทสฺสเนน ปหาตพฺพํ ธมฺมํ ปฏิจฺจ น-อกุสโล นทสฺสเนน ปหาตพฺโพ ธมฺโม อุปฺปชฺชติ เหตุปจฺจยา.\csromanspacer{}\csromanspacer{}อกุสลํ ทสฺสเนน ปหาตพฺพํ ธมฺมํ ปฏิจฺจ นกุสโล นทสฺสเนน ปหาตพฺโพ ธมฺโม อุปฺปชฺชติ เหตุปจฺจยา.\csromanspacer{}\csromanspacer{}อกุสลํ ทสฺสเนน ปหาตพฺพํ ธมฺมํ ปฏิจฺจ นกุสโล นทสฺสเนน ปหาตพฺโพ จ น-อกุสโล นทสฺสเนน ปหาตพฺโพ จ ธมฺมา อุปฺปชฺชนฺติ เหตุปจฺจยา.\csromanspacer{}\csromanspacer{}เหตุยา ตีณิ.}

\prose{อพฺยากตํ นทสฺสเนน ปหาตพฺพํ ธมฺมํ ปจฺจยา น-อพฺยากโต นนทสฺสเนน ปหาตพฺโพ ธมฺโม อุปฺปชฺชติ เหตุปจฺจยา.\csromanspacer{}\csromanspacer{}อพฺยากตํ นทสฺสเนน ปหาตพฺพํ ธมฺมํ ปจฺจยา นกุสโล นนทสฺสเนน ปหาตพฺโพ ธมฺโม อุปฺปชฺชติ เหตุปจฺจยา.\csromanspacer{}\csromanspacer{}อพฺยากตํ นทสฺสเนน ปหาตพฺพํ ธมฺมํ ปจฺจยา นกุสโล นนทสฺสเนน ปหาตพฺโพ จ น-อพฺยากโต นนทสฺสเนน ปหาตพฺโพ จ ธมฺมา อุปฺปชฺชนฺติ เหตุปจฺจยา.\csromanspacer{}\csromanspacer{}เหตุยา ตีณิ.}

\proseitem{48}{อกุสลํ ภาวนาย ปหาตพฺพํ ธมฺมํ ปฏิจฺจ น-อกุสโล นภาวนาย ปหาตพฺโพ ธมฺโม อุปฺปชฺชติ เหตุปจฺจยา.\csromanspacer{}\csromanspacer{}อกุสลํ ภาวนาย ปหาตพฺพํ ธมฺมํ ปฏิจฺจ นกุสโล นภาวนาย ปหาตพฺโพ ธมฺโม อุปฺปชฺชติ เหตุปจฺจยา.\csromanspacer{}\csromanspacer{}อกุสลํ ภาวนาย ปหาตพฺพํ ธมฺมํ ปฏิจฺจ นกุสโล นภาวนาย ปหาตพฺโพ จ น-อกุสโล นภาวนาย ปหาตพฺโพ จ ธมฺมา อุปฺปชฺชนฺติ เหตุปจฺจยา.\csromanspacer{}\csromanspacer{}เหตุยา ตีณิ.}

\prose{อพฺยากตํ นภาวนาย ปหาตพฺพํ ธมฺมํ ปจฺจยา น-อพฺยากโต นนภาวนาย ปหาตพฺโพ ธมฺโม อุปฺปชฺชติ เหตุปจฺจยา.\csromanspacer{}\csromanspacer{}เหตุยา ตีณิ.}

\proseitem{49}{อกุสลํ ทสฺสเนน ปหาตพฺพเหตุกํ ธมฺมํ ปฏิจฺจ น-อกุสโล นทสฺสเนน ปหาตพฺพเหตุโก ธมฺโม อุปฺปชฺชติ เหตุปจฺจยา.\csromanspacer{}\csromanspacer{}เหตุยา ตีณิ.}

\prose{อกุสลํ นทสฺสเนน ปหาตพฺพเหตุกํ ธมฺมํ ปฏิจฺจ นกุสโล นนทสฺสเนน ปหาตพฺพเหตุโก ธมฺโม อุปฺปชฺชติ เหตุปจฺจยา.\csromanspacer{}\csromanspacer{}เหตุยา ตีณิ.}

\csromanheader{2}{1. กุสลตฺติก 83. ปิฏฺฐิทุก}
\prose{\csromanfolio{263}อพฺยากตํ นทสฺสเนน ปหาตพฺพเหตุกํ ธมฺมํ ปจฺจยา น-อพฺยากโต นนทสฺสเนน ปหาตพฺพเหตุโก ธมฺโม อุปฺปชฺชติ เหตุปจฺจยา.\csromanspacer{}\csromanspacer{}เหตุยา ตีณิ.}

\proseitem{50}{อกุสลํ ภาวนาย ปหาตพฺพเหตุกํ ธมฺมํ ปฏิจฺจ น-อกุสโล นภาวนาย ปหาตพฺพเหตุโก ธมฺโม อุปฺปชฺชติ เหตุปจฺจยา.\csromanspacer{}\csromanspacer{}เหตุยา ตีณิ.}

\prose{อกุสลํ นภาวนาย ปหาตพฺพเหตุกํ ธมฺมํ ปฏิจฺจ นกุสโล นนภาวนาย ปหาตพฺพเหตุโก ธมฺโม อุปฺปชฺชติ เหตุปจฺจยา.\csromanspacer{}\csromanspacer{}เหตุยา ตีณิ.}

\proseitem{51}{กุสลํ สวิตกฺกํ ธมฺมํ ปฏิจฺจ นกุสโล นสวิตกฺโก ธมฺโม อุปฺปชฺชติ เหตุปจฺจยา.\csromanspacer{}\csromanspacer{}เหตุยา เตรส.}

\prose{กุสลํ อวิตกฺกํ ธมฺมํ ปฏิจฺจ น-อกุสโล น-อวิตกฺโก ธมฺโม อุปฺปชฺชติ เหตุปจฺจยา.\csromanspacer{}\csromanspacer{}เหตุยา นว.}

\proseitem{52}{กุสลํ สวิจารํ ธมฺมํ ปฏิจฺจ นกุสโล นสวิจาโร ธมฺโม อุปฺปชฺชติ เหตุปจฺจยา.\csromanspacer{}\csromanspacer{}เหตุยา เตรส.}

\prose{กุสลํ อวิจารํ ธมฺมํ ปฏิจฺจ น-อกุสโล น-อวิจาโร ธมฺโม อุปฺปชฺชติ เหตุปจฺจยา.\csromanspacer{}\csromanspacer{}เหตุยา นว.}

\proseitem{53}{กุสลํ สปฺปีติกํ ธมฺมํ ปฏิจฺจ นกุสโล นสปฺปีติโก ธมฺโม อุปฺปชฺชติ เหตุปจฺจยา.\csromanspacer{}\csromanspacer{}เหตุยา เตรส.}

\prose{กุสลํ อปฺปีติกํ ธมฺมํ ปฏิจฺจ น-อกุสโล น-อปฺปีติโก ธมฺโม อุปฺปชฺชติ เหตุปจฺจยา.\csromanspacer{}\csromanspacer{}เหตุยา นว.}

\proseitem{54}{กุสลํ ปีติสหคตํ ธมฺมํ ปฏิจฺจ นกุสโล นปีติสหคโต ธมฺโม อุปฺปชฺชติ เหตุปจฺจยา.\csromanspacer{}\csromanspacer{}เหตุยา เตรส.}

\prose{กุสลํ นปีติสหคตํ ธมฺมํ ปฏิจฺจ น-อกุสโล นนปีติสหคโต ธมฺโม อุปฺปชฺชติ เหตุปจฺจยา.\csromanspacer{}\csromanspacer{}เหตุยา นว.}

\proseitem{55}{กุสลํ สุขสหคตํ ธมฺมํ ปฏิจฺจ นกุสโล นสุขสหคโต ธมฺโม อุปฺปชฺชติ เหตุปจฺจยา.\csromanspacer{}\csromanspacer{}เหตุยา เตรส.}

\csromanheader{2}{ธมฺมานุโลมปจฺจนีย ติกทุกปฏฺฐานปาฬิ}
\prose{\csromanfolio{264}กุสลํ นสุขสหคตํ ธมฺมํ ปฏิจฺจ น-อกุสโล นนสุขสหคโต ธมฺโม อุปฺปชฺชติ เหตุปจฺจยา.\csromanspacer{}\csromanspacer{}เหตุยา นว.}

\proseitem{56}{กุสลํ อุเปกฺขาสหคตํ ธมฺมํ ปฏิจฺจ นกุสโล น-อุเปกฺขาสหคโต ธมฺโม อุปฺปชฺชติ เหตุปจฺจยา.\csromanspacer{}\csromanspacer{}เหตุยา เตรส.}

\prose{กุสลํ น-อุเปกฺขาสหคตํ ธมฺมํ ปฏิจฺจ นกุสโล นน-อุเปกฺขาสหคโต ธมฺโม อุปฺปชฺชติ เหตุปจฺจยา.\csromanspacer{}\csromanspacer{}เหตุยา นว.}

\proseitem{57}{อพฺยากตํ กามาวจรํ ธมฺมํ ปฏิจฺจ นกุสโล นกามาวจโร ธมฺโม อุปฺปชฺชติ เหตุปจฺจยา.\csromanspacer{}\csromanspacer{}เหตุยา ตีณิ.}

\prose{กุสลํ นกามาวจรํ ธมฺมํ ปฏิจฺจ นกุสโล นนกามาวจโร ธมฺโม อุปฺปชฺชติ เหตุปจฺจยา.\csromanspacer{}\csromanspacer{}เหตุยา ฉ.}

\proseitem{58}{กุสลํ รูปาวจรํ ธมฺมํ ปฏิจฺจ นกุสโล นรูปาวจโร ธมฺโม อุปฺปชฺชติ เหตุปจฺจยา.\csromanspacer{}\csromanspacer{}เหตุยา ฉ.}

\prose{อพฺยากตํ นรูปาวจรํ ธมฺมํ ปฏิจฺจ นกุสโล นนรูปาวจโร ธมฺโม อุปฺปชฺชติ เหตุปจฺจยา.\csromanspacer{}\csromanspacer{}เหตุยา ตีณิ.}

\proseitem{59}{กุสลํ อรูปาวจรํ ธมฺมํ ปฏิจฺจ นกุสโล น-อรูปาวจโร ธมฺโม อุปฺปชฺชติ เหตุปจฺจยา.\csromanspacer{}\csromanspacer{}เหตุยา ฉ.}

\prose{อพฺยากตํ น-อรูปาวจรํ ธมฺมํ ปจฺจยา น-อพฺยากโต นน-อรูปาวจโร ธมฺโม อุปฺปชฺชติ เหตุปจฺจยา.\csromanspacer{}\csromanspacer{}เหตุยา ปญฺจ.}

\proseitem{60}{อพฺยากตํ ปริยาปนฺนํ ธมฺมํ ปจฺจยา น-อพฺยากโต นปริยาปนฺโน ธมฺโม อุปฺปชฺชติ เหตุปจฺจยา.\csromanspacer{}\csromanspacer{}เหตุยา ปญฺจ.}

\prose{กุสลํ อปริยาปนฺนํ ธมฺมํ ปฏิจฺจ นกุสโล น-อปริยาปนฺโน ธมฺโม อุปฺปชฺชติ เหตุปจฺจยา.\csromanspacer{}\csromanspacer{}เหตุยา ฉ.}

\proseitem{61}{กุสลํ นิยฺยานิกํ ธมฺมํ ปฏิจฺจ นกุสโล นนิยฺยานิโก ธมฺโม อุปฺปชฺชติ เหตุปจฺจยา.\csromanspacer{}\csromanspacer{}เหตุยา ตีณิ.}

\prose{อพฺยากตํ อนิยฺยานิกํ ธมฺมํ ปจฺจยา น-อพฺยากโต น-อนิยฺยานิโก ธมฺโม อุปฺปชฺชติ เหตุปจฺจยา.\csromanspacer{}\csromanspacer{}เหตุยา ตีณิ.}

\csromanheader{2}{2. เวทนาตฺติก 1. เหตุทุก}
\proseitem{62}{\csromanfolio{265}กุสลํ นิยตํ ธมฺมํ ปฏิจฺจ นกุสโล นนิยโต ธมฺโม อุปฺปชฺชติ เหตุปจฺจยา.\csromanspacer{}\csromanspacer{}เหตุยา ฉ.}

\prose{อพฺยากตํ อนิยตํ ธมฺมํ ปจฺจยา น-อพฺยากโต น-อนิยโต ธมฺโม อุปฺปชฺชติ เหตุปจฺจยา.\csromanspacer{}\csromanspacer{}เหตุยา ปญฺจ.}

\proseitem{63}{อพฺยากตํ ส-อุตฺตรํ ธมฺมํ ปจฺจยา น-อพฺยากโต นส-อุตฺตโร ธมฺโม อุปฺปชฺชติ เหตุปจฺจยา.\csromanspacer{}\csromanspacer{}เหตุยา ปญฺจ.}

\prose{กุสลํ อนุตฺตรํ ธมฺมํ ปฏิจฺจ นกุสโล น-อนุตฺตโร ธมฺโม อุปฺปชฺชติ เหตุปจฺจยา.\csromanspacer{}\csromanspacer{}เหตุยา ฉ.}

\proseitem{64}{อกุสลํ สรณํ ธมฺมํ ปฏิจฺจ น-อกุสโล นสรโณ ธมฺโม อุปฺปชฺชติ เหตุปจฺจยา.\csromanspacer{}\csromanspacer{}อกุสลํ สรณํ ธมฺมํ ปฏิจฺจ นกุสโล นสรโณ ธมฺโม อุปฺปชฺชติ เหตุปจฺจยา.\csromanspacer{}\csromanspacer{}อกุสลํ สรณํ ธมฺมํ ปฏิจฺจ นกุสโล นสรโณ จ น-อกุสโล นสรโณ จ ธมฺมา อุปฺปชฺชนฺติ เหตุปจฺจยา.\csromanspacer{}\csromanspacer{}เหตุยา ตีณิ.}

\prose{อพฺยากตํ อรณํ ธมฺมํ ปจฺจยา น-อพฺยากโต น-อรโณ ธมฺโม อุปฺปชฺชติ เหตุปจฺจยา.\csromanspacer{}\csromanspacer{}อพฺยากตํ อรณํ ธมฺมํ ปจฺจยา นกุสโล น-อรโณ ธมฺโม อุปฺปชฺชติ เหตุปจฺจยา.\csromanspacer{}\csromanspacer{}อพฺยากตํ อรณํ ธมฺมํ ปจฺจยา นกุสโล น-อรโณ จ น-อพฺยากโต น-อรโณ จ ธมฺมา อุปฺปชฺชนฺติ เหตุปจฺจยา.\csromanspacer{}\csromanspacer{}เหตุยา ตีณิ.}

\csromanheader{2}{2. เวทนาตฺติก 1. เหตุทุก}
\proseitem{65}{สุขาย เวทนาย สมฺปยุตฺตํ เหตุํ ธมฺมํ ปฏิจฺจ นสุขาย เวทนาย สมฺปยุตฺโต นเหตุ ธมฺโม อุปฺปชฺชติ เหตุปจฺจยา.\csromanspacer{}\csromanspacer{}สุขาย เวทนาย สมฺปยุตฺตํ เหตุํ ธมฺมํ ปฏิจฺจ นทุกฺขาย เวทนาย สมฺปยุตฺโต นเหตุ ธมฺโม อุปฺปชฺชติ เหตุปจฺจยา.\csromanspacer{}\csromanspacer{}สุขาย เวทนาย สมฺปยุตฺตํ เหตุํ ธมฺมํ ปฏิจฺจ น-อทุกฺขมสุขาย เวทนาย สมฺปยุตฺโต นเหตุ ธมฺโม อุปฺปชฺชติ เหตุปจฺจยา ฯเปฯ.\csromanspacer{} สตฺต ปญฺหา.}

\prose{ทุกฺขาย เวทนาย สมฺปยุตฺตํ เหตุํ ธมฺมํ ปฏิจฺจ นทุกฺขาย เวทนาย สมฺปยุตฺโต นเหตุ ธมฺโม อุปฺปชฺชติ เหตุปจฺจยา.\csromanspacer{} สตฺต ปญฺหา.}

\csromanheader{2}{ธมฺมานุโลมปจฺจนีย ติกทุกปฏฺฐานปาฬิ}
\prose{\csromanfolio{266}อทุกฺขมสุขาย เวทนาย สมฺปยุตฺตํ เหตุํ ธมฺมํ ปฏิจฺจ น-อทุกฺขมสุขาย เวทนาย สมฺปยุตฺโต นเหตุ ธมฺโม อุปฺปชฺชติ เหตุปจฺจยา.\csromanspacer{} สตฺต ปญฺหา.\csromanspacer{}\csromanspacer{}เหตุยา เอกวีส ฯเปฯ อวิคเต เอกวีส. (สพฺพตฺถ เอกวีส.)}

\proseitem{66}{สุขาย เวทนาย สมฺปยุตฺตํ นเหตุํ ธมฺมํ ปฏิจฺจ นทุกฺขาย เวทนาย สมฺปยุตฺโต นนเหตุ ธมฺโม อุปฺปชฺชติ เหตุปจฺจยา.\csromanspacer{}\csromanspacer{}สุขาย เวทนาย สมฺปยุตฺตํ นเหตุํ ธมฺมํ ปฏิจฺจ น-อทุกฺขมสุขาย เวทนาย สมฺปยุตฺโต นนเหตุ ธมฺโม อุปฺปชฺชติ เหตุปจฺจยา.\csromanspacer{}\csromanspacer{}สุขาย เวทนาย สมฺปยุตฺตํ นเหตุํ ธมฺมํ ปฏิจฺจ นทุกฺขาย เวทนาย สมฺปยุตฺโต นนเหตุ จ น-อทุกฺขมสุขาย เวทนาย สมฺปยุตฺโต นนเหตุ จ ธมฺมา อุปฺปชฺชนฺติ เหตุปจฺจยา.\csromanspacer{} ตีณิ.}

\prose{ทุกฺขาย เวทนาย สมฺปยุตฺตํ นเหตุํ ธมฺมํ ปฏิจฺจ นสุขาย เวทนาย สมฺปยุตฺโต นนเหตุ ธมฺโม อุปฺปชฺชติ เหตุปจฺจยา.\csromanspacer{} ตีณิ.}

\prose{อทุกฺขมสุขาย เวทนาย สมฺปยุตฺตํ นเหตุํ ธมฺมํ ปฏิจฺจ นสุขาย เวทนาย สมฺปยุตฺโต นนเหตุ ธมฺโม อุปฺปชฺชติ เหตุปจฺจยา.\csromanspacer{} ตีณิ.\csromanspacer{}\csromanspacer{}เหตุยา นว.}

\csromanheader{2}{3. วิปากตฺติก 1. เหตุทุก}
\proseitem{67}{วิปากํ เหตุํ ธมฺมํ ปฏิจฺจ นวิปาโก นเหตุ ธมฺโม อุปฺปชฺชติ เหตุปจฺจยา.\csromanspacer{}\csromanspacer{}วิปากํ เหตุํ ธมฺมํ ปฏิจฺจ นวิปากธมฺมธมฺโม นเหตุ ธมฺโม อุปฺปชฺชติ เหตุปจฺจยา.\csromanspacer{}\csromanspacer{}วิปากํ เหตุํ ธมฺมํ ปฏิจฺจ นเนววิปากนวิปากธมฺมธมฺโม นเหตุ ธมฺโม อุปฺปชฺชติ เหตุปจฺจยา.\csromanspacer{}\csromanspacer{}วิปากํ เหตุํ ธมฺมํ ปฏิจฺจ นวิปากธมฺมธมฺโม นเหตุ จ นเนววิปากนวิปากธมฺมธมฺโม นเหตุ จ ธมฺมา อุปฺปชฺชนฺติ เหตุปจฺจยา.\csromanspacer{}\csromanspacer{}วิปากํ เหตุํ ธมฺมํ ปฏิจฺจ นวิปาโก นเหตุ จ นวิปากธมฺมธมฺโม นเหตุ จ ธมฺมา อุปฺปชฺชนฺติ เหตุปจฺจยา.\csromanspacer{} ปญฺจ.}

\prose{วิปากธมฺมธมฺมํ เหตุํ ธมฺมํ ปฏิจฺจ นวิปากธมฺมธมฺโม นเหตุ ธมฺโม อุปฺปชฺชติ เหตุปจฺจยา.\csromanspacer{}\csromanspacer{}วิปากธมฺมธมฺมํ เหตุํ ธมฺมํ ปฏิจฺจ นวิปาโก นเหตุ ธมฺโม อุปฺปชฺชติ เหตุปจฺจยา.\csromanspacer{}\csromanspacer{}วิปากธมฺมธมฺมํ เหตุํ ธมฺมํ ปฏิจฺจ นเนววิปากนวิปากธมฺมธมฺโม นเหตุ ธมฺโม อุปฺปชฺชติ เหตุปจฺจยา.\csromanspacer{}\csromanspacer{}วิปากธมฺมธมฺมํ เหตุํ ธมฺมํ ปฏิจฺจ นวิปาโก นเหตุ จ นเนววิปากนวิปากธมฺมธมฺโม นเหตุ จ ธมฺมา อุปฺปชฺชนฺติ เหตุปจฺจยา.\csromanspacer{}\csromanspacer{}วิปากธมฺมธมฺมํ เหตุํ ธมฺมํ ปฏิจฺจ}

\vspace{\tipitakaparskip}
\csromancenter{อุปาทินฺนตฺติก 1.\csromanspacer{} เหตุทุก นวิปาโก นเหตุ จ นวิปากธมฺมธมฺโม นเหตุ จ ธมฺมา อุปฺปชฺชนฺติ เหตุปจฺจยา.\csromanspacer{} ปญฺจ.}
\prose{\csromanfolio{267}เนววิปากนวิปากธมฺมธมฺมํ เหตุํ ธมฺมํ ปฏิจฺจ นวิปาโก นเหตุ ธมฺโม อุปฺปชฺชติ เหตุปจฺจยา.\csromanspacer{}\csromanspacer{}เนววิปากนวิปากธมฺมธมฺมํ เหตุํ ธมฺมํ ปฏิจฺจ นวิปากธมฺมธมฺโม นเหตุ ธมฺโม อุปฺปชฺชติ เหตุปจฺจยา.\csromanspacer{}\csromanspacer{}เนววิปากนวิปากธมฺมธมฺมํ เหตุํ ธมฺมํ ปฏิจฺจ นวิปาโก นเหตุ จ นวิปากธมฺมธมฺโม นเหตุ จ ธมฺมา อุปฺปชฺชนฺติ เหตุปจฺจยา.\csromanspacer{} ตีณิ.\csromanspacer{}\csromanspacer{}เหตุยา เตรส.}

\proseitem{68}{วิปากํ นเหตุํ ธมฺมํ ปฏิจฺจ นวิปากธมฺมธมฺโม นนเหตุ ธมฺโม อุปฺปชฺชติ เหตุปจฺจยา.\csromanspacer{}\csromanspacer{}วิปากํ นเหตุํ ธมฺมํ ปฏิจฺจ นเนววิปากนวิปากธมฺมธมฺโม นนเหตุ ธมฺโม อุปฺปชฺชติ เหตุปจฺจยา.\csromanspacer{}\csromanspacer{}วิปากํ นเหตุํ ธมฺมํ ปฏิจฺจ นวิปากธมฺมธมฺโม นนเหตุ จ นเนววิปากนวิปากธมฺมธมฺโม นนเหตุ จ ธมฺมา อุปฺปชฺชนฺติ เหตุปจฺจยา.\csromanspacer{} ตีณิ.}

\prose{วิปากธมฺมธมฺมํ นเหตุํ ธมฺมํ ปฏิจฺจ นวิปาโก นนเหตุ ธมฺโม อุปฺปชฺชติ เหตุปจฺจยา.\csromanspacer{}\csromanspacer{}วิปากธมฺมธมฺมํ นเหตุํ ธมฺมํ ปฏิจฺจ นเนววิปากนวิปากธมฺมธมฺโม นนเหตุ ธมฺโม อุปฺปชฺชติ เหตุปจฺจยา.\csromanspacer{}\csromanspacer{}วิปากธมฺมธมฺมํ นเหตุํ ธมฺมํ ปฏิจฺจ นวิปาโก นนเหตุ จ นเนววิปากนวิปากธมฺมธมฺโม นนเหตุ จ ธมฺมา อุปฺปชฺชนฺติ เหตุปจฺจยา.\csromanspacer{} ตีณิ.}

\prose{เนววิปากนวิปากธมฺมธมฺมํ นเหตุํ ธมฺมํ ปฏิจฺจ นเนววิปากนวิปากธมฺมธมฺโม นนเหตุ ธมฺโม อุปฺปชฺชติ เหตุปจฺจยา.\csromanspacer{}\csromanspacer{}เนววิปากนวิปากธมฺมธมฺมํ นเหตุํ ธมฺมํ ปฏิจฺจ นวิปาโก นนเหตุ ธมฺโม อุปฺปชฺชติ เหตุปจฺจยา ฯเปฯ.\csromanspacer{} เนววิปากนวิปากธมฺมธมฺมํ นเหตุํ ธมฺมํ ปฏิจฺจ นวิปาโก นนเหตุ จ นวิปากธมฺมธมฺโม นนเหตุ จ ธมฺมา อุปฺปชฺชนฺติ เหตุปจฺจยา.\csromanspacer{} ปญฺจ.}

\prose{วิปากํ นเหตุญฺจ เนววิปากนวิปากธมฺมธมฺมํ นเหตุญฺจ ธมฺมํ ปฏิจฺจ นวิปากธมฺมธมฺโม นนเหตุ ธมฺโม อุปฺปชฺชติ เหตุปจฺจยา.\csromanspacer{} ตีณิ.\csromanspacer{}\csromanspacer{}เหตุยา จุทฺทส.}

\csromanheader{2}{4. อุปาทินฺนตฺติก 1. เหตุทุก}
\proseitem{69}{อุปาทินฺนุปาทานิยํ เหตุํ ธมฺมํ ปฏิจฺจ น-อุปาทินฺนุปาทานิโย นเหตุ ธมฺโม อุปฺปชฺชติ เหตุปจฺจยา.\csromanspacer{}\csromanspacer{}อุปาทินฺนุปาทานิยํ เหตุํ ธมฺมํ ปฏิจฺจ น-อนุปาทินฺนุปาทานิโย นเหตุ ธมฺโม อุปฺปชฺชติ เหตุปจฺจยา.\csromanspacer{}\csromanspacer{}}

\vspace{\tipitakaparskip}
\csromancenter{ธมฺมานุโลมปจฺจนีย ติกทุกปฏฺฐานปาฬิ อุปาทินฺนุปาทานิยํ เหตุํ ธมฺมํ ปฏิจฺจ น-อนุปาทินฺน-อนุปาทานิโย นเหตุ ธมฺโม อุปฺปชฺชติ เหตุปจฺจยา.\csromanspacer{}\csromanspacer{}อุปาทินฺนุปาทานิยํ เหตุํ ธมฺมํ ปฏิจฺจ น-อุปาทินฺนุปาทานิโย นเหตุ จ น-อนุปาทินฺน-อนุปาทานิโย นเหตุ จ ธมฺมา อุปฺปชฺชนฺติ เหตุปจฺจยา.\csromanspacer{}\csromanspacer{}อุปาทินฺนุปาทานิยํ เหตุํ ธมฺมํ ปฏิจฺจ น-อนุปาทินฺนุปาทานิโย นเหตุ จ น-อนุปาทินฺน-อนุปาทานิโย นเหตุ จ ธมฺมา อุปฺปชฺชนฺติ เหตุปจฺจยา.\csromanspacer{} ปญฺจ.}
\prose{\csromanfolio{268}อนุปาทินฺนุปาทานิยํ เหตุํ ธมฺมํ ปฏิจฺจ น-อุปาทินฺนุปาทานิโย นเหตุ ธมฺโม อุปฺปชฺชติ เหตุปจฺจยา.\csromanspacer{} ตีณิ.}

\prose{อนุปาทินฺน-อนุปาทานิยํ เหตุํ ธมฺมํ ปฏิจฺจ น-อนุปาทินฺน-อนุปาทานิโย นเหตุ ธมฺโม อุปฺปชฺชติ เหตุปจฺจยา.\csromanspacer{} ปญฺจ ปญฺหา.\csromanspacer{}\csromanspacer{}เหตุยา เตรส.}

\proseitem{70}{อุปาทินฺนุปาทานิยํ นเหตุํ ธมฺมํ ปฏิจฺจ น-อนุปาทินฺนุปาทานิโย นนเหตุ ธมฺโม อุปฺปชฺชติ เหตุปจฺจยา.\csromanspacer{} ตีณิ.\csromanspacer{}\csromanspacer{}เหตุยา นว.}

\csromanheader{2}{5. สํกิลิฏฺฐตฺติก 1. เหตุทุก}
\proseitem{71}{สํกิลิฏฺฐสํกิเลสิกํ เหตุํ ธมฺมํ ปฏิจฺจ นสํกิลิฏฺฐสํกิเลสิโก นเหตุ ธมฺโม อุปฺปชฺชติ เหตุปจฺจยา.\csromanspacer{}\csromanspacer{}เหตุยา เตรส.}

\prose{สํกิลิฏฺฐสํกิเลสิกํ นเหตุํ ธมฺมํ ปฏิจฺจ น-อสํกิลิฏฺฐสํกิเลสิโก นนเหตุ ธมฺโม อุปฺปชฺชติ เหตุปจฺจยา.\csromanspacer{} ตีณิ.\csromanspacer{}\csromanspacer{}เหตุยา นว.}

\csromanheader{2}{6. วิตกฺกตฺติก 1. เหตุทุก}
\proseitem{72}{สวิตกฺกสวิจารํ เหตุํ ธมฺมํ ปฏิจฺจ นสวิตกฺกสวิจาโร นเหตุ ธมฺโม อุปฺปชฺชติ เหตุปจฺจยา.\csromanspacer{}\csromanspacer{}เหตุยา ปนฺนรส.}

\prose{สวิตกฺกสวิจารํ นเหตุํ ธมฺมํ ปฏิจฺจ น-อวิตกฺกวิจารมตฺโต นนเหตุ ธมฺโม อุปฺปชฺชติ เหตุปจฺจยา.\csromanspacer{}\csromanspacer{}เหตุยา อฏฺฐวีส.}

\csromanheader{2}{7. ปีติตฺติก 1. เหตุทุก}
\proseitem{73}{ปีติสหคตํ เหตุํ ธมฺมํ ปฏิจฺจ นปีติสหคโต นเหตุ ธมฺโม อุปฺปชฺชติ เหตุปจฺจยา.\csromanspacer{}\csromanspacer{}เหตุยา อฏฺฐวีส.}

\csromanheader{2}{11. เสกฺขตฺติก 1. เหตุทุก}
\prose{\csromanfolio{269}ปีติสหคตํ นเหตุํ ธมฺมํ ปฏิจฺจ น-อุเปกฺขาสหคโต นนเหตุ ธมฺโม อุปฺปชฺชติ เหตุปจฺจยา.\csromanspacer{}\csromanspacer{}เหตุยา อฏฺฐารส\footnote{เหตุยา อฏฺฐ?}.}

\csromanheader{2}{8. ทสฺสนตฺติก 1. เหตุทุก}
\proseitem{74}{ทสฺสเนน ปหาตพฺพํ เหตุํ ธมฺมํ ปฏิจฺจ นทสฺสเนน ปหาตพฺโพ นเหตุ ธมฺโม อุปฺปชฺชติ เหตุปจฺจยา.\csromanspacer{}\csromanspacer{}เหตุยา เตรส.}

\prose{ทสฺสเนน ปหาตพฺพํ นเหตุํ ธมฺมํ ปฏิจฺจ นภาวนาย ปหาตพฺโพ นนเหตุ ธมฺโม อุปฺปชฺชติ เหตุปจฺจยา.\csromanspacer{}\csromanspacer{}เหตุยา นว.}

\csromanheader{2}{9. ทสฺสนเหตุตฺติก 1. เหตุทุก}
\proseitem{75}{ทสฺสเนน ปหาตพฺพเหตุกํ เหตุํ ธมฺมํ ปฏิจฺจ นทสฺสเนน ปหาตพฺพเหตุโก นเหตุ ธมฺโม อุปฺปชฺชติ เหตุปจฺจยา.\csromanspacer{}\csromanspacer{}เหตุยา โสฬส.}

\prose{ทสฺสเนน ปหาตพฺพเหตุกํ นเหตุํ ธมฺมํ ปฏิจฺจ นภาวนาย ปหาตพฺพเหตุโก นนเหตุ ธมฺโม อุปฺปชฺชติ เหตุปจฺจยา.\csromanspacer{}\csromanspacer{}เหตุยา นว.}

\csromanheader{2}{10. อาจยคามิตฺติก 1. เหตุทุก}
\proseitem{76}{อาจยคามิํ เหตุํ ธมฺมํ ปฏิจฺจ น-อาจยคามี นเหตุ ธมฺโม อุปฺปชฺชติ เหตุปจฺจยา.\csromanspacer{}\csromanspacer{}เหตุยา เตรส.}

\prose{อาจยคามิํ นเหตุํ ธมฺมํ ปฏิจฺจ น-อปจยคามี นนเหตุ ธมฺโม อุปฺปชฺชติ เหตุปจฺจยา.\csromanspacer{}\csromanspacer{}เหตุยา นว.}

\csromanheader{2}{11. เสกฺขตฺติก 1. เหตุทุก}
\proseitem{77}{เสกฺขํ เหตุํ ธมฺมํ ปฏิจฺจ นเสกฺโข นเหตุ ธมฺโม อุปฺปชฺชติ เหตุปจฺจยา.\csromanspacer{}\csromanspacer{}เหตุยา เตรส.}

\prose{เสกฺขํ นเหตุํ ธมฺมํ ปฏิจฺจ น-อเสกฺโข นนเหตุ ธมฺโม อุปฺปชฺชติ เหตุปจฺจยา.\csromanspacer{}\csromanspacer{}เหตุยา นว.}

\csromanheader{2}{ธมฺมานุโลมปจฺจนีย ติกทุกปฏฺฐานปาฬิ \textbf{12. ปริตฺตตฺติก 1. เหตุทุก}}
\proseitem{78}{\csromanfolio{270}ปริตฺตํ เหตุํ ธมฺมํ ปฏิจฺจ นมหคฺคโต นเหตุ ธมฺโม อุปฺปชฺชติ เหตุปจฺจยา.\csromanspacer{}\csromanspacer{}เหตุยา เตรส.}

\prose{ปริตฺตํ นเหตุํ ธมฺมํ ปฏิจฺจ นปริตฺโต นนเหตุ ธมฺโม อุปฺปชฺชติ เหตุปจฺจยา.\csromanspacer{}\csromanspacer{}เหตุยา จุทฺทส.}

\csromanheader{2}{13. ปริตฺตารมฺมณตฺติก 1. เหตุทุก}
\proseitem{79}{ปริตฺตารมฺมณํ เหตุํ ธมฺมํ ปฏิจฺจ นปริตฺตารมฺมโณ นเหตุ ธมฺโม อุปฺปชฺชติ เหตุปจฺจยา.\csromanspacer{}\csromanspacer{}เหตุยา เอกวีส.}

\prose{ปริตฺตารมฺมณํ นเหตุํ ธมฺมํ ปฏิจฺจ นมหคฺคตารมฺมโณ นนเหตุ ธมฺโม อุปฺปชฺชติ เหตุปจฺจยา.\csromanspacer{}\csromanspacer{}เหตุยา นว.}

\csromanheader{2}{14. หีนตฺติก 1. เหตุทุก}
\proseitem{80}{หีนํ เหตุํ ธมฺมํ ปฏิจฺจ นหีโน นเหตุ ธมฺโม อุปฺปชฺชติ เหตุปจฺจยา.\csromanspacer{}\csromanspacer{}เหตุยา เตรส.}

\prose{หีนํ นเหตุํ ธมฺมํ ปฏิจฺจ นมชฺฌิโม นนเหตุ ธมฺโม อุปฺปชฺชติ เหตุปจฺจยา.\csromanspacer{}\csromanspacer{}เหตุยา นว.}

\csromanheader{2}{15. มิจฺฉตฺตนิยตตฺติก 1. เหตุทุก}
\proseitem{81}{มิจฺฉตฺตนิยตํ เหตุํ ธมฺมํ ปฏิจฺจ นมิจฺฉตฺตนิยโต นเหตุ ธมฺโม อุปฺปชฺชติ เหตุปจฺจยา.\csromanspacer{}\csromanspacer{}เหตุยา เตรส.}

\prose{มิจฺฉตฺตนิยตํ นเหตุํ ธมฺมํ ปฏิจฺจ นสมฺมตฺตนิยโต นนเหตุ ธมฺโม อุปฺปชฺชติ เหตุปจฺจยา.\csromanspacer{}\csromanspacer{}เหตุยา นว.}

\csromanheader{2}{16. มคฺคารมฺมณตฺติก 1. เหตุทุก}
\proseitem{82}{มคฺคารมฺมณํ เหตุํ ธมฺมํ ปฏิจฺจ นมคฺคารมฺมโณ นเหตุ ธมฺโม อุปฺปชฺชติ เหตุปจฺจยา.\csromanspacer{}\csromanspacer{}เหตุยา ปญฺจวีส.}

\prose{มคฺคารมฺมณํ นเหตุํ ธมฺมํ ปฏิจฺจ นมคฺคเหตุโก นนเหตุ ธมฺโม อุปฺปชฺชติ เหตุปจฺจยา.\csromanspacer{}\csromanspacer{}เหตุยา เตรส\footnote{เหตุยา อฏฺฐ?}.}

\vspace{\tipitakaparskip}
\csromancenter{อชฺฌตฺตารมฺมณตฺติก 1.\csromanspacer{} เหตุทุก \textbf{17.}\csromanspacer{}\textbf{ อุปฺปนฺนตฺติก 1.}\csromanspacer{}\textbf{ เหตุทุก}}
\proseitem{83}{\csromanfolio{271}อุปฺปนฺโน เหตุ ธมฺโม น-อนุปฺปนฺนสฺส นเหตุสฺส ธมฺมสฺส เหตุปจฺจเยน ปจฺจโย.\csromanspacer{}\csromanspacer{}อุปฺปนฺโน เหตุ ธมฺโม น-อุปฺปาทิสฺส นเหตุสฺส ธมฺมสฺส เหตุปจฺจเยน ปจฺจโย.\csromanspacer{}\csromanspacer{}อุปฺปนฺโน เหตุ ธมฺโม น-อนุปฺปนฺนสฺส นเหตุสฺส จ น-อุปฺปาทิสฺส นเหตุสฺส จ ธมฺมสฺส เหตุปจฺจเยน ปจฺจโย.\csromanspacer{}\csromanspacer{}เหตุยา ตีณิ.}

\csromanheader{2}{18. อตีตตฺติก 1. เหตุทุก}
\proseitem{84}{ปจฺจุปฺปนฺโน เหตุ ธมฺโม น-อตีตสฺส นเหตุสฺส ธมฺมสฺส เหตุปจฺจเยน ปจฺจโย.\csromanspacer{}\csromanspacer{}ปจฺจุปฺปนฺโน เหตุ ธมฺโม น-อนาคตสฺส นเหตุสฺส ธมฺมสฺส เหตุปจฺจเยน ปจฺจโย.\csromanspacer{}\csromanspacer{}ปจฺจุปฺปนฺโน เหตุ ธมฺโม น-อตีตสฺส นเหตุสฺส จ น-อนาคตสฺส นเหตุสฺส จ ธมฺมสฺส เหตุปจฺจเยน ปจฺจโย.\csromanspacer{}\csromanspacer{}เหตุยา ตีณิ.}

\csromanheader{2}{19. อตีตารมฺมณตฺติก 1. เหตุทุก}
\proseitem{85}{อตีตารมฺมณํ เหตุํ ธมฺมํ ปฏิจฺจ น-อตีตารมฺมโณ นเหตุ ธมฺโม อุปฺปชฺชติ เหตุปจฺจยา.\csromanspacer{}\csromanspacer{}เหตุยา เอกวีส.}

\prose{อตีตารมฺมณํ นเหตุํ ธมฺมํ ปฏิจฺจ น-อนาคตารมฺมโณ นนเหตุ ธมฺโม อุปฺปชฺชติ เหตุปจฺจยา.\csromanspacer{}\csromanspacer{}เหตุยา นว.}

\csromanheader{2}{20. อชฺฌตฺตตฺติก 1. เหตุทุก}
\proseitem{86}{อชฺฌตฺตํ เหตุํ ธมฺมํ ปฏิจฺจ นพหิทฺธา นเหตุ ธมฺโม อุปฺปชฺชติ เหตุปจฺจยา.\csromanspacer{}\csromanspacer{}พหิทฺธา เหตุํ ธมฺมํ ปฏิจฺจ น-อชฺฌตฺโต นเหตุ ธมฺโม อุปฺปชฺชติ เหตุปจฺจยา.\csromanspacer{}\csromanspacer{}เหตุยา เทฺว.}

\prose{อชฺฌตฺตํ นเหตุํ ธมฺมํ ปฏิจฺจ นพหิทฺธา นนเหตุ ธมฺโม อุปฺปชฺชติ เหตุปจฺจยา.\csromanspacer{}\csromanspacer{}พหิทฺธา นเหตุํ ธมฺมํ ปฏิจฺจ น-อชฺฌตฺโต นนเหตุ ธมฺโม อุปฺปชฺชติ เหตุปจฺจยา.\csromanspacer{}\csromanspacer{}เหตุยา เทฺว.}

\csromanheader{2}{21. อชฺฌตฺตารมฺมณตฺติก 1. เหตุทุก}
\proseitem{87}{อชฺฌตฺตารมฺมณํ เหตุํ ธมฺมํ ปฏิจฺจ น-อชฺฌตฺตารมฺมโณ นเหตุ ธมฺโม อุปฺปชฺชติ เหตุปจฺจยา.\csromanspacer{}\csromanspacer{}เหตุยา ฉ.}

\csromanheader{2}{ธมฺมานุโลมปจฺจนีย ติกทุกปฏฺฐานปาฬิ}
\prose{\csromanfolio{272}อชฺฌตฺตารมฺมณํ นเหตุํ ธมฺมํ ปฏิจฺจ นพหิทฺธารมฺมโณ นนเหตุ ธมฺโม อุปฺปชฺชติ เหตุปจฺจยา.\csromanspacer{}\csromanspacer{}เหตุยา เทฺว. \textbf{22.}\csromanspacer{}\textbf{ สนิทสฺสนตฺติก 1.}\csromanspacer{}\textbf{ เหตุทุก}}

\proseitem{88}{อนิทสฺสน-อปฺปฏิฆํ เหตุํ ธมฺมํ ปฏิจฺจ น-อนิทสฺสน-อปฺปฏิโฆ นเหตุ ธมฺโม อุปฺปชฺชติ เหตุปจฺจยา.\csromanspacer{}\csromanspacer{}อนิทสฺสน-อปฺปฏิฆํ เหตุํ ธมฺมํ ปฏิจฺจ นสนิทสฺสนสปฺปฏิโฆ นเหตุ ธมฺโม อุปฺปชฺชติ เหตุปจฺจยา ฯเปฯ.\csromanspacer{}\csromanspacer{}อนิทสฺสน-อปฺปฏิฆํ เหตุํ ธมฺมํ ปฏิจฺจ นสนิทสฺสนสปฺปฏิโฆ นเหตุ จ น-อนิทสฺสนสปฺปฏิโฆ นเหตุ จ ธมฺมา อุปฺปชฺชนฺติ เหตุปจฺจยา.\csromanspacer{}\csromanspacer{}เหตุยา ฉ.}

\proseitem{89}{อนิทสฺสน-อปฺปฏิฆํ นเหตุํ ธมฺมํ ปฏิจฺจ นสนิทสฺสนสปฺปฏิโฆ นนเหตุ ธมฺโม อุปฺปชฺชติ เหตุปจฺจยา.\csromanspacer{}\csromanspacer{}อนิทสฺสน-อปฺปฏิฆํ นเหตุํ ธมฺมํ ปฏิจฺจ น-อนิทสฺสนสปฺปฏิโฆ นนเหตุ ธมฺโม อุปฺปชฺชติ เหตุปจฺจยา.\csromanspacer{}\csromanspacer{}อนิทสฺสน-อปฺปฏิฆํ นเหตุํ ธมฺมํ ปฏิจฺจ นสนิทสฺสนสปฺปฏิโฆ นนเหตุ จ น-อนิทสฺสนสปฺปฏิโฆ นนเหตุ จ ธมฺมา อุปฺปชฺชนฺติ เหตุปจฺจยา.\csromanspacer{}\csromanspacer{}เหตุยา ตีณิ. \textbf{22.}\csromanspacer{}\textbf{ สนิทสฺสนตฺติก 2.}\csromanspacer{}\textbf{ สเหตุทุก}}

\proseitem{90}{อนิทสฺสน-อปฺปฏิฆํ สเหตุกํ ธมฺมํ ปฏิจฺจ น-อนิทสฺสน-อปฺปฏิโฆ นสเหตุโก ธมฺโม อุปฺปชฺชติ เหตุปจฺจยา.\csromanspacer{}\csromanspacer{}อนิทสฺสน-อปฺปฏิฆํ สเหตุกํ ธมฺมํ ปฏิจฺจ นสนิทสฺสนสปฺปฏิโฆ นสเหตุโก ธมฺโม อุปฺปชฺชติ เหตุปจฺจยา ฯเปฯ.\csromanspacer{}\csromanspacer{}อนิทสฺสน-อปฺปฏิฆํ สเหตุกํ ธมฺมํ ปฏิจฺจ นสนิทสฺสนสปฺปฏิโฆ นสเหตุโก จ น-อนิทสฺสนสปฺปฏิโฆ นสเหตุโก จ ธมฺมา อุปฺปชฺชนฺติ เหตุปจฺจยา.\csromanspacer{}\csromanspacer{}เหตุยา ฉ.}

\prose{อนิทสฺสน-อปฺปฏิฆํ อเหตุกํ ธมฺมํ ปฏิจฺจ นสนิทสฺสนสปฺปฏิโฆ น-อเหตุโก ธมฺโม อุปฺปชฺชติ เหตุปจฺจยา. (สํขิตฺตํ.) .\csromanspacer{} เหตุยา ตีณิ.}

\csromanheader{2}{22. สนิทสฺสนตฺติก 3. เหตุสมฺปยุตฺตทุก}
\proseitem{91}{อนิทสฺสน-อปฺปฏิฆํ เหตุสมฺปยุตฺตํ ธมฺมํ ปฏิจฺจ น-อนิทสฺสน-อปฺปฏิโฆ นเหตุสมฺปยุตฺโต ธมฺโม อุปฺปชฺชติ เหตุปจฺจยา.\csromanspacer{}\csromanspacer{}}

\vspace{\tipitakaparskip}
\csromancenter{สนิทสฺสนตฺติก 6.\csromanspacer{} นเหตุสเหตุกทุก อนิทสฺสน-อปฺปฏิฆํ เหตุสมฺปยุตฺตํ ธมฺมํ ปฏิจฺจ นสนิทสฺสนสปฺปฏิโฆ นเหตุสมฺปยุตฺโต ธมฺโม อุปฺปชฺชติ เหตุปจฺจยา ฯเปฯ.\csromanspacer{}\csromanspacer{}อนิทสฺสน-อปฺปฏิฆํ เหตุสมฺปยุตฺตํ ธมฺมํ ปฏิจฺจ นสนิทสฺสนสปฺปฏิโฆ นเหตุสมฺปยุตฺโต จ น-อนิทสฺสนสปฺปฏิโฆ นเหตุสมฺปยุตฺโต จ ธมฺมา อุปฺปชฺชนฺติ เหตุปจฺจยา.\csromanspacer{}\csromanspacer{}เหตุยา ฉ.}
\prose{\csromanfolio{273}อนิทสฺสน-อปฺปฏิฆํ เหตุวิปฺปยุตฺตํ ธมฺมํ ปฏิจฺจ นสนิทสฺสนสปฺปฏิโฆ นเหตุวิปฺปยุตฺโต ธมฺโม อุปฺปชฺชติ เหตุปจฺจยา.\csromanspacer{}\csromanspacer{}เหตุยา ตีณิ.}

\csromanheader{2}{22. สนิทสฺสนตฺติก 4. เหตุสเหตุกทุก}
\proseitem{92}{อนิทสฺสน-อปฺปฏิฆํ เหตุญฺเจว สเหตุกญฺจ ธมฺมํ ปฏิจฺจ นสนิทสฺสนสปฺปฏิโฆ นเหตุ เจว น-อเหตุโก จ ธมฺโม อุปฺปชฺชติ เหตุปจฺจยา.\csromanspacer{}\csromanspacer{}เหตุยา ตีณิ.}

\prose{อนิทสฺสน-อปฺปฏิฆํ สเหตุกญฺเจว น จ เหตุํ ธมฺมํ ปฏิจฺจ นสนิทสฺสนสปฺปฏิโฆ น-อเหตุโก เจว นนเหตุ จ ธมฺโม อุปฺปชฺชติ เหตุปจฺจยา.\csromanspacer{}\csromanspacer{}เหตุยา ตีณิ.}

\csromanheader{2}{22. สนิทสฺสนตฺติก 5. เหตุเหตุสมฺปยุตฺตทุก}
\proseitem{93}{อนิทสฺสน-อปฺปฏิฆํ เหตุญฺเจว เหตุสมฺปยุตฺตญฺจ ธมฺมํ ปฏิจฺจ นสนิทสฺสนสปฺปฏิโฆ นเหตุ เจว นเหตุวิปฺปยุตฺโต จ ธมฺโม อุปฺปชฺชติ เหตุปจฺจยา.\csromanspacer{}\csromanspacer{}เหตุยา ตีณิ. อนิทสฺสน-อปฺปฏิฆํ เหตุสมฺปยุตฺตญฺเจว น จ เหตุํ ธมฺมํ ปฏิจฺจ นสนิทสฺสนสปฺปฏิโฆ นเหตุวิปฺปยุตฺโต เจว นนเหตุ จ ธมฺโม อุปฺปชฺชติ เหตุปจฺจยา.\csromanspacer{}\csromanspacer{}เหตุยา ตีณิ.}

\csromanheader{2}{22. สนิทสฺสนตฺติก 6. นเหตุสเหตุกทุก}
\proseitem{94}{อนิทสฺสน-อปฺปฏิฆํ นเหตุํ สเหตุกํ ธมฺมํ ปฏิจฺจ น-อนิทสฺสน-อปฺปฏิโฆ นเหตุ นสเหตุโก ธมฺโม อุปฺปชฺชติ เหตุปจฺจยา.\csromanspacer{}\csromanspacer{}เหตุยา ฉ.}

\csromanheader{2}{ธมฺมานุโลมปจฺจนีย ติกทุกปฏฺฐานปาฬิ}
\prose{\csromanfolio{274}อนิทสฺสน-อปฺปฏิฆํ นเหตุํ อเหตุกํ ธมฺมํ ปฏิจฺจ นสนิทสฺสนสปฺปฏิโฆ นเหตุ น-อเหตุโก ธมฺโม อุปฺปชฺชติ เหตุปจฺจยา.\csromanspacer{}\csromanspacer{}เหตุยา ตีณิ.}

\csromanheader{2}{22. สนิทสฺสนตฺติก 7. สปฺปจฺจยทุก}
\proseitem{95}{อนิทสฺสน-อปฺปฏิโฆ อปฺปจฺจโย ธมฺโม นสนิทสฺสนสปฺปฏิฆสฺส น-อปฺปจฺจยสฺส ธมฺมสฺส อารมฺมณปจฺจเยน ปจฺจโย.}

\prose{อารมฺมเณ ตีณิ, อธิปติยา ตีณิ, อุปนิสฺสเย ตีณิ. (สํขิตฺตํ.)}

\csromanheader{2}{22. สนิทสฺสนตฺติก 9. สนิทสฺสนทุก}
\proseitem{96}{สนิทสฺสนสปฺปฏิโฆ สนิทสฺสโน ธมฺโม นสนิทสฺสนสปฺปฏิฆสฺส นสนิทสฺสนสฺส ธมฺมสฺส อารมฺมณปจฺจเยน ปจฺจโย. (สํขิตฺตํ.) .\csromanspacer{} อารมฺมเณ ตีณิ, อธิปติยา ตีณิ, อุปนิสฺสเย ปุเรชาเต อตฺถิยา อวิคเต ตีณิ.}

\prose{อนิทสฺสนสปฺปฏิฆํ อนิทสฺสนํ ธมฺมํ ปฏิจฺจ น-อนิทสฺสนสปฺปฏิโฆ น-อนิทสฺสโน ธมฺโม อุปฺปชฺชติ เหตุปจฺจยา.\csromanspacer{}\csromanspacer{}เหตุยา นว.}

\csromanheader{2}{22. สนิทสฺสนตฺติก 10. สปฺปฏิฆทุก}
\proseitem{97}{อนิทสฺสนสปฺปฏิฆํ สปฺปฏิฆํ ธมฺมํ ปฏิจฺจ น-อนิทสฺสนสปฺปฏิโฆ นสปฺปฏิโฆ ธมฺโม อุปฺปชฺชติ เหตุปจฺจยา.\csromanspacer{}\csromanspacer{}อนิทสฺสนสปฺปฏิฆํ สปฺปฏิฆํ ธมฺมํ ปฏิจฺจ นสนิทสฺสนสปฺปฏิโฆ นสปฺปฏิโฆ ธมฺโม อุปฺปชฺชติ เหตุปจฺจยา.\csromanspacer{}\csromanspacer{}อนิทสฺสนสปฺปฏิฆํ สปฺปฏิฆํ ธมฺมํ ปฏิจฺจ นสนิทสฺสนสปฺปฏิโฆ นสปฺปฏิโฆ จ น-อนิทสฺสนสปฺปฏิโฆ นสปฺปฏิโฆ จ ธมฺมา อุปฺปชฺชนฺติ เหตุปจฺจยา.\csromanspacer{}\csromanspacer{}เหตุยา ตีณิ.}

\prose{อนิทสฺสน-อปฺปฏิฆํ อปฺปฏิฆํ ธมฺมํ ปฏิจฺจ น-อนิทสฺสน-อปฺปฏิโฆ น-อปฺปฏิโฆ ธมฺโม อุปฺปชฺชติ เหตุปจฺจยา.\csromanspacer{}\csromanspacer{}อนิทสฺสน-อปฺปฏิฆํ อปฺปฏิฆํ ธมฺมํ ปฏิจฺจ นสนิทสฺสนสปฺปฏิโฆ น-อปฺปฏิโฆ ธมฺโม อุปฺปชฺชติ เหตุปจฺจยา ฯเปฯ.\csromanspacer{} อนิทสฺสน-อปฺปฏิฆํ อปฺปฏิฆํ ธมฺมํ ปฏิจฺจ น-อนิทสฺสนสปฺปฏิโฆ น-อปฺปฏิโฆ จ น-อนิทสฺสน-อปฺปฏิโฆ น-อปฺปฏิโฆ จ ธมฺมา อุปฺปชฺชนฺติ เหตุปจฺจยา.\csromanspacer{}\csromanspacer{}เหตุยา ปญฺจ.}

\vspace{\tipitakaparskip}
\csromancenter{สนิทสฺสนตฺติก 13.\csromanspacer{} เกนจิวิญฺเญยฺยทุก \textbf{22.}\csromanspacer{}\textbf{ สนิทสฺสนตฺติก 11.}\csromanspacer{}\textbf{ รูปีทุก}}
\proseitem{98}{\csromanfolio{275}อนิทสฺสน-อปฺปฏิฆํ รูปิํ ธมฺมํ ปฏิจฺจ นสนิทสฺสนสปฺปฏิโฆ นรูปี ธมฺโม อุปฺปชฺชติ เหตุปจฺจยา.\csromanspacer{}\csromanspacer{}เหตุยา ตีณิ.}

\prose{อนิทสฺสน-อปฺปฏิฆํ อรูปิํ ธมฺมํ ปฏิจฺจ น-อนิทสฺสน-อปฺปฏิโฆ น-อรูปี ธมฺโม อุปฺปชฺชติ เหตุปจฺจยา.\csromanspacer{}\csromanspacer{}อนิทสฺสน-อปฺปฏิฆํ อรูปิํ ธมฺมํ ปฏิจฺจ นสนิทสฺสนสปฺปฏิโฆ น-อรูปี ธมฺโม อุปฺปชฺชติ เหตุปจฺจยา ฯเปฯ.\csromanspacer{}\csromanspacer{}อนิทสฺสน-อปฺปฏิฆํ อรูปิํ ธมฺมํ ปฏิจฺจ นสนิทสฺสนสปฺปฏิโฆ น-อรูปี จ น-อนิทสฺสนสปฺปฏิโฆ น-อรูปี จ ธมฺมา อุปฺปชฺชนฺติ เหตุปจฺจยา.\csromanspacer{}\csromanspacer{}เหตุยา ฉ.}

\csromanheader{2}{22. สนิทสฺสนตฺติก 12. โลกิยทุก}
\proseitem{99}{อนิทสฺสน-อปฺปฏิฆํ โลกิยํ ธมฺมํ ปจฺจยา นสนิทสฺสนสปฺปฏิโฆ นโลกิโย ธมฺโม อุปฺปชฺชติ เหตุปจฺจยา. (สํขิตฺตํ.) .\csromanspacer{} เหตุยา ตีณิ ฯเปฯ อวิคเต ตีณิ.}

\prose{อนิทสฺสน-อปฺปฏิฆํ โลกุตฺตรํ ธมฺมํ ปฏิจฺจ น-อนิทสฺสน-อปฺปฏิโฆ นโลกุตฺตโร ธมฺโม อุปฺปชฺชติ เหตุปจฺจยา.\csromanspacer{}\csromanspacer{}เหตุยา ฉ.}

\csromanheader{2}{22. สนิทสฺสนตฺติก 13. เกนจิวิญฺเญยฺยทุก}
\proseitem{100}{อนิทสฺสนสปฺปฏิฆํ เกนจิ วิญฺเญยฺยํ ธมฺมํ ปฏิจฺจ น-อนิทสฺสนสปฺปฏิโฆ นเกนจิ วิญฺเญยฺโย ธมฺโม อุปฺปชฺชติ เหตุปจฺจยา.\csromanspacer{}\csromanspacer{}อนิทสฺสนสปฺปฏิฆํ เกนจิ วิญฺเญยฺยํ ธมฺมํ ปฏิจฺจ นสนิทสฺสนสปฺปฏิโฆ นเกนจิ วิญฺเญยฺโย ธมฺโม อุปฺปชฺชติ เหตุปจฺจยา.\csromanspacer{}\csromanspacer{}อนิทสฺสนสปฺปฏิฆํ เกนจิ วิญฺเญยฺยํ ธมฺมํ ปฏิจฺจ น-อนิทสฺสน-อปฺปฏิโฆ นเกนจิ วิญฺเญยฺโย ธมฺโม อุปฺปชฺชติ เหตุปจฺจยา ฯเปฯ.\csromanspacer{}\csromanspacer{}อนิทสฺสนสปฺปฏิฆํ เกนจิ วิญฺเญยฺยํ ธมฺมํ ปฏิจฺจ นสนิทสฺสนสปฺปฏิโฆ นเกนจิ วิญฺเญยฺโย จ น-อนิทสฺสนสปฺปฏิโฆ นเกนจิ วิญฺเญยฺโย จ ธมฺมา อุปฺปชฺชนฺติ เหตุปจฺจยา.\csromanspacer{} ฉ.}

\prose{อนิทสฺสน-อปฺปฏิฆํ เกนจิ วิญฺเญยฺยํ ธมฺมํ ปฏิจฺจ น-อนิทสฺสน-อปฺปฏิโฆ นเกนจิ วิญฺเญยฺโย ธมฺโม อุปฺปชฺชติ เหตุปจฺจยา.\csromanspacer{} ฉ.}

\prose{อนิทสฺสนสปฺปฏิฆํ เกนจิ วิญฺเญยฺยญฺจ อนิทสฺสน-อปฺปฏิฆํ เกนจิ วิญฺเญยฺยญฺจ ธมฺมํ ปฏิจฺจ นสนิทสฺสนสปฺปฏิโฆ นเกนจิ วิญฺเญยฺโย ธมฺโม อุปฺปชฺชติ เหตุปจฺจยา.\csromanspacer{} ฉ.\csromanspacer{}\csromanspacer{}เหตุยา อฏฺฐารส.}

\csromanheader{2}{ธมฺมานุโลมปจฺจนีย ติกทุกปฏฺฐานปาฬิ}
\prose{\csromanfolio{276}อนิทสฺสนสปฺปฏิฆํ นเกนจิ วิญฺเญยฺยํ ธมฺมํ ปฏิจฺจ นสนิทสฺสนสปฺปฏิโฆ นนเกนจิ วิญฺเญยฺโย ธมฺโม อุปฺปชฺชติ เหตุปจฺจยา.\csromanspacer{}\csromanspacer{}เหตุยา อฏฺฐารส.}

\csromanheader{2}{22. สนิทสฺสนตฺติก 14. อาสวทุก}
\proseitem{101}{อนิทสฺสน-อปฺปฏิฆํ อาสวํ ธมฺมํ ปฏิจฺจ น-อนิทสฺสน-อปฺปฏิโฆ โน-อาสโว ธมฺโม อุปฺปชฺชติ เหตุปจฺจยา.\csromanspacer{}\csromanspacer{}อนิทสฺสน-อปฺปฏิฆํ อาสวํ ธมฺมํ ปฏิจฺจ นสนิทสฺสนสปฺปฏิโฆ โน-อาสโว ธมฺโม อุปฺปชฺชติ เหตุปจฺจยา ฯเปฯ.\csromanspacer{}\csromanspacer{}อนิทสฺสน-อปฺปฏิฆํ อาสวํ ธมฺมํ ปฏิจฺจ นสนิทสฺสนสปฺปฏิโฆ โน-อาสโว จ น-อนิทสฺสนสปฺปฏิโฆ โน-อาสโว จ ธมฺมา อุปฺปชฺชนฺติ เหตุปจฺจยา.\csromanspacer{}\csromanspacer{}เหตุยา ฉ.}

\prose{อนิทสฺสน-อปฺปฏิฆํ โน-อาสวํ ธมฺมํ ปฏิจฺจ นสนิทสฺสนสปฺปฏิโฆ นโน-อาสโว ธมฺโม อุปฺปชฺชติ เหตุปจฺจยา.\csromanspacer{}\csromanspacer{}อนิทสฺสน-อปฺปฏิฆํ โน-อาสวํ ธมฺมํ ปฏิจฺจ น-อนิทสฺสนสปฺปฏิโฆ นโน-อาสโว ธมฺโม อุปฺปชฺชติ เหตุปจฺจยา.\csromanspacer{}\csromanspacer{}อนิทสฺสน-อปฺปฏิฆํ โน-อาสวํ ธมฺมํ ปฏิจฺจ นสนิทสฺสนสปฺปฏิโฆ นโน-อาสโว จ น-อนิทสฺสนสปฺปฏิโฆ นโน-อาสโว จ ธมฺมา อุปฺปชฺชนฺติ เหตุปจฺจยา.\csromanspacer{}\csromanspacer{}เหตุยา ตีณิ.}

\csromanheader{2}{22. สนิทสฺสนตฺติก 15. สาสวทุก}
\proseitem{102}{อนิทสฺสน-อปฺปฏิฆํ สาสวํ ธมฺมํ ปจฺจยา นสนิทสฺสนสปฺปฏิโฆ นสาสโว ธมฺโม อุปฺปชฺชติ เหตุปจฺจยา.\csromanspacer{}\csromanspacer{}อนิทสฺสน-อปฺปฏิฆํ สาสวํ ธมฺมํ ปจฺจยา น-อนิทสฺสนสปฺปฏิโฆ นสาสโว ธมฺโม อุปฺปชฺชติ เหตุปจฺจยา.\csromanspacer{}\csromanspacer{}อนิทสฺสน-อปฺปฏิฆํ สาสวํ ธมฺมํ ปจฺจยา นสนิทสฺสนสปฺปฏิโฆ นสาสโว จ น-อนิทสฺสนสปฺปฏิโฆ นสาสโว จ ธมฺมา อุปฺปชฺชนฺติ เหตุปจฺจยา.\csromanspacer{}\csromanspacer{}เหตุยา ตีณิ.}

\prose{อนิทสฺสน-อปฺปฏิฆํ อนาสวํ ธมฺมํ ปฏิจฺจ น-อนิทสฺสน-อปฺปฏิโฆ น-อนาสโว ธมฺโม อุปฺปชฺชติ เหตุปจฺจยา.\csromanspacer{}\csromanspacer{}อนิทสฺสน-อปฺปฏิฆํ อนาสวํ ธมฺมํ ปฏิจฺจ นสนิทสฺสนสปฺปฏิโฆ น-อนาสโว ธมฺโม อุปฺปชฺชติ เหตุปจฺจยา ฯเปฯ.\csromanspacer{}\csromanspacer{}อนิทสฺสน-อปฺปฏิฆํ อนาสวํ ธมฺมํ ปฏิจฺจ นสนิทสฺสนสปฺปฏิโฆ น-อนาสโว จ น-อนิทสฺสนสปฺปฏิโฆ น-อนาสโว จ ธมฺมา อุปฺปชฺชนฺติ เหตุปจฺจยา.\csromanspacer{}\csromanspacer{}เหตุยา ฉ.}

\vspace{\tipitakaparskip}
\csromancenter{สนิทสฺสนตฺติก 17.\csromanspacer{} อาสวสาสวทุก \textbf{22.}\csromanspacer{}\textbf{ สนิทสฺสนตฺติก 16.}\csromanspacer{}\textbf{ อาสวสมฺปยุตฺตทุก}}
\proseitem{103}{\csromanfolio{277}อนิทสฺสน-อปฺปฏิฆํ อาสวสมฺปยุตฺตํ ธมฺมํ ปฏิจฺจ น-อนิทสฺสน-อปฺปฏิโฆ น-อาสวสมฺปยุตฺโต ธมฺโม อุปฺปชฺชติ เหตุปจฺจยา.\csromanspacer{}\csromanspacer{}อนิทสฺสน-อปฺปฏิฆํ อาสวสมฺปยุตฺตํ ธมฺมํ ปฏิจฺจ นสนิทสฺสนสปฺปฏิโฆ น-อาสวสมฺปยุตฺโต ธมฺโม อุปฺปชฺชติ เหตุปจฺจยา ฯเปฯ.\csromanspacer{}\csromanspacer{}อนิทสฺสน-อปฺปฏิฆํ อาสวสมฺปยุตฺตํ ธมฺมํ ปฏิจฺจ นสนิทสฺสนสปฺปฏิโฆ น-อาสวสมฺปยุตฺโต จ น-อนิทสฺสนสปฺปฏิโฆ น-อาสวสมฺปยุตฺโต จ ธมฺมา อุปฺปชฺชนฺติ เหตุปจฺจยา.\csromanspacer{}\csromanspacer{}เหตุยา ฉ.}

\prose{อนิทสฺสน-อปฺปฏิฆํ อาสววิปฺปยุตฺตํ ธมฺมํ ปฏิจฺจ นสนิทสฺสนสปฺปฏิโฆ น-อาสววิปฺปยุตฺโต ธมฺโม อุปฺปชฺชติ เหตุปจฺจยา.\csromanspacer{}\csromanspacer{}เหตุยา ตีณิ.}

\csromanheader{2}{22. สนิทสฺสนตฺติก 17. อาสวสาสวทุก}
\proseitem{104}{อนิทสฺสน-อปฺปฏิฆํ อาสวญฺเจว สาสวญฺจ ธมฺมํ ปฏิจฺจ น-อนิทสฺสน-อปฺปฏิโฆ น-อาสโว เจว น-อนาสโว จ ธมฺโม อุปฺปชฺชติ เหตุปจฺจยา.\csromanspacer{}\csromanspacer{}อนิทสฺสน-อปฺปฏิฆํ อาสวญฺเจว สาสวญฺจ ธมฺมํ ปฏิจฺจ นสนิทสฺสนสปฺปฏิโฆ น-อาสโว เจว น-อนาสโว จ ธมฺโม อุปฺปชฺชติ เหตุปจฺจยา ฯเปฯ.\csromanspacer{}\csromanspacer{}อนิทสฺสน-อปฺปฏิฆํ อาสวญฺเจว สาสวญฺจ ธมฺมํ ปฏิจฺจ นสนิทสฺสนสปฺปฏิโฆ น-อาสโว เจว น-อนาสโว จ น-อนิทสฺสนสปฺปฏิโฆ น-อาสโว เจว น-อนาสโว จ ธมฺมา อุปฺปชฺชนฺติ เหตุปจฺจยา.\csromanspacer{}\csromanspacer{}เหตุยา ฉ.}

\prose{อนิทสฺสน-อปฺปฏิฆํ สาสวญฺเจว โน จ อาสวํ ธมฺมํ ปฏิจฺจ นสนิทสฺสนสปฺปฏิโฆ น-อนาสโว เจว นโน-อาสโว จ ธมฺโม อุปฺปชฺชติ เหตุปจฺจยา.\csromanspacer{}\csromanspacer{}อนิทสฺสน-อปฺปฏิฆํ สาสวญฺเจว โน จ อาสวํ ธมฺมํ ปฏิจฺจ น-อนิทสฺสนสปฺปฏิโฆ น-อนาสโว เจว นโน จ อาสโว ธมฺโม อุปฺปชฺชติ เหตุปจฺจยา.\csromanspacer{}\csromanspacer{}อนิทสฺสนสปฺปฏิฆํ สาสวญฺเจว โน จ อาสวํ ธมฺมํ ปฏิจฺจ นสนิทสฺสนสปฺปฏิโฆ น-อนาสโว เจว นโน จ อาสโว น-อนิทสฺสนสปฺปฏิโฆ น-อนาสโว เจว นโน อาสโว จ ธมฺมา อุปฺปชฺชนฺติ เหตุปจฺจยา.\csromanspacer{}\csromanspacer{}เหตุยา ตีณิ.}

\vspace{\tipitakaparskip}
\csromancenter{ธมฺมานุโลมปจฺจนีย ติกทุกปฏฺฐานปาฬิ \textbf{22.}\csromanspacer{}\textbf{ สนิทสฺสนตฺติก 18.}\csromanspacer{}\textbf{ อาสว-อาสวสมฺปยุตฺตทุก}}
\proseitem{105}{\csromanfolio{278}อนิทสฺสน-อปฺปฏิฆํ อาสวญฺเจว อาสวสมฺปยุตฺตญฺจ ธมฺมํ ปฏิจฺจ นสนิทสฺสนสปฺปฏิโฆ น-อาสโว เจว น-อาสววิปฺปยุตฺโต จ ธมฺโม อุปฺปชฺชติ เหตุปจฺจยา.\csromanspacer{}\csromanspacer{}อนิทสฺสน-อปฺปฏิฆํ อาสวญฺเจว อาสวสมฺปยุตฺตญฺจ ธมฺมํ ปฏิจฺจ น-อนิทสฺสนสปฺปฏิโฆ น-อาสโว เจว น-อาสววิปฺปยุตฺโต จ ธมฺโม อุปฺปชฺชติ เหตุปจฺจยา.\csromanspacer{}\csromanspacer{}อนิทสฺสน-อปฺปฏิฆํ อาสวญฺเจว อาสวสมฺปยุตฺตญฺจ ธมฺมํ ปฏิจฺจ นสนิทสฺสนสปฺปฏิโฆ น-อาสโว เจว น-อาสววิปฺปยุตฺโต จ น-อนิทสฺสนสปฺปฏิโฆ น-อาสโว เจว น-อาสววิปฺปยุตฺโต จ ธมฺมา อุปฺปชฺชนฺติ เหตุปจฺจยา.\csromanspacer{}\csromanspacer{}เหตุยา ตีณิ.}

\prose{อนิทสฺสน-อปฺปฏิฆํ อาสวสมฺปยุตฺตญฺเจว โน จ อาสวํ ธมฺมํ ปฏิจฺจ นสนิทสฺสนสปฺปฏิโฆ น-อาสววิปฺปยุตฺโต เจว นโน จ อาสโว ธมฺโม อุปฺปชฺชติ เหตุปจฺจยา.\csromanspacer{}\csromanspacer{}เหตุยา ตีณิ.}

\csromanheader{2}{22. สนิทสฺสนตฺติก 19. อาสววิปฺปยุตฺตสาสวทุก}
\proseitem{106}{อนิทสฺสน-อปฺปฏิฆํ อาสววิปฺปยุตฺตํ สาสวํ ธมฺมํ ปจฺจยา นสนิทสฺสนสปฺปฏิโฆ อาสววิปฺปยุตฺโต นสาสโว ธมฺโม อุปฺปชฺชติ เหตุปจฺจยา.\csromanspacer{}\csromanspacer{}อนิทสฺสน-อปฺปฏิฆํ อาสววิปฺปยุตฺตํ สาสวํ ธมฺมํ ปจฺจยา น-อนิทสฺสนสปฺปฏิโฆ อาสววิปฺปยุตฺโต นสาสโว ธมฺโม อุปฺปชฺชติ เหตุปจฺจยา.\csromanspacer{}\csromanspacer{}อนิทสฺสน-อปฺปฏิฆํ อาสววิปฺปยุตฺตํ สาสวํ ธมฺมํ ปจฺจยา นสนิทสฺสนสปฺปฏิโฆ อาสววิปฺปยุตฺโต นสาสโว จ น-อนิทสฺสนสปฺปฏิโฆ อาสววิปฺปยุตฺโต นสาสโว จ ธมฺมา อุปฺปชฺชนฺติ เหตุปจฺจยา.\csromanspacer{}\csromanspacer{}เหตุยา ตีณิ.}

\prose{อนิทสฺสน-อปฺปฏิฆํ อาสววิปฺปยุตฺตํ อนาสวํ ธมฺมํ ปฏิจฺจ น-อนิทสฺสน-อปฺปฏิโฆ อาสววิปฺปยุตฺโต น-อนาสโว ธมฺโม อุปฺปชฺชติ เหตุปจฺจยา.\csromanspacer{}\csromanspacer{}อนิทสฺสน-อปฺปฏิฆํ อาสววิปฺปยุตฺตํ อนาสวํ ธมฺมํ ปฏิจฺจ นสนิทสฺสนสปฺปฏิโฆ อาสววิปฺปยุตฺโต น-อนาสโว ธมฺโม อุปฺปชฺชติ เหตุปจฺจยา ฯเปฯ.\csromanspacer{}\csromanspacer{}อนิทสฺสน-อปฺปฏิฆํ อาสววิปฺปยุตฺตํ อนาสวํ ธมฺมํ ปฏิจฺจ นสนิทสฺสนสปฺปฏิโฆ อาสววิปฺปยุตฺโต น-อนาสโว จ น-อนิทสฺสนสปฺปฏิโฆ อาสววิปฺปยุตฺโต น-อนาสโว จ ธมฺมา อุปฺปชฺชนฺติ เหตุปจฺจยา.\csromanspacer{}\csromanspacer{}เหตุยา ฉ.}

\vspace{\tipitakaparskip}
\csromancenter{สนิทสฺสนตฺติก 56.\csromanspacer{} จิตฺตทุก \textbf{22.}\csromanspacer{}\textbf{ สนิทสฺสนตฺติก 20-54.}\csromanspacer{}\textbf{ สญฺโญชนาทิทุก}}
\proseitem{107}{\csromanfolio{279}อนิทสฺสน-อปฺปฏิฆํ สญฺโญชนํ ธมฺมํ ปฏิจฺจ น-อนิทสฺสน-อปฺปฏิโฆ โนสญฺโญชโน ธมฺโม อุปฺปชฺชติ เหตุปจฺจยา.}

\proseitem{108}{อนิทสฺสน-อปฺปฏิฆํ คนฺถํ ธมฺมํ ปฏิจฺจ น-อนิทสฺสน-อปฺปฏิโฆ โนคนฺโถ ธมฺโม อุปฺปชฺชติ เหตุปจฺจยา ฯเปฯ.}

\prose{อนิทสฺสน-อปฺปฏิฆํ โอฆํ ธมฺมํ ปฏิจฺจ น-อนิทสฺสน-อปฺปฏิโฆ โน-โอโฆ ธมฺโม อุปฺปชฺชติ เหตุปจฺจยา ฯเปฯ.}

\prose{อนิทสฺสน-อปฺปฏิฆํ โยคํ ธมฺมํ ปฏิจฺจ น-อนิทสฺสน-อปฺปฏิโฆ โนโยโค ธมฺโม อุปฺปชฺชติ เหตุปจฺจยา.}

\proseitem{109}{อนิทสฺสน-อปฺปฏิฆํ นีวรณํ ธมฺมํ ปฏิจฺจ น-อนิทสฺสน-อปฺปฏิโฆ โนนีวรโณ ธมฺโม อุปฺปชฺชติ เหตุปจฺจยา.}

\proseitem{110}{อนิทสฺสน-อปฺปฏิฆํ ปรามาสํ ธมฺมํ ปฏิจฺจ น-อนิทสฺสน-อปฺปฏิโฆ โนปรามาโส ธมฺโม อุปฺปชฺชติ เหตุปจฺจยา.}

\csromanheader{2}{22. สนิทสฺสนตฺติก 55. สารมฺมณทุก}
\proseitem{111}{อนิทสฺสน-อปฺปฏิฆํ สารมฺมณํ ธมฺมํ ปฏิจฺจ น-อนิทสฺสน-อปฺปฏิโฆ นสารมฺมโณ ธมฺโม อุปฺปชฺชติ เหตุปจฺจยา.\csromanspacer{}\csromanspacer{}อนิทสฺสน-อปฺปฏิฆํ สารมฺมณํ ธมฺมํ ปฏิจฺจ นสนิทสฺสนสปฺปฏิโฆ นสารมฺมโณ ธมฺโม อุปฺปชฺชติ เหตุปจฺจยา ฯเปฯ.\csromanspacer{}\csromanspacer{}อนิทสฺสน-อปฺปฏิฆํ สารมฺมณํ ธมฺมํ ปฏิจฺจ นสนิทสฺสนสปฺปฏิโฆ นสารมฺมโณ จ น-อนิทสฺสนสปฺปฏิโฆ นสารมฺมโณ จ ธมฺมา อุปฺปชฺชนฺติ เหตุปจฺจยา.}

\prose{อนิทสฺสน-อปฺปฏิฆํ อนารมฺมณํ ธมฺมํ ปฏิจฺจ นสนิทสฺสนสปฺปฏิโฆ น-อนารมฺมโณ ธมฺโม อุปฺปชฺชติ เหตุปจฺจยา.\csromanspacer{}\csromanspacer{}เหตุยา ตีณิ.}

\csromanheader{2}{22. สนิทสฺสนตฺติก 56. จิตฺตทุก}
\proseitem{112}{อนิทสฺสน-อปฺปฏิฆํ จิตฺตํ ธมฺมํ ปฏิจฺจ น-อนิทสฺสน-อปฺปฏิโฆ โนจิตฺโต ธมฺโม อุปฺปชฺชติ เหตุปจฺจยา.\csromanspacer{}\csromanspacer{}อนิทสฺสน-อปฺปฏิฆํ จิตฺตํ ธมฺมํ ปฏิจฺจ นสนิทสฺสนสปฺปฏิโฆ โนจิตฺโต ธมฺโม อุปฺปชฺชติ เหตุปจฺจยา ฯเปฯ.\csromanspacer{}\csromanspacer{}}

\vspace{\tipitakaparskip}
\csromancenter{ธมฺมานุโลมปจฺจนีย ติกทุกปฏฺฐานปาฬิ อนิทสฺสน-อปฺปฏิฆํ จิตฺตํ ธมฺมํ ปฏิจฺจ นสนิทสฺสนสปฺปฏิโฆ โนจิตฺโต จ น-อนิทสฺสนสปฺปฏิโฆ โนจิตฺโต จ ธมฺมา อุปฺปชฺชนฺติ เหตุปจฺจยา.\csromanspacer{}\csromanspacer{}เหตุยา ฉ.}
\prose{\csromanfolio{280}อนิทสฺสน-อปฺปฏิฆํ โนจิตฺตํ ธมฺมํ ปฏิจฺจ นสนิทสฺสนสปฺปฏิโฆ นโนจิตฺโต ธมฺโม อุปฺปชฺชติ เหตุปจฺจยา.\csromanspacer{}\csromanspacer{}เหตุยา ตีณิ.}

\csromanheader{2}{22. สนิทสฺสนตฺติก 57. เจตสิกทุก}
\proseitem{113}{อนิทสฺสน-อปฺปฏิฆํ เจตสิกํ ธมฺมํ ปฏิจฺจ น-อนิทสฺสน-อปฺปฏิโฆ นเจตสิโก ธมฺโม อุปฺปชฺชติ เหตุปจฺจยา.\csromanspacer{}\csromanspacer{}อนิทสฺสน-อปฺปฏิฆํ เจตสิกํ ธมฺมํ ปฏิจฺจ นสนิทสฺสนสปฺปฏิโฆ นเจตสิโก ธมฺโม อุปฺปชฺชติ เหตุปจฺจยา ฯเปฯ.\csromanspacer{}\csromanspacer{}อนิทสฺสน-อปฺปฏิฆํ เจตสิกํ ธมฺมํ ปฏิจฺจ นสนิทสฺสนสปฺปฏิโฆ นเจตสิโก จ น-อนิทสฺสนสปฺปฏิโฆ นเจตสิโก จ ธมฺมา อุปฺปชฺชนฺติ เหตุปจฺจยา.\csromanspacer{}\csromanspacer{}เหตุยา ฉ.}

\csromanheader{2}{22. สนิทสฺสนตฺติก 58. จิตฺตสมฺปยุตฺตทุก}
\proseitem{114}{อนิทสฺสน-อปฺปฏิฆํ จิตฺตสมฺปยุตฺตํ ธมฺมํ ปฏิจฺจ น-อนิทสฺสน-อปฺปฏิโฆ นจิตฺตสมฺปยุตฺโต ธมฺโม อุปฺปชฺชติ เหตุปจฺจยา.\csromanspacer{}\csromanspacer{}อนิทสฺสน-อปฺปฏิฆํ จิตฺตสมฺปยุตฺตํ ธมฺมํ ปฏิจฺจ นสนิทสฺสนสปฺปฏิโฆ นจิตฺตสมฺปยุตฺโต ธมฺโม อุปฺปชฺชติ เหตุปจฺจยา ฯเปฯ.\csromanspacer{}\csromanspacer{}อนิทสฺสน-อปฺปฏิฆํ จิตฺตสมฺปยุตฺตํ ธมฺมํ ปฏิจฺจ นสนิทสฺสนสปฺปฏิโฆ นจิตฺตสมฺปยุตฺโต จ น-อนิทสฺสนสปฺปฏิโฆ นจิตฺตสมฺปยุตฺโต จ ธมฺมา อุปฺปชฺชนฺติ เหตุปจฺจยา.\csromanspacer{}\csromanspacer{}เหตุยา ฉ.}

\csromanheader{2}{22. สนิทสฺสนตฺติก 59. จิตฺตสํสฏฺฐทุก}
\proseitem{115}{อนิทสฺสน-อปฺปฏิฆํ จิตฺตสํสฏฺฐํ ธมฺมํ ปฏิจฺจ น-อนิทสฺสน-อปฺปฏิโฆ นจิตฺตสํสฏฺโฐ ธมฺโม อุปฺปชฺชติ เหตุปจฺจยา.\csromanspacer{}\csromanspacer{}อนิทสฺสน-อปฺปฏิฆํ จิตฺตสํสฏฺฐํ ธมฺมํ ปฏิจฺจ นสนิทสฺสนสปฺปฏิโฆ นจิตฺตสํสฏฺโฐ ธมฺโม อุปฺปชฺชติ เหตุปจฺจยา ฯเปฯ.\csromanspacer{}\csromanspacer{}อนิทสฺสน-อปฺปฏิฆํ จิตฺตสํสฏฺฐํ ธมฺมํ ปฏิจฺจ นสนิทสฺสนสปฺปฏิโฆ นจิตฺตสํสฏฺโฐ จ น-อนิทสฺสนสปฺปฏิโฆ นจิตฺตสํสฏฺโฐ จ ธมฺมา อุปฺปชฺชนฺติ เหตุปจฺจยา.\csromanspacer{}\csromanspacer{}เหตุยา ฉ.}

\vspace{\tipitakaparskip}
\csromancenter{สนิทสฺสนตฺติก 66.\csromanspacer{} อชฺฌตฺตทุก \textbf{22.}\csromanspacer{}\textbf{ สนิทสฺสนตฺติก 60.}\csromanspacer{}\textbf{ จิตฺตสมุฏฺฐานทุก}}
\proseitem{116}{\csromanfolio{281}อนิทสฺสน-อปฺปฏิฆํ จิตฺตสมุฏฺฐานํ ธมฺมํ ปฏิจฺจ น-อนิทสฺสน-อปฺปฏิโฆ โนจิตฺตสมุฏฺฐาโน ธมฺโม อุปฺปชฺชติ เหตุปจฺจยา.\csromanspacer{}\csromanspacer{}เหตุยา ฉ.}

\csromanheader{2}{22. สนิทสฺสนตฺติก 61. จิตฺตสหภูทุก}
\proseitem{117}{อนิทสฺสน-อปฺปฏิฆํ จิตฺตสหภุํ ธมฺมํ ปฏิจฺจ น-อนิทสฺสน-อปฺปฏิโฆ โนจิตฺตสหภู ธมฺโม อุปฺปชฺชติ เหตุปจฺจยา.\csromanspacer{}\csromanspacer{}เหตุยา ฉ.}

\csromanheader{2}{22. สนิทสฺสนตฺติก 62. จิตฺตานุปริวตฺติทุก}
\proseitem{118}{อนิทสฺสน-อปฺปฏิฆํ จิตฺตานุปริวตฺติํ ธมฺมํ ปฏิจฺจ น-อนิทสฺสน-อปฺปฏิโฆ โนจิตฺตานุปริวตฺตี ธมฺโม อุปฺปชฺชติ เหตุปจฺจยา.\csromanspacer{}\csromanspacer{}เหตุยา ฉ.}

\csromanheader{2}{22. สนิทสฺสนตฺติก 63-65. จิตฺตสํสฏฺฐสมุฏฺฐานทุกาทิ}
\proseitem{119}{อนิทสฺสน-อปฺปฏิฆํ จิตฺตสํสฏฺฐสมุฏฺฐานํ ธมฺมํ ปฏิจฺจ น-อนิทสฺสน-อปฺปฏิโฆ โนจิตฺตสํสฏฺฐสมุฏฺฐาโน ธมฺโม อุปฺปชฺชติ เหตุปจฺจยา.\csromanspacer{}\csromanspacer{}เหตุยา ฉ.}

\proseitem{120}{อนิทสฺสน-อปฺปฏิฆํ จิตฺตสํสฏฺฐสมุฏฺฐานสหภุํ ธมฺมํ ปฏิจฺจ น-อนิทสฺสน-อปฺปฏิโฆ โนจิตฺตสํสฏฺฐสมุฏฺฐานสหภู ธมฺโม อุปฺปชฺชติ เหตุปจฺจยา.\csromanspacer{} เหตุยา ฉ.}

\proseitem{121}{อนิทสฺสน-อปฺปฏิฆํ จิตฺตสํสฏฺฐสมุฏฺฐานานุปริวตฺติํ ธมฺมํ ปฏิจฺจ น-อนิทสฺสน-อปฺปฏิโฆ โนจิตฺตสํสฏฺฐสมุฏฺฐานานุปริวตฺตี ธมฺโม อุปฺปชฺชติ เหตุปจฺจยา.\csromanspacer{}\csromanspacer{}เหตุยา ฉ.}

\csromanheader{2}{22. สนิทสฺสนตฺติก 66. อชฺฌตฺตทุก}
\proseitem{122}{อนิทสฺสน-อปฺปฏิฆํ อชฺฌตฺติกํ ธมฺมํ ปฏิจฺจ น-อนิทสฺสน-อปฺปฏิโฆ น-อชฺฌตฺติโก ธมฺโม อุปฺปชฺชติ เหตุปจฺจยา.\csromanspacer{}\csromanspacer{}อนิทสฺสน-อปฺปฏิฆํ อชฺฌตฺติกํ ธมฺมํ ปฏิจฺจ นสนิทสฺสนสปฺปฏิโฆ น-อชฺฌตฺติโก ธมฺโม อุปฺปชฺชติ เหตุปจฺจยา ฯเปฯ.\csromanspacer{}\csromanspacer{}อนิทสฺสน-อปฺปฏิฆํ อชฺฌตฺติกํ ธมฺมํ ปฏิจฺจ นสนิทสฺสนสปฺปฏิโฆ น-อชฺฌตฺติโก จ น-อนิทสฺสนสปฺปฏิโฆ น-อชฺฌตฺติโก จ ธมฺมา อุปฺปชฺชนฺติ เหตุปจฺจยา.\csromanspacer{}\csromanspacer{}เหตุยา ฉ, อารมฺมเณ ตีณิ ฯเปฯ อวิคเต ฉ.}

\csromanheader{2}{ธมฺมานุโลมปจฺจนีย ติกทุกปฏฺฐานปาฬิ}
\prose{\csromanfolio{282}อนิทสฺสนสปฺปฏิฆํ พาหิรํ ธมฺมํ ปฏิจฺจ นสนิทสฺสนสปฺปฏิโฆ นพาหิโร ธมฺโม อุปฺปชฺชติ เหตุปจฺจยา.\csromanspacer{}\csromanspacer{}เหตุยา เอกาทส.}

\csromanheader{2}{22. สนิทสฺสนตฺติก 67. อุปาทาทุก}
\proseitem{123}{อนิทสฺสน-อปฺปฏิฆํ อุปาทา ธมฺมํ ปฏิจฺจ นสนิทสฺสนสปฺปฏิโฆ โน-อุปาทา ธมฺโม อุปฺปชฺชติ เหตุปจฺจยา.\csromanspacer{}\csromanspacer{}เหตุยา ตีณิ.}

\prose{อนิทสฺสนสปฺปฏิฆํ โน-อุปาทา ธมฺมํ ปฏิจฺจ น-อนิทสฺสนสปฺปฏิโฆ นโน-อุปาทา ธมฺโม อุปฺปชฺชติ เหตุปจฺจยา.\csromanspacer{} ฉ.}

\prose{อนิทสฺสน-อปฺปฏิฆํ โน-อุปาทา ธมฺมํ ปฏิจฺจ น-อนิทสฺสน-อปฺปฏิโฆ นโน-อุปาทา ธมฺโม อุปฺปชฺชติ เหตุปจฺจยา.\csromanspacer{} ฉ.}

\prose{อนิทสฺสนสปฺปฏิฆํ โน-อุปาทา จ อนิทสฺสน-อปฺปฏิฆํ โน-อุปาทา จ ธมฺมํ ปฏิจฺจ นสนิทสฺสนสปฺปฏิโฆ นโน-อุปาทา ธมฺโม อุปฺปชฺชติ เหตุปจฺจยา.\csromanspacer{} ฉ.\csromanspacer{} เหตุยา อฏฺฐารส.}

\csromanheader{2}{22. สนิทสฺสนตฺติก 68. อุปาทินฺนทุก}
\proseitem{124}{อนิทสฺสน-อปฺปฏิฆํ อุปาทินฺนํ ธมฺมํ ปฏิจฺจ น-อนิทสฺสน-อปฺปฏิโฆ น-อุปาทินฺโน ธมฺโม อุปฺปชฺชติ เหตุปจฺจยา.\csromanspacer{}\csromanspacer{}อนิทสฺสน-อปฺปฏิฆํ อุปาทินฺนํ ธมฺมํ ปฏิจฺจ นสนิทสฺสนสปฺปฏิโฆ น-อุปาทินฺโน ธมฺโม อุปฺปชฺชติ เหตุปจฺจยา ฯเปฯ.\csromanspacer{}\csromanspacer{}อนิทสฺสน-อปฺปฏิฆํ อุปาทินฺนํ ธมฺมํ ปฏิจฺจ นสนิทสฺสนสปฺปฏิโฆ น-อุปาทินฺโน จ น-อนิทสฺสนสปฺปฏิโฆ น-อุปาทินฺโน จ ธมฺมา อุปฺปชฺชนฺติ เหตุปจฺจยา.\csromanspacer{}\csromanspacer{}เหตุยา ฉ.}

\proseitem{125}{สนิทสฺสนสปฺปฏิโฆ อนุปาทินฺโน ธมฺโม นสนิทสฺสนสปฺปฏิฆสฺส น-อนุปาทินฺนสฺส ธมฺมสฺส อารมฺมณปจฺจเยน ปจฺจโย. (สํขิตฺตํ.) .\csromanspacer{} อารมฺมเณ นว.}

\csromanheader{2}{22. สนิทสฺสนตฺติก 69-82. อุปาทานทุกาทิ}
\proseitem{126}{อนิทสฺสน-อปฺปฏิฆํ อุปาทานํ ธมฺมํ ปฏิจฺจ น-อนิทสฺสน-อปฺปฏิโฆ โน-อุปาทาโน ธมฺโม อุปฺปชฺชติ เหตุปจฺจยา.\csromanspacer{}\csromanspacer{}เหตุยา ฉ.}

\proseitem{127}{อนิทสฺสน-อปฺปฏิฆํ กิเลสํ ธมฺมํ ปฏิจฺจ น-อนิทสฺสน-อปฺปฏิโฆ โนกิเลโส ธมฺโม อุปฺปชฺชติ เหตุปจฺจยา.\csromanspacer{}\csromanspacer{}เหตุยา ฉ.}

\proseitem{22}{\csromanfolio{283}สนิทสฺสนตฺติก 85.\csromanspacer{} ทสฺสเนนปหาตพฺพเหตุกทุก \textbf{22.}\csromanspacer{}\textbf{ สนิทสฺสนตฺติก 83.}\csromanspacer{}\textbf{ ทสฺสเนนปหาตพฺพทุก}}

\proseitem{128}{อนิทสฺสน-อปฺปฏิฆํ ทสฺสเนน ปหาตพฺพํ ธมฺมํ ปฏิจฺจ น-อนิทสฺสน-อปฺปฏิโฆ นทสฺสเนน ปหาตพฺโพ ธมฺโม อุปฺปชฺชติ เหตุปจฺจยา.\csromanspacer{}\csromanspacer{}เหตุยา ฉ.}

\prose{อนิทสฺสน-อปฺปฏิฆํ นทสฺสเนน ปหาตพฺพํ ธมฺมํ ปจฺจยา นสนิทสฺสนสปฺปฏิโฆ นนทสฺสเนน ปหาตพฺโพ ธมฺโม อุปฺปชฺชติ เหตุปจฺจยา.\csromanspacer{}\csromanspacer{}อนิทสฺสน-อปฺปฏิฆํ นทสฺสเนน ปหาตพฺพํ ธมฺมํ ปจฺจยา น-อนิทสฺสนสปฺปฏิโฆ นนทสฺสเนน ปหาตพฺโพ ธมฺโม อุปฺปชฺชติ เหตุปจฺจยา.\csromanspacer{}\csromanspacer{}อนิทสฺสน-อปฺปฏิฆํ นทสฺสเนน ปหาตพฺพํ ธมฺมํ ปจฺจยา นสนิทสฺสนสปฺปฏิโฆ นนทสฺสเนน ปหาตพฺโพ จ น-อนิทสฺสนสปฺปฏิโฆ นนทสฺสเนน ปหาตพฺโพ จ ธมฺมา อุปฺปชฺชนฺติ เหตุปจฺจยา.\csromanspacer{}\csromanspacer{}เหตุยา ตีณิ.}

\csromanheader{2}{22. สนิทสฺสนตฺติก 84. ภาวนายปหาตพฺพทุก}
\proseitem{129}{อนิทสฺสน-อปฺปฏิฆํ ภาวนาย ปหาตพฺพํ ธมฺมํ ปฏิจฺจ น-อนิทสฺสน-อปฺปฏิโฆ นภาวนาย ปหาตพฺโพ ธมฺโม อุปฺปชฺชติ เหตุปจฺจยา.\csromanspacer{}\csromanspacer{}เหตุยา ฉ.}

\prose{อนิทสฺสน-อปฺปฏิฆํ นภาวนาย ปหาตพฺพํ ธมฺมํ ปจฺจยา นสนิทสฺสนสปฺปฏิโฆ นนภาวนาย ปหาตพฺโพ ธมฺโม อุปฺปชฺชติ เหตุปจฺจยา.\csromanspacer{}\csromanspacer{}เหตุยา ตีณิ.}

\csromanheader{2}{22. สนิทสฺสนตฺติก 85. ทสฺสเนนปหาตพฺพเหตุกทุก}
\proseitem{130}{อนิทสฺสน-อปฺปฏิฆํ ทสฺสเนน ปหาตพฺพเหตุกํ ธมฺมํ ปฏิจฺจ น-อนิทสฺสน-อปฺปฏิโฆ นทสฺสเนน ปหาตพฺพเหตุโก ธมฺโม อุปฺปชฺชติ เหตุปจฺจยา.\csromanspacer{}\csromanspacer{}เหตุยา ฉ.}

\prose{อนิทสฺสน-อปฺปฏิฆํ นทสฺสเนน ปหาตพฺพเหตุกํ ธมฺมํ ปฏิจฺจ นสนิทสฺสนสปฺปฏิโฆ นนทสฺสเนน ปหาตพฺพเหตุโก ธมฺโม อุปฺปชฺชติ เหตุปจฺจยา.\csromanspacer{}\csromanspacer{}เหตุยา ตีณิ.}

\vspace{\tipitakaparskip}
\csromancenter{ธมฺมานุโลมปจฺจนีย ติกทุกปฏฺฐานปาฬิ \textbf{22.}\csromanspacer{}\textbf{ สนิทสฺสนตฺติก 86.}\csromanspacer{}\textbf{ ภาวนายปหาตพฺพเหตุกทุก}}
\proseitem{131}{\csromanfolio{284}อนิทสฺสน-อปฺปฏิฆํ ภาวนาย ปหาตพฺพเหตุกํ ธมฺมํ ปฏิจฺจ น-อนิทสฺสน-อปฺปฏิโฆ นภาวนาย ปหาตพฺพเหตุโก ธมฺโม อุปฺปชฺชติ เหตุปจฺจยา.\csromanspacer{}\csromanspacer{}เหตุยา ฉ.}

\prose{อนิทสฺสน-อปฺปฏิฆํ นภาวนาย ปหาตพฺพเหตุกํ ธมฺมํ ปฏิจฺจ นสนิทสฺสนสปฺปฏิโฆ นนภาวนาย ปหาตพฺพเหตุโก ธมฺโม อุปฺปชฺชติ เหตุปจฺจยา.\csromanspacer{}\csromanspacer{}เหตุยา ตีณิ.}

\csromanheader{2}{22. สนิทสฺสนตฺติก 87. สวิตกฺกทุก}
\proseitem{132}{อนิทสฺสน-อปฺปฏิฆํ สวิตกฺกํ ธมฺมํ ปฏิจฺจ น-อนิทสฺสน-อปฺปฏิโฆ นสวิตกฺโก ธมฺโม อุปฺปชฺชติ เหตุปจฺจยา.\csromanspacer{}\csromanspacer{}อนิทสฺสน-อปฺปฏิฆํ สวิตกฺกํ ธมฺมํ ปฏิจฺจ นสนิทสฺสนสปฺปฏิโฆ นสวิตกฺโก ธมฺโม อุปฺปชฺชติ เหตุปจฺจยา ฯเปฯ.\csromanspacer{}\csromanspacer{}อนิทสฺสน-อปฺปฏิฆํ สวิตกฺกํ ธมฺมํ ปฏิจฺจ นสนิทสฺสนสปฺปฏิโฆ นสวิตกฺโก จ น-อนิทสฺสนสปฺปฏิโฆ นสวิตกฺโก จ ธมฺมา อุปฺปชฺชนฺติ เหตุปจฺจยา.\csromanspacer{}\csromanspacer{}เหตุยา ฉ.}

\proseitem{133}{อนิทสฺสน-อปฺปฏิฆํ อวิตกฺกํ ธมฺมํ ปฏิจฺจ นสนิทสฺสนสปฺปฏิโฆ น-อวิตกฺโก ธมฺโม อุปฺปชฺชติ เหตุปจฺจยา.\csromanspacer{}\csromanspacer{}อนิทสฺสน-อปฺปฏิฆํ อวิตกฺกํ ธมฺมํ ปฏิจฺจ น-อนิทสฺสนสปฺปฏิโฆ น-อวิตกฺโก ธมฺโม อุปฺปชฺชติ เหตุปจฺจยา.\csromanspacer{}\csromanspacer{}อนิทสฺสน-อปฺปฏิฆํ อวิตกฺกํ ธมฺมํ ปฏิจฺจ นสนิทสฺสนสปฺปฏิโฆ น-อวิตกฺโก จ น-อนิทสฺสนสปฺปฏิโฆ น-อวิตกฺโก จ ธมฺมา อุปฺปชฺชนฺติ เหตุปจฺจยา.\csromanspacer{}\csromanspacer{}เหตุยา ตีณิ.}

\csromanheader{2}{22. สนิทสฺสนตฺติก 88. สวิจารทุก}
\proseitem{134}{อนิทสฺสน-อปฺปฏิฆํ สวิจารํ ธมฺมํ ปฏิจฺจ น-อนิทสฺสน-อปฺปฏิโฆ นสวิจาโร ธมฺโม อุปฺปชฺชติ เหตุปจฺจยา.\csromanspacer{}\csromanspacer{}เหตุยา ฉ.}

\prose{อนิทสฺสน-อปฺปฏิฆํ อวิจารํ ธมฺมํ ปฏิจฺจ นสนิทสฺสนสปฺปฏิโฆ น-อวิจาโร ธมฺโม อุปฺปชฺชติ เหตุปจฺจยา.\csromanspacer{}\csromanspacer{}อนิทสฺสน-อปฺปฏิฆํ อวิจารํ ธมฺมํ ปฏิจฺจ น-อนิทสฺสนสปฺปฏิโฆ น-อวิจาโร ธมฺโม อุปฺปชฺชติ เหตุปจฺจยา.\csromanspacer{}\csromanspacer{}อนิทสฺสน-อปฺปฏิฆํ อวิจารํ ธมฺมํ ปฏิจฺจ นสนิทสฺสนสปฺปฏิโฆ น-อวิจาโร จ น-อนิทสฺสนสปฺปฏิโฆ น-อวิจาโร จ ธมฺมา อุปฺปชฺชนฺติ เหตุปจฺจยา.\csromanspacer{}\csromanspacer{}เหตุยา ตีณิ.}

\vspace{\tipitakaparskip}
\csromancenter{สนิทสฺสนตฺติก 91.\csromanspacer{} สุขสหคตทุก \textbf{22.}\csromanspacer{}\textbf{ สนิทสฺสนตฺติก 89.}\csromanspacer{}\textbf{ สปฺปีติกทุก}}
\proseitem{135}{\csromanfolio{285}อนิทสฺสน-อปฺปฏิฆํ สปฺปีติกํ ธมฺมํ ปฏิจฺจ น-อนิทสฺสน-อปฺปฏิโฆ นสปฺปีติโก ธมฺโม อุปฺปชฺชติ เหตุปจฺจยา.\csromanspacer{}\csromanspacer{}อนิทสฺสน-อปฺปฏิฆํ สปฺปีติกํ ธมฺมํ ปฏิจฺจ นสนิทสฺสนสปฺปฏิโฆ นสปฺปีติโก ธมฺโม อุปฺปชฺชติ เหตุปจฺจยา ฯเปฯ.\csromanspacer{}\csromanspacer{}อนิทสฺสน-อปฺปฏิฆํ สปฺปีติกํ ธมฺมํ ปฏิจฺจ นสนิทสฺสนสปฺปฏิโฆ นสปฺปีติโก จ น-อนิทสฺสนสปฺปฏิโฆ นสปฺปีติโก จ ธมฺมา อุปฺปชฺชนฺติ เหตุปจฺจยา.\csromanspacer{}\csromanspacer{}เหตุยา ฉ.}

\prose{อนิทสฺสน-อปฺปฏิฆํ อปฺปีติกํ ธมฺมํ ปฏิจฺจ นสนิทสฺสนสปฺปฏิโฆ น-อปฺปีติโก ธมฺโม อุปฺปชฺชติ เหตุปจฺจยา.\csromanspacer{}\csromanspacer{}อนิทสฺสน-อปฺปฏิฆํ อปฺปีติกํ ธมฺมํ ปฏิจฺจ น-อนิทสฺสนสปฺปฏิโฆ น-อปฺปีติโก ธมฺโม อุปฺปชฺชติ เหตุปจฺจยา.\csromanspacer{}\csromanspacer{}อนิทสฺสน-อปฺปฏิฆํ อปฺปีติกํ ธมฺมํ ปฏิจฺจ นสนิทสฺสนสปฺปฏิโฆ น-อปฺปีติโก จ น-อนิทสฺสนสปฺปฏิโฆ น-อปฺปีติโก จ ธมฺมา อุปฺปชฺชนฺติ เหตุปจฺจยา.\csromanspacer{}\csromanspacer{}เหตุยา ตีณิ.}

\csromanheader{2}{22. สนิทสฺสนตฺติก 90. ปีติสหคตทุก}
\proseitem{136}{อนิทสฺสน-อปฺปฏิฆํ ปีติสหคตํ ธมฺมํ ปฏิจฺจ น-อนิทสฺสน-อปฺปฏิโฆ นปีติสหคโต ธมฺโม อุปฺปชฺชติ เหตุปจฺจยา.\csromanspacer{}\csromanspacer{}อนิทสฺสน-อปฺปฏิฆํ ปีติสหคตํ ธมฺมํ ปฏิจฺจ นสนิทสฺสนสปฺปฏิโฆ นปีติสหคโต ธมฺโม อุปฺปชฺชติ เหตุปจฺจยา ฯเปฯ.\csromanspacer{}\csromanspacer{}อนิทสฺสน-อปฺปฏิฆํ ปีติสหคตํ ธมฺมํ ปฏิจฺจ นสนิทสฺสนสปฺปฏิโฆ นปีติสหคโต จ น-อนิทสฺสนสปฺปฏิโฆ นปีติสหคโต จ ธมฺมา อุปฺปชฺชนฺติ เหตุปจฺจยา.\csromanspacer{}\csromanspacer{}เหตุยา ฉ.}

\prose{อนิทสฺสน-อปฺปฏิฆํ นปีติสหคตํ ธมฺมํ ปฏิจฺจ นสนิทสฺสนสปฺปฏิโฆ นนปีติสหคโต ธมฺโม อุปฺปชฺชติ เหตุปจฺจยา.\csromanspacer{}\csromanspacer{}เหตุยา ตีณิ.}

\csromanheader{2}{22. สนิทสฺสนตฺติก 91. สุขสหคตทุก}
\proseitem{137}{อนิทสฺสน-อปฺปฏิฆํ สุขสหคตํ ธมฺมํ ปฏิจฺจ น-อนิทสฺสน-อปฺปฏิโฆ นสุขสหคโต ธมฺโม อุปฺปชฺชติ เหตุปจฺจยา.\csromanspacer{}\csromanspacer{}เหตุยา ฉ.}

\prose{อนิทสฺสน-อปฺปฏิฆํ นสุขสหคตํ ธมฺมํ ปฏิจฺจ นสนิทสฺสนสปฺปฏิโฆ นนสุขสหคโต ธมฺโม อุปฺปชฺชติ เหตุปจฺจยา.\csromanspacer{}\csromanspacer{}เหตุยา ตีณิ.}

\vspace{\tipitakaparskip}
\csromancenter{ธมฺมานุโลมปจฺจนีย ติกทุกปฏฺฐานปาฬิ \textbf{22.}\csromanspacer{}\textbf{ สนิทสฺสนตฺติก 92.}\csromanspacer{}\textbf{ อุเปกฺขาสหคตทุก}}
\proseitem{138}{\csromanfolio{286}อนิทสฺสน-อปฺปฏิฆํ อุเปกฺขาสหคตํ ธมฺมํ ปฏิจฺจ น-อนิทสฺสน-อปฺปฏิโฆ น-อุเปกฺขาสหคโต ธมฺโม อุปฺปชฺชติ เหตุปจฺจยา.\csromanspacer{}\csromanspacer{}อนิทสฺสน-อปฺปฏิฆํ อุเปกฺขาสหคตํ ธมฺมํ ปฏิจฺจ นสนิทสฺสนสปฺปฏิโฆ น-อุเปกฺขาสหคโต ธมฺโม อุปฺปชฺชติ เหตุปจฺจยา ฯเปฯ.\csromanspacer{}\csromanspacer{}อนิทสฺสน-อปฺปฏิฆํ อุเปกฺขาสหคตํ ธมฺมํ ปฏิจฺจ นสนิทสฺสนสปฺปฏิโฆ น-อุเปกฺขาสหคโต จ น-อนิทสฺสนสปฺปฏิโฆ น-อุเปกฺขาสหคโต จ ธมฺมา อุปฺปชฺชนฺติ เหตุปจฺจยา.\csromanspacer{}\csromanspacer{}เหตุยา ฉ.}

\prose{อนิทสฺสน-อปฺปฏิฆํ น-อุเปกฺขาสหคตํ ธมฺมํ ปฏิจฺจ นสนิทสฺสนสปฺปฏิโฆ นน-อุเปกฺขาสหคโต ธมฺโม อุปฺปชฺชติ เหตุปจฺจยา.\csromanspacer{}\csromanspacer{}เหตุยา ตีณิ.}

\csromanheader{2}{22. สนิทสฺสนตฺติก 93. กามาวจรทุก}
\proseitem{139}{อนิทสฺสน-อปฺปฏิฆํ กามาวจรํ ธมฺมํ ปฏิจฺจ นสนิทสฺสนสปฺปฏิโฆ นกามาวจโร ธมฺโม อุปฺปชฺชติ เหตุปจฺจยา.\csromanspacer{}\csromanspacer{}เหตุยา ตีณิ.}

\prose{อนิทสฺสน-อปฺปฏิฆํ นกามาวจรํ ธมฺมํ ปฏิจฺจ น-อนิทสฺสน-อปฺปฏิโฆ นนกามาวจโร ธมฺโม อุปฺปชฺชติ เหตุปจฺจยา ฯเปฯ.\csromanspacer{}\csromanspacer{}อนิทสฺสน-อปฺปฏิฆํ นกามาวจรํ ธมฺมํ ปฏิจฺจ นสนิทสฺสนสปฺปฏิโฆ นนกามาวจโร จ น-อนิทสฺสนสปฺปฏิโฆ นนกามาวจโร จ ธมฺมา อุปฺปชฺชนฺติ เหตุปจฺจยา.\csromanspacer{}\csromanspacer{}เหตุยา ฉ.}

\csromanheader{2}{22. สนิทสฺสนตฺติก 94. รูปาวจรทุก}
\proseitem{140}{อนิทสฺสน-อปฺปฏิฆํ รูปาวจรํ ธมฺมํ ปฏิจฺจ น-อนิทสฺสน-อปฺปฏิโฆ นรูปาวจโร ธมฺโม อุปฺปชฺชติ เหตุปจฺจยา.\csromanspacer{}\csromanspacer{}เหตุยา ฉ.}

\prose{อนิทสฺสน-อปฺปฏิฆํ นรูปาวจรํ ธมฺมํ ปฏิจฺจ นสนิทสฺสนสปฺปฏิโฆ นนรูปาวจโร ธมฺโม อุปฺปชฺชติ เหตุปจฺจยา.\csromanspacer{}\csromanspacer{}เหตุยา ตีณิ.}

\csromanheader{2}{22. สนิทสฺสนตฺติก 95. อรูปาวจรทุก}
\proseitem{141}{อนิทสฺสน-อปฺปฏิฆํ อรูปาวจรํ ธมฺมํ ปฏิจฺจ น-อนิทสฺสน-อปฺปฏิโฆ น-อรูปาวจโร ธมฺโม อุปฺปชฺชติ เหตุปจฺจยา.\csromanspacer{}\csromanspacer{}เหตุยา ฉ.}

\csromanheader{2}{22. สนิทสฺสนตฺติก 99. ส-อุตฺตรทุก}
\prose{\csromanfolio{287}อนิทสฺสน-อปฺปฏิฆํ น-อรูปาวจรํ ธมฺมํ ปจฺจยา นสนิทสฺสนสปฺปฏิโฆ นน-อรูปาวจโร ธมฺโม อุปฺปชฺชติ เหตุปจฺจยา.\csromanspacer{}\csromanspacer{}เหตุยา ตีณิ.}

\csromanheader{2}{22. สนิทสฺสนตฺติก 96. ปริยาปนฺนทุก}
\proseitem{142}{อนิทสฺสน-อปฺปฏิฆํ ปริยาปนฺนํ ธมฺมํ ปจฺจยา นสนิทสฺสนสปฺปฏิโฆ นปริยาปนฺโน ธมฺโม อุปฺปชฺชติ เหตุปจฺจยา.\csromanspacer{}\csromanspacer{}เหตุยา ตีณิ.}

\prose{อนิทสฺสน-อปฺปฏิฆํ อปริยาปนฺนํ ธมฺมํ ปฏิจฺจ น-อนิทสฺสน-อปฺปฏิโฆ น-อปริยาปนฺโน ธมฺโม อุปฺปชฺชติ เหตุปจฺจยา.\csromanspacer{}\csromanspacer{}อนิทสฺสน-อปฺปฏิฆํ อปริยาปนฺนํ ธมฺมํ ปฏิจฺจ นสนิทสฺสนสปฺปฏิโฆ น-อปริยาปนฺโน ธมฺโม อุปฺปชฺชติ เหตุปจฺจยา ฯเปฯ.\csromanspacer{}\csromanspacer{}อนิทสฺสน-อปฺปฏิฆํ อปริยาปนฺนํ ธมฺมํ ปฏิจฺจ นสนิทสฺสนสปฺปฏิโฆ น-อปริยาปนฺโน จ น-อนิทสฺสนสปฺปฏิโฆ น-อปริยาปนฺโน จ ธมฺมา อุปฺปชฺชนฺติ เหตุปจฺจยา.\csromanspacer{}\csromanspacer{}เหตุยา ฉ.}

\csromanheader{2}{22. สนิทสฺสนตฺติก 97. นิยฺยานิกทุก}
\proseitem{143}{อนิทสฺสน-อปฺปฏิฆํ นิยฺยานิกํ ธมฺมํ ปฏิจฺจ น-อนิทสฺสน-อปฺปฏิโฆ นนิยฺยานิโก ธมฺโม อุปฺปชฺชติ เหตุปจฺจยา.\csromanspacer{}\csromanspacer{}เหตุยา ฉ.}

\prose{อนิทสฺสน-อปฺปฏิฆํ อนิยฺยานิกํ ธมฺมํ ปจฺจยา นสนิทสฺสนสปฺปฏิโฆ น-อนิยฺยานิโก ธมฺโม อุปฺปชฺชติ เหตุปจฺจยา.\csromanspacer{}\csromanspacer{}เหตุยา ตีณิ.}

\csromanheader{2}{22. สนิทสฺสนตฺติก 98. นิยตทุก}
\proseitem{144}{อนิทสฺสน-อปฺปฏิฆํ นิยตํ ธมฺมํ ปฏิจฺจ น-อนิทสฺสน-อปฺปฏิโฆ นนิยโต ธมฺโม อุปฺปชฺชติ เหตุปจฺจยา.\csromanspacer{}\csromanspacer{}เหตุยา ฉ.}

\prose{อนิทสฺสน-อปฺปฏิฆํ อนิยตํ ธมฺมํ ปจฺจยา นสนิทสฺสนสปฺปฏิโฆ น-อนิยโต ธมฺโม อุปฺปชฺชติ เหตุปจฺจยา.\csromanspacer{}\csromanspacer{}เหตุยา ตีณิ.}

\csromanheader{2}{22. สนิทสฺสนตฺติก 99. ส-อุตฺตรทุก}
\proseitem{145}{อนิทสฺสน-อปฺปฏิฆํ ส-อุตฺตรํ ธมฺมํ ปจฺจยา นสนิทสฺสนสปฺปฏิโฆ นส-อุตฺตโร ธมฺโม อุปฺปชฺชติ เหตุปจฺจยา.\csromanspacer{}\csromanspacer{}เหตุยา ตีณิ.}

\prose{อนิทสฺสน-อปฺปฏิฆํ อนุตฺตรํ ธมฺมํ ปฏิจฺจ น-อนิทสฺสน-อปฺปฏิโฆ น-อนุตฺตโร ธมฺโม อุปฺปชฺชติ เหตุปจฺจยา.\csromanspacer{}\csromanspacer{}เหตุยา ฉ.}

\csromanheader{2}{ธมฺมานุโลมปจฺจนีย ติกทุกปฏฺฐานปาฬิ \textbf{22. สนิทสฺสนตฺติก 100. สรณทุก}}
\proseitem{146}{\csromanfolio{288}อนิทสฺสน-อปฺปฏิฆํ สรณํ ธมฺมํ ปฏิจฺจ น-อนิทสฺสน-อปฺปฏิโฆ นสรโณ ธมฺโม อุปฺปชฺชติ เหตุปจฺจยา.\csromanspacer{}\csromanspacer{}เหตุยา ฉ.}

\prose{อนิทสฺสน-อปฺปฏิฆํ อรณํ ธมฺมํ ปจฺจยา นสนิทสฺสนสปฺปฏิโฆ น-อรโณ ธมฺโม อุปฺปชฺชติ เหตุปจฺจยา.\csromanspacer{}\csromanspacer{}เหตุยา ตีณิ. (สหชาตวารมฺปิ ฯเปฯ สมฺปยุตฺตวารมฺปิ ปฏิจฺจวารสทิสํ.)}

\proseitem{147}{สนิทสฺสนสปฺปฏิโฆ อรโณ ธมฺโม นสนิทสฺสนสปฺปฏิฆสฺส น-อรณสฺส ธมฺมสฺส อารมฺมณปจฺจเยน ปจฺจโย.}

\prose{(ยถา กุสลตฺติเก ปญฺหาวารสฺส อนุโลมมฺปิ ปจฺจนียมฺปิ อนุโลมปจฺจนียมฺปิ ปจฺจนียานุโลมมฺปิ คณิตํ, เอวํ คเณตพฺพํ.)}

\nitthitam{\textbf{ธมฺมานุโลมปจฺจนีเย ติกทุกปฏฺฐานํ นิฏฺฐิตํ.}}
\pitaka{\textbf{อภิธมฺมปิฏก}}
\csromanheader{2}{\textbf{ธมฺมานุโลมปจฺจนีย}}
\gambhira{\textbf{ติกติกปฏฺฐานปาฬิ}}
\namakkaram{นโม ตสฺส ภควโต อรหโต สมฺมาสมฺพุทฺธสฺส.}
\csromanheader{2}{1. กุสลตฺติก 1. เวทนาตฺติก}
\proseitem{1}{\csromanfolio{289}กุสลํ สุขาย เวทนาย สมฺปยุตฺตํ ธมฺมํ ปฏิจฺจ นกุสโล นสุขาย เวทนาย สมฺปยุตฺโต ธมฺโม อุปฺปชฺชติ เหตุปจฺจยา.\csromanspacer{}\csromanspacer{}กุสลํ สุขาย เวทนาย สมฺปยุตฺตํ ธมฺมํ ปฏิจฺจ น-อกุสโล นสุขาย เวทนาย สมฺปยุตฺโต ธมฺโม อุปฺปชฺชติ เหตุปจฺจยา ฯเปฯ.\csromanspacer{}\csromanspacer{}กุสลํ สุขาย เวทนาย สมฺปยุตฺตํ ธมฺมํ ปฏิจฺจ นกุสโล นสุขาย เวทนาย สมฺปยุตฺโต จ น-อกุสโล นสุขาย เวทนาย สมฺปยุตฺโต จ ธมฺมา อุปฺปชฺชนฺติ เหตุปจฺจยา.\csromanspacer{} ปญฺจ.}

\prose{อกุสลํ สุขาย เวทนาย สมฺปยุตฺตํ ธมฺมํ ปฏิจฺจ น-อกุสโล นสุขาย เวทนาย สมฺปยุตฺโต ธมฺโม อุปฺปชฺชติ เหตุปจฺจยา.\csromanspacer{}\csromanspacer{}อกุสลํ สุขาย เวทนาย สมฺปยุตฺตํ ธมฺมํ ปฏิจฺจ นกุสโล นสุขาย เวทนาย สมฺปยุตฺโต ธมฺโม อุปฺปชฺชติ เหตุปจฺจยา ฯเปฯ.\csromanspacer{}\csromanspacer{}อกุสลํ สุขาย เวทนาย สมฺปยุตฺตํ ธมฺมํ ปฏิจฺจ นกุสโล นสุขาย เวทนาย สมฺปยุตฺโต จ น-อกุสโล นสุขาย เวทนาย สมฺปยุตฺโต จ ธมฺมา อุปฺปชฺชนฺติ เหตุปจฺจยา.\csromanspacer{} ปญฺจ.}

\prose{อพฺยากตํ สุขาย เวทนาย สมฺปยุตฺตํ ธมฺมํ ปฏิจฺจ นกุสโล นสุขาย เวทนาย สมฺปยุตฺโต ธมฺโม อุปฺปชฺชติ เหตุปจฺจยา.\csromanspacer{}\csromanspacer{}อพฺยากตํ สุขาย เวทนาย สมฺปยุตฺตํ ธมฺมํ ปฏิจฺจ น-อกุสโล นสุขาย}

\vspace{\tipitakaparskip}
\csromancenter{ธมฺมานุโลมปจฺจนีย ติกติกปฏฺฐานปาฬิ เวทนาย สมฺปยุตฺโต ธมฺโม อุปฺปชฺชติ เหตุปจฺจยา.\csromanspacer{}\csromanspacer{}อพฺยากตํ สุขาย เวทนาย สมฺปยุตฺตํ ธมฺมํ ปฏิจฺจ นกุสโล นสุขาย เวทนาย สมฺปยุตฺโต จ น-อกุสโล นสุขาย เวทนาย สมฺปยุตฺโต จ ธมฺมา อุปฺปชฺชนฺติ เหตุปจฺจยา.\csromanspacer{} ตีณิ. (สํขิตฺตํ.)}
\csromancenter{เหตุยา เตรส, อารมฺมเณ นว ฯเปฯ อวิคเต เตรส.}
\csromanheader{2}{\textbf{ปจฺจนีย-นเหตุ น-อารมฺมณ}}
\proseitem{2}{\csromanfolio{290}อพฺยากตํ สุขาย เวทนาย สมฺปยุตฺตํ ธมฺมํ ปฏิจฺจ นกุสโล นสุขาย เวทนาย สมฺปยุตฺโต ธมฺโม อุปฺปชฺชติ นเหตุปจฺจยา.\csromanspacer{} ตีณิ.}

\proseitem{3}{กุสลํ สุขาย เวทนาย สมฺปยุตฺตํ ธมฺมํ ปฏิจฺจ นกุสโล นสุขาย เวทนาย สมฺปยุตฺโต ธมฺโม อุปฺปชฺชติ น-อารมฺมณปจฺจยา.\csromanspacer{}\csromanspacer{}กุสลํ สุขาย เวทนาย สมฺปยุตฺตํ ธมฺมํ ปฏิจฺจ น-อกุสโล นสุขาย เวทนาย สมฺปยุตฺโต ธมฺโม อุปฺปชฺชติ น-อารมฺมณปจฺจยา.\csromanspacer{}\csromanspacer{}กุสลํ สุขาย เวทนาย สมฺปยุตฺตํ ธมฺมํ ปฏิจฺจ นกุสโล นสุขาย เวทนาย สมฺปยุตฺโต จ น-อกุสโล นสุขาย เวทนาย สมฺปยุตฺโต จ ธมฺมา อุปฺปชฺชนฺติ น-อารมฺมณปจฺจยา.\csromanspacer{} ตีณิ.}

\prose{อกุสลํ สุขาย เวทนาย สมฺปยุตฺตํ ธมฺมํ ปฏิจฺจ น-อกุสโล นสุขาย เวทนาย สมฺปยุตฺโต ธมฺโม อุปฺปชฺชติ น-อารมฺมณปจฺจยา.\csromanspacer{} ตีณิ.}

\prose{อพฺยากตํ สุขาย เวทนาย สมฺปยุตฺตํ ธมฺมํ ปฏิจฺจ นกุสโล นสุขาย เวทนาย สมฺปยุตฺโต ธมฺโม อุปฺปชฺชติ น-อารมฺมณปจฺจยา.\csromanspacer{}\csromanspacer{}อพฺยากตํ สุขาย เวทนาย สมฺปยุตฺตํ ธมฺมํ ปฏิจฺจ น-อกุสโล นสุขาย เวทนาย สมฺปยุตฺโต ธมฺโม อุปฺปชฺชติ น-อารมฺมณปจฺจยา.\csromanspacer{}\csromanspacer{}อพฺยากตํ สุขาย เวทนาย สมฺปยุตฺตํ ธมฺมํ ปฏิจฺจ นกุสโล นสุขาย เวทนาย สมฺปยุตฺโต จ น-อกุสโล นสุขาย เวทนาย สมฺปยุตฺโต จ ธมฺมา อุปฺปชฺชนฺติ น-อารมฺมณปจฺจยา.\csromanspacer{} ตีณิ.}

\proseitem{4}{นเหตุยา ตีณิ, น-อารมฺมเณ นว, น-อธิปติยา เตรส ฯเปฯ โนวิคเต นว.}

\prose{เหตุปจฺจยา น-อารมฺมเณ นว.}

\prose{นเหตุปจฺจยา อารมฺมเณ ตีณิ.}

\csromanheader{2}{1. กุสลตฺติก 1. เวทนาตฺติก}
\prose{\csromanfolio{291}(สหชาตวารมฺปิ ปจฺจยวารมฺปิ นิสฺสยวารมฺปิ สํสฏฺฐวารมฺปิ สมฺปยุตฺตวารมฺปิ ปฏิจฺจวารสทิสํ.)}

\proseitem{5}{กุสโล สุขาย เวทนาย สมฺปยุตฺโต ธมฺโม นกุสลสฺส นสุขาย เวทนาย สมฺปยุตฺตสฺส ธมฺมสฺส เหตุปจฺจเยน ปจฺจโย ฯเปฯ.}

\csromanheader{2}{อารมฺมเณ อฏฺฐารส. (ปญฺหาวารมฺปิ วิตฺถาเรตพฺพํ.)}
\csromanheader{2}{\textbf{อกุสลปท-เหตุ}}
\proseitem{6}{อกุสลํ ทุกฺขาย เวทนาย สมฺปยุตฺตํ ธมฺมํ ปฏิจฺจ น-อกุสโล นทุกฺขาย เวทนาย สมฺปยุตฺโต ธมฺโม อุปฺปชฺชติ เหตุปจฺจยา.\csromanspacer{}\csromanspacer{}อกุสลํ ทุกฺขาย เวทนาย สมฺปยุตฺตํ ธมฺมํ ปฏิจฺจ นกุสโล นทุกฺขาย เวทนาย สมฺปยุตฺโต ธมฺโม อุปฺปชฺชติ เหตุปจฺจยา ฯเปฯ.\csromanspacer{}\csromanspacer{}อกุสลํ ทุกฺขาย เวทนาย สมฺปยุตฺตํ ธมฺมํ ปฏิจฺจ นกุสโล นทุกฺขาย เวทนาย สมฺปยุตฺโต จ น-อกุสโล นทุกฺขาย เวทนาย สมฺปยุตฺโต จ ธมฺมา อุปฺปชฺชนฺติ เหตุปจฺจยา.\csromanspacer{}\csromanspacer{}เหตุยา ปญฺจ, อารมฺมเณ ฉ ฯเปฯ อวิคเต อฏฺฐ.}

\csromanheader{2}{\textbf{อพฺยากตปท-ปจฺจนีย}}
\proseitem{7}{อพฺยากตํ ทุกฺขาย เวทนาย สมฺปยุตฺตํ ธมฺมํ ปฏิจฺจ นกุสโล นทุกฺขาย เวทนาย สมฺปยุตฺโต ธมฺโม อุปฺปชฺชติ นเหตุปจฺจยา.\csromanspacer{}\csromanspacer{}อพฺยากตํ ทุกฺขาย เวทนาย สมฺปยุตฺตํ ธมฺมํ ปฏิจฺจ น-อกุสโล นทุกฺขาย เวทนาย สมฺปยุตฺโต ธมฺโม อุปฺปชฺชติ นเหตุปจฺจยา.\csromanspacer{}\csromanspacer{}อพฺยากตํ ทุกฺขาย เวทนาย สมฺปยุตฺตํ ธมฺมํ ปฏิจฺจ นกุสโล นทุกฺขาย เวทนาย สมฺปยุตฺโต จ น-อกุสโล นทุกฺขาย เวทนาย สมฺปยุตฺโต จ ธมฺมา อุปฺปชฺชนฺติ เหตุปจฺจยา.}

\prose{นเหตุยา ตีณิ. (สหชาตวารมฺปิ ฯเปฯ ปญฺหาวารมฺปิ วิตฺถาเรตพฺพํ.)}

\proseitem{8}{กุสลํ อทุกฺขมสุขาย เวทนาย สมฺปยุตฺตํ ธมฺมํ ปฏิจฺจ นกุสโล น-อทุกฺขมสุขาย เวทนาย สมฺปยุตฺโต ธมฺโม อุปฺปชฺชติ เหตุปจฺจยา. (สํขิตฺตํ.) .\csromanspacer{} เหตุยา เตรส. (สพฺพตฺถ วิตฺถาโร.)}

\csromanheader{2}{ธมฺมานุโลมปจฺจนีย ติกติกปฏฺฐานปาฬิ \textbf{1. กุสลตฺติก 2. วิปากตฺติก}}
\proseitem{9}{\csromanfolio{292}อพฺยากตํ วิปากํ ธมฺมํ ปฏิจฺจ นกุสโล นวิปาโก ธมฺโม อุปฺปชฺชติ เหตุปจฺจยา.\csromanspacer{}\csromanspacer{}เหตุยา ตีณิ.}

\proseitem{10}{กุสลํ วิปากธมฺมธมฺมํ ปฏิจฺจ นกุสโล นวิปากธมฺมธมฺโม อุปฺปชฺชติ เหตุปจฺจยา.\csromanspacer{}\csromanspacer{}กุสลํ วิปากธมฺมธมฺมํ ปฏิจฺจ น-อกุสโล นวิปากธมฺมธมฺโม อุปฺปชฺชติ เหตุปจฺจยา.\csromanspacer{}\csromanspacer{}กุสลํ วิปากธมฺมธมฺมํ ปฏิจฺจ นกุสโล นวิปากธมฺมธมฺโม จ น-อกุสโล นวิปากธมฺมธมฺโม จ ธมฺมา อุปฺปชฺชนฺติ เหตุปจฺจยา.\csromanspacer{} ตีณิ.}

\prose{อกุสลํ วิปากธมฺมธมฺมํ ปฏิจฺจ น-อกุสโล นวิปากธมฺมธมฺโม อุปฺปชฺชติ เหตุปจฺจยา.\csromanspacer{}\csromanspacer{}อกุสลํ วิปากธมฺมธมฺมํ ปฏิจฺจ นกุสโล นวิปากธมฺมธมฺโม อุปฺปชฺชติ เหตุปจฺจยา.\csromanspacer{}\csromanspacer{}อกุสลํ วิปากธมฺมธมฺมํ ปฏิจฺจ นกุสโล นวิปากธมฺมธมฺโม จ น-อกุสโล นวิปากธมฺมธมฺโม จ ธมฺมา อุปฺปชฺชนฺติ เหตุปจฺจยา. (3).\csromanspacer{} เหตุยา ฉ.}

\proseitem{11}{อพฺยากตํ เนววิปากนวิปากธมฺมธมฺมํ ปฏิจฺจ นกุสโล นเนววิปากนวิปากธมฺมธมฺโม อุปฺปชฺชติ เหตุปจฺจยา.\csromanspacer{}\csromanspacer{}อพฺยากตํ เนววิปากนวิปากธมฺมธมฺมํ ปฏิจฺจ น-อกุสโล นเนววิปากนวิปากธมฺมธมฺโม อุปฺปชฺชติ เหตุปจฺจยา.\csromanspacer{}\csromanspacer{}อพฺยากตํ เนววิปากนวิปากธมฺมธมฺมํ ปฏิจฺจ นกุสโล นเนววิปากนวิปากธมฺมธมฺโม จ น-อกุสโล เนววิปากนวิปากธมฺมธมฺโม จ ธมฺมา อุปฺปชฺชนฺติ เหตุปจฺจยา.\csromanspacer{}\csromanspacer{}เหตุยา ตีณิ.}

\csromanheader{2}{1. กุสลตฺติก 3. อุปาทินฺนตฺติก}
\proseitem{12}{อพฺยากตํ อุปาทินฺนุปาทานิยํ ธมฺมํ ปฏิจฺจ นกุสโล น-อุปาทินฺนุปาทานิโย ธมฺโม อุปฺปชฺชติ เหตุปจฺจยา.\csromanspacer{}\csromanspacer{}อพฺยากตํ อุปาทินฺนุปาทานิยํ ธมฺมํ ปฏิจฺจ น-อกุสโล น-อุปาทินฺนุปาทานิโย ธมฺโม อุปฺปชฺชติ เหตุปจฺจยา.\csromanspacer{}\csromanspacer{}อพฺยากตํ อุปาทินฺนุปาทานิยํ ธมฺมํ ปฏิจฺจ นกุสโล น-อุปาทินฺนุปาทานิโย จ น-อกุสโล น-อุปาทินฺนุปาทานิโย จ ธมฺมา อุปฺปชฺชนฺติ เหตุปจฺจยา.\csromanspacer{}\csromanspacer{}เหตุยา ตีณิ.}

\proseitem{13}{กุสโล อนุปาทินฺนุปาทานิโย ธมฺโม นกุสลสฺส น-อนุปาทินฺนุปาทานิยสฺส ธมฺมสฺส อารมฺมณปจฺจเยน ปจฺจโย.\csromanspacer{}\csromanspacer{}กุสโล}

\vspace{\tipitakaparskip}
\csromancenter{กุสลตฺติก 4.\csromanspacer{} สํกิลิฏฺฐตฺติก อนุปาทินฺนุปาทานิโย ธมฺโม น-อกุสลสฺส น-อนุปาทินฺนุปาทานิยสฺส ธมฺมสฺส อารมฺมณปจฺจเยน ปจฺจโย.\csromanspacer{}\csromanspacer{}กุสโล อนุปาทินฺนุปาทานิโย ธมฺโม นกุสลสฺส น-อนุปาทินฺนุปาทานิยสฺส จ น-อกุสลสฺส น-อนุปาทินฺนุปาทานิยสฺส จ ธมฺมสฺส อารมฺมณปจฺจเยน ปจฺจโย.\csromanspacer{} ตีณิ.}
\prose{\csromanfolio{293}อกุสเล ตีณิ.\csromanspacer{} อพฺยากตํ อนุปาทินฺนุปาทานิเย ตีณิเยว ฯเปฯ.}

\proseitem{14}{กุสลํ อนุปาทินฺน-อนุปาทานิยํ ธมฺมํ ปฏิจฺจ นกุสโล น-อนุปาทินฺน-อนุปาทานิโย ธมฺโม อุปฺปชฺชติ เหตุปจฺจยา.\csromanspacer{} ตีณิ.}

\prose{อพฺยากตํ อนุปาทินฺน-อนุปาทานิยํ ธมฺมํ ปฏิจฺจ นกุสโล น-อนุปาทินฺน-อนุปาทานิโย ธมฺโม อุปฺปชฺชติ เหตุปจฺจยา.\csromanspacer{} ตีณิ.\csromanspacer{}\csromanspacer{}เหตุยา ฉ.}

\csromanheader{2}{1. กุสลตฺติก 4. สํกิลิฏฺฐตฺติก}
\proseitem{15}{อกุสลํ สํกิลิฏฺฐสํกิเลสิกํ ธมฺมํ ปฏิจฺจ น-อกุสโล นสํกิลิฏฺฐสํกิเลสิโก ธมฺโม อุปฺปชฺชติ เหตุปจฺจยา.\csromanspacer{}\csromanspacer{}เหตุยา ตีณิ.}

\prose{อพฺยากตํ อสํกิลิฏฺฐสํกิเลสิกํ ธมฺมํ ปจฺจยา น-อพฺยากโต น-อสํกิลิฏฺฐสํกิเลสิโก ธมฺโม อุปฺปชฺชติ เหตุปจฺจยา.\csromanspacer{}\csromanspacer{}อพฺยากตํ อสํกิลิฏฺฐสํกิเลสิกํ ธมฺมํ ปจฺจยา นกุสโล น-อสํกิลิฏฺฐสํกิเลสิโก ธมฺโม อุปฺปชฺชติ เหตุปจฺจยา.\csromanspacer{}\csromanspacer{}อพฺยากตํ อสํกิลิฏฺฐสํกิเลสิกํ ธมฺมํ ปจฺจยา น-อกุสโล น-อสํกิลิฏฺฐสํกิเลสิโก ธมฺโม อุปฺปชฺชติ เหตุปจฺจยา.\csromanspacer{}\csromanspacer{}อพฺยากตํ อสํกิลิฏฺฐสํกิเลสิกํ ธมฺมํ ปจฺจยา นกุสโล น-อสํกิลิฏฺฐสํกิเลสิโก จ น-อพฺยากโต น-อสํกิลิฏฺฐสํกิเลสิโก จ ธมฺมา อุปฺปชฺชนฺติ เหตุปจฺจยา.\csromanspacer{}\csromanspacer{}อพฺยากตํ อสํกิลิฏฺฐสํกิเลสิกํ ธมฺมํ ปจฺจยา น-อกุสโล น-อสํกิลิฏฺฐสํกิเลโก จ น-อพฺยากโต น-อสํกิลิฏฺฐสํกิเลสิโก จ ธมฺมา อุปฺปชฺชนฺติ เหตุปจฺจยา.\csromanspacer{}\csromanspacer{}อพฺยากตํ อสํกิลิฏฺฐสํกิเลสิกํ ธมฺมํ ปจฺจยา นกุสโล น-อสํกิลิฏฺฐสํกิเลสิโก จ น-อกุสโล น-อสํกิลิฏฺฐสํกิเลสิโก จ ธมฺมา อุปฺปชฺชนฺติ เหตุปจฺจยา.\csromanspacer{} เหตุยา ฉ.}

\proseitem{16}{กุสลํ อสํกิลิฏฺฐ-อสํกิเลสิกํ ธมฺมํ ปฏิจฺจ นกุสโล น-อสํกิลิฏฺฐ-อสํกิเลสิโก ธมฺโม อุปฺปชฺชติ เหตุปจฺจยา.\csromanspacer{} ตีณิ.}

\csromanheader{2}{ธมฺมานุโลมปจฺจนีย ติกติกปฏฺฐานปาฬิ}
\prose{\csromanfolio{294}อพฺยากตํ อสํกิลิฏฺฐ-อสํกิเลสิกํ ธมฺมํ ปฏิจฺจ นกุสโล น-อสํกิลิฏฺฐ-อสํกิเลสิโก ธมฺโม อุปฺปชฺชติ เหตุปจฺจยา.\csromanspacer{} ตีณิ.\csromanspacer{}\csromanspacer{}เหตุยา ฉ.}

\csromanheader{2}{1. กุสลตฺติก 5. วิตกฺกตฺติก}
\proseitem{17}{กุสลํ สวิตกฺกสวิจารํ ธมฺมํ ปฏิจฺจ นกุสโล นสวิตกฺกสวิจาโร ธมฺโม อุปฺปชฺชติ เหตุปจฺจยา.\csromanspacer{}\csromanspacer{}เหตุยา เตรส. (กุสเล ปญฺจ.\csromanspacer{} อกุสเล ปญฺจ.\csromanspacer{} อพฺยากเต ตีณิ.)}

\prose{กุสลํ อวิตกฺกวิจารมตฺตํ ธมฺมํ ปฏิจฺจ นกุสโล น-อวิตกฺกวิจารมตฺโต ธมฺโม อุปฺปชฺชติ เหตุปจฺจยา.\csromanspacer{}\csromanspacer{}เหตุยา เตรส.}

\prose{กุสลํ อวิตกฺก-อวิจารํ ธมฺมํ ปฏิจฺจ นกุสโล น-อวิตกฺก-อวิจาโร ธมฺโม อุปฺปชฺชติ เหตุปจฺจยา.\csromanspacer{} ตีณิ.}

\prose{อพฺยากตํ อวิตกฺก-อวิจารํ ธมฺมํ ปฏิจฺจ นกุสโล น-อวิตกฺก-อวิจาโร ธมฺโม อุปฺปชฺชติ เหตุปจฺจยา.\csromanspacer{} ตีณิ.\csromanspacer{}\csromanspacer{}เหตุยา ฉ.}

\csromanheader{2}{1. กุสลตฺติก 6. ปีติตฺติก}
\proseitem{18}{กุสลํ ปีติสหคตํ ธมฺมํ ปฏิจฺจ นกุสโล นปีติสหคโต ธมฺโม อุปฺปชฺชติ เหตุปจฺจยา.\csromanspacer{} ปญฺจ.\csromanspacer{}\csromanspacer{}เหตุยา เตรส.}

\prose{กุสลํ สุขสหคตํ ธมฺมํ ปฏิจฺจ นกุสโล นสุขสหคโต ธมฺโม อุปฺปชฺชติ เหตุปจฺจยา.\csromanspacer{} ปญฺจ.\csromanspacer{}\csromanspacer{}เหตุยา เตรส.}

\prose{กุสลํ อุเปกฺขาสหคตํ ธมฺมํ ปฏิจฺจ นกุสโล น-อุเปกฺขาสหคโต ธมฺโม อุปฺปชฺชติ เหตุปจฺจยา.\csromanspacer{} ปญฺจ.\csromanspacer{}\csromanspacer{}เหตุยา เตรส.}

\csromanheader{2}{1. กุสลตฺติก 7. ทสฺสนตฺติก}
\proseitem{19}{อกุสลํ ทสฺสเนน ปหาตพฺพํ ธมฺมํ ปฏิจฺจ น-อกุสโล นทสฺสเนน ปหาตพฺโพ ธมฺโม อุปฺปชฺชติ เหตุปจฺจยา.\csromanspacer{}\csromanspacer{}เหตุยา ตีณิ.}

\prose{อกุสลํ ภาวนาย ปหาตพฺพํ ธมฺมํ ปฏิจฺจ น-อกุสโล นภาวนาย ปหาตพฺโพ ธมฺโม อุปฺปชฺชติ เหตุปจฺจยา.\csromanspacer{}\csromanspacer{}เหตุยา ตีณิ.}

\csromanheader{2}{1. กุสลตฺติก 10. เสกฺขตฺติก}
\prose{\csromanfolio{295}อพฺยากตํ เนวทสฺสเนน นภาวนาย ปหาตพฺพํ ธมฺมํ ปจฺจยา น-อพฺยากโต นเนวทสฺสเนน นภาวนาย ปหาตพฺโพ ธมฺโม อุปฺปชฺชติ เหตุปจฺจยา.\csromanspacer{}\csromanspacer{}เหตุยา ตีณิ.}

\csromanheader{2}{1. กุสลตฺติก 8. ทสฺสนเหตุตฺติก}
\proseitem{20}{อกุสลํ ทสฺสเนน ปหาตพฺพเหตุกํ ธมฺมํ ปฏิจฺจ น-อกุสโล นทสฺสเนน ปหาตพฺพเหตุโก ธมฺโม อุปฺปชฺชติ เหตุปจฺจยา.\csromanspacer{}\csromanspacer{}เหตุยา ตีณิ.}

\prose{อกุสลํ ภาวนาย ปหาตพฺพเหตุกํ ธมฺมํ ปฏิจฺจ น-อกุสโล นภาวนาย ปหาตพฺพเหตุโก ธมฺโม อุปฺปชฺชติ เหตุปจฺจยา.\csromanspacer{}\csromanspacer{}เหตุยา ตีณิ.}

\prose{อกุสลํ เนวทสฺสเนน นภาวนาย ปหาตพฺพเหตุกํ ธมฺมํ ปฏิจฺจ นกุสโล นเนวทสฺสเนน นภาวนาย ปหาตพฺพเหตุโก ธมฺโม อุปฺปชฺชติ เหตุปจฺจยา.\csromanspacer{}\csromanspacer{}เหตุยา ตีณิ.}

\prose{อพฺยากตํ เนวทสฺสเนน นภาวนาย ปหาตพฺพเหตุกํ ธมฺมํ ปจฺจยา น-อพฺยากโต นเนวทสฺสเนน นภาวนาย ปหาตพฺพเหตุโก ธมฺโม อุปฺปชฺชติ เหตุปจฺจยา.\csromanspacer{}\csromanspacer{}เหตุยา ตีณิ.}

\csromanheader{2}{1. กุสลตฺติก 9. อาจยคามิตฺติก}
\proseitem{21}{กุสลํ อาจยคามิํ ธมฺมํ ปฏิจฺจ นกุสโล น-อาจยคามี ธมฺโม อุปฺปชฺชติ เหตุปจฺจยา.\csromanspacer{}\csromanspacer{}เหตุยา ฉ.}

\prose{กุสลํ อปจยคามิํ ธมฺมํ ปฏิจฺจ นกุสโล น-อปจยคามี ธมฺโม อุปฺปชฺชติ เหตุปจฺจยา.\csromanspacer{}\csromanspacer{}เหตุยา ตีณิ.}

\prose{อพฺยากตํ เนวาจยคามินาปจยคามิํ ธมฺมํ ปจฺจยา น-อพฺยากโต นเนวาจยคามินาปจยคามี ธมฺโม อุปฺปชฺชติ เหตุปจฺจยา.\csromanspacer{}\csromanspacer{}เหตุยา ปญฺจ.}

\csromanheader{2}{1. กุสลตฺติก 10. เสกฺขตฺติก}
\proseitem{22}{กุสลํ เสกฺขํ ธมฺมํ ปฏิจฺจ นกุสโล นเสกฺโข ธมฺโม อุปฺปชฺชติ เหตุปจฺจยา. (เทฺว มูลานิ.) .\csromanspacer{} เหตุยา ฉ.}

\csromanheader{2}{ธมฺมานุโลมปจฺจนีย ติกติกปฏฺฐานปาฬิ}
\prose{\csromanfolio{296}อพฺยากตํ อเสกฺขํ ธมฺมํ ปฏิจฺจ นกุสโล น-อเสกฺโข ธมฺโม อุปฺปชฺชติ เหตุปจฺจยา.\csromanspacer{}\csromanspacer{}เหตุยา ตีณิ.}

\prose{อพฺยากตํ เนวเสกฺขนาเสกฺขํ ธมฺมํ ปจฺจยา น-อพฺยากโต นเนวเสกฺขนาเสกฺโข ธมฺโม อุปฺปชฺชติ เหตุปจฺจยา.\csromanspacer{}\csromanspacer{}เหตุยา ปญฺจ.}

\csromanheader{2}{1. กุสลตฺติก 11. ปริตฺตตฺติก}
\proseitem{23}{อพฺยากตํ ปริตฺตํ ธมฺมํ ปฏิจฺจ นกุสโล นปริตฺโต ธมฺโม อุปฺปชฺชติ เหตุปจฺจยา.\csromanspacer{}\csromanspacer{}เหตุยา ตีณิ.}

\prose{กุสลํ มหคฺคตํ ธมฺมํ ปฏิจฺจ นกุสโล นมหคฺคโต ธมฺโม อุปฺปชฺชติ เหตุปจฺจยา.\csromanspacer{}\csromanspacer{}เหตุยา ฉ.}

\prose{กุสลํ อปฺปมาณํ ธมฺมํ ปฏิจฺจ นกุสโล น-อปฺปมาโณ ธมฺโม อุปฺปชฺชติ เหตุปจฺจยา.\csromanspacer{}\csromanspacer{}เหตุยา ฉ.}

\csromanheader{2}{1. กุสลตฺติก 12. ปริตฺตารมฺมณตฺติก}
\proseitem{24}{กุสลํ ปริตฺตารมฺมณํ ธมฺมํ ปฏิจฺจ นกุสโล นปริตฺตารมฺมโณ ธมฺโม อุปฺปชฺชติ เหตุปจฺจยา.\csromanspacer{}\csromanspacer{}เหตุยา นว.}

\prose{กุสลํ มหคฺคตารมฺมณํ ธมฺมํ ปฏิจฺจ นกุสโล นมหคฺคตารมฺมโณ ธมฺโม อุปฺปชฺชติ เหตุปจฺจยา.\csromanspacer{}\csromanspacer{}เหตุยา นว.}

\prose{กุสลํ อปฺปมาณารมฺมณํ ธมฺมํ ปฏิจฺจ นกุสโล น-อปฺปมาณารมฺมโณ ธมฺโม อุปฺปชฺชติ เหตุปจฺจยา.\csromanspacer{}\csromanspacer{}เหตุยา ฉ.}

\csromanheader{2}{1. กุสลตฺติก 13. หีนตฺติก}
\proseitem{25}{อกุสลํ หีนํ ธมฺมํ ปฏิจฺจ น-อกุสโล นหีโน ธมฺโม อุปฺปชฺชติ เหตุปจฺจยา.\csromanspacer{}\csromanspacer{}เหตุยา ตีณิ.}

\prose{อพฺยากตํ มชฺฌิมํ ธมฺมํ ปจฺจยา น-อพฺยากโต นมชฺฌิโม ธมฺโม อุปฺปชฺชติ เหตุปจฺจยา.\csromanspacer{}\csromanspacer{}เหตุยา ฉ.}

\prose{กุสลํ ปณีตํ ธมฺมํ ปฏิจฺจ นกุสโล นปณีโต ธมฺโม อุปฺปชฺชติ เหตุปจฺจยา.\csromanspacer{}\csromanspacer{}เหตุยา ฉ.}

\vspace{\tipitakaparskip}
\csromancenter{กุสลตฺติก 18.\csromanspacer{} อตีตารมฺมณตฺติก \textbf{1.}\csromanspacer{}\textbf{ กุสลตฺติก 14.}\csromanspacer{}\textbf{ มิจฺฉตฺตนิยตฺตตฺติก}}
\proseitem{26}{\csromanfolio{297}อกุสลํ มิจฺฉตฺตนิยตํ ธมฺมํ ปฏิจฺจ น-อกุสโล นมิจฺฉตฺตนิยโต ธมฺโม อุปฺปชฺชติ เหตุปจฺจยา.\csromanspacer{}\csromanspacer{}เหตุยา ตีณิ.}

\prose{กุสลํ สมฺมตฺตนิยตํ ธมฺมํ ปฏิจฺจ นกุสโล นสมฺมตฺตนิยโต ธมฺโม อุปฺปชฺชติ เหตุปจฺจยา.\csromanspacer{}\csromanspacer{}เหตุยา ตีณิ.}

\prose{อพฺยากตํ อนิยตํ ธมฺมํ ปจฺจยา น-อพฺยากโต น-อนิยโต ธมฺโม อุปฺปชฺชติ เหตุปจฺจยา.\csromanspacer{}\csromanspacer{}เหตุยา ปญฺจ.}

\csromanheader{2}{1. กุสลตฺติก 15. มคฺคารมฺมณตฺติก}
\proseitem{27}{กุสลํ มคฺคารมฺมณํ ธมฺมํ ปฏิจฺจ นกุสโล นมคฺคารมฺมโณ ธมฺโม อุปฺปชฺชติ เหตุปจฺจยา.\csromanspacer{}\csromanspacer{}เหตุยา ฉ.}

\prose{กุสลํ มคฺคเหตุกํ ธมฺมํ ปฏิจฺจ นกุสโล นมคฺคเหตุโก ธมฺโม อุปฺปชฺชติ เหตุปจฺจยา.\csromanspacer{}\csromanspacer{}เหตุยา ตีณิ.}

\prose{กุสลํ มคฺคาธิปติํ ธมฺมํ ปฏิจฺจ นกุสโล นมคฺคาธิปติ ธมฺโม อุปฺปชฺชติ เหตุปจฺจยา.\csromanspacer{}\csromanspacer{}เหตุยา ฉ.}

\csromanheader{2}{1. กุสลตฺติก 16. อุปฺปนฺนตฺติก}
\proseitem{28}{กุสโล อนุปฺปนฺโน ธมฺโม นกุสลสฺส น-อนุปฺปนฺนสฺส ธมฺมสฺส อารมฺมณปจฺจเยน ปจฺจโย.\csromanspacer{}\csromanspacer{}อารมฺมเณ อฏฺฐารส. (อุปฺปาที อนุปฺปนฺนสทิสํ.)}

\csromanheader{2}{1. กุสลตฺติก 17. อตีตตฺติก}
\proseitem{29}{กุสโล อตีโต ธมฺโม นกุสลสฺส น-อตีตสฺส ธมฺมสฺส อารมฺมณปจฺจเยน ปจฺจโย.}

\csromanheader{2}{อารมฺมเณ อฏฺฐารส. (อนาคตํ อตีตสทิสํ.)}
\csromanheader{2}{1. กุสลตฺติก 18. อตีตารมฺมณตฺติก}
\proseitem{30}{กุสลํ อตีตารมฺมณํ ธมฺมํ ปฏิจฺจ นกุสโล น-อตีตารมฺมโณ ธมฺโม อุปฺปชฺชติ เหตุปจฺจยา.\csromanspacer{}\csromanspacer{}เหตุยา นว.}

\csromanheader{2}{ธมฺมานุโลมปจฺจนีย ติกติกปฏฺฐานปาฬิ}
\prose{\csromanfolio{298}กุสลํ อนาคตารมฺมณํ ธมฺมํ ปฏิจฺจ นกุสโล น-อนาคตารมฺมโณ ธมฺโม อุปฺปชฺชติ เหตุปจฺจยา.\csromanspacer{}\csromanspacer{}เหตุยา นว.}

\prose{กุสลํ ปจฺจุปฺปนฺนารมฺมณํ ธมฺมํ ปฏิจฺจ นกุสโล นปจฺจุปฺปนฺนารมฺมโณ ธมฺโม อุปฺปชฺชติ เหตุปจฺจยา.\csromanspacer{}\csromanspacer{}เหตุยา นว.}

\csromanheader{2}{1. กุสลตฺติก 20-19. อชฺฌตฺตตฺติกทฺวย}
\proseitem{31}{กุสโล อชฺฌตฺโต ธมฺโม น-อชฺฌตฺตสฺส นกุสลสฺส ธมฺมสฺส อารมฺมณปจฺจเยน ปจฺจโย. (สํขิตฺตํ.) .\csromanspacer{} อารมฺมเณ อฏฺฐารส.}

\prose{กุสโล พหิทฺธา ธมฺโม นพหิทฺธา นกุสลสฺส ธมฺมสฺส อารมฺมณปจฺจเยน ปจฺจโย. (สํขิตฺตํ.) .\csromanspacer{} อารมฺมเณ อฏฺฐารส.}

\proseitem{32}{กุสลํ อชฺฌตฺตารมฺมณํ ธมฺมํ ปฏิจฺจ นกุสโล น-อชฺฌตฺตารมฺมโณ ธมฺโม อุปฺปชฺชติ เหตุปจฺจยา.\csromanspacer{}\csromanspacer{}เหตุยา นว.}

\prose{กุสลํ พหิทฺธารมฺมณํ ธมฺมํ ปฏิจฺจ นกุสโล นพหิทฺธารมฺมโณ ธมฺโม อุปฺปชฺชติ เหตุปจฺจยา.\csromanspacer{}\csromanspacer{}เหตุยา นว.}

\csromanheader{2}{1. กุสลตฺติก 21. สนิทสฺสนตฺติก}
\proseitem{33}{อพฺยากโต สนิทสฺสนสปฺปฏิโฆ ธมฺโม น-อพฺยากตสฺส นสนิทสฺสนสปฺปฏิฆสฺส ธมฺมสฺส อารมฺมณปจฺจเยน ปจฺจโย.}

\csromanheader{2}{อารมฺมเณ ฉ. (ตีณิ เวทิตกํ กาตพฺพํ.)}
\proseitem{34}{อพฺยากตํ อนิทสฺสนสปฺปฏิฆํ ธมฺมํ ปฏิจฺจ กุสโล น-อนิทสฺสนสปฺปฏิโฆ ธมฺโม อุปฺปชฺชติ เหตุปจฺจยา.\csromanspacer{}\csromanspacer{}อพฺยากตํ อนิทสฺสนสปฺปฏิฆํ ธมฺมํ ปฏิจฺจ น-อกุสโล น-อนิทสฺสนสปฺปฏิโฆ ธมฺโม อุปฺปชฺชติ เหตุปจฺจยา.\csromanspacer{}\csromanspacer{}อพฺยากตํ อนิทสฺสนสปฺปฏิฆํ ธมฺมํ ปฏิจฺจ นกุสโล น-อนิทสฺสนสปฺปฏิโฆ จ น-อกุสโล น-อนิทสฺสนสปฺปฏิโฆ จ ธมฺมา อุปฺปชฺชนฺติ เหตุปจฺจยา.\csromanspacer{}\csromanspacer{}เหตุยา ตีณิ.}

\proseitem{35}{กุสลํ อนิทสฺสน-อปฺปฏิฆํ ธมฺมํ ปฏิจฺจ นกุสโล น-อนิทสฺสน-อปฺปฏิโฆ ธมฺโม อุปฺปชฺชติ เหตุปจฺจยา.\csromanspacer{}\csromanspacer{}กุสลํ อนิทสฺสน-อปฺปฏิฆํ ธมฺมํ ปฏิจฺจ น-อกุสโล น-อนิทสฺสน-อปฺปฏิโฆ ธมฺโม อุปฺปชฺชติ}

\vspace{\tipitakaparskip}
\csromancenter{เวทนาตฺติก 1.\csromanspacer{} กุสลตฺติก เหตุปจฺจยา.\csromanspacer{}\csromanspacer{}กุสลํ อนิทสฺสน-อปฺปฏิฆํ ธมฺมํ ปฏิจฺจ นกุสโล น-อนิทสฺสน-อปฺปฏิโฆ จ น-อกุสโล น-อนิทสฺสน-อปฺปฏิโฆ จ ธมฺมา อุปฺปชฺชนฺติ เหตุปจฺจยา.\csromanspacer{} ตีณิ.}
\prose{\csromanfolio{299}อกุสลํ อนิทสฺสน-อปฺปฏิฆํ ธมฺมํ ปฏิจฺจ น-อกุสโล น-อนิทสฺสน-อปฺปฏิโฆ ธมฺโม อุปฺปชฺชติ เหตุปจฺจยา.\csromanspacer{}\csromanspacer{}อกุสลํ อนิทสฺสน-อปฺปฏิฆํ ธมฺมํ ปฏิจฺจ นกุสโล น-อนิทสฺสน-อปฺปฏิโฆ ธมฺโม อุปฺปชฺชติ เหตุปจฺจยา.\csromanspacer{}\csromanspacer{}อกุสลํ อนิทสฺสน-อปฺปฏิฆํ ธมฺมํ ปฏิจฺจ นกุสโล น-อนิทสฺสน-อปฺปฏิโฆ จ น-อกุสโล น-อนิทสฺสน-อปฺปฏิโฆ จ ธมฺมา อุปฺปชฺชนฺติ เหตุปจฺจยา.\csromanspacer{} ตีณิ.}

\prose{อพฺยากตํ อนิทสฺสน-อปฺปฏิฆํ ธมฺมํ ปฏิจฺจ นกุสโล น-อนิทสฺสน-อปฺปฏิโฆ ธมฺโม อุปฺปชฺชติ เหตุปจฺจยา.\csromanspacer{}\csromanspacer{}อพฺยากตํ อนิทสฺสน-อปฺปฏิฆํ ธมฺมํ ปฏิจฺจ น-อกุสโล น-อนิทสฺสน-อปฺปฏิโฆ ธมฺโม อุปฺปชฺชติ เหตุปจฺจยา.\csromanspacer{}\csromanspacer{}อพฺยากตํ อนิทสฺสน-อปฺปฏิฆํ ธมฺมํ ปฏิจฺจ นกุสโล น-อนิทสฺสน-อปฺปฏิโฆ จ น-อกุสโล น-อนิทสฺสน-อปฺปฏิโฆ จ ธมฺมา อุปฺปชฺชนฺติ เหตุปจฺจยา.\csromanspacer{} ตีณิ. กุสลํ อนิทสฺสน-อปฺปฏิฆญฺจ อพฺยากตํ อนิทสฺสน-อปฺปฏิฆญฺจ ธมฺมํ ปฏิจฺจ นกุสโล น-อนิทสฺสน-อปฺปฏิโฆ ธมฺโม อุปฺปชฺชติ เหตุปจฺจยา.\csromanspacer{} ตีณิ.}

\prose{อกุสลํ อนิทสฺสน-อปฺปฏิฆญฺจ อพฺยากตํ อนิทสฺสน-อปฺปฏิฆญฺจ ธมฺมํ ปฏิจฺจ นกุสโล น-อนิทสฺสน-อปฺปฏิโฆ ธมฺโม อุปฺปชฺชติ เหตุปจฺจยา.\csromanspacer{} ตีณิ.\csromanspacer{} เหตุยา ปนฺนรส.}

\csromanheader{2}{2. เวทนาตฺติก 1. กุสลตฺติก}
\proseitem{36}{สุขาย เวทนาย สมฺปยุตฺตํ กุสลํ ธมฺมํ ปฏิจฺจ นสุขาย เวทนาย สมฺปยุตฺโต นกุสโล ธมฺโม อุปฺปชฺชติ เหตุปจฺจยา.\csromanspacer{}\csromanspacer{}สุขาย เวทนาย สมฺปยุตฺตํ กุสลํ ธมฺมํ ปฏิจฺจ นทุกฺขาย เวทนาย สมฺปยุตฺโต นกุสโล ธมฺโม อุปฺปชฺชติ เหตุปจฺจยา.\csromanspacer{}\csromanspacer{}สุขาย เวทนาย สมฺปยุตฺตํ กุสลํ ธมฺมํ ปฏิจฺจ น-อทุกฺขมสุขาย เวทนาย สมฺปยุตฺโต นกุสโล ธมฺโม อุปฺปชฺชติ เหตุปจฺจยา.\csromanspacer{}\csromanspacer{}สุขาย เวทนาย สมฺปยุตฺตํ กุสลํ ธมฺมํ ปฏิจฺจ นสุขาย เวทนาย สมฺปยุตฺโต นกุสโล จ น-อทุกฺขมสุขาย เวทนาย สมฺปยุตฺโต นกุสโล จ ธมฺมา อุปฺปชฺชนฺติ}

\vspace{\tipitakaparskip}
\csromancenter{ธมฺมานุโลมปจฺจนีย ติกติกปฏฺฐานปาฬิ เหตุปจฺจยา.\csromanspacer{}\csromanspacer{}สุขาย เวทนาย สมฺปยุตฺตํ กุสลํ ธมฺมํ ปฏิจฺจ นทุกฺขาย เวทนาย สมฺปยุตฺโต นกุสโล จ น-อทุกฺขมสุขาย เวทนาย สมฺปยุตฺโต นกุสโล จ ธมฺมา อุปฺปชฺชนฺติ เหตุปจฺจยา.\csromanspacer{}\csromanspacer{}สุขาย เวทนาย สมฺปยุตฺตํ กุสลํ ธมฺมํ ปฏิจฺจ นสุขาย เวทนาย สมฺปยุตฺโต นกุสโล จ นทุกฺขาย เวทนาย สมฺปยุตฺโต นกุสโล จ ธมฺมา อุปฺปชฺชนฺติ เหตุปจฺจยา.\csromanspacer{}\csromanspacer{}สุขาย เวทนาย สมฺปยุตฺตํ กุสลํ ธมฺมํ ปฏิจฺจ นสุขาย เวทนาย สมฺปยุตฺโต นกุสโล จ นทุกฺขาย เวทนาย สมฺปยุตฺโต นกุสโล จ น-อทุกฺขมสุขาย เวทนาย สมฺปยุตฺโต นกุสโล จ ธมฺมา อุปฺปชฺชนฺติ เหตุปจฺจยา.\csromanspacer{} สตฺต.}
\prose{\csromanfolio{300}อทุกฺขมสุขาย เวทนาย สมฺปยุตฺตํ กุสลํ ธมฺมํ ปฏิจฺจ น-อทุกฺขมสุขาย เวทนาย สมฺปยุตฺโต นกุสโล ธมฺโม อุปฺปชฺชติ เหตุปจฺจยา.\csromanspacer{} สตฺต.\csromanspacer{} เหตุยา จุทฺทส. สุขาย เวทนาย สมฺปยุตฺตํ อกุสลํ ธมฺมํ ปฏิจฺจ นสุขาย เวทนาย สมฺปยุตฺโต น-อกุสโล ธมฺโม อุปฺปชฺชติ เหตุปจฺจยา.\csromanspacer{} สตฺต.}

\proseitem{37}{ทุกฺขาย เวทนาย สมฺปยุตฺตํ อกุสลํ ธมฺมํ ปฏิจฺจ นทุกฺขาย เวทนาย สมฺปยุตฺโต น-อกุสโล ธมฺโม อุปฺปชฺชติ เหตุปจฺจยา.\csromanspacer{} สตฺต.}

\prose{อทุกฺขมสุขาย เวทนาย สมฺปยุตฺตํ อกุสลํ ธมฺมํ ปฏิจฺจ น-อทุกฺขมสุขาย เวทนาย สมฺปยุตฺโต น-อกุสโล ธมฺโม อุปฺปชฺชติ เหตุปจฺจยา.\csromanspacer{} สตฺต. (เอกวีสติ ปญฺหา.)}

\proseitem{38}{สุขาย เวทนาย สมฺปยุตฺโต อพฺยากโต ธมฺโม นสุขาย เวทนาย สมฺปยุตฺตสฺส น-อพฺยากตสฺส ธมฺมสฺส อารมฺมณปจฺจเยน ปจฺจโย. (สํขิตฺตํ.)}

\prose{อารมฺมเณ เอกวีสติ. (ทุกฺขาย เวทนาย สมฺปยุตฺต-อพฺยากตมูลํ อทุกฺขมสุขาย เวทนาย สมฺปยุตฺต-อพฺยากตมูลมฺปิ กาตพฺพํ.)}

\csromanheader{2}{3. วิปากตฺติก 1. กุสลตฺติก}
\proseitem{39}{วิปากธมฺมธมฺมํ กุสลํ ธมฺมํ ปฏิจฺจ นวิปากธมฺมธมฺโม นกุสโล ธมฺโม อุปฺปชฺชติ เหตุปจฺจยา.\csromanspacer{}\csromanspacer{}เหตุยา ตีณิ.}

\csromanheader{2}{6. วิตกฺกตฺติก 1. กุสลตฺติก}
\prose{\csromanfolio{301}วิปากธมฺมธมฺมํ อกุสลํ ธมฺมํ ปฏิจฺจ นวิปากธมฺมธมฺโม น-อกุสโล ธมฺโม อุปฺปชฺชติ เหตุปจฺจยา.\csromanspacer{}\csromanspacer{}เหตุยา ตีณิ.}

\prose{เนววิปากนวิปากธมฺมธมฺมํ อพฺยากตํ ธมฺมํ ปจฺจยา นเนววิปากนวิปากธมฺมธมฺโม น-อพฺยากโต ธมฺโม อุปฺปชฺชติ เหตุปจฺจยา.\csromanspacer{}\csromanspacer{}เหตุยา ตีณิ.}

\csromanheader{2}{4. อุปาทินฺนตฺติก 1. กุสลตฺติก}
\proseitem{40}{อนุปาทินฺนุปาทานิยํ กุสลํ ธมฺมํ ปฏิจฺจ น-อุปาทินฺนุปาทานิโย นกุสโล ธมฺโม อุปฺปชฺชติ เหตุปจฺจยา.\csromanspacer{} ตีณิ.}

\prose{อนุปาทินฺน-อนุปาทานิยํ กุสลํ ธมฺมํ ปฏิจฺจ น-อนุปาทินฺน-อนุปาทานิโย นกุสโล ธมฺโม อุปฺปชฺชติ เหตุปจฺจยา.\csromanspacer{} ตีณิ.\csromanspacer{} เหตุยา ฉ.}

\prose{อนุปาทินฺนุปาทานิยํ อกุสลํ ธมฺมํ ปฏิจฺจ น-อุปาทินฺนุปาทานิโย น-อกุสโล ธมฺโม อุปฺปชฺชติ เหตุปจฺจยา.\csromanspacer{}\csromanspacer{}เหตุยา ตีณิ.}

\prose{อุปาทินฺนุปาทานิยํ อพฺยากตํ ธมฺมํ ปจฺจยา น-อุปาทินฺนุปาทานิโย น-อพฺยากโต ธมฺโม อุปฺปชฺชติ เหตุปจฺจยา.\csromanspacer{}\csromanspacer{}เหตุยา ปญฺจ.}

\csromanheader{2}{5. สํกิลิฏฺฐตฺติก 1. กุสลตฺติก}
\proseitem{41}{อสํกิลิฏฺฐสํกิเลสิกํ กุสลํ ธมฺมํ ปฏิจฺจ นสํกิลิฏฺฐสํกิเลสิโก นกุสโล ธมฺโม อุปฺปชฺชติ เหตุปจฺจยา.\csromanspacer{} ตีณิ.}

\prose{อสํกิลิฏฺฐ-อสํกิเลสิกํ กุสลํ ธมฺมํ ปฏิจฺจ น-อสํกิลิฏฺฐ-อสํกิเลสิโก นกุสโล ธมฺโม อุปฺปชฺชติ เหตุปจฺจยา.\csromanspacer{} ตีณิ.\csromanspacer{} เหตุยา ฉ.}

\prose{สํกิลิฏฺฐสํกิเลสิกํ อกุสลํ ธมฺมํ ปฏิจฺจ นสํกิลิฏฺฐสํกิเลสิโก น-อกุสโล ธมฺโม อุปฺปชฺชติ เหตุปจฺจยา.\csromanspacer{} ตีณิ.}

\prose{อสํกิลิฏฺฐสํกิเลสิกํ อพฺยากตํ ธมฺมํ ปจฺจยา น-อสํกิลิฏฺฐสํกิเลสิโก น-อพฺยากโต ธมฺโม อุปฺปชฺชติ เหตุปจฺจยา. (สํขิตฺตํ.) .\csromanspacer{} เหตุยา ฉ ฯเปฯ อวิคเต ฉ.}

\csromanheader{2}{6. วิตกฺกตฺติก 1. กุสลตฺติก}
\proseitem{42}{สวิตกฺกสวิจารํ กุสลํ ธมฺมํ ปฏิจฺจ นสวิตกฺกสวิจาโร นกุสโล ธมฺโม อุปฺปชฺชติ เหตุปจฺจยา.\csromanspacer{} ตีณิ.}

\csromanheader{2}{ธมฺมานุโลมปจฺจนีย ติกติกปฏฺฐานปาฬิ}
\prose{\csromanfolio{302}อวิตกฺกวิจารมตฺตํ กุสลํ ธมฺมํ ปฏิจฺจ น-อวิตกฺกวิจารมตฺโต นกุสโล ธมฺโม อุปฺปชฺชติ เหตุปจฺจยา.\csromanspacer{}\csromanspacer{}เหตุยา ปนฺนรส.}

\prose{สวิตกฺกสวิจารํ อกุสลํ ธมฺมํ ปฏิจฺจ นสวิตกฺกสวิจาโร น-อกุสโล ธมฺโม อุปฺปชฺชติ เหตุปจฺจยา.\csromanspacer{}\csromanspacer{}เหตุยา นว.}

\csromanheader{2}{7. ปีติตฺติก 1. กุสลตฺติก}
\proseitem{43}{ปีติสหคตํ กุสลํ ธมฺมํ ปฏิจฺจ นปีติสหคโต นกุสโล ธมฺโม อุปฺปชฺชติ เหตุปจฺจยา.\csromanspacer{} สตฺต.}

\prose{สุขสหคตํ กุสลํ ธมฺมํ ปฏิจฺจ นสุขสหคโต นกุสโล ธมฺโม อุปฺปชฺชติ เหตุปจฺจยา.\csromanspacer{} สตฺต.}

\prose{อุเปกฺขาสหคตํ กุสลํ ธมฺมํ ปฏิจฺจ น-อุเปกฺขาสหคโต นกุสโล ธมฺโม อุปฺปชฺชติ เหตุปจฺจยา.\csromanspacer{} สตฺต.\csromanspacer{}\csromanspacer{}(เวทนาตฺติกสทิสํ.\csromanspacer{} สํขิตฺตํ.)}

\csromanheader{2}{22. สนิทสฺสนตฺติก 1. กุสลตฺติก}
\proseitem{44}{อนิทสฺสน-อปฺปฏิฆํ กุสลํ ธมฺมํ ปฏิจฺจ น-อนิทสฺสน-อปฺปฏิโฆ นกุสโล ธมฺโม อุปฺปชฺชติ เหตุปจฺจยา.\csromanspacer{}\csromanspacer{}อนิทสฺสน-อปฺปฏิฆํ กุสลํ ธมฺมํ ปฏิจฺจ นสนิทสฺสนสปฺปฏิโฆ นกุสโล ธมฺโม อุปฺปชฺชติ เหตุปจฺจยา ฯเปฯ.\csromanspacer{}\csromanspacer{}อนิทสฺสน-อปฺปฏิฆํ กุสลํ ธมฺมํ ปฏิจฺจ นสนิทสฺสนสปฺปฏิโฆ นกุสโล จ น-อนิทสฺสนสปฺปฏิโฆ นกุสโล จ ธมฺมา อุปฺปชฺชนฺติ เหตุปจฺจยา.\csromanspacer{}\csromanspacer{}เหตุยา ฉ.}

\prose{อนิทสฺสน-อปฺปฏิฆํ อกุสลํ ธมฺมํ ปฏิจฺจ น-อนิทสฺสน-อปฺปฏิโฆ น-อกุสโล ธมฺโม อุปฺปชฺชติ เหตุปจฺจยา.\csromanspacer{}\csromanspacer{}อนิทสฺสน-อปฺปฏิฆํ อกุสลํ ธมฺมํ ปฏิจฺจ นสนิทสฺสนสปฺปฏิโฆ นกุสโล ธมฺโม อุปฺปชฺชติ เหตุปจฺจยา ฯเปฯ.\csromanspacer{}\csromanspacer{}อนิทสฺสน-อปฺปฏิฆํ อกุสลํ ธมฺมํ ปฏิจฺจ นสนิทสฺสนสปฺปฏิโฆ นกุสโล จ น-อนิทสฺสนสปฺปฏิโฆ น-อกุสโล จ ธมฺมา อุปฺปชฺชนฺติ เหตุปจฺจยา. (สํขิตฺตํ.) .\csromanspacer{} เหตุยา ฉ ฯเปฯ อวิคเต ฉ.}

\prose{อนิทสฺสน-อปฺปฏิฆํ อพฺยากตํ ธมฺมํ ปจฺจยา นสนิทสฺสนสปฺปฏิโฆ น-อพฺยากโต ธมฺโม อุปฺปชฺชติ เหตุปจฺจยา.\csromanspacer{}\csromanspacer{}อนิทสฺสน-อปฺปฏิฆํ อพฺยากตํ ธมฺมํ ปจฺจยา น-อนิทสฺสนสปฺปฏิโฆ น-อพฺยากโต ธมฺโม อุปฺปชฺชติ}

\vspace{\tipitakaparskip}
\csromancenter{สนิทสฺสนตฺติก 4.\csromanspacer{} อุปาทินฺนตฺติก เหตุปจฺจยา.\csromanspacer{}\csromanspacer{}อนิทสฺสน-อปฺปฏิฆํ อพฺยากตํ ธมฺมํ ปจฺจยา นสนิทสฺสนสปฺปฏิโฆ น-อพฺยากโต จ น-อนิทสฺสนสปฺปฏิโฆ น-อพฺยากโต จ ธมฺมา อุปฺปชฺชนฺติ เหตุปจฺจยา.\csromanspacer{}\csromanspacer{}เหตุยา ตีณิ ฯเปฯ อวิคเต ตีณิ.}
\csromanheader{2}{22. สนิทสฺสนตฺติก 2. เวทนาตฺติก}
\proseitem{45}{\csromanfolio{303}อนิทสฺสน-อปฺปฏิฆํ สุขาย เวทนาย สมฺปยุตฺตํ ธมฺมํ ปฏิจฺจ น-อนิทสฺสน-อปฺปฏิโฆ นสุขาย เวทนาย สมฺปยุตฺโต ธมฺโม อุปฺปชฺชติ เหตุปจฺจยา.\csromanspacer{}\csromanspacer{}เหตุยา ฉ.}

\prose{อนิทสฺสน-อปฺปฏิฆํ ทุกฺขาย เวทนาย สมฺปยุตฺตํ ธมฺมํ ปฏิจฺจ น-อนิทสฺสน-อปฺปฏิโฆ นทุกฺขาย เวทนาย สมฺปยุตฺโต ธมฺโม อุปฺปชฺชติ เหตุปจฺจยา.\csromanspacer{}\csromanspacer{}เหตุยา ฉ.}

\prose{อนิทสฺสน-อปฺปฏิฆํ อทุกฺขมสุขาย เวทนาย สมฺปยุตฺตํ ธมฺมํ ปฏิจฺจ น-อนิทสฺสน-อปฺปฏิโฆ น-อทุกฺขมสุขาย เวทนาย สมฺปยุตฺโต ธมฺโม อุปฺปชฺชติ เหตุปจฺจยา.\csromanspacer{}\csromanspacer{}เหตุยา ฉ.}

\csromanheader{2}{22. สนิทสฺสนตฺติก 3. วิปากตฺติก}
\proseitem{46}{อนิทสฺสน-อปฺปฏิฆํ วิปากํ ธมฺมํ ปฏิจฺจ น-อนิทสฺสน-อปฺปฏิโฆ นวิปาโก ธมฺโม อุปฺปชฺชติ เหตุปจฺจยา.\csromanspacer{}\csromanspacer{}เหตุยา ฉ.}

\prose{อนิทสฺสน-อปฺปฏิฆํ วิปากธมฺมธมฺมํ ปฏิจฺจ น-อนิทสฺสน-อปฺปฏิโฆ นวิปากธมฺมธมฺโม อุปฺปชฺชติ เหตุปจฺจยา.\csromanspacer{}\csromanspacer{}เหตุยา ฉ.}

\prose{อนิทสฺสน-อปฺปฏิฆํ เนววิปากนวิปากธมฺมธมฺมํ ปฏิจฺจ นสนิทสฺสนสปฺปฏิโฆ นเนววิปากนวิปากธมฺมธมฺโม อุปฺปชฺชติ เหตุปจฺจยา.\csromanspacer{}\csromanspacer{}เหตุยา ตีณิ.}

\csromanheader{2}{22. สนิทสฺสนตฺติก 4. อุปาทินฺนตฺติก}
\proseitem{47}{อนิทสฺสน-อปฺปฏิฆํ อุปาทินฺนุปาทานิยํ ธมฺมํ ปฏิจฺจ น-อนิทสฺสน-อปฺปฏิโฆ น-อุปาทินฺนุปาทานิโย ธมฺโม อุปฺปชฺชติ เหตุปจฺจยา.\csromanspacer{}\csromanspacer{}เหตุยา ฉ.}

\prose{สนิทสฺสนสปฺปฏิโฆ อนุปาทินฺนุปาทานิโย ธมฺโม นสนิทสฺสนสปฺปฏิฆสฺส น-อนุปาทินฺนุปาทานิยสฺส ธมฺมสฺส อารมฺมณปจฺจเยน ปจฺจโย.}

\csromangathameasure{ธมฺมานุโลมปจฺจนีย ติกติกปฏฺฐานปาฬิ (สํขิตฺตํ.) . อารมฺมเณ นว, อนนฺตเร ตีณิ ฯเปฯ อุปนิสฺสเย ปุเรชาเต นว ฯเปฯ อวิคเต นว.}
\csromangathagroup{\csromangathastanza{\csromanfolio{304}ธมฺมานุโลมปจฺจนีย ติกติกปฏฺฐานปาฬิ (สํขิตฺตํ.) .\csromanspacer{} อารมฺมเณ นว, อนนฺตเร ตีณิ ฯเปฯ อุปนิสฺสเย ปุเรชาเต นว ฯเปฯ อวิคเต นว.}}

\prose{อนิทสฺสน-อปฺปฏิฆํ อนุปาทินฺน-อนุปาทานิยํ ธมฺมํ ปฏิจฺจ น-อนิทสฺสน-อปฺปฏิโฆ น-อนุปาทินฺน-อนุปาทานิโย ธมฺโม อุปฺปชฺชติ เหตุปจฺจยา.\csromanspacer{}\csromanspacer{}เหตุยา ฉ.}

\csromanheader{2}{22. สนิทสฺสนตฺติก 5. สํกิลิฏฺฐตฺติก}
\proseitem{48}{อนิทสฺสน-อปฺปฏิฆํ สํกิลิฏฺฐสํกิเลสิกํ ธมฺมํ ปฏิจฺจ น-อนิทสฺสน-อปฺปฏิโฆ นสํกิลิฏฺฐสํกิเลสิโก ธมฺโม อุปฺปชฺชติ เหตุปจฺจยา.\csromanspacer{}\csromanspacer{}เหตุยา ฉ.}

\prose{อนิทสฺสน-อปฺปฏิฆํ อสํกิลิฏฺฐสํกิเลสิกํ ธมฺมํ ปจฺจยา นสนิทสฺสนสปฺปฏิโฆ น-อสํกิลิฏฺฐสํกิเลสิโก ธมฺโม อุปฺปชฺชติ เหตุปจฺจยา.\csromanspacer{}\csromanspacer{}เหตุยา ตีณิ.}

\prose{อนิทสฺสน-อปฺปฏิฆํ อสํกิลิฏฺฐ-อสํกิเลสิกํ ธมฺมํ ปฏิจฺจ น-อนิทสฺสน-อปฺปฏิโฆ น-อสํกิลิฏฺฐ-อสํกิเลสิโก ธมฺโม อุปฺปชฺชติ เหตุปจฺจยา.\csromanspacer{}\csromanspacer{}เหตุยา ฉ.}

\csromanheader{2}{22. สนิทสฺสนตฺติก 6. วิตกฺกตฺติก}
\proseitem{49}{อนิทสฺสน-อปฺปฏิฆํ สวิตกฺกสวิจารํ ธมฺมํ ปฏิจฺจ น-อนิทสฺสน-อปฺปฏิโฆ นสวิตกฺกสวิจาโร ธมฺโม อุปฺปชฺชติ เหตุปจฺจยา.\csromanspacer{}\csromanspacer{}เหตุยา ฉ.}

\prose{อนิทสฺสน-อปฺปฏิฆํ อวิตกฺกวิจารมตฺตํ ธมฺมํ ปฏิจฺจ น-อนิทสฺสน-อปฺปฏิโฆ น-อวิตกฺกวิจารมตฺโต ธมฺโม อุปฺปชฺชติ เหตุปจฺจยา.\csromanspacer{}\csromanspacer{}เหตุยา ฉ.}

\prose{อนิทสฺสน-อปฺปฏิฆํ อวิตกฺก-อวิจารํ ธมฺมํ ปฏิจฺจ นสนิทสฺสนสปฺปฏิโฆ น-อวิตกฺก-อวิจาโร ธมฺโม อุปฺปชฺชติ เหตุปจฺจยา.\csromanspacer{}\csromanspacer{}เหตุยา ตีณิ.}

\csromanheader{2}{22. สนิทสฺสนตฺติก 21. อชฺฌตฺตารมฺมณตฺติก}
\proseitem{50}{อนิทสฺสน-อปฺปฏิฆํ อชฺฌตฺตารมฺมณํ ธมฺมํ ปฏิจฺจ น-อนิทสฺสน-อปฺปฏิโฆ น-อชฺฌตฺตารมฺมโณ ธมฺโม อุปฺปชฺชติ เหตุปจฺจยา.\csromanspacer{}\csromanspacer{}เหตุยา ฉ.}

\csromanheader{2}{22. สนิทสฺสนตฺติก 21. อชฺฌตฺตารมฺมณตฺติก}
\prose{\csromanfolio{305}อนิทสฺสน-อปฺปฏิฆํ พหิทฺธารมฺมณํ ธมฺมํ ปฏิจฺจ น-อนิทสฺสน-อปฺปฏิโฆ นพหิทฺธารมฺมโณ ธมฺโม อุปฺปชฺชติ เหตุปจฺจยา.\csromanspacer{}\csromanspacer{}อนิทสฺสน-อปฺปฏิฆํ พหิทฺธารมฺมณํ ธมฺมํ ปฏิจฺจ นสนิทสฺสนสปฺปฏิโฆ นพหิทฺธารมฺมโณ ธมฺโม อุปฺปชฺชติ เหตุปจฺจยา.\csromanspacer{}\csromanspacer{}อนิทสฺสน-อปฺปฏิฆํ พหิทฺธารมฺมณํ ธมฺมํ ปฏิจฺจ น-อนิทสฺสนสปฺปฏิโฆ นพหิทฺธารมฺมโณ ธมฺโม อุปฺปชฺชติ เหตุปจฺจยา ฯเปฯ.\csromanspacer{}\csromanspacer{}อนิทสฺสน-อปฺปฏิฆํ พหิทฺธารมฺมณํ ธมฺมํ ปฏิจฺจ นสนิทสฺสนสปฺปฏิโฆ นพหิทฺธารมฺมโณ จ น-อนิทสฺสนสปฺปฏิโฆ นพหิทฺธารมฺมโณ จ ธมฺมา อุปฺปชฺชนฺติ เหตุปจฺจยา.\csromanspacer{}\csromanspacer{}เหตุยา ฉ ฯเปฯ อวิคเต ฉ. (ปญฺหาวารํ วิตฺถาเรตพฺพํ.)}

\nitthitam{\textbf{ธมฺมานุโลมปจฺจนีเย ติกติกปฏฺฐานํ นิฏฺฐิตํ.}}
\pitaka{\textbf{อภิธมฺมปิฏก}}
\csromanheader{2}{\textbf{ธมฺมานุโลมปจฺจนีย}}
\gambhira{\textbf{ทุกทุกปฏฺฐานปาฬิ}}
\namakkaram{นโม ตสฺส ภควโต อรหโต สมฺมาสมฺพุทฺธสฺส.}
\csromanheader{2}{1. เหตุทุก 1. สเหตุกทุก}
\proseitem{1}{\csromanfolio{307}เหตุํ สเหตุกํ ธมฺมํ ปฏิจฺจ นเหตุ นสเหตุโก ธมฺโม อุปฺปชฺชติ เหตุปจฺจยา.\csromanspacer{}\csromanspacer{}นเหตุํ สเหตุกํ ธมฺมํ ปฏิจฺจ นเหตุ นสเหตุโก ธมฺโม อุปฺปชฺชติ เหตุปจฺจยา.\csromanspacer{}\csromanspacer{}เหตุํ สเหตุกญฺจ นเหตุํ สเหตุกญฺจ ธมฺมํ ปฏิจฺจ นเหตุ นสเหตุโก ธมฺโม อุปฺปชฺชติ เหตุปจฺจยา. (สํขิตฺตํ.) .\csromanspacer{} เหตุยา ตีณิ, อารมฺมเณ เอกํ ฯเปฯ อวิคเต ปญฺจ.}

\proseitem{2}{นเหตุํ สเหตุกํ ธมฺมํ ปฏิจฺจ นนเหตุ นสเหตุโก ธมฺโม อุปฺปชฺชติ นเหตุปจฺจยา. (สํขิตฺตํ.) .\csromanspacer{} นเหตุยา เอกํ, น-อารมฺมเณ ตีณิ ฯเปฯ โนวิคเต ตีณิ.}

\prose{(สหชาตวารมฺปิ ฯเปฯ นมฺปยุตฺตวารมฺปิ ปฏิจฺจวารสทิสํ.)}

\csromanheader{2}{\textbf{เหตุ-อารมฺมณ}}
\csromanheader{2}{3. เหตุ สเหตุโก ธมฺโม นเหตุสฺส นสเหตุกสฺส ธมฺมสฺส}
\prose{เหตุ สเหตุโก ธมฺโม นเหตุสฺส นสเหตุกสฺส ธมฺมสฺส อารมฺมณปจฺจเยน ปจฺจโย.\csromanspacer{}\csromanspacer{}เหตุ สเหตุโก ธมฺโม นนเหตุสฺส}

\vspace{\tipitakaparskip}
\csromancenter{ธมฺมานุโลมปจฺจนีย ทุกทุกปฏฺฐานปาฬิ นสเหตุกสฺส ธมฺมสฺส อารมฺมณปจฺจเยน ปจฺจโย.\csromanspacer{}\csromanspacer{}เหตุ สเหตุโก ธมฺโม นเหตุสฺส นสเหตุกสฺส จ นนเหตุสฺส นสเหตุกสฺส จ ธมฺมสฺส อารมฺมณปจฺจเยน ปจฺจโย. (3)}
\prose{\csromanfolio{308}นเหตุ สเหตุโก ธมฺโม นนเหตุสฺส นสเหตุกสฺส ธมฺมสฺส อารมฺมณปจฺจเยน ปจฺจโย.\csromanspacer{}\csromanspacer{}นเหตุ สเหตุโก ธมฺโม นเหตุสฺส นสเหตุกสฺส ธมฺมสฺส อารมฺมณปจฺจเยน ปจฺจโย.\csromanspacer{}\csromanspacer{}นเหตุ สเหตุโก ธมฺโม นเหตุสฺส นสเหตุกสฺส จ นนเหตุสฺส นสเหตุกสฺส จ ธมฺมสฺส อารมฺมณปจฺจเยน ปจฺจโย. (3) (สํขิตฺตํ.)}

\prose{เหตุยา เอกํ, อารมฺมเณ ฉ, อธิปติยา เทฺว ฯเปฯ อวิคเต ปญฺจ. (ปญฺหาวารํ วิตฺถาเรตพฺพํ.)}

\proseitem{4}{เหตุํ อเหตุกํ ธมฺมํ ปฏิจฺจ นเหตุ น-อเหตุโก ธมฺโม อุปฺปชฺชติ เหตุปจฺจยา. (1)}

\prose{นเหตุํ อเหตุกํ ธมฺมํ ปฏิจฺจ นนเหตุ น-อเหตุโก ธมฺโม อุปฺปชฺชติ เหตุปจฺจยา.\csromanspacer{} ตีณิ. (สํขิตฺตํ.) .\csromanspacer{} เหตุยา จตฺตาริ.}

\csromanheader{2}{1. เหตุทุก 2. เหตุสมฺปยุตฺตทุก}
\proseitem{5}{เหตุํ เหตุสมฺปยุตฺตํ ธมฺมํ ปฏิจฺจ นเหตุ นเหตุสมฺปยุตฺโต ธมฺโม อุปฺปชฺชติ เหตุปจฺจยา.\csromanspacer{} ตีณิ. (สเหตุกสทิสํ.)}

\prose{เหตุํ เหตุวิปฺปยุตฺตํ ธมฺมํ ปฏิจฺจ นเหตุ นเหตุวิปฺปยุตฺโต}

\prose{นเหตุํ เหตุวิปฺปยุตฺตํ ธมฺมํ ปฏิจฺจ นนเหตุ นเหตุวิปฺปยุตฺโต ธมฺโม อุปฺปชฺชติ เหตุปจฺจยา.\csromanspacer{} ตีณิ.\csromanspacer{}\csromanspacer{}เหตุยา จตฺตาริ.}

\csromanheader{2}{1. เหตุทุก 3-5. เหตุสเหตุกทุกาทิ}
\proseitem{6}{เหตุํ เหตุญฺเจว สเหตุกญฺจ ธมฺมํ ปฏิจฺจ นเหตุ นเหตุ เจว น-อเหตุโก จ ธมฺโม อุปฺปชฺชติ เหตุปจฺจยา.\csromanspacer{}\csromanspacer{}เหตุยา เอกํ.}

\prose{นเหตุํ สเหตุกญฺเจว น จ เหตุํ ธมฺมํ ปฏิจฺจ นนเหตุ น-อเหตุโก เจว นนเหตุ จ ธมฺโม อุปฺปชฺชติ เหตุปจฺจยา.\csromanspacer{}\csromanspacer{}เหตุยา เอกํ.}

\csromanheader{2}{1. เหตุทุก 6-12. จูฬนฺตรทุก}
\proseitem{7}{\csromanfolio{309}เหตุํ เหตุญฺเจว เหตุสมฺปยุตฺตญฺจ ธมฺมํ ปฏิจฺจ นเหตุ นเหตุ เจว นเหตุวิปฺปยุตฺโต จ ธมฺโม อุปฺปชฺชติ เหตุปจฺจยา.\csromanspacer{}\csromanspacer{}เหตุยา เอกํ.}

\prose{นเหตุํ เหตุสมฺปยุตฺตญฺเจว น จ เหตุํ ธมฺมํ ปฏิจฺจ นนเหตุ นเหตุ วิปฺปยุตฺโต เจว นนเหตุ จ ธมฺโม อุปฺปชฺชติ เหตุปจฺจยา.\csromanspacer{}\csromanspacer{}เหตุยา เอกํ.}

\proseitem{8}{นเหตุํ นเหตุํ สเหตุกํ ธมฺมํ ปฏิจฺจ นเหตุ นเหตุ นสเหตุโก ธมฺโม อุปฺปชฺชติ เหตุปจฺจยา.\csromanspacer{}\csromanspacer{}เหตุยา อธิปติยา เอกํ.}

\prose{นเหตุํ นเหตุํ อเหตุกํ ธมฺมํ ปฏิจฺจ นเหตุ นเหตุ น-อเหตุโก ธมฺโม อุปฺปชฺชติ เหตุปจฺจยา.\csromanspacer{}\csromanspacer{}เหตุยา เอกํ.}

\csromanheader{2}{1. เหตุทุก 6-12. จูฬนฺตรทุก}
\proseitem{9}{นเหตุ อปฺปจฺจโย ธมฺโม นนเหตุสฺส น-อปฺปจฺจยสฺส ธมฺมสฺส อารมฺมณปจฺจเยน ปจฺจโย ฯเปฯ.\csromanspacer{}\csromanspacer{}นเหตุ อปฺปจฺจโย ธมฺโม นเหตุสฺส น-อปฺปจฺจยสฺส ธมฺมสฺส อารมฺมณปจฺจเยน ปจฺจโย.\csromanspacer{}\csromanspacer{}นเหตุ อปฺปจฺจโย ธมฺโม นเหตุสฺส น-อปฺปจฺจยสฺส จ นนเหตุสฺส น-อปฺปจฺจยสฺส จ ธมฺมสฺส อารมฺมณปจฺจเยน ปจฺจโย ฯเปฯ. (อสงฺขตํ อปฺปจฺจยสทิสํ.)}

\proseitem{10}{นเหตุ สนิทสฺสโน ธมฺโม นนเหตุสฺส นสนิทสฺสนสฺส ธมฺมสฺส อารมฺมณปจฺจเยน ปจฺจโย.\csromanspacer{}\csromanspacer{}นเหตุ สนิทสฺสโน ธมฺโม นเหตุสฺส นสนิทสฺสนสฺส ธมฺมสฺส อารมฺมณปจฺจเยน ปจฺจโย.\csromanspacer{}\csromanspacer{}นเหตุ สนิทสฺสโน ธมฺโม นเหตุสฺส นสนิทสฺสนสฺส จ นนเหตุสฺส นสนิทสฺสนสฺส จ ธมฺมสฺส อารมฺมณปจฺจเยน ปจฺจโย.\csromanspacer{}\csromanspacer{}อารมฺมเณ ตีณิ.}

\prose{เหตุํ อนิทสฺสนํ ธมฺมํ ปฏิจฺจ นเหตุ น-อนิทสฺสโน ธมฺโม อุปฺปชฺชติ เหตุปจฺจยา.\csromanspacer{}\csromanspacer{}เหตุยา ตีณิ.}

\proseitem{11}{นเหตุํ สปฺปฏิฆํ ธมฺมํ ปฏิจฺจ นเหตุ นสปฺปฏิโฆ ธมฺโม อุปฺปชฺชติ เหตุปจฺจยา.\csromanspacer{}\csromanspacer{}เหตุยา เอกํ.}

\prose{เหตุํ อปฺปฏิฆํ ธมฺมํ ปฏิจฺจ นเหตุ น-อปฺปฏิโฆ ธมฺโม อุปฺปชฺชติ เหตุปจฺจยา.\csromanspacer{}\csromanspacer{}เหตุยา ตีณิ.}

\csromanheader{2}{ธมฺมานุโลมปจฺจนีย ทุกทุกปฏฺฐานปาฬิ}
\proseitem{12}{\csromanfolio{310}นเหตุํ รูปิํ ธมฺมํ ปฏิจฺจ นนเหตุ นรูปี ธมฺโม อุปฺปชฺชติ เหตุปจฺจยา.\csromanspacer{}\csromanspacer{}เหตุยา ตีณิ.}

\prose{เหตุํ อรูปิํ ธมฺมํ ปฏิจฺจ นเหตุ น-อรูปี ธมฺโม อุปฺปชฺชติ เหตุปจฺจยา.\csromanspacer{}\csromanspacer{}เหตุยา ตีณิ.}

\proseitem{13}{นเหตุํ โลกิยํ ธมฺมํ ปจฺจยา นนเหตุ นโลกิโย ธมฺโม อุปฺปชฺชติ เหตุปจฺจยา.\csromanspacer{} เหตุยา ตีณิ.}

\prose{เหตุํ โลกุตฺตรํ ธมฺมํ ปฏิจฺจ นเหตุ นโลกุตฺตโร ธมฺโม อุปฺปชฺชติ เหตุปจฺจยา.\csromanspacer{}\csromanspacer{}นเหตุํ โลกุตฺตรํ ธมฺมํ ปฏิจฺจ นเหตุ นโลกุตฺตโร ธมฺโม อุปฺปชฺชติ เหตุปจฺจยา.\csromanspacer{}\csromanspacer{}เหตุํ โลกุตฺตรญฺจ นเหตุํ โลกุตฺตรญฺจ ธมฺมํ ปฏิจฺจ นเหตุ นโลกุตฺตโร ธมฺโม อุปฺปชฺชติ เหตุปจฺจยา.\csromanspacer{}\csromanspacer{}เหตุยา ตีณิ.}

\proseitem{14}{เหตุํ เกนจิ วิญฺเญยฺยํ ธมฺมํ ปฏิจฺจ นเหตุ นเกนจิ วิญฺเญยฺโย ธมฺโม อุปฺปชฺชติ เหตุปจฺจยา.\csromanspacer{}\csromanspacer{}เหตุยา นว.}

\prose{เหตุํ นเกนจิ วิญฺเญยฺยํ ธมฺมํ ปฏิจฺจ นเหตุ นนเกนจิ วิญฺเญยฺโย ธมฺโม อุปฺปชฺชติ เหตุปจฺจยา.\csromanspacer{}\csromanspacer{}เหตุยา นว.}

\csromanheader{2}{1. เหตุทุก 13. อาสวโคจฺฉก}
\proseitem{15}{เหตุํ อาสวํ ธมฺมํ ปฏิจฺจ นเหตุ น-อาสโว ธมฺโม อุปฺปชฺชติ เหตุปจฺจยา.\csromanspacer{}\csromanspacer{}เหตุยา ปญฺจ.}

\prose{เหตุํ โน-อาสวํ ธมฺมํ ปฏิจฺจ นนเหตุ นโน-อาสโว ธมฺโม อุปฺปชฺชติ เหตุปจฺจยา.\csromanspacer{}\csromanspacer{}เหตุยา ปญฺจ.}

\proseitem{16}{นเหตุํ สาสวํ ธมฺมํ ปจฺจยา นนเหตุ นสาสโว ธมฺโม อุปฺปชฺชติ เหตุปจฺจยา.\csromanspacer{}\csromanspacer{}เหตุยา ตีณิ.}

\prose{เหตุํ อนาสวํ ธมฺมํ ปฏิจฺจ นเหตุ น-อนาสโว ธมฺโม อุปฺปชฺชติ เหตุปจฺจยา.\csromanspacer{}\csromanspacer{}เหตุยา ตีณิ.}

\proseitem{17}{เหตุํ อาสวสมฺปยุตฺตํ ธมฺมํ ปฏิจฺจ นเหตุ น-อาสวสมฺปยุตฺโต ธมฺโม อุปฺปชฺชติ เหตุปจฺจยา.\csromanspacer{}\csromanspacer{}เหตุยา นว.}

\csromanheader{2}{1. เหตุทุก 19. สญฺโญชนาทิโคจฺฉก}
\prose{\csromanfolio{311}เหตุํ อาสววิปฺปยุตฺตํ ธมฺมํ ปฏิจฺจ นเหตุ น-อาสววิปฺปยุตฺโต ธมฺโม อุปฺปชฺชติ เหตุปจฺจยา.\csromanspacer{}\csromanspacer{}เหตุยา ตีณิ.}

\proseitem{18}{เหตุํ อาสวญฺเจว สาสวญฺจ ธมฺมํ ปฏิจฺจ นเหตุ น-อาสโว เจว น-อนาสโว จ ธมฺโม อุปฺปชฺชติ เหตุปจฺจยา.\csromanspacer{}\csromanspacer{}เหตุยา ปญฺจ.}

\prose{เหตุํ สาสวญฺเจว โน จ อาสวํ ธมฺมํ ปฏิจฺจ นนเหตุ น-อนาสโว เจว นโน จ อาสโว ธมฺโม อุปฺปชฺชติ เหตุปจฺจยา.\csromanspacer{}\csromanspacer{}เหตุยา ปญฺจ.}

\proseitem{19}{เหตุํ อาสวญฺเจว อาสวสมฺปยุตฺตญฺจ ธมฺมํ ปฏิจฺจ นเหตุ น-อาสโว เจว น-อาสววิปฺปยุตฺโต จ ธมฺโม อุปฺปชฺชติ เหตุปจฺจยา.\csromanspacer{}\csromanspacer{}เหตุยา ตีณิ.}

\prose{นเหตุํ อาสวสมฺปยุตฺตญฺเจว โน จ อาสวํ ธมฺมํ ปฏิจฺจ นนเหตุ น-อาสววิปฺปยุตฺโต เจว นโน-อาสโว จ ธมฺโม อุปฺปชฺชติ เหตุปจฺจยา.\csromanspacer{}\csromanspacer{}เหตุยา ตีณิ.}

\proseitem{20}{นเหตุํ อาสววิปฺปยุตฺตํ สาสวํ ธมฺมํ ปจฺจยา นนเหตุ อาสววิปฺปยุตฺโต นสาสโว ธมฺโม อุปฺปชฺชติ เหตุปจฺจยา. (โลกิยสทิสํ.)}

\prose{เหตุํ อาสววิปฺปยุตฺตํ อนาสวํ ธมฺมํ ปฏิจฺจ นเหตุ อาสววิปฺปยุตฺโต น-อนาสโว ธมฺโม อุปฺปชฺชติ เหตุปจฺจยา.\csromanspacer{}\csromanspacer{}นเหตุํ อาสววิปฺปยุตฺตํ อนาสวํ ธมฺมํ ปฏิจฺจ นเหตุ อาสววิปฺปยุตฺโต น-อนาสโว ธมฺโม อุปฺปชฺชติ เหตุปจฺจยา.\csromanspacer{}\csromanspacer{}เหตุํ อาสววิปฺปยุตฺตํ อนาสวญฺจ นเหตุํ อาสววิปฺปยุตฺตํ อนาสวญฺจ ธมฺมํ ปฏิจฺจ นเหตุ อาสววิปฺปยุตฺโต น-อนาสโว ธมฺโม อุปฺปชฺชติ เหตุปจฺจยา.\csromanspacer{}\csromanspacer{}เหตุยา ตีณิ.}

\csromanheader{2}{1. เหตุทุก 19. สญฺโญชนาทิโคจฺฉก}
\proseitem{21}{เหตุํ สญฺโญชนํ ธมฺมํ ปฏิจฺจ นเหตุ โนสญฺโญชโน ธมฺโม อุปฺปชฺชติ เหตุปจฺจยา.\csromanspacer{}\csromanspacer{}เหตุยา ตีณิ.}

\proseitem{22}{เหตุํ คนฺถํ ธมฺมํ ปฏิจฺจ นเหตุ โนคนฺโถ ธมฺโม อุปฺปชฺชติ เหตุปจฺจยา.\csromanspacer{}\csromanspacer{}เหตุยา นว.}

\csromanheader{2}{ธมฺมานุโลมปจฺจนีย ทุกทุกปฏฺฐานปาฬิ}
\proseitem{23}{\csromanfolio{312}เหตุํ โอฆํ ธมฺมํ ปฏิจฺจ นเหตุ โน-โอโฆ ธมฺโม อุปฺปชฺชติ เหตุปจฺจยา.\csromanspacer{}\csromanspacer{}เหตุยา ปญฺจ.}

\proseitem{24}{เหตุํ โยคํ ธมฺมํ ปฏิจฺจ นเหตุ โนโยโค ธมฺโม อุปฺปชฺชติ เหตุปจฺจยา.\csromanspacer{}\csromanspacer{}เหตุยา ปญฺจ.}

\proseitem{25}{เหตุํ นีวรณํ ธมฺมํ ปฏิจฺจ นเหตุ โนนีวรโณ ธมฺโม อุปฺปชฺชติ เหตุปจฺจยา.\csromanspacer{}\csromanspacer{}เหตุยา ตีณิ.}

\proseitem{26}{นเหตุํ ปรามาสํ ธมฺมํ ปฏิจฺจ นนเหตุ โนปรามาโส ธมฺโม อุปฺปชฺชติ เหตุปจฺจยา.\csromanspacer{}\csromanspacer{}เหตุยา ตีณิ.}

\csromanheader{2}{1. เหตุทุก 54. มหนฺตรทุก}
\proseitem{27}{เหตุํ สารมฺมณํ ธมฺมํ ปฏิจฺจ นเหตุ นสารมฺมโณ ธมฺโม อุปฺปชฺชติ เหตุปจฺจยา.\csromanspacer{}\csromanspacer{}เหตุยา ตีณิ.}

\prose{นเหตุํ อนารมฺมณํ ธมฺมํ ปฏิจฺจ นนเหตุ น-อารมฺมโณ ธมฺโม อุปฺปชฺชติ เหตุปจฺจยา.\csromanspacer{}\csromanspacer{}เหตุยา ตีณิ.}

\proseitem{28}{นเหตุํ จิตฺตํ ธมฺมํ ปฏิจฺจ นนเหตุ โนจิตฺโต ธมฺโม อุปฺปชฺชติ เหตุปจฺจยา.\csromanspacer{}\csromanspacer{}เหตุยา ตีณิ.}

\prose{เหตุํ โนจิตฺตํ ธมฺมํ ปฏิจฺจ นเหตุ นโนจิตฺโต ธมฺโม อุปฺปชฺชติ เหตุปจฺจยา.\csromanspacer{}\csromanspacer{}เหตุยา ตีณิ.}

\proseitem{29}{เหตุํ เจตสิกํ ธมฺมํ ปฏิจฺจ นเหตุ นเจตสิโก ธมฺโม อุปฺปชฺชติ เหตุปจฺจยา.\csromanspacer{}\csromanspacer{}เหตุยา ตีณิ.}

\prose{นเหตุํ อเจตสิกํ ธมฺมํ ปฏิจฺจ นนเหตุ น-อเจตสิโก ธมฺโม อุปฺปชฺชติ เหตุปจฺจยา.\csromanspacer{}\csromanspacer{}เหตุยา ตีณิ.}

\proseitem{30}{เหตุํ จิตฺตสมฺปยุตฺตํ ธมฺมํ ปฏิจฺจ นเหตุ นจิตฺตสมฺปยุตฺโต ธมฺโม อุปฺปชฺชติ เหตุปจฺจยา.\csromanspacer{}\csromanspacer{}เหตุยา ตีณิ.}

\prose{นเหตุํ จิตฺตวิปฺปยุตฺตํ ธมฺมํ ปฏิจฺจ นนเหตุ นจิตฺตวิปฺปยุตฺโต ธมฺโม อุปฺปชฺชติ เหตุปจฺจยา.\csromanspacer{}\csromanspacer{}เหตุยา ตีณิ.}

\csromanheader{2}{1. เหตุทุก 54. มหนฺตรทุก}
\proseitem{31}{\csromanfolio{313}เหตุํ จิตฺตสํสฏฺฐํ ธมฺมํ ปฏิจฺจ นเหตุ นจิตฺตสํสฏฺโฐ ธมฺโม อุปฺปชฺชติ เหตุปจฺจยา.\csromanspacer{}\csromanspacer{}เหตุยา ตีณิ.}

\prose{นเหตุํ จิตฺตวิสํสฏฺฐํ ธมฺมํ ปฏิจฺจ นนเหตุ นจิตฺตวิสํสฏฺโฐ ธมฺโม อุปฺปชฺชติ เหตุปจฺจยา.\csromanspacer{}\csromanspacer{}เหตุยา ตีณิ.}

\proseitem{32}{เหตุํ จิตฺตสมุฏฺฐานํ ธมฺมํ ปฏิจฺจ นเหตุ โนจิตฺตสมุฏฺฐาโน ธมฺโม อุปฺปชฺชติ เหตุปจฺจยา.\csromanspacer{}\csromanspacer{}เหตุยา ตีณิ.}

\prose{นเหตุํ โนจิตฺตสมุฏฺฐานํ ธมฺมํ ปฏิจฺจ นนเหตุ นโนจิตฺตสมุฏฺฐาโน ธมฺโม อุปฺปชฺชติ เหตุปจฺจยา.\csromanspacer{}\csromanspacer{}เหตุยา ตีณิ.}

\proseitem{33}{เหตุํ จิตฺตสหภุํ ธมฺมํ ปฏิจฺจ นเหตุ โนจิตฺตสหภู ธมฺโม อุปฺปชฺชติ เหตุปจฺจยา.\csromanspacer{}\csromanspacer{}เหตุยา ตีณิ.}

\prose{นเหตุํ โนจิตฺตสหภุํ ธมฺมํ ปฏิจฺจ นนเหตุ นโนจิตฺตสหภู ธมฺโม อุปฺปชฺชติ เหตุปจฺจยา.\csromanspacer{}\csromanspacer{}เหตุยา ตีณิ.}

\proseitem{34}{เหตุํ จิตฺตานุปริวตฺติํ ธมฺมํ ปฏิจฺจ นเหตุ โนจิตฺตานุปริวตฺตี ธมฺโม อุปฺปชฺชติ เหตุปจฺจยา.\csromanspacer{}\csromanspacer{}เหตุยา ตีณิ.}

\prose{นเหตุํ โนจิตฺตานุปริวตฺติํ ธมฺมํ ปฏิจฺจ นนเหตุ นโนจิตฺตานุปริวตฺตี ธมฺโม อุปฺปชฺชติ เหตุปจฺจยา.\csromanspacer{} เหตุยา ตีณิ.}

\proseitem{35}{เหตุํ จิตฺตสํสฏฺฐสมุฏฺฐานํ ธมฺมํ ปฏิจฺจ นเหตุ โนจิตฺตสํสฏฺฐสมุฏฺฐาโน ธมฺโม อุปฺปชฺชติ เหตุปจฺจยา.\csromanspacer{}\csromanspacer{}เหตุยา ตีณิ.}

\prose{นเหตุํ โนจิตฺตสํสฏฺฐสมุฏฺฐานํ ธมฺมํ ปฏิจฺจ นนเหตุ นโนจิตฺตสํสฏฺฐสมุฏฺฐาโน ธมฺโม อุปฺปชฺชติ เหตุปจฺจยา.\csromanspacer{}\csromanspacer{}เหตุยา ตีณิ.}

\proseitem{36}{เหตุํ จิตฺตสํสฏฺฐสมุฏฺฐานสหภุํ ธมฺมํ ปฏิจฺจ นเหตุ โนจิตฺตสํสฏฺฐสมุฏฺฐานสหภู ธมฺโม อุปฺปชฺชติ เหตุปจฺจยา.\csromanspacer{}\csromanspacer{}เหตุยา ตีณิ.}

\prose{นเหตุํ โนจิตฺตสํสฏฺฐสมุฏฺฐานสหภุํ ธมฺมํ ปฏิจฺจ นนเหตุ นโนจิตฺตสํสฏฺฐสมุฏฺฐานสหภู ธมฺโม อุปฺปชฺชติ เหตุปจฺจยา.\csromanspacer{}\csromanspacer{}เหตุยา ตีณิ.}

\proseitem{37}{เหตุํ จิตฺตสํสฏฺฐสมุฏฺฐานานุปริวตฺติํ ธมฺมํ ปฏิจฺจ นเหตุ โนจิตฺตสํสฏฺฐสมุฏฺฐานานุปริวตฺตี ธมฺโม อุปฺปชฺชติ เหตุปจฺจยา.\csromanspacer{}\csromanspacer{}เหตุยา ตีณิ.}

\csromanheader{2}{ธมฺมานุโลมปจฺจนีย ทุกทุกปฏฺฐานปาฬิ}
\prose{\csromanfolio{314}นเหตุํ โนจิตฺตสํสฏฺฐสมุฏฺฐานานุปริวตฺติํ ธมฺมํ ปฏิจฺจ นนเหตุ นโนจิตฺตสํสฏฺฐสมุฏฺฐานานุปริวตฺตี ธมฺโม อุปฺปชฺชติ เหตุปจฺจยา.}

\csromanheader{2}{เหตุยา ตีณิ. (สํขิตฺตํ. นยวเสน วิตฺถาเรตพฺพํ.)}
\csromanheader{2}{1. เหตุทุก 82. ทสฺสเนนปหาตพฺพทุกาทิ}
\proseitem{38}{เหตุํ ทสฺสเนน ปหาตพฺพํ ธมฺมํ ปฏิจฺจ นเหตุ นทสฺสเนน ปหาตพฺโพ ธมฺโม อุปฺปชฺชติ เหตุปจฺจยา.\csromanspacer{}\csromanspacer{}เหตุยา ตีณิ.}

\prose{นเหตุํ นทสฺสเนน ปหาตพฺพํ ธมฺมํ ปจฺจยา นนเหตุ นนทสฺสเนน ปหาตพฺโพ ธมฺโม อุปฺปชฺชติ เหตุปจฺจยา.\csromanspacer{} ตีณิเมว.}

\proseitem{39}{เหตุํ ทสฺสเนน ปหาตพฺพเหตุกํ ธมฺมํ ปฏิจฺจ นเหตุ นทสฺสเนน ปหาตพฺพเหตุโก ธมฺโม อุปฺปชฺชติ เหตุปจฺจยา.\csromanspacer{}\csromanspacer{}เหตุยา ตีณิ.}

\prose{นเหตุํ นทสฺสเนน ปหาตพฺพเหตุกํ ธมฺมํ ปฏิจฺจ นนเหตุ นนทสฺสเนน ปหาตพฺพเหตุโก ธมฺโม อุปฺปชฺชติ เหตุปจฺจยา.\csromanspacer{}\csromanspacer{}เหตุยา เอกํ.}

\csromanheader{2}{1. เหตุทุก 99. สรณทุก}
\proseitem{40}{เหตุํ สรณํ ธมฺมํ ปฏิจฺจ นเหตุ นสรโณ ธมฺโม อุปฺปชฺชติ เหตุปจฺจยา.\csromanspacer{}\csromanspacer{}นเหตุํ สรณํ ธมฺมํ ปฏิจฺจ นเหตุ นสรโณ ธมฺโม อุปฺปชฺชติ เหตุปจฺจยา.\csromanspacer{}\csromanspacer{}เหตุํ สรณญฺจ นเหตุํ สรณญฺจ ธมฺมํ ปฏิจฺจ นเหตุ นสรโณ ธมฺโม อุปฺปชฺชติ เหตุปจฺจยา.\csromanspacer{}\csromanspacer{}เหตุยา ตีณิ.}

\prose{นเหตุํ อรณํ ธมฺมํ ปจฺจยา นนเหตุ น-อรโณ ธมฺโม อุปฺปชฺชติ เหตุปจฺจยา.\csromanspacer{}\csromanspacer{}นเหตุํ อรณํ ธมฺมํ ปจฺจยา นเหตุ น-อรโณ ธมฺโม อุปฺปชฺชติ เหตุปจฺจยา.\csromanspacer{}\csromanspacer{}นเหตุํ อรณํ ธมฺมํ ปจฺจยา นเหตุ น-อรโณ จ นนเหตุ น-อรโณ จ ธมฺมา อุปฺปชฺชนฺติ เหตุปจฺจยา.\csromanspacer{}\csromanspacer{}เหตุยา ตีณิ.}

\csromanheader{2}{2. สเหตุกทุก 1. เหตุทุก}
\proseitem{41}{สเหตุกํ เหตุํ ธมฺมํ ปฏิจฺจ นสเหตุโก นเหตุ ธมฺโม อุปฺปชฺชติ เหตุปจฺจยา.\csromanspacer{} ตีณิ.}

\csromanheader{2}{7-13. จูฬนฺตรทุก 1. เหตุทุก}
\prose{\csromanfolio{315}อเหตุกํ เหตุํ ธมฺมํ ปฏิจฺจ น-อเหตุโก นเหตุ ธมฺโม อุปฺปชฺชติ เหตุปจฺจยา.\csromanspacer{} ตีณิ.\csromanspacer{}\csromanspacer{}เหตุยา ฉ.}

\prose{สเหตุกํ นเหตุํ ธมฺมํ ปฏิจฺจ น-อเหตุโก นนเหตุ ธมฺโม อุปฺปชฺชติ เหตุปจฺจยา.\csromanspacer{}\csromanspacer{}เหตุยา ตีณิ.}

\csromanheader{2}{3. เหตุสมฺปยุตฺตทุก 1. เหตุทุก}
\proseitem{42}{เหตุสมฺปยุตฺตํ เหตุํ ธมฺมํ ปฏิจฺจ นเหตุสมฺปยุตฺโต นเหตุ ธมฺโม อุปฺปชฺชติ เหตุปจฺจยา.\csromanspacer{} ตีณิ.}

\prose{เหตุวิปฺปยุตฺตํ เหตุํ ธมฺมํ ปฏิจฺจ นเหตุวิปฺปยุตฺโต นเหตุ ธมฺโม อุปฺปชฺชติ เหตุปจฺจยา.\csromanspacer{} ตีณิ.\csromanspacer{}\csromanspacer{}เหตุยา ฉ.}

\prose{เหตุสมฺปยุตฺตํ นเหตุํ ธมฺมํ ปฏิจฺจ นเหตุวิปฺปยุตฺโต นนเหตุ ธมฺโม อุปฺปชฺชติ เหตุปจฺจยา.\csromanspacer{}\csromanspacer{}เหตุยา ตีณิ.}

\csromanheader{2}{4. เหตุสเหตุกทุก 1. เหตุทุก}
\proseitem{43}{เหตุญฺเจว สเหตุกญฺจ เหตุํ ธมฺมํ ปฏิจฺจ นเหตุ เจว น-อเหตุโก จ นเหตุ ธมฺโม อุปฺปชฺชติ เหตุปจฺจยา.\csromanspacer{}\csromanspacer{}เหตุยา เอกํ.}

\prose{สเหตุกญฺเจว น จ เหตุํ นเหตุํ ธมฺมํ ปฏิจฺจ น-อเหตุโก เจว นน จ เหตุ นนเหตุ ธมฺโม อุปฺปชฺชติ เหตุปจฺจยา.\csromanspacer{}\csromanspacer{}เหตุยา เอกํ.}

\csromanheader{2}{5. เหตุเหตุสมฺปยุตฺตทุก 1. เหตุทุก}
\proseitem{44}{เหตุญฺเจว เหตุสมฺปยุตฺตญฺจ เหตุํ ธมฺมํ ปฏิจฺจ นเหตุ เจว นเหตุวิปฺปยุตฺโต จ นเหตุ ธมฺโม อุปฺปชฺชติ เหตุปจฺจยา.\csromanspacer{}\csromanspacer{}เหตุยา เอกํ.}

\prose{เหตุสมฺปยุตฺตญฺเจว น จ เหตุํ นเหตุํ ธมฺมํ ปฏิจฺจ นเหตุวิปฺปยุตฺโต เจว นน จ เหตุ นนเหตุ ธมฺโม อุปฺปชฺชติ เหตุปจฺจยา.\csromanspacer{}\csromanspacer{}เหตุยา เอกํ.}

\csromanheader{2}{7-13. จูฬนฺตรทุก 1. เหตุทุก}
\proseitem{45}{สปฺปจฺจยํ เหตุํ ธมฺมํ ปฏิจฺจ น-อปฺปจฺจโย นเหตุ ธมฺโม อุปฺปชฺชติ เหตุปจฺจยา.\csromanspacer{}\csromanspacer{}เหตุยา เอกํ.}

\csromanheader{2}{ธมฺมานุโลมปจฺจนีย ทุกทุกปฏฺฐานปาฬิ}
\prose{\csromanfolio{316}สปฺปจฺจยํ นเหตุํ ธมฺมํ ปฏิจฺจ น-อปฺปจฺจโย นนเหตุ ธมฺโม อุปฺปชฺชติ เหตุปจฺจยา.\csromanspacer{}\csromanspacer{}เหตุยา เอกํ. (สงฺขตํ สปฺปจฺจยสทิสํ.)}

\proseitem{46}{อนิทสฺสนํ เหตุํ ธมฺมํ ปฏิจฺจ น-อนิทสฺสโน นเหตุ ธมฺโม อุปฺปชฺชติ เหตุปจฺจยา.\csromanspacer{}\csromanspacer{}เหตุยา ตีณิ.}

\prose{อนิทสฺสนํ นเหตุํ ธมฺมํ ปฏิจฺจ นสนิทสฺสโน นนเหตุ ธมฺโม อุปฺปชฺชติ เหตุปจฺจยา.\csromanspacer{}\csromanspacer{}เหตุยา เอกํ.}

\proseitem{47}{อปฺปฏิฆํ เหตุํ ธมฺมํ ปฏิจฺจ น-อปฺปฏิโฆ นเหตุ ธมฺโม อุปฺปชฺชติ เหตุปจฺจยา.\csromanspacer{}\csromanspacer{}เหตุยา ตีณิ.}

\prose{อปฺปฏิฆํ นเหตุํ ธมฺมํ ปฏิจฺจ นสปฺปฏิโฆ นนเหตุ ธมฺโม อุปฺปชฺชติ เหตุปจฺจยา.\csromanspacer{}\csromanspacer{}เหตุยา เอกํ.}

\proseitem{48}{อรูปิํ เหตุํ ธมฺมํ ปฏิจฺจ น-อรูปี นเหตุ ธมฺโม อุปฺปชฺชติ เหตุปจฺจยา.\csromanspacer{}\csromanspacer{}เหตุยา ตีณิ.}

\prose{รูปิํ นเหตุํ ธมฺมํ ปฏิจฺจ นรูปี นนเหตุ ธมฺโม อุปฺปชฺชติ เหตุปจฺจยา.}

\prose{อรูปิํ นเหตุํ ธมฺมํ ปฏิจฺจ น-อรูปี นนเหตุ ธมฺโม อุปฺปชฺชติ เหตุปจฺจยา.\csromanspacer{}\csromanspacer{}เหตุยา ตีณิ.}

\proseitem{49}{โลกิยํ เหตุํ ธมฺมํ ปฏิจฺจ นโลกุตฺตโร นเหตุ ธมฺโม อุปฺปชฺชติ เหตุปจฺจยา.\csromanspacer{} เอกํ.}

\prose{โลกุตฺตรํ เหตุํ ธมฺมํ ปฏิจฺจ นโลกุตฺตโร นเหตุ ธมฺโม อุปฺปชฺชติ เหตุปจฺจยา.\csromanspacer{} ตีณิ.\csromanspacer{}\csromanspacer{}เหตุยา จตฺตาริ.}

\prose{โลกิยํ นเหตุํ ธมฺมํ ปฏิจฺจ นโลกุตฺตโร นนเหตุ ธมฺโม อุปฺปชฺชติ เหตุปจฺจยา.\csromanspacer{} เอกํ.}

\prose{โลกุตฺตรํ นเหตุํ ธมฺมํ ปฏิจฺจ นโลกิโย นนเหตุ ธมฺโม อุปฺปชฺชติ เหตุปจฺจยา.\csromanspacer{} เอกํ.\csromanspacer{}\csromanspacer{}เหตุยา เทฺว.}

\proseitem{50}{เกนจิ วิญฺเญยฺยํ เหตุํ ธมฺมํ ปฏิจฺจ นเกนจิ วิญฺเญยฺโย นเหตุ ธมฺโม อุปฺปชฺชติ เหตุปจฺจยา.\csromanspacer{} ตีณิ.}

\prose{นเกนจิ วิญฺเญยฺยํ เหตุํ ธมฺมํ ปฏิจฺจ นนเกนจิ วิญฺเญยฺโย นเหตุ ธมฺโม อุปฺปชฺชติ เหตุปจฺจยา.\csromanspacer{} ตีณิ.}

\csromanheader{2}{16. อาสวสมฺปยุตฺตทุก 1. เหตุทุก}
\prose{\csromanfolio{317}เกนจิ วิญฺเญยฺยํ เหตุญฺจ นเกนจิ วิญฺเญยฺยํ เหตุญฺจ ธมฺมํ ปฏิจฺจ นเกนจิ วิญฺเญยฺโย นเหตุ ธมฺโม อุปฺปชฺชติ เหตุปจฺจยา.\csromanspacer{}\csromanspacer{}เหตุยา นว.}

\prose{เกนจิ วิญฺเญยฺยํ นเหตุํ ธมฺมํ ปฏิจฺจ นเกนจิ วิญฺเญยฺโย นนเหตุ ธมฺโม อุปฺปชฺชติ เหตุปจฺจยา.\csromanspacer{}\csromanspacer{}เหตุยา นว.}

\csromanheader{2}{14. อาสวทุก 1. เหตุทุก}
\proseitem{51}{อาสวํ เหตุํ ธมฺมํ ปฏิจฺจ โน-อาสโว นเหตุ ธมฺโม อุปฺปชฺชติ เหตุปจฺจยา.\csromanspacer{} ตีณิ.}

\prose{โน-อาสวํ เหตุํ ธมฺมํ ปฏิจฺจ นโน-อาสโว นเหตุ ธมฺโม อุปฺปชฺชติ เหตุปจฺจยา.\csromanspacer{} เอกํ.}

\prose{อาสวํ เหตุญฺจ โน-อาสวํ เหตุญฺจ ธมฺมํ ปฏิจฺจ โน-อาสโว นเหตุ ธมฺโม อุปฺปชฺชติ เหตุปจฺจยา.\csromanspacer{} เอกํ. (สพฺพตฺถ ปญฺจ.)}

\prose{โน-อาสวํ นเหตุํ ธมฺมํ ปฏิจฺจ นโน-อาสโว นนเหตุ ธมฺโม อุปฺปชฺชติ เหตุปจฺจยา.\csromanspacer{}\csromanspacer{}เหตุยา ปญฺจ.}

\csromanheader{2}{15. สาสวทุก 1. เหตุทุก}
\proseitem{52}{สาสวํ เหตุํ ธมฺมํ ปฏิจฺจ น-อนาสโว นเหตุ ธมฺโม อุปฺปชฺชติ เหตุปจฺจยา.\csromanspacer{} เอกํ.}

\prose{อนาสวํ เหตุํ ธมฺมํ ปฏิจฺจ น-อนาสโว นเหตุ ธมฺโม อุปฺปชฺชติ เหตุปจฺจยา.\csromanspacer{} ตีณิ.\csromanspacer{}\csromanspacer{}เหตุยา จตฺตาริ. (โลกิยทุกสทิสํ.)}

\csromanheader{2}{16. อาสวสมฺปยุตฺตทุก 1. เหตุทุก}
\proseitem{53}{อาสวสมฺปยุตฺตํ เหตุํ ธมฺมํ ปฏิจฺจ น-อาสวสมฺปยุตฺโต นเหตุ ธมฺโม อุปฺปชฺชติ เหตุปจฺจยา.\csromanspacer{} ตีณิ.}

\prose{อาสววิปฺปยุตฺตํ เหตุํ ธมฺมํ ปฏิจฺจ น-อาสววิปฺปยุตฺโต นเหตุ ธมฺโม อุปฺปชฺชติ เหตุปจฺจยา.\csromanspacer{} ตีณิ.}

\prose{อาสวสมฺปยุตฺตํ เหตุญฺจ อาสววิปฺปยุตฺตํ เหตุญฺจ ธมฺมํ ปฏิจฺจ น-อาสวสมฺปยุตฺโต นเหตุ ธมฺโม อุปฺปชฺชติ เหตุปจฺจยา.\csromanspacer{} ตีณิ.\csromanspacer{}\csromanspacer{}เหตุยา นว.}

\prose{อาสวสมฺปยุตฺตํ นเหตุํ ธมฺมํ ปฏิจฺจ น-อาสวสมฺปยุตฺโต นนเหตุ ธมฺโม อุปฺปชฺชติ เหตุปจฺจยา.\csromanspacer{} ตีณิ.}

\csromanheader{2}{ธมฺมานุโลมปจฺจนีย ทุกทุกปฏฺฐานปาฬิ}
\prose{\csromanfolio{318}อาสววิปฺปยุตฺตํ นเหตุํ ธมฺมํ ปฏิจฺจ น-อาสวสมฺปยุตฺโต นนเหตุ ธมฺโม อุปฺปชฺชติ เหตุปจฺจยา.\csromanspacer{} เอกํ.\csromanspacer{}\csromanspacer{}เหตุยา จตฺตาริ.}

\csromanheader{2}{17. อาสวสาสวทุก 1. เหตุทุก}
\proseitem{54}{อาสวญฺเจว สาสวญฺจ เหตุํ ธมฺมํ ปฏิจฺจ น-อาสโว เจว น-อนาสโว จ นเหตุ ธมฺโม อุปฺปชฺชติ เหตุปจฺจยา.\csromanspacer{} ตีณิ.}

\prose{สาสวญฺเจว โน จ อาสวํ เหตุํ ธมฺมํ ปฏิจฺจ น-อาสโว เจว น-อนาสโว จ นเหตุ ธมฺโม อุปฺปชฺชติ เหตุปจฺจยา.\csromanspacer{} เอกํ.}

\prose{อาสวญฺเจว สาสวํ เหตุญฺจ สาสวญฺเจว โน จ อาสวํ เหตุญฺจ ธมฺมํ ปฏิจฺจ น-อาสโว เจว น-อนาสโว จ นเหตุ ธมฺโม อุปฺปชฺชติ เหตุปจฺจยา.\csromanspacer{} เอกํ.\csromanspacer{}\csromanspacer{}เหตุยา ปญฺจ.}

\prose{อาสวญฺเจว สาสวญฺจ นเหตุํ ธมฺมํ ปฏิจฺจ น-อนาสโว เจว นโน-อาสโว จ นนเหตุ ธมฺโม อุปฺปชฺชติ เหตุปจฺจยา.\csromanspacer{}\csromanspacer{}เหตุยา ปญฺจ.\csromanspacer{}\csromanspacer{}}

\csromanheader{2}{18. อาสว-อาสวสมฺปยุตฺตทุก 1. เหตุทุก}
\proseitem{55}{อาสวญฺเจว อาสวสมฺปยุตฺตญฺจ เหตุํ ธมฺมํ ปฏิจฺจ น-อาสโว เจว น-อาสววิปฺปยุตฺโต จ นเหตุ ธมฺโม อุปฺปชฺชติ เหตุปจฺจยา.\csromanspacer{} ตีณิ.}

\prose{อาสวสมฺปยุตฺตญฺเจว โน จ อาสวํ เหตุํ ธมฺมํ ปฏิจฺจ น-อาสโว เจว น-อาสววิปฺปยุตฺโต จ นเหตุ ธมฺโม อุปฺปชฺชติ เหตุปจฺจยา.\csromanspacer{} เอกํ.\csromanspacer{}\csromanspacer{}เหตุยา จตฺตาริ.}

\prose{อาสวญฺเจว อาสวสมฺปยุตฺตญฺจ นเหตุํ ธมฺมํ ปฏิจฺจ น-อาสววิปฺปยุตฺโต เจว นโน จ อาสโว นนเหตุ ธมฺโม อุปฺปชฺชติ เหตุปจฺจยา.\csromanspacer{} เอกํ.}

\prose{อาสวสมฺปยุตฺตญฺเจว โน จ อาสวํ นเหตุํ ธมฺมํ ปฏิจฺจ น-อาสววิปฺปยุตฺโต เจว นโน จ อาสโว นนเหตุ ธมฺโม อุปฺปชฺชติ เหตุปจฺจยา.\csromanspacer{} ตีณิ.\csromanspacer{}\csromanspacer{}เหตุยา จตฺตาริ.}

\csromanheader{2}{19. อาสววิปฺปยุตฺตสาสวทุก 1. เหตุทุก}
\proseitem{56}{อาสววิปฺปยุตฺตํ สาสวํ เหตุํ ธมฺมํ ปฏิจฺจ อาสววิปฺปยุตฺโต น-อนาสโว นเหตุ ธมฺโม อุปฺปชฺชติ เหตุปจฺจยา.\csromanspacer{} เอกํ.}

\csromanheader{2}{55. มหนฺตรทุก 1. เหตุทุก}
\prose{\csromanfolio{319}อาสววิปฺปยุตฺตํ อนาสวํ เหตุํ ธมฺมํ ปฏิจฺจ อาสววิปฺปยุตฺโต น-อนาสโว นเหตุ ธมฺโม อุปฺปชฺชติ เหตุปจฺจยา.\csromanspacer{} ตีณิ.\csromanspacer{} เหตุยา จตฺตาริ.}

\prose{อาสววิปฺปยุตฺตํ สาสวํ นเหตุํ ธมฺมํ ปฏิจฺจ อาสววิปฺปยุตฺโต น-อนาสโว นนเหตุ ธมฺโม อุปฺปชฺชติ เหตุปจฺจยา.\csromanspacer{}\csromanspacer{}อาสววิปฺปยุตฺตํ อนาสวํ นเหตุํ ธมฺมํ ปฏิจฺจ อาสววิปฺปยุตฺโต นสาสโว นนเหตุ ธมฺโม อุปฺปชฺชติ เหตุปจฺจยา.\csromanspacer{}\csromanspacer{}เหตุยา เทฺว.}

\csromanheader{2}{20. สญฺโญชนโคจฺฉกาทิ 1. เหตุทุก}
\proseitem{57}{สญฺโญชนํ เหตุํ ธมฺมํ ปฏิจฺจ โนสญฺโญชโน นเหตุ ธมฺโม อุปฺปชฺชติ เหตุปจฺจยา.}

\proseitem{58}{คนฺถํ เหตุํ ธมฺมํ ปฏิจฺจ โนคนฺโถ นเหตุ ธมฺโม อุปฺปชฺชติ เหตุปจฺจยา.}

\proseitem{59}{โอฆํ เหตุํ ธมฺมํ ปฏิจฺจ โน-โอโฆ นเหตุ ธมฺโม อุปฺปชฺชติ เหตุปจฺจยา.}

\proseitem{60}{โยคํ เหตุํ ธมฺมํ ปฏิจฺจ โนโยโค นเหตุ ธมฺโม อุปฺปชฺชติ เหตุปจฺจยา.}

\proseitem{61}{นีวรณํ เหตุํ ธมฺมํ ปฏิจฺจ โนนีวรโณ นเหตุ ธมฺโม อุปฺปชฺชติ เหตุปจฺจยา.}

\proseitem{62}{โนปรามาสํ เหตุํ ธมฺมํ ปฏิจฺจ นโนปรามาโส นเหตุ ธมฺโม อุปฺปชฺชติ เหตุปจฺจยา.\csromanspacer{}\csromanspacer{}เหตุยา ตีณิ.}

\csromanheader{2}{55. มหนฺตรทุก 1. เหตุทุก}
\proseitem{63}{สารมฺมณํ เหตุํ ธมฺมํ ปฏิจฺจ นสารมฺมโณ นเหตุ ธมฺโม อุปฺปชฺชติ เหตุปจฺจยา.\csromanspacer{}\csromanspacer{}เหตุยา ตีณิ.}

\prose{สารมฺมณํ นเหตุํ ธมฺมํ ปฏิจฺจ น-อนารมฺมโณ นนเหตุ ธมฺโม อุปฺปชฺชติ เหตุปจฺจยา.\csromanspacer{}\csromanspacer{}อนารมฺมณํ นเหตุํ ธมฺมํ ปฏิจฺจ น-อนารมฺมโณ นนเหตุ}

\vspace{\tipitakaparskip}
\csromancenter{ธมฺมานุโลมปจฺจนีย ทุกทุกปฏฺฐานปาฬิ ธมฺโม อุปฺปชฺชติ เหตุปจฺจยา.\csromanspacer{}\csromanspacer{}สารมฺมณํ นเหตุญฺจ อนารมฺมณํ นเหตุญฺจ ธมฺมํ ปฏิจฺจ น-อนารมฺมโณ นนเหตุ ธมฺโม อุปฺปชฺชติ เหตุปจฺจยา.\csromanspacer{}\csromanspacer{}เหตุยา ตีณิ.}
\proseitem{64}{\csromanfolio{320}โนจิตฺตํ เหตุํ ธมฺมํ ปฏิจฺจ นโนจิตฺโต นเหตุ ธมฺโม อุปฺปชฺชติ เหตุปจฺจยา.\csromanspacer{}\csromanspacer{}เหตุยา ตีณิ.}

\prose{จิตฺตํ นเหตุํ ธมฺมํ ปฏิจฺจ โนจิตฺโต นนเหตุ ธมฺโม อุปฺปชฺชติ เหตุปจฺจยา.\csromanspacer{}\csromanspacer{}เหตุยา ตีณิ.}

\proseitem{65}{เจตสิกํ เหตุํ ธมฺมํ ปฏิจฺจ นเจตสิโก นเหตุ ธมฺโม อุปฺปชฺชติ เหตุปจฺจยา.\csromanspacer{}\csromanspacer{}เหตุยา ตีณิ.}

\prose{เจตสิกํ นเหตุํ ธมฺมํ ปฏิจฺจ น-อเจตสิโก นนเหตุ ธมฺโม อุปฺปชฺชติ เหตุปจฺจยา.\csromanspacer{}\csromanspacer{}เหตุยา ตีณิ.}

\proseitem{66}{จิตฺตสมฺปยุตฺตํ เหตุํ ธมฺมํ ปฏิจฺจ นจิตฺตสมฺปยุตฺโต นเหตุ ธมฺโม อุปฺปชฺชติ เหตุปจฺจยา.\csromanspacer{}\csromanspacer{}เหตุยา ตีณิ. (สํขิตฺตํ.)}

\proseitem{67}{จิตฺตสํสฏฺฐํ เหตุํ ธมฺมํ ปฏิจฺจ นจิตฺตสํสฏฺโฐ นเหตุ ธมฺโม อุปฺปชฺชติ เหตุปจฺจยา.\csromanspacer{}\csromanspacer{}เหตุยา ตีณิ.}

\proseitem{68}{จิตฺตสมุฏฺฐานํ เหตุํ ธมฺมํ ปฏิจฺจ โนจิตฺตสมุฏฺฐาโน นเหตุ ธมฺโม อุปฺปชฺชติ เหตุปจฺจยา.\csromanspacer{}\csromanspacer{}เหตุยา ตีณิ.}

\proseitem{69}{จิตฺตสหภุํ เหตุํ ธมฺมํ ปฏิจฺจ โนจิตฺตสหภู นเหตุ ธมฺโม อุปฺปชฺชติ เหตุปจฺจยา.\csromanspacer{}\csromanspacer{}เหตุยา ตีณิ.}

\proseitem{70}{จิตฺตานุปริวตฺติํ เหตุํ ธมฺมํ ปฏิจฺจ โนจิตฺตานุปริวตฺตี นเหตุ ธมฺโม อุปฺปชฺชติ เหตุปจฺจยา.\csromanspacer{}\csromanspacer{}เหตุยา ตีณิ.}

\proseitem{71}{จิตฺตสํสฏฺฐสมุฏฺฐานํ เหตุํ ธมฺมํ ปฏิจฺจ โนจิตฺตสํสฏฺฐสมุฏฺฐาโน นเหตุ ธมฺโม อุปฺปชฺชติ เหตุปจฺจยา.\csromanspacer{}\csromanspacer{}เหตุยา ตีณิ.}

\proseitem{72}{จิตฺตสํสฏฺฐสมุฏฺฐานสหภุํ เหตุํ ธมฺมํ ปฏิจฺจ โนจิตฺตสํสฏฺฐสมุฏฺฐานสหภู นเหตุ ธมฺโม อุปฺปชฺชติ เหตุปจฺจยา.\csromanspacer{}\csromanspacer{}เหตุยา ตีณิ.}

\csromanheader{2}{69. อุปาทานโคจฺฉกาทิ 1. เหตุทุก}
\proseitem{73}{\csromanfolio{321}จิตฺตสํสฏฺฐสมุฏฺฐานานุปริวตฺติํ เหตุํ ธมฺมํ ปฏิจฺจ โนจิตฺตสํสฏฺฐสมุฏฺฐานานุปริวตฺตี นเหตุ ธมฺโม อุปฺปชฺชติ เหตุปจฺจยา.\csromanspacer{}\csromanspacer{}เหตุยา ตีณิ.}

\proseitem{74}{พาหิรํ เหตุํ ธมฺมํ ปฏิจฺจ นพาหิโร นเหตุ ธมฺโม อุปฺปชฺชติ เหตุปจฺจยา.\csromanspacer{}\csromanspacer{}เหตุยา ตีณิ.}

\prose{อชฺฌตฺติกํ นเหตุํ ธมฺมํ ปฏิจฺจ น-อชฺฌตฺติโก นนเหตุ ธมฺโม อุปฺปชฺชติ เหตุปจฺจยา.\csromanspacer{}\csromanspacer{}เหตุยา ตีณิ.}

\proseitem{75}{โน-อุปาทา เหตุํ ธมฺมํ ปฏิจฺจ นโน-อุปาทา นเหตุ ธมฺโม อุปฺปชฺชติ เหตุปจฺจยา.\csromanspacer{}\csromanspacer{}เหตุยา ตีณิ.}

\prose{อุปาทา นเหตุํ ธมฺมํ ปฏิจฺจ โน-อุปาทา นนเหตุ ธมฺโม อุปฺปชฺชติ เหตุปจฺจยา.\csromanspacer{}\csromanspacer{}เหตุยา ตีณิ.}

\proseitem{76}{อุปาทินฺนํ เหตุํ ธมฺมํ ปฏิจฺจ น-อุปาทินฺโน นเหตุ ธมฺโม อุปฺปชฺชติ เหตุปจฺจยา.\csromanspacer{} ตีณิ.}

\prose{อนุปาทินฺนํ เหตุํ ธมฺมํ ปฏิจฺจ น-อุปาทินฺโน นเหตุ ธมฺโม อุปฺปชฺชติ เหตุปจฺจยา เอกํ.\csromanspacer{}\csromanspacer{}เหตุยา จตฺตาริ.}

\prose{อุปาทินฺนํ นเหตุํ ธมฺมํ ปฏิจฺจ น-อนุปาทินฺโน นนเหตุ ธมฺโม อุปฺปชฺชติ เหตุปจฺจยา.\csromanspacer{} เอกํ.}

\prose{อนุปาทินฺนํ นเหตุํ ธมฺมํ ปฏิจฺจ น-อุปาทินฺโน นนเหตุ ธมฺโม อุปฺปชฺชติ เหตุปจฺจยา.\csromanspacer{} เอกํ.\csromanspacer{}\csromanspacer{}เหตุยา เทฺว.}

\csromanheader{2}{69. อุปาทานโคจฺฉกาทิ 1. เหตุทุก}
\proseitem{77}{อุปาทานํ เหตุํ ธมฺมํ ปฏิจฺจ โน-อุปาทาโน นเหตุ ธมฺโม อุปฺปชฺชติ เหตุปจฺจยา.}

\proseitem{78}{กิเลสํ เหตุํ ธมฺมํ ปฏิจฺจ โนกิเลโส นเหตุ ธมฺโม อุปฺปชฺชติ เหตุปจฺจยา.}

\csromanheader{2}{ธมฺมานุโลมปจฺจนีย ทุกทุกปฏฺฐานปาฬิ \textbf{83. ปิฏฺฐิทุก 1. เหตุทุก}}
\proseitem{79}{\csromanfolio{322}ทสฺสเนน ปหาตพฺพํ เหตุํ ธมฺมํ ปฏิจฺจ นทสฺสเนน ปหาตพฺโพ นเหตุ ธมฺโม อุปฺปชฺชติ เหตุปจฺจยา.\csromanspacer{} ตีณิ.}

\prose{นทสฺสเนน ปหาตพฺพํ เหตุํ ธมฺมํ ปฏิจฺจ นทสฺสเนน ปหาตพฺโพ นเหตุ ธมฺโม อุปฺปชฺชติ เหตุปจฺจยา.\csromanspacer{} เอกํ.\csromanspacer{}\csromanspacer{}เหตุยา จตฺตาริ.}

\prose{ทสฺสเนน ปหาตพฺพํ นเหตุํ ธมฺมํ ปฏิจฺจ นนทสฺสเนน ปหาตพฺโพ นนเหตุ ธมฺโม เหตุปจฺจยา.\csromanspacer{} เอกํ.}

\prose{นทสฺสเนน ปหาตพฺพํ นเหตุํ ธมฺมํ ปฏิจฺจ นทสฺสเนน ปหาตพฺโพ นนเหตุ ธมฺโม อุปฺปชฺชติ เหตุปจฺจยา.\csromanspacer{} เอกํ.\csromanspacer{}\csromanspacer{}เหตุยา เทฺว.}

\proseitem{80}{ภาวนาย ปหาตพฺพํ เหตุํ ธมฺมํ ปฏิจฺจ นภาวนาย ปหาตพฺโพ นเหตุ ธมฺโม อุปฺปชฺชติ เหตุปจฺจยา.\csromanspacer{} ตีณิ.\csromanspacer{}\csromanspacer{}นภาวนาย ปหาตพฺพํ เหตุํ ธมฺมํ ปฏิจฺจ นภาวนาย ปหาตพฺโพ นเหตุ ธมฺโม อุปฺปชฺชติ เหตุปจฺจยา.\csromanspacer{} เอกํ.\csromanspacer{}\csromanspacer{}เหตุยา จตฺตาริ.}

\prose{ภาวนาย ปหาตพฺพํ นเหตุํ ธมฺมํ ปฏิจฺจ นนภาวนาย ปหาตพฺโพ นนเหตุ ธมฺโม อุปฺปชฺชติ เหตุปจฺจยา.\csromanspacer{} เอกํ.}

\prose{นภาวนาย ปหาตพฺพํ นเหตุํ ธมฺมํ ปฏิจฺจ นภาวนาย ปหาตพฺโพ นนเหตุ ธมฺโม อุปฺปชฺชติ เหตุปจฺจยา.\csromanspacer{} เอกํ.\csromanspacer{}\csromanspacer{}เหตุยา เทฺว.}

\proseitem{81}{ทสฺสเนน ปหาตพฺพเหตุกํ เหตุํ ธมฺมํ ปฏิจฺจ นทสฺสเนน ปหาตพฺพเหตุโก นเหตุ ธมฺโม อุปฺปชฺชติ เหตุปจฺจยา.\csromanspacer{} ตีณิ.}

\prose{นทสฺสเนน ปหาตพฺพเหตุกํ เหตุํ ธมฺมํ ปฏิจฺจ นนทสฺสเนน ปหาตพฺพเหตุโก นเหตุ ธมฺโม อุปฺปชฺชติ เหตุปจฺจยา.\csromanspacer{} ตีณิ.\csromanspacer{}\csromanspacer{}(สพฺพตฺถ ปเญฺห สํขิตฺตํ.)}

\prose{ทสฺสเนน ปหาตพฺพเหตุกํ นเหตุํ ธมฺมํ ปฏิจฺจ นนทสฺสเนน ปหาตพฺพเหตุโก นนเหตุ ธมฺโม อุปฺปชฺชติ เหตุปจฺจยา.\csromanspacer{}\csromanspacer{}นทสฺสเนน ปหาตพฺพเหตุกํ นเหตุํ ธมฺมํ ปฏิจฺจ นทสฺสเนน ปหาตพฺพเหตุโก นนเหตุ ธมฺโม อุปฺปชฺชติ เหตุปจฺจยา.\csromanspacer{}\csromanspacer{}เหตุยา เทฺว.}

\csromanheader{2}{100. สรณทุก 1. เหตุทุกาทิ \textbf{100. สรณทุก 1. เหตุทุกาทิ}}
\proseitem{82}{\csromanfolio{323}สรณํ เหตุํ ธมฺมํ ปฏิจฺจ นสรโณ นเหตุ ธมฺโม อุปฺปชฺชติ เหตุปจฺจยา.\csromanspacer{} ตีณิ.}

\prose{อรณํ เหตุํ ธมฺมํ ปฏิจฺจ นสรโณ นเหตุ ธมฺโม อุปฺปชฺชติ เหตุปจฺจยา.\csromanspacer{} เอกํ.\csromanspacer{}\csromanspacer{}เหตุยา จตฺตาริ.}

\prose{สรณํ นเหตุํ ธมฺมํ ปฏิจฺจ น-อรโณ นนเหตุ ธมฺโม อุปฺปชฺชติ เหตุปจฺจยา.\csromanspacer{} เอกํ.}

\prose{อรณํ นเหตุํ ธมฺมํ ปฏิจฺจ นสรโณ นนเหตุ ธมฺโม อุปฺปชฺชติ เหตุปจฺจยา.\csromanspacer{} เอกํ.\csromanspacer{}\csromanspacer{}เหตุยา เทฺว.}

\proseitem{83}{สรณํ สเหตุกํ ธมฺมํ ปฏิจฺจ นสรโณ นสเหตุโก ธมฺโม อุปฺปชฺชติ เหตุปจฺจยา.\csromanspacer{} เอกํ.}

\prose{อรณํ สเหตุกํ ธมฺมํ ปฏิจฺจ นสรโณ นสเหตุโก ธมฺโม อุปฺปชฺชติ เหตุปจฺจยา.\csromanspacer{} เอกํ.\csromanspacer{}\csromanspacer{}เหตุยา เทฺว.}

\prose{สรณํ อเหตุกํ ธมฺมํ ปฏิจฺจ น-อรโณ น-อเหตุโก ธมฺโม อุปฺปชฺชติ เหตุปจฺจยา.\csromanspacer{} เอกํ.}

\prose{อรณํ อเหตุกํ ธมฺมํ ปฏิจฺจ นสรโณ น-อเหตุโก ธมฺโม อุปฺปชฺชติ เหตุปจฺจยา.\csromanspacer{} เอกํ.\csromanspacer{}\csromanspacer{}เหตุยา เทฺว.}

\proseitem{84}{สรณํ เหตุสมฺปยุตฺตํ ธมฺมํ ปฏิจฺจ นสรโณ นเหตุสมฺปยุตฺโต ธมฺโม อุปฺปชฺชติ เหตุปจฺจยา.\csromanspacer{} เอกํ.\csromanspacer{}\csromanspacer{}อรณํ เหตุสมฺปยุตฺตํ ธมฺมํ ปฏิจฺจ นสรโณ นเหตุสมฺปยุตฺโต ธมฺโม อุปฺปชฺชติ เหตุปจฺจยา.\csromanspacer{} เอกํ.\csromanspacer{}\csromanspacer{}เหตุยา เทฺว. (สเหตุกทุกสทิสํ.)}

\proseitem{85}{สรณํ เหตุญฺเจว สเหตุกญฺจ ธมฺมํ ปฏิจฺจ น-อรโณ นเหตุ เจว น-อเหตุโก จ ธมฺโม อุปฺปชฺชติ เหตุปจฺจยา.\csromanspacer{} เอกํ.}

\prose{อรณํ เหตุญฺเจว สเหตุกญฺจ ธมฺมํ ปฏิจฺจ นสรโณ นเหตุ เจว น-อเหตุโก จ ธมฺโม อุปฺปชฺชติ เหตุปจฺจยา.\csromanspacer{} เอกํ.\csromanspacer{} เหตุยา เทฺว.}

\prose{สรณํ สเหตุกญฺเจว น จ เหตุํ ธมฺมํ ปฏิจฺจ น-อรโณ น-อเหตุโก เจว นน จ เหตุ ธมฺโม อุปฺปชฺชติ เหตุปจฺจยา.\csromanspacer{} เอกํ.}

\csromanheader{2}{ธมฺมานุโลมปจฺจนีย ทุกทุกปฏฺฐานปาฬิ}
\prose{\csromanfolio{324}อรณํ สเหตุกญฺเจว น จ เหตุํ ธมฺมํ ปฏิจฺจ นสรโณ น-อเหตุโก เจว นน จ เหตุ ธมฺโม อุปฺปชฺชติ เหตุปจฺจยา.\csromanspacer{} เอกํ.\csromanspacer{}\csromanspacer{}เหตุยา เทฺว.}

\proseitem{86}{สรณํ เหตุญฺเจว เหตุสมฺปยุตฺตญฺจ ธมฺมํ ปฏิจฺจ น-อรโณ นเหตุ เจว นเหตุวิปฺปยุตฺโต จ ธมฺโม อุปฺปชฺชติ เหตุปจฺจยา. (เหตุสเหตุกทุกสทิสํ.)}

\proseitem{87}{สรณํ นเหตุํ สเหตุกํ ธมฺมํ ปฏิจฺจ นสรโณ นเหตุ นสเหตุโก ธมฺโม อุปฺปชฺชติ เหตุปจฺจยา.\csromanspacer{} เอกํ.}

\prose{อรณํ นเหตุํ สเหตุกํ ธมฺมํ ปฏิจฺจ นสรโณ นเหตุ นสเหตุโก ธมฺโม อุปฺปชฺชติ เหตุปจฺจยา.\csromanspacer{} เอกํ.\csromanspacer{}\csromanspacer{}เหตุยา เทฺว.}

\prose{อรณํ นเหตุํ อเหตุกํ ธมฺมํ ปฏิจฺจ นสรโณ นเหตุ น-อเหตุโก ธมฺโม อุปฺปชฺชติ เหตุปจฺจยา.\csromanspacer{}\csromanspacer{}เหตุยา เอกํ.}

\csromanheader{2}{100. สรณทุก 7. จูฬนฺตรทุก}
\proseitem{88}{อรโณ อปฺปจฺจโย ธมฺโม นสรณสฺส น-อปฺปจฺจยสฺส ธมฺมสฺส อารมฺมณปจฺจเยน ปจฺจโย.}

\proseitem{89}{อรโณ อสงฺขโต ธมฺโม นสรณสฺส น-อสงฺขตสฺส ธมฺมสฺส อารมฺมณปจฺจเยน ปจฺจโย.}

\proseitem{90}{อรโณ สนิทสฺสโน ธมฺโม น-อรณสฺส นสนิทสฺสนสฺส ธมฺมสฺส อารมฺมณปจฺจเยน ปจฺจโย.\csromanspacer{}\csromanspacer{}อรโณ สนิทสฺสโน ธมฺโม นสรณสฺส นสนิทสฺสนสฺส ธมฺมสฺส อารมฺมณปจฺจเยน ปจฺจโย.}

\csromanheader{2}{100. สรณทุก 14. อาสวโคจฺฉก}
\proseitem{91}{สรณํ อาสวํ ธมฺมํ ปฏิจฺจ นสรโณ โน-อาสโว ธมฺโม อุปฺปชฺชติ เหตุปจฺจยา.\csromanspacer{}\csromanspacer{}สรณํ อาสวํ ธมฺมํ ปฏิจฺจ น-อรโณ โน-อาสโว ธมฺโม อุปฺปชฺชติ เหตุปจฺจยา.\csromanspacer{}\csromanspacer{}สรณํ อาสวํ ธมฺมํ ปฏิจฺจ นสรโณ โน-อาสโว จ น-อรโณ โน-อาสโว จ ธมฺมา อุปฺปชฺชนฺติ เหตุปจฺจยา.}

\csromanheader{2}{100. สรณทุก 83. ปิฏฺฐิทุก}
\prose{\csromanfolio{325}สรณํ โน-อาสวํ ธมฺมํ ปฏิจฺจ น-อรโณ นโน-อาสโว ธมฺโม อุปฺปชฺชติ เหตุปจฺจยา.\csromanspacer{}\csromanspacer{}เหตุยา เอกํ.}

\proseitem{92}{อรณํ สาสวํ ธมฺมํ ปจฺจยา นสรโณ นสาสโว ธมฺโม อุปฺปชฺชติ เหตุปจฺจยา.}

\prose{อรณํ อนาสวํ ธมฺมํ ปฏิจฺจ นสรโณ น-อนาสโว ธมฺโม อุปฺปชฺชติ เหตุปจฺจยา.\csromanspacer{}\csromanspacer{}เหตุยา เอกํ. (เอเตน อุปาเยน สพฺพตฺถ วิตฺถาเรตพฺพํ.)}

\csromanheader{2}{100. สรณทุก 55. มหนฺตรทุก}
\proseitem{93}{สรณํ สารมฺมณํ ธมฺมํ ปฏิจฺจ นสรโณ นสารมฺมโณ ธมฺโม อุปฺปชฺชติ เหตุปจฺจยา.\csromanspacer{}\csromanspacer{}อรณํ สารมฺมณํ ธมฺมํ ปฏิจฺจ นสรโณ นสารมฺมโณ ธมฺโม อุปฺปชฺชติ เหตุปจฺจยา. (สํขิตฺตํ.)}

\csromanheader{2}{100. สรณทุก 83. ปิฏฺฐิทุก}
\proseitem{94}{สรณํ ทสฺสเนน ปหาตพฺพํ ธมฺมํ ปฏิจฺจ นสรโณ นทสฺสเนน ปหาตพฺโพ ธมฺโม อุปฺปชฺชติ เหตุปจฺจยา.\csromanspacer{}\csromanspacer{}เหตุยา เอกํ.}

\prose{อรณํ นทสฺสเนน ปหาตพฺพํ ธมฺมํ ปจฺจยา น-อรโณ นนทสฺสเนน ปหาตพฺโพ ธมฺโม อุปฺปชฺชติ เหตุปจฺจยา. (สํขิตฺตํ.)}

\proseitem{95}{อรณํ ส-อุตฺตรํ ธมฺมํ ปจฺจยา นสรโณ นส-อุตฺตโร ธมฺโม อุปฺปชฺชติ เหตุปจฺจยา.}

\prose{อรณํ อนุตฺตรํ ธมฺมํ ปฏิจฺจ นสรโณ น-อนุตฺตโร ธมฺโม อุปฺปชฺชติ เหตุปจฺจยา. (สหชาตวารมฺปิ ฯเปฯ ปญฺหาวารมฺปิ วิตฺถาเรตพฺพํ.)}

\nitthitam{ธมฺมานุโลมปจฺจนีเย ทุกทุกปฏฺฐานํ นิฏฺฐิตํ.}
\nitthitam{\textbf{อนุโลมปจฺจนียปฏฺฐานํ นิฏฺฐิตํ.}}
\pitaka{\textbf{อภิธมฺมปิฏก}}
\csromanheader{2}{\textbf{ธมฺมปจฺจนียานุโลม}}
\gambhira{\textbf{ติกปฏฺฐานปาฬิ}}
\namakkaram{นโม ตสฺส ภควโต อรหโต สมฺมาสมฺพุทฺธสฺส.}
\csromanheader{2}{1. กุสลตฺติก 1. ปฏิจฺจวาร}
\csromanheader{2}{\textbf{ปจฺจยจตุกฺก-เหตุ-อารมฺมณ}}
\proseitem{1}{\csromanfolio{327}นกุสลํ ธมฺมํ ปฏิจฺจ อกุสโล ธมฺโม อุปฺปชฺชติ เหตุปจฺจยา.\csromanspacer{} อกุสลํ เอกํ ขนฺธํ ปฏิจฺจ ตโย ขนฺธา ฯเปฯ เทฺว ขนฺเธ ปฏิจฺจ เทฺว ขนฺธา.\csromanspacer{}\csromanspacer{}นกุสลํ ธมฺมํ ปฏิจฺจ อพฺยากโต ธมฺโม อุปฺปชฺชติ เหตุปจฺจยา.\csromanspacer{} วิปากาพฺยากตํ กิริยาพฺยากตํ เอกํ ขนฺธํ ปฏิจฺจ ตโย ขนฺธา จิตฺตสมุฏฺฐานญฺจ รูปํ ฯเปฯ.\csromanspacer{}\csromanspacer{}นกุสลํ ธมฺมํ ปฏิจฺจ อกุสโล จ อพฺยากโต จ ธมฺมา อุปฺปชฺชนฺติ เหตุปจฺจยา.\csromanspacer{} ตีณิ.}

\prose{น-อกุสลํ ธมฺมํ ปฏิจฺจ กุสโล ธมฺโม อุปฺปชฺชติ เหตุปจฺจยา.\csromanspacer{}\csromanspacer{}น-อกุสลํ ธมฺมํ ปฏิจฺจ อพฺยากโต ธมฺโม อุปฺปชฺชติ เหตุปจฺจยา.\csromanspacer{}\csromanspacer{}น-อกุสลํ ธมฺมํ ปฏิจฺจ กุสโล จ อพฺยากโต จ ธมฺมา อุปฺปชฺชนฺติ เหตุปจฺจยา.\csromanspacer{} ตีณิ.}

\prose{น-อพฺยากตํ ธมฺมํ ปฏิจฺจ อพฺยากโต ธมฺโม อุปฺปชฺชติ เหตุปจฺจยา.\csromanspacer{}\csromanspacer{}น-อพฺยากตํ ธมฺมํ ปฏิจฺจ กุสโล ธมฺโม อุปฺปชฺชติ เหตุปจฺจยา.\csromanspacer{}\csromanspacer{}น-อพฺยากตํ ธมฺมํ ปฏิจฺจ อกุสโล ธมฺโม อุปฺปชฺชติ เหตุปจฺจยา.\csromanspacer{}\csromanspacer{}น-อพฺยากตํ ธมฺมํ ปฏิจฺจ กุสโล จ อพฺยากโต จ ธมฺมา อุปฺปชฺชนฺติ เหตุปจฺจยา.\csromanspacer{}\csromanspacer{}น-อพฺยากตํ ธมฺมํ ปฏิจฺจ อกุสโล จ อพฺยากโต จ ธมฺมา อุปฺปชฺชนฺติ เหตุปจฺจยา.\csromanspacer{} ปญฺจ.}

\gambhira{ธมฺมปจฺจนียานุโลม ติกปฏฺฐานปาฬิ}
\prose{\csromanfolio{328}นกุสลญฺจ น-อพฺยากตญฺจ ธมฺมํ ปฏิจฺจ อกุสโล ธมฺโม อุปฺปชฺชติ เหตุปจฺจยา.\csromanspacer{}\csromanspacer{}นกุสลญฺจ น-อพฺยากตญฺจ ธมฺมํ ปฏิจฺจ อพฺยากโต ธมฺโม อุปฺปชฺชติ เหตุปจฺจยา.\csromanspacer{}\csromanspacer{}นกุสลญฺจ น-อพฺยากตญฺจ ธมฺมํ ปฏิจฺจ อกุสโล จ อพฺยากโต จ ธมฺมา อุปฺปชฺชนฺติ เหตุปจฺจยา.\csromanspacer{} ตีณิ.}

\prose{น-อกุสลญฺจ น-อพฺยากตญฺจ ธมฺมํ ปฏิจฺจ กุสโล ธมฺโม อุปฺปชฺชติ เหตุปจฺจยา.\csromanspacer{}\csromanspacer{}น-อกุสลญฺจ น-อพฺยากตญฺจ ธมฺมํ ปฏิจฺจ อพฺยากโต ธมฺโม อุปฺปชฺชติ เหตุปจฺจยา.\csromanspacer{}\csromanspacer{}น-อกุสลญฺจ น-อพฺยากตญฺจ ธมฺมํ ปฏิจฺจ กุสโล จ อพฺยากโต จ ธมฺมา อุปฺปชฺชนฺติ เหตุปจฺจยา.\csromanspacer{} ตีณิ.}

\prose{นกุสลญฺจ น-อกุสลญฺจ ธมฺมํ ปฏิจฺจ อพฺยากโต ธมฺโม อุปฺปชฺชติ เหตุปจฺจยา.\csromanspacer{} เอกํ.}

\proseitem{2}{นกุสลํ ธมฺมํ ปฏิจฺจ อกุสโล ธมฺโม อุปฺปชฺชติ อารมฺมณปจฺจยา.\csromanspacer{}\csromanspacer{}นกุสลํ ธมฺมํ ปฏิจฺจ อพฺยากโต ธมฺโม อุปฺปชฺชติ อารมฺมณปจฺจยา.\csromanspacer{} เทฺว.}

\prose{น-อกุสลํ ธมฺมํ ปฏิจฺจ กุสโล ธมฺโม อุปฺปชฺชติ อารมฺมณปจฺจยา.\csromanspacer{}\csromanspacer{}น-อกุสลํ ธมฺมํ ปฏิจฺจ อพฺยากโต ธมฺโม อุปฺปชฺชติ อารมฺมณปจฺจยา.\csromanspacer{} เทฺว.}

\prose{น-อพฺยากตํ ธมฺมํ ปฏิจฺจ กุสโล ธมฺโม อุปฺปชฺชติ อารมฺมณปจฺจยา.\csromanspacer{}\csromanspacer{}น-อพฺยากตํ ธมฺมํ ปฏิจฺจ อกุสโล ธมฺโม อุปฺปชฺชติ อารมฺมณปจฺจยา.\csromanspacer{} เทฺว.}

\prose{นกุสลญฺจ น-อพฺยากตญฺจ ธมฺมํ ปฏิจฺจ อกุสโล ธมฺโม อุปฺปชฺชติ อารมฺมณปจฺจยา.\csromanspacer{} เอกํ.}

\prose{น-อกุสลญฺจ น-อพฺยากตญฺจ ธมฺมํ ปฏิจฺจ กุสโล ธมฺโม อุปฺปชฺชติ อารมฺมณปจฺจยา.\csromanspacer{} เอกํ.}

\prose{นกุสลญฺจ น-อกุสลญฺจ ธมฺมํ ปฏิจฺจ อพฺยากโต ธมฺโม อุปฺปชฺชติ อารมฺมณปจฺจยา.\csromanspacer{} เอกํ.\csromanspacer{}\csromanspacer{}เหตุยา อฏฺฐารส, อารมฺมเณ นว, อธิปติยา อฏฺฐารส ฯเปฯ อวิคเต อฏฺฐารส.}

\csromanheader{2}{\textbf{ปจฺจนีย-นเหตุ น-อารมฺมณ}}
\proseitem{3}{นกุสลํ ธมฺมํ ปฏิจฺจ อกุสโล ธมฺโม อุปฺปชฺชติ นเหตุปจฺจยา.\csromanspacer{}\csromanspacer{}นกุสลํ ธมฺมํ ปฏิจฺจ อพฺยากโต ธมฺโม อุปฺปชฺชติ เหตุปจฺจยา ฯเปฯ.\csromanspacer{} นกุสลญฺจ น-อกุสลญฺจ ธมฺมํ ปฏิจฺจ อพฺยากโต ธมฺโม อุปฺปชฺชติ นเหตุปจฺจยา.}

\csromanheader{2}{1. กุสลตฺติก}
\proseitem{4}{\csromanfolio{329}นกุสลํ ธมฺมํ ปฏิจฺจ อพฺยากโต ธมฺโม อุปฺปชฺชติ น-อารมฺมณปจฺจยา ฯเปฯ.\csromanspacer{} นกุสลญฺจ น-อกุสลญฺจ ธมฺมํ ปฏิจฺจ อพฺยากโต ธมฺโม อุปฺปชฺชติ น-อารมฺมณปจฺจยา. (สํขิตฺตํ.)}

\proseitem{5}{นเหตุยา ฉ, น-อารมฺมเณ ฉ, น-อธิปติยา อฏฺฐารส ฯเปฯ โนวิคเต ฉ.}

\csromanheader{2}{เหตุปจฺจยา น-อารมฺมเณ ฉ. (สํขิตฺตํ.)}
\prose{นเหตุปจฺจยา อารมฺมเณ ฉ. (สํขิตฺตํ.) (สหชาตวาโร ปฏิจฺจวารสทิโส.)}

\csromanheader{2}{1. กุสลตฺติก 3. ปจฺจยวาร}
\proseitem{6}{นกุสลํ ธมฺมํ ปจฺจยา กุสโล ธมฺโม อุปฺปชฺชติ เหตุปจฺจยา.\csromanspacer{}\csromanspacer{}นกุสลํ ธมฺมํ ปจฺจยา อกุสโล ธมฺโม อุปฺปชฺชติ เหตุปจฺจยา.\csromanspacer{}\csromanspacer{}นกุสลํ ธมฺมํ ปจฺจยา อพฺยากโต ธมฺโม อุปฺปชฺชติ เหตุปจฺจยา.\csromanspacer{}\csromanspacer{}นกุสลํ ธมฺมํ ปจฺจยา กุสโล จ อพฺยากโต จ ธมฺมา อุปฺปชฺชนฺติ เหตุปจฺจยา.\csromanspacer{}\csromanspacer{}นกุสลํ ธมฺมํ ปจฺจยา อกุสโล จ อพฺยากโต จ ธมฺมา อุปฺปชฺชนฺติ เหตุปจฺจยา.\csromanspacer{} ปญฺจ.}

\prose{น-อกุสลํ ธมฺมํ ปจฺจยา อกุสโล ธมฺโม อุปฺปชฺชติ เหตุปจฺจยา.\csromanspacer{} ปญฺจ.}

\prose{น-อพฺยากตํ ธมฺมํ ปจฺจยา อพฺยากโต ธมฺโม อุปฺปชฺชติ เหตุปจฺจยา.\csromanspacer{} ปญฺจ.}

\prose{นกุสลญฺจ น-อพฺยากตญฺจ ธมฺมํ ปจฺจยา.\csromanspacer{} ตีณิ.}

\prose{น-อกุสลญฺจ น-อพฺยากตญฺจ ธมฺมํ ปจฺจยา.\csromanspacer{} ตีณิ.\csromanspacer{} นกุสลญฺจ น-อกุสลญฺจ ธมฺมํ ปจฺจยา กุสโล ธมฺโม อุปฺปชฺชติ เหตุปจฺจยา.\csromanspacer{} ปญฺจ. (สํขิตฺตํ.) .\csromanspacer{} เหตุยา ฉพฺพีสติ, อารมฺมเณ เตรส ฯเปฯ อวิคเต ฉพฺพีสติ.}

\csromanheader{2}{สํสฏฺฐวาเร เหตุยา นว ฯเปฯ อวิคเต นว. (สํขิตฺตํ.)}
\csromanheader{2}{ธมฺมปจฺจนียานุโลม ติกปฏฺฐานปาฬิ \textbf{1. กุสลตฺติก 7. ปญฺหาวาร}}
\proseitem{7}{\csromanfolio{330}นกุสโล ธมฺโม อกุสลสฺส ธมฺมสฺส เหตุปจฺจเยน ปจฺจโย.\csromanspacer{}\csromanspacer{}นกุสโล ธมฺโม อพฺยากตสฺส ธมฺมสฺส เหตุปจฺจเยน ปจฺจโย.\csromanspacer{}\csromanspacer{}นกุสโล ธมฺโม อกุสลสฺส จ อพฺยากตสฺส จ ธมฺมสฺส เหตุปจฺจเยน ปจฺจโย.\csromanspacer{} ตีณิ.}

\prose{น-อกุสโล ธมฺโม กุสลสฺส ธมฺมสฺส เหตุปจฺจเยน ปจฺจโย.\csromanspacer{} ตีณิ.}

\prose{น-อพฺยากโต ธมฺโม อพฺยากตสฺส ธมฺมสฺส เหตุปจฺจเยน ปจฺจโย.\csromanspacer{} ปญฺจ.}

\prose{นกุสโล จ น-อพฺยากโต จ ธมฺมา อกุสลสฺส ธมฺมสฺส เหตุปจฺจเยน ปจฺจโย.\csromanspacer{} ตีณิ.}

\prose{น-อกุสโล จ น-อพฺยากโต จ ธมฺมา กุสลสฺส ธมฺมสฺส เหตุปจฺจเยน ปจฺจโย.\csromanspacer{} ตีณิ.}

\prose{นกุสโล จ น-อกุสโล จ ธมฺมา อพฺยากตสฺส ธมฺมสฺส เหตุปจฺจเยน ปจฺจโย.\csromanspacer{} เอกํ.}

\csromanheader{2}{\textbf{อารมฺมณ}}
\proseitem{8}{นกุสโล ธมฺโม กุสลสฺส ธมฺมสฺส อารมฺมณปจฺจเยน ปจฺจโย.\csromanspacer{} ตีณิ.}

\prose{น-อกุสโล ธมฺโม อกุสลสฺส ธมฺมสฺส อารมฺมณปจฺจเยน ปจฺจโย.\csromanspacer{} ตีณิ.}

\prose{น-อพฺยากโต ธมฺโม อพฺยากตสฺส ธมฺมสฺส อารมฺมณปจฺจเยน ปจฺจโย.\csromanspacer{} ตีณิ.}

\prose{นกุสโล จ น-อพฺยากโต จ ธมฺมา กุสลสฺส ธมฺมสฺส อารมฺมณปจฺจเยน ปจฺจโย.\csromanspacer{}\csromanspacer{}นกุสโล จ น-อพฺยากโต จ ธมฺมา อกุสลสฺส ธมฺมสฺส อารมฺมณปจฺจเยน ปจฺจโย.\csromanspacer{}\csromanspacer{}นกุสโล จ น-อพฺยากโต จ ธมฺมา อพฺยากตสฺส ธมฺมสฺส อารมฺมณปจฺจเยน ปจฺจโย.\csromanspacer{} ตีณิ.}

\prose{น-อกุสโล จ น-อพฺยากโต จ ธมฺมา กุสลสฺส ธมฺมสฺส อารมฺมณปจฺจเยน ปจฺจโย.\csromanspacer{}\csromanspacer{}น-อกุสโล จ น-อพฺยากโต จ ธมฺมา อกุสลสฺส ธมฺมสฺส อารมฺมณปจฺจเยน ปจฺจโย.\csromanspacer{}\csromanspacer{}น-อกุสโล จ น-อพฺยากโต จ ธมฺมา อพฺยากตสฺส ธมฺมสฺส อารมฺมณปจฺจเยน ปจฺจโย.\csromanspacer{} ตีณิ.}

\csromanheader{2}{3. วิปากตฺติก}
\prose{\csromanfolio{331}นกุสโล จ น-อกุสโล จ ธมฺมา กุสลสฺส ธมฺมสฺส อารมฺมณปจฺจเยน ปจฺจโย.\csromanspacer{} ตีณิ. (สํขิตฺตํ.)}

\proseitem{9}{เหตุยา อฏฺฐารส, อารมฺมเณ อฏฺฐารส, อธิปติยา เตวีส ฯเปฯ อวิคเต ทฺวาวีส. (ปญฺหาวารํ วิตฺถาเรตพฺพํ.)}

\csromanheader{2}{2. เวทนาตฺติก 1. ปฏิจฺจวาร}
\proseitem{10}{นสุขาย เวทนาย สมฺปยุตฺตํ ธมฺมํ ปฏิจฺจ สุขาย เวทนาย สมฺปยุตฺโต ธมฺโม อุปฺปชฺชติ เหตุปจฺจยา.\csromanspacer{}\csromanspacer{}นสุขาย เวทนาย สมฺปยุตฺตํ ธมฺมํ ปฏิจฺจ ทุกฺขาย เวทนาย สมฺปยุตฺโต ธมฺโม อุปฺปชฺชติ เหตุปจฺจยา.\csromanspacer{}\csromanspacer{}นสุขาย เวทนาย สมฺปยุตฺตํ ธมฺมํ ปฏิจฺจ อทุกฺขมสุขาย เวทนาย สมฺปยุตฺโต ธมฺโม อุปฺปชฺชติ เหตุปจฺจยา. (3)}

\prose{นทุกฺขาย เวทนาย สมฺปยุตฺตํ ธมฺมํ ปฏิจฺจ ทุกฺขาย เวทนาย สมฺปยุตฺโต ธมฺโม อุปฺปชฺชติ เหตุปจฺจยา.\csromanspacer{}\csromanspacer{}นทุกฺขาย เวทนาย สมฺปยุตฺตํ ธมฺมํ ปฏิจฺจ สุขาย เวทนาย สมฺปยุตฺโต ธมฺโม อุปฺปชฺชติ เหตุปจฺจยา.\csromanspacer{}\csromanspacer{}นทุกฺขาย เวทนาย สมฺปยุตฺตํ ธมฺมํ ปฏิจฺจ อทุกฺขมสุขาย เวทนาย สมฺปยุตฺโต ธมฺโม อุปฺปชฺชติ เหตุปจฺจยา. (3)}

\prose{น-อทุกฺขมสุขาย เวทนาย สมฺปยุตฺตํ ธมฺมํ ปฏิจฺจ อทุกฺขมสุขาย เวทนาย สมฺปยุตฺโต ธมฺโม อุปฺปชฺชติ เหตุปจฺจยา.\csromanspacer{}\csromanspacer{}น-อทุกฺขมสุขาย เวทนาย สมฺปยุตฺตํ ธมฺมํ ปฏิจฺจ สุขาย เวทนาย สมฺปยุตฺโต ธมฺโม อุปฺปชฺชติ เหตุปจฺจยา.\csromanspacer{}\csromanspacer{}น-อทุกฺขมสุขาย เวทนาย สมฺปยุตฺตํ ธมฺมํ ปฏิจฺจ ทุกฺขาย เวทนาย สมฺปยุตฺโต ธมฺโม อุปฺปชฺชติ เหตุปจฺจยา. (3) (สํขิตฺตํ.) .\csromanspacer{} เหตุยา เอกวีส, อธิปติยา เอกวีส ฯเปฯ อวิคเต เอกวีส.}

\csromanheader{2}{3. วิปากตฺติก 1. ปฏิจฺจวาร}
\proseitem{11}{นวิปากํ ธมฺมํ ปฏิจฺจ วิปาโก ธมฺโม อุปฺปชฺชติ เหตุปจฺจยา.\csromanspacer{}\csromanspacer{}นวิปากํ ธมฺมํ ปฏิจฺจ วิปากธมฺมธมฺโม อุปฺปชฺชติ เหตุปจฺจยา.\csromanspacer{} ปญฺจ.}

\prose{นวิปากธมฺมธมฺมํ ปฏิจฺจ วิปาโก ธมฺโม อุปฺปชฺชติ เหตุปจฺจยา.\csromanspacer{}\csromanspacer{}นวิปากธมฺมธมฺมํ ปฏิจฺจ เนววิปากนวิปากธมฺมธมฺโม อุปฺปชฺชติ เหตุปจฺจยา.\csromanspacer{} ตีณิ.}

\gambhira{ธมฺมปจฺจนียานุโลม ติกปฏฺฐานปาฬิ}
\prose{\csromanfolio{332}นเนววิปากนวิปากธมฺมธมฺมํ ปฏิจฺจ เนววิปากนวิปากธมฺมธมฺโม อุปฺปชฺชติ เหตุปจฺจยา.\csromanspacer{}\csromanspacer{}นเนววิปากนวิปากธมฺมธมฺมํ ปฏิจฺจ วิปาโก ธมฺโม อุปฺปชฺชติ เหตุปจฺจยา.\csromanspacer{} ปญฺจ.\csromanspacer{}\csromanspacer{}เหตุยา ทฺวาวีส.}

\csromanheader{2}{4. อุปาทินฺนตฺติก 1. ปฏิจฺจวาร}
\proseitem{12}{น-อุปาทินฺนุปาทานิยํ ธมฺมํ ปฏิจฺจ อนุปาทินฺนุปาทานิโย ธมฺโม อุปฺปชฺชติ เหตุปจฺจยา.\csromanspacer{} ตีณิ.}

\prose{น-อนุปาทินฺนุปาทานิยํ ธมฺมํ ปฏิจฺจ อนุปาทินฺนุปาทานิโย ธมฺโม อุปฺปชฺชติ เหตุปจฺจยา.\csromanspacer{} ปญฺจ.}

\prose{น-อนุปาทินฺน-อนุปาทานิยํ ธมฺมํ ปฏิจฺจ อุปาทินฺนุปาทานิโย ธมฺโม อุปฺปชฺชติ เหตุปจฺจยา.\csromanspacer{} ตีณิ.}

\prose{น-อุปาทินฺนุปาทานิยญฺจ น-อนุปาทินฺน-อนุปาทานิยญฺจ ธมฺมํ ปฏิจฺจ อนุปาทินฺนุปาทานิโย ธมฺโม อุปฺปชฺชติ เหตุปจฺจยา.\csromanspacer{} เอกํ.}

\prose{น-อนุปาทินฺนุปาทานิยญฺจ น-อนุปาทินฺน-อนุปาทานิยญฺจ ธมฺมํ ปฏิจฺจ อุปาทินฺนุปาทานิโย ธมฺโม อุปฺปชฺชติ เหตุปจฺจยา.\csromanspacer{} ตีณิ.}

\prose{น-อุปาทินฺนุปาทานิยญฺจ น-อนุปาทินฺนุปาทานิยญฺจ ธมฺมํ ปฏิจฺจ อนุปาทินฺน-อนุปาทานิโย ธมฺโม อุปฺปชฺชติ เหตุปจฺจยา.\csromanspacer{} ตีณิ.\csromanspacer{}\csromanspacer{}เหตุยา อฏฺฐารส.}

\csromanheader{2}{5. สํกิลิฏฺฐตฺติก 1. ปฏิจฺจวาร}
\proseitem{13}{นสํกิลิฏฺฐสํกิเลสิกํ ธมฺมํ ปฏิจฺจ อสํกิลิฏฺฐสํกิเลสิโก ธมฺโม อุปฺปชฺชติ เหตุปจฺจยา.\csromanspacer{}\csromanspacer{}นสํกิลิฏฺฐสํกิเลสิกํ ธมฺมํ ปฏิจฺจ อสํกิลิฏฺฐ-อสํกิเลสิโก ธมฺโม อุปฺปชฺชติ เหตุปจฺจยา.\csromanspacer{}\csromanspacer{}นสํกิลิฏฺฐสํกิเลสิกํ ธมฺมํ ปฏิจฺจ อสํกิลิฏฺฐสํกิเลสิโก จ อสํกิลิฏฺฐ-อสํกิเลสิโก จ ธมฺมา อุปฺปชฺชนฺติ เหตุปจฺจยา.\csromanspacer{} ตีณิ.}

\prose{น-อสํกิลิฏฺฐสํกิเลสิกํ ธมฺมํ ปฏิจฺจ อสํกิลิฏฺฐสํกิเลสิโก ธมฺโม อุปฺปชฺชติ เหตุปจฺจยา.\csromanspacer{} ปญฺจ.}

\prose{น-อสํกิลิฏฺฐ-อสํกิเลสิกํ ธมฺมํ ปฏิจฺจ สํกิลิฏฺฐสํกิเลสิโก ธมฺโม อุปฺปชฺชติ เหตุปจฺจยา.\csromanspacer{}\csromanspacer{}น-อสํกิลิฏฺฐ-อสํกิเลสิกํ ธมฺมํ ปฏิจฺจ}

\vspace{\tipitakaparskip}
\csromancenter{ทสฺสนตฺติก อสํกิลิฏฺฐสํกิเลสิโก ธมฺโม อุปฺปชฺชติ เหตุปจฺจยา.\csromanspacer{}\csromanspacer{}น-อสํกิลิฏฺฐ-อสํกิเลสิกํ ธมฺมํ ปฏิจฺจ สํกิลิฏฺฐสํกิเลสิโก จ อสํกิลิฏฺฐสํกิเลสิโก จ ธมฺมา อุปฺปชฺชนฺติ เหตุปจฺจยา.\csromanspacer{} ตีณิ.\csromanspacer{} เหตุยา อฏฺฐารส. \textbf{6.}\csromanspacer{}\textbf{ วิตกฺกตฺติก 1.}\csromanspacer{}\textbf{ ปฏิจฺจวาร}}
\proseitem{14}{\csromanfolio{333}นสวิตกฺกสวิจารํ ธมฺมํ ปฏิจฺจ สวิตกฺกสวิจาโร ธมฺโม อุปฺปชฺชติ เหตุปจฺจยา.\csromanspacer{}\csromanspacer{}นสวิตกฺกสวิจารํ ธมฺมํ ปฏิจฺจ อวิตกฺกวิจารมตฺโต ธมฺโม อุปฺปชฺชติ เหตุปจฺจยา.\csromanspacer{}\csromanspacer{}นสวิตกฺกสวิจารํ ธมฺมํ ปฏิจฺจ อวิตกฺก-อวิจาโร ธมฺโม อุปฺปชฺชติ เหตุปจฺจยา.\csromanspacer{} สตฺต.}

\prose{น-อวิตกฺกวิจารมตฺตํ ธมฺมํ ปฏิจฺจ อวิตกฺกวิจารมตฺโต ธมฺโม อุปฺปชฺชติ เหตุปจฺจยา.\csromanspacer{} สตฺต.}

\prose{น-อวิตกฺก-อวิจารํ ธมฺมํ ปฏิจฺจ อวิตกฺก-อวิจาโร ธมฺโม อุปฺปชฺชติ เหตุปจฺจยา.\csromanspacer{} สตฺต.\csromanspacer{}\csromanspacer{}เหตุยา เอกูนปญฺญาส. \textbf{7.}\csromanspacer{}\textbf{ ปีติตฺติก 1.}\csromanspacer{}\textbf{ ปฏิจฺจวาร}}

\proseitem{15}{นปีติสหคตํ ธมฺมํ ปฏิจฺจ ปีติสหคโต ธมฺโม อุปฺปชฺชติ เหตุปจฺจยา.\csromanspacer{} จตฺตาริ.}

\prose{นสุขสหคตํ ธมฺมํ ปฏิจฺจ สุขสหคโต ธมฺโม อุปฺปชฺชติ เหตุปจฺจยา.\csromanspacer{} จตฺตาริ.}

\prose{น-อุเปกฺขาสหคตํ ธมฺมํ ปฏิจฺจ อุเปกฺขาสหคโต ธมฺโม อุปฺปชฺชติ เหตุปจฺจยา.\csromanspacer{}\csromanspacer{}น-อุเปกฺขาสหคตํ ธมฺมํ ปฏิจฺจ ปีติสหคโต ธมฺโม อุปฺปชฺชติ เหตุปจฺจยา.\csromanspacer{} จตฺตาริ.\csromanspacer{}\csromanspacer{}เหตุยา อฏฺฐวีส. \textbf{8.}\csromanspacer{}\textbf{ ทสฺสนตฺติก 1.}\csromanspacer{}\textbf{ ปฏิจฺจวาร}}

\proseitem{16}{นทสฺสเนน ปหาตพฺพํ ธมฺมํ ปฏิจฺจ ภาวนาย ปหาตพฺโพ ธมฺโม อุปฺปชฺชติ เหตุปจฺจยา.\csromanspacer{}\csromanspacer{}นทสฺสเนน ปหาตพฺพํ ธมฺมํ ปฏิจฺจ เนวทสฺสเนน นภาวนาย ปหาตพฺโพ ธมฺโม อุปฺปชฺชติ เหตุปจฺจยา.\csromanspacer{}\csromanspacer{}นทสฺสเนน ปหาตพฺพํ ธมฺมํ ปฏิจฺจ ภาวนาย ปหาตพฺโพ จ เนวทสฺสเนน นภาวนาย ปหาตพฺโพ จ ธมฺมา อุปฺปชฺชนฺติ เหตุปจฺจยา.\csromanspacer{} ตีณิ.}

\gambhira{ธมฺมปจฺจนียานุโลม ติกปฏฺฐานปาฬิ}
\prose{\csromanfolio{334}นภาวนาย ปหาตพฺพํ ธมฺมํ ปฏิจฺจ ทสฺสเนน ปหาตพฺโพ ธมฺโม อุปฺปชฺชติ เหตุปจฺจยา.\csromanspacer{}\csromanspacer{}นภาวนาย ปหาตพฺพํ ธมฺมํ ปฏิจฺจ เนวทสฺสเนน นภาวนาย ปหาตพฺโพ ธมฺโม อุปฺปชฺชติ เหตุปจฺจยา.\csromanspacer{}\csromanspacer{}นภาวนาย ปหาตพฺพํ ธมฺมํ ปฏิจฺจ ทสฺสเนน ปหาตพฺโพ จ เนวทสฺสเนน นภาวนาย ปหาตพฺโพ จ ธมฺมา อุปฺปชฺชนฺติ เหตุปจฺจยา.\csromanspacer{} ตีณิ.}

\prose{นเนวทสฺสเนน นภาวนาย ปหาตพฺพํ ธมฺมํ ปฏิจฺจ เนวทสฺสเนน นภาวนาย ปหาตพฺโพ ธมฺโม อุปฺปชฺชติ เหตุปจฺจยา.\csromanspacer{}\csromanspacer{}นเนวทสฺสเนน นภาวนาย ปหาตพฺพํ ธมฺมํ ปฏิจฺจ ทสฺสเนน ปหาตพฺโพ ธมฺโม อุปฺปชฺชติ เหตุปจฺจยา.\csromanspacer{} ปญฺจ. (สํขิตฺตํ.) .\csromanspacer{} เหตุยา อฏฺฐารส.}

\csromanheader{2}{9. ทสฺสนเหตุตฺติก 1. ปฏิจฺจวาร}
\proseitem{17}{นทสฺสเนน ปหาตพฺพเหตุกํ ธมฺมํ ปฏิจฺจ ทสฺสเนน ปหาตพฺพเหตุโก ธมฺโม อุปฺปชฺชติ เหตุปจฺจยา.\csromanspacer{}\csromanspacer{}เหตุยา ฉพฺพีส.}

\csromanheader{2}{10. อาจยคามิตฺติก 1. ปฏิจฺจวาร}
\proseitem{18}{น-อาจยคามิํ ธมฺมํ ปฏิจฺจ อปจยคามี ธมฺโม อุปฺปชฺชติ เหตุปจฺจยา.\csromanspacer{}\csromanspacer{}เหตุยา อฏฺฐารส.}

\csromanheader{2}{11. เสกฺขตฺติก 1. ปฏิจฺจวาร}
\proseitem{19}{นเสกฺขํ ธมฺมํ ปฏิจฺจ อเสกฺโข ธมฺโม อุปฺปชฺชติ เหตุปจฺจยา.\csromanspacer{}\csromanspacer{}เหตุยา อฏฺฐารส.}

\csromanheader{2}{12-16. ปริตฺตตฺติกาทิ 1. ปฏิจฺจวาร}
\proseitem{20}{นปริตฺตํ ธมฺมํ ปฏิจฺจ ปริตฺโต ธมฺโม อุปฺปชฺชติ เหตุปจฺจยา.\csromanspacer{} เหตุยา ทฺวาวีส.}

\proseitem{21}{นปริตฺตารมฺมณํ ธมฺมํ ปฏิจฺจ ปริตฺตารมฺมโณ ธมฺโม อุปฺปชฺชติ เหตุปจฺจยา.\csromanspacer{}\csromanspacer{}เหตุยา เตรส.}

\proseitem{22}{นหีนํ ธมฺมํ ปฏิจฺจ มชฺฌิโม ธมฺโม อุปฺปชฺชติ เหตุปจฺจยา.\csromanspacer{} เหตุยา อฏฺฐารส.}

\csromanheader{2}{20-21. อชฺฌตฺตตฺติกทฺวย}
\proseitem{23}{\csromanfolio{335}นมิจฺฉตฺตนิยตํ ธมฺมํ ปฏิจฺจ สมฺมตฺตนิยโต ธมฺโม อุปฺปชฺชติ เหตุปจฺจยา.\csromanspacer{}\csromanspacer{}เหตุยา อฏฺฐารส.}

\proseitem{24}{นมคฺคารมฺมณํ ธมฺมํ ปฏิจฺจ มคฺคเหตุโก ธมฺโม อุปฺปชฺชติ เหตุปจฺจยา.\csromanspacer{}\csromanspacer{}เหตุยา ทส\footnote{เหตุยา อฏฺฐ?}.}

\csromanheader{2}{17. อุปฺปนฺนตฺติก 7. ปญฺหาวาร}
\proseitem{25}{น-อนุปฺปนฺโน ธมฺโม อุปฺปนฺนสฺส ธมฺมสฺส เหตุปจฺจเยน ปจฺจโย.\csromanspacer{}\csromanspacer{}น-อุปฺปาที ธมฺโม อุปฺปนฺนสฺส ธมฺมสฺส เหตุปจฺจเยน ปจฺจโย.\csromanspacer{}\csromanspacer{}น-อนุปฺปนฺโน จ น-อุปฺปาที จ ธมฺมา อุปฺปนฺนสฺส ธมฺมสฺส เหตุปจฺจเยน ปจฺจโย.\csromanspacer{} เหตุยา ตีณิ.}

\csromanheader{2}{18-19. อตีตตฺติกทฺวย 1-7. ปฏิจฺจาทิวาร}
\proseitem{26}{น-อตีโต ธมฺโม ปจฺจุปฺปนฺนสฺส ธมฺมสฺส เหตุปจฺจเยน ปจฺจโย.\csromanspacer{}\csromanspacer{}น-อนาคโต ธมฺโม ปจฺจุปฺปนฺนสฺส ธมฺมสฺส เหตุปจฺจเยน ปจฺจโย.\csromanspacer{}\csromanspacer{}น-อตีโต จ น-อนาคโต จ ธมฺมา ปจฺจุปฺปนฺนสฺส ธมฺมสฺส เหตุปจฺจเยน ปจฺจโย.\csromanspacer{}\csromanspacer{}เหตุยา ตีณิ.}

\proseitem{27}{น-อตีตารมฺมณํ ธมฺมํ ปฏิจฺจ อตีตารมฺมโณ ธมฺโม อุปฺปชฺชติ เหตุปจฺจยา.\csromanspacer{}\csromanspacer{}น-อตีตารมฺมณํ ธมฺมํ ปฏิจฺจ อนาคตารมฺมโณ ธมฺโม อุปฺปชฺชติ เหตุปจฺจยา.\csromanspacer{}\csromanspacer{}น-อตีตารมฺมณํ ธมฺมํ ปฏิจฺจ ปจฺจุปฺปนฺนารมฺมโณ ธมฺโม อุปฺปชฺชติ เหตุปจฺจยา.\csromanspacer{} ตีณิ.}

\prose{น-อนาคตารมฺมณํ ธมฺมํ ปฏิจฺจ อตีตารมฺมโณ ธมฺโม อุปฺปชฺชติ เหตุปจฺจยา.\csromanspacer{}\csromanspacer{}น-อนาคตารมฺมณํ ธมฺมํ ปฏิจฺจ ปจฺจุปฺปนฺนารมฺมโณ ธมฺโม อุปฺปชฺชติ เหตุปจฺจยา.\csromanspacer{} เทฺว.}

\prose{นปจฺจุปฺปนฺนารมฺมณํ ธมฺมํ ปฏิจฺจ ปจฺจุปฺปนฺนารมฺมโณ ธมฺโม อุปฺปชฺชติ เหตุปจฺจยา.\csromanspacer{}\csromanspacer{}นปจฺจุปฺปนฺนารมฺมณํ ธมฺมํ ปฏิจฺจ อตีตารมฺมโณ ธมฺโม อุปฺปชฺชติ เหตุปจฺจยา.\csromanspacer{}\csromanspacer{}นปจฺจุปฺปนฺนารมฺมณํ ธมฺมํ ปฏิจฺจ อนาคตารมฺมโณ ธมฺโม อุปฺปชฺชติ เหตุปจฺจยา.\csromanspacer{} ตีณิ. (สํขิตฺตํ.) .\csromanspacer{} เหตุยา สตฺตรส.}

\csromanheader{2}{20-21. อชฺฌตฺตตฺติกทฺวย 1. ปฏิจฺจวาร}
\proseitem{28}{น-อชฺฌตฺตํ ธมฺมํ ปฏิจฺจ พหิทฺธา ธมฺโม อุปฺปชฺชติ เหตุปจฺจยา.\csromanspacer{} เอกํ.}

\gambhira{ธมฺมปจฺจนียานุโลม ติกปฏฺฐานปาฬิ}
\prose{\csromanfolio{336}นพหิทฺธา ธมฺมํ ปฏิจฺจ อชฺฌตฺโต ธมฺโม อุปฺปชฺชติ เหตุปจฺจยา.\csromanspacer{} เอกํ.\csromanspacer{} เหตุยา เทฺว.}

\proseitem{29}{น-อชฺฌตฺตารมฺมณํ ธมฺมํ ปฏิจฺจ อชฺฌตฺตารมฺมโณ ธมฺโม อุปฺปชฺชติ เหตุปจฺจยา.\csromanspacer{} เทฺว.}

\prose{นพหิทฺธารมฺมณํ ธมฺมํ ปฏิจฺจ พหิทฺธารมฺมโณ ธมฺโม อุปฺปชฺชติ เหตุปจฺจยา.\csromanspacer{} เทฺว.}

\prose{น-อชฺฌตฺตารมฺมณญฺจ นพหิทฺธารมฺมณญฺจ ธมฺมํ ปฏิจฺจ อชฺฌตฺตารมฺมโณ ธมฺโม อุปฺปชฺชติ เหตุปจฺจยา.\csromanspacer{} เทฺว.\csromanspacer{}\csromanspacer{}เหตุยา ฉ.}

\csromanheader{2}{22. สนิทสฺสนตฺติก 1. ปฏิจฺจวาร}
\proseitem{30}{นสนิทสฺสนสปฺปฏิฆํ ธมฺมํ ปฏิจฺจ สนิทสฺสนสปฺปฏิโฆ ธมฺโม อุปฺปชฺชติ เหตุปจฺจยา.\csromanspacer{}\csromanspacer{}นสนิทสฺสนสปฺปฏิฆํ ธมฺมํ ปฏิจฺจ อนิทสฺสนสปฺปฏิโฆ ธมฺโม อุปฺปชฺชติ เหตุปจฺจยา.\csromanspacer{}\csromanspacer{}นสนิทสฺสนสปฺปฏิฆํ ธมฺมํ ปฏิจฺจ อนิทสฺสน-อปฺปฏิโฆ ธมฺโม อุปฺปชฺชติ เหตุปจฺจยา.\csromanspacer{}\csromanspacer{}นสนิทสฺสนสปฺปฏิฆํ ธมฺมํ ปฏิจฺจ สนิทสฺสนสปฺปฏิโฆ จ อนิทสฺสน-อปฺปฏิโฆ จ ธมฺมา อุปฺปชฺชนฺติ เหตุปจฺจยา.\csromanspacer{}\csromanspacer{}นสนิทสฺสนสปฺปฏิฆํ ธมฺมํ ปฏิจฺจ อนิทสฺสนสปฺปฏิโฆ จ อนิทสฺสน-อปฺปฏิโฆ จ ธมฺมา อุปฺปชฺชนฺติ เหตุปจฺจยา.\csromanspacer{}\csromanspacer{}นสนิทสฺสนสปฺปฏิฆํ ธมฺมํ ปฏิจฺจ สนิทสฺสนสปฺปฏิโฆ จ อนิทสฺสนสปฺปฏิโฆ จ ธมฺมา อุปฺปชฺชนฺติ เหตุปจฺจยา.\csromanspacer{}\csromanspacer{}นสนิทสฺสนสปฺปฏิฆํ ธมฺมํ ปฏิจฺจ สนิทสฺสนสปฺปฏิโฆ จ อนิทสฺสนสปฺปฏิโฆ จ อนิทสฺสน-อปฺปฏิโฆ จ ธมฺมา อุปฺปชฺชนฺติ เหตุปจฺจยา.\csromanspacer{} สตฺต.}

\prose{น-อนิทสฺสนสปฺปฏิฆํ ธมฺมํ ปฏิจฺจ อนิทสฺสนสปฺปฏิโฆ ธมฺโม อุปฺปชฺชติ เหตุปจฺจยา.\csromanspacer{}\csromanspacer{}น-อนิทสฺสนสปฺปฏิฆํ ธมฺมํ ปฏิจฺจ สนิทสฺสนสปฺปฏิโฆ ธมฺโม อุปฺปชฺชติ เหตุปจฺจยา.\csromanspacer{} สตฺต.}

\prose{น-อนิทสฺสน-อปฺปฏิฆํ ธมฺมํ ปฏิจฺจ อนิทสฺสน-อปฺปฏิโฆ ธมฺโม อุปฺปชฺชติ เหตุปจฺจยา.\csromanspacer{} สตฺต.\csromanspacer{}\csromanspacer{}เหตุยา ปญฺจติํส. (ปญฺหาวารมฺปิ วิตฺถาเรตพฺพํ.)}

\nitthitam{\textbf{ธมฺมปจฺจนียานุโลเม ติกปฏฺฐานํ นิฏฺฐิตํ.}}
\pitaka{\textbf{อภิธมฺมปิฏก}}
\csromanheader{2}{\textbf{ธมฺมปจฺจนียานุโลม}}
\gambhira{\textbf{ทุกปฏฺฐานปาฬิ}}
\namakkaram{นโม ตสฺส ภควโต อรหโต สมฺมาสมฺพุทฺธสฺส.}
\csromanheader{2}{1-6. เหตุโคจฺฉก 1. ปฏิจฺจวาร}
\proseitem{1}{\csromanfolio{337}นเหตุํ ธมฺมํ ปฏิจฺจ เหตุ ธมฺโม อุปฺปชฺชติ เหตุปจฺจยา.\csromanspacer{}\csromanspacer{}นเหตุํ ธมฺมํ ปฏิจฺจ นเหตุ ธมฺโม อุปฺปชฺชติ เหตุปจฺจยา.\csromanspacer{}\csromanspacer{}นเหตุํ ธมฺมํ ปฏิจฺจ เหตุ จ นเหตุ จ ธมฺมา อุปฺปชฺชนฺติ เหตุปจฺจยา. (3)}

\prose{นนเหตุํ ธมฺมํ ปฏิจฺจ เหตุ ธมฺโม อุปฺปชฺชติ เหตุปจฺจยา.\csromanspacer{}\csromanspacer{}นนเหตุํ ธมฺมํ ปฏิจฺจ นเหตุ ธมฺโม อุปฺปชฺชติ เหตุปจฺจยา.\csromanspacer{}\csromanspacer{}นนเหตุํ ธมฺมํ ปฏิจฺจ เหตุ จ นเหตุ จ ธมฺมา อุปฺปชฺชนฺติ เหตุปจฺจยา. (3)}

\prose{นเหตุญฺจ นนเหตุญฺจ ธมฺมํ ปฏิจฺจ เหตุ ธมฺโม อุปฺปชฺชติ เหตุปจฺจยา.\csromanspacer{}\csromanspacer{}นเหตุญฺจ นนเหตุญฺจ ธมฺมํ ปฏิจฺจ นเหตุ ธมฺโม อุปฺปชฺชติ เหตุปจฺจยา.\csromanspacer{}\csromanspacer{}นเหตุญฺจ นนเหตุญฺจ ธมฺมํ ปฏิจฺจ เหตุ จ นเหตุ จ ธมฺมา อุปฺปชฺชนฺติ เหตุปจฺจยา. (3) (สํขิตฺตํ.) .\csromanspacer{} เหตุยา นว, อารมฺมเณ นว ฯเปฯ อวิคเต นว. (สหชาตวารมฺปิ ฯเปฯ ปญฺหาวารมฺปิ วิตฺถาเรตพฺพํ.)}

\proseitem{2}{นสเหตุกํ ธมฺมํ ปฏิจฺจ สเหตุโก ธมฺโม อุปฺปชฺชติ เหตุปจฺจยา.\csromanspacer{}\csromanspacer{}นสเหตุกํ ธมฺมํ ปฏิจฺจ อเหตุโก ธมฺโม อุปฺปชฺชติ เหตุปจฺจยา.\csromanspacer{}\csromanspacer{}นสเหตุกํ ธมฺมํ ปฏิจฺจ สเหตุโก จ อเหตุโก จ ธมฺมา อุปฺปชฺชนฺติ เหตุปจฺจยา.\csromanspacer{} ตีณิ.}

\gambhira{ธมฺมปจฺจนียานุโลม ทุกปฏฺฐานปาฬิ}
\prose{\csromanfolio{338}น-อเหตุกํ ธมฺมํ ปฏิจฺจ อเหตุโก ธมฺโม อุปฺปชฺชติ เหตุปจฺจยา.\csromanspacer{}\csromanspacer{}น-อเหตุกํ ธมฺมํ ปฏิจฺจ สเหตุโก ธมฺโม อุปฺปชฺชติ เหตุปจฺจยา.\csromanspacer{}\csromanspacer{}น-อเหตุกํ ธมฺมํ ปฏิจฺจ สเหตุโก จ อเหตุโก จ ธมฺมา อุปฺปชฺชนฺติ เหตุปจฺจยา.\csromanspacer{} ตีณิ.}

\prose{นสเหตุกญฺจ น-อเหตุกญฺจ ธมฺมํ ปฏิจฺจ สเหตุโก ธมฺโม อุปฺปชฺชติ เหตุปจฺจยา.\csromanspacer{}\csromanspacer{}นสเหตุกญฺจ น-อเหตุกญฺจ ธมฺมํ ปฏิจฺจ อเหตุโก ธมฺโม อุปฺปชฺชติ เหตุปจฺจยา.\csromanspacer{}\csromanspacer{}นสเหตุกญฺจ น-อเหตุกญฺจ ธมฺมํ ปฏิจฺจ สเหตุโก จ อเหตุโก จ ธมฺมา อุปฺปชฺชนฺติ เหตุปจฺจยา.\csromanspacer{} ตีณิ.\csromanspacer{}\csromanspacer{}เหตุยา นว.}

\proseitem{3}{นเหตุสมฺปยุตฺตํ ธมฺมํ ปฏิจฺจ เหตุสมฺปยุตฺโต ธมฺโม อุปฺปชฺชติ เหตุปจฺจยา.\csromanspacer{}\csromanspacer{}เหตุยา นว.}

\proseitem{4}{นเหตุญฺเจว น-อเหตุกญฺจ ธมฺมํ ปฏิจฺจ เหตุ เจว สเหตุโก ธมฺโม อุปฺปชฺชติ เหตุปจฺจยา.\csromanspacer{}\csromanspacer{}เหตุยา นว.}

\proseitem{5}{นเหตุญฺเจว นเหตุวิปฺปยุตฺตญฺจ ธมฺมํ ปฏิจฺจ เหตุ เจว เหตุสมฺปยุตฺโต จ ธมฺโม อุปฺปชฺชติ เหตุปจฺจยา.\csromanspacer{}\csromanspacer{}เหตุยา นว.}

\proseitem{6}{นเหตุํ นสเหตุกํ ธมฺมํ ปฏิจฺจ นเหตุ สเหตุโก ธมฺโม อุปฺปชฺชติ เหตุปจฺจยา.\csromanspacer{}\csromanspacer{}เหตุยา นว.}

\csromanheader{2}{7-13. จูฬนฺตรทุก}
\proseitem{7}{น-อปฺปจฺจยํ ธมฺมํ ปฏิจฺจ สปฺปจฺจโย ธมฺโม อุปฺปชฺชติ เหตุปจฺจยา.\csromanspacer{}\csromanspacer{}เหตุยา เอกํ.}

\proseitem{8}{น-อสงฺขตํ ธมฺมํ ปฏิจฺจ สงฺขโต ธมฺโม อุปฺปชฺชติ เหตุปจฺจยา.\csromanspacer{}\csromanspacer{}เหตุยา เอกํ.}

\proseitem{9}{นสนิทสฺสนํ ธมฺมํ ปฏิจฺจ สนิทสฺสโน ธมฺโม อุปฺปชฺชติ เหตุปจฺจยา.\csromanspacer{}\csromanspacer{}นสนิทสฺสนํ ธมฺมํ ปฏิจฺจ อนิทสฺสโน ธมฺโม อุปฺปชฺชติ เหตุปจฺจยา.\csromanspacer{}\csromanspacer{}นสนิทสฺสนํ ธมฺมํ ปฏิจฺจ สนิทสฺสโน จ อนิทสฺสโน จ ธมฺมา อุปฺปชฺชนฺติ เหตุปจฺจยา.\csromanspacer{}\csromanspacer{}เหตุยา ตีณิ.}

\csromanheader{2}{14-19. อาสวโคจฺฉก}
\prose{\csromanfolio{339}นสนิทสฺสโน ธมฺโม สนิทสฺสนสฺส ธมฺมสฺส เหตุปจฺจเยน ปจฺจโย.\csromanspacer{}\csromanspacer{}เหตุยา ตีณิ.}

\prose{นสนิทสฺสโน ธมฺโม อนิทสฺสนสฺส ธมฺมสฺส อารมฺมณปจฺจเยน ปจฺจโย.\csromanspacer{}\csromanspacer{}น-อนิทสฺสโน ธมฺโม อนิทสฺสนสฺส ธมฺมสฺส อารมฺมณปจฺจเยน ปจฺจโย.}

\proseitem{10}{นสปฺปฏิฆํ ธมฺมํ ปฏิจฺจ สปฺปฏิโฆ ธมฺโม อุปฺปชฺชติ เหตุปจฺจยา.\csromanspacer{}\csromanspacer{}นสปฺปฏิฆํ ธมฺมํ ปฏิจฺจ อปฺปฏิโฆ ธมฺโม อุปฺปชฺชติ เหตุปจฺจยา.\csromanspacer{}\csromanspacer{}นสปฺปฏิฆํ ธมฺมํ ปฏิจฺจ สปฺปฏิโฆ จ อปฺปฏิโฆ จ ธมฺมา อุปฺปชฺชนฺติ เหตุปจฺจยา.\csromanspacer{} ตีณิ.}

\prose{น-อปฺปฏิฆํ ธมฺมํ ปฏิจฺจ อปฺปฏิโฆ ธมฺโม อุปฺปชฺชติ เหตุปจฺจยา.\csromanspacer{}\csromanspacer{}น-อปฺปฏิฆํ ธมฺมํ ปฏิจฺจ สปฺปฏิโฆ ธมฺโม อุปฺปชฺชติ เหตุปจฺจยา.\csromanspacer{}\csromanspacer{}น-อปฺปฏิฆํ ธมฺมํ ปฏิจฺจ สปฺปฏิโฆ จ อปฺปฏิโฆ จ ธมฺมา อุปฺปชฺชนฺติ เหตุปจฺจยา.\csromanspacer{} ตีณิ.}

\prose{นสปฺปฏิฆญฺจ น-อปฺปฏิฆญฺจ ธมฺมํ ปฏิจฺจ สปฺปฏิโฆ ธมฺโม อุปฺปชฺชติ เหตุปจฺจยา.\csromanspacer{} ตีณิ.\csromanspacer{}\csromanspacer{}เหตุยา นว.}

\proseitem{11}{นรูปิํ ธมฺมํ ปฏิจฺจ รูปี ธมฺโม อุปฺปชฺชติ เหตุปจฺจยา.\csromanspacer{}\csromanspacer{}เหตุยา นว.}

\proseitem{12}{นโลกิยํ ธมฺมํ ปฏิจฺจ โลกิโย ธมฺโม อุปฺปชฺชติ เหตุปจฺจยา.\csromanspacer{}\csromanspacer{}นโลกิยํ ธมฺมํ ปฏิจฺจ โลกุตฺตโร ธมฺโม อุปฺปชฺชติ เหตุปจฺจยา.\csromanspacer{}\csromanspacer{}นโลกิยํ ธมฺมํ ปฏิจฺจ โลกิโย จ โลกุตฺตโร จ ธมฺมา อุปฺปชฺชนฺติ เหตุปจฺจยา. (3)}

\prose{นโลกุตฺตรํ ธมฺมํ ปฏิจฺจ โลกิโย ธมฺโม อุปฺปชฺชติ เหตุปจฺจยา. (1)}

\prose{นโลกิยญฺจ นโลกุตฺตรญฺจ ธมฺมํ ปฏิจฺจ โลกิโย ธมฺโม อุปฺปชฺชติ เหตุปจฺจยา. (1) เหตุยา ปญฺจ.}

\proseitem{13}{นเกนจิ วิญฺเญยฺยํ ธมฺมํ ปฏิจฺจ เกนจิ วิญฺเญยฺโย ธมฺโม อุปฺปชฺชติ เหตุปจฺจยา.\csromanspacer{}\csromanspacer{}เหตุยา นว.}

\csromanheader{2}{14-19. อาสวโคจฺฉก}
\proseitem{14}{น-อาสวํ ธมฺมํ ปฏิจฺจ อาสโว ธมฺโม อุปฺปชฺชติ เหตุปจฺจยา.\csromanspacer{}\csromanspacer{}น-อาสวํ ธมฺมํ ปฏิจฺจ โน-อาสโว ธมฺโม อุปฺปชฺชติ เหตุปจฺจยา.\csromanspacer{}\csromanspacer{}}

\vspace{\tipitakaparskip}
\csromancenter{ธมฺมปจฺจนียานุโลม ทุกปฏฺฐานปาฬิ น-อาสวํ ธมฺมํ ปฏิจฺจ อาสโว จ โน-อาสโว จ ธมฺมา อุปฺปชฺชนฺติ เหตุปจฺจยา.\csromanspacer{} ตีณิ.}
\prose{\csromanfolio{340}นโน-อาสวํ ธมฺมํ ปฏิจฺจ โน-อาสโว ธมฺโม อุปฺปชฺชติ เหตุปจฺจยา.\csromanspacer{}\csromanspacer{}นโน-อาสวํ ธมฺมํ ปฏิจฺจ อาสโว ธมฺโม อุปฺปชฺชติ เหตุปจฺจยา.\csromanspacer{}\csromanspacer{}นโน-อาสวํ ธมฺมํ ปฏิจฺจ อาสโว จ โน-อาสโว จ ธมฺมา อุปฺปชฺชนฺติ เหตุปจฺจยา.\csromanspacer{} ตีณิ.}

\prose{น-อาสวญฺจ นโน-อาสวญฺจ ธมฺมํ ปฏิจฺจ อาสโว ธมฺโม อุปฺปชฺชติ เหตุปจฺจยา.\csromanspacer{} ตีณิ.\csromanspacer{}\csromanspacer{}เหตุยา นว.}

\proseitem{15}{นสาสวํ ธมฺมํ ปฏิจฺจ สาสโว ธมฺโม อุปฺปชฺชติ เหตุปจฺจยา.\csromanspacer{}\csromanspacer{}นสาสวํ ธมฺมํ ปฏิจฺจ อนาสโว ธมฺโม อุปฺปชฺชติ เหตุปจฺจยา.\csromanspacer{}\csromanspacer{}นสาสวํ ธมฺมํ ปฏิจฺจ สาสโว จ อนาสโว จ ธมฺมา อุปฺปชฺชนฺติ เหตุปจฺจยา.\csromanspacer{} ตีณิ.}

\prose{น-อนาสวํ ธมฺมํ ปฏิจฺจ สาสโว ธมฺโม อุปฺปชฺชติ เหตุปจฺจยา.\csromanspacer{} เอกํ.}

\prose{นสาสวญฺจ น-อนาสวญฺจ ธมฺมํ ปฏิจฺจ สาสโว ธมฺโม อุปฺปชฺชติ เหตุปจฺจยา.\csromanspacer{} เอกํ (สํขิตฺตํ.) .\csromanspacer{} เหตุยา ปญฺจ.}

\proseitem{16}{น-อาสวสมฺปยุตฺตํ ธมฺมํ ปฏิจฺจ อาสวสมฺปยุตฺโต ธมฺโม อุปฺปชฺชติ เหตุปจฺจยา.\csromanspacer{}\csromanspacer{}เหตุยา นว.}

\proseitem{17}{น-อาสวญฺเจว อนาสวญฺจ ธมฺมํ ปฏิจฺจ อาสโว เจว สาสโว จ ธมฺโม อุปฺปชฺชติ เหตุปจฺจยา.\csromanspacer{}\csromanspacer{}เหตุยา นว.}

\proseitem{18}{น-อาสวญฺเจว น-อาสววิปฺปยุตฺตญฺจ ธมฺมํ ปฏิจฺจ อาสโว เจว อาสวสมฺปยุตฺโต จ ธมฺโม อุปฺปชฺชติ เหตุปจฺจยา.\csromanspacer{}\csromanspacer{}เหตุยา นว.}

\proseitem{19}{อาสววิปฺปยุตฺตํ นสาสวํ ธมฺมํ ปฏิจฺจ อาสววิปฺปยุตฺโต สาสโว ธมฺโม อุปฺปชฺชติ เหตุปจฺจยา.\csromanspacer{}\csromanspacer{}เหตุยา ปญฺจ.}

\csromanheader{2}{20-54. สญฺโญชนาทิฉโคจฺฉก}
\proseitem{20}{นสญฺโญชนํ ธมฺมํ ปฏิจฺจ สญฺโญชโน ธมฺโม อุปฺปชฺชติ เหตุปจฺจยา ฯเปฯ.}

\prose{นคนฺถํ ธมฺมํ ปฏิจฺจ คนฺโถ ธมฺโม อุปฺปชฺชติ เหตุปจฺจยา ฯเปฯ.}

\csromanheader{2}{55-68. มหนฺตรทุก}
\prose{\csromanfolio{341}น-โอฆํ ธมฺมํ ปฏิจฺจ โอโฆ ธมฺโม อุปฺปชฺชติ เหตุปจฺจยา ฯเปฯ.}

\prose{นโยคํ ธมฺมํ ปฏิจฺจ โยโค ธมฺโม อุปฺปชฺชติ เหตุปจฺจยา ฯเปฯ.}

\prose{นนีวรณํ ธมฺมํ ปฏิจฺจ นีวรโณ ธมฺโม อุปฺปชฺชติ เหตุปจฺจยา.}

\proseitem{21}{นปรามาสํ ธมฺมํ ปฏิจฺจ ปรามาโส ธมฺโม อุปฺปชฺชติ เหตุปจฺจยา.}

\csromanheader{2}{55-68. มหนฺตรทุก}
\proseitem{22}{นสารมฺมณํ ธมฺมํ ปฏิจฺจ สารมฺมโณ ธมฺโม อุปฺปชฺชติ เหตุปจฺจยา.\csromanspacer{}\csromanspacer{}นสารมฺมณํ ธมฺมํ ปฏิจฺจ อนารมฺมโณ ธมฺโม อุปฺปชฺชติ เหตุปจฺจยา.\csromanspacer{}\csromanspacer{}นสารมฺมณํ ธมฺมํ ปฏิจฺจ สารมฺมโณ จ อนารมฺมโณ จ ธมฺมา อุปฺปชฺชนฺติ เหตุปจฺจยา.\csromanspacer{} ตีณิ.}

\prose{น-อนารมฺมณํ ธมฺมํ ปฏิจฺจ อนารมฺมโณ ธมฺโม อุปฺปชฺชติ เหตุปจฺจยา.\csromanspacer{} ตีณิ.}

\prose{นสารมฺมณญฺจ น-อนารมฺมณญฺจ ธมฺมํ ปฏิจฺจ สารมฺมโณ ธมฺโม อุปฺปชฺชติ เหตุปจฺจยา.\csromanspacer{} ตีณิ.\csromanspacer{}\csromanspacer{}เหตุยา นว.}

\proseitem{23}{นจิตฺตํ ธมฺมํ ปฏิจฺจ จิตฺโต ธมฺโม อุปฺปชฺชติ เหตุปจฺจยา.\csromanspacer{}\csromanspacer{}นจิตฺตํ ธมฺมํ ปฏิจฺจ โนจิตฺโต ธมฺโม อุปฺปชฺชติ เหตุปจฺจยา.\csromanspacer{}\csromanspacer{}นจิตฺตํ ธมฺมํ ปฏิจฺจ จิตฺโต จ โนจิตฺโต จ ธมฺมา อุปฺปชฺชนฺติ เหตุปจฺจยา.\csromanspacer{} ตีณิ.}

\prose{นโนจิตฺตํ ธมฺมํ ปฏิจฺจ โนจิตฺโต ธมฺโม อุปฺปชฺชติ เหตุปจฺจยา.\csromanspacer{} เอกํ.}

\prose{นจิตฺตญฺจ นโนจิตฺตญฺจ ธมฺมํ ปฏิจฺจ โนจิตฺโต ธมฺโม อุปฺปชฺชติ เหตุปจฺจยา.\csromanspacer{} เอกํ.\csromanspacer{}\csromanspacer{}เหตุยา ปญฺจ.}

\proseitem{24}{นเจตสิกํ ธมฺมํ ปฏิจฺจ เจตสิโก ธมฺโม อุปฺปชฺชติ เหตุปจฺจยา.\csromanspacer{}\csromanspacer{}นเจตสิกํ ธมฺมํ ปฏิจฺจ อเจตสิโก ธมฺโม อุปฺปชฺชติ เหตุปจฺจยา.\csromanspacer{}\csromanspacer{}นเจตสิกํ ธมฺมํ ปฏิจฺจ เจตสิโก จ อเจตสิโก จ ธมฺมา อุปฺปชฺชนฺติ เหตุปจฺจยา.\csromanspacer{}\csromanspacer{}เหตุยา นว.}

\proseitem{25}{นจิตฺตสมฺปยุตฺตํ ธมฺมํ ปฏิจฺจ จิตฺตสมฺปยุตฺโต ธมฺโม อุปฺปชฺชติ เหตุปจฺจยา.\csromanspacer{}\csromanspacer{}นจิตฺตสมฺปยุตฺตํ ธมฺมํ ปฏิจฺจ จิตฺตวิปฺปยุตฺโต ธมฺโม อุปฺปชฺชติ เหตุปจฺจยา.\csromanspacer{}\csromanspacer{}นจิตฺตสมฺปยุตฺตํ ธมฺมํ ปฏิจฺจ จิตฺตสมฺปยุตฺโต จ จิตฺตวิปฺปยุตฺโต จ ธมฺมา อุปฺปชฺชนฺติ เหตุปจฺจยา. (3)}

\gambhira{ธมฺมปจฺจนียานุโลม ทุกปฏฺฐานปาฬิ}
\prose{\csromanfolio{342}นจิตฺตวิปฺปยุตฺตํ ธมฺมํ ปฏิจฺจ จิตฺตวิปฺปยุตฺโต ธมฺโม อุปฺปชฺชติ เหตุปจฺจยา.\csromanspacer{}\csromanspacer{}นจิตฺตวิปฺปยุตฺตํ ธมฺมํ ปฏิจฺจ จิตฺตสมฺปยุตฺโต ธมฺโม อุปฺปชฺชติ เหตุปจฺจยา.\csromanspacer{}\csromanspacer{}นจิตฺตวิปฺปยุตฺตํ ธมฺมํ ปฏิจฺจ จิตฺตสมฺปยุตฺโต จ จิตฺตวิปฺปยุตฺโต จ ธมฺมา อุปฺปชฺชนฺติ เหตุปจฺจยา. (3)}

\prose{นจิตฺตสมฺปยุตฺตญฺจ นจิตฺตวิปฺปยุตฺตญฺจ ธมฺมํ ปฏิจฺจ จิตฺตสมฺปยุตฺโต ธมฺโม อุปฺปชฺชติ เหตุปจฺจยา.\csromanspacer{} ตีณิ.\csromanspacer{}\csromanspacer{}เหตุยา นว.}

\proseitem{26}{นจิตฺตสํสฏฺฐํ ธมฺมํ ปฏิจฺจ จิตฺตสํสฏฺโฐ ธมฺโม อุปฺปชฺชติ เหตุปจฺจยา.\csromanspacer{}\csromanspacer{}นจิตฺตสํสฏฺฐํ ธมฺมํ ปฏิจฺจ จิตฺตวิสํสฏฺโฐ ธมฺโม อุปฺปชฺชติ เหตุปจฺจยา.\csromanspacer{}\csromanspacer{}นจิตฺตสํสฏฺฐํ ธมฺมํ ปฏิจฺจ จิตฺตสํสฏฺโฐ จ จิตฺตวิสํสฏฺโฐ จ ธมฺมา อุปฺปชฺชนฺติ เหตุปจฺจยา.\csromanspacer{}\csromanspacer{}เหตุยา นว.}

\proseitem{27}{นจิตฺตสมุฏฺฐานํ ธมฺมํ ปฏิจฺจ จิตฺตสมุฏฺฐาโน ธมฺโม อุปฺปชฺชติ เหตุปจฺจยา.\csromanspacer{}\csromanspacer{}นจิตฺตสมุฏฺฐานํ ธมฺมํ ปฏิจฺจ โนจิตฺตสมุฏฺฐาโน ธมฺโม อุปฺปชฺชติ เหตุปจฺจยา.\csromanspacer{}\csromanspacer{}นจิตฺตสมุฏฺฐานํ ธมฺมํ ปฏิจฺจ จิตฺตสมุฏฺฐาโน จ โนจิตฺตสมุฏฺฐาโน จ ธมฺมา อุปฺปชฺชนฺติ เหตุปจฺจยา.\csromanspacer{}\csromanspacer{}เหตุยา นว.}

\proseitem{28}{นจิตฺตสหภุํ ธมฺมํ ปฏิจฺจ จิตฺตสหภู ธมฺโม อุปฺปชฺชติ เหตุปจฺจยา.\csromanspacer{}\csromanspacer{}นจิตฺตสหภุํ ธมฺมํ ปฏิจฺจ โนจิตฺตสหภู ธมฺโม อุปฺปชฺชติ เหตุปจฺจยา.\csromanspacer{}\csromanspacer{}นจิตฺตสหภุํ ธมฺมํ ปฏิจฺจ จิตฺตสหภู จ โนจิตฺตสหภู จ ธมฺมา อุปฺปชฺชนฺติ เหตุปจฺจยา.\csromanspacer{}\csromanspacer{}เหตุยา นว.}

\proseitem{29}{นจิตฺตานุปริวตฺติํ ธมฺมํ ปฏิจฺจ จิตฺตานุปริวตฺตี ธมฺโม อุปฺปชฺชติ เหตุปจฺจยา.\csromanspacer{}\csromanspacer{}นจิตฺตานุปริวตฺติํ ธมฺมํ ปฏิจฺจ โนจิตฺตานุปริวตฺตี ธมฺโม อุปฺปชฺชติ เหตุปจฺจยา.\csromanspacer{}\csromanspacer{}นจิตฺตานุปริวตฺติํ ธมฺมํ ปฏิจฺจ จิตฺตานุปริวตฺตี จ โนจิตฺตานุปริวตฺตี จ ธมฺมา อุปฺปชฺชนฺติ เหตุปจฺจยา.\csromanspacer{}\csromanspacer{}เหตุยา นว.}

\proseitem{30}{นจิตฺตสํสฏฺฐสมุฏฺฐานํ ธมฺมํ ปฏิจฺจ จิตฺตสํสฏฺฐสมุฏฺฐาโน ธมฺโม อุปฺปชฺชติ เหตุปจฺจยา.\csromanspacer{}\csromanspacer{}เหตุยา นว.}

\proseitem{31}{นจิตฺตสํสฏฺฐสมุฏฺฐานสหภุํ ธมฺมํ ปฏิจฺจ จิตฺตสํสฏฺฐสมุฏฺฐานสหภู ธมฺโม อุปฺปชฺชติ เหตุปจฺจยา.\csromanspacer{}\csromanspacer{}เหตุยา นว.}

\proseitem{32}{นจิตฺตสํสฏฺฐสมุฏฺฐานานุปริวตฺติํ ธมฺมํ ปฏิจฺจ จิตฺตสํสฏฺฐสมุฏฺฐานานุปริวตฺตี ธมฺโม อุปฺปชฺชติ เหตุปจฺจยา.\csromanspacer{}\csromanspacer{}เหตุยา นว.}

\csromanheader{2}{83-99. ปิฏฺฐิทุก}
\proseitem{33}{\csromanfolio{343}น-อชฺฌตฺติกํ ธมฺมํ ปฏิจฺจ อชฺฌตฺติโก ธมฺโม อุปฺปชฺชติ เหตุปจฺจยา.\csromanspacer{} ตีณิ.}

\prose{นพาหิรํ ธมฺมํ ปฏิจฺจ พาหิโร ธมฺโม อุปฺปชฺชติ เหตุปจฺจยา.\csromanspacer{} ตีณิ.}

\prose{น-อชฺฌตฺติกญฺจ นพาหิรญฺจ ธมฺมํ ปฏิจฺจ อชฺฌตฺติโก ธมฺโม อุปฺปชฺชติ เหตุปจฺจยา.\csromanspacer{} ตีณิ.\csromanspacer{}\csromanspacer{}เหตุยา นว.}

\proseitem{34}{น-อุปาทา ธมฺมํ ปฏิจฺจ อุปาทา ธมฺโม อุปฺปชฺชติ เหตุปจฺจยา.\csromanspacer{}\csromanspacer{}เหตุยา ปญฺจ.}

\proseitem{35}{น-อุปาทินฺนํ ธมฺมํ ปฏิจฺจ อนุปาทินฺโน ธมฺโม อุปฺปชฺชติ เหตุปจฺจยา.\csromanspacer{}\csromanspacer{}เหตุยา ปญฺจ.}

\csromanheader{2}{69-74. อุปาทานโคจฺฉก}
\proseitem{36}{น-อุปาทานํ ธมฺมํ ปฏิจฺจ อุปาทาโน ธมฺโม อุปฺปชฺชติ เหตุปจฺจยา. (สํขิตฺตํ.)}

\csromanheader{2}{75-82. กิเลสโคจฺฉก}
\vspace{\tipitakaparskip}
\csromancenter{นกิเลสํ ธมฺมํ ปฏิจฺจ กิเลโส ธมฺโม อุปฺปชฺชติ เหตุปจฺจยา.}
\csromanheader{2}{83-99. ปิฏฺฐิทุก}
\proseitem{38}{นทสฺสเนน ปหาตพฺพํ ธมฺมํ ปฏิจฺจ นทสฺสเนน ปหาตพฺโพ ธมฺโม อุปฺปชฺชติ เหตุปจฺจยา.\csromanspacer{} เอกํ.}

\prose{นนทสฺสเนน ปหาตพฺพํ ธมฺมํ ปฏิจฺจ ทสฺสเนน ปหาตพฺโพ ธมฺโม อุปฺปชฺชติ เหตุปจฺจยา.\csromanspacer{}\csromanspacer{}นนทสฺสเนน ปหาตพฺพํ ธมฺมํ ปฏิจฺจ นทสฺสเนน ปหาตพฺโพ ธมฺโม อุปฺปชฺชติ เหตุปจฺจยา.\csromanspacer{}\csromanspacer{}นนทสฺสเนน ปหาตพฺพํ ธมฺมํ ปฏิจฺจ ทสฺสเนน ปหาตพฺโพ จ นทสฺสเนน ปหาตพฺโพ จ ธมฺมา อุปฺปชฺชนฺติ เหตุปจฺจยา.\csromanspacer{} ตีณิ.}

\prose{นทสฺสเนน ปหาตพฺพญฺจ นนทสฺสเนน ปหาตพฺพญฺจ ธมฺมํ ปฏิจฺจ นทสฺสเนน ปหาตพฺโพ ธมฺโม อุปฺปชฺชติ เหตุปจฺจยา.\csromanspacer{} เอกํ.\csromanspacer{}\csromanspacer{}เหตุยา ปญฺจ.}

\proseitem{39}{นภาวนาย ปหาตพฺพํ ธมฺมํ ปฏิจฺจ นภาวนาย ปหาตพฺโพ ธมฺโม อุปฺปชฺชติ เหตุปจฺจยา.\csromanspacer{}\csromanspacer{}เหตุยา ปญฺจ.}

\gambhira{ธมฺมปจฺจนียานุโลม ทุกปฏฺฐานปาฬิ}
\proseitem{40}{\csromanfolio{344}นทสฺสเนน ปหาตพฺพเหตุกํ ธมฺมํ ปฏิจฺจ ทสฺสเนน ปหาตพฺพเหตุโก ธมฺโม อุปฺปชฺชติ เหตุปจฺจยา.\csromanspacer{}\csromanspacer{}เหตุยา นว.}

\proseitem{41}{นภาวนาย ปหาตพฺพเหตุกํ ธมฺมํ ปฏิจฺจ ภาวนาย ปหาตพฺพเหตุโก ธมฺโม อุปฺปชฺชติ เหตุปจฺจยา.\csromanspacer{}\csromanspacer{}เหตุยา นว.}

\proseitem{42}{นสวิตกฺกํ ธมฺมํ ปฏิจฺจ สวิตกฺโก ธมฺโม อุปฺปชฺชติ เหตุปจฺจยา.\csromanspacer{}\csromanspacer{}นสวิตกฺกํ ธมฺมํ ปฏิจฺจ อวิตกฺโก ธมฺโม อุปฺปชฺชติ เหตุปจฺจยา.\csromanspacer{}\csromanspacer{}นสวิตกฺกํ ธมฺมํ ปฏิจฺจ สวิตกฺโก จ อวิตกฺโก จ ธมฺมา อุปฺปชฺชนฺติ เหตุปจฺจยา.\csromanspacer{} ตีณิ.}

\prose{น-อวิตกฺกํ ธมฺมํ ปฏิจฺจ อวิตกฺโก ธมฺโม อุปฺปชฺชติ เหตุปจฺจยา.\csromanspacer{}\csromanspacer{}น-อวิตกฺกํ ธมฺมํ ปฏิจฺจ สวิตกฺโก ธมฺโม อุปฺปชฺชติ เหตุปจฺจยา.\csromanspacer{}\csromanspacer{}น-อวิตกฺกํ ธมฺมํ ปฏิจฺจ สวิตกฺโก จ อวิตกฺโก จ ธมฺมา อุปฺปชฺชนฺติ เหตุปจฺจยา. (3)}

\prose{นสวิตกฺกญฺจ น-อวิตกฺกญฺจ ธมฺมํ ปฏิจฺจ สวิตกฺโก ธมฺโม อุปฺปชฺชติ เหตุปจฺจยา.\csromanspacer{}\csromanspacer{}นสวิตกฺกญฺจ น-อวิตกฺกญฺจ ธมฺมํ ปฏิจฺจ อวิตกฺโก ธมฺโม อุปฺปชฺชติ เหตุปจฺจยา.\csromanspacer{}\csromanspacer{}นสวิตกฺกญฺจ น-อวิตกฺกญฺจ ธมฺมํ ปฏิจฺจ สวิตกฺโก จ อวิตกฺโก จ ธมฺมา อุปฺปชฺชนฺติ เหตุปจฺจยา. (3).\csromanspacer{} เหตุยา นว.}

\proseitem{43}{นสวิจารํ ธมฺมํ ปฏิจฺจ สวิจาโร ธมฺโม อุปฺปชฺชติ เหตุปจฺจยา.\csromanspacer{}\csromanspacer{}เหตุยา นว.}

\proseitem{44}{นสปฺปีติกํ ธมฺมํ ปฏิจฺจ สปฺปีติโก ธมฺโม อุปฺปชฺชติ เหตุปจฺจยา.\csromanspacer{}\csromanspacer{}เหตุยา นว.}

\proseitem{45}{นปีติสหคตํ ธมฺมํ ปฏิจฺจ ปีติสหคโต ธมฺโม อุปฺปชฺชติ เหตุปจฺจยา.\csromanspacer{} เหตุยา นว.}

\proseitem{46}{นสุขสหคตํ ธมฺมํ ปฏิจฺจ สุขสหคโต ธมฺโม อุปฺปชฺชติ เหตุปจฺจยา.\csromanspacer{}\csromanspacer{}เหตุยา นว.}

\proseitem{47}{น-อุเปกฺขาสหคตํ ธมฺมํ ปฏิจฺจ อุเปกฺขาสหคโต ธมฺโม อุปฺปชฺชติ เหตุปจฺจยา.\csromanspacer{}\csromanspacer{}น-อุเปกฺขาสหคตํ ธมฺมํ ปฏิจฺจ น-อุเปกฺขาสหคโต ธมฺโม อุปฺปชฺชติ เหตุปจฺจยา.\csromanspacer{}\csromanspacer{}น-อุเปกฺขาสหคตํ ธมฺมํ ปฏิจฺจ อุเปกฺขาสหคโต จ น-อุเปกฺขาสหคโต จ ธมฺมา อุปฺปชฺชนฺติ เหตุปจฺจยา.\csromanspacer{} ตีณิ.}

\csromanheader{2}{83-99. ปิฏฺฐิทุก}
\prose{\csromanfolio{345}นน-อุเปกฺขาสหคตํ ธมฺมํ ปฏิจฺจ น-อุเปกฺขาสหคโต ธมฺโม อุปฺปชฺชติ เหตุปจฺจยา.\csromanspacer{} ตีณิ.}

\prose{น-อุเปกฺขาสหคตญฺจ นน-อุเปกฺขาสหคตญฺจ ธมฺมํ ปฏิจฺจ อุเปกฺขาสหคโต ธมฺโม อุปฺปชฺชติ เหตุปจฺจยา.\csromanspacer{}\csromanspacer{}น-อุเปกฺขาสหคตญฺจ นน-อุเปกฺขาสหคตญฺจ ธมฺมํ ปฏิจฺจ น-อุเปกฺขาสหคโต ธมฺโม อุปฺปชฺชติ เหตุปจฺจยา.\csromanspacer{}\csromanspacer{}น-อุเปกฺขาสหคตญฺจ นน-อุเปกฺขาสหคตญฺจ ธมฺมํ ปฏิจฺจ อุเปกฺขาสหคโต จ น-อุเปกฺขาสหคโต จ ธมฺมา อุปฺปชฺชนฺติ เหตุปจฺจยา.\csromanspacer{} ตีณิ.\csromanspacer{}\csromanspacer{}เหตุยา นว.}

\proseitem{48}{นกามาวจรํ ธมฺมํ ปฏิจฺจ กามาวจโร ธมฺโม อุปฺปชฺชติ เหตุปจฺจยา.\csromanspacer{}\csromanspacer{}นกามาวจรํ ธมฺมํ ปฏิจฺจ นกามาวจโร ธมฺโม อุปฺปชฺชติ เหตุปจฺจยา.\csromanspacer{}\csromanspacer{}นกามาวจรํ ธมฺมํ ปฏิจฺจ กามาวจโร จ นกามาวจโร จ ธมฺมา อุปฺปชฺชนฺติ เหตุปจฺจยา.\csromanspacer{} ตีณิ.}

\prose{นนกามาวจรํ ธมฺมํ ปฏิจฺจ นกามาวจโร ธมฺโม อุปฺปชฺชติ เหตุปจฺจยา.\csromanspacer{} ตีณิ.}

\prose{นกามาวจรญฺจ นนกามาวจรญฺจ ธมฺมํ ปฏิจฺจ กามาวจโร ธมฺโม อุปฺปชฺชติ เหตุปจฺจยา.\csromanspacer{} ตีณิ.\csromanspacer{}\csromanspacer{}เหตุยา นว.}

\proseitem{49}{นรูปาวจรํ ธมฺมํ ปฏิจฺจ รูปาวจโร ธมฺโม อุปฺปชฺชติ เหตุปจฺจยา.\csromanspacer{}\csromanspacer{}เหตุยา นว.}

\proseitem{50}{น-อรูปาวจรํ ธมฺมํ ปฏิจฺจ น-อรูปาวจโร ธมฺโม อุปฺปชฺชติ เหตุปจฺจยา.\csromanspacer{}\csromanspacer{}เหตุยา ปญฺจ.}

\proseitem{51}{นปริยาปนฺนํ ธมฺมํ ปฏิจฺจ ปริยาปนฺโน ธมฺโม อุปฺปชฺชติ เหตุปจฺจยา.\csromanspacer{} ตีณิ.}

\prose{น-อปริยาปนฺนํ ธมฺมํ ปฏิจฺจ ปริยาปนฺโน ธมฺโม อุปฺปชฺชติ เหตุปจฺจยา.\csromanspacer{} เอกํ.}

\prose{นปริยาปนฺนญฺจ น-อปริยาปนฺนญฺจ ธมฺมํ ปฏิจฺจ ปริยาปนฺโน ธมฺโม อุปฺปชฺชติ เหตุปจฺจยา.\csromanspacer{} เอกํ.\csromanspacer{}\csromanspacer{}เหตุยา ปญฺจ.}

\proseitem{52}{นนิยฺยานิกํ ธมฺมํ ปฏิจฺจ อนิยฺยานิโก ธมฺโม อุปฺปชฺชติ เหตุปจฺจยา.\csromanspacer{} เอกํ.}

\gambhira{ธมฺมปจฺจนียานุโลม ทุกปฏฺฐานปาฬิ}
\prose{\csromanfolio{346}น-อนิยฺยานิกํ ธมฺมํ ปฏิจฺจ อนิยฺยานิโก ธมฺโม อุปฺปชฺชติ เหตุปจฺจยา.\csromanspacer{} ตีณิ.}

\prose{นนิยฺยานิกญฺจ น-อนิยฺยานิกญฺจ ธมฺมํ ปฏิจฺจ อนิยฺยานิโก ธมฺโม อุปฺปชฺชติ เหตุปจฺจยา.\csromanspacer{} เอกํ.\csromanspacer{}\csromanspacer{}เหตุยา ปญฺจ.}

\proseitem{53}{นนิยตํ ธมฺมํ ปฏิจฺจ อนิยโต ธมฺโม อุปฺปชฺชติ เหตุปจฺจยา.\csromanspacer{} เอกํ.}

\prose{น-อนิยตํ ธมฺมํ ปฏิจฺจ อนิยโต ธมฺโม อุปฺปชฺชติ เหตุปจฺจยา.\csromanspacer{} ตีณิ.}

\prose{นนิยตญฺจ น-อนิยตญฺจ ธมฺมํ ปฏิจฺจ อนิยโต ธมฺโม อุปฺปชฺชติ เหตุปจฺจยา.\csromanspacer{} เอกํ.\csromanspacer{}\csromanspacer{}เหตุยา ปญฺจ.}

\proseitem{54}{นส-อุตฺตรํ ธมฺมํ ปฏิจฺจ ส-อุตฺตโร ธมฺโม อุปฺปชฺชติ เหตุปจฺจยา.\csromanspacer{}\csromanspacer{}นส-อุตฺตรํ ธมฺมํ ปฏิจฺจ อนุตฺตโร ธมฺโม อุปฺปชฺชติ เหตุปจฺจยา.\csromanspacer{}\csromanspacer{}นส-อุตฺตรํ ธมฺมํ ปฏิจฺจ ส-อุตฺตโร จ อนุตฺตโร จ ธมฺมา อุปฺปชฺชนฺติ เหตุปจฺจยา.\csromanspacer{} ตีณิ.}

\prose{น-อนุตฺตรํ ธมฺมํ ปฏิจฺจ ส-อุตฺตโร ธมฺโม อุปฺปชฺชติ เหตุปจฺจยา.\csromanspacer{} เอกํ.}

\prose{นส-อุตฺตรญฺจ น-อนุตฺตรญฺจ ธมฺมํ ปฏิจฺจ ส-อุตฺตโร ธมฺโม อุปฺปชฺชติ เหตุปจฺจยา.\csromanspacer{}\csromanspacer{}เหตุยา ปญฺจ.}

\csromanheader{2}{100. สรณทุก 1-7. ปฏิจฺจาทิวาร}
\proseitem{55}{นสรณํ ธมฺมํ ปฏิจฺจ อรโณ ธมฺโม อุปฺปชฺชติ เหตุปจฺจยา.\csromanspacer{} เอกํ.}

\prose{น-อรณํ ธมฺมํ ปฏิจฺจ อรโณ ธมฺโม อุปฺปชฺชติ เหตุปจฺจยา.\csromanspacer{} ตีณิ.}

\prose{นสรณญฺจ น-อรณญฺจ ธมฺมํ ปฏิจฺจ อรโณ ธมฺโม อุปฺปชฺชติ เหตุปจฺจยา.\csromanspacer{} เอกํ.\csromanspacer{}\csromanspacer{}เหตุยา ปญฺจ, อารมฺมเณ เทฺว ฯเปฯ อวิคเต ปญฺจ.}

\csromanheader{2}{\textbf{ปจฺจนีย-นเหตุ}}
\proseitem{56}{นสรณํ ธมฺมํ ปฏิจฺจ อรโณ ธมฺโม อุปฺปชฺชติ นเหตุปจฺจยา.\csromanspacer{}\csromanspacer{}น-อรณํ ธมฺมํ ปฏิจฺจ สรโณ ธมฺโม อุปฺปชฺชติ นเหตุปจฺจยา. (สํขิตฺตํ.)}

\prose{นเหตุยา เทฺว, น-อารมฺมเณ ตีณิ ฯเปฯ โนวิคเต ตีณิ.}

\prose{(สหชาตวารมฺปิ ฯเปฯ สมฺปยุตฺตวารมฺปิ ปฏิจฺจวารสทิสํ.)}

\csromanheader{2}{100. สรณทุก \textbf{เหตุ อารมฺมณ}}
\proseitem{57}{\csromanfolio{347}นสรโณ ธมฺโม อรณสฺส ธมฺมสฺส เหตุปจฺจเยน ปจฺจโย. (1)}

\prose{น-อรโณ ธมฺโม อรณสฺส ธมฺมสฺส เหตุปจฺจเยน ปจฺจโย.\csromanspacer{}\csromanspacer{}น-อรโณ ธมฺโม สรณสฺส ธมฺมสฺส เหตุปจฺจเยน ปจฺจโย.\csromanspacer{}\csromanspacer{}น-อรโณ ธมฺโม สรณสฺส จ อรณสฺส จ ธมฺมสฺส เหตุปจฺจเยน ปจฺจโย. (3)}

\proseitem{58}{นสรโณ ธมฺโม สรณสฺส ธมฺมสฺส อารมฺมณปจฺจเยน ปจฺจโย.\csromanspacer{}\csromanspacer{}นสรโณ ธมฺโม อรณสฺส ธมฺมสฺส อารมฺมณปจฺจเยน ปจฺจโย.\csromanspacer{} เทฺว.}

\prose{น-อรโณ ธมฺโม อรณสฺส ธมฺมสฺส อารมฺมณปจฺจเยน ปจฺจโย.\csromanspacer{}\csromanspacer{}น-อรโณ ธมฺโม สรณสฺส ธมฺมสฺส อารมฺมณปจฺจเยน ปจฺจโย.\csromanspacer{} เทฺว.}

\proseitem{59}{เหตุยา จตฺตาริ, อารมฺมเณ จตฺตาริ ฯเปฯ อวิคเต สตฺต.}

\csromanheader{2}{\textbf{ปจฺจนียุทฺธาร}}
\proseitem{60}{นสรโณ ธมฺโม สรณสฺส ธมฺมสฺส อารมฺมณปจฺจเยน ปจฺจโย.\csromanspacer{} อุปนิสฺสยปจฺจเยน ปจฺจโย.\csromanspacer{} ปุเรชาตปจฺจเยน ปจฺจโย.}

\prose{นเหตุยา สตฺต, น-อารมฺมเณ สตฺต ฯเปฯ โน-อวิคเต จตฺตาริ.}

\prose{(ยถา กุสลตฺติเก ปญฺหาวารสฺส อนุโลมมฺปิ ปจฺจนียมฺปิ อนุโลมปจฺจนียมฺปิ ปจฺจนียานุโลมมฺปิ คณิตํ, เอวํ คเณตพฺพํ.)}

\nitthitam{\textbf{ธมฺมปจฺจนียานุโลเม ทุกปฏฺฐานํ นิฏฺฐิตํ.}}
\pitaka{\textbf{อภิธมฺมปิฏก}}
\csromanheader{2}{\textbf{ธมฺมปจฺจนียานุโลม}}
\gambhira{\textbf{ทุกติกปฏฺฐานปาฬิ}}
\namakkaram{นโม ตสฺส ภควโต อรหโต สมฺมาสมฺพุทฺธสฺส.}
\csromanheader{2}{1. เหตุทุก 1. กุสลตฺติก}
\proseitem{1}{\csromanfolio{349}นเหตุํ นกุสลํ ธมฺมํ ปจฺจยา เหตุ กุสโล ธมฺโม อุปฺปชฺชติ เหตุปจฺจยา.\csromanspacer{}\csromanspacer{}นเหตุํ นกุสลํ ธมฺมํ ปจฺจยา นเหตุ กุสโล ธมฺโม อุปฺปชฺชติ เหตุปจฺจยา.\csromanspacer{}\csromanspacer{}นเหตุํ นกุสลํ ธมฺมํ ปจฺจยา เหตุ กุสโล จ นเหตุ กุสโล จ ธมฺมา อุปฺปชฺชนฺติ เหตุปจฺจยา.\csromanspacer{}\csromanspacer{}เหตุยา ตีณิ.}

\prose{นเหตุํ น-อกุสลํ ธมฺมํ ปจฺจยา เหตุ อกุสโล ธมฺโม อุปฺปชฺชติ เหตุปจฺจยา.\csromanspacer{}\csromanspacer{}เหตุยา ตีณิ.}

\proseitem{2}{นเหตุํ น-อพฺยากตํ ธมฺมํ ปฏิจฺจ นเหตุ อพฺยากโต ธมฺโม อุปฺปชฺชติ เหตุปจฺจยา.\csromanspacer{}\csromanspacer{}นนเหตุํ น-อพฺยากตํ ธมฺมํ ปฏิจฺจ นเหตุ อพฺยากโต ธมฺโม อุปฺปชฺชติ เหตุปจฺจยา.\csromanspacer{}\csromanspacer{}นเหตุํ น-อพฺยากตญฺจ นนเหตุํ น-อพฺยากตญฺจ ธมฺมํ ปฏิจฺจ นเหตุ อพฺยากโต ธมฺโม อุปฺปชฺชติ เหตุปจฺจยา.\csromanspacer{}\csromanspacer{}เหตุยา ตีณิ.}

\csromanheader{2}{2-4. สเหตุกทุกาทิ 1. กุสลตฺติก}
\proseitem{3}{นสเหตุกํ นกุสลํ ธมฺมํ ปจฺจยา สเหตุโก กุสโล ธมฺโม อุปฺปชฺชติ เหตุปจฺจยา.\csromanspacer{}\csromanspacer{}เหตุยา เอกํ.}

\prose{นสเหตุกํ น-อกุสลํ ธมฺมํ ปจฺจยา สเหตุโก อกุสโล ธมฺโม อุปฺปชฺชติ เหตุปจฺจยา.\csromanspacer{}\csromanspacer{}เหตุยา เอกํ.}

\csromanheader{2}{ธมฺมปจฺจนียานุโลม ทุกติกปฏฺฐานปาฬิ}
\proseitem{4}{\csromanfolio{350}นสเหตุกํ น-อพฺยากตํ ธมฺมํ ปฏิจฺจ อเหตุโก อพฺยากโต ธมฺโม อุปฺปชฺชติ เหตุปจฺจยา.\csromanspacer{}\csromanspacer{}น-อเหตุกํ น-อพฺยากตํ ธมฺมํ ปฏิจฺจ อเหตุโก อพฺยากโต ธมฺโม อุปฺปชฺชติ เหตุปจฺจยา.\csromanspacer{}\csromanspacer{}นสเหตุกํ น-อพฺยากตญฺจ น-อเหตุกํ น-อพฺยากตญฺจ ธมฺมํ ปฏิจฺจ อเหตุโก อพฺยากโต ธมฺโม อุปฺปชฺชติ เหตุปจฺจยา.\csromanspacer{}\csromanspacer{}เหตุยา ตีณิ.}

\proseitem{5}{นเหตุสมฺปยุตฺตํ นกุสลํ ธมฺมํ ปจฺจยา.}

\proseitem{6}{นเหตุ เจว น-อเหตุโก จ นกุสโล ธมฺโม เหตุสฺส เจว สเหตุกสฺส จ กุสลสฺส ธมฺมสฺส อารมฺมณปจฺจเยน ปจฺจโย.\csromanspacer{}\csromanspacer{}นเหตุ เจว น-อเหตุโก จ นกุสโล ธมฺโม สเหตุกสฺส เจว น จ เหตุสฺส กุสลสฺส ธมฺมสฺส อารมฺมณปจฺจเยน ปจฺจโย.\csromanspacer{}\csromanspacer{}นเหตุ เจว น-อเหตุโก จ นกุสโล ธมฺโม เหตุสฺส เจว สเหตุกสฺส กุสลสฺส จ สเหตุกสฺส เจว น จ เหตุสฺส กุสลสฺส จ ธมฺมสฺส อารมฺมณปจฺจเยน ปจฺจโย.\csromanspacer{} ตีณิ.}

\prose{น-อเหตุโก เจว นน จ เหตุ นกุสโล ธมฺโม เหตุสฺส เจว สเหตุกสฺส จ กุสลสฺส ธมฺมสฺส อารมฺมณปจฺจเยน ปจฺจโย.\csromanspacer{} ตีณิ.}

\prose{นเหตุ เจว น-อเหตุโก นกุสโล จ น-อเหตุโก เจว นนเหตุ นกุสโล จ ธมฺมา เหตุสฺส เจว สเหตุกสฺส จ กุสลสฺส ธมฺมสฺส อารมฺมณปจฺจเยน ปจฺจโย.\csromanspacer{} ตีณิ. (สพฺพตฺถ นว ปญฺหา.)}

\proseitem{7}{นเหตุ เจว น-อเหตุโก จ น-อกุสโล ธมฺโม เหตุสฺส เจว สเหตุกสฺส จ อกุสลสฺส ธมฺมสฺส อารมฺมณปจฺจเยน ปจฺจโย. (สํขิตฺตํ.\csromanspacer{} นว ปญฺหา.)}

\prose{นเหตุ เจว น-อเหตุโก จ น-อพฺยากโต ธมฺโม เหตุสฺส เจว สเหตุกสฺส จ อพฺยากตสฺส ธมฺมสฺส อารมฺมณปจฺจเยน ปจฺจโย.\csromanspacer{} ตีณิ. (นว ปญฺหา.)}

\csromanheader{2}{5. เหตุเหตุสมฺปยุตฺตทุก 1. กุสลตฺติก}
\proseitem{8}{นเหตุ เจว นเหตุวิปฺปยุตฺโต จ นกุสโล ธมฺโม เหตุสฺส เจว เหตุสมฺปยุตฺตสฺส จ กุสลสฺส ธมฺมสฺส อารมฺมณปจฺจเยน ปจฺจโย. (สํขิตฺตํ.\csromanspacer{} นว ปญฺหา กาตพฺพา.)}

\csromanheader{2}{7-13. จูฬนฺตรทุก 1. กุสลตฺติก}
\prose{\csromanfolio{351}นเหตุ เจว นเหตุวิปฺปยุตฺโต จ น-อกุสโล ธมฺโม เหตุสฺส เจว เหตุสมฺปยุตฺตสฺส จ อกุสลสฺส ธมฺมสฺส อารมฺมณปจฺจเยน ปจฺจโย. (สํขิตฺตํ.\csromanspacer{} นว ปญฺหา กาตพฺพา.)}

\prose{นเหตุ เจว นเหตุวิปฺปยุตฺโต จ น-อพฺยากโต ธมฺโม เหตุสฺส เจว เหตุสมฺปยุตฺตสฺส จ อพฺยากตสฺส ธมฺมสฺส อารมฺมณปจฺจเยน ปจฺจโย.\csromanspacer{} ตีณิ. (นว ปญฺหา กาตพฺพา.)}

\csromanheader{2}{6. นเหตุสเหตุกทุก 1. กุสลตฺติก}
\proseitem{9}{นเหตุํ นสเหตุกํ นกุสลํ ธมฺมํ ปจฺจยา นเหตุ สเหตุโก กุสโล ธมฺโม อุปฺปชฺชติ เหตุปจฺจยา.\csromanspacer{}\csromanspacer{}เหตุยา เอกํ.}

\prose{นเหตุํ นสเหตุกํ น-อกุสลํ ธมฺมํ ปจฺจยา นเหตุ สเหตุโก อกุสโล ธมฺโม อุปฺปชฺชติ เหตุปจฺจยา.\csromanspacer{}\csromanspacer{}เหตุยา เอกํ.}

\prose{นเหตุํ น-อเหตุกํ น-อพฺยากตํ ธมฺมํ ปฏิจฺจ นเหตุ อเหตุโก อพฺยากโต ธมฺโม อุปฺปชฺชติ เหตุปจฺจยา.\csromanspacer{}\csromanspacer{}เหตุยา เอกํ.}

\csromanheader{2}{7-13. จูฬนฺตรทุก 1. กุสลตฺติก}
\proseitem{10}{น-อปฺปจฺจยํ นกุสลํ ธมฺมํ ปจฺจยา สปฺปจฺจโย กุสโล ธมฺโม อุปฺปชฺชติ เหตุปจฺจยา.\csromanspacer{}\csromanspacer{}เหตุยา เอกํ.}

\prose{น-อปฺปจฺจยํ น-อกุสลํ ธมฺมํ ปจฺจยา สปฺปจฺจโย อกุสโล ธมฺโม อุปฺปชฺชติ เหตุปจฺจยา.\csromanspacer{}\csromanspacer{}เหตุยา เอกํ.}

\prose{น-อปฺปจฺจยํ น-อพฺยากตํ ธมฺมํ ปฏิจฺจ สปฺปจฺจโย อพฺยากโต ธมฺโม อุปฺปชฺชติ เหตุปจฺจยา.\csromanspacer{}\csromanspacer{}เหตุยา เอกํ.}

\csromanheader{2}{11. น-อสงฺขตํ ฯเปฯ. (สปฺปจฺจยทุกสทิสํ.)}
\proseitem{12}{นสนิทสฺสนํ นกุสลํ ธมฺมํ ปจฺจยา อนิทสฺสโน กุสโล ธมฺโม อุปฺปชฺชติ เหตุปจฺจยา.\csromanspacer{}\csromanspacer{}เหตุยา เอกํ.}

\prose{นสนิทสฺสนํ น-อกุสลํ ธมฺมํ ปจฺจยา อนิทสฺสโน อกุสโล ธมฺโม อุปฺปชฺชติ เหตุปจฺจยา.\csromanspacer{}\csromanspacer{}เหตุยา เอกํ.}

\csromanheader{2}{ธมฺมปจฺจนียานุโลม ทุกติกปฏฺฐานปาฬิ}
\prose{\csromanfolio{352}นสนิทสฺสนํ น-อพฺยากตํ ธมฺมํ ปฏิจฺจ สนิทสฺสโน อพฺยากโต ธมฺโม อุปฺปชฺชติ เหตุปจฺจยา.\csromanspacer{}\csromanspacer{}นสนิทสฺสนํ น-อพฺยากตํ ธมฺมํ ปฏิจฺจ อนิทสฺสโน อพฺยากโต ธมฺโม อุปฺปชฺชติ เหตุปจฺจยา.\csromanspacer{}\csromanspacer{}นสนิทสฺสนํ น-อพฺยากตํ ธมฺมํ ปฏิจฺจ สนิทสฺสโน อพฺยากโต จ อนิทสฺสโน อพฺยากโต จ ธมฺมา อุปฺปชฺชนฺติ เหตุปจฺจยา. (3).\csromanspacer{} เหตุยา ตีณิ.}

\proseitem{13}{นสปฺปฏิฆํ นกุสลํ ธมฺมํ ปจฺจยา อปฺปฏิโฆ กุสโล ธมฺโม อุปฺปชฺชติ เหตุปจฺจยา.}

\proseitem{14}{น-อรูปิํ นกุสลํ ธมฺมํ ปจฺจยา อรูปี กุสโล ธมฺโม อุปฺปชฺชติ เหตุปจฺจยา.\csromanspacer{}\csromanspacer{}เหตุยา เอกํ.}

\prose{น-อรูปิํ น-อกุสลํ ธมฺมํ ปจฺจยา อรูปี อกุสโล ธมฺโม อุปฺปชฺชติ เหตุปจฺจยา.\csromanspacer{}\csromanspacer{}เหตุยา เอกํ.}

\prose{นรูปิํ น-อพฺยากตํ ธมฺมํ ปฏิจฺจ รูปี อพฺยากโต ธมฺโม อุปฺปชฺชติ เหตุปจฺจยา.\csromanspacer{}\csromanspacer{}เหตุยา เอกํ.}

\proseitem{15}{นโลกุตฺตรํ นกุสลํ ธมฺมํ ปจฺจยา โลกิโย กุสโล ธมฺโม อุปฺปชฺชติ เหตุปจฺจยา.\csromanspacer{}\csromanspacer{}นโลกุตฺตรํ นกุสลํ ธมฺมํ ปจฺจยา โลกุตฺตโร กุสโล ธมฺโม อุปฺปชฺชติ เหตุปจฺจยา.\csromanspacer{}\csromanspacer{}เหตุยา เทฺว.}

\prose{นโลกุตฺตรํ น-อกุสลํ ธมฺมํ ปจฺจยา โลกิโย อกุสโล ธมฺโม อุปฺปชฺชติ เหตุปจฺจยา.\csromanspacer{}\csromanspacer{}เหตุยา เอกํ.}

\prose{นโลกิยํ น-อพฺยากตํ ธมฺมํ ปฏิจฺจ โลกิโย อพฺยากโต ธมฺโม อุปฺปชฺชติ เหตุปจฺจยา.\csromanspacer{}\csromanspacer{}นโลกุตฺตรํ น-อพฺยากตํ ธมฺมํ ปฏิจฺจ โลกิโย อพฺยากโต ธมฺโม อุปฺปชฺชติ เหตุปจฺจยา.\csromanspacer{}\csromanspacer{}เหตุยา เทฺว.}

\proseitem{16}{นเกนจิ วิญฺเญยฺยํ นกุสลํ ธมฺมํ ปจฺจยา เกนจิ วิญฺเญยฺโย กุสโล ธมฺโม อุปฺปชฺชติ เหตุปจฺจยา.\csromanspacer{} ตีณิ.}

\prose{นเกนจิ นวิญฺเญยฺยํ นกุสลํ ธมฺมํ ปจฺจยา เกนจิ นวิญฺเญยฺโย กุสโล ธมฺโม อุปฺปชฺชติ เหตุปจฺจยา.\csromanspacer{} ตีณิ.}

\csromanheader{2}{16. อาสวสมฺปยุตฺตทุก 1. กุสลตฺติก}
\prose{\csromanfolio{353}นเกนจิ วิญฺเญยฺยํ นกุสลญฺจ นเกนจิ นวิญฺเญยฺยํ นกุสลญฺจ ธมฺมํ ปจฺจยา เกนจิ วิญฺเญยฺโย กุสโล ธมฺโม อุปฺปชฺชติ เหตุปจฺจยา.\csromanspacer{} ตีณิ. (สพฺพตฺถ นว.)}

\prose{นเกนจิ วิญฺเญยฺยํ น-อกุสลํ ธมฺมํ ปจฺจยา เกนจิ วิญฺเญยฺโย อกุสโล ธมฺโม อุปฺปชฺชติ เหตุปจฺจยา.\csromanspacer{}\csromanspacer{}นเกนจิ นวิญฺเญยฺยํ น-อกุสลํ ธมฺมํ ปจฺจยา เกนจิ นวิญฺเญยฺโย อกุสโล ธมฺโม อุปฺปชฺชติ เหตุปจฺจยา.\csromanspacer{}\csromanspacer{}เหตุยา นว.}

\prose{นเกนจิ วิญฺเญยฺยํ น-อพฺยากตํ ธมฺมํ ปฏิจฺจ เกนจิ วิญฺเญยฺโย อพฺยากโต ธมฺโม อุปฺปชฺชติ เหตุปจฺจยา.\csromanspacer{}\csromanspacer{}เหตุยา นว.}

\csromanheader{2}{14. อาสวทุก 1. กุสลตฺติก}
\proseitem{17}{น-อาสวํ นกุสลํ ธมฺมํ ปจฺจยา โน-อาสโว กุสโล ธมฺโม อุปฺปชฺชติ เหตุปจฺจยา.\csromanspacer{}\csromanspacer{}เหตุยา เอกํ.}

\prose{น-อาสวํ น-อกุสลํ ธมฺมํ ปจฺจยา อาสโว อกุสโล ธมฺโม อุปฺปชฺชติ เหตุปจฺจยา.\csromanspacer{}\csromanspacer{}เหตุยา ตีณิ.}

\prose{น-อาสวํ น-อพฺยากตํ ธมฺมํ ปฏิจฺจ โน-อาสโว อพฺยากโต ธมฺโม อุปฺปชฺชติ เหตุปจฺจยา.\csromanspacer{}\csromanspacer{}นโน-อาสวํ น-อพฺยากตํ ธมฺมํ ปฏิจฺจ โน-อาสโว อพฺยากโต ธมฺโม อุปฺปชฺชติ เหตุปจฺจยา. (ทุกมูลกํ เอกํ.) .\csromanspacer{} เหตุยา ตีณิ.}

\csromanheader{2}{15. สาสวทุก 1. กุสลตฺติก}
\proseitem{18}{น-อนาสวํ นกุสลํ ธมฺมํ ปจฺจยา อนาสโว กุสโล ธมฺโม อุปฺปชฺชติ เหตุปจฺจยา.\csromanspacer{}\csromanspacer{}น-อนาสวํ นกุสลํ ธมฺมํ ปจฺจยา สาสโว กุสโล ธมฺโม อุปฺปชฺชติ เหตุปจฺจยา.\csromanspacer{}\csromanspacer{}เหตุยา เทฺว.}

\prose{น-อนาสวํ น-อกุสลํ ธมฺมํ ปจฺจยา สาสโว อกุสโล ธมฺโม อุปฺปชฺชติ เหตุปจฺจยา.\csromanspacer{}\csromanspacer{}เหตุยา เอกํ.}

\csromanheader{2}{16. อาสวสมฺปยุตฺตทุก 1. กุสลตฺติก}
\proseitem{19}{น-อาสวสมฺปยุตฺตํ นกุสลํ ธมฺมํ ปจฺจยา อาสววิปฺปยุตฺโต กุสโล ธมฺโม อุปฺปชฺชติ เหตุปจฺจยา.\csromanspacer{}\csromanspacer{}เหตุยา เอกํ.}

\csromanheader{2}{ธมฺมปจฺจนียานุโลม ทุกติกปฏฺฐานปาฬิ}
\prose{\csromanfolio{354}น-อาสวสมฺปยุตฺตํ น-อกุสลํ ธมฺมํ ปจฺจยา อาสวสมฺปยุตฺโต อกุสโล ธมฺโม อุปฺปชฺชติ เหตุปจฺจยา.\csromanspacer{}\csromanspacer{}เหตุยา ตีณิ.}

\prose{น-อาสวสมฺปยุตฺตํ น-อพฺยากตํ ธมฺมํ ปฏิจฺจ อาสววิปฺปยุตฺโต อพฺยากโต ธมฺโม อุปฺปชฺชติ เหตุปจฺจยา.\csromanspacer{}\csromanspacer{}เหตุยา ตีณิ.}

\csromanheader{2}{17. อาสวสาสวทุก 1. กุสลตฺติก}
\proseitem{20}{น-อาสวญฺเจว น-อนาสวญฺจ นกุสลํ ธมฺมํ ปจฺจยา สาสโว เจว โน-อาสโว จ กุสโล ธมฺโม อุปฺปชฺชติ เหตุปจฺจยา.\csromanspacer{}\csromanspacer{}เหตุยา เอกํ.}

\prose{น-อาสวญฺเจว น-อนาสวญฺจ น-อกุสลํ ธมฺมํ ปจฺจยา อาสโว เจว สาสโว จ อกุสโล ธมฺโม อุปฺปชฺชติ เหตุปจฺจยา.\csromanspacer{}\csromanspacer{}เหตุยา ตีณิ.}

\prose{น-อาสวญฺเจว น-อนาสวญฺจ น-อพฺยากตํ ธมฺมํ ปฏิจฺจ สาสโว เจว โน จ อาสโว อพฺยากโต ธมฺโม อุปฺปชฺชติ เหตุปจฺจยา.\csromanspacer{}\csromanspacer{}เหตุยา ตีณิ. (อาสว-อาสวสมฺปยุตฺตทุกํ นตฺถิ.)}

\csromanheader{2}{19. อาสววิปฺปยุตฺตสาสวทุก 1. กุสลตฺติก}
\proseitem{21}{อาสววิปฺปยุตฺตํ น-อนาสวํ นกุสลํ ธมฺมํ ปจฺจยา อาสววิปฺปยุตฺโต อนาสโว กุสโล ธมฺโม อุปฺปชฺชติ เหตุปจฺจยา.\csromanspacer{}\csromanspacer{}อาสววิปฺปยุตฺตํ น-อนาสวํ นกุสลํ ธมฺมํ ปจฺจยา อาสววิปฺปยุตฺโต สาสโว กุสโล ธมฺโม อุปฺปชฺชติ เหตุปจฺจยา.\csromanspacer{}\csromanspacer{}เหตุยา เทฺว.}

\prose{อาสววิปฺปยุตฺตํ น-อนาสวํ น-อกุสลํ ธมฺมํ ปจฺจยา อาสววิปฺปยุตฺโต สาสโว อกุสโล ธมฺโม อุปฺปชฺชติ เหตุปจฺจยา.\csromanspacer{}\csromanspacer{}เหตุยา เอกํ.}

\prose{อาสววิปฺปยุตฺตํ นสาสวํ น-อพฺยากตํ ธมฺมํ ปฏิจฺจ อาสววิปฺปยุตฺโต สาสโว อพฺยากโต ธมฺโม อุปฺปชฺชติ เหตุปจฺจยา.\csromanspacer{}\csromanspacer{}อาสววิปฺปยุตฺตํ น-อนาสวํ น-อพฺยากตํ ธมฺมํ ปฏิจฺจ อาสววิปฺปยุตฺโต สาสโว อพฺยากโต ธมฺโม อุปฺปชฺชติ เหตุปจฺจยา.\csromanspacer{}\csromanspacer{}เหตุยา เทฺว.}

\vspace{\tipitakaparskip}
\csromancenter{มหนฺตรทุก 1.\csromanspacer{} กุสลตฺติก \textbf{20-54.}\csromanspacer{}\textbf{ สญฺโญชนาทิฉโคจฺฉก 1.}\csromanspacer{}\textbf{ กุสลตฺติก}}
\proseitem{22}{\csromanfolio{355}นสญฺโญชนํ น-อกุสลํ ธมฺมํ ปจฺจยา สญฺโญชโน อกุสโล ธมฺโม อุปฺปชฺชติ เหตุปจฺจยา.\csromanspacer{}\csromanspacer{}เหตุยา ตีณิ.}

\proseitem{23}{นคนฺถํ น-อกุสลํ ธมฺมํ ปจฺจยา คนฺโถ อกุสโล ธมฺโม อุปฺปชฺชติ เหตุปจฺจยา.\csromanspacer{}\csromanspacer{}เหตุยา ตีณิ.}

\proseitem{24}{น-โอฆํ น-อกุสลํ ธมฺมํ ปจฺจยา โอโฆ อกุสโล ธมฺโม อุปฺปชฺชติ เหตุปจฺจยา.\csromanspacer{}\csromanspacer{}เหตุยา ตีณิ.}

\proseitem{25}{นโยคํ น-อกุสลํ ธมฺมํ ปจฺจยา โยโค อกุสโล ธมฺโม อุปฺปชฺชติ เหตุปจฺจยา.\csromanspacer{}\csromanspacer{}เหตุยา ตีณิ.}

\proseitem{26}{นนีวรณํ น-อกุสลํ ธมฺมํ ปจฺจยา นีวรโณ อกุสโล ธมฺโม อุปฺปชฺชติ เหตุปจฺจยา.\csromanspacer{}\csromanspacer{}เหตุยา ตีณิ.}

\proseitem{27}{นปรามาสํ น-อกุสลํ ธมฺมํ ปจฺจยา ปรามาโส อกุสโล ธมฺโม อุปฺปชฺชติ เหตุปจฺจยา.\csromanspacer{}\csromanspacer{}เหตุยา ตีณิ.}

\csromanheader{2}{55-68. มหนฺตรทุก 1. กุสลตฺติก}
\proseitem{28}{นสารมฺมณํ นกุสลํ ธมฺมํ ปจฺจยา สารมฺมโณ กุสโล ธมฺโม อุปฺปชฺชติ เหตุปจฺจยา.\csromanspacer{}\csromanspacer{}เหตุยา เอกํ.}

\prose{นสารมฺมณํ น-อกุสลํ ธมฺมํ ปจฺจยา สารมฺมโณ อกุสโล ธมฺโม อุปฺปชฺชติ เหตุปจฺจยา.\csromanspacer{}\csromanspacer{}เหตุยา เอกํ.}

\prose{น-อนารมฺมณํ น-อพฺยากตํ ธมฺมํ ปฏิจฺจ อนารมฺมโณ อพฺยากโต ธมฺโม อุปฺปชฺชติ เหตุปจฺจยา.\csromanspacer{}\csromanspacer{}เหตุยา เอกํ.}

\proseitem{29}{โนจิตฺตํ นกุสลํ ธมฺมํ ปจฺจยา จิตฺโต กุสโล ธมฺโม อุปฺปชฺชติ เหตุปจฺจยา.\csromanspacer{}\csromanspacer{}เหตุยา ตีณิ. (น-อกุสเล ตีณิ.\csromanspacer{} น-อพฺยากเต ตีณิ.\csromanspacer{} สํขิตฺตํ.)}

\proseitem{30}{นเจตสิกํ นกุสลํ ธมฺมํ ปจฺจยา เจตสิโก กุสโล ธมฺโม อุปฺปชฺชติ เหตุปจฺจยา.\csromanspacer{}\csromanspacer{}เหตุยา ตีณิ.}

\csromanheader{2}{ธมฺมปจฺจนียานุโลม ทุกติกปฏฺฐานปาฬิ}
\proseitem{31}{\csromanfolio{356}นจิตฺตสมฺปยุตฺตํ นกุสลํ ธมฺมํ ปจฺจยา จิตฺตสมฺปยุตฺโต กุสโล ธมฺโม อุปฺปชฺชติ เหตุปจฺจยา.\csromanspacer{}\csromanspacer{}เหตุยา เอกํ.}

\proseitem{32}{นจิตฺตสํสฏฺฐํ นกุสลํ ธมฺมํ ปจฺจยา จิตฺตสํสฏฺโฐ กุสโล ธมฺโม อุปฺปชฺชติ เหตุปจฺจยา.\csromanspacer{}\csromanspacer{}เหตุยา เอกํ.}

\proseitem{33}{นจิตฺตสมุฏฺฐานํ นกุสลํ ธมฺมํ ปจฺจยา จิตฺตสมุฏฺฐาโน กุสโล ธมฺโม อุปฺปชฺชติ เหตุปจฺจยา.\csromanspacer{}\csromanspacer{}เหตุยา ตีณิ.}

\proseitem{34}{นจิตฺตสหภุํ นกุสลํ ธมฺมํ ปจฺจยา จิตฺตสหภู กุสโล ธมฺโม อุปฺปชฺชติ เหตุปจฺจยา.\csromanspacer{}\csromanspacer{}เหตุยา ตีณิ.}

\proseitem{35}{นจิตฺตานุปริวตฺติํ นกุสลํ ธมฺมํ ปจฺจยา จิตฺตานุปริวตฺตี กุสโล ธมฺโม อุปฺปชฺชติ เหตุปจฺจยา.\csromanspacer{}\csromanspacer{}เหตุยา ตีณิ.}

\proseitem{36}{นจิตฺตสํสฏฺฐสมุฏฺฐานํ นกุสลํ ธมฺมํ ปจฺจยา จิตฺตสํสฏฺฐสมุฏฺฐาโน กุสโล ธมฺโม อุปฺปชฺชติ เหตุปจฺจยา.\csromanspacer{}\csromanspacer{}เหตุยา ตีณิ.}

\proseitem{37}{นจิตฺตสํสฏฺฐสมุฏฺฐานสหภุํ นกุสลํ ธมฺมํ ปจฺจยา จิตฺตสํสฏฺฐสมุฏฺฐานสหภู กุสโล ธมฺโม อุปฺปชฺชติ เหตุปจฺจยา.\csromanspacer{}\csromanspacer{}เหตุยา ตีณิ.}

\proseitem{38}{นจิตฺตสํสฏฺฐสมุฏฺฐานานุปริวตฺติํ นกุสลํ ธมฺมํ ปจฺจยา จิตฺตสํสฏฺฐสมุฏฺฐานานุปริวตฺตี กุสโล ธมฺโม อุปฺปชฺชติ เหตุปจฺจยา.\csromanspacer{}\csromanspacer{}เหตุยา ตีณิ.}

\proseitem{39}{น-อชฺฌตฺติกํ นกุสลํ ธมฺมํ ปจฺจยา อชฺฌตฺติโก กุสโล ธมฺโม อุปฺปชฺชติ เหตุปจฺจยา.\csromanspacer{}\csromanspacer{}เหตุยา ตีณิ.}

\proseitem{40}{นโน-อุปาทา นกุสลํ ธมฺมํ ปจฺจยา โน-อุปาทา กุสโล ธมฺโม อุปฺปชฺชติ เหตุปจฺจยา.\csromanspacer{}\csromanspacer{}เหตุยา เอกํ.}

\proseitem{41}{น-อนุปาทินฺนํ นกุสลํ ธมฺมํ ปจฺจยา อนุปาทินฺโน กุสโล ธมฺโม อุปฺปชฺชติ เหตุปจฺจยา.\csromanspacer{}\csromanspacer{}เหตุยา เอกํ.}

\csromanheader{2}{83-100. ปิฏฺฐิทุก 1. กุสลตฺติก \textbf{69. อุปาทานทุก 1. กุสลตฺติก}}
\proseitem{42}{\csromanfolio{357}น-อุปาทานํ น-อกุสลํ ธมฺมํ ปจฺจยา อุปาทาโน อกุสโล ธมฺโม อุปฺปชฺชติ เหตุปจฺจยา.\csromanspacer{}\csromanspacer{}เหตุยา ตีณิ.}

\csromanheader{2}{75. กิเลสทุก 1. กุสลตฺติก}
\proseitem{43}{นกิเลสํ น-อกุสลํ ธมฺมํ ปจฺจยา กิเลโส อกุสโล ธมฺโม อุปฺปชฺชติ เหตุปจฺจยา.\csromanspacer{}\csromanspacer{}เหตุยา ตีณิ.}

\csromanheader{2}{83-100. ปิฏฺฐิทุก 1. กุสลตฺติก}
\proseitem{44}{นทสฺสเนน ปหาตพฺพํ นกุสลํ ธมฺมํ ปจฺจยา นทสฺสเนน ปหาตพฺโพ กุสโล ธมฺโม อุปฺปชฺชติ เหตุปจฺจยา.\csromanspacer{}\csromanspacer{}เหตุยา เอกํ.}

\prose{นทสฺสเนน ปหาตพฺพํ น-อกุสลํ ธมฺมํ ปจฺจยา ทสฺสเนน ปหาตพฺโพ อกุสโล ธมฺโม อุปฺปชฺชติ เหตุปจฺจยา.\csromanspacer{}\csromanspacer{}นทสฺสเนน ปหาตพฺพํ น-อกุสลํ ธมฺมํ ปจฺจยา นทสฺสเนน ปหาตพฺโพ อกุสโล ธมฺโม อุปฺปชฺชติ เหตุปจฺจยา.\csromanspacer{}\csromanspacer{}เหตุยา เทฺว.}

\prose{นทสฺสเนน ปหาตพฺพํ น-อพฺยากตํ ธมฺมํ ปฏิจฺจ นทสฺสเนน ปหาตพฺโพ อพฺยากโต ธมฺโม อุปฺปชฺชติ เหตุปจฺจยา.\csromanspacer{}\csromanspacer{}นนทสฺสเนน ปหาตพฺพํ น-อพฺยากตํ ธมฺมํ ปฏิจฺจ นทสฺสเนน ปหาตพฺโพ อพฺยากโต ธมฺโม อุปฺปชฺชติ เหตุปจฺจยา.\csromanspacer{}\csromanspacer{}เหตุยา เทฺว.}

\proseitem{45}{นภาวนาย ปหาตพฺพํ นกุสลํ ธมฺมํ ปจฺจยา นภาวนาย ปหาตพฺโพ กุสโล ธมฺโม อุปฺปชฺชติ เหตุปจฺจยา.\csromanspacer{}\csromanspacer{}เหตุยา เอกํ.}

\proseitem{46}{นทสฺสเนน ปหาตพฺพเหตุกํ นกุสลํ ธมฺมํ ปจฺจยา นทสฺสเนน ปหาตพฺพเหตุโก กุสโล ธมฺโม อุปฺปชฺชติ เหตุปจฺจยา.\csromanspacer{}\csromanspacer{}เหตุยา เอกํ.}

\proseitem{47}{นภาวนาย ปหาตพฺพเหตุกํ นกุสลํ ธมฺมํ ปจฺจยา นภาวนาย ปหาตพฺพเหตุโก กุสโล ธมฺโม อุปฺปชฺชติ เหตุปจฺจยา.\csromanspacer{}\csromanspacer{}เหตุยา เอกํ.}

\proseitem{48}{นสวิตกฺกํ นกุสลํ ธมฺมํ ปจฺจยา สวิตกฺโก กุสโล ธมฺโม อุปฺปชฺชติ เหตุปจฺจยา.\csromanspacer{}\csromanspacer{}เหตุยา ตีณิ.}

\csromanheader{2}{ธมฺมปจฺจนียานุโลม ทุกติกปฏฺฐานปาฬิ}
\proseitem{49}{\csromanfolio{358}นสวิจารํ นกุสลํ ธมฺมํ ปจฺจยา สวิจาโร กุสโล ธมฺโม อุปฺปชฺชติ เหตุปจฺจยา.\csromanspacer{}\csromanspacer{}เหตุยา ตีณิ.}

\proseitem{50}{นสปฺปีติกํ นกุสลํ ธมฺมํ ปจฺจยา สปฺปีติโก กุสโล ธมฺโม อุปฺปชฺชติ เหตุปจฺจยา.\csromanspacer{}\csromanspacer{}เหตุยา ตีณิ.}

\proseitem{51}{นปีติสหคตํ นกุสลํ ธมฺมํ ปจฺจยา ปีติสหคโต กุสโล ธมฺโม อุปฺปชฺชติ เหตุปจฺจยา.\csromanspacer{}\csromanspacer{}เหตุยา ตีณิ.}

\proseitem{52}{นสุขสหคตํ นกุสลํ ธมฺมํ ปจฺจยา สุขสหคโต กุสโล ธมฺโม อุปฺปชฺชติ เหตุปจฺจยา.\csromanspacer{}\csromanspacer{}เหตุยา ตีณิ.}

\proseitem{53}{น-อุเปกฺขาสหคตํ นกุสลํ ธมฺมํ ปจฺจยา อุเปกฺขาสหคโต กุสโล ธมฺโม อุปฺปชฺชติ เหตุปจฺจยา.\csromanspacer{}\csromanspacer{}เหตุยา ตีณิ.}

\proseitem{54}{นนกามาวจรํ นกุสลํ ธมฺมํ ปจฺจยา นกามาวจโร กุสโล ธมฺโม อุปฺปชฺชติ เหตุปจฺจยา.\csromanspacer{}\csromanspacer{}นนกามาวจรํ นกุสลํ ธมฺมํ ปจฺจยา กามาวจโร กุสโล ธมฺโม อุปฺปชฺชติ เหตุปจฺจยา.\csromanspacer{}\csromanspacer{}เหตุยา เทฺว.}

\proseitem{55}{นรูปาวจรํ นกุสลํ ธมฺมํ ปจฺจยา รูปาวจโร กุสโล ธมฺโม อุปฺปชฺชติ เหตุปจฺจยา.\csromanspacer{}\csromanspacer{}นรูปาวจรํ นกุสลํ ธมฺมํ ปจฺจยา นรูปาวจโร กุสโล ธมฺโม อุปฺปชฺชติ เหตุปจฺจยา.\csromanspacer{}\csromanspacer{}เหตุยา เทฺว.}

\proseitem{56}{น-อรูปาวจรํ นกุสลํ ธมฺมํ ปจฺจยา อรูปาวจโร กุสโล ธมฺโม อุปฺปชฺชติ เหตุปจฺจยา.\csromanspacer{}\csromanspacer{}น-อรูปาวจรํ นกุสลํ ธมฺมํ ปจฺจยา น-อรูปาวจโร กุสโล ธมฺโม อุปฺปชฺชติ เหตุปจฺจยา.\csromanspacer{}\csromanspacer{}เหตุยา เทฺว.}

\proseitem{57}{น-อปริยาปนฺนํ นกุสลํ ธมฺมํ ปจฺจยา อปริยาปนฺโน กุสโล ธมฺโม อุปฺปชฺชติ เหตุปจฺจยา.\csromanspacer{}\csromanspacer{}น-อปริยาปนฺนํ นกุสลํ ธมฺมํ ปจฺจยา ปริยาปนฺโน กุสโล ธมฺโม อุปฺปชฺชติ เหตุปจฺจยา.\csromanspacer{}\csromanspacer{}เหตุยา เทฺว.}

\proseitem{58}{นนิยฺยานิกํ นกุสลํ ธมฺมํ ปจฺจยา นิยฺยานิโก กุสโล ธมฺโม อุปฺปชฺชติ เหตุปจฺจยา.\csromanspacer{}\csromanspacer{}นนิยฺยานิกํ นกุสลํ ธมฺมํ ปจฺจยา น-อนิยฺยานิโก กุสโล ธมฺโม อุปฺปชฺชติ เหตุปจฺจยา.\csromanspacer{}\csromanspacer{}เหตุยา เทฺว.}

\csromanheader{2}{1. เหตุทุก 3. วิปากตฺติก}
\proseitem{59}{\csromanfolio{359}นนิยตํ นกุสลํ ธมฺมํ ปจฺจยา นิยโต กุสโล ธมฺโม อุปฺปชฺชติ เหตุปจฺจยา.\csromanspacer{}\csromanspacer{}นนิยตํ นกุสลํ ธมฺมํ ปจฺจยา อนิยโต กุสโล ธมฺโม อุปฺปชฺชติ เหตุปจฺจยา.\csromanspacer{}\csromanspacer{}เหตุยา เทฺว.}

\proseitem{60}{น-อนุตฺตรํ นกุสลํ ธมฺมํ ปจฺจยา อนุตฺตโร กุสโล ธมฺโม อุปฺปชฺชติ เหตุปจฺจยา.\csromanspacer{}\csromanspacer{}น-อนุตฺตรํ นกุสลํ ธมฺมํ ปจฺจยา ส-อุตฺตโร กุสโล ธมฺโม อุปฺปชฺชติ เหตุปจฺจยา.\csromanspacer{}\csromanspacer{}เหตุยา เทฺว.}

\proseitem{61}{นสรณํ นกุสลํ ธมฺมํ ปจฺจยา อรโณ กุสโล ธมฺโม อุปฺปชฺชติ เหตุปจฺจยา.\csromanspacer{}\csromanspacer{}เหตุยา เอกํ.}

\prose{นสรณํ น-อกุสลํ ธมฺมํ ปจฺจยา สรโณ อกุสโล ธมฺโม อุปฺปชฺชติ เหตุปจฺจยา.\csromanspacer{}\csromanspacer{}เหตุยา เอกํ.}

\prose{นสรณํ น-อพฺยากตํ ธมฺมํ ปฏิจฺจ อรโณ อพฺยากโต ธมฺโม อุปฺปชฺชติ เหตุปจฺจยา.\csromanspacer{}\csromanspacer{}น-อรณํ น-อพฺยากตํ ธมฺมํ ปฏิจฺจ อรโณ อพฺยากโต ธมฺโม อุปฺปชฺชติ เหตุปจฺจยา.\csromanspacer{}\csromanspacer{}เหตุยา เทฺว.}

\csromanheader{2}{1. เหตุทุก 2. เวทนาตฺติก}
\proseitem{62}{นเหตุํ นสุขาย เวทนาย สมฺปยุตฺตํ ธมฺมํ ปฏิจฺจ เหตุ สุขาย เวทนาย สมฺปยุตฺโต ธมฺโม อุปฺปชฺชติ เหตุปจฺจยา.\csromanspacer{}\csromanspacer{}นเหตุํ นสุขาย เวทนาย สมฺปยุตฺตํ ธมฺมํ ปฏิจฺจ นเหตุ สุขาย เวทนาย สมฺปยุตฺโต ธมฺโม อุปฺปชฺชติ เหตุปจฺจยา.\csromanspacer{}\csromanspacer{}นเหตุํ นสุขาย เวทนาย สมฺปยุตฺตํ ธมฺมํ ปฏิจฺจ เหตุ สุขาย เวทนาย สมฺปยุตฺโต จ นเหตุ สุขาย เวทนาย สมฺปยุตฺโต จ ธมฺมา อุปฺปชฺชนฺติ เหตุปจฺจยา.\csromanspacer{}\csromanspacer{}เหตุยา ตีณิ.\csromanspacer{}\csromanspacer{}นเหตุ นทุกฺขาย เวทนาย สมฺปยุตฺตมูลํ ตีณิเยว.}

\prose{นเหตุํ น-อทุกฺขามสุขาย เวทนาย สมฺปยุตฺตํ ธมฺมํ ปฏิจฺจ เหตุ อทุกฺขมสุขาย เวทนาย สมฺปยุตฺโต ธมฺโม อุปฺปชฺชติ เหตุปจฺจยา.\csromanspacer{}\csromanspacer{}เหตุยา ตีณิ.}

\csromanheader{2}{1. เหตุทุก 3. วิปากตฺติก}
\proseitem{63}{นเหตุํ นวิปากํ ธมฺมํ ปฏิจฺจ เหตุ วิปาโก ธมฺโม อุปฺปชฺชติ เหตุปจฺจยา.\csromanspacer{}\csromanspacer{}เหตุยา ตีณิ.}

\csromanheader{2}{ธมฺมปจฺจนียานุโลม ทุกติกปฏฺฐานปาฬิ}
\prose{\csromanfolio{360}นเหตุํ นวิปากธมฺมธมฺมํ ปจฺจยา เหตุ วิปากธมฺมธมฺโม อุปฺปชฺชติ เหตุปจฺจยา.\csromanspacer{}\csromanspacer{}เหตุยา ตีณิ.}

\prose{นเหตุํ นเนววิปากนวิปากธมฺมธมฺมํ ปฏิจฺจ นเหตุ เนววิปากนวิปากธมฺมธมฺโม อุปฺปชฺชติ เหตุปจฺจยา.\csromanspacer{}\csromanspacer{}เหตุยา ตีณิ.}

\csromanheader{2}{1. เหตุทุก 4. อุปาทินฺนตฺติก}
\proseitem{64}{นเหตุ น-อุปาทินฺนุปาทานิโย ธมฺโม เหตุสฺส อุปาทินฺนุปาทานิยสฺส ธมฺมสฺส อารมฺมณปจฺจเยน ปจฺจโย.\csromanspacer{}\csromanspacer{}นเหตุ น-อุปาทินฺนุปาทานิโย ธมฺโม นเหตุสฺส อุปาทินฺนุปาทานิยสฺส ธมฺมสฺส อารมฺมณปจฺจเยน ปจฺจโย.\csromanspacer{}\csromanspacer{}นเหตุ น-อุปาทินฺนุปาทานิโย ธมฺโม เหตุสฺส อุปาทินฺนุปาทานิยสฺส จ นเหตุสฺส อุปาทินฺนุปาทานิยสฺส จ ธมฺมสฺส อารมฺมณปจฺจเยน ปจฺจโย. (3)}

\prose{นนเหตุ น-อุปาทินฺนุปาทานิโย ธมฺโม นเหตุสฺส อุปาทินฺนุปาทานิยสฺส ธมฺมสฺส อารมฺมณปจฺจเยน ปจฺจโย.\csromanspacer{} ตีณิ.}

\prose{นเหตุ น-อุปาทินฺนุปาทานิโย จ นนเหตุ น-อุปาทินฺนุปาทานิโย จ ธมฺมา เหตุสฺส อุปาทินฺนุปาทานิยสฺส ธมฺมสฺส อารมฺมณปจฺจเยน ปจฺจโย.\csromanspacer{} ตีณิ. (นว ปญฺหา.\csromanspacer{} สํขิตฺตํ.)}

\prose{นเหตุํ น-อนุปาทินฺนุปาทานิยํ ธมฺมํ ปจฺจยา เหตุ อนุปาทินฺนุปาทานิโย ธมฺโม อุปฺปชฺชติ เหตุปจฺจยา.\csromanspacer{}\csromanspacer{}เหตุยา ตีณิ.}

\prose{นเหตุํ น-อนุปาทินฺน-อนุปาทานิยํ ธมฺมํ ปจฺจยา เหตุ อนุปาทินฺน-อนุปาทานิโย ธมฺโม อุปฺปชฺชติ เหตุปจฺจยา.\csromanspacer{}\csromanspacer{}เหตุยา ตีณิ.}

\csromanheader{2}{1. เหตุทุก 5. สํกิลิฏฺฐตฺติก}
\proseitem{65}{นเหตุํ นสํกิลิฏฺฐสํกิเลสิกํ ธมฺมํ ปจฺจยา เหตุ สํกิลิฏฺฐสํกิเลสิโก ธมฺโม อุปฺปชฺชติ เหตุปจฺจยา.\csromanspacer{}\csromanspacer{}เหตุยา ตีณิ.}

\prose{นเหตุํ น-อสํกิลิฏฺฐสํกิเลสิกํ ธมฺมํ ปฏิจฺจ นเหตุ อสํกิลิฏฺฐสํกิเลสิโก ธมฺโม อุปฺปชฺชติ เหตุปจฺจยา.\csromanspacer{}\csromanspacer{}เหตุยา ตีณิ.}

\prose{นเหตุํ น-อสํกิลิฏฺฐ-อสํกิเลสิกํ ธมฺมํ ปจฺจยา เหตุ อสํกิลิฏฺฐ-อสํกิเลสิโก ธมฺโม อุปฺปชฺชติ เหตุปจฺจยา.\csromanspacer{}\csromanspacer{}เหตุยา ตีณิ.}

\csromanheader{2}{1. เหตุทุก 9. ทสฺสนเหตุตฺติก \textbf{1. เหตุทุก 6. วิตกฺกตฺติก}}
\proseitem{66}{\csromanfolio{361}นเหตุํ นสวิตกฺกสวิจารํ ธมฺมํ ปฏิจฺจ เหตุ สวิตกฺกสวิจาโร ธมฺโม อุปฺปชฺชติ เหตุปจฺจยา.\csromanspacer{}\csromanspacer{}เหตุยา ตีณิ.}

\prose{นเหตุํ น-อวิตกฺกวิจารมตฺตํ ธมฺมํ ปฏิจฺจ เหตุ อวิตกฺกวิจารมตฺโต ธมฺโม อุปฺปชฺชติ เหตุปจฺจยา.\csromanspacer{}\csromanspacer{}เหตุยา ตีณิ.}

\prose{นเหตุํ น-อวิตกฺก-อวิจารํ ธมฺมํ ปฏิจฺจ นเหตุ อวิตกฺก-อวิจาโร ธมฺโม อุปฺปชฺชติ เหตุปจฺจยา.\csromanspacer{}\csromanspacer{}เหตุยา ตีณิ.}

\csromanheader{2}{1. เหตุทุก 7. ปีติตฺติก}
\proseitem{67}{นเหตุํ นปีติสหคตํ ธมฺมํ ปฏิจฺจ เหตุ ปีติสหคโต ธมฺโม อุปฺปชฺชติ เหตุปจฺจยา.\csromanspacer{}\csromanspacer{}เหตุยา ตีณิ.}

\prose{นเหตุํ นสุขสหคตํ ธมฺมํ ปฏิจฺจ เหตุ สุขสหคโต ธมฺโม อุปฺปชฺชติ เหตุปจฺจยา.\csromanspacer{}\csromanspacer{}เหตุยา ตีณิ.}

\prose{นเหตุํ น-อุเปกฺขาสหคตํ ธมฺมํ ปฏิจฺจ เหตุ อุเปกฺขาสหคโต ธมฺโม อุปฺปชฺชติ เหตุปจฺจยา.\csromanspacer{}\csromanspacer{}เหตุยา ตีณิ.}

\csromanheader{2}{1. เหตุทุก 8. ทสฺสนตฺติก}
\proseitem{68}{นเหตุํ นทสฺสเนน ปหาตพฺพํ ธมฺมํ ปจฺจยา เหตุ ทสฺสเนน ปหาตพฺโพ ธมฺโม อุปฺปชฺชติ เหตุปจฺจยา.\csromanspacer{}\csromanspacer{}เหตุยา ตีณิ.}

\prose{นเหตุํ นภาวนาย ปหาตพฺพํ ธมฺมํ ปจฺจยา เหตุ ภาวนาย ปหาตพฺโพ ธมฺโม อุปฺปชฺชติ เหตุปจฺจยา.\csromanspacer{}\csromanspacer{}เหตุยา ตีณิ.}

\proseitem{69}{นเหตุํ นเนวทสฺสเนน นภาวนาย ปหาตพฺพํ ธมฺมํ ปฏิจฺจ นเหตุ เนวทสฺสเนน นภาวนาย ปหาตพฺโพ ธมฺโม อุปฺปชฺชติ เหตุปจฺจยา.\csromanspacer{}\csromanspacer{}นนเหตุํ นเนวทสฺสเนน นภาวนาย ปหาตพฺพํ ธมฺมํ ปฏิจฺจ นเหตุ เนวทสฺสเนน นภาวนาย ปหาตพฺโพ ธมฺโม อุปฺปชฺชติ เหตุปจฺจยา.\csromanspacer{}\csromanspacer{}(ทุกมูลํ เอกํ สํขิตฺตํ.\csromanspacer{} สพฺพตฺถ ตีณิ.)}

\csromanheader{2}{1. เหตุทุก 9. ทสฺสนเหตุตฺติก}
\proseitem{70}{นนเหตุํ นทสฺสเนน ปหาตพฺพเหตุกํ ธมฺมํ ปฏิจฺจ นเหตุ ทสฺสเนน ปหาตพฺพเหตุโก ธมฺโม อุปฺปชฺชติ เหตุปจฺจยา.\csromanspacer{}\csromanspacer{}เหตุยา เอกํ.}

\csromanheader{2}{ธมฺมปจฺจนียานุโลม ทุกติกปฏฺฐานปาฬิ}
\prose{\csromanfolio{362}นนเหตุํ นภาวนาย ปหาตพฺพเหตุกํ ธมฺมํ ปฏิจฺจ นเหตุ ภาวนาย ปหาตพฺพเหตุโก ธมฺโม อุปฺปชฺชติ เหตุปจฺจยา.\csromanspacer{}\csromanspacer{}เหตุยา เอกํ.}

\proseitem{71}{นเหตุํ นเนวทสฺสเนน นภาวนาย ปหาตพฺพเหตุกํ ธมฺมํ ปฏิจฺจ นเหตุ เนวทสฺสเนน นภาวนาย ปหาตพฺพเหตุโก ธมฺโม อุปฺปชฺชติ เหตุปจฺจยา.\csromanspacer{}\csromanspacer{}นนเหตุํ นเนวทสฺสเนน นภาวนาย ปหาตพฺพเหตุกํ ธมฺมํ ปฏิจฺจ นเหตุ เนวทสฺสเนน นภาวนาย ปหาตพฺพเหตุโก ธมฺโม อุปฺปชฺชติ เหตุปจฺจยา.\csromanspacer{} คณิตเกน ตีณิ.}

\csromanheader{2}{1. เหตุทุก 10. อาจยคามิตฺติก}
\proseitem{72}{นเหตุํ น-อาจยคามิํ ธมฺมํ ปจฺจยา เหตุ อาจยคามี ธมฺโม อุปฺปชฺชติ เหตุปจฺจยา.\csromanspacer{}\csromanspacer{}เหตุยา ตีณิ.}

\prose{นเหตุํ น-อปจยคามิํ ธมฺมํ ปจฺจยา เหตุ อปจยคามี ธมฺโม อุปฺปชฺชติ เหตุปจฺจยา.\csromanspacer{}\csromanspacer{}เหตุยา ตีณิ.}

\prose{นเหตุํ นเนวาจยคามินาปจยคามิํ ธมฺมํ ปฏิจฺจ นเหตุ เนวาจยคามินาปจยคามี ธมฺโม อุปฺปชฺชติ เหตุปจฺจยา.\csromanspacer{}\csromanspacer{}นนเหตุํ นเนวาจยคามินาปจยคามิํ ธมฺมํ ปฏิจฺจ นเหตุ เนวาจยคามินาปจยคามี ธมฺโม อุปฺปชฺชติ เหตุปจฺจยา.\csromanspacer{} คณิตเกน ตีณิ.}

\csromanheader{2}{1. เหตุทุก 11. เสกฺขตฺติก}
\proseitem{73}{นเหตุํ นเสกฺขํ ธมฺมํ ปจฺจยา เหตุ เสกฺโข ธมฺโม อุปฺปชฺชติ เหตุปจฺจยา.\csromanspacer{}\csromanspacer{}เหตุยา ตีณิ.}

\prose{นเหตุํ น-อเสกฺขํ ธมฺมํ ปจฺจยา เหตุ อเสกฺโข ธมฺโม อุปฺปชฺชติ เหตุปจฺจยา.\csromanspacer{}\csromanspacer{}เหตุยา ตีณิ.}

\prose{นเหตุํ นเนวเสกฺขนาเสกฺขํ ธมฺมํ ปฏิจฺจ นเหตุ เนวเสกฺขนาเสกฺโข ธมฺโม อุปฺปชฺชติ เหตุปจฺจยา.\csromanspacer{}\csromanspacer{}นนเหตุํ นเนวเสกฺขนาเสกฺขํ ธมฺมํ ปฏิจฺจ นเหตุ เนวเสกฺขนาเสกฺโข ธมฺโม อุปฺปชฺชติ เหตุปจฺจยา.\csromanspacer{} คณิตเกน ตีณิ.}

\vspace{\tipitakaparskip}
\csromancenter{เหตุทุก 15.\csromanspacer{} มิจฺฉตฺตนิยตตฺติก \textbf{1.}\csromanspacer{}\textbf{ เหตุทุก 12.}\csromanspacer{}\textbf{ ปริตฺตตฺติก}}
\proseitem{74}{\csromanfolio{363}นเหตุํ นปริตฺตํ ธมฺมํ ปฏิจฺจ นเหตุ ปริตฺโต ธมฺโม อุปฺปชฺชติ เหตุปจฺจยา.\csromanspacer{}\csromanspacer{}นนเหตุํ นปริตฺตํ ธมฺมํ ปฏิจฺจ นเหตุ ปริตฺโต ธมฺโม อุปฺปชฺชติ เหตุปจฺจยา.\csromanspacer{} คณิตเกน ตีณิ.}

\prose{นเหตุํ นมหคฺคตํ ธมฺมํ ปจฺจยา เหตุ มหคฺคโต ธมฺโม อุปฺปชฺชติ เหตุปจฺจยา.\csromanspacer{}\csromanspacer{}เหตุยา ตีณิ.}

\prose{นเหตุํ น-อปฺปมาณํ ธมฺมํ ปจฺจยา เหตุ อปฺปมาโณ ธมฺโม อุปฺปชฺชติ เหตุปจฺจยา.\csromanspacer{}\csromanspacer{}เหตุยา ตีณิ.}

\csromanheader{2}{1. เหตุทุก 13. ปริตฺตารมฺมณตฺติก}
\proseitem{75}{นเหตุํ นปริตฺตารมฺมณํ ธมฺมํ ปฏิจฺจ เหตุ ปริตฺตารมฺมโณ ธมฺโม อุปฺปชฺชติ เหตุปจฺจยา.\csromanspacer{}\csromanspacer{}เหตุยา ตีณิ.}

\prose{นเหตุํ นมหคฺคตารมฺมณํ ธมฺมํ ปจฺจยา เหตุ มหคฺคตารมฺมโณ ธมฺโม อุปฺปชฺชติ เหตุปจฺจยา.\csromanspacer{}\csromanspacer{}เหตุยา ตีณิ.}

\prose{นเหตุํ น-อปฺปมาณารมฺมณํ ธมฺมํ ปจฺจยา เหตุ อปฺปมาณารมฺมโณ ธมฺโม อุปฺปชฺชติ เหตุปจฺจยา.\csromanspacer{}\csromanspacer{}เหตุยา ตีณิ.}

\csromanheader{2}{1. เหตุทุก 14. หีนตฺติก}
\proseitem{76}{นเหตุํ นหีนํ ธมฺมํ ปจฺจยา เหตุ หีโน ธมฺโม อุปฺปชฺชติ เหตุปจฺจยา.\csromanspacer{}\csromanspacer{}เหตุยา ตีณิ.}

\prose{นเหตุํ นมชฺฌิมํ ธมฺมํ ปฏิจฺจ นเหตุ มชฺฌิโม ธมฺโม อุปฺปชฺชติ เหตุปจฺจยา.\csromanspacer{}\csromanspacer{}นนเหตุ นมชฺฌิมํ ธมฺมํ ปฏิจฺจ นเหตุ มชฺฌิโม ธมฺโม อุปฺปชฺชติ เหตุปจฺจยา.\csromanspacer{} คณิตเกน ตีณิ.}

\prose{นเหตุํ นปณีตํ ธมฺมํ ปจฺจยา เหตุ ปณีโต ธมฺโม อุปฺปชฺชติ เหตุปจฺจยา.\csromanspacer{}\csromanspacer{}เหตุยา ตีณิ.}

\csromanheader{2}{1. เหตุทุก 15. มิจฺฉตฺตนิยตตฺติก}
\proseitem{77}{นเหตุํ นมิจฺฉตฺตนิยตํ ธมฺมํ ปจฺจยา เหตุ มิจฺฉตฺตนิยโต ธมฺโม อุปฺปชฺชติ เหตุปจฺจยา.\csromanspacer{}\csromanspacer{}เหตุยา ตีณิ.}

\csromanheader{2}{ธมฺมปจฺจนียานุโลม ทุกติกปฏฺฐานปาฬิ}
\prose{\csromanfolio{364}นเหตุํ นสมฺมตฺตนิยตํ ธมฺมํ ปจฺจยา เหตุ สมฺมตฺตนิยโต ธมฺโม อุปฺปชฺชติ เหตุปจฺจยา.\csromanspacer{}\csromanspacer{}เหตุยา ตีณิ.}

\prose{นเหตุํ น-อนิยตํ ธมฺมํ ปฏิจฺจ นเหตุ อนิยโต ธมฺโม อุปฺปชฺชติ เหตุปจฺจยา.\csromanspacer{}\csromanspacer{}นนเหตุํ น-อนิยตํ ธมฺมํ ปฏิจฺจ นเหตุ อนิยโต ธมฺโม อุปฺปชฺชติ เหตุปจฺจยา.\csromanspacer{} คณิตเกน ตีณิ.}

\csromanheader{2}{1. เหตุทุก 16. มคฺคารมฺมณตฺติก}
\proseitem{78}{นเหตุํ นมคฺคารมฺมณํ ธมฺมํ ปจฺจยา เหตุ มคฺคารมฺมโณ ธมฺโม อุปฺปชฺชติ เหตุปจฺจยา.\csromanspacer{}\csromanspacer{}เหตุยา ตีณิ.}

\prose{นเหตุํ นมคฺคเหตุกํ ธมฺมํ ปจฺจยา เหตุ มคฺคเหตุโก ธมฺโม อุปฺปชฺชติ เหตุปจฺจยา.\csromanspacer{}\csromanspacer{}เหตุยา ตีณิ.}

\prose{นเหตุํ นมคฺคาธิปติํ ธมฺมํ ปจฺจยา เหตุ มคฺคาธิปติ ธมฺโม อุปฺปชฺชติ เหตุปจฺจยา.\csromanspacer{}\csromanspacer{}เหตุยา ตีณิ.}

\csromanheader{2}{1. เหตุทุก 17. อุปฺปนฺนตฺติก}
\proseitem{79}{นเหตุ น-อุปฺปนฺโน ธมฺโม เหตุสฺส อุปฺปนฺนสฺส ธมฺมสฺส อารมฺมณปจฺจเยน ปจฺจโย.\csromanspacer{}\csromanspacer{}อารมฺมเณ นว.}

\csromanheader{2}{1. เหตุทุก 18. อตีตตฺติก}
\proseitem{80}{นเหตุ นปจฺจุปฺปนฺโน ธมฺโม เหตุสฺส ปจฺจุปฺปนฺนสฺส ธมฺมสฺส อารมฺมณปจฺจเยน ปจฺจโย.\csromanspacer{}\csromanspacer{}อารมฺมเณ นว.}

\csromanheader{2}{1. เหตุทุก 19. อตีตารมฺมณตฺติก}
\proseitem{81}{นเหตุํ น-อตีตารมฺมณํ ธมฺมํ ปฏิจฺจ เหตุ อตีตารมฺมโณ ธมฺโม อุปฺปชฺชติ เหตุปจฺจยา.\csromanspacer{}\csromanspacer{}เหตุยา ตีณิ.}

\prose{นเหตุํ น-อนาคตารมฺมณํ ธมฺมํ ปจฺจยา เหตุ อนาคตารมฺมโณ ธมฺโม อุปฺปชฺชติ เหตุปจฺจยา.\csromanspacer{}\csromanspacer{}เหตุยา ตีณิ.}

\prose{นเหตุํ นปจฺจุปฺปนฺนารมฺมณํ ธมฺมํ ปฏิจฺจ เหตุ ปจฺจุปฺปนฺนารมฺมโณ ธมฺโม อุปฺปชฺชติ เหตุปจฺจยา.\csromanspacer{}\csromanspacer{}เหตุยา ตีณิ.}

\vspace{\tipitakaparskip}
\csromancenter{สเหตุกทุก 22.\csromanspacer{} สนิทสฺสนตฺติก \textbf{1.}\csromanspacer{}\textbf{ เหตุทุก 20.}\csromanspacer{}\textbf{ อชฺฌตฺตตฺติก}}
\proseitem{82}{\csromanfolio{365}นเหตุ น-อชฺฌตฺโต ธมฺโม เหตุสฺส อชฺฌตฺตสฺส ธมฺมสฺส อารมฺมณปจฺจเยน ปจฺจโย. (สํขิตฺตํ.) .\csromanspacer{} อารมฺมเณ นว, อธิปติยา อุปนิสฺสเย นว, ปุเรชาเต อตฺถิยา อวิคเต ตีณิ.}

\prose{นเหตุ นพหิทฺธา ธมฺโม เหตุสฺส พหิทฺธา ธมฺมสฺส อารมฺมณปจฺจเยน ปจฺจโย. (สํขิตฺตํ.) .\csromanspacer{} อารมฺมเณ นว, อธิปติยา อุปนิสฺสเย นว, ปุเรชาเต อตฺถิยา อวิคเต ตีณิ. (อชฺฌตฺตตฺติโก น ลพฺภติ ปฏิจฺจวาราทีสุ.)}

\csromanheader{2}{1. เหตุทุก 21. อชฺฌตฺตารมฺมณตฺติก}
\proseitem{83}{นเหตุํ น-อชฺฌตฺตารมฺมณํ ธมฺมํ ปจฺจยา เหตุ อชฺฌตฺตารมฺมโณ ธมฺโม อุปฺปชฺชติ เหตุปจฺจยา.\csromanspacer{}\csromanspacer{}เหตุยา ตีณิ.}

\prose{นเหตุํ นพหิทฺธารมฺมณํ ธมฺมํ ปจฺจยา เหตุ พหิทฺธารมฺมโณ ธมฺโม อุปฺปชฺชติ เหตุปจฺจยา.\csromanspacer{}\csromanspacer{}เหตุยา ตีณิ. (อชฺฌตฺตพหิทฺธารมฺมณํ นตฺถิ.)}

\csromanheader{2}{1. เหตุทุก 22. สนิทสฺสนตฺติก}
\proseitem{84}{นเหตุํ นสนิทสฺสนสปฺปฏิฆํ ธมฺมํ ปฏิจฺจ นเหตุ สนิทสฺสนสปฺปฏิโฆ ธมฺโม อุปฺปชฺชติ เหตุปจฺจยา.\csromanspacer{} เอกํ.}

\prose{นนเหตุํ นสนิทสฺสนสปฺปฏิฆํ ธมฺมํ ปฏิจฺจ นเหตุ สนิทสฺสนสปฺปฏิโฆ ธมฺโม อุปฺปชฺชติ เหตุปจฺจยา.\csromanspacer{} เอกํ. (นเหตุน-อนิทสฺสนสปฺปฏิฆมูเลปิ ตีณิเยว.)}

\proseitem{85}{นเหตุํ น-อนิทสฺสน-อปฺปฏิฆํ ธมฺมํ ปฏิจฺจ นเหตุ อนิทสฺสน-อปฺปฏิโฆ ธมฺโม อุปฺปชฺชติ เหตุปจฺจยา.\csromanspacer{}\csromanspacer{}เหตุยา เอกํ.}

\csromanheader{2}{2. สเหตุกทุก 22. สนิทสฺสนตฺติก}
\proseitem{86}{นสเหตุกํ นสนิทสฺสนสปฺปฏิฆํ ธมฺมํ ปฏิจฺจ อเหตุโก สนิทสฺสนสปฺปฏิโฆ ธมฺโม อุปฺปชฺชติ เหตุปจฺจยา.\csromanspacer{} เอกํ.}

\prose{น-อเหตุกํ นสนิทสฺสนสปฺปฏิฆํ ธมฺมํ ปฏิจฺจ อเหตุโก สนิทสฺสนสปฺปฏิโฆ ธมฺโม อุปฺปชฺชติ เหตุปจฺจยา.\csromanspacer{} เอกํ.\csromanspacer{} คณิตเกน ตีณิ ปญฺหา.}

\vspace{\tipitakaparskip}
\csromancenter{(นสเหตุกน-อนิทสฺสนสปฺปฏิฆมูเลปิ ตีณิเยว.)}
\csromanheader{2}{ธมฺมปจฺจนียานุโลม ทุกติกปฏฺฐานปาฬิ}
\prose{\csromanfolio{366}นสเหตุกํ น-อนิทสฺสน-อปฺปฏิฆํ ธมฺมํ ปฏิจฺจ อเหตุโก อนิทสฺสน-อปฺปฏิโฆ ธมฺโม อุปฺปชฺชติ เหตุปจฺจยา.\csromanspacer{}\csromanspacer{}เหตุยา เอกํ.}

\prose{(นเหตุสมฺปยุตฺตํ นสนิทสฺสนสปฺปฏิฆํ สเหตุกทุกสทิสํ.\csromanspacer{} ตีณิ ปญฺหา.\csromanspacer{} เหตุสเหตุกทุเก จ เหตุเหตุสมฺปยุตฺตทุเก จ ปญฺหา โน ลพฺภติ.) \textbf{6.}\csromanspacer{}\textbf{ นเหตุสเหตุกทุก 22.}\csromanspacer{}\textbf{ สนิทสฺสนตฺติก}}

\proseitem{87}{นเหตุํ นสเหตุกํ นสนิทสฺสนสปฺปฏิฆํ ธมฺมํ ปฏิจฺจ นเหตุ อเหตุโก สนิทสฺสนสปฺปฏิโฆ ธมฺโม อุปฺปชฺชติ เหตุปจฺจยา.\csromanspacer{}\csromanspacer{}เหตุยา เอกํ.}

\prose{นเหตุํ น-อเหตุกํ นสนิทสฺสนสปฺปฏิฆํ ธมฺมํ ปฏิจฺจ นเหตุ อเหตุโก สนิทสฺสนสปฺปฏิโฆ ธมฺโม อุปฺปชฺชติ เหตุปจฺจยา.\csromanspacer{} เอกํ.\csromanspacer{}\csromanspacer{}คณิตเกน ตีณิ.}

\prose{นเหตุํ นสเหตุกํ น-อนิทสฺสนสปฺปฏิฆํ ธมฺมํ ปฏิจฺจ.\csromanspacer{} ตีณิเยว.}

\prose{นเหตุํ นสเหตุกํ น-อนิทสฺสน-อปฺปฏิฆํ ธมฺมํ ปฏิจฺจ นเหตุ อเหตุโก อนิทสฺสน-อปฺปฏิโฆ ธมฺโม อุปฺปชฺชติ เหตุปจฺจยา.\csromanspacer{}\csromanspacer{}เหตุยา เอกํ. \textbf{7.}\csromanspacer{}\textbf{ สปฺปจฺจยทุก 22.}\csromanspacer{}\textbf{ สนิทสฺสนตฺติก}}

\proseitem{88}{น-อปฺปจฺจยํ นสนิทสฺสนสปฺปฏิฆํ ธมฺมํ ปฏิจฺจ สปฺปจฺจโย สนิทสฺสนสปฺปฏิโฆ ธมฺโม อุปฺปชฺชติ เหตุปจฺจยา.\csromanspacer{}\csromanspacer{}เหตุยา เอกํ.}

\prose{น-อปฺปจฺจยํ น-อนิทสฺสนสปฺปฏิฆํ ธมฺมํ ปฏิจฺจ สปฺปจฺจโย อนิทสฺสนสปฺปฏิโฆ ธมฺโม อุปฺปชฺชติ เหตุปจฺจยา.\csromanspacer{}\csromanspacer{}เหตุยา เอกํ.}

\prose{(น-อปฺปจฺจยน-อนิทสฺสน-อปฺปฏิฆมูเลปิ เอกเมว.) .\csromanspacer{} เหตุยา เอกํ, อธิปติยา เอกํ ฯเปฯ อวิคเต เอกํ.\csromanspacer{}\csromanspacer{}(นสงฺขตํ นสปฺปจฺจยสทิสํ.) \textbf{9.}\csromanspacer{}\textbf{ สนิทสฺสนทุก 22.}\csromanspacer{}\textbf{ สนิทสฺสนตฺติก}}

\proseitem{89}{นสนิทสฺสนํ นสนิทสฺสนสปฺปฏิฆํ ธมฺมํ ปฏิจฺจ สนิทสฺสโน สนิทสฺสนสปฺปฏิโฆ ธมฺโม อุปฺปชฺชติ เหตุปจฺจยา.\csromanspacer{}\csromanspacer{}เหตุยา เอกํ.}

\csromanheader{2}{12. โลกิยทุก 22. สนิทสฺสนตฺติก}
\prose{\csromanfolio{367}นสนิทสฺสนํ น-อนิทสฺสนสปฺปฏิฆํ ธมฺมํ ปฏิจฺจ อนิทสฺสโน อนิทสฺสนสปฺปฏิโฆ ธมฺโม อุปฺปชฺชติ เหตุปจฺจยา.\csromanspacer{}\csromanspacer{}เหตุยา เอกํ.}

\prose{นสนิทสฺสนํ น-อนิทสฺสน-อปฺปฏิฆํ ธมฺมํ ปฏิจฺจ อนิทสฺสโน อนิทสฺสน-อปฺปฏิโฆ ธมฺโม อุปฺปชฺชติ เหตุปจฺจยา.\csromanspacer{}\csromanspacer{}เหตุยา เอกํ.}

\csromanheader{2}{10. สปฺปฏิฆทุก 22. สนิทสฺสนตฺติก}
\proseitem{90}{นสปฺปฏิฆํ นสนิทสฺสนสปฺปฏิฆํ ธมฺมํ ปฏิจฺจ สปฺปฏิโฆ สนิทสฺสนสปฺปฏิโฆ ธมฺโม อุปฺปชฺชติ เหตุปจฺจยา.\csromanspacer{} เอกํ.}

\prose{น-อปฺปฏิฆํ นสนิทสฺสนสปฺปฏิฆํ ธมฺมํ ปฏิจฺจ สปฺปฏิโฆ สนิทสฺสนสปฺปฏิโฆ ธมฺโม อุปฺปชฺชติ เหตุปจฺจยา.\csromanspacer{} เอกํ.}

\prose{นสปฺปฏิฆํ นสนิทสฺสนสปฺปฏิฆญฺจ น-อปฺปฏิฆํ นสนิทสฺสนสปฺปฏิฆญฺจ ธมฺมํ ปฏิจฺจ สปฺปฏิโฆ สนิทสฺสนสปฺปฏิโฆ ธมฺโม อุปฺปชฺชติ เหตุปจฺจยา.\csromanspacer{} เอกํ. (สพฺพตฺถ ตีณิ.)}

\prose{นสปฺปฏิฆํ น-อนิทสฺสน-อปฺปฏิฆํ ธมฺมํ ปฏิจฺจ สปฺปฏิโฆ อนิทสฺสนสปฺปฏิโฆ ธมฺโม อุปฺปชฺชติ เหตุปจฺจยา.\csromanspacer{}\csromanspacer{}เหตุยา เอกํ.}

\prose{น-อปฺปฏิฆํ น-อนิทสฺสน-อปฺปฏิฆํ ธมฺมํ ปฏิจฺจ อปฺปฏิโฆ อนิทสฺสน-อปฺปฏิโฆ ธมฺโม อุปฺปชฺชติ เหตุปจฺจยา.\csromanspacer{}\csromanspacer{}เหตุยา เอกํ.}

\csromanheader{2}{11. รูปีทุก 22. สนิทสฺสนตฺติก}
\proseitem{91}{นรูปิํ นสนิทสฺสนสปฺปฏิฆํ ธมฺมํ ปฏิจฺจ รูปี สนิทสฺสนสปฺปฏิโฆ ธมฺโม อุปฺปชฺชติ เหตุปจฺจยา.\csromanspacer{}\csromanspacer{}เหตุยา ตีณิ.}

\prose{นรูปิํ น-อนิทสฺสนสปฺปฏิฆํ ธมฺมํ ปฏิจฺจ รูปี อนิทสฺสนสปฺปฏิโฆ ธมฺโม อุปฺปชฺชติ เหตุปจฺจยา.\csromanspacer{}\csromanspacer{}เหตุยา ตีณิ.}

\prose{น-อรูปิํ น-อนิทสฺสน-อปฺปฏิฆํ ธมฺมํ ปฏิจฺจ รูปี อนิทสฺสน-อปฺปฏิโฆ ธมฺโม อุปฺปชฺชติ เหตุปจฺจยา.\csromanspacer{}\csromanspacer{}เหตุยา เอกํ.}

\csromanheader{2}{12. โลกิยทุก 22. สนิทสฺสนตฺติก}
\proseitem{92}{นโลกิยํ นสนิทสฺสนสปฺปฏิฆํ ธมฺมํ ปฏิจฺจ โลกิโย สนิทสฺสนสปฺปฏิโฆ ธมฺโม อุปฺปชฺชติ เหตุปจฺจยา.\csromanspacer{} เอกํ.}

\csromanheader{2}{ธมฺมปจฺจนียานุโลม ทุกติกปฏฺฐานปาฬิ}
\prose{\csromanfolio{368}นโลกุตฺตรํ นสนิทสฺสนสปฺปฏิฆํ ธมฺมํ ปฏิจฺจ โลกิโย สนิทสฺสนสปฺปฏิโฆ ธมฺโม อุปฺปชฺชติ เหตุปจฺจยา.\csromanspacer{} เอกํ.}

\prose{(คณิตเกน คเณตพฺพา ตีณิ ปญฺหา.)}

\prose{นโลกิยํ น-อนิทสฺสนสปฺปฏิฆํ ธมฺมํ ปฏิจฺจ โลกิโย อนิทสฺสนสปฺปฏิโฆ. (ตีณิเยว ปญฺหา กาตพฺพา.)}

\prose{นโลกุตฺตรํ น-อนิทสฺสน-อปฺปฏิฆํ ธมฺมํ ปฏิจฺจ โลกิโย อนิทสฺสน-อปฺปฏิโฆ ธมฺโม อุปฺปชฺชติ เหตุปจฺจยา.\csromanspacer{}\csromanspacer{}เหตุยา เอกํ.}

\csromanheader{2}{13. เกนจิวิญฺเญยฺยทุก 22. สนิทสฺสนตฺติก}
\proseitem{93}{นเกนจิ วิญฺเญยฺยํ นสนิทสฺสนสปฺปฏิฆํ ธมฺมํ ปฏิจฺจ เกนจิ วิญฺเญยฺโย สนิทสฺสนสปฺปฏิโฆ ธมฺโม อุปฺปชฺชติ เหตุปจฺจยา.\csromanspacer{}\csromanspacer{}นเกนจิ วิญฺเญยฺยํ นสนิทสฺสนสปฺปฏิฆํ ธมฺมํ ปฏิจฺจ เกนจิ นวิญฺเญยฺโย สนิทสฺสนสปฺปฏิโฆ ธมฺโม อุปฺปชฺชติ เหตุปจฺจยา.\csromanspacer{}\csromanspacer{}นเกนจิ วิญฺเญยฺยํ นสนิทสฺสนสปฺปฏิฆํ ธมฺมํ ปฏิจฺจ เกนจิ วิญฺเญยฺโย สนิทสฺสนสปฺปฏิโฆ จ เกนจิ นวิญฺเญยฺโย สนิทสฺสนสปฺปฏิโฆ จ ธมฺมา อุปฺปชฺชนฺติ เหตุปจฺจยา. (3)}

\prose{นเกนจิ นวิญฺเญยฺยํ นสนิทสฺสนสปฺปฏิฆํ ธมฺมํ ปฏิจฺจ.\csromanspacer{} ตีณิเยว.}

\prose{นเกนจิ วิญฺเญยฺยํ นสนิทสฺสนสปฺปฏิฆญฺจ นเกนจิ นวิญฺเญยฺยํ นสนิทสฺสนสปฺปฏิฆญฺจ ธมฺมํ ปฏิจฺจ.\csromanspacer{} ตีณิเยว. (สพฺพตฺถ นว.)}

\prose{นเกนจิ วิญฺเญยฺยํ น-อนิทสฺสนสปฺปฏิฆสฺส ปุริมสทิสํ นวปญฺหํ กาตพฺพํ.\csromanspacer{} นเกนจิ วิญฺเญยฺยํ น-อนิทสฺสน-อปฺปฏิฆมูลสฺส นวปญฺหเมว.\csromanspacer{}\csromanspacer{}เหตุยา นว, อธิปติยา นว ฯเปฯ อวิคเต นว.}

\csromanheader{2}{14. อาสวทุก 22. สนิทสฺสนตฺติก}
\proseitem{94}{น-อาสวํ นสนิทสฺสนสปฺปฏิฆํ ธมฺมํ ปฏิจฺจ โน-อาสโว สนิทสฺสนสปฺปฏิโฆ ธมฺโม อุปฺปชฺชติ เหตุปจฺจยา.\csromanspacer{} เอกํ.}

\prose{นโน-อาสวํ นสนิทสฺสนสปฺปฏิฆํ ธมฺมํ ปฏิจฺจ โน-อาสโว สนิทสฺสนสปฺปฏิโฆ ธมฺโม อุปฺปชฺชติ เหตุปจฺจยา.\csromanspacer{} เอกํ.\csromanspacer{} คณิตเกน ตีณิ.}

\csromanheader{2}{17. อาสวสาสวทุก 22. สนิทสฺสนตฺติก}
\prose{\csromanfolio{369}โน-อาสวํ น-อนิทสฺสนสปฺปฏิฆํ ธมฺมํ ปฏิจฺจ. (ปุริมนเยน ตีณิ ปญฺหา.)}

\prose{โน-อาสวํ น-อนิทสฺสน-อปฺปฏิฆํ ธมฺมํ ปฏิจฺจ โน-อาสโว อนิทสฺสน-อปฺปฏิโฆ ธมฺโม อุปฺปชฺชติ เหตุปจฺจยา.\csromanspacer{}\csromanspacer{}เหตุยา เอกํ.}

\csromanheader{2}{15. สาสวทุก 22. สนิทสฺสนตฺติก}
\proseitem{95}{นสาสวํ นสนิทสฺสนสปฺปฏิฆํ ธมฺมํ ปฏิจฺจ สาสโว สนิทสฺสนสปฺปฏิโฆ ธมฺโม อุปฺปชฺชติ เหตุปจฺจยา.\csromanspacer{} เอกํ.}

\prose{น-อนาสวํ นสนิทสฺสนสปฺปฏิฆํ ธมฺมํ ปฏิจฺจ สาสโว สนิทสฺสนสปฺปฏิโฆ ธมฺโม อุปฺปชฺชติ เหตุปจฺจยา.\csromanspacer{} เอกํ.\csromanspacer{} คณิตเกน ตีณิ.}

\prose{(นสาสวํ น-อนิทสฺสนสปฺปฏิฆมูลสฺส ปุริมนเยน ตีณิ ปญฺหา.)}

\prose{น-อนาสวํ น-อนิทสฺสน-อปฺปฏิฆํ ธมฺมํ ปฏิจฺจ สาสโว อนิทสฺสน-อปฺปฏิโฆ ธมฺโม อุปฺปชฺชติ เหตุปจฺจยา.\csromanspacer{}\csromanspacer{}เหตุยา เอกํ.}

\csromanheader{2}{16. อาสวสมฺปยุตฺตทุก 22. สนิทสฺสนตฺติก}
\proseitem{96}{น-อาสวสมฺปยุตฺตํ นสนิทสฺสนสปฺปฏิฆํ ธมฺมํ ปฏิจฺจ อาสววิปฺปยุตฺโต สนิทสฺสนสปฺปฏิโฆ ธมฺโม อุปฺปชฺชติ เหตุปจฺจยา.\csromanspacer{} เอกํ.}

\prose{น-อาสววิปฺปยุตฺตํ นสนิทสฺสนสปฺปฏิฆํ ธมฺมํ ปฏิจฺจ อาสววิปฺปยุตฺโต สนิทสฺสนสปฺปฏิโฆ ธมฺโม อุปฺปชฺชติ เหตุปจฺจยา.\csromanspacer{} เอกํ.\csromanspacer{} คณิตเกน ตีณิ.}

\prose{(น-อาสวสมฺปยุตฺตน-อนิทสฺสนสปฺปฏิฆมูเล ตีณิเยว.)}

\prose{น-อาสวสมฺปยุตฺตํ น-อนิทสฺสน-อปฺปฏิฆํ ธมฺมํ ปฏิจฺจ อาสววิปฺปยุตฺโต อนิทสฺสน-อปฺปฏิโฆ ธมฺโม อุปฺปชฺชติ เหตุปจฺจยา.\csromanspacer{}\csromanspacer{}เหตุยา เอกํ.}

\csromanheader{2}{17. อาสวสาสวทุก 22. สนิทสฺสนตฺติก}
\proseitem{97}{น-อาสวญฺเจว น-อนาสวญฺจ นสนิทสฺสนสปฺปฏิฆํ ธมฺมํ ปฏิจฺจ สาสโว เจว โน จ อาสโว สนิทสฺสนสปฺปฏิโฆ ธมฺโม อุปฺปชฺชติ เหตุปจฺจยา.}

\csromanheader{2}{ธมฺมปจฺจนียานุโลม ทุกติกปฏฺฐานปาฬิ}
\prose{\csromanfolio{370}น-อนาสวญฺเจว นโน จ อาสวํ นสนิทสฺสนสปฺปฏิฆํ ธมฺมํ ปฏิจฺจ สาสโว เจว โน จ อาสโว สนิทสฺสนสปฺปฏิโฆ ธมฺโม อุปฺปชฺชติ เหตุปจฺจยา.\csromanspacer{} คณิตเกน ตีณิ.}

\prose{(น-อาสวญฺเจว น-อนาสวํ น-อนิทสฺสนสปฺปฏิฆมูเลปิ ปุริมนเยน ตีณิเยว.)}

\prose{น-อาสวญฺเจว น-อนาสวญฺจ น-อนิทสฺสน-อปฺปฏิฆํ ธมฺมํ ปฏิจฺจ สาสโว เจว โน จ อาสโว อนิทสฺสน-อปฺปฏิโฆ ธมฺโม อุปฺปชฺชติ เหตุปจฺจยา.\csromanspacer{}\csromanspacer{}เหตุยา เอกํ. (อาสว-อาสวสมฺปยุตฺตทุเก ปญฺหา นลพฺภติ.)}

\csromanheader{2}{19. อาสววิปฺปยุตฺตสาสวทุก 22. สนิทสฺสนตฺติก}
\proseitem{98}{อาสววิปฺปยุตฺตํ นสาสวํ นสนิทสฺสนสปฺปฏิฆํ ธมฺมํ ปฏิจฺจ อาสววิปฺปยุตฺโต สาสโว สนิทสฺสนสปฺปฏิโฆ ธมฺโม อุปฺปชฺชติ เหตุปจฺจยา.\csromanspacer{} เอกํ.}

\prose{อาสววิปฺปยุตฺตํ น-อนาสวํ นสนิทสฺสนสปฺปฏิฆํ ธมฺมํ ปฏิจฺจ อาสววิปฺปยุตฺโต สาสโว สนิทสฺสนสปฺปฏิโฆ ธมฺโม อุปฺปชฺชติ เหตุปจฺจยา.\csromanspacer{} เอกํ.\csromanspacer{} คณิตเกน ตีณิ.}

\prose{(อาสววิปฺปยุตฺตนสาสวน-อนิทสฺสนสปฺปฏิฆมูเลปิ ปุริมนเยเนว ตีณิ ปญฺหา.)}

\prose{อาสววิปฺปยุตฺตํ น-อนาสวํ น-อนิทสฺสน-อปฺปฏิฆํ ธมฺมํ ปฏิจฺจ อาสววิปฺปยุตฺโต สาสโว อนิทสฺสน-อปฺปฏิโฆ ธมฺโม อุปฺปชฺชติ เหตุปจฺจยา.\csromanspacer{}\csromanspacer{}เหตุยา เอกํ.}

\csromanheader{2}{20-54. สญฺโญชนาทิฉโคจฺฉก 22. สนิทสฺสนตฺติก}
\proseitem{99}{โนสญฺโญชนํ นสนิทสฺสนสปฺปฏิฆํ ธมฺมํ ปฏิจฺจ โนสญฺโญชโน สนิทสฺสนสปฺปฏิโฆ ธมฺโม อุปฺปชฺชติ เหตุปจฺจยา. (ตีณิ ปญฺหา.)}

\proseitem{100}{โนคนฺถํ นสนิทสฺสนสปฺปฏิฆํ ธมฺมํ ปฏิจฺจ โนคนฺโถ สนิทสฺสนสปฺปฏิโฆ ธมฺโม อุปฺปชฺชติ เหตุปจฺจยา.\csromanspacer{}\csromanspacer{}เหตุยา ตีณิ.}

\csromanheader{2}{55-68. มหนฺตรทุก 22. สนิทสฺสนตฺติก}
\proseitem{101}{\csromanfolio{371}โน-โอฆํ นสนิทสฺสนสปฺปฏิฆํ ธมฺมํ ปฏิจฺจ โน-โอโฆ สนิทสฺสนสปฺปฏิโฆ ธมฺโม อุปฺปชฺชติ เหตุปจฺจยา.\csromanspacer{}\csromanspacer{}เหตุยา ตีณิ.}

\proseitem{102}{โนโยคํ นสนิทสฺสนสปฺปฏิฆํ ธมฺมํ ปฏิจฺจ โนโยโค สนิทสฺสนสปฺปฏิโฆ ธมฺโม อุปฺปชฺชติ เหตุปจฺจยา.\csromanspacer{}\csromanspacer{}เหตุยา ตีณิ.}

\proseitem{103}{โนนีวรณํ นสนิทสฺสนสปฺปฏิฆํ ธมฺมํ ปฏิจฺจ โนนีวรโณ สนิทสฺสนสปฺปฏิโฆ ธมฺโม อุปฺปชฺชติ เหตุปจฺจยา.\csromanspacer{}\csromanspacer{}เหตุยา ตีณิ.}

\proseitem{104}{โนปรามาสํ นสนิทสฺสนสปฺปฏิฆํ ธมฺมํ ปฏิจฺจ โนปรามาโส สนิทสฺสนสปฺปฏิโฆ ธมฺโม อุปฺปชฺชติ เหตุปจฺจยา.\csromanspacer{}\csromanspacer{}เหตุยา ตีณิ.}

\csromanheader{2}{55-68. มหนฺตรทุก 22. สนิทสฺสนตฺติก}
\proseitem{105}{นสารมฺมณํ นสนิทสฺสนสปฺปฏิฆํ ธมฺมํ ปฏิจฺจ อนารมฺมโณ สนิทสฺสนสปฺปฏิโฆ ธมฺโม อุปฺปชฺชติ เหตุปจฺจยา.\csromanspacer{} เอกํ.}

\prose{น-อนารมฺมณํ นสนิทสฺสนสปฺปฏิฆํ ธมฺมํ ปฏิจฺจ อนารมฺมโณ สนิทสฺสนสปฺปฏิโฆ ธมฺโม อุปฺปชฺชติ เหตุปจฺจยา.\csromanspacer{} เอกํ.\csromanspacer{} คณิตเกน ตีณิ.}

\prose{(อนิทสฺสนสปฺปฏิเฆ ตีณิเยว.)}

\prose{นสารมฺมณํ น-อนิทสฺสน-อปฺปฏิฆํ ธมฺมํ ปฏิจฺจ อนารมฺมโณ อนิทสฺสน-อปฺปฏิโฆ ธมฺโม อุปฺปชฺชติ เหตุปจฺจยา.\csromanspacer{}\csromanspacer{}เหตุยา เอกํ.}

\proseitem{106}{โนจิตฺตํ นสนิทสฺสนสปฺปฏิฆํ ธมฺมํ ปฏิจฺจ โนจิตฺโต สนิทสฺสนสปฺปฏิโฆ ธมฺโม อุปฺปชฺชติ เหตุปจฺจยา.\csromanspacer{} เอกํ.}

\prose{นโนจิตฺตํ นสนิทสฺสนสปฺปฏิฆํ ธมฺมํ ปฏิจฺจ โนจิตฺโต สนิทสฺสนสปฺปฏิโฆ ธมฺโม อุปฺปชฺชติ เหตุปจฺจยา.\csromanspacer{} เอกํ.\csromanspacer{} คณิตเกน ตีณิ.}

\proseitem{107}{นเจตสิกํ นสนิทสฺสนสปฺปฏิฆํ ธมฺมํ ปฏิจฺจ นเจตสิโก สนิทสฺสนสปฺปฏิโฆ ธมฺโม อุปฺปชฺชติ เหตุปจฺจยา.\csromanspacer{}\csromanspacer{}เหตุยา ตีณิ.}

\proseitem{108}{นจิตฺตสมฺปยุตฺตํ นสนิทสฺสนสปฺปฏิฆํ ธมฺมํ ปฏิจฺจ จิตฺตวิปฺปยุตฺโต สนิทสฺสนสปฺปฏิโฆ ธมฺโม อุปฺปชฺชติ เหตุปจฺจยา.\csromanspacer{}\csromanspacer{}เหตุยา ตีณิ.}

\csromanheader{2}{ธมฺมปจฺจนียานุโลม ทุกติกปฏฺฐานปาฬิ}
\proseitem{109}{\csromanfolio{372}นจิตฺตสํสฏฺฐํ นสนิทสฺสนสปฺปฏิฆํ ธมฺมํ ปฏิจฺจ จิตฺตวิสํสฏฺโฐ สนิทสฺสนสปฺปฏิโฆ ธมฺโม อุปฺปชฺชติ เหตุปจฺจยา.\csromanspacer{}\csromanspacer{}เหตุยา ตีณิ.}

\proseitem{110}{โนจิตฺตสมุฏฺฐานํ นสนิทสฺสนสปฺปฏิฆํ ธมฺมํ ปฏิจฺจ จิตฺตสมุฏฺฐาโน สนิทสฺสนสปฺปฏิโฆ ธมฺโม อุปฺปชฺชติ เหตุปจฺจยา. (จิตฺตสมุฏฺฐานรูเปเนว ตีณิ.\csromanspacer{} สํขิตฺตํ.)}

\proseitem{111}{โนจิตฺตสหภุํ นสนิทสฺสนสปฺปฏิฆํ ธมฺมํ ปฏิจฺจ โนจิตฺตสหภู สนิทสฺสนสปฺปฏิโฆ ธมฺโม อุปฺปชฺชติ เหตุปจฺจยา.\csromanspacer{}\csromanspacer{}เหตุยา ตีณิ.}

\proseitem{112}{โนจิตฺตานุปริวตฺติํ นสนิทสฺสนสปฺปฏิฆํ ธมฺมํ ปฏิจฺจ โนจิตฺตานุปริวตฺตี สนิทสฺสนสปฺปฏิโฆ ธมฺโม อุปฺปชฺชติ เหตุปจฺจยา.\csromanspacer{}\csromanspacer{}เหตุยา ตีณิ.}

\proseitem{113}{โนจิตฺตสํสฏฺฐสมุฏฺฐานํ นสนิทสฺสนสปฺปฏิฆํ ธมฺมํ ปฏิจฺจ โนจิตฺตสํสฏฺฐสมุฏฺฐาโน สนิทสฺสนสปฺปฏิโฆ ธมฺโม อุปฺปชฺชติ เหตุปจฺจยา.\csromanspacer{}\csromanspacer{}เหตุยา ตีณิ.}

\proseitem{114}{โนจิตฺตสํสฏฺฐสมุฏฺฐานสหภุํ นสนิทสฺสนสปฺปฏิฆํ ธมฺมํ ปฏิจฺจ โนจิตฺตสํสฏฺฐสมุฏฺฐานสหภู สนิทสฺสนสปฺปฏิโฆ ธมฺโม อุปฺปชฺชติ เหตุปจฺจยา.\csromanspacer{}\csromanspacer{}เหตุยา ตีณิ.}

\proseitem{115}{โนจิตฺตสํสฏฺฐสมุฏฺฐานานุปริวตฺติํ นสนิทสฺสนสปฺปฏิฆํ ธมฺมํ ปฏิจฺจ โนจิตฺตสํสฏฺฐสมุฏฺฐานานุปริวตฺตี สนิทสฺสนสปฺปฏิโฆ ธมฺโม อุปฺปชฺชติ เหตุปจฺจยา.\csromanspacer{}\csromanspacer{}เหตุยา ตีณิ.}

\proseitem{116}{น-อชฺฌตฺติกํ นสนิทสฺสนสปฺปฏิฆํ ธมฺมํ ปฏิจฺจ พาหิโร สนิทสฺสนสปฺปฏิโฆ ธมฺโม อุปฺปชฺชติ เหตุปจฺจยา.\csromanspacer{} เอกํ.}

\prose{นพาหิรํ นสนิทสฺสนสปฺปฏิฆํ ธมฺมํ ปฏิจฺจ พาหิโร สนิทสฺสนสปฺปฏิโฆ ธมฺโม อุปฺปชฺชติ เหตุปจฺจยา.\csromanspacer{} เอกํ.\csromanspacer{} คณิตเกน ตีณิ.}

\proseitem{117}{โน-อุปาทา นสนิทสฺสนสปฺปฏิฆํ ธมฺมํ ปฏิจฺจ อุปาทา สนิทสฺสนสปฺปฏิโฆ ธมฺโม อุปฺปชฺชติ เหตุปจฺจยา.\csromanspacer{}\csromanspacer{}เหตุยา เอกํ.}

\csromanheader{2}{83-100. ปิฏฺฐิทุก 22. สนิทสฺสนตฺติก}
\proseitem{118}{\csromanfolio{373}น-อุปาทินฺนํ นสนิทสฺสนสปฺปฏิฆํ ธมฺมํ ปฏิจฺจ อนุปาทินฺโน สนิทสฺสนสปฺปฏิโฆ ธมฺโม อุปฺปชฺชติ เหตุปจฺจยา.\csromanspacer{}\csromanspacer{}เหตุยา ตีณิ.}

\csromanheader{2}{69. อุปาทานทุก 22. สนิทสฺสนตฺติก}
\proseitem{119}{โน-อุปาทานํ นสนิทสฺสนสปฺปฏิฆํ ธมฺมํ ปฏิจฺจ โน-อุปาทาโน สนิทสฺสนสปฺปฏิโฆ ธมฺโม อุปฺปชฺชติ เหตุปจฺจยา.\csromanspacer{}\csromanspacer{}เหตุยา ตีณิ.}

\csromanheader{2}{75. กิเลสทุก 22. สนิทสฺสนตฺติก}
\proseitem{120}{โนกิเลสํ นสนิทสฺสนสปฺปฏิฆํ ธมฺมํ ปฏิจฺจ โนกิเลโส สนิทสฺสนสปฺปฏิโฆ ธมฺโม อุปฺปชฺชติ เหตุปจฺจยา.\csromanspacer{}\csromanspacer{}เหตุยา ตีณิ.}

\csromanheader{2}{83-100. ปิฏฺฐิทุก 22. สนิทสฺสนตฺติก}
\proseitem{121}{นทสฺสเนน ปหาตพฺพํ นสนิทสฺสนสปฺปฏิฆํ ธมฺมํ ปฏิจฺจ นทสฺสเนน ปหาตพฺโพ สนิทสฺสนสปฺปฏิโฆ ธมฺโม อุปฺปชฺชติ เหตุปจฺจยา.\csromanspacer{}\csromanspacer{}เหตุยา ตีณิ.}

\proseitem{122}{นภาวนาย ปหาตพฺพํ นสนิทสฺสนสปฺปฏิฆํ ธมฺมํ ปฏิจฺจ นภาวนาย ปหาตพฺโพ สนิทสฺสนสปฺปฏิโฆ ธมฺโม อุปฺปชฺชติ เหตุปจฺจยา.\csromanspacer{}\csromanspacer{}เหตุยา ตีณิ.}

\proseitem{123}{นทสฺสเนน ปหาตพฺพเหตุกํ นสนิทสฺสนสปฺปฏิฆํ ธมฺมํ ปฏิจฺจ นทสฺสเนน ปหาตพฺพเหตุโก สนิทสฺสนสปฺปฏิโฆ ธมฺโม อุปฺปชฺชติ เหตุปจฺจยา.\csromanspacer{}\csromanspacer{}เหตุยา ตีณิ.}

\proseitem{124}{นภาวนาย ปหาตพฺพเหตุกํ นสนิทสฺสนสปฺปฏิฆํ ธมฺมํ ปฏิจฺจ นภาวนาย ปหาตพฺพเหตุโก สนิทสฺสนสปฺปฏิโฆ ธมฺโม อุปฺปชฺชติ เหตุปจฺจยา.\csromanspacer{}\csromanspacer{}เหตุยา ตีณิ.}

\proseitem{125}{นสวิตกฺกํ นสนิทสฺสนสปฺปฏิฆํ ธมฺมํ ปฏิจฺจ อวิตกฺโก สนิทสฺสนสปฺปฏิโฆ ธมฺโม อุปฺปชฺชติ เหตุปจฺจยา.\csromanspacer{}\csromanspacer{}เหตุยา ตีณิ.}

\proseitem{126}{นสวิจารํ นสนิทสฺสนสปฺปฏิฆํ ธมฺมํ ปฏิจฺจ อวิจาโร สนิทสฺสนสปฺปฏิโฆ ธมฺโม อุปฺปชฺชติ เหตุปจฺจยา.\csromanspacer{}\csromanspacer{}เหตุยา ตีณิ.}

\csromanheader{2}{ธมฺมปจฺจนียานุโลม ทุกติกปฏฺฐานปาฬิ}
\proseitem{127}{\csromanfolio{374}นสปฺปีติกํ นสนิทสฺสนสปฺปฏิฆํ ธมฺมํ ปฏิจฺจ นสปฺปีติโก สนิทสฺสนสปฺปฏิโฆ ธมฺโม อุปฺปชฺชติ เหตุปจฺจยา.\csromanspacer{}\csromanspacer{}เหตุยา ตีณิ.}

\proseitem{128}{นปีติสหคตํ นสนิทสฺสนสปฺปฏิฆํ ธมฺมํ ปฏิจฺจ นปีติสหคโต สนิทสฺสนสปฺปฏิโฆ ธมฺโม อุปฺปชฺชติ เหตุปจฺจยา.\csromanspacer{}\csromanspacer{}เหตุยา ตีณิ.}

\proseitem{129}{นสุขสหคตํ นสนิทสฺสนสปฺปฏิฆํ ธมฺมํ ปฏิจฺจ นสุขสหคโต สนิทสฺสนสปฺปฏิโฆ ธมฺโม อุปฺปชฺชติ เหตุปจฺจยา.\csromanspacer{}\csromanspacer{}เหตุยา ตีณิ.}

\proseitem{130}{น-อุเปกฺขาสหคตํ นสนิทสฺสนสปฺปฏิฆํ ธมฺมํ ปฏิจฺจ น-อุเปกฺขาสหคโต สนิทสฺสนสปฺปฏิโฆ ธมฺโม อุปฺปชฺชติ เหตุปจฺจยา.\csromanspacer{}\csromanspacer{}เหตุยา ตีณิ.}

\proseitem{131}{นกามาวจรํ นสนิทสฺสนสปฺปฏิฆํ ธมฺมํ ปฏิจฺจ กามาวจโร สนิทสฺสนสปฺปฏิโฆ ธมฺโม อุปฺปชฺชติ เหตุปจฺจยา.\csromanspacer{}\csromanspacer{}เหตุยา ตีณิ.}

\proseitem{132}{นรูปาวจรํ นสนิทสฺสนสปฺปฏิฆํ ธมฺมํ ปฏิจฺจ นรูปาวจโร สนิทสฺสนสปฺปฏิโฆ ธมฺโม อุปฺปชฺชติ เหตุปจฺจยา.\csromanspacer{}\csromanspacer{}เหตุยา ตีณิ.}

\proseitem{133}{น-อรูปาวจรํ นสนิทสฺสนสปฺปฏิฆํ ธมฺมํ ปฏิจฺจ น-อรูปาวจโร สนิทสฺสนสปฺปฏิโฆ ธมฺโม อุปฺปชฺชติ เหตุปจฺจยา.\csromanspacer{}\csromanspacer{}เหตุยา ตีณิ.}

\proseitem{134}{นปริยาปนฺนํ นสนิทสฺสนสปฺปฏิฆํ ธมฺมํ ปฏิจฺจ ปริยาปนฺโน สนิทสฺสนสปฺปฏิโฆ ธมฺโม อุปฺปชฺชติ เหตุปจฺจยา.\csromanspacer{}\csromanspacer{}เหตุยา ตีณิ.}

\proseitem{135}{นนิยฺยานิกํ นสนิทสฺสนสปฺปฏิฆํ ธมฺมํ ปฏิจฺจ อนิยฺยานิโก สนิทสฺสนสปฺปฏิโฆ ธมฺโม อุปฺปชฺชติ เหตุปจฺจยา.\csromanspacer{}\csromanspacer{}เหตุยา ตีณิ.}

\proseitem{136}{นนิยตํ นสนิทสฺสนสปฺปฏิฆํ ธมฺมํ ปฏิจฺจ อนิยโต สนิทสฺสนสปฺปฏิโฆ ธมฺโม อุปฺปชฺชติ เหตุปจฺจยา.\csromanspacer{}\csromanspacer{}เหตุยา ตีณิ.}

\proseitem{137}{นส-อุตฺตรํ นสนิทสฺสนสปฺปฏิฆํ ธมฺมํ ปฏิจฺจ ส-อุตฺตโร สนิทสฺสนสปฺปฏิโฆ ธมฺโม อุปฺปชฺชติ เหตุปจฺจยา.\csromanspacer{} เหตุยา ตีณิ.}

\proseitem{138}{นสรณํ นสนิทสฺสนสปฺปฏิฆํ ธมฺมํ ปฏิจฺจ อรโณ สนิทสฺสนสปฺปฏิโฆ ธมฺโม อุปฺปชฺชติ เหตุปจฺจยา.\csromanspacer{} เอกํ.}

\csromanheader{2}{83-100. ปิฏฺฐิทุก 22. สนิทสฺสนตฺติก}
\prose{\csromanfolio{375}น-อรณํ นสนิทสฺสนสปฺปฏิฆํ ธมฺมํ ปฏิจฺจ อรโณ สนิทสฺสนสปฺปฏิโฆ ธมฺโม อุปฺปชฺชติ เหตุปจฺจยา.\csromanspacer{} เอกํ.}

\prose{นสรณํ นสนิทสฺสนสปฺปฏิฆญฺจ น-อรณํ นสนิทสฺสนสปฺปฏิฆญฺจ ธมฺมํ ปฏิจฺจ อรโณ สนิทสฺสนสปฺปฏิโฆ ธมฺโม อุปฺปชฺชติ เหตุปจฺจยา.\csromanspacer{}\csromanspacer{}เหตุยา ตีณิ.\csromanspacer{}\csromanspacer{}(อนิทสฺสนสปฺปฏิเฆ ตีณิเยว.)}

\prose{นสรณํ น-อนิทสฺสน-อปฺปฏิฆํ ธมฺมํ ปฏิจฺจ อรโณ อนิทสฺสน-อปฺปฏิโฆ ธมฺโม อุปฺปชฺชติ เหตุปจฺจยา.\csromanspacer{}\csromanspacer{}เหตุยา เอกํ.}

\csromanheader{2}{\textbf{ปจฺจนีย-นเหตุ น-อารมฺมณ}}
\proseitem{139}{นสรณํ นสนิทสฺสนสปฺปฏิฆํ ธมฺมํ ปฏิจฺจ อรโณ สนิทสฺสนสปฺปฏิโฆ ธมฺโม อุปฺปชฺชติ นเหตุปจฺจยา.\csromanspacer{} เอกํ.}

\prose{นสรณํ นสนิทสฺสนสปฺปฏิฆํ ธมฺมํ ปฏิจฺจ อรโณ สนิทสฺสนสปฺปฏิโฆ ธมฺโม อุปฺปชฺชติ น-อารมฺมณปจฺจยา.\csromanspacer{} ตีณิ.\csromanspacer{}\csromanspacer{}นเหตุยา เอกํ, น-อารมฺมเณ ตีณิ, น-อธิปติยา ตีณิ ฯเปฯ โนวิคเต ตีณิ.}

\prose{เหตุปจฺจยา น-อารมฺมเณ ตีณิ.\csromanspacer{}\csromanspacer{}น-อารมฺมณปจฺจยา เหตุยา ตีณิ.}

\prose{(สหชาตวารมฺปิ ปจฺจยวารมฺปิ นิสฺสยวารมฺปิ สํสฏฺฐวารมฺปิ สมฺปยุตฺตวารมฺปิ วิตฺถาเรตพฺพํ.)}

\csromanheader{2}{\textbf{ปญฺหาวาร-เหตุ}}
\proseitem{140}{นสรโณ นสนิทสฺสนสปฺปฏิโฆ ธมฺโม อรณสฺส สนิทสฺสนสปฺปฏิฆสฺส ธมฺมสฺส เหตุปจฺจเยน ปจฺจโย.\csromanspacer{}\csromanspacer{}น-อรโณ นสนิทสฺสนสปฺปฏิโฆ ธมฺโม อรณสฺส สนิทสฺสนสปฺปฏิฆสฺส ธมฺมสฺส เหตุปจฺจเยน ปจฺจโย.}

\prose{เหตุยา เทฺว, อธิปติยา เทฺว ฯเปฯ อวิคเต ตีณิ.}

\prose{(ยถา กุสลตฺติเก ปญฺหาวารสฺส อนุโลมมฺปิ ปจฺจนียมฺปิ อนุโลมปจฺจนียมฺปิ ปจฺจนียานุโลมมฺปิ คณิตํ, เอวํ คเณตพฺพํ.)}

\nitthitam{\textbf{ธมฺมปจฺจนียานุโลเม ทุกติกปฏฺฐานํ นิฏฺฐิตํ.}}
\pitaka{\textbf{อภิธมฺมปิฏก}}
\csromanheader{2}{\textbf{ธมฺมปจฺจนียานุโลม}}
\gambhira{\textbf{ติกทุกปฏฺฐานปาฬิ}}
\namakkaram{นโม ตสฺส ภควโต อรหโต สมฺมาสมฺพุทฺธสฺส.}
\csromanheader{2}{1. กุสลตฺติก 1. เหตุทุก}
\proseitem{1}{\csromanfolio{377}นกุสลํ นเหตุํ ธมฺมํ ปฏิจฺจ อกุสโล เหตุ ธมฺโม อุปฺปชฺชติ เหตุปจฺจยา.\csromanspacer{}\csromanspacer{}นกุสลํ นเหตุํ ธมฺมํ ปฏิจฺจ อพฺยากโต เหตุ ธมฺโม อุปฺปชฺชติ เหตุปจฺจยา. (2)}

\prose{น-อกุสลํ นเหตุํ ธมฺมํ ปฏิจฺจ กุสโล เหตุ ธมฺโม อุปฺปชฺชติ เหตุปจฺจยา.\csromanspacer{}\csromanspacer{}น-อกุสลํ นเหตุํ ธมฺมํ ปฏิจฺจ อพฺยากโต เหตุ ธมฺโม อุปฺปชฺชติ เหตุปจฺจยา. (2)}

\prose{น-อพฺยากตํ นเหตุํ ธมฺมํ ปฏิจฺจ กุสโล เหตุ ธมฺโม อุปฺปชฺชติ เหตุปจฺจยา.\csromanspacer{}\csromanspacer{}น-อพฺยากตํ นเหตุํ ธมฺมํ ปฏิจฺจ อกุสโล เหตุ ธมฺโม อุปฺปชฺชติ เหตุปจฺจยา. (2)}

\prose{นกุสลํ นเหตุญฺจ น-อพฺยากตํ นเหตุญฺจ ธมฺมํ ปฏิจฺจ อกุสโล เหตุ ธมฺโม อุปฺปชฺชติ เหตุปจฺจยา.\csromanspacer{}\csromanspacer{}น-อกุสลํ นเหตุญฺจ น-อพฺยากตํ นเหตุญฺจ ธมฺมํ ปฏิจฺจ กุสโล เหตุ ธมฺโม อุปฺปชฺชติ เหตุปจฺจยา.\csromanspacer{}\csromanspacer{}นกุสลํ นเหตุญฺจ น-อกุสลํ นเหตุญฺจ ธมฺมํ ปฏิจฺจ อพฺยากโต เหตุ ธมฺโม อุปฺปชฺชติ เหตุปจฺจยา.\csromanspacer{} ตีณิ. (สํขิตฺตํ.) .\csromanspacer{} เหตุยา นว.}

\proseitem{2}{นกุสลํ นนเหตุํ ธมฺมํ ปฏิจฺจ อกุสโล นเหตุ ธมฺโม อุปฺปชฺชติ เหตุปจฺจยา.\csromanspacer{}\csromanspacer{}นกุสลํ นนเหตุํ ธมฺมํ ปฏิจฺจ อพฺยากโต นเหตุ}

\vspace{\tipitakaparskip}
\csromancenter{ธมฺมปจฺจนียานุโลม ติกทุกปฏฺฐานปาฬิ ธมฺโม อุปฺปชฺชติ เหตุปจฺจยา.\csromanspacer{}\csromanspacer{}นกุสลํ นนเหตุํ ธมฺมํ ปฏิจฺจ อกุสโล นเหตุ จ อพฺยากโต นเหตุ จ ธมฺมา อุปฺปชฺชนฺติ เหตุปจฺจยา.\csromanspacer{} ตีณิ.}
\prose{\csromanfolio{378}น-อกุสลํ นนเหตุํ ธมฺมํ ปฏิจฺจ กุสโล นเหตุ ธมฺโม อุปฺปชฺชติ เหตุปจฺจยา.\csromanspacer{}\csromanspacer{}น-อกุสลํ นนเหตุํ ธมฺมํ ปฏิจฺจ อพฺยากโต นเหตุ ธมฺโม อุปฺปชฺชติ เหตุปจฺจยา.\csromanspacer{}\csromanspacer{}น-อกุสลํ นนเหตุํ ธมฺมํ ปฏิจฺจ กุสโล นเหตุ จ อพฺยากโต นเหตุ จ ธมฺมา อุปฺปชฺชนฺติ เหตุปจฺจยา.\csromanspacer{} ตีณิ.}

\prose{น-อพฺยากตํ นนเหตุํ ธมฺมํ ปฏิจฺจ อพฺยากโต นเหตุ ธมฺโม อุปฺปชฺชติ เหตุปจฺจยา.\csromanspacer{} ปญฺจ.}

\prose{นกุสลํ นนเหตุญฺจ น-อพฺยากตํ นนเหตุญฺจ ธมฺมํ ปฏิจฺจ อกุสโล นเหตุ ธมฺโม อุปฺปชฺชติ เหตุปจฺจยา.\csromanspacer{} ตีณิ.}

\prose{น-อกุสลํ นนเหตุญฺจ น-อพฺยากตํ นนเหตุญฺจ ธมฺมํ ปฏิจฺจ กุสโล นเหตุ ธมฺโม อุปฺปชฺชติ เหตุปจฺจยา.\csromanspacer{} ตีณิ.}

\prose{นกุสลํ นนเหตุญฺจ น-อกุสลํ นนเหตุญฺจ ธมฺมํ ปฏิจฺจ อพฺยากโต นเหตุ ธมฺโม อุปฺปชฺชติ เหตุปจฺจยา.\csromanspacer{} เอกํ.\csromanspacer{}\csromanspacer{}เหตุยา อฏฺฐารส.}

\csromanheader{2}{1. กุสลตฺติก 2. สเหตุกทุก}
\proseitem{3}{นกุสลํ นสเหตุกํ ธมฺมํ ปฏิจฺจ อกุสโล สเหตุโก ธมฺโม อุปฺปชฺชติ เหตุปจฺจยา.\csromanspacer{}\csromanspacer{}นกุสลํ นสเหตุกํ ธมฺมํ ปฏิจฺจ อพฺยากโต สเหตุโก ธมฺโม อุปฺปชฺชติ เหตุปจฺจยา.\csromanspacer{}\csromanspacer{}น-อกุสลํ นสเหตุกํ ธมฺมํ ปฏิจฺจ อพฺยากโต สเหตุโก ธมฺโม อุปฺปชฺชติ เหตุปจฺจยา.\csromanspacer{}\csromanspacer{}น-อพฺยากตํ นสเหตุกํ ธมฺมํ ปฏิจฺจ อกุสโล สเหตุโก ธมฺโม อุปฺปชฺชติ เหตุปจฺจยา.\csromanspacer{}\csromanspacer{}นกุสลํ นสเหตุกญฺจ น-อพฺยากตํ นสเหตุกญฺจ ธมฺมํ ปฏิจฺจ อกุสโล สเหตุโก ธมฺโม อุปฺปชฺชติ เหตุปจฺจยา.\csromanspacer{}\csromanspacer{}นกุสลํ นสเหตุกญฺจ น-อกุสลํ นสเหตุกญฺจ ธมฺมํ ปฏิจฺจ อพฺยากโต สเหตุโก ธมฺโม อุปฺปชฺชติ เหตุปจฺจยา.\csromanspacer{}\csromanspacer{}เหตุยา ฉ.}

\proseitem{4}{นกุสลํ น-อเหตุกํ ธมฺมํ ปฏิจฺจ อพฺยากโต อเหตุโก ธมฺโม อุปฺปชฺชติ เหตุปจฺจยา.\csromanspacer{}\csromanspacer{}น-อกุสลํ น-อเหตุกํ ธมฺมํ ปฏิจฺจ อพฺยากโต อเหตุโก ธมฺโม อุปฺปชฺชติ เหตุปจฺจยา.}

\vspace{\tipitakaparskip}
\csromancenter{กุสลตฺติก 4-6.\csromanspacer{} เหตุสเหตุกทุกาทิ น-อพฺยากตํ น-อเหตุกํ ธมฺมํ ปฏิจฺจ อพฺยากโต อเหตุโก ธมฺโม อุปฺปชฺชติ เหตุปจฺจยา.\csromanspacer{}\csromanspacer{}นกุสลํ น-อเหตุกญฺจ น-อพฺยากตํ น-อเหตุกญฺจ ธมฺมํ ปฏิจฺจ อพฺยากโต อเหตุโก ธมฺโม อุปฺปชฺชติ เหตุปจฺจยา.\csromanspacer{}\csromanspacer{}น-อกุสลํ น-อเหตุกญฺจ น-อพฺยากตํ น-อเหตุกญฺจ ธมฺมํ ปฏิจฺจ อพฺยากโต อเหตุโก ธมฺโม อุปฺปชฺชติ เหตุปจฺจยา.\csromanspacer{}\csromanspacer{}นกุสลํ น-อเหตุกญฺจ น-อกุสลํ น-อเหตุกญฺจ ธมฺมํ ปฏิจฺจ อพฺยากโต อเหตุโก ธมฺโม อุปฺปชฺชติ เหตุปจฺจยา.\csromanspacer{}\csromanspacer{}เหตุยา ฉ.}
\csromanheader{2}{1. กุสลตฺติก 3. เหตุสมฺปยุตฺตทุก}
\proseitem{5}{\csromanfolio{379}นกุสลํ นเหตุสมฺปยุตฺตํ ธมฺมํ ปฏิจฺจ อกุสโล เหตุสมฺปยุตฺโต ธมฺโม อุปฺปชฺชติ เหตุปจฺจยา.\csromanspacer{}\csromanspacer{}นกุสลํ นเหตุสมฺปยุตฺตํ ธมฺมํ ปฏิจฺจ อพฺยากโต เหตุสมฺปยุตฺโต ธมฺโม อุปฺปชฺชติ เหตุปจฺจยา.\csromanspacer{}\csromanspacer{}น-อกุสลํ นเหตุสมฺปยุตฺตํ ธมฺมํ ปฏิจฺจ อพฺยากโต เหตุสมฺปยุตฺโต ธมฺโม อุปฺปชฺชติ เหตุปจฺจยา.\csromanspacer{} ฉ. (สเหตุกสทิสํ.)}

\prose{นกุสลํ นเหตุวิปฺปยุตฺตํ ธมฺมํ ปฏิจฺจ อพฺยากโต เหตุวิปฺปยุตฺโต ธมฺโม อุปฺปชฺชติ เหตุปจฺจยา.\csromanspacer{} ฉ. (อเหตุกสทิสํ.\csromanspacer{} สํขิตฺตํ.)}

\csromanheader{2}{1. กุสลตฺติก 4-6. เหตุสเหตุกทุกาทิ}
\proseitem{6}{นกุสลํ นเหตุญฺเจว น-อเหตุกญฺจ ธมฺมํ ปฏิจฺจ อกุสโล เหตุ เจว สเหตุโก จ ธมฺโม อุปฺปชฺชติ เหตุปจฺจยา.\csromanspacer{}\csromanspacer{}นกุสลํ นเหตุญฺเจว น-อเหตุกญฺจ ธมฺมํ ปฏิจฺจ อพฺยากโต เหตุ เจว สเหตุโก จ ธมฺโม อุปฺปชฺชติ เหตุปจฺจยา.\csromanspacer{} เทฺว.}

\prose{(น-อกุสเล เทฺว.\csromanspacer{} น-อพฺยากเต เทฺว.\csromanspacer{} ปฐมํ คณิตเกน เอกํ, ทุติยํ คณิตเกน เอกํ, ตติยํ คณิตเกน เอกํ, สพฺพตฺถ นว ปญฺหา.)}

\prose{นกุสลํ น-อเหตุกญฺเจว น จ เหตุํ ธมฺมํ ปฏิจฺจ อกุสโล สเหตุโก เจว น จ เหตุ ธมฺโม อุปฺปชฺชติ เหตุปจฺจยา.\csromanspacer{}\csromanspacer{}เหตุยา นว.}

\proseitem{7}{นกุสลํ นเหตุญฺเจว นเหตุวิปฺปยุตฺตญฺจ ธมฺมํ ปฏิจฺจ อกุสโล เหตุ เจว เหตุสมฺปยุตฺโต จ ธมฺโม อุปฺปชฺชติ เหตุปจฺจยา.\csromanspacer{}\csromanspacer{}เหตุยา นว.}

\csromanheader{2}{ธมฺมปจฺจนียานุโลม ติกทุกปฏฺฐานปาฬิ}
\prose{\csromanfolio{380}นกุสลํ นเหตุวิปฺปยุตฺตญฺเจว นนเหตุญฺจ ธมฺมํ ปฏิจฺจ อกุสโล เหตุสมฺปยุตฺโต เจว น จ เหตุ ธมฺโม อุปฺปชฺชติ เหตุปจฺจยา.\csromanspacer{}\csromanspacer{}เหตุยา นว.}

\proseitem{8}{นกุสลํ นเหตุํ นสเหตุกํ ธมฺมํ ปฏิจฺจ อพฺยากโต นเหตุ สเหตุโก ธมฺโม อุปฺปชฺชติ เหตุปจฺจยา.\csromanspacer{}\csromanspacer{}เหตุยา ตีณิ.}

\proseitem{9}{นกุสลํ นเหตุํ น-อเหตุกํ ธมฺมํ ปฏิจฺจ อกุสโล นเหตุ อเหตุโก ธมฺโม อุปฺปชฺชติ เหตุปจฺจยา. (1)}

\prose{น-อกุสลํ นเหตุํ น-อเหตุกํ ธมฺมํ ปฏิจฺจ อพฺยากโต นเหตุ อเหตุโก ธมฺโม อุปฺปชฺชติ เหตุปจฺจยา. (1)}

\prose{น-อพฺยากตํ นเหตุํ น-อเหตุกํ ธมฺมํ ปฏิจฺจ อพฺยากโต นเหตุ อเหตุโก ธมฺโม อุปฺปชฺชติ เหตุปจฺจยา. (1) (คณิตเก ตีณิ ปญฺหา.) .\csromanspacer{} เหตุยา ฉ.}

\csromanheader{2}{1. กุสลตฺติก 7-13. จูฬนฺตรทุก}
\proseitem{10}{นกุสโล นสปฺปจฺจโย ธมฺโม กุสลสฺส สปฺปจฺจยสฺส ธมฺมสฺส อารมฺมณปจฺจเยน ปจฺจโย.\csromanspacer{}\csromanspacer{}อารมฺมเณ ฉ. (สงฺขตํ สปฺปจฺจยสทิสํ.)}

\proseitem{11}{นกุสลํ นสนิทสฺสนํ ธมฺมํ ปฏิจฺจ อพฺยากโต สนิทสฺสโน ธมฺโม อุปฺปชฺชติ เหตุปจฺจยา.\csromanspacer{}\csromanspacer{}เหตุยา ฉ.}

\prose{นกุสโล น-อนิทสฺสโน ธมฺโม กุสลสฺส อนิทสฺสนสฺส ธมฺมสฺส อารมฺมณปจฺจเยน ปจฺจโย. (นว ปญฺหา กาตพฺพา.)}

\proseitem{12}{นกุสลํ นสปฺปฏิฆํ ธมฺมํ ปฏิจฺจ อพฺยากโต สปฺปฏิโฆ ธมฺโม อุปฺปชฺชติ เหตุปจฺจยา.\csromanspacer{}\csromanspacer{}เหตุยา ฉ.}

\prose{นกุสลํ น-อปฺปฏิฆํ ธมฺมํ ปฏิจฺจ อพฺยากโต อปฺปฏิโฆ ธมฺโม อุปฺปชฺชติ เหตุปจฺจยา.\csromanspacer{}\csromanspacer{}เหตุยา ตีณิ.}

\proseitem{13}{นกุสลํ นรูปิํ ธมฺมํ ปฏิจฺจ อพฺยากโต รูปี ธมฺโม อุปฺปชฺชติ เหตุปจฺจยา.\csromanspacer{}\csromanspacer{}เหตุยา ฉ.}

\csromanheader{2}{1. กุสลตฺติก 14. อาสวโคจฺฉก}
\prose{\csromanfolio{381}นกุสลํ น-อรูปิํ ธมฺมํ ปฏิจฺจ อพฺยากโต อรูปี ธมฺโม อุปฺปชฺชติ เหตุปจฺจยา.\csromanspacer{}\csromanspacer{}เหตุยา ตีณิ.}

\proseitem{14}{นกุสลํ นโลกิยํ ธมฺมํ ปฏิจฺจ อพฺยากโต โลกิโย ธมฺโม อุปฺปชฺชติ เหตุปจฺจยา.\csromanspacer{}\csromanspacer{}เหตุยา ปญฺจ.}

\prose{นกุสลํ นโลกุตฺตรํ ธมฺมํ ปจฺจยา กุสโล โลกุตฺตโร ธมฺโม อุปฺปชฺชติ เหตุปจฺจยา.\csromanspacer{}\csromanspacer{}นกุสลํ นโลกุตฺตรํ ธมฺมํ ปจฺจยา อพฺยากโต โลกุตฺตโร ธมฺโม อุปฺปชฺชติ เหตุปจฺจยา. (ฉ ปญฺหา กาตพฺพา.)}

\proseitem{15}{นกุสลํ นเกนจิ วิญฺเญยฺยํ ธมฺมํ ปฏิจฺจ อกุสโล เกนจิ วิญฺเญยฺโย ธมฺโม อุปฺปชฺชติ เหตุปจฺจยา.\csromanspacer{}\csromanspacer{}เหตุยา อฏฺฐารส.}

\prose{นกุสลํ นเกนจิ นวิญฺเญยฺยํ ธมฺมํ ปฏิจฺจ อกุสโล เกนจิ วิญฺเญยฺโย ธมฺโม อุปฺปชฺชติ เหตุปจฺจยา.\csromanspacer{}\csromanspacer{}เหตุยา อฏฺฐารส.}

\csromanheader{2}{1. กุสลตฺติก 14. อาสวโคจฺฉก}
\proseitem{16}{นกุสลํ โน-อาสวํ ธมฺมํ ปฏิจฺจ อกุสโล อาสโว ธมฺโม อุปฺปชฺชติ เหตุปจฺจยา.\csromanspacer{}\csromanspacer{}เหตุยา ตีณิ.}

\prose{นกุสลํ นโน-อาสวํ ธมฺมํ ปฏิจฺจ อกุสโล โน-อาสโว ธมฺโม อุปฺปชฺชติ เหตุปจฺจยา.\csromanspacer{} เหตุยา นว.}

\proseitem{17}{นกุสลํ นสาสวํ ธมฺมํ ปฏิจฺจ อพฺยากโต สาสโว ธมฺโม อุปฺปชฺชติ เหตุปจฺจยา.\csromanspacer{}\csromanspacer{}เหตุยา ปญฺจ.}

\prose{นกุสลํ น-อนาสวํ ธมฺมํ ปจฺจยา กุสโล อนาสโว ธมฺโม อุปฺปชฺชติ เหตุปจฺจยา. (สํขิตฺตํ.) .\csromanspacer{} เหตุยา ฉ ฯเปฯ วิปาเก ตีณิ ฯเปฯ อวิคเต ฉ.}

\proseitem{18}{นกุสลํ น-อาสวสมฺปยุตฺตํ ธมฺมํ ปฏิจฺจ อกุสโล อาสวสมฺปยุตฺโต ธมฺโม อุปฺปชฺชติ เหตุปจฺจยา.\csromanspacer{}\csromanspacer{}เหตุยา ตีณิ.}

\prose{นกุสลํ น-อาสววิปฺปยุตฺตํ ธมฺมํ ปฏิจฺจ อกุสโล อาสววิปฺปยุตฺโต ธมฺโม อุปฺปชฺชติ เหตุปจฺจยา.\csromanspacer{}\csromanspacer{}เหตุยา นว.}

\csromanheader{2}{ธมฺมปจฺจนียานุโลม ติกทุกปฏฺฐานปาฬิ}
\proseitem{19}{\csromanfolio{382}นกุสลํ น-อาสวญฺเจว น-อนาสวญฺจ ธมฺมํ ปฏิจฺจ อกุสโล อาสโว เจว สาสโว จ ธมฺโม อุปฺปชฺชติ เหตุปจฺจยา.\csromanspacer{}\csromanspacer{}เหตุยา ติณิ.}

\prose{นกุสลํ น-อนาสวญฺเจว นโน จ อาสวํ ธมฺมํ ปฏิจฺจ อกุสโล สาสโว เจว โน จ อาสโว ธมฺโม อุปฺปชฺชติ เหตุปจฺจยา.\csromanspacer{}\csromanspacer{}เหตุยา นว.}

\proseitem{20}{นกุสลํ น-อาสวญฺเจว น-อาสววิปฺปยุตฺตญฺจ ธมฺมํ ปฏิจฺจ อกุสโล อาสโว เจว อาสวสมฺปยุตฺโต จ ธมฺโม อุปฺปชฺชติ เหตุปจฺจยา.\csromanspacer{}\csromanspacer{}เหตุยา ตีณิ.}

\prose{นกุสลํ น-อาสววิปฺปยุตฺตญฺเจว นโน จ อาสวํ ธมฺมํ ปฏิจฺจ อกุสโล อาสวสมฺปยุตฺโต เจว โน อาสโว ธมฺโม อุปฺปชฺชติ เหตุปจฺจยา.\csromanspacer{}\csromanspacer{}เหตุยา ตีณิ.}

\proseitem{21}{นกุสลํ อาสววิปฺปยุตฺตํ นสาสวํ ธมฺมํ ปฏิจฺจ อพฺยากโต อาสววิปฺปยุตฺโต สาสโว ธมฺโม อุปฺปชฺชติ เหตุปจฺจยา.\csromanspacer{}\csromanspacer{}เหตุยา ปญฺจ.}

\csromanheader{2}{1. กุสลตฺติก 20. สญฺโญชนาทิโคจฺฉก}
\proseitem{22}{นกุสลํ โนสญฺโญชนํ ธมฺมํ ปฏิจฺจ อกุสโล สญฺโญชโน ธมฺโม อุปฺปชฺชติ เหตุปจฺจยา.}

\proseitem{23}{นกุสลํ โนคนฺถํ ธมฺมํ ปฏิจฺจ อกุสโล คนฺโถ ธมฺโม อุปฺปชฺชติ เหตุปจฺจยา.}

\proseitem{24}{นกุสลํ โน-โอฆํ ธมฺมํ ปฏิจฺจ อกุสโล โอโฆ ธมฺโม อุปฺปชฺชติ เหตุปจฺจยา.}

\proseitem{25}{นกุสลํ โนโยคํ ธมฺมํ ปฏิจฺจ อกุสโล โยโค ธมฺโม อุปฺปชฺชติ เหตุปจฺจยา.}

\proseitem{26}{นกุสลํ โนนีวรณํ ธมฺมํ ปฏิจฺจ อกุสโล นีวรโณ ธมฺโม อุปฺปชฺชติ เหตุปจฺจยา.}

\csromanheader{2}{1. กุสลตฺติก 96. ปริยาปนฺนทุก}
\proseitem{27}{\csromanfolio{383}นกุสลํ โนปรามาสํ ธมฺมํ ปฏิจฺจ อกุสโล ปรามาโส ธมฺโม อุปฺปชฺชติ เหตุปจฺจยา.}

\csromanheader{2}{1. กุสลตฺติก 55. มหนฺตรทุก}
\proseitem{28}{นกุสลํ นสารมฺมณํ ธมฺมํ ปฏิจฺจ อพฺยากโต สารมฺมโณ ธมฺโม อุปฺปชฺชติ เหตุปจฺจยา ฯเปฯ.}

\csromanheader{2}{1. กุสลตฺติก 93. กามาวจรทุก}
\proseitem{29}{นกุสลํ นกามาวจรํ ธมฺมํ ปฏิจฺจ อพฺยากโต กามาวจโร ธมฺโม อุปฺปชฺชติ เหตุปจฺจยา.\csromanspacer{}\csromanspacer{}เหตุยา ปญฺจ.}

\prose{นกุสลํ นนกามาวจรํ ธมฺมํ ปฏิจฺจ อพฺยากโต นกามาวจโร ธมฺโม อุปฺปชฺชติ เหตุปจฺจยา.\csromanspacer{}\csromanspacer{}เหตุยา ตีณิ.}

\csromanheader{2}{1. กุสลตฺติก 94. รูปาวจรทุก}
\proseitem{30}{นกุสลํ นรูปาวจรํ ธมฺมํ ปฏิจฺจ อพฺยากโต รูปาวจโร ธมฺโม อุปฺปชฺชติ เหตุปจฺจยา.\csromanspacer{}\csromanspacer{}เหตุยา ตีณิ.}

\prose{นกุสลํ นนรูปาวจรํ ธมฺมํ ปฏิจฺจ อพฺยากโต นรูปาวจโร ธมฺโม อุปฺปชฺชติ เหตุปจฺจยา.\csromanspacer{}\csromanspacer{}เหตุยา ปญฺจ.}

\csromanheader{2}{1. กุสลตฺติก 95. อรูปาวจรทุก}
\proseitem{31}{นกุสลํ น-อรูปาวจรํ ธมฺมํ ปจฺจยา กุสโล อรูปาวจโร ธมฺโม อุปฺปชฺชติ เหตุปจฺจยา.\csromanspacer{}\csromanspacer{}เหตุยา ฉ.}

\prose{นกุสลํ นน-อรูปาวจรํ ธมฺมํ ปฏิจฺจ อพฺยากโต น อรูปาวจโร ธมฺโม อุปฺปชฺชติ เหตุปจฺจยา.\csromanspacer{}\csromanspacer{}เหตุยา ปญฺจ.}

\csromanheader{2}{1. กุสลตฺติก 96. ปริยาปนฺนทุก}
\proseitem{32}{นกุสลํ นปริยาปนฺนํ ธมฺมํ ปฏิจฺจ อพฺยากโต ปริยาปนฺโน ธมฺโม อุปฺปชฺชติ เหตุปจฺจยา.\csromanspacer{}\csromanspacer{}น-อกุสลํ นปริยาปนฺนํ ธมฺมํ ปฏิจฺจ อพฺยากโต ปริยาปนฺโน ธมฺโม อุปฺปชฺชติ เหตุปจฺจยา.\csromanspacer{}\csromanspacer{}น-อพฺยากตํ นปริยาปนฺนํ ธมฺมํ ปฏิจฺจ อพฺยากโต ปริยาปนฺโน ธมฺโม อุปฺปชฺชติ เหตุปจฺจยา.\csromanspacer{} เหตุยา ปญฺจ.}

\csromanheader{2}{ธมฺมปจฺจนียานุโลม ติกทุกปฏฺฐานปาฬิ}
\prose{\csromanfolio{384}นกุสลํ น-อปริยาปนฺนํ ธมฺมํ ปจฺจยา กุสโล อปริยาปนฺโน ธมฺโม อุปฺปชฺชติ เหตุปจฺจยา.\csromanspacer{}\csromanspacer{}นกุสลํ น-อปริยาปนฺนํ ธมฺมํ ปจฺจยา อพฺยากโต อปริยาปนฺโน ธมฺโม อุปฺปชฺชติ เหตุปจฺจยา.\csromanspacer{} เทฺว. (อกุสลมูเล เทฺว.\csromanspacer{} ทุกมูเล เทฺว.) .\csromanspacer{} เหตุยา ฉ.}

\csromanheader{2}{1. กุสลตฺติก 97. นิยฺยานิกทุก}
\proseitem{33}{นกุสลํ นนิยฺยานิกํ ธมฺมํ ปจฺจยา กุสโล นิยฺยานิโก ธมฺโม อุปฺปชฺชติ เหตุปจฺจยา.\csromanspacer{}\csromanspacer{}เหตุยา ตีณิ.}

\prose{น-อกุสลํ น-อนิยฺยานิกํ ธมฺมํ ปฏิจฺจ อพฺยากโต อนิยฺยานิโก ธมฺโม อุปฺปชฺชติ เหตุปจฺจยา.\csromanspacer{}\csromanspacer{}น-อพฺยากตํ น-อนิยฺยานิกํ ธมฺมํ ปฏิจฺจ อพฺยากโต อนิยฺยานิโก ธมฺโม อุปฺปชฺชติ เหตุปจฺจยา.\csromanspacer{}\csromanspacer{}น-อกุสลํ น-อนิยฺยานิกญฺจ น-อพฺยากตํ น-อนิยฺยานิกญฺจ ธมฺมํ ปฏิจฺจ อพฺยากโต อนิยฺยานิโก ธมฺโม อุปฺปชฺชติ เหตุปจฺจยา.\csromanspacer{}\csromanspacer{}เหตุยา ตีณิ.}

\csromanheader{2}{1. กุสลตฺติก 98. นิยตทุก}
\proseitem{34}{นกุสลํ นนิยตํ ธมฺมํ ปจฺจยา กุสโล นิยโต ธมฺโม อุปฺปชฺชติ เหตุปจฺจยา. (อปริยาปนฺนสทิสํ.\csromanspacer{} ฉ ปญฺหา.)}

\prose{นกุสลํ น-อนิยตํ ธมฺมํ ปฏิจฺจ อพฺยากโต อนิยโต ธมฺโม อุปฺปชฺชติ เหตุปจฺจยา.\csromanspacer{}\csromanspacer{}น-อกุสลํ น-อนิยตํ ธมฺมํ ปฏิจฺจ อพฺยากโต อนิยโต ธมฺโม อุปฺปชฺชติ เหตุปจฺจยา.\csromanspacer{}\csromanspacer{}น-อพฺยากตํ น-อนิยตํ ธมฺมํ ปฏิจฺจ อพฺยากโต อนิยโต ธมฺโม อุปฺปชฺชติ เหตุปจฺจยา.\csromanspacer{}\csromanspacer{}นกุสลํ น-อนิยตญฺจ น-อพฺยากตํ น-อนิยตญฺจ ธมฺมํ ปฏิจฺจ อพฺยากโต อนิยโต ธมฺโม อุปฺปชฺชติ เหตุปจฺจยา.\csromanspacer{}\csromanspacer{}นกุสลํ น-อนิยตญฺจ น-อพฺยากตํ น-อนิยตญฺจ ธมฺมํ ปฏิจฺจ อพฺยากโต อนิยโต ธมฺโม อุปฺปชฺชติ เหตุปจฺจยา.\csromanspacer{}\csromanspacer{}เหตุยา ปญฺจ.}

\csromanheader{2}{1. กุสลตฺติก 99. ส-อุตฺตรทุก}
\proseitem{35}{นกุสลํ นส-อุตฺตรํ ธมฺมํ ปฏิจฺจ อพฺยากโต ส-อุตฺตโร ธมฺโม อุปฺปชฺชติ เหตุปจฺจยา.\csromanspacer{}\csromanspacer{}น-อกุสลํ นส-อุตฺตรํ ธมฺมํ ปฏิจฺจ อพฺยากโต ส-อุตฺตโร ธมฺโม อุปฺปชฺชติ เหตุปจฺจยา.\csromanspacer{}\csromanspacer{}น-อพฺยากตํ นส-อุตฺตรํ ธมฺมํ}

\vspace{\tipitakaparskip}
\csromancenter{เวทนาตฺติก 1.\csromanspacer{} เหตุทุก ปฏิจฺจ อพฺยากโต ส-อุตฺตโร ธมฺโม อุปฺปชฺชติ เหตุปจฺจยา.\csromanspacer{}\csromanspacer{}น-อกุสลํ นส-อุตฺตรญฺจ น-อพฺยากตํ นส-อุตฺตรญฺจ ธมฺมํ ปฏิจฺจ อพฺยากโต ส-อุตฺตโร ธมฺโม อุปฺปชฺชติ เหตุปจฺจยา.\csromanspacer{}\csromanspacer{}นกุสลํ นส-อุตฺตรญฺจ น-อกุสลํ นส-อุตฺตรญฺจ ธมฺมํ ปฏิจฺจ อพฺยากโต ส-อุตฺตโร ธมฺโม อุปฺปชฺชติ เหตุปจฺจยา.\csromanspacer{}\csromanspacer{}เหตุยา ปญฺจ.}
\prose{\csromanfolio{385}นกุสลํ น-อนุตฺตรํ ธมฺมํ ปจฺจยา กุสโล อนุตฺตโร ธมฺโม อุปฺปชฺชติ เหตุปจฺจยา. (สํขิตฺตํ.) .\csromanspacer{} เหตุยา ฉ ฯเปฯ วิปาเก ตีณิ ฯเปฯ อวิคเต ฉ.}

\csromanheader{2}{1. กุสลตฺติก 100. สรณทุก}
\proseitem{36}{นกุสลํ นสรณํ ธมฺมํ ปจฺจยา อกุสโล สรโณ ธมฺโม อุปฺปชฺชติ เหตุปจฺจยา.\csromanspacer{}\csromanspacer{}เหตุยา ตีณิ.}

\prose{นกุสลํ น-อรณํ ธมฺมํ ปฏิจฺจ อพฺยากโต อรโณ ธมฺโม อุปฺปชฺชติ เหตุปจฺจยา.\csromanspacer{}\csromanspacer{}น-อพฺยากตํ น-อรณํ ธมฺมํ ปฏิจฺจ อพฺยากโต อรโณ ธมฺโม อุปฺปชฺชติ เหตุปจฺจยา.\csromanspacer{}\csromanspacer{}นกุสลํ น-อรณญฺจ น-อพฺยากตํ น-อรณญฺจ ธมฺมํ ปฏิจฺจ อพฺยากโต อรโณ ธมฺโม อุปฺปชฺชติ เหตุปจฺจยา.\csromanspacer{}\csromanspacer{}เหตุยา ตีณิ.}

\csromanheader{2}{2. เวทนาตฺติก 1. เหตุทุก}
\proseitem{37}{นสุขาย เวทนาย สมฺปยุตฺตํ นเหตุํ ธมฺมํ ปฏิจฺจ สุขาย เวทนาย สมฺปยุตฺโต เหตุ ธมฺโม อุปฺปชฺชติ เหตุปจฺจยา.\csromanspacer{}\csromanspacer{}นสุขาย เวทนาย สมฺปยุตฺตํ นเหตุํ ธมฺมํ ปฏิจฺจ ทุกฺขาย เวทนาย สมฺปยุตฺโต เหตุ ธมฺโม อุปฺปชฺชติ เหตุปจฺจยา.\csromanspacer{}\csromanspacer{}นสุขาย เวทนาย สมฺปยุตฺตํ นเหตุํ ธมฺมํ ปฏิจฺจ อทุกฺขมสุขาย เวทนาย สมฺปยุตฺโต เหตุ ธมฺโม อุปฺปชฺชติ เหตุปจฺจยา.\csromanspacer{} ตีณิ.}

\prose{นทุกฺขาย เวทนาย สมฺปยุตฺตํ นเหตุํ ธมฺมํ ปฏิจฺจ ทุกฺขาย เวทนาย สมฺปยุตฺโต เหตุ ธมฺโม อุปฺปชฺชติ เหตุปจฺจยา.\csromanspacer{} ตีณิ.}

\prose{น-อทุกฺขมสุขาย เวทนาย สมฺปยุตฺตํ นเหตุํ ธมฺมํ ปฏิจฺจ อทุกฺขมสุขาย เวทนาย สมฺปยุตฺโต เหตุ ธมฺโม อุปฺปชฺชติ เหตุปจฺจยา.\csromanspacer{} ตีณิ.}

\csromanheader{2}{ธมฺมปจฺจนียานุโลม ติกทุกปฏฺฐานปาฬิ}
\prose{\csromanfolio{386}นสุขาย เวทนาย สมฺปยุตฺตํ นเหตุญฺจ น-อทุกฺขมสุขาย เวทนาย สมฺปยุตฺตํ นเหตุญฺจ ธมฺมํ ปฏิจฺจ สุขาย เวทนาย สมฺปยุตฺโต เหตุ ธมฺโม อุปฺปชฺชติ เหตุปจฺจยา.\csromanspacer{} ตีณิ. (ทุติยํ คณิตเกน ตีณิ.) .\csromanspacer{} เหตุยา เอกวีส.}

\prose{นสุขาย เวทนาย สมฺปยุตฺตํ นนเหตุํ ธมฺมํ ปฏิจฺจ ทุกฺขาย เวทนาย สมฺปยุตฺโต นเหตุ ธมฺโม อุปฺปชฺชติ เหตุปจฺจยา.\csromanspacer{} เทฺว.}

\prose{นทุกฺขาย เวทนาย สมฺปยุตฺตํ นนเหตุํ ธมฺมํ ปฏิจฺจ. (เทฺว ปญฺหาเยว.)}

\prose{น-อทุกฺขมสุขาย เวทนาย สมฺปยุตฺตํ นนเหตุํ ธมฺมํ ปฏิจฺจ. (เทฺวเยว.\csromanspacer{} ปฐมํ คณิตเกน เอกํ, ทุติยํ คณิตเกน เอกํ, ตติยํ คณิตเกน เอกํ กาตพฺพํ.) .\csromanspacer{} เหตุยา นว.}

\csromanheader{2}{3. วิปากตฺติก 1. เหตุทุก}
\proseitem{38}{นวิปากํ นเหตุํ ธมฺมํ ปฏิจฺจ วิปาโก เหตุ ธมฺโม อุปฺปชฺชติ เหตุปจฺจยา.\csromanspacer{}\csromanspacer{}นวิปากํ นเหตุํ ธมฺมํ ปฏิจฺจ วิปากธมฺมธมฺโม เหตุ ธมฺโม อุปฺปชฺชติ เหตุปจฺจยา.\csromanspacer{}\csromanspacer{}นวิปากํ นเหตุํ ธมฺมํ ปฏิจฺจ เนววิปากนวิปากธมฺมธมฺโม เหตุ ธมฺโม อุปฺปชฺชติ เหตุปจฺจยา.\csromanspacer{} ตีณิ.}

\prose{นวิปากธมฺมธมฺมํ นเหตุํ ธมฺมํ ปฏิจฺจ วิปาโก เหตุ ธมฺโม อุปฺปชฺชติ เหตุปจฺจยา.\csromanspacer{}\csromanspacer{}นวิปากธมฺมธมฺมํ นเหตุํ ธมฺมํ ปฏิจฺจ เนววิปากนวิปากธมฺมธมฺโม เหตุ ธมฺโม อุปฺปชฺชติ เหตุปจฺจยา.\csromanspacer{} เทฺว.}

\prose{นเนววิปากนวิปากธมฺมธมฺมํ นเหตุํ ธมฺมํ ปฏิจฺจ วิปาโก เหตุ ธมฺโม อุปฺปชฺชติ เหตุปจฺจยา.\csromanspacer{}\csromanspacer{}นเนววิปากนวิปากธมฺมธมฺมํ นเหตุํ ธมฺมํ ปฏิจฺจ วิปากธมฺมธมฺโม เหตุ ธมฺโม อุปฺปชฺชติ เหตุปจฺจยา.\csromanspacer{} เทฺว.}

\prose{นวิปากํ นเหตุญฺจ นเนววิปากนวิปากธมฺมธมฺมํ นเหตุญฺจ ธมฺมํ ปฏิจฺจ วิปากธมฺมธมฺโม เหตุ ธมฺโม อุปฺปชฺชติ เหตุปจฺจยา.\csromanspacer{} เอกํ.}

\prose{นวิปากธมฺมธมฺมํ นเหตุญฺจ นเนววิปากนวิปากธมฺมธมฺมํ นเหตุญฺจ ธมฺมํ ปฏิจฺจ วิปาโก เหตุ ธมฺโม อุปฺปชฺชติ เหตุปจฺจยา.\csromanspacer{} เอกํ.}

\prose{นวิปากํ นเหตุญฺจ นวิปากธมฺมธมฺมํ นเหตุญฺจ ธมฺมํ ปฏิจฺจ วิปาโก เหตุ ธมฺโม อุปฺปชฺชติ เหตุปจฺจยา.\csromanspacer{}\csromanspacer{}นวิปากํ นเหตุญฺจ นวิปากธมฺมธมฺมํ}

\vspace{\tipitakaparskip}
\csromancenter{วิตกฺกตฺติกาทิ 1.\csromanspacer{} เหตุทุก นเหตุญฺจ ธมฺมํ ปฏิจฺจ เนววิปากนวิปากธมฺมธมฺโม เหตุ ธมฺโม อุปฺปชฺชติ เหตุปจฺจยา.\csromanspacer{} เทฺว. (สํขิตฺตํ.) .\csromanspacer{} เหตุยา เอกาทส.}
\prose{\csromanfolio{387}นวิปากํ นนเหตุํ ธมฺมํ ปฏิจฺจ วิปากธมฺมธมฺโม นเหตุ ธมฺโม อุปฺปชฺชติ เหตุปจฺจยา.\csromanspacer{}\csromanspacer{}เหตุยา อฏฺฐารส. (สํขิตฺตํ.)}

\csromanheader{2}{4. อุปาทินฺนตฺติก 1. เหตุทุก}
\proseitem{39}{น-อุปาทินฺนุปาทานิยํ นเหตุํ ธมฺมํ ปฏิจฺจ อนุปาทินฺนุปาทานิโย เหตุ ธมฺโม อุปฺปชฺชติ เหตุปจฺจยา.\csromanspacer{}\csromanspacer{}เหตุยา นว.}

\prose{น-อุปาทินฺนุปาทานิยํ นนเหตุํ ธมฺมํ ปฏิจฺจ อนุปาทินฺนุปาทานิโย นเหตุ ธมฺโม อุปฺปชฺชติ เหตุปจฺจยา.\csromanspacer{}\csromanspacer{}เหตุยา อฏฺฐารส.}

\csromanheader{2}{5. สํกิลิฏฺฐตฺติก 1. เหตุทุก}
\proseitem{40}{นสํกิลิฏฺฐสํกิเลสิกํ นเหตุํ ธมฺมํ ปฏิจฺจ อสํกิลิฏฺฐสํกิเลสิโก เหตุ ธมฺโม อุปฺปชฺชติ เหตุปจฺจยา.\csromanspacer{}\csromanspacer{}เหตุยา นว.}

\csromanheader{2}{6-11. วิตกฺกตฺติกาทิ 1. เหตุทุก}
\proseitem{41}{นสวิตกฺกสวิจารํ นเหตุํ ธมฺมํ ปฏิจฺจ สวิตกฺกสวิจาโร เหตุ ธมฺโม อุปฺปชฺชติ เหตุปจฺจยา.\csromanspacer{}\csromanspacer{}เหตุยา ปนฺนรส.}

\proseitem{42}{นปีติสหคตํ นเหตุํ ธมฺมํ ปฏิจฺจ ปีติสหคโต เหตุ ธมฺโม อุปฺปชฺชติ เหตุปจฺจยา.\csromanspacer{}\csromanspacer{}เหตุยา อฏฺฐวีส.}

\proseitem{43}{นทสฺสเนน ปหาตพฺพํ นเหตุํ ธมฺมํ ปฏิจฺจ ภาวนาย ปหาตพฺโพ เหตุ ธมฺโม อุปฺปชฺชติ เหตุปจฺจยา.\csromanspacer{}\csromanspacer{}เหตุยา นว.}

\proseitem{44}{นทสฺสเนน ปหาตพฺพเหตุกํ นเหตุํ ธมฺมํ ปฏิจฺจ ภาวนาย ปหาตพฺพเหตุโก เหตุ ธมฺโม อุปฺปชฺชติ เหตุปจฺจยา.\csromanspacer{}\csromanspacer{}เหตุยา นว.}

\proseitem{45}{น-อาจยคามิํ นเหตุํ ธมฺมํ ปฏิจฺจ อปจยคามี เหตุ ธมฺโม อุปฺปชฺชติ เหตุปจฺจยา.\csromanspacer{}\csromanspacer{}เหตุยา นว.}

\csromanheader{2}{ธมฺมปจฺจนียานุโลม ติกทุกปฏฺฐานปาฬิ}
\proseitem{46}{\csromanfolio{388}นเสกฺขํ นเหตุํ ธมฺมํ ปฏิจฺจ อเสกฺโข เหตุ ธมฺโม อุปฺปชฺชติ เหตุปจฺจยา.\csromanspacer{}\csromanspacer{}เหตุยา นว.}

\csromanheader{2}{12-13. ปริตฺตตฺติกทฺวย 1. เหตุทุก}
\proseitem{47}{นปริตฺตํ นเหตุํ ธมฺมํ ปฏิจฺจ มหคฺคโต เหตุ ธมฺโม อุปฺปชฺชติ เหตุปจฺจยา.\csromanspacer{}\csromanspacer{}เหตุยา เอกาทส.}

\proseitem{48}{นปริตฺตารมฺมณํ นเหตุํ ธมฺมํ ปฏิจฺจ ปริตฺตารมฺมโณ เหตุ ธมฺโม อุปฺปชฺชติ เหตุปจฺจยา.\csromanspacer{}\csromanspacer{}เหตุยา เตรส.}

\csromanheader{2}{14-17. หีนตฺติกาทิ 1. เหตุทุก}
\proseitem{49}{นหีนํ นเหตุํ ธมฺมํ ปฏิจฺจ มชฺฌิโม เหตุ ธมฺโม อุปฺปชฺชติ เหตุปจฺจยา.\csromanspacer{}\csromanspacer{}เหตุยา นว.}

\proseitem{50}{นมิจฺฉตฺตนิยตํ นเหตุํ ธมฺมํ ปฏิจฺจ สมฺมตฺตนิยโต เหตุ ธมฺโม อุปฺปชฺชติ เหตุปจฺจยา.\csromanspacer{}\csromanspacer{}เหตุยา นว.}

\proseitem{51}{นมคฺคารมฺมณํ นเหตุํ ธมฺมํ ปฏิจฺจ มคฺคเหตุโก เหตุ ธมฺโม อุปฺปชฺชติ เหตุปจฺจยา. (สํขิตฺตํ.) .\csromanspacer{} เหตุยา ทส.}

\proseitem{52}{น-อนุปฺปนฺโน นนเหตุ ธมฺโม อุปฺปนฺนสฺส นเหตุสฺส ธมฺมสฺส เหตุปจฺจเยน ปจฺจโย.\csromanspacer{}\csromanspacer{}น-อุปฺปาที นนเหตุ ธมฺโม อุปฺปนฺนสฺส นเหตุสฺส ธมฺมสฺส เหตุปจฺจเยน ปจฺจโย.\csromanspacer{}\csromanspacer{}น-อนุปฺปนฺโน นนเหตุ จ น-อุปฺปาที นนเหตุ จ ธมฺมา อุปฺปนฺนสฺส นเหตุสฺส ธมฺมสฺส เหตุปจฺจเยน ปจฺจโย.\csromanspacer{}\csromanspacer{}เหตุยา ตีณิ.}

\csromanheader{2}{18-19. อตีตตฺติกทฺวย 1. เหตุทุก}
\proseitem{53}{น-อตีโต นนเหตุ ธมฺโม ปจฺจุปฺปนฺนสฺส นเหตุสฺส ธมฺมสฺส เหตุปจฺจเยน ปจฺจโย.\csromanspacer{} น-อนาคโต นนเหตุ ธมฺโม ปจฺจุปฺปนฺนสฺส นเหตุสฺส ธมฺมสฺส เหตุปจฺจเยน ปจฺจโย.\csromanspacer{}\csromanspacer{}น-อตีโต นนเหตุ จ น-อนาคโต นนเหตุ จ ธมฺมา ปจฺจุปฺปนฺนสฺส นเหตุสฺส ธมฺมสฺส เหตุปจฺจเยน ปจฺจโย.\csromanspacer{}\csromanspacer{}เหตุยา ตีณิ.}

\csromanheader{2}{22. สนิทสฺสนตฺติก 2. สเหตุกทุก}
\proseitem{54}{\csromanfolio{389}น-อตีตารมฺมณํ นเหตุํ ธมฺมํ ปฏิจฺจ อตีตารมฺมโณ เหตุ ธมฺโม อุปฺปชฺชติ เหตุปจฺจยา.\csromanspacer{}\csromanspacer{}เหตุยา สตฺตรส.}

\csromanheader{2}{20-21. อชฺฌตฺตตฺติกทฺวย 1. เหตุทุก}
\proseitem{55}{น-อชฺฌตฺตํ นเหตุํ ธมฺมํ ปฏิจฺจ พหิทฺธา เหตุ ธมฺโม อุปฺปชฺชติ เหตุปจฺจยา.}

\prose{นพหิทฺธา นเหตุํ ธมฺมํ ปฏิจฺจ อชฺฌตฺโต เหตุ ธมฺโม อุปฺปชฺชติ เหตุปจฺจยา.\csromanspacer{}\csromanspacer{}เหตุยา เทฺว.}

\proseitem{56}{น-อชฺฌตฺตารมฺมณํ นเหตุํ ธมฺมํ ปฏิจฺจ อชฺฌตฺตารมฺมโณ เหตุ ธมฺโม อุปฺปชฺชติ เหตุปจฺจยา.\csromanspacer{} เทฺว.}

\prose{นพหิทฺธารมฺมณํ นเหตุํ ธมฺมํ ปฏิจฺจ พหิทฺธารมฺมโณ เหตุ ธมฺโม อุปฺปชฺชติ เหตุปจฺจยา.\csromanspacer{}\csromanspacer{}เหตุยา ฉ.}

\csromanheader{2}{22. สนิทสฺสนตฺติก 1. เหตุทุก}
\proseitem{57}{นสนิทสฺสนสปฺปฏิฆํ นเหตุํ ธมฺมํ ปฏิจฺจ อนิทสฺสน-อปฺปฏิโฆ เหตุ ธมฺโม อุปฺปชฺชติ เหตุปจฺจยา.\csromanspacer{}\csromanspacer{}น-อนิทสฺสนสปฺปฏิฆํ นเหตุํ ธมฺมํ ปฏิจฺจ อนิทสฺสน-อปฺปฏิโฆ เหตุ ธมฺโม อุปฺปชฺชติ เหตุปจฺจยา.\csromanspacer{}\csromanspacer{}นสนิทสฺสนสปฺปฏิฆํ นเหตุญฺจ น-อนิทสฺสนสปฺปฏิฆํ นเหตุญฺจ ธมฺมํ ปฏิจฺจ อนิทสฺสน-อปฺปฏิโฆ เหตุ ธมฺโม อุปฺปชฺชติ เหตุปจฺจยา.\csromanspacer{}\csromanspacer{}เหตุยา ตีณิ.}

\prose{นสนิทสฺสนสปฺปฏิฆํ นนเหตุํ ธมฺมํ ปฏิจฺจ สนิทสฺสนสปฺปฏิโฆ นเหตุ ธมฺโม อุปฺปชฺชติ เหตุปจฺจยา.\csromanspacer{}\csromanspacer{}นสนิทสฺสนสปฺปฏิฆํ นนเหตุํ ธมฺมํ ปฏิจฺจ อนิทสฺสนสปฺปฏิโฆ นเหตุ ธมฺโม อุปฺปชฺชติ เหตุปจฺจยา. (นเหตุ นนเหตุ โอวตฺตา.\csromanspacer{} ตีณิ มูลานิ.\csromanspacer{} เอกวีสติ ปญฺหา กาตพฺพา.)}

\csromanheader{2}{22. สนิทสฺสนตฺติก 2. สเหตุกทุก}
\proseitem{58}{นสนิทสฺสนสปฺปฏิฆํ นสเหตุกํ ธมฺมํ ปฏิจฺจ อนิทสฺสน-อปฺปฏิโฆ สเหตุโก ธมฺโม อุปฺปชฺชติ เหตุปจฺจยา.\csromanspacer{}\csromanspacer{}น-อนิทสฺสนสปฺปฏิฆํ นสเหตุกํ ธมฺมํ ปฏิจฺจ อนิทสฺสน-อปฺปฏิโฆ สเหตุโก ธมฺโม อุปฺปชฺชติ}

\vspace{\tipitakaparskip}
\csromancenter{ธมฺมปจฺจนียานุโลม ติกทุกปฏฺฐานปาฬิ เหตุปจฺจยา.\csromanspacer{}\csromanspacer{}นสนิทสฺสนสปฺปฏิฆํ นสเหตุกญฺจ น-อนิทสฺสนสปฺปฏิฆํ นสเหตุกญฺจ ธมฺมํ ปฏิจฺจ อนิทสฺสน-อปฺปฏิโฆ สเหตุโก ธมฺโม อุปฺปชฺชติ เหตุปจฺจยา.\csromanspacer{}\csromanspacer{}เหตุยา ตีณิ.}
\prose{\csromanfolio{390}นสนิทสฺสนสปฺปฏิฆํ น-อเหตุกํ ธมฺมํ ปฏิจฺจ สนิทสฺสนสปฺปฏิโฆ อเหตุโก ธมฺโม อุปฺปชฺชติ เหตุปจฺจยา.\csromanspacer{} สตฺต.}

\prose{น-อนิทสฺสนสปฺปฏิฆํ น-อเหตุกํ ธมฺมํ ปฏิจฺจ อนิทสฺสนสปฺปฏิโฆ อเหตุโก ธมฺโม อุปฺปชฺชติ เหตุปจฺจยา.\csromanspacer{} สตฺต.}

\prose{นสนิทสฺสนสปฺปฏิฆํ น-อเหตุกญฺจ น-อนิทสฺสนสปฺปฏิฆํ น-อเหตุกญฺจ ธมฺมํ ปฏิจฺจ อนิทสฺสนสปฺปฏิโฆ อเหตุโก ธมฺโม อุปฺปชฺชติ เหตุปจฺจยา.\csromanspacer{} สตฺต.\csromanspacer{}\csromanspacer{}เหตุยา เอกวีส.}

\csromanheader{2}{22. สนิทสฺสนตฺติก 3-6. เหตุสมฺปยุตฺตทุกาทิ}
\proseitem{59}{นสนิทสฺสนสปฺปฏิฆํ นเหตุสมฺปยุตฺตํ ธมฺมํ ปฏิจฺจ อนิทสฺสน-อปฺปฏิโฆ เหตุสมฺปยุตฺโต ธมฺโม อุปฺปชฺชติ เหตุปจฺจยา.\csromanspacer{}\csromanspacer{}เหตุยา ตีณิ.}

\prose{นสนิทสฺสนสปฺปฏิฆํ นเหตุวิปฺปยุตฺตํ ธมฺมํ ปฏิจฺจ สนิทสฺสนสปฺปฏิโฆ เหตุวิปฺปยุตฺโต ธมฺโม อุปฺปชฺชติ เหตุปจฺจยา.\csromanspacer{}\csromanspacer{}เหตุยา เอกวีส.}

\proseitem{60}{นสนิทสฺสนสปฺปฏิฆํ นเหตุญฺเจว น-อเหตุกญฺจ ธมฺมํ ปฏิจฺจ อนิทสฺสน-อปฺปฏิโฆ เหตุ เจว สเหตุโก จ ธมฺโม อุปฺปชฺชติ เหตุปจฺจยา.\csromanspacer{}\csromanspacer{}เหตุยา ตีณิ.}

\prose{นสนิทสฺสนสปฺปฏิฆํ น-อเหตุกญฺเจว นน จ เหตุํ ธมฺมํ ปฏิจฺจ อนิทสฺสน-อปฺปฏิโฆ สเหตุโก เจว น จ เหตุ ธมฺโม อุปฺปชฺชติ เหตุปจฺจยา.\csromanspacer{}\csromanspacer{}เหตุยา ตีณิ.}

\proseitem{61}{นสนิทสฺสนสปฺปฏิฆํ นเหตุญฺเจว นเหตุวิปฺปยุตฺตญฺจ ธมฺมํ ปฏิจฺจ อนิทสฺสน-อปฺปฏิโฆ เหตุ เจว เหตุสมฺปยุตฺโต จ ธมฺโม อุปฺปชฺชติ เหตุปจฺจยา.\csromanspacer{}\csromanspacer{}เหตุยา ตีณิ.}

\prose{นสนิทสฺสนสปฺปฏิฆํ นเหตุวิปฺปยุตฺตญฺเจว นน จ เหตุํ ธมฺมํ ปฏิจฺจ อนิทสฺสน-อปฺปฏิโฆ เหตุสมฺปยุตฺโต เจว น จ เหตุ ธมฺโม อุปฺปชฺชติ เหตุปจฺจยา.\csromanspacer{}\csromanspacer{}เหตุยา ตีณิ.}

\csromanheader{2}{22. สนิทสฺสนตฺติก 7. จูฬนฺตรทุก}
\proseitem{62}{\csromanfolio{391}นสนิทสฺสนสปฺปฏิฆํ นเหตุํ นสเหตุกํ ธมฺมํ ปฏิจฺจ อนิทสฺสน-อปฺปฏิโฆ นเหตุ สเหตุโก ธมฺโม อุปฺปชฺชติ เหตุปจฺจยา.\csromanspacer{}\csromanspacer{}เหตุยา ตีณิ.}

\prose{นสนิทสฺสนสปฺปฏิฆํ นเหตุํ น-อเหตุกํ ธมฺมํ ปฏิจฺจ สนิทสฺสนสปฺปฏิโฆ นเหตุ อเหตุโก ธมฺโม อุปฺปชฺชติ เหตุปจฺจยา.\csromanspacer{}\csromanspacer{}เหตุยา เอกวีส.}

\csromanheader{2}{22. สนิทสฺสนตฺติก 7. จูฬนฺตรทุก}
\proseitem{63}{นสนิทสฺสนสปฺปฏิโฆ นสปฺปจฺจโย ธมฺโม อนิทสฺสน-อปฺปฏิฆสฺส สปฺปจฺจยสฺส ธมฺมสฺส อารมฺมณปจฺจเยน ปจฺจโย. (สํขิตฺตํ.) .\csromanspacer{} อารมฺมเณ ตีณิ, อธิปติยา อุปนิสฺสเย ตีณิ.}

\proseitem{64}{นสนิทสฺสนสปฺปฏิฆํ นสนิทสฺสนํ ธมฺมํ ปฏิจฺจ สนิทสฺสนสปฺปฏิโฆ สนิทสฺสโน ธมฺโม อุปฺปชฺชติ เหตุปจฺจยา.\csromanspacer{}\csromanspacer{}เหตุยา ปญฺจ.}

\prose{น-อนิทสฺสน-อปฺปฏิโฆ นนสนิทสฺสโน ธมฺโม อนิทสฺสน-อปฺปฏิฆสฺส นสนิทสฺสนสฺส ธมฺมสฺส อารมฺมณปจฺจเยน ปจฺจโย. (สํขิตฺตํ.) .\csromanspacer{} อารมฺมเณ ตีณิ, อธิปติยา อุปนิสฺสเย ปุเรชาเต อตฺถิยา อวิปเต ตีณิ.}

\proseitem{65}{นสนิทสฺสนสปฺปฏิฆํ นสปฺปฏิฆํ ธมฺมํ ปฏิจฺจ สนิทสฺสนสปฺปฏิโฆ สปฺปฏิโฆ ธมฺโม อุปฺปชฺชติ เหตุปจฺจยา.\csromanspacer{}\csromanspacer{}เหตุยา นว.}

\prose{นสนิทสฺสนสปฺปฏิฆํ น-อปฺปฏิฆํ ธมฺมํ ปฏิจฺจ อนิทสฺสน-อปฺปฏิโฆ อปฺปฏิโฆ ธมฺโม อุปฺปชฺชติ เหตุปจฺจยา.\csromanspacer{}\csromanspacer{}เหตุยา ตีณิ.}

\proseitem{66}{นสนิทสฺสนสปฺปฏิฆํ นรูปิํ ธมฺมํ ปฏิจฺจ สนิทสฺสนสปฺปฏิโฆ รูปี ธมฺโม อุปฺปชฺชติ เหตุปจฺจยา.\csromanspacer{}\csromanspacer{}เหตุยา เอกวีส.}

\prose{นสนิทสฺสนสปฺปฏิฆํ น-อรูปิํ ธมฺมํ ปฏิจฺจ อนิทสฺสน-อปฺปฏิโฆ อรูปี ธมฺโม อุปฺปชฺชติ เหตุปจฺจยา.\csromanspacer{}\csromanspacer{}เหตุยา ตีณิ.}

\proseitem{67}{นสนิทสฺสนสปฺปฏิฆํ นโลกิยํ ธมฺมํ ปฏิจฺจ สนิทสฺสนสปฺปฏิโฆ โลกิโย ธมฺโม อุปฺปชฺชติ เหตุปจฺจยา.\csromanspacer{} สตฺต.}

\csromanheader{2}{ธมฺมปจฺจนียานุโลม ติกทุกปฏฺฐานปาฬิ}
\prose{\csromanfolio{392}น-อนิทสฺสนสปฺปฏิฆํ นโลกิยํ ธมฺมํ ปฏิจฺจ อนิทสฺสนสปฺปฏิโฆ โลกิโย ธมฺโม อุปฺปชฺชติ เหตุปจฺจยา.\csromanspacer{} สตฺต.\csromanspacer{}\csromanspacer{}เหตุยา เอกวีส.}

\prose{นสนิทสฺสนสปฺปฏิฆํ นโลกุตฺตรํ ธมฺมํ ปจฺจยา อนิทสฺสน-อปฺปฏิโฆ โลกุตฺตโร ธมฺโม อุปฺปชฺชติ เหตุปจฺจยา. (สํขิตฺตํ.) .\csromanspacer{} เหตุยา ตีณิ ฯเปฯ อวิคเต ตีณิ.}

\prose{นสนิทสฺสนสปฺปฏิฆํ นโลกุตฺตรํ ธมฺมํ ปจฺจยา อนิทสฺสน-อปฺปฏิโฆ โลกุตฺตโร ธมฺโม อุปฺปชฺชติ เหตุปจฺจยา. (สํขิตฺตํ.) .\csromanspacer{} เหตุยา ตีณิ ฯเปฯ อวิคเต ตีณิ.}

\proseitem{68}{นสนิทสฺสนสปฺปฏิฆํ นเกนจิ วิญฺเญยฺยํ ธมฺมํ ปฏิจฺจ สนิทสฺสนสปฺปฏิโฆ เกนจิ วิญฺเญยฺโย ธมฺโม อุปฺปชฺชติ เหตุปจฺจยา.\csromanspacer{}\csromanspacer{}เหตุยา ปญฺจติํส.}

\prose{นสนิทสฺสนสปฺปฏิฆํ นนเกนจิ วิญฺเญยฺยํ ธมฺมํ ปฏิจฺจ สนิทสฺสนสปฺปฏิโฆ นเกนจิ วิญฺเญยฺโย ธมฺโม อุปฺปชฺชติ เหตุปจฺจยา.\csromanspacer{}\csromanspacer{}เหตุยา ปญฺจติํส.}

\csromanheader{2}{22. สนิทสฺสนตฺติก 14. อาสวโคจฺฉก}
\proseitem{69}{นสนิทสฺสนสปฺปฏิฆํ โน-อาสวํ ธมฺมํ ปฏิจฺจ อนิทสฺสน-อปฺปฏิโฆ อาสโว ธมฺโม อุปฺปชฺชติ เหตุปจฺจยา.\csromanspacer{}\csromanspacer{}เหตุยา ตีณิ.}

\prose{นสนิทสฺสนสปฺปฏิฆํ นโน-อาสวํ ธมฺมํ ปฏิจฺจ สนิทสฺสนสปฺปฏิโฆ โน-อาสโว ธมฺโม อุปฺปชฺชติ เหตุปจฺจยา.\csromanspacer{}\csromanspacer{}เหตุยา เอกวีส.}

\proseitem{70}{นสนิทสฺสนสปฺปฏิฆํ นสาสวํ ธมฺมํ ปฏิจฺจ สนิทสฺสนสปฺปฏิโฆ สาสโว ธมฺโม อุปฺปชฺชติ เหตุปจฺจยา.\csromanspacer{}\csromanspacer{}เหตุยา เอกวีส.}

\prose{นสนิทสฺสนสปฺปฏิฆํ น-อนาสวํ ธมฺมํ ปจฺจยา อนิทสฺสน-อปฺปฏิโฆ อนาสโว ธมฺโม อุปฺปชฺชติ เหตุปจฺจยา. (สํขิตฺตํ.) .\csromanspacer{} เหตุยา ตีณิ ฯเปฯ อวิคเต ตีณิ.}

\proseitem{71}{นสนิทสฺสนสปฺปฏิฆํ น-อาสวสมฺปยุตฺตํ ธมฺมํ ปฏิจฺจ อนิทสฺสน-อปฺปฏิโฆ อาสวสมฺปยุตฺโต ธมฺโม อุปฺปชฺชติ เหตุปจฺจยา.\csromanspacer{}\csromanspacer{}เหตุยา ตีณิ.}

\csromanheader{2}{22. สนิทสฺสนตฺติก 20. สญฺโญชนาทิโคจฺฉก}
\prose{\csromanfolio{393}นสนิทสฺสนสปฺปฏิฆํ น-อาสววิปฺปยุตฺตํ ธมฺมํ ปฏิจฺจ สนิทสฺสนสปฺปฏิโฆ อาสววิปฺปยุตฺโต ธมฺโม อุปฺปชฺชติ เหตุปจฺจยา.\csromanspacer{}\csromanspacer{}เหตุยา เอกวีส.}

\proseitem{72}{นสนิทสฺสนสปฺปฏิฆํ น-อาสวญฺเจว น-อนาสวญฺจ ธมฺมํ ปฏิจฺจ อนิทสฺสน-อปฺปฏิโฆ อาสโว เจว สาสโว จ ธมฺโม อุปฺปชฺชติ เหตุปจฺจยา.\csromanspacer{}\csromanspacer{}เหตุยา ตีณิ.}

\prose{นสนิทสฺสนสปฺปฏิฆํ น-อนาสวญฺเจว นโน จ อาสวํ ธมฺมํ ปฏิจฺจ สนิทสฺสนสปฺปฏิโฆ สาสโว เจว โน จ อาสโว ธมฺโม อุปฺปชฺชติ เหตุปจฺจยา.\csromanspacer{}\csromanspacer{}เหตุยา เอกวีส.}

\proseitem{73}{นสนิทสฺสนสปฺปฏิฆํ น-อาสวญฺเจว น-อาสววิปฺปยุตฺตญฺจ ธมฺมํ ปฏิจฺจ อนิทสฺสน-อปฺปฏิโฆ อาสโว เจว อาสวสมฺปยุตฺโต จ ธมฺโม อุปฺปชฺชติ เหตุปจฺจยา.\csromanspacer{}\csromanspacer{}เหตุยา ตีณิ.}

\prose{นสนิทสฺสนสปฺปฏิฆํ น-อาสววิปฺปยุตฺตญฺเจว นโน จ อาสวํ ธมฺมํ ปฏิจฺจ อนิทสฺสน-อปฺปฏิโฆ อาสวสมฺปยุตฺโต เจว โน จ อาสโว ธมฺโม อุปฺปชฺชติ เหตุปจฺจยา.\csromanspacer{}\csromanspacer{}เหตุยา ตีณิ.}

\proseitem{74}{นสนิทสฺสนสปฺปฏิฆํ อาสววิปฺปยุตฺตํ นสาสวํ ธมฺมํ ปฏิจฺจ สนิทสฺสนสปฺปฏิโฆ อาสววิปฺปยุตฺโต สาสโว ธมฺโม อุปฺปชฺชติ เหตุปจฺจยา.\csromanspacer{}\csromanspacer{}เหตุยา เอกวีส.}

\prose{นสนิทสฺสนสปฺปฏิฆํ อาสววิปฺปยุตฺตํ น-อนาสวํ ธมฺมํ ปจฺจยา อนิทสฺสน-อปฺปฏิโฆ อาสววิปฺปยุตฺโต อนาสโว ธมฺโม อุปฺปชฺชติ เหตุปจฺจยา.\csromanspacer{}\csromanspacer{}เหตุยา ตีณิ.}

\csromanheader{2}{22. สนิทสฺสนตฺติก 20. สญฺโญชนาทิโคจฺฉก}
\proseitem{75}{นสนิทสฺสนสปฺปฏิฆํ โนสญฺโญชนํ ธมฺมํ ปฏิจฺจ อนิทสฺสน-อปฺปฏิโฆ สญฺโญชโน ธมฺโม อุปฺปชฺชติ เหตุปจฺจยา.\csromanspacer{}\csromanspacer{}เหตุยา ตีณิ.}

\proseitem{76}{นสนิทสฺสนสปฺปฏิฆํ โนคนฺถํ ธมฺมํ ปฏิจฺจ อนิทสฺสน-อปฺปฏิโฆ คนฺโถ ธมฺโม อุปฺปชฺชติ เหตุปจฺจยา.\csromanspacer{}\csromanspacer{}เหตุยา ตีณิ.}

\csromanheader{2}{ธมฺมปจฺจนียานุโลม ติกทุกปฏฺฐานปาฬิ}
\proseitem{77}{\csromanfolio{394}นสนิทสฺสนสปฺปฏิฆํ โน-โอฆํ ธมฺมํ ปฏิจฺจ อนิทสฺสน-อปฺปฏิโฆ โอโฆ ธมฺโม อุปฺปชฺชติ เหตุปจฺจยา.\csromanspacer{}\csromanspacer{}เหตุยา ตีณิ.}

\proseitem{78}{นสนิทสฺสนสปฺปฏิฆํ โนโยคํ ธมฺมํ ปฏิจฺจ อนิทสฺสน-อปฺปฏิโฆ โยโค ธมฺโม อุปฺปชฺชติ เหตุปจฺจยา.\csromanspacer{}\csromanspacer{}เหตุยา ตีณิ.}

\proseitem{79}{นสนิทสฺสนสปฺปฏิฆํ โนนีวรณํ ธมฺมํ ปฏิจฺจ อนิทสฺสน-อปฺปฏิโฆ นีวรโณ ธมฺโม อุปฺปชฺชติ เหตุปจฺจยา.\csromanspacer{}\csromanspacer{}เหตุยา ตีณิ.}

\proseitem{80}{นสนิทสฺสนสปฺปฏิฆํ โนปรามาสํ ธมฺมํ ปฏิจฺจ อนิทสฺสน-อปฺปฏิโฆ ปรามาโส ธมฺโม อุปฺปชฺชติ เหตุปจฺจยา.\csromanspacer{}\csromanspacer{}เหตุยา ตีณิ. \textbf{22.}\csromanspacer{}\textbf{ สนิทสฺสนตฺติก 55.}\csromanspacer{}\textbf{ มหนฺตรทุก}}

\proseitem{81}{นสนิทสฺสนสปฺปฏิฆํ นสารมฺมณํ ธมฺมํ ปฏิจฺจ อนิทสฺสน-อปฺปฏิโฆ สารมฺมโณ ธมฺโม อุปฺปชฺชติ เหตุปจฺจยา.\csromanspacer{}\csromanspacer{}เหตุยา ตีณิ.}

\prose{นสนิทสฺสนสปฺปฏิฆํ น-อนารมฺมณํ ธมฺมํ ปฏิจฺจ สนิทสฺสนสปฺปฏิโฆ อนารมฺมโณ ธมฺโม อุปฺปชฺชติ เหตุปจฺจยา.\csromanspacer{}\csromanspacer{}เหตุยา เอกวีส.}

\proseitem{82}{นสนิทสฺสนสปฺปฏิฆํ โนจิตฺตํ ธมฺมํ ปฏิจฺจ อนิทสฺสน-อปฺปฏิโฆ จิตฺโต ธมฺโม อุปฺปชฺชติ เหตุปจฺจยา.\csromanspacer{}\csromanspacer{}เหตุยา ตีณิ.}

\prose{นสนิทสฺสนสปฺปฏิฆํ นโนจิตฺตํ ธมฺมํ ปฏิจฺจ สนิทสฺสนสปฺปฏิโฆ โนจิตฺโต ธมฺโม อุปฺปชฺชติ เหตุปจฺจยา.\csromanspacer{}\csromanspacer{}เหตุยา เอกวีส.}

\proseitem{83}{นสนิทสฺสนสปฺปฏิฆํ นเจตสิกํ ธมฺมํ ปฏิจฺจ อนิทสฺสน-อปฺปฏิโฆ เจตสิโก ธมฺโม อุปฺปชฺชติ เหตุปจฺจยา.\csromanspacer{}\csromanspacer{}เหตุยา ตีณิ.}

\prose{นสนิทสฺสนสปฺปฏิฆํ น-อเจตสิกํ ธมฺมํ ปฏิจฺจ สนิทสฺสนสปฺปฏิโฆ อเจตสิโก ธมฺโม อุปฺปชฺชติ เหตุปจฺจยา.\csromanspacer{}\csromanspacer{}เหตุยา เอกวีส.}

\proseitem{84}{นสนิทสฺสนสปฺปฏิฆํ นจิตฺตสมฺปยุตฺตํ ธมฺมํ ปฏิจฺจ อนิทสฺสน-อปฺปฏิโฆ จิตฺตสมฺปยุตฺโต ธมฺโม อุปฺปชฺชติ เหตุปจฺจยา.\csromanspacer{}\csromanspacer{}เหตุยา ตีณิ.}

\csromanheader{2}{22. สนิทสฺสนตฺติก 55. มหนฺตรทุก}
\prose{\csromanfolio{395}นสนิทสฺสนสปฺปฏิฆํ นจิตฺตวิปฺปยุตฺตํ ธมฺมํ ปฏิจฺจ สนิทสฺสนสปฺปฏิโฆ จิตฺตวิปฺปยุตฺโต ธมฺโม อุปฺปชฺชติ เหตุปจฺจยา.\csromanspacer{}\csromanspacer{}เหตุยา เอกวีส.}

\proseitem{85}{นสนิทสฺสนสปฺปฏิฆํ นจิตฺตสํสฏฺฐํ ธมฺมํ ปฏิจฺจ อนิทสฺสน-อปฺปฏิโฆ จิตฺตสํสฏฺโฐ ธมฺโม อุปฺปชฺชติ เหตุปจฺจยา.\csromanspacer{}\csromanspacer{}เหตุยา ตีณิ.}

\prose{นสนิทสฺสนสปฺปฏิฆํ นจิตฺตวิสํสฏฺฐํ ธมฺมํ ปฏิจฺจ สนิทสฺสนสปฺปฏิโฆ จิตฺตวิสํสฏฺโฐ ธมฺโม อุปฺปชฺชติ เหตุปจฺจยา.\csromanspacer{}\csromanspacer{}เหตุยา เอกวีส.}

\proseitem{86}{นสนิทสฺสนสปฺปฏิฆํ โนจิตฺตสมุฏฺฐานํ ธมฺมํ ปฏิจฺจ สนิทสฺสนสปฺปฏิโฆ จิตฺตสมุฏฺฐาโน ธมฺโม อุปฺปชฺชติ เหตุปจฺจยา.\csromanspacer{}\csromanspacer{}เหตุยา เอกวีส.}

\prose{นสนิทสฺสนสปฺปฏิฆํ นโนจิตฺตสมุฏฺฐานํ ธมฺมํ ปฏิจฺจ อนิทสฺสน-อปฺปฏิโฆ โนจิตฺตสมุฏฺฐาโน ธมฺโม อุปฺปชฺชติ เหตุปจฺจยา.\csromanspacer{}\csromanspacer{}เหตุยา ตีณิ.}

\proseitem{87}{นสนิทสฺสนสปฺปฏิฆํ โนจิตฺตสหภุํ ธมฺมํ ปฏิจฺจ อนิทสฺสน-อปฺปฏิโฆ จิตฺตสหภู ธมฺโม อุปฺปชฺชติ เหตุปจฺจยา.\csromanspacer{}\csromanspacer{}เหตุยา ปญฺจ.}

\prose{นสนิทสฺสนสปฺปฏิฆํ นโนจิตฺตสหภุํ ธมฺมํ ปฏิจฺจ สนิทสฺสนสปฺปฏิโฆ โนจิตฺตสหภู ธมฺโม อุปฺปชฺชติ เหตุปจฺจยา.\csromanspacer{}\csromanspacer{}เหตุยา เอกวีส.}

\proseitem{88}{นสนิทสฺสนสปฺปฏิฆํ โนจิตฺตานุปริวตฺติํ ธมฺมํ ปฏิจฺจ อนิทสฺสน-อปฺปฏิโฆ จิตฺตานุปริวตฺตี ธมฺโม อุปฺปชฺชติ เหตุปจฺจยา.\csromanspacer{}\csromanspacer{}เหตุยา ปญฺจ.}

\prose{นสนิทสฺสนสปฺปฏิฆํ นโนจิตฺตนุปริวตฺติํ ธมฺมํ ปฏิจฺจ สนิทสฺสนสปฺปฏิโฆ โนจิตฺตานุปริวตฺตี ธมฺโม อุปฺปชฺชติ เหตุปจฺจยา.\csromanspacer{}\csromanspacer{}เหตุยา เอกวีส.}

\proseitem{89}{นสนิทสฺสนสปฺปฏิฆํ โนจิตฺตสํสฏฺฐสมุฏฺฐานํ ธมฺมํ ปฏิจฺจ อนิทสฺสน-อปฺปฏิโฆ จิตฺตสํสฏฺฐสมุฏฺฐาโน ธมฺโม อุปฺปชฺชติ เหตุปจฺจยา.\csromanspacer{}\csromanspacer{}เหตุยา ตีณิ.}

\prose{นสนิทสฺสนสปฺปฏิฆํ นโนจิตฺตสํสฏฺฐสมุฏฺฐานํ ธมฺมํ ปฏิจฺจ สนิทสฺสนสปฺปฏิโฆ โนจิตฺตสํสฏฺฐสมุฏฺฐาโน ธมฺโม อุปฺปชฺชติ เหตุปจฺจยา.\csromanspacer{}\csromanspacer{}เหตุยา เอกวีส.}

\csromanheader{2}{ธมฺมปจฺจนียานุโลม ติกทุกปฏฺฐานปาฬิ}
\proseitem{90}{\csromanfolio{396}นสนิทสฺสนสปฺปฏิฆํ โนจิตฺตสํสฏฺฐสมุฏฺฐานสหภุํ ธมฺมํ ปฏิจฺจ อนิทสฺสน-อปฺปฏิโฆ จิตฺตสํสฏฺฐสมุฏฺฐานสหภู ธมฺโม อุปฺปชฺชติ เหตุปจฺจยา.\csromanspacer{}\csromanspacer{}เหตุยา ตีณิ.}

\prose{นสนิทสฺสนสปฺปฏิฆํ นโนจิตฺตสํสฏฺฐสมุฏฺฐานสหภุํ ธมฺมํ ปฏิจฺจ สนิทสฺสนสปฺปฏิโฆ โนจิตฺตสํสฏฺฐสมุฏฺฐานสหภู ธมฺโม อุปฺปชฺชติ เหตุปจฺจยา.\csromanspacer{}\csromanspacer{}เหตุยา เอกวีส.}

\proseitem{91}{นสนิทสฺสนสปฺปฏิฆํ โนจิตฺตสํสฏฺฐสมุฏฺฐานานุปริวตฺติํ ธมฺมํ ปฏิจฺจ อนิทสฺสน-อปฺปฏิโฆ จิตฺตสํสฏฺฐสมุฏฺฐานานุปริวตฺตี ธมฺโม อุปฺปชฺชติ เหตุปจฺจยา.\csromanspacer{}\csromanspacer{}เหตุยา ตีณิ.}

\prose{นสนิทสฺสนสปฺปฏิฆํ นโนจิตฺตสํสฏฺฐสมุฏฺฐานานุปริวตฺติํ ธมฺมํ ปฏิจฺจ สนิทสฺสนสปฺปฏิโฆ โนจิตฺตสํสฏฺฐสมุฏฺฐานานุปริวตฺตี ธมฺโม อุปฺปชฺชติ เหตุปจฺจยา.\csromanspacer{}\csromanspacer{}เหตุยา เอกวีส.}

\proseitem{92}{นสนิทสฺสนสปฺปฏิฆํ น-อชฺฌตฺติกํ ธมฺมํ ปฏิจฺจ อนิทสฺสนสปฺปฏิโฆ อชฺฌตฺติโก ธมฺโม อุปฺปชฺชติ เหตุปจฺจยา.\csromanspacer{}\csromanspacer{}เหตุยา เอกาทส.}

\prose{นสนิทสฺสนสปฺปฏิฆํ นพาหิรํ ธมฺมํ ปฏิจฺจ สนิทสฺสนสปฺปฏิโฆ พาหิโร ธมฺโม อุปฺปชฺชติ เหตุปจฺจยา.\csromanspacer{}\csromanspacer{}เหตุยา เอกวีส.}

\proseitem{93}{นสนิทสฺสนสปฺปฏิฆํ โน-อุปาทา ธมฺมํ ปฏิจฺจ สนิทสฺสนสปฺปฏิโฆ อุปาทา ธมฺโม อุปฺปชฺชติ เหตุปจฺจยา.\csromanspacer{}\csromanspacer{}เหตุยา ปญฺจติํส.}

\prose{นสนิทสฺสนสปฺปฏิฆํ นโน-อุปาทา ธมฺมํ ปฏิจฺจ อนิทสฺสน-อปฺปฏิโฆ โน-อุปาทา ธมฺโม อุปฺปชฺชติ เหตุปจฺจยา.\csromanspacer{}\csromanspacer{}เหตุยา ตีณิ.}

\proseitem{94}{นสนิทสฺสนสปฺปฏิฆํ น-อนุปาทินฺนํ ธมฺมํ ปฏิจฺจ สนิทสฺสนสปฺปฏิโฆ อนุปาทินฺโน ธมฺโม อุปฺปชฺชติ เหตุปจฺจยา.\csromanspacer{}\csromanspacer{}เหตุยา เอกวีส.}

\csromanheader{2}{22. สนิทสฺสนตฺติก 69. อุปาทานทุก}
\proseitem{95}{นสนิทสฺสนสปฺปฏิฆํ โน-อุปาทานํ ธมฺมํ ปฏิจฺจ อนิทสฺสน-อปฺปฏิโฆ อุปาทาโน ธมฺโม อุปฺปชฺชติ เหตุปจฺจยา.\csromanspacer{}\csromanspacer{}เหตุยา ตีณิ.}

\vspace{\tipitakaparskip}
\csromancenter{สนิทสฺสนตฺติก 83-99.\csromanspacer{} ปิฏฺฐิทุก \textbf{22.}\csromanspacer{}\textbf{ สนิทสฺสนตฺติก 75.}\csromanspacer{}\textbf{ กิเลสทุก}}
\proseitem{96}{\csromanfolio{397}นสนิทสฺสนสปฺปฏิฆํ โนกิเลสํ ธมฺมํ ปฏิจฺจ อนิทสฺสน-อปฺปฏิโฆ กิเลโส ธมฺโม อุปฺปชฺชติ เหตุปจฺจยา.\csromanspacer{}\csromanspacer{}เหตุยา ตีณิ.}

\csromanheader{2}{22. สนิทสฺสนตฺติก 83-99. ปิฏฺฐิทุก}
\proseitem{97}{นสนิทสฺสนสปฺปฏิฆํ นทสฺสเนน ปหาตพฺพํ ธมฺมํ ปจฺจยา อนิทสฺสน-อปฺปฏิโฆ ทสฺสเนน ปหาตพฺโพ ธมฺโม อุปฺปชฺชติ เหตุปจฺจยา.\csromanspacer{}\csromanspacer{}เหตุยา ตีณิ.}

\prose{นสนิทสฺสนสปฺปฏิฆํ นนทสฺสเนน ปหาตพฺพํ ธมฺมํ ปฏิจฺจ สนิทสฺสนสปฺปฏิโฆ นทสฺสเนน ปหาตพฺโพ ธมฺโม อุปฺปชฺชติ เหตุปจฺจยา.\csromanspacer{}\csromanspacer{}เหตุยา เอกวีส.}

\proseitem{98}{นสนิทสฺสนสปฺปฏิฆํ นภาวนาย ปหาตพฺพํ ธมฺมํ ปจฺจยา อนิทสฺสน-อปฺปฏิโฆ ภาวนาย ปหาตพฺโพ ธมฺโม อุปฺปชฺชติ เหตุปจฺจยา.\csromanspacer{}\csromanspacer{}เหตุยา ตีณิ.}

\prose{นสนิทสฺสนสปฺปฏิฆํ นนภาวนาย ปหาตพฺพํ ธมฺมํ ปฏิจฺจ สนิทสฺสนสปฺปฏิโฆ นภาวนาย ปหาตพฺโพ ธมฺโม อุปฺปชฺชติ เหตุปจฺจยา.\csromanspacer{}\csromanspacer{}เหตุยา เอกวีส.}

\proseitem{99}{นสนิทสฺสนสปฺปฏิฆํ นทสฺสเนน ปหาตพฺพเหตุกํ ธมฺมํ ปฏิจฺจ อนิทสฺสน-อปฺปฏิโฆ ทสฺสเนน ปหาตพฺพเหตุโก ธมฺโม อุปฺปชฺชติ เหตุปจฺจยา.\csromanspacer{} เหตุยา ตีณิ.}

\prose{นสนิทสฺสนสปฺปฏิฆํ นนทสฺสเนน ปหาตพฺพเหตุกํ ธมฺมํ ปฏิจฺจ สนิทสฺสนสปฺปฏิโฆ นทสฺสเนน ปหาตพฺพเหตุโก ธมฺโม อุปฺปชฺชติ เหตุปจฺจยา.\csromanspacer{}\csromanspacer{}เหตุยา เอกวีส.}

\proseitem{100}{นสนิทสฺสนสปฺปฏิฆํ นภาวนาย ปหาตพฺพเหตุกํ ธมฺมํ ปจฺจยา อนิทสฺสน-อปฺปฏิโฆ ภาวนาย ปหาตพฺพเหตุโก ธมฺโม อุปฺปชฺชติ เหตุปจฺจยา.\csromanspacer{}\csromanspacer{}เหตุยา ตีณิ.}

\prose{นสนิทสฺสนสปฺปฏิฆํ นนภาวนาย ปหาตพฺพเหตุกํ ธมฺมํ ปฏิจฺจ สนิทสฺสนสปฺปฏิโฆ นภาวนาย ปหาตพฺพเหตุโก ธมฺโม อุปฺปชฺชติ เหตุปจฺจยา.\csromanspacer{}\csromanspacer{}เหตุยา เอกวีส.}

\csromanheader{2}{ธมฺมปจฺจนียานุโลม ติกทุกปฏฺฐานปาฬิ}
\proseitem{101}{\csromanfolio{398}นสนิทสฺสนสปฺปฏิฆํ นสวิตกฺกํ ธมฺมํ ปฏิจฺจ อนิทสฺสน-อปฺปฏิโฆ สวิตกฺโก ธมฺโม อุปฺปชฺชติ เหตุปจฺจยา.\csromanspacer{}\csromanspacer{}เหตุยา ตีณิ.}

\prose{นสนิทสฺสนสปฺปฏิฆํ น-อวิตกฺกํ ธมฺมํ ปฏิจฺจ สนิทสฺสนสปฺปฏิโฆ อวิตกฺโก ธมฺโม อุปฺปชฺชติ เหตุปจฺจยา.\csromanspacer{}\csromanspacer{}เหตุยา เอกวีส.}

\proseitem{102}{นสนิทสฺสนสปฺปฏิฆํ นสวิจารํ ธมฺมํ ปฏิจฺจ อนิทสฺสน-อปฺปฏิโฆ สวิจาโร ธมฺโม อุปฺปชฺชติ เหตุปจฺจยา.\csromanspacer{}\csromanspacer{}เหตุยา ตีณิ.}

\prose{นสนิทสฺสนสปฺปฏิฆํ น-อวิจารํ ธมฺมํ ปฏิจฺจ สนิทสฺสนสปฺปฏิโฆ อวิจาโร ธมฺโม อุปฺปชฺชติ เหตุปจฺจยา.\csromanspacer{}\csromanspacer{}เหตุยา เอกวีส.}

\proseitem{103}{นสนิทสฺสนสปฺปฏิฆํ นสปฺปีติกํ ธมฺมํ ปฏิจฺจ อนิทสฺสน-อปฺปฏิโฆ สปฺปีติโก ธมฺโม อุปฺปชฺชติ เหตุปจฺจยา.\csromanspacer{}\csromanspacer{}เหตุยา ตีณิ.}

\prose{นสนิทสฺสนสปฺปฏิฆํ น-อปฺปีติกํ ธมฺมํ ปฏิจฺจ สนิทสฺสนสปฺปฏิโฆ อปฺปีติโก ธมฺโม อุปฺปชฺชติ เหตุปจฺจยา.\csromanspacer{}\csromanspacer{}เหตุยา เอกวีส.}

\proseitem{104}{นสนิทสฺสนสปฺปฏิฆํ นปีติสหคตํ ธมฺมํ ปฏิจฺจ อนิทสฺสน-อปฺปฏิโฆ ปีติสหคโต ธมฺโม อุปฺปชฺชติ เหตุปจฺจยา.\csromanspacer{}\csromanspacer{}เหตุยา ตีณิ.}

\prose{นสนิทสฺสนสปฺปฏิฆํ นนปีติสหคตํ ธมฺมํ ปฏิจฺจ สนิทสฺสนสปฺปฏิโฆ นปีติสหคโต ธมฺโม อุปฺปชฺชติ เหตุปจฺจยา.\csromanspacer{}\csromanspacer{}เหตุยา เอกวีส.}

\proseitem{105}{นสนิทสฺสนสปฺปฏิฆํ นสุขสหคตํ ธมฺมํ ปฏิจฺจ อนิทสฺสน-อปฺปฏิโฆ สุขสหคโต ธมฺโม อุปฺปชฺชติ เหตุปจฺจยา.\csromanspacer{}\csromanspacer{}เหตุยา ตีณิ.}

\prose{นสนิทสฺสนสปฺปฏิฆํ นนสุขสหคตํ ธมฺมํ ปฏิจฺจ สนิทสฺสนสปฺปฏิโฆ นสุขสหคโต ธมฺโม อุปฺปชฺชติ เหตุปจฺจยา.\csromanspacer{}\csromanspacer{}เหตุยา เอกวีส.}

\proseitem{106}{นสนิทสฺสนสปฺปฏิฆํ น-อุเปกฺขาสหคตํ ธมฺมํ ปฏิจฺจ อนิทสฺสน-อปฺปฏิโฆ อุเปกฺขาสหคโต ธมฺโม อุปฺปชฺชติ เหตุปจฺจยา.\csromanspacer{}\csromanspacer{}เหตุยา ตีณิ.}

\csromanheader{2}{22. สนิทสฺสนตฺติก 83-99. ปิฏฺฐิทุก}
\prose{\csromanfolio{399}นสนิทสฺสนสปฺปฏิฆํ นน-อุเปกฺขาสหคตํ ธมฺมํ ปฏิจฺจ สนิทสฺสนสปฺปฏิโฆ น-อุเปกฺขาสหคโต ธมฺโม อุปฺปชฺชติ เหตุปจฺจยา.\csromanspacer{}\csromanspacer{}เหตุยา เอกวีส.}

\proseitem{107}{นสนิทสฺสนสปฺปฏิฆํ นกามาวจรํ ธมฺมํ ปฏิจฺจ สนิทสฺสนสปฺปฏิโฆ กามาวจโร ธมฺโม อุปฺปชฺชติ เหตุปจฺจยา.\csromanspacer{}\csromanspacer{}เหตุยา เอกวีส.}

\prose{นสนิทสฺสนสปฺปฏิฆํ นนกามาวจรํ ธมฺมํ ปฏิจฺจ อนิทสฺสน-อปฺปฏิโฆ นกามาวจโร ธมฺโม อุปฺปชฺชติ เหตุปจฺจยา.\csromanspacer{}\csromanspacer{}เหตุยา ตีณิ.}

\proseitem{108}{นสนิทสฺสนสปฺปฏิฆํ นรูปาวจรํ ธมฺมํ ปฏิจฺจ อนิทสฺสน-อปฺปฏิโฆ รูปาวจโร ธมฺโม อุปฺปชฺชติ เหตุปจฺจยา.\csromanspacer{}\csromanspacer{}เหตุยา ติณิ.}

\prose{นสนิทสฺสนสปฺปฏิฆํ นนรูปาวจรํ ธมฺมํ ปฏิจฺจ สนิทสฺสนสปฺปฏิโฆ นรูปาวจโร ธมฺโม อุปฺปชฺชติ เหตุปจฺจยา.\csromanspacer{}\csromanspacer{}เหตุยา เอกวีส.}

\proseitem{109}{นสนิทสฺสนสปฺปฏิฆํ น-อรูปาวจรํ ธมฺมํ ปจฺจยา อนิทสฺสน-อปฺปฏิโฆ อรูปาวจโร ธมฺโม อุปฺปชฺชติ เหตุปจฺจยา.\csromanspacer{}\csromanspacer{}เหตุยา ตีณิ.}

\prose{นสนิทสฺสนสปฺปฏิฆํ นน-อรูปาวจรํ ธมฺมํ ปฏิจฺจ สนิทสฺสนสปฺปฏิโฆ น-อรูปาวจโร ธมฺโม อุปฺปชฺชติ เหตุปจฺจยา.\csromanspacer{}\csromanspacer{}เหตุยา เอกวีส.}

\proseitem{110}{นสนิทสฺสนสปฺปฏิฆํ นปริยาปนฺนํ ธมฺมํ ปฏิจฺจ สนิทสฺสนสปฺปฏิโฆ ปริยาปนฺโน ธมฺโม อุปฺปชฺชติ เหตุปจฺจยา.\csromanspacer{}\csromanspacer{}เหตุยา เอกวีส.}

\prose{นสนิทสฺสนสปฺปฏิฆํ น-อปริยาปนฺนํ ธมฺมํ ปจฺจยา อนิทสฺสน-อปฺปฏิโฆ อปริยาปนฺโน ธมฺโม อุปฺปชฺชติ เหตุปจฺจยา.\csromanspacer{}\csromanspacer{}เหตุยา ตีณิ.}

\proseitem{111}{นสนิทสฺสนสปฺปฏิฆํ นนิยฺยานิกํ ธมฺมํ ปจฺจยา อนิทสฺสน-อปฺปฏิโฆ นิยฺยานิโก ธมฺโม อุปฺปชฺชติ เหตุปจฺจยา.\csromanspacer{}\csromanspacer{}เหตุยา ตีณิ.}

\prose{นสนิทสฺสนสปฺปฏิฆํ น-อนิยฺยานิกํ ธมฺมํ ปฏิจฺจ สนิทสฺสนสปฺปฏิโฆ อนิยฺยานิโก ธมฺโม อุปฺปชฺชติ เหตุปจฺจยา.\csromanspacer{}\csromanspacer{}เหตุยา เอกวีส.}

\csromanheader{2}{ธมฺมปจฺจนียานุโลม ติกทุกปฏฺฐานปาฬิ}
\proseitem{112}{\csromanfolio{400}นสนิทสฺสนสปฺปฏิฆํ นนิยตํ ธมฺมํ ปจฺจยา อนิทสฺสน-อปฺปฏิโฆ นิยโต ธมฺโม อุปฺปชฺชติ เหตุปจฺจยา.\csromanspacer{}\csromanspacer{}เหตุยา ตีณิ.}

\prose{นสนิทสฺสนสปฺปฏิฆํ น-อนิยตํ ธมฺมํ ปฏิจฺจ สนิทสฺสนสปฺปฏิโฆ อนิยโต ธมฺโม อุปฺปชฺชติ เหตุปจฺจยา.\csromanspacer{}\csromanspacer{}เหตุยา เอกวีส.}

\proseitem{113}{นสนิทสฺสนสปฺปฏิฆํ นส-อุตฺตรํ ธมฺมํ ปฏิจฺจ สนิทสฺสนสปฺปฏิโฆ ส-อุตฺตโร ธมฺโม อุปฺปชฺชติ เหตุปจฺจยา.\csromanspacer{}\csromanspacer{}เหตุยา เอกวีส.}

\prose{นสนิทสฺสนสปฺปฏิฆํ น-อนุตฺตรํ ธมฺมํ ปจฺจยา อนิทสฺสน-อปฺปฏิโฆ อนุตฺตโร ธมฺโม อุปฺปชฺชติ เหตุปจฺจยา.\csromanspacer{}\csromanspacer{}เหตุยา ตีณิ.}

\csromanheader{2}{22. สนิทสฺสนตฺติก 100. สรณทุก}
\proseitem{114}{นสนิทสฺสนสปฺปฏิฆํ นสรณํ ธมฺมํ ปจฺจยา อนิทสฺสน-อปฺปฏิโฆ สรโณ ธมฺโม อุปฺปชฺชติ เหตุปจฺจยา.\csromanspacer{}\csromanspacer{}เหตุยา ตีณิ.}

\prose{นสนิทสฺสนสปฺปฏิฆํ น-อรณํ ธมฺมํ ปฏิจฺจ สนิทสฺสนสปฺปฏิโฆ อรโณ ธมฺโม อุปฺปชฺชติ เหตุปจฺจยา.\csromanspacer{} สตฺต.}

\prose{น-อนิทสฺสนสปฺปฏิฆํ น-อรณํ ธมฺมํ ปฏิจฺจ สนิทสฺสนสปฺปฏิโฆ อรโณ ธมฺโม อุปฺปชฺชติ เหตุปจฺจยา.\csromanspacer{} สตฺต.}

\prose{นสนิทสฺสนสปฺปฏิฆํ น-อรณญฺจ น-อนิทสฺสนสปฺปฏิฆํ น-อรณญฺจ ธมฺมํ ปฏิจฺจ สนิทสฺสนสปฺปฏิโฆ อรโณ ธมฺโม อุปฺปชฺชติ เหตุปจฺจยา.\csromanspacer{} สตฺต.\csromanspacer{}\csromanspacer{}เหตุยา เอกวีส. (สหชาตวารมฺปิ ปฏิจฺจวารสทิสํ วิตฺถาเรตพฺพํ.)}

\csromanheader{2}{\textbf{เหตุ อารมฺมณ}}
\proseitem{115}{นสนิทสฺสนสปฺปฏิโฆ น-อรโณ ธมฺโม สนิทสฺสนสปฺปฏิฆสฺส อรณสฺส ธมฺมสฺส เหตุปจฺจเยน ปจฺจโย.\csromanspacer{} สตฺต.}

\prose{น-อนิทสฺสนสปฺปฏิโฆ น-อรโณ ธมฺโม สนิทสฺสนสปฺปฏิฆสฺส อรณสฺส ธมฺมสฺส เหตุปจฺจเยน ปจฺจโย.\csromanspacer{} สตฺต.}

\csromanheader{2}{22. สนิทสฺสนตฺติก 100. สรณทุก}
\prose{\csromanfolio{401}นสนิทสฺสนสปฺปฏิโฆ น-อรโณ จ น-อนิทสฺสนสปฺปฏิโฆ น-อรโณ จ ธมฺมา สนิทสฺสนสปฺปฏิฆสฺส อรณสฺส ธมฺมสฺส เหตุปจฺจเยน ปจฺจโย.\csromanspacer{} สตฺต.}

\proseitem{116}{นสนิทสฺสนสปฺปฏิโฆ น-อรโณ ธมฺโม อนิทสฺสน-อปฺปฏิฆสฺส อรณสฺส ธมฺมสฺส อารมฺมณปจฺจเยน ปจฺจโย.\csromanspacer{}\csromanspacer{}น-อนิทสฺสนสปฺปฏิโฆ น-อรโณ ธมฺโม อนิทสฺสน-อปฺปฏิฆสฺส อรณสฺส ธมฺมสฺส อารมฺมณปจฺจเยน ปจฺจโย.\csromanspacer{}\csromanspacer{}นสนิทสฺสนสปฺปฏิโฆ น-อรโณ จ น-อนิทสฺสนสปฺปฏิโฆ น-อรโณ จ ธมฺมา อนิทสฺสน-อปฺปฏิฆสฺส อรณสฺส ธมฺมสฺส อารมฺมณปจฺจเยน ปจฺจโย.}

\proseitem{117}{เหตุยา เอกวีส, อารมฺมเณ ตีณิ ฯเปฯ อวิคเต เอกวีส.}

\prose{(ยถา กุสลตฺติเก ปญฺหาวารสฺส อนุโลมมฺปิ ปจฺจนียมฺปิ อนุโลมปจฺจนียมฺปิ ปจฺจนียานุโลมมฺปิ คณิตํ, เอวํ คเณตพฺพํ.)}

\nitthitam{\textbf{ธมฺมปจฺจนียานุโลเม ติกทุกปฏฺฐานํ นิฏฺฐิตํ.}}
\pitaka{\textbf{อภิธมฺมปิฏก}}
\csromanheader{2}{\textbf{ธมฺมปจฺจนียานุโลม}}
\gambhira{\textbf{ติกติกปฏฺฐานปาฬิ}}
\namakkaram{นโม ตสฺส ภควโต อรหโต สมฺมาสมฺพุทฺธสฺส.}
\csromanheader{2}{1. กุสลตฺติก 1. เวทนาตฺติก}
\proseitem{1}{\csromanfolio{403}นกุสลํ นสุขาย เวทนาย สมฺปยุตฺตํ ธมฺมํ ปฏิจฺจ อกุสโล สุขาย เวทนาย สมฺปยุตฺโต ธมฺโม อุปฺปชฺชติ เหตุปจฺจยา.\csromanspacer{}\csromanspacer{}นกุสลํ นสุขาย เวทนาย สมฺปยุตฺตํ ธมฺมํ ปฏิจฺจ อพฺยากโต สุขาย เวทนาย สมฺปยุตฺโต ธมฺโม อุปฺปชฺชติ เหตุปจฺจยา.\csromanspacer{} เทฺว.}

\prose{น-อกุสลํ นสุขาย เวทนาย สมฺปยุตฺตํ ธมฺมํ ปฏิจฺจ กุสโล สุขาย เวทนาย สมฺปยุตฺโต ธมฺโม อุปฺปชฺชติ เหตุปจฺจยา.\csromanspacer{}\csromanspacer{}น-อกุสลํ นสุขาย เวทนาย สมฺปยุตฺตํ ธมฺมํ ปฏิจฺจ อพฺยากโต สุขาย เวทนาย สมฺปยุตฺโต ธมฺโม อุปฺปชฺชติ เหตุปจฺจยา.\csromanspacer{} เทฺว.}

\prose{น-อพฺยากตํ นสุขาย เวทนาย สมฺปยุตฺตํ ธมฺมํ ปฏิจฺจ กุสโล สุขาย เวทนาย สมฺปยุตฺโต ธมฺโม อุปฺปชฺชติ เหตุปจฺจยา.\csromanspacer{}\csromanspacer{}น-อพฺยากตํ นสุขาย เวทนาย สมฺปยุตฺตํ ธมฺมํ ปฏิจฺจ อกุสโล สุขาย เวทนาย สมฺปยุตฺโต ธมฺโม อุปฺปชฺชติ เหตุปจฺจยา.\csromanspacer{} เทฺว.}

\prose{นกุสลํ นสุขาย เวทนาย สมฺปยุตฺตญฺจ น-อพฺยากตํ นสุขาย เวทนาย สมฺปยุตฺตญฺจ ธมฺมํ ปฏิจฺจ อกุสโล สุขาย เวทนาย สมฺปยุตฺโต ธมฺโม อุปฺปชฺชติ เหตุปจฺจยา.\csromanspacer{} เอกํ.}

\csromanheader{2}{ธมฺมปจฺจนียานุโลม ติกติกปฏฺฐานปาฬิ}
\prose{\csromanfolio{404}น-อกุสลํ นสุขาย เวทนาย สมฺปยุตฺตญฺจ น-อพฺยากตํ นสุขาย เวทนาย สมฺปยุตฺตญฺจ ธมฺมํ ปฏิจฺจ กุสโล สุขาย เวทนาย สมฺปยุตฺโต ธมฺโม อุปฺปชฺชติ เหตุปจฺจยา.\csromanspacer{} เอกํ.}

\prose{นกุสลํ นสุขาย เวทนาย สมฺปยุตฺตญฺจ น-อกุสลํ นสุขาย เวทนาย สมฺปยุตฺตญฺจ ธมฺมํ ปฏิจฺจ อพฺยากโต สุขาย เวทนาย สมฺปยุตฺโต ธมฺโม อุปฺปชฺชติ เหตุปจฺจยา.\csromanspacer{} เอกํ.\csromanspacer{}\csromanspacer{}เหตุยา นว.}

\proseitem{2}{นกุสลํ นทุกฺขาย เวทนาย สมฺปยุตฺตํ ธมฺมํ ปฏิจฺจ อกุสโล ทุกฺขาย เวทนาย สมฺปยุตฺโต ธมฺโม อุปฺปชฺชติ เหตุปจฺจยา.\csromanspacer{} เอกํ.}

\prose{น-อพฺยากตํ นทุกฺขาย เวทนาย สมฺปยุตฺตํ ธมฺมํ ปฏิจฺจ อกุสโล ทุกฺขาย เวทนาย สมฺปยุตฺโต ธมฺโม อุปฺปชฺชติ เหตุปจฺจยา.\csromanspacer{} เอกํ.}

\prose{นกุสลํ นทุกฺขาย เวทนาย สมฺปยุตฺตญฺจ น-อพฺยากตํ นทุกฺขาย เวทนาย สมฺปยุตฺตญฺจ ธมฺมํ ปฏิจฺจ อกุสโล ทุกฺขาย เวทนาย สมฺปยุตฺโต ธมฺโม อุปฺปชฺชติ เหตุปจฺจยา.\csromanspacer{} เอกํ.\csromanspacer{} เหตุยา ตีณิ.}

\proseitem{3}{นกุสลํ น-อทุกฺขมสุขาย เวทนาย สมฺปยุตฺตํ ธมฺมํ ปฏิจฺจ อกุสโล อทุกฺขมสุขาย เวทนาย สมฺปยุตฺโต ธมฺโม อุปฺปชฺชติ เหตุปจฺจยา.\csromanspacer{}\csromanspacer{}นกุสลํ น-อทุกฺขมสุขาย เวทนาย สมฺปยุตฺตํ ธมฺมํ ปฏิจฺจ อพฺยากโต อทุกฺขมสุขาย เวทนาย สมฺปยุตฺโต ธมฺโม อุปฺปชฺชติ เหตุปจฺจยา.\csromanspacer{} เทฺว.}

\prose{น-อกุสลํ น-อทุกฺขมสุขาย เวทนาย สมฺปยุตฺตํ ธมฺมํ ปฏิจฺจ กุสโล อทุกฺขมสุขาย เวทนาย สมฺปยุตฺโต ธมฺโม อุปฺปชฺชติ เหตุปจฺจยา.\csromanspacer{}\csromanspacer{}น-อกุสลํ น-อทุกฺขมสุขาย เวทนาย สมฺปยุตฺตํ ธมฺมํ ปฏิจฺจ อพฺยากโต อทุกฺขมสุขาย เวทนาย สมฺปยุตฺโต ธมฺโม อุปฺปชฺชติ เหตุปจฺจยา.\csromanspacer{} เทฺว.}

\prose{น-อพฺยากตํ น-อทุกฺขมสุขาย เวทนาย สมฺปยุตฺตํ ธมฺมํ ปฏิจฺจ กุสโล อทุกฺขมสุขาย เวทนาย สมฺปยุตฺโต ธมฺโม อุปฺปชฺชติ เหตุปจฺจยา.\csromanspacer{}\csromanspacer{}น-อพฺยากตํ น-อทุกฺขมสุขาย เวทนาย สมฺปยุตฺตํ ธมฺมํ ปฏิจฺจ อกุสโล อทุกฺขมสุขาย เวทนาย สมฺปยุตฺโต ธมฺโม อุปฺปชฺชติ เหตุปจฺจยา.\csromanspacer{} เทฺว.}

\prose{นกุสลํ น-อทุกฺขมสุขาย เวทนาย สมฺปยุตฺตญฺจ น-อพฺยากตํ น-อทุกฺขมสุขาย เวทนาย สมฺปยุตฺตญฺจ ธมฺมํ ปฏิจฺจ อกุสโล อทุกฺขมสุขาย เวทนาย สมฺปยุตฺโต ธมฺโม อุปฺปชฺชติ เหตุปจฺจยา.\csromanspacer{} เอกํ.}

\csromanheader{2}{1. กุสลตฺติก 3. อุปาทินฺนตฺติก}
\prose{\csromanfolio{405}น-อกุสลํ น-อทุกฺขมสุขาย เวทนาย สมฺปยุตฺตญฺจ น-อพฺยากตํ น-อทุกฺขมสุขาย เวทนาย สมฺปยุตฺตญฺจ ธมฺมํ ปฏิจฺจ กุสโล อทุกฺขมสุขาย เวทนาย สมฺปยุตฺโต ธมฺโม อุปฺปชฺชติ เหตุปจฺจยา.\csromanspacer{} เอกํ.}

\prose{นกุสลํ น-อทุกฺขมสุขาย เวทนาย สมฺปยุตฺตญฺจ น-อกุสลํ น-อทุกฺขมสุขาย เวทนาย สมฺปยุตฺตญฺจ ธมฺมํ ปฏิจฺจ อพฺยากโต อทุกฺขมสุขาย เวทนาย สมฺปยุตฺโต ธมฺโม อุปฺปชฺชติ เหตุปจฺจยา.\csromanspacer{} เอกํ.\csromanspacer{}\csromanspacer{}เหตุยา นว.}

\csromanheader{2}{1. กุสลตฺติก 2. วิปากตฺติก}
\proseitem{4}{นกุสลํ นวิปากํ ธมฺมํ ปฏิจฺจ อพฺยากโต วิปาโก ธมฺโม อุปฺปชฺชติ เหตุปจฺจยา.\csromanspacer{} เอกํ.}

\prose{น-อกุสลํ นวิปากํ ธมฺมํ ปฏิจฺจ อพฺยากโต วิปาโก ธมฺโม อุปฺปชฺชติ เหตุปจฺจยา.\csromanspacer{} เอกํ.}

\prose{นกุสลํ นวิปากญฺจ น-อกุสลํ นวิปากญฺจ ธมฺมํ ปฏิจฺจ อพฺยากโต วิปาโก ธมฺโม อุปฺปชฺชติ เหตุปจฺจยา.\csromanspacer{} เอกํ.\csromanspacer{}\csromanspacer{}เหตุยา ตีณิ.}

\proseitem{5}{นกุสลํ นวิปากธมฺมธมฺมํ ปจฺจยา กุสโล วิปากธมฺมธมฺโม อุปฺปชฺชติ เหตุปจฺจยา.\csromanspacer{}\csromanspacer{}นกุสลํ นวิปากธมฺมธมฺมํ ปจฺจยา อกุสโล วิปากธมฺมธมฺโม อุปฺปชฺชติ เหตุปจฺจยา.\csromanspacer{} เทฺว.}

\prose{น-อกุสลํ นวิปากธมฺมธมฺมํ ปจฺจยา อกุสโล วิปากธมฺมธมฺโม อุปฺปชฺชติ เหตุปจฺจยา.\csromanspacer{} เทฺว.}

\prose{นกุสลํ นวิปากธมฺมธมฺมญฺจ น-อกุสลํ นวิปากธมฺมธมฺมญฺจ ธมฺมํ ปจฺจยา กุสโล วิปากธมฺมธมฺโม อุปฺปชฺชติ เหตุปจฺจยา.\csromanspacer{} เทฺว. (สํขิตฺตํ.) .\csromanspacer{} เหตุยา ฉ ปญฺหา.}

\proseitem{6}{นกุสลํ นเนววิปากนวิปากธมฺมธมฺมํ ปฏิจฺจ อพฺยากโต เนววิปากนวิปากธมฺมธมฺโม อุปฺปชฺชติ เหตุปจฺจยา.\csromanspacer{}\csromanspacer{}เหตุยา ฉ.}

\csromanheader{2}{1. กุสลตฺติก 3. อุปาทินฺนตฺติก}
\proseitem{7}{นกุสโล น-อุปาทินฺนุปาทานิโย ธมฺโม อพฺยากตสฺส อุปาทินฺนุปาทานิยสฺส ธมฺมสฺส อารมฺมณปจฺจเยน ปจฺจโย.\csromanspacer{} ฉ ปญฺหา.}

\csromanheader{2}{ธมฺมปจฺจนียานุโลม ติกติกปฏฺฐานปาฬิ}
\prose{\csromanfolio{406}นกุสลํ น-อนุปาทินฺนุปาทานิยํ ธมฺมํ ปฏิจฺจ อพฺยากโต อนุปาทินฺนุปาทานิโย ธมฺโม อุปฺปชฺชติ เหตุปจฺจยา.\csromanspacer{}\csromanspacer{}เหตุยา ปญฺจ.}

\prose{นกุสลํ น-อนุปาทินฺน-อนุปาทานิยํ ธมฺมํ ปจฺจยา กุสโล อนุปาทินฺน-อนุปาทานิโย ธมฺโม อุปฺปชฺชติ เหตุปจฺจยา.\csromanspacer{}\csromanspacer{}เหตุยา ฉ.}

\csromanheader{2}{1. กุสลตฺติก 4. สํกิลิฏฺฐตฺติก}
\proseitem{8}{นกุสลํ นสํกิลิฏฺฐสํกิเลสิกํ ธมฺมํ ปจฺจยา อกุสโล สํกิลิฏฺฐสํกิเลสิโก ธมฺโม อุปฺปชฺชติ เหตุปจฺจยา. (อกุสลาเนว ตีณิ.)}

\prose{นกุสลํ น-อสํกิลิฏฺฐสํกิเลสิกํ ธมฺมํ ปฏิจฺจ อพฺยากโต อสํกิลิฏฺฐ-อสํกิเลสิโก ธมฺโม อุปฺปชฺชติ เหตุปจฺจยา.\csromanspacer{}\csromanspacer{}เหตุยา ฉ.}

\prose{นกุสลํ น-อสํกิลิฏฺฐ-อสํกิเลสิกํ ธมฺมํ ปจฺจยา กุสโล อสํกิลิฏฺฐ-อสํกิเลสิโก ธมฺโม อุปฺปชฺชติ เหตุปจฺจยา.\csromanspacer{}\csromanspacer{}เหตุยา ฉ.}

\csromanheader{2}{1. กุสลตฺติก 5. วิตกฺกตฺติก}
\proseitem{9}{นกุสลํ นสวิตกฺกสวิจารํ ธมฺมํ ปฏิจฺจ อกุสโล สวิตกฺกสวิจาโร ธมฺโม อุปฺปชฺชติ เหตุปจฺจยา.\csromanspacer{}\csromanspacer{}เหตุยา นว.}

\prose{นกุสลํ น-อวิตกฺกวิจารมตฺตํ ธมฺมํ ปฏิจฺจ อกุสโล อวิตกฺกวิจารมตฺโต ธมฺโม อุปฺปชฺชติ เหตุปจฺจยา.\csromanspacer{}\csromanspacer{}เหตุยา นว.}

\prose{นกุสลํ น-อวิตกฺก-อวิจารํ ธมฺมํ ปฏิจฺจ อพฺยากโต อวิตกฺก-อวิจาโร ธมฺโม อุปฺปชฺชติ เหตุปจฺจยา.\csromanspacer{}\csromanspacer{}เหตุยา ทฺวาทส.}

\csromanheader{2}{1. กุสลตฺติก 6. ปีติตฺติก}
\proseitem{10}{นกุสลํ นปีติสหคตํ ธมฺมํ ปฏิจฺจ อกุสโล ปีติสหคโต ธมฺโม อุปฺปชฺชติ เหตุปจฺจยา.\csromanspacer{}\csromanspacer{}เหตุยา นว.}

\prose{นกุสลํ นสุขสหคตํ ธมฺมํ ปฏิจฺจ อกุสโล สุขสหคโต ธมฺโม อุปฺปชฺชติ เหตุปจฺจยา.\csromanspacer{}\csromanspacer{}เหตุยา นว.}

\prose{นกุสลํ น-อุเปกฺขาสหคตํ ธมฺมํ ปฏิจฺจ อกุสโล อุเปกฺขาสหคโต ธมฺโม อุปฺปชฺชติ เหตุปจฺจยา.\csromanspacer{}\csromanspacer{}เหตุยา นว.}

\csromanheader{2}{1. กุสลตฺติก 9. อาจยคามิตฺติก \textbf{1. กุสลตฺติก 7. ทสฺสนตฺติก}}
\proseitem{11}{\csromanfolio{407}นกุสลํ นทสฺสเนน ปหาตพฺพํ ธมฺมํ ปจฺจยา อกุสโล ทสฺสเนน ปหาตพฺโพ ธมฺโม อุปฺปชฺชติ เหตุปจฺจยา.\csromanspacer{}\csromanspacer{}เหตุยา ตีณิ.}

\prose{นกุสลํ นภาวนาย ปหาตพฺพํ ธมฺมํ ปจฺจยา อกุสโล ภาวนาย ปหาตพฺโพ ธมฺโม อุปฺปชฺชติ เหตุปจฺจยา.\csromanspacer{}\csromanspacer{}เหตุยา ตีณิ.}

\prose{นกุสลํ นเนวทสฺสเนน นภาวนาย ปหาตพฺพํ ธมฺมํ ปฏิจฺจ อพฺยากโต เนวทสฺสเนน นภาวนาย ปหาตพฺโพ ธมฺโม อุปฺปชฺชติ เหตุปจฺจยา.\csromanspacer{}\csromanspacer{}เหตุยา ตีณิ.}

\csromanheader{2}{1. กุสลตฺติก 8. ทสฺสนเหตุตฺติก}
\proseitem{12}{นกุสลํ นทสฺสเนน ปหาตพฺพเหตุกํ ธมฺมํ ปฏิจฺจ อกุสโล ทสฺสเนน ปหาตพฺพเหตุโก ธมฺโม อุปฺปชฺชติ เหตุปจฺจยา.\csromanspacer{}\csromanspacer{}เหตุยา ตีณิ.}

\prose{นกุสลํ นภาวนาย ปหาตพฺพเหตุกํ ธมฺมํ ปฏิจฺจ อกุสโล ภาวนาย ปหาตพฺพเหตุโก ธมฺโม อุปฺปชฺชติ เหตุปจฺจยา.\csromanspacer{}\csromanspacer{}เหตุยา ตีณิ.}

\prose{นกุสลํ นเนวทสฺสเนน นภาวนาย ปหาตพฺพเหตุกํ ธมฺมํ ปฏิจฺจ อพฺยากโต เนวทสฺสเนน นภาวนาย ปหาตพฺพเหตุโก ธมฺโม อุปฺปชฺชติ เหตุปจฺจยา.\csromanspacer{}\csromanspacer{}เหตุยา ตีณิ.}

\csromanheader{2}{1. กุสลตฺติก 9. อาจยคามิตฺติก}
\proseitem{13}{นกุสลํ น-อาจยคามิํ ธมฺมํ ปจฺจยา กุสโล อาจยคามี ธมฺโม อุปฺปชฺชติ เหตุปจฺจยา.\csromanspacer{}\csromanspacer{}เหตุยา ฉ.}

\prose{นกุสลํ น-อปจยคามิํ ธมฺมํ ปจฺจยา กุสโล อปจยคามี ธมฺโม อุปฺปชฺชติ เหตุปจฺจยา.\csromanspacer{}\csromanspacer{}เหตุยา ตีณิ.}

\prose{นกุสลํ นเนวาจยคามินาปจยคามิํ ธมฺมํ ปฏิจฺจ อพฺยากโต เนวาจยคามินาปจยคามี ธมฺโม อุปฺปชฺชติ เหตุปจฺจยา.\csromanspacer{}\csromanspacer{}เหตุยา ปญฺจ.}

\csromanheader{2}{ธมฺมปจฺจนียานุโลม ติกติกปฏฺฐานปาฬิ \textbf{1. กุสลตฺติก 10. เสกฺขตฺติก}}
\proseitem{14}{\csromanfolio{408}นกุสลํ นเสกฺขํ ธมฺมํ ปจฺจยา กุสโล เสกฺโข ธมฺโม อุปฺปชฺชติ เหตุปจฺจยา.\csromanspacer{}\csromanspacer{}เหตุยา ฉ.}

\prose{นกุสลํ น-อเสกฺขํ ธมฺมํ ปจฺจยา อพฺยากโต อเสกฺโข ธมฺโม อุปฺปชฺชติ เหตุปจฺจยา.\csromanspacer{}\csromanspacer{}เหตุยา ตีณิ.}

\prose{นกุสลํ นเนวเสกฺขนาเสกฺขํ ธมฺมํ ปฏิจฺจ อพฺยากโต เนวเสกฺขนาเสกฺโข ธมฺโม อุปฺปชฺชติ เหตุปจฺจยา.\csromanspacer{}\csromanspacer{}เหตุยา ปญฺจ.}

\csromanheader{2}{1. กุสลตฺติก 11. ปริตฺตตฺติก}
\proseitem{15}{นกุสลํ นปริตฺตํ ธมฺมํ ปฏิจฺจ อพฺยากโต ปริตฺโต ธมฺโม อุปฺปชฺชติ เหตุปจฺจยา.\csromanspacer{}\csromanspacer{}เหตุยา ปญฺจ.}

\prose{นกุสลํ นมหคฺคตํ ธมฺมํ ปฏิจฺจ อพฺยากโต มหคฺคโต ธมฺโม อุปฺปชฺชติ เหตุปจฺจยา.\csromanspacer{}\csromanspacer{}เหตุยา ตีณิ.}

\prose{นกุสลํ น-อปฺปมาณํ ธมฺมํ ปจฺจยา กุสโล อปฺปมาโณ ธมฺโม อุปฺปชฺชติ เหตุปจฺจยา.\csromanspacer{}\csromanspacer{}เหตุยา ฉ.}

\csromanheader{2}{1. กุสลตฺติก 12. ปริตฺตารมฺมณตฺติก}
\proseitem{16}{นกุสลํ นปริตฺตารมฺมณํ ธมฺมํ ปฏิจฺจ อพฺยากโต ปริตฺตารมฺมโณ ธมฺโม อุปฺปชฺชติ เหตุปจฺจยา.\csromanspacer{}\csromanspacer{}เหตุยา ตีณิ.}

\prose{นกุสลํ นมหคฺคตารมฺมณํ ธมฺมํ ปจฺจยา กุสโล มหคฺคตารมฺมโณ ธมฺโม อุปฺปชฺชติ เหตุปจฺจยา.\csromanspacer{}\csromanspacer{}เหตุยา นว.}

\prose{นกุสลํ น-อปฺปมาณารมฺมณํ ธมฺมํ ปจฺจยา กุสโล อปฺปมาณารมฺมโณ ธมฺโม อุปฺปชฺชติ เหตุปจฺจยา.\csromanspacer{}\csromanspacer{}เหตุยา ฉ.}

\csromanheader{2}{1. กุสลตฺติก 13. หีนตฺติก}
\proseitem{17}{นกุสลํ นหีนํ ธมฺมํ ปจฺจยา อกุสโล หีโน ธมฺโม อุปฺปชฺชติ เหตุปจฺจยา.\csromanspacer{}\csromanspacer{}เหตุยา ตีณิ.}

\prose{นกุสลํ นมชฺฌิมํ ธมฺมํ ปฏิจฺจ อพฺยากโต มชฺฌิโม ธมฺโม อุปฺปชฺชติ เหตุปจฺจยา.\csromanspacer{}\csromanspacer{}เหตุยา ฉ.}

\csromanheader{2}{1. กุสลตฺติก 17-18. อตีตตฺติกทฺวย}
\prose{\csromanfolio{409}นกุสลํ นปณีตํ ธมฺมํ ปจฺจยา กุสโล ปณีโต ธมฺโม อุปฺปชฺชติ เหตุปจฺจยา.\csromanspacer{}\csromanspacer{}เหตุยา ฉ.}

\csromanheader{2}{1. กุสลตฺติก 14. มิจฺฉตฺตนิยตตฺติก}
\proseitem{18}{นกุสลํ นมิจฺฉตฺตนิยตํ ธมฺมํ ปจฺจยา อกุสโล มิจฺฉตฺตนิยโต ธมฺโม อุปฺปชฺชติ เหตุปจฺจยา.\csromanspacer{}\csromanspacer{}เหตุยา ตีณิ.}

\prose{นกุสลํ นสมฺมตฺตนิยตํ ธมฺมํ ปจฺจยา กุสโล สมฺมตฺตนิยโต ธมฺโม อุปฺปชฺชติ เหตุปจฺจยา.\csromanspacer{}\csromanspacer{}เหตุยา ตีณิ.}

\prose{นกุสลํ น-อนิยตํ ธมฺมํ ปฏิจฺจ อพฺยากโต อนิยโต ธมฺโม อุปฺปชฺชติ เหตุปจฺจยา.\csromanspacer{}\csromanspacer{}เหตุยา ปญฺจ.}

\csromanheader{2}{1. กุสลตฺติก 15. มคฺคารมฺมณตฺติก}
\proseitem{19}{นกุสลํ นมคฺคารมฺมณํ ธมฺมํ ปจฺจยา กุสโล มคฺคารมฺมโณ ธมฺโม อุปฺปชฺชติ เหตุปจฺจยา.\csromanspacer{}\csromanspacer{}เหตุยา ฉ.}

\prose{นกุสลํ นมคฺคเหตุกํ ธมฺมํ ปจฺจยา กุสโล มคฺคเหตุโก ธมฺโม อุปฺปชฺชติ เหตุปจฺจยา.\csromanspacer{}\csromanspacer{}เหตุยา ตีณิ.}

\prose{นกุสลํ นมคฺคาธิปติํ ธมฺมํ ปจฺจยา กุสโล มคฺคาธิปติ ธมฺโม อุปฺปชฺชติ เหตุปจฺจยา.\csromanspacer{}\csromanspacer{}เหตุยา ฉ.}

\csromanheader{2}{1. กุสลตฺติก 16. อุปฺปนฺนตฺติก}
\proseitem{20}{นกุสโล น-อุปฺปนฺโน ธมฺโม กุสลสฺส อุปฺปนฺนสฺส ธมฺมสฺส อารมฺมณปจฺจเยน ปจฺจโย.\csromanspacer{}\csromanspacer{}อารมฺมเณ อฏฺฐารส.}

\csromanheader{2}{1. กุสลตฺติก 17-18. อตีตตฺติกทฺวย}
\proseitem{21}{นกุสโล นปจฺจุปฺปนฺโน ธมฺโม กุสลสฺส ปจฺจุปฺปนฺนสฺส ธมฺมสฺส อารมฺมณปจฺจเยน ปจฺจโย.\csromanspacer{}\csromanspacer{}อารมฺมเณ อฏฺฐารส.}

\proseitem{22}{นกุสลํ น-อตีตารมฺมณํ ธมฺมํ ปฏิจฺจ อพฺยากโต อตีตารมฺมโณ ธมฺโม อุปฺปชฺชติ เหตุปจฺจยา.\csromanspacer{}\csromanspacer{}เหตุยา ตีณิ.}

\csromanheader{2}{ธมฺมปจฺจนียานุโลม ติกติกปฏฺฐานปาฬิ}
\prose{\csromanfolio{410}นกุสลํ น-อนาคตารมฺมณํ ธมฺมํ ปจฺจยา กุสโล อนาคตารมฺมโณ ธมฺโม อุปฺปชฺชติ เหตุปจฺจยา.\csromanspacer{}\csromanspacer{}เหตุยา นว.}

\prose{นกุสลํ นปจฺจุปฺปนฺนารมฺมณํ ธมฺมํ ปฏิจฺจ อพฺยากโต ปจฺจุปฺปนฺนารมฺมโณ ธมฺโม อุปฺปชฺชติ เหตุปจฺจยา.\csromanspacer{}\csromanspacer{}เหตุยา ตีณิ.}

\csromanheader{2}{1. กุสลตฺติก 19-20. อชฺฌตฺตตฺติกทฺวย}
\proseitem{23}{นกุสโล น-อชฺฌตฺโต ธมฺโม อชฺฌตฺตสฺส กุสลสฺส ธมฺมสฺส อารมฺมณปจฺจเยน ปจฺจโย. (สํขิตฺตํ.) .\csromanspacer{} อารมฺมเณ อฏฺฐารส, อธิปติยา โสฬส, อุปนิสฺสเย อฏฺฐารส, ปุเรชาเต อตฺถิยา อวิคเต นว.}

\prose{นกุสโล นพหิทฺธา ธมฺโม พหิทฺธา กุสลสฺส ธมฺมสฺส อารมฺมณปจฺจเยน ปจฺจโย. (สํขิตฺตํ.) .\csromanspacer{} อารมฺมเณ อฏฺฐารส, อธิปติยา ฉ, อุปนิสฺสเย อฏฺฐารส, ปุเรชาเต อตฺถิยา อวิคเต นว.}

\proseitem{24}{นกุสลํ น-อชฺฌตฺตารมฺมณํ ธมฺมํ ปฏิจฺจ อพฺยากโต อชฺฌตฺตารมฺมโณ ธมฺโม อุปฺปชฺชติ เหตุปจฺจยา.\csromanspacer{}\csromanspacer{}เหตุยา ตีณิ.}

\prose{นกุสลํ นพหิทฺธารมฺมณํ ธมฺมํ ปฏิจฺจ อพฺยากโต พหิทฺธารมฺมโณ ธมฺโม อุปฺปชฺชติ เหตุปจฺจยา.\csromanspacer{}\csromanspacer{}เหตุยา ตีณิ.}

\csromanheader{2}{1. กุสลตฺติก 21. สนิทสฺสนตฺติก}
\proseitem{25}{นกุสลํ นสนิทสฺสนสปฺปฏิฆํ ธมฺมํ ปฏิจฺจ อพฺยากโต สนิทสฺสนสปฺปฏิโฆ ธมฺโม อุปฺปชฺชติ เหตุปจฺจยา.\csromanspacer{} เอกํ.}

\prose{น-อกุสลํ นสนิทสฺสนสปฺปฏิฆํ ธมฺมํ ปฏิจฺจ อพฺยากโต สนิทสฺสนสปฺปฏิโฆ ธมฺโม อุปฺปชฺชติ เหตุปจฺจยา.\csromanspacer{} เอกํ.}

\prose{น-อพฺยากตํ นสนิทสฺสนสปฺปฏิฆํ ธมฺมํ ปฏิจฺจ อพฺยากโต สนิทสฺสนสปฺปฏิโฆ ธมฺโม อุปฺปชฺชติ เหตุปจฺจยา.\csromanspacer{} เอกํ.}

\prose{นกุสลํ นสนิทสฺสนสปฺปฏิฆญฺจ น-อพฺยากตํ นสนิทสฺสนสปฺปฏิฆญฺจ ธมฺมํ ปฏิจฺจ อพฺยากโต สนิทสฺสนสปฺปฏิโฆ ธมฺโม อุปฺปชฺชติ เหตุปจฺจยา.\csromanspacer{} เอกํ.}

\csromanheader{2}{2. เวทนาตฺติก 1. กุสลตฺติก}
\prose{\csromanfolio{411}น-อกุสลํ นสนิทสฺสนสปฺปฏิฆญฺจ น-อพฺยากตํ นสนิทสฺสนสปฺปฏิฆญฺจ ธมฺมํ ปฏิจฺจ อพฺยากโต สนิทสฺสนสปฺปฏิโฆ ธมฺโม อุปฺปชฺชติ เหตุปจฺจยา.\csromanspacer{} เอกํ.}

\prose{นกุสลํ นสนิทสฺสนสปฺปฏิฆญฺจ น-อกุสลํ นสนิทสฺสนสปฺปฏิฆญฺจ ธมฺมํ ปฏิจฺจ อพฺยากโต สนิทสฺสนสปฺปฏิโฆ ธมฺโม อุปฺปชฺชติ เหตุปจฺจยา.\csromanspacer{} เอกํ.\csromanspacer{}\csromanspacer{}เหตุยา ฉ.}

\proseitem{26}{นกุสลํ น-อนิทสฺสนสปฺปฏิฆํ ธมฺมํ ปฏิจฺจ อพฺยากโต อนิทสฺสนสปฺปฏิโฆ ธมฺโม อุปฺปชฺชติ เหตุปจฺจยา.\csromanspacer{}\csromanspacer{}เหตุยา ฉ.}

\prose{นกุสลํ น-อนิทสฺสน-อปฺปฏิฆํ ธมฺมํ ปฏิจฺจ อพฺยากโต อนิทสฺสนสปฺปฏิโฆ ธมฺโม อุปฺปชฺชติ เหตุปจฺจยา.\csromanspacer{}\csromanspacer{}เหตุยา ตีณิ.}

\csromanheader{2}{2. เวทนาตฺติก 1. กุสลตฺติก}
\proseitem{27}{นสุขาย เวทนาย สมฺปยุตฺตํ นกุสลํ ธมฺมํ ปจฺจยา สุขาย เวทนาย สมฺปยุตฺโต กุสโล ธมฺโม อุปฺปชฺชติ เหตุปจฺจยา.\csromanspacer{}\csromanspacer{}นสุขาย เวทนาย สมฺปยุตฺตํ นกุสลํ ธมฺมํ ปจฺจยา อทุกฺขมสุขาย เวทนาย สมฺปยุตฺโต กุสโล ธมฺโม อุปฺปชฺชติ เหตุปจฺจยา.\csromanspacer{} เทฺว.}

\prose{นทุกฺขาย เวทนาย สมฺปยุตฺตํ นกุสลํ ธมฺมํ ปจฺจยา สุขาย เวทนาย สมฺปยุตฺโต กุสโล ธมฺโม อุปฺปชฺชติ เหตุปจฺจยา.\csromanspacer{}\csromanspacer{}นทุกฺขาย เวทนาย สมฺปยุตฺตํ นกุสลํ ธมฺมํ ปจฺจยา อทุกฺขมสุขาย เวทนาย สมฺปยุตฺโต กุสโล ธมฺโม อุปฺปชฺชติ เหตุปจฺจยา.\csromanspacer{} เทฺว.}

\prose{น-อทุกฺขมสุขาย เวทนาย สมฺปยุตฺตํ นกุสลํ ธมฺมํ ปจฺจยา อทุกฺขมสุขาย เวทนาย สมฺปยุตฺโต กุสโล ธมฺโม อุปฺปชฺชติ เหตุปจฺจยา.\csromanspacer{}\csromanspacer{}น-อทุกฺขมสุขาย เวทนาย สมฺปยุตฺตํ นกุสลํ ธมฺมํ ปจฺจยา สุขาย เวทนาย สมฺปยุตฺโต กุสโล ธมฺโม อุปฺปชฺชติ เหตุปจฺจยา.\csromanspacer{} เทฺว. (จตฺตาริ คณิตเกน เทฺว เทฺว ปญฺหา กาตพฺพา.) .\csromanspacer{} เหตุยา จุทฺทส.}

\proseitem{28}{นสุขาย เวทนาย สมฺปยุตฺตํ น-อกุสลํ ธมฺมํ ปจฺจยา สุขาย เวทนาย สมฺปยุตฺโต อกุสโล ธมฺโม อุปฺปชฺชติ เหตุปจฺจยา.\csromanspacer{}\csromanspacer{}นสุขาย เวทนาย สมฺปยุตฺตํ น-อกุสลํ ธมฺมํ ปจฺจยา ทุกฺขาย เวทนาย สมฺปยุตฺโต อกุสโล ธมฺโม อุปฺปชฺชติ เหตุปจฺจยา.\csromanspacer{}\csromanspacer{}เหตุยา เอกวีส.}

\csromanheader{2}{ธมฺมปจฺจนียานุโลม ติกติกปฏฺฐานปาฬิ}
\proseitem{29}{\csromanfolio{412}นสุขาย เวทนาย สมฺปยุตฺโต น-อพฺยากโต ธมฺโม สุขาย เวทนาย สมฺปยุตฺตสฺส อพฺยากตสฺส ธมฺมสฺส อารมฺมณปจฺจเยน ปจฺจโย.\csromanspacer{} อารมฺมเณ จุทฺทส.}

\csromanheader{2}{3. วิปากตฺติก 1. กุสลตฺติก}
\proseitem{30}{นวิปากํ นกุสลํ ธมฺมํ ปจฺจยา วิปากธมฺมธมฺโม กุสโล ธมฺโม อุปฺปชฺชติ เหตุปจฺจยา.\csromanspacer{}\csromanspacer{}เหตุยา ตีณิ.}

\prose{นวิปากํ น-อกุสลํ ธมฺมํ ปจฺจยา วิปากธมฺมธมฺโม อกุสโล ธมฺโม อุปฺปชฺชติ เหตุปจฺจยา.\csromanspacer{}\csromanspacer{}เหตุยา ตีณิ.}

\prose{นวิปากํ น-อพฺยากตํ ธมฺมํ ปฏิจฺจ เนววิปากนวิปากธมฺมธมฺโม อพฺยากโต ธมฺโม อุปฺปชฺชติ เหตุปจฺจยา.\csromanspacer{}\csromanspacer{}เหตุยา ตีณิ.}

\csromanheader{2}{4. อุปาทินฺนตฺติก 1. กุสลตฺติก}
\proseitem{31}{น-อนุปาทินฺนุปาทานิยํ นกุสลํ ธมฺมํ ปจฺจยา อนุปาทินฺนุปาทานิโย กุสโล ธมฺโม อุปฺปชฺชติ เหตุปจฺจยา.\csromanspacer{}\csromanspacer{}เหตุยา ฉ.}

\prose{น-อนุปาทินฺนุปาทานิยํ น-อกุสลํ ธมฺมํ ปจฺจยา อนุปาทินฺนุปาทานิโย อกุสโล ธมฺโม อุปฺปชฺชติ เหตุปจฺจยา.\csromanspacer{}\csromanspacer{}เหตุยา ตีณิ.}

\prose{น-อุปาทินฺนุปาทานิยํ น-อพฺยากตํ ธมฺมํ ปฏิจฺจ อนุปาทินฺนุปาทานิโย อพฺยากโต ธมฺโม อุปฺปชฺชติ เหตุปจฺจยา.\csromanspacer{}\csromanspacer{}เหตุยา ปญฺจ.}

\csromanheader{2}{5. สํกิลิฏฺฐตฺติก 1. กุสลตฺติก}
\proseitem{32}{นสํกิลิฏฺฐสํกิเลสิกํ นกุสลํ ธมฺมํ ปจฺจยา อสํกิลิฏฺฐสํกิเลสิโก กุสโล ธมฺโม อุปฺปชฺชติ เหตุปจฺจยา.\csromanspacer{}\csromanspacer{}เหตุยา ฉ.}

\prose{นสํกิลิฏฺฐสํกิเลสิกํ น-อกุสลํ ธมฺมํ ปจฺจยา สํกิลิฏฺฐสํกิเลสิโก อกุสโล ธมฺโม อุปฺปชฺชติ เหตุปจฺจยา.\csromanspacer{}\csromanspacer{}เหตุยา ตีณิ.}

\prose{นสํกิลิฏฺฐสํกิเลสิกํ น-อพฺยากตํ ธมฺมํ ปฏิจฺจ อสํกิลิฏฺฐสํกิเลสิโก อพฺยากโต ธมฺโม อุปฺปชฺชติ เหตุปจฺจยา.\csromanspacer{}\csromanspacer{}เหตุยา ฉ.}

\csromanheader{2}{8. ทสฺสนตฺติก 1. กุสลตฺติก \textbf{6. วิตกฺกตฺติก 1. กุสลตฺติก}}
\proseitem{33}{\csromanfolio{413}นสวิตกฺกสวิจารํ นกุสลํ ธมฺมํ ปจฺจยา สวิตกฺกสวิจาโร กุสโล ธมฺโม อุปฺปชฺชติ เหตุปจฺจยา.\csromanspacer{}\csromanspacer{}เหตุยา ปนฺนรส.}

\prose{นสวิตกฺกสวิจารํ น-อกุสลํ ธมฺมํ ปจฺจยา สวิตกฺกสวิจาโร อกุสโล ธมฺโม อุปฺปชฺชติ เหตุปจฺจยา.\csromanspacer{}\csromanspacer{}เหตุยา นว.}

\prose{นสวิตกฺกสวิจารํ น-อพฺยากตํ ธมฺมํ ปฏิจฺจ อวิตกฺก-อวิจาโร อพฺยากโต ธมฺโม อุปฺปชฺชติ เหตุปจฺจยา.\csromanspacer{}\csromanspacer{}เหตุยา สตฺต.}

\csromanheader{2}{7. ปีติตฺติก 1. กุสลตฺติก}
\proseitem{34}{นปีติสหคตํ นกุสลํ ธมฺมํ ปจฺจยา ปีติสหคโต กุสโล ธมฺโม อุปฺปชฺชติ เหตุปจฺจยา.\csromanspacer{}\csromanspacer{}เหตุยา อฏฺฐวีส.}

\prose{นปีติสหคตํ น-อกุสลํ ธมฺมํ ปจฺจยา ปีติสหคโต อกุสโล ธมฺโม อุปฺปชฺชติ เหตุปจฺจยา.\csromanspacer{}\csromanspacer{}เหตุยา อฏฺฐวีส.}

\prose{นปีติสหคโต น-อพฺยากโต ธมฺโม ปีติสหคตสฺส อพฺยากตสฺส ธมฺมสฺส อารมฺมณปจฺจเยน ปจฺจโย. (สํขิตฺตํ.) .\csromanspacer{} อารมฺมเณ อฏฺฐวีส, อธิปติยา อนนฺตเร อฏฺฐวีส ฯเปฯ อุปนิสฺสเย อฏฺฐวีส, กมฺเม จตุวีส, นตฺถิยา วิคเต อฏฺฐวีส.}

\csromanheader{2}{8. ทสฺสนตฺติก 1. กุสลตฺติก}
\proseitem{35}{นทสฺสเนน ปหาตพฺพํ นกุสลํ ธมฺมํ ปจฺจยา เนวทสฺสเนน นภาวนาย ปหาตพฺโพ กุสโล ธมฺโม อุปฺปชฺชติ เหตุปจฺจยา.\csromanspacer{}\csromanspacer{}เหตุยา ตีณิ.}

\prose{นทสฺสเนน ปหาตพฺพํ น-อกุสลํ ธมฺมํ ปจฺจยา ทสฺสเนน ปหาตพฺโพ อกุสโล ธมฺโม อุปฺปชฺชติ เหตุปจฺจยา.\csromanspacer{}\csromanspacer{}เหตุยา ฉ.}

\prose{นทสฺสเนน ปหาตพฺพํ น-อพฺยากตํ ธมฺมํ ปฏิจฺจ เนวทสฺสเนน นภาวนาย ปหาตพฺโพ อพฺยากโต ธมฺโม อุปฺปชฺชติ เหตุปจฺจยา.\csromanspacer{}\csromanspacer{}เหตุยา ฉ.}

\csromanheader{2}{ธมฺมปจฺจนียานุโลม ติกติกปฏฺฐานปาฬิ \textbf{9. ทสฺสนเหตุตฺติก 1. กุสลตฺติก}}
\proseitem{36}{\csromanfolio{414}นทสฺสเนน ปหาตพฺพเหตุกํ นกุสลํ ธมฺมํ ปจฺจยา เนวทสฺสเนน นภาวนาย ปหาตพฺพเหตุโก กุสโล ธมฺโม อุปฺปชฺชติ เหตุปจฺจยา.\csromanspacer{}\csromanspacer{}เหตุยา ตีณิ.}

\prose{นทสฺสเนน ปหาตพฺพเหตุกํ น-อกุสลํ ธมฺมํ ปจฺจยา ทสฺสเนน ปหาตพฺพเหตุโก อกุสโล ธมฺโม อุปฺปชฺชติ เหตุปจฺจยา.\csromanspacer{}\csromanspacer{}เหตุยา ฉ.}

\prose{นเนวทสฺสเนน นภาวนาย ปหาตพฺพเหตุกํ น-อพฺยากตํ ธมฺมํ ปฏิจฺจ เนวทสฺสเนน นภาวนาย ปหาตพฺพเหตุโก อพฺยากโต ธมฺโม อุปฺปชฺชติ เหตุปจฺจยา.\csromanspacer{}\csromanspacer{}เหตุยา ฉ.}

\csromanheader{2}{10. อาจยคามิตฺติก 1. กุสลตฺติก}
\proseitem{37}{น-อาจยคามิํ นกุสลํ ธมฺมํ ปจฺจยา อาจยคามี กุสโล ธมฺโม อุปฺปชฺชติ เหตุปจฺจยา.\csromanspacer{}\csromanspacer{}เหตุยา ฉ.}

\prose{น-อาจยคามิํ น-อกุสลํ ธมฺมํ ปจฺจยา อาจยคามี อกุสโล ธมฺโม อุปฺปชฺชติ เหตุปจฺจยา.\csromanspacer{}\csromanspacer{}เหตุยา ตีณิ.}

\prose{น-อาจยคามิํ น-อพฺยากตํ ธมฺมํ ปฏิจฺจ เนวาจยคามินาปจยคามี อพฺยากโต ธมฺโม อุปฺปชฺชติ เหตุปจฺจยา.\csromanspacer{}\csromanspacer{}เหตุยา ปญฺจ.}

\csromanheader{2}{11. เสกฺขตฺติก 1. กุสลตฺติก}
\proseitem{38}{นเสกฺขํ นกุสลํ ธมฺมํ ปจฺจยา เสกฺโข กุสโล ธมฺโม อุปฺปชฺชติ เหตุปจฺจยา.\csromanspacer{}\csromanspacer{}เหตุยา ฉ.}

\prose{นเสกฺขํ น-อกุสลํ ธมฺมํ ปจฺจยา เนวเสกฺขนาเสกฺโข อกุสโล ธมฺโม อุปฺปชฺชติ เหตุปจฺจยา.\csromanspacer{}\csromanspacer{}เหตุยา ตีณิ.}

\prose{นเสกฺขํ น-อพฺยากตํ ธมฺมํ ปฏิจฺจ เนวเสกฺขนาเสกฺโข อพฺยากโต ธมฺโม อุปฺปชฺชติ เหตุปจฺจยา.\csromanspacer{}\csromanspacer{}เหตุยา ปญฺจ.}

\csromanheader{2}{12. ปริตฺตตฺติก 1. กุสลตฺติก}
\proseitem{39}{นมหคฺคตํ นกุสลํ ธมฺมํ ปจฺจยา มหคฺคโต กุสโล ธมฺโม อุปฺปชฺชติ เหตุปจฺจยา.\csromanspacer{}\csromanspacer{}เหตุยา นว.}

\csromanheader{2}{15. มิจฺฉตฺตนิยตตฺติก 1. กุสลตฺติก}
\prose{\csromanfolio{415}นมหคฺคตํ น-อกุสลํ ธมฺมํ ปจฺจยา ปริตฺโต อกุสโล ธมฺโม อุปฺปชฺชติ เหตุปจฺจยา.\csromanspacer{}\csromanspacer{}เหตุยา ตีณิ.}

\prose{นปริตฺตํ น-อพฺยากตํ ธมฺมํ ปฏิจฺจ ปริตฺโต อพฺยากโต ธมฺโม อุปฺปชฺชติ เหตุปจฺจยา.\csromanspacer{}\csromanspacer{}เหตุยา ฉ.}

\csromanheader{2}{13. ปริตฺตารมฺมณตฺติก 1. กุสลตฺติก}
\proseitem{40}{นปริตฺตารมฺมณํ นกุสลํ ธมฺมํ ปจฺจยา ปริตฺตารมฺมโณ กุสโล ธมฺโม อุปฺปชฺชติ เหตุปจฺจยา.\csromanspacer{}\csromanspacer{}เหตุยา เอกวีส.}

\prose{นปริตฺตารมฺมณํ น-อกุสลํ ธมฺมํ ปจฺจยา ปริตฺตารมฺมโณ อกุสโล ธมฺโม อุปฺปชฺชติ เหตุปจฺจยา.\csromanspacer{}\csromanspacer{}เหตุยา จุทฺทส.}

\prose{นปริตฺตารมฺมโณ น-อพฺยากโต ธมฺโม ปริตฺตารมฺมณสฺส อพฺยากตสฺส ธมฺมสฺส อารมฺมณปจฺจเยน ปจฺจโย. (สํขิตฺตํ.) .\csromanspacer{} อารมฺมเณ เอกวีส.}

\csromanheader{2}{14. หีนตฺติก 1. กุสลตฺติก}
\proseitem{41}{นหีนํ นกุสลํ ธมฺมํ ปจฺจยา มชฺฌิโม กุสโล ธมฺโม อุปฺปชฺชติ เหตุปจฺจยา.\csromanspacer{}\csromanspacer{}เหตุยา ฉ.}

\prose{นหีนํ น-อกุสลํ ธมฺมํ ปจฺจยา หีโน อกุสโล ธมฺโม อุปฺปชฺชติ เหตุปจฺจยา.\csromanspacer{}\csromanspacer{}เหตุยา ตีณิ.}

\prose{นหีนํ น-อพฺยากตํ ธมฺมํ ปฏิจฺจ มชฺฌิโม อพฺยากโต ธมฺโม อุปฺปชฺชติ เหตุปจฺจยา.\csromanspacer{}\csromanspacer{}เหตุยา ฉ.}

\csromanheader{2}{15. มิจฺฉตฺตนิยตตฺติก 1. กุสลตฺติก}
\proseitem{42}{นมิจฺฉตฺตนิยตํ นกุสลํ ธมฺมํ ปจฺจยา สมฺมตฺตนิยโต กุสโล ธมฺโม อุปฺปชฺชติ เหตุปจฺจยา.\csromanspacer{}\csromanspacer{}เหตุยา ฉ.}

\prose{นมิจฺฉตฺตนิยตํ น-อกุสลํ ธมฺมํ ปจฺจยา มิจฺฉตฺตนิยโต อกุสโล ธมฺโม อุปฺปชฺชติ เหตุปจฺจยา.\csromanspacer{}\csromanspacer{}เหตุยา ฉ.}

\prose{นมิจฺฉตฺตนิยตํ น-อพฺยากตํ ธมฺมํ ปฏิจฺจ อนิยโต อพฺยากโต ธมฺโม อุปฺปชฺชติ เหตุปจฺจยา.\csromanspacer{}\csromanspacer{}เหตุยา ฉ.}

\csromanheader{2}{ธมฺมปจฺจนียานุโลม ติกติกปฏฺฐานปาฬิ \textbf{16. มคฺคารมฺมณตฺติก 1. กุสลตฺติก}}
\proseitem{43}{\csromanfolio{416}นมคฺคารมฺมณํ นกุสลํ ธมฺมํ ปจฺจยา มคฺคารมฺมโณ กุสโล ธมฺโม อุปฺปชฺชติ เหตุปจฺจยา.\csromanspacer{}\csromanspacer{}เหตุยา ปญฺจติํส. (น-อกุสลํ น-อพฺยากตํ นตฺถิ.)}

\csromanheader{2}{17. อุปฺปนฺนตฺติก 1. กุสลตฺติก}
\proseitem{44}{น-อุปฺปนฺโน นกุสโล ธมฺโม อุปฺปนฺนสฺส กุสลสฺส ธมฺมสฺส อารมฺมณปจฺจเยน ปจฺจโย.\csromanspacer{}\csromanspacer{}อารมฺมเณ สตฺต.}

\prose{น-อุปฺปนฺโน น-อกุสโล ธมฺโม อุปฺปนฺนสฺส อกุสลสฺส ธมฺมสฺส อารมฺมณปจฺจเยน ปจฺจโย.\csromanspacer{}\csromanspacer{}อารมฺมเณ ฉ.}

\prose{น-อุปฺปนฺโน น-อพฺยากโต ธมฺโม อุปฺปนฺนสฺส อพฺยากตสฺส ธมฺมสฺส อารมฺมณปจฺจเยน ปจฺจโย.\csromanspacer{}\csromanspacer{}อารมฺมเณ สตฺต. (อตีตตฺติกํ อุปฺปนฺนตฺติกสทิสํ.)}

\csromanheader{2}{19. อตีตารมฺมณตฺติก 1. กุสลตฺติก}
\proseitem{45}{น-อตีตารมฺมณํ นกุสลํ ธมฺมํ ปจฺจยา อตีตารมฺมโณ กุสโล ธมฺโม อุปฺปชฺชติ เหตุปจฺจยา.\csromanspacer{}\csromanspacer{}เหตุยา เอกวีส.}

\prose{น-อตีตารมฺมณํ น-อกุสลํ ธมฺมํ ปจฺจยา อตีตารมฺมโณ อกุสโล ธมฺโม อุปฺปชฺชติ เหตุปจฺจยา.\csromanspacer{}\csromanspacer{}เหตุยา เอกวีส.}

\csromanheader{2}{20. อชฺฌตฺตตฺติก 1. กุสลตฺติก}
\proseitem{46}{น-อชฺฌตฺตํ นกุสลํ ธมฺมํ ปจฺจยา พหิทฺธา กุสโล ธมฺโม อุปฺปชฺชติ เหตุปจฺจยา.\csromanspacer{}\csromanspacer{}นพหิทฺธา นกุสลํ ธมฺมํ ปจฺจยา อชฺฌตฺโต กุสโล ธมฺโม อุปฺปชฺชติ เหตุปจฺจยา.\csromanspacer{}\csromanspacer{}เหตุยา เทฺว.}

\prose{น-อชฺฌตฺตํ น-อกุสลํ ธมฺมํ ปจฺจยา พหิทฺธา อกุสโล ธมฺโม อุปฺปชฺชติ เหตุปจฺจยา.\csromanspacer{} นพหิทฺธา น-อกุสลํ ธมฺมํ ปจฺจยา อชฺฌตฺโต อกุสโล ธมฺโม อุปฺปชฺชติ เหตุปจฺจยา.\csromanspacer{}\csromanspacer{}เหตุยา เทฺว.}

\vspace{\tipitakaparskip}
\csromancenter{สนิทสฺสนตฺติก 2.\csromanspacer{} เวทนาตฺติก \textbf{21.}\csromanspacer{}\textbf{ อชฺฌตฺตารมฺมณตฺติก 1.}\csromanspacer{}\textbf{ กุสลตฺติก}}
\proseitem{47}{\csromanfolio{417}น-อชฺฌตฺตารมฺมณํ นกุสลํ ธมฺมํ ปจฺจยา อชฺฌตฺตารมฺมโณ กุสโล ธมฺโม อุปฺปชฺชติ เหตุปจฺจยา.\csromanspacer{}\csromanspacer{}เหตุยา ฉ.}

\prose{น-อชฺฌตฺตารมฺมณํ น-อกุสลํ ธมฺมํ ปจฺจยา อชฺฌตฺตารมฺมโณ อกุสโล ธมฺโม อุปฺปชฺชติ เหตุปจฺจยา.\csromanspacer{}\csromanspacer{}เหตุยา ฉ.}

\csromanheader{2}{22. สนิทสฺสนตฺติก 1. กุสลตฺติก}
\proseitem{48}{นสนิทสฺสนสปฺปฏิฆํ นกุสลํ ธมฺมํ ปจฺจยา อนิทสฺสน-อปฺปฏิโฆ กุสโล ธมฺโม อุปฺปชฺชติ เหตุปจฺจยา.\csromanspacer{}\csromanspacer{}น-อนิทสฺสนสปฺปฏิฆํ นกุสลํ ธมฺมํ ปจฺจยา อนิทสฺสน-อปฺปฏิโฆ กุสโล ธมฺโม อุปฺปชฺชติ เหตุปจฺจยา.\csromanspacer{}\csromanspacer{}นสนิทสฺสนสปฺปฏิฆํ นกุสลญฺจ น-อนิทสฺสนสปฺปฏิฆํ นกุสลญฺจ ธมฺมํ ปจฺจยา อนิทสฺสน-อปฺปฏิโฆ กุสโล ธมฺโม อุปฺปชฺชติ เหตุปจฺจยา.\csromanspacer{}\csromanspacer{}เหตุยา ตีณิ.}

\prose{นสนิทสฺสนสปฺปฏิฆํ น-อกุสลํ ธมฺมํ ปจฺจยา อนิทสฺสน-อปฺปฏิโฆ อกุสโล ธมฺโม อุปฺปชฺชติ เหตุปจฺจยา.\csromanspacer{}\csromanspacer{}เหตุยา ตีณิ.}

\prose{นสนิทสฺสนสปฺปฏิฆํ น-อพฺยากตํ ธมฺมํ ปฏิจฺจ สนิทสฺสนสปฺปฏิโฆ อพฺยากโต ธมฺโม อุปฺปชฺชติ เหตุปจฺจยา.\csromanspacer{} สตฺต.}

\prose{(น-อนิทสฺสนสปฺปฏิฆน-อพฺยากตมูลานิ สตฺตเมว.\csromanspacer{} ทุกมูลานิ สตฺตเมว.\csromanspacer{} สพฺพํ เอกวีสติเมว.)}

\csromanheader{2}{22. สนิทสฺสนตฺติก 2. เวทนาตฺติก}
\proseitem{49}{นสนิทสฺสนสปฺปฏิฆํ นสุขาย เวทนาย สมฺปยุตฺตํ ธมฺมํ ปฏิจฺจ อนิทสฺสน-อปฺปฏิโฆ สุขาย เวทนาย สมฺปยุตฺโต ธมฺโม อุปฺปชฺชติ เหตุปจฺจยา.\csromanspacer{}\csromanspacer{}เหตุยา ตีณิ.}

\prose{นสนิทสฺสนสปฺปฏิฆํ นทุกฺขาย เวทนาย สมฺปยุตฺตํ ธมฺมํ ปฏิจฺจ อนิทสฺสน-อปฺปฏิโฆ ทุกฺขาย เวทนาย สมฺปยุตฺโต ธมฺโม อุปฺปชฺชติ เหตุปจฺจยา.\csromanspacer{}\csromanspacer{}เหตุยา ตีณิ.}

\prose{นสนิทสฺสนสปฺปฏิฆํ น-อทุกฺขมสุขาย เวทนาย สมฺปยุตฺตํ ธมฺมํ ปฏิจฺจ อนิทสฺสน-อปฺปฏิโฆ อทุกฺขมสุขาย เวทนาย สมฺปยุตฺโต ธมฺโม อุปฺปชฺชติ เหตุปจฺจยา.\csromanspacer{}\csromanspacer{}เหตุยา ตีณิ.}

\csromanheader{2}{ธมฺมปจฺจนียานุโลม ติกติกปฏฺฐานปาฬิ \textbf{22. สนิทสฺสนตฺติก 3. วิปากตฺติก}}
\proseitem{50}{\csromanfolio{418}นสนิทสฺสนสปฺปฏิฆํ นวิปากํ ธมฺมํ ปฏิจฺจ อนิทสฺสน-อปฺปฏิโฆ วิปาโก ธมฺโม อุปฺปชฺชติ เหตุปจฺจยา.\csromanspacer{}\csromanspacer{}เหตุยา ตีณิ.}

\prose{นสนิทสฺสนสปฺปฏิฆํ นวิปากธมฺมธมฺมํ ปจฺจยา อนิทสฺสน-อปฺปฏิโฆ วิปากธมฺมธมฺโม อุปฺปชฺชติ เหตุปจฺจยา.\csromanspacer{}\csromanspacer{}เหตุยา ตีณิ.}

\prose{นสนิทสฺสนสปฺปฏิฆํ นเนววิปากนวิปากธมฺมธมฺมํ ปฏิจฺจ สนิทสฺสนสปฺปฏิโฆ เนววิปากนวิปากธมฺมธมฺโม อุปฺปชฺชติ เหตุปจฺจยา.\csromanspacer{}\csromanspacer{}เหตุยา เอกวีส.}

\csromanheader{2}{22. สนิทสฺสนตฺติก 4. อุปาทินฺนตฺติก}
\proseitem{51}{นสนิทสฺสนสปฺปฏิโฆ น-อุปาทินฺนุปาทานิโย ธมฺโม อนิทสฺสน-อปฺปฏิฆสฺส อุปาทินฺนุปาทานิยสฺส ธมฺมสฺส อารมฺมณปจฺจเยน ปจฺจโย.\csromanspacer{}\csromanspacer{}อารมฺมเณ ฉ.}

\prose{นสนิทสฺสนสปฺปฏิฆํ น-อนุปาทินฺนุปาทานิยํ ธมฺมํ ปฏิจฺจ สนิทสฺสนสปฺปฏิโฆ อนุปาทินฺนุปาทานิโย ธมฺโม อุปฺปชฺชติ เหตุปจฺจยา.\csromanspacer{}\csromanspacer{}เหตุยา เอกวีส.}

\prose{นสนิทสฺสนสปฺปฏิฆํ น-อนุปาทินฺน-อนุปาทานิยํ ธมฺมํ ปจฺจยา อนิทสฺสน-อปฺปฏิโฆ อนุปาทินฺน-อนุปาทานิโย ธมฺโม อุปฺปชฺชติ เหตุปจฺจยา.\csromanspacer{}\csromanspacer{}เหตุยา ตีณิ.}

\csromanheader{2}{22. สนิทสฺสนตฺติก 5. สํกิลิฏฺฐตฺติก}
\proseitem{52}{นสนิทสฺสนสปฺปฏิฆํ นสํกิลิฏฺฐสํกิเลสิกํ ธมฺมํ ปจฺจยา อนิทสฺสน-อปฺปฏิโฆ สํกิลิฏฺฐสํกิเลสิโก ธมฺโม อุปฺปชฺชติ เหตุปจฺจยา.\csromanspacer{}\csromanspacer{}เหตุยา ตีณิ.}

\prose{นสนิทสฺสนสปฺปฏิฆํ น-อสํกิลิฏฺฐสํกิเลสิกํ ธมฺมํ ปฏิจฺจ สนิทสฺสนสปฺปฏิโฆ อสํกิลิฏฺฐสํกิเลสิโก ธมฺโม อุปฺปชฺชติ เหตุปจฺจยา.\csromanspacer{}\csromanspacer{}เหตุยา เอกวีส.}

\prose{นสนิทสฺสนสปฺปฏิฆํ น-อสํกิลิฏฺฐ-อสํกิเลสิกํ ธมฺมํ ปจฺจยา อนิทสฺสน-อปฺปฏิโฆ อสํกิลิฏฺฐ-อสํกิเลสิโก ธมฺโม อุปฺปชฺชติ เหตุปจฺจยา.\csromanspacer{}\csromanspacer{}เหตุยา ตีณิ.}

\vspace{\tipitakaparskip}
\csromancenter{สนิทสฺสนตฺติก 8.\csromanspacer{} ทสฺสนตฺติก \textbf{22.}\csromanspacer{}\textbf{ สนิทสฺสนตฺติก 6.}\csromanspacer{}\textbf{ วิตกฺกตฺติก}}
\proseitem{53}{\csromanfolio{419}นสนิทสฺสนสปฺปฏิฆํ นสวิตกฺกสวิจารํ ธมฺมํ ปฏิจฺจ อนิทสฺสน-อปฺปฏิโฆ สวิตกฺกสวิจาโร ธมฺโม อุปฺปชฺชติ เหตุปจฺจยา.\csromanspacer{}\csromanspacer{}เหตุยา ตีณิ.}

\prose{นสนิทสฺสนสปฺปฏิฆํ น-อวิตกฺกวิจารมตฺตํ ธมฺมํ ปฏิจฺจ อนิทสฺสน-อปฺปฏิโฆ อวิตกฺกวิจารมตฺโต ธมฺโม อุปฺปชฺชติ เหตุปจฺจยา.\csromanspacer{}\csromanspacer{}เหตุยา ตีณิ.}

\prose{นสนิทสฺสนสปฺปฏิฆํ น-อวิตกฺก-อวิจารํ ธมฺมํ ปฏิจฺจ สนิทสฺสนสปฺปฏิโฆ อวิตกฺก-อวิจาโร ธมฺโม อุปฺปชฺชติ เหตุปจฺจยา.\csromanspacer{}\csromanspacer{}เหตุยา เอกวีส.}

\csromanheader{2}{22. สนิทสฺสนตฺติก 7. ปีติตฺติก}
\proseitem{54}{นสนิทสฺสนสปฺปฏิฆํ นปีติสหคตํ ธมฺมํ ปฏิจฺจ อนิทสฺสน-อปฺปฏิโฆ ปีติสหคโต ธมฺโม อุปฺปชฺชติ เหตุปจฺจยา.\csromanspacer{}\csromanspacer{}เหตุยา ตีณิ.}

\prose{นสนิทสฺสนสปฺปฏิฆํ นสุขสหคตํ ธมฺมํ ปฏิจฺจ อนิทสฺสน-อปฺปฏิโฆ สุขสหคโต ธมฺโม อุปฺปชฺชติ เหตุปจฺจยา.\csromanspacer{}\csromanspacer{}เหตุยา ตีณิ.}

\prose{นสนิทสฺสนสปฺปฏิฆํ น-อุเปกฺขาสหคตํ ธมฺมํ ปฏิจฺจ อนิทสฺสน-อปฺปฏิโฆ อุเปกฺขาสหคโต ธมฺโม อุปฺปชฺชติ เหตุปจฺจยา.\csromanspacer{}\csromanspacer{}เหตุยา ตีณิ.}

\csromanheader{2}{22. สนิทสฺสนตฺติก 8. ทสฺสนตฺติก}
\proseitem{55}{นสนิทสฺสนสปฺปฏิฆํ นทสฺสเนน ปหาตพฺพํ ธมฺมํ ปจฺจยา อนิทสฺสน-อปฺปฏิโฆ ทสฺสเนน ปหาตพฺโพ ธมฺโม อุปฺปชฺชติ เหตุปจฺจยา.\csromanspacer{}\csromanspacer{}เหตุยา ตีณิ.}

\prose{นสนิทสฺสนสปฺปฏิฆํ นภาวนาย ปหาตพฺพํ ธมฺมํ ปจฺจยา อนิทสฺสน-อปฺปฏิโฆ ภาวนาย ปหาตพฺโพ ธมฺโม อุปฺปชฺชติ เหตุปจฺจยา.\csromanspacer{}\csromanspacer{}เหตุยา ตีณิ.}

\prose{นสนิทสฺสนสปฺปฏิฆํ นเนวทสฺสเนน นภาวนาย ปหาตพฺพํ ธมฺมํ ปฏิจฺจ สนิทสฺสนสปฺปฏิโฆ เนวทสฺสเนน นภาวนาย ปหาตพฺโพ ธมฺโม อุปฺปชฺชติ เหตุปจฺจยา.\csromanspacer{}\csromanspacer{}เหตุยา เอกวีส.}

\vspace{\tipitakaparskip}
\csromancenter{ธมฺมปจฺจนียานุโลม ติกติกปฏฺฐานปาฬิ \textbf{22.}\csromanspacer{}\textbf{ สนิทสฺสนตฺติก 9.}\csromanspacer{}\textbf{ ทสฺสนเหตุตฺติก}}
\proseitem{56}{\csromanfolio{420}นสนิทสฺสนสปฺปฏิฆํ นทสฺสเนน ปหาตพฺพเหตุกํ ธมฺมํ ปฏิจฺจ อนิทสฺสน-อปฺปฏิโฆ ทสฺสเนน ปหาตพฺพเหตุโก ธมฺโม อุปฺปชฺชติ เหตุปจฺจยา.\csromanspacer{}\csromanspacer{}เหตุยา ตีณิ.}

\prose{นสนิทสฺสนสปฺปฏิฆํ นภาวนาย ปหาตพฺพเหตุกํ ธมฺมํ ปฏิจฺจ อนิทสฺสน-อปฺปฏิโฆ ภาวนาย ปหาตพฺพเหตุโก ธมฺโม อุปฺปชฺชติ เหตุปจฺจยา.\csromanspacer{}\csromanspacer{}เหตุยา ตีณิ.}

\prose{นสนิทสฺสนสปฺปฏิฆํ นเนวทสฺสเนน นภาวนาย ปหาตพฺพเหตุกํ ธมฺมํ ปฏิจฺจ สนิทสฺสนสปฺปฏิโฆ เนวทสฺสเนน นภาวนาย ปหาตพฺพเหตุโก ธมฺโม อุปฺปชฺชติ เหตุปจฺจยา.\csromanspacer{}\csromanspacer{}เหตุยา เอกวีส. \textbf{22.}\csromanspacer{}\textbf{ สนิทสฺสนตฺติก 10.}\csromanspacer{}\textbf{ อาจยคามิตฺติก}}

\proseitem{57}{นสนิทสฺสนสปฺปฏิฆํ น-อาจยคามิํ ธมฺมํ ปจฺจยา อนิทสฺสน-อปฺปฏิโฆ อาจยคามี ธมฺโม อุปฺปชฺชติ เหตุปจฺจยา.\csromanspacer{}\csromanspacer{}เหตุยา ตีณิ.}

\prose{นสนิทสฺสนสปฺปฏิฆํ น-อปจยคามิํ ธมฺมํ ปจฺจยา อนิทสฺสน-อปฺปฏิโฆ อปจยคามี ธมฺโม อุปฺปชฺชติ เหตุปจฺจยา.\csromanspacer{}\csromanspacer{}เหตุยา ตีณิ.}

\prose{นสนิทสฺสนสปฺปฏิฆํ นเนวาจยคามินาปจยคามิํ ธมฺมํ ปฏิจฺจ สนิทสฺสนสปฺปฏิโฆ เนวาจยคามินาปจยคามี ธมฺโม อุปฺปชฺชติ เหตุปจฺจยา.\csromanspacer{}\csromanspacer{}เหตุยา เอกวีส. \textbf{22.}\csromanspacer{}\textbf{ สนิทสฺสนตฺติก 11.}\csromanspacer{}\textbf{ เสกฺขตฺติก}}

\proseitem{58}{นสนิทสฺสนสปฺปฏิฆํ นเสกฺขํ ธมฺมํ ปจฺจยา อนิทสฺสน-อปฺปฏิโฆ เสกฺโข ธมฺโม อุปฺปชฺชติ เหตุปจฺจยา.\csromanspacer{}\csromanspacer{}เหตุยา ตีณิ.}

\prose{นสนิทสฺสนสปฺปฏิฆํ น-อเสกฺขํ ธมฺมํ ปจฺจยา อนิทสฺสน-อปฺปฏิโฆ อเสกฺโข ธมฺโม อุปฺปชฺชติ เหตุปจฺจยา.\csromanspacer{}\csromanspacer{}เหตุยา ตีณิ.}

\prose{นสนิทสฺสนสปฺปฏิฆํ นเนวเสกฺขนาเสกฺขํ ธมฺมํ ปฏิจฺจ สนิทสฺสนสปฺปฏิโฆ เนวเสกฺขนาเสกฺโข ธมฺโม อุปฺปชฺชติ เหตุปจฺจยา.\csromanspacer{}\csromanspacer{}เหตุยา เอกวีส.}

\vspace{\tipitakaparskip}
\csromancenter{สนิทสฺสนตฺติก 15.\csromanspacer{} มิจฺฉตฺตนิยตตฺติก \textbf{22.}\csromanspacer{}\textbf{ สนิทสฺสนตฺติก 12.}\csromanspacer{}\textbf{ ปริตฺตตฺติก}}
\proseitem{59}{\csromanfolio{421}นสนิทสฺสนสปฺปฏิฆํ นปริตฺตํ ธมฺมํ ปฏิจฺจ สนิทสฺสนสปฺปฏิโฆ ปริตฺโต ธมฺโม อุปฺปชฺชติ เหตุปจฺจยา.\csromanspacer{}\csromanspacer{}เหตุยา เอกวีส.}

\prose{นสนิทสฺสนสปฺปฏิฆํ นมหคฺคตํ ธมฺมํ ปฏิจฺจ อนิทสฺสน-อปฺปฏิโฆ มหคฺคโต ธมฺโม อุปฺปชฺชติ เหตุปจฺจยา.\csromanspacer{}\csromanspacer{}เหตุยา ตีณิ.}

\prose{นสนิทสฺสนสปฺปฏิฆํ น-อปฺปมาณํ ธมฺมํ ปจฺจยา อนิทสฺสน-อปฺปฏิโฆ อปฺปมาโณ ธมฺโม อุปฺปชฺชติ เหตุปจฺจยา.\csromanspacer{}\csromanspacer{}เหตุยา ตีณิ.}

\csromanheader{2}{22. สนิทสฺสนตฺติก 13. ปริตฺตารมฺมณตฺติก}
\proseitem{60}{นสนิทสฺสนสปฺปฏิฆํ นปริตฺตารมฺมณํ ธมฺมํ ปฏิจฺจ อนิทสฺสน-อปฺปฏิโฆ ปริตฺตารมฺมโณ ธมฺโม อุปฺปชฺชติ เหตุปจฺจยา.\csromanspacer{}\csromanspacer{}เหตุยา ตีณิ.}

\prose{นสนิทสฺสนสปฺปฏิฆํ นมหคฺคตารมฺมณํ ธมฺมํ ปจฺจยา อนิทสฺสน-อปฺปฏิโฆ มหคฺคตารมฺมโณ ธมฺโม อุปฺปชฺชติ เหตุปจฺจยา.\csromanspacer{}\csromanspacer{}เหตุยา ตีณิ.}

\prose{นสนิทสฺสนสปฺปฏิฆํ น-อปฺปมาณารมฺมณํ ธมฺมํ ปจฺจยา อนิทสฺสน-อปฺปฏิโฆ อปฺปมาณารมฺมโณ ธมฺโม อุปฺปชฺชติ เหตุปจฺจยา.\csromanspacer{}\csromanspacer{}เหตุยา ตีณิ.}

\csromanheader{2}{22. สนิทสฺสนตฺติก 14. หีนตฺติก}
\proseitem{61}{นสนิทสฺสนสปฺปฏิฆํ นหีนํ ธมฺมํ ปจฺจยา อนิทสฺสน-อปฺปฏิโฆ หีโน ธมฺโม อุปฺปชฺชติ เหตุปจฺจยา.\csromanspacer{}\csromanspacer{}เหตุยา ตีณิ.}

\prose{นสนิทสฺสนสปฺปฏิฆํ นมชฺฌิมํ ธมฺมํ ปฏิจฺจ สนิทสฺสนสปฺปฏิโฆ มชฺฌิโม ธมฺโม อุปฺปชฺชติ เหตุปจฺจยา.\csromanspacer{}\csromanspacer{}เหตุยา เอกวีส.}

\prose{นสนิทสฺสนสปฺปฏิฆํ นปณีตํ ธมฺมํ ปจฺจยา อนิทสฺสน-อปฺปฏิโฆ ปณีโต ธมฺโม อุปฺปชฺชติ เหตุปจฺจยา.\csromanspacer{}\csromanspacer{}เหตุยา ตีณิ.}

\csromanheader{2}{22. สนิทสฺสนตฺติก 15. มิจฺฉตฺตนิยตตฺติก}
\proseitem{62}{นสนิทสฺสนสปฺปฏิฆํ นมิจฺฉตฺตนิยตํ ธมฺมํ ปจฺจยา อนิทสฺสน-อปฺปฏิโฆ มิจฺฉตฺตนิยโต ธมฺโม อุปฺปชฺชติ เหตุปจฺจยา.\csromanspacer{}\csromanspacer{}เหตุยา ตีณิ.}

\csromanheader{2}{ธมฺมปจฺจนียานุโลม ติกติกปฏฺฐานปาฬิ}
\prose{\csromanfolio{422}นสนิทสฺสนสปฺปฏิฆํ นสมฺมตฺตนิยตํ ธมฺมํ ปจฺจยา อนิทสฺสน-อปฺปฏิโฆ สมฺมตฺตนิยโต ธมฺโม อุปฺปชฺชติ เหตุปจฺจยา.\csromanspacer{}\csromanspacer{}เหตุยา ตีณิ.}

\prose{นสนิทสฺสนสปฺปฏิฆํ น-อนิยตํ ธมฺมํ ปฏิจฺจ สนิทสฺสนสปฺปฏิโฆ อนิยโต ธมฺโม อุปฺปชฺชติ เหตุปจฺจยา.\csromanspacer{}\csromanspacer{}เหตุยา เอกวีส.}

\csromanheader{2}{22. สนิทสฺสนตฺติก 16. มคฺคารมฺมณตฺติก}
\proseitem{63}{นสนิทสฺสนสปฺปฏิฆํ นมคฺคารมฺมณํ ธมฺมํ ปจฺจยา อนิทสฺสน-อปฺปฏิโฆ มคฺคารมฺมโณ ธมฺโม อุปฺปชฺชติ เหตุปจฺจยา.\csromanspacer{}\csromanspacer{}เหตุยา ตีณิ.}

\prose{นสนิทสฺสนสปฺปฏิฆํ นมคฺคเหตุกํ ธมฺมํ ปจฺจยา อนิทสฺสน-อปฺปฏิโฆ มคฺคเหตุโก ธมฺโม อุปฺปชฺชติ เหตุปจฺจยา.\csromanspacer{}\csromanspacer{}เหตุยา ตีณิ.}

\prose{นสนิทสฺสนสปฺปฏิฆํ นมคฺคาธิปติํ ธมฺมํ ปจฺจยา อนิทสฺสน-อปฺปฏิโฆ มคฺคาธิปติ ธมฺโม อุปฺปชฺชติ เหตุปจฺจยา.\csromanspacer{}\csromanspacer{}เหตุยา ตีณิ.}

\csromanheader{2}{22. สนิทสฺสนตฺติก 17. อุปฺปนฺนตฺติก}
\proseitem{64}{นสนิทสฺสนสปฺปฏิโฆ น-อุปฺปนฺโน ธมฺโม อนิทสฺสน-อปฺปฏิฆสฺส อุปฺปนฺนสฺส ธมฺมสฺส อารมฺมณปจฺจเยน ปจฺจโย.\csromanspacer{}\csromanspacer{}อารมฺมเณ ฉ.}

\csromanheader{2}{22. สนิทสฺสนตฺติก 18. อตีตตฺติก}
\proseitem{65}{นสนิทสฺสนสปฺปฏิโฆ นปจฺจุปฺปนฺโน ธมฺโม อนิทสฺสน-อปฺปฏิฆสฺส ปจฺจุปฺปนฺนสฺส ธมฺมสฺส อารมฺมณปจฺจเยน ปจฺจโย.\csromanspacer{}\csromanspacer{}อารมฺมเณ ฉ.}

\csromanheader{2}{22. สนิทสฺสนตฺติก 19. อตีตารมฺมณตฺติก}
\proseitem{66}{นสนิทสฺสนสปฺปฏิฆํ น-อตีตารมฺมณํ ธมฺมํ ปฏิจฺจ อนิทสฺสน-อปฺปฏิโฆ อตีตารมฺมโณ ธมฺโม อุปฺปชฺชติ เหตุปจฺจยา.\csromanspacer{}\csromanspacer{}เหตุยา ตีณิ.}

\prose{นสนิทสฺสนสปฺปฏิฆํ น-อนาคตารมฺมณํ ธมฺมํ ปจฺจยา อนิทสฺสน-อปฺปฏิโฆ อนาคตารมฺมโณ ธมฺโม อุปฺปชฺชติ เหตุปจฺจยา.\csromanspacer{}\csromanspacer{}เหตุยา ตีณิ.}

\prose{นสนิทสฺสนสปฺปฏิฆํ นปจฺจุปฺปนฺนารมฺมณํ ธมฺมํ ปฏิจฺจ อนิทสฺสน-อปฺปฏิโฆ ปจฺจุปฺปนฺนารมฺมโณ ธมฺโม อุปฺปชฺชติ เหตุปจฺจยา.\csromanspacer{}\csromanspacer{}เหตุยา ตีณิ.}

\vspace{\tipitakaparskip}
\csromancenter{สนิทสฺสนตฺติก 20-21.\csromanspacer{} อชฺฌตฺตตฺติกทฺวย \textbf{22.}\csromanspacer{}\textbf{ สนิทสฺสนตฺติก 20-21.}\csromanspacer{}\textbf{ อชฺฌตฺตตฺติกทฺวย}}
\proseitem{67}{\csromanfolio{423}นสนิทสฺสนสปฺปฏิโฆ น-อชฺฌตฺโต ธมฺโม อนิทสฺสน-อปฺปฏิฆสฺส อชฺฌตฺตสฺส ธมฺมสฺส อารมฺมณปจฺจเยน ปจฺจโย. (สํขิตฺตํ.) .\csromanspacer{} อารมฺมเณ ฉ, อธิปติยา ฉ, ปุเรชาเต อตฺถิยา อวิคเต ฉ.}

\prose{นสนิทสฺสนสปฺปฏิโฆ นพหิทฺธา ธมฺโม อนิทสฺสน-อปฺปฏิฆสฺส พหิทฺธา ธมฺมสฺส อารมฺมณปจฺจเยน ปจฺจโย. (สํขิตฺตํ.) .\csromanspacer{} อารมฺมเณ ฉ, อธิปติยา ฉ, ปุเรชาเต อตฺถิยา อวิคเต ฉ.}

\proseitem{68}{นสนิทสฺสนสปฺปฏิฆํ น-อชฺฌตฺตารมฺมณํ ธมฺมํ ปฏิจฺจ อนิทสฺสน-อปฺปฏิโฆ อชฺฌตฺตารมฺมโณ ธมฺโม อุปฺปชฺชติ เหตุปจฺจยา.\csromanspacer{}\csromanspacer{}เหตุยา ตีณิ.}

\prose{นสนิทสฺสนสปฺปฏิฆํ นพหิทฺธารมฺมณํ ธมฺมํ ปฏิจฺจ อนิทสฺสน-อปฺปฏิโฆ พหิทฺธารมฺมโณ ธมฺโม อุปฺปชฺชติ เหตุปจฺจยา.\csromanspacer{}\csromanspacer{}เหตุยา ตีณิ ฯเปฯ อวิคเต ตีณิ.}

\prose{(สหชาตวารมฺปิ ปจฺจยวารมฺปิ นิสฺสยวารมฺปิ สํสฏฺฐวารมฺปิ สมฺปยุตฺตวารมฺปิ ปญฺหาวารมฺปิ วิตฺถาเรตพฺพํ.)}

\nitthitam{\textbf{ธมฺมปจฺจนียานุโลเม ติกติกปฏฺฐานํ นิฏฺฐิตํ.}}
\pitaka{\textbf{อภิธมฺมปิฏก}}
\csromanheader{2}{\textbf{ธมฺมปจฺจนียานุโลม}}
\gambhira{\textbf{ทุกทุกปฏฺฐานปาฬิ}}
\namakkaram{นโม ตสฺส ภควโต อรหโต สมฺมาสมฺพุทฺธสฺส.}
\csromanheader{2}{1. เหตุทุก 1. สเหตุกทุก}
\proseitem{1}{\csromanfolio{425}นเหตุํ นสเหตุกํ ธมฺมํ ปฏิจฺจ เหตุ สเหตุโก ธมฺโม อุปฺปชฺชติ เหตุปจฺจยา.\csromanspacer{}\csromanspacer{}นเหตุํ นสเหตุกํ ธมฺมํ ปฏิจฺจ นเหตุ สเหตุโก ธมฺโม อุปฺปชฺชติ เหตุปจฺจยา.\csromanspacer{}\csromanspacer{}นเหตุํ นสเหตุกํ ธมฺมํ ปฏิจฺจ เหตุ สเหตุโก จ นเหตุ สเหตุโก จ ธมฺมา อุปฺปชฺชนฺติ เหตุปจฺจยา.\csromanspacer{} ตีณิ.}

\prose{นนเหตุํ นสเหตุกํ ธมฺมํ ปฏิจฺจ นเหตุ สเหตุโก ธมฺโม อุปฺปชฺชติ เหตุปจฺจยา.\csromanspacer{} เอกํ.}

\prose{เหตุยา จตฺตาริ, อารมฺมเณ จตฺตาริ ฯเปฯ อวิคเต จตฺตาริ.}

\proseitem{2}{นนเหตุ น-อเหตุโก ธมฺโม นเหตุสฺส อเหตุกสฺส ธมฺมสฺส เหตุปจฺจเยน ปจฺจโย.}

\prose{เหตุยา เอกํ, อารมฺมเณ ฉ ฯเปฯ อวิคเต ปญฺจ. (ปญฺหาวารํ วิตฺถาเรตพฺพํ.)}

\proseitem{3}{นเหตุํ น-อเหตุกํ ธมฺมํ ปฏิจฺจ นเหตุ อเหตุโก ธมฺโม อุปฺปชฺชติ เหตุปจฺจยา.\csromanspacer{} เอกํ.}

\csromanheader{2}{ธมฺมปจฺจนียานุโลม ทุกทุกปฏฺฐานปาฬิ}
\prose{\csromanfolio{426}นนเหตุํ น-อเหตุกํ ธมฺมํ ปฏิจฺจ นเหตุ อเหตุโก ธมฺโม อุปฺปชฺชติ เหตุปจฺจยา.\csromanspacer{} เอกํ.}

\prose{นเหตุํ น-อเหตุกญฺจ นนเหตุํ น-อเหตุกญฺจ ธมฺมํ ปฏิจฺจ นเหตุ อเหตุโก ธมฺโม อุปฺปชฺชติ เหตุปจฺจยา.\csromanspacer{} เอกํ.\csromanspacer{}\csromanspacer{}เหตุยา ตีณิ.}

\csromanheader{2}{1. เหตุทุก 2-5. เหตุสมฺปยุตฺตทุกาทิ}
\proseitem{4}{นเหตุํ นเหตุสมฺปยุตฺตํ ธมฺมํ ปฏิจฺจ เหตุ เหตุสมฺปยุตฺโต ธมฺโม อุปฺปชฺชติ เหตุปจฺจยา.\csromanspacer{}\csromanspacer{}เหตุยา จตฺตาริ.}

\prose{นเหตุํ นเหตุวิปฺปยุตฺตํ ธมฺมํ ปฏิจฺจ นเหตุ เหตุวิปฺปยุตฺโต ธมฺโม อุปฺปชฺชติ เหตุปจฺจยา.\csromanspacer{}\csromanspacer{}เหตุยา ตีณิ.}

\proseitem{5}{นเหตุํ นเหตุญฺเจว น-อเหตุกญฺจ ธมฺมํ ปฏิจฺจ เหตุ เหตุ เจว สเหตุโก จ ธมฺโม อุปฺปชฺชติ เหตุปจฺจยา.\csromanspacer{}\csromanspacer{}เหตุยา เอกํ.}

\prose{นนเหตุํ น-อเหตุกญฺเจว นนเหตุญฺจ ธมฺมํ ปฏิจฺจ นเหตุ สเหตุโก เจว น จ เหตุ ธมฺโม อุปฺปชฺชติ เหตุปจฺจยา.\csromanspacer{}\csromanspacer{}เหตุยา เอกํ.}

\proseitem{6}{นเหตุํ นเหตุญฺเจว นเหตุวิปฺปยุตฺตญฺจ ธมฺมํ ปฏิจฺจ เหตุ เหตุ เจว เหตุสมฺปยุตฺโต จ ธมฺโม อุปฺปชฺชติ เหตุปจฺจยา.\csromanspacer{}\csromanspacer{}เหตุยา เอกํ.}

\prose{นนเหตุํ นเหตุวิปฺปยุตฺตญฺเจว นนเหตุญฺจ ธมฺมํ ปฏิจฺจ นเหตุ เหตุสมฺปยุตฺโต เจว น จ เหตุ ธมฺโม อุปฺปชฺชติ เหตุปจฺจยา.\csromanspacer{}\csromanspacer{}เหตุยา เอกํ.}

\proseitem{7}{นเหตุํ นเหตุํ นสเหตุกํ ธมฺมํ ปฏิจฺจ นเหตุ นเหตุ สเหตุโก ธมฺโม อุปฺปชฺชติ เหตุปจฺจยา.\csromanspacer{}\csromanspacer{}เหตุยา เอกํ.}

\prose{นเหตุํ นเหตุํ น-อเหตุกํ ธมฺมํ ปฏิจฺจ นเหตุ นเหตุ อเหตุโก ธมฺโม อุปฺปชฺชติ เหตุปจฺจยา.\csromanspacer{}\csromanspacer{}เหตุยา เอกํ.}

\csromanheader{2}{1. เหตุทุก 6-12. จูฬนฺตรทุก}
\proseitem{8}{นเหตุ นสปฺปจฺจโย ธมฺโม เหตุสฺส สปฺปจฺจยสฺส ธมฺมสฺส อารมฺมณปจฺจเยน ปจฺจโย.\csromanspacer{}\csromanspacer{}อารมฺมเณ ตีณิ. (สงฺขตํ สปฺปจฺจยสทิสํ.)}

\csromanheader{2}{1. เหตุทุก 13-18. อาสวโคจฺฉก}
\proseitem{9}{\csromanfolio{427}นเหตุํ นสนิทสฺสนํ ธมฺมํ ปฏิจฺจ นเหตุ สนิทสฺสโน ธมฺโม อุปฺปชฺชติ เหตุปจฺจยา.\csromanspacer{}\csromanspacer{}เหตุยา ตีณิ.}

\prose{นเหตุ น-อนิทสฺสโน ธมฺโม เหตุสฺส อนิทสฺสนสฺส ธมฺมสฺส อารมฺมณปจฺจเยน ปจฺจโย.\csromanspacer{}\csromanspacer{}อารมฺมเณ ตีณิ.}

\proseitem{10}{นเหตุํ นสปฺปฏิฆํ ธมฺมํ ปฏิจฺจ นเหตุ สปฺปฏิโฆ ธมฺโม อุปฺปชฺชติ เหตุปจฺจยา.\csromanspacer{}\csromanspacer{}เหตุยา ตีณิ.}

\prose{นเหตุํ น-อปฺปฏิฆํ ธมฺมํ ปฏิจฺจ นเหตุ อปฺปฏิโฆ ธมฺโม อุปฺปชฺชติ เหตุปจฺจยา.\csromanspacer{}\csromanspacer{}เหตุยา เอกํ.}

\proseitem{11}{นเหตุํ นรูปิํ ธมฺมํ ปฏิจฺจ เหตุ รูปี ธมฺโม อุปฺปชฺชติ เหตุปจฺจยา.\csromanspacer{}\csromanspacer{}เหตุยา ตีณิ.}

\prose{นเหตุํ น-อรูปิํ ธมฺมํ ปฏิจฺจ เหตุ อรูปี ธมฺโม อุปฺปชฺชติ เหตุปจฺจยา.\csromanspacer{}\csromanspacer{}เหตุยา ตีณิ.}

\proseitem{12}{นเหตุํ นโลกิยํ ธมฺมํ ปฏิจฺจ นเหตุ โลกิโย ธมฺโม อุปฺปชฺชติ เหตุปจฺจยา.\csromanspacer{}\csromanspacer{}นนเหตุํ นโลกิยํ ธมฺมํ ปฏิจฺจ นเหตุ โลกิโย ธมฺโม อุปฺปชฺชติ เหตุปจฺจยา. (คณิตเกน ตีณิ.)}

\prose{นเหตุํ นโลกุตฺตรํ ธมฺมํ ปจฺจยา เหตุ โลกุตฺตโร ธมฺโม อุปฺปชฺชติ เหตุปจฺจยา.\csromanspacer{}\csromanspacer{}เหตุยา ตีณิ.}

\proseitem{13}{นเหตุํ นเกนจิ วิญฺเญยฺยํ ธมฺมํ ปฏิจฺจ เหตุ เกนจิ วิญฺเญยฺโย ธมฺโม อุปฺปชฺชติ เหตุปจฺจยา.\csromanspacer{}\csromanspacer{}เหตุยา นว.}

\prose{นเหตุํ นเกนจิ นวิญฺเญยฺยํ ธมฺมํ ปฏิจฺจ เหตุ เกนจิ นวิญฺเญยฺโย ธมฺโม อุปฺปชฺชติ เหตุปจฺจยา.\csromanspacer{}\csromanspacer{}เหตุยา นว.}

\csromanheader{2}{1. เหตุทุก 13-18. อาสวโคจฺฉก}
\proseitem{14}{นเหตุํ โน-อาสวํ ธมฺมํ ปฏิจฺจ เหตุ อาสโว ธมฺโม อุปฺปชฺชติ เหตุปจฺจยา.\csromanspacer{}\csromanspacer{}นเหตุํ โน-อาสวํ ธมฺมํ ปฏิจฺจ นเหตุ อาสโว ธมฺโม อุปฺปชฺชติ เหตุปจฺจยา.\csromanspacer{}\csromanspacer{}นเหตุํ โน-อาสวํ ธมฺมํ ปฏิจฺจ เหตุ อาสโว จ นเหตุ อาสโว จ ธมฺมา อุปฺปชฺชนฺติ เหตุปจฺจยา.\csromanspacer{} ตีณิ.}

\csromanheader{2}{ธมฺมปจฺจนียานุโลม ทุกทุกปฏฺฐานปาฬิ}
\prose{\csromanfolio{428}นนเหตุํ โน-อาสวํ ธมฺมํ ปฏิจฺจ เหตุ อาสโว ธมฺโม อุปฺปชฺชติ เหตุปจฺจยา.\csromanspacer{} เอกํ.}

\prose{นเหตุํ โน-อาสวญฺจ นนเหตุํ โน-อาสวญฺจ ธมฺมํ ปฏิจฺจ เหตุ อาสโว ธมฺโม อุปฺปชฺชติ เหตุปจฺจยา.\csromanspacer{} เอกํ.\csromanspacer{}\csromanspacer{}เหตุยา ปญฺจ.}

\prose{นเหตุํ นโน-อาสวํ ธมฺมํ ปฏิจฺจ นเหตุ โน-อาสโว ธมฺโม อุปฺปชฺชติ เหตุปจฺจยา.\csromanspacer{} เอกํ.}

\prose{นนเหตุํ นโน-อาสวํ ธมฺมํ ปฏิจฺจ นเหตุ โน-อาสโว ธมฺโม อุปฺปชฺชติ เหตุปจฺจยา.\csromanspacer{}\csromanspacer{}นนเหตุํ นโน-อาสวํ ธมฺมํ ปฏิจฺจ เหตุ โน-อาสโว ธมฺโม อุปฺปชฺชติ เหตุปจฺจยา.\csromanspacer{}\csromanspacer{}นนเหตุํ นโน-อาสวํ ธมฺมํ ปฏิจฺจ เหตุ โน-อาสโว จ นเหตุ โน-อาสโว จ ธมฺมา อุปฺปชฺชนฺติ เหตุปจฺจยา.\csromanspacer{} ตีณิ.}

\prose{นเหตุํ นโน-อาสวญฺจ นนเหตุํ นโน-อาสวญฺจ ธมฺมํ ปฏิจฺจ นเหตุ โน-อาสโว ธมฺโม อุปฺปชฺชติ เหตุปจฺจยา.\csromanspacer{} เอกํ.\csromanspacer{}\csromanspacer{}เหตุยา ปญฺจ.}

\proseitem{15}{นเหตุํ นสาสวํ ธมฺมํ ปฏิจฺจ นเหตุ สาสโว ธมฺโม อุปฺปชฺชติ เหตุปจฺจยา.\csromanspacer{}\csromanspacer{}เหตุยา ตีณิ.}

\prose{นเหตุํ น-อนาสวํ ธมฺมํ ปจฺจยา เหตุ อนาสโว ธมฺโม อุปฺปชฺชติ เหตุปจฺจยา.\csromanspacer{}\csromanspacer{}เหตุยา ตีณิ.}

\proseitem{16}{นนเหตุํ น-อาสวสมฺปยุตฺตํ ธมฺมํ ปฏิจฺจ นเหตุ อาสวสมฺปยุตฺโต ธมฺโม อุปฺปชฺชติ เหตุปจฺจยา.\csromanspacer{}\csromanspacer{}เหตุยา ตีณิ.}

\prose{นเหตุํ น-อาสววิปฺปยุตฺตํ ธมฺมํ ปฏิจฺจ เหตุ อาสววิปฺปยุตฺโต ธมฺโม อุปฺปชฺชติ เหตุปจฺจยา.\csromanspacer{}\csromanspacer{}เหตุยา นว.}

\proseitem{17}{นเหตุํ น-อาสวญฺเจว น-อนาสวญฺจ ธมฺมํ ปฏิจฺจ เหตุ อาสโว เจว สาสโว จ ธมฺโม อุปฺปชฺชติ เหตุปจฺจยา.\csromanspacer{}\csromanspacer{}เหตุยา ปญฺจ.}

\prose{นเหตุํ น-อนาสวญฺเจว นโน จ อาสวํ ธมฺมํ ปฏิจฺจ นเหตุ สาสโว เจว โน จ อาสโว ธมฺโม อุปฺปชฺชติ เหตุปจฺจยา.\csromanspacer{}\csromanspacer{}เหตุยา ปญฺจ.}

\csromanheader{2}{1. เหตุทุก 19. สญฺโญชนาทิโคจฺฉก}
\proseitem{18}{\csromanfolio{429}นเหตุํ น-อาสวญฺเจว น-อาสววิปฺปยุตฺตญฺจ ธมฺมํ ปฏิจฺจ เหตุ อาสโว เจว อาสวสมฺปยุตฺโต จ ธมฺโม อุปฺปชฺชติ เหตุปจฺจยา.\csromanspacer{}\csromanspacer{}เหตุยา ตีณิ.}

\prose{นเหตุํ น-อาสววิปฺปยุตฺตญฺเจว นโน จ อาสวํ ธมฺมํ ปฏิจฺจ นเหตุ อาสวสมฺปยุตฺโต เจว โน จ อาสโว ธมฺโม อุปฺปชฺชติ เหตุปจฺจยา.\csromanspacer{}\csromanspacer{}เหตุยา ตีณิ.}

\proseitem{19}{นเหตุํ อาสววิปฺปยุตฺตํ นสาสวํ ธมฺมํ ปฏิจฺจ นเหตุ อาสววิปฺปยุตฺโต สาสโว ธมฺโม อุปฺปชฺชติ เหตุปจฺจยา.\csromanspacer{}\csromanspacer{}เหตุยา ตีณิ.}

\prose{นเหตุํ อาสววิปฺปยุตฺตํ น-อนาสวํ ธมฺมํ ปจฺจยา เหตุ อาสววิปฺปยุตฺโต อนาสโว ธมฺโม อุปฺปชฺชติ เหตุปจฺจยา.\csromanspacer{}\csromanspacer{}เหตุยา ตีณิ.}

\csromanheader{2}{1. เหตุทุก 19. สญฺโญชนาทิโคจฺฉก}
\proseitem{20}{นเหตุํ โนสญฺโญชนํ ธมฺมํ ปฏิจฺจ เหตุ สญฺโญชโน ธมฺโม อุปฺปชฺชติ เหตุปจฺจยา.\csromanspacer{}\csromanspacer{}เหตุยา ตีณิ.}

\proseitem{21}{นเหตุํ โนคนฺถํ ธมฺมํ ปฏิจฺจ เหตุ คนฺโถ ธมฺโม อุปฺปชฺชติ เหตุปจฺจยา.\csromanspacer{}\csromanspacer{}เหตุยา นว.}

\proseitem{22}{นเหตุํ โน-โอฆํ ธมฺมํ ปฏิจฺจ เหตุ โอโฆ ธมฺโม อุปฺปชฺชติ เหตุปจฺจยา.\csromanspacer{}\csromanspacer{}เหตุยา ปญฺจ.}

\proseitem{23}{นเหตุํ โนโยคํ ธมฺมํ ปฏิจฺจ เหตุ โยโค ธมฺโม อุปฺปชฺชติ เหตุปจฺจยา.\csromanspacer{}\csromanspacer{}เหตุยา ปญฺจ.}

\proseitem{24}{นเหตุํ โนนีวรณํ ธมฺมํ ปฏิจฺจ เหตุ นีวรโณ ธมฺโม อุปฺปชฺชติ เหตุปจฺจยา.\csromanspacer{}\csromanspacer{}เหตุยา ตีณิ.}

\proseitem{25}{นเหตุํ โนปรามาสํ ธมฺมํ ปฏิจฺจ นเหตุ ปรามาโส ธมฺโม อุปฺปชฺชติ เหตุปจฺจยา.\csromanspacer{}\csromanspacer{}เหตุยา ตีณิ.}

\csromanheader{2}{ธมฺมปจฺจนียานุโลม ทุกทุกปฏฺฐานปาฬิ \textbf{1. เหตุทุก 54. มหนฺตรทุก}}
\proseitem{26}{\csromanfolio{430}นเหตุํ นสารมฺมณํ ธมฺมํ ปฏิจฺจ เหตุ สารมฺมโณ ธมฺโม อุปฺปชฺชติ เหตุปจฺจยา.\csromanspacer{}\csromanspacer{}เหตุยา ตีณิ.}

\prose{นเหตุํ น-อนารมฺมณํ ธมฺมํ ปฏิจฺจ นเหตุ อนารมฺมโณ ธมฺโม อุปฺปชฺชติ เหตุปจฺจยา.\csromanspacer{}\csromanspacer{}เหตุยา ตีณิ.}

\proseitem{27}{นเหตุํ โนจิตฺตํ ธมฺมํ ปฏิจฺจ นเหตุ จิตฺโต ธมฺโม อุปฺปชฺชติ เหตุปจฺจยา.\csromanspacer{}\csromanspacer{}เหตุยา ตีณิ.}

\prose{นเหตุํ นโนจิตฺตํ ธมฺมํ ปฏิจฺจ เหตุ โนจิตฺโต ธมฺโม อุปฺปชฺชติ เหตุปจฺจยา.\csromanspacer{}\csromanspacer{}เหตุยา ตีณิ.}

\proseitem{28}{นเหตุํ นเจตสิกํ ธมฺมํ ปฏิจฺจ เหตุ เจตสิโก ธมฺโม อุปฺปชฺชติ เหตุปจฺจยา.\csromanspacer{}\csromanspacer{}เหตุยา ตีณิ.}

\prose{นเหตุํ น-อเจตสิกํ ธมฺมํ ปฏิจฺจ นเหตุ อเจตสิโก ธมฺโม อุปฺปชฺชติ เหตุปจฺจยา.\csromanspacer{}\csromanspacer{}เหตุยา ตีณิ.}

\proseitem{29}{นเหตุํ นจิตฺตสมฺปยุตฺตํ ธมฺมํ ปฏิจฺจ เหตุ จิตฺตสมฺปยุตฺโต ธมฺโม อุปฺปชฺชติ เหตุปจฺจยา.\csromanspacer{}\csromanspacer{}เหตุยา ตีณิ.}

\prose{นเหตุํ นจิตฺตวิปฺปยุตฺตํ ธมฺมํ ปฏิจฺจ นเหตุ จิตฺตวิปฺปยุตฺโต ธมฺโม อุปฺปชฺชติ เหตุปจฺจยา.\csromanspacer{}\csromanspacer{}เหตุยา ตีณิ.}

\proseitem{30}{นเหตุํ นจิตฺตสํสฏฺฐํ ธมฺมํ ปฏิจฺจ เหตุ จิตฺตสํสฏฺโฐ ธมฺโม อุปฺปชฺชติ เหตุปจฺจยา.\csromanspacer{}\csromanspacer{}เหตุยา ตีณิ.}

\prose{นเหตุํ นจิตฺตวิสํสฏฺฐํ ธมฺมํ ปฏิจฺจ นเหตุ จิตฺตวิสํสฏฺโฐ ธมฺโม อุปฺปชฺชติ เหตุปจฺจยา.\csromanspacer{}\csromanspacer{}เหตุยา ตีณิ.}

\proseitem{31}{นเหตุํ โนจิตฺตสมุฏฺฐานํ ธมฺมํ ปฏิจฺจ เหตุ จิตฺตสมุฏฺฐาโน ธมฺโม อุปฺปชฺชติ เหตุปจฺจยา.\csromanspacer{}\csromanspacer{}เหตุยา ตีณิ.}

\prose{นเหตุํ นโนจิตฺตสมุฏฺฐานํ ธมฺมํ ปฏิจฺจ เหตุ โนจิตฺตสมุฏฺฐาโน ธมฺโม อุปฺปชฺชติ เหตุปจฺจยา.\csromanspacer{}\csromanspacer{}เหตุยา ตีณิ.}

\csromanheader{2}{1. เหตุทุก 54. มหนฺตรทุก}
\proseitem{32}{\csromanfolio{431}นเหตุํ โนจิตฺตสหภุํ ธมฺมํ ปฏิจฺจ เหตุ จิตฺตสหภู ธมฺโม อุปฺปชฺชติ เหตุปจฺจยา.\csromanspacer{}\csromanspacer{}เหตุยา ตีณิ.}

\prose{นเหตุํ นโนจิตฺตสหภุํ ธมฺมํ ปฏิจฺจ นเหตุ โนจิตฺตสหภู ธมฺโม อุปฺปชฺชติ เหตุปจฺจยา.\csromanspacer{}\csromanspacer{}เหตุยา ตีณิ.}

\proseitem{33}{นเหตุํ นจิตฺตานุปริวตฺติํ ธมฺมํ ปฏิจฺจ เหตุ จิตฺตานุปริวตฺตี ธมฺโม อุปฺปชฺชติ เหตุปจฺจยา.\csromanspacer{}\csromanspacer{}เหตุยา ตีณิ.}

\prose{นเหตุํ นโนจิตฺตานุปริวตฺติํ ธมฺมํ ปฏิจฺจ นเหตุ โนจิตฺตานุปริวตฺตี ธมฺโม อุปฺปชฺชติ เหตุปจฺจยา.\csromanspacer{}\csromanspacer{}เหตุยา ตีณิ.}

\proseitem{34}{นเหตุํ นจิตฺตสํสฏฺฐสมุฏฺฐานํ ธมฺมํ ปฏิจฺจ เหตุ จิตฺตสํสฏฺฐสมุฏฺฐาโน ธมฺโม อุปฺปชฺชติ เหตุปจฺจยา.\csromanspacer{}\csromanspacer{}เหตุยา ตีณิ.}

\prose{นเหตุํ นโนจิตฺตสํสฏฺฐสมุฏฺฐานํ ธมฺมํ ปฏิจฺจ นเหตุ โนจิตฺตสํสฏฺฐสมุฏฺฐาโน ธมฺโม อุปฺปชฺชติ เหตุปจฺจยา.\csromanspacer{}\csromanspacer{}เหตุยา ตีณิ.}

\proseitem{35}{นเหตุํ โนจิตฺตสํสฏฺฐสมุฏฺฐานสหภุํ ธมฺมํ ปฏิจฺจ เหตุ จิตฺตสํสฏฺฐสมุฏฺฐานสหภู ธมฺโม อุปฺปชฺชติ เหตุปจฺจยา.\csromanspacer{}\csromanspacer{}เหตุยา ตีณิ.}

\prose{นเหตุํ นโนจิตฺตสํสฏฺฐสมุฏฺฐานสหภุํ ธมฺมํ ปฏิจฺจ นเหตุ โนจิตฺตสํสฏฺฐสมุฏฺฐานสหภู ธมฺโม อุปฺปชฺชติ เหตุปจฺจยา.\csromanspacer{}\csromanspacer{}เหตุยา ตีณิ.}

\proseitem{36}{นเหตุํ โนจิตฺตสํสฏฺฐสมุฏฺฐานานุปริวตฺติํ ธมฺมํ ปฏิจฺจ เหตุ จิตฺตสํสฏฺฐสมุฏฺฐานานุปริวตฺตี ธมฺโม อุปฺปชฺชติ เหตุปจฺจยา.\csromanspacer{}\csromanspacer{}เหตุยา ตีณิ.}

\prose{นเหตุํ นโนจิตฺตสํสฏฺฐสมุฏฺฐานานุปริวตฺติํ ธมฺมํ ปฏิจฺจ นเหตุ โนจิตฺตสํสฏฺฐสมุฏฺฐานานุปริวตฺตี ธมฺโม อุปฺปชฺชติ เหตุปจฺจยา.\csromanspacer{}\csromanspacer{}เหตุยา ตีณิ.}

\proseitem{37}{นเหตุํ น-อชฺฌตฺติกํ ธมฺมํ ปฏิจฺจ นเหตุ อชฺฌตฺติโก ธมฺโม อุปฺปชฺชติ เหตุปจฺจยา.\csromanspacer{}\csromanspacer{}เหตุยา ตีณิ.}

\csromanheader{2}{ธมฺมปจฺจนียานุโลม ทุกทุกปฏฺฐานปาฬิ}
\prose{\csromanfolio{432}นเหตุํ นพาหิรํ ธมฺมํ ปฏิจฺจ เหตุ พาหิโร ธมฺโม อุปฺปชฺชติ เหตุปจฺจยา.\csromanspacer{}\csromanspacer{}เหตุยา ตีณิ.}

\proseitem{38}{นเหตุํ น-อุปาทา ธมฺมํ ปฏิจฺจ นเหตุ อุปาทา ธมฺโม อุปฺปชฺชติ เหตุปจฺจยา.\csromanspacer{}\csromanspacer{}เหตุยา ตีณิ.}

\prose{นเหตุํ นโน-อุปาทา ธมฺมํ ปฏิจฺจ เหตุ โน-อุปาทา ธมฺโม อุปฺปชฺชติ เหตุปจฺจยา.\csromanspacer{}\csromanspacer{}เหตุยา ตีณิ.}

\proseitem{39}{นเหตุํ น-อนุปาทินฺนํ ธมฺมํ ปฏิจฺจ นเหตุ อนุปาทินฺโน ธมฺโม อุปฺปชฺชติ เหตุปจฺจยา.\csromanspacer{}\csromanspacer{}เหตุยา ตีณิ.}

\csromanheader{2}{1. เหตุทุก 68. อุปาทานโคจฺฉก}
\proseitem{40}{นเหตุํ น-อุปาทานํ ธมฺมํ ปฏิจฺจ เหตุ อุปาทาโน ธมฺโม อุปฺปชฺชติ เหตุปจฺจยา.\csromanspacer{}\csromanspacer{}เหตุยา นว.}

\csromanheader{2}{1. เหตุทุก 74. กิเลสโคจฺฉก}
\proseitem{41}{นเหตุํ นกิเลสํ ธมฺมํ ปฏิจฺจ เหตุ กิเลโส ธมฺโม อุปฺปชฺชติ เหตุปจฺจยา.\csromanspacer{}\csromanspacer{}เหตุยา ตีณิ.}

\csromanheader{2}{1. เหตุทุก 82. ปิฏฺฐิทุก}
\proseitem{42}{นเหตุํ นทสฺสเนน ปหาตพฺพํ ธมฺมํ ปจฺจยา เหตุ ทสฺสเนน ปหาตพฺโพ ธมฺโม อุปฺปชฺชติ เหตุปจฺจยา.\csromanspacer{}\csromanspacer{}เหตุยา ตีณิ.}

\prose{นเหตุํ นนทสฺสเนน ปหาตพฺพํ ธมฺมํ ปฏิจฺจ นเหตุ นทสฺสเนน ปหาตพฺโพ ธมฺโม อุปฺปชฺชติ เหตุปจฺจยา.\csromanspacer{}\csromanspacer{}เหตุยา ตีณิ.}

\proseitem{43}{นเหตุํ นภาวนาย ปหาตพฺพํ ธมฺมํ ปจฺจยา เหตุ ภาวนาย ปหาตพฺโพ ธมฺโม อุปฺปชฺชติ เหตุปจฺจยา.\csromanspacer{}\csromanspacer{}เหตุยา ตีณิ.}

\prose{นเหตุํ นนภาวนาย ปหาตพฺพํ ธมฺมํ ปฏิจฺจ นเหตุ นภาวนาย ปหาตพฺโพ ธมฺโม อุปฺปชฺชติ เหตุปจฺจยา.\csromanspacer{}\csromanspacer{}เหตุยา ตีณิ.}

\proseitem{44}{นนเหตุํ นทสฺสเนน ปหาตพฺพเหตุกํ ธมฺมํ ปฏิจฺจ นเหตุ ทสฺสเนน ปหาตพฺพเหตุโก ธมฺโม อุปฺปชฺชติ เหตุปจฺจยา.\csromanspacer{}\csromanspacer{}เหตุยา เอกํ.}

\csromanheader{2}{2. สเหตุกาทิทุก 1. เหตุทุก}
\prose{\csromanfolio{433}นเหตุํ นนทสฺสเนน ปหาตพฺพเหตุกํ ธมฺมํ ปฏิจฺจ นเหตุ นทสฺสเนน ปหาตพฺพเหตุโก ธมฺโม อุปฺปชฺชติ เหตุปจฺจยา.\csromanspacer{}\csromanspacer{}เหตุยา ตีณิ.}

\proseitem{45}{นนเหตุํ นภาวนาย ปหาตพฺพเหตุกํ ธมฺมํ ปฏิจฺจ นเหตุ ภาวนาย ปหาตพฺพเหตุโก ธมฺโม อุปฺปชฺชติ เหตุปจฺจยา.\csromanspacer{}\csromanspacer{}เหตุยา เอกํ.}

\prose{นเหตุํ นนภาวนาย ปหาตพฺพเหตุกํ ธมฺมํ ปฏิจฺจ นเหตุ นภาวนาย ปหาตพฺพเหตุโก ธมฺโม อุปฺปชฺชติ เหตุปจฺจยา.\csromanspacer{}\csromanspacer{}เหตุยา ตีณิ.}

\proseitem{46}{นเหตุํ นสวิตกฺกํ ธมฺมํ ปฏิจฺจ เหตุ สวิตกฺโก ธมฺโม อุปฺปชฺชติ เหตุปจฺจยา.\csromanspacer{}\csromanspacer{}เหตุยา ตีณิ. (สพฺพตฺถ สํขิตฺตํ.)}

\proseitem{47}{นเหตุํ นสรณํ ธมฺมํ ปจฺจยา เหตุ สรโณ ธมฺโม อุปฺปชฺชติ เหตุปจฺจยา.\csromanspacer{}\csromanspacer{}เหตุยา ตีณิ.}

\prose{นเหตุํ น-อรณํ ธมฺมํ ปฏิจฺจ นเหตุ อรโณ ธมฺโม อุปฺปชฺชติ เหตุปจฺจยา.\csromanspacer{}\csromanspacer{}นนเหตุํ น-อรณํ ธมฺมํ ปฏิจฺจ นเหตุ อรโณ ธมฺโม อุปฺปชฺชติ เหตุปจฺจยา.\csromanspacer{}\csromanspacer{}นเหตุํ น-อรณญฺจ นนเหตุํ น-อรณญฺจ ธมฺมํ ปฏิจฺจ นเหตุ อรโณ ธมฺโม อุปฺปชฺชติ เหตุปจฺจยา.\csromanspacer{}\csromanspacer{}เหตุยา ตีณิ. \textbf{2.}\csromanspacer{}\textbf{ สเหตุกาทิทุก 1.}\csromanspacer{}\textbf{ เหตุทุก}}

\proseitem{48}{นสเหตุกํ นเหตุํ ธมฺมํ ปฏิจฺจ สเหตุโก เหตุ ธมฺโม อุปฺปชฺชติ เหตุปจฺจยา.\csromanspacer{}\csromanspacer{}น-อเหตุกํ นเหตุํ ธมฺมํ ปฏิจฺจ สเหตุโก เหตุ ธมฺโม อุปฺปชฺชติ เหตุปจฺจยา.\csromanspacer{}\csromanspacer{}นสเหตุกํ นเหตุญฺจ น-อเหตุกํ นเหตุญฺจ ธมฺมํ ปฏิจฺจ สเหตุโก เหตุ ธมฺโม อุปฺปชฺชติ เหตุปจฺจยา.\csromanspacer{}\csromanspacer{}เหตุยา ตีณิ.}

\prose{นสเหตุกํ นนเหตุํ ธมฺมํ ปฏิจฺจ สเหตุโก นเหตุ ธมฺโม อุปฺปชฺชติ เหตุปจฺจยา.\csromanspacer{} ตีณิ.}

\prose{น-อเหตุกํ นนเหตุํ ธมฺมํ ปฏิจฺจ อเหตุโก นเหตุ ธมฺโม อุปฺปชฺชติ เหตุปจฺจยา.\csromanspacer{} ตีณิ.\csromanspacer{}\csromanspacer{}เหตุยา ฉ.}

\proseitem{49}{นเหตุสมฺปยุตฺตํ นเหตุํ ธมฺมํ ปฏิจฺจ เหตุสมฺปยุตฺโต เหตุ ธมฺโม อุปฺปชฺชติ เหตุปจฺจยา.\csromanspacer{} ตีณิ. (สเหตุกทุกสทิสํ.)}

\csromanheader{2}{ธมฺมปจฺจนียานุโลม ทุกทุกปฏฺฐานปาฬิ}
\proseitem{50}{\csromanfolio{434}นเหตุญฺเจว น-อเหตุกญฺจ นเหตุํ ธมฺมํ ปฏิจฺจ เหตุ เจว สเหตุโก จ เหตุ ธมฺโม อุปฺปชฺชติ เหตุปจฺจยา.\csromanspacer{}\csromanspacer{}เหตุยา เอกํ.}

\prose{น-อเหตุกญฺเจว นน จ เหตุํ นนเหตุํ ธมฺมํ ปฏิจฺจ สเหตุโก เจว น จ เหตุ นเหตุ ธมฺโม อุปฺปชฺชติ เหตุปจฺจยา.\csromanspacer{}\csromanspacer{}เหตุยา เอกํ.}

\proseitem{51}{นเหตุญฺเจว นเหตุวิปฺปยุตฺตญฺจ นเหตุํ ธมฺมํ ปฏิจฺจ เหตุ เจว เหตุสมฺปยุตฺโต จ เหตุ ธมฺโม อุปฺปชฺชติ เหตุปจฺจยา.\csromanspacer{}\csromanspacer{}เหตุยา เอกํ.}

\prose{นเหตุวิปฺปยุตฺตญฺเจว นนเหตุญฺจ นนเหตุํ ธมฺมํ ปฏิจฺจ เหตุสมฺปยุตฺโต เจว น จ เหตุ นเหตุ ธมฺโม อุปฺปชฺชติ เหตุปจฺจยา.\csromanspacer{}\csromanspacer{}เหตุยา เอกํ. (อนฺติมทุกํ น ลพฺภติ.)}

\csromanheader{2}{7-13. จูฬนฺตรทุก 1. เหตุทุก}
\proseitem{52}{น-อปฺปจฺจยํ นเหตุํ ธมฺมํ ปฏิจฺจ สปฺปจฺจโย เหตุ ธมฺโม อุปฺปชฺชติ เหตุปจฺจยา.\csromanspacer{}\csromanspacer{}เหตุยา เอกํ.}

\prose{น-อปฺปจฺจยํ นนเหตุํ ธมฺมํ ปฏิจฺจ สปฺปจฺจโย นเหตุ ธมฺโม อุปฺปชฺชติ เหตุปจฺจยา.\csromanspacer{}\csromanspacer{}เหตุยา เอกํ. (สงฺขตํ สปฺปจฺจยสทิสํ.)}

\proseitem{53}{นสนิทสฺสนํ นเหตุํ ธมฺมํ ปฏิจฺจ อนิทสฺสโน เหตุ ธมฺโม อุปฺปชฺชติ เหตุปจฺจยา.\csromanspacer{}\csromanspacer{}เหตุยา เอกํ.}

\prose{นสนิทสฺสนํ นนเหตุํ ธมฺมํ ปฏิจฺจ สนิทสฺสโน นเหตุ ธมฺโม อุปฺปชฺชติ เหตุปจฺจยา.\csromanspacer{}\csromanspacer{}เหตุยา ตีณิ.}

\proseitem{54}{นสปฺปฏิฆํ นเหตุํ ธมฺมํ ปฏิจฺจ อปฺปฏิโฆ เหตุ ธมฺโม อุปฺปชฺชติ เหตุปจฺจยา.\csromanspacer{}\csromanspacer{}เหตุยา เอกํ.}

\prose{นสปฺปฏิฆํ นนเหตุํ ธมฺมํ ปฏิจฺจ สปฺปฏิโฆ นเหตุ ธมฺโม อุปฺปชฺชติ เหตุปจฺจยา.\csromanspacer{}\csromanspacer{}เหตุยา ตีณิ.}

\proseitem{55}{นรูปิํ นเหตุํ ธมฺมํ ปฏิจฺจ อรูปี เหตุ ธมฺโม อุปฺปชฺชติ เหตุปจฺจยา.}

\csromanheader{2}{14. อาสวโคจฺฉก 1. เหตุทุก}
\prose{\csromanfolio{435}นรูปิํ นนเหตุํ ธมฺมํ ปฏิจฺจ อรูปี นเหตุ ธมฺโม อุปฺปชฺชติ เหตุปจฺจยา.\csromanspacer{} คณิตเกน ตีณิ.}

\proseitem{56}{นโลกิยํ นเหตุํ ธมฺมํ ปฏิจฺจ โลกุตฺตโร เหตุ ธมฺโม อุปฺปชฺชติ เหตุปจฺจยา.\csromanspacer{} เอกํ.}

\prose{นโลกุตฺตรํ นเหตุํ ธมฺมํ ปฏิจฺจ โลกิโย เหตุ ธมฺโม อุปฺปชฺชติ เหตุปจฺจยา.\csromanspacer{} เอกํ.\csromanspacer{}\csromanspacer{}เหตุยา เทฺว.}

\prose{นโลกิยํ นนเหตุํ ธมฺมํ ปฏิจฺจ โลกิโย นเหตุ ธมฺโม อุปฺปชฺชติ เหตุปจฺจยา.\csromanspacer{} ตีณิ.}

\prose{นโลกุตฺตรํ นนเหตุํ ธมฺมํ ปฏิจฺจ โลกิโย นเหตุ ธมฺโม อุปฺปชฺชติ เหตุปจฺจยา.\csromanspacer{} เอกํ.\csromanspacer{}\csromanspacer{}เหตุยา จตฺตาริ.}

\proseitem{57}{นเกนจิ วิญฺเญยฺยํ นเหตุํ ธมฺมํ ปฏิจฺจ เกนจิ วิญฺเญยฺโย เหตุ ธมฺโม อุปฺปชฺชติ เหตุปจฺจยา.\csromanspacer{}\csromanspacer{}เหตุยา นว.}

\prose{นเกนจิ นวิญฺเญยฺยํ นนเหตุํ ธมฺมํ ปฏิจฺจ เกนจิ นวิญฺเญยฺโย นเหตุ ธมฺโม อุปฺปชฺชติ เหตุปจฺจยา.\csromanspacer{}\csromanspacer{}เหตุยา นว.}

\csromanheader{2}{14. อาสวโคจฺฉก 1. เหตุทุก}
\proseitem{58}{โน-อาสวํ นเหตุํ ธมฺมํ ปฏิจฺจ อาสโว เหตุ ธมฺโม อุปฺปชฺชติ เหตุปจฺจยา.\csromanspacer{} ตีณิ.}

\prose{นโน-อาสวํ นเหตุํ ธมฺมํ ปฏิจฺจ อาสโว เหตุ ธมฺโม อุปฺปชฺชติ เหตุปจฺจยา.\csromanspacer{} เอกํ.}

\prose{โน-อาสวํ นเหตุญฺจ นโน-อาสวํ นเหตุญฺจ ธมฺมํ ปฏิจฺจ อาสโว เหตุ ธมฺโม อุปฺปชฺชติ เหตุปจฺจยา.\csromanspacer{} เอกํ.\csromanspacer{}\csromanspacer{}เหตุยา ปญฺจ.}

\prose{โน-อาสวํ นนเหตุํ ธมฺมํ ปฏิจฺจ โน-อาสโว นเหตุ ธมฺโม อุปฺปชฺชติ เหตุปจฺจยา.\csromanspacer{} เอกํ.}

\prose{นโน-อาสวํ นนเหตุํ ธมฺมํ ปฏิจฺจ โน-อาสโว นเหตุ ธมฺโม อุปฺปชฺชติ เหตุปจฺจยา.\csromanspacer{} ตีณิ.}

\csromanheader{2}{ธมฺมปจฺจนียานุโลม ทุกทุกปฏฺฐานปาฬิ}
\prose{\csromanfolio{436}โน-อาสวํ นนเหตุญฺจ นโน-อาสวํ นนเหตุญฺจ ธมฺมํ ปฏิจฺจ โน-อาสโว นเหตุ ธมฺโม อุปฺปชฺชติ เหตุปจฺจยา.\csromanspacer{} เอกํ.\csromanspacer{}\csromanspacer{}เหตุยา ปญฺจ.}

\csromanheader{2}{55. มหนฺตรทุก 1. เหตุทุก}
\proseitem{59}{นสารมฺมณํ นเหตุํ ธมฺมํ ปฏิจฺจ สารมฺมโณ เหตุ ธมฺโม อุปฺปชฺชติ เหตุปจฺจยา.\csromanspacer{}\csromanspacer{}เหตุยา ตีณิ.}

\prose{น-อนารมฺมณํ นนเหตุํ ธมฺมํ ปฏิจฺจ สารมฺมโณ นเหตุ ธมฺโม อุปฺปชฺชติ เหตุปจฺจยา.\csromanspacer{}\csromanspacer{}เหตุยา ตีณิ. (สํขิตฺตํ.)}

\csromanheader{2}{100. สรณทุก 1. เหตุทุก}
\proseitem{60}{นสรณํ นเหตุํ ธมฺมํ ปฏิจฺจ อรโณ เหตุ ธมฺโม อุปฺปชฺชติ เหตุปจฺจยา.\csromanspacer{}\csromanspacer{}น-อรณํ นเหตุํ ธมฺมํ ปฏิจฺจ สรโณ เหตุ ธมฺโม อุปฺปชฺชติ เหตุปจฺจยา.\csromanspacer{}\csromanspacer{}เหตุยา เทฺว.}

\prose{นสรณํ นนเหตุํ ธมฺมํ ปฏิจฺจ อรโณ นเหตุ ธมฺโม อุปฺปชฺชติ เหตุปจฺจยา. (1)}

\prose{น-อรณํ นนเหตุํ ธมฺมํ ปฏิจฺจ อรโณ นเหตุ ธมฺโม อุปฺปชฺชติ เหตุปจฺจยา.\csromanspacer{}\csromanspacer{}น-อรณํ นนเหตุํ ธมฺมํ ปฏิจฺจ สรโณ นเหตุ ธมฺโม อุปฺปชฺชติ เหตุปจฺจยา.\csromanspacer{}\csromanspacer{}น-อรณํ นนเหตุํ ธมฺมํ ปฏิจฺจ สรโณ นเหตุ จ อรโณ นเหตุ จ ธมฺมา อุปฺปชฺชนฺติ เหตุปจฺจยา.\csromanspacer{} ตีณิ. (สํขิตฺตํ.)}

\csromanheader{2}{100. สรณทุก 2. สเหตุกทุก}
\proseitem{61}{นสรณํ นสเหตุกํ ธมฺมํ ปฏิจฺจ อรโณ สเหตุโก ธมฺโม อุปฺปชฺชติ เหตุปจฺจยา.\csromanspacer{}\csromanspacer{}น-อรณํ นสเหตุกํ ธมฺมํ ปฏิจฺจ สรโณ สเหตุโก ธมฺโม อุปฺปชฺชติ เหตุปจฺจยา.\csromanspacer{}\csromanspacer{}เหตุยา เทฺว.}

\proseitem{62}{นสรณํ น-อเหตุกํ ธมฺมํ ปฏิจฺจ อรโณ อเหตุโก ธมฺโม อุปฺปชฺชติ เหตุปจฺจยา.\csromanspacer{}\csromanspacer{}น-อรณํ น-อเหตุกํ ธมฺมํ ปฏิจฺจ อรโณ อเหตุโก ธมฺโม อุปฺปชฺชติ เหตุปจฺจยา.\csromanspacer{}\csromanspacer{}เหตุยา เทฺว.}

\vspace{\tipitakaparskip}
\csromancenter{สรณทุก 6.\csromanspacer{} นเหตุสเหตุกทุก \textbf{100.}\csromanspacer{}\textbf{ สรณทุก 3.}\csromanspacer{}\textbf{ เหตุสมฺปยุตฺตทุก}}
\proseitem{63}{\csromanfolio{437}นสรณํ นเหตุสมฺปยุตฺตํ ธมฺมํ ปฏิจฺจ อรโณ เหตุสมฺปยุตฺโต ธมฺโม อุปฺปชฺชติ เหตุปจฺจยา.\csromanspacer{}\csromanspacer{}น-อรณํ นเหตุสมฺปยุตฺตํ ธมฺมํ ปฏิจฺจ สรโณ เหตุสมฺปยุตฺโต ธมฺโม อุปฺปชฺชติ เหตุปจฺจยา.\csromanspacer{}\csromanspacer{}เหตุยา เทฺว.}

\proseitem{64}{นสรณํ นเหตุวิปฺปยุตฺตํ ธมฺมํ ปฏิจฺจ อรโณ เหตุวิปฺปยุตฺโต ธมฺโม อุปฺปชฺชติ เหตุปจฺจยา.\csromanspacer{}\csromanspacer{}น-อรณํ นเหตุวิปฺปยุตฺตํ ธมฺมํ ปฏิจฺจ อรโณ เหตุวิปฺปยุตฺโต ธมฺโม อุปฺปชฺชติ เหตุปจฺจยา.\csromanspacer{}\csromanspacer{}เหตุยา เทฺว.}

\csromanheader{2}{100. สรณทุก 4. เหตุสเหตุกทุก}
\proseitem{65}{นสรณํ นเหตุญฺเจว น-อเหตุกญฺจ ธมฺมํ ปฏิจฺจ อรโณ เหตุ เจว สเหตุโก จ ธมฺโม อุปฺปชฺชติ เหตุปจฺจยา.\csromanspacer{}\csromanspacer{}น-อรณํ นเหตุญฺเจว น-อเหตุกญฺจ ธมฺมํ ปฏิจฺจ สรโณ เหตุ เจว สเหตุโก จ ธมฺโม อุปฺปชฺชติ เหตุปจฺจยา.\csromanspacer{}\csromanspacer{}เหตุยา เทฺว.}

\prose{นสรณํ น-อเหตุกญฺเจว นนเหตุญฺจ ธมฺมํ ปฏิจฺจ อรโณ สเหตุโก เจว น จ เหตุ ธมฺโม อุปฺปชฺชติ เหตุปจฺจยา.\csromanspacer{}\csromanspacer{}น-อรณํ น-อเหตุกญฺเจว นนเหตุญฺจ ธมฺมํ ปฏิจฺจ สรโณ สเหตุโก เจว น จ เหตุ ธมฺโม อุปฺปชฺชติ เหตุปจฺจยา.\csromanspacer{}\csromanspacer{}เหตุยา เทฺว.\csromanspacer{}\csromanspacer{}(เหตุเหตุสมฺปยุตฺตทุกํ สํขิตฺตํ.)}

\csromanheader{2}{100. สรณทุก 6. นเหตุสเหตุกทุก}
\proseitem{66}{นสรณํ นเหตุํ นสเหตุกํ ธมฺมํ ปฏิจฺจ อรโณ นเหตุ สเหตุโก ธมฺโม อุปฺปชฺชติ เหตุปจฺจยา.\csromanspacer{}\csromanspacer{}เหตุยา เอกํ.}

\prose{นสรณํ นเหตุํ น-อเหตุกํ ธมฺมํ ปฏิจฺจ อรโณ นเหตุ อเหตุโก ธมฺโม อุปฺปชฺชติ เหตุปจฺจยา.\csromanspacer{} เอกํ.}

\prose{น-อรณํ นเหตุํ น-อเหตุกํ ธมฺมํ ปฏิจฺจ อรโณ นเหตุ อเหตุโก ธมฺโม อุปฺปชฺชติ เหตุปจฺจยา.\csromanspacer{} เอกํ.\csromanspacer{}\csromanspacer{}เหตุยา เทฺว.}

\csromanheader{2}{ธมฺมปจฺจนียานุโลม ทุกทุกปฏฺฐานปาฬิ \textbf{100. สรณทุก 7. จูฬนฺตรทุก}}
\proseitem{67}{\csromanfolio{438}นสรโณ นสปฺปจฺจโย ธมฺโม อรณสฺส สปฺปจฺจยสฺส ธมฺมสฺส อารมฺมณปจฺจเยน ปจฺจโย.\csromanspacer{}\csromanspacer{}อารมฺมเณ เอกํ.}

\csromanheader{2}{68. นสรโณ นสงฺขโต ธมฺโม. (สํขิตฺตํ.)}
\proseitem{69}{นสรณํ นสนิทสฺสนํ ธมฺมํ ปฏิจฺจ อรโณ สนิทสฺสโน ธมฺโม อุปฺปชฺชติ เหตุปจฺจยา.\csromanspacer{}\csromanspacer{}เหตุยา ตีณิ.}

\prose{นสรโณ น-อนิทสฺสโน ธมฺโม สรณสฺส อนิทสฺสนสฺส ธมฺมสฺส อารมฺมณปจฺจเยน ปจฺจโย.\csromanspacer{}\csromanspacer{}นสรโณ น-อนิทสฺสโน ธมฺโม อรณสฺส อนิทสฺสนสฺส ธมฺมสฺส อารมฺมณปจฺจเยน ปจฺจโย.\csromanspacer{}\csromanspacer{}อารมฺมเณ เทฺว.}

\proseitem{70}{นสรณํ นสปฺปฏิฆํ ธมฺมํ ปฏิจฺจ อรโณ สปฺปฏิโฆ ธมฺโม อุปฺปชฺชติ เหตุปจฺจยา.\csromanspacer{}\csromanspacer{}เหตุยา ตีณิ.}

\proseitem{71}{นสรณํ นรูปิํ ธมฺมํ ปฏิจฺจ อรโณ รูปี ธมฺโม อุปฺปชฺชติ เหตุปจฺจยา.\csromanspacer{}\csromanspacer{}น-อรณํ นรูปิํ ธมฺมํ ปฏิจฺจ อรโณ รูปี ธมฺโม อุปฺปชฺชติ เหตุปจฺจยา.\csromanspacer{}\csromanspacer{}เหตุยา เทฺว.}

\prose{นสรณํ น-อรูปิํ ธมฺมํ ปจฺจยา สรโณ อรูปี ธมฺโม อุปฺปชฺชติ เหตุปจฺจยา.\csromanspacer{}\csromanspacer{}นสรณํ น-อรูปิํ ธมฺมํ ปจฺจยา อรโณ อรูปี ธมฺโม อุปฺปชฺชติ เหตุปจฺจยา.\csromanspacer{}\csromanspacer{}เหตุยา เทฺว.}

\proseitem{72}{นสรณํ นโลกิยํ ธมฺมํ ปฏิจฺจ อรโณ โลกิโย ธมฺโม อุปฺปชฺชติ เหตุปจฺจยา.\csromanspacer{}\csromanspacer{}เหตุยา เอกํ.}

\prose{นสรณํ นโลกุตฺตรํ ธมฺมํ ปจฺจยา อรโณ โลกุตฺตโร ธมฺโม อุปฺปชฺชติ เหตุปจฺจยา.\csromanspacer{}\csromanspacer{}เหตุยา เอกํ.}

\proseitem{73}{นสรณํ นเกนจิ วิญฺเญยฺยํ ธมฺมํ ปฏิจฺจ อรโณ เกนจิ วิญฺเญยฺโย ธมฺโม อุปฺปชฺชติ เหตุปจฺจยา.\csromanspacer{}\csromanspacer{}เหตุยา ปญฺจ.}

\prose{นสรณํ นเกนจิ นวิญฺเญยฺยํ ธมฺมํ ปฏิจฺจ อรโณ เกนจิ นวิญฺเญยฺโย ธมฺโม อุปฺปชฺชติ เหตุปจฺจยา.\csromanspacer{}\csromanspacer{}เหตุยา ปญฺจ.}

\vspace{\tipitakaparskip}
\csromancenter{สรณทุก 55.\csromanspacer{} มหนฺตรทุก \textbf{100.}\csromanspacer{}\textbf{ สรณทุก 14.}\csromanspacer{}\textbf{ อาสวโคจฺฉกาทิ}}
\proseitem{74}{\csromanfolio{439}น-อรณํ น-อาสวํ ธมฺมํ ปฏิจฺจ สรโณ อาสโว ธมฺโม อุปฺปชฺชติ เหตุปจฺจยา.\csromanspacer{}\csromanspacer{}เหตุยา เอกํ.}

\proseitem{75}{น-อรณํ นสญฺโญชนํ ธมฺมํ ปฏิจฺจ สรโณ สญฺโญชโน ธมฺโม อุปฺปชฺชติ เหตุปจฺจยา.\csromanspacer{}\csromanspacer{}เหตุยา เอกํ.}

\proseitem{76}{น-อรณํ นคนฺถํ ธมฺมํ ปฏิจฺจ สรโณ คนฺโถ ธมฺโม อุปฺปชฺชติ เหตุปจฺจยา.\csromanspacer{}\csromanspacer{}เหตุยา เอกํ.}

\proseitem{77}{น-อรณํ น-โอฆํ ธมฺมํ ปฏิจฺจ สรโณ โอโฆ ธมฺโม อุปฺปชฺชติ เหตุปจฺจยา.\csromanspacer{}\csromanspacer{}เหตุยา เอกํ.}

\proseitem{78}{น-อรณํ โนโยคํ ธมฺมํ ปฏิจฺจ สรโณ โยโค ธมฺโม อุปฺปชฺชติ เหตุปจฺจยา.\csromanspacer{}\csromanspacer{}เหตุยา เอกํ.}

\proseitem{79}{น-อรณํ นนีวรณํ ธมฺมํ ปฏิจฺจ สรโณ นีวรโณ ธมฺโม อุปฺปชฺชติ เหตุปจฺจยา.\csromanspacer{}\csromanspacer{}เหตุยา เอกํ.}

\proseitem{80}{น-อรณํ นปรามาสํ ธมฺมํ ปฏิจฺจ สรโณ ปรามาโส ธมฺโม อุปฺปชฺชติ เหตุปจฺจยา.\csromanspacer{}\csromanspacer{}เหตุยา เอกํ.}

\csromanheader{2}{100. สรณทุก 55. มหนฺตรทุก}
\proseitem{81}{นสรณํ นสารมฺมณํ ธมฺมํ ปฏิจฺจ อรโณ สารมฺมโณ ธมฺโม อุปฺปชฺชติ เหตุปจฺจยา.\csromanspacer{}\csromanspacer{}เหตุยา เอกํ.}

\prose{นสรณํ น-อนารมฺมณํ ธมฺมํ ปฏิจฺจ อรโณ อนารมฺมโณ ธมฺโม อุปฺปชฺชติ เหตุปจฺจยา.\csromanspacer{}\csromanspacer{}น-อรณํ น-อนารมฺมณํ ธมฺมํ ปฏิจฺจ อรโณ อนารมฺมโณ ธมฺโม อุปฺปชฺชติ เหตุปจฺจยา.\csromanspacer{}\csromanspacer{}เหตุยา เทฺว.}

\proseitem{82}{นสรณํ นจิตฺตํ ธมฺมํ ปฏิจฺจ อรโณ จิตฺโต ธมฺโม อุปฺปชฺชติ เหตุปจฺจยา.\csromanspacer{}\csromanspacer{}น-อรณํ นจิตฺตํ ธมฺมํ ปฏิจฺจ สรโณ จิตฺโต ธมฺโม อุปฺปชฺชติ เหตุปจฺจยา.\csromanspacer{}\csromanspacer{}เหตุยา เทฺว. (สํขิตฺตํ.)}

\csromanheader{2}{ธมฺมปจฺจนียานุโลม ทุกทุกปฏฺฐานปาฬิ}
\proseitem{83}{\csromanfolio{440}นสรณํ นเจตสิกํ ธมฺมํ ปฏิจฺจ อรโณ เจตสิโก ธมฺโม อุปฺปชฺชติ เหตุปจฺจยา.\csromanspacer{}\csromanspacer{}น-อรณํ นเจตสิกํ ธมฺมํ ปฏิจฺจ สรโณ เจตสิโก ธมฺโม อุปฺปชฺชติ เหตุปจฺจยา.\csromanspacer{}\csromanspacer{}เหตุยา เทฺว.}

\proseitem{84}{นสรณํ นจิตฺตสมฺปยุตฺตํ ธมฺมํ ปฏิจฺจ อรโณ จิตฺตสมฺปยุตฺโต ธมฺโม อุปฺปชฺชติ เหตุปจฺจยา.\csromanspacer{}\csromanspacer{}น-อรณํ นจิตฺตสมฺปยุตฺตํ ธมฺมํ ปฏิจฺจ สรโณ จิตฺตสมฺปยุตฺโต ธมฺโม อุปฺปชฺชติ เหตุปจฺจยา.\csromanspacer{}\csromanspacer{}เหตุยา เทฺว.}

\proseitem{85}{นสรณํ นจิตฺตสํสฏฺฐํ ธมฺมํ ปฏิจฺจ อรโณ จิตฺตสํสฏฺโฐ ธมฺโม อุปฺปชฺชติ เหตุปจฺจยา.\csromanspacer{}\csromanspacer{}น-อรณํ นจิตฺตสํสฏฺฐํ ธมฺมํ ปฏิจฺจ สรโณ จิตฺตสํสฏฺโฐ ธมฺโม อุปฺปชฺชติ เหตุปจฺจยา.\csromanspacer{}\csromanspacer{}เหตุยา เทฺว.}

\csromanheader{2}{100. สรณทุก 83. ปิฏฺฐิทุก}
\proseitem{86}{นสรณํ นทสฺสเนน ปหาตพฺพํ ธมฺมํ ปฏิจฺจ สรโณ ทสฺสเนน ปหาตพฺโพ ธมฺโม อุปฺปชฺชติ เหตุปจฺจยา.\csromanspacer{}\csromanspacer{}เหตุยา เอกํ.}

\prose{น-อรณํ นนทสฺสเนน ปหาตพฺพํ ธมฺมํ ปฏิจฺจ อรโณ นทสฺสเนน ปหาตพฺโพ ธมฺโม อุปฺปชฺชติ เหตุปจฺจยา.\csromanspacer{}\csromanspacer{}เหตุยา เอกํ.}

\proseitem{87}{นสรณํ นส-อุตฺตรํ ธมฺมํ ปฏิจฺจ อรโณ ส-อุตฺตโร ธมฺโม อุปฺปชฺชติ เหตุปจฺจยา.\csromanspacer{}\csromanspacer{}เหตุยา เอกํ.}

\prose{(สหชาตวารมฺปิ ปจฺจยวารมฺปิ นิสฺสยวารมฺปิ สํสฏฺฐวารมฺปิ สมฺปยุตฺตวารมฺปิ ปฏิจฺจวารสทิสํ.)}

\csromanheader{2}{\textbf{ปญฺหาวาร เหตุ-อารมฺมณ}}
\proseitem{88}{นสรโณ นส-อุตฺตโร ธมฺโม อรณสฺส ส-อุตฺตรสฺส ธมฺมสฺส เหตุปจฺจเยน ปจฺจโย.\csromanspacer{} เอกํ.}

\prose{นสรโณ นส-อุตฺตโร ธมฺโม อรณสฺส ส-อุตฺตรสฺส ธมฺมสฺส อารมฺมณปจฺจเยน ปจฺจโย.\csromanspacer{} เอกํ.\csromanspacer{}\csromanspacer{}เหตุยา เอกํ.}

\csromanheader{2}{\textbf{ปจฺจนียุทฺธาร} \textbf{ปจฺจนียุทฺธาร}}
\proseitem{89}{\csromanfolio{441}นสรโณ นส-อุตฺตโร ธมฺโม อรณสฺส ส-อุตฺตรสฺส ธมฺมสฺส อารมฺมณปจฺจเยน ปจฺจโย.\csromanspacer{} สหชาตปจฺจเยน ปจฺจโย.\csromanspacer{} อุปนิสฺสยปจฺจเยน ปจฺจโย.\csromanspacer{} ปจฺฉาชาตปจฺจเยน ปจฺจโย. (สํขิตฺตํ.) .\csromanspacer{} นเหตุยา เอกํ, น-อารมฺมเณ เอกํ.}

\csromanheader{2}{เหตุปจฺจยา น-อารมฺมเณ เอกํ. (สํขิตฺตํ.)}
\prose{นเหตุปจฺจยา อารมฺมเณ เอกํ. (สํขิตฺตํ.) (ยถา กุสลตฺติเก ปญฺหาวารํ เอวํ วิตฺถาเรตพฺพํ.)}

\csromanheader{2}{\textbf{อนุตฺตรปท-เหตุ-อนนฺตร}}
\proseitem{90}{นสรณํ น-อนุตฺตรํ ธมฺมํ ปจฺจยา อรโณ อนุตฺตโร ธมฺโม อุปฺปชฺชติ เหตุปจฺจยา.\csromanspacer{}\csromanspacer{}เหตุยา เอกํ ฯเปฯ อวิคเต เอกํ.}

\proseitem{91}{นสรโณ น-อนุตฺตโร ธมฺโม อรณสฺส อนุตฺตรสฺส ธมฺมสฺส อนนฺตรปจฺจเยน ปจฺจโย.\csromanspacer{}\csromanspacer{}อนนฺตเร เอกํ, สมนนฺตเร เอกํ, อุปนิสฺสเย เทฺว, ปุเรชาเต เอกํ อาเสวเน เอกํ, วิปฺปยุตฺเต เอกํ, อตฺถิยา เอกํ, นตฺถิยา เอกํ, วิคเต เอกํ, อวิคเต เอกํ.}

\csromanheader{2}{\textbf{ปจฺจนียุทฺธาร}}
\proseitem{92}{นสรโณ น-อนุตฺตโร ธมฺโม อรณสฺส อนุตฺตรสฺส ธมฺมสฺส อุปนิสฺสยปจฺจเยน ปจฺจโย.\csromanspacer{} ปุเรชาตปจฺจเยน ปจฺจโย.}

\prose{น-อรโณ น-อนุตฺตโร ธมฺโม อรณสฺส อนุตฺตรสฺส ธมฺมสฺส อุปนิสฺสยปจฺจเยน ปจฺจโย.\csromanspacer{}\csromanspacer{}นเหตุยา เทฺว, น-อารมฺมเณ เทฺว ฯเปฯ น-อุปนิสฺสเย เอกํ, นปุเรชาเต เทฺว ฯเปฯ โน-อวิคเต เทฺว.}

\csromanheader{2}{อุปนิสฺสยปจฺจยา นเหตุยา เทฺว. (สํขิตฺตํ.)}
\csromanheader{2}{ธมฺมปจฺจนียานุโลม ทุกทุกปฏฺฐานปาฬิ}
\prose{\csromanfolio{442}นเหตุปจฺจยา อุปนิสฺสเย เทฺว, ปุเรชาเต เอกํ ฯเปฯ อตฺถิยา เอกํ ฯเปฯ อวิคเต เอกํ. (สํขิตฺตํ.) (ยถา กุสลตฺติเก ปญฺหาวารํ เอวํ วิตฺถาเรตพฺพํ.)}

\prose{อนุโลมทุกติกปฏฺฐานโต ปฏฺฐาย ยาว ปริโยสานา ติํสมตฺเตหิ ภาณวาเรหิ ปฏฺฐานํ.}

\nitthitam{ธมฺมปจฺจนียานุโลเม ทุกทุกปฏฺฐานํ นิฏฺฐิตํ.}
\nitthitam{ปจฺจนียานุโลมปฏฺฐานํ นิฏฺฐิตํ.}
\nitthitam{\textbf{ปฏฺฐานปกรณํ นิฏฺฐิตํ.}}
\orphannote{( ) * อยํ สงฺขฺยา วิจาเรตพฺพา, นโลกิยนโลกุตฺตรธมฺโม นาม นตฺถิ.}

